\sekciono{A}
\begin{multicols}{2}
\artiklo{a}%
\vorto{a}\fako{muziko} La noto sisesma di la « do »-gamo\lingui{DE}
\artiklo{-a}%
\vorto{-a} La litero lasta di adjektivo
\artiklo{a(d)}%
\vorto{a(d)} Prepoziciono quan onu uzas por indikar la skopo di la ago, la loko quan onu volas atingar, la destinario, la persono a qua onu donas, donacas od atribuas ulo. Konseque, lu indikas logikale la objekto di ta o ca penso o sentimento, opoze a « subjekto » qua recevas li
\artiklo{-ab-}%
\vorto{-ab-} Sufixo qua indikas, che verbo, la anteeso kompare a tempo-punto di epoko pasinta o futura
\artiklo{abado}%
\vorto{abado} Superioro di monakeyo (katolika) de viri, di olqua la revenui esis benefico por ilta qua havis la ofico\lingui{DEFIRS}
\artiklo{abako}%
\vorto{abako}\senco{1}{\textit{(arkitekt.)} Parto supra di la kapitelo di kolono, qua havas la formo di tableto}\senco{2}{\textit{(matem.)} Aparato por kalkulo, de origino Greka, qua konsistas ye tabelo, dividita per ula quanto de ranuri paralela reciproke, en qui povis glitar buli movebla}\lingui{DEFIRS}
\artiklo{abandonar}%
\vorto{abandonar}\fako{trans.} Cesar sua sorgo (di persono, ofico, posteno) sive nevolunte o kontre-vole, sive violacante ula devo. (ex. : abandonar sua filii, sua gepatri, sua posteno)\lingui{DEFIS}
\artiklo{abaniko}%
\vorto{abaniko} Mi-cirklo de stofo, de papero, de plumi, muntita sur lameli movebla, quan onu desfaldas por agitar la aero cirkum su e tale produktar minvarmeso\lingui{SP}
\artiklo{abasar}%
\vorto{abasar}\fako{trans.} Pozar (plu) infre. – \textit{(metaf.)} \textit{(aludante persono)} Retroduktar (lu) a rango infra, a situeso sociala min-supra. (Onu abasas su per situar su a rango infra, per adoptar posturo humila)\lingui{DEFI}
\artiklo{abatar}%
\vorto{abatar}\fako{trans.} Faligar, per plura frapi (ulu od ulo qua stacas)\lingui{DFIS}
\artiklo{abatiso}%
\vorto{abatiso} Parti acesora di animali (precipue pultro) quin onu buchas por manjo
\artiklo{abceso}%
\vorto{abceso}\fako{patol.} Kavajo pusoza en la korpo homala\lingui{DEFIS}
\artiklo{abciso}%
\vorto{abciso}\fako{geom.} Un ek la du koordinati rekta-linea qui uzesas por determinar la situeso di punto en spaco\lingui{DEFIRS}
\artiklo{abdikar}%
\vorto{abdikar}\fako{trans.) (pri rejo od altra suvereno} Cesar sua regno per transmisar la autoritato imperala a sucedanto, o naturala, o selektita\lingui{DEFIRS}
\artiklo{abdomino}%
\vorto{abdomino}\fako{anat.} Korpo-kavajo qua kontenas la viceri digestala. (Ta vorto ciencala esas plu preciza kam « ventro »)\lingui{DEFIS}
\artiklo{abduktar}%
\vorto{abduktar}\fako{trans.} Movigar parto di la homo-korpo tale ke la movajo eskartesas de la plano mediana di la korpo\lingui{DEFIS}
\artiklo{abelo}%
\vorto{abelo}\fako{zool.} Insekto himenoptera qua vivas esame e produktas mielo e vaxo\latina{apis mellifica}\lingui{DEFIS}
\artiklo{aberaco}%
\vorto{aberaco}\senco{1}{\textit{(astron.)} Deviaceso semblanta di la radii luma qui venas de astro}\senco{2}{\textit{(fiziko)} Disperseso di la radii luma qui trairas lenso}\lingui{DEFIRS}
\artiklo{abieto}%
\vorto{abieto}\fako{bot.} Arboro rezinifera, di qua la folii nula-sezone divenas flavatra e sika\latina{abies pectinata}\lingui{EFIS}
\artiklo{abismo}%
\vorto{abismo}\senco{1}{Profundegajo qua ne esas mezurebla pro la grandeso di lua dimensioni}\senco{2}{\textit{(metaf.)}(a) Profundegajo quan spirito ne povas mezurar konceptale. (b) Profundegajo en qua onu esas quaze nihiligata, su-egaranta}\lingui{EFIS}
\artiklo{abiso}%
\vorto{abiso} Profundegajo submara
\artiklo{abjekta}%
\vorto{abjekta} Qua esas ye la grado maxim infra di desnobleso, pri la hierarkio sociala e pri lua karaktero personala\lingui{DEFIS}
\artiklo{abjurar}%
\vorto{abjurar}\fako{trans.}\senco{1}{Renuncar, en su-obligo solena, la religio quan onu profesis til-lore}\senco{2}{\textit{(metaf.)} Renuncar publike to quon onu profesis kredar, amar, prizar}\lingui{DEFIS}
\artiklo{ablacionar}%
\vorto{ablacionar}\fako{trans.) (kirurg.} Deprenar de la korpo homala (per operaco kirurgiala) parto morboza, malada\lingui{DEFIS}
\artiklo{ablaktar}%
\vorto{ablaktar}\fako{trans.} Separar infanteto de lua patrino o de la salariato qua alaktis lu, kande oportas komencar donar a lu alimente nutrivi altra kam lakto\lingui{DEFIS}
\artiklo{ablativo}%
\vorto{ablativo}\fako{gram.} Kazo di la deklino-sistemo Greka, Latina, e c., qua indikas ke substantivo esas la departo-punto, o la instrumento, di ago\lingui{DEFIRS}
\artiklo{ablegato}%
\vorto{ablegato}\fako{religio katolika} Delegito qua komisesis por portar la bireto a prelato jus nominita « kardinalo »\lingui{DEFI}
\artiklo{ableto}%
\vorto{ableto}\fako{zool.} Mikra fisho ek la familio « kiprinoidi », di qua la squami perlomatrea uzesas por facar la perli artificala\latina{alburnus ludicus}\lingui{EF}
\artiklo{ablucionar}%
\vorto{ablucionar}\fako{trans.}\senco{1}{\textit{(liturgio katolika)} Varsar sur la fingri di la sacerdoto vino ed aquo qui transfalas aden la kalico, e quan lu drinkas pose}\senco{2}{\textit{(medic.)} Lavar parto di la homo-korpo segun la preskripti medikala}\lingui{DEFIS}
\artiklo{abnegar}%
\vorto{abnegar}\fako{trans.} Renuncar omna profito, omna interesto propra, e mem plezuro, por konsakrar su ad altru od a verko sociala, religiala. Renuncar su ipsa : sua feliceso propra\lingui{DEFIS}
\artiklo{abolisar}%
\vorto{abolisar}\fako{trans.} \textit{(pri institucuro o kustumo : rejeso, privileji, sklaveso-rejimo)} Destruktar (li) en maniero tala ke li ne povos ri-existar\lingui{DEF}
\artiklo{abominar}%
\vorto{abominar}\fako{trans.} Hororar (ulo, ulu) pro lua despieso o maledikeso, o pro ke onu regardas lu kom maledikenda od odienda\lingui{EFIS}
\artiklo{abonar}%
\vorto{abonar}\fako{trans.} Kontratar forfete e kom valida dum quanto de tempo konvencionata, por recevar jurnalo, revuo, od asistar serio de spektakli en teatro\lingui{DFIRS}
\artiklo{abordar}%
\vorto{abordar}\fako{trans.} Irar sur ulo, o – per extenso di la senco – venar flanke a, bordo an bordo kun ulo (por irar sur olu)\lingui{EFIRS}
\artiklo{aborijena}%
\vorto{aborijena} Qua apartenas a populo qua habitas la lando de la origino (historiala)\lingui{EFIS}
\artiklo{abortar}%
\vorto{abortar}\fako{netrans.}\senco{1}{\textit{(aludante feto)} Mortar, od ocidesar, dum la jermeso-periodo}\senco{2}{\textit{(aludante muliero gravida)} Parturar, ante la instanto normala, filio qua mortis od ocidesis dum la jermeso-periodo}\lingui{DEFIS}
\artiklo{aboyar}%
\vorto{aboyar}\fako{netrans.} \textit{(aludante hundo o kelka animali altra ek la sama genero)} Kriar\lingui{FI}
\artiklo{abreviar}%
\vorto{abreviar}\fako{trans.} Igar (ulo) plu kurta spacale (verko literatura, historia, e c.) per uzar min multa vorti\lingui{DEFIS}
\artiklo{abrikoto}%
\vorto{abrikoto}\fako{bot.} Frukto grosa, flava e kelke reda, di qua la pulpo esas sukroza e saporoza – de arboro ek la genero « pruniero », familio de la « rozacei », qua devenis (onu konjektas lo) de Armenia\latina{prunus armeniaca}\lingui{DEFIRS}
\artiklo{abrogar}%
\vorto{abrogar}\fako{trans.} Deklarar « nevalida » to quo esas establisita, institucita. (Supresar lego per altra lego, dekreto per altra dekreto qua nevalidigas la antea)\lingui{DEFIS}
\artiklo{abrotano}%
\vorto{abrotano}\fako{bot.} Speco di artemizio di qua la folii, kande onu manuagas li, lasas sur la fingri odoro citronala\latina{artemisia abrotanum}
\artiklo{abrupta}%
\vorto{abrupta}\senco{1}{Di qua la pento esas ruptita bruske, ne-egale}\senco{2}{\textit{(metaf.)} \textit{(pri stilo)} Di qua la aluro esas neharmonioza e ne-egala}\lingui{DEFIS}
\artiklo{absenta}%
\vorto{absenta} Qua ne esas en la loko ube lu povas o devas esar\lingui{DEFIRS}
\artiklo{absinto}%
\vorto{absinto}\fako{bot.} Planto aromata tre bitra, ek la familio « kompozaji »\latina{artemisia absinthium}\lingui{DEFIRS}
\artiklo{absoluta}%
\vorto{absoluta} Qua ne submisesas a kondiciono, a restrikto\lingui{DEFIRS}
\artiklo{absolvar}%
\vorto{absolvar}\fako{trans.} Dispensar de la puniso, de la responso a qua lu su expozis pro ago pri qua lego preskriptas puniso\lingui{DEFIS}
\artiklo{absorbar}%
\vorto{absorbar}\fako{trans.} Igar penetrar, retenar en sua substanco, liquido, gaso\lingui{DEFIS}
\artiklo{abstenar}%
\vorto{abstenar}\fako{trans.}\senco{1}{Interdiktar a su ipsa agar (ulo)}\senco{2}{Interdiktar a su ipsa uzar, juar (ulo)}\lingui{DEFI}
\artiklo{abstersar}%
\vorto{abstersar}\fako{trans.) (medic.} Netigar de la puso, de la sanio qua adheras a la surfaco di plago, di ulcero\lingui{DEFIS}
\artiklo{abstinencar}%
\vorto{abstinencar}\fako{netrans.) (religio} Ne manjar karno en la dii quin indikas la Eklezio katolika\lingui{DEFIS}
\artiklo{abstraktar}%
\vorto{abstraktar}\fako{trans.}\senco{1}{Izolar per penso}\senco{2}{\textit{(filoz.)} Konsiderar parto, izolante lu ek la toto.}\lingui{DEFIRS}\subvorto{}{nombro abstraktita}{Qua ne indikas *grandoro reala, objekti mezurebla o kontebla}{}{}
\artiklo{abstruza}%
\vorto{abstruza} Desfacile komprenebla e tedanta\lingui{DEFIS}
\artiklo{absurda}%
\vorto{absurda} Kontrea ye la raciono komuna\lingui{DEFIRS}
\artiklo{abulio}%
\vorto{abulio}\fako{medic.} Morbo mentala quan karakterizas la savo di to quon onu devas agar e la nepovo agar lo\lingui{DEFIS}
\artiklo{abundar}%
\vorto{abundar}\fako{netrans.} Esar plu multa kam suficante\lingui{EFIS}
\artiklo{abutar}%
\vorto{abutar}\fako{netrans.} Arivar per sua extremajo. (Ex. : strado qua abutas an placo; nervi qui abutas an la cerebro)\lingui{EFIS}
\artiklo{abutmento}%
\vorto{abutmento}\fako{teknol.} Konstrukturo, masivo petra, destinita a suportar pulso. (Ex. : sur la rivi di fluvio, por rezistar la pulsi da la arki)\lingui{F}
\artiklo{acefala}%
\vorto{acefala}\fako{biol.} Sen-kapa. (Ex. : la ostri esas acefala)
\artiklo{acelerar}%
\vorto{acelerar}\fako{trans.} Igar plu rapida\lingui{EFIS}
\artiklo{acendar}%
\vorto{acendar}\fako{trans.} Igar lumifanta per igar (lu) produktar flami, mem nur mikra. \textit{(Anke metaf.)}\lingui{ISL}
\artiklo{acensar}%
\vorto{acensar}\fako{netrans., sur} Irar adsupre, hisar su adsupre\lingui{DEFIS}
\artiklo{acento}%
\vorto{acento}\fako{fonetiko} Plu-lautigo di la voco sur ta o ca silabo di vorto\lingui{DEFIRS}
\artiklo{aceptar}%
\vorto{aceptar}\senco{1}{Konsentar prenar, recevar to quon ulo ofras, donas, donacas, prizentas}\senco{2}{Recevar volunte (ulu) qua venas kom vizitanto}\lingui{DEFIRS}
\artiklo{acerba}%
\vorto{acerba} Pikanta ed atakema \textit{(metaf.)}\lingui{EFIS}
\artiklo{acero}%
\vorto{acero}\fako{bot.} Arboro di qua la ligno, veinoza, uzesas por la faco di mobli\latina{acer campestre}
\artiklo{acesar}%
\vorto{acesar}\fako{trans.} Povar arivar an loko, an persono\lingui{EFIS}
\artiklo{acesito}%
\vorto{acesito} Nomino quan onu grantas, en konkurso, a la personi qui proximeskis maxime a la premio\lingui{DEFIS}
\artiklo{acesora}%
\vorto{acesora} Qua adjuntesas a la kozo precipua\lingui{DEFIRS}
\artiklo{acetato}%
\vorto{acetato}\fako{kemio} CH₃CO₂H
\artiklo{acetileno}%
\vorto{acetileno}\fako{kemio} CH : CH
\artiklo{acetometro}%
\vorto{acetometro} Acidometro specala, uzata por mezurar la koncentreso-grado di vinagro
\artiklo{acetono}%
\vorto{acetono}\fako{kemio} CH₃CO.CH₃
\artiklo{-acho}%
\vorto{-acho} Sufixo pejorativa qua expresas qualeso tre infra, nuanco di desestimo, antipatio, repugneso
\artiklo{aciano}%
\vorto{aciano}\fako{bot.} Varietato di centaurei, kun floro blua, qua kreskas precipue en la frumento-agri\latina{centaurea cyanus}
\artiklo{acida}%
\vorto{acida} Di qua la saporo esas pikanta.\lingui{EFIRS}\subvorto{kemio}{acido}{Oxo-kompozajo qua redigas la tornasolo blua, e qua, kande onu deskompozas lua sali per baterio, fluas a la polo pozitiva}{}{}
\artiklo{acidentar}%
\vorto{acidentar}\fako{trans.}\senco{1}{Ruptar hazarde, ne-expektite e (generale) desfortunoze la iro normala di la kozi}\senco{2}{\textit{(filoz.)} Modifikar nedurive, efemere, la eso}\lingui{EFIRS}
\artiklo{acidometro}%
\vorto{acidometro} Instrumento per qua onu mezuras la koncentreso-grado di acido\lingui{EFIRS}
\artiklo{aciono}%
\vorto{aciono}\fako{financo}\senco{1}{Parto di la kapitalo suskriptenda, qua yurizas pri parto di la trezoro societala e di la profito di societo}\senco{2}{La akto qua reprezentas singla di ta parti}\lingui{DFIRS}
\artiklo{acipitro}%
\vorto{acipitro}\fako{zool.} Rapt-ucelo jornala, qua ne flugas en la alta strati di atmosfero\latina{astur palumbarius}
\artiklo{acizo}%
\vorto{acizo} Imposto di ula vari (nutrivi, e c.) ye la enireyo di urbo\lingui{DEFIRS}
\artiklo{adaptar}%
\vorto{adaptar}\fako{trans., ad} Unionar (kozo) ad altra kozo qua konvenas a lu\lingui{DEFIRS}
\artiklo{adenino}%
\vorto{adenino}\fako{kemio} C₅H₅N₅
\artiklo{adenito}%
\vorto{adenito}\fako{patol.} Inflamuro di la glandi o ganglioni limfala
\artiklo{adepto}%
\vorto{adepto} Persono quan onu iniciis pri ula arto o doktrino\lingui{DEFIRS}
\artiklo{adequata}%
\vorto{adequata}\fako{filoz.} Qua equivalas, esas adaptita, korespondas perfekte a la objekto (aludante penso, cienco, e c.)\lingui{DEFS}
\artiklo{adherar}%
\vorto{adherar}\fako{netrans., ad}\senco{1}{Quaze glutinesar an ulo}\senco{2}{\textit{(metaf.)} Ligar su komplete ad opiniono, a partiso}\lingui{DEFIS}
\artiklo{adiabata}%
\vorto{adiabata}\fako{fiziko} Quan kaloro povas permear
\artiklo{adiar}%
\vorto{adiar}\fako{trans.}\senco{1}{Pronuncar la formulo di politeso kande onu su separas (de ulu)}\senco{2}{Pronuncar la formulo di renunco od abandono}\lingui{DEFIS}
\artiklo{adicionar}%
\vorto{adicionar}\fako{matem.} Adjuntar nombro ad un o plura nombri di la sama naturo, por unionar li ad una\lingui{DEFIS}
\artiklo{adipocero}%
\vorto{adipocero}\fako{kemio} Nomo olima di tri substanci quin onu kredis inter-identa, ma qui interdiferas esencale : spermaceto, la grasajo di la kadavri, e kolesterino (C₂₇H₄₅OH)
\artiklo{adjektivo}%
\vorto{adjektivo}\fako{gram.} Vorto-speco qua expresas la eso-maniero\lingui{DEFIS}
\artiklo{adjudikar}%
\vorto{adjudikar}\fako{trans., ad}\senco{1}{\textit{(yuro-cienco)} Atribuar, per judicio, a la autoritato o persono qui esas kompetenta}\senco{2}{\textit{(lor vendo)} Atribuar a la persono qua ofras po lo la preco maxima}\lingui{DEFIS}
\artiklo{adjurar}%
\vorto{adjurar}\fako{trans., pri} Sumnar (ulu), ye la nomo di sakrajo ke lu agez ulo\lingui{EF}
\artiklo{adjutanto}%
\vorto{adjutanto}\fako{armeo} Oficiro qua servas generalo\lingui{DEFIRS}
\artiklo{administrar}%
\vorto{administrar}\fako{trans.} Direktar, surveyar la jero di la aferi di societo, di stato, di privato\lingui{DEFIRS}
\artiklo{admiralo}%
\vorto{admiralo}\fako{armeo navala} Oficiro di qua la grado esas maxima en la hierarkio di la navaro militala\lingui{DEFIRS}
\artiklo{admirar}%
\vorto{admirar}\fako{trans., pri, pro} Kontemplar astonate\lingui{EFIS}
\artiklo{admisar}%
\vorto{admisar}\fako{trans., aden} Aceptar aden loko ulu qua havas la qualesi postulata por lo\lingui{EFIS}
\artiklo{adolecar}%
\vorto{adolecar}\fako{netrans.} Esar en ta evo-periodo qua sequas puereso e preiras adulteso (t.e. de 15 til 20 ev-yari)\lingui{EFIS}
\artiklo{adoptar}%
\vorto{adoptar}\fako{trans.}\senco{1}{\textit{(ulu)} Prenar legale kom filio. \textit{(Anke metafore)}}\senco{2}{Igar sua la judiko-maniero o la ago-maniero di altru}\lingui{DEFIRS}
\artiklo{adorar}%
\vorto{adorar}\fako{trans.} Honorizar quale se lu esus deo\lingui{EFIS}
\artiklo{adosar}%
\vorto{adosar}\fako{trans., an ulo} Apogar per la dorso\lingui{FIS}
\artiklo{adraganto}%
\vorto{adraganto}\fako{bot.} Gumo qua defluas, forme di paceti vermatra, de tragakanto\latina{Astragalus}\lingui{DEFIRS}
\artiklo{adreso}%
\vorto{adreso}\senco{1}{Domicilo di persono o di organismo}\senco{2}{Surskriburo di letro-kuverto, e c., qua indikas ube habitas, ube lojas la destinario}\lingui{DEFRS}
\artiklo{adsorbar}%
\vorto{adsorbar}\fako{trans.} Koncentrar, sur sua surfaco, substanci dissolvita\lingui{DEFI}
\artiklo{adulta}%
\vorto{adulta} Qua atingis la fino di sua kresko. (Dicesas precipue pri la homo qua evas plu kam 20 yari)\lingui{EFIS}
\artiklo{adulterar}%
\vorto{adulterar}\fako{netrans., kun} Violacar sua promiso solena di fideleso sexuala a sua spozo\lingui{EFIS}
\artiklo{adurolo}%
\vorto{adurolo}\fako{kemio} Revelilo fotografala, ek klorhidoquinono
\artiklo{adventiva}%
\vorto{adventiva}\senco{1}{Qua devenas de ago acidenta}\subvorto{biol.}{planto adventiva}{Quan la homo ne semis}{2}{}
\artiklo{advento}%
\vorto{advento}\fako{religio katolika} La quar semani qui preiras la Kristo-nasko – tempo determinita da la Eklezio katolika, dum qua la religiani devas preparar su pri la arivo di Yesu-Kristo\lingui{DEFIS}
\artiklo{adverbo}%
\vorto{adverbo}\fako{gram.} Vorto-speco nevariiva qua modifikas la senco di adjektivo, di verbo, di adverbo\lingui{DEFIRS}
\artiklo{adversa}%
\vorto{adversa} Kontrea. \textit{(en lukto)} La persono qua opozesas ad altra\lingui{EFIS}
\artiklo{advokato}%
\vorto{advokato} La homo di qua la profesiono konsistas ye pledar por la personi qui procesas koram tribunalo\lingui{DEFIRS}
\artiklo{adyunto}%
\vorto{adyunto} Qua juntesas ad ulu kom lua helpanto\lingui{DEFIRS}
\artiklo{aerenkimo}%
\vorto{aerenkimo}\fako{biol.} Tisuo ek lakuni respirala, quan onu renkontras en la organi (precipue la radiki) qui kreskas en slamo
\artiklo{aero}%
\vorto{aero} Fluido gasa, diafana, ek oxo, azoto e kelka gasi skarsa, difuzita en omna regioni di Tero, utila por respiro e kombusto\lingui{DEFIS}
\artiklo{aerobia}%
\vorto{aerobia}\fako{biol.} Qua bezonas la oxo di aero por respirar\lingui{DEFIS}
\artiklo{aerolito}%
\vorto{aerolito}\fako{astron.} Maso minerala qua falas adsur la surfaco di Tero, de la regioni superega di atmosfero\lingui{DEFIRS}
\artiklo{aerometrio}%
\vorto{aerometrio} La parto di fiziko qua studias la denseso-grado di aero, di gasi\lingui{DEFIRS}
\artiklo{aerometro}%
\vorto{aerometro} Instrumento per qua onu povas mezurar la denseso-grado di aero, di gasi\lingui{DEFIRS}
\artiklo{aeroplano}%
\vorto{aeroplano}\fako{aviaco} Mashino qua povas sustenar su super-sule, en atmosfero, sen esar plu lejera kam ica, efike da la preso da vento sur la surfaci inklinita di lua quaza ali\lingui{DEFIRS}
\artiklo{aeroskopo}%
\vorto{aeroskopo} Instrumento per qua onu kolektas (por studiar li) la polvuni mikroskopala qui suspendesas en atmosfero\lingui{DEFIRS}
\artiklo{aerostato}%
\vorto{aerostato} Balonego di qua la susteneso en atmosfero obtenesas per la uzo di gaso plu lejera kam aero\lingui{DEFIRS}
\artiklo{afabla}%
\vorto{afabla} Qua aceptas sua kunhomi benigne, aminde\lingui{IS}
\artiklo{afazio}%
\vorto{afazio}\fako{patol.} Privaceso di la parol-fakultato, konjekteble pro lezuro cerebrala\lingui{DEFIS}
\artiklo{afecionar}%
\vorto{afecionar}\fako{trans.} Unionesar ad (ulu, ulo) per amo o prizo duranta\lingui{EFIS}
\artiklo{afektacar}%
\vorto{afektacar}\fako{trans.} Esar, agar, kondutar en maniero qua esas nur la semblo (di ulo)\lingui{DEFIRS}
\artiklo{afektar}%
\vorto{afektar}\fako{trans.}\senco{1}{\textit{(medic.)} \textit{(pri morbo)} Efikar a la homo-korpo}\senco{2}{\textit{(psikol.)} Efikar a la psiko}\lingui{DEFIRS}
\artiklo{afelio}%
\vorto{afelio}\fako{astron.} La punto di la orbito di planeto, ube la planeto distas maxime de Suno\lingui{DEFIRS}
\artiklo{afero}%
\vorto{afero}\senco{1}{To quon onu havas kom agenda (en maniero generala)}\senco{2}{(1) \textit{(politiko)} To quo havas kom objekto la interesti di la stato. (2) To quo havas kom objekto la interesti privata. (3) \textit{komerco} Vendo-kontrato. (4)\textit{(judicio)} To quo esas la objekto di debato judiciala}\lingui{DEFIR}
\artiklo{afidio}%
\vorto{afidio}\fako{zool.} Insekto tre mikra qua adheras a la folii ed a la rami di la planti, por sugar li\lingui{EFIS}
\artiklo{afina}%
\vorto{afina}\fako{ad} Karakterizata da la tendenco a proximigar su (ad) ed ad unionar su (kun), a kombinar su (kun)\lingui{DEFIS}
\artiklo{afirmar}%
\vorto{afirmar}\fako{trans.} Deklarar ke to o co esas vera\lingui{DEFIS}
\artiklo{afisho}%
\vorto{afisho} Anunco a publiko per manuskriburo od imprimuro quan onu adherigas sur pordo, sur muro\lingui{DEFIRS}
\artiklo{afixo}%
\vorto{afixo}\fako{gram.} Partikulo quan onu adjuntas a la vorti por modifikar lia senco\lingui{DEFIRS}
\artiklo{afliktar}%
\vorto{afliktar}\fako{trans.} Frapar dolorigante (lua) korpo o psiko\lingui{EFIS}
\artiklo{aforismo}%
\vorto{aforismo} Propoziciono pri ula cienco, arto, enuncita forme di sentenco\lingui{DEFIRS}
\artiklo{afrankar}%
\vorto{afrankar}\fako{trans.} Dispensar la destinario de la pago di la transporto-kusto, per agar ta pago lor la expedio\lingui{DEFIRS}
\artiklo{afrontar}%
\vorto{afrontar}\fako{trans.}\senco{1}{Opozar su, fronto an fronto}\senco{2}{\textit{(metaf.)} Abordar (danjero, tempesto) audacoze}\lingui{EFIS}
\artiklo{afto}%
\vorto{afto}\fako{patol.} Mikra ulcero blankatra, qua aparas e kreskas sur la membrani mukoza di la boko o di faringo\lingui{DEFIRS}
\artiklo{afusto}%
\vorto{afusto}\fako{artilerio} Suportilo di kanono\lingui{FIS}
\artiklo{afuzionar}%
\vorto{afuzionar}\fako{trans.} Igar liquido fluar strate sur korpo-parto malada\lingui{DEFIS}
\artiklo{agacar}%
\vorto{agacar}\fako{trans.} Produktar mikra iriteso nervala o psikala – exemple per la froto di kultelo pozita vertikale sur plado o sur vitroplako, e c.\lingui{F}
\artiklo{agama}%
\vorto{agama}\fako{biol.} Sen uniono intersexua
\artiklo{agapo}%
\vorto{agapo}\fako{religio} Repasto quan, origine, la Kristani agis komune\lingui{DEFIRS}
\artiklo{agar}%
\vorto{agar}\fako{trans.} Exekutar la gesti, la movi di la korpo o di parti di olu, la movi psikala, qui esas necesa\lingui{DFIS}
\artiklo{-agar}%
\vorto{-agar} Uzesas kun vorti di instrumento por formacar verbi qui signifikas « agar per… ». Exemple : \textit{krucagar} : agar per kruco (por mortigar ulu); \textit{klovagar} : agar per klovi (por fixigar ulo : planki, e c.); \textit{agrafagar}, \textit{butonagar}, e c.
\artiklo{agar-agaro}%
\vorto{agar-agaro} Gelatino vejetantala, obtenita de mar-algi
\artiklo{agariko}%
\vorto{agariko}\fako{bot.} Nomo di fungi diversa-sorta, manjebla o venenoza\latina{agaricus}\lingui{EFIRS}
\artiklo{agato}%
\vorto{agato}\fako{mineral.} Petro harda, varietato di quarco, qua divenas briloza per poliseso\lingui{DEFIRS}
\artiklo{agavo}%
\vorto{agavo}\fako{bot.} Genero de amarilidacei, qua devenas de Amerika, e konsistas ye cirkume cent speci Amerikana; olua sapo, sukroza e tre abundanta, uzesas da la Mexikiani por preparar drinkaji fermentacita\latina{agave}\lingui{DEFIRS}
\artiklo{agendo}%
\vorto{agendo} Registro, kayereto, di qua singla folio, qua indikas la dio-quantesmo di la monato, uzesas por notar sur lu to quon onu havas kom agenda\lingui{DEFIS}
\artiklo{agento}%
\vorto{agento} La persono qua esas komisita pri agor vice la komisinto\lingui{DEFIRS}
\artiklo{agitar}%
\vorto{agitar}\fako{trans.} Dis-movigar (liquido, la spiriti di publiko), sukusar omna-sinse\lingui{DEFIRS}
\artiklo{aglo}%
\vorto{aglo}\senco{1}{\textit{(zool.)} Rapt-ucelo jornala, di qua la regardo esas tre perceptiva; la flugo, rapida; la beko, tre forta, bazo-rekta, e kurvigita nur ye lua pinto}\senco{2}{To quo *weras aglo kom emblemo (standardi, blazono, orden-insigno)}\latina{aquila}\lingui{EFIS}
\artiklo{aglomerar}%
\vorto{aglomerar}\fako{trans.} Asemblar por formacar maso kompakta\lingui{DEFIRS}
\artiklo{aglosio}%
\vorto{aglosio}\fako{patol.} Absenteso di la lango
\artiklo{aglutinar}%
\vorto{aglutinar}\fako{trans.} Interglutinar plura kozi tale ke li divenas una\lingui{DEFIS}
\artiklo{agnelo}%
\vorto{agnelo} Muton-yuno, en la senco religiala e poeziala\lingui{FIR}
\artiklo{agnoskar}%
\vorto{agnoskar}\fako{trans.} Lasar enirar sua spirito kom vera, valida, yusta, legale konstitucita\lingui{L}
\artiklo{agoniar}%
\vorto{agoniar}\fako{netrans.} Esar la objekto di la lukto lasta da naturo kontre morto, en la korpo di vivanto\lingui{DEFIRS}
\artiklo{agorafobio}%
\vorto{agorafobio} Pavoro morba, da ula personi, kande li devas trans-irar placo, strado, ponto\lingui{EFIRS}
\artiklo{agosto}%
\vorto{agosto} La monato okesma di la yaro kalendarala
\artiklo{agrafo}%
\vorto{agrafo}\senco{1}{\textit{(arkitekt.)} Krampono fera qua inter-adherigas la quadri di muro}\senco{2}{\textit{(vesto)} Roketo metala qua uzesas por interjuntar la du bordi inter-opozanta di vesto, fixigita sur ica, ed eniranta ringo, fixigita sur ita}\lingui{DEFIS}
\artiklo{agreabla}%
\vorto{agreabla}\fako{ad} Qua donas plezuro (ad ulu)\lingui{EFIS}
\artiklo{agregar}%
\vorto{agregar}\fako{trans.}\senco{1}{Interunionar – por formacar ek li ensemblo – molekuli, partikuli materio}\senco{2}{Unionar ad ensemblo, a korpo}\lingui{DEFIRS}\subvorto{en la universitato di Francia}{agregito}{La persono qua, pro sucesir lor la konkurso pri lo, deklaresas kapabla pri docar en liceo, en fakultato. (Ta grado esas min alta kam doktoreso)}{3}{}
\artiklo{agrimonio}%
\vorto{agrimonio}\fako{bot.} Planto rozacea, kun mikra flori qui esas dispozita spike, olim uzita kom astriktivo\latina{agrimonia eupatoria}\lingui{DEFIS}
\artiklo{agro}%
\vorto{agro}\senco{1}{Tereno sen-tekta e plata, limitizita por uzeso partikulara}\senco{2}{Tereno por kultivado, ne cirkondata da muri}\lingui{DEFIRS}
\artiklo{agronomio}%
\vorto{agronomio} La cienco qua donas la reguli di agrokultivo\lingui{DEFIRS}
\artiklo{agronomo}%
\vorto{agronomo} Persono experta pri la cienco « agrokultivo »\lingui{DEFIRS}
\artiklo{agulo}%
\vorto{agulo}\senco{1}{Stangeto ek stalo polisita, di qua un extremajo esas pinta, e la altra, obtuzigita, ed un-trua, por recevar filo}\senco{2}{Mikra stango di qua un extremajo esas pinta. (Agulo \textit{di relvoyo} : parti di reli movebla, uzata por igar la treni pasar bifurke de ca a ta relvoyo)}\lingui{FIS}
\artiklo{ah!}%
\vorto{ah!} Interjeciono qua indikas ke onu quik expresos emoco forta qua kaptis la psiko (joyo, doloro, iraco, e c.)
\artiklo{ailanto}%
\vorto{ailanto}\fako{bot.} Planto ek la familio « nutacei »\latina{Ailanthus}
\artiklo{ajila}%
\vorto{ajila} Di qua la movi esas flexebla, rapida\lingui{EFIS}
\artiklo{ajio}%
\vorto{ajio}\fako{banko-komerco} Profito de la difero inter la valoro nominala e la valoro reala di la moneto-peci – inter la valoro nuna e la preco lor la pagoteso di la valor-paperi, di la vari di qui la kurso varias\lingui{DEFIRS}
\artiklo{ajiotar}%
\vorto{ajiotar}\fako{netrans., pri} Spekular pri la valor-paperi, pri la vari, di qui la kursi esas variiva\lingui{DEFIRS}
\artiklo{-ajo}%
\vorto{-ajo}\senco{1}{Sufixo qua indikas : I. Kozo facita ek ula materio, o qua posedas ta o ca karaktero (ex. : lanajo, kotonajo, bonajo, bitrajo)}\senco{2}{\textit{(per extenso di la senco)} Ago, parolo, procedo di… (amikalajo, pueralajo)}\senco{3}{\textit{(kun verbo transitiva)} Indikas la objekto pasiva di la ago (cakaze, lu equivalas « -ata ») (ex. : manjajo : karno, pano, legumi); vidajo : kozo vidata}\senco{4}{\textit{(kun verbo mixita)} -ajo = -(at)ajo (ex. : chanjajo : kozo chanjata.)}\senco{5}{\textit{(kun verbo netransitiva)} Indikas la subjekto, ed equivalas \textit{-(ant)ajo} (ex. : rezultajo : to quo rezultas; eventajo : to quo eventas)}
\artiklo{ajornar}%
\vorto{ajornar}\fako{trans., til} Rezolvar ke (ulo) agesos o facesos erste en ta o ca dio, vice nun\lingui{EFI}
\artiklo{ajuro}%
\vorto{ajuro}\senco{1}{\textit{(arkitekt.)} Singla ek la mikra aperturi di ornamentajo}\senco{2}{\textit{(stofo)} Quaza truo en broduro o dentelo}\lingui{DEF}
\artiklo{ajustar}%
\vorto{ajustar}\fako{trans.} Intermuntar ed inter-harmoniigar la parti di mashino, di mekanismo, por atingar la skopo determinita\lingui{EF}
\artiklo{ajuto}%
\vorto{ajuto} Tubo adaptita a la orifico di fonteno, di recipiento, e c. por modifikar (pri lua volumo o lua formo) la fluo di la liquido\lingui{EF}
\artiklo{akacio}%
\vorto{akacio}\fako{bot.} Genero de planti mimozea, di qui ula speci donas kachuo, glutino-gumo, e c.\latina{acacia vera}\lingui{DEFIRS}
\artiklo{akademio}%
\vorto{akademio} Societo literaturala, ciencala, artala, di qua la membri esas, generale, elektita e fixa-nombra, e qua ne docas\lingui{DEFIRS}
\artiklo{akantino}%
\vorto{akantino}\fako{biol.} Substanco organika, albuminoza, solvebla en sal-aquo – ye qua konsistas la skeleto di ula radiolarii
\artiklo{akanto}%
\vorto{akanto}\fako{bot.} Planto bikotiledona, di qua la folii esas dentoza\latina{acanthus}\lingui{DEFIRS}
\artiklo{akantonar}%
\vorto{akantonar}\fako{trans., en) (milit.} Instalar la armeani en la loki diversa di teritorio ube li devas vivar dum kelka tempo\lingui{DEFIS}
\artiklo{akaparar}%
\vorto{akaparar}\fako{trans.}\senco{1}{\textit{(ulo)} Rezervigar por su, komprar la tota quanto de ta o ca varo qua ofresas vende en merkato. – Prenar por su ipsa, detrimente di la kunhomi}\senco{2}{\textit{(ulu)} \textit{(metaf.)} Kaptar (lu) e ne lasar (lu) relatar kun la personi cetera}
\artiklo{akaro}%
\vorto{akaro}\fako{zool.} Familio de insekti di qui la tipo esas L. \textit{acarus}, kun korpo, generale, disko-forma e globatra, sen distingo tre preciza di abdomino de cefalotorako\latina{acarus}\lingui{EFIRSL}
\artiklo{aklamar}%
\vorto{aklamar}\fako{trans.} Aceptar per klami e krii de entuziasmo – agita da plura, multa, personi kune\lingui{DEFIS}
\artiklo{aklimatar}%
\vorto{aklimatar}\fako{trans.} Igar (lu) adaptar su a klimato altra kam ta di lua nasko-lando. – \textit{(anke metaf.)}\lingui{DEFIRS}
\artiklo{akneo}%
\vorto{akneo}\fako{patol.} Erupto di papuli\lingui{EFIS}
\artiklo{akoluto}%
\vorto{akoluto}\fako{eklezio katolika} Kleriko qua, en kirko, exekutis la taski grosiera\lingui{DEFIRS}
\artiklo{akomodar}%
\vorto{akomodar}\fako{trans.} Aranjar kozo tale ke lu konvenas por ulu\lingui{DEFIRS}
\artiklo{akompanar}%
\vorto{akompanar}\senco{1}{Juntar su (ad ulu) por irar adube lu iras}\senco{2}{\textit{(muziko)} Adjuntar, a la parto precipua quan kantas o pleas ulu, partin acesora por igar olta saliar}\lingui{EFIRS}
\artiklo{akonito}%
\vorto{akonito}\fako{bot.} Planto venenoza, ek la familio « renonkulacei »\latina{aconitum}\lingui{DEFIRS}
\artiklo{akonto}%
\vorto{akonto} Pago di parto de debo\lingui{DEFIS}
\artiklo{akordar}%
\vorto{akordar}\fako{trans.}\senco{1}{\textit{(muziko)} Existigar relato preciza inter la soni diversa di muzik-instrumento segun diapazono. Regular plura instrumenti muzikala segun la sama diapazonosono}\senco{2}{\textit{(gram.)} Igar (per flexioni) vorto relatar kun altra vorto di qua lu devas sequar la varii}\lingui{DEFIRS}
\artiklo{akordeono}%
\vorto{akordeono} Muzik-instrumento, kun klavaro, langeti metala, libere movebla, quin vibrigas suflilo\lingui{DEFIRS}
\artiklo{akoro}%
\vorto{akoro}\fako{bot.} Genero de aroidei, kun flori hermafrodita, qui enkovras la spadico\latina{acorus calamus}\lingui{L}
\artiklo{akra}%
\vorto{akra}\senco{1}{Di qua la saporo, la odoro, esas forta, quaze acida, ed iritas la respir-organi e la guturo}\senco{2}{\textit{(metaf.)} \textit{(aludante persono)} Di qua la karaktero esas bitre vundiva}\lingui{EFIS}
\artiklo{akreditar}%
\vorto{akreditar}\fako{trans., ye}\senco{1}{Igar (ulu) egardata kom fidinda (anke financale)}\senco{2}{Igar (ulu) agnoskesar kom ambasadisto ye guvernistaro stranjera}\lingui{DEFIRS}
\artiklo{akrido}%
\vorto{akrido}\fako{zool.} Insekto saltera, qua devastas la agri trupe, e quan onu konfundas ofte a la lokusto\latina{acridium}\lingui{I}
\artiklo{akrobato}%
\vorto{akrobato} La homo qua agas publike prodaji di forteso, di equilibro, di ajileso; gimnastikisto, klauno, herkulatro, equilibristo, jonglisto, dislokito, desostizito, serpentatro, e c.\lingui{DEFIRS}
\artiklo{akrochar}%
\vorto{akrochar}\fako{trans., an}\senco{1}{Pendigar de hoko}\senco{2}{Retenar per hoko o per ulo qua efikas quale hoko}\lingui{F}
\artiklo{akromatino}%
\vorto{akromatino}\fako{biol.} Substanco liquida di la celulo-nukleo qua esas kolorizebla per reaktivi, ma multe plu desfacile kam kromatino\lingui{DEFIRS}
\artiklo{akropolo}%
\vorto{akropolo}\fako{che la Greki antiqua} Citadelo, en la parto maxim alte-situita di la urbo\lingui{DEFIRS}
\artiklo{akrostiko}%
\vorto{akrostiko} Serio de versi di qui la literi unesma di singla lineo riproduktas, segun lia ordino, la literi ye qui konsistas la nomo di persono o di kozo\lingui{DEFIRS}
\artiklo{akroterio}%
\vorto{akroterio}\fako{arkitekt.} Soklo di statuo, di vazo, instalita sur frontono\lingui{DEFIRS}
\artiklo{aktinika}%
\vorto{aktinika} Dicesas pri ta radii per qui la suno-lumo efikas kemiale a substanci diversa
\artiklo{aktinio}%
\vorto{aktinio}\fako{zool.} Zoofito kun korpo pulpatra, quaze muskuloza, ornita diverse, adheranta a sulo per bazo cirkla qua efikas quale ventuzo\latina{actinia}\lingui{DEFIRS}
\artiklo{aktiva}%
\vorto{aktiva}\lingui{DEFIRS}\subvorto{gram.}{voco aktiva}{Qua expresas la ago}{1}{}\subvorto{milit.}{servado aktiva}{Servado da la armeanigeblo qua esas reale en armeo, sub la standardo}{2}{}\subvorto{financo}{debaji aktiva}{Pri qui onu esas la debario}{3}{}\subvorto{komerco}{aktivo}{To quon ulu expektas recevor kom debata a lu – to quon lu posedas, proprietas}{4}{}
\artiklo{akto}%
\vorto{akto}\senco{1}{\textit{(yuro)} Folio skriburoza qua konstatas fakto, konvenciono, obligeso}\senco{2}{\textit{(teatro)} Parto di teatrajo, separita de la cetera da intervalo}\lingui{DEFIRS}
\artiklo{aktoro}%
\vorto{aktoro}\fako{teatro, cinemo} La persono qua reprezentas un ek la roli, en teatrajo o cinemajo\lingui{DEFIRS}
\artiklo{aktuala}%
\vorto{aktuala}\senco{1}{\textit{(filoz.)} Qua realeskas – qua pasas de virtualeso ad ago}\senco{2}{Qua interesas en la periodo nuna}\lingui{DEFIS}
\artiklo{akular}%
\vorto{akular}\fako{trans., ad}\senco{1}{Igar butar kontre ulo, per la parto dopa}\senco{2}{Pulsar aden loko ube ne plus esas posibla retroirar. – \textit{(anke metaf.)}}\lingui{FIS}
\artiklo{akumular}%
\vorto{akumular}\fako{trans.} Pozar kozi, ici sur iti, plenigante, ecese, sen-limite. – \textit{(anke metaf.)}\lingui{DEFIRS}
\artiklo{akumulatoro}%
\vorto{akumulatoro}\fako{elektro} Aparato baterio, fonto di elektro, quan onu povas transportar ed uzar en sua lojeyo\lingui{DEFIRS}
\artiklo{akurata}%
\vorto{akurata} Qua arivas (persono) precize ye la kloko konvencionita\lingui{DIRS}
\artiklo{akushar}%
\vorto{akushar}\fako{trans.} Igar la feto homala ekirar- per ekpulso naturala, o per extrakto – de la organo en qua lu su developis\lingui{DEFR}
\artiklo{akustiko}%
\vorto{akustiko} La parto di fiziko qua koncernas la soni\lingui{DEFIRS}
\artiklo{akuta}%
\vorto{akuta}\senco{1}{Qua finas per pinto longa.}\subvorto{}{angulo akuta}{Min apertita kam angulo orta}{}{}\senco{2}{Qua igas sentar su quaze kom penetranta. \textit{(metaf.)}}\lingui{DEFIS}
\artiklo{akuzar}%
\vorto{akuzar}\fako{trans., pri}\senco{1}{Signalar (lu) kom kulpinta}\senco{2}{\textit{(yuro-cienco)} (a) Signalar (lu) kom kulpinta, koram tribunalo; (b) Sendar a tribunalo}\senco{3}{Signalar kom kriminalo la homo quan la judiciisto inquestala *inkulpas, t.e. judikas kom tre versemble kulpinta}\lingui{EFIS}
\artiklo{akuzativo}%
\vorto{akuzativo}\fako{gram.} \textit{(en la deklino-sistemo di la lingui Greka, Latina, Germana, Rusa, e c.)} La kazo di la komplemento direta\lingui{DEFIRS}
\artiklo{-ala}%
\vorto{-ala} Sufixo qua signifikas « qua relatas, koncernas, apartenas a, dependas de, konvenas por », ed anke « segun la maniero di, analoga a la verki da, digna de »
\artiklo{alabastro}%
\vorto{alabastro}\senco{1}{Varietato di gipso, diafane-blanka, mola e facile laborebla}\senco{2}{Varietato di marmoro, lakto-blanka, diafana, qua devenas precipue de la stalaktiti e stalagmiti}\lingui{DEFIRS}
\artiklo{alaktar}%
\vorto{alaktar}\fako{trans.) (aludante muliero} Nutrar (infanteto) per sua lakto\lingui{EFIS}
\artiklo{alambiko}%
\vorto{alambiko} Aparato por distilar\lingui{DEFIRS}
\artiklo{alarmar}%
\vorto{alarmar}\fako{trans.} Advokar a lia armi korpo de soldati\lingui{DEFIS}
\artiklo{alaudo}%
\vorto{alaudo}\fako{zool.} Ucelo di la familio « konirostri », qua vivas en la agri\latina{alauda}\lingui{FISL}
\artiklo{albatroso}%
\vorto{albatroso}\fako{zool.} Mar-ucelo di la mi-sfero australa di Tero, ek la familio « longa-pluma », di qua la beko esas forta e hokatra, e la pedi, sen-polexa\latina{diomedea}\lingui{DEFIRS}
\artiklo{albergo}%
\vorto{albergo} Hotelo, situita (generale) en la pre-urbi di urbo, an ruro-voyo od en vilajo\lingui{DFIS}
\artiklo{albina}%
\vorto{albina}\fako{patol.} Anomaleso korpala quan karakterizas, che la homo, la senkoloreso di la pelo, di la hari, di la pili, la paleso rozea di la iriso, la redatreso di la pupilo, e quan onu vidas anke che ula animali (kunikli furo-blanka) od elefanti pelo-blanka\lingui{DEFIRS}
\artiklo{albo}%
\vorto{albo}\fako{liturgio katolika} Tuniko ek telo blanka, kunligita super la reni per zono o kordono, quan la episkopi, la sacerdoti, e c. *weras sur sua sutano por celebrar la meso\lingui{DEFIS}
\artiklo{albumeno}%
\vorto{albumeno} Substanco qua envelopas la embriono en ula grani e nutras lu dum la jermifo. \textit{(anke pri la ovi)}\lingui{DEFIRS}
\artiklo{albumino}%
\vorto{albumino} Substanco quan onu renkontras en la sero di sango, en la albumeno di la ovo, e qua koagulas efike da kaloro\lingui{DEFIRS}
\artiklo{albuminoido}%
\vorto{albuminoido} To di qua la naturo esas olta di albumino; dicesas pri la substanci azotoza (nitroza) di animali o vejetanti\lingui{DEFIRS}
\artiklo{albuminurio}%
\vorto{albuminurio}\fako{patol.} Formo di maladeso dum qua urino esas albuminoza\lingui{DEFIRS}
\artiklo{albumo}%
\vorto{albumo} Mikra registro, di qua la folii esas destinita a recevar noturi, autografi, skisuri, versi, muzikaji, e c.\lingui{DEFIRS}
\artiklo{alburno}%
\vorto{alburno}\fako{bot.} En la arbori bikotiledona, la strato unesma de ligno, ordinare blanka, qua jacas nemediate sub la kortico\lingui{EFISL}
\artiklo{alcedo}%
\vorto{alcedo}\fako{zool.} Mar-hirundo\latina{alcedo}
\artiklo{alciono}%
\vorto{alciono}\fako{mitol.) (en la epoki antiqua} La alcedo, egardata kom ucelo qua facas sua nesto sur maro kalma, e (da la navigisti) kom ucelo bon-augura\lingui{DEFIS}
\artiklo{alegar}%
\vorto{alegar}\fako{trans.} Mencionar o citar kom autoritato, donar kom justifikivo, kom punto di sua argumento\lingui{DEFIS}
\artiklo{alegorio}%
\vorto{alegorio} Vorti figura qui prizentas a la spirito senco qua esas celita sub la senco literala\lingui{DEFIRS}
\artiklo{alejar}%
\vorto{alejar}\fako{trans.} Igar min pezoza, Igar min kargoza. \textit{(anke metaf.)}\lingui{EFIS}
\artiklo{aleno}%
\vorto{aleno} Puncilo quan uzas la shuifisti, la selifisti, por borar trui en la ledro-peci quin li volas inter-sutar\lingui{F}
\artiklo{aleo}%
\vorto{aleo}\senco{1}{Voyo quan facas bordumo de arbori, de arbusti folioza o de herbi longa-stipa}\senco{2}{\textit{(arkeol.)} Paseyo tektoza, unionuro ek du rangi inter-paralela de dolmeni}\lingui{DEFR}
\artiklo{alerta}%
\vorto{alerta} Ago-rapida\lingui{F}
\artiklo{alexandrino}%
\vorto{alexandrino} Verso de dek-e-du silabi, kun cezuro an-pos la sisesma\lingui{DEFIRS}
\artiklo{alezar}%
\vorto{alezar}\fako{trans.) (teknol.} Polisar, glatigar la parti kava di objekto borita o konkavigita\lingui{FIS}
\artiklo{alfabeto}%
\vorto{alfabeto} La serio de la literi qui reprezentas la soni ye qui konsistas la vorti\lingui{DEFIRS}
\artiklo{alfenido}%
\vorto{alfenido} Aloyuro de nikelo e de kupro\lingui{DEFIRS}
\artiklo{algebro}%
\vorto{algebro}\fako{matem.} Cienco di qua la skopo esas simpligar e solvar la problemi pri la *grandori, reprezentante li per literi e signi, e determinante la relati qui interligas la elementi diversa di la problemo, per formulo quan onu povas aplikar ad omna kazi simila; e qua, preterirante la skopo, studias la relati kom ipsajo, sen egardar lia origino, formi, kombinuri e transformesi\lingui{DEFIRS}
\artiklo{algo}%
\vorto{algo}\fako{bot.} Vejetanto kriptogama qua kreskas en aquo sen-sala o saloza en la loki tre humida\latina{alga}\lingui{DEFIRS}
\artiklo{algoritmo}%
\vorto{algoritmo}\fako{matem.} Tipo di noto-sistemo qua esas propra di genero partikulara di kalkulado\lingui{DEFIRS}
\artiklo{aliancar}%
\vorto{aliancar}\fako{trans., kun} Unionar personi, familii, per mariajo. (Ta vorto ne esas sinonimo di \textit{mariajar} : \textit{mariajo} relatas nur la gespozi; \textit{aliancar} relatas lia familii.)\lingui{DEFIS}
\artiklo{alibio}%
\vorto{alibio}\fako{yuro kriminala} La fakto esir ye la instanto kande krimino eventis, en loko altra\lingui{DEFRS}
\artiklo{alidado}%
\vorto{alidado} Linealo movebla por la relevo di mapi, qua havas, ye singla di sua extremaji, dioptri tra qui onu vizas la objekto di qua onu volas determinar la direciono\lingui{DEFRS}
\artiklo{alienacar}%
\vorto{alienacar}\fako{netrans.} Divenar quaze stranjera ye su ipsa, perdar sua rezono-fakultato\lingui{EFIS}
\artiklo{alienar}%
\vorto{alienar}\fako{trans.) (yuro-cienco} Igar la proprieto di ulu pasar ad altru, per vendo, cedo, e c.\lingui{DEFI}
\artiklo{aligatoro}%
\vorto{aligatoro}\fako{zool.} Sorto di krokodilo en la fluvii di Amerika\latina{alligator}\lingui{DEFIRSL}
\artiklo{alimentar}%
\vorto{alimentar}\fako{trans., ye} Provizar ye to quo esas necesa pri la duro di lua funciono, vivo, e c. (ex. : alimentar urbo, kaldiero, mashino)\lingui{DEFIS}
\artiklo{alineo}%
\vorto{alineo} Separo quan onu agas inter frazo e la preiranta, per igar lu spaco-komencar sur la lineo sequanta, dop mikra interspaco lakuna\lingui{DFI}
\artiklo{aliniar}%
\vorto{aliniar}\fako{trans.} Rangizar (soldati, planti) tale ke la ensemblo formacas quaza lineo
\artiklo{alio}%
\vorto{alio}\fako{bot.} Varietato di planti liliacea, di qua la bulbo odoras e saporas forte e pikante : shaloto, oniono, cibolo, ciboleto\latina{allium sativum}\lingui{FISL}
\artiklo{aliquanta}%
\vorto{aliquanta}\fako{matem.} Dicesas pri *grandoro qua ne kontenesas ulaquantafoye sen restajo, en toto. (Ex. : 2 esas parto aliquanta de 13)\lingui{EFIS}
\artiklo{aliquota}%
\vorto{aliquota}\fako{matem.} Dicesas pri *grandoro qua kontenesas ulaquantafoye sen restajo. (Ex. : 2, 3, 4… esas parti aliquota di 12)\lingui{EFIS}
\artiklo{aliteracar}%
\vorto{aliteracar}\fako{netrans. e trans.} Komencar (o finar) per la sama sono o soni\lingui{DEFIS}
\artiklo{alizeo}%
\vorto{alizeo}\fako{meteorol.} Vento qua suflas uniforme e konstante, de la komenco til la fino di la yaro, inter la tropiki\lingui{FIS}
\artiklo{alkado}%
\vorto{alkado}\fako{en Hispania} Judiciisto qua esforcas konciliar la interprocesonti qui intencas demandar la judicieso di sua afero koram la tribunalo di la yuro civila\lingui{DEFIS}
\artiklo{alkalio}%
\vorto{alkalio}\fako{kemio} Kompozajo kemiala qua efikas kom bazo kande lu kontaktas acido, en la sali\lingui{DEFIRS}
\artiklo{alkaloido}%
\vorto{alkaloido}\fako{kemio} Kompozajo organika di qua la propraji esas alkaliatra\lingui{DEFIRS}
\artiklo{alkemio}%
\vorto{alkemio} Savo arkana qua mixas, a procedi kemiala, la formuli di pseudo-« arto sakra » por deskovrar la moyeno fabrikar oro e panaceo\lingui{DEFIRS}
\artiklo{alko}%
\vorto{alko}\fako{zool.} Speco de cervi, qua habitas la regioni norda di Tero\latina{cervus alces}\lingui{DEISL}
\artiklo{alkoholo}%
\vorto{alkoholo} Liquido quan onu obtenas per la distilo di la drinkaji fermentacita\lingui{DEFIRS}
\artiklo{alkoholometrio}%
\vorto{alkoholometrio} Ta parto di fiziko qua traktas la procedi quin onu uzas por determinar la quanto de alkoholo quan kontenas la liquori alkoholoza\lingui{DEFIRS}
\artiklo{alkoholometro}%
\vorto{alkoholometro} Instrumento per qua onu determinas la quanto di alkoholo quan kontenas vino, brandio, e c.\lingui{DEFIRS}
\artiklo{alkovo}%
\vorto{alkovo} Kavajo en dormo-chambro por lokizar ibe un o plura liti\lingui{DEFIRS}
\artiklo{almanako}%
\vorto{almanako} Kalendaro qua kontenas anke indiki astronomiala, konjekti e predici pri la vetero dum la tota yaro qua komencas\lingui{DEFIRS}
\artiklo{almeo}%
\vorto{almeo} Dansistino e kantistino, che la Orientani\lingui{DEFIS}
\artiklo{almonar}%
\vorto{almonar}\fako{netrans., ad} Donacar karitatale a la povri\lingui{DEF}
\artiklo{alno}%
\vorto{alno}\fako{bot.} Arboro ek la familio « betulacei », qua kreskas an la rivo di la aquo-flui, lagi, lageti\latina{alnus}\lingui{FIL}
\artiklo{alo}%
\vorto{alo}\senco{1}{Apendico quan havas, an singla flanko di sua korpo, la animali qui povas flugar, e quan li desfaldas ed agitas por sustenar su en atmosfero}\senco{2}{\textit{(pri vento-muelilo)} La granda kadri telizita qui, per moveso da vento, igas la muelo rotacar}\senco{3}{Irgo qua su extensas o desfaldas ye singla di la du flanki di kozo : (a) alo di armeo : la parto sinistra e la parto dextra di la batalio-korpo; (b) alo di edifico, di aeroplano, e c.}\lingui{EFIS}
\artiklo{alodio}%
\vorto{alodio}\fako{feudismo} Domeno qua, donacite kom proprietajo di la donacario, esis heredebla e nedependanta, ecepte pri la homajo di vasaleso\lingui{DEFIRS}
\artiklo{alonge}%
\vorto{alonge} Paralele ye la bordo, ye la rivo, ye la litoro di…, ma kelka-diste (de lu)\lingui{EFI}
\artiklo{alonjo}%
\vorto{alonjo}\senco{1}{\textit{(tekn.)} Tubo quan onu adjuntas ad altra por igar lu plu longa o por igar lu komunikar kun altra tubo o vazo; anke : prolonguro di tablo}\senco{2}{\textit{(komerco)} \textit{(pri trato)} Papero-peco quan onu glutinas an lu kande, sur la trato ipsa, ne plus esas sat multa spaco pri la plusa indosi}\lingui{DEF}
\artiklo{aloo}%
\vorto{aloo}\fako{bot.} Planto dika-folia, liliacea, di qua la folii kontenas suko bitra\latina{aloe socotrina}\lingui{DEFIRSL}
\artiklo{alopatio}%
\vorto{alopatio}\fako{medic.} La medicino qua kuracas per la inversaji (kontraste kun homeopatio), konforme a la formulo Latina \textit{contraria contrariis curantur}\lingui{DEFIRS}
\artiklo{alopecio}%
\vorto{alopecio}\fako{patol.} Desaparo, tota o parta, di la hari, di la pili, e c. (che homo od animalo)\lingui{DEFIS}
\artiklo{aloso}%
\vorto{aloso}\fako{zool.} Fisho ek la familio « Kupi », di qua la karno prizesas multe\latina{alosa vulgaris}\lingui{DEFSL}
\artiklo{alotropa}%
\vorto{alotropa}\fako{kemio} Dicesas pri la korpi simpla qui, per altrigar sua stando, aquiras propraji anke altra\lingui{DEFIRS}
\artiklo{aloyar}%
\vorto{aloyar}\fako{trans.} Kombinar, mixar metali (per fuzo)\lingui{EF}
\artiklo{alpako}%
\vorto{alpako}\fako{zool.} Ruminero ek la genero « lamai », di qua la lan-felo esas silkatra\latina{auchenia pacos}
\artiklo{alpari}%
\vorto{« alpari »}\senco{1}{\textit{(komerco)} Inter-egalesi di la kambio-kursi monetala pri du merkati financala}\senco{2}{\textit{(financo)} Interegaleso di la kapitalo nominala di aciono od obligaciono, e la kurso aktuala}\senco{3}{\textit{(laboro)} Employato qua, po sua laboro, esas lojata e repastas gratuite, che sua patrono, ma ne recevas, de ica, salariopekunio}
\artiklo{alpo}%
\vorto{alpo} Pastureyo ye altitudo tre granda sur la monti Alpa\lingui{DEFIRS}
\artiklo{alta}%
\vorto{alta}\senco{1}{Qua preter-iras la nivelo ordinara}\senco{2}{Qua preter-iras la nivelo di altra objekto}\senco{3}{Superiora ye la kozi sama-speca}\senco{4}{Qua havas (ta o ca dimensiono) de lua pedo til lua somito}\lingui{EFIS}
\artiklo{altaro}%
\vorto{altaro}\fako{religio}\senco{1}{Tablo gazona, petra, e c. sur qua onu ofris olim sakrifikaji a la dei}\senco{2}{\textit{(che la kriatani)} Konstrukturo petra, marmora, tablo-forma, sur qua onu celebras la meso}\lingui{DEIRS}
\artiklo{alteo}%
\vorto{alteo}\fako{bot.} Planto mucilajatra, ek la familio « malvacei »\latina{althaea}\lingui{DEFIRSL}
\artiklo{alterar}%
\vorto{alterar}\fako{trans.} Modifikar pri lua naturo. Chanjar la naturo e qualeso di kozo (generale a min bona)\lingui{DEFIS}
\artiklo{alternar}%
\vorto{alternar}\fako{netrans.} Venar, ica pos ita, singla ye sua foyo\lingui{DEFIS}
\artiklo{alternativa}%
\vorto{alternativa} Dicesas pri la elektro-flui qui parkuras cirkuito, lore ta-direcione, lore ca-direcione\lingui{DEFIS}
\artiklo{alternativo}%
\vorto{alternativo} Obligeso optar inter du selektebli\lingui{DEFIS}
\artiklo{alternatoro}%
\vorto{alternatoro} Elektro-mashino qua produktas elektrofluo alternativa\lingui{DEFIS}
\artiklo{altitudo}%
\vorto{altitudo} Alteso di loko, kompare a la nivelo di maro\lingui{EFIS}
\artiklo{alto}%
\vorto{alto}\fako{muziko} Instrumento muzikala, kun kordi ed arketo, meza inter violino e violoncelo\lingui{DEFIRS}
\artiklo{altra}%
\vorto{altra} Qua ne esas sama kam ulu, ulo, qua diferas de lu, qua esas diversa\lingui{DEFIRS}
\artiklo{aludar}%
\vorto{aludar}\fako{trans.} Uzar vorto, frazo, qua donas ideo pri persono, pri kozo, pri eventajo\lingui{EFIS}
\artiklo{alumeto}%
\vorto{alumeto} Stipeto de ligno, o de altra substanco, inflamebla, uzata por acendar\lingui{F}
\artiklo{aluminio}%
\vorto{aluminio}\fako{kemio} Korpo simpla, metala, forjebla e duktila, tenaca, poke densa, ofte uzata aloye kun kupro (por facar juveli, vazi, e c.)\simbolo{Al}\lingui{DEFIRS}
\artiklo{alumino}%
\vorto{alumino}\fako{kemio} Oxido di aluminio, uzata por facar porcelano. Al₂O₃\lingui{EFIRS}
\artiklo{aluno}%
\vorto{aluno}\fako{kemio} SO₄M₂, (SO₄)₃, M'₂, 24 H₂O. (M : natro, o kalio, amonio, rubidio, cezio, arjento; M' = aluminio, fero, magnezo, kromo o galio)\lingui{DEFIS}
\artiklo{aluro}%
\vorto{aluro} La maniero marchar (pazope, trote, galope, e c.)\lingui{DFR}
\artiklo{aluviono}%
\vorto{aluviono} Sedimento quan lasas la aqui kande li retroiras\lingui{DEF}
\artiklo{alveolo}%
\vorto{alveolo}\senco{1}{Celulo vaxa quan facas la abeli}\senco{2}{Kavajo di la maxil-osto en qua singla dento esas inkastrita}\lingui{DEFIS}
\artiklo{amalgamo}%
\vorto{amalgamo}\senco{1}{Aloyuro de merkurio kun metalo}\senco{2}{\textit{(metaf.)} Mixuro de elementi heterogena}\lingui{DEFIRS}
\artiklo{amansar}%
\vorto{amansar}\fako{trans.}\senco{1}{\textit{(ulo)} Igar min sovaja la animalo}\senco{2}{\textit{(ulu)} \textit{(metaf.)} Igar (persono) plu sociema}\lingui{IS}
\artiklo{amar}%
\vorto{amar}\fako{trans.} Ligesar, per sua kordio, ad (ulu)\lingui{EFIS}
\artiklo{amaranto}%
\vorto{amaranto}\fako{bot.} Planto autunala, kun flori reda bele, purpuree, veluratre\latina{amarantus}\lingui{DEFIRS}
\artiklo{amaro}%
\vorto{amaro}\fako{navig.} Kordego uzata por retenar navo, fixigar objekto en navo, e c.\lingui{FIS}
\artiklo{amaso}%
\vorto{amaso} To quo esas akumulita tale ke lu formacas maso, granda-quanta, per adjunti intersequanta\lingui{DEFIR}
\artiklo{amatoro}%
\vorto{amatoro} La persono qua prizas e kultivas, por sua plezuro propra, ta o ca arto, cienco, e c.\lingui{DEFIRS}
\artiklo{amauroso}%
\vorto{amauroso}\fako{patol.} Deskresko, o desaparo kompleta, di la vido-fakultato, sen kauzo detektebla en la vid-organaro\lingui{DEFIS}
\artiklo{amazono}%
\vorto{amazono}\senco{1}{\textit{(mitol.)} Muliero qua apartenis a tribuo militema qua ne aceptis viri kun su}\senco{2}{Muliero qua kustumas kavalkar}\lingui{DEFIRS}
\artiklo{ambajo}%
\vorto{ambajo} Singla ek la perifrazi qui cirkondas la pens-enunco\lingui{EFIS}
\artiklo{ambasado}%
\vorto{ambasado} Misiono permananta, establisita da guvernistaro apud guvernistaro stranjera, por igar plu facilan lia inter-relati\lingui{DEFIS}
\artiklo{ambiciar}%
\vorto{ambiciar}\fako{trans.}\senco{1}{Dezirar pasionoze la honori, la granda ofici, e solicitar li}\senco{2}{Solicitar to quon onu egardas kom honorizivan realigar, exekutar}\lingui{DEFIRS}
\artiklo{ambidextra}%
\vorto{ambidextra} Qua uzas egale-bone sua manui (dextra e sinistra)\lingui{DEFIS}
\artiklo{ambigua}%
\vorto{ambigua} Qua havas plura senci inter qui la spirito esas hezitanta\lingui{DEFIS}
\artiklo{amblar}%
\vorto{amblar}\fako{netrans.} Aluro di quadrupedo qua marchas, trotas, levante alterne sua gambi di la sama latero\lingui{EFIS}
\artiklo{ambliopa}%
\vorto{ambliopa}\fako{patol.} Feblesko di la vido-fakultato\lingui{EFIRS}
\artiklo{ambono}%
\vorto{ambono}\fako{religio} Tribuno petra, marmora, en la koreyo di la baziliki, de ube onu lektis laute la epistolo o la evangelio\lingui{DEFIRS}
\artiklo{amboso}%
\vorto{amboso} Maso de fero stalo-parkovrita, sur qua onu forjas la metali kolda o varmigita ad-rede od ad-inkandece\lingui{D}
\artiklo{ambro}%
\vorto{ambro} Substanco vaxatra, mosk-odoranta, retrojetata da maro\lingui{DEFIRS}
\artiklo{ambrozio}%
\vorto{ambrozio}\fako{mitol.} Substanco de quo delektesis la nemortivi. \textit{(Anke metaf.)}\lingui{DEFIRS}
\artiklo{ambulanco}%
\vorto{ambulanco} Establisuro provizora, destinita a furnisar flego\lingui{DEFIRS}
\artiklo{amebo}%
\vorto{amebo} Protozoario mikroskopala, qua vivas en la aquo sen-sala o saloza, ed havas pseudopodi por movar su\lingui{DEFIRS}
\artiklo{amebocito}%
\vorto{amebocito}\fako{biol.} Celulo di la organismi komplexa, animali o planti, qua su movas quale la amebi\lingui{DEFIS}
\artiklo{ameboida}%
\vorto{ameboida} Dicesas pri la movi analoga ad olti di la amebi\lingui{DEFIS}
\artiklo{amen}%
\vorto{« amen »} « Tale eventez ». Formulo per qua finas la pregi Latina\lingui{DEFIRS}
\artiklo{amendo}%
\vorto{amendo} To per quo (ordinare : pekunio) onu reparas sua kulpo\lingui{EFIS}
\artiklo{ametisto}%
\vorto{ametisto} Lapido violea\lingui{DEFIRS}
\artiklo{amfiastero}%
\vorto{amfiastero}\fako{biol.} Radiaco di la citoplasmo qua, dum mitoso, eventas cirkum la du poli, ecepte en la regiono di la spindelo\lingui{DEFIS}
\artiklo{amfibia}%
\vorto{amfibia} Qua vivas en aquo o sur tero\lingui{DEFIRS}
\artiklo{amfibola}%
\vorto{amfibola}\fako{medic.} Qualeso di febro di qua la intenseso-grado esas tre variiva\lingui{DEFIS}
\artiklo{amfigenezo}%
\vorto{amfigenezo}\fako{biol.} Vario nedeterminita, produktita (segun la dicipuli di Darwin) da la medio qua, pro modifikar su, produktas chanji celulala qui ulakaze esas tre poka\lingui{DEFIS}
\artiklo{amfimixio}%
\vorto{amfimixio}\fako{biol.} Riprodukto che qua la nova ento havas kom departo-punto celulo produktita da la fuzo di du celuli ne-kompleta (riprodukto inter-sexua) o di lia parti esencala (konjugo)\lingui{DEFIS}
\artiklo{amfioxo}%
\vorto{amfioxo} Mikra animalo marala, fishoforma, qua reprezentas la grado unesma di la skalo de la animali vertebroza, e quan onu renkontras sur litori sabloza\latina{amphioxus lanceolatus}\lingui{EFL}
\artiklo{amfipirenino}%
\vorto{amfipirenino}\fako{biol.} Substanco akromata de qua konsistas la membrano di la tisuo celulara\lingui{DEFIS}
\artiklo{amfiteatro}%
\vorto{amfiteatro}\senco{1}{Edifico cirklo-kurva, provizita ye gradi superpozita, ube eventis kombati di gladiatori}\senco{2}{Parto di teatro moderna, provizita per gradi superpozita}\senco{3}{Chambrego provizita ye gradi, ube profesoro diskursas doce}\lingui{DEFIRS}
\artiklo{amforo}%
\vorto{amforo}\fako{che la antiqui} Vazo de terakoto, destinita kontenor oleo, vino, e c.\lingui{DEFIRS}
\artiklo{amianto}%
\vorto{amianto} Varietato blanka di asbesto (SiO₃)₂n, di qua la filamenti esas brilanta, nekombustebla\lingui{DEFIRS}
\artiklo{amido}%
\vorto{amido}\fako{kemio} Nomo generala di la kompozaji qui derivas de amoniako o de amino primara o sekundara per la substituco di la radiki acida ad hidro\lingui{DEFIS}
\artiklo{amidolo}%
\vorto{amidolo}\fako{kemio} Diamidofenolo L.2.4., uzata kom revelilo fotografala\lingui{DEFIS}
\artiklo{amiko}%
\vorto{amiko} Ta persono qua unionesas ad altra persono per afeciono altra kam olta qua devenas de la parenteso o de la atrakto sexuala\lingui{EFIS}
\artiklo{amilo}%
\vorto{amilo}\fako{kemio} (CH₃) CH. CH₂CH₂NH₂
\artiklo{amiloplasto}%
\vorto{amiloplasto}\fako{biol.} Nomo di la granuli diferenciita, che protoplasmo, qui esas sen-kolora ed en qui naskas amilo\lingui{DEFIS}
\artiklo{amino}%
\vorto{amino} CH₃NH₂ o (CH₃)₃ NH o C₆H₅N (CH₃)₂
\artiklo{amitoso}%
\vorto{amitoso}\fako{biol.} Su-divido nemediata di la celuli, sen su-divido di la nukleo-peloto. (Antonimo : « mitoso »)\lingui{DEFIS}
\artiklo{amnestiar}%
\vorto{amnestiar}\fako{trans.} Efacar, kom suvereno, tota kategorio de krimini, de delikti, e c.\lingui{DEFIRS}
\artiklo{amniono}%
\vorto{amniono}\fako{anat.} La membrano maxim interna ek ti qui tegas la feto\lingui{DEFIS}
\artiklo{amnioto}%
\vorto{amnioto}\fako{anat.} Liquido, gelatinatra, qua, en planti, cirkondas la embriono, e divenas videbla erste pos la fekundigo\lingui{DEFIS}
\artiklo{amoniako}%
\vorto{amoniako}\fako{kemio} Amonio-salo\lingui{DEFIRS}
\artiklo{amonio}%
\vorto{amonio} NH₄\lingui{DEFIRS}
\artiklo{amonito}%
\vorto{amonito}\fako{paleont.} Konko fosila\lingui{DEFIRS}
\artiklo{amono}%
\vorto{amono} NH₃\lingui{DEFIRS}
\artiklo{amorar}%
\vorto{amorar}\fako{trans.} Afecionar pasionoze persono di la altra sexuo, (Ta vorto expresas specale la amo di la viri e di la mulieri, ne kom relato sexuala e pure fiziologiala, ma kom sentimento specala)\lingui{DEFIRS}
\artiklo{amorcar}%
\vorto{amorcar}\fako{trans.} Preparar por lua funciono. (Ex. : amorcar pumpilo)\lingui{F}
\artiklo{amorfa}%
\vorto{amorfa} Sen formo fixa\lingui{DEFIRS}
\artiklo{amortisar}%
\vorto{amortisar}\fako{trans.}\senco{1}{Igar ulo min violentoza; febligar}\senco{2}{Extingar debajo per pagar la kapitalo}\senco{3}{Igar plu mola}\lingui{DEFIRS}
\artiklo{ampelopso}%
\vorto{ampelopso}\fako{bot.} Speco de vito ornamentala, qua klimas sur muri\latina{ampelopsis}\lingui{FSL}
\artiklo{ampère}%
\vorto{« ampère »}\fako{elektro} Intenseso-grado di elektro-fluo qua povas deskompozar 0,001118 gramo di arjento singla-sekunde, per elektrolizo, ek aquo-dissolvuro de arjento-nitrato\lingui{DEFIRS}
\artiklo{ampèremetro}%
\vorto{ampèremetro}\fako{elektro} Instrumento por mezurar la intenseso-grado di elektro-fluo, e di qua la cifro-plako koncernas le « ampère »\lingui{DEFIRS}
\artiklo{ampla}%
\vorto{ampla} Vorto qua expresas relato inter du kozi de speci diversa, maxim-multa-kaze la relato di kontenanto e kontenajo. « \textit{Amplan} » qualifikesas : \textit{mantelo} qua fitas tre laxe olua *weranto; \textit{chambro} od \textit{agro} ube granda asemblitaro esos prezenta : \textit{provizuro} qua satisfacas la demandi. Stato qua havas granda mar-komerco mustas havar navaro « ampla »\lingui{EFIS}
\artiklo{amplitudo}%
\vorto{amplitudo} Grandeso-grado di arko\lingui{DEFIRS}
\artiklo{ampulo}%
\vorto{ampulo}\senco{1}{Mikra fialo, pasable sfera, qua finas pinte o quan onu klozas per weldo-lampo pos enduktir en olu liquido}\senco{2}{Vitra vazo qua kontenas la fili di elektro-lampo}\senco{3}{Veziketo produktata da la su-akumulo di seraji sub la epidermo}\lingui{EFIS}
\artiklo{amputar}%
\vorto{amputar}\fako{trans., de} De-sekar per tranchilo, per skalpelo (membro, parto de membro, organo, e c.) – \textit{(anke metaf.)}\lingui{DEFIRS}
\artiklo{amuleto}%
\vorto{amuleto} Objekto a qua onu atribuas supersticoze kuraciveso\lingui{DEFIRS}
\artiklo{amuzar}%
\vorto{amuzar}\fako{trans.} Okupar per kozi qui disipas tempo\lingui{DEF}
\artiklo{-an-}%
\vorto{-an-} Sufixo qua indikas individuo qua apartenas a klaso (urbo, lando, ensemblo) – homo qua esas « membro di »; plu precize : qua apartenas ad ula domeno (lando, socio, e c.). Konseque, lu ne uzesas nur pri homi, ma anke pri kozi
\artiklo{an}%
\vorto{an} Prepoziciono qua indikas relato di kontigueso o di apogo, tale ke la kozo kontaktas o preske\lingui{D}
\artiklo{anabaptismo}%
\vorto{anabaptismo}\fako{religio kristana} Herezio qua konsistas ye negar la efiko di la bapto kande olca eventas ante la racion-evo, e ye baptar itere la adulti\lingui{DEFIRS}
\artiklo{anabaptisto}%
\vorto{anabaptisto} Adheranto di anabaptismo\lingui{DEFIRS}
\artiklo{anabioso}%
\vorto{anabioso}\fako{biol.} Itero di la vivo videbla, pos longa periodo de vivo latenta\lingui{DEFIS}
\artiklo{anado}%
\vorto{anado}\fako{zool.} Ucelo palmipeda, di la familio « lamelorostri », qua vivas o ne-amansite od amansite\latina{anas, anatis}\lingui{FISL}
\artiklo{anaerobia}%
\vorto{anaerobia}\fako{biol.) (pri la organismi mikroskopala} Qua povas vivar en medio sen-oxa e sen-aera\lingui{DEFIS}
\artiklo{anafazo}%
\vorto{anafazo}\fako{biol.} Fazo di la konstituco di la nuklei filio lor la divido mitosala di la nuklei, divido per qua finas mitoso\lingui{DEFIS}
\artiklo{anagalo}%
\vorto{anagalo}\fako{bot.} Planto singla-yara, ek la familio « primulacei »\latina{anagallis}\lingui{DEFIS}
\artiklo{anagramo}%
\vorto{anagramo} Vorto formacita ek la transpozo di la literi di altra vorto\lingui{DEFIS}
\artiklo{anakoluto}%
\vorto{anakoluto} Elipso en qua onu omisas, en frazo, la korelativo gramatikala di vorto expresata. (Ex. : « Qua servas bone sua patrio / Ne bezonas ancestri » (Voltaire). La omiso di \textit{ta} avan \textit{qua} esas anakoluto)\lingui{DEFIRS}
\artiklo{anakoreto}%
\vorto{anakoreto} Religiano qua vivas sole\lingui{DEFIRS}
\artiklo{anakronismo}%
\vorto{anakronismo} Eroro qua konsistas ye konsiderar ke la fakto eventis ante sua dato eventala o (per extenso di la senco), ye dato altra kam la lua\lingui{DEFIRS}
\artiklo{analgeziar}%
\vorto{analgeziar}\fako{trans.} Produktar anestezio partala di la pelo o di la mukozi\lingui{EFIS}
\artiklo{analitika}%
\vorto{analitika}\senco{1}{Qua procedas per analizado}\lingui{DEFIRS}\subvorto{}{geometrio analitika}{Apliko di algebro a geometrio, qua, retroduktante la nociono « figuro » ad olta di « situeso », ed olta di « situeso » ad olta di « dimensiono », reprezentas la loki geometriala per equacioni, o deduktas, de la analizo di la equacioni, la karakterizivi di la figuri}{2}{}\subvorto{}{linguo analitika}{Qua tendencas deskompozar la moyeni per qui onu expresas la pensaji}{3}{}
\artiklo{analizar}%
\vorto{analizar}\fako{trans.} Dividar ensemblo ad olua parti\lingui{DEFIRS}
\artiklo{analo}%
\vorto{analo} Singla de la naracuri pri la eventaji di la yaro koncernata\lingui{DEFIS}
\artiklo{analoga}%
\vorto{analoga} Qua havas kelka karakteri komuna kun altra kozo\lingui{DEFIRS}
\artiklo{anamorfoso}%
\vorto{anamorfoso}\fako{matem.} Fenomeno fizikala qua eventas kande la desegnuro deformita de objekto, pozate perpendikle ye la axo di spegulo cilindra o kona, donas, per reflekto, imajo de la objekto sen deformo\lingui{DEFIRS}
\artiklo{ananaso}%
\vorto{ananaso}\fako{bot.} Frukto saporoza e parfumoza, di planto dornoza, ek la familio « bromeliacei », di qua la formo esas ovatra\latina{Ananasa sativa}\lingui{DFIRS}
\artiklo{anapesto}%
\vorto{anapesto} Verso en qua la silabi alternas yene : 2 kurta, 1 longa\lingui{DEFIRS}
\artiklo{anarkio}%
\vorto{anarkio} Stando sociala qua karakterizesas da la ne-existo di guvernanto\lingui{DEFIRS}
\artiklo{anastigmata}%
\vorto{anastigmata}\fako{optiko} Vitro, od okulo, sen astigmateso\lingui{DEFIRS}
\artiklo{anastomoso}%
\vorto{anastomoso}\fako{anat.} Junturo di du vaskuli\lingui{DEFIRS}
\artiklo{anatemar}%
\vorto{anatemar}\fako{trans.) (religio kristana} Frapar per la malediko per qua la Eklezio katolika forpulsas ulu de sua medio\lingui{DEFIRS}
\artiklo{anatomio}%
\vorto{anatomio} Cienco qua traktas la strukturo di la organi di nia korpo per dissekar li\lingui{DEFIRS}
\artiklo{ancestro}%
\vorto{ancestro} Un persono ek la serio de patri, avi, preavi, e c., retroire til epoko nedefinita\lingui{EF}
\artiklo{anchovo}%
\vorto{anchovo}\fako{zool.} Mikra maro-fisho qua, konservate en salo, manjesas kom acesorajo repastala\latina{engraulis encrasicholus}\lingui{DEFIRS}
\artiklo{anciena}%
\vorto{anciena} Qua esis tala ja multa tempo ante la instanto konsiderata\lingui{DEFIS}
\artiklo{andrinoplo}%
\vorto{andrinoplo} Sorto di stofo kotona, tintita rede\lingui{FIS}
\artiklo{androceo}%
\vorto{androceo}\fako{bot.} Ensemblo de la stamini, de la organi maskula di floro\lingui{DEF}
\artiklo{androdioika}%
\vorto{androdioika}\fako{bot.} Qualeso di planto che qua la flori esas : ici, maskulo – iti, hermafrodita, la du sorti esante sur stipiti diversa\lingui{DEFIS}
\artiklo{androgina}%
\vorto{androgina}\fako{bot.} Qua posedas e la organi maskulala e la organi feminala, ma en du flori separita\lingui{DEFIRS}
\artiklo{andromonoika}%
\vorto{andromonoika}\fako{bot.} Qualeso di planto che qua la flori maskula e la flori hermafrodita esas sur la sama stipito\lingui{DEFIS}
\artiklo{anekdoto}%
\vorto{anekdoto} Kurta naracuro pri mikra eventajo savinda\lingui{DEFIRS}
\artiklo{anemio}%
\vorto{anemio}\fako{patol.} Deperiso produktata da la povresko di sango\lingui{DEFIRS}
\artiklo{anemofila}%
\vorto{anemofila}\fako{biol.} Dicesas pri la planti che qui la dissemo di poleno eventas per vento\lingui{DEFIS}
\artiklo{anemometro}%
\vorto{anemometro} Instrumento per qua onu povas mezurar la rapideso-grado di vento\lingui{DEFIRS}
\artiklo{anemono}%
\vorto{anemono}\fako{bot.} Planto de la familio « ranunkulei », grupo de la « anemonei »\latina{anemone}\lingui{DEFIRS}
\artiklo{aneroido}%
\vorto{aneroido} Barometro sen liquido\lingui{DEFIRS}
\artiklo{anestezio}%
\vorto{anestezio}\fako{trans.) (kirurg.} Igar ne-sentanta, per igar la malado absorbar ula substanci (kloroformo, etero, e c.)\lingui{DEFIRS}
\artiklo{aneto}%
\vorto{aneto}\fako{bot.} Planto umbelifera di qua la semini esas toniziva ed ecitiva\latina{anethum graveolens}\lingui{EFIRS}
\artiklo{anevrismo}%
\vorto{anevrismo}\fako{patol.} Tumoro produktita sur vaskulo arteria o sur kordio-parieto, da la distenseso di la tuniki\lingui{DEFIRS}
\artiklo{anexar}%
\vorto{anexar}\fako{trans.} Juntar ad ensemblo ulo qua divenas dependanta de ta ensemblo\lingui{DEFIRS}
\artiklo{angeliko}%
\vorto{angeliko}\fako{bot.} Planto aromata, umbelifera, di qua onu konfitas la kosti verda\latina{angelica archangelica}\lingui{DEFIRS}
\artiklo{angelo}%
\vorto{angelo} Fera hoko, kun pinti, quan onu muntas ye la extremajo di pesko-filo e quan onu provizas ye lurivo por kaptar fishi\lingui{DS}
\artiklo{angelus}%
\vorto{« angelus »}\fako{liturg. katol.} Prego quan onu agas en la matino, ye dimezo ed en la vespero, memore di la anjelo-saluto\lingui{DEFIRS}
\artiklo{angino}%
\vorto{angino}\fako{patol.} Inflamo-morbo di la guturo, kun o sen produkto di membrani acesora\lingui{DEFIRS}
\artiklo{angiosperma}%
\vorto{angiosperma}\fako{bot.} Dicesas pri la planti di qui la grani kontenesas interne di la frukto\lingui{DEFIS}
\artiklo{anglikana}%
\vorto{anglikana}\fako{religio kristana} Qua apartenas a ta formi di protestantismo qua esas la Stato-religio di Britania\lingui{DEFIRS}
\artiklo{angora}%
\vorto{« angora »} Nomo di raso di kato e di kapro, di qui la pili esas longa e silkatra\lingui{DEFIRS}
\artiklo{angorar}%
\vorto{angorar}\fako{netrans., pri, pro} Esar la objekto di ta sento fiziologiala di fauco-konstrikto e di quaza sufoko qua akompanas desquieteso od expekto pasionoza\lingui{EFIS}
\artiklo{anguilo}%
\vorto{anguilo}\fako{zool.} Genero de fisho reptera, di qua la korpo esas tenua, la pelo eskapanta, e quin onu trovas en la aqui sen-sala\latina{anguilla}\lingui{FIS}
\artiklo{angulo}%
\vorto{angulo}\fako{geom.} Inter-eskarto-grado di du rekta linei qui inter-sekas sur plano\lingui{EFIS}
\artiklo{anhelar}%
\vorto{anhelar}\fako{netrans., pro} Respirar en maniero penoza, hastoza. (Ex. : Me anheligis me, pro acensir ta eskalero tante rapide.)\lingui{EFIS}
\artiklo{anhidra}%
\vorto{anhidra}\fako{kemio} Qua ne kontenas aquo\lingui{DEFIRS}
\artiklo{anilino}%
\vorto{anilino}\fako{kemio} CHNH₂\lingui{DEFIRS}
\artiklo{animalo}%
\vorto{animalo} Ento organizita qua povas sentar e movar su\lingui{DEFIRS}
\artiklo{animar}%
\vorto{animar}\fako{trans.} Igar plu vivoza en lua agado, en lua afecioni\lingui{DEFIRS}
\artiklo{animismo}%
\vorto{animismo}\fako{biol.} Doktrino qua konsideras la anmo kom la kauzo di la fenomeni psikala e fiziologiala\lingui{DEF}
\artiklo{aniono}%
\vorto{aniono}\fako{elektro} Iono qua portas elektro negativa, e qua, en dissolvuro elektrolitika, movas su ad la anodo di la fluo\lingui{DEFIRS}
\artiklo{aniversario}%
\vorto{aniversario} Dato qua korespondas a la dio kande ulo eventis, un o plura yari antee\lingui{DEFIS}
\artiklo{anizo}%
\vorto{anizo}\fako{bot.} Planto umbelifera di qua la semini esas aromata\latina{pimpinella anisum}\lingui{DEFIRS}
\artiklo{anizotropa}%
\vorto{anizotropa}\fako{biol.} Dicesas pri la korpi organika quan karakterizas bi-refraktiveso. (Ex. : la fragmento obskura di la muskulo-fibro strioza)\lingui{DEFIS}
\artiklo{anjelo}%
\vorto{anjelo}\fako{religio krist.} Pura spirito qua mediacas Deo e la homo\lingui{DEFIRS}
\artiklo{anke}%
\vorto{anke} Simile, egale, same – aludante ago o stando qua riproduktesas\lingui{I}
\artiklo{ankilosar}%
\vorto{ankilosar}\fako{trans.) (patol.} Paralizar per la weldo di artiko movebla\lingui{DEFIRS}
\artiklo{ankore}%
\vorto{ankore} Expresas la ne-ceso di la ago o di la stando (quin indikas la verbo)\lingui{FI}
\artiklo{ankro}%
\vorto{ankro} Peco de fero, provizita ye du pintegi, quan onu imersas til la fundo di maro ube olu akroche fixigas su tale ke la navo, sur la korpo di qua lu esas muntita, retenesas\lingui{DEFIRS}
\artiklo{anmo}%
\vorto{anmo}\fako{religio o filoz.} Ento metafizikala, principo di vivo e di penso\lingui{DEFIS}
\artiklo{anobio}%
\vorto{anobio}\fako{zool.} Genero de insekto koleoptera, familio de la « anobiidi »\latina{anobium}\lingui{FL}
\artiklo{anodo}%
\vorto{anodo}\fako{elektro} Elektrodo pozitiva di elektro-baterio, di ampulo radiografala\lingui{DEFIRS}
\artiklo{anomala}%
\vorto{anomala} Qua aspektas stranje, baroke\lingui{DEFIRS}
\artiklo{anonima}%
\vorto{anonima}\senco{1}{Qua ne esas provizita ye la nomo di la autoro}\senco{2}{Qua ne konocigas sua nomo}\lingui{DEFIRS}
\artiklo{anso}%
\vorto{anso} Parto salianta, ringa od arko, ye qua esas provizita utensilo od implemento, por plufaciligar la sizo di olca
\artiklo{-ant-}%
\vorto{-ant-} Sufixo di la participo aktiva prezenta
\artiklo{antagonismo}%
\vorto{antagonismo} Stando di interluktado\lingui{DEFIRS}
\artiklo{antagonisto}%
\vorto{antagonisto} Ta qua luktas kontre altra\lingui{DEFIRS}
\artiklo{antarktika}%
\vorto{antarktika} Latero qua ne direktesas vers la Urso-stelaro, ma vers la sudo-polo di la Terglobo\lingui{DEFIRS}
\artiklo{ante}%
\vorto{ante} En epoko qua pre-iras la evento quan onu mencionas\lingui{DEFIRS}
\artiklo{antecedento}%
\vorto{antecedento}\senco{1}{Antea agi di persono}\senco{2}{Termino di propoziciono qua pre-iras logikale to quo relative « konsequas »}\senco{3}{Termino unesma di relato aritmetikala o geometriala}\lingui{DEFIS}
\artiklo{anteno}%
\vorto{anteno}\senco{1}{\textit{(zool.)} Apendico kornatra, artikoza e movebla, situita sur la kapo di la maxim multa ek la insekti}\senco{2}{\textit{(elektro, radio e televiziono)} Longa vergo, fixigita vertikale relate sulo}\lingui{DEFIS}
\artiklo{anteridio}%
\vorto{anteridio}\fako{biol.} Organo riproduktala maskulala di la maxim multa kriptogami, en qua naskas la anterozoidi\lingui{DEFIS}
\artiklo{antero}%
\vorto{antero}\fako{bot.} Mikra posho qua situesas sur la extremajo di la stamino e qua kontenas la poleno\lingui{DEFIS}
\artiklo{anterozoido}%
\vorto{anterozoido}\fako{bot.} Korpi riproduktala maskulala di ula kriptogami\lingui{DEFIS}
\artiklo{antezo}%
\vorto{antezo}\fako{biol.} Desklozo di la floro, qua liberigas la poleno\lingui{DEFIS}
\artiklo{anti-}%
\vorto{anti-} Prefixo (rezervita a la terminaro ciencala e teknikala) qua signifikas « kontre »\lingui{DEFIRS}
\artiklo{anticipar}%
\vorto{anticipar}\fako{trans.}\senco{1}{Prenar, proprigar a su ante la dato}\senco{2}{Facar, agar, exekutar ante la dato}\lingui{DEFIS}
\artiklo{antidoto}%
\vorto{antidoto} Substanco quan onu absorbigas da ulu, por kombatar en lu la efiko di veneno, di viruso\lingui{DEFIRS}
\artiklo{antifermento}%
\vorto{antifermento}\fako{kemio, biol.} Korpo qua su opozas a la efiko di fermento\lingui{DEFIS}
\artiklo{antifono}%
\vorto{antifono}\fako{liturgio katol.} Versaro quan onu kantas, tote o parte, ante psalmo o kantiko, e quan onu repetas tote, pose\lingui{DEFIRS}
\artiklo{antifrazo}%
\vorto{antifrazo} Frazeto uzata – eufemisme od ironie – por expresar la kontreajo di la senco vera\lingui{DEFIRS}
\artiklo{antiklina}%
\vorto{antiklina}\fako{geol.} Dicesas pri la plisuro di strato en qua la parto konkava regardas adinfre
\artiklo{antilogio}%
\vorto{antilogio} Des-kohero di la idei, di la dico-maniero\lingui{DEFIS}
\artiklo{antilopo}%
\vorto{antilopo}\fako{zool.} Ruminero di formo gracila, gracoza, kun korni ne-kaduka, kava, cirkondanta, nukleo ostoza, bloka\latina{antilope}\lingui{DEFIRS}
\artiklo{antimonio}%
\vorto{antimonio}\fako{kemio} Korpo simpla, metala, qua uzesas por kompozar emetiko, e qua pulverigite uzesas nomo di « kohl »\lingui{DEFIRS}
\artiklo{antinomio}%
\vorto{antinomio} Des-kohero, reala o semblanta, di du legi naturala (*leyi) o yurala\lingui{DEFIRS}
\artiklo{antipatiar}%
\vorto{antipatiar}\fako{trans.} Instinte eskartar su de (ulo o de ulu)\lingui{DEFIRS}
\artiklo{antipirino}%
\vorto{antipirino} Medikamento obtenita de un ek la kompozanti di karbono, e qua kombatas febro, nevralgii\lingui{DEFIRS}
\artiklo{antipodo}%
\vorto{antipodo}\senco{1}{Persono qua okupas sur la Terglobo loko diametre opozata}\senco{2}{Sur la Terglobo, loko diametre opozata ad altra}\lingui{DEFIRS}
\artiklo{antiqua}%
\vorto{antiqua} Qua apartenas a la vivo-maniero, a la kustumi di epoki tre antea (partikulare, qua apartenas a la civilizeso-sistemo di la Greki e di la Romani lor la quaresma e lor la unesma yarcenti respektive ante nia ero)\lingui{DEFIRS}
\artiklo{antisepsiar}%
\vorto{antisepsiar}\fako{trans.} Uzar produkturo apta kombatar la putro morbala\lingui{DEFIRS}
\artiklo{antitezo}%
\vorto{antitezo} Kontrasto qua produktesas da la apudigo di du idei, di du expresuri, qui inter-opozesas\lingui{DEFIRS}
\artiklo{antitoxino}%
\vorto{antitoxino} Substanco apta nihiligar la efiko di toxino\lingui{DEF}
\artiklo{antologio}%
\vorto{antologio} Kolektajo de mikra versari\lingui{DEFIRS}
\artiklo{antonimo}%
\vorto{antonimo} Vorto qua, konsiderate relate ad altra vorto, havas senco absolute opozata\lingui{DEFIRS}
\artiklo{antonomazio}%
\vorto{antonomazio} Metaforo qua indikas persono per la karakterizivo di qua lu esas la tipo; od individuo qua posedas ula karakterizivi per la persono qua esas lua tipo\lingui{DEFIRS}
\artiklo{antracito}%
\vorto{antracito}\fako{mineral.} Karbono fosila, qua brulas nur desfacile, kun flamo kurta, sen odoro nek fumo\lingui{DEFIRS}
\artiklo{antraxo}%
\vorto{antraxo}\fako{patol.} Inflamo-tumoro di la tisuo celulara sub-pela\lingui{DEFIS}
\artiklo{antropoida}%
\vorto{antropoida} Qua similesas la homo, di qua la vizajo esas formo-vicina a la vizajo homala\lingui{DEFIRS}
\artiklo{antropologio}%
\vorto{antropologio} Ta parto di la naturo-historio qua traktas la homo-speco\lingui{DEFIRS}
\artiklo{antropometrio}%
\vorto{antropometrio}\senco{1}{Studio di la proporcioni di la parti di la homo-korpo}\senco{2}{Ta studio aplikata a la krimininti por rikonocar li}\lingui{DEFIRS}
\artiklo{antropomorfa}%
\vorto{antropomorfa} Qua atribuas a la deajo la vizajo, la sentimenti, la pasioni di la homo\lingui{DEFIRS}
\artiklo{anuncar}%
\vorto{anuncar}\fako{trans., ad} Konocigar da publiko per avizo vocala o skribura\lingui{DEFIS}
\artiklo{anuso}%
\vorto{anuso}\fako{anat.} Orifico extera di la rektumo, per qua ekiras feko\lingui{DEFIS}
\artiklo{anxiar}%
\vorto{anxiar}\fako{netrans., pri, pro, de} Esar tale desquieta ke la kordio quaze konstriktesas\lingui{DEFIS}
\artiklo{aoristo}%
\vorto{aoristo}\fako{gram.} Tempo di la konjugado Greka, qua expresas la ago pasinta en maniero nepreciza (same kam « -is » che Ido)\lingui{DEFIRS}
\artiklo{aortito}%
\vorto{aortito}\fako{patol.} Inflamuro di la tuniko extera di aorto\lingui{DEFIRS}
\artiklo{aorto}%
\vorto{aorto}\fako{anat.} Arterio precipua qua departas de la kordio e qua esas la fundamento di omna vaskuli arteria di la « granda cirkulo-kanalaro »\lingui{DEFIRS}
\artiklo{apanajo}%
\vorto{apanajo} Domeno quan la reji di Francia grantis a sua filiuli juniora, a sua fratuli, e qua retrovenis a la trono kande omna decendanti di la grantarii mortis\lingui{DEFIRS}
\artiklo{aparar}%
\vorto{aparar} Videbleskar subite\lingui{DEFIS}
\artiklo{aparato}%
\vorto{aparato} Ensemblo de la objekti, instrumenti, qui uzesas por exekutar o facar ulo\lingui{DEFIRS}
\artiklo{aparta}%
\vorto{aparta} Quaze separita, izolita, sola, quankam permananta kom parto di ensemblo\lingui{DEFIS}
\artiklo{apartamento}%
\vorto{apartamento} Parto de domo, qua konsistas ye ula quanto de chambri e qua esas habiteyo komfortoza\lingui{DEFIRS}
\artiklo{apartenar}%
\vorto{apartenar}\fako{netrans., ad} Esar la proprietajo di ulu. \textit{(Anke metaf.)}\lingui{EFIS}
\artiklo{apatio}%
\vorto{apatio} Stando di anmo o di mento o di psiko quin nulo povas emocigar\lingui{DEFIRS}
\artiklo{apelar}%
\vorto{apelar}\fako{netrans., ad ulu, pri ulo} Rekursar a la judiciistaro kompetenta di la grado superiora, por igar reformar verdikto qua enuncesis da tribunalo inferiora\lingui{DEFIRS}
\artiklo{apendicito}%
\vorto{apendicito}\fako{patol.} Inflamuro di la apendico cekala vermoforma\lingui{DEFIS}
\artiklo{apendico}%
\vorto{apendico} Parto qua servas kom parto di bloko precipua\lingui{DEFIS}
\artiklo{apene}%
\vorto{apene} Preske ne\lingui{FIS}
\artiklo{aperceptar}%
\vorto{aperceptar}\fako{trans.) (filoz.} Intuico, fakultato od ago qua konsistas ye intelektar quik, per koncio, ta o ca ideo o verajo. (En la filozofio da Leibniz, ta vorto signifikis la percepto juntata a reflekto)
\artiklo{apertar}%
\vorto{apertar}\fako{trans., por} Igar acesebla per deprenar to quo kluzigas\lingui{DEFIS}
\artiklo{apetitar}%
\vorto{apetitar}\fako{trans.} Tendencar a to quo satisfacos sua bezoni. \textit{(Anke metaf.)}\lingui{DEFIRS}
\artiklo{api-pomo}%
\vorto{api-pomo}\fako{bot.} Mikra pomo krakanta, di qua un latero esas blankatra, e la altra, reda\latina{malum appianum}
\artiklo{apikala}%
\vorto{apikala} Dicesas pri la parto qua esas la somito di organo
\artiklo{apio}%
\vorto{apio}\fako{bot.} Planto umbelifera, di qua un speco divenis, per kultivado, celerio\latina{apium}\lingui{DEFIRS}
\artiklo{aplanata}%
\vorto{aplanata}\fako{fotogr.} Tipo di objektali rektalinea e simetra\lingui{DEF}
\artiklo{aplanatika}%
\vorto{aplanatika}\fako{optiko} En qua ne eventas aberaco sferesala\lingui{DEF}
\artiklo{aplastar}%
\vorto{aplastar}\fako{trans.}\senco{1}{Platigar e deformar (korpo) per sinkar ula parto per shoko violentoza o per granda preso}\senco{2}{\textit{(metaf.)} Igar sukombar ad pezozajo tro grava, a kargajo tro grava; igar sukombar en lukto per efikigar forco a qua ne esas posibla rezistar}\lingui{FIS}
\artiklo{aplaudar}%
\vorto{aplaudar}\fako{trans.} Intershokigar la parto interna di la manui kom manifesto di sua granda admiro\lingui{DEFIRS}
\artiklo{aplikar}%
\vorto{aplikar}\fako{trans., ad}\senco{1}{Pozar ulo sur ulo, tale ke la pozajo kovras la pozario komplete, adherante a lu}\senco{2}{Igar efikar (ad ulu, ad ulo) agado od esforco}\senco{3}{\textit{(metaf.)} Uzar procedo, igar efikar regulo, ad ulo, ad ulu}\lingui{DEFIRS}
\artiklo{aplombo}%
\vorto{aplombo}\senco{1}{Equilibro stabila di korpo kande vertikalo, quan indikas aplombo-filo, koincidas kun olua baricentro ed abutas a la fundamento qua subtenas lu}\senco{2}{\textit{(metaf.)} Audaco neperturbebla}\lingui{DEFIRS}
\artiklo{apodiktika}%
\vorto{apodiktika}\fako{filoz.} Qua enuncas verajo necesa. (Ex. : Omna cirklo havas centro)\lingui{DEFIS}
\artiklo{apofizo}%
\vorto{apofizo}\fako{anat.} Saliajo naturala di la osti\lingui{DEFIS}
\artiklo{apogar}%
\vorto{apogar}\fako{trans., ad, an, sur} Suportar od igar suportar per forco laterala od obliqua. Suportar, sustenar per apogilo\lingui{FIS}
\artiklo{apogeo}%
\vorto{apogeo}\fako{astron.} Punto ube Luno e Suno distas de Tero maxime\lingui{DEFIRS}
\artiklo{apokalipso}%
\vorto{apokalipso}\fako{religio} Libro de la Nova Testamento en qua esas naracita la revelo agita a Santa Ioannes en la insulo Patmos\lingui{DEFIRS}
\artiklo{apokrifa}%
\vorto{apokrifa}\fako{religio}\senco{1}{Ne aceptata kom vera da la Eklezio katolika}\senco{2}{\textit{(pri libro)} Suspektinda, substitucita fraude, dubitinda}\lingui{DEFIRS}
\artiklo{apokromata}%
\vorto{apokromata}\fako{fotogr.} Dicesas pri objektalo fotografala o mikroskopala, en qua la akromateso realigesas, ne nur pri la kolori (objektalo akromata ordinara), ma pri tri kolori adminime\lingui{DEFIRS}
\artiklo{apologiar}%
\vorto{apologiar}\fako{trans.} Kompozuro, diskurso justifikiva\lingui{DEFIRS}
\artiklo{apologo}%
\vorto{apologo} Kurta rakonto, proza o versa, qua kontenas leciono di etiko aplikanta, maxim-multa-kaze per simboli ek animali od sen-vivaji\lingui{DEFIRS}
\artiklo{aponevroso}%
\vorto{aponevroso}\fako{anat.} Membrano fibroza blanka, tre rezistiva, qua tegas komplete e mantenas la muskuli, inter-separas la faski de muskuli, o ligas li an la osti\lingui{DEF}
\artiklo{apoplexio}%
\vorto{apoplexio}\fako{patol.} Paralizeso plu o min kompleta, produktita da lezuro cerebrala\lingui{DEFIRS}
\artiklo{apostata}%
\vorto{apostata}\fako{religio} Qua abandonis sua religio (la religio en qua lu naskis) po altra\lingui{DEFIRS}
\artiklo{a posteriori}%
\vorto{« a posteriori »} Segun la fakti observita
\artiklo{apostilo}%
\vorto{apostilo} Noto sur peticiono, por rekomendar lu\lingui{EFIS}
\artiklo{apostolo}%
\vorto{apostolo}\fako{religio katol.} Singla ek la dek-e-du dicipuli qui recevis de Iesu-Kristo la komiso predikor evangelio. – Ta qua propagas la doktrino Kristana\lingui{DEFIR}
\artiklo{apostrofo}%
\vorto{apostrofo}\senco{1}{\textit{(gram.)} Signo ortografiala qua remplasas litero elizionita}\senco{2}{\textit{(retor.)} Procedo oratorala per qua onu interpelas subite persono prezenta, ento nevidebla, kozo quan onu traktas quale se lu esus persono. – Interpelo vigoroza di ulu}\lingui{DEFIRS}
\artiklo{apoteko}%
\vorto{apoteko} Butiko ube la medikamenti preparesas, konservesas, vendesas. a konsumeri\lingui{DEFIRS}
\artiklo{apotemo}%
\vorto{apotemo}\fako{geom.}\senco{1}{Radio di la cirklo enskribata en poligono regulala}\senco{2}{Alteso di un ek la edri di piramido regulala}\lingui{DEFIRS}
\artiklo{apoteoso}%
\vorto{apoteoso}\senco{1}{Granto di honori deala}\senco{2}{Granto di honori extraordinara}\lingui{DEFIRS}
\artiklo{apoziciono}%
\vorto{apoziciono}\fako{gram.} Qualifiko quan onu adjuntas a substantivo per altra substantivo apud-pozata\lingui{DEFIRS}
\artiklo{aprentiso}%
\vorto{aprentiso}\fako{pri} Ta qua lernas mestiero\lingui{DEFIS}
\artiklo{apretar}%
\vorto{apretar}\fako{trans.) (tekn.}\senco{1}{Submisar ula texuri a preparajo qua igas li koheranta, solida, glata, briloza}\senco{2}{Indutar telo sur qua onu esas piktonta}\lingui{DFIR}
\artiklo{aprilo}%
\vorto{aprilo} La monato quaresma di la kalendar-yaro\lingui{DEFIRS}
\artiklo{a priori}%
\vorto{a priori} Segun donaji qui pre-iras experimentado\lingui{DEFIRS}
\artiklo{apro}%
\vorto{apro}\fako{zool.} Mamifero, ek la pakidermi, ek genero qua vivas sovaja, ed egardesas kom la tronko di la porko domestika\latina{aper}\lingui{DRL}
\artiklo{aprobar}%
\vorto{aprobar}\fako{trans.}\senco{1}{Judikar bona, laudinda}\senco{2}{Deklarar (per formulo) ke akto, konto, esas exakta}\lingui{DEFIRS}
\artiklo{aprocho}%
\vorto{aprocho}\fako{fortifik.} Terlaboruro per qua militisti *asiejanta avancas a la fortifikuri e rempari enemikala\lingui{DEFIRS}
\artiklo{aproximar}%
\vorto{aproximar}\fako{trans.) (tekn.} Evaluar tale ke onu saciesas per nura su-proximigo a lo exakta
\artiklo{apsido}%
\vorto{apsido}\senco{1}{\textit{(arkeol.)} Mi-cirklo qua formacas la fundo di la baziliki pagana, e divenis la santuario di la baziliki Kristana}\senco{2}{\textit{(astron.)} Punto di la orbito di planeto ube lu distas maxime o minime de la astro cirkum qua lu su movas}\lingui{DEFIRS}
\artiklo{apta}%
\vorto{apta} Qua havas la karakteri necesa por (agar, facar, efikar)\lingui{EFIS}
\artiklo{apud}%
\vorto{apud} Tre proxime, ma ne tam grande kam indikas « an »\lingui{IL}
\artiklo{apuntar}%
\vorto{apuntar}\fako{fusilo, kanono, lorno, teleskopo} Direktar, mantenar fixigita vers to quon onu vizas\lingui{DEFIS}
\artiklo{aquaforto}%
\vorto{aquaforto} Azot-acido (AzO₃H) dilutita per aquo, quan la grabisto uzas por korodar la kupro-plako\lingui{DEFIRS}
\artiklo{aquamarino}%
\vorto{aquamarino} Varietato di smeraldo, verda quale la mar-aquo\lingui{DEFIRS}
\artiklo{aquarelo}%
\vorto{aquarelo} Pikturo lejera, en preske omna kazi sur papero, facita ek farbi diafana, dilutita per aquo\lingui{DEFIRS}
\artiklo{aquario}%
\vorto{aquario} Tanko en qua onu entratenas animali o planti aquala\lingui{DEFIRS}
\artiklo{aquatinto}%
\vorto{aquatinto} Graburo per aquaforto, qua imitas aquopikturo\lingui{DEFIRS}
\artiklo{aquedukto}%
\vorto{aquedukto} Tubego subsula od elevita sur arki, qua adduktas aquo de loko ad altra\lingui{DEFIRS}
\artiklo{aquilegio}%
\vorto{aquilegio}\fako{bot.} Planto renunkulacea\latina{aquilegia}\lingui{L}
\artiklo{aquilono}%
\vorto{aquilono}\senco{1}{Vento de nordo}\senco{2}{\textit{(literaturo)} Vento kolda e violentoza}\lingui{DEFIRS}
\artiklo{aquirar}%
\vorto{aquirar}\fako{trans.) (per, po} Prokurar a su la posedo di ulo\lingui{DEFIS}
\artiklo{aquo}%
\vorto{aquo}\fako{kemio} Liquido qua konsistas ye du parti de hidrogeno ed un parto de oxigeno. H₂O\lingui{DEFIRS}
\artiklo{-ar}%
\vorto{-ar} Finalo di la verbi ye infinitivo di prezento e di neprecizeso epokala
\artiklo{arabesko}%
\vorto{arabesko} Ornamento di pikturo o di skulturo, facita ek planti, animali, fantazio-figuri interplektita kapricoze\lingui{DEFIRS}
\artiklo{arachar}%
\vorto{arachar}\fako{trans., de, ek} Deprenar per violento, koakto, de ulu to quon lu retenas. \textit{(Anke metaf.)}\lingui{FS}
\artiklo{araknoido}%
\vorto{araknoido}\fako{anat.} Un ek la tri envelopili di encefalo, membrano seroza tre dina, situita inter du altra\lingui{DEF}
\artiklo{arako}%
\vorto{arako} Liquoro alkoholoza quan onu obtenas de rizo o de la suko di kokoso\lingui{DEFIRS}
\artiklo{araneo}%
\vorto{araneo}\fako{zool.} Artropodo qua filifas e texas telo reta, per qua lu kaptas la insekti quin lu devoros\latina{aranea}\lingui{EFIS}
\artiklo{aranjar}%
\vorto{aranjar}\fako{trans.}\senco{1}{Dispozar la kozi tale ke (la mariajo, la partio di chaso, di ruro-promeno, e c.) povos eventar}\senco{2}{Agar tale ke (la afero) abutos ad interkonsento}\lingui{DEFI}
\artiklo{arbaleto}%
\vorto{arbaleto} Quaza arm-arko, muntita an afusto-ligna, qua direktas la projektilo\lingui{DEFIRS}
\artiklo{arbitrajar}%
\vorto{arbitrajar}\fako{netrans., pri) (financo} Bank-operaco qua konsistas ye (1) selektar la borso-urbo ube la kambio-koto esas la maxim avantajoza pri tratar o livrar valor-paperi; (2) igar valor-papero remplasar altra qua aspektas kom min avantajoza\lingui{DEFIRS}
\artiklo{arbitrar}%
\vorto{arbitrar}\fako{trans.} (\textit{aludante la persono qua esas selektita da} \textit{la litij-partoprenanti o nominita da tribunalo por aranjar} \textit{la afero}) Decidar, evaluar\lingui{DEFIRS}
\artiklo{arbitrio}%
\vorto{arbitrio} Volo libera\lingui{DEFIS}
\artiklo{arboro}%
\vorto{arboro}\senco{1}{\textit{(bot.)} Vejetanto perena, kun trunko alta e ligna}\senco{2}{To quo kelke similesas forme arboro}\lingui{EFIS}
\artiklo{arbusto}%
\vorto{arbusto}\fako{bot.} Planto di qua la stipo similesas olta di arboreto, ma qua kreskas quale la planti herbatra\lingui{FIS}
\artiklo{arbuto}%
\vorto{arbuto}\fako{bot.} Frukto globatra di planto di qua la foliaro esas sempre verda\latina{arbutus}\lingui{EFIL}
\artiklo{archo}%
\vorto{archo}\senco{1}{\textit{(biblo)} Sorto di navo klozita (*kluza) quan konstruktis Noe (Noah) por prezervar de la diluvio sua familio e la animali diversa}\senco{2}{Kofro qua kontenis la lego-tabeli qui esis donita a Mozes}\lingui{DEFIS}
\artiklo{ardezo}%
\vorto{ardezo} Argilo-skisto, nebulatre griza, facile separebla a folii dina\lingui{FI}
\artiklo{ardorar}%
\vorto{ardorar}\fako{netrans., pri} Sentar deziro tre granda\lingui{DEFIS}
\artiklo{ardua}%
\vorto{ardua}\senco{1}{Ube onu pasas nur desfacile pro la asperaji}\senco{2}{\textit{(metaf.)} Ube onu renkontras obstakli embarasanta}\lingui{EFIS}
\artiklo{areko}%
\vorto{areko}\fako{bot.} Genero de planto, de la familio « palmiero »\latina{areca}\lingui{DEFIRS}
\artiklo{areno}%
\vorto{areno} Parto sabloza di cirko, di amfiteatro, ube eventas lukti\lingui{DEFIRS}
\artiklo{areo}%
\vorto{areo} Spaco inkluzata en surfaco plana\lingui{DEFIS}
\artiklo{areolo}%
\vorto{areolo}\fako{anat.}\senco{1}{Mikra spaco inter la fibri di la tisuo celulara}\senco{2}{Cirklo brunatra qua cirkondas la mamo-pinto; cirklo redatra qua cirkondas punto inflamura}\lingui{DEFIS}
\artiklo{areometrio}%
\vorto{areometrio} Cienco qua havas kom objekto la mezuro di la denseso-grado di korpi, e precipue di liquidi\lingui{DEFIRS}
\artiklo{areometro}%
\vorto{areometro} Instrumento uzata por mezurar la denseso-grado di korpi, precipue di liquidi\lingui{DEFIRS}
\artiklo{areopago}%
\vorto{areopago} Tribunalo suprega, che la Ateniani antiqua, qua judiciis sur la Marso-kolino. \textit{(anke metaf.)}\lingui{DEFIRS}
\artiklo{arestar}%
\vorto{arestar}\fako{trans.} Kaptar ulu por enkarcerigar lu\lingui{DEFIRS}
\artiklo{argalo}%
\vorto{argalo}\fako{zool.} Ucelo di India, sur la ali di qua esas plumi dina, tenua, de qui onu facas kapvesti hominala eleganta\latina{leptoptilus argala}\lingui{EF}
\artiklo{argano}%
\vorto{argano} Aparato qua konsistas ye peci ordinare desmuntebla, e quan onu uzas por elevar kargaji\lingui{FIS}
\artiklo{argilo}%
\vorto{argilo}\fako{geol.} Tero mola e grasa, de qua onu fabrikas vazi multa-sorta, modeli por skulto, e c.\lingui{DEFIS}
\artiklo{argonauto}%
\vorto{argonauto} Nomo di la heroi qui akompanis Jazon en Kolkidia, sur la navo « Argo »\lingui{DEFIRS}
\artiklo{argono}%
\vorto{argono}\fako{kemio} Gaso simpla, senkolora, sensapora, senodoro, ye qua konsistas cirkume la centimo di nia atmosfero\simbolo{Ar}\lingui{DEF}
\artiklo{argumentar}%
\vorto{argumentar}\fako{trans., pri} Esforcar pruvar per la prizento di rezono\lingui{DEFIRS}
\artiklo{argus}%
\vorto{« argus »} Persono mitologiala qua havis cent okuli ed a qua Junono donabis kom instrucioni surveyor la nimfo Io\lingui{DEFIRS}
\artiklo{arida}%
\vorto{arida} Sen vejetanti, sterila\lingui{DEFIS}
\artiklo{arieto}%
\vorto{arieto}\fako{milit-arto} En la milit-arto di la antiqui e di la mezepokani : milito-mashino kun, ye un extremajo, mutono-kapo ek fero, per qua onu frapis la muregi di urbo *asiejata (*siejata)\lingui{S}
\artiklo{-ario}%
\vorto{-ario} Sufixo qua indikas « persono a qua nia ago direktesas, od ula kozo destinesas »
\artiklo{ario}%
\vorto{ario} Ensemblo de noti muzikala e de toni ye qui konsistas kanto\lingui{DEFIRS}
\artiklo{aristo}%
\vorto{aristo}\senco{1}{Stangeto osta en la skeleto di la fishi}\senco{2}{Interseko-lineo di du plani}\senco{3}{Barbo di la spiko di ula planti graminea}\lingui{EFIS}
\artiklo{aristokrata}%
\vorto{aristokrata} Qua apartenas a ta formo di guvernado en olqua la autoritato guvernala esas la privilejo di la nobeli\lingui{DEFIRS}
\artiklo{aristolokio}%
\vorto{aristolokio}\fako{bot.} Planto lignoza di qua la radiko uzesas kom tonizivo ed apertivo\latina{aristolochia}\lingui{DEFIS}
\artiklo{aritmetiko}%
\vorto{aritmetiko}\fako{matem.} Cienco di la nombri, di kalkulado\lingui{DEFIRS}
\artiklo{aritmografo}%
\vorto{aritmografo} Aparato por kalkular mekanike; kalkulo-mashino\lingui{DEFIRS}
\artiklo{aritmologio}%
\vorto{aritmologio} Nomo di la cienco generala di la nombri, di la mezuro di la *grandori, irgequalan esas oli\lingui{DEFIRS}
\artiklo{aritmomanciar}%
\vorto{aritmomanciar}\fako{netrans.} Asertar savar (per procedi super natura) to quo esas celita en la pasinto, la prezento e la futuro – per nombri\lingui{DEFIRS}
\artiklo{aritmometro}%
\vorto{aritmometro} Kalkulo-linealo\lingui{DEFIRS}
\artiklo{arivar}%
\vorto{arivar}\fako{netrans., an} Abutar an la fino-punto di sua iro\lingui{EFIS}
\artiklo{arjentano}%
\vorto{arjentano} Aloyuro de nikelo kun kupro o zinko\lingui{DEFIS}
\artiklo{arjento}%
\vorto{arjento}\fako{kemio} Metalo blanka, briloza, quan onu uzas por fabrikar moneto-peci, vazi, juveli, adjuntante ad olu aloyuro de kupro\simbolo{Ag}\lingui{EFIS}
\artiklo{arkabuzo}%
\vorto{arkabuzo} Olima paf-armo portebla\lingui{DEFIRS}
\artiklo{arkado}%
\vorto{arkado} Konstrukturo qua havas la formo di arko, qua su apogas a pilastri o koloni qui posibligas subpaso\lingui{DEFIRS}
\artiklo{arkaika}%
\vorto{arkaika} Qua evas de la komenco di arto\lingui{DEFIRS}
\artiklo{arkaismo}%
\vorto{arkaismo} Expresuro, vort-aranjo qua esas apene uzata pro olua quaza antiqueso\lingui{DEFIRS}
\artiklo{arkano}%
\vorto{arkano} Sekreto ciencala (o pseudo-ciencala) quan onu revelas o docas nur ad iniciati. (Ex. : la arkani di alkemio)\lingui{DEFIRS}
\artiklo{arkegonio}%
\vorto{arkegonio}\fako{bot.} Genit-organo feminala di la planti hepatikala, di la mosi e di la kriptogami vaskulala e gimnosperma\lingui{DEFIS}
\artiklo{arkeologio}%
\vorto{arkeologio} Cienco pri la kozi antiqua\lingui{DEFIRS}
\artiklo{arkeologo}%
\vorto{arkeologo} La persono qua su okupas pri arkeologio
\artiklo{arkesporo}%
\vorto{arkesporo}\fako{bot.} Parto interna di sporuyo che la muski, filiki\lingui{DEFIS}
\artiklo{arki-}%
\vorto{arki-} Prefixo qua indikas « en grado eminenta »; lu, do, indikas grado supera\lingui{DEFIRS}
\artiklo{arkimandrito}%
\vorto{arkimandrito}\fako{eklezio Greka} Superioro di monakeyo unesmaklasa\lingui{DEFIRS}
\artiklo{arkipelago}%
\vorto{arkipelago} Grupo de insuli\lingui{DEFIRS}
\artiklo{arkiplasmo}%
\vorto{arkiplasmo}\fako{biol.} Substanco ye qua konsistas la direktosferi en kariokinezo\lingui{DEFIS}
\artiklo{arkitekto}%
\vorto{arkitekto} La persono qua facas la mapo di edifico e direktas la konstrukto di olca. \textit{(anke metaf.)}\lingui{DEFIRS}
\artiklo{arkitekturo}%
\vorto{arkitekturo} La arto konstruktar edifici\lingui{DEFIRS}
\artiklo{arkitravo}%
\vorto{arkitravo} Infra parto di la entablamento qua su apogas sur la kapiteli di la koloni\lingui{DEFIRS}
\artiklo{arkivo}%
\vorto{arkivo} Singla kompozanto ek la ensemblo de la akti, dokumenti pri la historio di populo, di urbo, di instituto publika o privata, di komerco-societo\lingui{DEFIRS}
\artiklo{arkivolto}%
\vorto{arkivolto} Kadro di la arkado di pordo, di fenestro\lingui{DEFIRS}
\artiklo{arko}%
\vorto{arko}\senco{1}{Parto limitoza di kurva lineo irgaspeca}\senco{2}{Konstrukturo qua havas la formo di kurvo}\senco{3}{Lumo dazlanta qua spricas inter du peci de karbono ligita an la du elektrodi di generatoro elektrala}\lingui{DEFIRS}
\artiklo{arkonto}%
\vorto{arkonto} Un ek la non oficieri elektita singlayare, qui guvernis la republiko di la urbo Atena, antique\lingui{DEFIRS}
\artiklo{arkoplasmo}%
\vorto{arkoplasmo}\fako{biol.} Atrakto-sfero o direkto-vezikulo, judikata kom un ek la organi prima di la celulo\lingui{DEFIS}
\artiklo{arktika}%
\vorto{arktika} Situita en la mi-sfero di la Terglobo, qua regardas vers la stelaro « Granda Urso », vers la nordo-polo di la Terglobo\lingui{DEFIRS}
\artiklo{arlekino}%
\vorto{arlekino} Persono komika di la farso Italiana; lu aparas sur la ceneyo, kun vesto konsistanta ye peci triangula, diversa-kolora, e kun ligna sabro en la manuo\lingui{DEFIRS}
\artiklo{armadilo}%
\vorto{armadilo}\fako{zool.} Mamifero sen-denta, kelke simila a la aselo\latina{dasypus}\lingui{DEFIS}
\artiklo{armatoro}%
\vorto{armatoro} Posedanto ed equipanto di navo\lingui{FIRS}
\artiklo{armaturo}%
\vorto{armaturo} Asembluro de peci de fero, por mantenar kunligite la diversa parti di karpentajo o di masonajo\lingui{DEFIRS}
\artiklo{armeo}%
\vorto{armeo} Trupo de armieri diversa-sorta, destinita a milito\lingui{DEFIRS}
\artiklo{armistico}%
\vorto{armistico} Ceso provizora di milito\lingui{DEFIS}
\artiklo{armo}%
\vorto{armo} Instrumento uzata da la homo por atakar o por su defensar\lingui{DEFIS}
\artiklo{armoracio}%
\vorto{armoracio} Speco di krucifero di la genero « koklearia »\latina{cochlearia armoracia}\lingui{FIS}
\artiklo{armoro}%
\vorto{armoro} Alta moblo, klozebla per klapi, destinita a kontenar objekti kelka-valora\lingui{FIS}
\artiklo{arniko}%
\vorto{arniko}\senco{1}{Planto aromata di la familio « kompozaji », qua havas povo stimulala e tonizala}\senco{2}{Medikamento liquida, ek ta planto}\latina{arnica montana}\lingui{DEFIS}
\artiklo{-aro}%
\vorto{-aro} Sufixo qua, soldita a radiko nomala, indikas « la maxim extensita kolektajo, ensemblo » ek la kozi o personi nomata da la radiko
\artiklo{arogar}%
\vorto{arogar}\fako{trans., ad} Atribuar (a su) sen havar yuro ye olu (titulo, povo)\lingui{DEFIS}
\artiklo{aromato}%
\vorto{aromato} Substanco odoranta, quan onu uzas kom parfumo, medikamento, kondimento, e c.\lingui{DEFIRS}
\artiklo{aromo}%
\vorto{aromo} Emanajo de substanci odoranta, precipue de olti qui efikas a flarado\lingui{DEFIRS}
\artiklo{aroruto}%
\vorto{aroruto} Farino manjebla quan onu extraktas de la radiko di L. « \textit{maranta indica} »\lingui{DEFIRS}
\artiklo{arozar}%
\vorto{arozar}\fako{trans.} Humidigar per varsar liquido gutope, imitante roso\lingui{F}
\artiklo{arpejar}%
\vorto{arpejar}\fako{netrans.) (muziko} Plear muzikajo-parto tre rapide, igante audeblan intersucedante la son-uni di akordo\lingui{DEFIRS}
\artiklo{arsenalo}%
\vorto{arsenalo} Loko ube onu fabrikas o depozas la armi, la municioni, la milit-instrumenti quin onu uzas tere, mare od aere\lingui{DEFIRS}
\artiklo{arseno}%
\vorto{arseno}\fako{kemio} Korpo simpla, metaloida, stalo-griza, facile ruptebla, di qui la kompozaji (la aloyuri eceptate) esas venenoza\simbolo{As}\lingui{DEFIRS}
\artiklo{artemizio}%
\vorto{artemizio}\fako{bot.} Planto di la familio « kompozaji », toniziva\latina{artemisia}\lingui{FL}
\artiklo{arterio}%
\vorto{arterio}\fako{anat.} Vaskulo qua portas la sango de la kordio ad omna parti di la korpo\lingui{DEFIRS}
\artiklo{arteriologio}%
\vorto{arteriologio} Ta parto di anatomio, qua traktas la arterii\lingui{DEFIRS}
\artiklo{arteriotomio}%
\vorto{arteriotomio} Aperto di arterio per lanceto\lingui{DEFIRS}
\artiklo{arterito}%
\vorto{arterito}\fako{patol.} Inflamuro di arterio\lingui{DEF}
\artiklo{artezo-puteo}%
\vorto{artezo-puteo} Puteo borita til aquo-strato subsula qua lore spriceskas\lingui{DEFIRS}
\artiklo{artichoko}%
\vorto{artichoko}\fako{bot.} Sorto di kardono di qua la floro konsistas ye folii imbrikita, kun fundamento pulpoza\latina{cynara scolymus}\lingui{DEFR}
\artiklo{artifico}%
\vorto{artifico} To quo uzesas por travestiar naturalajo, verajo, por trompar\lingui{DEFIS}
\artiklo{artiklo}%
\vorto{artiklo}\senco{1}{\textit{(gram.)} Partikulo qua, en ula lingui nacionala, juntesas a substantivo quan lu determinas}\senco{2}{\textit{(jurnalo, revuo, e c.)} Singla de la parti di edituro kolektiva}\lingui{DEFIRS}
\artiklo{artiko}%
\vorto{artiko}\fako{anat.} Junturo naturala inter du parti di korpo, movebla ica sur ita\lingui{DEFIS}
\artiklo{artikular}%
\vorto{artikular}\fako{trans.} Emisar son-uni vocala, donante a li formi diversa, per ula movi da la labii, da la lango\lingui{DEFIRS}
\artiklo{artilrio}%
\vorto{artilrio}\fako{milit-arto} Armeo-korpo uzata por la funcionigo di la kanoni\lingui{DEFIRS}
\artiklo{arto}%
\vorto{arto} Singla de la generi en qui la homo o la animalo produktas verki segun ula reguli. – Aplikado di la savi por realigar koncepto\lingui{DEFIRS}
\artiklo{artrito}%
\vorto{artrito}\fako{patol.} Inflamuro di la tisui fibroza e seroza di la artiki\lingui{DEFIRS}
\artiklo{artro}%
\vorto{artro}\fako{anat.} Korpo-parto movebla sur altra a qua lu esas juntita\lingui{FS}
\artiklo{artrodio}%
\vorto{artrodio}\senco{1}{\textit{(anat.)} Artiko en qua esas inkastrita la extremajo di altra osto}\senco{2}{\textit{(mekan.)} Junturo qua su movas quale genuo}\lingui{EFS}
\artiklo{artropodo}%
\vorto{artropodo}\fako{zool.} Brancho di la animalo-regno, qua inkluzas la animali chitinoza (krabi, kankri, e c.)\lingui{DEFIS}
\artiklo{arufar}%
\vorto{arufar}\fako{trans.} Manuagar kontre la senso normala la hari sur la kapo\lingui{I}
\artiklo{arumo}%
\vorto{arumo}\fako{bot.} Planto herbatra di qua ula speci kontenas nutrofekulo\latina{arum}\lingui{DEFIRS}
\artiklo{aruspico}%
\vorto{aruspico}\fako{che la Romani antiqua} Sacerdoto qua exploris la internaji di la viktimi por obtenar de ta internaji auguri\lingui{DEFIRS}
\artiklo{-as}%
\vorto{-as}\fako{gram.} Finalo di la indikativo prezenta (e sen-epoko) di la verbi
\artiklo{asafetido}%
\vorto{asafetido}\fako{bot.} Gumo odoranta fetide, qua venas de speco di ferulo de Persia ed uzesas medicine\latina{ferula asa foetida}\lingui{DEFIS}
\artiklo{asaltar}%
\vorto{asaltar}\fako{trans.) (milit-arto} Atakar bruske por violentovinkar loko fortifikita\lingui{EFIRS}
\artiklo{asasinar}%
\vorto{asasinar}\fako{trans.} Mortigar kun preparo, embuske o surprize\lingui{DEFIS}
\artiklo{asbesto}%
\vorto{asbesto} Substanco minerala amfibola, qua esas fuzebla nur ye temperaturo extreme alta-grada\lingui{DEFIRS}
\artiklo{asekurar}%
\vorto{asekurar}\fako{trans., kontre} Garantiar pagor ula quanto de pekunio se ta o ca risko (objekto di la asekuro) realigesos – po pekunio-quanto quan la asekurito pagas un-foye, o singla-yare (akonte)\lingui{DEFIS}
\artiklo{aselo}%
\vorto{aselo}\fako{zool.} Mikra krustaceo izopoda qua prizas vivar en loki obskura e humida\latina{oniscus asellus}\lingui{DIL}
\artiklo{asemblar}%
\vorto{asemblar}\senco{1}{Juntigar provizore plura personi ad un loko por ula skopo komuna}\senco{2}{Kunigar plura kozi por ajustar li}\lingui{EFI}
\artiklo{asentar}%
\vorto{asentar}\fako{trans.} Adoptar la koncepto-maniero quan onu aprobas. (\textit{asento} koncernas opinioni pure teoriala; \textit{konsento} relatas prefere opinioni praktikala, t.e. propozi e rezolvi)\lingui{DEFIS}
\artiklo{asepta}%
\vorto{asepta} Qua ne kontenas mikrobi\lingui{DEFIRS}
\artiklo{asesoro}%
\vorto{asesoro} La persono qua adjuntesas a judiciisto precipua\lingui{DEFIRS}
\artiklo{asfalto}%
\vorto{asfalto} Bitumo solida, brile-nigra, renkontrebla surface di ula lagi, precipue di la lago Asphaltita\lingui{DEFIRS}
\artiklo{asfixiar}%
\vorto{asfixiar}\fako{trans.} Haltigar la arterio-pulsado; asfixio esas la stando qua sucedas la oxigenizo neperfekta di sango, e finas per morto kande lu tro duras\lingui{DEFIRS}
\artiklo{asfodelo}%
\vorto{asfodelo}\fako{bot.} Genero de planti perena, kun bela flori, dispozita grape\latina{asphodelus}\lingui{DEFIRS}
\artiklo{asidua}%
\vorto{asidua} Qua esas konstante prezenta che ulu, en ula loko, por exekutar ibe sua taski\lingui{EFIS}
\artiklo{asiejar}%
\vorto{asiejar}\fako{trans.} Instalar su kun armeo an-cirkum (urbo, fortifikajo) por agar lo necesa por atako sucesonta, e por sizar (lu)\lingui{EFIS}
\artiklo{asignar}%
\vorto{asignar}\fako{trans.} Sumnar ulu ke lu su prizentez, ye la dio indikata, a la judiciisto\lingui{DEFIRS}
\artiklo{asimilar}%
\vorto{asimilar}\fako{trans.}\senco{1}{\textit{(fiziol.)} Transformar la nutrivi tale ke li divenas identa a sua substanco}\senco{2}{\textit{(metaf.)} Igar suan la idei di altru}\lingui{DEFIRS}
\artiklo{asimptoto}%
\vorto{asimptoto}\fako{geom.} Lineo rekta, ligata konstanta-relate a kurvo di qua lu proximigas su sen atingar lu\lingui{DEFIRS}
\artiklo{asinektika}%
\vorto{asinektika}\fako{matem.) (aludante anguli e figuri geometriala} Ne-homologa\lingui{DEFIS}
\artiklo{asistar}%
\vorto{asistar}\fako{trans.} Esar prezenta ye… e regardar\lingui{DFIRS}
\artiklo{askarido}%
\vorto{askarido}\fako{zool.} Vermo nematoida qua vivas en la intestini di la vertebrozi, precipue di la homo\latina{ascaridae genus : ascaris}\lingui{DEFIRS}
\artiklo{asketo}%
\vorto{asketo}\senco{1}{Persono qua, en la komenco di la Eklezio kristana, retroiris de la socio a loko ube lu esis sola, por ibe konsakrar su ad exerci di pieso ed ad auto-punisi}\senco{2}{Persono qua impozas a su vivo-maniero austera}\lingui{DEFIRS}
\artiklo{asklepiado}%
\vorto{asklepiado} Verso qua konsistas ye un spondeo, du koriambi ed un iambo\lingui{DEFIRS}
\artiklo{askoltar}%
\vorto{askoltar}\fako{trans.} Atencar to quo dicesas; esforcar por audar
\artiklo{asno}%
\vorto{asno}\fako{zool.} Animalo di la genero « solipedi », kun oreli longa, spino salianta, krio desharmonioza (bramo), e quan onu uzas precipue kom kargajo-bestio\latina{asinus}\lingui{EFIS}
\artiklo{aso}%
\vorto{aso} Ludo-karto o ludo-kubo di nur un punto\lingui{DEFIS}
\artiklo{asociar}%
\vorto{asociar}\fako{trans.} Igar partoprenar to quon agas ulu. – Unionar ad ulu kom partoprenonto di ta quon lu agas\lingui{DEFIRS}
\artiklo{asomar}%
\vorto{asomar}\fako{trans.} Ocidar per frapo de ulo tre pezoza\lingui{F}
\artiklo{asonancar}%
\vorto{asonancar}\fako{netrans., kun}\senco{1}{Rivenar \textit{(aludante la sama son-uno)}}\senco{2}{Konsistar ye rimo nekompleta, quan karakterizas nur la identeso di la vokalo acentoza inter du finali maskula e du finali femina. \textit{(Ex. ek la Franca linguo : perte e peste)}}\lingui{DEFIS}
\artiklo{asortar}%
\vorto{asortar}\fako{trans., kun} Kunpozar kozi qui inter-harmonias\lingui{DEFIR}
\artiklo{asparagino}%
\vorto{asparagino} Substanco quan onu renkontras precipue en la yuna sprosi di asparago e qua uzesas medicine kom urinifiva\simbolo{C₈H₈AZ₂O₆}\lingui{DEFIRS}
\artiklo{asparago}%
\vorto{asparago}\fako{bot.} Planto gardenala, di qua la yuna sprosi esas manjebla\latina{asparagus}\lingui{DEFIRS}
\artiklo{aspektar}%
\vorto{aspektar}\fako{netrans.} Vidatesar kun karakterizivi qui igas judikar lu kom… (ed esas ya tala)\lingui{DEFIRS}
\artiklo{aspera}%
\vorto{aspera} Qua shokas desagreable la takto-senso\lingui{EFIS}
\artiklo{aspersar}%
\vorto{aspersar}\fako{trans., ye} Humidigar ye kelka guti\lingui{DEFIRS}
\artiklo{asperulo}%
\vorto{asperulo}\fako{bot.} Genero de rubiacei de la regioni ye temperaturo meza\latina{asperula}\lingui{EFIS}
\artiklo{aspikjeleo}%
\vorto{aspikjeleo} Karno-disho o fisho-disho kun jeleo\lingui{DEFIS}
\artiklo{aspiko}%
\vorto{aspiko}\fako{zool.} Varietato di vipero, serpento venenoza di Egiptia e di Libia, quan karakterizas influro remarkinda di lua kolo\latina{naja haje}\lingui{EFIRS}
\artiklo{aspirar}%
\vorto{aspirar}\fako{trans., per, ad}\senco{1}{Atraktar per inspiro di aero}\senco{2}{Direktar sua deziro ad ulu od ulo}\lingui{DEFIRS}
\artiklo{aspiratoro}%
\vorto{aspiratoro} Nomo di aparati diversa qui havas kom rolo aspirar la fluidi (aero, gaso, vapori, liquidi, e c.) pura o qui kontenas korpi solida suspende, la polvi (segala, limala, e c.)\lingui{DEFIRS}
\artiklo{aspirino}%
\vorto{aspirino} Acido acetil-salicila : CH₃COOC₆H₄COOH. Korpo blanka, solida, fuzo-punto : 133° Celsius. Uzata medicine kom kontre-febro, kontre-doloro\lingui{DEFIRS}
\artiklo{astatika}%
\vorto{astatika}\fako{magnet.} Qua permanas movebla kande deprenata de sua cirkonstanci naturala\lingui{DEFIRS}
\artiklo{astenopa}%
\vorto{astenopa}\fako{patol.} Dicesas pri persono di qua la vidiveso ne esas konstanta. (Komence, la malado vidas bone; balde, lu vidas quaze tra nebulo; fine, lu ne plus povas dicernar la objekti)\lingui{DEFIS}
\artiklo{asterio}%
\vorto{asterio} Zoofito radiera\latina{asterias}\lingui{DEFIRS}
\artiklo{asterismo}%
\vorto{asterismo} Fenomeno optikala : reflekto radiacanta di lumo sur ula lapidi\lingui{DEFIRS}
\artiklo{astero}%
\vorto{astero}\fako{bot.} Planto kompozaja di qua la flori radioza similesas steli\latina{aster}\lingui{DEFIRS}
\artiklo{asteroido}%
\vorto{asteroido}\fako{astron.} Mikra korpo qua cirkulas en cielo\lingui{DEFIRS}
\artiklo{astigmata}%
\vorto{astigmata}\fako{patol.} Di qua la okulo havas kom karakterizivo kurveso-neegalaji en la medii refraktiva\lingui{DEF}
\artiklo{astmo}%
\vorto{astmo}\fako{patol.} Sufoko intermitanta, atribuata a nevroso di la respir-organi\lingui{DEFIRS}
\artiklo{astonar}%
\vorto{astonar}\fako{trans.} Shokar la psiko per ulo extraordinara\lingui{EF}
\artiklo{astragalo}%
\vorto{astragalo}\senco{1}{\textit{(anat.)} Osto di la tarso, di formo kuboida, inkluzita inter la maleoli di la tibio e di la peroneo}\senco{2}{\textit{(arkitekt.)} Muluro pozita sur kolono, inter la fusto e la kapitelo}\lingui{DEFIRS}
\artiklo{astrakano}%
\vorto{astrakano} Felo di qua la pili esas lokloza, devenanta de bovido nasko-mortinta\lingui{DEFIRS}
\artiklo{astriktar}%
\vorto{astriktar}\fako{trans.} Densigar ula tisui organika. (Aluno esas astriktiva)\lingui{DEFIS}
\artiklo{astro}%
\vorto{astro}\fako{astron.} Ciel-korpo\lingui{DEFIS}
\artiklo{astrocentro}%
\vorto{astrocentro}\fako{biol.} Fuzelo akromata di la nukleo\lingui{DEFIRS}
\artiklo{astrofiziko}%
\vorto{astrofiziko} Studio di la astri per la metodi diversa di fiziko\lingui{DEFIRS}
\artiklo{astroida}%
\vorto{astroida} Astro-forma, stelo-forma\lingui{DEF}
\artiklo{astrolabo}%
\vorto{astrolabo} Instrumento uzata por mezurar la alteso-grado di la astri super horizonto\lingui{DEFIRS}
\artiklo{astrologio}%
\vorto{astrologio} Exploro di la astri pri lia influo asertita di la homo-fato\lingui{DEFIRS}
\artiklo{astrologo}%
\vorto{astrologo} La persono qua praktikas astrologio\lingui{DEFIRS}
\artiklo{astronomio}%
\vorto{astronomio} Cienco pri la astri\lingui{DEFIRS}
\artiklo{astronomo}%
\vorto{astronomo} La persono qua konsakras su ad astronomio\lingui{DEFIRS}
\artiklo{astrosfero}%
\vorto{astrosfero}\fako{biol.} Sfero atraktiva\lingui{DEFIRS}
\artiklo{asumar}%
\vorto{asumar}\fako{trans.} Aceptar responsor pri (ulo)\lingui{EFIS}
\artiklo{asunciono}%
\vorto{asunciono}\fako{religio katolika} Festo-dio institucita da la Eklezio katolika, por honorizar la transporto mirakla di la korpo e di la anmo di la Virgino aden cielo\lingui{EFIS}
\artiklo{-at-}%
\vorto{-at-} Sufixo qua indikas la participo prezenta pasiva
\artiklo{atakar}%
\vorto{atakar}\fako{trans., per, ye} Komencar la frapo (di enemiko, di adverso)\lingui{DEFIRS}
\artiklo{ataraxio}%
\vorto{ataraxio} Quieteso di la psiko quan trublas nula deziro, nula timo\lingui{DEFIRS}
\artiklo{atasheo}%
\vorto{atasheo}\fako{diplomac.} La persono qua, en ambasadeyo, esas, sur la skalo hierarkiala, nemediate dop la sekretario\lingui{DEFIR}
\artiklo{atavismo}%
\vorto{atavismo} Herediveso di ula karakteri korpala o mentala qui, en ula kazi, devenas de plura generacioni\lingui{DEFIRS}
\artiklo{ataxio}%
\vorto{ataxio}\fako{patol.} Neregulaleso di la movi di la korpo, produktata da la perturbeso di la funcioni nervala\lingui{DEFIS}
\artiklo{ateismo}%
\vorto{ateismo} Opiniono di la persono qua ne kredas ke existas Deo\lingui{DEFIRS}
\artiklo{ateisto}%
\vorto{ateisto} La persono qua profesas ateismo\lingui{DEFIRS}
\artiklo{ateliero}%
\vorto{ateliero} Loko en qua laboras piktisto, skultisto, fotografisto, *artizano\lingui{DF}
\artiklo{atencar}%
\vorto{atencar}\fako{trans.} Direktar sua tota penso ad\lingui{DEFIS}
\artiklo{atentar}%
\vorto{atentar}\fako{trans.} Krimino-probar\lingui{DEFIS}
\artiklo{ateromo}%
\vorto{ateromo}\fako{patol.} Degenero di la tuniko interna di la arterii, qua implikas lokale la existo di amasi de kolesterino e di kalko-sali\lingui{DEFIRS}
\artiklo{atestar}%
\vorto{atestar}\fako{trans.} Deklarar ke ulo esas exakta, reala\lingui{DEFIRS}
\artiklo{atiko}%
\vorto{atiko} Supera parto di edifico, ordinare ornita per pilastri, e di qua la rolo esas disimular la tekto\lingui{EF}
\artiklo{atingar}%
\vorto{atingar}\fako{trans.} Sucesar tushar to quo distas\lingui{EFIS}
\artiklo{atipa}%
\vorto{atipa}\fako{biol.} Qua esas plu o min fora de la tipo normala\lingui{DEFIS}
\artiklo{atlanto}%
\vorto{atlanto}\fako{arkitekt.} Statuo di viro, nuda e drapirita, qua subtenas entablamento o kornico\lingui{DEFIRS}
\artiklo{atlaso}%
\vorto{atlaso} Kolekturo de mapi, de imaji anatomiala, anexita a libro\lingui{DEFIRS}
\artiklo{atleto}%
\vorto{atleto} La homo qua, en la ludi publika, kombatis pri luktado, e pri boxado (en Grekia antiqua)\lingui{DEFIRS}
\artiklo{atmosfero}%
\vorto{atmosfero} Aero-zono qua cirkondas la Terglobo\lingui{DEFIRS}
\artiklo{atolo}%
\vorto{atolo} Insulo ek koralio\lingui{DFS}
\artiklo{atomo}%
\vorto{atomo} Ta parteto di korpo simpla, quan onu konsideris kom nepartigebla, e qua esas la elemento-quanto maxim mikra apta divenar parto de kombinuro\lingui{DEFIRS}
\artiklo{atonio}%
\vorto{atonio} Ne-existo di vivozeso\lingui{DEFIRS}
\artiklo{-atr-}%
\vorto{-atr-} Sufixo qua indikas simileso
\artiklo{atraktar}%
\vorto{atraktar}\fako{trans.} Kaptar per ruzo-stroko\lingui{DEFS}
\artiklo{atribuar}%
\vorto{atribuar}\fako{trans., ad} Donar (ulo) ad ulu kom lua parto\lingui{DEFIS}
\artiklo{atributo}%
\vorto{atributo}\fako{gram.} Termino di la propoziciono, qua expresas la eso-maniero quan onu afirmas pri la subjekto\lingui{DEFIRS}
\artiklo{atricar}%
\vorto{atricar}\fako{netrans.) (teol.} Repentar pri sua peko sive pro shamo, sive pro timo a puniseso\lingui{DEFIS}
\artiklo{atriumo}%
\vorto{atriumo}\fako{che la Romani antiqua} Korto interna, generale cirkondata da portiko apertata\lingui{DEF}
\artiklo{atrofiar}%
\vorto{atrofiar}\fako{trans.} Igar deperisar parto di la korpo o di organo qua, pro ne funcionar, ne plus alimentesas da la fluo sangala\lingui{DEFIRS}
\artiklo{atropino}%
\vorto{atropino}\fako{kemio} C₁₇H₂₃NO₃. Alkaloido, extraktata de beladono, qua dilatas la pupilo\lingui{DEFIRS}
\artiklo{atuno}%
\vorto{atuno}\fako{zool.} Mar-fisho manjebla, di la genero « scomber »\latina{thynnus}\lingui{DEFIRS}
\artiklo{aturdar}%
\vorto{aturdar}\fako{trans.} Cesir konciar o preske ne plus konciar, konseque (kom exemplo) de frapo, de granda shoko, e c.\lingui{FIS}
\artiklo{aucionar}%
\vorto{aucionar}\fako{trans.} Ofrar preco plu alta kam altru\lingui{DEFRS}
\artiklo{audacar}%
\vorto{audacar}\fako{trans.} Probar en maniero ne-timigebla\lingui{EFIS}
\artiklo{audar}%
\vorto{audar}\fako{trans.} Perceptar per la orel-organaro\lingui{DEFIRS}
\artiklo{audienco}%
\vorto{audienco} Acepto di la viziteri, demanderi\lingui{DEFIRS}
\artiklo{auditoro}%
\vorto{auditoro} Titulo di employato en tribunalo\lingui{DEFIRS}
\artiklo{augias-stablo}%
\vorto{augias-stablo} Loko tre sordida\lingui{DEFIS}
\artiklo{augmentar}%
\vorto{augmentar}\fako{trans.} Igar plu granda, per adjuntar parti\lingui{DEFIS}
\artiklo{augurar}%
\vorto{augurar}\fako{trans.} Igar konsiderar kom esonte bona o mala (pri la eventi)\lingui{DEFIRS}
\artiklo{augusta}%
\vorto{augusta} Titulo honorizala, uzata a la membri di familio suverena\lingui{DEFIRS}
\artiklo{aula}%
\vorto{« aula »}\senco{1}{En la universitati Germana, Usana, Suisa : la chambrego di la akti e di la festi skolala}\senco{2}{Palaco en qua rezidas la kortanaro di rejio}\lingui{DEFIRS}
\artiklo{aureolo}%
\vorto{aureolo}\senco{1}{Cirklo lumala ye qua la piktisti cirkondas la kapo di Iesu-Kristo, di la Virgino e di la santi}\senco{2}{Omna cirklo lumala o koloroza quan la okulo vidas cirkum objekto. \textit{(anke metaf.)}}\lingui{DEFIRS}
\artiklo{aurikulo}%
\vorto{aurikulo} Singla de la du supra kavaji di la kordio\lingui{DEFIS}
\artiklo{auroro}%
\vorto{auroro} Koloro oranjea qua venas de la radii di Suno levanta, qui frapas horizonto. – \textit{Metaf}. : La komenco di ulo\lingui{DEFIRS}
\artiklo{auskultar}%
\vorto{auskultar}\fako{trans.) (medic.} Askoltar, per pozar la orelo an la parti vicina di la kordio, di la pulmoni, e c., la bruisi normala od anormala qui eventas en ta organi\lingui{DEFIRS}
\artiklo{auspico}%
\vorto{auspico}\fako{en Roma antiqua} Predico ek la flugo od ek la kanto di la uceli, quan onu agis jus ante komencar agado, entraprezo publika. \textit{(anke metaf.)}\lingui{DEFIS}
\artiklo{austera}%
\vorto{austera} Tre severa por su pri la mori, kun nuanco di rigoro ed asketeso\lingui{DEFIS}
\artiklo{austro}%
\vorto{austro} Vento de la regioni suda\lingui{DEFIRS}
\artiklo{autentika}%
\vorto{autentika}\senco{1}{Di qua la exakteso, la vereso, ne esas kontestebla}\senco{2}{Di qua la exakteso, la vereso, garantiesas per akto statala, municipala, per formalajo legala}\lingui{DEFIRS}
\artiklo{auto-}%
\vorto{auto-} Prefixo rezervenda a la terminaro ciencala e teknikala, e di qua la ideo povas tradukesar per « su ». \textit{(ex. : su-defenso)} o per « \textit{propra} » \textit{(ex. propr-amo)}\lingui{DEFIRS}
\artiklo{autobiografio}%
\vorto{autobiografio} Biografio di persono, kompozita da lu ipsa\lingui{DEFIRS}
\artiklo{autoblasto}%
\vorto{autoblasto}\fako{bot.} Organismo vejetala maxim simpla\lingui{DEFIS}
\artiklo{autobuso}%
\vorto{autobuso} Omnibuso movata per motoro vice per kavali\lingui{DEFIRS}
\artiklo{autodafe}%
\vorto{« autodafe »} Lekto ed exekuto publika di la sentenco per qua Inquiziciono, en Hispania, kondamnis la hereziani a la suplicieso per kombusto\lingui{DEFIRS}
\artiklo{autodidakto}%
\vorto{autodidakto} Persono qua obtenis sua savaji ipse\lingui{DEFIRS}
\artiklo{autogama}%
\vorto{autogama}\fako{biol.) (aludante ula protisti} Qua riproduktas su sen divido di la celulo. (Nur la nuklei diferencias su, e pleas la rolo di gameti)\lingui{DEFIRS}
\artiklo{autografo}%
\vorto{autografo}\fako{aludante letro, libro} Skribita da la autoro propra-manue\lingui{DEFIRS}
\artiklo{autoklavo}%
\vorto{autoklavo}\senco{1}{Dicesas pri klozilo di tanko qua kontenas ulo presata, e qua mantenesas klozata da la ipsa presajo qua efikas de interne}\senco{2}{Tanko klozata per kovrilo mantenata da la preso interna, uzata por steriligar la substanci enpozita}\lingui{DEFIRS}
\artiklo{autokrato}%
\vorto{autokrato} Suvereno di qua la povo esas suprega e nekontrolata\lingui{DEFIRS}
\artiklo{autolizo}%
\vorto{autolizo} Digesteso spontana di la tisui vejetala od animalala\lingui{DEFIRS}
\artiklo{automata}%
\vorto{automata} Qua semble movas su, pro la efiko da mekanismo celata\lingui{DEFIRS}
\artiklo{automobilo}%
\vorto{automobilo} Veturo movata per motoro qua alimentesas da explozosubstanci (kerozeno, benzino, e c.)\lingui{DEFIRS}
\artiklo{automorfa}%
\vorto{automorfa} Dicesas pri la minerali di la roki qui limitizesas da formi kristalatra e plana\lingui{DEFIRS}
\artiklo{automorfismo}%
\vorto{automorfismo}\fako{biol.} Rejimo di automorfoso\lingui{DEFIS}
\artiklo{automorfoso}%
\vorto{automorfoso}\fako{biol.} Morfoso kompozata da faktori interna, che planto\lingui{DEFIS}
\artiklo{autonoma}%
\vorto{autonoma}\senco{1}{Qua regnesas da sua legi propra}\senco{2}{\textit{(biol.)} Dicesas pri la formaco di ovo quan genitis du gameti sama-rasa}\lingui{DEFIRS}
\artiklo{autopsiar}%
\vorto{autopsiar}\fako{trans.} Studiar anatomiale parti interna di kadavro, por informar la judiciistaro o la ciencisti\lingui{DEFIRS}
\artiklo{autoritato}%
\vorto{autoritato}\senco{1}{Yuro, povo, darfo pri koaktar ad obedio}\senco{2}{\textit{(metaf.)} Fakultato impozar su a la opiniono kunhomala per sua merito}\lingui{DEFIRS}
\artiklo{autoro}%
\vorto{autoro}\senco{1}{La kauzo prima di ulo}\senco{2}{La persono qua kompozis verko pensala}\lingui{I}
\artiklo{autosporo}%
\vorto{autosporo}\fako{biol.} Sporo adulteskinta interne di la sporangio, o jermifanta, che algi\lingui{DEFIS}
\artiklo{autotipar}%
\vorto{autotipar}\fako{trans.} Fotografar sur plako metala, lumo-sentiva, qua pose gravesos korode per acidi\lingui{DEFS}
\artiklo{autotomio}%
\vorto{autotomio}\fako{biol.} Rupturo reflexa di parto korpala, ex. che akrido, lacerto\lingui{DEFIS}
\artiklo{autuno}%
\vorto{autuno} Sezono qua sequas somero e pre-iras vintro\lingui{EFIS}
\artiklo{auxosporo}%
\vorto{auxosporo}\fako{biol.} Ovo qua rezultas de la kopulaco, che diatomei\lingui{DEFIS}
\artiklo{avalanchar}%
\vorto{avalanchar}\fako{netrans.) (aludante blokego de nivo-floki} Separar su de monto e precipitar su aden la valo\lingui{EFIS}
\artiklo{avalo}%
\vorto{avalo}\fako{banko e komerco} Garantio quan onu donas pri la pagoteso, da ula kunhomo, di pago-promisilo, per skribar infre di olca « Aceptate avale » e signatar ca formulo\lingui{DEFIS}
\artiklo{avan}%
\vorto{avan} Prepoziciono qua relatas la loko, la plaso, quan okupas, en spaco, la enti o kozi, kontre ke \textit{ante} relatas tempo. « Avan » indikas : situita en la areo quan povas vidar persono\lingui{DEFIRS}
\artiklo{avancar}%
\vorto{avancar}\fako{netrans.} Lokizar su plu fore ad-avan. \textit{(trans.)} Lokizar plu fore ad-avan\lingui{DEFIS}
\artiklo{avantajo}%
\vorto{avantajo} To quo donas ula superioreso generale\lingui{DEFIS}
\artiklo{avantalo}%
\vorto{avantalo}\senco{1}{Peco de stofo, de ledro, quan onu ligas sur la parto avana di la korpo, super sua vesti, por evitar la domajeso di olci kande onu agas o facas ulo sordidigiva}\senco{2}{Peco de ledro, ligita sur la parto avana di veturo sen-tekta, por protektar de la pluv-aquo, de la spricizuri}
\artiklo{avara}%
\vorto{avara} Qua prizas pekunio por la plezuro de akumular olu\lingui{EFIS}
\artiklo{avariar}%
\vorto{avariar}\fako{trans.} Domajar, ne nur navo, ma anke kozi transportata per navo, fervoyo, veturo\lingui{DEFIRS}
\artiklo{avelano}%
\vorto{avelano}\fako{bot.} Frukto kun shelo un-lojia, ne-dehiscenta, inkastrita per involukro\latina{corylus evellana}\lingui{FISL}
\artiklo{aveno}%
\vorto{aveno}\fako{bot.} Planto cereala nutrala, ek la familio « graminei »\latina{avena}\lingui{EFISL}
\artiklo{aventurar}%
\vorto{aventurar}\fako{netrans.} Irar renkontre di ulo di qua la rezulto ne esas pre-savebla\lingui{DEFIRS}
\artiklo{avenuo}%
\vorto{avenuo} Larja voyo, ordinare arbor-lineoza, per qua onu abutas a palaco, a kastelo, a placo, e c.\lingui{DEFIRS}
\artiklo{averso}%
\vorto{averso} La facio di medalio, di moneto-peco : \textit{(pri la pagini di kayero, di libro, e c.)} Singla ek la pagini kun numero ne-para\lingui{DF}
\artiklo{avertar}%
\vorto{avertar}\fako{trans., pri} Informar (ulu) pri ulo futura, bona o mala\lingui{DEFIS}
\artiklo{aviacar}%
\vorto{aviacar}\fako{netrans.) (aludante homo} Cirkular en atmosfero, per moyeno altra kam aerostato od aeronavo\lingui{DEFIRS}
\artiklo{avida}%
\vorto{avida} Qua deziras ulo ne-represeble\lingui{DEFIS}
\artiklo{avizar}%
\vorto{avizar}\fako{trans., pri} Donar informo praktikala, konsilo od impero\lingui{DFIRS}
\artiklo{avo}%
\vorto{avo} La patro di la patro di homo\lingui{ISL}
\artiklo{avoué}%
\vorto{« avoué »} En Francia : oficiero judiciala, institucita dum la 1789-revoluciono por remplasar la prokuratori, e di qua la rolo konsistas ye okupar su pri la proceduro judiciala e pri la inquesti di la procesi civila, ma sen pledar (la pledo resortisas a la advokati)\lingui{F}
\artiklo{axelo}%
\vorto{axelo}\fako{anat.} Korpo-parto qua esas sub la junturo di la brakio kun la shultro\lingui{DEFIS}
\artiklo{axiomo}%
\vorto{axiomo} Verajo generala, nedemonstrebla, qua impozas su a la spirito kom evidentajo\lingui{DEFIRS}
\artiklo{axo}%
\vorto{axo}\senco{1}{\textit{(geom.)} Lineo matematikala cirkum qua rotacas figuro plana, por produktar rotaco-solido}\senco{2}{\textit{(astron.)} Lineo cirkum qua eventas la rotaco di korpo cielala}\senco{3}{\textit{(mekan.)} Omna lineo cirkum qua korpo povas movar su}\lingui{DEFIS}
\artiklo{axoido}%
\vorto{axoido}\fako{anat.} Vertebro duesma di la kolo, qua esas insertita ye la unesma\lingui{DEFIS}
\artiklo{axometro}%
\vorto{axometro} Indexo qua sequas la movi di la guvernilo-stango di navo ed indikas lua direciono\lingui{DF}
\artiklo{axonometrio}%
\vorto{axonometrio} Projekturo ortogonala sur plano obliqua ye la tri dimensioni di la korpo reprezentenda. La axi di la perspektivo esas la projekturi di la edri di la triedro tri-rektangula quan determinas la tri dimensioni\lingui{DEFIS}
\artiklo{aye!}%
\vorto{aye!} Interjeciono; kurta klameto de doloro
\artiklo{azaleo}%
\vorto{azaleo}\fako{bot.} Planto ek la familio « erikacei », kultivata kom apartamento-planto\latina{azalea}\lingui{DEFIRS}
\artiklo{azigosporo}%
\vorto{azigosporo}\fako{biol.} Gameto kreskinta, ma qua permanas en la stando di simpla sporo\lingui{DEFIS}
\artiklo{azilo}%
\vorto{azilo}\senco{1}{Loko konsiderata kom neviolacebla, ube la krimininti darfis refujar}\senco{2}{Instituto di karitato, en qua onu gardas la puereti dum ke lia gepatri okupas su pri la taski di sua mestiero}\lingui{DEFIS}
\artiklo{azimo}%
\vorto{azimo} Pano sen-hefa\lingui{DEFIS}
\artiklo{azimuto}%
\vorto{azimuto} Angulo formacata da la plano meridiana di loko, kun la plano vertikala qua pasas tra punto quan onu konsideras\lingui{DEFIRS}
\artiklo{azoto}%
\vorto{azoto}\fako{kemio} Korpo simpla, gasa, sen odoro, sen saporo, neapta por respiro o kombusto, qua esas elemento di aero, ed esas un ek la elementi fundamentala di la tisui vejetala od organika\simbolo{Az (N)}\lingui{DEFIRS}
\artiklo{azuro}%
\vorto{azuro} Koloro ciel-blua\lingui{DEFIRS}
\end{multicols}
\sekciono{B}
\begin{multicols}{2}
\artiklo{b}%
\vorto{b}\fako{muziko} La noto sepesma di la « do » gamo
\artiklo{ba!}%
\vorto{ba!} \textit{(interjeciono)} Expresas la ne-sucio pri la konsequaji de ulo, la quaza desprizo\lingui{DEFRS}
\artiklo{baba}%
\vorto{« baba »} Kuko qua kontenas sukro, korinto e rumo\lingui{F}
\artiklo{babilar}%
\vorto{babilar}\fako{netrans.} Parolar ne-modere e pri frivolaji\lingui{DEF}
\artiklo{babucho}%
\vorto{babucho} Pantoflo de ledro, sen la una o dua ledro-peci ye qui konsistas la talono, e kun suolo plata\lingui{DEFIRS}
\artiklo{babuino}%
\vorto{babuino}\fako{zool.} Sorto di simio cinocefala, kun labii salianta\latina{Cynocephalus babuin}\lingui{DEFIRS}
\artiklo{bachelero}%
\vorto{bachelero} La grado universitatala unesma (minim alta)\lingui{DEFIRS}
\artiklo{bacilo}%
\vorto{bacilo}\fako{biol.} Mikrobo di naturo vejetala, longa, bastonatra\lingui{DEFIRS}
\artiklo{baderno}%
\vorto{baderno} Treso de olda kordegi navala, uzata por garnisar la masti, la yardi, e c., por prezervar li de la fricionado, o por impedar ke la kargajo renversesos da la navo-rulado\lingui{FIS}
\artiklo{badiano}%
\vorto{badiano}\fako{bot.} Planto ek la familio « magnoliacei », di qua la frukto esas aromata\latina{Illicium anisatum}\lingui{DEFRS}
\artiklo{badijonar}%
\vorto{badijonar}\fako{trans.} Indutar (parto malada) ye medikamento liquida\lingui{EF}
\artiklo{bagajo}%
\vorto{bagajo} To quon onu kunportas dum voyajo, expediciono\lingui{DEFIRS}
\artiklo{bagatelo}%
\vorto{bagatelo} Objekto qua valoras tre poke\lingui{DEFIS}
\artiklo{bagno}%
\vorto{bagno} Loko ube onu retenis la forcati en ula mar-portui, pos la supreso di la galeri\lingui{DEFIS}
\artiklo{bakanalo}%
\vorto{bakanalo} Orgio bruisoza\lingui{DEFIRS}
\artiklo{bakar}%
\vorto{bakar}\fako{trans., en} Igar apta por nutro, per la efiko da fairo, en forno\lingui{DE}
\artiklo{bakterio}%
\vorto{bakterio}\fako{biol.} Mikrobo di naturo vejetala, moviva o nemoviva. longa\lingui{DEFIRS}
\artiklo{bakteriologio}%
\vorto{bakteriologio} La parto di la biologio, qua havas kom temo di studio la bakterii\lingui{DEFIRS}
\artiklo{bakteriologo}%
\vorto{bakteriologo} Ciencisto pri bakteriologio\lingui{DEFIRS}
\artiklo{balado}%
\vorto{balado} Mikra poemo ek versi inter-egala, qua konsistas ek tri strofi simetra ed un kupleto uzata kom homajo-versi\lingui{DEFIRS}
\artiklo{balanciero}%
\vorto{balanciero} Maso (generale roto-forma) qua ocilas cirkum sua baricentro per resorto\lingui{DEFIRS}
\artiklo{balanco}%
\vorto{balanco} Instrumento por ponderar korpo, equilibriganta olca per maso de metalo uzata kom pez-unajo\lingui{DEFIRS}
\artiklo{balas-rubino}%
\vorto{balas-rubino} Lapido ek aluminato di magnezio, varietato rozea\lingui{DEFIRS}
\artiklo{balasto}%
\vorto{balasto} Grosa sabl-uni o fer-skorio, uzata : 1. en la marnavigado, por igar la navi plu equilibroza, per abasar la baricentro; 2. sur la relvoyi, por kovrar la traversi qui suportas la reli\lingui{DEFRS}
\artiklo{balayar}%
\vorto{balayar}\fako{trans., per} Pulsar de avan su, sur sulo – polvuni, rezidui, aquo, nivo-floki, e c., per fasko de mikra stipi o per brosilo longa-krina, fixigita a mancho longa
\artiklo{balbutar}%
\vorto{balbutar}\fako{trans.} Artikular la vorti en maniero poke dicernebla\lingui{FIS}
\artiklo{baldakino}%
\vorto{baldakino}\senco{1}{Doselo, subtenata da koloni, e garnisita ye tapeti, qua aspektas quale krono super la altaro, en la kirki}\senco{2}{Lito-doselo de ube pendas kurteni}\lingui{DEFIRS}
\artiklo{balde}%
\vorto{balde} Poka tempo pose\lingui{D}
\artiklo{baldrio}%
\vorto{baldrio} Larja bendo de ledro o de stofo qua, suportante la espado o la sabro, portesas da la persono diagonale, de la shultro dextra a la hancho sinistra\lingui{DEFI}
\artiklo{baleno}%
\vorto{baleno}\fako{zool.} Granda mamifero, ek la klaso « cetacei », di qua la boko sen-denta, esas garnisita – ye la du lateri, an la maxilo ye granda lami dina, flexebla (barti)\latina{balaena}\lingui{FISL}
\artiklo{baleto}%
\vorto{baleto} Danso alegoria, exekutata sive en festo, sive en teatro, kom intermezo di opero, di komedio\lingui{DEFIRS}
\artiklo{balifo}%
\vorto{balifo} Oficiro nobela o yurala qua, en Francia, olim judiciis ye la nomo di la rejo o di sinioro\lingui{EFIS}
\artiklo{baliso}%
\vorto{baliso}\fako{navig.} Signalilo por guidar la naviganto en paseyo danjeroza\lingui{FS}
\artiklo{balistiko}%
\vorto{balistiko} Cienco pri la rapideso-grado e la direciono di la projektili\lingui{DEFIRS}
\artiklo{balisto}%
\vorto{balisto} Milito-mashino quan la antiqui uzis por lansar stonegi tre pezoza\lingui{DEFIRS}
\artiklo{balkono}%
\vorto{balkono}\senco{1}{Platformo plu o min larja qua salias ek la fasado di edifico, ye la sama nivelo kam un o plura fenestri di etajo, e qua komunikas kun la internajo per un o plura aperturi}\senco{2}{Galerio en la dopa parto di granda navo, kom ornivo o komodajo}\senco{3}{Parto di la galerio unesma qua, en teatro, apudesas la lojii avan-cena}\lingui{DEFIRS}
\artiklo{balnar}%
\vorto{balnar}\senco{1}{\textit{(trans.)} Pozar e mantenar en aquo (la korpo o parto di la korpo) por netigar, koldigar, e c.}\senco{2}{\textit{(netrans.)} Humidigar abunde}\lingui{DEFIRS}
\artiklo{balo}%
\vorto{balo} Asemblitaro grandanombra de personi qui dansas\lingui{DEFIRS}
\artiklo{balono}%
\vorto{balono} Globo o bulo kava, mola, plena ye aero (kompresita) e, konseque, elastika, por ludar\lingui{DEFIRS}
\artiklo{balotar}%
\vorto{balotar}\fako{trans.} Votar duesma-foye por selektar definitive ek la kandidati qui, lor elekto, preske atingis la majoritato\lingui{DEFIRS}
\artiklo{balotear}%
\vorto{balotear}\fako{netrans.} Movar su, ocilante o tremante \textit{(aludante maso solida, ma elastika o viskoza, quale gelatino)}\lingui{FIS}
\artiklo{balustrado}%
\vorto{balustrado} Rango formacita ek balustri, generale horizontala\lingui{DEFIRS}
\artiklo{balustro}%
\vorto{balustro} Mikra pilastro fasonita\lingui{DEFIRS}
\artiklo{balzamino}%
\vorto{balzamino} Planto gardenala, di qua la semini kontenesas en kapsulo qua desklozesas (quale movata da resorto) kande onu tushas olu\latina{impatiens balsamina}\lingui{DEFIRS}
\artiklo{balzamo}%
\vorto{balzamo}\fako{anke metaf.}\senco{1}{Substanco odoroza, rezinatra, qua fluas de ula vejetanti en formo di liquido}\senco{2}{Nomo di preparaji farmaciala diversa, oleoza e balzamoza}\lingui{DEFIRS}
\artiklo{bambuo}%
\vorto{bambuo}\fako{bot.} Planto gramina di qua la kortico esas plektebla aden korbi, mati, e c., e la yuna stipi, transformebla a mobli, promeno-bastoni, e c.\latina{bambusa}\lingui{DEFIRS}
\artiklo{banano}%
\vorto{banano}\fako{bot.} Frukto de planto herbacea, ek la familio « musacei », veninta de India, di qua la stipo esas dika, kona, ed havas quaza krono ek folii larja, de la centro di qua naskas beri saporoza, dispozita grape\latina{musa sapientum}\lingui{DEFIRS}
\artiklo{bandajo}%
\vorto{bandajo}\senco{1}{Aparato qua uzesas por fixigar pansilo, mantenar organo qua tendencas eskartar de lua plaso normala}\senco{2}{Sistemo de bendi fera o stala, qua uzesas por mantenar la felgi di roto, la mulduri di giserio}\lingui{DEFIRS}
\artiklo{bandear}%
\vorto{bandear}\fako{netrans.) (pri navo} Inklinar su adflanke\lingui{FIS}
\artiklo{banderolo}%
\vorto{banderolo} Bendo longa ek stofo qua flotacas en la aero cirkondanta; finas ofte per du pinti; en pikturi, portas surskriburi\lingui{DEFIRS}
\artiklo{bandito}%
\vorto{bandito} Deliktanto (kriminanto) proskriptita qua furtas, asasinas sur la granda voyi\lingui{DEFIRS}
\artiklo{bando}%
\vorto{bando} Asemblitaro de personi qui kuniras ed agas interkonsente\lingui{DEFIRS}
\artiklo{bandoliero}%
\vorto{bandoliero} Rimeno, o sharpo, pozita sur un shultro e qua suportas ula kozo (ex. : espado) sur la altra latero, do pasante oblique sur la korpo di la persono\lingui{DEFIRS}
\artiklo{banero}%
\vorto{banero} Speco di insigno rektangula (ek stofo) fixigita per latero horizontala a stango di sama longeso, qua pendas per du kordi de la stango per qua onu portas la insigno\lingui{DEFIS}
\artiklo{banko}%
\vorto{banko} Firmo qua komercas pri pekunio, indosas e transmisas la kredit-akti ye pago-dato konvencionita, pag-aceptita, diskontas olci, plufaciligas (per trati) la kambio di moneto inter la landi, recevas la pekunio-depozaji, e c.\lingui{DEFIRS}
\artiklo{bankrotar}%
\vorto{bankrotar}\fako{netrans.) (pri komercisto} Faliar – falio qua punisesas da la lego pro eventir lua-kulpe\lingui{DEFIRS}
\artiklo{baobabo}%
\vorto{baobabo}\fako{bot.} Arboro di Afrika, di qua la grandeso atingas giganteso\latina{adansonia digitata}\lingui{DEFIRS}
\artiklo{baptar}%
\vorto{baptar}\fako{trans.) (en la religio kristana} Igar ulu Kristano per la sakramento destinita ad efacar la peko originala, e qua konsistas ye varsar aquo sur la kapo, pronuncante vorti sakramentala\lingui{EFIS}
\artiklo{barakano}%
\vorto{barakano} Dika stofo de lano, texita tre dense, qua fabrikesis en la XVIII-esma yarcento, precipue en la nordo-regiono di Francia\lingui{DEFIRS}
\artiklo{barako}%
\vorto{barako} Konstrukturo de planki, por shirmar peskisti, muton-gardisti, komercisti en aero publika, e c.\lingui{DEFIRS}
\artiklo{baraktar}%
\vorto{baraktar}\fako{netrans.} Luktar per esforcar por su desintrikar\lingui{R}
\artiklo{barar}%
\vorto{barar}\fako{trans.} Obstruktar, klozar per stango.\lingui{DEFIRS}\subvorto{}{baro-fluxo}{Granda bar-ondo qua avancas quaza-monto en ula estuarii}{}{}
\artiklo{baratrio}%
\vorto{baratrio} Fraudo da la navestro kontre armatori od asekuristi\lingui{DEFIRS}
\artiklo{barbakano}%
\vorto{barbakano} Fenduro facita, por la fluo de la pluv-aquo, en muro qua subtenas tero\lingui{DEFIS}
\artiklo{barbara}%
\vorto{barbara} Stranjera pri civilizeso\lingui{DEFIRS}
\artiklo{barbelo}%
\vorto{barbelo}\fako{zool.} River-fisho, ek la familio « ciprinoidi », di qua la maxilo esas garnisita ye filamento simila a ta qua existas cirkum la boko di preske omna insekti – e di qua la karno nutro-prizesas\lingui{DEFIRS}
\artiklo{barbo}%
\vorto{barbo} Pili qui garnisas la mentono, la vangi, la labio supera di adultulo. \textit{(anke metaf.)}\lingui{DEFIRS}
\artiklo{barbotar}%
\vorto{barbotar}\fako{netrans.} Agitar su en aquo, igante lu spricar, trublante lu\lingui{FIS}
\artiklo{barbuliar}%
\vorto{barbuliar}\fako{netrans.} Parolar tro rapide e, konseque, neklare e konfuze, nekompreneble\lingui{FIS}
\artiklo{bardano}%
\vorto{bardano}\fako{bot.} Planto ek la familio « kompozaji », di qua la radiko uzesis olim medicinale kom sudorifigivo, e la (larja) folii kom medikamento kontre tinio\latina{lappa}\lingui{FIS}
\artiklo{bardo}%
\vorto{bardo} Poeto Keltika qua kompozis e recitis kansoni militala, religiala\lingui{DEFIS}
\artiklo{barejo}%
\vorto{barejo} Stofo de lano lejera, texita nedense\lingui{DEFIRS}
\artiklo{barelo}%
\vorto{barelo} Granda vazo de ligno, cilindra, inflita en la mezo, konsistante ye duvi asemblita e cirkondita da ringi fera o ligna, klozita per du fundi ligna, uzata por transportar liquidi ed ula vari altra\lingui{DEFIS}
\artiklo{baricentro}%
\vorto{baricentro} Punto di korpo per qua pasas konstante la rezultanto di la gravito-forci qui efikas ad omna ek lua molekuli\lingui{FIS}
\artiklo{barikado}%
\vorto{barikado} Fortifikuro improvizita, ek pavi amasigita, arbori abatita, veturi renversita, e c., por impedar la aceso di strado\lingui{DEFIRS}
\artiklo{bario}%
\vorto{bario}\fako{kemio} Korpo simpla, metala, flavatre blanka e brilanta, tre oxidigebla, qua deskompozas la aquo ye la temperaturo normala\simbolo{Ba}\lingui{DEFIS}
\artiklo{barito}%
\vorto{barito}\fako{kemio} Oxido di bario, substanco alkalioza venenoza, saporanta korodive, quarople plu densa kam aero\simbolo{BaO}
\artiklo{baritono}%
\vorto{baritono}\senco{1}{Muzik-instrumento antiqua, quaze viulo, qua havas, sur sua mancho, dek-e-sis kordi de latuno quin onu pinchis per la fingri, ultre sep tripo-kordi a qui onu efikis per arketo}\senco{2}{Homo-voco, inter tenoro e basvoco}\lingui{DEFIRS}
\artiklo{barkarolo}%
\vorto{barkarolo} Kansono ritmoza di la gondolisti di la urbo Venezia (Italia)\lingui{DEFIRS}
\artiklo{barkaso}%
\vorto{barkaso} Granda barko uzata por la desembarko di la pasajeri e di la vari kande ne existas portui ube la navi povas an-kayeskar\lingui{DFIRS}
\artiklo{barko}%
\vorto{barko} Mikra kanoto, sen-masta, movebla per remili, e generale ne-uzata en aquo marala\lingui{DEFIRS}
\artiklo{barografo}%
\vorto{barografo} Barometro enrejistrera, konstruktita specale por donar, sur grafiko kontinua, la kurvo di la altitudi atingita da aernavigero, aviacisto\lingui{DEFIRS}
\artiklo{baroka}%
\vorto{baroka} Qua aspektas nereguloza o stranja (ex. : perlo, stilo)\lingui{DEFIS}
\artiklo{barometro}%
\vorto{barometro} Instrumento fizikala qua uzesas : I. por mezurar la pezo di aero, per recipiento apertita, plena ye merkurio qua, presata da atmosfero, su elevas plu o min alte en tubo gradizita, aero-vakua; II. por mezurar la altitudi, pro ke la aero-kolono deskreskas altesale, e ke desdenseskas atmosfero segun ke onu acensas; III. por prognozar vetero segun la denseso-grado di atmosfero\lingui{DEFIRS}
\artiklo{barono}%
\vorto{barono} En la hierarkio di la tituli di nobeleso (grantita da la suvereno, od heredala), olta qua venas pos la viskomto ed ante la *chevaliere\lingui{DEFIRS}
\artiklo{baroskopo}%
\vorto{baroskopo} Aparato por konstatar ke korpo pozita en aero perdas parto de sua pezo, egala a la pezo di la aero-volumino quan olu diplasas\lingui{DEFIRS}
\artiklo{barsho}%
\vorto{barsho}\fako{zool.} Fisho qua tre similesas perko\latina{labrax lupus}\lingui{DFI}
\artiklo{barto}%
\vorto{barto} Omna ek la lami korno-harda qui garnisas transverse la maxilo di la baleno e di ula cetacei\lingui{D}
\artiklo{basa}%
\vorto{basa}\senco{1}{Qua ne atingas la nivelo ordinara}\senco{2}{Qua ne atingas la nivelo di altra objekto}\lingui{F}
\artiklo{basamento}%
\vorto{basamento} Infra parto (bazo), videbla, di edifico, qua jacas sur la fundamento\lingui{DEFIS}
\artiklo{baseno}%
\vorto{baseno}\senco{1}{Recipiento portebla, cirklatra, di qua la bordi esas alta, destinita a kontenar aquo}\subvorto{}{mar-baseno}{La regiono quan trafluas omna fluvii qui adiras ta maro}{2}{}\subvorto{}{fluvio-baseno}{La teritorio quan trairas fluvio kun omna lua riveri}{3}{}\senco{4}{Geologio-baseno : la teritorio ek masivo komuna de strati, de tereni geologiala}\lingui{DEFIRS}
\artiklo{basko}%
\vorto{basko} Parto de stofo taliita, qua pendas cirkum la zono (infre di la korsajo di la muliero-robi, o di la viro-fraki.)\lingui{EFS}
\artiklo{baskular}%
\vorto{baskular}\fako{netrans.} Falar de ta a ca alteso-punto, influe da pezo o forco, e retrovenar a la alteso-grado antea se altra (kontrea) pezo o forco ne efikas\lingui{DEFS}
\artiklo{basono}%
\vorto{basono}\fako{muziko} Sufl-instrumento muzikala ek ligno, kun langeto qua esas quaze la basa voco di la hoboyo\lingui{DEFIS}
\artiklo{basoto}%
\vorto{basoto}\fako{zool.} Hundo kun gambi tre kurta e tordita\lingui{EFIS}
\artiklo{bastardo}%
\vorto{bastardo}\senco{1}{Qua naskis de patrulo e de patrino ne unionita per mariajo}\senco{2}{Qua ne esas raso-pura, speco-pura}\lingui{DEFIRS}
\artiklo{bastingo}%
\vorto{bastingo}\fako{navig.} Garnisuro cirkum ferdeko, ek reti futerizita ye dika tolo, por stivar la hamaki e por amortisar la projektili lor kombato\lingui{FI}
\artiklo{bastiono}%
\vorto{bastiono} Masivo de tero, ordinare kovrita per gazono o per masonuro, konstruktita exter la anguli salianta di la remparo-cirkuito di fortreso por protektar lu\lingui{DEFIRS}
\artiklo{basto}%
\vorto{basto} Parto interna di la kortico di arboro, konsistante ye la strati recenta maxim vicina a la alburno, e dispozita ye dina folieti, ita sur ica\lingui{DE}
\artiklo{bastono}%
\vorto{bastono} Longa peco de ligno, generale prenita de arboro-brancho o de dina stipo, ed uzata kom su-apogilo lor marcho\lingui{DEFIRS}
\artiklo{bataliar}%
\vorto{bataliar}\fako{netrans.} Kombatar (aludante individui, ma ne armeo)\lingui{DEFIRS}
\artiklo{bataliono}%
\vorto{bataliono} Trupo de batalieri. – en la armei moderna : korpo de infantriani, qua subdividesas a plura kompanii\lingui{DEFIRS}
\artiklo{batar}%
\vorto{batar}\fako{trans.} Frapar (ulu od ulo) plura-foye, intersucedante, generale per utensilo : ad ulu, por dolorigar lu; ad ulo, por netigar lu (batar stuli, tapisi). \textit{(Bato esas do esence ago volata e duranta, sempre korpala)}\lingui{DEFIS}
\artiklo{batelo}%
\vorto{batelo} Nomo generala, donita a naveti omnaspeca, ma precipue a naveti de dimensioni pasable granda, quale gabari, e c.\lingui{DFIS}
\artiklo{baterio}%
\vorto{baterio}\senco{1}{\textit{(milit.)} Asemblajo de plura kanoni por destruktar poziciono enemika}\senco{2}{\textit{(elektro)} Asemblajo de plura elektro-kondensatori por produktar granda elektro-descharji}\lingui{DEFIRS}
\artiklo{batisto}%
\vorto{batisto} Dina stofo, extreme dina\lingui{DEFIRS}
\artiklo{batrako}%
\vorto{batrako}\fako{zool.} Amfibio, animalo sango-kolda, di qua la tipo esas la rano, e qua konstitucas la klaso quaresma di la vertebrozi\lingui{DEFIS}
\artiklo{bavar}%
\vorto{bavar}\fako{netrans.} Lasar fluar la salivo ek la boko\lingui{FIS}
\artiklo{bayadero}%
\vorto{bayadero} Dansistino Indiana\lingui{DEFIRS}
\artiklo{bayo}%
\vorto{bayo} Mikra golfo di qua la aperturo esas streta\lingui{DEFIRS}
\artiklo{bayoneto}%
\vorto{bayoneto} Lamo de stalo qua, ajustebla segunvole ye la pinto di la fusilo, posibligas uzar olca kom armo di la tipo « sabro, espado »\lingui{DEFIRS}
\artiklo{bazalto}%
\vorto{bazalto}\fako{geol.} Roko nigratra, de origino fairala, qua ofte prizentas fragmenti prisma inter-paralela\lingui{DEFIRS}
\artiklo{bazaro}%
\vorto{bazaro}\senco{1}{Merkato publika, che la Orientani}\senco{2}{Edifico plur-etajo qua kontenas butiki diversa}\senco{3}{Loko tektoza, galerio ube vendesas objekti omna-sorta}\lingui{DEFIRS}
\artiklo{baziliko}%
\vorto{baziliko}\senco{1}{Che la Romani antiqua : edifico rektangula, dividita aden plura navi da rangi de koloni – uzita kom *seanceyo di la tribunali o kom rendevueyo}\senco{2}{Baziliko Romana antiqua, konsakrita a la kulto kristana, e qua ordinare divenis la kirko precipua di la urbo}\lingui{DEFIRS}
\artiklo{bazilisko}%
\vorto{bazilisko}\fako{mitol.} Reptero mitologiala di qua la regardo mortigis\lingui{DEFIRS}
\artiklo{bazilo}%
\vorto{bazilo}\fako{bot.} Genero de planti odoroza, de la familio « labiacei », di qua un varietato kultivesas en la gardeni kom plezur-planto\latina{ocynum basilicum}\lingui{DEFIRS}
\artiklo{bazo}%
\vorto{bazo}\senco{1}{Parto infra di korpo per qua olca su apogas sur to quo suportas lu}\senco{2}{\textit{(geom.)} Lineo o plano qua limitizas korpo, figuro geometriala, e quan onu prenas kom nivelo por mezurar de ibe la alteso-grado di la korpo di ta figuro}\lingui{DEFIRS}
\artiklo{bear}%
\vorto{bear}\fako{netrans.} Apertar la boko. \textit{(anke en la senco figurala : « truo beanta »)}\lingui{F}
\artiklo{beata}%
\vorto{beata} Qua ravisesas en Deo. – \textit{(metaf.)} Qua simulas beato; hipokrita\lingui{DEFIRSL}
\artiklo{beatifikar}%
\vorto{beatifikar}\fako{trans.) (aludante la papo} Pozar inter la beati (homo qua vivis sante)\lingui{DEFIRS}
\artiklo{bechika}%
\vorto{bechika} Efikiva kontre tusado\lingui{EFIS}
\artiklo{bedelo}%
\vorto{bedelo} Adulto qua preiras o sequas la procesioni (generale religiala)\lingui{DEFIRS}
\artiklo{bedo}%
\vorto{bedo} Bendo de tero quan onu garnisas ye flori, arbusti, en gardeno\lingui{DE}
\artiklo{begino}%
\vorto{begino} Sorto de regulierini qui vivis kongregacione segun reguli monakeyala, sen vovir\lingui{DEFIRS}
\artiklo{begonio}%
\vorto{begonio}\fako{bot.} Tipo di la familio « begoniacei », quan onu kultivas kom plezur-planto\latina{begonia}\lingui{DEFIRS}
\artiklo{bekafigo}%
\vorto{bekafigo}\fako{zool.} Ula mikra ucelo\latina{sulvia hortensis}\lingui{DEFIS}
\artiklo{bekaso}%
\vorto{bekaso}\fako{zool.} Ucelo migrera, ek la klaso « steltieri », de la familio « longa-beka »\latina{scolopax rusticola}\lingui{DEFIRS}
\artiklo{beko}%
\vorto{beko}\senco{1}{Boko di la uceli, qua konsistas ye du lami korno-harda (mandibuli) ed uzebla kom armo por atakar o por su defensar}\senco{2}{La extremajo di ula objekti qui finas pintatre}\lingui{EFIS}
\artiklo{bela}%
\vorto{bela} Qua vekigas la admiro-sentimento, per igar senteblan o per manifestar ta o ca perfektajo korpala od etikala\lingui{EFI}
\artiklo{beladono}%
\vorto{beladono}\fako{bot.} Planto herbacea, perena, ek la familio « solanei », venenoza en omna lua parti, e di qua la toxikeso uzesas medicinale\latina{atropa belladonna}\lingui{DEFIRS}
\artiklo{belemnito}%
\vorto{belemnito} Konko fosila, longa, cilindra, finanta per pinto\lingui{DEFIRS}
\artiklo{belvedero}%
\vorto{belvedero} Paviliono konstruktita en loko de ube la peizajo esas agreabla e vasta, od an la somito di edifico de ube onu vidas vaste\lingui{DEFIR}
\artiklo{bemolo}%
\vorto{bemolo}\fako{muziko} Signo pozita avan muzik-noto por savigar ke onu devas abasar ye un mi-tono ta noto ed olti di la sama grado en la sama mezuro. – Signo an la klefo por savigar ke onu devas abasar ye un mi-tono omna noti skribita sur la sama lineo, en la muzikajo\lingui{FIRS}
\artiklo{bendo}%
\vorto{bendo} Areo limitizita per du linei paralela\lingui{DEFIRS}
\artiklo{benedikar}%
\vorto{benedikar}\fako{trans.} Dicar vorti qui *wishas ad ulu lua feliceso\lingui{DEFIS}
\artiklo{benediktano}%
\vorto{benediktano} Reguliero en la ordeno di santa Benediktus\lingui{DEFIRS}
\artiklo{benefico}%
\vorto{benefico}\fako{religio katolika} Ofico ekleziala qua implikis la granto di revenuo\lingui{DEFIRS}
\artiklo{benigna}%
\vorto{benigna} Bonvoloza ad omna personi – indulgema ad omni\lingui{DEFIS}
\artiklo{benko}%
\vorto{benko} Sidilo longa (ligna, kom exemplo) kun o sen dors-apogilo, streta, sur qua plura personi povas sidar kune\lingui{DEFIRS}
\artiklo{benzino}%
\vorto{benzino} Hidrokarbido, liquido senkolora, extraktita de rezini e de gudro minkarbonala, qua uzesas por preparar granda quanto de kompozaji aromatoza, por netigar la stofi, e c.\simbolo{C₆H₆}\lingui{DEFIRS}
\artiklo{benzolo}%
\vorto{benzolo} Kompozanto di la karbongudri qua distilo pasas sub 170° C. e konsistas ye mixuro qua kontenas benzeno, tolueno, xilono\lingui{DEFIRS}
\artiklo{benzoo}%
\vorto{benzoo} Substanco rezinoza qua defluas, tra incizuro, de varietato de L. \textit{styrax benzoin}\latina{styrax benzoin}\lingui{DEFIRS}
\artiklo{bequadro}%
\vorto{bequadro}\fako{muziko} Signo analoga a la b kun ventro quadrata, uzata por nihiligar la bemoli o la diezi, e retroduktar noto a lua tono naturala (de kande la diezi uzesis)\lingui{FIRS}
\artiklo{berberiso}%
\vorto{berberiso}\fako{bot.} Planto lignoza, kovrita ye pikili, e qua produktas grapi de beri reda tre acida, de qui onu facas sorto di vino\latina{berberis vulgaris}\lingui{DEFIRS}
\artiklo{beredo}%
\vorto{beredo} Toko de lano, di qua la parto supra esas plata e cirkla\lingui{DEFIR}
\artiklo{berenjeno}%
\vorto{berenjeno}\fako{bot.} Varietato di solano, di qua la frukto (qua havas la formo di granda bero longa) manjesas kom legumo\latina{solanum melongena}\lingui{FS}
\artiklo{bergamoto}%
\vorto{bergamoto}\fako{bot.} Varietato di piro tre parfumoza, qua quaze liquideskas lor manjeso\latina{citrus bergamia}\lingui{DEFIRS}
\artiklo{beriberio}%
\vorto{beriberio}\fako{patol.} Sorto di anemio, partikulara ye la habitanti di insulo Ceylon\lingui{F}
\artiklo{berilo}%
\vorto{berilo} Varietato di smeraldo, diafana, maro-verda\lingui{DEFIRS}
\artiklo{berjero}%
\vorto{berjero} Larja fotelo, di qua la sido-plako esas garnisita ye kuseno mola\lingui{DFR}
\artiklo{berlino}%
\vorto{berlino} Granda karoso klozita, quar-rota, di qua la parti avana e posa esas simetra\lingui{DEFIRS}
\artiklo{berloko}%
\vorto{berloko} Siglili, klefi, mikra juveli, qui pendas de rubando, de kateno di posh-horlojo\lingui{DFIR}
\artiklo{bermo}%
\vorto{bermo} Parto horizontala en taluso\lingui{DEFRS}
\artiklo{bernaklo}%
\vorto{bernaklo}\fako{zool.} Genero de krustacei kun pedi apendicoza, di qui la konko subtenesas da pedikulo kontraktebla\latina{branta bernicla}\lingui{DEFIS}
\artiklo{bero}%
\vorto{bero}\fako{bot.} Frukto tre pulpoza, di qua la semini situesas meze di la pulpo\lingui{DE}
\artiklo{bersar}%
\vorto{bersar}\fako{trans.} Ociligar dum kelka tempo, dolce (infanteto) en lua lito o sur la gremio, por dormigar o kalmigar lu\lingui{F}
\artiklo{bestio}%
\vorto{bestio} Omna animalo, eceptante la homo\lingui{DEFIRS}
\artiklo{betelo}%
\vorto{betelo}\fako{bot.} Arbusto sarmentoza de India orientala, de la genero « pipro »\latina{piper betel}\lingui{DEFIRS}
\artiklo{beto}%
\vorto{beto}\fako{bot.} Planto herbacea, kun radiko pulpoza, genero di la familio « kenopodei »\latina{beta}\lingui{DEFIS}
\artiklo{betonio}%
\vorto{betonio}\fako{bot.} Planto herbacea, perena, kun bela flori spikedispozite – genero di la familio « labiacei »\latina{betonica officinalis}\lingui{DEFIS}
\artiklo{betono}%
\vorto{betono} Mixuro de mortero hidraulikala e de stoni triturita, qua hardeskas en aquo\lingui{DEFIRS}
\artiklo{betravo}%
\vorto{betravo}\fako{bot.} Speco di beto kun radiko rotacanta quale la rapo, blanka o reda, qua donas nutrivo por la homo e la bestii, e di qua onu extraktas sukro\latina{beta rapa}\lingui{ES}
\artiklo{bezigo}%
\vorto{bezigo} Karto-ludo. – Nomo, en ta ludo, di la unioneso di la karelo-pajo e di la piquo-damo\lingui{DEFIS}
\artiklo{bezoaro}%
\vorto{bezoaro} Agreguro kalkoloza quan onu renkontras en la intestino ed en la stomako di ula animali, e quan onu uzis kom antidoto, amuleto, e c.\lingui{DEFR}
\artiklo{bezonar}%
\vorto{bezonar}\fako{trans.} Indijar ulo quon postulas la naturo o la cirkonstanci o la skopo vizata\lingui{FI}
\artiklo{bi-}%
\vorto{bi-} Du-… \textit{(uzata por la termini ciencala : bikapa, biloba, e c.)}
\artiklo{bibliofilo}%
\vorto{bibliofilo} Persono qua prizas e kolektas la libri editita pri ta o ca temo, la diversa edituri, lia valoro, lia skarseso-grado, e c.\lingui{DEFIRS}
\artiklo{bibliografio}%
\vorto{bibliografio} Enumero di la libri publikigita pri ta o ca temo, lia imprimo-dati diversa, lia valoro, lia skarseso-grado e c.\lingui{DEFIRS}
\artiklo{bibliomanio}%
\vorto{bibliomanio} Prizego pasionoza di la libri\lingui{DEFIRS}
\artiklo{biblioteko}%
\vorto{biblioteko}\senco{1}{Moblo garnisita ye tabuli sur qui onu ordinas libri}\senco{2}{Edifico qua kontenas kolektajo de libri}\lingui{DEFIRS}
\artiklo{biblo}%
\vorto{biblo} Asemblajo ek la libri « Olda Testamento » e « Nova Testamento » – che la Kristani; la « Olda Testamento », che la Judi\lingui{DEFIRS}
\artiklo{bicepso}%
\vorto{bicepso}\fako{anat.} Muskulo qua havas du ligo-punti ye lua parto supra (kom ex. : la muskulo avana di la avanbrakio)\lingui{DEFIS}
\artiklo{biciklo}%
\vorto{biciklo} Velocipedo kun du roti inter-egala, dispozita ica dop ita\lingui{DEFIRS}
\artiklo{bideto}%
\vorto{bideto} Tualeto-moblo, kun kuveto sidilatra, por ablucioni intergamba (por la mulieri)\lingui{DEFIRS}
\artiklo{bidono}%
\vorto{bidono} Lada vazo klozita, por oleo, petrolo, benzino, e c. qua kontenas de 1 til 10 litri; gurdatra kontenilo ek lado, quan la soldato portas\lingui{FIRS}
\artiklo{bielo}%
\vorto{bielo}\fako{tekn.} Stango qua ligas un peco ad altra (en mashino), generale stango di pistono a manivelo, e transformas movo rektalinea a movo rotaca\lingui{DEFIRS}
\artiklo{bifsteko}%
\vorto{bifsteko} Loncho de bovo-karno destinita a koqueso sur grililo\lingui{DEFIRS}
\artiklo{bifurkar}%
\vorto{bifurkar}\senco{1}{\textit{(trans.)} Dividar a du branchi}\senco{2}{\textit{(netrans.)} Dividar su quale la gambi di homo bifurkas de trunko}\lingui{DEFIS}
\artiklo{bigama}%
\vorto{bigama}\senco{1}{\textit{(yuro ekleziala)} Mariajita dufoye; mariajita kun vidvino}\senco{2}{Qua su mariajis duesmafoye, la mariajo unesma ne esante dissolvita}\lingui{DEFIRS}
\artiklo{bigoto}%
\vorto{bigoto} Qua agas devocaji streta-mente, minucioze e kelke hipokrite\lingui{DEFI}
\artiklo{bikursalo}%
\vorto{bikursalo}\fako{geom.} Kurvo sur qua selektesis punto de ube duktesas rekto; ta kurvo renkontresas da ta rekto en nur du punti; la duopla punto esinte selektita origine. (Ex. : lemniskato)\lingui{DEFI}
\artiklo{bikuspida}%
\vorto{bikuspida}\fako{natur-historio} Qua finas per du pinti\lingui{DEFIS}
\artiklo{bilanco}%
\vorto{bilanco}\fako{komerco} Equilibrigo di aktivo e di pasivo di konto komercala\lingui{DEFIRS}
\artiklo{bilboketo}%
\vorto{bilboketo} Ludilo qua konsistas ye bastoneto qua finas, ca-latere per pinto, talatere per pleto, mezi di qua esas suspendita bulo quan onu devas lansar tale ke onu retrorecevos lu sur la pinto o sur la pleto\lingui{DFIR}
\artiklo{bilgo}%
\vorto{bilgo} Kolektuyo (en la fondo di la holdo di navo) di la aqui, rezidui ed exkrementi, quan onu vakuigas per pumpili\lingui{DE}
\artiklo{bilieto}%
\vorto{bilieto} Paper-folieto, karto, per qua konstatesas, favore di ulu, yuro tempala, komprita, od obtenita per favoreso\lingui{DEFIRS}
\artiklo{biliono}%
\vorto{biliono} 1.000.000²\lingui{DEFIRS}
\artiklo{biljardo}%
\vorto{biljardo}\senco{1}{Pleo qua konsistas ye pulsar ivoro-buli per longa bastono}\senco{2}{Tablo sur qua esas tensita tapiso verda, garnisita ye rebordi, e sur qua onu pleas biliardo}\lingui{DEFIRS}
\artiklo{bilo}%
\vorto{bilo} Liquido bitra, flava verdatre, sekrecita da hepato, qua penetras en la duodenumo segun ke la nutro-materio decensas aden lu, e pri qua onu konjektas ke lu helpas digestado\lingui{DEFIS}
\artiklo{binara}%
\vorto{binara}\fako{pri sistemo di nombrifado} En qua onu expresas omna nombri per du cifri (1 e 0)\lingui{DEFIS}
\artiklo{bindar}%
\vorto{bindar}\fako{trans.} Sutar ad-ensemble la kayeri di libro, inkastrar li en kovrilo kartona di qua la dorso e la platajo kovresas da telo, felo, e c.\lingui{DE}
\artiklo{bineto}%
\vorto{bineto}\fako{agrokultivo} Implemento kun mancho longa, di qua la ferajo prizentas ye ca latero di la mufo, lamo streta e plata; talatere, du (e mem tri) denti – e qua uzesas por kultive modifikar duesmafoye sulo o planto\lingui{FS}
\artiklo{binoklo}%
\vorto{binoklo} Du okulo-vitrodiski solidara\lingui{DEFIRS}
\artiklo{binomio}%
\vorto{binomio}\fako{algebro} Quanto qua konsistas ye du termini unionita per la signo « plus » o la signo « minus »\lingui{DEFIRS}
\artiklo{biodinamiko}%
\vorto{biodinamiko} Cienco qua relatas la teorio (tre anciena) segun qua existas vivo-forci qui diferas de la forci fiziko-kemiala, e qui esus la fonto di la manifesti specifika : sentiveso, moviveso, genetiveso, morto, e c., di la enti organoza\lingui{DEFIRS}
\artiklo{biogenezo}%
\vorto{biogenezo} Teorio segun qua omna ento vivanta venis de ento qua genitis lu\lingui{DE}
\artiklo{biografio}%
\vorto{biografio} Redakturo qua naracas la viv-eventaji di persono\lingui{DEFIRS}
\artiklo{biokemio}%
\vorto{biokemio} Parto di biologio, qua traktas la relati di la strukturo kemiala di la elementi nemediata, a la mekanismo funcionala di la celuli, tisui ed organi di la organismi vivanta\lingui{DEFIRS}
\artiklo{biologio}%
\vorto{biologio} Cienco generala pri la vivo, lua kondicioni e lua formi diversa (organizeso di la enti vivanta, medii ube li vivas, funcioni di la vivo, e c.)\lingui{DEFIRS}
\artiklo{biomekaniko}%
\vorto{biomekaniko} Parto di biologio, qua traktas la expliko mekanista di la fenomeni vivala\lingui{DEFIRS}
\artiklo{biremo}%
\vorto{biremo} Galero qua havas singla-latere du rangi de remisti\lingui{DEFIRS}
\artiklo{bireto}%
\vorto{bireto} Mikra boneto; kun quar edri quadrata, qua esas faldebla\lingui{DEFIRS}
\artiklo{birko}%
\vorto{birko}\fako{bot.} Arboro qua esas un genero di la familio « betulacei »\latina{betula}\lingui{E}
\artiklo{biro}%
\vorto{biro} Drinkajo fermentacinta, preparita ek hordeo ed ek jermifinta frumento e lupulo\lingui{DEFI}
\artiklo{bis}%
\vorto{bis} Duesmafoye\lingui{DFIRS}
\artiklo{bisko}%
\vorto{bisko} Kankro-supo\lingui{DEF}
\artiklo{biskoto}%
\vorto{biskoto} Loncho de pano preparita per lakto (vice per aquo) e sikigita e bakita en forno\lingui{DEF}
\artiklo{bismuto}%
\vorto{bismuto}\fako{kemio} Metalo grize-blanka kun nuanci redatra, lameloza, frajila, qua, amalgamita kun merkurio, uzesas por stanizar la speguli\simbolo{Bi}\lingui{DEFIRS}
\artiklo{biso}%
\vorto{biso} Materio texebla, sorto di lino flavatra, de qua la Antiqui fabrikis stofi tre granda-preco\lingui{DEFIRS}
\artiklo{bisquito}%
\vorto{bisquito} Kukajo lejera, qua konsistas ye ovi, farino e sukro\lingui{DEFIRS}
\artiklo{bistro}%
\vorto{bistro} Farbo bruna flavatre, quan onu obtenas de fuligino koquita e dilutita, e quan onu uzas por aquo-piktar\lingui{DEFR}
\artiklo{bisturio}%
\vorto{bisturio} Mikra kultelo kirurgiala, di qua la lamo, movebla, mantenesas da resorto kande onu apertas lu\lingui{DEFIRS}
\artiklo{bito}%
\vorto{bito}\fako{navig.} Bloko pozita en la parto avana di la navo, e qua konsistas ye du montanti ed un traverso sur qua spulesis e fixigesis per amaro la kabli\lingui{DEFIRS}
\artiklo{bitra}%
\vorto{bitra} Qua saporas repugnante. -\textit{(metaf.)} Qua impresas penigante, chagrenigante\lingui{DE}
\artiklo{bitumo}%
\vorto{bitumo} Substanco minerala, liquida (nafto) o solida (asfalto), bruna\lingui{DEFIRS}
\artiklo{bivakar}%
\vorto{bivakar}\fako{netrans.} Kampar sen tendi, havante cielo kom plafono, generale cirkum fairi, nur dum un nokto\lingui{DEFIRS}
\artiklo{bixestila}%
\vorto{bixestila} En qua inkluzesis la dio quan onu adjuntas a la monato februaro singla-quaresma-yare\lingui{DEFIS}
\artiklo{bizantina}%
\vorto{bizantina} Qua prizas disipar sua tempo en diskutar subtilaji ne-utila\lingui{DEFIRS}
\artiklo{bizelo}%
\vorto{bizelo} Bordo (di planko, cizelilo, e c.) tranchita oblique\lingui{FS}
\artiklo{bizono}%
\vorto{bizono}\fako{zool.} Bovo sovaja di Usa, di qua la dorso esas giboza, la korni kurta e rondatra, kun longa barbo\latina{bos americanus}\lingui{DEFIRS}
\artiklo{blamar}%
\vorto{blamar}\fako{trans., pro, pri} Desaprobar ulu quan onu judikas kom aginta male\lingui{DEFIS}
\artiklo{blanka}%
\vorto{blanka} Di qua la koloro esas pasable simila a lakto, a nivo-floki, e c.\lingui{DEFIRS}
\artiklo{blasfemar}%
\vorto{blasfemar}\fako{netrans.} Pronuncar vorti qui ofensas la deajo\lingui{DEFIS}
\artiklo{blato}%
\vorto{blato}\fako{zool.} Insekto ortoptera bruna, qua odoras fetide, ajila, devorema, qua devoras la nutrivo-provizuri, rodas la stofi, e c.\latina{blatta}\lingui{DFI}
\artiklo{blazar}%
\vorto{blazar}\fako{trans.} Obtuzigar per impresi granda ed ofte iterata\lingui{DEF}
\artiklo{blazono}%
\vorto{blazono} Ensemblo de emblemi konsiderata kom signo distingiva di familio nobela, di urbo, di korporaciono, e c., quan onu reprezentis sur la skudi, la baneri, la siglili, e c.\lingui{DEFIS}
\artiklo{blendo}%
\vorto{blendo}\fako{kemio} Erco de zinko-sulfuro, ordinare unionita a la mineralo-strati de plombo, de arjento\lingui{DEFIRS}
\artiklo{blezar}%
\vorto{blezar}\fako{netrans.} Pronuncar « s » quale « th » Angla o « c » (avan « a » e « i ») Hispana (t.e. kun la lango-pinto inter la denti.)\lingui{F}
\artiklo{blinda}%
\vorto{blinda} Sen vidiveso\lingui{DE}
\artiklo{blokhauso}%
\vorto{blokhauso}\fako{milit.} Fuorteto regiono-sola, ek ligno, ordinare juntita a fortifikajo precipua per koridoro subsula\lingui{DEF}
\artiklo{bloko}%
\vorto{bloko} La toto di korpo; amaso granda e pezoza\lingui{DEFIRS}
\artiklo{blokusar}%
\vorto{blokusar}\fako{trans.} *Asiejar fortreso, kampeyo, portuo, por impedar omna komuniko da la asiejiti kun la exteresanti\lingui{DEFIRS}
\artiklo{blonda}%
\vorto{blonda} \textit{(aludante la hari, la brovi, la barbo, e c.)} Di qua la koloro esas inter or-flava e klar-kastanea\lingui{DEFIRS}
\artiklo{blotisar}%
\vorto{blotisar}\fako{netrans.} Lokizar su squate, kontraktante su, an poka spaco\lingui{F}
\artiklo{blua}%
\vorto{blua} Di qua la koloro pasable similesas olta di cielo sen-nuba, jorne\lingui{DEFI}
\artiklo{blutar}%
\vorto{blutar}\fako{trans.} Pasigar tra la dika telo (burelo) di sivo cilindra, adube falas la mueluro e qua – retenante la brano – lasas pasar nur la farino\lingui{F}
\artiklo{bluzo}%
\vorto{bluzo} Supervesto, ek telo dika, lano, e c., kun aperturo avana, e di qua la formo similesas olta di kamizo – quan *weras la rurani, la laboristi brakiala, e c., e la infanti e pueri, kom streno-vesto\lingui{DEFIRS}
\artiklo{bo-}%
\vorto{bo-} Prefixo qua indikas la parenteso rezulte de mariajo
\artiklo{boao}%
\vorto{boao} Serpento di Sud-Amerika, di qua la grandeso e la forteso esas timinda – qua trituras per sua ringi la animali mem la maxim granda, ed englutas li pos kovrir li ye sua salivo\latina{boa constrictor}\lingui{DEFIRS}
\artiklo{bobino}%
\vorto{bobino} Mikra cilindro ligna, kun rebordi, sur qua la filo lina o silka volvesas o desvolvesas reguloze, spulesis\lingui{DEFRS}
\artiklo{bodmeriar}%
\vorto{bodmeriar}\fako{netrans.} Prestar pekunio di qua la rimborseso garantiesas da vari qui mar-transportesas, ma qua rimborsesos nur se ta vari atingas lia destino-portuo, o proporcione a la quanto di la vari qui tale arivos\lingui{N}
\artiklo{bodrucho}%
\vorto{bodrucho} Dina membrano de la grosa intestino di la bovi, di la mutoni, desgrasizita, preparita a pelikulo preske diafana, quan onu pozas sur la vunduri – ek qua onu fabrikas baloneti – e per qua onu kovras la or-folii lor lia bateso\lingui{F}
\artiklo{bogio}%
\vorto{bogio} Ensemblo de quar (o du) roti qui formacas vagoneto e suportas la extremajo di vagono o di lokomotivo\lingui{DEFS}
\artiklo{boikotar}%
\vorto{boikotar}\fako{trans.} Interkonsentar (tacite od expresite) por punisar, per domajo pekuniala od etikala, individuo o lando, per refuzar exekutar ula labori, komprar ula vari, e c.\lingui{DEFIRS}
\artiklo{bokalo}%
\vorto{bokalo} Vazo cilindra o cilindratra, ek vitro, sen pedo, kun fundo plata e boko larja, uzata precipue en kemio, farmacio (e por konfituri)\lingui{DFIRS}
\artiklo{boko}%
\vorto{boko}\senco{1}{Che la homo e kelka animali : kavajo ye la parto infra di la vizajo, qua komunikas interne kun la kanalo respirala e la kanalo nutrivala}\senco{2}{Aperturo qua esas la enireyo di kavajo}\senco{3}{Parto di fluvio o di rivero, ube lia aquo enfluas rispektive la maro od altra aquofluo}\lingui{FIS}
\artiklo{bolero}%
\vorto{« bolero »} Kansono, dans-ario Hispana, generale ye modo minora, tri-tempa, e ye movo rapida\lingui{DEFIRS}
\artiklo{boleto}%
\vorto{boleto}\fako{bot.} Fungo kun chapelo pedikuloza, mi-sfera, di qua la surfaco suba konsistas ye tubi cilindra – genero di qua plura speci esas manjebla (cepo, e c.) e di qua un altra (agariko di la querko) esas la materio di tindro\latina{boletus}\lingui{DEFIS}
\artiklo{boliar}%
\vorto{boliar}\fako{netrans.) (aludante liquido} Subisar ta transformeso rapida ye ula temperaturo (qua dependas de la liquido e de la aero-preso) per qua parto de la liquido vaporeskas, forme di buli qui acensas groseskante til la surfaco ube li krevesas, disjetante en la atmosfero cirkondanta la aero o la vaporo quin li kontenis\lingui{EFIS}
\artiklo{bolido}%
\vorto{bolido}\fako{astron.} Meteoro fairatra, akompanata da traco-linei, lumoza, fumuro, detono ed adterfalo di aeroliti\lingui{DEFIRS}
\artiklo{bolo}%
\vorto{bolo} Kalico mi-sfera ek fayenco, porcelano, generale destinita a kontenor nutrivi liquida\lingui{DEFS}
\artiklo{bolto}%
\vorto{bolto}\fako{tekn.} Stifto de fero qua uzesas por interfixigar du peci de ligno, de fero, e c., finanta capinte per kapo ronda qua limitizas lua extremajo supera, e tapinte per skrubino o kolo\lingui{DER}
\artiklo{bombardar}%
\vorto{bombardar}\fako{trans.} Asaltar per la lanso di bombi\lingui{DEFIRS}
\artiklo{bombinatoro}%
\vorto{bombinatoro}\fako{zool.} Genero de batraki sen-kauda, analoga a rospi, karakterizata da lia lango ronda, ligita tota-aree a la boko-sulo, e da lia ventro qua esas veinizita ye makuli nigra ed oranjo-flava\latina{bombinator}\lingui{FSL}
\artiklo{bombo}%
\vorto{bombo} Projektilo ek fero, kava, plenigita ye substanco detoniva qua explozas lor lua falo\lingui{DEFIRS}
\artiklo{bombono}%
\vorto{bombono} Sukrajo (drajeo, e c.) fabrikita da la sukrajisti\lingui{DEFIRS}
\artiklo{bona}%
\vorto{bona} Qua prokuras avantajo, satisfaco; qua havas en su perfekteso\lingui{DEFIRSL}
\artiklo{boneto}%
\vorto{boneto} Kapvesto mola, senborda, ek stofo, trikoturo, furo, e c.\lingui{EFIS}
\artiklo{bonimentar}%
\vorto{bonimentar}\fako{netrans.) (aludante la sharlatani e feriala aktori} Diskursar fanfaronere por laudar sua vari od atraktar publiko\lingui{F}
\artiklo{bonzo}%
\vorto{bonzo} Budha-sacerdoto en la Oriento-regioni di Azia\lingui{DEFIRS}
\artiklo{boracho}%
\vorto{boracho}\fako{bot.} Planto herbacea, tipo di la familio « boragacei », di qua la stipo e la folii kovresas ye pili pikanta – e la folii, dispozita grape\latina{borago}\lingui{DEFIRS}
\artiklo{borar}%
\vorto{borar}\fako{trans.} Truifar rotace. (Borilo portesas e movesas da drilo)\lingui{DER}
\artiklo{boraxo}%
\vorto{boraxo}\fako{kemio} Antea nomo di la sodo-borato\simbolo{B₄O₇Na₂}\lingui{DEFIRS}
\artiklo{borborigmo}%
\vorto{borborigmo} Bruiso da la diplaseso di la gasi en la abdomino, meze di la liquidi quin kontenas la intestini\lingui{DEFIS}
\artiklo{bordear}%
\vorto{bordear}\fako{netrans.} Avancar \textit{(aludante navo)} zigzage kontre vento\lingui{FIS}
\artiklo{bordo}%
\vorto{bordo} Parto extrema qua limitizas la konturo di objekto\lingui{DEFIRS}
\artiklo{boreo}%
\vorto{boreo} Nomo mitologiala di la nordo-vento\lingui{DEFIRS}
\artiklo{borgezo}%
\vorto{borgezo} La persono qua apartenas a la meza klaso di la habitanti di urbo\lingui{DEFIRS}
\artiklo{boro}%
\vorto{boro}\fako{kemio} Korpo simpla, metaloida, bruna verdatre, sen saporo nek odoro\simbolo{B}\lingui{DEFIRS}
\artiklo{borso}%
\vorto{borso} Loko publika ube su asemblas la kurtajisti, monetokambiisti, e c., por operaci komercala, financala pri valoraji (acioni, obligacioni) di qui la kurso esas variiva\lingui{DEFIRS}
\artiklo{bosajo}%
\vorto{bosajo} Saliajo sur la surfaco di muro, destinita a talieso, skulteso figurala od ornamentala\lingui{DEFI}
\artiklo{boselar}%
\vorto{boselar}\fako{trans.} Formacar reliefi en metalo per martelago (de interne)\lingui{DEF}
\artiklo{bosko}%
\vorto{bosko} Grupo de arbori qui kovras areo pasable granda\lingui{DEFIS}
\artiklo{bostono}%
\vorto{bostono}\senco{1}{Ludo quar-partnera per kinadekedu karti}\senco{2}{Sorto di valso, deveninta de Usa, qua exekutesas per glitar segun mezuro tri-tempa, en ritmo « moderato »}\lingui{DEFIRS}
\artiklo{botaniko}%
\vorto{botaniko} Parto di la natur-historio qua havas kom temo di studio la vejetanti\lingui{DEFIRS}
\artiklo{botanizar}%
\vorto{botanizar}\fako{netrans.} Kolektar por la studio botanikala. (ne « praktikar o studiar botaniko »)\lingui{DEFIRS}
\artiklo{botelo}%
\vorto{botelo} Vazo portebla, di qua la ventro esas cilindra, la kolo streta e longa, destinita a kontenar liquidi, precipue vino\lingui{DEFIRS}
\artiklo{boto}%
\vorto{boto} Pedo-vesto ledra qua kontenas la gambo. (\textit{Boteto} kontenas nur la infra parto di la gambo; \textit{shuo} kontenas nur la pedo)\lingui{EFRS}
\artiklo{bovo}%
\vorto{bovo}\fako{zool.} Mamifero ruminera quan onu uzas prefere por la agrokultivo se lu esas bovulo, e por obtenar lua lakto se lu esas bovino – o quan onu grasigas por bucheso. (Kande bovulo esas destinita a la tauro-kombati, lua nomo esas « tauro »)\lingui{EFIS}
\artiklo{boxar}%
\vorto{boxar}\fako{trans. e netrans.} Interbatar per pugno-frapi, en la maniero Angla\lingui{DEFR}
\artiklo{boyardo}%
\vorto{boyardo} Sinioro, personego, en Rusia caroza, en Transilvania\lingui{DEFIRS}
\artiklo{boyo}%
\vorto{boyo}\fako{navig.} Omna korpo flotacanta (peco de ligno o korko, barelo, e c.) destinita a markizar, sur la surfaco di aquo, la punto ube lansesis ankro – ed informar pri rifo od obstaklo od indikar la direciono di paseyo o di kanaleto\lingui{DEFIRS}
\artiklo{braceleto}%
\vorto{braceleto} Ornivo ringo-forma qua cirkondas la ponyeto\lingui{DEFIRS}
\artiklo{bracho}%
\vorto{bracho} Viro-vesto, qua kovras la persono, de la hanchi til sub la genui, e bifurkas tale ke olu envelopas la kruri\lingui{EFIS}
\artiklo{braco}%
\vorto{braco} Kordego an la du extremaji di yardo\lingui{DEFIS}
\artiklo{bradipo}%
\vorto{bradipo}\fako{zool.} Mamifero de la tribuo « sen-denti »\latina{bradypus}\lingui{EFIS}
\artiklo{brakicefala}%
\vorto{brakicefala} Di qua la kapo esas kurta\lingui{DEFIRS}
\artiklo{brakio}%
\vorto{brakio}\senco{1}{Membro en la parto supra di la korpo homala, qua artikizesas ye la shultro e qua finas per la manuo; che ula animali, membri qui korespondas a la brakii homala pro lia situeso o funciono}\senco{2}{Parto de la maro qua konstriktesas da du litori}\lingui{DEFIRS}
\artiklo{brakiopodo}%
\vorto{brakiopodo}\fako{zool.} Klaso de moluski kun konko bivalva, sen-movila, di qua la boko situesas inter du brakii tre karnoza ed esas provizita ye multa filamenti quin oli povas igar ekirar od enirar per spular olci spirale\lingui{DEFIS}
\artiklo{brakistokrono}%
\vorto{brakistokrono}\fako{mekan.} Kurvo quan devas sequar korpo pezoza por irar de ca a ta punto en tempo minima, la rapideso-grado lor la departo esante nula\lingui{DEFI}
\artiklo{brako}%
\vorto{brako} Speco di chas-hundo, reza-pila, e kun oreli pendanta\lingui{DEFIS}
\artiklo{brakteo}%
\vorto{brakteo}\fako{bot.} Folieto qua diferas de la folio ordinara qua, che ula planti, situite ye la inserto-punto di la floro, akompanas o parte kovras olca\lingui{DEFIS}
\artiklo{bramano}%
\vorto{bramano} Sacerdoto di la bramanismo, qua apartenas a la unesma ek la kasti Indiana, e qua dicis naskir de la kapo di Brama\lingui{DEFIRS}
\artiklo{bramar}%
\vorto{bramar}\fako{netrans.} Vorto generala por omna krii di bestii\lingui{FIS}
\artiklo{brancho}%
\vorto{brancho}\senco{1}{\textit{(bot.)} Sproso ligna qua kreskas sur la trunko di arboro}\senco{2}{\textit{(geom.)} Singla de la linei intersimetra, di qui la uniono formacas ula kurvi geometriala (hiperbolo, parabolo)}\senco{3}{To omna quo, departante de stango, aspektas analoge kam la arboro-branchi}\senco{4}{La familii diversa, genitita de trunko komuna. (Ex. : familio-branchi)}\lingui{DEFIR}
\artiklo{brandio}%
\vorto{brandio} Kompozajo ek aquo ed alkoholo, obtenita per la distilo di vino, cidro, grani, betravo, e c.\lingui{DEF}
\artiklo{brandisar}%
\vorto{brandisar}\fako{trans.} Ociligar en sua manuo minacoze (espado, lanco, javelino)\lingui{EFIS}
\artiklo{brando}%
\vorto{brando} Fasko de palio flamoza per qua onu dis-tushas por incendiar, o quan onu kunportas por lumizar\lingui{DEFI}
\artiklo{brankardo}%
\vorto{brankardo} Singla de la du branchi (ligna), plu o min lejera, maxim ofte kurvigita por su adaptar a la korpo di la kavalo, qui artike juntesas a la avana parto di veturo, ed inter qui onu jungas la kavalo\lingui{F}
\artiklo{brankio}%
\vorto{brankio} Lami dispozita quale la denti di pektilo, ye qua konsistas la respir-aparato di la animali qui vivas en aquo e respiras la aero quan olca kontenas dissolvite (fishi, krustacei)\lingui{EFIS}
\artiklo{brankiopodo}%
\vorto{brankiopodo}\fako{zool.} Mikra krustaceo, di qua la pedi, apta por natado, esas garnisita ye apendici brankia\lingui{DEFIS}
\artiklo{brano}%
\vorto{brano} Reziduo de la muelo, asemblajo ek la eskombri de la shelo di la frumento-grani\lingui{EF}
\artiklo{brasar}%
\vorto{brasar}\fako{trans.} Movar, agitar por operaco. (Ex. : brasar oro ed arjento en kruzelo)\lingui{FIS}
\artiklo{brava}%
\vorto{brava} Pronta afrontar la danjero (precipue en la kombati)\lingui{DEFIRS}
\artiklo{bravur-ario}%
\vorto{« bravur-ario »} Ario briloza, destinita ad eminentigar la kantisto\lingui{DEF}
\artiklo{brazar}%
\vorto{brazar}\fako{trans.} Unionar (du metali interdiferanta) per aloyuro quan onu fuzas en fairego\lingui{EF}
\artiklo{brechio}%
\vorto{brechio} Roko di qua la strukturo esas fragmentoza, ek peceti anguloza, diversa-kolora, aglomerita en cemento naturala
\artiklo{brecho}%
\vorto{brecho} Rupturo qua posibligas enirar remparo-cirkuito\lingui{DEFIRS}
\artiklo{breko}%
\vorto{breko} Veturo quar-rota, generale sen-tekta, qua havas, ye sua parto avana, sidili situita alte, e du benki longesale\lingui{DEFIS}
\artiklo{brelano}%
\vorto{brelano} Hazardo-ludo, quan onu pleis per distributar tri karti ad omna ludonto\lingui{DFI}
\artiklo{bremo}%
\vorto{bremo}\fako{zool.} Fisho river-aquala, qua similesas karpo, e di qua la karno nutro-prizesas\latina{abramis brama}\lingui{EFS}
\artiklo{bretelo}%
\vorto{bretelo} Bendo de ledro, quan onu pasigas sur la shultro, por suportar sako, dorso-korbo, kamilo; bendo de texuro elastika, bifurkoza, por subtenar la pantalono virala\lingui{FI}
\artiklo{breviario}%
\vorto{breviario}\fako{religio katolika} Asemblajo ek la prego-texti qui, en la religio katolika, esas recitenda da la kleriki ye ula kloki di la jorno\lingui{DEFIS}
\artiklo{brevo}%
\vorto{brevo} Reskripto pontifikala, ordinare en linguo Latina; papo-sigluroza\lingui{DEFIRS}
\artiklo{brezo}%
\vorto{brezo} Ligno transformita a karbono per kombusto, sive ardoranta, sive extingita\lingui{EFIS}
\artiklo{brido}%
\vorto{brido} La du rimeni qui, fixigita ye singla lateri di la morso, uzesas por haltigar o direktar sel-kavalo o veturkavalo\lingui{EFIS}
\artiklo{brigado}%
\vorto{brigado} Korpo de armeani mi-diviziona\lingui{DEFIRS}
\artiklo{brigantino}%
\vorto{brigantino}\fako{navig.} Mikra navo, rigoza quale brigo, ma senponta\lingui{DEFIRS}
\artiklo{brigo}%
\vorto{brigo}\fako{navig.} Navo qua havas nur un granda masto, inklinita addope ad un mezeno-masto\lingui{DEFIRS}
\artiklo{briko}%
\vorto{briko} Karelo ek argilo hardigita per fairo o suno-radii, quan onu uzas por la konstrukto di muri, di kameni, por kovrar sulo, e c.\lingui{DEFR}
\artiklo{brilar}%
\vorto{brilar}\fako{netrans.} Difuzar lumo forta\lingui{DEFIS}
\artiklo{brilianto}%
\vorto{brilianto} Diamanto, facetoza sure e sube, quan onu povas inkastrar ajure por ke lu esez plu briloza\lingui{DEFRS}
\artiklo{brino}%
\vorto{brino} Parto (seciono o segmento) di kordo\lingui{FS}
\artiklo{briocho}%
\vorto{briocho} Kuko ek pasto fermentacinta, konsistanta ye farino, butro ed ovi, e qua ordinare havas la formo di kaloto inflita, quaze-kronizita da kapo\lingui{F}
\artiklo{briska}%
\vorto{briska} Di qua la gayeso esas kelke tro libera\lingui{E}
\artiklo{brizanto}%
\vorto{brizanto} Rokajo reze aquo-surfaco (en maro) sur qua la ondi ruptesas\lingui{F}
\artiklo{brizo}%
\vorto{brizo} Vento fresha qua naskas sur la litoro ye ula kloki\lingui{DEFIRS}
\artiklo{brocho}%
\vorto{brocho} Muliero-juvelo, kun pinglo, uzata por ligar shalo, kolumo, e c.\lingui{DEFRS}
\artiklo{brodar}%
\vorto{brodar}\fako{trans.} Plubeligar (texuro) per ornivi reliefa, manuo-facita\lingui{EFS}
\artiklo{brokantar}%
\vorto{brokantar}\fako{trans.} Komprar e pose vendar o kambiar vari ne- *nuva (ne-nova), ne-fresha\lingui{EF}
\artiklo{brokato}%
\vorto{brokato} Stofo de silko rekamita ye flori, ornivi ek oro, arjento, silko\lingui{DEFIRS}
\artiklo{«brokoli»-kaulo}%
\vorto{«brokoli»-kaulo}\fako{bot.} Italia-kaulo, varietato di kaulo, kun pedunkli min grosa e plu longa kam ti di la florkaulo ordinara\latina{brassica oleracea botrytis}\lingui{DEFIRS}
\artiklo{bromo}%
\vorto{bromo}\fako{kemio} Korpo simpla, metaloida, obskur-reda, qua tre odoras quale kloro, di qua lu esas vicina pro la ensemblo de sua propraji\simbolo{Br}\lingui{DEFIRS}
\artiklo{bronkio}%
\vorto{bronkio}\fako{anat.} Singla ek la du tubi kartilagoza quin formacas la bifurko di la trakeo e qui subdividesas, en pulmono, a rami sencese plu mikra\lingui{DEFIRS}
\artiklo{bronkito}%
\vorto{bronkito}\fako{patol.} Inflameso di la mukozo qua tapete kovras la bronki\lingui{DEFIRS}
\artiklo{bronzo}%
\vorto{bronzo}\fako{kemio} Metalo harda e sonora, obskur-kolora, qua konsistas ye aloyuro de kupro e stano, segun proporcioni variiva ed a qua onu adjuntas olakaze zinko, plombo e mem arjento – quan onu uzas kom materio por gisar statui, medalii, kloshi, kanoni\lingui{DEFIRS}
\artiklo{brosar}%
\vorto{brosar}\fako{trans.} Frotar, per asemblajo ek faski de krini, de porko-pili, ek fragmenti de hundo-herbo e c., ajustita inter-nivele sur plako de ligno, de ivoro, e c. por netigar, sen-polvigar, e c.\latina{triticum caninum}\lingui{EFS}
\artiklo{broshar}%
\vorto{broshar}\fako{trans.} Pasigar filo tra folii o kayeri de folii por sutar li ad-ensemble\lingui{DEFR}
\artiklo{brovo}%
\vorto{brovo}\fako{anat.} Saliajo arko-forma, garnisita ye pili, qua esas situita super singla okulo\lingui{DER}
\artiklo{brueto}%
\vorto{brueto} Sorto di mikra baskul-veturo, quan onu pulsas o tranas per du brankardi, e qua rulas per mikra roto an la dopa parto\lingui{F}
\artiklo{bruisar}%
\vorto{bruisar}\fako{netrans.} Sono od asemblajo de soni qui rezultas de vibri nereguloza, e ne relatas skalo muzikala\lingui{F}
\artiklo{brular}%
\vorto{brular}\fako{trans. e netrans.}\senco{1}{Parkonsumar, destruktar per la efiko da fairo}\senco{2}{Domajar (kom ex. : la dermo) per la efiko da fairo o de kaloro intensa}\lingui{F}
\artiklo{bruna}%
\vorto{bruna} Di qua la koloro tendencas esor nigra\lingui{DEFIR}
\artiklo{brunisar}%
\vorto{brunisar}\fako{trans.} Polisar metali per instrumento ek stalo, agato, hematito, sanguino, e c.\lingui{DEFIS}
\artiklo{bruo}%
\vorto{bruo}\fako{bot.} Kovrilo verda, fibroza, ledratra, qua protektas la shelo di la nuco, di la mandelo, di la avelano, e c.\lingui{F}
\artiklo{bruska}%
\vorto{bruska} Qua agas per movo subita e violentoza\lingui{DEFIS}
\artiklo{bruta}%
\vorto{bruta} De qua la instinto grosiera ne esas fasonita da civilizeso\lingui{EFIS}
\artiklo{bubalo}%
\vorto{bubalo} Antilopo di Afrika, kun korni ringoza e kurvigita addope\latina{bubalis}\lingui{EFIRS}
\artiklo{bubo}%
\vorto{bubo} Puero mizeroza e malicoza qua vagas en la stradi\lingui{D}
\artiklo{bubono}%
\vorto{bubono}\fako{patol.} Tumoro glandatra, produktita da la inflameso di la ganglioni limfala sub-pela, specale ye la inguino, la axelo, la kolo\lingui{DEFIRS}
\artiklo{buchar}%
\vorto{buchar}\fako{trans.} Ocidar la bestii destinita a la homo-nutro\lingui{EF}
\artiklo{budar}%
\vorto{budar}\fako{netrans.} Manifestar nekontenteso, mala humoro, per silenco, agado, expreso di la fizionomio\lingui{F}
\artiklo{budjeto}%
\vorto{budjeto}\fako{financo} Etato komparala di la spensi e di la revenui publika, pri omna yaro\lingui{DEFR}
\artiklo{budmakero}%
\vorto{budmakero} Profesionisto di la pariado sur la agri di kavalkonkurso-kuri\lingui{EF}
\artiklo{buduaro}%
\vorto{buduaro} Mikra chambro di la apartamento di damo, refujeyo eleganta ube el su lokizas kande elu esas sola\lingui{DEFR}
\artiklo{bufalo}%
\vorto{bufalo}\fako{zool.} Speco di bovo di la regioni tropikala, aklimatita en Grekia, Italia, e c., plu granda e plu forta kam la bovo ordinara, plu sovaja kam lu, e di qua la muzelo esas plu larja kam la lua, la korni plu larja e kurvigita addope\latina{bubalus}\lingui{DEFIS}
\artiklo{bufar}%
\vorto{bufar}\fako{netrans.} Kreskar volumine ad omna direcioni\lingui{DF}
\artiklo{bufeto}%
\vorto{bufeto} Tablo, plad-etajero ube esas servita dishi; kukaji, drinkaji, ye la dispono di la invititi, voyajanti\lingui{DEFR}
\artiklo{bufono}%
\vorto{bufono} La persono qua esforcas ridigar per joki grosiera\lingui{DEFRS}
\artiklo{bufro}%
\vorto{bufro} Maso polsterizita (o resortoza) destinita ad amortisar la shoki\lingui{DER}
\artiklo{buglo}%
\vorto{buglo} Trumpeto kun klavi\lingui{DEFS}
\artiklo{bugrano}%
\vorto{bugrano} Dika telo apretita, quan la taliori uzas kom futero interna por rigidigar la infra parto di la pantalono, la kolumo, la surfaldaji di vesto\lingui{EFIS}
\artiklo{buho}%
\vorto{buho}\fako{zool.} Nokt-ucelo di la familio « strigi » (forsan pro la plum-penacheti qui ornas lua kapo)\latina{buho maximus}\lingui{L}
\artiklo{bujio}%
\vorto{bujio}\senco{1}{Kandelo de stearino}\senco{2}{Organo uzata por acendar la gas-mixuro en la explozo-motori}\lingui{DEFS}
\artiklo{buketo}%
\vorto{buketo} Asemblajo de flori koliita e kunligita; fasko analoga a la aspekto di flor-buketo\lingui{DEFR}
\artiklo{buklo}%
\vorto{buklo} Sorto di ringo o mi-ringo de metalo tra qua pasas axo kun halto-pinti, uzata por fixigar la extremajo di rimeno, di gurto, di klapo, e c.\lingui{EFIRS}
\artiklo{bukoliko}%
\vorto{bukoliko} Qua apartenas a la genero « dramato pastorala »\lingui{DEFIS}
\artiklo{bulbo}%
\vorto{bulbo} Influro tuberkulala, ovoida, globatra, formacita da la asembleso di squami o di folii rudimenta\lingui{EFIS}
\artiklo{buldogo}%
\vorto{buldogo}\fako{zool.} Varietato di hundo, plu mikra, di qua la nazo esas min aplastita, la kaudo tordita, la pili falvea e nigre strioza\lingui{DEF}
\artiklo{buletino}%
\vorto{buletino} Mikra jurnalo, revuo o periodala raporto\lingui{DEFIS}
\artiklo{bulimio}%
\vorto{bulimio}\fako{patol.} Apetito ne-naturala, ne-saturebla, quan produktas forsan tenio od afekto hipokondra, histeria, e c.\lingui{DFIS}
\artiklo{bulino}%
\vorto{bulino}\fako{navig.} Kordo qua tensas la seglo per la infra angulo\lingui{DEF}
\artiklo{buliono}%
\vorto{buliono} Liquido en qua onu boliigis ula substanci\lingui{DEFRS}
\artiklo{bulo}%
\vorto{bulo}\senco{1}{Korpo sfera, generale destinita a rular, a rotacar}\senco{2}{Globeto quan la aero, la gaso quan kontenas liquido, formacas kande lu acensas til la surfaco}\lingui{EFS}
\artiklo{bulvardo}%
\vorto{bulvardo}\senco{1}{En la urbi di qui la fortifikaji esas supresita : promeno-voyo qua -.arbor-lineoza – cirkondas la urbo en la loko quan okupis la terplenajo avan la rempari}\senco{2}{Larja voyo arbor-lineoza, interne di urbo}\lingui{DEFIR}
\artiklo{bumerango}%
\vorto{bumerango} Jet-armo uzata da la aborijeni di Australia, e di qua onu konstatis la existo anke che ula populi di Dekkan e la Egiptiani antiqua\lingui{DEFIRS}
\artiklo{bunta}%
\vorto{bunta} Igita disparata per la asemblo di kolori male asortita
\artiklo{burasko}%
\vorto{burasko} To ne esas ventego, nek tempesto, ma subita e neduranta vento-frapo\lingui{FIS}
\artiklo{burbiliono}%
\vorto{burbiliono} Grumelo blankatra, filamentoza, tenaca, quan onu renkontras funde di la furunkli pusifanta\lingui{F}
\artiklo{burdono}%
\vorto{burdono}\fako{zool.} Speco di abelo, di qua la korpo esas tre voluminoza, tre piloza, qua vivas en galerii subsula, en socii mikra-nombra, remarkinda pro lua bruiso per-rostra, e quan onu konfundas ofte a la « burdono falsa »\latina{bombus}\lingui{F}
\artiklo{burelo}%
\vorto{burelo} Stofo grosiera de lano\lingui{FIS}
\artiklo{bureto}%
\vorto{bureto} Mikra vazo de lado, de kupro, e c., qua kontenas oleo destinita a lubrifikar la mashini, la velocipedi, e c., quan lu lasas ekirar nur gutope\lingui{DEFS}
\artiklo{burgo}%
\vorto{burgo} Granda vilajo ube, ordinare, lokizesas la merkato di la vilaji cirkondanta\lingui{EFS}
\artiklo{burgravo}%
\vorto{burgravo} Komto, sinioro kasteloza, a la judiciado di qua, ordinare, resortisis urbo\lingui{DEFIS}
\artiklo{burjono}%
\vorto{burjono}\fako{bot.} Sproso, rudimenta, qua kontenas jerme la stipi, branchi, folii, flori, e frukti di la vejetanto\lingui{EF}
\artiklo{burleska}%
\vorto{burleska} Di qua la komikeso esas extravago\lingui{DEFIS}
\artiklo{burnuso}%
\vorto{burnuso} Mantelo de lano blanka o bruna, kun kapuco quan *weras la Arabi\lingui{DEFR}
\artiklo{buro}%
\vorto{buro} Amaso de pili, sive de animali, sive de ula stofo lana o silka; to quon onu pozis sur la kargajo di la pafili\lingui{FIS}
\artiklo{burokrato}%
\vorto{burokrato} La persono qua molestoze uzas, a publiko, la povo quan donas a lu lua grado en la kontori di stato-administreyo\lingui{DEFIS}
\artiklo{burso}%
\vorto{burso} Saketo por portar moneto-peci en la posho\lingui{DEFIS}
\artiklo{bushelo}%
\vorto{bushelo} En Francia, olima mezurilo di kapaceso (cirkume 13 litri); nun, la bushelo kontenas -segun la regioni 1 o 2 dekalitri. – En Usa, la kapaceso varias segun la stati\lingui{EFR}
\artiklo{busho}%
\vorto{busho} Grupo, asemblajo de arboreti (nefruktifera)\lingui{DEFI}
\artiklo{busko}%
\vorto{busko} Lamo flexebla ek barto-materio, stalo, e c., qua, per aplikeso muldoza a la kurveso di la pektoro, mantenas rigide la parto avana di korsajo, di jupo-korpo\lingui{EF}
\artiklo{buso}%
\vorto{buso}\fako{bot.} Arboreto di la familio « euforbiacei », di qua la foliaro esas sempre verda\latina{buxus}\lingui{DEFIRS}
\artiklo{busolo}%
\vorto{busolo} Instrumento facita segun la duopla proprajo di la magnet-agulo : direktar sua pinto ad (vers) la nordo-polo, e proximigar su ad vertikalo segun ke lu portesas ad la poli\lingui{DFI}
\artiklo{busprito}%
\vorto{busprito}\fako{navig.} Masto avana di navo, inklinita a la aquo-surfaco, qua portas la seglo avana e fixigita an kordego\lingui{DEFIRS}
\artiklo{busto}%
\vorto{busto}\senco{1}{Che la homo : la kapo e la parto supra di la korpo}\senco{2}{Skulturo qua reprezentas la kapo e la parto supra di la korpo, maxim-multa-kaze sen la brakii}\lingui{DEFIRS}
\artiklo{bustrofedono}%
\vorto{bustrofedono} Skribo-maniero che la Greki antiqua : la linei de literi iras alterne de dextre a sinistre, e de sinistre a dextre. Onu renkontras olu en ula epigrafi, subskriburi, medalii\lingui{DEF}
\artiklo{butar}%
\vorto{butar}\fako{netrans., kontre}\senco{1}{Marchante, pedo-shokar ne-expektite kontre korpo qua salias sur la voyo}\senco{2}{Apogar ulo kontre trabo, arko, e c.}\lingui{EF}
\artiklo{butiko}%
\vorto{butiko} Chambrego, ter-etaje, qua debushas an la strado, ed en qua la komercisti expozas e vendas endetale sua vari\lingui{FI}
\artiklo{butono}%
\vorto{butono} Mikra peco cirkla, plata o konvexa, ek metalo, ivoro, perlomatro, ligno, ulakaze kovrita ye stofo, e qua, fixigite sur parto di vesto, ligas olca ad altra sur qua esas buklo kordona o truo\lingui{EFIS}\subvorto{elektro}{*botono}{Mikra cilindro, ek materio elektronekonduktiva, mantenata en glito-tubeto, kun un resorto, e qua uzesas por igar interkontaktar la extremaji di elektro-filo fluidoza, o por ruptar ta kontakto.}{}{}
\artiklo{butoro}%
\vorto{butoro}\fako{zool.} Speco di herono qua ne esas dresebla por chasado\latina{botaurus}\lingui{F}
\artiklo{butro}%
\vorto{butro} Substanco grasa, manjebla, flava pale, quan onu separas de lakto per batado\lingui{DEFIR}
\artiklo{buvrelo}%
\vorto{buvrelo}\fako{zool.} Pasero di qua la beko esas kona, grosa, kurta e konvexa, la kapo e la ali nigra, la ventro reda, quan onu amansas ed exercas pri siflar\latina{pyrrula}\lingui{F}
\artiklo{buxo}%
\vorto{buxo} Kolektuyo di qua la formo e la dimensioni esas variiva, ek ligno, kartono, metalo, e c., generale portebla, qua havas kovrilo, e destinita a kunportar, transportar od kontenar objekti prezerve\lingui{DE}
\artiklo{buzardo}%
\vorto{buzardo}\fako{zool.} Rapt-ucelo, genero « avidi desnobla », di qua la tarso esas tenua ed alta, havante mi-kolumo de plumi, e qua similesas la buzo\latina{buteo vulgaris}\lingui{DEFS}
\artiklo{buzo}%
\vorto{buzo}\fako{zool.} Raptucelo, genero « avidi desnobla », kun korpo grosiera e quan lua posturo sen-mova dum plura hori (vartante sua raptotajo) igas aspektar stupida\latina{falco buteo}\lingui{DEFIRS}
\end{multicols}
\sekciono{C}
\begin{multicols}{2}
\artiklo{c}%
\vorto{c}\fako{muziko} Nomo di la noto unesma di la « do »-gamo
\artiklo{ca}%
\vorto{ca}\fako{ica} Pronomo qua uzesas por ri-nomar kozo qua jus nomesis o quik nomesos. – Generale, ta \textit{(ita)} uzesas por korespondar a ta senco – ca \textit{(ica)} esante rezervata por indikar : 1. objekto plu proxima; 2. la duesma ek du nomaji (ta \textit{(ita)} koncernas lore la unesma)\lingui{F}
\artiklo{cambium}%
\vorto{« cambium »}\fako{bot.} Parto di la sapo qua, pro solideskar pokope, augmentas la substanco di la basto e di la alburno inter qui lu cirkulas\lingui{DEFIRSL}
\artiklo{caro}%
\vorto{caro}\fako{olim, ante 1915} La suvereno di Rusia\lingui{DEFIRS}
\artiklo{cecidio}%
\vorto{cecidio}\fako{bot.} Nomo generala di la gali quin produktas, sur la vejetanti, ula paraziti animala o vejetanta (insekti, akari, vermi, bakterii)
\artiklo{cedar}%
\vorto{cedar}\fako{trans., ad} Abandonar (ulo ad ulu), renuncante sua yuro\lingui{EFIS}
\artiklo{cedilio}%
\vorto{cedilio} Mikra signo (kelke simila a komo) qua – en la linguo Franca (kom exemplo), pozite sub la litero « c » kande olca sequesas nemediate da la vokali « a », « o », « u », indikas ke ta « c » esas pronuncenda quale « s » che Ido\lingui{DEFIRS}
\artiklo{cedrato}%
\vorto{cedrato}\fako{bot.} Frukto de varietato di citroniero, deveninta de Persia\latina{citrus medica}\lingui{DEFIS}
\artiklo{cedro}%
\vorto{cedro}\fako{bot.} Granda arboro verda-folia, de la familio « koniferi », deveninta de Avan-Azia – di qua la ligno odoras ed esas poke putriva\latina{cedrus}\lingui{DEFIRS}
\artiklo{cedulo}%
\vorto{cedulo} Dokumento per qua onu notifikas ulo\lingui{EFIRS}
\artiklo{cefalopodo}%
\vorto{cefalopodo}\fako{zool.} Molusko di qua la korpo konsistas ye sako membrana, qua kontenas la viceri, e di qua la kapo esas garnisita ye palpili\lingui{DEFIS}
\artiklo{ceko}%
\vorto{ceko}\fako{anat.} Komenco-parto di la grosa intestino, klozita sua-baze per quaza sako-voyo, e qua recevas, de la intestino tenua, la reziduo de la digesto\lingui{EFIRS}
\artiklo{celar}%
\vorto{celar}\fako{trans.} Igar eskapar videso\lingui{EFIS}
\artiklo{celebrar}%
\vorto{celebrar}\fako{trans.}\senco{1}{Exekutar solene (mariajo, meso, e c.)}\senco{2}{Laudar, publikigar solene}\lingui{DEFIS}
\artiklo{celerio}%
\vorto{celerio}\fako{bot.} Apio odoranta, dolcigita per kultivado, e di qua onu manjas la yuna stipi e la bazo di la petioli\latina{apium graveolens}\lingui{DEFR}
\artiklo{celiba}%
\vorto{celiba} Qua vivas sen mariajir su\lingui{DEFIS}
\artiklo{celulo}%
\vorto{celulo}\senco{1}{Mikra chambro (ex. : di monako, di karcerano)}\senco{2}{Alveolo quan konstruktas la abeli}\senco{3}{Mikra kavajo en ula organi di la animali e di la vejetanti}\senco{4}{Elemento anatomiala embriona che la planti o la animali}\lingui{DEFIRS}
\artiklo{celuloido}%
\vorto{celuloido} Kompozajo kemiala, qua imitas skalio, koralio, e c.\lingui{DEFIRS}
\artiklo{celuloso}%
\vorto{celuloso} Korpo neutra ye qua konsistas la tisui di la vejetanti e specale la tisuo celulara ipsa\lingui{DEFIRS}
\artiklo{cembro}%
\vorto{cembro}\fako{bot.} Nomo di speco di pino, di qua la grani esas manjebla\latina{cembro}\lingui{F}
\artiklo{cementacar}%
\vorto{cementacar}\fako{trans.} Pozar stratope, sur metalo, substanco qua, efike da temperaturo tre alta, modifikas lu, adportante a lu un ek sua elementi\lingui{DEFIRS}
\artiklo{cemento}%
\vorto{cemento} Mixuro de kalko e de briki pistita, por interjuntar la petro-bloki en masonuro\lingui{DEFIRS}
\artiklo{ceno}%
\vorto{ceno}\senco{1}{Parto di teatro ube pleas la aktori}\senco{2}{Loko ube eventas la pleajo}\senco{3}{Divido di akto teatrala, determinita da la ekiro di la protagonisti}\lingui{DEFIRS}
\artiklo{cenobito}%
\vorto{cenobito} En la komenco di la historio di la Eklezio Kristana, la personulo qua vivis komune kun altra regulieri (opoze ad « anakoreto »)\lingui{DEFIS}
\artiklo{cenotafio}%
\vorto{cenotafio} Imituro di tombo, erektita memorige di mortinto, e qua ne kontenas lua reliquii\lingui{DEFIS}
\artiklo{censo}%
\vorto{censo} La revenuo evaluata da la autoritati statala, sive por taxo, imposto, sive por ula yuro (elektiveso)\lingui{DEFIRS}
\artiklo{censoro}%
\vorto{censoro} En Roma antiqua, un ek la du judiciisti di qui la tasko konsistis ye okupar su pri la censo, reprimandar e punisar la personi qui violacis la mori\lingui{DEFIRS}
\artiklo{censurar}%
\vorto{censurar}\fako{trans.}\senco{1}{Reprimandar ulu publike}\senco{2}{\textit{(aludante la Eklezio Kristana o korpo di ta Eklezio)} Kondamnar publike doktrino}\lingui{DEFIRS}
\artiklo{cent}%
\vorto{cent} 100\lingui{DEFIRS}
\artiklo{centaureo}%
\vorto{centaureo}\fako{bot.} Planto de la familio « kompozaji », di qua plura speci uzesas medicinale\latina{chlora perfoliata centaurea}\lingui{EFIS}
\artiklo{centauro}%
\vorto{centauro} En la mitologio antiqua, ento mi-vira, mi-kavala\lingui{DEFIRS}
\artiklo{centezimala}%
\vorto{centezimala} Di qua la parti esas centimi\lingui{DEFIRS}
\artiklo{centiaro}%
\vorto{centiaro} La parto centima di aro\lingui{DEFIRS}
\artiklo{centifolio}%
\vorto{centifolio}\fako{bot.} Rozo di qua la petali esas tre multa\latina{rosa centifolia}\lingui{DEF}
\artiklo{centigrado}%
\vorto{centigrado} Dividita a 100 gradi\lingui{DEFIS}
\artiklo{centigramo}%
\vorto{centigramo} Parto centima di gramo\lingui{DEFIRS}
\artiklo{centilitro}%
\vorto{centilitro} Parto centima di litro\lingui{DEFIRS}
\artiklo{centimetro}%
\vorto{centimetro} Parto centima di metro\lingui{DEFIRS}
\artiklo{centimo}%
\vorto{centimo} Parto centima di franko (la unajo monetala di Francia)\lingui{DEFIRS}
\artiklo{centono}%
\vorto{centono} Versajo ek versi o fragmenti di versi pruntita hike ed ibe\lingui{DEF}
\artiklo{centrifugar}%
\vorto{centrifugar}\fako{trans.} Efikigar la forco per qua moveblo qua cirkulas segun kurvo retropulsesas de la centro di ta kurvo\lingui{DEFIRS}
\artiklo{centripetar}%
\vorto{centripetar}\fako{trans.} Efikigar la forco per qua moveblo qua cirkulas segun kurvo atraktesas ad la centro di ta kurvo\lingui{DEFIRS}
\artiklo{centro}%
\vorto{centro}\senco{1}{\textit{(geom.)} Punto interna, situita egaldiste de omna punti di cirklokurvo o di la surfaco di sfero}\senco{2}{To quo esas la centro, la mezo di extensajo}\senco{3}{Apliko-punto di la rezultanto de plura forci}\senco{4}{Loko adube konvergas forci, agi diversa}\lingui{DEFIRS}
\artiklo{centrolecito}%
\vorto{centrolecito} Ovo di qua la plasmo nutriva akumulesas en la centro di la celulo\lingui{DEFIRS}
\artiklo{centrosfero}%
\vorto{centrosfero}\fako{biol.} Maso de protoplasmo qua cirkondas la centrosomo e prizentas ulakaze du zoni sama-centra e neagale koloroza\lingui{DEFIRS}
\artiklo{centrosomo}%
\vorto{centrosomo}\fako{biol.} Mikra korpuskulo qua okupas ordinare la centro di la celulo, apud la parieto nukleala, e qua konjekteble pleas granda rolo en mitoso\lingui{DEFIRS}
\artiklo{centurio}%
\vorto{centurio}\fako{en Roma antiqua}\senco{1}{Asemblitaro de cent civitani ye qui konsistis un ek la dividuri politikala di la populo}\senco{2}{Dividuro de kohorto a cent homi}\lingui{DEFIRS}
\artiklo{centuriono}%
\vorto{centuriono} En epoki antiqua : oficiro di la armeo Romana, qua komandis cent soldati\lingui{DEFIS}
\artiklo{cepo}%
\vorto{cepo}\fako{bot.} Varietato di grosa fungo tre pulpoza, manjebla\latina{boletus edulis}\lingui{F}
\artiklo{ceptro}%
\vorto{ceptro} Komando-bastono, impero-bastono, un ek la insigni rejala, imperiestrala\lingui{DEFIRS}
\artiklo{ceramiko}%
\vorto{ceramiko} Arto (ed industrio) di la fabriko di objekti ek terakoto, fayenco, porcelano, emalio\lingui{DEFIS}
\artiklo{cerastio}%
\vorto{cerastio}\fako{bot.} Genero de « kariofilacei » qua inkluzas plu kam cent speci qui kreskas en omna parti di la Terglobo
\artiklo{cerato}%
\vorto{cerato} Medikamento uzenda korpo-surface, kom unguento, konsistanta ye vaxo dissolvita en oleo\lingui{DEFIS}
\artiklo{cerbero}%
\vorto{cerbero}\fako{mitol.} Hundo tri-kapa qua facionis an la pordo di inferno\lingui{DEFIRS}
\artiklo{cereala}%
\vorto{cereala}\fako{bot.} Di qua la grani, farinoza, uzesas por la nutro di la homo\lingui{DEFIS}
\artiklo{cerebelo}%
\vorto{cerebelo}\fako{anat.} Parto dopa ed infra di la encefalo, qua unionas la cerebro a la spino\lingui{EFIS}
\artiklo{cerebro}%
\vorto{cerebro}\fako{anat.} La encefalo, amaso nervoza, ye qua esas plena la kranio-kavajo, e precipue la parti avana e supra di la encefalo\lingui{DEFIS}
\artiklo{ceremoniar}%
\vorto{ceremoniar}\fako{netrans.} Konformigar su a la formo di la pompo qua akompanas la celebro di solenajo\lingui{DEFIRS}
\artiklo{cerfolio}%
\vorto{cerfolio}\fako{bot.} Planto umbelifera, aromata, di qua la folii uzesas kom kondimenti\latina{anthriacus cerefolium}\lingui{FIL}
\artiklo{cerio}%
\vorto{cerio}\fako{kemio} Korpo simpla metala, ferea, harda, denseso-grado : 6,7\simbolo{Ce}
\artiklo{ceriz-lauro}%
\vorto{ceriz-lauro}\fako{bot.} Arboreto de la familio « apocinacei ». (prunus) laurocerasus\latina{cerasus}\lingui{F}
\artiklo{cerizo}%
\vorto{cerizo}\fako{bot.} Mikra frukto, reda plu o min forte, globa, pulpoza, di qua la kerno esas glata, quan produktas arboro de la familio « rozacei » (di qua la ligno uzesas por moblifado)\latina{prunus cerasus}\lingui{EFS}
\artiklo{cernar}%
\vorto{cernar}\fako{trans.} Enemike cirkondar ed enklozar por kaptar\lingui{DF}
\artiklo{cero}%
\vorto{cero} Granda kandelo de vaxo\lingui{EFIS}
\artiklo{certa}%
\vorto{certa} Qua konsideresas, egardesas, kom nedubitebla, kom vera\lingui{EFIS}\subvorto{}{*certena}{Qua ne dubitas pri la vereso, la exakteso, di ulo}{}{}
\artiklo{certifikar}%
\vorto{certifikar}\fako{trans.} Garantiar kom vera, kom exakta\lingui{EF}
\artiklo{cerumeno}%
\vorto{cerumeno}\fako{anat.} Materio grasatra, flavatra, qua, sekrecite da la folikuli aden la tubeto di la orelo extera, lubrifikas la membrano di olca, ed ibe su amasigas ofte, hardeskante\lingui{DEFIS}
\artiklo{ceruzo}%
\vorto{ceruzo}\fako{kemio} Subkarbonato di plombo, substanco blanka, friabla, quan onu uzas kom fardo, e kom farbo por induto edificala, e de qua onu facas la emalio di porcelano\lingui{EFIS}
\artiklo{cervo}%
\vorto{cervo}\fako{zool.} Mamifero ruminera, di qua la kapo esas garnisita ye korni branchatra\latina{cervus elaphus}\lingui{FIRS}
\artiklo{ces}%
\vorto{« ces »}\fako{muziko} La noto « do » bemoloza
\artiklo{cesar}%
\vorto{cesar}\fako{trans. e netrans.} Ne durar\lingui{EFIS}
\artiklo{cesio}%
\vorto{cesio}\fako{kemio} Metalo alkalioza, korpo simpla\simbolo{Cs}\lingui{DEFIRS}
\artiklo{cesto}%
\vorto{cesto} En la epoki antiqua : rimeno ulakaze garnisita ye plombo, per qua la atleti cirkondis sua manui por la pugno-kombato\lingui{DEFI}
\artiklo{cetaceo}%
\vorto{cetaceo}\fako{zool.} Granda mamifero marala, fishatra (baleno, kashaloto, e c.)\lingui{DEFIS}
\artiklo{cetera}%
\vorto{cetera} Ta qua restas kande finis la enumero, la depreno\lingui{FL}
\artiklo{cetonio}%
\vorto{cetonio}\fako{zool.} Koleoptero di qua la kolori esas metalatra quale la kantarido\latina{cetonia}
\artiklo{cezaro}%
\vorto{cezaro} Imperiestro\lingui{DEFIRS}
\artiklo{cezuro}%
\vorto{cezuro} Divido ritmala di verso, indikita da silabo qua finas vorto sur qua onu devas plufortigar la pronunco\lingui{DEFIS}
\artiklo{chagrenar}%
\vorto{chagrenar}\fako{netrans.} Sentar bitre doloro psikala\lingui{F}
\artiklo{chaketo}%
\vorto{chaketo} Ludo analoga a la triktrako\lingui{FS}
\artiklo{chako}%
\vorto{chako} Kapvesto kun viziero. ek materio rigida, ordinare cilindratra\lingui{DEF}
\artiklo{chalazo}%
\vorto{chalazo}\senco{1}{Hilo interna di la vejetanto}\senco{2}{Punto jermala di la ucel-ovo}\lingui{DF}
\artiklo{chalenjo}%
\vorto{chalenjo}\fako{premio} En konkurso en qua la premio ne donesas, ma prestesas dum tempo a la vinkinto, qua havas lu til kande altra persono vinkas e konquestas la premio\lingui{F}
\artiklo{chamado}%
\vorto{chamado} Tambur-agar, trumpet-agar por avertar la enemiko ke onu deziras parlementar o kapitulacar\lingui{DEFRS}
\artiklo{chamaro}%
\vorto{chamaro} Speco di redingoto kun kordoni transversa an-avane, vesto Polona de epoko revolucionala di la XIX-esma yarcento
\artiklo{chambelano}%
\vorto{chambelano} Oficiro komisita por okupar su pri la chambro di rejo, di princo, qua *weris klefo kom insigno\lingui{EFIS}
\artiklo{chambro}%
\vorto{chambro}\senco{1}{Parto di lojeyo}\senco{2}{Salono ube onu su asemblas por deliberar}\senco{3}{En la guverno-sistemi reprezentala : singla de la du asemblitari qui posedas la yuro legifar same kam la chefo di la exekutigo-autoritato}\senco{4}{\textit{(teknologio)} Nomo di kavaji diversa qui plu o min similesas chambro}\lingui{EFI}
\artiklo{chamo}%
\vorto{chamo}\fako{zool.} Quadripedo ruminera, de la genero « antilopi », kun korni kava e glata, quan onu renkontras en la alta monti di Alpi\latina{rupicapra tragus}\lingui{EF}
\artiklo{champanio}%
\vorto{champanio} Vino de Champania (Francia)\lingui{DEFIRS}
\artiklo{champiniono}%
\vorto{champiniono}\fako{bot.} Genero de fungi di la familio « agaricinei », karakterizata da la prezenteso di ringo ye la supro di la pedo, lami libera rozea, pose nigra, t.e. departanta de la bordo di la chapelo e qui ne arivas til la pedo, e spori bruna purpurea\lingui{F}
\artiklo{championo}%
\vorto{championo} Singla de la adversi qui kombatis en lico. \textit{(metaf.)} La persono qua konsakras su a la defenso di ideo, di idealo\lingui{DEF}
\artiklo{chanco}%
\vorto{chanco} Probableso favoroza o desfavoroza\lingui{DEFI}
\artiklo{chanfreno}%
\vorto{chanfreno} Parto avana di la kapo di kavalo, de la fronto til la naztrui\lingui{EF}
\artiklo{chanjar}%
\vorto{chanjar}\fako{trans. e netrans.}\senco{1}{Substitucar ula kozo ad altra en un loko o funciono}\senco{2}{Prenar kozo vice altra, o pos altra}\senco{3}{Divenar altra}\lingui{EFL}
\artiklo{chankro}%
\vorto{chankro}\fako{patol.} Ulcero qua tendencas rodar la parti qui cirkondas lu\lingui{DEFR}
\artiklo{chanto}%
\vorto{chanto} Edro laterala, streta, di korpo larja e plata\lingui{DF}
\artiklo{chapelo}%
\vorto{chapelo}\fako{I. Kapovesto -virala}\senco{1}{Qua konsistas ordinare ye korpo rigida o mi-rigida por recevar la kapo, kun rebordo cirklo-kurva plu o min larja}\senco{2}{Kapovesto (mulierala) qua konsistas generale ye korpo por recevar la kapo, o kaloto, kun rebordo}\lingui{FI}
\artiklo{chapitro}%
\vorto{chapitro} Singla de la parti ye qua konsistas la dividuri di libro, kodexo, kontrato, budjeto, e c.\lingui{DEFIRS}
\artiklo{chapleto}%
\vorto{chapleto} Asemblajo de kin deki de grani filizita, separita per grano plu grosa\lingui{EFL}
\artiklo{chapo}%
\vorto{chapo}\senco{1}{Mantelo ceremoniala quan *weras la episkopi, la kantori, la sacerdoti por ula oficii solena}\senco{2}{\textit{(tekn.)} Peco qua parkovras ulo}\lingui{DFS}
\artiklo{charioto}%
\vorto{charioto} Kargo-veturo kun quar roti, rideloza ed un-timona, uzata en la ruro por transportar la rekoltaji, la vari vendenda, e c. – en la armei, por transportar la bagaji, e c.\lingui{F}
\artiklo{charjar}%
\vorto{charjar}\fako{trans.}\senco{1}{Lokizar (veturo, e c.) sub la pezo di la objekti (kargajo) portenda ulube, transportenda}\senco{2}{Garnisar (ulo) ye pezajo, ye ula quanto, destinita a ta o ca uzo (fusilo, pipo, bobino, e c.)}\lingui{EFIS}
\artiklo{charmar}%
\vorto{charmar}\fako{trans.} Ravisar, seduktar per atraktivo forta\lingui{EF}
\artiklo{charniro}%
\vorto{charniro} Ligilo ek du fero-plaketi, artikoza, de qui adminime un povas movar cirkum axo\lingui{EFIRS}
\artiklo{charo}%
\vorto{charo}\senco{1}{Veturo lejera di la antiqui, kun du roti, aperturoza dope, tirata da du o plura kavali jungita tale ke li avancas fronte, e sur la parto avana di qua la duktanto stacas}\senco{2}{Veturo dekorita, en la festi, procesioni, masko-festi, e qua portas personegi simbola, historiala, maskiti, e c.}\lingui{DEFIS}
\artiklo{charto}%
\vorto{charto} En la mezepoko, omna akto en qua enrejistresis la akti di proprieto, vendo, privileji grantita, e c.\lingui{DEFIRS}
\artiklo{chasar}%
\vorto{chasar}\fako{trans.} Proxime-sequar la animali por kaptar o destruktar li\lingui{EFIS}
\artiklo{chasta}%
\vorto{chasta} Qua konservas sexuo-pura sua korpo e sua psiko\lingui{EFIS}
\artiklo{chazublo}%
\vorto{chazublo} Vesto sacerdotala per qua kovresas la sacerdoto dum ke lu celebras la santa sakrifiko, qua pendas sur la albo, ed havas kom ornivo precipua la figuro di la Kruco (kristana)\lingui{EF}
\artiklo{che}%
\vorto{che} En la domo, habiteyo, lando, domeno (materiala o spiritala – inter ta o ca sorto de grupani)\lingui{F}
\artiklo{chefo}%
\vorto{chefo} La persono qua esas ye la kapo di grupo quan lu duktas ed a qua lu imperas\lingui{DEFRS}
\artiklo{cheko}%
\vorto{cheko}\fako{banko} Folieto separita de kayero talonoza, kompozita ek bilieti pagebla ye vido \textit{(da banko)} e qua, pro pagendeso a la portanto, esas indosebla\lingui{DEFRS}
\artiklo{chelato}%
\vorto{chelato} Kasko di qua la kaloto esas mi-sferatra, la viziero kurta e fixa, la nukumo granda, quan (en la mezepoko ed en la yarcento XVI-esma) *weris la militisti kavaliera\lingui{EFIS}
\artiklo{chera}%
\vorto{chera} Qua kustas plu multe kam lo analoga\lingui{FIS}
\artiklo{cherpar}%
\vorto{cherpar}\fako{trans., ek, de} Prenar ek puteo, fonto, barelo, e c. kelke de la liquido quan lu kontenas. – \textit{(anke metaf.)}
\artiklo{*chevaliero}%
\vorto{*chevaliero}\fako{en Roma, antique}\senco{1}{Civitano di un ek la tri stato-klasi, inter la patrico e la plebeyo}\senco{2}{\textit{(mezepoko)} Membro di ta institucuro milito-tendencema, kun karaktero religiala, establisita che la nobelaro feudala, qua impozis a sua membri la desegardo di danjero, loyaleso, la protekto gratuita di la febli, e politegeso a la dami}\senco{3}{\textit{(nun)} Membro di ordeno honorumala, institucita da suvereno}\senco{4}{\textit{(nun)} En ordeno honorumala plura-klasa, membro di la klaso minim-alta}\lingui{I}
\artiklo{chevrono}%
\vorto{chevrono}\senco{1}{Peco ek ligno, fixigita an la pento di tekto, e suportanta la lati qui subtenas la tekto-kovruro, la teguli, la ardezo-plaki}\senco{2}{Bendi plata, dispozita V-forme sur la skudo, la blazoni}\lingui{EFRS}
\artiklo{chifrar}%
\vorto{chifrar}\fako{trans.} Kompozar kripto-grafuro, t.e. texto sekreta, qua tradukas la vera texto per ula moyeno o regulo (« klefo »)\lingui{DEF}
\artiklo{chigo}%
\vorto{chigo}\fako{zool.} Insekto di Sud-Amerika\latina{sarcopsylla penetrans}\lingui{EFS}
\artiklo{chika}%
\vorto{chika} Qua havas gracio eleganta, audacoza\lingui{DEFR}
\artiklo{chilo}%
\vorto{chilo}\fako{fiziol.} Suko qua konsistas (en la intestino tenua) ye la parto nutriva di la bulo nutrala, qua divenas sango\lingui{DEFIRS}
\artiklo{chimo}%
\vorto{chimo} Sorto di paplo quan produktas, en la stomako, la elaboro komencala di la nutrivi\lingui{DEFIRS}
\artiklo{chimpanzeo}%
\vorto{chimpanzeo}\fako{zool.} Simio alta-statura\latina{troglodytea niger}\lingui{DEFIRS}
\artiklo{chiniono}%
\vorto{chiniono} Che la muliero : la parto di la hararo qua amasigesas ed erektesas dop la kapo\lingui{DEF}
\artiklo{chipa}%
\vorto{chipa} Qua kustas min multe kam lo analoga\lingui{E}
\artiklo{chitino}%
\vorto{chitino}\fako{kemio} Elemento precipua di la kuraso di la « artropodi »\simbolo{C₁₆, H₂₆, N₂, O₁₀}\lingui{DEFIS}
\artiklo{chitozamino}%
\vorto{chitozamino}\fako{kemio} Bazo qua formacesas per hidrolizo di la chitino e di multa albuminoidi
\artiklo{choano}%
\vorto{choano}\fako{anat.} Orifico dopa di la naztrui\lingui{DEFIS}
\artiklo{chokolado}%
\vorto{chokolado} Substanco nutrala, ek kakao-grani, pistita kun sukro\lingui{DEFIRS}
\artiklo{chomar}%
\vorto{chomar}\fako{netrans.} Nelaborar pro neemployateso nevolata\lingui{F}
\artiklo{chopino}%
\vorto{chopino} Birglaso qua kontenas cirkume un mi-litro\lingui{DEF}
\artiklo{choseo}%
\vorto{choseo}\senco{1}{Bendo de tereno, ofte stonizita, qua dominacas rivero, lageto, marsho, quan lu transiras, ed uzata kom voyo, paseyo}\senco{2}{La mezo (generale pavizita o stonizita) di voyo, strado}\lingui{DF}
\artiklo{chuvo}%
\vorto{chuvo}\fako{zool.} Mikra korniko quan onu renkontras en la kloshturmi. F. ci. \textit{(gram.)} Pluralo di ca\latina{corvus monedula}
\artiklo{cianato}%
\vorto{cianato}\fako{kemio} Salo od etero di la cianacido
\artiklo{cianhidra}%
\vorto{cianhidra}\fako{kemio} Produktita da la kombino di ciano kun oxo\simbolo{H. CN}\lingui{EF}
\artiklo{ciano}%
\vorto{ciano}\fako{kemio} Gaso senkolora, produktita da la kombino di azoto kun karbo, qua kondutas quale korpo simpla en la reakti kemiala\simbolo{N : C, C}\lingui{DEF}
\artiklo{ciatika}%
\vorto{ciatika}\fako{anat.} Qua relatas la grosa nervo qua iras de la plexuso sakrumala til la muskuli dopa di la kruro, di la suro e di la pedo\lingui{EF}
\artiklo{cibolo}%
\vorto{cibolo}\fako{bot.} Planto legum-gardenala, analoga a la oniono, quan onu uzas koquarte kom kondimento\latina{allium fistulosum}\lingui{DEFIRS}
\artiklo{ciborio}%
\vorto{ciborio}\fako{liturgio katol.} Vazo kalico-forma, kun kovrilo, qua kontenas la hostii konsakrita por la komuniono di la Kristo-Ekleziani\lingui{DEFI}
\artiklo{cicisbeo}%
\vorto{cicisbeo} Amoranto di damo, qua esas quaze oficale agnoskata e servas elu publike\lingui{DEFIRS}
\artiklo{cidro}%
\vorto{cidro} Drinkajo ek pomo-suko fermentacita\lingui{DEFIRS}
\artiklo{cielo}%
\vorto{cielo}\senco{1}{Spaco en qua jiras omna astri}\senco{2}{Parto de la spaco qua aspektas quale vulto mi-sfera di qua la horizonto esas la limito}\lingui{FIS}
\artiklo{cienco}%
\vorto{cienco}\senco{1}{Savo exakta e sistemigita pri ula temo o grupo de temi, qua relatas la mondo fizikala od un ek la arti liberala (gramatiko, logiko, muziko, e c.)}\senco{2}{Sistemo de savaji, en qua fako de fakti preciza koordinesas e reduktesas a legi (*leyi)}\lingui{EFIS}
\artiklo{cifro}%
\vorto{cifro} Tipo o signo qua reprezentas nombro (Araba o Romana)\lingui{DEFIRS}
\artiklo{cigano}%
\vorto{cigano} Membro di tribui vaganta qui venis de Oriento\lingui{DFIS}
\artiklo{cigno}%
\vorto{cigno}\fako{zool.} Ucelo palmipeda aquala, analoga a ganso, di qua la kolo e la beko esas longa, remarkinda pro la gracio e la eleganteso di lua konturi e movi\latina{cygnus}\lingui{EFIS}
\artiklo{cikado}%
\vorto{cikado}\fako{zool.} Insekto hemiptera qua emisas bruisado akra e monotona produktita da la interfroto di du membrani situita en la abdomino\latina{cicada}\lingui{DEFIRS}
\artiklo{cikatro}%
\vorto{cikatro} Marko poslasita sur la pelo da vunduro, bruluro, pos risanesko\lingui{EFIS}
\artiklo{ciklameno}%
\vorto{ciklameno}\fako{bot.} Planto herbacea, perena, de la familio « primulacei », di qua la folii esas ronda kordio-forme e purpurea, la floro pendanta blanka e purpura\latina{cyclamen}\lingui{DEFIS}
\artiklo{ciklido}%
\vorto{ciklido}\fako{geom.} Areo di la klaso quaresma, qua admisas kom lineo duopla la cirklo di la infinito, e posedas dek serii de secioni cirkla\lingui{DEFIRS}
\artiklo{ciklito}%
\vorto{ciklito}\fako{patol.} Varietato di iriso-koroidito qua afektas la parto dopa ed extera di la koroido (parto qua havas prolonguri pasable multa)\lingui{DEFIS}
\artiklo{ciklo}%
\vorto{ciklo}\senco{1}{Periodo kontinua ek ula quanto de yari, dum qua ula fenomeni astronomiala eventas sam-ordine}\senco{2}{Serio de modifikesi di korpo qua pasas tra estadi diversa, e retrovenas a sua departo-punto}\senco{3}{Serio de poemi pri un epoko, un temo, e c.}\senco{4}{Lineo spirala quan formacas, cirkum stipo o rameto, la origino-punto di folii qui inter-korespondas}\lingui{DEFIRS}
\artiklo{cikloido}%
\vorto{cikloido}\fako{geom.} Kurvo trasata da la par-jiro di punto qua apartenas a cirklo-kurvo qua rulas sur rekto fixa\lingui{DEFIRS}
\artiklo{ciklometrio}%
\vorto{ciklometrio} Arto mezurar la cirkli\lingui{DEF}
\artiklo{ciklono}%
\vorto{ciklono} Sturmo-vento qua vorticas, atraktante (per forco nerezistebla) irgo quon lu renkontras sur tero e sur maro\lingui{DEFIRS}
\artiklo{ciklopo}%
\vorto{ciklopo} Giganto mitologiala, reprezentata kom havanta nur un okulo (meze di sua fronto)\lingui{DEFIRS}
\artiklo{cikloso}%
\vorto{cikloso}\fako{biol.} Cirkulado en-celula che la planti (kontraste kun la cirkulado generala di la planti)
\artiklo{cikonio}%
\vorto{cikonio}\fako{zool.} Ucelo migrera, de la klaso « steltieri », di qua la kolo, la beko, la gambi esas tre longa\latina{cyconia}\lingui{FISL}
\artiklo{cikorio}%
\vorto{cikorio}\fako{bot.} Planto legum-gardenala, di qua la folii esas mikra e frizita, quan onu manjas kom legumo o salado, de qua onu facas infuzuri\latina{cichorium}\lingui{DEFIRS}
\artiklo{cikuto}%
\vorto{cikuto}\fako{bot.} Planto perena, venenoza, de la familio « umbeliferi »\latina{conium maculatum, e speco di cicuta}\lingui{I}
\artiklo{ciliaro}%
\vorto{ciliaro}\fako{anat.} Parto avana ed extera di la koroido, qua havas prolonguri pasable multa
\artiklo{cilico}%
\vorto{cilico}\fako{religio katolika} Kamizo, zono de krini, *werita sur la pelo, por penitencar o su mortifikar\lingui{DEFIS}
\artiklo{cilindro}%
\vorto{cilindro}\senco{1}{Solido produktita da rekto movebla qua rotacas cirkum sua axo a qua lu esas paralela}\senco{2}{Rolero de petro, de metalo, de ligno, quan onu igas rotacar cirkum lua axo, e sub qua onu pasigas to quon onu volas submisar a preso uniforma}\lingui{DEFIRS}
\artiklo{cilindroido}%
\vorto{cilindroido}\fako{geom.} Solido qua similesas cilindro, ma di qua la bazi esas elipsi vice cirkli\lingui{DEFIRS}
\artiklo{cilio}%
\vorto{cilio}\fako{anat.}\senco{1}{Pilo qua garnisas la bordo libera di la palpebro (che la mamiferi)}\senco{2}{\textit{(bot.)} Sorto di pilo silkatra qua bordumas ula parti di la vejetanti}\lingui{F}
\artiklo{cimaso}%
\vorto{cimaso}\fako{arkitekt.} Muluro ye la supro di kornico\lingui{DEFIS}
\artiklo{cimbalo}%
\vorto{cimbalo}\fako{muziko}\senco{1}{\textit{(che la antiqui)} Instrumento ek du miglobi bronza, kava, quan onu tenis per ringo fixigita ye la suprajo, e quin onu intershokis}\senco{2}{\textit{(nun)} Instrumento ek du diski inter-egala, kupra, quin onu interfrapas}\lingui{DEFRS}
\artiklo{cimitaro}%
\vorto{cimitaro} Sabro di qua la lamo esas larja, kurva, uzata che la Orientani\lingui{EFIS}
\artiklo{cimo}%
\vorto{cimo}\fako{zool.} Insekto di qua la korpo esas plata, la odoro fetida, qua sugas la sango di la personi dormanta\latina{cimex lectularium}\lingui{IL}
\artiklo{cinabro}%
\vorto{cinabro}\fako{kemio} Merkurio-sulfuro, reda\simbolo{HgS}\lingui{DEFIRS}
\artiklo{cinamo}%
\vorto{cinamo} Che la antiqui : substanco aromata quan onu dicas esir miro\lingui{DER}
\artiklo{cindro}%
\vorto{cindro} Reziduo pulveratra, de la kombusto di ula materii : ligno, terkarbono, e c.\lingui{EFI}
\artiklo{cinematiko}%
\vorto{cinematiko} La parto di mekaniko qua koncernas la legi (*leyi) di movo konsiderata pure, sole\lingui{DEFIRS}
\artiklo{cinematografar}%
\vorto{cinematografar}\fako{trans.} Enrejistrar cinemo-filmo\lingui{EF}
\artiklo{cinemo}%
\vorto{cinemo} Salono, chambrego, publika ube projektesas, sur *skrino, fotografuri di qui la figuri esas quaze-vivanta (ek filmo)\lingui{EF}
\artiklo{cinerario}%
\vorto{cinerario}\fako{bot.} Planto de la familio « kompozaji », deveninta de Kabo di Bona Espero\latina{cineraria}\lingui{DEFIS}
\artiklo{cinetiko}%
\vorto{cinetiko}\fako{fiziko} Teorio pri ensemblo de fenomeni, fondita exkluzive sur la movo di materio\lingui{DEFIRS}
\artiklo{cinezo}%
\vorto{cinezo}\fako{biol.} Divideso mediata di la nukleo, en qua la kromosomo e la centrosomi displasesas\lingui{DEFIRS}
\artiklo{cinika}%
\vorto{cinika}\senco{1}{Qua apartenas a la skolo filozofiala quan fondis Antistenes, tale nomizita pro ke olu esis establisita komence en Cinosargos, pre-urbo di la urbo Athinai}\senco{2}{Qua su lokizas super deco, la konveno-reguli, e pudoro}\lingui{DEFIS}
\artiklo{cinocefalo}%
\vorto{cinocefalo}\fako{zool.} Genero de simio, di qua la muzelo esas longa quale che la hundo\latina{cynocephalus}\lingui{EFIS}
\artiklo{cintilo}%
\vorto{cintilo} Peceto flama, lumoza, qua su separas de korpo kombustata\lingui{EFIS}
\artiklo{cipo}%
\vorto{cipo}\fako{arkitekt.} Mi-kolono sen kapitelo, qua simulas kolono ruptita, quan la antiqui erektis kom monumento funerala, quan li lokizis sur la voyi por indikar la direcioni, od en la loki konsakrita da ta o ca eventaji memorinda\lingui{DEFIRS}
\artiklo{cipreso}%
\vorto{cipreso}\fako{bot.} Arboro rezinala, de la familio « koniferi », quan lua foliaro obskura igas selektar por plantaceso en la tombeyi, apud la tombi\latina{cypressus}\lingui{DEFIRS}
\artiklo{ciprino}%
\vorto{ciprino}\fako{zool.} Genero de fisho qua vivas en aquo sen-sala, di qua la tipo esas la karpo, e di qua la koloro esas pasable oranjea o redatra\latina{cyprinus}\lingui{EFIS}
\artiklo{cirajo}%
\vorto{cirajo} Kompozajo nigra, quan onu extensas sur la ledro di la pedovesti, di la harnesi, e quan onu frotas per brosilo por igar lu brilanta\lingui{FIS}
\artiklo{circinata}%
\vorto{circinata}\fako{dermatologio} \textit{(lezuri di pelo)} , grupigita cirklatre o cirklo-fragmentatre, di qua la centro esas sendomaja\lingui{EFS}
\artiklo{cirklo}%
\vorto{cirklo}\fako{geom.} Parto de plano, di qua la limito esas kurvo klozita di qua omna punti distas egale de punto interna (centro)\lingui{EFIRS}
\artiklo{cirko}%
\vorto{cirko}\senco{1}{Klozajo cirkla ube onu celebris (che la antiqui) la ludi publika}\senco{2}{Sorto di teatro, cirkoforma, ube la spektajo konsistas ye prodalaji di kavalko, di voltijo}\lingui{DEFIRS}
\artiklo{cirkoncizar}%
\vorto{cirkoncizar}\fako{trans.} Distranchar la prepuco\lingui{EFIS}
\artiklo{cirkondar}%
\vorto{cirkondar}\fako{trans.}\senco{1}{Esar cirkum}\senco{2}{Garnisar cirkume}
\artiklo{cirkonflexo}%
\vorto{cirkonflexo} En la linguo Franca, Portugalana, e c. : signo ortografiala (\textasciicircum{}) quan onu lokizas super ula vokali longa (âsuê\{\textbackslash\{\}îexôasciicircum\{\}\}\{a\}, ê, î, ô)\lingui{DEFIRS}
\artiklo{cirkonspekta}%
\vorto{cirkonspekta} Qua surveyas prudente, e desfideme, to quon lu dicas od agas, o to quon dicas od agas la kunhomi\lingui{EFIS}
\artiklo{cirkonstanco}%
\vorto{cirkonstanco}\senco{1}{Singla de la fakti di eventajo o di situeso}\senco{2}{Fakto acesora, qua akompanas eventajo precipua}\lingui{EFIS}
\artiklo{cirkonvalaciono}%
\vorto{cirkonvalaciono}\fako{milit.} Trancheo kun palisadi, parapeti, fortifikajo tempala per qua armeo envelopas la placo *asiejata, por cesigar la relati di olca kun la mondo extera, o per qua lu su protektas kontre la ataki de la trupi qui venas por sokursar la placo *asiejata\lingui{EFIRS}
\artiklo{cirkonvoluciono}%
\vorto{cirkonvoluciono} Envolveso – sinui cirkloforma\lingui{EFIS}
\artiklo{cirkuito}%
\vorto{cirkuito}\senco{1}{Spaco parkurenda por cirkumirar komplete loko}\senco{2}{Iro cirkuma, vice iro tra la voyo rekta}\lingui{EFIS}
\artiklo{cirkular}%
\vorto{cirkular}\fako{netrans.}\senco{1}{Irar e retroirar omna-direcione sinue}\senco{2}{\textit{(pri moneto)} Aceptesar da omna homi di la regiono a qua lu esas destinita, kom valoroza}\lingui{DEFIRS}
\artiklo{cirkulero}%
\vorto{cirkulero} Letro de qua un kopiuro uniforma sendesas sama-instante a plura (e, generale : a tre multa) personi\lingui{DEFIRS}
\artiklo{cirkum}%
\vorto{cirkum} En la loko, tempo o quanto qua esas vicina, apuda – kompare a la objekto, quaze ringe exter lu\lingui{DL}
\artiklo{cirkumnutaciono}%
\vorto{cirkumnutaciono}\fako{bot.} Parkuro helicoida da la extremajo di la radiko, o di la stipo, di vejetanto
\artiklo{cis}%
\vorto{cis} Ca-latere di – ne trans lu\lingui{L}
\artiklo{cisipara}%
\vorto{cisipara}\fako{biol.} Qua su riproduktas per sua divideso a du parti, generale segun la mikra axo di la korpo\lingui{DEFIRS}
\artiklo{cisoido}%
\vorto{cisoido}\fako{geom.} Kurvo di la grado triesma, konsistanta ye du branchi simetra qui departas de un punto e tendencas a tangento komuna, sen atingar lu ul-instante\lingui{DEF}
\artiklo{cisterciano}%
\vorto{cisterciano}\fako{religio kristana} Reguliero qua resortias a la ordeno religiala qua kreesis en la Franca urbo Citeaux\lingui{DEFIRS}
\artiklo{cisterno}%
\vorto{cisterno} Granda tanko a qua duktesas e rekoliesas la aquo pluvala\lingui{DEFIRS}
\artiklo{citadelo}%
\vorto{citadelo} Kastelo forta qua protektas urbo\lingui{DEFIRS}
\artiklo{citar}%
\vorto{citar}\fako{trans.} Adportar voce o skribe extraktaji de texto por sekondar to quon onu asertas\lingui{DEFIRS}
\artiklo{citaro}%
\vorto{citaro}\fako{muziko} Sorto di muzik-instrumento kordoza, quan uzis la antiqui\lingui{DEFIRS}
\artiklo{citizo}%
\vorto{citizo}\fako{bot.} Genero de planti di qua la tipo esas la Alpo-citizo o « falsa ebeniero »\latina{cytisus}\lingui{DEFIRS}
\artiklo{citodo}%
\vorto{citodo}\fako{biol.} Maso protoplasma nuda, sen-nuklea (ex. : la monero di Haeckel)
\artiklo{citologio}%
\vorto{citologio}\fako{biol.} Parto di histologio di qua la studio-temo esas la celuli, de la vidopunti konstituceso, morfologio ed evoluciono\lingui{DEFIRS}
\artiklo{citoplasmo}%
\vorto{citoplasmo}\fako{biol.} Protoplasmo, situita exter la nukleo di la celulo vivanta\lingui{DEFIRS}
\artiklo{citotropismo}%
\vorto{citotropismo}\fako{biol.} Atrakto di celulo ad altra celuli vivanta, od ad ula substanci kemiala sen-viva\lingui{DEFIRS}
\artiklo{citrono}%
\vorto{citrono}\fako{bot.} Frukto flava klare, saporanta acide, de arboro di la genero « oranjo »\latina{citrus}\lingui{DEF}
\artiklo{civeto}%
\vorto{civeto}\fako{zool.} Mamifero karnivora, analoga a la martro, qua sekrecas quaza mosko\latina{viverra zibetha}\lingui{DEFIRS}
\artiklo{civila}%
\vorto{civila} Qua relatas la civitani. (antonimo di armeala, religiala, kriminala, politikala).\lingui{DEFIRS}\subvorto{}{estado civila}{Ensemblo de la dokumenti relativa a la estado sociala di singla civitano : naskala akto, mariajala akto, e c.}{}{}
\artiklo{civilizar}%
\vorto{civilizar}\fako{trans.} Pasigar de la estado primitiva, naturala, ad estado plu progresoza pri la kultiveso etikala, morala, pensala e sociala\lingui{DEFIRS}
\artiklo{civismo}%
\vorto{civismo} Devoteso di la civitano a la Stato\lingui{DEFIS}
\artiklo{civito}%
\vorto{civito}\senco{1}{Urbo, konsiderata kom korpo politikala}\senco{2}{La korpo ek la civitanaro}\lingui{IS}
\artiklo{cizelar}%
\vorto{cizelar}\fako{trans.} Laborar metalajo (bronzo, oro, arjento, e c.) per instrumento ek stalo, longa, plata, tranchiva ye un extremajo e quan onu uzas por entaliar e nochizar\lingui{DEFIRS}
\artiklo{cizo}%
\vorto{cizo} Instrumento qua konsistas ye du lami di qui la bordi tranchanta interkrucumas, ed uzata por tranchar la kozi dina (stofo, papero, filo, e c.)\lingui{EFI}
\artiklo{co}%
\vorto{co} Ca kozo
\end{multicols}
\sekciono{D}
\begin{multicols}{2}
\artiklo{d}%
\vorto{d}\fako{muziko} Noto duesma di la « do »-gamo
\artiklo{da}%
\vorto{da} Indikas la individuo, objekto, kozo o kauzo qua igas, facas od efektigas ulo
\artiklo{dafno}%
\vorto{dafno}\fako{bot.} Planto, genero de la familio « timelei », analoga a lauro
\artiklo{dago}%
\vorto{dago} Poniardo olima, di qua la lamo akuta povis penetrar tra la lakuni di la kuraso o tra la mashkuti\lingui{EFIS}
\artiklo{dahlia}%
\vorto{« dahlia »}\fako{bot.} Planto exotika, de la familio « kompozaji », kultivata en nia gardeni pro lua flori di qui la kapituli esas nereguloza, la kolori brilanta e diversa\latina{Dahlia}\lingui{DEF}
\artiklo{daimo}%
\vorto{daimo}\fako{zool.} Genero di animalo, analoga a cervo, ma plu mikra, e di qua la supra parto di la kornobranchi esas aplastita e palmoida\latina{Cervus dama}\lingui{F}
\artiklo{daktilo}%
\vorto{daktilo} Verso-pedo ek un silabo longa, sequata da du kurta\lingui{DEFIS}
\artiklo{dalmatiko}%
\vorto{dalmatiko}\fako{liturgio} Ornivo kirko-vestala, sorto di chazublo sen-kruca, quan metas la diakono e la subdiakoni qui helpas la sacerdoto\lingui{DEFIRS}
\artiklo{daltonismo}%
\vorto{daltonismo}\fako{patol.} Stando anormala di la vid-organo, qua impedas dicernar ula kolori, precipue la kolori komplementala\lingui{DEFIRS}
\artiklo{damaskinar}%
\vorto{damaskinar}\fako{trans.} Inkrustar oro, arjento (en peco de fero o de stalo)\lingui{DEFIRS}
\artiklo{damasko}%
\vorto{damasko} Stofo de silko, di qua la texuro prizentas flori, desegnuri; stofo, linjo, qua prizentas ornamenti analoga\lingui{DEFIRS}
\artiklo{damijano}%
\vorto{damijano} Tre grosa botelo fera, gresa, terakota, di qua la ventro esas larja, la kolo kurta, maxim-multa-kaze envelopita ye ozier-plektajo, ed ansoza, uzata por transportar liquidi\lingui{EFIS}
\artiklo{damnar}%
\vorto{damnar}\fako{trans.} Kondamnar a la punisi pos-morta ed eterna (religio kristana)\lingui{DEF}
\artiklo{damo}%
\vorto{damo}\senco{1}{Muliero spozoza}\senco{2}{Muliero vidva}\senco{3}{En la karto-ludo : figuro qua reprezentas rejino}\senco{4}{En la shako-ludo : ta rejino qua unika darfas serchar e kaptar omna-direcione ed irga-diste}\senco{5}{En la piono-ludo : piono qua arivinte a la lasta lineo di la piono-plako, duopligesas por distingar lu de la ceteri, e darfas marchar e kaptar diagonale irga-diste}\lingui{DEFIS}
\artiklo{damzelo}%
\vorto{damzelo} Titulo qua aplikesas a la celibini\lingui{EF}
\artiklo{dandio}%
\vorto{dandio} Viro qua ostentas sua eleganteso segun-moda pri sua vesti e gesti\lingui{DEFR}
\artiklo{danjero}%
\vorto{danjero} To quo minacas la interesti, la sekureso di ulu\lingui{EF}
\artiklo{dankar}%
\vorto{dankar}\fako{trans.) (pri, pro} Agnoskar recevir bonfaco\lingui{DE}
\artiklo{dansar}%
\vorto{dansar}\fako{netrans.} Agar serio de pazi, de salti, de posturi, ritmoza e kadencoza, maxim-multa-kaze segun la soni de muzik-instrumento\lingui{DEFRS}
\artiklo{dardo}%
\vorto{dardo} Lans-armo antiqua, bastono kun fera pinto triangula, quan onu lansis pos brakio-ociligir por augmentar la impulso-forco\lingui{EFIS}
\artiklo{darfar}%
\vorto{darfar}\fako{trans.} Havar la yuro o permiso. (Ta vorto havas la senco legala, kontrea a la ideo di interdikto)\lingui{D}
\artiklo{darnar}%
\vorto{darnar}\fako{trans.} Modo specala di rapecar stofi, vesti, ne per pozar peco de stofo, ma imitante la texuro en la loko truoza\lingui{E}
\artiklo{datelo}%
\vorto{datelo}\fako{bot.} Frukto sukroza, ovala, di qua la grano esas longa, e quan produktas ula sorto di palmiero\latina{phoenix dactylifera}\lingui{DEFIS}
\artiklo{dativo}%
\vorto{dativo}\fako{gram.} Kazo di la deklino (Latina, Greka, Sanskrita, e c.) per qua expresesas ke kozo atribuesas a la komplemento\lingui{DEFIRSL}
\artiklo{dato}%
\vorto{dato} Indiko pri la yaro, la monato, la dio en qua fakto eventis od eventos\lingui{DEFIS}
\artiklo{daturo}%
\vorto{daturo}\fako{bot.} Genero de planti ek la familio « solanei », di qua un speco esas narkotigiva e venenoza\latina{datura}\lingui{EFIRS}
\artiklo{daxo}%
\vorto{daxo}\fako{zool.} Mamifero plantigrada, karnavida, ma anke fruktivora e granivora, kurta-gamba, di qua la ungli esas forta e hokatra, exkav-apta\latina{meles taxus}\lingui{DL}
\artiklo{de}%
\vorto{de}\fako{gram.} Prepoziciono qua indikas la punto di departo (en spaco e tempo) la origino, dependo, punto departala; uzesas anke : 1. kun la substantivi qui signifikas mezuro, quanto, kontenanto; 2. kun la adjektivi « plena, longa, larja, alta, profunda, dika, e c. » qui, fakte, relatas mezuro, dimensiono; 3. kun titulo di nobeleso
\artiklo{debar}%
\vorto{debar}\fako{trans.} Devar pagor per pekunio, kozo o servo\lingui{EFIRS}
\artiklo{debatar}%
\vorto{debatar}\fako{trans.} Diskutar kun un o plura konversanti, ye maniero ordinoza ed atencoza – singlu prizentante la argumenti qui servas sua propra interesto o prefero\lingui{DEFIRS}
\artiklo{debeto}%
\vorto{debeto}\fako{komerco} To quon ulu debas\lingui{DEFIRS}
\artiklo{debila}%
\vorto{debila} Febla naturale, temperamentale\lingui{EFIS}
\artiklo{debitar}%
\vorto{debitar}\fako{trans.}\senco{1}{Vendar endetale}\senco{2}{Lasar fluar ula quanto de liquido, dum ula quanto de tempo}\lingui{FIS}
\artiklo{debochar}%
\vorto{debochar}\fako{netrans.} Kondutar ne-etikale pri la relati sexuala\lingui{EFR}
\artiklo{debushar}%
\vorto{debushar}\fako{netrans.} Ekirar de loko streta, por extensar su en loko min konstriktita\lingui{DEF}
\artiklo{debutar}%
\vorto{debutar}\fako{netrans.} Agar sua probo dat-unesma\lingui{DEFIS}
\artiklo{decar}%
\vorto{decar}\fako{netrans.} Respektar, egardar, extere la konveno-reguli, la bona mori\lingui{DEFIS}
\artiklo{decembro}%
\vorto{decembro} Monato quan, nun, esas la dek-e-duesma (e la dato-lasta) di la yaro\lingui{DEFIRS}
\artiklo{decendar}%
\vorto{decendar}\fako{netrans.} Venar, segun lineo qua iras nemediate de gepatri a gefilii, de raso, de trunko prima\lingui{DEFIS}
\artiklo{decensar}%
\vorto{decensar}\fako{netrans.} Irar de supre ad infre. Portesar de supre ad infre\lingui{EFIS}
\artiklo{deceptar}%
\vorto{deceptar}\fako{trans.}\senco{1}{Trompar per la aspekto luranta}\senco{2}{Trompar per ne realigar la esperajo}\lingui{EFIS}
\artiklo{dechifrar}%
\vorto{dechifrar}\fako{trans.} Deskovrar la klefo e tradukar kriptografuro\lingui{DEFIRS}
\artiklo{decidar}%
\vorto{decidar}\fako{trans.} Pri temo qua genitis deskonkordi, enuncar sua judiko pos studiir la faktori di la problemo\lingui{EFIS}
\artiklo{decidua}%
\vorto{decidua}\fako{bot.} Qua su separas e falas naturale, kande lua vivo cesis\lingui{DEFIS}
\artiklo{decigramo}%
\vorto{decigramo} La parto dekima di gramo\lingui{DEFIRS}
\artiklo{deciliono}%
\vorto{deciliono} 10⁶⁰
\artiklo{decilitro}%
\vorto{decilitro} La parto dekima di litro\lingui{DEFIRS}
\artiklo{decimacar}%
\vorto{decimacar}\fako{trans.} Ek deko de personi, ocidar una, indikita lotrie\lingui{DEFIS}
\artiklo{decimala}%
\vorto{decimala} Qua procedas per deki\lingui{DEFIRS}
\artiklo{decimetro}%
\vorto{decimetro} La parto dekima di metro\lingui{DEFIRS}
\artiklo{decimo}%
\vorto{decimo} La parto dekima di franko\lingui{DEFIRS}
\artiklo{dedikar}%
\vorto{dedikar}\fako{trans.) (ad}\senco{1}{Pozar (templo, kirko) sub la protekto deala}\senco{2}{Pozar (libro, generale) sub la protekto di ulu per enskribar en olu la nomo di ta persono}\lingui{DEFIS}
\artiklo{dedolar}%
\vorto{dedolar}\fako{trans.) (kirurg.} Deprenar la parto surfacala di la tisui\lingui{F}
\artiklo{deduktar}%
\vorto{deduktar}\fako{trans.} Ekirigar (logike) de propoziciono (la serio de konsequaji quin lu kontenas implicite)\lingui{DEFIRS}
\artiklo{defektiva}%
\vorto{defektiva}\fako{gram.} Che qua mankas ula formi di la deklinotipo, di la konjugo-tipo a qua lu apartenas\lingui{DEFIS}
\artiklo{defekto}%
\vorto{defekto} To quo, en persono od en kozo, ne esas tala qualan lu devas esar\lingui{EFIS}
\artiklo{defensar}%
\vorto{defensar}\fako{trans., kontre, de}\senco{1}{Helpar (ulu) kontre ti qui atakas lu}\senco{2}{Shirmar (ulo) de to quo minacas lu}\lingui{DEFIS}
\artiklo{deferencar}%
\vorto{deferencar}\fako{trans.} Manifestar egardo ad ulu per exekutar lua volo o deziro\lingui{DEFIS}
\artiklo{deferenta}%
\vorto{deferenta}\senco{1}{\textit{(anatomio)} Kanalo deferenta : kanalo ek-sekrecala di la testikuli}\senco{2}{\textit{(astronomio)} Cirklo deferenta : cirklo imaginita da la olima astronomi por explikar la ecentrikeso, la perigeo e la apogeo di la planeti}\senco{3}{\textit{(bonatiko)} Kanalo deferenta : organo precipua por la cirkulado di sapo}
\artiklo{defervecar}%
\vorto{defervecar}\fako{netrans.) (medic.} Deskresko gradoza di febro, qua anuncas la konvaleco pri la morbi akuta\lingui{DEFIS}
\artiklo{defiar}%
\vorto{defiar}\fako{trans.} Instigar ulu ad agor o facor ulo, per enuncar dubito pri lua audaco o kapableso agar o facar lo\lingui{EFIS}
\artiklo{deficienta}%
\vorto{deficienta}\fako{matem.} \textit{(nombro)} tala ke la sumo di lua parti quantima ne equivalas lu\lingui{EF}
\artiklo{deficito}%
\vorto{deficito} To pro la manko di quo ne inter-equilibrigas la revenui e la spensi, la havajo e la debeto, la quanto de produkturi e la quanto bezonata da la konsumeri\lingui{DEFIRS}
\artiklo{defilar}%
\vorto{defilar}\fako{netrans.) (aludante trupo de homi} Pazar ganso-marche, seriope o kolonope, koram superioro\lingui{DEFIS}
\artiklo{defileo}%
\vorto{defileo} Paseyo streta inter du monti\lingui{DEFIS}
\artiklo{definar}%
\vorto{definar}\fako{trans.} Konocigar, per formulo preciza, to quon esas kozo, o la senco di termino\lingui{DEFIS}
\artiklo{definitiva}%
\vorto{definitiva} Fixigita tale ke la kozo ne plus esas ri-agenda o rifacenda\lingui{DEFIS}
\artiklo{deflacar}%
\vorto{deflacar}\fako{trans.) (financo} Reduktar la quanto de moneto-papero qua cirkulas\lingui{DEFIRS}\subvorto{}{antonimo}{Inflaco}{}{}
\artiklo{deflegmar}%
\vorto{deflegmar}\fako{trans.} Deprenar de la alkoholo la flegmi
\artiklo{deformar}%
\vorto{deformar}\fako{trans.} Alterigar pri lua formo\lingui{DEF}
\artiklo{degenerar}%
\vorto{degenerar}\fako{netrans.} Perdar la bona qualesi di sua raso\lingui{DEFIS}
\artiklo{degnar}%
\vorto{degnar}\fako{trans.} Voluntar (agar ulo) pro egardo di ulu quan onu judikas digna ye lo\lingui{EFIS}
\artiklo{degradar}%
\vorto{degradar}\fako{trans.}\senco{1}{Decensigar (ulu) de la grado dignesala, honorumala, qua esas lua}\senco{2}{Tre abasar, etikale}\lingui{DEFIS}
\artiklo{dehicenta}%
\vorto{dehicenta}\fako{bot.} Dicesas pri organo klozita qua su apertas ipse ye ula periodo di lua kresko\lingui{F}
\artiklo{dejunar}%
\vorto{dejunar}\fako{netrans.} Repastar, ye dimezo\lingui{F}
\artiklo{dejurar}%
\vorto{dejurar}\fako{netrans.} Esar en la periodo dum qua onu devas laborar, servar, agar, obedie di employanto od exekute di lo konvencionita kun organismo. (Anke pri oficiro en armeo)\lingui{DR}
\artiklo{dek}%
\vorto{dek} 10
\artiklo{dekadar}%
\vorto{dekadar}\fako{netrans.} Irar a ruineso\lingui{DEFIRS}
\artiklo{dekaedro}%
\vorto{dekaedro}\fako{geom.} Figuro qua havas dek edri
\artiklo{dekagono}%
\vorto{dekagono}\fako{geom.} Poligono qua havas dek anguli e dek lateri
\artiklo{dekagramo}%
\vorto{dekagramo} 10 grami
\artiklo{dekalitro}%
\vorto{dekalitro} 10 litri
\artiklo{dekalkomanio}%
\vorto{dekalkomanio} Procedo per qua onu transportas imaji kolorizita adsur surfaco (porcelano, papero, e c.) qua retenas lia stampuro
\artiklo{dekalogo}%
\vorto{dekalogo}\fako{religio} Asemblajo de la dek imperi da Deo a Mozes (sur la monto Sinaya)\lingui{DEFIS}
\artiklo{dekametro}%
\vorto{dekametro} 10 metri
\artiklo{dekano}%
\vorto{dekano}\senco{1}{Alta-rangiero di kirko (en kirko katedrala, la prezidero di la kanonikaro)}\senco{2}{Alta-rangiero di universitato (la profesoro selektita kom lua chefo dum ula quanto de tempo)}\lingui{DEFIRS}
\artiklo{dekantar}%
\vorto{dekantar}\fako{trans.} Transvarsar (liquido) tale ke en la vazo unesma restez la lizo\lingui{DEFIS}
\artiklo{dekapar}%
\vorto{dekapar}\fako{trans.} Deprenar la strato de oxido qua kovras la surfaco di metalo\lingui{DFIS}
\artiklo{dekastero}%
\vorto{dekastero} 10 steri
\artiklo{deklamar}%
\vorto{deklamar}\fako{trans.} Recitar, saliigante la senco per la toni e la gesti\lingui{DEFIRS}
\artiklo{deklarar}%
\vorto{deklarar}\fako{trans.}\senco{1}{Konocigar aperte (ulo)}\senco{2}{Enuncar aperte fakto}\lingui{DEFIRS}
\artiklo{deklinar}%
\vorto{deklinar}\fako{trans.}\senco{1}{Enuncar la serio de la dezinenci (di kazi, di genri, di nombri) tra qua pasas nomo, pronomo, adjektivo qualifikala en ula lingui}\senco{2}{\textit{(pri astro)} Distar de la equatoro di la sfero cielala, supere o sube}\senco{3}{\textit{(pri la magnet-agulo)} Deviacar de la meridiano terglobala, per ula angulo}\lingui{DEFIS}
\artiklo{dekoktar}%
\vorto{dekoktar}\fako{trans.} Boliigar substanco en liquido, por ke la utilajo quan lu kontenas dissolvesez en ta liquido\lingui{DEFIRS}
\artiklo{dekoltar}%
\vorto{dekoltar}\fako{trans.} Abatar addope o sesgar vesto por videbligar la kolo, e c.\lingui{DEFR}
\artiklo{dekorar}%
\vorto{dekorar}\fako{trans.} Garnisar ye acesoraji destinita a plubeligar\lingui{DEFIRS}
\artiklo{dekremento}%
\vorto{dekremento}\fako{radiotekniko} Deskresko di la ocili o di elektro-fluo
\artiklo{dekremetro}%
\vorto{dekremetro}\fako{radiotekniko} Aparato uzata por mezurar nemediate la dekremento logaritmala di cirkuito de elektro-fluo alternanta o di ondo-serio elektromagnetala
\artiklo{dekretar}%
\vorto{dekretar}\fako{trans.) (aludante autoritatozo suprega} Decidar, rezolvar\lingui{DEFIRS}
\artiklo{dekubitar}%
\vorto{dekubitar}\fako{netrans.) (aludante homo-korpo, kushita} Pozesar (flankale, dorsale)\lingui{DEFIS}
\artiklo{delegar}%
\vorto{delegar}\fako{trans., ad ulu} Charjar (ulu) per ofico, transmisante a lu omna, o parto de, sua autoritato\lingui{DEFIRS}
\artiklo{delektar}%
\vorto{delektar}\fako{trans.} Igar (ulu) savurar juo\lingui{DEFIS}
\artiklo{deletera}%
\vorto{deletera}\fako{vorto ciencala} Qua endanjerigas la vivo
\artiklo{delfinelo}%
\vorto{delfinelo}\fako{bot.} Nomo di ranunkulaceo, planto astriktiva\latina{delphinium}\lingui{F}
\artiklo{delfino}%
\vorto{delfino}\senco{1}{\textit{(historio di Francia)} Titulo di la filiulo seniora di la reji di Francia}\senco{2}{\textit{(zool.)} Cetaceo suflera, mikra, di qua la kranio esas konvexa, quan la fabli antiqua ed olima reprezentas kom amika a la homo}\latina{delphinus}\lingui{DEFIRS}
\artiklo{deliberar}%
\vorto{deliberar}\fako{netrans., pri} Ponderar en su ipsa, o kun altra personi, pri problemo solvenda o rezolvo agenda, la motivin « por » e la motivi « kontre » la decido o rezolvo\lingui{DEFIS}
\artiklo{delico}%
\vorto{delico} Juo delikatega, bonega, ecelanta\lingui{EFIS}
\artiklo{delikata}%
\vorto{delikata} Qua esas extreme dina, tenua, pura : di qua la gustosentiveso igas nefacila la satisfaco\lingui{DEFIRS}
\artiklo{deliktar}%
\vorto{deliktar}\fako{netrans.} Agar tale ke la lego esas violacita. – \textit{(en Francia)} Opoze a krimino ed a kontravenco : violaco di la lego qua efektigas puniseso segun la delikto-kodexo\lingui{DEFIRS}
\artiklo{deliquecar}%
\vorto{deliquecar}\fako{netrans.) (kemio} Qua aquiris la proprajo absorbar la vaporo di aero humida e, pro lo, dissolvesar\lingui{DEFIS}
\artiklo{delirar}%
\vorto{delirar}\fako{netrans.}\senco{1}{Standar tale ke, pro febro, ebrieso, e c., la idei esas acidente nekoheranta ed ecitegita}\senco{2}{Havar sua spirito en ta stando di exalteso plu o min violentoza qua karakterizas alienaco}\senco{3}{Havar sua spirito en tala stando di exalteso violentoza, ke lu cesas obediar raciono}\lingui{DEFIRS}
\artiklo{delto}%
\vorto{delto}\fako{geogr.} Extensajo de tero ferma, maxim-multa-kaze triangulo-forma, formacita da aluvioni qui, per amaseskar meze la boko di fluvio, partigis ica aden du o plura boki\lingui{DEFIRS}
\artiklo{deltoido}%
\vorto{deltoido}\fako{anat.} Muskulo triangulatra qua, per juntar la humero a la klavikulo ed a la skapulo, servas por elevar ed abasar la brakio
\artiklo{demagogio}%
\vorto{demagogio}\senco{1}{Sistemo di politiko segun qua onu flatas turbo}\senco{2}{Stato ube la povo statala abandonesas a turbo}\lingui{DEFIRS}
\artiklo{demagogo}%
\vorto{demagogo} La homo qua esforcas igar su povoza a turbo, per flatar lu\lingui{DEFIS}
\artiklo{demandar}%
\vorto{demandar}\fako{trans., ulo, ad ulu}\senco{1}{Savigar ke onu volas havar (ulo)}\senco{2}{Savigar da ulu ke onu expektas (ulo) de lu}\lingui{FIS}
\artiklo{demarshar}%
\vorto{demarshar}\fako{netrans.} Ago sociala per qua persono volas influar altra por obtenar ulo. Demarsho implikas sempre la ideo di « pazi » en la socio por atingar ula skopo, por obtenar irga helpo o favoro\lingui{F}
\artiklo{dementa}%
\vorto{dementa}\fako{patol.} Malada de grava desordineso di sua raciono\lingui{DEFIS}
\artiklo{dementiar}%
\vorto{dementiar}\fako{trans.}\senco{1}{Kontredicar (ulu) kom nedicinta lo exakta}\senco{2}{Kontredicar (to quon dicas ulu) kom neexakta}\lingui{FIS}
\artiklo{demisionar}%
\vorto{demisionar}\fako{netrans.} Retretar volunte de ofico\lingui{DFIS}
\artiklo{demiurgo}%
\vorto{demiurgo}\fako{en la filozofio olima} Deo, konsiderata ne kom kreinto, ma kom ordininto, di materio\lingui{DEFIRS}
\artiklo{demokrata}%
\vorto{demokrata} Esar partisano di demokratio, e (plu generale) havar la mento demokratiala o demokratismala\lingui{DEFIRS}
\artiklo{demokratio}%
\vorto{demokratio} La rejimo politikala en qua la populo guvernesas da su ipsa (t.e. per sua elektiti) : lua principo esas la suvereneso di la majoritato\lingui{DEFIRS}
\artiklo{demolisar}%
\vorto{demolisar}\fako{trans.} Deskonstruktar, per faligar, sucedante, omna parti qui kompozas lu\lingui{DEFIS}
\artiklo{demono}%
\vorto{demono}\senco{1}{\textit{(che la Antiqui)} Feo bonfacera o malfacera qua-segun la kredo lora – prezide direktis la fato di homo, di populo, ed inspiris a lu bonaji o malaji}\senco{2}{\textit{(nun)} Diablo}\senco{3}{\textit{(metaf.)} Personiguro di defekto, di vicio}\lingui{DEFIRS}
\artiklo{demonstrar}%
\vorto{demonstrar}\fako{trans.}\senco{1}{Pruvar, per rezono, la exakteso di propoziciono}\senco{2}{Docar per montrar la kozi, per pozar li avan la regardo}\lingui{DEFIRS}
\artiklo{demonstrativa}%
\vorto{demonstrativa}\fako{gram.} Qua indikas persono o kozo\lingui{DEFIRS}
\artiklo{demotika}%
\vorto{demotika}\fako{skribo-formo} Kursiva, vulgara, che la Egiptiani antiqua, simpliguro di la skriboformo hieratika (hieroglifi)\lingui{DEF}
\artiklo{denaro}%
\vorto{denaro} En Roma (antique), monet-uno ek arjento qua valoris cirkume dek « as-i »\lingui{DEFIRS}
\artiklo{denaturar}%
\vorto{denaturar}\fako{trans.} Igar neapta kom nutrivo, ma lasar sat apta pri uzeso en industrio\lingui{DEFIRS}
\artiklo{dendrito}%
\vorto{dendrito} Petro sur qua esas videbla desegnuri naturala qui reprezentas arboro-rami : ta ipsa rami
\artiklo{denigrar}%
\vorto{denigrar}\fako{trans.} Inspirar mala opiniono pri ulu (od ulo) per aserti ed opinioni severa o nejusta\lingui{F}
\artiklo{denominatoro}%
\vorto{denominatoro}\fako{matem.} Ta ek la du termini di fraciono qua indikas a quanta parti inter-egala dividesis la unajo\lingui{EFIS}
\artiklo{densa}%
\vorto{densa}\senco{1}{Di qua la maso esas granda kompare a la volumino}\senco{2}{Relato di la pezo di korpo (solida o liquida) a la pezo di volumino egala de aquo distilita ye 4° C. Relato di la pezo di gaso a la pezo di volumino egala di aero sama-faktore}\lingui{EFIS}
\artiklo{densimetro}%
\vorto{densimetro} Aparato por mezurar nemediate la denseso-grado di la korpi (precipue la liquidi)\lingui{DEFIRS}
\artiklo{dentelo}%
\vorto{dentelo} Texuro ajuroza, ek reto tenua de filo o de silko, di qua la bordo esas generale dentoza, uzata kom fundo por flori, ornivi facita per spindelo o per agulo\lingui{F}
\artiklo{dentifrico}%
\vorto{dentifrico} Preparajo qua uzesas por frotar, por netigar la denti\lingui{EFIS}
\artiklo{dento}%
\vorto{dento}\senco{1}{\textit{(anat.)} Singla ek la mikra osti emalioza qui garnisas la maxilo e la mandibulo di la homo, di ula animali, ed uzesas da li por tranchar, arachar, pistar la nutrivi, por mordar, e c.}\senco{2}{\textit{(tekn.)} To quo esas dentatra}\lingui{EFIRS}
\artiklo{denuncar}%
\vorto{denuncar}\fako{trans.} Indikar (ulu) kom kulpinta, precipue a la tribunali od a la policisti\lingui{DEFIRS}
\artiklo{deo}%
\vorto{deo}\fako{mitol., religio} Ento superega, objekto di la kulto da la personi qui kredas ke lu existas, e ke li e mondo submisesis a lua povo\lingui{EFIS}
\artiklo{departar}%
\vorto{departar}\fako{netrans., de, ad} Forirar de loko\lingui{EFIS}
\artiklo{departmento}%
\vorto{departmento} Dividuro teritoriala, o resortiso teritoriala\lingui{DEFIRS}
\artiklo{dependar}%
\vorto{dependar}\fako{netrans., de} Ligesar necese a persono od a kozo, kom nepovanta esor o realigesor sen lu\lingui{EFIS}
\artiklo{deperisar}%
\vorto{deperisar}\fako{netrans.} Tale febleskar ke morto probable eventos balde\lingui{FI}
\artiklo{depesho}%
\vorto{depesho} Informo (publika o privata) sendita (ad ulu) per moyeno rapida\lingui{DEFIRS}
\artiklo{deplorar}%
\vorto{deplorar}\fako{trans.} Manifestar, expresar, profunda chagreno pro (ulo)\lingui{EFIRS}
\artiklo{deponenta}%
\vorto{deponenta}\fako{gram. Latina} Qua havas la formo pasiva e la senco aktiva\lingui{DEFIS}
\artiklo{deportar}%
\vorto{deportar}\fako{trans.} Duktar (kondamnito) adexter la teritorio ad-an loko de ube lu devos ne ekirar\lingui{DEFIRS}
\artiklo{depozar}%
\vorto{depozar}\fako{trans.} Pozar (kozo) en loko ube lu esos sekura\lingui{DEFIRS}
\artiklo{depresar}%
\vorto{depresar}\fako{trans.} Abasar la nivelo per preso de alte ad infre\lingui{DEFIS}
\artiklo{deputar}%
\vorto{deputar}\fako{trans.} Elektar (ulu) por reprezentar la naciono en asemblitaro deliberanta\lingui{DEFIRS}
\artiklo{derivar}%
\vorto{derivar}\fako{trans.}\senco{1}{Ekirigar de lua lito aquo fluanta, por direktar lu ad altra direciono}\senco{2}{Atraktar (sango, humoro) de parto di la korpo ad altra, ube lia adfluo esas min danjeroza}\senco{3}{Ekirigar (vorto) de altra per modifikar lua formo; (senco) de altra senco per extenso, e c.}\lingui{EFIS}
\artiklo{dermatogeno}%
\vorto{dermatogeno}\fako{anat.} Tisuo yuna qua, ye la somito di la stipo o di la radiko, produktas la epidermo o la « kofio »
\artiklo{dermo}%
\vorto{dermo}\fako{anat.} Tisuo-fibroza ye qua konsistas la strato maxim dika e maxim profunde-situita di pelo, e quan kovras membrano surfacala (epidermo)\lingui{DEFIS}
\artiklo{derogar}%
\vorto{derogar}\fako{netrans. de) (yuro-cienco} Eceptar, violacar provizore, ula lego o regulo\lingui{DEFIS}
\artiklo{dervisho}%
\vorto{dervisho}\fako{che la Mohamedisti} Quaza monako\lingui{DEFIRS}
\artiklo{des}%
\vorto{« des »}\fako{muziko} La noto duesma di la « do »-gamo bemoloza
\artiklo{des-}%
\vorto{des-} La antonimo di (to quon indikas la radiko)
\artiklo{desegnar}%
\vorto{desegnar}\fako{trans., per, sur, e c}\senco{1}{Riproduktar, per krayono o plumo, la formo di la objekti}\senco{2}{Igar videbla la konturi}\lingui{EFIS}
\artiklo{desero}%
\vorto{desero} La kozi servata laste, lor repasto : fromajo, frukti, kuko, e c.\lingui{DEFR}
\artiklo{desertar}%
\vorto{desertar}\fako{trans., e netrans.} Abandonar loko (precipue aludante armeano) o partiso, e c. quin onu devas ne livar\lingui{DEFIRS}
\artiklo{deservar}%
\vorto{deservar}\fako{trans.) (en templo, kirko, sinagogo} Agar la servado religiala\lingui{DEFS}
\artiklo{desikatoro}%
\vorto{desikatoro}\senco{1}{\textit{(kemio)} Aparato sikigera}\senco{2}{\textit{(tekn.)} Aparato en qua onu pozas diversa texo-materii por determinar la pezo vera di la vari en paki}\lingui{DEFIRS}
\artiklo{deskriptar}%
\vorto{deskriptar}\fako{trans.} Reprezentar skribe o parole (persono, kozo) pri lua traiti extera\lingui{EFIS}
\artiklo{despenso}%
\vorto{despenso} Loko ube onu inkluzas la provizuri (en kolegio, farmodomo, e c.)\lingui{FIS}
\artiklo{despitar}%
\vorto{despitar}\fako{netrans., pri, pro, de} Iracar pro desprizeso, desegardeso, o pro ke altru preferesas a ni, e chagrenas pri lo\lingui{EFIS}
\artiklo{despoto}%
\vorto{despoto} Monarko absolutista qua guvernas segun sua arbitrio\lingui{DEFIRS}
\artiklo{destinar}%
\vorto{destinar}\fako{trans., ad, por}\senco{1}{Fixigar anticipe (ad ulu) kom esonta lua loto}\senco{2}{Fixigar anticipe kom uzota por ta o ca skopo}\lingui{DEFIS}
\artiklo{destruktar}%
\vorto{destruktar}\fako{trans.}\senco{1}{Desfacar (lo facita)}\senco{2}{Ruinar tote}\lingui{DEFIRS}
\artiklo{detachar}%
\vorto{detachar}\fako{trans.) (armeo} Sendar (ulu, trupo) separite, aparte\lingui{DEF}
\artiklo{detachmento}%
\vorto{detachmento}\fako{armeo} Mikra korpo de armeani quan onu detachas de korpo plu granda\lingui{DEF}
\artiklo{detalo}%
\vorto{detalo} La partikularajo di objekto\lingui{DEFIRS}
\artiklo{detektivo}%
\vorto{detektivo} Policisto komisita por explorar\lingui{DEFR}
\artiklo{detektoro}%
\vorto{detektoro}\fako{elektro} Aparato por transformar elektro-magnetala ondi a formo di energio perceptebla da la orelo o da la okulo\lingui{DEFIS}
\artiklo{determinar}%
\vorto{determinar}\fako{trans.}\senco{1}{Fixigar (ulo) quo esas ankore dubitebla}\senco{2}{\textit{(gram.)} Fixigar la rolo, en la frazo, di propoziciono chefa \textit{(aludante propoziciono acesora)}}\lingui{DEFIRS}
\artiklo{detersar}%
\vorto{detersar}\fako{trans., per) (trans.} Netigar per eventigar la evakueso di la humoro, di la puso, e c.\lingui{DEFI}
\artiklo{detonar}%
\vorto{detonar}\fako{netrans.} Eventigar bruiso subita per bruska expanso di gaso\lingui{DEFIS}
\artiklo{detrimentar}%
\vorto{detrimentar}\fako{trans.} Efektigar ad ulu (ad ulo), perdo, mem ne-vole e mediate\lingui{EFIS}
\artiklo{deutoplasmo}%
\vorto{deutoplasmo}\fako{biol.} Rezervo-materii inkluzita en la citoplasmo e ne vivoza
\artiklo{devalorar}%
\vorto{devalorar}\fako{trans.) (financo} Diminutar la valoro di moneto papera relate oro
\artiklo{devaluar}%
\vorto{devaluar}\fako{trans.) (financo} Substitucar a monet-unajo, altra monet-unajo qua kontenas min multa oro. (Ta operaco konsakras legale la devaloreso di moneto, por sekurigar la stabileso di ica)
\artiklo{devancar}%
\vorto{devancar}\fako{trans.} Pasar avan sua rivali, pri rango, en hierarkio\lingui{F}
\artiklo{devar}%
\vorto{devar}\fako{trans.} Obligesar agor ulo. (Ta vorto konvenas ad obligesi omna ed omnaspeca; uzenda en omna imperi, preskripti e konsili. Obligeso dependas de volado)\lingui{FIS}
\artiklo{devastar}%
\vorto{devastar}\fako{trans.} Ruinar (regiono per destruktar arbori, cereali, lojeyi, e c., tale ke la sulo igesas nuda)\lingui{EFIS}
\artiklo{developar}%
\vorto{developar}\fako{trans.}\senco{1}{Igar parkreskar to quo kontenesas en jermo}\senco{2}{Igar par-kreskar ulo – to quo esas virtuala, potenciala}\lingui{EFI}
\artiklo{deviacar}%
\vorto{deviacar}\fako{netrans., de} Eskartesar de la voyo rekta; eskartesar de la direciono normala\lingui{EFIS}
\artiklo{devizo}%
\vorto{devizo} Formulo quan onu skribis super, sub od apud la skudo en la blazono, olim\lingui{DEFIRS}
\artiklo{devoco}%
\vorto{devoco} Zelo por religio, por la praktikado religiala\lingui{DEFIS}
\artiklo{devorar}%
\vorto{devorar}\fako{trans.}\senco{1}{\textit{(aludante animalo)} Saturar su per sua raptajo}\senco{2}{Manjar avide}\lingui{EFIS}
\artiklo{devota}%
\vorto{devota} Konsakrita a servar od a sorgar la interesti di ulu\lingui{DEFI}
\artiklo{dextra}%
\vorto{dextra} Situita ye la latero opozita ad olta ube trovesas la pinto di la homo-kordio\lingui{EFIS}
\artiklo{dextrino}%
\vorto{dextrino}\fako{kemio} Substanco analoga a gumo, di qua la kompozanti kemiala esas sama kam olti di amilo, e quan onu uzas por apretar, bluizar\simbolo{(C₆H₁₀O₅) n}\lingui{DEFIRS}
\artiklo{dezerta}%
\vorto{dezerta} Qua ne esas sovaja (od arida), ma ne-habitata\lingui{EFIS}
\artiklo{dezinenco}%
\vorto{dezinenco} Fino-literi di vorto, qui expresas la flexioni\lingui{DFIS}
\artiklo{dezirar}%
\vorto{dezirar}\fako{trans., por} Tendencar a (kozo quan onu volas posedar, od ago quan onu volas efektigar – segun quante ta voli esas satisfacebla)\lingui{EFIS}
\artiklo{dezolar}%
\vorto{dezolar}\fako{trans.} Frapar (ulu) per aflikto troa, pro qua judikeble omno mankas a lu\lingui{EFIS}
\artiklo{di}%
\vorto{di}\fako{gram.} Prepoziciono qua expresas la relato di persono o kozo : 1. kun termino a qua ta persono o kozo apartenas, o de qua lu dependas; 2. kun termino submisita a lua efiko
\artiklo{diabaso}%
\vorto{diabaso}\fako{geol.} Roko eruptala, qua havas la sama kompozanto mineralogiala kam diorito, e strukturo ek intrikajo de prismi aplastita, feldspata, aglomerita en plaji de piroxeno\lingui{DEF}
\artiklo{diabeto}%
\vorto{diabeto}\fako{patol.} Morbo karakterizata da ek-sekreco abundanta di urino qua kontenas plu o min multa elementi sukra – e da deperiso gradoza\lingui{DEFIRS}
\artiklo{diablo}%
\vorto{diablo}\senco{1}{\textit{(en la religio Kristana)} La demono di lo mala, Satanas, princo di la anjeli rebela}\senco{2}{Persono o kozo quan onu komparas a diablo}\lingui{DEFIRS}
\artiklo{diademo}%
\vorto{diademo}\senco{1}{Bendo ornizita ye broduri, lapidi, quan la suvereni *weris zone sur sua fronto}\senco{2}{Bendo, kapo-ringa, ye qua konsistas la krono di la suvereni}\lingui{DEFIRS}
\artiklo{diado}%
\vorto{diado}\senco{1}{Asemblajo de du principi, de du elementi}\senco{2}{\textit{(filoz.)} Duo ek la principo unesma (monado) e la principo inferiora, qua devenas de ica}
\artiklo{diafana}%
\vorto{diafana} Qua lasas pasar la lumo-radii, tale ke onu povas dicernar la objekti tra lu\lingui{EFIS}
\artiklo{diafizo}%
\vorto{diafizo}\fako{anat.} Korpo di la osti longa\lingui{DEF}
\artiklo{diafragmo}%
\vorto{diafragmo}\senco{1}{\textit{(anat.)} Muskulo parieta inter la pektoro e la abdomino. – Parieto qua dividas ula organi}\senco{2}{\textit{(optiko)} Cirklo opaka, muntita en teleskopo o lorno, por haltigar la radii qui ne kuniras a la foko. – Disko kun truo ronda, quan la fotografisto muntas en la objektalo samaskope}\lingui{EFIRS}
\artiklo{diagnozar}%
\vorto{diagnozar}\fako{trans.} Determinar (morbo) per la simptomi qui esas lua propraji\lingui{DEFIRS}
\artiklo{diagonala}%
\vorto{diagonala} Qua iras de angulo ad angulo opozita\lingui{DEFIRS}
\artiklo{diagramo}%
\vorto{diagramo} Expresuro grafika, tabelo lineala, cifrala, ek la parti di ensemblo ed ek lia inter-relati\lingui{DEFIRS}
\artiklo{diakaustika}%
\vorto{diakaustika}\fako{geom.} Produktita da la inter-seki di la radii refraktita
\artiklo{diakilono}%
\vorto{diakilono}\senco{1}{Plastro ek dekokturo de gladiolo-radiki, de oleo e de litargiro, a qui onu adjuntas ulakaze gumo-rezini, pecho blanka, terpentin-oleo, vaxo}\senco{2}{Sparadrapo ek ta plastro, extensita sur telo peco}\lingui{DEF}
\artiklo{diakoniso}%
\vorto{diakoniso}\fako{religio}\senco{1}{En la komenco-periodo di la eklezio Kristana : filiino o vidvino qua recevis la impozo di la manui, e su konsakris a la entrateno di la templi, di la povri, e c.}\senco{2}{En ula sekti di protestantismo : damo qua su konsakris a verki piesala, karitatala}\lingui{DEFIRS}
\artiklo{diakono}%
\vorto{diakono}\fako{en la eklezio katolika} Kleriko qua havas la duesma ek la ordi superiora\lingui{DEFIRS}
\artiklo{diakritika}%
\vorto{diakritika} Qua esas signo adjuntita a litero (supersigno, sub-signo, e c.) – kom exempli : ç, â, ĝ, ë, e c.\lingui{DEFIRS}
\artiklo{dialektiko}%
\vorto{dialektiko} Metodo per qua onu deduktas rezoni, uzata por demonstrar o refutar\lingui{DEFIRS}
\artiklo{dialekto}%
\vorto{dialekto}\fako{linguistiko} Varietato regionala di linguo\lingui{DEFIRS}
\artiklo{dializar}%
\vorto{dializar}\fako{trans.} Separar du substanci dissolvura, per difuzo tra parieto poroza, quan un de li povas trairar\lingui{DEF}
\artiklo{dialogar}%
\vorto{dialogar}\fako{netrans., kun} Konversar kun un persono\lingui{DEFIRS}
\artiklo{diamanto}%
\vorto{diamanto}\fako{mineral.} Lapido ek karbo kristaligita, un sorto di qua uzesas kom juvelo pos facetizo (se lu esas pura)\lingui{DEFIRS}
\artiklo{diametro}%
\vorto{diametro}\fako{geom.} Lineo rekta qua pasas tra la centro di cirklo, di sfero, di sferoido, e di qua la extremaji tushas la periferio\lingui{DEFIRS}
\artiklo{dianto}%
\vorto{dianto}\fako{bot.} Planto gardenala, de la familio « kariofilei »\latina{dianthus}\lingui{L}
\artiklo{diapazono}%
\vorto{diapazono}\fako{muziko} Instrumento qua donas noto uzata kom departo-punto por regular la tono\lingui{DEFIRS}
\artiklo{diapedezo}%
\vorto{diapedezo}\fako{patol.} Hemoragio tra la pelo-pori\lingui{DEFIRS}
\artiklo{diarear}%
\vorto{diarear}\fako{netrans.) (patol.} Ofte evakuar feko liquida ek la intestino\lingui{DEFIS}
\artiklo{diaso}%
\vorto{diaso}\fako{geol.} Nomo donita olim, da Macon (1839) a la strato geologiala permiana\lingui{DE}
\artiklo{diastazo}%
\vorto{diastazo}\fako{kemio} Substanco azotoza, neutra, extraktita de hordeo, de frumento, e c., qua efikas quale fermento a la amilo di la grani, tale ke lu separas la sukro e la dextrino de la substanci nesolvebla kun qui li esas mixita\lingui{DEFIRS}
\artiklo{diastolo}%
\vorto{diastolo}\fako{anat.} Movo di dilato di la kordio e di la arterii, qua alternas kun la movo di konstrikto (sistolo)\lingui{DEFIS}
\artiklo{diatezo}%
\vorto{diatezo}\fako{fiziol.} Ta stando generala di la organismo di homo, maxim-multa-kaze atavismala od inata, qua, pos esir latenta, su manifestas per ula morbi\lingui{DEFIS}
\artiklo{diatomeo}%
\vorto{diatomeo}\fako{bot.} Familio de algi feoficea, od algi bruna\lingui{DEFIRS}
\artiklo{diatonika}%
\vorto{diatonika}\fako{muziko} Qua procedas segun la sucedo naturala di la toni e la mi-toni di la gamo\lingui{DEFIRS}
\artiklo{diatribo}%
\vorto{diatribo}\senco{1}{Diserturo kritikoza}\senco{2}{Kritiko akra}\lingui{DEFIRS}
\artiklo{dicar}%
\vorto{dicar}\fako{trans., ad, pri, pro} Savigar per vorti e signi (to quon onu havas komunikenda)\lingui{FIS}
\artiklo{dicernar}%
\vorto{dicernar}\fako{trans.} Vidar o komprenar la difero inter du kozi (t. e. ke li interdiferas)\lingui{EFIS}
\artiklo{dicionario}%
\vorto{dicionario} Vorto-libro en qua esas definita la senco di la vorti klasifikita alfabetale (quale en ica radikaro)\lingui{DEFIRS}
\artiklo{diciono}%
\vorto{diciono} Arto, maniero parolar o deklamar\lingui{DEFIRS}
\artiklo{diciplinar}%
\vorto{diciplinar}\fako{trans.} Kustumigar a la regulo di subordino e di ordino, impozita a la membri di korpo, di organismo\lingui{DEFIRS}
\artiklo{dicipulo}%
\vorto{dicipulo} Persono qua asimilas la lecioni di maestro o la doktrino di ulu\lingui{EFIS}
\artiklo{didaktiko}%
\vorto{didaktiko} La doco-cienco\lingui{DEFIRS}
\artiklo{didelfo}%
\vorto{didelfo}\fako{zool.} Grup-uno mamifera (marsupiali), di qua la femino havas abdomino-posho en qua elu shirmas sua yuni dum la kelka monati qui sequas lia nasko\latina{didelphus}\lingui{L}
\artiklo{diedra}%
\vorto{diedra}\fako{geom.} Qua konsistas ye du plani\lingui{DEF}
\artiklo{dielektriko}%
\vorto{dielektriko}\fako{elektro} Omna medio qua posibligas nur tre febla o desegardebla kondukto di elektro
\artiklo{dietetiko}%
\vorto{dietetiko} Parto di medicino, qua traktas la nutro-rejimi\lingui{FIRS}
\artiklo{dieto}%
\vorto{dieto}\senco{1}{Nutro-rejimo qua konsistas ye la absteno tota o parta de nutrivi}\senco{2}{En ula stati di Europa (ante 1914) : Germania, Suedia, Polonia, Suisia, Hungaria – asemblitaro politikala suprega, ek elementi diversa segun la charto politikala di la stato}\lingui{DEFIRS}
\artiklo{diezo}%
\vorto{diezo}\fako{muziko}\senco{1}{Intervalo de un mi-tono per qua onu elevas noto}\senco{2}{Signo (\#) quan onu pozas acidente avan noto od an la klefo por savigar ke onu devas elevar ye un mi-tono omna noti skribita sur la sama lineo, en la muzikajo}
\artiklo{difamar}%
\vorto{difamar}\fako{trans.} Dicar ulo, exakta o ne-exakta, qua atakas la honoro di ulu o lua bona reputeso; ago punisebla da la lego civila\lingui{EFIRS}
\artiklo{diferar}%
\vorto{diferar}\fako{netrans.} Esar des-simila (kam)\lingui{DEFIS}
\artiklo{diferencialo}%
\vorto{diferencialo}\fako{tekn.} Mekanismo muntita sur la arboro di la movo-roti di automobilo, por posibligar ke, lor jiro, la roto exter-jira aquirez la rotaco-rapideso di la roto interna-jira\lingui{DEFIRS}
\artiklo{diferenciar}%
\vorto{diferenciar}\fako{trans.} Eventigar la ensemblo de fenomeni qui, de la ovo, produktas la tisui diversa ye qui konsistas organismo komplexa\lingui{DEFIS}
\artiklo{difraktar}%
\vorto{difraktar}\fako{trans.} Produktar deviaco semblanta di la lumo-radii kande li frolas la bordi di korpo opaka\lingui{DEFIS}
\artiklo{difterio}%
\vorto{difterio}\fako{patol.} Morbo quan karakterizas la formaco di pseudo-membrani en la guturo\lingui{DEFIRS}
\artiklo{diftongo}%
\vorto{diftongo}\fako{linguistiko} Asemblajo de du vokali qui pronuncesas per un voc-emiso (kom ex. : au, eu, ie…)\lingui{DEFIRS}
\artiklo{difuzar}%
\vorto{difuzar}\fako{trans.} Extensar omna-direcione, omna-sinse\lingui{DEFIRS}
\artiklo{digestar}%
\vorto{digestar}\fako{trans.) (fiziol.} Elaborar (nutrivo) por igar asimilebla la parto nutriva, e separar ica de olta qua esas eliminenda. – \textit{(anke metaf.)}\lingui{DEFIS}
\artiklo{digitalino}%
\vorto{digitalino} Substanco efikera di la digitalo purpura\lingui{DEF}
\artiklo{digitalo}%
\vorto{digitalo}\fako{bot.} Planto di qua la floro esas simila a ganto-fingro, e la floro uzata por plulentigar la movi di la homo-kordio\latina{digitalis}\lingui{DEFISL}
\artiklo{digitigrado}%
\vorto{digitigrado}\fako{zool.} Qua, pro havar la tarso e la metatarso plu alta kam normale, apogas su sur sua fingri exkluzive, lor marcho\lingui{EFIS}
\artiklo{digna}%
\vorto{digna} Qua meritas (favoroze o desfavoroze) ulo\lingui{FIS}
\artiklo{digo}%
\vorto{digo}\senco{1}{Konstrukturo por preventar inundi; longa konstrukturo destinita a retenar la aquo, sur litoro o sur rivo, alonge lageto, an la enireyo di portuo}\senco{2}{\textit{(metaf.)} Obstaklo opozita a to quo tendencas ekirar de sua limiti}\lingui{FIS}
\artiklo{digramo}%
\vorto{digramo}\fako{linguistiko} Grupo de du literi; sonuno simpla, reprezentata skribe per grupo de du literi. (Kom ex. : sh, ch, e c.)\lingui{DEFIS}
\artiklo{digresar}%
\vorto{digresar}\fako{netrans.) (pri diserto} Traktar, eskartante su, deviacante, de la temo. – \textit{(astron.)} Elongaco maxima di planeto inferiora\lingui{EFIS}
\artiklo{dika}%
\vorto{dika} Konsiderata segun la dimensiono opozata a longeso ed a larjeso\lingui{DE}
\artiklo{diklina}%
\vorto{diklina}\fako{bot.} Di qua la floruli separesas de la florini\lingui{DEFIS}
\artiklo{dikogamo}%
\vorto{dikogamo}\fako{bot.} Dicesas pri la flori hermafrodita, en qui ne eventas sam-instante la developeso di la stamini e di la pistilo
\artiklo{dikotiledona}%
\vorto{dikotiledona}\fako{bot.} Di qua la embriono havas du lobi (kotiledoni)\lingui{DEFIS}
\artiklo{dikotoma}%
\vorto{dikotoma}\fako{bot.} Di qua la stipi di la rameti e di la pedunkli esas dividita – subdividita duope\lingui{DEFIRS}
\artiklo{dikotomio}%
\vorto{dikotomio}\fako{logiko} Metodo di divido, per subdivido du-unajo\lingui{DEFIRS}
\artiklo{diktamo}%
\vorto{diktamo}\fako{bot.} Planto aromata de la familio « labiei », konsiderita olim kom remediilo kontre la vunduri\latina{dictamnus}
\artiklo{diktar}%
\vorto{diktar}\fako{trans.} Dicar (ulo) koram ulu, por ke lu skribez olu segun ke la vorti pronuncesas\lingui{DEFIRS}
\artiklo{diktatoro}%
\vorto{diktatoro}\senco{1}{\textit{(en Roma, antique)} Stat-oficiero extraordinara e provizora, di qua la autoritato esis sen restrikto}\senco{2}{\textit{(nun)} Ta homo qua recevas od arogas a su la yuro koncentrar en su omna autoritato, omna yuri}\lingui{DEFIRS}
\artiklo{dilatar}%
\vorto{dilatar}\fako{trans.}\senco{1}{Extensar korpo elastika}\senco{2}{\textit{(fiziko)} Augmentar la volumino di korpo (per la efiko di kaloro) sen augmentar lua maso}\lingui{EFIS}
\artiklo{dilemo}%
\vorto{dilemo}\fako{logiko} Rezono en qua onu reduktas omna kazi a du alternativo-branchi interkontrea, inter qui absolute oportas selektar, ica esante exakta se la altra esas ne-exakta, e qui du duktas a la konkluzajo quan onu volas demonstrar\lingui{DEFIRS}
\artiklo{diletar}%
\vorto{diletar}\fako{trans.} Prizar e kultivar, pro la plezuro de lo, arto, cienco, e c.\lingui{DEFIRS}
\artiklo{diligenta}%
\vorto{diligenta}\senco{1}{Qua manifestas sorgo zeloza}\senco{2}{Qua manifestas agereso konstanta en la exekuto di ulo}\lingui{EFIS}
\artiklo{dilijenco}%
\vorto{dilijenco} Granda veturo publika, kun plura faki, por transportar la voyajanti\lingui{DEFIRS}
\artiklo{dilutar}%
\vorto{dilutar}\fako{trans.}\senco{1}{Dissolvar substanco en liquido; (homeopatio) tenuigar la dozo de medikamento, per adjuntar a lu liquido}\senco{2}{Pozar en aquo ula substanci solvebla, por transformar li a pulvero nepalpebla, per dekanto}\lingui{FIS}
\artiklo{diluvio}%
\vorto{diluvio} Granda inundo qua submersis la terglobo e perisigis preske omna homi (segun Biblo)\lingui{DEFIS}
\artiklo{diminutar}%
\vorto{diminutar}\fako{trans., deprene di} Igar min granda, min multa, per la depreno di parto\lingui{DEFIS}
\artiklo{dimitio}%
\vorto{dimitio} Stofo mixura, ek lino-warpo e kotono-wefto\lingui{DE}
\artiklo{dimorfa}%
\vorto{dimorfa}\fako{kristalografio} Qua havas kom proprajo kristaleskar ye du formi qui apartenas a du sistemi geometriala interdiferanta
\artiklo{dina}%
\vorto{dina} Tre poke dika\lingui{DE}
\artiklo{dinamiko}%
\vorto{dinamiko} La cienco pri la kalkulo di la movi di la korpi qui submisesas a la efiko da la forci qui movas li\lingui{DEFIRS}
\artiklo{dinamismo}%
\vorto{dinamismo}\fako{filoz.} Sistemo en qua materio konsideresas kom facita ek forci qui esas ipse efikiva\lingui{DEFIRS}
\artiklo{dinamisto}%
\vorto{dinamisto} Partisano di dinamismo
\artiklo{dinamito}%
\vorto{dinamito} Explozivo ek mixuro de nitroglicerino e substanco inerta (sablo quarcoza, greso polvo, e c.) – inventita da Alfred Nobel\lingui{DEFIRS}
\artiklo{dinamo}%
\vorto{dinamo}\fako{tekn.} Generatoro di elektro-fluo kontinua, kun magneto-feldo auto-ecitata\lingui{DEFIRS}
\artiklo{dinamometrio}%
\vorto{dinamometrio} Evaluo e komparo di la fortesi, per dinamometro
\artiklo{dinamometro}%
\vorto{dinamometro} Instrumento por mezurar la povo di motoro, la forteso muskulala di la homi, di la animali\lingui{DEFIRS}
\artiklo{dinastio}%
\vorto{dinastio} Serio de suvereni, decendinta de un ancestro\lingui{DEFIRS}
\artiklo{dindo}%
\vorto{dindo}\fako{zool.} Galinaceo di qua la kaudo, kurvatra abanike, esas extensebla rotatre quale olta di pavono\latina{meleagris gallopavo}\lingui{FI}
\artiklo{dinear}%
\vorto{dinear}\fako{netrans.} Repastar, en la vespero (cirkum 19 o 20 kloki)\lingui{DEF}
\artiklo{dinosaurio}%
\vorto{dinosaurio}\fako{paleont.} Klaso de repteri, ek nur formi desaparinta, qui vivis sur la tersurfaco dum la epoki prehistoriala sekundara, ed inkluzanta : animali ek grandesi tre diversa, de kelka centimetri til 25 metri de longeso\lingui{DEF}
\artiklo{dinoterio}%
\vorto{dinoterio}\fako{zool.} Genero de mamiferi rostroza, ye qua konsistas la familio « dinoteridi »\lingui{DEF}
\artiklo{dio}%
\vorto{dio} Tempo-quanto qua fluas (inter du pasi da Suno ye la sama meridiano, de dimezo a dimezo), e qua varias segun la epoko di la yaro (t.e. segun la situeso di Tero sur ekliptiko)\lingui{EFIS}
\artiklo{diocezo}%
\vorto{diocezo}\fako{eklezio katolika} Distrikto di qua la judicieso religiala resortisas ad episkopo od arkiepiskopo\lingui{DEF}
\artiklo{dioika}%
\vorto{dioika}\fako{bot.} Di qua la floruli e la florini esas sur stipi diversa
\artiklo{dioptriko}%
\vorto{dioptriko}\fako{optiko} Teorio di la vitro-diski optikala\lingui{DEF}
\artiklo{dioptrio}%
\vorto{dioptrio}\fako{optiko} Mezur-unajo por la vitro-diski di la binokli\lingui{DEFIRS}
\artiklo{dioptro}%
\vorto{dioptro} La fenestreto tra qua onu regardas por vizar punto per la instrumento di qua la nomo esas « alidado »\lingui{DEIRS}
\artiklo{dioramo}%
\vorto{dioramo} Pikturo granda-dimensiona, facita ek farbi diafana sur telo sen bordi videbla, di qua onu variigas la lumizo por prizentar diversa iluzioni optikala a la personi qui spektas lu de loko obskura\lingui{DEFIRS}
\artiklo{diorito}%
\vorto{diorito}\fako{geol.} Roko fairatra, ek feldspato e verda amfibolo\lingui{DEFIRS}
\artiklo{diplasar}%
\vorto{diplasar}\fako{trans.} Deprenar (ulo, ulu) de lua plaso\lingui{EF}
\artiklo{diplomaco}%
\vorto{diplomaco} Fako di politiko, qua koncernas la relati di populo, di guvernistaro, kun altra, la negocii, la kontrati internaciona\lingui{DEFIRS}
\artiklo{diplomatiko}%
\vorto{diplomatiko} Cienco di qua la studiajo esas la diplomi, charti, ed altra tituli, la verifiko di la autentikeso, di la integreso di la akti, e c.\lingui{DEFIRS}
\artiklo{diplomo}%
\vorto{diplomo}\senco{1}{\textit{(olim)} Akto oficala (faldita e siglita) qua emanis de suvereno ed establisis o konfirmis yuro, privilejo}\senco{2}{\textit{(nun)} Akto emananta de universitato, fakultato, skolo publika, qua grantas grado, titulo}\lingui{DEFIRS}
\artiklo{diploo}%
\vorto{diploo}\fako{anat.}\senco{1}{Asemblajo de du lami osta, kompakta, ye qui konsistas la surfaci interna ed extera di la kranio}\senco{2}{Tisuo sponjatra quan onu renkontras inter la surfaci di ta osti, e, generale, en la osti plata}
\artiklo{diplostemono}%
\vorto{diplostemono}\fako{bot.} Floro o planto di qua la quanto di la stamini esas duopla ye olta di la sepali e di la petali
\artiklo{diporto}%
\vorto{diporto}\fako{financo} Ajorno, til liquidaco borsala proxima, di la livro di valor-akti (acioni, obligacioni), vendita, pagante premio; la premio pagata por ta ajorno\lingui{DFIS}
\artiklo{diptero}%
\vorto{diptero}\fako{zool.} Insekto sugera, di qua nur la du ali supra esas developita, le infra esante remplasigita da organi rudimenta\lingui{DEFIS}
\artiklo{diptiko}%
\vorto{diptiko}\senco{1}{\textit{(en la epoki antiqua)} Asemblajo de du tabuleti ivora sur qua onu skribis (en Roma) la nomo di la konsuli e di la stat-oficieri alta-ranga; en la Eklezio katolika, la nomi di la episkopi, di la bonfaceri di la Eklezio, di la martiri, e c.}\senco{2}{\textit{(nun)} Pikturo basreliefa, kovrita da quaza shutro di qua la surfaco interna esas anke pikturoza e skulturoza}\lingui{DEFIS}
\artiklo{direciono}%
\vorto{direciono} Lineo rekta segun qua korpo su movas o pozesas\lingui{DEFIRS}
\artiklo{direktar}%
\vorto{direktar}\fako{trans.}\senco{1}{Funcionigar segun ta o ca konduto-lineo; (aludante kleriko) trasar ad ulu la konduto-lineo quan lu devas sequar pri lua koncienco}\senco{2}{Movigar ad skopo indikita}\lingui{DEFIRS}
\artiklo{Direktorio}%
\vorto{Direktorio}\fako{historio di Francia} Komisitaro de kin membri, qua investesis per la povo exekutala, de 1795 til 1799\lingui{DEFIRS}
\artiklo{direta}%
\vorto{direta}\fako{pri treno}\senco{1}{Qua, segun lua itinerario, atingas sua destino-staciono sen haltir ye la multa stacioni di la interspaco}\senco{2}{\textit{(pri voyo)} Rekta; sen-sinua o poke-sinua}\senco{3}{\textit{(pri cito)} En qua onu raportas, uzante la pronomo « me », la vorti pronuncita (reale o konjektate) da ulu}\senco{4}{\textit{(pri metodo)} Qua iras rekte a la skopo}\lingui{DEFIS}
\artiklo{dis}%
\vorto{« dis »}\fako{muziko} La noto duesma di la « do »-gamo, diezoza
\artiklo{dis-}%
\vorto{dis-} Prefixo qua indikas ideo di semeso hike ed ibe, hazarde
\artiklo{disenterio}%
\vorto{disenterio}\fako{patol.} Inflamuro di qua la fonto esas en la grosa intestino, qua karakterizesas da ofta kaki di materio mukatra, sangoza, akompanata da koliki forta e da doloriganta tenseso di la regiono anusala qua bezonigas kako-probo (olqua esas sen-sucesa)\lingui{DEFIRS}
\artiklo{disertar}%
\vorto{disertar}\fako{trans., pri} Prizentar ideo-developo pri temo, pri doktrino-punto\lingui{DEFIRS}
\artiklo{disho}%
\vorto{disho} Nutrivo destinita a prizentesar lor repasto\lingui{E}
\artiklo{disidenta}%
\vorto{disidenta} Qua ne konkordas kun la religiani maxim multa, pri ula punto, dogmato\lingui{DEFIRS}
\artiklo{disimular}%
\vorto{disimular}\fako{trans.} Celar sua pensi, opinioni, sentimenti, o fingar altrai\lingui{EFIS}
\artiklo{disipar}%
\vorto{disipar}\fako{trans.} Ofte spensar abunde ed ecese de sua havajo, po kozi frivola e sterila. \textit{(metaf.)} Onu disipas sua energio, sua atenco ye multa objekti e skopi\lingui{EFIS}
\artiklo{disjuntiva}%
\vorto{disjuntiva}\senco{1}{\textit{(logiko)} Qua indikas la separeso di la du termini di propoziciono}\senco{2}{\textit{(gram.)} Qua indikas ke la vorto, pri lua senco, de la vorto qua sequas lu}
\artiklo{disko}%
\vorto{disko}\senco{1}{\textit{(en la epoki antiqua)} Cirklajo pezoza, fera o petra, quan onu probis lansar maxim adfore lor la exerci, la ludi atletala}\senco{2}{Nomo di omna korpo plata e cirkla (kom ex. : disko fonografala, e c.)}\lingui{DEFIRS}
\artiklo{diskontar}%
\vorto{diskontar}\fako{trans.) (banko} Anticipe pagar ad ulu la pekunio-quanto qua debesas a lu, prece di retenajo kalkulita segun la interesti di la pekunio-quanto dum la tempo qua fluos til la pago-dato normala, konvencionita (ye procento qua varias segun la kurso di la merkato)\lingui{DEFIRS}
\artiklo{diskreta}%
\vorto{diskreta} Qua evitas jenar o tedar la kunhomi per sua demandi, o nocar li per agi neprudenta (ex. : revelo di sekretaji)\lingui{DEFIRS}
\artiklo{diskriminanto}%
\vorto{diskriminanto}\fako{algebro}\senco{1}{Funciono di la koeficiento di equaciono di la grado duesma, qua posibligas dicernar ka la kurvo reprezentata da la equaciono en koordinati ortangula esas kurvo vera, o sistemo de du rekti, – en la kazo di elipso – kad ica esas reala od imaginala}\senco{2}{Dicesas anke pri la equaciono B² – E ac – a, b, c, esante la koeficienti di ax² + bx + c = 0}\lingui{DEF}
\artiklo{diskursar}%
\vorto{diskursar}\fako{netrans., pri}\senco{1}{Developar temo oratorala koram un o plura personi, animante la pronunco per gesti}\senco{2}{Developar docive temo (parole o skribe)}\lingui{DEFIRS}
\artiklo{diskursiva}%
\vorto{diskursiva}\fako{logiko} Qua pasas de ideo ad altra, qua procedas per serchi inter-sucedanta (opoze ad intuicala)\lingui{DEF}
\artiklo{diskutar}%
\vorto{diskutar}\fako{trans.} Studiar problemo, questiono, ponderante la « por » e la « kontre ». \textit{(aludante personi qui ne havas la sama opiniono)}. Kambiar argumento-punti pri ula temo\lingui{DEFIRS}
\artiklo{dislokar}%
\vorto{dislokar}\fako{trans.} Separar violentoze (la parti di ensemblo)\lingui{DEFIS}
\artiklo{dismo}%
\vorto{dismo}\senco{1}{Dekimo de la rekoltajo (che la Judi, olim) presustracionata por ofrar a la Sinioro o donata a la levidi}\senco{2}{Parto (variiva) de la rekoltajo, presustracionata da la eklezio katolika e da la siniori feudala}\lingui{FIS}
\artiklo{disociar}%
\vorto{disociar}\fako{trans.) (kemio} Separar la elementi unionita. \textit{(anke metaf.)}\lingui{DEFIRS}
\artiklo{disonancar}%
\vorto{disonancar}\fako{netrans.} Formacar harmonio poke agreabla audale\lingui{DEFIRS}
\artiklo{dispacho}%
\vorto{dispacho}\fako{komerco} Aranjo, konvenciono pri la perdi, la avarii, inter kompanio di asekuro navala e la asekurito\lingui{DFS}
\artiklo{disparata}%
\vorto{disparata} Qua shokante dessimilesas to quo cirkondas lu\lingui{DEFIS}
\artiklo{dispensar}%
\vorto{dispensar}\fako{trans., de}\senco{1}{Permisar specale (ad ulu) agar o facar ulo quon interdiktas la regulo di la eklezio katolika, o ne agar o facar to quon lu preskriptas}\senco{2}{Permisar (ad ulu) ne agar o facar ulo}\lingui{DEFIRS}
\artiklo{dispensario}%
\vorto{dispensario} Loko en qua la povri povas konsultar mediki e recevar flegi o medikamenti (sen lojar en olu)\lingui{EFIS}
\artiklo{dispepsio}%
\vorto{dispepsio} Desfacileso pri digestar\lingui{DEFIS}
\artiklo{dispersar}%
\vorto{dispersar}\fako{trans.} Separar per pulsar li, per igar li dis-irar (kozi, personi, qui asemblesis)\lingui{DEFIS}
\artiklo{dispiremo}%
\vorto{dispiremo}\fako{histol.} Spiremo duopla qua, dum mitoso, aparas an la poli di la celulo, kande la kromosomi rigrupigas su por formacar la elementi nuklea di la du celuli nova
\artiklo{disponar}%
\vorto{disponar}\fako{trans.} Agar segun sua arbitrio pri ulo, pri ulu; darfar agar, uzar, ulo quale onu volas lo\lingui{DEFIRS}
\artiklo{disputar}%
\vorto{disputar}\fako{netrans., pri, pro} Komencar kontre ulu lukto per vorti violentoza, grosiera\lingui{DEFIRS}
\artiklo{distar}%
\vorto{distar}\fako{netrans., de, per} Interseparesar per ula quanto de spaco o de tempo\lingui{DEFIRS}
\artiklo{distiko}%
\vorto{distiko}\senco{1}{\textit{(metriko antiqua)} Asemblajo de du versi; asemblajo de un hexametro ed un pentametro}\senco{2}{\textit{(prozodio moderna. Poemeto de du versi)}}\lingui{DEFIS}
\artiklo{distilar}%
\vorto{distilar}\fako{trans.} Purigar (substanco) per vaporigar la parto volatila e kondensar ica, per (kom exemplo) alambiko\lingui{DEFIRS}
\artiklo{distingar}%
\vorto{distingar}\fako{trans., per} Igar rikonocar (persono, kozo) per atribuar a lu karakterizivo, traito, marko, tale ke onu povas dicernar lu\lingui{DEFIS}
\artiklo{distinta}%
\vorto{distinta} Markizita ye karaktero partikulara, qua impedas konfundar lu kun ulo analoga
\artiklo{distraktar}%
\vorto{distraktar}\fako{trans.} Deturnar la mento, la atenco de lua objekto, sive per amuzivi, sive per irga moyeno altra. Existas menti nature distraktita o distraktebla, qui distraktesas da irga kauzo extera, o mem da nula\lingui{EFIS}
\artiklo{distributar}%
\vorto{distributar}\fako{trans., inter, che} Dividar, donante un parto a singlu, inter plura personi o plura loki, kun dicerno, ne hazarde\lingui{DEFIRS}
\artiklo{distrikto}%
\vorto{distrikto} Teritorio limitoza qua, en ula landi, resortisas a ta o ca judiciistaro\lingui{DEFIRS}
\artiklo{ditenar}%
\vorto{ditenar}\fako{trans.} Konservar (ulo) kom posedajo til ula evento\lingui{DEFIS}
\artiklo{ditirambo}%
\vorto{ditirambo}\senco{1}{\textit{(en la epoki antiqua)} Poemo lirika por honorizar Bako}\senco{2}{\textit{(metaf.)} Laudo entuziasmoza}\lingui{DEFIRS}
\artiklo{ditreso}%
\vorto{ditreso}\senco{1}{Kordio-preseso laceranta}\senco{2}{Situeso kordio-laceranta}\senco{3}{Situeso di ta qua balde perisos se ne helpota quik}\lingui{EF}
\artiklo{divagar}%
\vorto{divagar}\fako{netrans.} Lasar sua penso irar exter la limiti di raciono\lingui{EFIS}
\artiklo{divano}%
\vorto{divano} Kanapeo sen dors-apogilo, ma kun kuseni movebla quin onu apogas kontre muro\lingui{DEFIRS}
\artiklo{divenar}%
\vorto{divenar}\fako{netrans.} Komencar esar (ulo) ye ta o ca instanto\lingui{FI}
\artiklo{divergar}%
\vorto{divergar}\fako{netrans.} \textit{(aludante elementi plu o min inter-apuda an la departo-punto)} Inter-eskartar sempre plu multe, segun ke li plulongeskas\lingui{DEFIS}
\artiklo{diversa}%
\vorto{diversa}\senco{1}{\textit{(aludante kozi quin onu komparas)} De qua singla havas karaktero diferanta de olta di la ceteri}\senco{2}{Plura (personi o kozi) di qui la naturo, la karaktero esas interdiferanta}\lingui{DEFIRS}
\artiklo{diversionar}%
\vorto{diversionar}\fako{netrans.}\senco{1}{\textit{(pri operaco militala)} Deturnar la enemiko de la punto quan lu atakas, per koaktar lu a su-defenso en punto altra}\senco{2}{Ago qua deturnas la spirito, la kordio, de to quo fatigas lu, suciigas lu, ad objekto altra}\lingui{DEFIRS}
\artiklo{divertikulo}%
\vorto{divertikulo}\fako{anat.} Apendico kava, analoga a sako-strado, qua debushas en kavajo naturala di la homo-korpo
\artiklo{dividar}%
\vorto{dividar}\fako{trans.}\senco{1}{\textit{(matem.)} Donite la produto de du faktori ed un ek ici, determinar la faktoro nekonocata}\senco{2}{Igar, de toto, plura parti}\lingui{DEFIRS}
\artiklo{divinar}%
\vorto{divinar}\fako{trans.}\senco{1}{Savar, per moyeni supernatura, to quo celesas en la tempo pasinta, prezenta, futura}\senco{2}{Deskovrar, per konjekti, supozi, interpreti, to quon onu ne savas}\lingui{EFIS}
\artiklo{diviziono}%
\vorto{diviziono}\senco{1}{En armeo, grupo de armeani qua kompozesas ek adminime du brigadi}\senco{2}{En navaro, parto de eskadro qua kompozesas ek adminime tri milito-navi komandata da un chefo}\lingui{DEFIRS}
\artiklo{divorcar}%
\vorto{divorcar}\fako{netrans., de} Dissolvar legale la mariajo per qua onu esas unionita a la spozo\lingui{EFIS}
\artiklo{divulgar}%
\vorto{divulgar}\fako{trans.} Dissavigar da multa personi\lingui{DEFIS}
\artiklo{dizastro}%
\vorto{dizastro} Desfortuno tre granda e subita, qua ruinas multa homi o tota socio, exemple : grava vinkeso, granda incendio, naufrajo\lingui{EFIS}
\artiklo{djin}%
\vorto{« djin »}\fako{che la Arabi} Demono
\artiklo{do}%
\vorto{« do »}\fako{muziko} Noto unesma di la gamo ordinara, sen acidento
\artiklo{do}%
\vorto{do} Konjunciono, uzata por kunduktar la konsequanto, la konkluzo di lo preiranta\lingui{FI}
\artiklo{doario}%
\vorto{doario}\fako{yuro olima}\senco{1}{Parto de la havaji donita da la spozulo a sua spozino, quan elu posedeskis se elu divenis vidva, e qua transpasis ad ilua filii}\senco{2}{Revenuo atribuata a rejino, a princino salika}\lingui{FI}
\artiklo{*doblar}%
\vorto{*doblar}\fako{trans.} \textit{(cinemo-enrejistro)} \textit{(pri parol-filmo)} Tradukar lo dicata, preske vortope (remplase di la pronuncatajo originala ye altra linguo), sorgante selektar vorti qui korespondas ne nur a la senco e sentimentonuanci, ma anke (same-longa) a la formo lor-pronunca di la boko di la aktoro
\artiklo{docar}%
\vorto{docar}\fako{trans., ad} Komunikar ad ulu (cienco, arto) per lecioni intersequanta\lingui{L}
\artiklo{docento}%
\vorto{docento} En la universitati di Germania, profesoro da qua la lecioni pagesas, ne da la stato, ma da la studenti ipsa\lingui{DEFR}
\artiklo{docila}%
\vorto{docila} Qua havas bona tendenci pri lasar su instruktesar e duktesar\lingui{EFIS}
\artiklo{dodekaedro}%
\vorto{dodekaedro}\fako{geom.} Poliedro kun dek-e-du edri\lingui{DEFIRS}
\artiklo{dodekagina}%
\vorto{dodekagina}\fako{bot.} Qua havas dek-e-du stamini\lingui{DEFIRS}
\artiklo{dogano}%
\vorto{dogano} Administrerio qua impostas la vari taxita, kande oli eniras od ekiras la teritorio\lingui{FIS}
\artiklo{dogmatiko}%
\vorto{dogmatiko} La teologio qua traktas la dogmato Kristana\lingui{DEFIRS}
\artiklo{dogmato}%
\vorto{dogmato} Kredendajo o punto, en doktrino religiala o filozofiala\lingui{DEFIRS}
\artiklo{dogo}%
\vorto{dogo}\fako{zool.} Hundo grosa e kurta, di qua la muzelo esas kurta, la maxilo e la mandibulo forta, la labii pendanta\lingui{DF}
\artiklo{dogro}%
\vorto{dogro} Mikra navo pontoza, quan onu uzas por peskar la haringi e la makreli\lingui{DEFIS}
\artiklo{dojo}%
\vorto{dojo} Duko elektita, chefo nominala di la republiki olima di Venezia e di Genova (Italia)\lingui{DEFIR}
\artiklo{doko}%
\vorto{doko} Maro-baseno por konstruktar o reparar navi; baseno cirkondata da kayi e warfi por charjar o descharjar navi\lingui{DFR}
\artiklo{doktoro}%
\vorto{doktoro} La persono qua obtenis la grado universitatala maxim alta en fakultato (teologio, yuro-cienco, medic., cienci, literaturo)\lingui{DEFIRS}
\artiklo{doktrino}%
\vorto{doktrino} Ensemblo de nocioni propozata da ulu kom docenda pri temo\lingui{DEFIRS}
\artiklo{dokumento}%
\vorto{dokumento} Texto skribura, raporto-texto, akto, qua uzesas por klarigar fakti historiala, judiciala, e c.\lingui{DEFIRS}
\artiklo{dolca}%
\vorto{dolca}\senco{1}{Qua plezas a la gustado pro saporo agreabla}\senco{2}{\textit{(metaf.)} Qua donas juo subtila}\lingui{FIS}
\artiklo{dolikocefala}%
\vorto{dolikocefala}\fako{homo} Di qua la kranio esas longa\lingui{DEFRS}
\artiklo{dolio}%
\vorto{dolio}\fako{tekn.} Kavajo qua esas ye la extremajo di la instrumento specala qua uzesas en la rivetago por impedar la riveto ekbatesar per la martelagi\lingui{E}
\artiklo{dollar}%
\vorto{« dollar »}\fako{ekonomiko} Monetuno arjenta od ora, en Usa e Kanada
\artiklo{dolmano}%
\vorto{dolmano} Vestio longa, ornita per frogi, quan *weris ula armeani (kom exemplo : en Francia, la husari, la kavalchaseri)\lingui{DEFIS}
\artiklo{dolmeno}%
\vorto{dolmeno}\fako{arkeol.} Monumento megalita, qua konsistas ye un petro plata, pozita sur un o plura petri erektita vertikale\lingui{DEFRS}
\artiklo{dolo}%
\vorto{dolo}\fako{yuro-cienco} Trompo (per tacar pri ula fakti konocinda da la koncernato; o per kredigar, da ica, fakti qui ne existas) qua detrimentas la interesti di ulu\lingui{DFIS}
\artiklo{dolomito}%
\vorto{dolomito} Karbonato de kalko o de magnezio, quan onu renkontras, kom roki, en masi quaze-marmora od en kristali di la sistemo romboedrala\lingui{DEFIRS}
\artiklo{dolorar}%
\vorto{dolorar}\fako{netrans., pro, de, pri} Subisar sento laceranta de ula lezuri, de ula standi morbala qui afektas la korpo\lingui{FIS}
\artiklo{domajar}%
\vorto{domajar}\fako{trans., per, pro} Facar altero materiala, qua minbonigas, minaptigas, kozo (o la korpo di persono)\lingui{EFI}
\artiklo{domeno}%
\vorto{domeno}\senco{1}{Sulo-peco quan ulu proprietas}\senco{2}{Yuro di proprieto e di sinioreso ye ulu od ulo. – \textit{(anke metaf.)}}\lingui{DEFIRS}
\artiklo{domestika}%
\vorto{domestika}\fako{pri animali} Qua vivas apud la homo e servas lu por satisfacar lua bezoni o por lua plezuro\lingui{EFIS}
\artiklo{domicilo}%
\vorto{domicilo}\senco{1}{La habiteyo legala, oficala, di ulu}\senco{2}{\textit{(astrologio)} Zodiako-signo sub qua astro konsideresis kom influanta maxime}\lingui{DEFIS}
\artiklo{dominacar}%
\vorto{dominacar}\fako{trans.}\senco{1}{Havar sub sua hegemonio, prepondero}\senco{2}{Havar apud-sub su}\lingui{DEFIRS}
\artiklo{dominanto}%
\vorto{dominanto}\senco{1}{\textit{(muziko)} En la gamo ordinara, la noto kinesma, ta qua (pro formacar kun la unesma la konsonanco maxim perfekta) determinas la tono ed uzesas kom fundamento di la harmonio}\senco{2}{En la plankanto, noto sur qua onu kantas la korpo di verso}\senco{3}{\textit{(kristalografio)} Tipo geometriala normala segun qua korpo kristaleskas}\lingui{DEF}
\artiklo{dominikano}%
\vorto{dominikano} Reguliero de la ordeno di la frati predikera, fondita da Santa Dominikus\lingui{DEFIRS}
\artiklo{domino}%
\vorto{domino}\senco{1}{Robo flotacanta, kapucoza, per qua onu su vestizas kande onu iras por masko-balo}\senco{2}{Ludilo ek duadek-ed-ok mikra paralelipipedi osta od ivora, dividita a du parti, per qui onu obtenas la kombinuri diversa : blanka-blanka, blanka-aso, blanka-du, blanka-tri, e c. til sis-sis}\lingui{DEFIRS}
\artiklo{domo}%
\vorto{domo}\senco{1}{Konstrukturo destinita ad uzesor kom habiteyo}\senco{2}{Konstrukturo destinita ad uzo specala, publika o privata}\lingui{EFRL}
\artiklo{domtar}%
\vorto{domtar}\fako{trans.}\senco{1}{Koaktar ad obedio (animalo ne-amansita)}\senco{2}{Koaktar ad obedio (ti qui refuzas submisar su)}\lingui{DFIS}
\artiklo{donacar}%
\vorto{donacar}\fako{trans.} Abandonar ad ulu la proprieto di ulo, sen recevar ulo po olu\lingui{EFIRS}
\artiklo{donar}%
\vorto{donar}\fako{trans., ad}\senco{1}{Igar disponebla da ulu}\senco{2}{Propozar kondiciono determinita en la enunco di problemo}\lingui{EFIS}
\artiklo{dop}%
\vorto{dop} An la latero opozita ad olta ube situesas la vizajo di persono o la facio di kozo\lingui{I}
\artiklo{dorado}%
\vorto{dorado}\fako{zool.}\senco{1}{Mar-fisho de la familio « komberidi »}\senco{2}{Mar-fisho de la familio « sparoidi »}\latina{Coryphaena hippurus; Chrysophrys aurata}\lingui{DFIRS}
\artiklo{dorika}%
\vorto{dorika}\fako{arkitekt.} La duesma ek la kin grupi arkitekturala che la Antiqui, karakterizata da koloni kaneloza, sen-baza, kapitelo ek granda muluro kalicatra, fuzo interruptata da triglifi\lingui{DEFIRS}
\artiklo{dorlotar}%
\vorto{dorlotar}\fako{ulu} Traktar tenere, afablege\lingui{F}
\artiklo{dormar}%
\vorto{dormar}\fako{netrans.} Esar en ta stando dum qua che homo od animalo, haltigesas ula funcioni di la vivo (precipue dum nokto), pro la neceseso di lia repozo\lingui{EFIS}
\artiklo{dorno}%
\vorto{dorno}\fako{bot.} Pikanto qua kreskas sur ula planti\lingui{DE}
\artiklo{dorso}%
\vorto{dorso}\fako{anat.} Facio dopa di la homo-korpo; facio supera di la animalo-korpo\lingui{EFIS}
\artiklo{doselo}%
\vorto{doselo} Kronatro qua superesas altaro, trono, e quan onu portas en la procesioni super la eukaristio\lingui{S}
\artiklo{dotar}%
\vorto{dotar}\fako{trans., per, ye} Provizar (homo) per la havajo quan lu adportas mariajo a sua spozo. – Provizar per ula qualesi\lingui{EFIS}
\artiklo{dozo}%
\vorto{dozo}\senco{1}{Quanto de medikamento, quan la malado devas absorbar en un foyo; quanto quan onu devas inkluzar, en medikamento kompozaja, de singla ek la substanci ye qui lu konsistas}\senco{2}{Quanto de to quo formacas kompozajo}\lingui{DEFIRS}
\artiklo{draceno}%
\vorto{draceno}\fako{bot.} Planto ornamentala\latina{dracaena}\lingui{DFS}
\artiklo{dragar}%
\vorto{dragar}\senco{1}{Skurar (portuo, rivero, e c.) per deprenar de la fundo di la skurajo la gravio, la slamo, la rezidui}\senco{2}{Arachar (ostri, musli) per skrapilo fera, adaptita a reto sakatra}\senco{3}{Probar peskar (ankro di qua la boyo) esas perdita per kordego garnisita ye grapino quan onu tranas sur la fundo di portuo, rivero, e c.}\lingui{EFS}
\artiklo{dragomano}%
\vorto{dragomano} Interpretisto en Oriento\lingui{DEFIRS}
\artiklo{dragono}%
\vorto{dragono}\senco{1}{\textit{(olim)} Soldato de la kavalrio lejera}\senco{2}{\textit{(nun)} Soldato de la kombato-infantrio, qua havas kasko stala kun krinaro pendanta}\lingui{DEFIRS}
\artiklo{drajeo}%
\vorto{drajeo} Mandelo, avelano, ed altra kozo kovrita per sukrajo harda\lingui{F}
\artiklo{drakmo}%
\vorto{drakmo}\senco{1}{Unajo di pondero e di moneto che la Greki antiqua}\senco{2}{Moneto-peco arjenta, che la Greki moderna, qua valoras cirkume un franko ora}\lingui{DEFIRS}
\artiklo{drako}%
\vorto{drako}\senco{1}{Animalo mitologiala, quan la Antiqui reprezentis kom havanta ali, skrach-ungli, serpento-kaudo, e c.}\senco{2}{\textit{(blaz.)} Reptero reprezentata en la eskudo kom havanta du pedi ed un kaudo}\senco{3}{\textit{(zool.)} Lacerto di India, qua vivas sur la arbori, e che qua la flanko-pelo povas extensar su tale ke olu formacas quaza ali e parafalo}\senco{4}{\textit{(astron.)} Stelaro en la mi-sfero boreala}\latina{draco}\lingui{DEFIRS}
\artiklo{drakonala}%
\vorto{drakonala} Di qua la severeso esas extrema\lingui{DEFIRS}
\artiklo{dramato}%
\vorto{dramato}\senco{1}{Teatrajo, en qua lo komika mixesas a lo tragediala}\senco{2}{\textit{(aludante eventajo ek la vivo reala)} Ensemblo de cirkonstanci violentoza, solvata per katastrofo}\lingui{DEFIRS}
\artiklo{drapirar}%
\vorto{drapirar}\fako{trans.} Dispozar (tapeto-stofo) tale ke lu formacas ondi harmonioza\lingui{DEFIR}
\artiklo{drapo}%
\vorto{drapo} Stofo ek warpo e wefto lana, di qua la texuro (submisite a lanugizo ed a fulo) esas lanugoza\lingui{F}
\artiklo{drashar}%
\vorto{drashar}\fako{trans.} Frapar repetoze sur la garbi (de cerealo) rekoltita e dispozita sur sulo-peco plana ed harda, por separar la grani de la spiki\lingui{DE}
\artiklo{drastika}%
\vorto{drastika}\fako{medic.} Qua purgas forte\lingui{DEFIRS}
\artiklo{drenar}%
\vorto{drenar}\fako{trans.) (agrokultivo} Liberigar (tereno humida) de la aquo ecesanta, plufaciligante la ekfluo per kanaleti ek briki, teguli, e c., quin onu dispozas ye ula profundeso-grado\lingui{DEFIRS}
\artiklo{dresar}%
\vorto{dresar} \textit{(trans.)} Kustumigar animalo pri agar docile e reguloze . (Onu dresas la animali por igar li exekutar ula movi determinita – ofte ne-naturala – o por servar la homo – kom ex. kavalkado – o por agar irga prodalaji)\lingui{DFR}
\artiklo{drezino}%
\vorto{drezino} Nomo di la biciklo lor lua inventeso\lingui{F}
\artiklo{driado}%
\vorto{driado}\fako{mitol.} Deino di la boski\lingui{DEFIS}
\artiklo{driftar}%
\vorto{driftar}\fako{netrans.} Esar pulsata da vento o marfluo deviace de sua voyo\lingui{DE}
\artiklo{drilo}%
\vorto{drilo}\fako{tekn.} La instrumento qua portas e movas la borilo en metalo\lingui{DER}
\artiklo{drinkar}%
\vorto{drinkar}\fako{trans.} Englutar liquido, aspirante lu per la boko\lingui{DE}
\artiklo{drogo}%
\vorto{drogo} Ingrediento uzata en kemio, farmacio, tinto, ed en la ekonomio hemala\lingui{DEFIS}
\artiklo{drola}%
\vorto{drola} Qua havas en su ulo stranja e plezanta\lingui{DEF}
\artiklo{dromedaro}%
\vorto{dromedaro}\fako{zool.} Kamelo un-giba\latina{Camelus dromedarius}\lingui{DEFIRS}
\artiklo{dronar}%
\vorto{dronar}\fako{trans.} Mortigar ulu, su, per asfixio en liquido\lingui{E}
\artiklo{dronto}%
\vorto{dronto}\fako{zool.} Granda ucelo kurta-pluma, di la insuli Maurico e Reunion, di qua la speco ne plus existas\latina{Didus ineptus}\lingui{DFS}
\artiklo{drugeto}%
\vorto{drugeto} Stofo de lano, de lino e lano, de lano e silko, de silko, sur qua quaze-semesas desegnuri rekamita qui ne texesas en la fundo di la stofo\lingui{DEF}
\artiklo{druido}%
\vorto{druido} Sacerdoto di la Galli (Kelti)\lingui{DEFIRS}
\artiklo{drupo}%
\vorto{drupo}\fako{bot.} Frukto tre pulpoza, simpla, nedehicenta, kernoza (abrikoto, persiko, pruno, e c.)\lingui{EFIS}
\artiklo{druzo}%
\vorto{druzo}\fako{mineral.} En mineralo, inkrustajo de kristali di naturo interdiferanta, mineralo en qua trovesas ta kristali\lingui{DEFS}
\artiklo{dubitar}%
\vorto{dubitar}\fako{netrans.}\senco{1}{Ne savar quan, ek du opinioni interkontrea, esas la exakta}\senco{2}{Ne savar kad existas ulo exakta (pri)}\lingui{EFIS}
\artiklo{*dublar}%
\vorto{*dublar}\fako{trans.} Preter-vehar (generale per automobilo) altra duktanto, pasante latere di ica lineo-dop qua onu avancis
\artiklo{duelar}%
\vorto{duelar}\fako{netrans.} Interkombatar \textit{(aludante nur du adversi.)} Interkombatar \textit{aludante du privati, de qui la una postulas} de la altra la honor-satisfaco (pos ofenseso) per armo-lukto\lingui{DEFIRS}
\artiklo{dueno}%
\vorto{dueno}\senco{1}{En Hispania, muliero pasable olda, quaza edukistino, adjuntita ad adolecantino od a yuna muliero di rango nobela, por moro-surveyar elu sencese}\senco{2}{Oldino qua moro-surveyas adolecantino o yuna muliero}\lingui{DEFS}
\artiklo{dueto}%
\vorto{dueto} Muzikajo por du voci o du instrumenti koncertanta, kun o sen akompanuro\lingui{DEFIRS}
\artiklo{dukato}%
\vorto{dukato} Olima monet-uno de oro pura, analoga a la florino, di qua la valoro equivalas 5 o 6 franki ora\lingui{DEFIRS}
\artiklo{duko}%
\vorto{duko}\senco{1}{La homo qua havis la guvernereso, la sinioreso, di teritorio qua kovris plu kam un komtio}\senco{2}{Suvereno di ula stati (duko di Brunswick, Granda Duko di Baden)}\senco{3}{Sinioro ye la grado maxim alta sub la rejo}\senco{4}{Princo di suvereno-lineo qua recevis la titulo (granda Duko di Rusia)}\lingui{EFIS}
\artiklo{duktar}%
\vorto{duktar}\fako{trans.} Igar (persono, kozo) agar od irar en ta o ca maniero, per direktar ed akompanar lu\lingui{EFIS}
\artiklo{duktila}%
\vorto{duktila} Di qua onu povas extensar, longigar (filoforme) la materio sen ruptar lu\lingui{DEFIS}
\artiklo{dulkamaro}%
\vorto{dulkamaro}\fako{bot.} Planto de la genero « solani », olim tre uzata kontre ula pelo-morbi\latina{Solanum dulcamara}\lingui{DEFIS}
\artiklo{dum}%
\vorto{dum} En la tempo-quanto kande ulo esas eventanta, agata o facata\lingui{L}
\artiklo{dungo}%
\vorto{dungo} Omno quon onu mixas a la kultivo-tero por igar lu plu fertila\lingui{DE}
\artiklo{duno}%
\vorto{duno} Kolino sabla, sur litoro\lingui{DEFIRS}
\artiklo{duodenumo}%
\vorto{duodenumo}\fako{anat.} Parto komencala di la intestino tenua\lingui{DEFIS}
\artiklo{duplexa}%
\vorto{duplexa}\fako{telegrafado} Qua posibligas sendar sam-instante, sur un filo, mesaji en la du sinsi
\artiklo{duplikato}%
\vorto{duplikato} Kopiuro, exemplero duesma di akto, kontrato-texto, quitigo-texto, epistolo, e c.\lingui{EFL}
\artiklo{dupliko}%
\vorto{dupliko}\fako{yuro-cienco} Respondo a repliko\lingui{DFIS}
\artiklo{dupo}%
\vorto{dupo} Persono quan onu trompas per misuzar lua naiveso\lingui{DEF}
\artiklo{duramatro}%
\vorto{duramatro}\fako{anat.} Ex la tri membrani qui envelopas la maso cerebra e la myelo, la maxim dika e solida, harda\lingui{DEFIS}
\artiklo{durar}%
\vorto{durar}\fako{trans. e netrans.} Ne-cesar existar, permanar plu o min longe.- Ne-cesar (sua ago)\lingui{DEFIRS}
\artiklo{durstar}%
\vorto{durstar}\fako{netrans. e trans.}\senco{1}{Bezonar drinko}\senco{2}{Dezirar nepaciente}\lingui{DE}
\artiklo{dushar}%
\vorto{dushar}\fako{trans.}\senco{1}{Arozar per aquo-sprico direktata a ta o ca parto di la korpo, por efikar medicinale}\senco{2}{Lavar (la homo-korpo) per aquo qua lansesas pluve o cirkle}\lingui{DEFIRS}
\artiklo{duvo}%
\vorto{duvo} Singla ek la planki kurvatra qui formacas la korpo di kuvo ligna, di barelo, e qui juntesas per ringi fera\lingui{DF}
\artiklo{duyongo}%
\vorto{duyongo}\fako{zool.} Mar-hundo de la insularo Malayala (oceano Indiana)\latina{halicore dugong}\lingui{DEFIRS}
\end{multicols}
\sekciono{E}
\begin{multicols}{2}
\artiklo{e}%
\vorto{e}\fako{muziko} Noto sepesma di la « do »-gamo
\artiklo{-e-}%
\vorto{-e-} Infixo qua signifikas « samakolora kam » o « sam-aspekta kam »
\artiklo{-e}%
\vorto{-e} Sufixo di la adverbi derivita
\artiklo{e(d)}%
\vorto{e(d)} Konjunciono qua uzesas por interligar la frazi o la propozicioni koordinita
\artiklo{ebeno}%
\vorto{ebeno}\fako{bot.} Ligno de la ebeniero, tre pezoza, tre harda, di qua la grano esas tre densa, e remarkinda pro la beleso di lua obskur-nigreso\latina{dicapyros ebenum}\lingui{DEFIS}
\artiklo{-ebl-}%
\vorto{-ebl-} Sufixo qua signifikas « quan onu povas… »
\artiklo{ebonito}%
\vorto{ebonito} Kauchuko hardigita per vulkanizo, de qua onu facas pektili, juveli, izolivi por la aparati elektrala\lingui{DEFIRS}
\artiklo{ebria}%
\vorto{ebria}\senco{1}{Di qua la cerebro-funciono trublesas da drinkir tro multe vino od altra drinkajo fermentacita}\senco{2}{Di qua la spirito trublesas da la deliro de pasiono}\lingui{EFIS}
\artiklo{ebuliar}%
\vorto{ebuliar}\fako{netrans.} Agitesar \textit{(aludante la surfaco di liquido)} kande lu bolieskas\lingui{EFIS}
\artiklo{ebulioskopo}%
\vorto{ebulioskopo} Aparato per qua onu determinas, per bolio, la quanto de alkoholo quan kontenas alkoholajo
\artiklo{ecelar}%
\vorto{ecelar}\fako{netrans.) (en, pri} Esar, en sua genero, ye grado eminenta\lingui{DEFI}
\artiklo{ecelenco}%
\vorto{ecelenco} Titulo honorala, donata a la ministri, a la ambasadisti\lingui{DEFIS}
\artiklo{ecentrika}%
\vorto{ecentrika}\senco{1}{\textit{(geom.)} Di qua la centro eskartesas de punto donita}\lingui{DEFIS}\subvorto{}{ecentriko}{Disko uzata por transformar movo cirklo-kurva, per ranuro kalkulita segun la movo produktenda rektalinea (iro e retroveno), ed en qua cirkulas roteto qua komandas la mekanismo}{2}{}
\artiklo{eceptar}%
\vorto{eceptar}\fako{trans.} Lasar extere persono, objekto, pri qua onu afirmas ulo\lingui{EFIS}
\artiklo{ecesar}%
\vorto{ecesar}\fako{trans.}\senco{1}{Transirar limito fixigita}\senco{2}{Trans-irar to quon ulu povas tolerar}\lingui{DEFIRS}
\artiklo{echekar}%
\vorto{echekar}\fako{trans.} Kontreagar vinkoze sua adverso, e minacar lu
\artiklo{echopo}%
\vorto{echopo} Mikra butiko, qua konsistas, maxim-multa-kaze, ye shirmo-tekto apogita an domo, an la muro di qua lu fixigesis en angulo konkava\lingui{F}
\artiklo{ecitar}%
\vorto{ecitar}\fako{trans.}\senco{1}{Produktar (movo) en la psiko, en la organismo}\senco{2}{Advokar (ulu) a psiko-movo (sentimento, rezolvo, e c.)}\senco{3}{Igar (movo psikala) plu vivaca}\senco{4}{Igar (ulo, la organismo) plu vivaca, plu ardoroza pri sentar, volar, e c.}\lingui{EFIS}
\artiklo{ecizar}%
\vorto{ecizar}\fako{trans.) (kirurg.} Ablacionar parti mikra-volumina, per instrumento tre tranchiva
\artiklo{-ed-}%
\vorto{-ed-} Sufixo qua signifikas « la kontenajo di »
\artiklo{ed}%
\vorto{ed} (videz « e »)
\artiklo{edelweiss}%
\vorto{« edelweiss »}\fako{bot.} Nomo vulgara di planto alpala ek la familio « kompozaji », remarkinda pro la lanugo blanka e lanatra qua kovras lu, precipue lua foliin\latina{Leontopodium}
\artiklo{edemo}%
\vorto{edemo}\fako{patol.} Influro qua konsequas de infiltruro seroza en tisuo celulara\lingui{DEFIS}
\artiklo{edeno}%
\vorto{edeno}\fako{religio}\senco{1}{En Biblo : la paradizo terala}\senco{2}{\textit{(metaf.)} Loko di delici}\lingui{DEFIS}
\artiklo{ederduno}%
\vorto{ederduno} Mikra plumi dina e dolca qui kreskas en la vivokomenco che la edero, la cigno, la ganso, ed ek qui onu facas lito-kovrili, matraci, kapkuseni, apogo-kuseni\lingui{DEFS}
\artiklo{edero}%
\vorto{edero}\fako{zool.} Anado de la regioni tre Europa-norda, de qua onu obtenas la lanugo « ederduno »\latina{somateria}\lingui{DEFI}
\artiklo{edifico}%
\vorto{edifico} Ulo plu generala kam domo, kastelo, palaco. Quaze meza inter ta tri vorti\lingui{EFIRS}
\artiklo{edifikar}%
\vorto{edifikar}\fako{trans.}\senco{1}{Firmigar pri lua pieso}\senco{2}{Informar komplete pri persono, kozo}\lingui{EFIRS}
\artiklo{ediktar}%
\vorto{ediktar}\fako{trans.} Establisar per dekreto suverenala\lingui{DEFIR}
\artiklo{edipus}%
\vorto{« edipus »}\senco{1}{Nomo di la princo famoza Greka qua divinis la enigmato propozita da la sfinxo}\senco{2}{Ta qua divinas facile la enigmati}\lingui{I}
\artiklo{editar}%
\vorto{editar}\fako{trans.} Imprimar ed ofrar ad kompreri eventuala (libro, jurnal, e c.); imprimar la texto da ula autoro\lingui{DFIRS}
\artiklo{edro}%
\vorto{edro}\fako{geom.} Singla de la plani qui limitizas solido\lingui{G}
\artiklo{edukar}%
\vorto{edukar}\fako{trans.}\senco{1}{Formacar infanto, puero, adolecanto, per developar e direktar lua fakultati korpala, morala (etikala), mentala, religiala}\senco{2}{Kustumigar ulu pri la ceremonialo mondumala (politeso, konveni, deci, e c.)}\senco{3}{Exercar la animali pri lia kapablesi inata}\senco{4}{Nutrar, plumultigar, entratenar ula animali por lia manjeso, lia laboro e la plezuro quan onu recevas de li; same pri planti}\lingui{EFIS}
\artiklo{efacar}%
\vorto{efacar}\fako{trans.} Desaparigar to quo tracesis\lingui{EF}
\artiklo{efektivo}%
\vorto{efektivo}\senco{1}{\textit{(armeo, administrantaro)} La nombro de la individui qui fakte servas}\senco{2}{La nombro de la individui di qui la nomi enrejistresis en la etati di societo, e c.}\lingui{DEFIS}
\artiklo{efekto}%
\vorto{efekto} Rezultajo de la efiko di kauzo\lingui{DEFIRS}
\artiklo{efemera}%
\vorto{efemera} Qua duras dum nur un dio\lingui{DEFIS}
\artiklo{efemerido}%
\vorto{efemerido}\senco{1}{Tabeli astronomiala qui indikas, pri singla dio di la yaro, la situeso di la astri}\senco{2}{Libro qua indikas anticipe la fakti astronomiala o meteorologiala, kalkulebla e predicebla pri ta o ca epoko}\lingui{DEFIS}
\artiklo{efervecar}%
\vorto{efervecar}\fako{netrans.}\senco{1}{\textit{(aludante gaso)} Ebuliar tra liquido}\senco{2}{Emocar vivace, ma neduronte}\lingui{EFIS}
\artiklo{eficienta}%
\vorto{eficienta} Qua produktas efekto maxima po laboro, esforco, energio, spenso, minima\lingui{EF}
\artiklo{efikar}%
\vorto{efikar}\fako{netrans.} Produktar la efekto quan onu expektas de lu\lingui{EFIS}
\artiklo{eflorecar}%
\vorto{eflorecar}\fako{netrans.}\senco{1}{\textit{(kemio)} La efekto qua konsistas ye la formaco di kristali ye la surfaco di solido o liquido}\senco{2}{\textit{(fiziol.)} Fenomeno qua karakterizesas da mikra influro di la pelo}\lingui{DEFIS}
\artiklo{-eg-}%
\vorto{-eg-} Sufixo qua indikas « la grado maxim alta »
\artiklo{egala}%
\vorto{egala} Same-quanta o same-quala kam to a quo onu komparas lu\lingui{DEFIRS}
\artiklo{egardar}%
\vorto{egardar}\fako{trans.} Atencar kozo o persono, pro lua importo\lingui{EF}
\artiklo{egarir}%
\vorto{egarir}\fako{trans.} Lokizir distrakte (ulo) ube onu ne plus trovas lu\lingui{F}
\artiklo{egido}%
\vorto{egido}\senco{1}{\textit{(en la epoki antiqua)} Shildo di Pallas, kovrita per la felo di la kaprino Amalte, e sur qua reprezentesabis la kapo di Meduzo}\senco{2}{To quo protektas}\lingui{DEFIS}
\artiklo{eglefino}%
\vorto{eglefino}\fako{zool.} Fisho de la genero « gadi »\latina{gadus aeglefinus}\lingui{FL}
\artiklo{ego}%
\vorto{ego} To ye quo konsistas la individueso, la personeso\lingui{L}
\artiklo{egorjar}%
\vorto{egorjar}\fako{trans.} Ocidar per kultelo-frapo en la kolo (fauco)\lingui{F}
\artiklo{egotismo}%
\vorto{egotismo} Manio parolar pri su\lingui{EFIRS}
\artiklo{egotisto}%
\vorto{egotisto} Maniiko pri parolar pri su\lingui{DEFIRS}
\artiklo{eis}%
\vorto{« eis »}\fako{muziko} La noto triesma di la « do »-gamo diezoza
\artiklo{ejektar}%
\vorto{ejektar}\fako{trans.}\senco{1}{\textit{(mekan.)} Retrojetar adextere aero per aparato qua evakuas ta fluido per altra fluid-amaso qua movas tre rapide}\senco{2}{\textit{(pafarmo)} Ekjetar (la mufo di kartocho)}\lingui{DEFIS}
\artiklo{ek}%
\vorto{ek}\senco{1}{De interne ad extere…}\senco{2}{Indikas la materio di kozo (de qua onu imaginas ke lu esas extraktita)}\senco{3}{\textit{(metaf.)} Indikas objekto, o mem ento, qua apartenas ad ensemblo, kolektitaro (e quin onu supozas extraktita de olu)}\lingui{I}
\artiklo{ekarte}%
\vorto{« ekarte »} Kartoludo quan partoprenas du personi, e segun la regulo di qua la ludanti darfas forjetar omna o kelka karti, e prenar altri kom remplasanta\lingui{F}
\artiklo{ekidno}%
\vorto{ekidno}\fako{zool.} Mamifero Australiana, de la familio « sen-denti », di qua la muzelo esas longa e tenua, la korpo kovrata da pikanti quale la herisono, qua ovifas e nutras su ek insekti. (hystrix)\latina{echidna}\lingui{DEFIS}
\artiklo{ekimoso}%
\vorto{ekimoso}\fako{patol.} Makulo livida sur la pelo, produktita da la ekfluo di la sango en la tisuo celulara\lingui{DEFIS}
\artiklo{ekino}%
\vorto{ekino}\fako{zool.} Ekinodermo globatra, herisata ye pikanti (nomata anke « mar-herisono ».)\latina{echinus}
\artiklo{ekinodermo}%
\vorto{ekinodermo}\fako{zool.} Klaso de animali radioza, di qui la korpo esas garnisita ye multega pinti artikoza e movebla\lingui{DEFIS}
\artiklo{eklampsio}%
\vorto{eklampsio}\fako{patol.} Konvulsi qui eventas maxim-multa-kaze pro albuminurio. Konvulsi che infanteti\lingui{DEFIS}
\artiklo{eklektika}%
\vorto{eklektika} Qua asemblas en un doktrino-korpo la exaktaji quaze dissemita en sistemi diversa di filozofio e di cienco\lingui{DEFIRS}
\artiklo{eklezio}%
\vorto{eklezio} Asemblitaro de ti qui adoras la deo di la kristani : la eklezio kristana, asemblitaro de ti qui profesas la lego di Kristo\lingui{EFIRS}
\artiklo{eklipsar}%
\vorto{eklipsar}\fako{trans.}\senco{1}{\textit{(pri astro)} Cesigar lua videbleso}\senco{2}{\textit{(metaf.)} Obskurigar per brilar supera-grade}\lingui{EFIRS}
\artiklo{ekliptiko}%
\vorto{ekliptiko}\fako{astron.} Orbito quan Tero parkuras cirkum Suno\lingui{DEFIRS}
\artiklo{eklogo}%
\vorto{eklogo} Mikra poemo en qua encenigesas muton-pastori\lingui{DEFIRS}
\artiklo{eko}%
\vorto{eko}\fako{akust.} Reflekto de la son-ondi retrosendita a la oreli da surfaco qua produktas retroshoko\lingui{DEFIRS}
\artiklo{ekologio}%
\vorto{ekologio} Cienco di la relati da la organismi kun la medio en qua li vivas, kun la cirkonstanci organika (biologio) e neorganika (kosmala) di la existo
\artiklo{ekonomiko}%
\vorto{ekonomiko} La cienco pri la produktado, la cirkulado, la kambiado e la repartiso di la kozi valoroza\lingui{DEFIRS}
\artiklo{ekonomio}%
\vorto{ekonomio}\senco{1}{Administro financala di domo, societo o stato}\senco{2}{\textit{(metaf.)} \textit{(pri organismo)} Gani (nutrado) e spensi (perdi di materio e di energio)}\lingui{DEFIRS}
\artiklo{ekonomo}%
\vorto{ekonomo} En granda dom-edo od en institucuro : la persono di qua la rolo konsistas ye jerar la spensi\lingui{DFIRS}
\artiklo{ekritar}%
\vorto{ekritar}\fako{trans.) (komerco} Notar en la konto-libri la fluktui di la aktivo e di la pasivo di omna konto\lingui{F}
\artiklo{ektodermo}%
\vorto{ektodermo}\fako{embriol.} Strato extera di la blastodermo, qua pose divenos la epidermo e la nervaro
\artiklo{ektoplasmo}%
\vorto{ektoplasmo}\senco{1}{\textit{(biol.)} Parto periferia di la celulo, diferenciajo de la citoplasmo, tale ke la kambii da ica kun la medio eventas segun la naturo di la celulo}\senco{2}{\textit{(cienci okulta)} Plasmo, da origino psikala, qua (segun-dice) emanas de mediumo}\lingui{I}
\artiklo{ekumenika}%
\vorto{ekumenika}\fako{religio} Universala\lingui{DEFIS}
\artiklo{ekzemo}%
\vorto{ekzemo}\fako{patol.} Morbo di la pelo, karakterizata da la erupto di mikra vezikuli, sequata de exkorii surfacala\lingui{DEFIRS}
\artiklo{el(u)}%
\vorto{el(u)} Pronomo qua remplasas personino kande onu volas emfazar ke lu esas femina
\artiklo{elaborar}%
\vorto{elaborar}\fako{trans.}\senco{1}{Transformar, produktar per laborado. \textit{(anke metaf.)}}\senco{2}{\textit{(fiziol.)} Igar absorbebla od asimilebla, substanco od igar ica tala ke lu esas ekpulsebla de la organismo o de la korpo}\lingui{EF}
\artiklo{elastika}%
\vorto{elastika} Qua esas kontraktebla, extensebla, deformebla per preso, tiro, e c. e retrohavas sua volumino, sua formo komencala kande ta efiko cesas\lingui{DEFIRS}
\artiklo{elastino}%
\vorto{elastino} Materio ek la klaso « albuminoidi », quan onu renkontras en la tegumenti ed en la subten-organi
\artiklo{elefantiazo}%
\vorto{elefantiazo}\fako{patol.} Quaza lepro, karakterizata da tuberkuli di la pelo qui igas olca rugoza quale ta di la elefanto\lingui{DEFIRS}
\artiklo{elefanto}%
\vorto{elefanto}\fako{zool.} Mamifero de la klaso « pakidermi », remarkinda pro la maso enorma di lua korpo, lua pelo sen-pila e rugoza, la dentegi qui garnisas lua maxilo, e lua nazo longigita a rostro\latina{elephas}\lingui{DEFIRS}
\artiklo{eleganta}%
\vorto{eleganta}\senco{1}{Karakterizata da distingateso plena ye gracio e ye sen-jeneso}\senco{2}{Di qua la manieri esas gracioza e senpene tala}\lingui{DEFIRS}
\artiklo{elegio}%
\vorto{elegio}\fako{en la epoki antiqua} Mikra poemo di karaktero melankolioza e tenera\lingui{DEFIRS}
\artiklo{elektar}%
\vorto{elektar}\fako{trans., por… -eso) (aludante plura personi} Selektar per voto (por ofico, plaso, posteno)\lingui{EFIS}
\artiklo{elektro}%
\vorto{elektro}\fako{fiziko} Proprajo quan ula korpi posedas od aquiras, kande li frotesas, kompresesas, e c., atraktar o retropulsar altra korpi, produktar cintili, deskompozar kemiale, sukusar la nervi, e c.\lingui{DEFIRS}
\artiklo{elektrodo}%
\vorto{elektrodo}\fako{fiziko} Konduktivo quan onu igas komunikar kun calatere la baterio, talatere medio a qua la fluo elektra efikas kemiale\lingui{DEFIRS}
\artiklo{elektroforo}%
\vorto{elektroforo} Aparato qua uzesas por akumular elektro\lingui{DEFIRS}
\artiklo{elektrolito}%
\vorto{elektrolito} Korpo qua, liquida, esas deskompozebla da la elektro-fluo\lingui{DEFIRS}
\artiklo{elektrolizar}%
\vorto{elektrolizar}\fako{trans.} Deskompozar kemiale per la fluo elektra\lingui{DEFIRS}
\artiklo{elektrologio}%
\vorto{elektrologio} Parto di fiziko qua traktas la fenomeni elektrala e magnetala (elektro-magnetala), e lia relati kun altra fenomeni\lingui{DEFIRS}
\artiklo{elektrometrio}%
\vorto{elektrometrio} Ensemblo de la metodi pri la mezuro di la elektro-grandori (intenseso, rezistiveso, e c.)\lingui{DEFIRS}
\artiklo{elektrometro}%
\vorto{elektrometro} Instrumento uzata por mezurar la intenseso o la naturo (pozitiva o negativa) di la fluo elektra ye qua kargesas korpo\lingui{DEFIRS}
\artiklo{elektromobilo}%
\vorto{elektromobilo} Veturo movata per elektro-motoro (vice per kavali o per gas-motoro)\lingui{DEFIRS}
\artiklo{elektrono}%
\vorto{elektrono}\fako{fiziko} Korpuskulo qua portas la maxim mikra kargajo izolebla ek elektro generale negativa\lingui{DEFIRS}
\artiklo{elektroskopio}%
\vorto{elektroskopio} Studiado di la elektroskopi e di la aplikokazi elektroskopala\lingui{DEFIRS}
\artiklo{elektroskopo}%
\vorto{elektroskopo} Aparato per qua onu povas konstatar la elektrozeso\lingui{DEFIRS}
\artiklo{elektuario}%
\vorto{elektuario}\fako{medic.} Medikamento koheranta mola, ek substanci diversa (pulvero, pulpi) interligita per siropo o vino\lingui{DEFIRS}
\artiklo{elemento}%
\vorto{elemento}\senco{1}{\textit{(che la antiqui)} Parto konstitucanta di kozo}\senco{2}{\textit{(che la moderni)} Korpo quan onu kredas simpla, kemiale}\senco{3}{En la linguo ordinara : singla ek la kozi di qui la asemblo e la kombino formacas kozo altra}\senco{4}{\textit{(pri cienco)} La principi prima sur qui lu fondesas, e per qui onu komencas lernar lu}\lingui{DEFIRS}
\artiklo{elevaciono}%
\vorto{elevaciono} Desegnuro di la formi extera di ula korpo vidata en direciono horizontala\lingui{EFS}
\artiklo{elevar}%
\vorto{elevar}\fako{trans.} Transportar a nivelo supera. \textit{(anke metafore)}\lingui{EFIRS}
\artiklo{elevatoro}%
\vorto{elevatoro}\senco{1}{\textit{(anat.)} Muskulo qua levas ulo}\senco{2}{\textit{(mekan.)} Aparato uzata por levar la navi en la reparo-doki}\lingui{DEFIRS}
\artiklo{elfo}%
\vorto{elfo}\fako{mitol. Skandinava} Koboldo di atmosfero\lingui{DEFIRS}
\artiklo{eliminar}%
\vorto{eliminar}\fako{trans.) (tekn., matem.} Ekirigar, eskartar\lingui{DEFIS}
\artiklo{elipso}%
\vorto{elipso}\senco{1}{\textit{(gramatiko)} Figuro di retoriko qua konsistas ye supresar un o plura vorti quin la lektero od askoltero devos suplear}\senco{2}{\textit{(geom.)} Kurvo qua rezultas de la seko di kono rekta per plano obliqua ye la axo e ye omna aristi}\lingui{DEFIRS}
\artiklo{elipsoida}%
\vorto{elipsoida} Elipso-forma\lingui{DEF}
\artiklo{elito}%
\vorto{elito} To quo selektesas ek korpo kom lo maxim bona\lingui{DEF}
\artiklo{elitro}%
\vorto{elitro}\fako{zool.} Alo rezistiva qua kovras, che ula insekti, ali plu dina e plu tenua\lingui{DEF}
\artiklo{elixiro}%
\vorto{elixiro} Preparajo farmaciala, qua rezultas de la mixo de ula siropi kun alkoholaji\lingui{DEFIRS}
\artiklo{elizeo}%
\vorto{elizeo}\senco{1}{\textit{(che la antiqui)} Regioni di inferno ube, segun la antiqui, la homi vertuoza, la heroi, sejornis pos lia morto ed ube printempo esis eterna}\senco{2}{\textit{(metaf.)} Loko delicoza}\lingui{DEFIRS}
\artiklo{elongaco}%
\vorto{elongaco} Angulo-disto, mezurata de la centro di Tero, inter Suno e ta o ca planeto
\artiklo{eloquenta}%
\vorto{eloquenta}\senco{1}{Tre talentoza pri parolar}\senco{2}{\textit{(metaf.)} Qua efikas quale parolo eloquenta}\lingui{DEFIS}
\artiklo{eludar}%
\vorto{eludar}\fako{trans.} Evitar habile, rapide e ruze\lingui{EFIS}
\artiklo{elzevir}%
\vorto{« elzevir »}\fako{imprim-arto} Tipo gracioza, derivita de ta tipo karakterizata da linei perpendikla (substitucita, en 1466, da du imprimisti Romana, a la tipi gotika)\lingui{DEFIRS}
\artiklo{-em-}%
\vorto{-em-} Sufixo qua signifikas « qua tendencas a… »
\artiklo{emalio}%
\vorto{emalio}\senco{1}{Verniso qua rezultas de la vitrigo di substanci fuzebla (silico, plomb-oxido, natroxo, potaso) quan onu aplikas, per fuzo, sur la terakotajo, la fayencajo, la metalajo, ed a qua onu adjuntas metal-oxidi pulverigita, por kolorizar lu diverse}\senco{2}{Verniso formacita da la vitrigo ye la surfaco di fayencajo}\lingui{DFIS}
\artiklo{emanar}%
\vorto{emanar}\fako{netrans., de}\senco{1}{\textit{(aludante partikuli nepalpebla)} Eskapar de korpo sen ke la korpo percepteble deskreskas substancale}\senco{2}{Devenar de ulo, de ulu}\senco{3}{\textit{(filoz.)} En la doktrino panteista pri emano, aludante la modi di substanco : ekirar, per libereski intersucedanta, de la substanco una ed universala qua esas Deo}\lingui{DEFIS}
\artiklo{emancipar}%
\vorto{emancipar}\fako{trans.}\senco{1}{\textit{(legaro di la Roma antiqua)} Liberigar de la autoritato patrulala}\senco{2}{\textit{(legaro moderna)} Liberigar (minoro) de la tutelo}\senco{3}{\textit{(en la Mezepoko)} Liberigar (serfo, vasalo) de lua devi a lua sinioro}\senco{4}{\textit{(metafore)} Liberigar de la tutelo di autoritato supera}\lingui{DEFIRS}
\artiklo{embalmar}%
\vorto{embalmar}\fako{trans.} Prezervar (kadavro) de putro, per enpozar substanci aromata, sikigiva ed antisepsiiva\lingui{EF}
\artiklo{embarasar}%
\vorto{embarasar}\fako{trans., per, pro}\senco{1}{Jenar (pri agar, movar)}\senco{2}{Duktar a perplexeso, en qua lu ne savas quon decidor o rezolvor}\lingui{EFIS}
\artiklo{embargar}%
\vorto{embargar}\fako{trans.}\senco{1}{\textit{(yuro marala)} Sorto di sequestro da guvernantaro sur la navi di naciono stranjera kun qua lu militas}\senco{2}{\textit{(yuro civila)} Posedeskar, ye la nomo di la tribunalo, mobli, imobli (di debanto por garantiar a la debario to quo debesas a lu), objekti interdiktita, kontrabando-vari, e c.}\lingui{DEFIRS}
\artiklo{embarkar}%
\vorto{embarkar}\senco{1}{\textit{(trans.)} Pozar en barko, en navo; enirar navo por voyajo sur aquo. \textit{(anke : enirar treno, veturo)}}\senco{2}{\textit{(netrans.)} \textit{(aludante maraquo.)} Penetrar violentoze super la bordo di la navo}\lingui{EFIS}
\artiklo{emberizo}%
\vorto{emberizo}\fako{zool.} Genero de uceli, di qua la speco tipa esas pasero « kornirostra », analoga a la verdrono\latina{emberiza}
\artiklo{emblemo}%
\vorto{emblemo}\senco{1}{Signo videbla, selektita konvencione por repreesentar ideo, kozo abstraktita, ed ulakaze akompanata da subskriburo explika}\senco{2}{Figuro qua memorigas kozo abstraktita per ula aludo}\senco{3}{Atributo di kozo, uzata por reprezentar olca}\lingui{DEFIRS}
\artiklo{embolio}%
\vorto{embolio}\fako{patol.}\senco{1}{Sanga koagulajo fibrinoza, kuntirata da la sango-cirkulado aden le arterii plu mikra, quin lu povas obstruktar}\senco{2}{Sango koagulajo, naskinta en la veini, e kuntirata da la sango-cirkulado ad la kordio}
\artiklo{embolo}%
\vorto{embolo}\fako{patol.} Korpo stranjera qua produktas la fenomeno « embolio »
\artiklo{embotar}%
\vorto{embotar}\fako{trans.) (tekn.}\senco{1}{Fasonar, per martelago, metalo por formacar sur olca la reliefo de puncuro ube la reproduktendo esas kava}\senco{2}{Fasonar, per martelago o saliigilo, metalo-plako, fero-tolo, por rondigar lu a baseno, kasrolo, por igar lu ornivo arkitekturala}\lingui{FS}
\artiklo{embracar}%
\vorto{embracar}\fako{trans.} Prenar e klemar en sua brakii. – \textit{(aludante kozo)}. Cirkondar per sua korpo – quaze kovrar\lingui{EFIS}
\artiklo{embragar}%
\vorto{embragar}\fako{trans.} Igar komunikar la motoro di mekanismo kun la parti di olca quin onu volas movigar, per glitigar kulise rimeno de libera pulio adsur fixa pulio (samaxo e sama-periferia), dento-roto ad piniono, e c.\lingui{F}
\artiklo{embrazuro}%
\vorto{embrazuro}\senco{1}{Vakuo en parieto}\senco{2}{Vakuo en muro, ube ajustesas la klapi di pordo, di fenestro}\senco{3}{Vakuo en la muro di parapeto, qua recevas la boko di kanono}\senco{4}{Aperturo di furnelo tra qua pasas la kolo di retorto}\senco{5}{Vakuo en la masivo di fornego}\lingui{EFR}
\artiklo{embriogenio}%
\vorto{embriogenio} Formaceso, developeso di la embriono che la animali\lingui{DEFIRS}
\artiklo{embriologio}%
\vorto{embriologio} Parto di la natur-historio qua traktas la formaceso e la developeso di la embriono
\artiklo{embriono}%
\vorto{embriono}\fako{fiziol.}\senco{1}{Jermo naskanta, produktita da la fekundigo}\senco{2}{\textit{(metafore)} Intencajo}\senco{3}{Jermo fekundigita, qua komencas developesar en la semino}\lingui{DEFIS}
\artiklo{embrokar}%
\vorto{embrokar}\fako{medic.) (netrans.} Fomentar, per varsar lente liquido grasa sur korpo-parto malada\lingui{DEFIS}
\artiklo{embuskar}%
\vorto{embuskar}\fako{trans.} Postenizar en loko de ube lu povas surprizar la enemiko (od ulu altra)\lingui{EFIS}
\artiklo{emendar}%
\vorto{emendar}\fako{trans.}\senco{1}{Plubonigar per modifikar to quo esas mala, defektoza}\senco{2}{\textit{(yuro)} Reformar verdikto}\lingui{DEFIS}
\artiklo{emerita}%
\vorto{emerita} Qua pensionesas kom retretinta pos sua tota kariero\lingui{DEFIS}
\artiklo{emersar}%
\vorto{emersar}\fako{netrans.} Ekirar de loko ube lu submersesis ed aparar trans la surfaco\lingui{EFIS}
\artiklo{emetiko}%
\vorto{emetiko} Substanco, preparajo por vomigar\lingui{DEFIS}
\artiklo{emetropa}%
\vorto{emetropa}\fako{oftalm.} Dicesas pri okulo homala qua vidas normale. (La okulo emetropa esas, do, nek miopa, nek presbita)\lingui{DEF}
\artiklo{emfazar}%
\vorto{emfazar}\fako{trans.}\senco{1}{Exajerar la tono, la termini, por importigar la kozi}\senco{2}{Aplikar ta figuro di retoriko per qua igesas importoza, termino qua ne esas tala ordinare}\lingui{DEFIRS}
\artiklo{eminenco}%
\vorto{eminenco} Titulo quan onu uzas a la kardinali ed a la chefa sacerdoto di Malta\lingui{DEFIS}
\artiklo{eminenta}%
\vorto{eminenta} Qua situesas aparte pro ta o ca traito di supereso. (propr. e metaf.)\lingui{DEFIRS}
\artiklo{emiro}%
\vorto{emiro}\senco{1}{Che la Arabi, guvernisto di provinco; chefo di tribuo}\senco{2}{Titulo di la viri qui decendas de Mohamed muliero-linee}\lingui{DEFIRS}
\artiklo{emisar}%
\vorto{emisar}\fako{trans.}\senco{1}{\textit{(fiziko)} Produktar adextere}\senco{2}{\textit{(financo)} Livrar por la cirkulado la valoraji (banko-bilieti, monet-uni) quin onu jus kreis}\lingui{DEFIS}
\artiklo{emocar}%
\vorto{emocar}\fako{netrans.} Ekirar de quieteso\lingui{EFIRS}
\artiklo{empalar}%
\vorto{empalar}\fako{trans.} Tormente trapikar per paliso, sinkante ica aden la rektumo\lingui{EFIS}
\artiklo{emperialo}%
\vorto{emperialo} La suprajo di veturo, aranjita por recevar voyajeri\lingui{FIS}
\artiklo{empireo}%
\vorto{empireo}\senco{1}{Segun ula teorii kosmologiala di la antiqui, la sfero ciela maxim granda, maxim alta, sama-centra, ta qua kontenis la fairo eterna, o la astri}\senco{2}{En la teologio Kristana : cielo, kom rezideyo di la beati}\lingui{DEFIRS}
\artiklo{empiriko}%
\vorto{empiriko} Metodo fondita sur la experimentado; metodo neciencala fondita sur experimentado nekompleta\lingui{DEFIRS}
\artiklo{employar}%
\vorto{employar}\fako{trans.} Uzar (sua-profite) ulu per igar lu laborar\lingui{EFS}
\artiklo{emposto}%
\vorto{emposto}\senco{1}{Strato de petro qua kronatrizas fosto o pilastro sur qua su apogas la vulto di arkado}\senco{2}{Parto supra di pordo, di fenestro, quan onu lasas fixa, por diminutar la alteso di la klapi}\senco{3}{Parto vitro-kareloza, fixa, di pordo, di parieto, tra qua kelka lumo-radii eniras chambro obskura}\lingui{DEFIS}
\artiklo{emula}%
\vorto{emula} Qua penas por egalesar o superesar la ceteri en kozi laudinda od en konkurso loyala\lingui{EFIS}
\artiklo{emulsar}%
\vorto{emulsar}\fako{trans.} Adjuntar a drinkajo liquido laktatra quan onu extraktas de la semini ek qui onu povas ekpresar oleo\lingui{DEFIRS}
\artiklo{emulsino}%
\vorto{emulsino} Substanco albuminoida quan onu extraktas de la mandeli
\artiklo{emundar}%
\vorto{emundar}\fako{trans.}\senco{1}{Purigar (arboro) de lua branchi mortinta, de lua planti parazita}\senco{2}{Netigar, beligar o preparar kozo per depreno di omna parti o peci neutila, nebona o nebela}\lingui{FIS}
\artiklo{emuntorio}%
\vorto{emuntorio}\fako{fiziol.} Tubo, organo, per qua evakuesas la humori superflua
\artiklo{en}%
\vorto{en} Qua esas inter limiti, spacale o tempale
\artiklo{enartroso}%
\vorto{enartroso}\fako{anat.} Artiko movebla\lingui{DEF}
\artiklo{encefalito}%
\vorto{encefalito}\fako{patol.} Influro di la encefalo\lingui{DEF}
\artiklo{encefalo}%
\vorto{encefalo}\fako{anat.} La ensemblo ek cerebro e cerebelo\lingui{EFIS}
\artiklo{encikliko}%
\vorto{encikliko}\fako{eklezio katolika} Letro cirkulera da la papo, pri punto di doktrino, di dogmato\lingui{DEF}
\artiklo{enciklopedio}%
\vorto{enciklopedio} Libro qua kontenas expozita, segun ordino metodala, La toto di la savaji homala\lingui{DEFIRS}
\artiklo{-end-}%
\vorto{-end-} Sufixo qua signifikas « quan onu devas », « quan oportas »
\artiklo{endemika}%
\vorto{endemika} Dicesas pri morbo qua esas specala ye regiono, ube lu regnas preske sencese\lingui{DEFIRS}
\artiklo{endemio}%
\vorto{endemio}\fako{patol.} Morbo, di qua ula kazi renkontresas intermite en ca o ta regiono\lingui{DEFIRS}
\artiklo{endermika}%
\vorto{endermika}\senco{1}{\textit{(medic.)} Di qua la fonto esas en la dermo}\senco{2}{\textit{(pri metodo)} Segun qua onu efikas a la dermo per fricioni, locioni, uncioni, e c.}
\artiklo{endetala}%
\vorto{endetala}\fako{pri vendo o kompro di vari} Mikra-quantope\lingui{FD}
\artiklo{endivio}%
\vorto{endivio}\fako{bot.} Cikorio gardenala\latina{cichorium endivia}\lingui{F}
\artiklo{endodermo}%
\vorto{endodermo}\fako{embriol.} Strato interna di la blastodermo qua, pose, divenos la digesto-tubo, e la kovro-epitelio e la glandi anexa di ica
\artiklo{endogena}%
\vorto{endogena}\fako{bot.} Qua kreskas interne di su ipsa
\artiklo{endokardio}%
\vorto{endokardio}\fako{anat.} Membrano qua tapetatre kovras la internajo, la kavaji di la kordio
\artiklo{endokarpo}%
\vorto{endokarpo}\fako{bot.} Un ek la tri strati o membrani ye qui konsistas la envelopilo (perikarpo) di la frukto, olta qua konstitucas la surfaco interna
\artiklo{endoparazito}%
\vorto{endoparazito} Parazito qua habitas interne di sua hosto
\artiklo{endoplasmo}%
\vorto{endoplasmo} Parto centrala od interna di la korpo celula di la enti monoplastida
\artiklo{endospermo}%
\vorto{endospermo} Mikra organo, funcionanta kom kotiledono en la graminei, la liliacei, e c.
\artiklo{eneagono}%
\vorto{eneagono}\fako{geom.} Poligono qua havas non edri e non anguli
\artiklo{enemiko}%
\vorto{enemiko} Homo qua havas sentimenti des-amikala\lingui{EFIS}
\artiklo{energetiko}%
\vorto{energetiko} Teorio fizikala, defensita da Duhem e Wilhelm Ostwald, qua konsideras, en la fenomeni, nur la manifesti mezurebla, abstraktante la naturo propra di la kozi\lingui{DEFIRS}
\artiklo{energio}%
\vorto{energio}\senco{1}{Forco vivaca, agera, di la organismo o di la psiko}\senco{2}{Lo apta efektigar laboro}\lingui{DEFIRS}
\artiklo{enfilar}%
\vorto{enfilar}\fako{trans.) (milit.} Kanonagar enemiki trupa o nava, tale ke li trairesas longesale\lingui{DEFIS}
\artiklo{engajar}%
\vorto{engajar}\fako{trans.} Ligar ulu kom employoto per kontrato\lingui{DEFIRS}
\artiklo{engrosa}%
\vorto{engrosa}\fako{pri vendo o kompro di vari} Granda-quantope
\artiklo{enigmato}%
\vorto{enigmato} Ulo divinenda segun defino-frazi obskura\lingui{EFIS}
\artiklo{enkaustiko}%
\vorto{enkaustiko}\senco{1}{\textit{(che la antiqui)} Farbo facita ek vaxi koloroza, liquidigita per fairo}\senco{2}{Mixuro de vaxo e de terpentino, quan onu uzas por briligar la parqueti, la mobli ligna, e c.}\senco{3}{Induturo vaxo-fundamenta, ye qua onu kovras la statui marmora por igar plu matida lia blankeso}\lingui{DEFIS}
\artiklo{enkilemo}%
\vorto{enkilemo}\fako{citologio} Liquido o suko nukleala en qua balnas la grani kromatika
\artiklo{enklitiko}%
\vorto{enklitiko}\fako{gram.} Vorto qua, en la skribo, unionesas a la vorto preiranta tale ke, aspekte, lu esas quaze una kun ita, e qua ne acentizesas. Kom ex. : L. \textit{que} en \textit{neque}; \textit{ne} en \textit{venisne}; F. \textit{ce} en \textit{est-ce}; \textit{je} en \textit{sais-je}\lingui{DEFIS}
\artiklo{enluminar}%
\vorto{enluminar}\fako{trans.} Piktar ye kolori vivaca, aplikata plata-nuance\lingui{DEFIS}
\artiklo{enorma}%
\vorto{enorma} Qua ecesas omna grandeso o quanteso mezurebla\lingui{DEF}
\artiklo{enoyar}%
\vorto{enoyar}\fako{netrans.} Esar la objekto di ta deskomforto quan sentas la psiko kande nulo interesas, okupas lu\lingui{EFI}
\artiklo{enrejistrar}%
\vorto{enrejistrar}\fako{trans.}\senco{1}{Enskribar o transkribar oficale en registri}\senco{2}{Dicesas anke pri ula aparati qui notas ocili, movi diversa (kom ex. la sfigmografo enrejistras la pulso arteriala)}\lingui{F}
\artiklo{ensemblo}%
\vorto{ensemblo} Asemblajo de personi, de kozi, qui omna formacas toto\lingui{DEFR}
\artiklo{entablamento}%
\vorto{entablamento}\fako{arkitekt.} Saliajo an la suprajo di la muregi di edifico, sur qua su apogas la karpenturo di la tekto\lingui{DEFIRS}
\artiklo{entamar}%
\vorto{entamar}\fako{trans.}\senco{1}{\textit{(pri kozo til lore netushita)} Agar la trancho, frapo, brecho, nocho ye lu, dat-unesme}\senco{2}{\textit{(metaf.)} Komencar (ulo), iniciar, inaugurar}\lingui{F}
\artiklo{enterito}%
\vorto{enterito}\fako{patol.} Influro di la mukozo qua tapete kovras la kanalo intestina\lingui{EF}
\artiklo{enternar}%
\vorto{enternar}\fako{trans.} Koaktar a rezido en ta o ca urbo, regiono, la ekiro esante interdiktita\lingui{DEFIS}
\artiklo{entimemo}%
\vorto{entimemo}\fako{logiko} Silogismo en qua un ek la premisi esas tacita, pro ke lua evidenteso posibligas suplear lu\lingui{DEFIRS}
\artiklo{ento}%
\vorto{ento} To quo existas\lingui{EFI}
\artiklo{entomofilo}%
\vorto{entomofilo}\fako{planto} Di qua la fekundigeso devenas de la transporteso, da la insekti, di la poleno
\artiklo{entomologio}%
\vorto{entomologio} Parto di la zoologio, qua traktas la insekti\lingui{DEFIRS}
\artiklo{entraprezar}%
\vorto{entraprezar}\fako{trans.}\senco{1}{Asumar la direkto, la exekuto (di afero, di tasko)}\senco{2}{Obligar su exekutor ula labori, furnisor ulo, po preco o segun kondicioni determinita}\lingui{FIS}
\artiklo{entratenar}%
\vorto{entratenar}\fako{trans.} Igar provizata ye la kozi qui esas necesa por vivo e por komforto\lingui{F}
\artiklo{entravar}%
\vorto{entravar}\fako{trans.}\senco{1}{Impedar o jenar la marcho di animalo per irga moyeno}\senco{2}{Impedar o jenar per ulo qua haltigas, embarasas}\lingui{FS}
\artiklo{entrechato}%
\vorto{entrechato}\fako{danso} Salto dum qua la pedi intershokas rapide\lingui{DFIR}
\artiklo{entremeso}%
\vorto{entremeso} Disho qua venas inter la rostajo e la desero\lingui{DEFIS}
\artiklo{entrenar}%
\vorto{entrenar}\fako{trans., pri, por, per} Preparar, per exerci e per rejimo specala, ad ago tre energio-konsumanta\lingui{EFS}
\artiklo{entreo}%
\vorto{entreo} Disho qua venas pos la fisho-disho ed ante la rostajo\lingui{DFS}
\artiklo{entresolo}%
\vorto{entresolo}\fako{arkitekt.} Lojeyo inter la ter-etajo e la etajo unesma\lingui{DEFS}
\artiklo{entuziasmar}%
\vorto{entuziasmar}\fako{netrans.}\senco{1}{Sentar sua psiko exaltata}\senco{2}{Ravisesar pro joyo}\lingui{DEFIRS}
\artiklo{enuklear}%
\vorto{enuklear}\fako{trans.) (kirurg.} Extirpar tumoro quan, per presar, onu ekirigas tra incizuro\lingui{EF}
\artiklo{enumerar}%
\vorto{enumerar}\fako{trans.} Enuncar unope omna parti di ensemblo\lingui{EFIS}
\artiklo{enuncar}%
\vorto{enuncar}\fako{trans.} Adexterigar, donante a lu formo determinita, to quon onu havas en sua spirito\lingui{EFIS}
\artiklo{envelopar}%
\vorto{envelopar}\fako{trans.}\senco{1}{Cirkondar ye ulo qua kovras (lu) omnapunte}\senco{2}{\textit{(milit.)} Cirkondar tale ke nula ekireyo posibligas eskapo}\lingui{DEFI}
\artiklo{enverguro}%
\vorto{enverguro}\senco{1}{Larjeso di seglo par-desvolvita}\senco{2}{Spaco quan embracas la ali di ucelo (si aeroplano) kande li esas extensita}\lingui{FS}
\artiklo{envidiar}%
\vorto{envidiar}\fako{trans.}\senco{1}{Dezirar (kozo, quan altru posedas)}\senco{2}{Dezirar (esar en la situeso di ulu)}\lingui{EFIS}
\artiklo{envio}%
\vorto{envio} Fragmenteto qua su separas de la pelo cirkum la ungli\lingui{F}
\artiklo{enzimo}%
\vorto{enzimo} Fermento solvebla
\artiklo{eogena}%
\vorto{eogena}\fako{geol.} Dicesas pri la grupo maxim anciena ek la sulo-strati terciara
\artiklo{eolipilo}%
\vorto{eolipilo}\fako{tekn.} Sfero kava, kun orifico streta, qua – plenigite ye aquo varmigita til bolio – igas olca ekirar per sprico kontinua, forta; se instalite sur quar roti, la shoko da vaporo kontre la aero extera igas la aparato avancar\lingui{DEFIS}
\artiklo{eozimo}%
\vorto{eozimo}\fako{kemio} Materio reda koloriziva (tetrabromofluoreceino)\lingui{DEF}
\artiklo{eozinofila}%
\vorto{eozinofila}\fako{citologio} Leukocito tre grosa, di qua la nukleo esas polimorfa, maxim ofte duopla, mem triopla
\artiklo{epakto}%
\vorto{epakto}\fako{kosmol.} Nombro qua indikas quanta dii esas adjuntenda a la Luno-yaro por ke lu egalesez la Suno-yaro, tale indikanta la evo di Luno kande finas yaro\lingui{DEFIRS}
\artiklo{eparko}%
\vorto{eparko}\senco{1}{\textit{(che la Romani antiqua)} Guvernisto di provinco en lando Greka}\senco{2}{\textit{(lor la imperio Greka)} Prefekto di Konstantinoplo}\senco{3}{\textit{(en Grekia)} Subprefekto}\lingui{EFIRS}
\artiklo{eperlano}%
\vorto{eperlano}\fako{zool.} Mikra marfisho, briloze perlomatrea, di qua la karno esas delikata\latina{Osmerus eperlanus}\lingui{FIL}
\artiklo{epicena}%
\vorto{epicena} Qua indikas sen-distinge la ento maskula, e la ento femina qua korespondas\lingui{DEFIS}
\artiklo{epiciklo}%
\vorto{epiciklo}\fako{astron.} En la astronomio olima : mikra cirklo qua (lore-konjekte) parkuris la periferio di cirklo plu granda, tale justifikante la neregulalaji videbla quin onu deskovris en la movi di la astri\lingui{DEFIRS}
\artiklo{epicikloido}%
\vorto{epicikloido}\fako{geom.} Kurvo produktata da la rotaco di punto selektita sur kurvo movanta qua rulas sen glito cirkum kurvo fixa\lingui{DEFIRS}
\artiklo{epidemio}%
\vorto{epidemio} Morbo qua atakas, en un loko, en la sama instanto, granda quanto de personi, ed aspektas kom dependanta de kauzo qua efikas generale\lingui{DEFIRS}
\artiklo{epidermo}%
\vorto{epidermo}\fako{anat.} Strato surfacala di la pelo, membrano dina e diafana qua kovras la dermo\lingui{DEFIRS}
\artiklo{epifanio}%
\vorto{epifanio}\fako{religio kristana}\senco{1}{Aparo di Yesu a la Reji Maga qui venis adorar lu}\senco{2}{Festo-dio Ekleziala (katolika) en la dio sisesma di januaro, por celebrar ta adoro}\lingui{DEFIRS}
\artiklo{epifita}%
\vorto{epifita} Qua vivas sur altra planti, sen tirar de li sua nutrivi\lingui{DE}
\artiklo{epifizo}%
\vorto{epifizo}\fako{anat.} Saliajo ostatra qua, ligita – che la infanto – a la osto per kartilago, divenas pose apofizo\lingui{DE}
\artiklo{epigastro}%
\vorto{epigastro}\fako{anat.} Parto supera di la abdomino, situita an-super la umbiliko\lingui{DEF}
\artiklo{epigenezo}%
\vorto{epigenezo} Doktrino segun qua onu konceptas la formaceso di la korpi vivanta per la adjunto sucedanta di parti qui ne preexistis en la jermo
\artiklo{epigloto}%
\vorto{epigloto}\fako{anat.} Valveto ye la parto supra di laringo, di qua lu klozas (per sua abaseso) la orifico supra lor la engluto, por impedar la eniro di la nutrivi aden la tubi respirala\lingui{DEF}
\artiklo{epigono}%
\vorto{epigono}\senco{1}{\textit{(mitol.)} Singla de la heroi di la milito duesma di Tebo}\senco{2}{\textit{(metaf.)} La homo qua apartenas a la generaciono dato-duesma}\lingui{DEF}
\artiklo{epigrafio}%
\vorto{epigrafio} La cienco di la epigrafi\lingui{DEFIRS}
\artiklo{epigrafo}%
\vorto{epigrafo}\senco{1}{La skriburo quan onu renkontras o pozas sur la monumenti}\senco{2}{Kurta cito-texto quan onu lokizas kape di libro, di chapitro, por indikar lua tendenco}\lingui{DEFIRS}
\artiklo{epigramo}%
\vorto{epigramo} Mikra versajo qua kontenas espritajo pikanta\lingui{DEFIRS}
\artiklo{epika}%
\vorto{epika} Qua naracas per versi ago heroala\lingui{DEFIRS}
\artiklo{epikarpo}%
\vorto{epikarpo} Shelo extera di la frukto
\artiklo{epilepsio}%
\vorto{epilepsio}\fako{patol.} Morbo di la encefalo, karakterizata da acesi dum qui la malado esvanas subite e kaptesas da konvulsi\lingui{DEFIRS}
\artiklo{epilogo}%
\vorto{epilogo}\senco{1}{\textit{(che la antiqui)} Kurta diskurso recitita da la aktoro lor la fino di la pleo, e per qua lu demandis de la spektinti aplaudi}\senco{2}{Rezumo, konkluzo, ye la fino di poemo, di diskurso, di libro}\lingui{DEFIRS}
\artiklo{episkopo}%
\vorto{episkopo} Alt-oficiero di la Eklezio kristana, chefo e pastoro di diocezo\lingui{DEFIRS}
\artiklo{epistemologio}%
\vorto{epistemologio} Kritikado di la cienci\lingui{DEFIS}
\artiklo{epistolo}%
\vorto{epistolo}\senco{1}{\textit{(literaturo)} Letro versa}\senco{2}{\textit{(liturgio katolika)} Extraktajo de la letri kanonala o de la cetera libri santa, quan onu lektas o kantas lor la meso, ante la evangelio}\lingui{DEFIRS}
\artiklo{epistologismo}%
\vorto{epistologismo}\fako{filoz.} Sorito du-silogisma, rezono dum qua onu prenas kom un ek la premisi lo konkluzita de silogismo preiranta
\artiklo{epitafo}%
\vorto{epitafo} Skriburo sur tombo\lingui{DEFIRS}
\artiklo{epitalamio}%
\vorto{epitalamio} Mikra poemo, kompozita okazione mariajo, honore di la jusa spozi\lingui{DEFIS}
\artiklo{epitelio}%
\vorto{epitelio}\fako{anat.} Quaza epidermo qua kovras la membrani mukoza, seroza, vaskulala e glandala\lingui{DEFIRS}
\artiklo{epiteto}%
\vorto{epiteto}\fako{gram.} Vorto quan onu adjuntas a substantivo por plusaliigar la ideo quan ica expresas\lingui{DEFIRS}
\artiklo{epitomo}%
\vorto{epitomo} Abreviuro de libro\lingui{DEFIRS}
\artiklo{epizodo}%
\vorto{epizodo}\senco{1}{Agado acesora qua ne ligesas rigoroze a la temo}\senco{2}{Fakto acesora qua ligesas plu o min rigoroze ad ensemblo de fakti}\lingui{DEFIRS}
\artiklo{epizootio}%
\vorto{epizootio} Morbo epidemiala, o kontagiala, qua afektas tota klaso de animali domestika en un foyo\lingui{DEFIRS}
\artiklo{epodo}%
\vorto{epodo}\senco{1}{En la poezio Greka e Latina : la dividuro triesma e lasta di la odo, qua venas pos la strofo e la anti-strofo}\senco{2}{Verso duesma (iambo quar-silaba) di distiko di qua la verso unesma esas iambo sis-silaba}\senco{3}{Mikra poemo ek distiki tala}\lingui{DEFIRS}
\artiklo{epoko}%
\vorto{epoko} Tempo-parto qua havas kom distingivo ulo historio-importoza\lingui{DEFIRS}
\artiklo{epoleto}%
\vorto{epoleto} Insigno armeanala, an la shultro, ek bendo di qua la extremajo mi-cirkla garnisesas ye franjo\lingui{DEFIRS}
\artiklo{epruveto}%
\vorto{epruveto} Tubo ek vitro o kristalo, di qua un extremajo esas klozita, qua uzesas por mezurar la volumino, la elastikeso di la gasi, por agar experimenti kemiala diversa, por determinar la koncentreso-grado di ula liquidi, e c.\lingui{FIS}
\artiklo{epuliso}%
\vorto{epuliso}\fako{patol.} Surkreskajo di la jenjivo
\artiklo{equaciono}%
\vorto{equaciono}\fako{algebro} Formulo qua expresas egaleso inter du quanti algebrala, e qua kontenas ordinare un o plura nekonocaji\lingui{EFIS}
\artiklo{equatoro}%
\vorto{equatoro}\senco{1}{Granda cirklo di la Ter-sfero, perpendikla ye la rotac-axo, e qua dividas la globo a du mi-sferi}\senco{2}{Granda cirklo di Suno e di la planeti, perpendikla ye lia rotac-axo}\senco{3}{Granda cirklo qua rezultas de la interseko di la plano di la Ter-equatoro e di la cielo-sfero}\lingui{DEFIRS}
\artiklo{equi-}%
\vorto{equi-} Prefixo uzata exkluzive en la vorti ciencala, qua signifikas « egala »
\artiklo{equilibrar}%
\vorto{equilibrar}\fako{netrans.}\senco{1}{Esar en la situeso di korpo, di punto materiala, qua – atraktate da forci inter-egala ed inter-kontrea – permanas sen movo}\senco{2}{Esar egale repartisita (pri fluido en ta o ca medio)}\senco{3}{\textit{(pri la kozi)} Esar distributita saje}\lingui{EFIRS}
\artiklo{equinoxo}%
\vorto{equinoxo} Singla ek la du periodi di la yaro dum qua (pro ke Suno pasas en la equatoro), la dio-duro egalesas la noktoduro sur la tota Ter-globo\lingui{DEFIS}
\artiklo{equipar}%
\vorto{equipar}\fako{trans.}\senco{1}{Provizar (navo) per omno necesa por la manovro, e por la nutro e la armizo di la homi embarkota}\senco{2}{Provizar (armeani) per la kozi necesa, vesti, armi, e c.}\senco{3}{Provizar (mashino, aparato) per omna kozi necesa por lua funciono}\lingui{DEFIR}
\artiklo{equipolar}%
\vorto{equipolar}\fako{trans.) (matem.} Equilibrigar segun la metodo da Giusto Bellavitis
\artiklo{equiseto}%
\vorto{equiseto}\fako{bot.} Planto kriptogama, di qua la stipo esas rugoza\latina{equisetum}\lingui{EL}
\artiklo{equitato}%
\vorto{equitato} Yusteso naturala, segun la komuna raciono, quan onu opozas a la yusteso legala, a la legaro civila\lingui{EFIS}
\artiklo{equivalar}%
\vorto{equivalar}\fako{trans.} Esar sama-valora (kam)\lingui{DEFIRS}
\artiklo{-er-}%
\vorto{-er-} Sufixo qua signifikas « qua su okupas ofte pri » (ma sen agar lo profesionale)
\artiklo{eratika}%
\vorto{eratika}\fako{geol.} Stranjera en la loko ube onu observas lu, e qua konjekteble adportesis ibe de fore (kom ex. grosa, enorma fragmenti de roki, e c.)\lingui{DEFIS}
\artiklo{erco}%
\vorto{erco} Substanco minerala qua kontenas un o plura metali\lingui{D}
\artiklo{erektar}%
\vorto{erektar}\fako{trans.} Stacigar – per equilibrigar vertikale\lingui{DEFIS}
\artiklo{ergatoplasmo}%
\vorto{ergatoplasmo}\fako{biol.} Parto di la plasmo qua efektigas la laboro celulifala – opoze a la trofoplasmo, di qua la rolo esas nutrala
\artiklo{ergotino}%
\vorto{ergotino}\fako{kemio} Bazo extraktita de la sekal-ergoto, ed uzata medicinale kontre la hemoragii\lingui{DEF}
\artiklo{ergotismo}%
\vorto{ergotismo} Venenago per la sekalo ergotoza\lingui{DEF}
\artiklo{ergoto}%
\vorto{ergoto}\fako{bot.} Produkturo vejetala parazita, venanta forme di mikra pinto, di sporno, sur la spiki di ula graminei, e qua esas venenoza\lingui{EFIS}
\artiklo{-eri-}%
\vorto{-eri-} Sufixo qua signifikas « butiko, butikego en qua onu ofras kom varo… », « establisuro ube onu (agas to o co). »
\artiklo{eriko}%
\vorto{eriko}\fako{bot.} Arbusto de la familio « erikacei », qua kreskas en la suli sabloza\latina{erica}\lingui{DIL}
\artiklo{eringo}%
\vorto{eringo}\fako{bot.} Genero de planto umbelifera, pikantoza\latina{eryngium}
\artiklo{eritroblasto}%
\vorto{eritroblasto}\fako{histol.} Globulo reda prima, kun nukleo ed hemoglobino, quan onu renkontras dum la periodi di regenero di la anemii grava, en la kakexio, e mem quik pos la nasko
\artiklo{eritrocito}%
\vorto{eritrocito}\fako{histol.} Globulo reda nukleoza, di qua la divideso produktas la globuli reda sen-nuklea (hematii)
\artiklo{eritrofila}%
\vorto{eritrofila}\fako{biol.} Qua prizas la koloro reda
\artiklo{erizipelo}%
\vorto{erizipelo}\fako{patol.} Inflamuro surfacala di la pelo, kun tenseso\lingui{EFIS}
\artiklo{ermino}%
\vorto{ermino}\fako{zool.} Martro blanka, mamifero digitigrada di qua la pelo, kovrita subventre da pili blanka tre tenua, esas furo tre chera\latina{putorius ermineus}\lingui{DEFIS}
\artiklo{ermito}%
\vorto{ermito}\fako{religio} Solitaro, retretinta a loko dezerta, por ibe exercar su a pieso\lingui{DEFIRS}
\artiklo{ero}%
\vorto{ero}\fako{kronologio}\senco{1}{Epoko konvencionita, de kande onu komencas kontar la yari di naciono}\senco{2}{Epoko remarkinda, de qua komencas koz-ordino nova}\lingui{DEFIRS}
\artiklo{erodar}%
\vorto{erodar}\fako{trans.} Rodar per efiko korodiva\lingui{DEFIS}
\artiklo{erorar}%
\vorto{erorar}\fako{netrans.}\senco{1}{Forirar de lo exakta}\senco{2}{\textit{(pri doktrino, opiniono)} Ne esar konforma a lo exakta, a lo vera}\lingui{EFIS}
\artiklo{erotika}%
\vorto{erotika} Qua koncernas amoro\lingui{DEFIRS}
\artiklo{erste}%
\vorto{erste} Ne ante (ta o ca dato, ta o ca evento)\lingui{D}
\artiklo{erudar}%
\vorto{erudar}\fako{trans., pri} Exercar pri la cienco di la dokumenti qui koncernas ta o ca ek la savaji homala\lingui{EFIS}
\artiklo{eruptar}%
\vorto{eruptar}\fako{netrans.) (pri kozo} Ekirar, desintrikar su bruske de to quo kontenas lu\lingui{DEFIS}
\artiklo{ervo}%
\vorto{ervo}\fako{bot.} Genero de leguminosi, di qua la speco tipa esas la lentilio\latina{ervum}\lingui{DEFISL}
\artiklo{-es-}%
\vorto{-es-} Sufixo qua indikas la qualeso, la estado, la stando
\artiklo{esamo}%
\vorto{esamo}\senco{1}{Grupo de abeli, de vespi, qui vivas socie}\senco{2}{Grupo multa-membro de insekti. – \textit{(anke metafore)}}\lingui{FL}
\artiklo{esar}%
\vorto{esar}\fako{netrans.} Verbo qua ligas la atributo a la subjekto di la propoziciono, afirmante la qualeso di olca\lingui{EFIRS}
\artiklo{esayo}%
\vorto{esayo} Libro en qua la autoro traktas surfacale temo, sen vizir a traktar lu profunde\lingui{DEFS}
\artiklo{esbosar}%
\vorto{esbosar}\fako{trans.} Komencar laboro, verko, indikante lua formi, kolori, e c., sen par-facar ula ek li\lingui{FI}
\artiklo{esenco}%
\vorto{esenco}\senco{1}{\textit{(filoz. e teol.)} Quo konstitucas la fundo di la ento}\senco{2}{Quo konstitucas la naturo propra di kozo}\senco{3}{Quo konstitucas la parto maxim importoza di kozo}\senco{4}{Quo kontenas la efiko partikulara ye substanco}\lingui{DEFIRS}
\artiklo{esforcar}%
\vorto{esforcar}\fako{netrans.} Uzar sua vigoro por rezistar o por vinkar rezistanto\lingui{FIS}
\artiklo{eshafodo}%
\vorto{eshafodo}\senco{1}{Konstrukturo qua subtenas la gradi por la spektanti}\senco{2}{Konstrukturo sur placo publika, ube onu exekutas la morto-verdikti}\senco{3}{Konstrukturo karpentura sur qua laboristi okupesas da la konstrukto, la reparo o la dekoro di la edifici}\lingui{DFIR}
\artiklo{eshelono}%
\vorto{eshelono}\fako{armeo} Personi, kozi, dispozita ye intervali\lingui{F}
\artiklo{-esk-}%
\vorto{-esk-} Sufixo qua indikas « divenar », « komencar… »
\artiklo{eskadro}%
\vorto{eskadro}\fako{navaro} Asemblajo de milito-navi komande da admiralo\lingui{DEFIRS}
\artiklo{eskadrono}%
\vorto{eskadrono}\fako{armeo} Trupo de armeani kavaliera; dividuro de regimento di la kavalrio, komandata da oficiro supra-grada\lingui{DEFIRS}
\artiklo{eskalar}%
\vorto{eskalar}\fako{netrans.) (pri navo} Haltar en portuo, dum voyajo, por embarkar o desembarkar pasajanti o vari\lingui{FIS}
\artiklo{eskaldar}%
\vorto{eskaldar}\fako{trans.}\senco{1}{Brular per aquo varmega}\senco{2}{Imersar en aquo bolianta}\lingui{EFS}
\artiklo{eskalero}%
\vorto{eskalero}\fako{en edifico} Serio de gradi sequanta, fixigita, uzata por acensar o decensar de etajo ad altra\lingui{FIS}
\artiklo{eskamotar}%
\vorto{eskamotar}\fako{trans.} Habile desaparigar objekto, per subtila procedo qua eskapas la regardo de la spektanto\lingui{DFS}
\artiklo{eskapar}%
\vorto{eskapar}\fako{trans.}\senco{1}{Tirar su de to, en quo onu retenesas}\senco{2}{Cesar retenesar}\lingui{EFIS}
\artiklo{eskarbilo}%
\vorto{eskarbilo} Fragmento de terkarbono, ne par-kombustita e mixita a cindri\lingui{FS}
\artiklo{eskaro}%
\vorto{eskaro}\fako{medic.} Krusto nigratra, quan formacas la desorganizeso di la tisui en parto atakata da gangreno o da la efiko di korodivo\lingui{EFIS}
\artiklo{eskarolo}%
\vorto{eskarolo}\fako{bot.} Varietato di cikorio kultivata, quan onu manjas kom salado\latina{lactuca scariola}\lingui{DFIS}
\artiklo{eskarpa}%
\vorto{eskarpa} Di qua la pento esas preske vertikala\lingui{EFIS}
\artiklo{eskartar}%
\vorto{eskartar}\fako{trans.}\senco{1}{Pozar (la parti di kozo) tale ke li kelke interdistas}\senco{2}{Pozar kelka-diste de kozo, de persono}\senco{3}{Forigar de la direciono quan onu devas sequar}\lingui{FIS}
\artiklo{eskombro}%
\vorto{eskombro} To quo restas de objekto ruptita\lingui{FIS}
\artiklo{eskopeto}%
\vorto{eskopeto} Speco di funel-fusilo, quan onu *weris bandolieratre\lingui{FS}
\artiklo{eskorto}%
\vorto{eskorto} Trupo qua akompanas persono de loko ad altra, por surveyar lua sekureso o por honorizar lu per lo; kaptiton, por impedar lua eskapo; konvoyon, por protektar lu\lingui{DEFIRS}
\artiklo{eskoto}%
\vorto{eskoto}\fako{navo} Grosa kordego manuo-movebla, fixigita an la angulo di basa segli por desvolvar od extensar li\lingui{DFS}
\artiklo{eskrokar}%
\vorto{eskrokar}\fako{trans.} Furtar (ulo de ulu) per dupigar (lu)\lingui{FI}
\artiklo{eskudelo}%
\vorto{eskudelo} Vazo kava, ligna, terakota o metala, e c., en qua onu servas buliono, supo, por manjar li\lingui{FIS}
\artiklo{eskupar}%
\vorto{eskupar}\fako{trans.} Vakuigar per ligna shovelo kava (la aquo qua eniras la navi)\lingui{EF}
\artiklo{-esm-}%
\vorto{-esm-} Sufixo qua signifikas la numero
\artiklo{esmerilio}%
\vorto{esmerilio}\fako{zool.} Speco di falkono, remarkinda pro la mi-kreso di lua staturo, la lejereso di lua flugo e la rapideso di lua movi\latina{falco subbuteo}\lingui{FIS}
\artiklo{esotera}%
\vorto{esotera}\fako{filoz.} Qua esas la objekto di docado partikulara, intima\lingui{DEFIRS}
\artiklo{espado}%
\vorto{espado} Atak-armo, ek lamo akuta ek stalo, muntita aden mancho konkoza, e quan onu *weras an la flanko, en gaino\lingui{FIS}
\artiklo{espalero}%
\vorto{espalero} Muro garnisita maxim-multa-kaze ye lato-greto alonge qua onu plantacas frukt-arbori di qui la branchi esas aplikita a lu\lingui{DEFRS}
\artiklo{esparo}%
\vorto{esparo}\fako{navo} Stango, min forta kam la masteto, qua suportas la paviliono, ek qua onu facas la stangi per qui onu impedas la altra navi venar kolizionar od abordar atakoze\lingui{DEF}
\artiklo{esparseto}%
\vorto{esparseto}\fako{bot.} Planto foreja, analoga a la luzerno, di qua la floro esas granda, purpuratra e flavatra, e dispozita spike od axo-grape\latina{onobrychis sativa}\lingui{DEF}
\artiklo{esparto}%
\vorto{esparto}\fako{bot.} Gramineo qua kreskas en la nordo-parto di Afrika, di qua la stipeti uzesas por facar objekti diversa (chapelo-frami, e c.), papero, e c.\latina{stipa tenacissima}\lingui{DEFIS}
\artiklo{espelar}%
\vorto{espelar}\fako{trans.} Lektar la vorti literope, e formacante la silabi\lingui{EF}
\artiklo{esperar}%
\vorto{esperar}\fako{trans.} Konsiderar (to quon onu deziras) kom realeskonta\lingui{FIS}
\artiklo{espineto}%
\vorto{espineto}\fako{muziko}\senco{1}{Muzik-instrumento portebla, kordoza, quan onu pleis per plumo-pinto}\senco{2}{Sorto di klavikordo, di qua la kordi agesis per pinto di korvo-plumo}\lingui{DEFIRS}
\artiklo{espino}%
\vorto{espino}\fako{bot.} Arboreto dornoza\lingui{FIS}
\artiklo{esplanado}%
\vorto{esplanado}\fako{armeo}\senco{1}{Tereno plata e nuda qua, en la fortifikaji, esas situita de la rempari a la domi di la urbo}\senco{2}{Sulo-peco nuda, qua esas situita exter e portee di fortreso}\lingui{DEFIS}
\artiklo{espreso}%
\vorto{espreso}\fako{treno} Qua iras rapide a sua destino-staciono, haltante intervale nur ye kelka stacioni\lingui{DEFIRS}
\artiklo{esprito}%
\vorto{esprito} Vivaceso pikanta di la spirito\lingui{DF}
\artiklo{espruva}%
\vorto{espruva} Qua ne lasas … efikar a su\lingui{EF}
\artiklo{esquado}%
\vorto{esquado}\senco{1}{\textit{(armeo)} Parto di kompanio de infantriani o de kavalriani}\senco{2}{Trupo de laboristi qui okupesas da tasko specala}\lingui{EFIS}
\artiklo{establisar}%
\vorto{establisar}\fako{trans.}\senco{1}{Fondar kom duronta multa-tempe}\senco{2}{Institucar}\lingui{DEFIRS}
\artiklo{estado}%
\vorto{estado} Eso mentala, anmala, psikala, od aferala di la individuo koncernata\subvorto{}{estado civila}{La situeso di persono kom filio legitima o ne-legitima, mariajita o celiba}{}{}
\artiklo{estafeto}%
\vorto{estafeto} Kuriero qua portas la depeshi de ca a ta posto-kontoro\lingui{DEFIRS}
\artiklo{estakado}%
\vorto{estakado} Asemblajo de palisi, interligita da kateni, da ferstangi, e c., en la enireyo di portuo, di kanalo, por barar lu\lingui{FIS}
\artiklo{estalar}%
\vorto{estalar}\fako{trans.}\senco{1}{Expozar, sur surfaco plu o min granda, vari vend-ofrata}\senco{2}{Extensar komplete sur surfaco}\senco{3}{\textit{(metaf.)} Expozar komplete kom regardenda}\lingui{F}
\artiklo{estampar}%
\vorto{estampar}\fako{trans.} Markizar (metalo, ledro, petro-kartono) per la traco (reliefe o kave) di mulduro, di modelo, di matrico, di qua lu subisas la preso\lingui{DEFIRS}
\artiklo{estanchar}%
\vorto{estanchar}\fako{trans.} Haltigar la fluo di liquido\lingui{EFS}
\artiklo{estayo}%
\vorto{estayo}\fako{navo} Grosa kordego, tensita de la parto avana di la navo til la kapo di omna masto, por fermigar ica avane, quale la vanti fermigas lu dope\lingui{EFS}
\artiklo{estetiko}%
\vorto{estetiko} Parto di filozofio, qua traktas lo bela en naturo ed en arto\lingui{DEFIRS}
\artiklo{estimar}%
\vorto{estimar}\fako{trans.} Opinionar favoroze pri la valoro (etikala, morala, karakterala) di persono\lingui{DEFIS}
\artiklo{estivar}%
\vorto{estivar}\fako{netrans.) (pri la bestio-trupi} Vivar dum somero en la pastureyi di la monti\lingui{DEFIRS}
\artiklo{esto}%
\vorto{esto} Un ek la quar punti kardinala, olta qua esas ube Suno su levas\lingui{DEFS}
\artiklo{estompar}%
\vorto{estompar}\fako{trans.}\senco{1}{Dilutar la traci de krayono e pastelo, sur desegnuro, per rolero fela o papera, di qua la extremaji esas ordinare pinta}\senco{2}{Kovrar ye nuanco subtila}\lingui{DEF}
\artiklo{-estr-}%
\vorto{-estr-} Sufixo qua signifikas « chefo di », « mastro di »
\artiklo{estrado}%
\vorto{estrado} Parto elevita super la planko-sulo di chambro, di chambrego, di salono, sur qua onu lokizas lito, trono, katedro, e c.\lingui{DEFRS}
\artiklo{estragono}%
\vorto{estragono} Planto herba, aromata, uzata kom kondimento\latina{artemisia dracuncules}\lingui{DEFIRS}
\artiklo{estribo}%
\vorto{estribo} Quaza triangulo de fero, pinta-baza, qua, suspendate an singla latero di la selo per rimeno, subtenas la pedo di la kavalkanto\lingui{FS}
\artiklo{estuario}%
\vorto{estuario}\fako{geogr.} Larja ekflueyo di fluvio, gulfatra\lingui{DEFIS}
\artiklo{estublo}%
\vorto{estublo} Parto di la stipo qua restas stacanta kande onu falchas la frumento, la sekalo\lingui{DEFI}
\artiklo{esturjono}%
\vorto{esturjono}\fako{zool.} Grosa fisho qua iras de la mari a la granda fluvii, e di qua la ovi uzesas por facar kaviaro\latina{acipenser Huso}\lingui{EFIS}
\artiklo{esvanar}%
\vorto{esvanar}\fako{netrans.} Ne plus konciar sua vivo\lingui{FI}
\artiklo{etajero}%
\vorto{etajero} Moblo ek montanti qui portas tabuleti horizontala, dispozita etaje, sur qui onu pozas libri, art-objekti, e c.\lingui{DFR}
\artiklo{etajo}%
\vorto{etajo}\fako{arkitekt.} En domo, en edifico, konsistanta ye plura apartmenti superpozita : omna parto sama-nivela\lingui{DEFR}
\artiklo{etalono}%
\vorto{etalono}\senco{1}{Tipo legale di la mezurili e ponderili}\senco{2}{\textit{(financo)} Metalo selektita por la fabriko di la moneto-peco qua uzesas kom tipa unajo di omna moneto-peci di la lando koncernata}\lingui{F}
\artiklo{etamino}%
\vorto{etamino} Stofo ek lano lejera. – Lanajo, krinajo, quan onu uzas por blutar\lingui{DEFIRS}
\artiklo{etapar}%
\vorto{etapar}\fako{netrans.) (pri armeo} Haltar en loko ube la trupi marchanta repozos dum nokto\lingui{DFIRS}
\artiklo{etato}%
\vorto{etato} To quo konstatas la estado di la kozi en ta o ca instanto\lingui{DFIS}
\artiklo{eterna}%
\vorto{eterna}\senco{1}{Qua nek komencis nek finos}\senco{2}{Qua ne finos}\senco{3}{Di qua la fino-instanto ne esas previdebla}\lingui{EFIS}
\artiklo{etero}%
\vorto{etero}\senco{1}{La parto maxim subtila di atmosfero, ye qua (krede da la antiqui) konsistas fairo}\senco{2}{\textit{(fiziko)} Fluido neponderebla, di qua ula fizikisti konjektis la existo por explikar la fenomeni « foto » e « kaloro »}\senco{3}{\textit{(kemio)} Liquido volatila qua produktas la alkoholi kande onu deprenas de li equivalanto ek aquo}\lingui{DEFIRS}
\artiklo{etiketo}%
\vorto{etiketo} Indiko-texto fixigita sur objekto, por savigar la naturo, la kontenajo, la preco, la destineso, e c. di ica
\artiklo{etiko}%
\vorto{etiko}\senco{1}{La cienco pri la mori}\senco{2}{Traktato pri ta cienco}\lingui{DEFIRS}
\artiklo{etimologio}%
\vorto{etimologio} Decendo de vorto relate un o plura vorti altra de qui lu derivesas\lingui{DEFIRS}
\artiklo{etiolar}%
\vorto{etiolar}\fako{trans.}\senco{1}{Igar (planto) febla e senkolora per kreskigar lu en loko obskura}\senco{2}{Igar (persono) febla e pala, per privacar lu de aero libera, de korp-exercado}\lingui{EFL}
\artiklo{etiologio}%
\vorto{etiologio}\senco{1}{\textit{(filoz.)} Exploro di la kauzi}\senco{2}{\textit{(medic.)} La parto di medicino, en qua onu exploras la morbi}\lingui{DEFIRS}
\artiklo{etmoido}%
\vorto{etmoido}\fako{anat.} \textit{(osto kraniala)} di qua la lamo supra esas truoza, inkastrita en la sesguro infra di la osto frontala, e qua kontributas a formacar la naztrui\lingui{DEF}
\artiklo{etnografio}%
\vorto{etnografio} Deskripto di la populi\lingui{DEFIRS}
\artiklo{etnologio}%
\vorto{etnologio} Cienco di la origino e di la deveno di la populi\lingui{DEFIRS}
\artiklo{etologio}%
\vorto{etologio} Cienco di la mori. – Traktato pri la mori
\artiklo{etuyo}%
\vorto{etuyo} Sorto di buxo adaptita a la formo di la objekto quan lu kontenos\lingui{DFIS}
\artiklo{eudiometrio}%
\vorto{eudiometrio} Analizo di aero, di mixuri de gasi, per eudiometro\lingui{DEFIRS}
\artiklo{eudiometro}%
\vorto{eudiometro} Instrumento uzata por determinar la proporcion di la parti qui kompozas aero od irgaltra mixuro de gasi\lingui{DEFIRS}
\artiklo{eufemismo}%
\vorto{eufemismo} Figuro de vorti, qua konsistas ye diminutar – per la termino o la frazoformo quin onu uzas – to quo esas shokiva en la vorto propra\lingui{DEFIS}
\artiklo{eufonio}%
\vorto{eufonio} Sorgo por evitar la soni orelo-harda en la pronunco\lingui{DEFIS}
\artiklo{euforbio}%
\vorto{euforbio}\fako{bot.} Planto di qua la suko esas laktatra, akra e korodiva\latina{euphorbia}
\artiklo{eukalipto}%
\vorto{eukalipto}\fako{bot.} Arboro balzamoza, de la familio « mirtacei », deveninta de Oceania, di qua la ligno esas veinoza, e la folii uzata kom astriktivo e kontre-febro\latina{eucalyptus}\lingui{DEFIS}
\artiklo{eukaristio}%
\vorto{eukaristio}\fako{religio kristana} Sakramento di la religio katolika, qua kontenas substancale la korpo, la sango, la anmo e la deeso di Iesu-Kristo, en la formo di pano e di vino, por servar a la fideli kom nutrivo anmala\lingui{DEFIRS}
\artiklo{eunuko}%
\vorto{eunuko}\senco{1}{Viro kastrita, di qua la ofico, en Oriento, konsistis ye surveyar la mulieri en la haremi}\senco{2}{Viro igita nekoitiva}\lingui{DEFIRS}
\artiklo{evakuar}%
\vorto{evakuar}\fako{trans.}\senco{1}{\textit{(medic.)} Ekpulsar (materii akumulita en parto di la korpo)}\senco{2}{Forsendar amase (la personi qui okupas loko)}\lingui{DEFIS}
\artiklo{evaluar}%
\vorto{evaluar}\fako{trans., ye} Determinar aproxime la valoro o la quanto, per kalkulo o konjekto\lingui{EFIS}
\artiklo{evangelio}%
\vorto{evangelio}\fako{religio kristana}\senco{1}{Doktrino e lego da Iesu-Kristo}\senco{2}{Singla de la libri (4) ube raportesas la biografio di Iesu-Kristo, lua doktrino e lua lego}\lingui{DEFIRS}
\artiklo{evar}%
\vorto{evar}\fako{trans.} Existar de (ta o ca quanto de tempo)\lingui{L}
\artiklo{eventar}%
\vorto{eventar}\fako{netrans.} Divenar fakto\lingui{EFI}
\artiklo{eventuala}%
\vorto{eventuala} Qua povas eventar, realeskar, se ula cirkonstanci existas\lingui{DEFIRS}
\artiklo{evidenta}%
\vorto{evidenta} Qua posedas ta klareso quan la exaktaji prizentas a la spirito\lingui{DEFIS}
\artiklo{eviero}%
\vorto{eviero} Speco di tablo kava, o kuvo ek petro o ceramikajo, sur qua onu lavas (generale en koqueyo) e qua havas truo por la ekfluo di la aquo reziduo\lingui{F}
\artiklo{eviktar}%
\vorto{eviktar}\fako{trans.) (yuro-cienco} Expropriar legale, kun la yuro a rekurso kontre la vendinto, ta qua aquiris sincere to quon onu ne havis la yuro vendar a lu\lingui{F}
\artiklo{evitar}%
\vorto{evitar}\fako{trans.} Probar ne-renkontror (ulu) o ne atingesor (da ulo)\lingui{EFIS}
\artiklo{evolucionar}%
\vorto{evolucionar}\fako{netrans.}\senco{1}{Manovrar turnante cirkum su}\senco{2}{Parkurar la ciklo di sua developeso}\lingui{DEFIRS}
\artiklo{ex-}%
\vorto{ex-} Prefixo qua signifikas « qua ne plus esas »
\artiklo{exajerar}%
\vorto{exajerar}\fako{trans.} Igar ulo transirar la limiti qui konvenas, parole, pense, age\lingui{EFIS}
\artiklo{exakta}%
\vorto{exakta} Rigoroze konforma a la fakti\lingui{DEFIRS}
\artiklo{exaltar}%
\vorto{exaltar}\fako{trans.}\senco{1}{Laudar (la merito di ulu) til alta grado di glorio}\senco{2}{Developar (sentimenti di la psiko) ad alta grado di intenseso}\lingui{DEFIS}
\artiklo{examenar}%
\vorto{examenar}\fako{trans.} Submisar (kandidato) ad experimenti por saveskar lua valoro\lingui{DEFIRS}
\artiklo{exantemo}%
\vorto{exantemo} Erupto pelala\lingui{DEFIS}
\artiklo{exarko}%
\vorto{exarko}\fako{epoki antiqua}\senco{1}{Reprezentero di la imperiestro di Oriento en provinco fora}\senco{2}{\textit{(nun)} Reprezentero di la patriarko di Konstantinoplo (Istanbul) komisita por vizitar la provinci}\lingui{DEFIRS}
\artiklo{exasperar}%
\vorto{exasperar}\fako{trans.}\senco{1}{\textit{(medic.)} Igar (morbo) atingar sua kolmo}\senco{2}{Igar (ulu) atingar la kolmo di la iriteso psikala}\senco{3}{Bitrigar extreme la humoro per irito troa}\lingui{EF}
\artiklo{exaucar}%
\vorto{exaucar}\fako{trans.}\senco{1}{Satisfacar (ulu) per grantar a lu to quo esas la objekto di lua deziri, di lua pregi}\senco{2}{Efektigar (la deziro, la prego, difAu)}\lingui{F}
\artiklo{exegezo}%
\vorto{exegezo} Interpreto di texto (historiala, yurala, e c.)\lingui{DEFIS}
\artiklo{exekutar}%
\vorto{exekutar}\fako{trans.} Efektigar to quo konsequas (de rezolvajo, de projetajo, de verdikto)\lingui{DEFIRS}
\artiklo{exemplero}%
\vorto{exemplero}\senco{1}{Singla de la objekti formacita per tipo unika quan onu *reproduktis}\senco{2}{Singla de la specimeni di un speco o di un varietato en kolektajo zoologiala, botanikala, mineralogiala}\lingui{DEFIRS}
\artiklo{exemplo}%
\vorto{exemplo}\senco{1}{Maniero esar, agar, di ulu, konsiderata kom imitinda}\senco{2}{La persono qua esas, agas, tale}\senco{3}{To quon subisas ulu (desfortuno, puniseso, e c.), egardata kom apta donar leciono}\senco{4}{La persono qua subisas ta desfortuno, ta puniso}\senco{5}{To quo eventis, konsiderata kom komparo-termino pri to simila quo povas eventar}\senco{6}{Fakto quan onu mencionas sekonde di aserto}\senco{7}{Extraktajo di texto quan onu citas, sekonde di expliko gramatikala}\lingui{DEFIRS}
\artiklo{exequatur}%
\vorto{« exequatur »}\senco{1}{\textit{(yuro-cienco)} Permiso legala necesa por ke verdikto da arbitri, da tribunali, stranjera divenez exekutebla}\senco{2}{Permiso quan bezonas aganto stranjera por exekutar en la lando la misiono quan konfidis a lu lua guvernistaro}
\artiklo{exercar}%
\vorto{exercar}\fako{trans.} Formacar, fasonar (ulu) por igar lu praktikar\lingui{DEFIRS}
\artiklo{exergo}%
\vorto{exergo}\fako{numismatiko} Spaco rezervita maxim-multa-kaze en la infra parto di medalio, por recevar surskriburo (moto, devizo, dato, e c.)\lingui{DEFIS}
\artiklo{exfoliar}%
\vorto{exfoliar}\fako{trans.} Deprenar foliope, lamelope (ula parti di substanco)\lingui{DEFIS}
\artiklo{exhalar}%
\vorto{exhalar}\fako{trans.}\senco{1}{Emisar (gaso, vaporo, odoro)}\senco{2}{\textit{(pri animali, vejetanti)} Lasar eskapar tra la tisui ula elementi eliminata de la organismo}\senco{3}{\textit{(metaf.)} Lasar eskapar (sentimenti)}\lingui{EFIS}
\artiklo{exhaustar}%
\vorto{exhaustar}\fako{trans.}\senco{1}{Sikigar per par-cherpo}\senco{2}{Vendar la lasta exemplero (di libro, e c.)}\senco{3}{Igar febla komplete}\senco{4}{\textit{(logiko)} Enumerar omna kazi, omna hipotezi posibla di temo}\senco{5}{\textit{(matem.)} Uzar la procedo qua posibligas neglijar la difero inter du *grandori, per pruvar ke la difero esas infinite mikra}\lingui{EFI}
\artiklo{exhortar}%
\vorto{exhortar}\fako{trans., pri} Probar igar (ulu) volar ulo, per diskursi persuadiva\lingui{EFIS}
\artiklo{exilar}%
\vorto{exilar}\fako{trans.} Koaktar (ulu) a livar lua patrio, interdiktante a lu retrovenor en olu\lingui{DEFI}
\artiklo{exino}%
\vorto{exino}\fako{biol.} Strato extera, transformita a kutino (C₆H₁₀O) e neegale dikeskinta (pinti e kresti, pori e plisuri), di poleno-grano
\artiklo{existar}%
\vorto{existar}\fako{netrans.} Apartenar a la nuna realaji\lingui{DEFIS}
\artiklo{exkavar}%
\vorto{exkavar}\fako{trans.} Facar, kavigante\lingui{EFIS}
\artiklo{exkluzar}%
\vorto{exkluzar}\fako{trans.}\senco{1}{Pozar (ulu) exter kozo kom ne esante plu partoprenonta olu}\senco{2}{Retropulsar (kozo) kom nekonciliebla kun altro}\lingui{DEFIRS}
\artiklo{exkomunikar}%
\vorto{exkomunikar}\fako{trans.) (religio kristana} Forjetar (ulu) de la komuniono (katolika)\lingui{DEFIRS}
\artiklo{exkoriar}%
\vorto{exkoriar}\fako{trans.) (medic.} Skrachar surfacale, entamante nur la pelo\lingui{EFIS}
\artiklo{exkremento}%
\vorto{exkremento}\fako{fiziol.} Materio liquida o solida (sudoro, muko nazala, urino, feko) quan ula organi ekpulsas de la korpo\lingui{DEFIRS}
\artiklo{exkursar}%
\vorto{exkursar}\fako{netrans.} Cirkular explorante ula quanto de la surfaco, de la areo di lando\lingui{DEFIRS}
\artiklo{exkuzar}%
\vorto{exkuzar}\fako{trans.} Alegar, favore di ulu, motivi por mingravigar to quon onu reprochas a lu, o por justifikar lu pri lo\lingui{DEFIRS}
\artiklo{ex libris}%
\vorto{« ex libris »} Enskriburo da ulu sur la libri di sua biblioteko, por indikar ke li apartenas a lu; ta enskriburo, maxim-multa-kaze, agesas per la glutino di desegnuro artala sur papero-folio, kun indiko di la nomo di la proprietanto
\artiklo{exodo}%
\vorto{exodo}\senco{1}{\textit{(pri la Hebrei antiqua)} Libro duesma di Biblo, qua kontenas la historio di la ekiro de Egiptia}\senco{2}{\textit{(metaf.)} Ekmigro grand-amase de populo}\lingui{DEFIS}
\artiklo{exogena}%
\vorto{exogena}\fako{bot.} Di qua la kresko eventas adexter su ipsa, per la strati extera
\artiklo{exorcisar}%
\vorto{exorcisar}\fako{trans.) (teol. kristana} Agar a la demoni (por expulsar li de la korpo di demonozi), od a la demonozo, la pregi, la ceremonii di la eklezio katolika\lingui{DEFIS}
\artiklo{exordiar}%
\vorto{exordiar}\fako{trans.) (retor.} Pri diskurso : komencar, preparante la atenco e la afableso-sentimenti\lingui{DEFIS}
\artiklo{exostoso}%
\vorto{exostoso}\senco{1}{\textit{(medic.)} Surkreskajo morbala sur la surfaco od interne di osto}\senco{2}{\textit{(bot.)} Surkreskajo sur la trunko, la branchi di ula arbori}\lingui{DEF}
\artiklo{exotera}%
\vorto{exotera}\senco{1}{\textit{(filoz.)} Dicesis pri la dicipuli qui vartis sua inicieso pri ula doktrini neprofunda – pri ula libri kompozita tale ke lia kontenajo esis komprenebla da la plebeyo}\senco{2}{\textit{(okultismo)} Dicesas pri la docajo quan onu darfas divulgar ad omna homi}\lingui{DEFIRS}
\artiklo{exotika}%
\vorto{exotika}\senco{1}{\textit{(natur-historio)} Qua ne esas produkturo naturala di nia klimati}\senco{2}{\textit{(metaf.)} Qua ne apartenas a la lando ipsa, nek a la landi relative vicina (pri la mori, la kustumi, e c.)}\lingui{DEFIS}
\artiklo{expansar}%
\vorto{expansar}\fako{netrans.}\senco{1}{\textit{(cienco)} Agar ta movo per qua gaso, vaporo dilatesas; agar ta movo per qua parto organa developesas, divenas plu voluminoza, o transformesas ad altra parto organa}\senco{2}{\textit{(metaf.)} Agar ta movo per qua la kordio, vice konstriktar sua sentimenti, difuzas li adextere. -Movo per qua ula doktrini, ula idei, kreskas ed atraktas adheranti kontinue plu multa}\lingui{DEFIS}
\artiklo{expediar}%
\vorto{expediar}\fako{trans.} Sendar ulo (ordinare vari) a destinario\lingui{DEFIRS}
\artiklo{expedicionar}%
\vorto{expedicionar}\fako{trans.} Sendar homi, komisar li, por militar o por explorar regiono\lingui{DEFIRS}
\artiklo{expektar}%
\vorto{expektar}\fako{trans.} Mente anticipar ula evento\lingui{EFIS}
\artiklo{expektorar}%
\vorto{expektorar}\fako{trans.} Ekjetar de la tubi respirala, per la boko\lingui{DEFIS}
\artiklo{experiencar}%
\vorto{experiencar}\fako{trans.} Konocar (ulo), savar (pri ulo) per uzir lu praktikale\lingui{EFIS}
\artiklo{experimentar}%
\vorto{experimentar}\fako{trans.} E \textit{(netrans., pri, per)}. En la observo-cienci : operaco quan onu agas por verifikar o por demonstrar ulo per praktiko\lingui{DEFIRS}
\artiklo{experta}%
\vorto{experta} Qua savas tre multe (pri ulo) per praktikado\lingui{DEFIRS}
\artiklo{expertizar}%
\vorto{expertizar}\fako{trans.} Submisar a la evaluo da experti. \textit{(pri experti)} Evaluar la preco, la qualesi di (ulo)\lingui{DEFIRS}
\artiklo{expiacar}%
\vorto{expiacar}\fako{trans.}\senco{1}{Lavar per ceremonio religiala, per sakrifiko (krimino pro qua onu esas despurigita.)}\senco{2}{Reparar (peko, krimino) per puniseso}\lingui{DEFIS}
\artiklo{expirar}%
\vorto{expirar}\fako{netrans.}\senco{1}{Ekpulsar de la pulmoni (la aero qua uzesis por respiro)}\senco{2}{\textit{(pri vejetanti)} Restitucar ad aero (pos prenir de lu ula elementi gasa) ti ek la elementi qui esas eliminenda de la organismo}\senco{3}{\textit{(metaf.)} Cesar existar o cesar validesar}\lingui{DEFIS}
\artiklo{explicita}%
\vorto{explicita} Par-enuncita\lingui{DEFIS}
\artiklo{explikar}%
\vorto{explikar}\fako{trans.}\senco{1}{Klarigar per developo (la senco di ulo)}\senco{2}{Justifikar fakto}\lingui{DEFIS}
\artiklo{explorar}%
\vorto{explorar}\fako{trans.}\senco{1}{Observar minucioze e detaloze, esforcante penetrar til la esenco di la kozi}\senco{2}{Parkurar (lando) studiante sorgoze, por skopo ciencala, strategiala, e c.}\lingui{DEFIRS}
\artiklo{explotar}%
\vorto{explotar}\fako{trans.} Valorigar (kozo) per extraktar profito de lua produkturo\lingui{DEFRS}
\artiklo{explozar}%
\vorto{explozar}\fako{netrans.} Spliteskar violentoze\lingui{DEFIS}
\artiklo{exponento}%
\vorto{exponento}\fako{matem.} Mikra cifro dextre e kelke super nombro, por indikar ye qua potenco lu egardesas\lingui{DEFIRS}
\artiklo{exportacar}%
\vorto{exportacar}\fako{trans.} Transportar a lando stranjera ed ibe vendar (la produkturi da la sulo o da la laboro nacionala)\lingui{DEFIRS}
\artiklo{expozar}%
\vorto{expozar}\fako{trans.}\senco{1}{Prizentar, lokizar tale ke (la kozo) esas facile videbla}\senco{2}{Prizentar, lokizar tale ke (la kozo) submisesas a la efiko da ulo}\lingui{DEFIS}
\artiklo{expresar}%
\vorto{expresar}\fako{trans.} Manifestar (penso, sentimento) per parolo, gesto od irgaltra signo extera\lingui{DEFIS}
\artiklo{expropriar}%
\vorto{expropriar}\fako{trans.} Desposedigar (ulu) legale de la proprieto di ulo\lingui{DEFIS}
\artiklo{extazar}%
\vorto{extazar}\fako{netrans.}\senco{1}{Havar sua anmo, sua psiko, ravisata tale ke dum kelka tempo, ica ne plus koncias su}\senco{2}{Ravisesar pro admiro, pro joyo}\lingui{DEFIRS}
\artiklo{extensar}%
\vorto{extensar}\fako{trans.}\senco{1}{Developar longesale, larjesale}\senco{2}{Estalar (su) sur spaco por kovrar ica, sur areo por kovrar ica}\senco{3}{\textit{(metaf.)} Aplikar a plura kozi tale ke ici kontenesas}\lingui{EFIS}
\artiklo{extenuar}%
\vorto{extenuar}\fako{trans.} Igar magra extreme, per exhausto di la vigoro\lingui{EFIS}
\artiklo{exter}%
\vorto{exter} Esar en la spaco-parto qua komencas ube finas la korpo\lingui{EFIS}
\artiklo{exterminar}%
\vorto{exterminar}\fako{trans.} Destruktar til la lasta uno (di raso, di speco)\lingui{EFIS}
\artiklo{externo}%
\vorto{externo} Skolano (lernanto) qua ne habitas la lerneyo\lingui{DEFIS}
\artiklo{extingar}%
\vorto{extingar}\fako{trans.}\senco{1}{Cesigar la kombusteso o la lumizo}\senco{2}{Kalmigar, quietigar}\lingui{EFIS}
\artiklo{extirpar}%
\vorto{extirpar}\fako{trans.}\senco{1}{Arachar kun lua radiki (planto) por ke lu ne povez ri-kreskar}\senco{2}{Deprenar radikale parto malada}\senco{3}{\textit{(metaf.)} Destruktar radikale (kozo establisita che la homi)}\lingui{DEF}
\artiklo{extorsar}%
\vorto{extorsar}\fako{trans.} Koaktar ulu cedar ulo (generale pekunio) per minaco, per violento nemateriala\lingui{DEFIS}
\artiklo{extra}%
\vorto{extra} Di qua la qualeso esas extraordinara\lingui{DEFIRS}
\artiklo{extradar}%
\vorto{extradar}\fako{trans.) (aludante guvernerio} Livrar persono akuzata ye krimino, refujinta sur olua teritorio, a guvernerio stranjera qua demandas lu\lingui{DEFIS}
\artiklo{extradoso}%
\vorto{extradoso}\fako{arkitekt.} Surfaco konvexa quan vulto prizentas ad-extere (opoze a « intradoso »)
\artiklo{extraktar}%
\vorto{extraktar}\fako{trans.}\senco{1}{Retrotirar (kozo) de loko ube lu esas sinkita, stekita, nature od acidente}\senco{2}{Tirar de libro, de manuskripto (un o plura texto-peci quin onu kopias)}\senco{3}{Deprenar de substanco, per deskompozar lu (elemento quan onu separas de lu)}\senco{4}{\textit{(aritm.)} Extraktar radiko quadratala, kubala : ektirar ta radiko de ta nombro per des-kompozar ica}\senco{5}{\textit{(kemio)} Extraktar salo : per analizo kemiala}\lingui{DEFIRS}
\artiklo{extraordinara}%
\vorto{extraordinara}\senco{1}{Qua ne resortisas a la regulo, a la normo}\senco{2}{Qua diferas de to quon onu vidas ordinare}\senco{3}{Qua esas super la nivelo ordinara di la kozi, di la personi}\lingui{DEF}
\artiklo{extravagar}%
\vorto{extravagar}\fako{netrans.} Agar, dicar kozi qui ne korespondas a la raciono komuna\lingui{DEFIS}
\artiklo{extrema}%
\vorto{extrema}\senco{1}{Qua esas tote ye la fino}\senco{2}{Qua esas tote ye la lasta grado}\senco{3}{Qua esas ye grado tre alta}\senco{4}{Qua agas til la lasta punto}\lingui{DEFIS}
\artiklo{extrinseka}%
\vorto{extrinseka}\fako{filoz.} Qua tiresas, cherpesas, ne de la kozo ipsa, ma de cirkonstanci acesora, extera\lingui{EFIS}
\artiklo{exudar}%
\vorto{exudar}\fako{netrans.} Ekfluar preske nesenteble, quale sudoro\lingui{DEFIS}
\artiklo{exulcero}%
\vorto{exulcero}\fako{medic.} Ulcero neprofunda
\artiklo{exultar}%
\vorto{exultar}\fako{netrans.} Kaptesar da joyo profunda o da sentimento (amo, amoro) qua su expresas per movi, kanti, e c.\lingui{EFI}
\artiklo{exutorio}%
\vorto{exutorio}\senco{1}{\textit{(medic.)} Ulcero artificala, destinita a plear la rolo di deturnivo, por produktar ed entratenar pusifo lokala}\senco{2}{Moyeno uzata por eskartar ulo mala, por des-embarasar su ye lu}\lingui{FIS}
\artiklo{-eyo}%
\vorto{-eyo} Sufixo qua signifikas « loko destinita por »
\artiklo{-ez}%
\vorto{-ez}\fako{gram.} Dezinenco di la imperativo
\artiklo{ezofago}%
\vorto{ezofago}\fako{anat.} Tubo membrana qua iras de la faringo an la stomako, e duktas ad ica la nutrivi\lingui{EFIS}
\end{multicols}
\sekciono{F}
\begin{multicols}{2}
\artiklo{fa}%
\vorto{fa}\fako{muziko} La noto quaresma di la « do »-gamo
\artiklo{fablo}%
\vorto{fablo} Naraco kurta, proza o versa, di qua la protagonisti esas, maxim-multa-kaze, od enti imaginita, od animali, e qua uzesas kom pruviva pri leciono di etiko praktikala\lingui{DEFIS}
\artiklo{fabo}%
\vorto{fabo}\fako{bot.} Planto ek la familio « leguminosi », kun shelo grosa, qua kontenas grano oblonga, uzata kom nutrivo\latina{vicia faba}\lingui{FIS}
\artiklo{fabrikar}%
\vorto{fabrikar}\fako{trans.} Facar, per materii prima quin onu laboras, ula kategorii de objekti destinita a livresor engrose a la komercisti, segun procedi mashinala o mekanikala\lingui{DEFIRS}
\artiklo{facar}%
\vorto{facar}\fako{trans.} Tirar kozo ek altra, adaptante lu a destineso nova\lingui{FIS}
\artiklo{faceto}%
\vorto{faceto} Latero di objekto plura-angula\lingui{DEFIRS}
\artiklo{facies}%
\vorto{« facies »}\fako{cienco} Aspekto, fizionomio
\artiklo{facila}%
\vorto{facila} Facebla od agebla sen peno\lingui{EFIS}
\artiklo{facinar}%
\vorto{facinar}\fako{trans.}\senco{1}{Seduktar nerezisteble per la povo di la regardo}\senco{2}{\textit{(metaf.)} Dazlar per prestijo}\lingui{DEFIS}
\artiklo{facio}%
\vorto{facio} La parto « vizaja » di irga korpo\lingui{EFIS}
\artiklo{facionar}%
\vorto{facionar}\fako{netrans.) (pri soldato} Komisesir sorgor por la sekureso di posteno, di institucuro publika avan qua lu stacas armoze dum plura hori til remplaseso da altru\lingui{F}
\artiklo{fago}%
\vorto{fago}\fako{bot.} Arboro granda, ek la familio « amentacei », di qua la kortico esas grizatra, la foliaro densa\latina{fagus}\lingui{ISL}
\artiklo{fagocito}%
\vorto{fagocito}\fako{biol.} Celulo qua povas inglobar e digestar la partikuli neorganika ed organika
\artiklo{fagocitoso}%
\vorto{fagocitoso}\fako{biol.} Rolo di la fagocito, deskovrita da Mechnikof
\artiklo{fairo}%
\vorto{fairo}\senco{1}{Emanajo de kaloro e de lumo, quan produktas la kombusteso di ula korpi, quan la Antiqui konsideris kom un ek la quar elementi}\senco{2}{\textit{(metaf.)} Ardoro, saporo varma ed ecitiva di la drinkaji alkoholoza; ardoro di la psiko, di la pasioni}\lingui{DE}
\artiklo{fakiro}%
\vorto{fakiro} Nomo (en India) di la asketi mendikera di omna sekti (Mohamedisti, Civaisti, Vishnuisti) qui probas aquirar santeso per kontemplo e su-mutili\lingui{DEFIRS}
\artiklo{faksimilo}%
\vorto{faksimilo} Reprodukturo exakta de skriburo, de desegnuro, sive manuale, sive per imprimo, grabo, e c.\lingui{DEFIRS}
\artiklo{faktitivo}%
\vorto{faktitivo}\fako{gram.} Qua indikas ke la subjekto di la verbo eventigas la ago
\artiklo{fakto}%
\vorto{fakto}\senco{1}{To quo existas reale}\senco{2}{\textit{(yuro-cienco)} To quo eventas en afero, to quo judiciesas}\senco{3}{To quo eventis}\lingui{DEFIR}
\artiklo{faktorialo}%
\vorto{faktorialo}\fako{matem.) (algebro} Produto di qua la faktori esas segun progresiono geometriala\lingui{DE}
\artiklo{faktorio}%
\vorto{faktorio} Kontoro di agenti di kompanio komercala en lando stranjera\lingui{DEFIS}
\artiklo{faktoro}%
\vorto{faktoro}\senco{1}{\textit{(matem.)} En la multipliko, la multiplikanto e la multiplikendo}\senco{2}{\textit{(metaf.)} To (o : ta) quo (o = qua) efikas por la rezultajo}\lingui{DEFIRS}
\artiklo{faktotum}%
\vorto{« faktotum »} La homo qua, serve di ulu, agas irgo por plezar a lu\lingui{DEF}
\artiklo{fakturo}%
\vorto{fakturo}\fako{komerco} Noto quan la vendisto furnisas a la komprinto di la vari quin ita livras, sur qua indikesas la preco di ta vari\lingui{DEFIS}
\artiklo{fakultativa}%
\vorto{fakultativa} Quan onu darfas agar o ne agar\lingui{DEFIS}
\artiklo{fakultato}%
\vorto{fakultato}\senco{1}{\textit{(filoz.)} Povo di ento libera pri uzar sua agiveso. funcionala}\senco{2}{\textit{(universitato)} Korpo de profesori komisita por kurso-docar ed examenar, en omna fako di la domeno « doco supra-grada »}\lingui{DEFIRS}
\artiklo{falangino}%
\vorto{falangino}\fako{anat.} Falango duesma di la fingri trifalanga
\artiklo{falango}%
\vorto{falango}\senco{1}{\textit{(anat.)} Singla ek la mikra osti longa, ye qui konsistas la fingri di la manui e di la pedi}\senco{2}{\textit{(milit.)} \textit{(che la Greki antiqua)} Korpo de infantriani}\lingui{DEFIRS}
\artiklo{falansterio}%
\vorto{falansterio}\fako{sociologio} En la sistemo sociala da la Franco Fourier : habiteyo di singla grupo de cent familii\lingui{DEFIRS}
\artiklo{falar}%
\vorto{falar}\fako{netrans.} Irar de supre ad infre, pro la efiko da gravito\lingui{DE}
\artiklo{falbalo}%
\vorto{falbalo} Bendo de stofo plisita, frunsita, e c., qua uzesas por ornar jupo, manteleto, shultro-tuko, qua fixigesas sur ici per lua parto supra\lingui{DFI}
\artiklo{falchar}%
\vorto{falchar}\fako{trans.) (agrokultivo} Tranchar (la feno, la cereali) per instrumento ek longa mancho ligna e granda lamo arkatra\lingui{FI}
\artiklo{falco}%
\vorto{falco} Entaliuro en la kadro di pordo, di fenestro, por ke la bordo di la klapo qua klozas lu esez kontre-aplikebla sen lasar lum-lakuno\lingui{DR}
\artiklo{faldar}%
\vorto{faldar}\fako{trans.} Igar duopla unfoye o plurafoye (materio flexebla, stofo, papero, e c.) per abatar ica parto sur ita\lingui{DE}
\artiklo{faleno}%
\vorto{faleno}\fako{zool.} Nokto-papiliono\latina{phalaena}\lingui{DEFIS}
\artiklo{faliar}%
\vorto{faliar}\senco{1}{\textit{(trans.)} Ne atingar to quon onu vizis}\senco{2}{\textit{(netrans.)} Ne agar to quon onu devas etikale}\senco{3}{\textit{(netrans.)} \textit{(komerco)} Deklarar su nesolventa, cesar la pagado di sua debaji}\lingui{DEFIS}
\artiklo{falko}%
\vorto{falko}\fako{navo} Plankaro extera per qua onu plualtigas la bordo di navo kande la ondi esas tro granda\lingui{FIS}
\artiklo{falkono}%
\vorto{falkono}\fako{zool.} Raptucelo – genero di la klaso « rapteri » – quan onu dresis olim por chasado\latina{falco}\lingui{DEFIS}
\artiklo{falsa}%
\vorto{falsa}\senco{1}{Qua havas, di kozo, nur la semblo}\senco{2}{Qua prenas, por trompar, la semblo di kozo}\lingui{DEFIRS}
\artiklo{falseo}%
\vorto{falseo} Voco akuta, produktita da la vibro di la kordi supra di la laringo\lingui{DEFIRS}
\artiklo{falva}%
\vorto{falva} Flava rufe\lingui{DF}
\artiklo{familiara}%
\vorto{familiara}\senco{1}{Qua relatas ad ulu libere, sen jeneso, sen ceremonialo}\senco{2}{Qua kustumas la uzo di ulo; di qua la praktiko la uzo, esas ordinara ye ulu}\lingui{DEFIS}
\artiklo{familio}%
\vorto{familio}\senco{1}{Asemblitaro de la personi quin unionas la sango o la alianco, qui vivas en un hemo, partikulare la patrulo, la patrino, e lia filii}\senco{2}{Asemblitaro de la personi unionita per la sango o la alianco}\senco{3}{La serio de la personi qui devenis de la sama familio, generacionope}\senco{4}{Asemblitaro de enti di qui la origino esas komuna, di qui la interesti esas komuna}\senco{5}{\textit{(zool. e bot.)} Kolekturo de generi intervicina pro ula karakterizivi generala quin li posedas komune}\lingui{DEFIRS}
\artiklo{famino}%
\vorto{famino} Sufrado generala, produktita da la indijo di nutrivi\lingui{EFI}
\artiklo{famo}%
\vorto{famo} Favoreso, segun-modeso, quan la nomo di ulu (o di ulo) obtenis\lingui{DEFI}
\artiklo{fanar}%
\vorto{fanar}\fako{trans.} Netigar, purigar (frumento) per sukusar lu sur sorto di korbo plata, di qua la bordo esas basa, por forjetar la polvo, la brano, la palio-peceti\lingui{DEF}
\artiklo{fanatika}%
\vorto{fanatika}\senco{1}{Qua kredas inspiresar deale}\senco{2}{Ta quan zelo religiala blindigas, incitas ad agi neyusta, violentoza, qui transiras la limiti di la raciono komuna}\lingui{DEFIRS}
\artiklo{fanerogamo}%
\vorto{fanerogamo}\fako{bot.} Planto di qua esas videbla la organi di la fruktifado\lingui{DEFIS}
\artiklo{fanfaro}%
\vorto{fanfaro}\fako{muziko} Ario armeala kurta, tre ritmoza, pleata per trumpeti od altra muzik-instrumenti kupra\lingui{DEFIS}
\artiklo{fanfaronar}%
\vorto{fanfaronar}\fako{netrans., pri, pro} Exajerar sua avantaji, sua forteso, e c., per parolo od altre, por trompar sua kun-homi\lingui{FIS}
\artiklo{fango}%
\vorto{fango} Tero, polvo, dilutita dense da aquo sur sulo, en la stradi, voyi\lingui{DFI}
\artiklo{fantasmagorio}%
\vorto{fantasmagorio} Produkto, en loko obskura, di figuri lumoza, di qui la imaji grandeskas e semblas marchar ad la spektanti\lingui{DEFIRS}
\artiklo{fantastika}%
\vorto{fantastika} Kreita da fantazio, da kaprico, da la imagino-fakultato exter la posiblaji e la realaji\lingui{DEFIS}
\artiklo{fantazia}%
\vorto{« fantazia »} Sorto di distrakto, di amuzo, en qua la kavalieri Araba agas galope evolucioni diversa, pafante
\artiklo{fantaziar}%
\vorto{fantaziar}\fako{trans., pri} Konceptar, vidar la kozi, havar idei qui naskas de imaginado-kaprico\lingui{DEFIRS}
\artiklo{fantomo}%
\vorto{fantomo}\senco{1}{Aparajo qua aspektas quale la imajo de persono qua ne plus vivas, o quale la imituro di persono reala}\senco{2}{\textit{(metaf.)} To quo nur aspektas quale persono, kozo}\lingui{DEFIRS}
\artiklo{farad}%
\vorto{« farad »}\fako{elektrotekniko} Kapaceso di konduktivo di qua la potencialo kreskas per un \textit{volt} po charjuro de 1 \textit{coulomb}. (La farad equivalas 10. kapaceso-unaji C.G.S.)\lingui{DEF}
\artiklo{faradizar}%
\vorto{faradizar}\fako{trans.} Elektrizar por skopi medicinala o terapiala, per elektrofluo induktita ed interruptita\lingui{DEFIS}
\artiklo{farandolo}%
\vorto{farandolo} Danso Provenciana, quaza marcho ritmoza, segun movado vivaca, agita da longa lineo de danseri qui intertenas per la manui\lingui{DEFIS}
\artiklo{faraono}%
\vorto{faraono}\senco{1}{Titulo di la suvereno, en Egiptia antiqua}\senco{2}{Hazardo-ludo}\lingui{DEFIRS}
\artiklo{farbo}%
\vorto{farbo} Substanco quan onu aplikas sur ula objekti por produktar la sento di koloro\lingui{D}
\artiklo{farcino}%
\vorto{farcino}\fako{patol.} Inflamuro, ofte kontagiala, akompanata da molesko di la ganglioni e di la vaskuli limfala, qua afektas la kavali, la muli\lingui{EFI}
\artiklo{fardo}%
\vorto{fardo}\senco{1}{Kompozajo quan onu aplikas sur la pelo por plubeligar la karnaciono}\senco{2}{\textit{(metaf.)} Aspekto versemblanta qua kolorizas o travestias la fakti}\lingui{F}
\artiklo{faringito}%
\vorto{faringito}\fako{patol.} Inflamuro di la faringo
\artiklo{faringo}%
\vorto{faringo}\fako{anat.} La dopa parto di la boko
\artiklo{farino}%
\vorto{farino} Substanco nutriva di ula grani de cereali, pulverigita\lingui{EFIS}
\artiklo{farizeo}%
\vorto{farizeo} La homo qua agas la gesti religiala, ma di qua la pieso esas plu semblanta kam reala\lingui{DEFIRS}
\artiklo{farmacio}%
\vorto{farmacio} La arto preparar la medikamenti\lingui{DEFIRS}
\artiklo{farmakopeo}%
\vorto{farmakopeo} Libro qua traktas la arto preparar la medikamenti\lingui{DEFIRS}
\artiklo{farmar}%
\vorto{farmar}\fako{trans.} Signatar la konvenciono per qua proprietero lugas ter-domeno ad ulu qua explotos olu po preco yarala\lingui{DEF}
\artiklo{faro}%
\vorto{faro} Turmo erektita sur litoro, ye la somito di qua esas fairo lumo-fonto qua uzesas por guidar (dum nokto) la iro di la navi\lingui{FIS}
\artiklo{farso}%
\vorto{farso}\senco{1}{Hachajo de karni spicizita ye qua onu garnisas la parto interna di pultro, di pasteto, e c.}\senco{2}{Mikra teatrajo bufonala}\senco{3}{Kozo bufonala quan onu dicas od agas}\lingui{DEFIRS}
\artiklo{fasado}%
\vorto{fasado} Parto avana di konstrukturo, ube esas la enireyo precipua\lingui{DEFIS}
\artiklo{fascia}%
\vorto{« fascia »}\fako{histol.} Termino per qua indikesas formacuri aponevrosala, veli sive fibra sive celulara, qui kovras od envelopas muskuli o regioni
\artiklo{fashino}%
\vorto{fashino} Branchari asemblita a faski, quin onu uzas por plenigar la fosati di fortifikajo, shirm-apogar kanono-baterii, impedar la krulo di tero, e c.\lingui{DEFIRS}
\artiklo{fashismo}%
\vorto{fashismo} Sistemo politikala, qua havas karaktero nacionalista, antisocialista, antiparlamentista, instaurita (en 1922) en Italia, da Benito Mussolini\lingui{DEFIRS}
\artiklo{fasiacar}%
\vorto{fasiacar}\fako{netrans.) (bot.} Deformesar teratologie per aplasto (stipi, rami, pedunkli)\lingui{EF}
\artiklo{fasko}%
\vorto{fasko}\senco{1}{Asemblajo de ula quanto de kozi simila, longa, interligita}\senco{2}{\textit{(optiko)} Ensemblo ek radii o de lumo-pinseli qui emanas de un fonto}\lingui{FIS}
\artiklo{fasonar}%
\vorto{fasonar}\fako{trans.} Laborar ula materio por fabrikar objekto definita (fasono di vesto, di juvelo)\lingui{DEFR}
\artiklo{fastar}%
\vorto{fastar}\fako{netrans., dum} Abstenar plu o min komplete nutrivi\lingui{DE}
\artiklo{fatigar}%
\vorto{fatigar}\fako{trans.}\senco{1}{Abatar per spensigar lua energio}\senco{2}{Repugnar per tedo, obsedo}\lingui{EFIS}
\artiklo{fato}%
\vorto{fato}\senco{1}{Rango atribuita ad homo da potento qua (imaginate) decidas segun sua arbitrio pri la destineso di singla homo}\senco{2}{Ta ipsa potento}\lingui{DEFIRS}
\artiklo{fatua}%
\vorto{fatua} Qua stulte kontentesas pri su ipsa\lingui{EFIS}
\artiklo{fauco}%
\vorto{fauco}\fako{anat.}\senco{1}{Parto extera ed avana di la kolo (do, anke la laringo, ma vidata extere)}\senco{2}{\textit{(metaf.)} Formi kava (tekn., geografio)}\lingui{IL}
\artiklo{fauno}%
\vorto{fauno}\fako{mitol.} Deo agrala, ordinare reprezentata kom kornoza, e kapro-pedoza\lingui{DEFIRS}
\artiklo{favo}%
\vorto{favo}\fako{patol.} Tinio produktita da fungo. \textit{(achorion Schoenleini)}\lingui{L}
\artiklo{favorar}%
\vorto{favorar}\fako{trans.} Avantajizar ulu, prefere kam la ceteri\lingui{DEFIRS}
\artiklo{fayenco}%
\vorto{fayenco} Terakotajo vernisita od emaliizita\lingui{DEFR}
\artiklo{fayo}%
\vorto{fayo} Silko-stofo gros-grana\lingui{DEFS}
\artiklo{fazano}%
\vorto{fazano}\fako{zool.} Ucelo, genero ek la familio « galinacei », sen kresto sur la kapo, di qua la kaudo esas longa, etajoza, di qua la plumaro esas, che la maskulo, extreme brilanta\latina{phasianus}\lingui{DEFIRS}
\artiklo{fazeolo}%
\vorto{fazeolo}\fako{bot.} Planto leguminosa\latina{phaseolus vulgaris}\lingui{DI}
\artiklo{fazo}%
\vorto{fazo}\senco{1}{Singla de la aspekti sucedanta di Luno e di la planeti dum lia rotaco, e segun la situeso (relate Tero) di la parto lumizata da Suno}\senco{2}{Singla de la estadi sucedanta ye qui pasas kozo dum lua evoluciono}\lingui{DEFIRS}
\artiklo{febla}%
\vorto{febla} Qua ne plus povas tre efikar. \textit{(anke metafore)}\lingui{EF}
\artiklo{febrar}%
\vorto{febrar}\fako{netrans.} Esar en la stando di maladatreso quan karakterizas pulsi tre, tro frequa, plumulteso di la kaloro korpala, e generala deskomforto\lingui{DEFIRS}
\artiklo{februaro}%
\vorto{februaro} Monato duesma di la yaro, qua kontenas 28 dii en la yari ordinara, e 29 en la yari bisextila\lingui{DEFIRS}
\artiklo{federar}%
\vorto{federar}\fako{trans.} Unionar (populi, societi, sindikati, e c.) per obligo reciproka\lingui{DEFIRS}
\artiklo{feko}%
\vorto{feko} Exkrementi anusala\lingui{DEFIS}
\artiklo{fekulo}%
\vorto{fekulo} Sedimento amilatra quan depozas la suki di ula materii vejetala, o quan formacas la farino lavita di la grani de cereali, leguminosi, di la tuberkuli de terpomi (= *kartofli), di la rizomi de manioko\lingui{EFIS}
\artiklo{fekunda}%
\vorto{fekunda} Apta produktar enti simila a su. \textit{(anke metaf.)}\lingui{EFIS}
\artiklo{feldo}%
\vorto{feldo}\fako{tekn.} Spaco, areo, en qua eventas ula distributo di forci (elektrala, magnetala, e c.)\lingui{EF}
\artiklo{felgo}%
\vorto{felgo}\fako{tekn.}\senco{1}{Singla ek la sis peci de ligno kurvigita qua, en roto, formacas la cirklo-kurvo extera en qua inkastresis la spoki}\senco{2}{Singla ek la quar peci de ligno, plu granda, qui formacas cirklo-kurvo duesma, tenono-asemblita en la unesma}\lingui{DE}
\artiklo{felica}%
\vorto{felica} Di qua la anmo, la psiko, esas satisfacita komplete\lingui{EFIS}
\artiklo{felina}%
\vorto{felina} Di qua la movi e gesti esas dolca e flexebla quale olti di la kato\lingui{EFIS}
\artiklo{felo}%
\vorto{felo}\fako{tekn.} Pelo deprenita de animalo mortinta, e preparita por skopi diversa\lingui{DE}
\artiklo{felpo}%
\vorto{felpo} Stofo de lano, de kotono, de silko qua imitas plusho, ma di qua la pili esas plu longa e min densa\lingui{DI}
\artiklo{felto}%
\vorto{felto} Texuro forta, densa, ek lano e pili, aglutinita e fulita\lingui{DEFIS}
\artiklo{femina}%
\vorto{femina} Qua apartenas a la sexuo organizita por konceptar e parturar, o por pondar ovi\lingui{DEFIRS}
\artiklo{femuro}%
\vorto{femuro}\fako{anat.} Osto di la kruro\lingui{EFIS}
\artiklo{fenacetino}%
\vorto{fenacetino} C₂H₅O . C₆H₄ . NH . CO . CH₃\lingui{DEF}
\artiklo{fenco}%
\vorto{fenco} Barilo qua klozas\lingui{DE}
\artiklo{fendar}%
\vorto{fendar}\fako{trans.}\senco{1}{Dividar longesale}\senco{2}{Dividar per eskartar la parti poke adheranta}\lingui{FIS}
\artiklo{fenestro}%
\vorto{fenestro}\senco{1}{Aperturo di qua la formo e la grandeso esas diversa, facita en la muri di konstrukturo por lasar aero e jorno penetrar}\senco{2}{Kadro vitroplakoza qua klozas ta aperturo}\lingui{DEFI}
\artiklo{fenikulo}%
\vorto{fenikulo}\fako{bot.} Planto aromata de la familio « umbeliferi »\latina{foeniculum vulgare}\lingui{DEFIL}
\artiklo{feno}%
\vorto{feno} Herbo quan onu falchas e lasas sikeskar, uzata kom nutrivo por la kavali, la brutaro\lingui{FIS}
\artiklo{fenolo}%
\vorto{fenolo}\fako{kemio} Substanco desinfektiva, extraktita de la olei tre pezoza, devenanta de la distilo di minkarbono e di gudro\simbolo{C₆H₅OH}\lingui{DEFIRS}
\artiklo{fenomeno}%
\vorto{fenomeno} Modifikeso di korpo quan perceptas nia sensi\lingui{DEFIRS}
\artiklo{fenomenologio}%
\vorto{fenomenologio}\senco{1}{Traktato, diserturo, pri la fenomeni, pri to quo povas frapar la sensi}\senco{2}{\textit{(filoz.)} En la Hegelsistemo, cienco pri la idei qui naskas de la percepti da la sensi}\lingui{DEF}
\artiklo{feo}%
\vorto{feo} Ento fantastika, quan onu reprezentas kom muliero dotita ye povo supernatura\lingui{DEFIR}
\artiklo{ferdeko}%
\vorto{ferdeko}\fako{navo} Plankaro fixigita sense di la longeso di navo, por kovrar la holdo o por formacar etaji\lingui{DE}
\artiklo{ferio}%
\vorto{ferio} Granda merkato ube vendesas vari omnasorta, qua eventas unfoye o plurafoye singlayare, en regiono, ordinare lor ed en la loko ube celebresas la festo lokala\lingui{EF}
\artiklo{ferma}%
\vorto{ferma}\senco{1}{Qua ne flexas, qua ne shancelas}\senco{2}{\textit{(metaf.)} Qua ne ocilas}\senco{3}{\textit{(komerco)} Di qua la kondicioni esas definitiva e ne dependas de ta o ca eventualajo}\lingui{DEFIS}
\artiklo{fermentacar}%
\vorto{fermentacar}\fako{netrans.} Subisar la transformo molekulala quan produktas fermento. \textit{(anke metaf.)}\lingui{DEFIRS}
\artiklo{fermento}%
\vorto{fermento}\senco{1}{To quo produktas ula transformi molekulala (diastazi, e c.)}\senco{2}{\textit{(metaf.)} To quo naskigas ula agitesi psikala}\lingui{DEFIRS}
\artiklo{fero}%
\vorto{fero} Metalo blanka, grizatra, fuzebla per alta temperaturo, tenaca, duktila, forjebla, magnetala, uzata tre diverse en industrio\simbolo{Fe}\lingui{EFIS}
\artiklo{feroca}%
\vorto{feroca} Di qua la karaktero esas kruela, sango-durstanta\lingui{EFIS}
\artiklo{fertila}%
\vorto{fertila} Qua produktas multo (exemple : agro)\lingui{DEFIS}
\artiklo{fervorar}%
\vorto{fervorar}\fako{netrans., pri, pro, por} Qua ardoras.vivace\lingui{EFIS}
\artiklo{fes}%
\vorto{« fes »} Noto quaresma di la « do » gamo, bemoloza
\artiklo{festar}%
\vorto{festar}\senco{1}{\textit{(trans.)} Agar la solenaji religiala, memorigala, quin onu celebras en ula dii di la yaro; joyo-manifestar honorize di eventajo memorinda}\senco{2}{\textit{(netrans.)} Partoprenar gala-festino, gala-balo, e c., donata da rejo, princo, urbo, privato}\lingui{DEFIS}
\artiklo{festinar}%
\vorto{festinar}\fako{netrans.} Agar repasto festala, pompoza, honorize di ulo o di ulu\lingui{EFS}
\artiklo{festono}%
\vorto{festono}\senco{1}{Envolvajo ek folii e flori, arkoforma, quan onu suspendas intervalope, kom festo-ornivi}\senco{2}{Bordo di vesto, di linjo, taliita formacante denti kurva (o quadrata, o tri-angula)}\senco{3}{Broduro qua kovras la bordo}\lingui{DEFIRS}
\artiklo{festuko}%
\vorto{festuko} Stipeto de palio\lingui{F}
\artiklo{fetida}%
\vorto{fetida} Mal-odoranta repugnante\lingui{DEFIS}
\artiklo{fetisho}%
\vorto{fetisho} Objekto materiala, deigita da ula populi sovaja\lingui{DEFIRS}
\artiklo{feto}%
\vorto{feto} Che la animali vivipara : la produkturo de la koncepto, kande lu ne plus esas embriono e komencas formacesar\lingui{DEFIS}
\artiklo{feudo}%
\vorto{feudo} Terdomeno di nobelo, di qua la posedanto resortisas a la sinioro di altra terdomeno, debas ad ilca sincereso e homajo, ed obligesas a lu pri ula servi e debi\lingui{DEFIRS}
\artiklo{fezo}%
\vorto{fezo} Boneto reda kun grosa fasketo ek lano o blua silko, uzata en la landi Mohamedista\lingui{DEFIRS}
\artiklo{fi!}%
\vorto{fi!} Interjeciono qua expresas la desprizo di ulo\lingui{DEFR}
\artiklo{fiakro}%
\vorto{fiakro} Lugo-veturo qua halte-vartas sur la placi, e quan onu lokacas iro-forfete o segun-dure\lingui{DEFR}
\artiklo{fialo}%
\vorto{fialo} Recipiento analoga a flakono, ma min granda, e kun streta « boko », quan onu uzas precipue por la liquidi farmaciala\lingui{DEFI}
\artiklo{fiancar}%
\vorto{fiancar}\fako{trans.} Obligar per promiso reciproka di mariajo, agita solene koram sacerdoto, e koram parenti, amiki\lingui{EFI}
\artiklo{fiasko}%
\vorto{fiasko} Nesuceso di autoro, di artisto, ye lua publiko\lingui{DEFIRS}
\artiklo{fibrilo}%
\vorto{fibrilo}\fako{bot.} Filamenti tenua ye qui konsistas la radiketi
\artiklo{fibrino}%
\vorto{fibrino} Substanco organika, liquida, sen-odora, sen-sapora, grizatra-blanka, quan onu renkontras partikulare en la sango, en la chilo, de qua lu separas su per koagulo kande ta liquidi esas ekpasinta\lingui{DEFIRS}
\artiklo{fibro}%
\vorto{fibro} Singla de la elementi tenua, longa, flexebla, di qui la interplekteso konstitucas ula substanco animala, vejetala o minerala\lingui{DEFIRS}
\artiklo{fibroino}%
\vorto{fibroino}\fako{kemio} Substanco albuminoida, homogena e diafana, ye qua konsistas parte la silko di la kokono di la insekti silkifera\lingui{DE}
\artiklo{fibulo}%
\vorto{fibulo}\fako{en la epoki antiqua} Brocho, agrafo ek metalo, por retenar la extremaji di vesto\lingui{DEFIS}
\artiklo{ficho}%
\vorto{ficho} Lamo de ivoro, de osto, de perlomatro, qua uzesas, dum karto-ludo, e c., por la konto di la punti quin singlu ganas, e qui pagesas lor la fino di la ludo\lingui{DFRS}
\artiklo{fidar}%
\vorto{fidar}\fako{netrans., ad} Tale estimar (ulu) ke onu konfidas a lu la sorgo okupor su senkontrole pri ulo\lingui{EFIS}
\artiklo{fideikomiso}%
\vorto{fideikomiso}\fako{yuro-cienco} Donaco o legaco ad ulu, konfidita a lu kom honesto, por transmisesor ad altru\lingui{DFIS}
\artiklo{fidela}%
\vorto{fidela} Qua ne falias pri to quon expektas de lu persono a qua lu ligesas per su-obligir\lingui{EFIS}
\artiklo{fiera}%
\vorto{fiera} Ta qua havas alta (ma generale justa) opiniono pri sua valoro, merito, povo, e qua manifestas lu extere\lingui{FIS}
\artiklo{fifro}%
\vorto{fifro}\fako{muziko} Mikra fluto sis-trua, di qua la sono esas tre akuta, uzita olim kom akompananto di la tamburo dum marchi armeala\lingui{DEF}
\artiklo{figo}%
\vorto{figo}\fako{bot.} Frukto di qua la pulpo esas mola e sukroza, produktata da arboro de la familio « urticei », di qui la flori esas asemblita en recevuyo komuna\latina{ficus}\lingui{DEFIRS}
\artiklo{figuranto}%
\vorto{figuranto} La homo qua, en teatro, engajesis por plear la rolo di la protagonisti qui ne parolas od esas tre acesora\lingui{DEFS}
\artiklo{figuro}%
\vorto{figuro}\senco{1}{Formo videbla di korpo}\senco{2}{\textit{(retor.)} Formo di la parolo, qua modifikas la expreso di la pensado}\senco{3}{\textit{(geom.)} Reprezenturo grafika di la formi quin geometrio konceptas en spaco}
\artiklo{fijo}%
\vorto{fijo} Karto sur qua onu skribas la tituli di la libri di biblioteko, noti, dokumenti, quan onu klasifikas alfabetale o numerale, en buxi, por rekursar a lu kande onu bezonas ta o ca libro, dokumento, e c.
\artiklo{fikario}%
\vorto{fikario}\fako{bot.} Planto herbacea di qua la flori esas flava, di la genero « ranunkuli »\latina{ranunculus ficaria}
\artiklo{fiko}%
\vorto{fiko}\fako{patol.} Surkreskajo karnoza, kun pedunklo streta, e somito inflata figatre
\artiklo{fiktiva}%
\vorto{fiktiva} Di qua la valoro, la realeso, esas nur konvencionala\lingui{DEFIRS}
\artiklo{filamento}%
\vorto{filamento} Elemento longa e tenua ye qua konsistas ula tisui vejetala, animalala o minerala\lingui{DEFIS}
\artiklo{filandro}%
\vorto{filandro} Filo lejera qua flugetas en aero\lingui{F}
\artiklo{filantropo}%
\vorto{filantropo} La homo qua amas homaro\lingui{DEFIRS}
\artiklo{filaso}%
\vorto{filaso} Materio texebla (kanabo, lino, e c.) ne ja filigita
\artiklo{filatelio}%
\vorto{filatelio} Cienco, studio, sercho di la posto-marki (*timbri postala)\lingui{DEFIS}
\artiklo{filetika}%
\vorto{filetika} Qua apartenas a, o relatas, sive la ancestraro di speco, sive la procedi qui abutis a la formaceso di la speci, sive la diversa lineaji di la grupi di la klasifiko naturala
\artiklo{fileto}%
\vorto{fileto}\senco{1}{Spiralo qua cirkondas skrubo}\senco{2}{Parto di la animalo, qua, pro havar nek tendini nek osti, esas distranchebla a lonchi dina}\lingui{DFIS}
\artiklo{filharmonio}%
\vorto{filharmonio} Amo pasionoza pri muziko\lingui{DEFIRS}
\artiklo{filheleno}%
\vorto{filheleno}\senco{1}{Partisano di la Greki antiqua, di la arti e di la civilizeso-sistemo di Grekia antiqua}\senco{2}{Amiko di la Greki moderna, partisano di lia nedependo politikala}\lingui{DEFIS}
\artiklo{filialo}%
\vorto{filialo} Firmo fondita da firmo « patra »
\artiklo{filigrano}%
\vorto{filigrano} Laboruro ek fili de oro, de arjento, de vitro, interplektita ed interligita per solduri preske neperceptebla\lingui{DEFIRS}
\artiklo{filiko}%
\vorto{filiko}\fako{bot.} Planto senkotiledona, kriptogama, herbacea, di qua la semineti kontenesas da kapsuli an la surfaco reversa di la folio, di qua la folii volvesas krosatre ante lia desklozeso, e di qua ula speci divenas arboratra en la regioni tropikala\latina{filicea}\lingui{IL}
\artiklo{filio}%
\vorto{filio} Persono konsiderata pri la ligilo qua unionas lu a ti de qui lu naskis\lingui{EFIRS}
\artiklo{filipiko}%
\vorto{filipiko} Diskurso violentoza kontre ulu qua guvernas la Stato\lingui{DEFIRS}
\artiklo{filistro}%
\vorto{filistro} Nomo donata da la kulturozi, da la artisti, a la personi quin li konsideras kom nekapabla komprenar to quo esas spiritala, artala
\artiklo{filmo}%
\vorto{filmo} Bendo fotografala celuloida, uzata en la kameri obskura di la cinematografilo\lingui{DEF}
\artiklo{filo}%
\vorto{filo} Brino tenua, longa, ek texo-materio, metalo\lingui{FIS}
\artiklo{filologo}%
\vorto{filologo} La cienco qua studias ta o ca linguo nacionala kom organo di la vivo mentala, spiritala, di la populo qua uzas olu\lingui{DEFIRS}
\artiklo{filoxero}%
\vorto{filoxero}\fako{zool.} Insekto qua rodas la radiko di la vito\latina{phylloxera vastatrix}\lingui{DEFIRS}
\artiklo{filozofiar}%
\vorto{filozofiar}\fako{netrans.}\senco{1}{Rezonar pri la origini e la kauzi}\senco{2}{Rezonar ciencoze pri ulo}\lingui{DEFIRS}
\artiklo{filozofo}%
\vorto{filozofo} La persono qua su konsakras a la cienco pri la origini e la kauzi
\artiklo{filtrar}%
\vorto{filtrar}\fako{trans.} Pasigar (liquido) tra korpo poroza (felto, stofo, papero, karbono, petro sponjatra, e c.) por purigar lu\lingui{DEFIRS}
\artiklo{financo}%
\vorto{financo} Provizuro de pekunio quan ulu (o Stato) disponas\lingui{DEFIRS}
\artiklo{finar}%
\vorto{finar}\senco{1}{\textit{(trans.)} Par-facar, par-agar, per traktar til (ed inkluzite) la lasta parto}\senco{2}{\textit{(netrans.)} Arivar an la punto ube onu haltas (en spaco, en tempo)}\lingui{EFIS}
\artiklo{fingar}%
\vorto{fingar}\fako{trans.} Prizentar kom reala (eso-maniero di qua onu prenas la aspekto)\lingui{DEFIRS}
\artiklo{fingro}%
\vorto{fingro}\fako{anat.} Singla de la prolonguri distinta, ek falangi artikoza, maxim-multa-kaze moviva, ube finas la pedi di la animali, la manui e la pedi di la homo e di la simio\lingui{DE}
\artiklo{finisar}%
\vorto{finisar}\fako{trans.) (tekn.} Sorgoze perfektigar la detali maxim mikra\lingui{F}
\artiklo{finita}%
\vorto{finita}\senco{1}{\textit{(matem.)} \textit{(pri *grandoro)} Qua esas quanto determinita; \textit{(pri progresiono)} konsistanta ye nombro determinita de termini}\senco{2}{\textit{(filoz.)} Di qua la eso esas limitoza}\lingui{EF}
\artiklo{finko}%
\vorto{finko}\fako{zool.} Ucelo kantera, pasero konobeka\latina{fringilla coelebs}\lingui{D}
\artiklo{finto}%
\vorto{finto}\fako{skerm-arto} Stroko per qua onu minacas ca latero por ke la adverso deskovrez ta latero quan onu volas atingar\lingui{F}
\artiklo{fiorituro}%
\vorto{fiorituro}\fako{muziko} Orno-traiti quin kantisto, exekutante, adjuntas a la frazo muzikala di la kompozisto\lingui{DFIRS}
\artiklo{firmamento}%
\vorto{firmamento} La vulto ciela\lingui{F}
\artiklo{firmo}%
\vorto{firmo} Nomo komercala, adoptita da komerco-societo o da privato komercista\lingui{DEFRS}
\artiklo{fis}%
\vorto{« fis »} La noto quaresma di la « do »-gamo, diezoza
\artiklo{fisho}%
\vorto{fisho}\fako{zool.} Animalo vertebroza qua vivas en aquo, respirante la aero qua esas dissolvita en ica, per organi partikulara (brankii)\lingui{DE}
\artiklo{fisko}%
\vorto{fisko} Administrantaro di la trezoro statala, di qua la tasko konsistas ye kolektar la imposto-revenuo\lingui{DEFIRS}
\artiklo{fistulo}%
\vorto{fistulo}\fako{patol.} Kanalo acidentala, formacita da ulcero, kavajo naturala\lingui{DEFIRS}
\artiklo{fitar}%
\vorto{fitar}\fako{netrans.} Tale adjustigita (peci mashinala), ke li ne plus povas movar en ula direciono exter la axala. – \textit{(pri vesto)}. Adaptita perfekte a la formo di la korpo di la persono qua *weras lu – od a parto di lua korpo (manui, pedi)\lingui{EI}
\artiklo{fiziko}%
\vorto{fiziko} La cienco di qua la studiajo esas la propraji generala di la korpi e la legi (*leyi) qui koncernas lia efiki extera, sen egardar la kompozanti\lingui{DEFIRS}
\artiklo{fiziokratio}%
\vorto{fiziokratio} Doktrino ekonomikala qua konsideras ke tero (sulo) esas la fonto unika di richeso\lingui{DEFIRS}
\artiklo{fiziologio}%
\vorto{fiziologio} La parto di la natur-historio qua traktas la roli di la organi che la enti vivanta en lia stando normala\lingui{DEFIRS}
\artiklo{fizionomio}%
\vorto{fizionomio} Expresuro che la vizajo\lingui{DEFIRS}
\artiklo{fizognomonio}%
\vorto{fizognomonio} La arto determinar la karaktero di persono segun la traiti di lua vizajo\lingui{DEFIRS}
\artiklo{flado}%
\vorto{flado} Tarto moleta, ek kremo, farino ed ovi\lingui{D}
\artiklo{flagelo}%
\vorto{flagelo}\fako{histol.} Filamento moviva qua (che ula protozoi, e che la gameti maskula, ed anke che la celuli epiteliala) funcionas kom organo lokomocala\lingui{DEFIRS}
\artiklo{flago}%
\vorto{flago} Insigno (generale min granda kam standardo) qua indikas la lojeyo, navo o loko, di ula chefo o di ula establisuro\lingui{DER}
\artiklo{flagrar}%
\vorto{flagrar}\fako{netrans.}\senco{1}{Kombustesar rapide kun flami granda e subita. \textit{(anke metaf.)}}\senco{2}{\textit{(pri delikto)} Eventar vidate da la persono ipsa qua konstatas lu}\lingui{DEFIS}
\artiklo{flajoleto}%
\vorto{flajoleto}\fako{muziko} Fluto bekoza, sis-tono, quan onu pluperfektigas per adjuntar a lu klavi\lingui{DEFS}
\artiklo{flako}%
\vorto{flako} Mikra marsheto de aquo, neprofunda\lingui{F}
\artiklo{flakono}%
\vorto{flakono} Mikra botelo, ek vitro, kristalo, porcelano, kun stopilo generale vitra o metala\lingui{DEFR}
\artiklo{flamingo}%
\vorto{flamingo}\fako{zool.} Ucelo de la klaso « steltieri », di qua la suba parto di la ali esas fairo-reda\lingui{DEFIRS}
\artiklo{flamo}%
\vorto{flamo} Kombinuro de la oxo di aero kun la partikuli o la gasi qui eskapas de la materii kombustata\lingui{DEFIS}
\artiklo{flanar}%
\vorto{flanar}\fako{netrans.} Irar sen skopo, lasante su distraktesar da ta kozo, da ca kozo, por pasigar tempo\lingui{DF}
\artiklo{flanelo}%
\vorto{flanelo} Stofo dolca, lanugoza, ek lano pektita o kardita, di qua la texuro esas kelke desdensa, simpla o krucumita\lingui{DEFIRS}
\artiklo{flanjo}%
\vorto{flanjo}\fako{tekn.} Parto diskatra, salianta, di tubi, stangi, e c., qua uzesas por ligar li per skrubi e bolti\lingui{DEIR}
\artiklo{flanko}%
\vorto{flanko}\senco{1}{\textit{(anat.)} Parto latera di la ventro, situita, an singla latero di la korpo, inter la lakuno kostala e la loko ube komencas la hancho}\senco{2}{\textit{(fortifikajo)} Parto di la bastiono, inter la kurteno e la shultro-angulo}\lingui{DEFIR}
\artiklo{flarar}%
\vorto{flarar}\fako{trans.} Probar dicerno segun la odoro. \textit{(anke metafore)}\lingui{F}
\artiklo{flatar}%
\vorto{flatar}\fako{trans.} Probar sedukto per laudi, nemeritita od exajerita\lingui{DEF}
\artiklo{flatuar}%
\vorto{flatuar}\fako{netrans.} Evakuar gaso intestinala, tra la sfinktero anusala\lingui{DEFIS}
\artiklo{flatulenta}%
\vorto{flatulenta}\fako{medic.} Che qua gasi, venti, akumulesis en la kanalo digestala
\artiklo{flava}%
\vorto{flava} Di qua la koloro esas olta di oro, di safrano, di citrono, e c.\lingui{FL}
\artiklo{flebito}%
\vorto{flebito}\fako{patol.} Inflamuro di la membrano interna di la veini\lingui{DEFIS}
\artiklo{flecho}%
\vorto{flecho}\senco{1}{Armo, ek stango ligna, garnisita, an un ek la extremaji, ye ferpeco, pintoza, an la altra ye aleti ek plumi o metalo, e quan onu lansas per arko od arbalesto}\senco{2}{\textit{(arkitekt.)} Alteso vertikala, de la punto maxim alta di vulto, di arko, super la komenco-punti}\senco{3}{\textit{(fortifikajo)} Masivo salianta angule}\senco{4}{\textit{(geom.)} Perpendiklo de la mezo di arko an la kordo qua klozas lu; o disto di la kordo di arko an la tangento paralela}\lingui{FIS}
\artiklo{flegar}%
\vorto{flegar}\fako{trans.}\senco{1}{Sorgar atencoze por la komforto di malado, di infanteto o di infanto}\senco{2}{Striglar (kavalo)}\lingui{D}
\artiklo{flegmo}%
\vorto{flegmo} Karaktero di la persono kalma, qua su dominacas\lingui{DEFIRS}
\artiklo{flegmono}%
\vorto{flegmono}\fako{patol.} Centro pusifanta qua rezultas de la inflameso di la tisuo celuloza subpela\lingui{DEFIRS}
\artiklo{flereto}%
\vorto{flereto} Speco di espado, di qua la lamo esas quadrato-seciona, tenua e flexebla, generale kun buleto ye la pinto, garnisita ye felo, uzata por exercar su pri skermo\lingui{DFIS}
\artiklo{flexar}%
\vorto{flexar}\fako{trans.}\senco{1}{Kurvigar per preso qua interproximigas la du extremaji}\senco{2}{\textit{(metaf.)} Cesar pokope sua rezisto}\lingui{EFIS}
\artiklo{flexiono}%
\vorto{flexiono}\fako{gram.} Modifikeso di vorto, da la chanjo di lua dezinenci (en la deklino, la konjugo, che ula lingui nacionala)\lingui{DEFIS}
\artiklo{flintglaso}%
\vorto{flintglaso} Kristalo pura, por la lensi akromata, kompozita ek plumbo-silikato e kalio-silikato, e tre refraktiva\lingui{DEFRS}
\artiklo{flirtar}%
\vorto{flirtar}\fako{netrans.) (aludante personi di sexui interdiferanta} Manovrar (per vorti o per konduto pleziva) por seduktar\lingui{DEFR}
\artiklo{flitro}%
\vorto{flitro} Lamelo de oro, de arjento, o de kupro orizita od arjentizita, per qua onu plubriligas broduri, galoni, rubandi, e c.\lingui{D}
\artiklo{flogar}%
\vorto{flogar}\fako{trans.} Batar per ledro-bendo o kordeto, agitante la mancho ye qua fixigesas un ek la extremaji\lingui{E}
\artiklo{floko}%
\vorto{floko}\senco{1}{Mikra tufo lejera ek lano, silko, kotono}\senco{2}{Mikra maso, poke densa}\lingui{DEFIR}
\artiklo{florino}%
\vorto{florino} Moneto-peco (olim ora, nun arjenta) uzata en Nederlando, e c.\lingui{DEFIRS}
\artiklo{floro}%
\vorto{floro}\fako{bot.} Parto di la planto, qua su developas pos la folii, desklozas sua korolo ofte odoroza, ornita ye kolori plu o min briloza, e, pos existo provizora, remplasesas da la frukto\lingui{DEFIRS}
\artiklo{florono}%
\vorto{florono}\senco{1}{Ornivo imitanta la formo di floro}\senco{2}{\textit{(arkitekt.)} Ornivo skultura qua reprezentas floro, folio}\lingui{FIS}
\artiklo{floso}%
\vorto{floso} Apendico membranoza per qua su movas en aquo la fishi e kelka altra animali aquala\lingui{D}
\artiklo{flotacar}%
\vorto{flotacar}\fako{netrans.} Portesar sur liquido segun la arbitrio di la ondi, di la fluo\lingui{DEFIS}
\artiklo{flotilio}%
\vorto{flotilio} Floto de mikra navi\lingui{DEFIRS}
\artiklo{floto}%
\vorto{floto} Asemblajo plu o min granda de navi militala o komercala qui navigas kune, ed havas un iro-skopo (generale komandata da un chefo)\lingui{EF}
\artiklo{floxo}%
\vorto{floxo}\fako{bot.} Planto di qua la infloresenco esas piramidatra\latina{phlox}
\artiklo{flozelo}%
\vorto{flozelo} Buro de silko, reziduo de la kokoni desvolvita\lingui{EF}
\artiklo{fluar}%
\vorto{fluar}\fako{netrans.) (aludante liquido} Irar de loko ad altra loko, per movo kontinua. – \textit{(anke metaf.)}\lingui{DEF}
\artiklo{flugar}%
\vorto{flugar}\fako{netrans.} Sustenigar su e movar su en aero atmosfero per ali (naturala od artificala). (ex. : la uceli e la maxim multa insekti; la aviacili : aeroplani, helikopteri)\lingui{DE}
\artiklo{fluida}%
\vorto{fluida} Di qua la molekuli, poke inter-adheranta, ne rezistas mem nur mikra esforco qua tendencas a diplasar li\lingui{DEFIRS}
\artiklo{fluktuar}%
\vorto{fluktuar}\fako{netrans.} Quaze ondifar en aero atmosfero, segun la arbitrio di vento o di altra impulso variiva. (kom ex. : standardo, hari.) – \textit{(anke pri objekto imersita)}\lingui{DEFIRS}
\artiklo{flundro}%
\vorto{flundro}\fako{zool.} Nomo di fisho de la genero « pleuronektoidi »
\artiklo{fluorecar}%
\vorto{fluorecar}\fako{netrans.) (pri ula dissolvuri} Kovreskar ye envelopilo lumoza qua aspektas analoge kam opalo kande onu lumizas li per suno-radii en chambro obskura\lingui{DEFIRS}
\artiklo{fluorino}%
\vorto{fluorino}\fako{kemio} Kombinuro de fluoro e kalcio\simbolo{CaF₂}\lingui{F}
\artiklo{fluoro}%
\vorto{fluoro}\fako{kemio} Korpo simpla, metaloida, qua deskompozas aquo ye la temperaturo ordinara, e korodas preske omna metali\simbolo{F}\lingui{DEFIRS}
\artiklo{fluro}%
\vorto{fluro}\fako{arkitekt.} Platformo en eskalero, ye la punto ube finas etajo\lingui{DE}
\artiklo{fluto}%
\vorto{fluto}\fako{muziko} Muzik-instrumento ventala, cilindro apertata supre, klozata infre, qua havas, apud la parto supra, aperturo qua uzesas kom boko-peco ube onu suflas, e longesale provizita ye trui sur qui onu pozas la fingri por apertar e klozar ta trui segun la noto quan onu volas produktar\lingui{DEFIRS}
\artiklo{fluvio}%
\vorto{fluvio}\senco{1}{Granda aquofluo qua konservas sua nomo til la maro ube lu debushas}\senco{2}{\textit{(metaf.)} To quo fluas abunde quale fluvio}\lingui{EFIRS}
\artiklo{fluxiono}%
\vorto{fluxiono}\senco{1}{\textit{(patol.)} Adfluo di sango o di altra liquido ad ula tisui, ad ula organi, konseque de stando inflamala}\senco{2}{\textit{(matem.)} La quanti finita quin produktas movo o fluo kontinua (la lineo, per la movo di la punto : la surfaco, per la movo di la lineo, e c.)}\lingui{DEFIS}
\artiklo{fluxo}%
\vorto{fluxo} La mareo adlitora\lingui{EFIRS}
\artiklo{foko}%
\vorto{foko}\senco{1}{\textit{(optiko)} Punto a qua konvergas la lumo-radii reflektita o refraktita}\senco{2}{\textit{(geom.)} Punto o punti de ube departas ed adube su asemblas la vektori}\lingui{DEFIR}
\artiklo{fola}%
\vorto{fola} Termino vulgara qua aplikesas precipue a la formi maxim videbla e shokiva di alienaco, a ti qui su manifestas per parolo sencesa, agiteso e violento\lingui{EFI}
\artiklo{folado}%
\vorto{folado}\fako{zool.} Molusko bivalva qua perforas la roki\latina{pholadidea}
\artiklo{folietono}%
\vorto{folietono} Parto rezervita, infre de jurnalo-pagino (cirkume la triimo di la alteso) por la publikigo di artikli kritikala, di romani, e c.\lingui{DEFRS}
\artiklo{folikulo}%
\vorto{folikulo}\senco{1}{\textit{(anat.)} Mikra sako membrana en la tegumento}\senco{2}{(bot.Frukto kapsula, ek folio volvita longesale, ed apertita per quaza suturo)}\lingui{DEFIRS}
\artiklo{folio}%
\vorto{folio}\senco{1}{\textit{(bot.)} Lamelo verda qua naskas de la stipi, de la rami}\senco{2}{To quo esas dina e plata quale arboro-folio (papero, e c.)}\lingui{EFIRS}
\artiklo{folium}%
\vorto{« folium »}\fako{matem.} Kurvo di qua parto esas pasable analoga a la formo di arboro-folio
\artiklo{folkloro}%
\vorto{folkloro} Civilizeso-formo materiala, strukturi sociala, tradicionaji e literaturo-sorti parolala, metodologio (ca fako inkluzante la tekniki yena : bibliografio, kartografio, enrejistro sonala, inquesti, muzeografio)\lingui{DEFR}
\artiklo{fomentar}%
\vorto{fomentar}\fako{trans.) (medic.} Remediar (parto di la korpo) per aplikar e mantenar sur lu topiko varma\lingui{DEFIS}
\artiklo{fondar}%
\vorto{fondar}\fako{trans.}\senco{1}{Apogar sur ulu, sur ulo (to quon onu entraprezas, esperas, kredas, e c.)}\senco{2}{Establisar (ulo) ye maniero solida, ferma}\lingui{DEFIS}
\artiklo{fonetiko}%
\vorto{fonetiko} La parto di linguistiko, qua traktas la son-uni (foni)\lingui{DEFIRS}
\artiklo{fonetismo}%
\vorto{fonetismo} Reprezento exakta, per skribo, di la foni vocala
\artiklo{fongo}%
\vorto{fongo}\fako{patol.} Surkreskajo karnoza, sponjatra, quale fungo (en ulcero, vunduri)
\artiklo{fono}%
\vorto{fono}\fako{fiziko} Un ek la serio de movi alternanta, tre rapida, da korpo solida, liquido o gaso, quin la nervo audala sentas per la orelo
\artiklo{fonografar}%
\vorto{fonografar}\fako{trans.} Enrejistrar voco per aparato qua konsistas ye disko vibranta, provizita ye stalo-stilo qua imprimas sur cilindro o disko rotacanta la vibri recevita, tale ke esas reproduktebla segun-vole, kande onu rotacigas la cilindro o la disko, la vorti dicita avan la aparato\lingui{DEFIRS}
\artiklo{fonolito}%
\vorto{fonolito}\fako{mineral.} Roko volkanala qua resonas sub martelofrapi\lingui{DEF}
\artiklo{fonometrio}%
\vorto{fonometrio} La arto mezurar la intenseso di la voco-foni
\artiklo{fonometro}%
\vorto{fonometro}\fako{tekn.} Aparato por mezurar la intenseso-grado di la foni\lingui{DEFIRS}
\artiklo{fontanelo}%
\vorto{fontanelo}\fako{anat.} Che la infanteti, lakuno inter la osti kraniala, plenigita ye tisuo mola, flexebla, qua osteskas erste kande la homo esas plu evoza\lingui{DEFIS}
\artiklo{fonteno}%
\vorto{fonteno} Aparato qua varsas, per robineto, la aquo ye qua onu plenigis lu o quan onu adduktas per tubi\lingui{DEFIR}
\artiklo{fonto}%
\vorto{fonto}\senco{1}{Origino di aquofluo, di fonteno – aquo-veino qua produktas lu en la loko ube lu komencas ekirar de sulo}\senco{2}{\textit{(metaf.)} Origino di kozo, to de quo olca devenas, derivesas}\lingui{EFIRS}
\artiklo{for}%
\vorto{for} Prepoziciono qua signifikas eskarteso, plu o min granda, de…\lingui{EF}
\artiklo{forcato}%
\vorto{forcato} Persono qua subisas puniso afliktiva ed infamiganta qua konsistis, olim, ye remar sur la galeri rejala; nun, ye deportesar en remparo-cirkuito fortifikita\lingui{FIS}
\artiklo{forcepso}%
\vorto{forcepso}\fako{medic.} Instrumento ek fero, du-brancha quan uzas la akushisto en la kazi di desfacileso, por sizar du-latere la kapo di la infanti ed atraktar lu ad su\lingui{DEFIS}
\artiklo{forco}%
\vorto{forco}\fako{fiziko}\senco{1}{La kauzo di movo}\senco{2}{Kauzo nekonocata di ula movi di materio (gravito, kaloro, elektro, e c.)}\lingui{EFIS}
\artiklo{forejo}%
\vorto{forejo} Feno di la prati naturala od artificala, herbi, vejetanti qui uzesas por nutrar la brutaro\lingui{DEFIRS}
\artiklo{foresto}%
\vorto{foresto} Vasta surfaco, areo, de tereno arboroza\lingui{DEFIS}
\artiklo{forfeto}%
\vorto{forfeto}\fako{yuro-cienco} Kontrato inter du personi, ek qui ica su obligas prenor – ita, livror, po preco determinita anticipe, ula laboruri, ula servi\lingui{FI}
\artiklo{forfikulo}%
\vorto{forfikulo}\fako{zool.} Insekto ortoptera, di qua la ventro finas per du hoketi pincatra\latina{forficula}\lingui{IL}
\artiklo{forjar}%
\vorto{forjar}\fako{trans.} Laborar (metalo) sur amboso, per fairo e martelo\lingui{EFS}
\artiklo{forko}%
\vorto{forko}\senco{1}{Instrumento kun longa mancho ligna, finanta per du o tri branchi pintoza, ek ligno o fero (por inversige agitar feno, palio, movar sterko, deprenar o kargar garbi)}\senco{2}{Objekto qua kelke similesas forko}\lingui{DEFI}
\artiklo{formacar}%
\vorto{formacar}\fako{trans.} Facar (ulu, ulo) donante a lu la eso e la formo\lingui{DEFIRS}
\artiklo{formalino}%
\vorto{formalino}\fako{kemio} H. CO₂H\lingui{DEF}
\artiklo{formato}%
\vorto{formato} Dimensiono quan onu donas a libro, a jurnalo, segun ke onu faldas la imprimo-folio unfoye, dufoye, trifoye\lingui{DFIRS}
\artiklo{formiko}%
\vorto{formiko}\fako{zool.} Mikra insekto ek la klaso « himenopteri », qua facas por su (en la kavaji di la olda arbori) lojeyo ube lu vivas societane\latina{formica}\lingui{FIS}
\artiklo{formo}%
\vorto{formo}\senco{1}{Aspekto extera di korpo}\senco{2}{\textit{(metaf.)} Karaktero sub qua kozo su manifestas}\senco{3}{Maniero extera agar segun ula reguli}\senco{4}{Maniero expresar penso}\lingui{DEFIRS}
\artiklo{formulario}%
\vorto{formulario} Kolektajo de formuli
\artiklo{formulo}%
\vorto{formulo} Formo determinita segun qua onu expresas kozo konvencionita\lingui{DEFIRS}
\artiklo{fornikar}%
\vorto{fornikar}\fako{netrans.) (religio} Pekar sexuale\lingui{EFIS}
\artiklo{forno}%
\vorto{forno} Masonuro vultoza, cirklatra, kun un aperturo avana, ube onu bakas pano, kuki, e c.\lingui{FI}
\artiklo{foronomio}%
\vorto{foronomio} Geometrio di movo, o cinematiko geometriala\lingui{DEFIRS}
\artiklo{forsan}%
\vorto{forsan} Adverbo qua signifikas « povas eventar ke… ma lo ne esas musto »
\artiklo{forsar}%
\vorto{forsar}\fako{trans.} Kurvigar, torsar, apertar (ulo) per violento\lingui{DEF}
\artiklo{forta}%
\vorto{forta}\senco{1}{\textit{(aludante korpo o parto di korpo)} Di qua la povo di impulso, di rezisto, esas tre efikiva}\senco{2}{\textit{(aludante la anmo, e fakultato di la anmo)} di la povo volala, sentimentala, inteligentesala esas kapabla efikar multe}\senco{3}{\textit{(aludante persono, o asemblitaro de personi)} Qua disponas la moyeni necesa por impozar sua autoritato}\lingui{EFIRS}
\artiklo{fortifikar}%
\vorto{fortifikar}\fako{trans.} Provizar ye defensili ek masonuri, teramasigi, e c.\lingui{DEFIRS}
\artiklo{fortreso}%
\vorto{fortreso} Placo fortifikita\lingui{EFIRS}
\artiklo{fortuno}%
\vorto{fortuno}\senco{1}{Deino quan la pagani reprezentis forme di muliero qua tenas en sua manuo abundo-korno, di qua la okuli esis linjo-bendizita, sidinta sur roto, por indikar ke elu distributas la bonaji e la malaji blinde e kapricoze}\senco{2}{Chanco bona}\lingui{DEFIS}
\artiklo{forumo}%
\vorto{forumo}\senco{1}{\textit{(che la Romani antiqua)} Placo ube la populo diskutis pri la aferi publika}\senco{2}{\textit{(metaf.)} Loko ube onu diskutas la aferi publika}\lingui{DEFIRS}
\artiklo{fosato}%
\vorto{fosato} Kavajo linea, en sulo, uzata por limitizar, por separar tereni, por plufaciligar la ekfluo di la aquo pluvala o drenala\lingui{FI}
\artiklo{fosfato}%
\vorto{fosfato}\fako{kemio} Kombinuro de la \textit{fosf}-acido kun korpo qua igas lu salo\lingui{DEF}
\artiklo{fosfeno}%
\vorto{fosfeno}\fako{fiziol.} Fenomeno lumala, quan onu povas produktar sur la retino per kompresar la okul-globo\lingui{DEF}
\artiklo{fosfo}%
\vorto{fosfo}\fako{kemio} Korpo simpla, tre inflamebla, venena, qua lasas eskapar, ye la temperaturo ordinara di atmosfero, vapori alie-odoranta e lumetas en loko obskura\simbolo{P}\lingui{DEFIRS}
\artiklo{fosforecar}%
\vorto{fosforecar}\fako{netrans.) (kemio} Lumar en loko obskura. (ex. : fosfo)\lingui{DEFIRS}
\artiklo{fosilo}%
\vorto{fosilo} Qua permanis enfosigita en la strati sedimenta anciena di la Terkortico\lingui{DEFIS}
\artiklo{foso}%
\vorto{foso} Kavajo od exkavajo plu o min larja e profunda (en sulo), uzata kom recevuyo\lingui{FIS}
\artiklo{fosto}%
\vorto{fosto}\senco{1}{Peco de karpenturo, pozita vertikale}\senco{2}{Longa ligno-peco stekita en sulo}\lingui{D}
\artiklo{fotelo}%
\vorto{fotelo} Granda stulo, kun dors-apogilo e brakii
\artiklo{foto}%
\vorto{foto}\fako{fiziko} Radiado de ula korpi, qua videbligas la objekti
\artiklo{fotofono}%
\vorto{fotofono} Aparato de ube la son-ondi, recevite sur spegulo vibranta qua transformas li ad lum-ondi, departas por re-produktar su kelka-diste\lingui{DEFIRS}
\artiklo{fotogena}%
\vorto{fotogena} Dicesas pri irga parto di organismo qua produktas ed emisas lumo-radii fiziologiala
\artiklo{fotograbar}%
\vorto{fotograbar}\fako{trans.} Aplikar la fotografado e la grabado per vitro-plako lumo-impresebligita\lingui{DEFIRS}
\artiklo{fotografar}%
\vorto{fotografar}\fako{trans.} Fixigar la imajo de korpi, obtenita per kamero obskura, sur vitroplako (o filmo) garnisita ye substanco impresebla da lumo\lingui{DEFIRS}
\artiklo{fotokemio}%
\vorto{fotokemio} La cienco di la reakti kemiala quin produktas la lumo-radii
\artiklo{fotolitografar}%
\vorto{fotolitografar}\fako{trans.} Aplikar la fotografado a la litografado per fotografuro dekalquebla sur petro
\artiklo{fotometrio}%
\vorto{fotometrio} Parto di fiziko, qua traktas la mezuro di la intenseso-gradi lumala
\artiklo{fotometro}%
\vorto{fotometro} Instrumento por mezurar la intenseso-grado di lumo\lingui{DEFIRS}
\artiklo{fotosfero}%
\vorto{fotosfero}\fako{astron.} La atmosfero lumoza qua cirkondas Suno\lingui{DEFIRS}
\artiklo{fotosintezo}%
\vorto{fotosintezo}\fako{bot.} Sintezo qua eventas en la folii di la planti verda, influe da la lumo-radii
\artiklo{fototaktismo}%
\vorto{fototaktismo} Movo (influe da lumo-radii) da korpi moviva (zoospori, e c.)
\artiklo{fototerapio}%
\vorto{fototerapio} Uzo di la lumo-radii kom agento terapiala\lingui{DEF}
\artiklo{fototipio}%
\vorto{fototipio} Procedo di *reprodukto fotomekanikala per imprim-inki, en qua uzesas substanci koloida (gelatino, albumino, bitumo) extensita sur fundi diversa (vitro, kupro, petro, zinko) ed igita inkizebla per la lumo-radii\lingui{DEF}
\artiklo{fototropismo}%
\vorto{fototropismo} Efiko da la lumo-radii a la kresko di la planti
\artiklo{foxo}%
\vorto{foxo}\fako{zool.} Quadripedo karnavida, ek la genero « hundi », di qua la muzelo esas dina, la kaudo longa e tufa\latina{canis vulpes}\lingui{DE}
\artiklo{foxtroto}%
\vorto{foxtroto} Danso (zoologiala, ek Usa) quar-tempa, segun ritmo inter marcho e polko\lingui{E}
\artiklo{foyero}%
\vorto{foyero}\fako{teatro}\senco{1}{Loko ube la spektinti iris olim por varmigar su dum la inter-akti, en la epoki dum qui la pleo-salono ne esis kalorizata}\senco{2}{Salono, galerio, ube la spektinti promenas dum la inter-akti}\lingui{DEFR}
\artiklo{foyo}%
\vorto{foyo}\senco{1}{Singla de la kazi ye qui fakto eventas}\senco{2}{Singla de la kazi ye qui quanto kontenesas kom elemento di toto}\lingui{F}
\artiklo{fraciono}%
\vorto{fraciono}\fako{matem.} Un o plura parti di la uno dividita o subdividita a parti inter-egala\lingui{DEFIRS}
\artiklo{fragmento}%
\vorto{fragmento}\senco{1}{Peco de kozo ruptita o dislacerita}\senco{2}{\textit{(metaf.)} To quo restas de libro di qua la maxim multo esas perdita o permanis nefinita. Peco citita ek libro}\lingui{DEFIS}
\artiklo{frago}%
\vorto{frago}\fako{bot.} Frukto multa-bera, qua devenas de planto herbacea de la familio « rozacei »\latina{fragaria}\lingui{FISL}
\artiklo{frajila}%
\vorto{frajila}\senco{1}{Ruptebla facile}\senco{2}{\textit{(metaf.)} Olqua facile cedas a la tentesi}\lingui{DEFIS}
\artiklo{frakasar}%
\vorto{frakasar}\fako{trans.} Ruptar splitige\lingui{FI}
\artiklo{frako}%
\vorto{frako} Viro-vesto nigra kun baski moruo-kaudatra, uzata por la vesper-festi mondumala, la ceremonii\lingui{DFRS}
\artiklo{framasono}%
\vorto{framasono} Persono qua iniciesis en asociitaro (olim sekreta) qua profesas principi di interfrateso, havas kom emblemi ula instrumenti di masonisti, e di qua la membraro konsistas ye grupi diversa (« logii ».)\lingui{EFIRS}
\artiklo{frambo}%
\vorto{frambo}\fako{bot.} Frukto multa-bera, tre parfumoza, devenanta de arboreto dornoza ek la familio « rozacei »\latina{rubus idaeus}\lingui{FS}
\artiklo{framo}%
\vorto{framo} Menuzuro quaza-kadra; asemblajo de montanti e de traversi; envelopilo buxatra\lingui{E}
\artiklo{franciskano}%
\vorto{franciskano} Reguliero di la ordeno (religiala, katolika) quan fondis santa Franciscus di Assisi\lingui{DEFIRS}
\artiklo{francisko}%
\vorto{francisko}\fako{arkeol.} Hakilo militala, kun du tranchivi, uzita da la Franki antiqua\lingui{DEFS}
\artiklo{frangulo}%
\vorto{frangulo}\fako{bot.} Arboreto qua kreskas inter la boski e la hegi, en la tereni humida, di qua la ligno esas transformebla a karbono lejera quan onu uzas por fabrikar la kanonpulvero, di qua la kortico interna esas purgiva, e di qua la beri esas verda o reda, segun la grado di matureso\latina{rhamnus frangula}\lingui{IL}
\artiklo{franjipano}%
\vorto{franjipano} Sorto di aromo\lingui{DEFIS}
\artiklo{franjo}%
\vorto{franjo}\senco{1}{Ornivo ek serio de brini, de torsadi pendanta (kotona, lana, silka, orfila, arjenta, e c.)}\senco{2}{\textit{(fiziko)} Franji, linei alterne briloza ed obskura, en la fenomeno « interfero »}\lingui{DEFIS}
\artiklo{franko}%
\vorto{franko} Unajo monetala di la sistemo decimala qua korespondas a moneto-peco ek arjento aloyita a 0 fr.165 de kupro, qua pezas 5 grami, dividita a 10 dekimi o 100 centimi\lingui{F}
\artiklo{frankolino}%
\vorto{frankolino}\fako{zool.} Ucelo de la genero « perdriki », samastatura kam la fazano, di qua la gambi esas longa, la beko forta e longa\latina{pternistes vulgaris}\lingui{DEFIRS}
\artiklo{frapar}%
\vorto{frapar}\fako{trans.} Shokar unfoye o plurafoye per apliko violentoza\lingui{F}
\artiklo{frato}%
\vorto{frato}\senco{1}{Persono konsiderata relate lua parenteso kun altru qua naskis de la sama gepatri}\senco{2}{\textit{(metaf.)} Persono konsiderata relate la ligilo qua unionas lu a la membri cetera di la familio homara}\lingui{EFIS}
\artiklo{fraudar}%
\vorto{fraudar}\fako{trans.}\senco{1}{Trompar (ulu) por prokurar a su, detrimento di ta ulu, ca o ta avantajo}\senco{2}{Frustrar la administrantaro doganala o fiskala}\lingui{DEFIS}
\artiklo{fraxino}%
\vorto{fraxino}\fako{bot.} Arboro de la familio « oleacei », di qua la ligno, kompakta, blanka e veinoza, uzesas en industrio\latina{fraxinus}\lingui{FISL}
\artiklo{frayo}%
\vorto{frayo} Fish-ovi fekundigita\lingui{FIS}
\artiklo{frazeologio}%
\vorto{frazeologio} Maniero konstruktar sua frazi, qua karakterizas linguo o skriptisto\lingui{DEFIRS}
\artiklo{frazo}%
\vorto{frazo}\senco{1}{\textit{(gram.)} Propoziciono simpla od asemblajo de propozicioni qua facas senco kompleta}\senco{2}{\textit{(muziko)} Serio de son-uni qui prizentas a la oreli developeso reguloza e kompleta}\lingui{DEFIRS}
\artiklo{fregato}%
\vorto{fregato}\fako{navig.} La maxim granda de la milito-navi kun un ferdeko o tota baterio, qua, importale, situesas dop la kombatonavo\lingui{DEFIRS}
\artiklo{fremisar}%
\vorto{fremisar}\fako{netrans., de, pro}\senco{1}{Agitar su bruisetoze}\senco{2}{Agitar su konvulse}\lingui{FS}
\artiklo{*frendo}%
\vorto{*frendo} Kun-partoprenanto di *muvmento
\artiklo{frenezio}%
\vorto{frenezio}\senco{1}{Deliro furioza}\senco{2}{\textit{(metaf.)} Eceso en ula pasiono}\lingui{DEFIS}
\artiklo{freno}%
\vorto{freno} Mekanismo uzata por lentigar o haltigar la movo di mashino\lingui{FIS}
\artiklo{frenologio}%
\vorto{frenologio} Doktrino segun qua singla de nia fakultati od instinti havas kom situeso ta o ca parto di la cerebro, ed esas konocebla per la observo di la protuberanci o di la kavaji korespondanta che la kranio\lingui{DEFIS}
\artiklo{frequa}%
\vorto{frequa} Qua eventas multafoye, tre mikra-intervale\lingui{DEFIS}
\artiklo{frequenco}%
\vorto{frequenco}\fako{elektrotekniko) (pri la elektro-flui alternanta} La nombro de ocili dum un sekundo
\artiklo{frequentar}%
\vorto{frequentar}\fako{trans.}\senco{1}{Venar (a loko) granda-nombre}\senco{2}{Venar (a loko) ofte}\senco{3}{Venar freque (kun persono) kom akompananto}\lingui{FIS}
\artiklo{fresha}%
\vorto{fresha} Qua ne esas velkinta\lingui{DEFIRS}
\artiklo{fresko}%
\vorto{fresko}\senco{1}{Pikturo sur muro, qua, facite per aquo-farbi sur induturo tre recenta, sikeskas e hardeskas kun lu}\senco{2}{To quo esas piktita per ta procedo}\lingui{DEFIRS}
\artiklo{fretar}%
\vorto{fretar}\fako{trans.} Lokacar navo por transportar vari\lingui{DEFS}
\artiklo{frezar}%
\vorto{frezar}\fako{trans.) (tekn.}\senco{1}{Mi-bizelizar la orifico di truo en qua esas insertota skrubo od irgaltra objekto}\senco{2}{Laborar, entaliar la ligno o la metali per solido rotaciva qua esas garnisita ye denti, per qua – kande la solido rotacas – onu povas deprenar spani de la surfaco kun qua onu igas lu kontaktar tre preso}\lingui{F}
\artiklo{friabla}%
\vorto{friabla} Qua esas facile reduktebla a peceti\lingui{EFIS}
\artiklo{frianda}%
\vorto{frianda}\fako{kozo} Qua luras per ulo delikata segun quante koncernesas la palato\lingui{F}
\artiklo{fricionar}%
\vorto{fricionar}\fako{trans., per} Frotar parto di la homo-korpo por igar plu rapida la cirkulo di sango, por kuracar doloro, e c.\lingui{DEFIS}
\artiklo{frigoro}%
\vorto{frigoro} To quo rezultas de la basigo di temperaturo; absenteso kompleta o parta di kaloro\lingui{EF}
\artiklo{frikasar}%
\vorto{frikasar}\fako{trans.} Koquar en sauco (legumi, o karno pecigita)\lingui{DEFIRS}
\artiklo{frikaseo}%
\vorto{frikaseo} Raguto de karno pecigita, e koquita en sauco\lingui{DEFIRS}
\artiklo{fripo}%
\vorto{fripo} Vesti, linjo, mobli olda, de qui onu su desembarasas per vendar li a brokantisti\lingui{EF}
\artiklo{fripono}%
\vorto{fripono}\senco{1}{Persono qua habile furtas mikra kozi}\senco{2}{Persono qua agas mala stroki malicoze}\lingui{F}
\artiklo{friskar}%
\vorto{friskar}\fako{netrans.} Agitar (su) per movi vivaca e kurta\lingui{E}
\artiklo{friso}%
\vorto{friso}\senco{1}{En la monumenti di stilo Greka : parto di la entablamento inter la arkitravo e la kornico}\senco{2}{Bendo piktita o skultita cirkum la parto supra di chambro, di vazo, super la kadro di pordo, di kameno, e c.}\lingui{DEFIRS}
\artiklo{fristo}%
\vorto{fristo} Tempo-quanto grantenda o grantita por efektigar ulo\lingui{D}
\artiklo{fritar}%
\vorto{fritar}\fako{trans.} Koquar per padelo, en graso, oleo od en butro bolianta\lingui{EFIS}
\artiklo{fritilario}%
\vorto{fritilario}\fako{bot.} Planto liliacea, di qua la floro havas la formo di korneto subversita, e di qua la bulbo pulpoza kontenas elemento akra, drastika\latina{fritillus}
\artiklo{frivola}%
\vorto{frivola}\senco{1}{Tante vana ke lu ne meritas ke onu fervorez por lu}\senco{2}{Qua fervoras por la kozi vana}\lingui{DEFIRS}
\artiklo{frizar}%
\vorto{frizar}\fako{trans.}\senco{1}{Volvar (hari, pili, e c.) cirkum li, preske lokle}\senco{2}{\textit{(netrans.)}\textit{(aludante hari, pili, e c.)} Ipse volvar cirkum su, preske lokle}\lingui{DEFS}
\artiklo{frogo}%
\vorto{frogo} Sorto di kazako kun longa maniki, di qua la uzo-kustumo venis de Brandeburgia, lor la rejo Ludovikus XIV. di Francia, ordinare garnisita ye butoni olivatra, interligita per galoni\lingui{F}
\artiklo{froko}%
\vorto{froko} Vesto di monako katolika\lingui{EF}
\artiklo{frolar}%
\vorto{frolar}\fako{trans.} Tushar lejere la bordo, la extremajo di ulo\lingui{F}
\artiklo{fromajo}%
\vorto{fromajo} Substanco nutrala quan onu obtenas de preparar di-verse butro, la parto kazeatra di lakto, o ta du materii kombinita\lingui{FI}
\artiklo{frondo}%
\vorto{frondo} Lans-armo, posho, quan onu igas jirar per du kordi qui lu suspendesas, e qua kontenas stono o globeto quan onu lasas eskapar bruske\lingui{FI}
\artiklo{frontiero}%
\vorto{frontiero}\senco{1}{Limito qua separas la teritorio di stato de olta di stato kontigua}\senco{2}{Limito qua impedas expanso}\lingui{EFIS}
\artiklo{frontispico}%
\vorto{frontispico}\senco{1}{\textit{(arkitekt.)} Facio precipua di edifico}\senco{2}{Titulo di libro, akompanata da vinyeti, e c.; graburo lokizita opoze a la pagino qua kontenas la titulo}\lingui{DEFIS}
\artiklo{fronto}%
\vorto{fronto}\senco{1}{La parto supra di la vizajo homala, inter la brovi e la lineo ube komencas la hararo}\senco{2}{La parto supra di la vizajo di la animali}\senco{3}{La parto supra di edifico, di monto, di arboro}\senco{4}{La facio avana di ula kozi (armeo qua marchas, e c.)}\lingui{DEFIRS}
\artiklo{frontono}%
\vorto{frontono} Ornivo qua quaze kronizas la enireyo precipua di edifico, e qua havas kom bazo la kornico di la entablamento\lingui{EFIRS}
\artiklo{frostar}%
\vorto{frostar}\fako{netrans.} Transformesar a glacio\lingui{DE}
\artiklo{frotar}%
\vorto{frotar}\fako{trans.} Igar (korpo) kontaktar kun altra korpo, quan onu igas pasar sur ita presante\lingui{DFS}
\artiklo{frua}%
\vorto{frua} Ante la instanto kande onu kustumas agar lo ordinare; ante la instanto kande onu expektas lua evento. \textit{(*anticipa)}\lingui{D}
\artiklo{*frua}%
\vorto{*frua} Nemulta tempo pos la komenco di la periodo aludata
\artiklo{frugala}%
\vorto{frugala} Qua saciesas per nutrivi simpla\lingui{DEFIS}
\artiklo{frugilego}%
\vorto{frugilego}\fako{zool.} Korvo di qua la beko esas rekta, akuta, ne garnisita ye plumi\latina{corvus frugilegus}
\artiklo{frukto}%
\vorto{frukto}\fako{bot.}\senco{1}{Produkturo de la vejetanto, qua sucedas a floro}\senco{2}{Organo di la vejetanti (ovario, fekundigita) qua kreskas en la floro, ube lu existas en estado rudimenta e kontenas la semino (ovulo) destinita a *reproduktar la planto}\senco{3}{\textit{(metaf.)} a. \textit{(aludante patrino)}. La filio quan elu konceptis, quan elu portas en sua ventro. b. \textit{(aludante mariajo)}. La filio qua naskas de la mariajo. c. Rezultajo avantajoza de ulo}\lingui{DEFIRS}
\artiklo{frumento}%
\vorto{frumento}\fako{bot.} Planto cereala, de la familio « cereali », di qua la spiko esas kompozaja, kontenanta grano manjebla de qua onu facas farino\latina{triticum}\lingui{DEFIRS}
\artiklo{frunsar}%
\vorto{frunsar}\fako{trans.}\senco{1}{Plisar (la brovi) per kontrakto}\senco{2}{Igar (stofo) plisita per pasar en lu filo o tre adheranta rubando}\lingui{FS}
\artiklo{frusta}%
\vorto{frusta} Di qua la reliefo esas kruda, grosiera, vulgara\lingui{FIS}
\artiklo{frustrar}%
\vorto{frustrar}\fako{trans., de} Privacar (ulu) de havo, de avantajo quan lu havas la yuro posedar\lingui{FIS}
\artiklo{ftizio}%
\vorto{ftizio} Deperiso pro morbo\lingui{DEFIS}
\artiklo{fudro}%
\vorto{fudro} Tre granda barelo (de 800 litri e mem plu multe)\lingui{DF}
\artiklo{fugar}%
\vorto{fugar}\senco{1}{\textit{(netrans.)} Forirar hastoze por evitar (ulu, ulo)}\senco{2}{\textit{(trans.)} Probar evitor (persono, kozo) per forirar de lu}\lingui{FIR}
\artiklo{fugaso}%
\vorto{fugaso}\fako{milit.} Min-truo provizora quan onu facas dum ula *asieji, por destruktar exploze peco de remparo-murego, e c.
\artiklo{fuino}%
\vorto{fuino}\fako{zool.} Mamifero karn-avida, de la genero « martri », di qua la korpo esas tenua e la muzelo longa\latina{mustela foina}\lingui{FISL}
\artiklo{fuko}%
\vorto{fuko}\fako{bot.} Algi marala quin la ondi jetas sur la litoro\latina{fucus}\lingui{EFIS}
\artiklo{fular}%
\vorto{fular}\fako{trans.}\senco{1}{Presar iteroze per la pedi, la manui, rolero, e c. (kom ex. : drapo) por igar lu plu densa}\senco{2}{Marchar sur ulo}\lingui{EFI}
\artiklo{fulardo}%
\vorto{fulardo} Silka texuro lejera\lingui{DEFR}
\artiklo{fulgurar}%
\vorto{fulgurar}\fako{netrans.} Produktar la lumesko elektrala qua eventas en la alta regioni di atmosfero, sen akompaneso da tondro
\artiklo{fulgurito}%
\vorto{fulgurito} Vitriguro da fulmino qua trapasas sablo-strato\lingui{DEFIRS}
\artiklo{fuligino}%
\vorto{fuligino} Materio nigra qua devenas de la kombusto nekompleta di materii organika, quan la fumuro depozas sur la surfaco di korpo kun qua lu kontaktas\lingui{EFIS}
\artiklo{fulko}%
\vorto{fulko}\fako{zool.} Marsh-ucelo, de la klaso « steltieri », analoga a la galinelo (gallinula chloropus)\latina{fulica}
\artiklo{fulminar}%
\vorto{fulminar}\fako{trans.} Frapar per la fairo elektrala qua explozas en spaco, aspekte di sulko lumoza kun tondro, ofte konsiderata (che la Antiqui e la Moderni) kom cielo-fairo lansita da deo iracoza\lingui{EFIS}
\artiklo{fumar}%
\vorto{fumar}\senco{1}{\textit{(trans.)} \textit{(aludante korpo kombustata)} Emisar gasatrajo ek gaso, aquo-vaporo partikuli plu o min tenua (cindro, fuligino, e c.), grizatra o nigra, akra, plu o min densa}\senco{2}{\textit{(trans.)} Kombustar per aspirar la fumuro (ex. tabako). (*smokar)}\lingui{EFIS}
\artiklo{fumario}%
\vorto{fumario}\fako{bot.} Genero de planti, tipo di la familio « fumariacei », planto farmaciala, sen-pila, kun flori purpurea e di qua la suko bitra uzesas kom kontrefebro urinifigiva\lingui{EFIS}
\artiklo{fumarolo}%
\vorto{fumarolo}\fako{geol.} Krevisuro di sulo volkanala, de qua eskapas fumuro, vapori\lingui{DEFIS}
\artiklo{funcionar}%
\vorto{funcionar} \textit{(aludante organo, aparato, mashino, administrantaro)} Agar to quo resortisas a lu\lingui{DEF}
\artiklo{fundamento}%
\vorto{fundamento}\senco{1}{Masonuro quan onu establisas, en sulo, kom bazo solida di edifico}\senco{2}{\textit{(metaf.)} En ensemblo, la elemento esencala sur qua su apogas lo cetera. – To sur quo onu apogas sua judiko}\lingui{DEFIRS}
\artiklo{fundo}%
\vorto{fundo}\senco{1}{La loko maxim infra en to quo esas kava}\senco{2}{\textit{(metaf.)} La parto interna, nevidebla, opoze a lo videbla extera}\senco{3}{To sur quo ulo su apogas}\senco{4}{La loko maxim fora (en lo konsiderata.)}\lingui{FIS}
\artiklo{funelo}%
\vorto{funelo}\senco{1}{Implemento ek tubo an qua esas muntita la somito di la kono qua formacas lua parto supra, per qua onu varsas liquido aden barelo, e c.}\senco{2}{La parto extera di la orelo}\lingui{E}
\artiklo{funero}%
\vorto{funero} Ensemblo de la ceremonii qui relatas la sepulto di homo jus-mortinta\lingui{DEFIS}
\artiklo{funesta}%
\vorto{funesta} Qua adportas desfortuno\lingui{FIS}
\artiklo{fungo}%
\vorto{fungo}\fako{bot.} Planto kriptogama qua, en medio apta, kreskas e su *reproduktas rapide, e di qua ula speci esas manjebla, la ceteri esante venenoza\lingui{EIL}
\artiklo{funikulara}%
\vorto{funikulara}\senco{1}{\textit{(aludante relvoyo)} Di qua la vagoni, por acensar rampo grand-angula, movesas per kordego spulita sur vindilo}\senco{2}{\textit{(matem.)} Qualeso di la kurvo quan produktas kordo flexebla, ne-extensebla, suspendita per lua extremaji a du punti fixa ed abandonita a la efiko di la gravito; qualeso di la poligono quan formacas kordo sur la punti diversa di qua aplikesas forci}\lingui{DEFIS}
\artiklo{fuorto}%
\vorto{fuorto}\fako{milit.} Laboruro de tero o de masonuro por protektar urbo, paseyo, defileo, portuo, e c.\lingui{DEFIRS}
\artiklo{fureto}%
\vorto{fureto}\fako{zool.} Mikra mamifero karnivora, de la genero « martri », quan onu povas dresar por chasar la kunikli, en la tertruo di qua onu enirigas lu\lingui{EFI}
\artiklo{furgono}%
\vorto{furgono}\fako{armeo} Longa veturo, *kluza (klozita), per qua onu transportas la bagaji, la provizuri de nutrivi\lingui{F}
\artiklo{furiar}%
\vorto{furiar}\fako{netrans.} Kaptesar da iraco dum qua onu ne plus dominacas su. – \textit{(metaf.)} Esar violenta extraordinare\lingui{DEFIRS}
\artiklo{furiero}%
\vorto{furiero}\senco{1}{Oficiro qua, dum voyajo, preiras princo ed agas lo necesa pri la lojeyi}\senco{2}{\textit{(armeo)} Sub-oficiro qua su okupas prokurar la lojeyi necesa por la armeani cirkulanta, repartisar la nutrivi, e c.}\lingui{DFIS}
\artiklo{furnazo}%
\vorto{furnazo} Forno inkandecanta. \textit{(anke metaf.)}
\artiklo{furnelo}%
\vorto{furnelo} Aparato en qua onu facas fairo, por koquar la nutrivi; ulakaze lu kontenas forno\lingui{FIS}
\artiklo{furnisar}%
\vorto{furnisar}\fako{trans.} Livrar sufico-quante, sive vende, sive donace\lingui{EFI}
\artiklo{furo}%
\vorto{furo} Pelo de ula animali, kun lua pili, preparita e fasonita por garnisar od uzesar kom futero di korpo-vesti, kapo-vesti, mufi, e c.\lingui{EFS}
\artiklo{furtar}%
\vorto{furtar}\fako{trans.} Proprigar a su, per ruzo o per violento, to quo savate esas la proprajo, la proprietajo, di altru\lingui{ISL}
\artiklo{furunklo}%
\vorto{furunklo}\fako{patol.} Tumoro inflamala limitoza, kun saliajo en lua centro, qua tempo-finas per pusifo\lingui{DEFIS}
\artiklo{fushar}%
\vorto{fushar}\fako{trans.} Exekutar tasko neglijoze\lingui{D}
\artiklo{fusilo}%
\vorto{fusilo} Paf-armo portebla, ek fer-tubo muntita sur ligna stango krosoza, provizita ye mekanismo di qua la deskuplo acendas la kargajo pulvero di la armo\lingui{FIRS}
\artiklo{fusteno}%
\vorto{fusteno} Stofo di qua la fili esas krucumita – la fili longesala ek lino, la fili transversa ek kotono, uzata por facar subjupi, kamizoli, e c.\lingui{EFIS}
\artiklo{fusto}%
\vorto{fusto}\senco{1}{\textit{(arkitekt.)} Stango di kolono, inter la bazo e la kapitelo}\senco{2}{\textit{(artilerio)} Suportilo, di kanono, di lorneto, di teleskopo, e c.}\lingui{EFIS}
\artiklo{futero}%
\vorto{futero} Stofo ye qua onu garnisas la parto interna di vesti\lingui{DIRS}
\artiklo{futuro}%
\vorto{futuro}\senco{1}{\textit{(gram.)} Tempo di verbo, qua expresas to quo esos}\senco{2}{To quo esas eventonta, existonta}\lingui{DEF}
\artiklo{fuxino}%
\vorto{fuxino}\fako{kemio} Farbo reda, preparita ek anilino\simbolo{C₂₀H₂₀N₃Cl, EH₂O}
\artiklo{fuxio}%
\vorto{fuxio}\fako{bot.} Orno-planto di qua la flori esas pendanta, genero de la familio « enoterei »\latina{fuchsia}
\artiklo{fuzar}%
\vorto{fuzar}\fako{trans.} Liquidigar (korpo solida) per la efiko di kaloro\lingui{EFIS}
\artiklo{fuzelo}%
\vorto{fuzelo}\fako{geom.} Parto di surfaco sferala, inter du duima granda cirkli – di surfaco cilindrala o konala, inter du plani facita segun la axo\lingui{F}
\artiklo{fuzeno}%
\vorto{fuzeno}\fako{bot.} Arboreto di qua la ligno uzesas por fabrikar spindeli lardo-spici, e c., e di qua la rameti karbonigita donas krayoni lejera por skisado\latina{evonymus}\lingui{FIS}
\artiklo{fuzeo}%
\vorto{fuzeo} Pirotekno-kozo inkluzita en kartono (kartocho) qua acendesas stratope, e lansas spricaji ek partikuli kombustata\lingui{F}
\artiklo{fyordo}%
\vorto{fyordo} Singla de la sesguri streta e profunda qui festonizas la litoro westala di Skandinavia e di Skotia\lingui{DEFIRS}
\end{multicols}
\sekciono{G}
\begin{multicols}{2}
\artiklo{g}%
\vorto{g}\fako{muziko} Noto kinesma di la « do »-gamo
\artiklo{gabano}%
\vorto{gabano} Kapoto di navano, vesto de lano, kurta, kun maniki e kapuco kovrita per gudro-telo\lingui{FIS}
\artiklo{gabaro}%
\vorto{gabaro}\fako{navig.} Navo segloza o remiloza, por charjar o descharjar la navi\lingui{DEFIS}
\artiklo{gabiono}%
\vorto{gabiono}\fako{milit.} Granda korbo cilindratra, plenigita ye tero, uzata por protektar la armeani e la laboranti en la tranchei\lingui{EFIRS}
\artiklo{gabro}%
\vorto{gabro}\fako{geol.} Roko granoza, kun plagioklazi qui kontenas egalquante minerali koloroza e minerali senkolora, od en qui iti esas la plu multa\lingui{DEFIS}
\artiklo{gado}%
\vorto{gado}\fako{zool.} Genero de fishi (moruo, merlano, e c.)\latina{gadus}\lingui{FL}
\artiklo{gafo}%
\vorto{gafo}\senco{1}{Stango kun pinto fera, garnisita ye hoko laterala, por pulsar barko, adtirar ulo adsur altra navo, sondar, e c.}\senco{2}{Hoko por adtirar adtere la grosa fishi}\lingui{EFI}
\artiklo{gagato}%
\vorto{gagato} Lignito brilo-nigra quan onu talias a tabuleti, ornamenti traurala\lingui{DL}
\artiklo{gaino}%
\vorto{gaino} Etuyo (di la lamo di instrumento tranchiva od akuta) tre longa, streta, en qua la objekto eniras tra la extremajo\lingui{FI}
\artiklo{gajo}%
\vorto{gajo}\senco{1}{Objekto depozita por garantiar la rimborsoteso di debajo}\senco{2}{En la ludi mondumala, objekto depozita da la ludanto qua eroris e quan lu povos retrohavar nur per submisar su a « puniso »}\lingui{EF}
\artiklo{gala-}%
\vorto{gala-} Prefixo qua signifikas « kun la pompo di festo »\lingui{DEFIS}
\artiklo{galanta}%
\vorto{galanta} Zeloza pri plezar a la mulieri\lingui{DEFIRS}
\artiklo{galanteno}%
\vorto{galanteno} Disho ek karno de pultro, de bovyuno, e c., koquita kun spici, quan onu servas kolda kun jeleo\lingui{DEFIRS}
\artiklo{galbano}%
\vorto{galbano}\fako{farmacio} Gum-rezino balzamoza, devenanta de arboro di Siria\lingui{DEFIRS}
\artiklo{galeno}%
\vorto{galeno}\fako{kemio} Plombo-sulfuro naturala\simbolo{PbS}\lingui{DEFIS}
\artiklo{galeo}%
\vorto{galeo}\fako{imprim-arto} Tabuleto rektangula e rebordoza, sur qua la kompostisto pozas la lineedi asemblita en la kompostilo\lingui{F}
\artiklo{galerio}%
\vorto{galerio}\senco{1}{Spaco tektoza cirkum edifico, apartmento, chambrego, o longesale, uzata kom promeneyo, paseyo}\senco{2}{Quaza balkono salianta qua quaze kronizas la dopa parto di navo}\lingui{DEFIRS}
\artiklo{galero}%
\vorto{galero}\senco{1}{\textit{(che la antiqui)} Milito-navo kun un o plura rangi de remili e segloza}\senco{2}{Milito-navo ramoza, ferdekoza, kun du masti e du segli Latina, uzita precipue en Mediteraneo}\lingui{DEFIS}
\artiklo{galeto}%
\vorto{galeto} Sorto di kuko, ordinare ronda e plata, ek pasto kompakta e foliatra, aplastita per rolero, e bakita\lingui{FS}
\artiklo{galicismo}%
\vorto{galicismo} Idiotismo di la linguo Franca\lingui{DEFIRS}
\artiklo{galikana}%
\vorto{galikana}\fako{historio religiala} Qua koncernas la eklezio katolika di Francia, egardata kom nedependanta (ula-punte) de la autoritato di la papo\lingui{DEFIS}
\artiklo{galimatiaso}%
\vorto{galimatiaso} Diskurso, skriburo, qua esas korekta gramatikale, ma konfuza ed obskura de la vidopunto « stilo »\lingui{DEFS}
\artiklo{galinelo}%
\vorto{galinelo}\fako{zool.} Speco di ucelo ek la « steltieri »\latina{gallinula chloropus}\lingui{EFISL}
\artiklo{galio}%
\vorto{galio}\fako{kemio} Metalo analoga a zinko, blanka, harda, poke forjebla\simbolo{Ga}\lingui{DEF}
\artiklo{galiono}%
\vorto{galiono}\fako{navig.} Charjo-navo uzita precipue olim da la Hispani por lia komerco kun Amerika, e specale por querar la oro di Peru\lingui{DEFIS}
\artiklo{galioto}%
\vorto{galioto}\senco{1}{Mikra galero remiloza e segloza}\senco{2}{Traverso ligna qua suportas la paneli qui klozas luko}\lingui{DEFIS}
\artiklo{galo}%
\vorto{galo}\fako{bot.} Surkreskajo produktita sur la folii, sur la stipi di la vejetanti (pro la ekfluo di sapo) da la pikuro da ula insekti qui venas pondar ibe sua ovi\latina{cynips}\lingui{DEFIRS}
\artiklo{galono}%
\vorto{galono} Texajo (lana, silka, ora), rubandatra, per qua onu ornas la uniformi, la kostumi, la chapeli, la kurteni, la mobli\lingui{DEFIRS}
\artiklo{galopar}%
\vorto{galopar}\fako{netrans.) (pri kavalo} Agar serio de salti longesala qui exekutesas sucede en tri tempo-parti, pozante unesme un ek la gambi dopa, pose un ek la gambi avana e la altra gambo dopa, fine la gambo avana qua restas – por obtenar la irorapideso maxima che kavalo\lingui{DEFIRS}
\artiklo{galosho}%
\vorto{galosho} Pedovesto, di qua la parto supra esas ledra, e la suolo ligna, uzata por prezervar la pedi de la humideso mediala\lingui{DEFIRS}
\artiklo{galvanala}%
\vorto{galvanala}\fako{fiziko} Qua relatas la elektro qua naskas de la kontakto di du korpi heterogena\lingui{DEFIRS}
\artiklo{galvanizar}%
\vorto{galvanizar}\fako{trans.} Igar (muskuli inerta) kontraktar per la elektro quan produktas en li la kontakto di du korpi heterogena\lingui{DEFIRS}
\artiklo{galvanometro}%
\vorto{galvanometro}\fako{fiziko} Instrumento qua uzesas por mezurar la intenseso-grado di la fluo galvanala\lingui{DEFIRS}
\artiklo{galvanoplastar}%
\vorto{galvanoplastar}\fako{trans.} Aplikar strato de oro, arjento, zinko, e c., per la efiko di baterio galvanala\lingui{DEFIRS}
\artiklo{galvanotropismo}%
\vorto{galvanotropismo}\fako{biol.} Efiko nemediata, da la elektro (fluo kontinua), a la protoplasmo di ula celuli e, partikulare, di la protozoi\lingui{DEFIRS}
\artiklo{gambitar}%
\vorto{gambitar}\fako{trans., ulu} Pasar sua gambo cirkum olta di persono, por renversar ica\lingui{FI}
\artiklo{gambo}%
\vorto{gambo}\senco{1}{\textit{(anat.)} Parto di la membro infra di la homo, qua prolongas la kruro, de la genuo til la pedo. – La tota membro}\senco{2}{Che ula quadripedi : singla de la quar membri qui subtenas la korpo e finas per hufi}\senco{3}{Membro di animalo quan lu uzas por marchar, sizar, kaptar}\lingui{FI}
\artiklo{gambolar}%
\vorto{gambolar}\fako{netrans.} Saltar agitante sua gambi\lingui{EF}
\artiklo{gamelo}%
\vorto{gamelo}\senco{1}{Granda skudelo en qua plura navani infra-grada, plura soldati, kunmanjas}\senco{2}{Lada vazo en qua singla soldato recevas sua porciono}\senco{3}{\textit{(milito-navo)} Gasto-tablo komuna por la oficiri, la dicipuli, la kirurgiisti}\lingui{FIS}
\artiklo{gameto}%
\vorto{gameto}\fako{biol.} Celulo reproduktera, maskula o femina
\artiklo{gametofito}%
\vorto{gametofito}\fako{biol.} Gameto di la planti, en la periodo en qua la du pronuklei maskula e femina ne ja esas aglutinita (sporofito) o kunfuzita (maxio)
\artiklo{gamo}%
\vorto{gamo}\fako{muziko} Serio de la sep noti di la skalo muzikala, dispozita segun lia sucedo naturala, acensanta o decensanta, per toni e mi-toni, intervale di un oktavo\lingui{EFIRS}
\artiklo{ganar}%
\vorto{ganar}\fako{trans.} Aquirar (profito-quanto)\lingui{EFIS}
\artiklo{gangliono}%
\vorto{gangliono}\fako{anat.} Organo globatra ek tubero de fibri nervala o de faskuli limfala\lingui{DEFIRS}
\artiklo{gango}%
\vorto{gango}\senco{1}{\textit{(teknol.)} Materio stranjera solida en qua kontenesas erco o lapido}\senco{2}{\textit{(anat.)} Substanco amorfa qua envelopas elemento anatomiala}\lingui{DEFIS}
\artiklo{gangrenar}%
\vorto{gangrenar}\fako{netrans.) (aludante la tisui animalala} Dissolvesar putrante. \textit{(anke metafore)}\lingui{DEFIRS}
\artiklo{ganso}%
\vorto{ganso}\fako{zool.} Ucelo palmipeda, de la sama familio kam olta di la anado\latina{anser}\lingui{DE}
\artiklo{ganto}%
\vorto{ganto} Envelopilo de felo o de texajo lina, lana, silka, e c., qua kovras la palmi di la manuo ed omna fingri single, tale ke la manui protektesas sen cesar esar libera\lingui{EFIS}
\artiklo{gapar}%
\vorto{gapar}\fako{netrans.} Cirkular en stradi, haltante (pro kuriozeso) mem avan kozi qui ne meritas regardeso\lingui{DE}
\artiklo{garantiar}%
\vorto{garantiar}\fako{trans.} *Certenigar ulu pri (ulo) sua-response; certifikar sua-response\lingui{DEFIRS}
\artiklo{garar}%
\vorto{garar}\fako{trans.}\senco{1}{Lokizar (navo) shirme en doko}\senco{2}{Venigar (treno) shirme en loko adaptita alonge la relvoyo normala}\senco{3}{Lokizar (en hangaro) veturo protektenda kontre la domaji veterala}\lingui{DEFIS}
\artiklo{garbanzo}%
\vorto{garbanzo}\fako{bot.} Nomo di varietato di pizi\latina{cicer}\lingui{S}
\artiklo{garbo}%
\vorto{garbo} Fasko de spiki falchita, en qua omna kapi dispozesas samapunte, qua plularjeskas inter la pedo e la kapo. \textit{(anke metafore, pri flori, pirotekno, aquo-sprico)}\lingui{DF}
\artiklo{gardar}%
\vorto{gardar}\fako{trans.} Surveyar por ke ulu, ulo, ne departez, ne egaresez, ne perdesez, ne domajesez\lingui{EFIS}
\artiklo{gardenio}%
\vorto{gardenio}\fako{bot.} Arboreto de la familio « rubiacei »\latina{gardenia}
\artiklo{gardeno}%
\vorto{gardeno} Tereno sur qua esas plantacita vejetanti utila o plezurala\lingui{DEF}
\artiklo{gardio}%
\vorto{gardio}\fako{zool.} Mikra fisho qua vivas en la aquo sen-sala, e di qua la karno tre manjo-prizesas\latina{cyprinus rutilus}\lingui{FS}
\artiklo{gareno}%
\vorto{gareno}\senco{1}{Loko rezervita por la pesko, la chaso}\senco{2}{Loko cirkondata da muri, fosati, treliso, ube vivas multa kunikli}\senco{3}{Bosko, erikeyo, ube la kunikli abundas}\lingui{FI}
\artiklo{gargarar}%
\vorto{gargarar}\fako{trans.} Humidetigar (la enireyo di la guturo) per liquido quan onu lasas ibe sejornar dum kelka instanti, agitante ica per aspiro-movo\lingui{DEFIS}
\artiklo{garito}%
\vorto{garito} Lojeyeto de ligno o de petro ube sentinelo shirmas su kontre la domaji veterala\lingui{FIS}
\artiklo{garnisar}%
\vorto{garnisar}\fako{trans.}\senco{1}{Cirkondar per ulo qua protektas}\senco{2}{Kompletigar (kozo) per pozar en lu to quon lu destinesis kontenor}\senco{3}{Kompletigar (kozo) per adjuntar ula parti kom acesora, orniva}\lingui{DEFIRS}
\artiklo{garnituro}%
\vorto{garnituro} Ensemblo de la acesoraji qui kompletigas kozo, od ornas, dekoras lu\lingui{DEFIR}
\artiklo{garnizono}%
\vorto{garnizono} Ensemblo de la armeani qui okupas fortifikajo por defensar lu\lingui{DEFIRS}
\artiklo{garoto}%
\vorto{garoto}\senco{1}{Ligno-peco quan onu pasigas en kordo por klemar ica per tordo}\senco{2}{Instrumento quan uzas la kirurgiisto por kompresar arterio}\senco{3}{\textit{(olim, en Hispania)} La fer-koliaro per qua onu tormente strangulis la karcerani}\lingui{EFS}
\artiklo{garsono}%
\vorto{garsono} Employato infra-klasa qua servas en balneyo, kafe-drinkerio, restorerio, kontoro, apartmento\lingui{FI}
\artiklo{gartero}%
\vorto{gartero} Kordono, rubando, e c., per qua onu fixigas la kalzi an-super od an-sub la genuo\lingui{EFS}
\artiklo{gaso}%
\vorto{gaso}\fako{kemio} Fluido aeroforma (ye la temperaturo ordinara)\lingui{DEFIRS}
\artiklo{gasometro}%
\vorto{gasometro}\senco{1}{Aparato per qua onu mezuras la volumino di gaso}\senco{2}{Aparato per qua onu regulas la ekfluo di lumizo-gaso, mezurante la quanto furnisata}
\artiklo{gasteropodo}%
\vorto{gasteropodo}\fako{zool.} Klaso de moluski, di qui la ventro esas diskoida e sur qua li reptas o su tranas\lingui{DEF}
\artiklo{gasto}%
\vorto{gasto}\senco{1}{\textit{(antique)} La persono qua profitis la yuro reciproka pri lojar, kande onu voyajis, iti che ici, yuro aplikita segun konvencioni inter privati, familii, urbi}\senco{2}{La persono qua, kom o quale amiko, lojesas e nutresas gratuite}\lingui{DER}
\artiklo{gastralgio}%
\vorto{gastralgio}\fako{patol.} Morbo inflamala di la stomako\lingui{DEFIS}
\artiklo{gastronomio}%
\vorto{gastronomio} Quaza arto di la gastronomo\lingui{DEFIRS}
\artiklo{gastronomo}%
\vorto{gastronomo} Persono qua havas gusto subtila e savas dicernar la qualesi di la manjaji e drinkaji\lingui{DEFIRS}
\artiklo{gastrula}%
\vorto{« gastrula »}\fako{biol.} Fazo duesma di la evoluciono di la ovo, qua donas formo larva ad omna metazoi
\artiklo{gaujar}%
\vorto{gaujar}\fako{trans.} Mezurar la kapaceso-grado\lingui{EF}
\artiklo{gavar}%
\vorto{gavar}\fako{trans.} Saturar ye nutrivi (la pultruni quin onu volas repletigar)\lingui{F}
\artiklo{gavialo}%
\vorto{gavialo}\fako{zool.} Speco di krokodilo, di qua la muzelo esas longa, qua, en la regiono di la fluvio Gange (India), su nutras ye fishi\latina{gavialis gangeticus}
\artiklo{gavoto}%
\vorto{gavoto} Danso du-tempa, kun du iteri de 4 ad 8 tempi
\artiklo{gaya}%
\vorto{gaya} Di qua la humoro esas ridemeso\lingui{EFIS}
\artiklo{gayako}%
\vorto{gayako}\fako{bot.} Arboro exotika, di qua la ligno esas extreme harda e rezinoza\latina{guiacum}\lingui{DEFIRS}
\artiklo{gazelo}%
\vorto{gazelo}\fako{zool.} Speco di antilopo qua habitas Afrika ed Azia\latina{antilope dorca}\lingui{DEFIRS}
\artiklo{gazo}%
\vorto{gazo} Stofo ek texajo lejera e diafana\lingui{DEFIRS}
\artiklo{gazolino}%
\vorto{gazolino}\fako{kemio} Petrol-etero, la elemento maxim volatila di la petrolo kruda
\artiklo{gazono}%
\vorto{gazono}\fako{bot.} Herbo tre tenua ek graminei diversa, qua, sur sulo, formacas quaza tapiso de verdajo\lingui{F}
\artiklo{ge-}%
\vorto{ge-} Prefixo qua signifikas « maskulo e femino, samasorta egardata pare »
\artiklo{geblo}%
\vorto{geblo}\fako{arkitekt.} Parto supra di la muro, qua finas per pinto qua portas la extremajo di la kolmo-trabo di la tekto-karpenturo\lingui{DE}
\artiklo{geheno}%
\vorto{geheno}\fako{biblo} Inferno\lingui{DEFIRS}
\artiklo{geisha}%
\vorto{« geisha »} Dansistino e kantistino Japoniana
\artiklo{gelatino}%
\vorto{gelatino} Substanco azotoza quan onu extraktas de la osti, de la kartilagi di la animali, e qua formacas, kun aquo, jeleo\lingui{DEFIRS}
\artiklo{gelto}%
\vorto{gelto} Parto proporcionala de la pekunio qua rezultas de la vendo di la vari, grantata a sua komizi (ultre la salario-quanto fixa di ici) da ula butikegi\lingui{DF}
\artiklo{gemo}%
\vorto{gemo} Lapido tre precoza\lingui{DEFIS}
\artiklo{gemulo}%
\vorto{gemulo}\fako{bot.}\senco{1}{Rudimento di la stipo, qua aparas lor la jermifo}\senco{2}{Mikra korpo ovoida qua, separante su de la planto, destinesas a *reproduktar ica}
\artiklo{genciano}%
\vorto{genciano}\fako{bot.} Planto de la familio « genciani », di qua la suko esas bitra, qua kreskas en la monti\latina{gentiana}\lingui{DEFIRS}
\artiklo{genealogio}%
\vorto{genealogio} Decendo-lineo pri un o plura personi, quan produktas la suceso da lua o lia ancestri\lingui{DEFIRS}
\artiklo{generaciono}%
\vorto{generaciono} La homi qui vivas en la sama periodo (evaluata segun la meza duro di la homo-vivo)\lingui{DEFIRS}
\artiklo{generalisimo}%
\vorto{generalisimo}\fako{armeo} La generalo qua esas la chefo di omna generali cetera\lingui{DEFIS}
\artiklo{generalo}%
\vorto{generalo}\fako{armeo} Milito-chefo qua komandas armeo, diviziono o brigado\lingui{DEFIRS}
\artiklo{generatoro}%
\vorto{generatoro}\fako{tekn.} Termino teknikala, kelke sinonima ye « kaldiero », qua tante plu uzesas, quante plu la tipi difereskas de la tanki cilindra, tre kapaca, quin onu nomizas « kaldieri »; la generatori esas generale aquo-tuboza e rezistas a presi normala ye 80 kilogrami sur omna centimetro quadrata\lingui{DEF}
\artiklo{genero}%
\vorto{genero}\senco{1}{Grupo naturala de enti qui intersimilesas pri ula traiti esencala}\senco{2}{Ideo generala sub qua la spirito asemblas to quo esas komuna ye termini diversa, per eliminar la difero-punti}\senco{3}{Ensemblo de la karakteri esencala di kozo}\lingui{DEFIS}
\artiklo{genezo}%
\vorto{genezo} Produktado di la enti organizita\lingui{DEFIS}
\artiklo{genio}%
\vorto{genio}\senco{1}{Ento supernatura, benigna o maligna, qua (krede da la pagani) direktis la fato di la individui, familii, nacioni}\senco{2}{Ento supernatura qua (en la fe-rakonti) dotesis ye povo magiala}\senco{3}{Ento alegoriala qua enkorpigas (en poezio, skultarto, piktarto) ideo abstraktita}\senco{4}{Talento naturala}\senco{5}{Kapableso supera quan Naturo grantis a la spiriti kreera}\lingui{DEFIRS}
\artiklo{genitar}%
\vorto{genitar}\senco{1}{\textit{(trans.)} Produktar per la fekundigo di la femino}\senco{2}{\textit{(geom.)} Produktar per movar su (ex. : « Sfero genitesas per la rotaco di mi-cirklo cirkum sua diametro »)}\lingui{DEFIRS}
\artiklo{genitivo}%
\vorto{genitivo}\fako{gram.} Kazo di vorto deklinebla, qua expresas lua relato ad altra vorto qua indikas kozo qua apartenas ad ita\lingui{DEFIS}
\artiklo{genitori}%
\vorto{genitori} Gepatri
\artiklo{genro}%
\vorto{genro}\fako{gram.} Formo specala, uzata da ula lingui, por distingar la enti maskula, la enti femina e la kozi sensexua
\artiklo{gento}%
\vorto{gento}\senco{1}{\textit{(che la Romani antiqua)} Familio o grupo de familii qua havis origino komuna, patrica}\senco{2}{La sama-rasani}\lingui{L}
\artiklo{genuo}%
\vorto{genuo}\fako{anat.} Artiko di la gambo e di la kruro, ye la parto avana\lingui{EFI}
\artiklo{geocentrala}%
\vorto{geocentrala}\fako{astron.} Mezurata relate Tero quan onu konsideras kom centro di la cielo-sfero\lingui{DEF}
\artiklo{geodezio}%
\vorto{geodezio} Cienco pri la mezuro di la Terglobo o di lua parti\lingui{DEFIRS}
\artiklo{geognozio}%
\vorto{geognozio} Cienco pri la kompozeso mineralogiala di la Terglobo
\artiklo{geografio}%
\vorto{geografio} Cienco di qua la studiajo esas la deskripto di la Teresurfaco\lingui{DEFIRS}
\artiklo{geografo}%
\vorto{geografo} La persono qua su okupas pri geografio\lingui{DEFIRS}
\artiklo{geologio}%
\vorto{geologio} Cienco qua havas kom studiajo la strati ye qui konsistas la Ter-kortico, lia estado nuna e lia formaceso\lingui{DEFIRS}
\artiklo{geologo}%
\vorto{geologo} La persono qua su okupas pri geologio\lingui{DEFIRS}
\artiklo{geomanciar}%
\vorto{geomanciar}\fako{netrans.} Divinar lo futura segun la figuri quin formacas manuedo de tero jetita sen-vice\lingui{DEF}
\artiklo{geometro}%
\vorto{geometro} Persono experta pri geometrio\lingui{DEFIRS}
\artiklo{georgismo}%
\vorto{georgismo}\fako{ekonomiko} Doktrino ekonomikala (da Henry George) qua kontestas la yuro a la proprieto privata di la Ter-peci, e propozas, kom moyeno praktikala por remediar la nuna estado qua naskis de ta yuro, taxar la valoro di la sulo-peci « nuda », t.e. sen inkluzar en la evaluo la valoro di la pluboniguri da la proprietero koncernata pro ke oli rezultas de lua laboro, kontre ke la sulo-peci « nuda » facesis da Naturo ipsa – ed impostar ta valoro – komence ye kelka-\%, e pokope plumulta-\%, til atingar 100\%. Lo posibligus la tota supreso di omna imposti cetera\lingui{DEFIRS}
\artiklo{geotaktismo}%
\vorto{geotaktismo}\fako{biol.} Movi da ula protoplasmi (ex. : la plasmodo di la mixomiceti) kande li influesas da gravito
\artiklo{geotropa}%
\vorto{geotropa}\fako{biol.} Qua koncernas la proprajo di ula korpi (precipue la radiki e la stipi) kreskar uladirecione pro influeso da gravito
\artiklo{gepardo}%
\vorto{gepardo}\fako{zool.} Animalo karnavida, qua vivas en India, de la genero « kati », kun krinaro, e di qua la ungli esas ne-retrotirebla\latina{cynailurus}\lingui{DFRS}
\artiklo{geranio}%
\vorto{geranio}\fako{bot.} Gardeno-planto, di qua la folii esas rondatra e la flori umbela\latina{geranium}
\artiklo{gerfalko}%
\vorto{gerfalko}\fako{zool.} Rapt-ucelo, quaza falko, granda-statura, audacoza e tre ajila\latina{falco gyrofalco}\lingui{DEFIS}
\artiklo{germanio}%
\vorto{germanio}\fako{kemio} Metalo griza, metal-briloza, volatila kande varmigita til redegeso\simbolo{Ge}\lingui{DEF}
\artiklo{ges}%
\vorto{« ges »} Noto kinesma di la « do »-gamo, bemoloza
\artiklo{gestar}%
\vorto{gestar}\fako{netrans.} Movar brakio, manuo, kapo, e c., por expresar ula pensi, ula sentimenti, e por igar plu expresiva la dicatajo\lingui{DEFIRS}
\artiklo{getro}%
\vorto{getro} Envelopilo drapa, ledra, klozebla per bukli o per butoni super la shuo e la infra parto di la gambo\lingui{EF}
\artiklo{geyser}%
\vorto{« geyser »}\fako{geol.} Fonto di aquo naturale bolianta qua spricas tre alte\lingui{DEF}
\artiklo{ghetto}%
\vorto{« ghetto »} Olim, en Italia : quartero ube la Judi koaktesis rezidar\lingui{DEF}
\artiklo{gibelino}%
\vorto{gibelino}\fako{historio di Italia} Dum la Mezepoko, partisano di la imperiestro Germana, opoze a la guelfi, partisani di la papo\lingui{DEFIRS}
\artiklo{gibo}%
\vorto{gibo}\senco{1}{Che la animali : protuberanco dorsala naturala, ek amaso grasa}\senco{2}{Che la homo : protuberanco produktita da deviaceso di la spino o da salio di la sternumo}\senco{3}{\textit{Anat}. Saliajo mi-sferatra di ula osti (kraniala, frontala, e c.)}\senco{4}{\textit{(pri tereno)} Neplatajo di la surfaco}\lingui{EFS}
\artiklo{gibono}%
\vorto{gibono} Granda simio di India\lingui{DEF}
\artiklo{gicheto}%
\vorto{gicheto} Mikra aperturo tra qua onu komunikas kun la oficieri en la kontori postala, kun la kasieri en la firmi banka e komercala\lingui{F}
\artiklo{giganto}%
\vorto{giganto} Persono di qua la staturo esas extraordinare longa\lingui{DEFIRS}
\artiklo{gildo}%
\vorto{gildo} Ensemblo de sama-speca mestieristi\lingui{DER}
\artiklo{gilochar}%
\vorto{gilochar}\fako{trans.) (tekn.} Ornar per interplekturo de traiti grabita\lingui{DEFR}
\artiklo{gilotino}%
\vorto{gilotino} Instrumento por tranchar la kapo di la personi kondamnita ad ocideso legala – ek lamo tre pezoza qua, glitante inter du fosti, falas bruske sur la kolo di la ocidoto sternita sub olu\lingui{DEFIRS}
\artiklo{gimbardo}%
\vorto{gimbardo}\fako{muziko} Muzik-instrumento, ek mi-cirklo stala o latuna qua finas per du branchi paralela, inter qui esas langeto stala quan onu agitas per la fingro dum ke onu tenas la instrumento inter la maxilo e la mandibulo
\artiklo{gimnastikar}%
\vorto{gimnastikar}\fako{netrans.}\senco{1}{Exekutar la exerco-movi apta plufortigar e plu-flexebligar la korpo}\senco{2}{Exekutar la exercotaski apta plu-fortigar e plu-flexebligar la fakultati pensala}\lingui{DEFIRS}
\artiklo{gimnazio}%
\vorto{gimnazio} Nomo donata, en Germania, a la kolegii ube onu edukesas klasikale\lingui{D}
\artiklo{gimnospermo}%
\vorto{gimnospermo} Di qua la semini aspektas quaze « nuda »\lingui{DEFIS}
\artiklo{gimnoto}%
\vorto{gimnoto}\fako{zool.} Fisho ek la regioni tropikala di Amerika, sen floso dorsala, e qua havas, an omna flanki di la kaudo, lami membrana qui emisas elektro\latina{gymnotua electricus}\lingui{DEFIS}
\artiklo{ginandra}%
\vorto{ginandra}\fako{bot.} Di qua la stamini insertesas sur la pistilo\lingui{DEFIRS}
\artiklo{gineceo}%
\vorto{gineceo}\senco{1}{\textit{(antique)} Parto de la domo qua rezervesis a la homini}\senco{2}{Loko ube permanis ordinare la homini, la mulieri, ube li laboris}\lingui{DEFIRS}
\artiklo{ginesto}%
\vorto{ginesto}\fako{bot.} Arboreto, di qua la floro esas flava, de la familio « leguminosi »\latina{genista}\lingui{DFIS}
\artiklo{ginodioika}%
\vorto{ginodioika}\fako{bot.} Planto qua havas flori hermafrodita e florini sur du stipi ne-sama
\artiklo{gipaeto}%
\vorto{gipaeto}\fako{zool.} La maxim granda ek la raptuceli di Europa\latina{gypaetus barbatus}\lingui{EFISL}
\artiklo{gipso}%
\vorto{gipso}\fako{kemio}\senco{1}{Kali-sulfato, maxim-multa-kaze hidratoza}\senco{2}{Ta materio, kalcinite e pulverigite, uzata por masono : onu humidigas lu por igar lu pasta, e, kande ri-sika, lu ri-esas tam harda kam ante kalcineso}\simbolo{SO₄Ca 2H₂O}\lingui{DEFIRS}
\artiklo{gipuro}%
\vorto{gipuro} Dentelo ek filo e tordita, di qua la desegnuri interplektita formacas la texajo, sen reto distinta da la ornivi\lingui{DEFIR}
\artiklo{girlando}%
\vorto{girlando} Kateno de flori, de folii tresita, quan onu suspendas kom ornivo\lingui{DEFIRS}
\artiklo{girometro}%
\vorto{girometro} Aparato uzata por mezurar la rotaco-rapideso di la mashini\lingui{DEFIRS}
\artiklo{giroskopo}%
\vorto{giroskopo}\fako{fiziko} Aparato por demonstrar la rotaco jornala di Tero cirkum sua axo, per la deviaceso, relate punti fixa selektita sur la Terglobo, di korpo libere suspendita per lua baricentro e qua rotacas cirkum axo\lingui{DEFIS}
\artiklo{gis}%
\vorto{« gis »} Noto kinesma di la « do »-gamo, diezoza
\artiklo{gisar}%
\vorto{gisar}\fako{trans.} Fabrikar (ulo : klosho, e c.) per substanco fuzanta\lingui{DI}
\artiklo{gitago}%
\vorto{gitago}\fako{bot.} Nomo di mesono-kariofilaceo di qua la semini, se mixita hazarde a la frumento, igas pano saporar bitre ed aspektar nigra\latina{agrostemma githago; nigella}
\artiklo{gitaro}%
\vorto{gitaro}\fako{muziko} Instrumento analoga a violino, kun sis kordi, quan onu pleas per pincar ici per la manuo dextra, e kun mancho dividita a mi-toni per mikra klavi qui indikas la loko ube esas pozenda la fingri di la manuo sinistra\lingui{DEFIRS}
\artiklo{gizardo}%
\vorto{gizardo}\fako{zool.} Stomako duesma di la uceli, qua venas dop la kropo, ed en qua la nutrivi trituresas e reduktesas a pasto\lingui{E}
\artiklo{glacar}%
\vorto{glacar}\fako{trans.} Indutar per verniso glata, diafana, qua aspektas quale glacio\lingui{DEFR}
\artiklo{glaciero}%
\vorto{glaciero} Amaso permananta de glacio, qua formacesas en la vali grand-altituda di la monti\lingui{DEFS}
\artiklo{glacio}%
\vorto{glacio}\senco{1}{Aquo konjelita da frigoro}\senco{2}{Kremo o siropo konjelita en glacio}\senco{3}{\textit{(metaf.)} Neardoroza extreme}\lingui{EFIS}
\artiklo{glacizo}%
\vorto{glacizo}\fako{milit.} Pento mikr-angula qua, en fortifikajo, abutas a la ruro cirkondanta\lingui{DEFRS}
\artiklo{gladiatoro}%
\vorto{gladiatoro}\fako{en la Roma antiqua} Viro quan onu igis kombatar, en cirko, kontre altra viri o kontre bestii feroca, por amuzar la plebeyo spektanta\lingui{DEFIRS}
\artiklo{gladiolo}%
\vorto{gladiolo}\fako{bot.} Planto analoga ad irido\latina{gladiolus}\lingui{DEFIL}
\artiklo{glando}%
\vorto{glando}\fako{anat.} Organo ek mikra utrikuli, qua sekrecas ula liquidi aden la korpo animalala\lingui{DEFIS}
\artiklo{glano}%
\vorto{glano}\senco{1}{\textit{(bot.)} Frukto di la querko, di qua la perikarpo, unionita a la semino, kovresas da involukro squamoza}\senco{2}{\textit{(anat.)} Extremajo di la peniso, di la klitorido}\senco{3}{Tufo de silko, lano, e c. uzata kom ornivo en la pasmenti}\lingui{DEFIS}
\artiklo{glaso}%
\vorto{glaso} La vazo ordinara, pedoza, (generale ek vitro), en qua onu drinkas vino od aquo\lingui{DE}
\artiklo{glata}%
\vorto{glata} Qua ne sentigas asperaji kande onu igas lu glitar sur sua pelo\lingui{DR}
\artiklo{glaukomo}%
\vorto{glaukomo}\fako{patol.} Okulo-morbo, pro qua la humoro vitratra divenas opaka, e la okul-fundo verdatra\lingui{DEF}
\artiklo{glaukonito}%
\vorto{glaukonito}\fako{mineral.} Hidroza silikato de fero, de alumino e de kalio, quaza kreto verdatra\lingui{DEFIS}
\artiklo{glavo}%
\vorto{glavo} Espado tranchiva, kun lamo kurta e rekta, di qua la extremajo esas du-bizela\lingui{FL}
\artiklo{glebo}%
\vorto{glebo} Mikra maso de tero kompakta\lingui{FI}
\artiklo{gleno}%
\vorto{gleno}\fako{anat.} Kavajo poke profunda di osto en qua esas inkastrita altra osto
\artiklo{glezar}%
\vorto{glezar}\fako{trans.} Aplikar (sur terakotajo, ceramikajo) induturo vitrigebla\lingui{DE}
\artiklo{glicerino}%
\vorto{glicerino}\fako{kemio} Korpo neutra, siropatra, quan onu obtenas per la duimigo di grasajo (stearino, oleino, e c.)\lingui{DEFIRS}
\artiklo{glicerio}%
\vorto{glicerio}\fako{bot.} Genero de planti perena, ek la familio « graminei »\latina{glyceria}
\artiklo{glicino}%
\vorto{glicino}\fako{bot.} Arbusto klimera, de la familio « papilionacei »\latina{glycine}\lingui{DEFIS}
\artiklo{glikogena}%
\vorto{glikogena}\fako{fiziol.} Qua produktas sukro
\artiklo{glikogenezo}%
\vorto{glikogenezo}\fako{kemio biologiala} Formaceso di la sukro (ex. : en la hepato)
\artiklo{glikolo}%
\vorto{glikolo}\fako{kemio} Substanco sukroza, kristaligebla, quan onu obtenas per duimigar (per saponigo) la acidi qui kompozas bilo, o per traktar la gelatino per sulfacido\lingui{DEFIRS}
\artiklo{glikoproteido}%
\vorto{glikoproteido}\fako{kemio biologiala} Asociuro molekulara ek proteino e komplexo hidrokarboza
\artiklo{glikosido}%
\vorto{glikosido} Nomo generala di la kompozaji naturala quin onu renkontras che la vejetanti, e qui per hidrolizo, influe da fermenti, produktas glikosi e substanci diversa-natura (acidi, fenoli, alkoholi, e c.)
\artiklo{glikoso}%
\vorto{glikoso}\fako{kemio} Elemento sukroza di la vito, di amilo, e c.\lingui{DEFIRS}
\artiklo{glinar}%
\vorto{glinar}\fako{trans.}\senco{1}{Rekoltar en agro (la spiki qui restas pos la mesono)}\senco{2}{\textit{(metaf.)} Rekoltar to quo lasesis da altru}\lingui{EF}
\artiklo{gliomo}%
\vorto{gliomo}\fako{patol.} Tumoro di la tisuo nervala
\artiklo{gliptar}%
\vorto{gliptar}\fako{trans.} Skultar medalie lapidi, kave (intalii) o reliefe (kamei)\lingui{DEFIS}
\artiklo{gliro}%
\vorto{gliro}\fako{zool.} Mikra mamifero rodera, qua torporas dum vintro\latina{myoxus glis}
\artiklo{glitar}%
\vorto{glitar}\fako{netrans.} Movesar, movar, kontinue sur la surfaco di korpo glata, konseque de pulso\lingui{DEF}
\artiklo{globo}%
\vorto{globo} Korpo sfera o sferoida\lingui{DEFIRS}
\artiklo{globoido}%
\vorto{globoido}\fako{bot.} Korpo globatra od ovatra, qua ulakaze esas inkluzita en grano di aleurono, e qua judikeble konsistas ye glicerofosfato di kalko e di magnezio
\artiklo{globulino}%
\vorto{globulino}\fako{kemio biologiala} Materio albuminoida, qua diferas de la albumini pro esar nedissolvebla en aquo distilita, ma dissolvebla en medii alkali-saloza
\artiklo{globulo}%
\vorto{globulo}\fako{anat.} Korpuskulo plu o min ronda, suspendita en sango\lingui{EFIS}
\artiklo{glorio}%
\vorto{glorio} Brilo de famo, de grandeso\lingui{DEFIRS}
\artiklo{glosario}%
\vorto{glosario}\senco{1}{Vortolibro kontenanta la vorti di qui la senco divenis obskura pro lia obsoleteso, e pri qui expliko esas nekareebla}\senco{2}{Nomenklaturo de la vorti ye qui konsistas linguo}\lingui{DEF}
\artiklo{gloso}%
\vorto{gloso} Expliko di la vorti di autoro literaturala, di qui la senco divenis obskura pro lia obsoleteso\lingui{DEFIS}
\artiklo{gloto}%
\vorto{gloto}\fako{anat.} Fenduro an la parto supra di laringo, qua uzesas por la emiso di la voco\lingui{EFIRS}
\artiklo{gluglar}%
\vorto{gluglar}\fako{netrans.} Bruisar quale la vino qua ekiras ek la kolo di botelo\lingui{DFIS}
\artiklo{glumo}%
\vorto{glumo}\fako{bot.} Envelopilo florala di singla ek la spiketi qui konstitucas la spiko di la gramini. \textit{(En la linguo ne-ciencala : brano)}\lingui{EF}
\artiklo{gluo}%
\vorto{gluo} Materio viskoza, de qua strato extensita inter du objekti, igas li inter-adherar\lingui{EFI}
\artiklo{glutar}%
\vorto{glutar}\fako{trans.}\senco{1}{Decensigar en la guturo}\senco{2}{\textit{(metaf.)} Desaparigar subite (en abismo, en abiso)}\lingui{EFIRS}
\artiklo{gluteino}%
\vorto{gluteino} Materio viskoza qua restas en la farino de la cereali, kande onu deprenis de ici la amilo
\artiklo{gluteno}%
\vorto{gluteno} Substanco tenaca qua glutinas, interligas ad ensemblo, la parti dividita di la korpi solida\lingui{DEFIS}
\artiklo{gluteo}%
\vorto{gluteo}\fako{anat.} Parto mi-sferatra, karnoza, qua formacas singla latero di la regiono dopa di la korpo di homo o di animalo\lingui{EL}
\artiklo{glutinar}%
\vorto{glutinar}\fako{trans.} Adherigar per gluo\lingui{EF}
\artiklo{gneiso}%
\vorto{gneiso}\fako{mineral.} Roko ek feldspato, mikao, e c., ma di naturo skistatra\lingui{DEFIRS}
\artiklo{gnomiko}%
\vorto{gnomiko}\fako{antique} Skriptisto Greka qua versigis sentenci etikala o morala\lingui{DEFIS}
\artiklo{gnomo}%
\vorto{gnomo}\fako{mitol.} Mikra genio qua regnas tero e la kontenajo di ica\lingui{DEFIRS}
\artiklo{gnomono}%
\vorto{gnomono} Stango vertikala di qua la ombro, projektate sur plano, igas savar, per sua longeso, la alteso di Suno, ed indikas la kloko per sua situeso\lingui{DEFIRS}
\artiklo{gnoso}%
\vorto{gnoso} Doktrino (religio-hereziala) segun qua la enti spiritala, ekirinte de la sino di Deo, retroeniros lu\lingui{DEFIS}
\artiklo{gnostiko}%
\vorto{gnostiko} Adepto di gnoso\lingui{DEFIRS}
\artiklo{gnuo}%
\vorto{gnuo}\fako{zool.} Speco di granda antilopo di Afrika\latina{catoblepas gnu}\lingui{DEFIRS}
\artiklo{gobio}%
\vorto{gobio}\fako{zool.} Mikra fisho manjebla, de la genero « ciprini », qua vivas trupe apud la fundo di la riveri\latina{cyprinus gobio}\lingui{DSL}
\artiklo{gobleto}%
\vorto{gobleto} Drinko-vazo, alta, cilindratra, sen anso, ed ordinare sen pedo\lingui{EFS}
\artiklo{godrono}%
\vorto{godrono}\senco{1}{\textit{(tekn.)} Muluro ye seciono ovala, borde di la repastovazaro arjenta}\senco{2}{\textit{(stofo)} Plisuro ye seciono cirkla en la kolumi plisita, en la jaboti. en la mansheti}\lingui{EF}
\artiklo{goeleto}%
\vorto{goeleto}\fako{navig.} Navo lejera, du-masta\lingui{DFIS}
\artiklo{*golfo}%
\vorto{*golfo} Ludo quan partoprenas 2 o 4 personi, en qua li uzas quaza biliardo-buli, mikra e tre harda, quin onu duktas, per bastono-frapi, rulige sur sulo, ad 9 o 18 trui en ica\lingui{DEFIRS}
\artiklo{gonagro}%
\vorto{gonagro}\fako{patol.} Goto di la genuo\lingui{DEF}
\artiklo{gondo}%
\vorto{gondo} Fer-peco axala, kudatra e pivotoza, qua sustenas la bendo sur qua rotacas la pordo, e c.\lingui{F}
\artiklo{gondolo}%
\vorto{gondolo}\senco{1}{Vehilo suspendita sub aerostato (direktebla o ne)}\senco{2}{Speco di batelo uzata en la laguni, precipue en la urbo Venezia (Italia)}\lingui{DEFIRS}
\artiklo{gonfalono}%
\vorto{gonfalono}\senco{1}{Banero militala, suspendita an la supra parto di lanco, di standardo}\senco{2}{Banero ekleziala (katolika) sub qua (dum la Mezepoko) su asemblis la vasali di la Eklezio}\lingui{F}
\artiklo{goniometrio}%
\vorto{goniometrio} Cienco pri la mezuro di la anguli\lingui{DEF}
\artiklo{goniometro}%
\vorto{goniometro}\fako{geom.} Instrumento por mezurar la anguli (di la kristali, precipue)\lingui{DEF}
\artiklo{gonokocio}%
\vorto{gonokocio}\fako{patol.} Nomo di la ensemblo de la morbi quin produktas la gonokoko od en qui la gonokoko pleas rolo certa : blenoragii, metriti, oftalmii blenoragiala, e c.
\artiklo{gonoreo}%
\vorto{gonoreo}\fako{patol.} Ekfluajo de la membrano genitala ed urinifala\lingui{DEFIRS}
\artiklo{gorgear}%
\vorto{gorgear}\fako{netrans.} Bruisar quale aquo qua falas de tekto-pluvkanalo\lingui{DEFI}
\artiklo{gorilo}%
\vorto{gorilo}\fako{zool.} Simio di Afrika, analoga a chimpanzeo\latina{gorilla gina}\lingui{DEFIRS}
\artiklo{gorjereto}%
\vorto{gorjereto}\fako{kirurg.} Instrumento kanaloza, kun kanalo streta, por litotomiar od operacar ulu pri fistulo\lingui{DEFIRS}
\artiklo{gotika}%
\vorto{gotika}\senco{1}{\textit{(arkitekt.)} Ogivala, qua sucedis la Romanala}\senco{2}{\textit{(pri literi)} Perpendikla e di qua la konturi esas quaze facetoza ed ornita ye pinti, hoki, ankore uzata en la libri Germana}\lingui{DEFIRS}
\artiklo{goto}%
\vorto{goto}\fako{patol.} Morbo qua atakas ordinare la artiki, iras de ica ad ita, ed ulakaze a la viceri\lingui{EFIS}
\artiklo{grabar}%
\vorto{grabar}\fako{trans.}\senco{1}{Trasar sur materio harda (per stilo, cizelilo) (figuro, epigrafo)}\senco{2}{Trasar (per stilo, cizelilo) sur plako de metalo o de ligno, la kopiuro de pikturo, de desegnuro, de muziko-texto, por *reproduktar lu mult-exemplere}\lingui{DEFIRS}
\artiklo{gracila}%
\vorto{gracila} Gracioze tenua ed alta-statura\lingui{EFIS}
\artiklo{gracio}%
\vorto{gracio} Agreablajo, plezivajo, en persono, en la kozi\lingui{DEFIRS}
\artiklo{graco}%
\vorto{graco}\fako{teol.} Sokurso supernatura quan Deo donas a la homo por helpar lu salvar su\lingui{EFIRS}
\artiklo{gradiento}%
\vorto{gradiento}\fako{meteorol.} Difero di la preso atmosferala, evaluata per milimetri e per grado geografiala, inter punto e la centro maxim apuda di ciklono o di anticiklono
\artiklo{grado}%
\vorto{grado}\senco{1}{La parto di eskalero, di skalo, sur qua onu pozas la pedo por acensar o decensar}\senco{2}{Singla de la mediataji qui duktas sucede de ca a ta estado}\senco{3}{En serio de termini di qui la importo, la valoro, kreskas o deskreskas, plaso di un ek la termini relate la ceteri}\senco{4}{Singla ek la dividuri inter-egala di serio markoza}\lingui{DEFIRS}
\artiklo{grafika}%
\vorto{grafika}\senco{1}{Qua trasas per desegno}\senco{2}{Qualeso di la trasuro da aparato enrejistrala}\senco{3}{Qua relatas la reprezento di la son-uni vocala per la skribado}\lingui{DEFIRS}
\artiklo{grafismo}%
\vorto{grafismo}\senco{1}{Maniero skribar la vorti di linguo}\senco{2}{Skribo-stilo di individuo, e qua (segun-dice) informas pri lua karaktero}
\artiklo{grafito}%
\vorto{grafito}\fako{mineral.} Substanco \textit{(ferkarburo)} de qua onu facas la krayoni ordinara\lingui{DEFIRS}
\artiklo{grafologio}%
\vorto{grafologio} Studiado di la karaktero di la homi segun lia skribo-stilo, di lia grafismo\lingui{DEFIRS}
\artiklo{grafometro}%
\vorto{grafometro} Mi-cirklo kun dioptro ed alidado, por mezurar la anguli lor agromezuro, mapo-relevo
\artiklo{gramatiko}%
\vorto{gramatiko} Cienco pri la reguli di linguo\lingui{DEFIRS}
\artiklo{gramino}%
\vorto{gramino}\fako{bot.} Familio de planti monokotiledona, qui havas kom stipo palio kava, nodoza, kun flori dispozita spike
\artiklo{gramo}%
\vorto{gramo} Ponder-unajo di la metro-sistemo, to quon pezas un centimetro kuba de aquo distilita, ye 4 gradi Celsius\lingui{DEFIRS}
\artiklo{granario}%
\vorto{granario} La parto interna maxim supra di domo, uzata kom depozeyo desembarasala\lingui{FIS}
\artiklo{granato}%
\vorto{granato} Lapido, vino-reda, qua povas striizar quarco\lingui{DEFIRS}
\artiklo{granda}%
\vorto{granda}\senco{1}{\textit{(ecepte pri la staturo homala)} Qua transiras la dimensioni ordinara (partikulare pri alteso o longeso)}\senco{2}{Qua transiras la quanto ordinara pri quanto o qualeso}\senco{3}{Qua transiras la nivelo, la grado ordinara pri la rango, la situeso sociala}\senco{4}{Qua transiras la nivelo, la grado ordinara, pri la qualeso di la spirito (mento) o di la kordio}\lingui{EFIS}
\artiklo{grandioza}%
\vorto{grandioza} Qua havas karaktero impozanta di grandeso\lingui{DEFIRS}
\artiklo{*grandoro}%
\vorto{*grandoro}\fako{matem.} Quanto kontinua, qua eventuale kreskas o deskreskas
\artiklo{granito}%
\vorto{granito}\fako{geol.} Roko fairala, ek feldspato, mikao, quarco, asemblita en maso kompakta\lingui{DEFIRS}
\artiklo{granivora}%
\vorto{granivora} Qua su nutras per cerealo-grani\lingui{DEFIS}
\artiklo{grano}%
\vorto{grano}\senco{1}{Singla de la frukti kontenita en la spiko di la cerealo}\senco{2}{Frukto o semino globatra di ula planti (vito, grozeliero)}\senco{3}{Asperajo granatra di surfaco (petro, stofo, e c.)}\senco{4}{Karaktero granoza di ca surfaco}\lingui{DEFIS}
\artiklo{grantar}%
\vorto{grantar}\fako{trans., ad} Donar o koncesar ulo kom favoro (do demandita o dezirita)\lingui{E}
\artiklo{granulo}%
\vorto{granulo}\fako{farmacio} Pilulo extreme mikra\lingui{DEFIRS}
\artiklo{granulomo}%
\vorto{granulomo}\fako{patol.} Mikra tumoro globatra, irgequala : sifilisala, tuberklosala, inflamala, e c.
\artiklo{grapino}%
\vorto{grapino}\senco{1}{\textit{(navig.)} Abordo-hoko}\senco{2}{Ankro di shalupo kun plura pinti hokatra}\senco{3}{\textit{(tekn.)} Instrumento por separar, en la vitbero-presilo, la grani de la raflo}\senco{4}{Instrumento por deprenar la sordidaji de la vitro fuzesanta}\senco{5}{Ferajo hokatra di la kamen-skrapisto por skrapar la fuligino}\senco{6}{Hoko quan onu fixigas an sua pedi por klimar sur la arbori}\lingui{EFI}
\artiklo{grapo}%
\vorto{grapo}\fako{bot.} Asemblajo de flori, de frukti, dispozita etajatre sur pedunklo komuna\lingui{FI}
\artiklo{graso}%
\vorto{graso}\senco{1}{Substanco oleoza, uncion-apta, dispersita en la tisuo celulara di la homo-korpo e di la animali}\senco{2}{Ta substanco, deprenita de la korpo di la animali, ed uzata por koquo od en la industrii}\lingui{EFIS}
\artiklo{gratar}%
\vorto{gratar}\fako{trans.} Frotar, skrapante la surfaco\lingui{FI}
\artiklo{gratifikar}%
\vorto{gratifikar}\fako{trans.} Donacar pekunio ad ulu kom atesto pri sua satisfaceso, ultre to quo debesas a lu po lua laboro\lingui{DFIS}
\artiklo{gratino}%
\vorto{gratino}\senco{1}{Parto di ula disho, qua adheras e risolas sur la parieto di la vazo lor koqueso}\senco{2}{Maniero koquar ula dishi kovrita de pan-krusto raspita, e pose risoligita}\lingui{DEF}
\artiklo{gratitudar}%
\vorto{gratitudar}\fako{trans., pri, pro} Atestar ke onu etiko-debas ad ulu pro ke onu recevis de lu bonfaco od helpo\lingui{EFIS}
\artiklo{gratuita}%
\vorto{gratuita}\senco{1}{Quan onu profitas sen pagar}\senco{2}{Gratuita : sen demandar pago kompense}\lingui{DEFIRS}
\artiklo{gratular}%
\vorto{gratular}\fako{trans., pri, pro} Dicar politeso-vorti (ad ulu) pri la fortuno quan lu profitas\lingui{DEFIS}
\artiklo{grava}%
\vorto{grava}\senco{1}{Desfacile movebla, levebla, pro lua pezo}\senco{2}{\textit{(muziko)} Qua apartenas a la gradi infra di la muziko-skalo}\senco{3}{\textit{(gram.)} En ula lingui (kom ex. la Franca) Qualeso di la supersigno analoga a komo inversigita qua donas a la vokalo « e » son-uno apertita, o distingas ula homonimi \textit{(kom ex.en F. : la e là)}}\senco{4}{\textit{(afero)} ne-vana, nefrivola, egardenda kom importoza}\lingui{EFIS}
\artiklo{gravelo}%
\vorto{gravelo}\fako{patol.} Konkreciono qua naskas en la reni\lingui{EF}
\artiklo{gravida}%
\vorto{gravida} Qua portas filio (yuno) en sua ventro\lingui{IL}
\artiklo{gravio}%
\vorto{gravio}\senco{1}{Sablo grosa-grana, qua devenas de la desagrego di roki stonoza}\senco{2}{Sabl-uni grosa per qui onu kovras la alei di gardeno}\lingui{EFR}
\artiklo{gravitar}%
\vorto{gravitar}\fako{netrans.} Obediar la forco pro qua la molekuli interatraktas proporcionale ye lia maso\lingui{DEFIS}
\artiklo{greftar}%
\vorto{greftar}\fako{trans., sur} Insertar sproso di planto sur altra planto, por ke ica portez la floro, la frukto dicita\lingui{EF}
\artiklo{grelar}%
\vorto{grelar}\fako{netrans.} Falo di pluvo-guteti konjelita, grano-forma, partikulare dum sturmo-vetero\lingui{F}
\artiklo{gremio}%
\vorto{gremio}\senco{1}{Che persono sidanta : parto di lua korpo qua iras de la tayo an la genui}\senco{2}{\textit{(metaf.)} La sino}\lingui{DIS}
\artiklo{grenado}%
\vorto{grenado}\senco{1}{\textit{(bot.)} Frukto sferoida qua kontenas semini reda, enklozita en mikra celuli, produktita da arbusto de la familio « mirtacei »}\senco{2}{\textit{(milit-arto)} Projektilo grenadatra, bulo de fero, plena ye stupo e de pulvero, di qua onu acendis la mecho por lansar lu manue}\latina{punica granatum}\lingui{DEFIRS}
\artiklo{greso}%
\vorto{greso}\fako{mineral.}\senco{1}{Roko ek sabl-uni quarcoza}\senco{2}{Rekuro-sablo, facita ek ta petro pulverigita}\lingui{FI}
\artiklo{greto}%
\vorto{greto} Asemblajo ajuroza ek stangi de fero o de ligno, uzata kom klozilo di gardeno, parko, e c.\lingui{EFS}
\artiklo{gribo}%
\vorto{gribo} Peco de panikulo o di lardo de porko, fritita en padelo
\artiklo{grifono}%
\vorto{grifono}\senco{1}{\textit{(mitol.)} Animalo mi-agla, mi-leona}\senco{2}{\textit{(zool.)} Mikra hundo pudala, kun longa pili herisita sur la kapo e la parto avana di la korpo}\lingui{DEFIS}
\artiklo{grilar}%
\vorto{grilar}\fako{trans.}\senco{1}{\textit{(koquarto)} Bruletar, per fairo vivaca, nutrivo (karno, fisho, pano) pozita sur koquo-utensilo ek dina stangi de fero}\senco{2}{Submisar erco a fairo por igar lu plu friabla}\senco{3}{Submisar a la efiko di flamo kotonajo (pos texto) o pultruno por supresar la lanugo}\lingui{EF}
\artiklo{grilio}%
\vorto{grilio}\fako{zool.} Insekto de la klaso « ortopteri », che qua la maskulo produktas bruiso akuta per interfrotar sua elitri, e qua tre prizas la loki suno-radioza, varma\latina{gryllus}\lingui{DFIS}
\artiklo{griliotalpo}%
\vorto{griliotalpo}\fako{zool.} Insekto qua devastas la gardeni, vivas sub la tersurfaco, e bruisas analoge a grilio\latina{gryllus gryllotalpa}
\artiklo{grimar}%
\vorto{grimar}\fako{trans.) (teatro} Rugizar (ulu) por igar lu aspektar plu olda\lingui{FR}
\artiklo{grimasar}%
\vorto{grimasar}\fako{netrans.} Quaze tordar sua vizajo\lingui{DEFR}
\artiklo{grincar}%
\vorto{grincar}\fako{netrans.) (pri mekanismo} Bruisar akre pro tre streta kontakto sen interpozeso di oleo-strato\lingui{FI}
\artiklo{grindar}%
\vorto{grindar}\fako{trans.} Submisar (korpo generale harda : torn-utensilo, frezo-cilindro) a la efiko di grosa disko qua rotacas rapide (ordinare ek petro tre harda : karborundo, smerilo, diamanto-polvo) por quaze limar lu
\artiklo{gripo}%
\vorto{gripo}\fako{patol.} Quaza kataro epidemiala\lingui{DEFIRS}
\artiklo{griza}%
\vorto{griza} Koloro meza inter blanka e nigra\lingui{EFIS}
\artiklo{grogo}%
\vorto{grogo} Drinkajo ek aquo kolda o varma, en qua onu varsas sukro, brandio od altra alkoholajo, adjuntante citronolonchi\lingui{DEFIR}
\artiklo{grondar}%
\vorto{grondar}\fako{netrans.} Audigar bruiso duroza e nelauta, quale la tondro fora od ula animali iracoza\lingui{DEFIRS}
\artiklo{gropo}%
\vorto{gropo}\fako{anat.} Che ula animali, parto dopa, sferatra, de la hancho til la loko ube komencas la kaudo\lingui{DEFIRS}
\artiklo{grosa}%
\vorto{grosa}\senco{1}{Qua transiras la volumino ordinara}\senco{2}{Qua transiras la volumino di altra kozo sama-speca}\lingui{EFIS}
\artiklo{grosdekeduo}%
\vorto{grosdekeduo} Dekeduo de dekedui. (t. e. 144)
\artiklo{grosiera}%
\vorto{grosiera}\senco{1}{\textit{(pri kozo)} Facita ye maniero vulgara, o fasonita krude}\senco{2}{\textit{(pri persono)} Di qua la rudeso ne pludolcigesas da civilizeso}\lingui{EFIS}
\artiklo{groteska}%
\vorto{groteska} Rid-ecitanta pro lua aspekto baroka, stranja\lingui{DEFIRS}
\artiklo{groto}%
\vorto{groto} Exkavajo pitoreska, naturala o manuo-facita\lingui{DEFIRS}
\artiklo{grozelo}%
\vorto{grozelo}\fako{bot.} Frukto redatra o verda, plu grosa kam la ribo (uzata, kande lu esas ankore nematura ed acido-saporanta, por facar sauco qua akompanas makrelo-disho) qua kreskas sur arbusto qua formacas genero de la familio « grosulariei »\latina{Ribes grossularia}\lingui{FS}
\artiklo{grucho}%
\vorto{grucho} Bastono suportila di qua la extremajo supra esas garnisita ye traverso, sur qua la infirmi su apogas (per la manuo o per la axelo) por marchar\lingui{DEI}
\artiklo{gruelo}%
\vorto{gruelo} La parto di la frumento maxim glutenoza, qua envelopas la jermo di la grano e – pro esar plu harda kam lo cetera – trituresas neperfekte sub la muelilo\lingui{F}
\artiklo{grumelo}%
\vorto{grumelo} Mikra amaso de substanco pulveratra aglomerita, o de substanco liquida koagulita\lingui{EFIS}
\artiklo{grumo}%
\vorto{grumo}\senco{1}{Yuna lakeo da qua onu sequigas su, kavalkante, ed a qua onu donas la reini kande onu decensas de kaval-movata veturo quan onu duktas}\senco{2}{Garsono puera qua exekutas la komisi (en hotelo)}\lingui{DEFR}
\artiklo{grunar}%
\vorto{grunar}\fako{netrans.} Murmurar kom manifesto di nekontenteso\lingui{DEFIS}
\artiklo{gruo}%
\vorto{gruo}\fako{zool.} Granda ucelo voyajera, de la klaso « steltieri »\latina{gruidae grua}\lingui{FL}
\artiklo{grupo}%
\vorto{grupo}\senco{1}{Ensemblo de personi, kozi, statuo}\senco{2}{Ula quanto de personi, de kozi, quin interproximigas ulo komuna}\lingui{DEFIRS}
\artiklo{gruzo}%
\vorto{gruzo}\fako{zool.} Ucelo galinacea, de la genero « tetrasi », qua vivas en la monti alta\latina{lagopus scoticus}\lingui{DEF}
\artiklo{guacho}%
\vorto{guacho} Induturo ek farbi dilutita kun gumo, ed igita pastatra per mielo od altra substanci\lingui{DEFI}
\artiklo{guanino}%
\vorto{guanino}\fako{kemio} C₅N₅H₅O
\artiklo{guano}%
\vorto{guano} Dungo tre efikiva, ek la amasi de mar-ucelo-feko sur la litoro, precipue di Chili\lingui{DEFIRS}
\artiklo{guardo}%
\vorto{guardo} Korpo ek soldati di qua la rolo konsistas precipue ye servar apud la suvereno\lingui{DEFIRS}
\artiklo{guatar}%
\vorto{guatar}\fako{trans.} Surveyar paciente por surprizar (ulu, ulo)\lingui{FI}
\artiklo{guayakolo}%
\vorto{guayakolo}\fako{kemio} Etero extraktita de la gayak-rezino\simbolo{CH₃O. C₆H₄. OH}
\artiklo{gubernio}%
\vorto{gubernio} Provinco direktata politikale e militale ye la nomo di la suvereno\lingui{DEFIRS}
\artiklo{gudro}%
\vorto{gudro}\senco{1}{Mixuro de sapo e de suki rezinoza, quan onu extraktas de la olda abieto-ligno per kombustar lu en furnelo}\senco{2}{Ta substanco a qua onu mixas sebo, oleo, uzata por indutar la kareno di navo, la kordegi, e c.}\lingui{F}
\artiklo{guelfo}%
\vorto{guelfo}\fako{historio di Italia} La persono, la partisano qua, dum la milito civila lor la Mezepoko, favoris la papo (kontraste kun la gibelini, qui favoris la imperiestro Germana)\lingui{DEFIRS}
\artiklo{gufo}%
\vorto{gufo}\fako{zool.} Nokt-ucelo de la familio « strigi »\latina{Otis vulgaris}\lingui{I}
\artiklo{guidar}%
\vorto{guidar}\fako{trans., ad, pri} Montrar la justa voyo ad ulu (akompanante lu) qua ne konocas olu, e qua povas hezitar od erorar\lingui{EFIS}
\artiklo{guindo}%
\vorto{guindo}\fako{bot.} Sorto di cerizo, nigra-reda\latina{Cerasus juliana}
\artiklo{*gulfo}%
\vorto{*gulfo}\fako{geogr.} Maro-parto qua formacas larja konkavajo en la tero di litoro\lingui{DEFIRS}
\artiklo{gulo}%
\vorto{gulo}\fako{mitol. orientala} Genio qua devoras la korpi di la mortinti\lingui{DEFIS}
\artiklo{gulyasho}%
\vorto{gulyasho} Bovo-karno, kun onioni hachita, papriko, alio, kumino-grani, majorano, muskado, e c., koquinta dum tri o quar hori : a la sauco adjuntesas kelka buliono e farino
\artiklo{gumiguto}%
\vorto{gumiguto} Gumo ek arboro di India\latina{garcinia}\lingui{DFIR}
\artiklo{gumlako}%
\vorto{gumlako} Gumo ek planti leguminosa di India\lingui{DEFIRS}
\artiklo{gumo}%
\vorto{gumo} Substanco qua defluas de ula arbori, uzata por glutinar, apretar, briligar, o kom dolcigivo en siropi, pasti, pastili, e c.\lingui{DEFIRS}
\artiklo{gurdo}%
\vorto{gurdo} Botelo plata, sur la parto extera di qua aplikesis oziero-plektajo, en qua onu varsas vino, brandio, e c.\lingui{EF}
\artiklo{gurmando}%
\vorto{gurmando}\fako{persono} Qua prizas la dishi bona (bone-saporoza) e manjas multo de li\lingui{DEF}
\artiklo{gurmo}%
\vorto{gurmo}\senco{1}{Flegmazio di la mukozo nazala, che la yuna kavalo}\senco{2}{Exantemo di la vizajo e di la pelo haroza (che la homo)}\lingui{F}
\artiklo{gurto}%
\vorto{gurto} Bendo de ledro o de telo, quan onu tensas por mantenar ulo\lingui{DE}
\artiklo{gustar}%
\vorto{gustar}\fako{trans.}\senco{1}{Perceptar la saporo (di ulo)}\senco{2}{Perceptar la saporo (di ulo) por judikar lo. \textit{(Anke metaf.)}}\lingui{DEFIS}
\artiklo{gutaperko}%
\vorto{gutaperko} Materio gumoza, flexebla, plastika, uzata kom substanco nepermeebla\lingui{DEFIRS}
\artiklo{guto}%
\vorto{guto} Mikra quanto de liquido qua su separas en la formo di globeto\lingui{FS}
\artiklo{guturo}%
\vorto{guturo}\fako{anat.}\senco{1}{Prolonguro di la faringo (dopa boko) qua komunikas kun la ezofago}\senco{2}{Prolonguro di la laringo qua komunikas kun la faringo}\lingui{EFIS}
\artiklo{guvernar}%
\vorto{guvernar}\fako{trans.}\senco{1}{Direktar la aferi di la stato}\senco{2}{Direktar la eduko di infanto, puero, adolecanto}\lingui{DEFIRS}
\artiklo{guyavo}%
\vorto{guyavo}\fako{bot.} Frukto (piroforma, di qua la pulpo esas sukroza e parfumoza) qua kreskas sur granda arboro exotika, qua formacas genero di la familio « mirtacei »\latina{psidium guyaba}\lingui{L}
\end{multicols}
\sekciono{H}
\begin{multicols}{2}
\artiklo{ha!}%
\vorto{ha!} Interjeciono qua indikas ke onu quik expresos ula emoco tre vivaca di psiko (joyo, doloro, iraco, e c.)
\artiklo{habila}%
\vorto{habila}\fako{pri persono} Apta sucesar en to quon lu agas o facas\lingui{DEFIS}
\artiklo{habilitar}%
\vorto{habilitar}\fako{trans.) (yuro-cienco} Igar kapabla legale pri agor ulo\lingui{DEFIS}
\artiklo{habitaklo}%
\vorto{habitaklo}\fako{navo} Armoro, sur la kastelo avana, avan la guvernilo-roto, en qua enklozesas la busolo, la horlojo e la lanternego\lingui{F}
\artiklo{habitar}%
\vorto{habitar}\fako{trans.} Rezidar permanante en domo propra o fixa\lingui{EFIS}
\artiklo{hachar}%
\vorto{hachar}\fako{trans.}\senco{1}{Tranchadar a peceti, per hakilo, hakokultelo, e c. \textit{(anke metaf.)}}\senco{2}{Provizar (desegnuro, e c. per mikra streki, paralela o qui interkrucumas)}\lingui{EF}
\artiklo{hagiografio}%
\vorto{hagiografio} Cienco qua traktas la vivo-maniero e la agi di la santi (di la eklezio katolika)
\artiklo{hagiografo}%
\vorto{hagiografo} Autoro qua naracas la vivo-maniero e la agi di la santi (di la eklezio katolika)\lingui{DEFIS}
\artiklo{hakar}%
\vorto{hakar}\fako{trans., per} Abatar, tranchar, fendar, per instrumento fera, lamo larja ye formo di triangulo kurvolinea; konvexa, ye lua bazo (qua esas tranchiva); konkava, ye lua lateri, e qua esas fixigita an mancho\lingui{DEFIS}
\artiklo{halbardo}%
\vorto{halbardo} Shaft-armo, kun mancho longa, ferlamo tranchiva e pintoza, plus du ferlami laterala (ali), de qui un esas krecentatra, la altra pinta\lingui{DEFIRS}
\artiklo{halo}%
\vorto{halo}\senco{1}{Chambro, salono, extreme vasta}\senco{2}{Granda placo tektoza, ube esas la merkato di la manjebli}\lingui{DEFS}
\artiklo{halofito}%
\vorto{halofito}\fako{biol.} Nomo generala di la vejetanti qui kreskas naturale en sulo tre mar-saloza, sive an maro, sive apud salini
\artiklo{halogena}%
\vorto{halogena} Nomo donita da Berzelius ad irga ek la korpi ye qui konsistas la familio di kloro : fluoro, kloro, bromo, iodo, qui esas apta formacar sali per kombinar su kun la metali
\artiklo{halono}%
\vorto{halono}\senco{1}{\textit{(astron.)} Cirklo lumoza, maxim-multa-kaze koloroza, qua aparas ul-instante cirkum la disko Suna, Luna o planeta kande li brilas tra atmosfero vaporoza}\senco{2}{\textit{(anat.)} Mikra cirklo reda qua cirkondas la mamo-pinto}
\artiklo{haltar}%
\vorto{haltar}\fako{netrans.} Cesar sua avanco, sua marcho, ne irar plu-fore; cesar funcionar\lingui{DEFI}
\artiklo{haltero}%
\vorto{haltero}\fako{gimnastiko} Kurta stango di qua la du extremaji finas per fer-bulo, quan onu elevas ed abasas por exercar la muskuli di la brakii\lingui{FIS}
\artiklo{haltiko}%
\vorto{haltiko}\fako{zool.} Insekto koleoptera, plantivora, saltera, ek la familio « krizomelidei », qua devoras la planti leguma e la vito\latina{altica}
\artiklo{halucinar}%
\vorto{halucinar}\fako{netrans.) (patol.} Subisar quaza alienaco neduriva dum qua onu perceptas sento, kontre ke ne existas extera objekto apta produktar ta sento\lingui{DEFIRS}
\artiklo{haluxo}%
\vorto{haluxo} La grosa pedo-fingro
\artiklo{hamadriado}%
\vorto{hamadriado}\fako{mitol.} Nimfo di la boski, di qua la fato esas ligita ad olca di arboro, e qua naskas e mortas kun ica
\artiklo{hamako}%
\vorto{hamako} Telo o reto suspendita horizontale per lua du extremaji, tale ke lu formacas lito portebla en qua onu su kushas o su ociligas\lingui{EFRS}
\artiklo{hamstro}%
\vorto{hamstro}\fako{zool.} Mikra mamifero rodera, di Rusia e di la nordo-regioni di Germania\latina{cricetus}\lingui{DFRS}
\artiklo{hancho}%
\vorto{hancho}\fako{anat.} Singla de la du parti simetra di la homo-korpo, qua salias adextere, inter la kruro e la kosti, e produktita da la expansuro di la iliak-osto, ube lu formacas artiko kun la femuro\lingui{DEFIS}
\artiklo{handikapar}%
\vorto{handikapar}\fako{trans.}\senco{1}{\textit{(kavalkur-konkurso)} (en kur-konkurso pri qua onu admisas kavali irge-evoza ed irge-quala) kargar ula ek li per pez-masi qui esas diversa segun lia evo, forteso (grado, por egaligar la suceso-chanci)}\senco{2}{\textit{(metaf.)} Impozar desavantajo ad ulu}\lingui{DEFIRS}
\artiklo{hangaro}%
\vorto{hangaro} Tekto pozita sur fosti o pilastri, sub qua onu shirmas veturi, ligno, e c.\lingui{DEF}
\artiklo{hano}%
\vorto{hano}\fako{zool.} Pultro de la klaso « galinacei », di qua la maskulo havas sur la kapo kresto reda, kun tarsi forta e garnisita ye sporno hokatra, e kantas per sorto di krio klara e penetranta\lingui{DE}
\artiklo{hanso}%
\vorto{hanso}\fako{olim} Korporaciono de la komercisti (exemplo : la hanso Teutonika)\lingui{DEFIRS}
\artiklo{haptiko}%
\vorto{haptiko} La cienco qua traktas la fenomeni pri tushado
\artiklo{haraso}%
\vorto{haraso}\senco{1}{Trupo ek kavaluli e kavalini di raso la maxim bona, asemblita skope di genito}\senco{2}{Establisuro specala por la eduko e produkto di kavali}\lingui{F}
\artiklo{harda}%
\vorto{harda}\senco{1}{\textit{(aludante kozo)} Qua tre rezistas lor tusheso}\senco{2}{Qua tre rezistas kontre entameso}\senco{3}{\textit{(aludante persono)} Qua ne apertesas por la sentimenti di benigneso, di humaneso, e c.}\lingui{DE}
\artiklo{haremo}%
\vorto{haremo}\senco{1}{Che la Mohamedisti : apartmento di la homini}\senco{2}{Asemblitaro ek ta homini}\lingui{DEFIRS}
\artiklo{haringo}%
\vorto{haringo}\fako{zool.} Mar-fisho de la familio « klupi » qua decensas singla-yare, en trupi nekontebla, de la Oceano Glaciala Arktika\latina{Clupea harengus}\lingui{DEFIS}
\artiklo{harmoniar}%
\vorto{harmoniar}\fako{netrans., kun} Havar sua parti tala, ke singla de li konkordas kun la ceteri. \textit{(anke metaf.)}\lingui{DEFRS}
\artiklo{harmoniko}%
\vorto{harmoniko}\fako{matem.} Proporciono ek tri nombri tala ke la eceso di la unesma ye la duesma relatas la eceso di la duesma sur la triesma quale la duesma ye la unesma
\artiklo{harmoniko}%
\vorto{harmoniko}\fako{muziko}\senco{1}{Muzik-instrumento antiqua, de kalici vitra, neegale plenigita ye aquo, e rotacanta, e qua, per la froto di la bordo di la kalici da la fingri, bruisas vibre}\senco{2}{Muzik-instrumento en qua la kordi remplasesas da lami vitra di qui la longeso ne esas inter-egala}\lingui{DEFIRS}
\artiklo{harneso}%
\vorto{harneso} Equipo-kozi di sel-kavalo o jungo-kavalo\lingui{DEFS}
\artiklo{haro}%
\vorto{haro}\fako{anat.} Singla de la pili qui garnisas la kranio-pelo (che la homo)\lingui{DE}
\artiklo{harpio}%
\vorto{harpio}\senco{1}{\textit{(mitol.)} Ento quan onu reprezentis kom havanta la vizajo di muliero e la korpo di vulturo}\senco{2}{\textit{(metaf.)} Persono kapt-avida, muliero quaze rabiika}\lingui{DEFIRS}
\artiklo{harpo}%
\vorto{harpo} Muzik-instrumento tri-angula, kun kordi vertikala, quin onu pincas per la du manui e quin onu vibrigas\lingui{DEFIRS}
\artiklo{harpuno}%
\vorto{harpuno}\senco{1}{Instrumento fera, uzata por akrochar, pikar}\senco{2}{Metal-squadro por ligar du peci ek konstrukturo}\lingui{DEFRS}
\artiklo{hashisho}%
\vorto{hashisho} Kanabo di India, di qua onu mastikas o fumas la folii sekigita\lingui{DEFIRS}
\artiklo{hastar}%
\vorto{hastar}\fako{netrans., pri} Ne disipar sua tempo (lor agar o facar ulo)\lingui{DEF}
\artiklo{haular}%
\vorto{haular}\fako{trans.} Tirar (batelo) ad su per kordo (por igar lu avancar alonge rivero, kanalo)\lingui{DEFIS}
\artiklo{hauo}%
\vorto{hauo} Plug-instrumento fera kun ligno-mancho, por exkavar tero; la fera parto esas larja, plata e tranchiva, kun aristo en plano perpendikla ye la mancho\lingui{DEF}
\artiklo{havar}%
\vorto{havar}\fako{trans.} Expresas la relato di ento kun to quo apartenas a lu\lingui{DEFIS}
\artiklo{hazardo}%
\vorto{hazardo} Kauzo blinda atribuita a la fakti di qui la kauzo reala eskapas ni\lingui{DEFRS}
\artiklo{he!}%
\vorto{he!} Interjeciono por vokar, avertar, atraktar la atenco, expresar la kompato
\artiklo{hederaceo}%
\vorto{hederaceo} Genero de « labiei », mikra herbo perena, kun stipi reptera, flori mikra klar-violea, qua vivas en la loki koldeta\latina{glechoma hederacea}
\artiklo{hedero}%
\vorto{hedero}\fako{bot.} Arbusto konsakrita da la Antiqui a Bako, planto klimera qua su ligas per radiketi acesora a la arbori, a la muri, quin lu kovras ye sua foliaro briloza e sempre verda\latina{hedera}\lingui{ISL}
\artiklo{hefo}%
\vorto{hefo} Parto de pasto quan onu lasis fermentacar, e quan onu mixas a la pasto fresha de qua onu facas pano, kuki, por ke ica (pos fermentacir kelke) inflez e divenez plu lejera\lingui{D}
\artiklo{hegemonio}%
\vorto{hegemonio} Speco di dominaco o di prepondero, da un naciono inter altri, de la vidopunto historiala e politikala\lingui{DEFIRS}
\artiklo{hego}%
\vorto{hego} Barilo o klozilo (cirkum agro, gardeno), formacita ek arbusti qui interplektas\lingui{DEF}
\artiklo{hejiro}%
\vorto{hejiro} Ero di la Mohamedisti, qua evas de kande Mohamed fugis de Mekka (yaro 622 pos I. K.)\lingui{DEFIRS}
\artiklo{hekatombo}%
\vorto{hekatombo}\senco{1}{\textit{(antique)} Sakrifiko di cent viktimi, di granda quanto de personi}\senco{2}{Masakro}\lingui{DEFIRS}
\artiklo{hektaro}%
\vorto{hektaro} Cent ari\lingui{DEFIRS}
\artiklo{hektika}%
\vorto{hektika}\fako{patol.} \textit{(febro)} karakterizata da deperiso progresanta\lingui{DEFIS}
\artiklo{hektografar}%
\vorto{hektografar}\fako{trans.} Imprimar (per aparato) skriburo, desegnuro, muziko-noti quin onu preliminare trasis, per perforo, sur *stencilo ek vaxofolio o parafino-folio quan onu aplikas sur la inkizorolero di ta aparato; onu pasigas sub la rolero la paper-folii virga : la inko trairas la perforuro- e lasas la traco necesa
\artiklo{hektogramo}%
\vorto{hektogramo} Cent grami\lingui{DEFIRS}
\artiklo{hektolitro}%
\vorto{hektolitro} Cent litri\lingui{DEFIRS}
\artiklo{hektometro}%
\vorto{hektometro} Cent metri\lingui{DEFIRS}
\artiklo{hektowato}%
\vorto{hektowato} Cent « watt »\lingui{DEFIRS}
\artiklo{heleboro}%
\vorto{heleboro}\fako{bot.} Planto herbacea de la familio « ranunkulacei », di qua la stipi esas rametoza, la flori verdatra, e la radiko uzata medicinale por la morbi nervala\latina{helleborus}\lingui{EFISL}
\artiklo{heliantemo}%
\vorto{heliantemo}\fako{bot.} Genero de planti ek la familio « cistacei »\latina{helianthemum}
\artiklo{helianto}%
\vorto{helianto}\fako{bot.} Genero de planti de la familio « kompozaji », a qua apartenas la sunfloro e la topinamburo\latina{Helianthus}\lingui{EFIS}
\artiklo{helico}%
\vorto{helico}\senco{1}{\textit{(geom.)} Lineo du-kurva, trasita filetatre cirkum cilindro}\senco{2}{\textit{(mekan.)} Aparato helicatra qua, movata da vaporo o motoro, propulsas la navi, la aeroplani, e c.}\lingui{EFIS}
\artiklo{helicoido}%
\vorto{helicoido}\fako{geom.} Surfaco produktita da rekto perpendikla ye la axo vertikala di cilindro rekta, e turnanta alonge la helico trasita sur la cilindro
\artiklo{heliko}%
\vorto{heliko}\fako{zool.} Molusko gasteropoda, reptera, di qua la konko esas helicatra\latina{cochlea}
\artiklo{helikoptero}%
\vorto{helikoptero} Aviac-aparato qua su elevas preske vertikale en atmosfero (e decensas de ibe) per surfaci rotacanta quin movas motoro\lingui{EFIRS}
\artiklo{helio}%
\vorto{helio}\fako{kemio} Gaso di atmosfero, deskovrita sur Suno (per la analizo spektrala) e pose sur Tero (en uranio-erco)\simbolo{He}\lingui{DEFIRS}
\artiklo{heliocentrala}%
\vorto{heliocentrala}\fako{astron.} Qua relatas la centro di Suno\lingui{DEF}
\artiklo{heliofilo}%
\vorto{heliofilo} Qua prizas la lumo-radii
\artiklo{heliofoba}%
\vorto{heliofoba} Qua patologie timas la lumo-radii
\artiklo{heliograbar}%
\vorto{heliograbar}\fako{trans.} Reproduktar imaji fotografura per imprimo de sika graburo
\artiklo{heliografar}%
\vorto{heliografar}\fako{trans.} Aplikar la procedi per qui reproduktesas, sur paperfolii fotografala, kalquo-desegnuri di qui la linei esas nigra sur fundo blanka
\artiklo{heliometro}%
\vorto{heliometro}\fako{astron.} Aparato uzata por mezurar la diametro semblanta di Suno e di la planeti\lingui{DEFIRS}
\artiklo{helioskopo}%
\vorto{helioskopo}\fako{astron.} Vitro-peco koloroza per qua onu povas regardar Suno\lingui{DEFIRS}
\artiklo{heliostato}%
\vorto{heliostato}\fako{tekn.} Aparato optikala qua mantenas en konstanta direciono sunoradio enduktita en kamero obskura\lingui{DEFIRS}
\artiklo{heliotropo}%
\vorto{heliotropo}\fako{bot.} Planto de la familio « boraginei », di qua un speco havas flori qui odoras dolcege\latina{Heliotropium europeum}
\artiklo{helixo}%
\vorto{helixo}\fako{anat.} La rebordo di la orel-funelo\lingui{EFIS}
\artiklo{heller}%
\vorto{« heller »}\fako{ante la milito universala di 1914-18} Moneto-peco ek bronzo e nikelo, adoptita da Austria-Hungaria, centima parto di la « krono »
\artiklo{helmo}%
\vorto{helmo}\senco{1}{\textit{(olim)} Kasko pintoza qua kovras la kapo e la vizajo, kun aperturo por la okuli}\senco{2}{\textit{(blazon-arto)} Peco qua indikas, per sua situeso, la gradi diversa di nobeleso}\lingui{DEFIS}
\artiklo{helpar}%
\vorto{helpar}\fako{trans.} Partoprenar lo facata od agata da ulu qua ne povas facar od agar lo sole\lingui{DE}
\artiklo{hem!}%
\vorto{hem!} Interjeciono qua expresas dubito, hezito
\artiklo{hematito}%
\vorto{hematito}\fako{mineral.} Sanguino, erco reda brunatre\lingui{DEFIRS}
\artiklo{hemerokalo}%
\vorto{hemerokalo}\fako{bot.} Planto liliacea, di qua la speco precipua esas la konvolvulo\latina{Hemerocallis}\lingui{EFIS}
\artiklo{hemiopa}%
\vorto{hemiopa}\fako{patol.} Malada de ta perturbeso vidala, pro qua onu dicernas nur la duimo sinistra o dextra di la objekto\lingui{EF}
\artiklo{hemiplegio}%
\vorto{hemiplegio}\fako{patol.} Paralizeso di duimo laterala di la homo-korpo\lingui{DEF}
\artiklo{hemiptero}%
\vorto{hemiptero}\fako{zool.} Insekto karakterizata da mi-elitri qui kovras la parto avana di la ali\lingui{DEFIS}
\artiklo{hemistiko}%
\vorto{hemistiko}\fako{prozodio} Duimo de verso, separita de la altra duimo per silenceto (cezuro)\lingui{DEFIS}
\artiklo{hemo}%
\vorto{hemo} La loko intima ube onu vivas (kontraste kun loko ube onu vivas, en hotelo o lojeyo)\lingui{DEF}
\artiklo{hemoglobino}%
\vorto{hemoglobino}\fako{kemio biol.} Pigmento nitroza e feroza, reda, quan onu renkontras en la hematii di la vertebrozi ed en la hemo-limfo di ula senvertebri\lingui{DEF}
\artiklo{hemolizo}%
\vorto{hemolizo}\fako{biol.} Destrukto di la sango-globuli reda
\artiklo{hemoptizio}%
\vorto{hemoptizio}\fako{patol.} Sputado di sango\lingui{DEF}
\artiklo{hemoragio}%
\vorto{hemoragio}\fako{patol.} Fluo di sango ek la vaskuli, kun o sen rupteso di la parieti\lingui{DEF}
\artiklo{hemoroido}%
\vorto{hemoroido}\fako{patol.} Tumoro di la veini di la anuso, de qua sango eskapas intervalope\lingui{DEFIRS}
\artiklo{henry}%
\vorto{« henry »}\fako{elektrotekniko} Unajo di indukto reciproka o di autoindukto, di qua la valoro egalesas 109 g unaji elektromagnetala C\lingui{S}
\artiklo{hepatiko}%
\vorto{hepatiko}\fako{bot.} Genero de « ranunkulei » de qua un speco rekomendesas kontre la hepato-morbi\latina{hepatica}\lingui{EFIS}
\artiklo{hepatito}%
\vorto{hepatito}\senco{1}{\textit{(mineral.)} Lapido hepatea}\senco{2}{\textit{(patol.)} Inflamuro di la hepato}\lingui{DEF}
\artiklo{hepato}%
\vorto{hepato}\fako{anat.} Vicero situita, che la homo e la precipua vertebrozi, an la parto supra di la latero rekta di la abdomino, organo qua sekrecas bilo, e produktas materio sukroza de qua lu richigas la sango\lingui{DEFIS}
\artiklo{heptaedro}%
\vorto{heptaedro}\fako{geom.} Solido sep-edra\lingui{DEP}
\artiklo{heptagono}%
\vorto{heptagono}\fako{geom.} Poligono qua havas sep lateri e sep anguli\lingui{DEF}
\artiklo{heraldo}%
\vorto{heraldo}\fako{olim} Oficiro di qua la rolo konsistis ye portar la mesaji (defii, sumni, milito-deklari, e c.), agar la proklami, regulizar la festi chevalierala, sorgar pri la blazoni\lingui{DEFIRS}
\artiklo{herbario}%
\vorto{herbario} Kolektajo ek planti sikigita e konservata inter paper-folii\lingui{DEFIRS}
\artiklo{herbivora}%
\vorto{herbivora} Qua su nutras per herbo\lingui{DEFIS}
\artiklo{herbo}%
\vorto{herbo}\fako{bot.}\senco{1}{Planto di qua la stipo esas ne-ligna e maxim-multa-kaze verda}\senco{2}{Vejetanti verda qui kovras la prati e quin onu uzas por nutrar la bestiaro}\lingui{EFIS}
\artiklo{herboro}%
\vorto{herboro}\fako{bot.} Planto medicinala, sikigita o preparita (ye la formo ye qua olu vendesas da herboristo)\lingui{DEFIS}
\artiklo{herdo}%
\vorto{herdo} Parto petro-plakizita di la kameno-sulo, sur qua onu produktas fairo\lingui{DE}
\artiklo{heredar}%
\vorto{heredar}\fako{trans.) (yuro-cienco} Posedeskar ulo per sucedo di persono (generale ek la gepatri o parenti) qua mortis. – \textit{(anke metaf.)}
\artiklo{herezio}%
\vorto{herezio}\fako{religio katolika} Doktrino di ula dogmati Kristana qui esas kontrea a la kredo katolika e, do, kondamnita da la eklezio katolika\lingui{DEFIRS}
\artiklo{herisar}%
\vorto{herisar}\fako{trans.}\senco{1}{\textit{(aludante homo od animalo)} Igar rekta e rigida (sua hari, sua pili, sua plumi)}\senco{2}{Garnisar (objekto) ye kozi pikiva qui impedas tushar lu}\lingui{FS}
\artiklo{herisono}%
\vorto{herisono}\fako{zool.} Mamifero insektivora, di qua la dorso-pelo esas kovrita ye pikivi longa e rigida ek pili aglomerita quin lu herisas per bulatreskar kande onu proximigas su a lu\latina{erinaceus}\lingui{FS}
\artiklo{herkogamio}%
\vorto{herkogamio}\fako{bot.} Nomo donita da Axell por signifikar la tipo de flori qui esas organizita tale ke la autopolenizo impedesas e ke la transporto di la poleno adsur la stigmato povas eventar nur per la insekti
\artiklo{herkulo}%
\vorto{herkulo}\senco{1}{Mi-deo di la Panteono Romana, simbolo di forteso e vigoro}\senco{2}{\textit{(metaf.)} Persono (generale viro) di qua la forteso o vigoro korpala esas tre granda}\lingui{DEFIRS}
\artiklo{hermafrodita}%
\vorto{hermafrodita}\fako{homo} A qua onu atribuas la du sexui; (planto) qua unionas en su la organi di la du sexui\lingui{DEFIRS}
\artiklo{hermetika}%
\vorto{hermetika} Dicesas pri la klozeso di recipiento, klozeso perfekta qua (ideale) obtenesas per kunfuzar la bordi di la orifico kun olti di la stopilo\lingui{DEFIRS}
\artiklo{hernio}%
\vorto{hernio}\fako{patol.}\senco{1}{Tumoro quan produktas parto de la viceri abdominala qua ekiras tra orifico irga, kunportante parto di la peritoneo, ica formacante quaze sako en qua lu su lokizas}\senco{2}{Tumoro quan produktas la ekiro di vicero o di parto de vicero, de lua kavajo naturala tra aperturo irgequala, acidentala o naturala}\lingui{DEFIS}
\artiklo{herono}%
\vorto{herono}\fako{zool.} Granda ucelo ek la klaso « steltieri », di qua la beko esas tre longa, la gambi magra e tre longa, e qua su nutras generale ek fishi\lingui{DEFI}
\artiklo{heroo}%
\vorto{heroo}\senco{1}{\textit{(mitol.)} Mi-deo}\senco{2}{La viro qua su distingas per sua prodaji, sua granda valoro, sua devoteso}\senco{3}{La viro qua, ye ta o ca instanto, akaparas la atenco, la admiro, da omna kunhomi}\lingui{DEFIRS}
\artiklo{herpeto}%
\vorto{herpeto}\fako{patol.} Morbo konstitucionala ed heredebla, nekontagiala, quan karakterizas desordini nervala\lingui{DEFIS}
\artiklo{herso}%
\vorto{herso}\senco{1}{\textit{(agrokultivo)} Instrumento provizita ye denti fera o ligna, quan onu tranas sur sulo, pos plugo, por ruptar la glebi – pos semo, por kovrar la semini}\senco{2}{\textit{(fortifikajo)} Greto fera o ligna, garnisita infre ye pinti grosa, quan onu elevas od abasas segun-arbitrie an la pordo di fortifikajo}\lingui{F}
\artiklo{heteroblastika}%
\vorto{heteroblastika}\fako{biol.} \textit{(fenomeno)} qua konsistas ye la formaceso, en la sama punti, che animali inter-parenta, di organi interequivalanta, ma di qui la origino embriologiala esas tamen inter-diferanta
\artiklo{heterodoxo}%
\vorto{heterodoxo} Qua deviacas de la ortodoxeso\lingui{DEF}
\artiklo{heterofila}%
\vorto{heterofila}\fako{bot.} \textit{(planto)} di qua la folii o la kalico-folio esas kelke polimorfa
\artiklo{heterogama}%
\vorto{heterogama}\fako{biol.} Di qua la du elementi sexuala interdiferanta esas fuzita
\artiklo{heterogena}%
\vorto{heterogena} Qua esas di naturo altra kam to a quo lu esas unionita\lingui{DEFIRS}
\artiklo{heterogenezo}%
\vorto{heterogenezo}\fako{biol.} Produkto di ento vivoza, da enti vivoza pre-existanta, ma di qui la naturo diferas
\artiklo{heterokromosomo}%
\vorto{heterokromosomo}\fako{biol.} Kromosomo qua diferas de la cetera kromosomi normala di la spermatozoido per sua formo e sua dimensioni (plugrandeso o plumikreso)
\artiklo{heteromorfoso}%
\vorto{heteromorfoso}\fako{biol.} Regenero anormala di parto di la korpo
\artiklo{heterosexuala}%
\vorto{heterosexuala} Di qua la sexui esas nesama
\artiklo{heterostila}%
\vorto{heterostila}\fako{bot.} Di qua la stili esas longa ne-egale
\artiklo{heterotipa}%
\vorto{heterotipa}\fako{teratol.} Monstro duopla en qua un ek la unaji, acesora e parazita, nekomplete developita, esas enplantacita en la parto avana di la korpo di la unajo precipua
\artiklo{heuristiko}%
\vorto{heuristiko} Parto di la historio-cienco qua traktas la serchado di la dokumenti
\artiklo{hexaedra}%
\vorto{hexaedra}\fako{geom.} Sis-edra\lingui{DEF}
\artiklo{hexagono}%
\vorto{hexagono}\fako{geom.} Figuro geometriala kun sis anguli e sis lateri\lingui{DEF}
\artiklo{hexagramo}%
\vorto{hexagramo}\senco{1}{\textit{(filoz.)} Singla ek la 64 divino-kombinuri di la Chiniani, kombinuri ek du de la 18 trigrami dispozita vento-roze, e di qui singla konsistas ye la kombinuro ek tri linei}\senco{2}{Asemblajo ek sis literi}
\artiklo{hexametro}%
\vorto{hexametro}\fako{metriko} Sis-peda verso\lingui{DEFIRS}
\artiklo{hezitar}%
\vorto{hezitar}\fako{netrans., pri, pro} Ne-decidar pri la adopto di selekto, di rezolvo\lingui{DEFI}
\artiklo{hiacinto}%
\vorto{hiacinto}\senco{1}{\textit{(bot.)} Planto ek la familio « liliacei »}\senco{2}{\textit{(mineral.)} Lapido, varietato di topazo}\latina{hyacinthus}\lingui{DEFIRS}
\artiklo{hialoplasmo}%
\vorto{hialoplasmo}\fako{biol.} Parto homogena, amorfa e fluida di protoplasmo
\artiklo{hiato}%
\vorto{hiato}\senco{1}{\textit{(gram.)} Son-uno quan produktas la renkontro di du vokali, ica finanta vorto, ita komencanta la vorto sequanta}\senco{2}{\textit{(metaf.)} Nekontinueso inter la ceni di teatrajo, inter la gradi di genealogio-arboro, e c.}\lingui{DEFIRS}
\artiklo{hibernar}%
\vorto{hibernar}\fako{netrans.) (pri ula animali} Vivar dum vintro, partorporante\lingui{EF}
\artiklo{hibrida}%
\vorto{hibrida}\senco{1}{\textit{(planto od animalo)} qua devenas de du speci}\senco{2}{\textit{(gram.)} (vorto) qua konsistas ye elementi, di qui ica apartenas a linguo, ita a linguo diferanta}\lingui{DEFIRS}
\artiklo{hidalgo}%
\vorto{« hidalgo »} Nobelo Hispana qua asertas decendar de familio Kristana anciena\lingui{DEFIRS}
\artiklo{hidatodo}%
\vorto{hidatodo}\fako{biol.} Nomo kreita da Haberlandt por signifikar organo mikroskopala qua ekpulsas la aquo quan kontenas la planto
\artiklo{hido}%
\vorto{hido}\fako{kemio} Korpo simpla, gaso, di qua la kombino kun oxo produktas aquo\simbolo{H}\lingui{DEFIRS}
\artiklo{hidramnioto}%
\vorto{hidramnioto}\fako{anat.} Hidropso di la amniono; hiper-abundo di la liquido amnionala
\artiklo{hidranto}%
\vorto{hidranto} Buxo ek gisfero, qua kontenas junto-tubo e robineto por aquo-cherpo, instalita sive borde di la choseo, sive an muro, ed uzata da la extingisti di incendii\lingui{DE}
\artiklo{hidrato}%
\vorto{hidrato}\fako{kemio} Kompozajo derivita ek aquo, per substituco, o formacita per la kombino nemediata di substanco kun aquo\lingui{DEFIS}
\artiklo{hidrauliko}%
\vorto{hidrauliko} La parto di fiziko qua traktas la *leyi (legi) pri la movo da la liquidi\lingui{DEFIRS}
\artiklo{hidro}%
\vorto{hidro}\senco{1}{\textit{(mitol.)} Serpento sep-kapa, qua produktis sep altra kapi kande onu tranchis un ek la sep}\senco{2}{\textit{(zool.)} Aquo-serpento reptera; anke : mikra polipo qua habitas aquo sen-sala e su nutras ek mikra krustacei e dafni, quin lu kaptas per sua palpili}\latina{hydra}\lingui{DEFIRS}
\artiklo{hidrocefala}%
\vorto{hidrocefala}\fako{anat.} Hidropso di la kapo\lingui{DEF}
\artiklo{hidrodinamiko}%
\vorto{hidrodinamiko} La parto di hidrauliko qua traktas la movi e la presi da la fluidi\lingui{DEFIRS}
\artiklo{hidrofila}%
\vorto{hidrofila}\senco{1}{Qua esas avida ye aquo, qua absorbas aquo}\senco{2}{\textit{(materio koloida)} qua inflesas en aquo}
\artiklo{hidrofobio}%
\vorto{hidrofobio}\fako{patol.} Hororo di la liquidi, qua karakterizas ula morbi (kom ex. : rabio)\lingui{DEFIRS}
\artiklo{hidrogeno}%
\vorto{hidrogeno}\fako{kemio} Sinonimo di « hido »\lingui{DEFIRS}
\artiklo{hidrografio}%
\vorto{hidrografio} Cienco qua havas kom temo la studio e la deskripto di la riveri, fluvii e mari\lingui{DEFIS}
\artiklo{hidrografo}%
\vorto{hidrografo} La persono qua esas kompetenta pri hidrografio
\artiklo{hidroleucito}%
\vorto{hidroleucito}\fako{biol.} Vakuolo kontenata en la celuli di la protoplasmo di la vejetanti, provizita ye membrano propra (*tonoplasto) e qua su *reproduktas per dupartigo
\artiklo{hidrolizo}%
\vorto{hidrolizo} Reakto di duimesko di molekulo per la efiko da aquo\lingui{DEFIS}
\artiklo{hidrologio}%
\vorto{hidrologio} Parto di la natur-historio qua traktas aquo, e, precipue, la aqui mineraloza\lingui{DEFIS}
\artiklo{hidromancio}%
\vorto{hidromancio} Divinado per aquo\lingui{DEFIS}
\artiklo{hidromekaniko}%
\vorto{hidromekaniko} Cienco pri la aparati en qui aquo uzesas kom forco\lingui{DEFIS}
\artiklo{hidrometrio}%
\vorto{hidrometrio} La cienco qua traktas la mezuro di la denseso e di la fluo-rapideso di la liquidi\lingui{DEFIRS}
\artiklo{hidrometro}%
\vorto{hidrometro} Instrumento per qua onu mezuras la dikeso di la aquo-strato qua pluve falis dum un yaro en ta o ca loko\lingui{DEFIRS}
\artiklo{hidropatio}%
\vorto{hidropatio} Sistemo terapiala di la mediki qui probas kuracar nur per la uzo (sur ed en la korpo) di aquo pura\lingui{DEFIRS}
\artiklo{hidropso}%
\vorto{hidropso}\fako{patol.} Akumulajo morbala ek liquido seroza en kavajo di la korpo, od en la tisuo celulara\lingui{DEFIS}
\artiklo{hidroquinono}%
\vorto{hidroquinono}\fako{kemio} C₆H₄ (OH)₂
\artiklo{hidroskopio}%
\vorto{hidroskopio} La arto deskovrar la aquo-fonti\lingui{DFIRS}
\artiklo{hidroskopo}%
\vorto{hidroskopo} La persono qua profesione okupesas per la arto deskovrar la aquo-fonti\lingui{DEFIRS}
\artiklo{hidrostatiko}%
\vorto{hidrostatiko} La parto di hidrauliko qua traktas la equilibro di la gasi e di la liquidi, e la presi da li sur la parieti di la recipienti qui kontenas li\lingui{DEFIRS}
\artiklo{hidroterapio}%
\vorto{hidroterapio}\fako{medic.} Kuraco di ula morbi per apliko di aquo sur la korpo (balni, dushi, vapor-ago, e c.)\lingui{DEFIRS}
\artiklo{hieno}%
\vorto{hieno}\fako{zool.} Mamifero karnavida, digitigrada, analoga a volfo\latina{hyaena}\lingui{DEFIRS}
\artiklo{hierarkio}%
\vorto{hierarkio}\senco{1}{\textit{(kristanismo)} Ordino di la subordineso di la ciel-kori, di la anjeli}\senco{2}{Ordino di la subordineso di ti qui esas situita en rangi ne-interegala}\lingui{DEFIRS}
\artiklo{hieratika}%
\vorto{hieratika} Qua koncernas la kozi sakra. – \textit{(che la Egiptiani antiqua)}. (skribo-stilo) trasebla rapide, por la linguo hieroglifala\lingui{DEFIRS}
\artiklo{hiero}%
\vorto{hiero} La dio qua preiras nemediate olta en qua onu vivas\lingui{FIS}
\artiklo{hieroglifo}%
\vorto{hieroglifo}\senco{1}{\textit{(arkeol.)} En la skribo-stilo di la Egiptiani antiqua, signo sakra qua expresis simbole la idei, per la reprezento di ula objekti naturala}\senco{2}{\textit{(metaf.)} Kozo enigmata}\lingui{DEFIRS}
\artiklo{high life}%
\vorto{« high life »} Granda mondumo, ek la aristokrati di la socio
\artiklo{higieno}%
\vorto{higieno} Parto de medicino, qua traktas la rejimo adoptenda da la homo por konservar lua saneso\lingui{DEFIRS}
\artiklo{higrometrio}%
\vorto{higrometrio} Parto di fiziko, qua mezuras la quanto de aquo-vaporo quan kontenas atmosfero\lingui{DEFIRS}
\artiklo{higrometro}%
\vorto{higrometro} Aparato per qua onu mezuras la quanto de aquo-vaporo qua existas en atmosfero, per korpo quan humideso igas modifikesar pezale, voluminale\lingui{DEFIRS}
\artiklo{higroskopo}%
\vorto{higroskopo} Aparato per qua onu konstatas la humideso-grado di atmosfero\lingui{DEF}
\artiklo{hike}%
\vorto{hike}\senco{1}{En la loko ube esas la persono qua pronuncas ca vorto}\senco{2}{En ica loko di la dico}\lingui{L}
\artiklo{hilo}%
\vorto{hilo}\senco{1}{\textit{(bot.)} Marko di la ligo-punto di la grano sur la *funikulo}\senco{2}{\textit{(anat.)} Ligo-punto ube vaskulo juntesas a vicero}\lingui{EFIS}
\artiklo{hiloto}%
\vorto{hiloto}\senco{1}{\textit{(epoki antiqua)} Sklavo, che la Spartani}\senco{2}{La viro, en socio, qua esas sur la grado maxim infra di la abjekteso-skalo}\lingui{DEFIRS}
\artiklo{himeneo}%
\vorto{himeneo} Kanto mariajala, che la Greki antiqua, e che la Romani lora\lingui{DEFIRS}
\artiklo{himenio}%
\vorto{himenio}\fako{bot.} Che la fungo, la mikra membrano en qua esas la spori
\artiklo{himeno}%
\vorto{himeno}\senco{1}{\textit{(anat.)} Membrano qua klozas parte la vagino, che la virgini}\senco{2}{\textit{(bot.)} Pelikulo di la korolo qua laceresas lor la desklozeso}\lingui{DEFIS}
\artiklo{himenoptero}%
\vorto{himenoptero}\fako{zool.} Insekto di qua la ali esas membrana (abeli, vespi, formiki, e c.)\lingui{DEFIRS}
\artiklo{himno}%
\vorto{himno}\senco{1}{Poemo religiala, honore di la dei, di la heroi}\senco{2}{Poemo lirika, honore di la deajo}\senco{3}{Kantiko Latina quan onu kantas o recitas en kirko}\lingui{DEFIRS}
\artiklo{hioido}%
\vorto{hioido}\fako{anat.} \textit{(osto)} qua esas horizontale inter la bazo di la lango e la laringo\lingui{EFIS}
\artiklo{hipalago}%
\vorto{hipalago}\fako{retor.} Stilo-figuro qua atribuas a vorto to quo konvenas ad altra vorto. (kom ex. : retrodonar ulu a vivo, vice : retrodonar vivo ad ulu)\lingui{DEFIS}
\artiklo{hipar}%
\vorto{hipar}\fako{netrans.} Sentar ta kontrakto intervalopa qua sukusas la diafragmo ed akompanesas da bruiso ne-artikulata partikulara\lingui{S}
\artiklo{hiper-}%
\vorto{hiper-} Prefixo qua signifikas « super, trans, ecese »
\artiklo{hiperbato}%
\vorto{hiperbato}\fako{retor.} Figuro qua konsistas ye permutar la ordino naturala di ula vorti en propoziciono\lingui{DEFIS}
\artiklo{hiperbolo}%
\vorto{hiperbolo}\senco{1}{\textit{(geom.)} Kurvo qua esas la loko di la punti qui distas de du punti fixa konstante}\senco{2}{\textit{(retor.)} Stilo-figuro qua konsistas ye exajerar la expreso-termino}\lingui{DEFIRS}
\artiklo{hiperboloido}%
\vorto{hiperboloido}\fako{geom.} Surfaco di la grado duesma, kun centro unika, produktita da hiperbolo\lingui{DEFIRS}
\artiklo{hiperboreo}%
\vorto{hiperboreo} Habitanto di la regioni qui esas extrem-norde\lingui{DEF}
\artiklo{hiperemio}%
\vorto{hiperemio}\fako{patol.} Prezenteso di tro multa sango en organo o parto de organo
\artiklo{hiperestezio}%
\vorto{hiperestezio}\fako{patol.} Hiper-sentiveso
\artiklo{hipermetropo}%
\vorto{hipermetropo}\fako{patol.} Qua afektesas da ta anomaleso di refrakto pro qua la lumo-radii paralela, vice konvergar a la retino, formacas sua foko trans ica kande la binokl-efiko ne intervenas\lingui{DEFS}
\artiklo{hipertonika}%
\vorto{hipertonika}\fako{biol.} Dicesas pri la *solutei injektebla di qui la koncentreso-grado molekulala superesas olta di la sero sangala
\artiklo{hipertrofiar}%
\vorto{hipertrofiar}\fako{trans.} Produktar hiper-kresko che organo o parto de organo, sen altero minboniganta\lingui{DEFIRS}
\artiklo{hipnotar}%
\vorto{hipnotar}\fako{netrans.} Submisesar a dormo artificala, produktita da la vido longa-dura di objekto briloza\lingui{DEFIRS}
\artiklo{hipo-}%
\vorto{hipo-} Prefixo qua signifikas « sub-, diminutanta, nesuficanta »
\artiklo{hipocikloido}%
\vorto{hipocikloido}\fako{geom.} Kazo en qua la cirkli movanta di la epicikloido esas extera\lingui{DEFIRS}
\artiklo{hipodermo}%
\vorto{hipodermo}\fako{anat.} Tisuo sub-derma
\artiklo{hipodromo}%
\vorto{hipodromo}\senco{1}{\textit{(en la epoki antiqua)} Cirko por la kavalo-konkursi e la charo-konkursi}\senco{2}{Agro por la kaval-konkursi}\lingui{DEFIRS}
\artiklo{hipofizo}%
\vorto{hipofizo}\fako{anat.} Organo glanda, an la edro infra di la cerebro ed inkluzita en la kavajo di la osto sfenoida, infre\lingui{DEFIRS}
\artiklo{hipogastro}%
\vorto{hipogastro}\fako{anat.} Infra parto di la ventro\lingui{DEFIS}
\artiklo{hipogeo}%
\vorto{hipogeo}\fako{arkeol.} Sepulteyo masonura sub-sula\lingui{DEFI}
\artiklo{hipogrifo}%
\vorto{hipogrifo}\fako{mitol.} Animalo fantastika, mi-kavala, mi-grifona, rezultinta de la kopulaco da grifono kun kavalino\lingui{DEFIRS}
\artiklo{hipokampo}%
\vorto{hipokampo}\fako{zool.} Genero de fishi, tipo de la familio « hipokampidei », kun formo longa, korpo dividita ad-ringe, e poligonala-secione, kovrita ye plaki harda\latina{hippocampus}\lingui{EFIS}
\artiklo{hipokondrio}%
\vorto{hipokondrio}\fako{patol.} Stando morbala qua igas melankoloza, trista\lingui{DEFIRS}
\artiklo{hipokondro}%
\vorto{hipokondro}\fako{anat.} Singla de la du parti laterala di la regiono epigastrala, sub la « falsa kosti »\lingui{DEFIRS}
\artiklo{hipokrito}%
\vorto{hipokrito} Persono qua fingas havar sentimenti religiala, vertuala\lingui{DEFIRS}
\artiklo{hipopiono}%
\vorto{hipopiono}\fako{patol.} Pusifo en la chambro avana di la okulo
\artiklo{hipopotamo}%
\vorto{hipopotamo}\fako{zool.} Mamifero pakiderma, kun korpo enorma, kovrita da pelo senpila tre dika, kurta-gamba, quan onu renkontras en la granda fluvii di la sudo-regioni di Amerika. Afrika\latina{hippopotamus}\lingui{DEFIRS}
\artiklo{hipostazo}%
\vorto{hipostazo}\fako{filoz.}\senco{1}{Persono. (Segun la teologiisti katolika, en Deo existas tri hipostazi ed un naturo, ed en Iesu Kristo un hipostazo e du naturi)}\senco{2}{\textit{(segun la filozofio-sistemi Alexandriana o neoplatonala)} La principo deala}\lingui{DEFIR}
\artiklo{hipoteko}%
\vorto{hipoteko}\fako{yuro-cienco} Yuro di prestinto, ad imoblo quan la prestario destinis a la pago di sua debajo – qua sequas ta imoblo che omna olua proprieteri\lingui{DEFIRS}
\artiklo{hipotenuzo}%
\vorto{hipotenuzo}\fako{geom.} Latero di triangulo ort-angula, opozita a la orta angulo\lingui{DEFIRS}
\artiklo{hipotezo}%
\vorto{hipotezo} Supozo, de qua onu deduktas konsequi verifikenda\lingui{DEFIRS}
\artiklo{hipsometrio}%
\vorto{hipsometrio} Mezuro di la altitudo di loko, per observi barometrala o geodeziala\lingui{DEFIS}
\artiklo{hipsometro}%
\vorto{hipsometro}\fako{fiziko} Aparato per qua onu mezuras la altitudo di loko, per la temperaturo ye qua bolieskas aquo ibe\lingui{DEFIS}
\artiklo{hirundo}%
\vorto{hirundo}\fako{zool.} Migr-ucelo, de la klaso « paseri », qua arivas en printempo, nestifas en la kameno-tubi o sub la tekto-karpenturi, ed autune migras a landi varma\lingui{FISL}
\artiklo{hisar}%
\vorto{hisar}\fako{trans.}\senco{1}{Suprigar per tiro (seglo, barelo, e c.)}\senco{2}{\textit{(metaf.)} Hisar su : irar adsupre od ad-sure penoze}\lingui{DEFIS}
\artiklo{hiskiamo}%
\vorto{hiskiamo}\fako{bot.} Planto de la familio « solanei », di qua la speco maxim ordinara esas la hiskiamo nigra, veneno tre efikiva qua esas narkotigiva\latina{hyosciamus niger}\lingui{FI}
\artiklo{histerezo}%
\vorto{histerezo}\fako{elektrotekniko} Proprajo karaktera di la substanci fero-magnetala : lia indukteso, en ca o ta instanto, dependas (en la specimeno konsiderata) ne nur de la feldo magnetala qua esas efikanta, ma anke de la estadi magnetala antea di la specimeno\lingui{DEFIS}
\artiklo{histerezometro}%
\vorto{histerezometro}\fako{elektrotekniko} Aparato per qua onu mezuras la qualeso di la fero uzata elektrale, de la vidopunto « histerezo »\lingui{DEFIS}
\artiklo{histerio}%
\vorto{histerio}\fako{patol.} Morbo nervala quan karakterizas sufoki, konvulsi, e c., judikita (nejuste) kom partikulara ye la homini\lingui{DEFIRS}
\artiklo{histerotomio}%
\vorto{histerotomio}\fako{kirurg.} Disseko di la utero\lingui{DEF}
\artiklo{histogenezo}%
\vorto{histogenezo}\fako{biol.} Formaceso e kresko di la tisui diversa di la embriono, de la blastomeri komencala\lingui{DEFIRS}
\artiklo{histolizo}%
\vorto{histolizo}\fako{biol.} Destrukto di la tisui vivanta, precipue en la komenco di la metamorfosi
\artiklo{histologio}%
\vorto{histologio} La parto di anatomio qua studias la formaco, la evoluciono e la kompozeso di la tisui vivanta\lingui{DEFIRS}
\artiklo{historio}%
\vorto{historio}\senco{1}{Biografio di populo, di naciono}\senco{2}{Studio di la enti diversa qui existas en naturo}\lingui{EFIRS}
\artiklo{historiografio}%
\vorto{historiografio} La arto di la historiografo\lingui{DEFIRS}
\artiklo{historiografo}%
\vorto{historiografo} La persono qua komisesis oficale da princo, pri kompozar la historio di la regno-duro di ica\lingui{DEFIRS}
\artiklo{histriono}%
\vorto{histriono}\senco{1}{\textit{(en la epoki antiqua)} Aktoro qua pleis farsi grosiera}\senco{2}{Bufono en la ferii}\senco{3}{Komediopleistacho}\lingui{DEFIS}
\artiklo{hivernar}%
\vorto{hivernar}\fako{netrans.} Travivar la vintro-tempo shirmate\lingui{EFIS}
\artiklo{ho!}%
\vorto{ho!}\senco{1}{Interjeciono per qua onu advokas}\senco{2}{Interjeciono qua expresas sprico di joyo, di doloro, di timo}\senco{3}{Interjeciono qua expresas surprizeso, advoko subita, sprico da la psiko kande lu emocas forte}\senco{4}{Interjeciono per qua onu expresas indigneso od astoneso}
\artiklo{*ho-}%
\vorto{*ho-} Prefixo qua signifikas « (radiko) dum qua ni vivas nun » (kom ex. : hodie, homonate, hosemane)
\artiklo{*hobio}%
\vorto{*hobio} Ideo quan onu preferas traktar, o pri qua onu preferas informesar; kozo pri qua onu su okupas prefere e kom distraktivo, amuzivo\lingui{E}
\artiklo{hoboyo}%
\vorto{hoboyo} Muzik-instrumento vento-funcionanta, langetoza, sen beko, kono-forma, qua finas per mikra funelo, e qua emisas foni klara e tre dolca\lingui{DEFIRS}
\artiklo{hodografo}%
\vorto{hodografo}\fako{matem.} De punto fixa O, onu trasas singlainstante vektoro qua equipolas la rapideso-vektoro. La extremajo di ta vektoro deskriptas kurvo qua esas la « hodografo » di ta movado
\artiklo{hodometro}%
\vorto{hodometro}\senco{1}{Instrumento horlojo-forma, di qua la indexo quik kande lu movesas, indikas la quanto de pazi ed anke la disto parkurita da ped-iranto dum ta o ca quanto de tempo}\senco{2}{Instrumento uzata por kontar la quanto de rotaci di roto, di mashin-arboro, en ta o ca tempo}\lingui{DEFIS}
\artiklo{hoko}%
\vorto{hoko}\senco{1}{Instrumento kun la extremajo kurva, per qua onu povas tirar ulo ad su}\senco{2}{Mikra peco kurvoza an qua onu povas suspendar ulo}\senco{3}{\textit{(metaf.)} Ulo kurvoza quale hoko}\lingui{DE}
\artiklo{hola!}%
\vorto{hola!} Interjeciono por advokar, por incitar ad halto
\artiklo{holdo}%
\vorto{holdo}\fako{navig.} La parto maxim infra di la internajo di navo\lingui{E}
\artiklo{holoblastika}%
\vorto{holoblastika}\fako{biol.} Dicesas pri la ovo qua su dividas tote ed en parti inter-egala (ovo *alecita, sen nutro-provizuri.)
\artiklo{holografa}%
\vorto{holografa}\fako{pri testamento} Tote skribita per la manuo di la testamentanto\lingui{DEFIS}
\artiklo{holokausto}%
\vorto{holokausto} Che la Hebrei antiqua : sakrifiko en qua la viktimo esis tote kombustata\lingui{EFIS}
\artiklo{holometro}%
\vorto{holometro}\fako{matem.} Instrumento per qua onu mezuras la altesoangulo di punto super horizonto
\artiklo{holoturio}%
\vorto{holoturio}\fako{zool.} Genero de *radiarii ekinoderma\latina{Holo thuria}\lingui{DEFIRS}
\artiklo{homajo}%
\vorto{homajo}\senco{1}{\textit{(epoki feudala)} Ago da la vasalo qua su deklaras devota a lua sinioro e promisas ad ica servor lu fidele e sen restrikto}\senco{2}{Ago submisa e respekta}\senco{3}{Respektosigno humilesala}\lingui{EFI}
\artiklo{homardo}%
\vorto{homardo}\fako{zool.} Genero de krustacei dek-peda, di qua la du gambi avana esas grosega e similesas pinci\latina{homarus vulgaris}\lingui{DFI}
\artiklo{homeopatio}%
\vorto{homeopatio}\fako{medic.} Sistemo terapiala – da S. Hahnemann, 1755-1843 qua konsistas ye traktar la morbi per dozi medikamentala infinitezima, qui produktas, en la individuo sana, la simptomi di morbo analoga ad olta quan onu probas kombatar\lingui{DEFIRS}
\artiklo{homeosexuala}%
\vorto{homeosexuala}\fako{patol.} Individuo qua sentas afineso sexuala nur kun la personi di sua propra sexuo\lingui{DEF}
\artiklo{homerala}%
\vorto{homerala}\fako{rido} Nerepresebla (quan Homeros, en la fino di la kanto unesma di \textit{Iliado}, atribuas a la dei)
\artiklo{homilio}%
\vorto{homilio}\senco{1}{\textit{(liturgio katolika)} Instruktado familiara di la populo pri evangelio e pri la materii religiala}\senco{2}{Breviario-leciono extraktita de la homilii di la patri ekleziala}\senco{3}{\textit{(desprizoze)} Etiko-prediko, leciono enoyiganta}\lingui{DEFIS}
\artiklo{homo}%
\vorto{homo}\fako{zool.} Mamifero bimanua, qua stacas vertikale, dotita ye la fakultato parolar e ye raciono\latina{homo}\lingui{EFIS}
\artiklo{homocentro}%
\vorto{homocentro}\fako{matem.} Centro komuna di plura centri
\artiklo{homogena}%
\vorto{homogena} Di qua la elementi konstitucanta esas sama-natura\lingui{DEFIRS}
\artiklo{homografio}%
\vorto{homografio}\fako{geom.} Interdependo partikulara di du figuri geometriala
\artiklo{homologa}%
\vorto{homologa}\fako{aludante un o plura termini} Qui interkorespondas\lingui{DEFIRS}
\artiklo{homologio}%
\vorto{homologio}\fako{geom.} Karaktero di la figuri tala (segun Poncelet) ke la punti interkorespondanta di ica e di ita esas duopa sur la rekti qui kunkuras ad un punto (centro di homologio) e ke la rekti qui juntas du punti ek ica e la du korespondanti di la altra iras krucumar sur rekto una (axo di homologio) se la figuri esas plana, o sur un plano (plano di homologio) se la figuri esas irgequala en spaco
\artiklo{homonima}%
\vorto{homonima}\fako{gram.} Vorti qui esas intersama ortografiale ed pronuncale, ma interdiferas pri la senco.\lingui{DEFIRS}\subvorto{}{homonimo}{Persono di qua la nomo esas sama kam olta di altra persono}{}{}
\artiklo{homoteta}%
\vorto{homoteta}\fako{geom.} Donite un punto fixa O ed un konstanto k, se, ad un punto M ek la spaco onu korespondantigas un punto M', tala ke \textit{OM'/OM} egalesas k, onu dicas ke la punti M e M' esas « homoteta » pri la punto O
\artiklo{homotipa}%
\vorto{homotipa}\fako{anat.} Di qua la organi, en un individuo, esas la analogi di altra organi. (kom ex. : la pedo-fingri kompare kun la manuo-fingri)
\artiklo{honesta}%
\vorto{honesta} Qua esas konforma a la devo, a vertuo\lingui{EFIS}
\artiklo{honorario}%
\vorto{honorario} Pago-pekunio po la servi da persono qua havas profesiono honorizanta (advokato, mediko, sacerdoto, e c.)\lingui{DEFIRS}
\artiklo{honoro}%
\vorto{honoro}\senco{1}{Digneso etikala qua naskas de la bezono di estimeso da la kunhomi e da ni ipsa}\senco{2}{Distingo-manifesto por ulu, flatema}\lingui{EFIS}
\artiklo{hop!}%
\vorto{hop!} Interjeciono uzata por stimular, ecitar a kuresko, o saltigar trans obstakli
\artiklo{hoplito}%
\vorto{hoplito}\fako{en la epoki antiqua} Infantriano di qua la armi esis tre pezoza\lingui{DEFIRS}
\artiklo{*horario}%
\vorto{*horario} Tabelo qua indikas la kloki di la departo ed arivo pri la treni, tramveturi, autobusi, aeroplani, e c. – o la tempopunti ye qui ulo agesos, facesos (segun programo)\lingui{F}
\artiklo{hordeo}%
\vorto{hordeo}\fako{bot.} Planto cereala singlayara, herbatra, di qua la stipo perpendikla esas garnisita ye folii alterna, lineatra. e la flori dispozita spiko-forme. – Grano quan produktas ta planto, uzata por facar pano-sorto grosiera, nutrar la kavali, la brutaro, fabrikar biro, e c.\latina{hordeum}\lingui{L}
\artiklo{hordo}%
\vorto{hordo}\senco{1}{\textit{(che la Tatari)} Homo-trupo vaganta, nomada}\senco{2}{Trupo vaganta de homi qui vivas socie, ma sen rezideyo fixa}\senco{3}{\textit{(metaf.)} Trupo de homi qui perturbas la ordino sociala, omnaforme}\lingui{DEFIRS}
\artiklo{horizonto}%
\vorto{horizonto}\senco{1}{Limito di la vido omna-direcione, che la observanto; lineo cirklo-kurva di qua lu esas la centro, ed ube cielo semblas juntar su a la Terglobo}\senco{2}{Parto di la Tersurfaco e di la sfero ciela quan limitizas ta lineo cirklo-kurva}\senco{3}{\textit{(metaf.)} La cirklo en qua su movas la mento, portee di la regardo da ica}\lingui{DEFIRS}
\artiklo{horlojo}%
\vorto{horlojo} Instrumento uzata por indikar qua kloko esas\lingui{EFIS}
\artiklo{hornblendo}%
\vorto{hornblendo} Silikato naturala di aluminio, di kalio, di fero e di magnezo, apartenanta a la genero « amfibola »; ta mineralo esas opaka, obskur-verda, od obskur-bruna o perfekte nigra\lingui{DEF}
\artiklo{horniso}%
\vorto{horniso}\fako{zool.} Granda vespo ferea, qua militas kontre la abeli por furtar de ici lia mielo\latina{vespa crabro}\lingui{DE}
\artiklo{horo}%
\vorto{horo} La parto 24-esma di un dio; « horo » mezuras duro o tempo-intervalo\lingui{EFIS}
\artiklo{hororar}%
\vorto{hororar}\fako{trans.}\senco{1}{Fremisar pro repulseso e pavoro, di qua la kauzo esas la vido di objekto abomininda}\senco{2}{\textit{(aludante psiko)} Repulsesar profunde, da persono o da kozo}\senco{3}{Movesar da sentimento di respekto e pavoro koram spektaklo}\lingui{EFIS}
\artiklo{horoskopo}%
\vorto{horoskopo} Observo di la situeso di ula planeti, di ula stelari, lor la nasko di infanteto, segun qua la astrologi asertis pre-anuncar la futuro di ta jus-naskinto\lingui{DEFIS}
\artiklo{hortensio}%
\vorto{hortensio}\fako{bot.} Arboreto de la familio « saxifragei », kultivata kom plezur-planto\latina{hydrangea hortensia}\lingui{DFRS}
\artiklo{hospico}%
\vorto{hospico}\senco{1}{Gastigo-habiteyo kreita da regulieri por la regulieri stranjera, la pilgrimanti, la voyajanti}\senco{2}{Gastigo-domo por infirmi, infanti, oldi}\lingui{DEFIS}
\artiklo{hospitalo}%
\vorto{hospitalo} Gastigo-domo ube aceptesas la maladi mizeroza por kuracar li\lingui{DEFIRS}
\artiklo{hostio}%
\vorto{hostio}\fako{liturgio katolika} Pano-peco dina, sen-hefa, destinita a la meso-sakrifiko\lingui{EF}\subvorto{}{la santa hostio}{La pano konsakrita e chanjita aden la korpo di Iesu Kristo}{}{}
\artiklo{hosto}%
\vorto{hosto} La persono qua gastigas\lingui{EFIRS}
\artiklo{hotelo}%
\vorto{hotelo} Domo di qua onu lugas la chambri, la apartmenti, mobloza ed ube onu restoras la voyajeri\lingui{DEFR}
\artiklo{hozanna}%
\vorto{« hozanna »}\senco{1}{\textit{(liturgio Hebrea)} Refreno di himno di la sinagogo}\senco{2}{\textit{(liturgio katolika)} Himno quan onu kantas en la Palmo-sundio, qua komencas per la vorto « Hozanna »}
\artiklo{hu!}%
\vorto{hu!} Interjeciono uzata por shamigar, por mokar, ulu
\artiklo{hufo}%
\vorto{hufo} Envelopilo korno-harda qua cirkondas la pedo di la rumineri, pakidermi, solipedi\lingui{DE}
\artiklo{hugenoto}%
\vorto{hugenoto}\fako{uzita desprizoze da la katoliki} Partisano di la herezio da Calvin\lingui{DEFIRS}
\artiklo{hum!}%
\vorto{hum!} Interjeciono por indikar desfido, dubito
\artiklo{*humanismo}%
\vorto{*humanismo} Doktrino di la *humanisti di Renesanco, qui igis tre prizata itere la lingui e la literaturo-verki di la Grekia e la Roma antiqua\lingui{DEFIRS}
\artiklo{*humanisto}%
\vorto{*humanisto} Persono qua par-savas la lingui, ed esas erudita pri la literaturo-verki, di la Grekia e la Roma antiqua\lingui{DEFIRS}
\artiklo{humero}%
\vorto{humero}\fako{anat.} Osto di la brakio, de la shultro til la kudo\lingui{EFIS}
\artiklo{humida}%
\vorto{humida} Qua kontenas en sua materio substanco aqua od aquatra\lingui{EFIS}
\artiklo{humila}%
\vorto{humila} Qua su abasas volunte koram la kunhomi\lingui{EFIS}
\artiklo{humoro}%
\vorto{humoro}\senco{1}{Liquido qua renkontresas en korpo organika}\senco{2}{Ta humoro egardata kom alterita e kom la kauzo di ula morbi}\senco{3}{Temperamento, karaktero}\lingui{EFIRS}
\artiklo{humuro}%
\vorto{humuro} Gayeso flegmoza, ulakaze akompanata da ironio\lingui{DEFIRS}
\artiklo{humuso}%
\vorto{humuso} Materio bruna o nigratra, qua produktesas en sulo, da la deskompozeso di la materii organika\lingui{DEFIS}
\artiklo{hundo}%
\vorto{hundo}\fako{zool.} Genero de mamifero digitigrada, domestika, qua aboyas; e di qua existas multa varietati\latina{canis familiaris}\lingui{DE}
\artiklo{hungrar}%
\vorto{hungrar}\fako{trans., ulo} Bezonar manjo\lingui{DF}
\artiklo{huo}%
\vorto{huo}\fako{zool.} Rapt-ucelo noktala quan onu dresis por chasado\latina{athene noctua}\lingui{F}
\artiklo{hura!}%
\vorto{hura!}\senco{1}{Klamo da la navani Britaniana honorize di chefo, di personego qua eniras la navo}\senco{2}{Aklamo-klamo honorize di ulu}
\artiklo{hurdo}%
\vorto{hurdo} Ozier-treliso ajuroza (sur qua onu sikigas frukti, e c., od uzata kom greto por la brutaro-parki)\lingui{DE}
\artiklo{huronala}%
\vorto{huronala}\fako{geol.} Nomo di la strati ek graniti e gneisi ye qui konsistas la somito di la arkeano
\artiklo{husaro}%
\vorto{husaro}\senco{1}{Armeo-kavaliero Hungariana}\senco{2}{Soldato ek la kavalrio lejera, di qua la uniformo (en Francia) kelke similesas olta di la armeo-kavalieri Hungariana}\lingui{DEFIS}
\artiklo{hutingo}%
\vorto{hutingo}\fako{zool.} Fisho qua vivas en aquo sen-sala, de la familio « salmonidei », qua habitas la granda lagi di Mez-Europa (Suisia, Germania, nordo-regioni di Italia)\latina{coregonus oxyrhynchus}
\end{multicols}
\sekciono{I}
\begin{multicols}{2}
\artiklo{-i-}%
\vorto{-i-} Sufixo qua indikas « lando, regiono, domeno qua dependas de… e, konseque, la resortiso di, la autoritato di… »
\artiklo{-i}%
\vorto{-i} Dezinenco plurala di la nomi e di la pronomi
\artiklo{iambo}%
\vorto{iambo}\senco{1}{\textit{(metriko antiqua)} Pedo ek du silabi, la unesma, kurta; la duesma, longa}\senco{2}{Verso satirala, qua konsistis prime ek sis iambi}\senco{3}{Versajo satira}\lingui{DEFIRS}
\artiklo{ibe}%
\vorto{ibe} En ta loko (quan onu indikas), kontraste kun olta en qua onu esas
\artiklo{ibiso}%
\vorto{ibiso}\fako{zool.} Ucelo steltiera, longa-beka, di qua un speco esis sakra, che la Egiptiani antiqua\latina{ibis}\lingui{DEFIRS}
\artiklo{ica}%
\vorto{ica} ca
\artiklo{-id-}%
\vorto{-id-} Sufixo qua signifikas « homo, decendanto di… »
\artiklo{idanto}%
\vorto{idanto}\fako{biol.} Segun la teorio di Waismann : partikuli qui reprezentas e transmisas atavisme la partikularaji di la strukturo di la enti vivoza
\artiklo{identa}%
\vorto{identa} Di qua la naturo esas absolute sama kam olta di altra kozo\lingui{DEFIRS}
\artiklo{*identifikar}%
\vorto{*identifikar}\fako{trans.} Komparar persono kun la deskripto o fotografuro donata da lua pasporto od altra dokumento, por divenar *certena ke lu esas ya la persono quan lu afirmas esar
\artiklo{ideo}%
\vorto{ideo} Imajo, en spirito, de ula ento o de ula eso-maniero\lingui{DEFIRS}
\artiklo{ideografio}%
\vorto{ideografio} Reprezento di la idei per signi qui esas la imajo de la objekto\lingui{DEFIRS}
\artiklo{ideologio}%
\vorto{ideologio}\fako{filoz.} La cienco pri la origino, la formac(es)o di la idei\lingui{DEFIRS}
\artiklo{idilio}%
\vorto{idilio}\senco{1}{Mikra poemo pastorala}\senco{2}{Amoro-naraco pastorala}\lingui{DEFIRS}
\artiklo{idioblasto}%
\vorto{idioblasto}\fako{bot.} \textit{(segun Sachs)}, celuli neregulale branchigita, izolita meze di tisuo tre diferanta
\artiklo{idiomo}%
\vorto{idiomo}\senco{1}{Linguo di naciono}\senco{2}{Dialekto di provinco}\lingui{DEFIRS}
\artiklo{idioplasmo}%
\vorto{idioplasmo}\fako{biol.} Parto di la plasmo, qua kontenas omna partikuli qui reprezentas la karakterizivi di la speco
\artiklo{idiota}%
\vorto{idiota}\fako{patol.} Di qua la cerebro esas nur hipo-kreskinta; di qua la kompreniveso esas tre poka\lingui{DEFIRS}
\artiklo{idiotismo}%
\vorto{idiotismo} Frazo-formo di qua la senco ne esas deskovrebla per la analizo, pro ke olu ne rezultas logikale de la senco e di la kombinuro de la elementi uzata; anke pri vorto e lua elementi. \textit{(Louis Couturat)}\lingui{DEFIRS}
\artiklo{idolo}%
\vorto{idolo} Figuro, statuo, qua reprezentas deo o deajo, e qua kultesas. – \textit{(metaf.)} Persono o kozo por qua onu profesas quaza kulto, amo, amoro pasionoza\lingui{DEFIRS}
\artiklo{-ier-}%
\vorto{-ier-}\senco{1}{Sufixo qua signifikas : I. Etuyo qua kontenas o sustenas la kozo indikata da la radiko}\senco{2}{Karakterizata da… (ulo extera)}\senco{3}{Arboro, arboreto, arbusto, qua produktas la frukto di qua la nomo indikesas da la radiko}
\artiklo{-if-}%
\vorto{-if-} Sufixo qua signifikas « produktar, genitar, sekrecar » (to quon indikas la radiko)
\artiklo{-ig-}%
\vorto{-ig-}\senco{1}{Sufixo qua signifikas : I. \textit{kun radiko nomala} : « donar la qualeso quan expresas la radiko »; « transformar a… »}\senco{2}{Kun \textit{radiko di verbo netransitiva} : « esar la kauzo di la… -ar ». – « Igar » uzesas anke kom radiko, kun la senco : « esar la kauzo di la… (ago, od ageso, indikata da la verbo). »}\lingui{D}
\artiklo{ignorar}%
\vorto{ignorar} Refuzar saveskar, konoceskar; vole nesavar, vole nekonocar\lingui{DEFIRS}
\artiklo{iguano}%
\vorto{iguano}\fako{zool.} Granda saurio di la tropiko-regioni di Amerika\latina{iguana}\lingui{DEFIRS}
\artiklo{-ik-}%
\vorto{-ik-} Sufixo qua signifikas « qua esas malada de…, pro… »
\artiklo{ikneumono}%
\vorto{ikneumono}\fako{zool.} Insekto himenoptera qua perforas la pelo di la raupi por depozar en oli sua ovi\latina{ichneumon}\lingui{EFIS}
\artiklo{ikono}%
\vorto{ikono}\fako{en Rusia caroza, ed en la eklezio Greka} Imajo kolorizita qua reprezentas la virgino Maria (di la katoliki) e la santi\lingui{EFIRS}
\artiklo{ikonogeno}%
\vorto{ikonogeno}\fako{fotogr.} Natro-salo ek la acido amino-1-naftolo-2, sulfonala-6, uzata kom revelilo fotografala\lingui{DE}
\artiklo{ikonografio}%
\vorto{ikonografio}\senco{1}{La cienco pri la monumenti figuroza, pikturi, statui, medalii, e c.}\senco{2}{Kolektajo ek portreti de famozi}\lingui{DEFIRS}
\artiklo{ikonoklasto}%
\vorto{ikonoklasto}\fako{historio religiala} Hereziano qua ruptas la santa imaji, qua ne admisas la reprezento figurala de la personi dea(la)\lingui{DEFIRS}
\artiklo{ikonometro}%
\vorto{ikonometro}\fako{fotogr.} Vizilo per qua onu povas determinar la disto fokala di la objektalo fotografala, apta donar (ye dimensioni donita) la imajo de fotografoto o de fotografotajo\lingui{DEFIS}
\artiklo{ikosaedro}%
\vorto{ikosaedro}\fako{geom.} Solido 20-edra\lingui{DEFIS}
\artiklo{iktero}%
\vorto{iktero}\fako{patol.} Morbo qua flavigas la pelo, qua devenas de perturbeso di la sekreco e di la fluo di la bilo\lingui{DEFIS}
\artiklo{iktiokolo}%
\vorto{iktiokolo} Substanco diafana, solida, preparita ek la natoveziko di la esturjoni, uzata por facar kapsuli gelatinatra, o la gumo-tafto\lingui{DEFI}
\artiklo{iktiologio}%
\vorto{iktiologio} La parto di la natur-historio, qua traktas la fishi
\artiklo{iktiosauro}%
\vorto{iktiosauro} Reptero fosila, kun fisho-vertebri\latina{ichtyosaurus}\lingui{DEFIRS}
\artiklo{iktus}%
\vorto{« iktus »}\fako{metriko} Frapo qua perceptigas la tempo forta di verso-pedo
\artiklo{iktuso}%
\vorto{iktuso}\fako{patol.} Afekto subita, qua frapas quale stroko. (kom ex. : iktuso apoplexiala, epilepsiala)
\artiklo{-il-}%
\vorto{-il-} Sufixo qua signifikas « instrumento por… » (ne « mashino por… »)
\artiklo{il(u)}%
\vorto{il(u)} Pronomo personala maskulala, persono triesma (uzata kande onu volas emfazar la sexuo)
\artiklo{ileono}%
\vorto{ileono}\fako{anat.} La lasta e maxim longa parto di la intestino tenua\lingui{EFIS}
\artiklo{ilexo}%
\vorto{ilexo}\fako{bot.} Arboro di la familio « ilicinei », sempre verda, di qua la folii esas briloza e dornoza\latina{ilex aquifolium}
\artiklo{iliako}%
\vorto{iliako}\fako{anat.} La maxim granda ek la tri osti flankala\lingui{EFISL}
\artiklo{iliterata}%
\vorto{iliterata} Qua ne savas lektar
\artiklo{iluminar}%
\vorto{iluminar}\fako{trans.} Ornar (manuskripto, libro imprimita) ye miniaturi, graburi, e c.\lingui{DEFIRS}
\artiklo{iluzionar}%
\vorto{iluzionar}\fako{netrans.} Subisar ta eroro sensala o spiritala, qua igas aceptar lo semblanta kom lo reala\lingui{DEFIRS}
\artiklo{-im-}%
\vorto{-im-} Sufixo qua signifikas nombro fracionala
\artiklo{imaginar}%
\vorto{imaginar}\fako{trans.} Reprezentar ulo en sua spirito\lingui{EFIS}
\artiklo{imaginara}%
\vorto{imaginara}\fako{matem.} Simbolo algebrala, ek radikalo duesma-grada pri nombro negativa\lingui{DEFIRS}
\artiklo{imajo}%
\vorto{imajo}\senco{1}{Aparajo videbla ek objekto lumoza, konsistanta ye la radii lumala qui spricas de olca}\senco{2}{Aparajo videbla ek objekto, imitita per desegnuro, pikturo, skulturo}\senco{3}{Aparajo videbla di objekto, konceptita per la imagino-fakultato}\lingui{EFIS}
\artiklo{imanar}%
\vorto{imanar}\fako{netrans.} Efikar konstante, vice efikar provizore \textit{(aludante Deo)}\lingui{DEFIS}
\artiklo{imbastar}%
\vorto{imbastar}\fako{trans.} Sutar per longa steburi ed (ordinare) per filo blanka, por pose probar (la vesto, la robo, e c.)\lingui{EFI}
\artiklo{imbecila}%
\vorto{imbecila} Di qua la spirito esas febla\lingui{DEFIS}
\artiklo{imbibar}%
\vorto{imbibar}\fako{trans., ulo, ye} Humidigar, per igar absorbar liquido maxima-quante\lingui{DEFIS}
\artiklo{imbrikar}%
\vorto{imbrikar}\fako{trans.} Dispozar, aranjar quale la teguli di tekto\lingui{EFIS}
\artiklo{imensa}%
\vorto{imensa} Di qua la grandeso-grado impedas la mezuro\lingui{EF}
\artiklo{imersar}%
\vorto{imersar}\fako{trans., aden} Parte sinkar (korpo) aden liquido\lingui{DEFIS}
\artiklo{imitar}%
\vorto{imitar}\fako{trans.} Agar quale ulu; facar (ulo) segun, (kom modelo), la verko da ulu\lingui{DEFIRS}
\artiklo{imoblo}%
\vorto{imoblo}\fako{yuro-cienco} Havajo nemovebla (domo, agro, e c.)\lingui{DEFIS}
\artiklo{imolar}%
\vorto{imolar}\fako{trans.}\senco{1}{Ocidar sakrifike por satisfacar la deajo}\senco{2}{Perisigar, ruinar, ulu por satisfacar interesto, pasiono, e c.}\lingui{EFI}
\artiklo{imortelo}%
\vorto{imortelo}\fako{bot.} Planto ek la familio « kompozaji », di qua la floro ne velkas, ed uzata por ornar la tombi\latina{helichrysum}\lingui{DEFR}
\artiklo{impedar}%
\vorto{impedar}\fako{trans., pri, per} Entravar, embarasar (ulu) pri lua ago od agado; entravar la evento (di ulo)\lingui{EFIS}
\artiklo{imperar}%
\vorto{imperar}\fako{trans.} Dicar (ad ulu) ke lu agez (to o co), kun implicita nuanco di puniso kaze di neobedio. (Ne uzesas pri la armeo-autoritato. Komp. « komandar »)\lingui{EFIS}
\artiklo{imperativo}%
\vorto{imperativo}\senco{1}{\textit{(gram.)} Modo qua expresas la impero}\senco{2}{\textit{(filoz.)} Ago-regulo absoluta e ne-kondicionoza}\lingui{DEFIS}
\artiklo{imperatoro}%
\vorto{imperatoro}\fako{historio} Komandanto di la armeo Romana antiqua\lingui{DEFIRS}
\artiklo{imperfekto}%
\vorto{imperfekto}\fako{gram.} En ula lingui (kom ex, en la linguo Franca), tempo di la modo « indikativo » e di la modo « subjuntivo », qua indikas ago pasinta, egardata kom eventinta sama-instante kam altra ago anke pasinta; o la kustumo pri ago, la ofta itero di olca\lingui{DEFIS}
\artiklo{imperio}%
\vorto{imperio} Naciono, o socio politikala, qua uzas sua autoritato a populi konquestinta da lu\lingui{DEFIRS}
\artiklo{impertinenta}%
\vorto{impertinenta} Qua manifestas desrespekto nekonvenanta\lingui{DEFIS}
\artiklo{impetuar}%
\vorto{impetuar}\fako{netrans.}\senco{1}{\textit{(pri kavalo)} Eskapar de manuo-teneso, per movi violentoza e rapida}\senco{2}{\textit{(pri persono)} Iracar subite}\lingui{EFIS}
\artiklo{implemento}%
\vorto{implemento} Utensilo manuala (por suto, skribo, e c.)\lingui{E}
\artiklo{implicita}%
\vorto{implicita} Kontenata en la senco, sen esar formale expresata\lingui{DEFIS}
\artiklo{implikar}%
\vorto{implikar}\fako{trans.) (logiko} Kontenar\lingui{EFIRS}
\artiklo{implorar}%
\vorto{implorar}\fako{trans.) (ulu, pri ulo}\senco{1}{Suplikar emocigante}\senco{2}{Demandar, suplikante ye maniero emociganta}\lingui{EFIS}
\artiklo{importacar}%
\vorto{importacar}\fako{trans.) (komerco} Igar (la produkturi da lando altra) enirar la lando\lingui{DEFIRS}
\artiklo{importar}%
\vorto{importar}\fako{netrans., pri} Implikar konsequaji grava pri ulu, pri ulo\lingui{EFIRS}
\artiklo{impostar}%
\vorto{impostar}\fako{trans., ye) (financo} Determinar la parto di la spensi nacionala publika, quan la stato impozas a singla ek la civitani\lingui{DEFI}
\artiklo{impotenta}%
\vorto{impotenta}\fako{fiziol.} Quan naturo impedas uzar un ek sua membri\lingui{DEFIRS}
\artiklo{impozar}%
\vorto{impozar}\fako{trans., ulo, ad ulu} Igar (ulu) subisar ne-voloze\lingui{DEFIS}
\artiklo{impregnar}%
\vorto{impregnar}\fako{trans., ye ulo} Igar penetrar en omna parti (di korpo)\lingui{DEFIS}
\artiklo{imprekar}%
\vorto{imprekar}\fako{netrans.} Dezirar ke ulu subisez desfortuno, mala chanco\lingui{EFIS}
\artiklo{impresar}%
\vorto{impresar}\fako{trans.} Influar spiritale, mentale. \textit{(anke metafore)}\lingui{DEFIS}
\artiklo{impresario}%
\vorto{impresario} Chefo di entraprezo teatrala\lingui{DEFIRS}
\artiklo{impresiono}%
\vorto{impresiono}\fako{pikt., e literat.} Verko artala qua tendencas expresar o reprezentar la impresi recevita da la artisto, de la naturo o de la temo\lingui{DEFIS}
\artiklo{imprimar}%
\vorto{imprimar}\fako{trans., per}\senco{1}{Reproduktar per graburo-planki, per tipi inkizita}\senco{2}{\textit{(fotogr.)} Obtenar imajo pozitiva de vitro-plako o filmo negativa}\lingui{EFIS}
\artiklo{improvizar}%
\vorto{improvizar}\fako{trans.} Produktar quik e sen preparo (diskurso, versi, e c.)\lingui{DEFIRS}
\artiklo{impulsar}%
\vorto{impulsar}\fako{trans., ad) (anke metaf.} Komencar pulsar (korpo); pulsar (ulu) ad agar ulo\lingui{DEFIRS}
\artiklo{imputar}%
\vorto{imputar}\fako{trans., ulo, ad ulu} Atribuar (la responso di ulo, ad ulu)\lingui{DEFIS}
\artiklo{imuna}%
\vorto{imuna}\senco{1}{Liberigita de devo, de obligeso, pro yuro (naturala od aquirita)}\senco{2}{Qua profitas grantajo (naturala od aquirita) pro qua lua korpo shirmesas de ula morbi}\lingui{DEFIRS}
\artiklo{-in-}%
\vorto{-in-} Sufixo qua indikas « ento femina », kande onu volas emfazar la sexuo
\artiklo{inaniciono}%
\vorto{inaniciono}\fako{fiziol.} Exhausteso pro ne-nutreso, o pro hiponutreso\lingui{DEFIS}
\artiklo{inata}%
\vorto{inata}\senco{1}{Qua esis en ni lor nia nasko}\senco{2}{\textit{(filoz.)} Qua inheras spirito, vice aquiresir per experienco}\lingui{EFIS}
\artiklo{inaugurar}%
\vorto{inaugurar}\fako{trans.}\senco{1}{Konsakrar per ceremonio oficala}\senco{2}{Markizar (la komenco di estado nova)}\lingui{DEFIRS}
\artiklo{incendiar}%
\vorto{incendiar}\fako{trans.} Facar granda fairo qua su extensas e ruinas\lingui{F}
\artiklo{incenso}%
\vorto{incenso} Rezino odoroza, di qua la odoro exhalesas precipue lor kombusteso, e quan onu extraktas de multa arbori ek la familio « bursaracei »\lingui{EFIS}
\artiklo{incestar}%
\vorto{incestar}\fako{netrans., kun} Koitar kun sua patro, filio, parento, o kun alianco-familiano, ye grado interdiktita da la lego\lingui{DEFIRS}
\artiklo{inch}%
\vorto{« inch »} Mezur-unajo Britaniana ed Usana : 2,53995 centimetri
\artiklo{incidar}%
\vorto{incidar}\fako{netrans., kontre) (fiziko} Arivar e butar kontre \textit{(ulo)}\lingui{EFIS}
\artiklo{incidento}%
\vorto{incidento} Mikra eventajo qua intervenas acesore\lingui{DEFIS}
\artiklo{incitar}%
\vorto{incitar}\fako{trans., ulu ad ulo} Konsilar, exhortar vivace ad agor ulo\lingui{DEFIS}
\artiklo{incizar}%
\vorto{incizar}\fako{trans.} Fendar butono-truatre, per instrumento sekiva\lingui{EFIS}
\artiklo{-ind-}%
\vorto{-ind-} Sufixo qua signifikas « qua meritas esar… -ata »
\artiklo{indemno}%
\vorto{indemno}\fako{yuro-cienco} To quon onu donas ad ulu por kompensar lua perdi, lua spensi\lingui{DEFIS}
\artiklo{indentar}%
\vorto{indentar}\fako{trans.) (imprim-arto} Retropulsar (lineo tipografiala) aden la longeso-kadro, pozante avan olu « lakuno- »- tipi
\artiklo{indexo}%
\vorto{indexo}\senco{1}{Tabelo alfabetala, fine di libro, en qua esas mencionita la vorti precipua e la propra nomi, uzita da la autoro, kun indiko pri la pagino ube esas singla de oli}\senco{2}{Mikra vergo ek metalo qua, per cifro-plako, indikas la kloki, la kompreseso-grado, e c.}\lingui{DEFIS}
\artiklo{indiciono}%
\vorto{indiciono}\fako{kronologio} Determino di la yaro di la cikli di la periodo Julianala, por datizar la akti\lingui{DEFIS}
\artiklo{indico}%
\vorto{indico}\senco{1}{Signo qua trovigas ulo}\senco{2}{\textit{(algebro)} Signo distingiva quan onu donas a litero qua reprezentas, en kalkulo, plura *grandori analoga}\lingui{DEFIS}\subvorto{fiziol.}{refrakto-indico}{To quo indikas la relato da la sinuso di la incido-angulo, a la sinuso di la angulo di refrakto}{3}{}
\artiklo{indieno}%
\vorto{indieno} Kotono-stofo kun imaji imprimita, qua origine fabrikesis en India\lingui{DFS}
\artiklo{indiferenta}%
\vorto{indiferenta}\senco{1}{Qua ne tendencas ad ulo prefere kam ad altro}\senco{2}{Qua ne interesesas da persono o da kozo, prefere kam da altra persono od altra kozo}\senco{3}{Qua ne esas plu interesiva kam altra kozo}\lingui{DEFIRS}
\artiklo{indignar}%
\vorto{indignar}\fako{netrans., pri, pro, kontre} Etiko-revoltar pro la spekto di konduto tre nekonvenanta\lingui{DEFIRS}
\artiklo{indigo}%
\vorto{indigo} Farbo blua quan onu extraktas de la folii e de la stipi di la L. \textit{indigofera}\latina{indigofera}\lingui{DEFIRS}
\artiklo{indijar}%
\vorto{indijar}\fako{trans.} Ne-havar ulo quon onu bezonas\lingui{EFIS}
\artiklo{indijena}%
\vorto{indijena} Qua naskis en la lando quan lu habitas\lingui{EFIS}
\artiklo{indikar}%
\vorto{indikar}\fako{trans., ad} Montrar ube esas (ulu, ulo)\lingui{EFIS}
\artiklo{indikativo}%
\vorto{indikativo}\fako{gram.} Modo di la verbo, qua expresas la ago ye maniero absoluta\lingui{DEFIS}
\artiklo{indikatoro}%
\vorto{indikatoro}\fako{tekn.} Aparato qua indikas (la rapideso-grado, la direcioni, la situeso di la bifurk-aguli fervoyala, la aquo-quanto di kaldiero, e c.)\lingui{DEFI}
\artiklo{individuo}%
\vorto{individuo} Ento qua esas unajo distingebla ek speco, genero, irga membro ek societo o socio de homi\lingui{DEFIRS}
\artiklo{indolenta}%
\vorto{indolenta} Qua evitas penar\lingui{DEFIS}
\artiklo{indosar}%
\vorto{indosar}\fako{trans.) (komerco} Skribar e signatar la transferoformulo reverse di mandato, di kambio-letro\lingui{DEFIS}
\artiklo{induktar}%
\vorto{induktar}\fako{trans.}\senco{1}{\textit{(logiko)} Rezonar, retroirante de la fakti observita, a la lego (*leyo) quan oli konjektigas, de la efekti a la kauzo}\senco{2}{\textit{(fiziko)} Produktar ta elektro-fluo qua rezultas de la efiko da magneti o da konduktili trairata da elektro-fluo (*korento)}\lingui{DEFIRS}
\artiklo{induktoro}%
\vorto{induktoro}\fako{elektro) (en mashino elektrifala} La aparato qua produktas la feldo induktera\lingui{DEFIRS}
\artiklo{indulgar}%
\vorto{indulgar}\fako{trans., ulu, pri ulo} Pardonar facile.sen postular rigoroze exkuzi\lingui{EFIS}
\artiklo{induljenco}%
\vorto{induljenco}\fako{teol. katol.} Remiso, tota o parta, di la punisi pro pekir, quan la eklezio katolika grantas ula-kondicione (fasto, prego, almono, e c.)\lingui{DEFIS}
\artiklo{indulto}%
\vorto{indulto}\fako{yuro ekleziala} Povo, autoritato, grantita da la papo, pri nominar ye ula benefici, o ye retenar ta benefici, kontre la yuro komuna\lingui{DEFIS}
\artiklo{industrio}%
\vorto{industrio} Ensemblo ek la arti, mestieri, qui transformas la materii prima\lingui{DEFIRS}
\artiklo{indutar}%
\vorto{indutar}\fako{trans.) (ye ulo} Kovrar ye farbo-strato\lingui{F}
\artiklo{inerta}%
\vorto{inerta}\senco{1}{Qua ne esas propre agiva}\senco{2}{\textit{(metaf.)} Sen-energia}\lingui{EFIRS}
\artiklo{infalibla}%
\vorto{infalibla} Qua ne povas ne-eventor\lingui{DEFIS}
\artiklo{infama}%
\vorto{infama} Stigmatizita da la opiniono sociala, da la lego\lingui{DEFIS}
\artiklo{infante}%
\vorto{« infante »} En Hispania ed en Portugal, nomo di ta ek la princi qua naskis nemediate pos la senioro
\artiklo{infanto}%
\vorto{infanto} Homo qua evas inter 1 e 7 yari\lingui{EFIS}
\artiklo{infantrio}%
\vorto{infantrio}\fako{milit-arto} Ensemblo ek armeani qui marchas, e kombatas pedo-stacante\lingui{DEFIRS}
\artiklo{infarkto}%
\vorto{infarkto}\fako{patol.} Infiltreso da tisuo ye ekfluo da sango\lingui{DEFIS}
\artiklo{infektar}%
\vorto{infektar}\fako{trans.} Impregnar ye jermi nociva\lingui{DEFIRS}
\artiklo{inferar}%
\vorto{inferar}\fako{trans., de, ek} Ektirar (konsequajo)\lingui{EFIS}
\artiklo{inferiora}%
\vorto{inferiora}\fako{en hierarkio} Qua okupas rango suba, infra\lingui{DEFIS}
\artiklo{inferno}%
\vorto{inferno}\senco{1}{\textit{(mitol.)} Loko subsula, habitata da la « ombri » di la mortinti, e di qua un regiono (Tartaro) esas rezervita por la maligni, ed altra (Elizeo) esas la lojeyo di la yusti}\senco{2}{\textit{(kristanismo)} Loko ube supliciesas la damniti (opoze a purgatorio ed a paradizo)}\lingui{EFIS}
\artiklo{infiltrar}%
\vorto{infiltrar}\fako{netrans., aden}\senco{1}{Insinuar su pokope aden la pori, la interstici, la fenduri, di korpo solida}\senco{2}{\textit{(metaf.)} Penetrar pokope aden la spirito, aden la mento}\lingui{DEFIS}
\artiklo{infinita}%
\vorto{infinita}\fako{matem.} To quon spirito ne konceptas kom limitoza; to quo esas plu granda kam irga quanto evaluebla. – \textit{(anke metaf.)}\lingui{DEFIS}
\artiklo{infinitezima}%
\vorto{infinitezima}\fako{matem.} Qua relatas la quanti infinite mikra\lingui{EFIS}
\artiklo{infinitivo}%
\vorto{infinitivo}\fako{gram.} Verbo-modo qua expresas la ago verbala en maniero nepreciza\lingui{DEFIS}
\artiklo{infirma}%
\vorto{infirma} Qua ne povas uzar (libere, normale) irga membro od organo\lingui{EFIS}
\artiklo{infixo}%
\vorto{infixo}\fako{gram.} Elemento qua insertesas interne di la foni ye qui konsistas radiko, por modifikar la senco di olca\lingui{DEFIS}
\artiklo{inflacar}%
\vorto{inflacar}\fako{netrans.) (financo} Emisar exajeroze moneto-papero\lingui{DEFIS}
\artiklo{inflamar}%
\vorto{inflamar}\fako{trans.} Produktar ta estado di maladeso, quan karakterizas varmeso, doloro ed (ulakaze) la tumefakteso di la parto malada\lingui{DEFIS}
\artiklo{inflar}%
\vorto{inflar}\fako{trans., per ulo, ye ulo} Distensar omna-sinse (korpo elastika) per preso interna. – \textit{(anke metaf.)}\lingui{EFIS}
\artiklo{inflexar}%
\vorto{inflexar}\fako{trans.}\senco{1}{\textit{(gram.)} Modifikar la pronunco-valoro di vokalo, influe da altra vokalo qua esas en la silabo qua sequas nemediate}\lingui{EFIS}\subvorto{geom.}{arko inflexita}{Qua konsistas ye du taloni, tangenta per olia somiti (quale la embracilo tipografiala)}{2}{}
\artiklo{infloresenco}%
\vorto{infloresenco}\fako{bot.} Dispozeso di la floro sur la stipito\lingui{EFIS}
\artiklo{influar}%
\vorto{influar}\fako{trans., per, pro, pri) (aludante kauzo nemateriala} Efikar ad ulu, ad ulo, tale ke ica modifikesas\lingui{DEFIS}
\artiklo{influenzo}%
\vorto{influenzo}\fako{patol.} Gripo tre grava\lingui{DEFIRS}
\artiklo{influxo}%
\vorto{influxo}\fako{biol.} Forco, di qua la naturo esas nekonocata, qua cirkulas alonge la nervi\lingui{DEFIS}
\artiklo{informar}%
\vorto{informar}\fako{trans., ulu, pri ulo} Igar (ulu) ne plus nesavar (fakti)\lingui{DEFIRS}
\artiklo{infre}%
\vorto{infre}\fako{pri objekto vertikala} Maxim apud-sule. – \textit{(pri objekto horizontala)}. Normale, maxim apud la regardanto\lingui{L}
\artiklo{infuzar}%
\vorto{infuzar}\fako{trans.} Varsar sur (kafeo, teo, e c.) aquo, od altra liquido, bolianta por ke olca prenez la substanci esenca quin olta kontenas\lingui{DEFIS}
\artiklo{infuzorio}%
\vorto{infuzorio}\fako{zool.} Animaleto qua kreskas en la infuzuri vejetala ed animalala, en la aqui koruptita\lingui{DEFIRS}
\artiklo{ingestar}%
\vorto{ingestar}\fako{trans.) (fiziol.} Enduktar aden sua stomako\lingui{DEFL}
\artiklo{inglobar}%
\vorto{inglobar}\fako{trans.} Ne simple « unionar », ma : mixar en un substanco; inkluzar substanco (farmaciala) en altra (de altra naturo)\lingui{FS}
\artiklo{ingranar}%
\vorto{ingranar}\fako{trans., ye) (tekn.} Igar (la denti di roto) enirar la lakuni qui interseparas la denti di altra roto, por transmisar ad olca la movo di olta\lingui{FIS}
\artiklo{ingrediento}%
\vorto{ingrediento} Elemento qua uzesas por la kompozo di mixajo o mixuro, di preparajo o preparuro\lingui{DEFIS}
\artiklo{inguino}%
\vorto{inguino}\fako{anat.} Parto di la homo-korpo, inter la infra ventro e la kruro\lingui{EFIS}
\artiklo{inhalar}%
\vorto{inhalar}\fako{trans.} Aspirar por englutar (vapori, gasi)\lingui{DEFIS}
\artiklo{inherar}%
\vorto{inherar}\fako{trans.} Tenar profunde an la ento di persono, di kozo\lingui{DEFIS}
\artiklo{inhibar}%
\vorto{inhibar}\fako{trans.) (biol.} Diminutar od haltigar la funciono di organo, per la irito di fora punto di la organismo\lingui{DEFIS}
\artiklo{iniciar}%
\vorto{iniciar}\fako{trans.}\senco{1}{Admisar a la konoco, a la partopreno di la misterii religiala}\senco{2}{\textit{(metaf.)} Admisar a la konoco di kozi sekreta, arkana}\lingui{DEFIRS}
\artiklo{iniciativar}%
\vorto{iniciativar}\fako{trans.} Esar la dat-unesma qua spontane propozas, organizas, exekutas (ulo)\lingui{DEFIRS}
\artiklo{injektar}%
\vorto{injektar}\fako{trans., en, aden} Lansar (liquido) aden substanco\lingui{DEFIRS}
\artiklo{injeniar}%
\vorto{injeniar}\fako{netrans., por, pri} Tre laborigar sua cerebro por abutar adulo\lingui{EFIS}
\artiklo{injenioro}%
\vorto{injenioro} La persono qua konstruktas mashini, qua direktas la konstrukto di ponti, voyi, e c.\lingui{DEFIRS}
\artiklo{inkandecar}%
\vorto{inkandecar}\fako{netrans.} Qua esas tante varmigita ke lu aspektas blanka\lingui{EFIS}
\artiklo{inkardinar}%
\vorto{inkardinar}\fako{trans.} Destinar, perpetue ed absolute, ula kleriko a la servo di diocezo determinata
\artiklo{inkarnacar}%
\vorto{inkarnacar}\fako{trans.) (teol.} Glitar su aden korpo ek karno homatra\lingui{DEFIS}
\artiklo{inkasar}%
\vorto{inkasar}\fako{trans.} Recevar e pozar aden la kaso (di la afero) pekunio-quanto\lingui{DFI}
\artiklo{inkastrar}%
\vorto{inkastrar}\fako{trans.} Inserar (objekto) aden altra objekto qua entaliesis por recevar olu; fixigar (lapido) aden la kastono, per abatar la rebordo di olca adcirkum la lapido\lingui{FIS}
\artiklo{inklinar}%
\vorto{inklinar}\fako{trans., ad} Igar lua axo kelke deviacar de lua vertikalo\lingui{EFIS}
\artiklo{inkluzar}%
\vorto{inkluzar} (trans., (ad) en.) Enpozar, cirkumklozar\lingui{DEFIRS}
\artiklo{inko}%
\vorto{inko} Liquido quan onu uzas kande onu skribas per plumo\lingui{EFI}
\artiklo{inkoativa}%
\vorto{inkoativa}\fako{gram.} Qua expresas la komenco di ago\lingui{DEFIS}
\artiklo{inkognito}%
\vorto{inkognito} Sen esar konocata, remarkata\lingui{DEFIRS}
\artiklo{inkombrar}%
\vorto{inkombrar}\fako{trans., ye) (aludante amaso de kozi} Embarasar per esar obstaklo pri la cirkulo\lingui{EFIS}
\artiklo{inkrustar}%
\vorto{inkrustar}\fako{trans.} Kovrar (objekto) ye strato petra e krustatra\lingui{DEFIRS}
\artiklo{inkubacar}%
\vorto{inkubacar}\fako{trans.} Havar en su la jermo di morbo, qua kreskas til kande onu maladeskas\lingui{DEFIS}
\artiklo{inkubo}%
\vorto{inkubo}\fako{teol.} Demono qua kopulacis, lokizinte su sur la dormanto\lingui{DEFIS}
\artiklo{inkulkar}%
\vorto{inkulkar}\fako{trans.} Igar par-enirar en la spirito\lingui{EFIS}
\artiklo{inkunablo}%
\vorto{inkunablo} Libro qua evas de la epoko kande onu jus inventabis la imprim-arto (generale til la yaro 1500)\lingui{DEFIS}
\artiklo{inkurar}%
\vorto{inkurar}\fako{trans.} Igar su en la kazo subisor\lingui{F}
\artiklo{inkursar}%
\vorto{inkursar}\fako{netrans.}\senco{1}{Devast-irupto da militanti aden la lando enemika}\senco{2}{\textit{(metaf.)} Trakto acesora di temo qua ne apartenas a la domeno precipua}\lingui{EFIRS}
\artiklo{inocenta}%
\vorto{inocenta}\fako{pri} Qua ne agis ulo mala\lingui{EFIS}
\artiklo{inokular}%
\vorto{inokular}\fako{trans.) (medic.} Transdonar ad ulu (viruso, jermo di morbo kontagiala); o transdonar artificale la viruso (febligita) di morbo, kom prezervivo kontre la ipsa morbo\lingui{DEFIRS}
\artiklo{inquestar}%
\vorto{inquestar}\fako{netrans., pri, che} Serchadar por savar ulo, per questioni, per askolto di testi, e c.\lingui{EFIS}
\artiklo{inquiziciono}%
\vorto{inquiziciono}\fako{historio} Tribunalo di la eklezio katolika, instalita por serchar e persequar la personi di qui la sentimenti esis enemika a la kredo katolika\lingui{DEFIRS}
\artiklo{inquizitoro}%
\vorto{inquizitoro}\fako{historio} Judiciisto ek la inquiziciono\lingui{DEFIRS}
\artiklo{insektivora}%
\vorto{insektivora} Qua su nutras ye insekti\lingui{DEFIRS}
\artiklo{insekto}%
\vorto{insekto}\fako{zool.} Mikra animalo sen-vertebra, ek la klaso « artropodi »\lingui{DEFIRS}
\artiklo{insertar}%
\vorto{insertar}\fako{trans., ulo, aden} Enpozar (kozo aden altra kozo) tale ke oli formacas quaze un korpo; enskribar (artiklo, klauzo) aden jurnalo, akto, e c.\lingui{DEFIS}
\artiklo{insidiar}%
\vorto{insidiar}\fako{trans.} Ruzar kontre ulu, por venkar lu, per travestiar naturo, per simular\lingui{EFIS}
\artiklo{insigno}%
\vorto{insigno}\senco{1}{Tabelo quan komercisto muntas sur sua butiko (maxim-multa-kaze kun emblemo o moto) por igar olca *rekonocebla}\senco{2}{Marko distingiva pri la autoritato, pri la rango, di persono}\senco{3}{Mikra objekto butonatra, quan onu fixigas an la surfaldajo di paltoto, di vestio, por indikar ke onu esas membro od adepto di ulo}\lingui{DEFIS}
\artiklo{insinuar}%
\vorto{insinuar}\fako{trans., su, ulo, aden}\senco{1}{Enduktar (ulo) dolce, gradoze; enduktar (ulu) habile, ruzoze}\senco{2}{\textit{(metaf.)} Igar (ulo) penetrar habile aden la spirito, aden la anmo, la psiko}\lingui{EFIRS}
\artiklo{insipida}%
\vorto{insipida} Qua repugnas gustado, pro olua sensaporeso. \textit{(anke metaf.)}\lingui{DEFIS}
\artiklo{insistar}%
\vorto{insistar}\fako{netrans., pri, pro, por} Presar \textit{(metaf.)} vigoroze lor demando\lingui{EFIS}
\artiklo{insolenta}%
\vorto{insolenta} Qua manifestas nerespekto insultoza, ofensoza\lingui{DEFIRS}
\artiklo{inspektar}%
\vorto{inspektar}\fako{trans.} Explorar atencoze (to quon onu devas surveyar)\lingui{DEFIRS}
\artiklo{inspirar}%
\vorto{inspirar}\fako{trans., ad} Naskigar en la kordio, en la spirito, (penso, rezolvo)\lingui{DEFIS}
\artiklo{instalar}%
\vorto{instalar}\fako{trans. ulu} Establisar solene en lua klapsidilo (episkopo, paroko); establisar solene en lua ofico (judiciisto, prezidero); establisar (ulu) en la lojeyo, en la loko, en la plaso qua destinesis por lu\lingui{DEFIS}
\artiklo{instanco}%
\vorto{instanco}\fako{yuro-cienco} Ago di persequo judiciala inter pledanto e defensanto\lingui{DEFIRS}
\artiklo{instanto}%
\vorto{instanto} Mikra tempo-fragmento\lingui{DEFIRS}
\artiklo{instigar}%
\vorto{instigar}\fako{trans., ad} Pulsar ad ago\lingui{EFIS}
\artiklo{instilar}%
\vorto{instilar}\fako{trans. aden) (medic.} Varsar gutope\lingui{DEFI}
\artiklo{instinto}%
\vorto{instinto} Impulsivo naturala; impulsivo interna qua instigas ento vivanta (homo o bestio) ad agi nerezonita, por la konservo di la individuo, di la speco\lingui{DEFIRS}
\artiklo{institucar}%
\vorto{institucar}\fako{trans.} Establisar durive\lingui{DEFIRS}
\artiklo{instituto}%
\vorto{instituto} Korpo institucura ek literaturisti, ciencisti, artisti, e c.\lingui{DEFIRS}
\artiklo{instruciono}%
\vorto{instruciono} Expliko por la direkto di afero, por la uzo di ulo\lingui{DEFIRS}
\artiklo{instruktar}%
\vorto{instruktar}\fako{trans.} Formacar la spirito, la mento, la psiko di (ulu) per lecioni, precepti\lingui{DEFIRS}
\artiklo{instrumento}%
\vorto{instrumento} To quon onu uzas por produktar efekto materiala\lingui{DEFIRS}
\artiklo{insulo}%
\vorto{insulo} Extensajo ek tero qua, en maro, cirkondesas da aquo omna-latere, tota-cirkume\lingui{DEFIS}
\artiklo{insultar}%
\vorto{insultar}\fako{trans.} Ofensar extreme (generale per vorti)\lingui{DEFIS}
\artiklo{-int-}%
\vorto{-int-}\fako{gram.} Dezinenco di la participo pasintala aktiva
\artiklo{intalio}%
\vorto{intalio} Petro harda, gravita kave\lingui{DEFI}
\artiklo{integra}%
\vorto{integra} Di qua la toto ne subisis diminuto\lingui{EFIRS}
\artiklo{integralo}%
\vorto{integralo}\fako{matem.} La integralo di diferencialo : quanto di qua la diferencialo esas infinitezima\lingui{DEFIRS}
\artiklo{intelektar}%
\vorto{intelektar}\fako{trans.} Konoceskar per la koncepto-fakultato pura\lingui{DEFIRS}
\artiklo{*intelektuelo}%
\vorto{*intelektuelo} Persono qua prizas prepondere lo pensala, lo spiritala\lingui{DEFIS}
\artiklo{inteligenta}%
\vorto{inteligenta} Qua komprenas facile\lingui{DEFIRS}
\artiklo{intencar}%
\vorto{intencar}\fako{trans.} Volar ad skopo\lingui{FIS}
\artiklo{intendanto}%
\vorto{intendanto}\fako{armeo} Ensemblo ek la oficiri qui su okupas pri la servicii administrala di la armearo\lingui{DFIRS}
\artiklo{intendanto}%
\vorto{intendanto}\senco{1}{La persono di qua la tasko konsistas ye jerar la aferi, administrar la menajo di richo, di sinioro}\senco{2}{Stat-oficiero, chefo di administrantaro guvernantarala}\lingui{DEFIRS}
\artiklo{intensa}%
\vorto{intensa} Qua efikas energioze\lingui{DEFIRS}
\artiklo{inter-}%
\vorto{inter-}\senco{1}{En la tempo o spaco qua separas du o multa personi, kozi, e c.}\senco{2}{\textit{(metaf.)} Signifikas partigo, kambio, reciprokeso}\lingui{EFRS}
\artiklo{interceptar}%
\vorto{interceptar}\fako{trans.} Sizar, kande olu pasas, kozo qua esas destinita ad ulu altra\lingui{EFIS}
\artiklo{intercesar}%
\vorto{intercesar}\fako{netrans., por, favore di} Intervenar por ke ulu obtenez lua pardoneso, lua indulgeso\lingui{DEFIRS}
\artiklo{interdiktar}%
\vorto{interdiktar}\fako{trans., ulo, ad ulu} Imperar ke ulu ne (agez ulo)\lingui{DEFIRS}
\artiklo{interesar}%
\vorto{interesar}\fako{trans., pri, ad} Ecitar en ulu lua kuriozeso, lua emoci\lingui{DEFIRS}
\artiklo{interesto}%
\vorto{interesto}\senco{1}{Ensemblo de la kozi qui esas avantajoza ye ulu}\senco{2}{\textit{(financo)} Profito-quanto de la pekunio, quan onu prestis, kalkulata (maxim-multa-kaze) procente}\lingui{DEFIRS}
\artiklo{interferar}%
\vorto{interferar}\fako{netrans.} \textit{(fiziko)} \textit{(aludante lumo-radii)} Renkontrar altra lumo-radii, de quo konsequas deskresko plu o min granda da lia lumo. – Fenomeno analoga pri la fonoondi\lingui{DEFIRS}
\artiklo{interimo}%
\vorto{interimo} Intervalo dum qua ofico o plaso (en hierarkio) vakas\lingui{DEFIRS}
\artiklo{interjeciono}%
\vorto{interjeciono}\fako{gram.} Klameto qua expresas la emoco subita di la psiko\lingui{DEFIS}
\artiklo{intermezo}%
\vorto{intermezo} Distraktivo (baleto, danso, kanto) quan onu lokizas inter la akti di teatrajo\lingui{DEFIRS}
\artiklo{intermitar}%
\vorto{intermitar}\fako{netrans.} Cesar e rieventar pos intervali ne-samadura\lingui{DEFIS}
\artiklo{interna}%
\vorto{interna} Qua esas en, o qua apartenas a, la spaco qua existas inter la limiti di korpo\lingui{DEFIRS}
\artiklo{internacionalo}%
\vorto{internacionalo}\senco{1}{Asociitaro generala ek salariati di landi diversa, por revendikar ula yuri qui esas komuna ye li}\senco{2}{La ralio-himno di la salariati revolucionista di omna landi, unionita por luktar kontre kapitalismo (poemo da Eugène Pottier 1871; muzik-ario da Adolphe Degeyter)}\lingui{DEFIRS}
\artiklo{interogativo}%
\vorto{interogativo}\fako{gram.} Formo qua expresas la ideo di questiono\lingui{DEFIS}
\artiklo{interpelar}%
\vorto{interpelar}\fako{trans., pri, pro}\senco{1}{Paroleskar (vers ulu, ad ulu) por questionar lu pri ulo, o demandar ulo de lu}\senco{2}{\textit{(yuro-cienco)} Sumnar (ulu) justifikar su pri fakto}\senco{3}{\textit{(politiko)} \textit{(aludante parlamentano)} Demandar a (ministro) ke ica justifikez, koram la parlamento, sua agi o ne-agi}\lingui{DEFIRS}
\artiklo{interpolar}%
\vorto{interpolar}\fako{trans.}\senco{1}{Insertar aden texto (frazo qua ne esis parto di olu)}\senco{2}{\textit{(fiziko)} Insertar aden observorezultaji terminin determinita da kalkulo}\senco{3}{\textit{(algebro)} Trovar formulo qua, pro satisfacar ula quanto de kazi observita, povas remplasar provizore la *leyo (lego) di la fenomeno}\lingui{DEFIRS}
\artiklo{interpretar}%
\vorto{interpretar}\fako{trans.}\senco{1}{Explikar to quo esas obskura e ambigua en texto}\senco{2}{Dicar to quon signifikas (oraklo, sonjajo)}\senco{3}{\textit{(teatro)} Reprezentar segun la intenci da autoro}\senco{4}{Rezumar, en la linguo quan komprenas la askoltanti, to quon dicas persono en linguo stranjera quan la askoltanti ne komprenas}\lingui{DEFIS}
\artiklo{interstico}%
\vorto{interstico} Tre mikra spaco vakua inter la parti di korpo\lingui{EFIS}
\artiklo{intervalo}%
\vorto{intervalo}\senco{1}{Disto-quanto inter du loki}\senco{2}{Disto-quanto inter du instanti, du epoki}\senco{3}{\textit{(muziko)} Disto-quanto qua separas son-uno de la sono qua sequas lu nemediate, acensante de grava ad akuta, o decensante de akuta a grava}\lingui{DEFIRS}
\artiklo{interviuvar}%
\vorto{interviuvar}\fako{trans.) (aludante la jurnalisti} Vizitar personego famoza (aktuala), por questionar lu pri lua vivo-maniero, lua agi, lua idei, lua intenci pri lua agi balda, e c.\lingui{DEF}
\artiklo{intestat}%
\vorto{« intestat »} Qua ne facis sua testamento
\artiklo{intestino}%
\vorto{intestino}\fako{anat.} Vicero en la kavajo abdomina, tubo digestala ek mukozi e muskuli, qua iras de la stomako til la anuso\lingui{EFIRS}
\artiklo{intima}%
\vorto{intima}\senco{1}{Qua esas tote interna}\senco{2}{Qua atingas la esenco di la kozi}\senco{3}{\textit{(aludante personi)} Qua unionesas (ad altra) strete}\lingui{DEFIRS}
\artiklo{intonar}%
\vorto{intonar}\fako{trans.} Komencar kantar (ario) por donar la tono a la personi cetera\lingui{DEFIS}
\artiklo{intoxikar}%
\vorto{intoxikar}\fako{trans.} Igar (ulu) absorbar substanco venenoza\lingui{DEFIRS}
\artiklo{intracelula}%
\vorto{intracelula}\fako{biol.} Qua esas interne di celulo (opoze ad « intercelula » : qua esas inter la celuli)
\artiklo{intradoso}%
\vorto{intradoso}\fako{arkitekt.} Kavajo interna, parto konkava di vulto\lingui{DEFS}
\artiklo{intramolekula}%
\vorto{intramolekula}\fako{biol.} Qua esas interne di molekulo
\artiklo{*intrenar}%
\vorto{*intrenar}\fako{trans.} La « shutro » centrala di la triptiko « komencar », « intrenar », « finar », uzata kande « -as », « -is », « -os » ne esas sat preciza
\artiklo{intrigar}%
\vorto{intrigar}\fako{netrans.} Mashinacar por igar sucesar\lingui{DEFIRS}
\artiklo{intrikar}%
\vorto{intrikar}\fako{trans.} Intermixar kozi tale ke onu povas nur tre desfacile liberigar o riordinar oli\lingui{EI}
\artiklo{intrinseka}%
\vorto{intrinseka} Qua ne dependas de konvenciono. \textit{(antonimo : « extrinseka »)}\lingui{EFIS}
\artiklo{introduktar}%
\vorto{introduktar}\fako{trans.} Igar ulu darfar eniror e frequentor societo, klubo, mondumo, per prizentar lu kom aceptinda\lingui{DEFIRS}
\artiklo{introitus}%
\vorto{« introitus »}\fako{liturgio katolika} Prego da la sacerdoto, lor celebro di meso, pos acensir la altaro, e kande la koristaro kantas, komence di meso granda
\artiklo{intruzar}%
\vorto{intruzar}\fako{netrans.} Enirar nedarfante\lingui{EFIS}
\artiklo{intuicar}%
\vorto{intuicar}\fako{trans.) (filoz.} Konocar nemediate, direte.- \textit{(metaf.)} Komprenar la kozi rapide\lingui{DEFIRS}
\artiklo{intumecar}%
\vorto{intumecar}\fako{netrans.) (patol.} Infleskar\lingui{DEFIS}
\artiklo{intusucepciono}%
\vorto{intusucepciono}\fako{biol.} Endukto, aden korpo organizita, di nutrivi quin olu absorbas ed asimilas\lingui{DEFIS}
\artiklo{inulino}%
\vorto{inulino}\fako{kemio} Sorto di amilo quan onu extraktas ek la radiko di inulo
\artiklo{inulo}%
\vorto{inulo}\fako{bot.} Planto perena, ek la familio « kompozaji », di qua la radiko uzesas medicine\latina{inula}\lingui{FIL}
\artiklo{inundar}%
\vorto{inundar}\fako{trans., per, ye} Kovrar per aquo. – \textit{(metaf.)} Kovrar per homi multega\lingui{EFIS}
\artiklo{invadar}%
\vorto{invadar}\fako{trans.}\senco{1}{Okupar bruske, per violento, (teritorio)}\senco{2}{Devastar pos iruptir}\lingui{DEFIS}
\artiklo{invaginar}%
\vorto{invaginar}\fako{trans.} Retrofaldar lamo ek tisuo adinterne di altra tisui\lingui{DEFIS}
\artiklo{invalida}%
\vorto{invalida}\senco{1}{Qua ne esas tote vigoroza}\senco{2}{\textit{(aludante ex-militisto)} Impedata da lua granda evo, o da vunduri, durar sua armeaneso}\lingui{DEFIRS}
\artiklo{invektivar}%
\vorto{invektivar}\fako{netrans., kontre} Diskursar violente, furioze, quankam mem decoze, kontre ulu od ulo\lingui{DEFIS}
\artiklo{inventar}%
\vorto{inventar}\fako{trans.}\senco{1}{Krear ulo nova}\senco{2}{\textit{(retor.)} Serchar e selektar la argumento-punti quin onu devas uzar, la idei quin la temo furnisas}\lingui{EFIRS}
\artiklo{inventariar}%
\vorto{inventariar}\fako{trans.} Kontar ed enrejistrar specope la havaji transportebla, vari, komerco-trati, mandati, acioni, obligacioni di firmo, di persono, e c.\lingui{DEFIRS}
\artiklo{inversa}%
\vorto{inversa} Qua esas, o qua venas, opozo-sinse\lingui{EFIS}
\artiklo{invertazo}%
\vorto{invertazo}\fako{kemio} Fermento solvebla, qua transformas sakaroso aden glikoso, e qua existas en vejetanti diversa
\artiklo{investar}%
\vorto{investar}\fako{trans.}\senco{1}{Solene igar posedeskar feudo, granda ofico, per ceremonio alegoriala}\senco{2}{Igar posedeskar povo, autoritato}\lingui{DEFIS}
\artiklo{invetera}%
\vorto{invetera} Qua lasis kreskar en su – til divenir neextirpebla pro lua ancieneso – kustumo mala. (ex. : lu esas fumero, hiper-drinkero, invetera)\lingui{EF}
\artiklo{invitar}%
\vorto{invitar}\fako{trans., ulu, ad ulo, pri} Pregar (ulu) ke lu irez a loko determinita\lingui{EFIS}
\artiklo{invokar}%
\vorto{invokar}\fako{trans.} Advokar per prego; advokar (ulu) por ke lu atestez favore di la advokanto\lingui{EF}
\artiklo{involuciono}%
\vorto{involuciono}\senco{1}{\textit{(geom.)} Dui de punti AA'+BB'+BB'+CC' – selektita sur rekto, esas involuciono kande li dividas harmonioze un segmento MN di la rekto}\senco{2}{\textit{(bot.)} Eso-maniero di parto volvita adinterne di su}\lingui{DEFIRS}
\artiklo{involukro}%
\vorto{involukro}\fako{bot.} Asemblajo ek braktei qui formacas, cirkum floro o kapitulo, quaze kalico\lingui{DEFIS}
\artiklo{inyamo}%
\vorto{inyamo}\fako{bot.} Planto tropikala di qua la stipito esas klimera, ek la familio « amilacei », e di qua la radikaro esas nutriva\latina{discorea}\lingui{DEFIRS}
\artiklo{iodo}%
\vorto{iodo}\fako{kemio} Korpo simpla, ek flitri briloza, bluatre-griza, metalea, saporanta feble quale kloro\simbolo{I}\lingui{DEFIRS}
\artiklo{iodoformo}%
\vorto{iodoformo}\fako{kemio} Kompozajo (CHI₃) deskovrita en 1824 da Serullas
\artiklo{ionika}%
\vorto{ionika}\fako{arkitekt.} \textit{(aludante la triesma ek la kin klasi arkitekturala antiqua)} Karakterizata da la voluti di la kapiteli\lingui{DEFIS}
\artiklo{iono}%
\vorto{iono}\fako{fiziko} Partikulo qua portas elektro pozitiva (iono pozitiva) o negativa (iono negativa)\lingui{DEF}
\artiklo{iota}%
\vorto{« iota »} La litero nonesma, e la maxim mikra, ek la alfabeto Greka, qua korespondas ad « y »
\artiklo{ipsa}%
\vorto{ipsa} Qua esas ya, rigoroze, to (o : ta) quon (o : quan) onu jus nomis, mencionis\lingui{L}
\artiklo{-ir}%
\vorto{-ir}\fako{gram.} Dezinenco di la infinitivo pasintala
\artiklo{iracar}%
\vorto{iracar}\fako{netrans.} Subisar violentoza ecito, psikala, desagreabla, efike da qua onu quaze eruptas (kontre ulu)\lingui{EFIS}
\artiklo{irade}%
\vorto{« irade »} Reskripto, da sultano di Konstantinopolo, pri aferi min importoza kam olti qui bezonigas « hatt-i-hou-maloun » o « hatti-i-cheriff »
\artiklo{iradiar}%
\vorto{iradiar}\fako{netrans.) (fiziko} Difuzar su radie\lingui{DEFIS}
\artiklo{irar}%
\vorto{irar}\fako{netrans.}\senco{1}{\textit{(aludante vivanto)} Movar su ad ula loko (opoze a « venar »- « irar » signifikas for-esko di la spaco-punto ube esas la persono qua parolas, ad qua onu parolas, pri qua onu parolas.)}\senco{2}{\textit{(aludante nevivanto)} Movar su (o movesar) ad ula skopo}\lingui{FIS}
\artiklo{irga}%
\vorto{irga} Signifikas la nedetermineso maxima
\artiklo{iridio}%
\vorto{iridio}\fako{kemio} Korpo simpla, metalo ruptebla, arjento-blanka  quan onu renkontras en ula platin-erci, e di qua la dissolvuri esas irisea\simbolo{Ir}
\artiklo{irido}%
\vorto{irido}\fako{bot.} Genero de planti monokotiledona, qua konsistas ye quanto grandega de speci\latina{iris}
\artiklo{irigacar}%
\vorto{irigacar}\fako{trans.} Arozar granda-quante, per kanaleti, tubi, e c.\lingui{EFI}
\artiklo{irisito}%
\vorto{irisito}\fako{patol.} Inflameso di la iriso
\artiklo{iriso}%
\vorto{iriso}\senco{1}{\textit{(anat.)} Parto koloroza di la okulo, membrano cirkla, avan la kristalino}\senco{2}{\textit{(optiko)} Arko luma, ye la kolori ek la prismo, e produktita da la refrakteso e la reflekteso di la lumo-radii da la pluvo-guti}\lingui{DEFIRS}
\artiklo{iritar}%
\vorto{iritar}\fako{trans.} Eventigar mikra inflameso\lingui{DEFIS}
\artiklo{ironio}%
\vorto{ironio}\senco{1}{Questiono quan Sokrates agis a la sofisti, por incitar li a su-kontredici, per qui lu demonstris a li lia eroro}\senco{2}{Moko-formo qua konsistas ye dicar lo inversa di to quon onu volas komprenigar}\lingui{DEFIRS}
\artiklo{iruptar}%
\vorto{iruptar}\fako{netrans.} Enirar subite e violentoze\lingui{EFIS}
\artiklo{-is}%
\vorto{-is}\fako{gram.} Dezinenco di la indikativo pasintala
\artiklo{iskiono}%
\vorto{iskiono}\fako{anat.} Parto infra di la osti hanchala, ube inkastresas la osto di la kruro
\artiklo{islamo}%
\vorto{islamo}\senco{1}{Religio e civilizeso-formo di la Mohamedisti}\senco{2}{Ensemblo ek la landi quin karakterizas ta religio e ta civilizeso-formo}
\artiklo{-ismo}%
\vorto{-ismo} Sufixo qua signifikas « doktrino », « partiso », « sistemo »
\artiklo{istmo}%
\vorto{istmo}\fako{geogr.} Streta lango de tero, inter du mari, e qua unionas peninsulo a la kontinento cetera\lingui{DEFIS}
\artiklo{-isto}%
\vorto{-isto}\senco{1}{Sufixo qua signifikas : I. « la persono qua profesione su okupas pri… »}\senco{2}{Adepto, adherinto di partiso, di skolo doktrinala}
\artiklo{-it}%
\vorto{-it} Dezinenco di la participo pasintala pasiva
\artiklo{ita}%
\vorto{ita} La persono, la kozo, qua maxime distas de ta qua parolas; la persono, la kozo, quan onu nomis enumer-unesme
\artiklo{*itemo}%
\vorto{*itemo} Singla ek la frazi (generale kun nombri) en enumero, etato, budjeto, fakturo, kalkulo\lingui{EFI}
\artiklo{iterar}%
\vorto{iterar}\fako{trans.} Ri-agar lo sama\lingui{EFIS}
\artiklo{itinerario}%
\vorto{itinerario}\senco{1}{Indiko pri la voyo uzenda}\senco{2}{Indiko pri omna loki ube onu pasas por irar de lando ad altra, qua kontenas ulakaze deskripti di la loki, o noti e voyajo-memori}\senco{3}{Indiko di omna stacioni qui existas sur la parkuro di relvoyo, di la diversa treni qui pasas en oli; di la transporto-preco, e c.}\lingui{DEFIS}
\artiklo{-iv}%
\vorto{-iv} Sufixo qua signifikas (soldite a radiko verbala) « kapabla, povanta, apta por… »
\artiklo{-ivoro}%
\vorto{-ivoro} Sufixo qua signifikas « qua su nutras ye… »
\artiklo{ivoro}%
\vorto{ivoro}\senco{1}{Materio ye qua konsistas la dentegi di elefanto}\senco{2}{Materio ye qua konsistas la denti di ula pakidermi}\lingui{EFI}
\artiklo{-iz}%
\vorto{-iz} Sufixo qua signifikas « provizar, garnisar, indutar, impregnar ye… »
\artiklo{izabela}%
\vorto{izabela} Pale-flava\lingui{DEFIS}
\artiklo{izobaro}%
\vorto{izobaro}\fako{fiziko} Lineo ek la punti di la Terglobo ube la preso atmosferala esas sama ye instanto determinita\lingui{DEFIRS}
\artiklo{izocela}%
\vorto{izocela}\fako{geom.} Di qua la lateri esas inter-egala\lingui{EFIS}
\artiklo{izodinama}%
\vorto{izodinama}\fako{fiziko} Di qua la forco esas inter-egala sur la du lateri\lingui{DEFIS}
\artiklo{izogama}%
\vorto{izogama}\fako{biol.} Dicesas pri la vejetanti inferiora, che qui la genito-elementi, qui interunionas por formacar la ovo, esas intersimila\lingui{DEFIS}
\artiklo{izogona}%
\vorto{izogona}\fako{geom.} Di qua la anguli esas inter-egala\lingui{DEFIS}
\artiklo{izoklino}%
\vorto{izoklino}\fako{magnet.} Lineo ek la punti di la Terglobo ube la deviaco da la magnet-indexo esas intersama\lingui{DEFIS}
\artiklo{izokrona}%
\vorto{izokrona}\fako{fiziko} Qua eventas sama-dure\lingui{DEFIS}
\artiklo{izolar}%
\vorto{izolar}\fako{trans., de} Separar de omna kozo : (a) igar (korpo) ne plus kontaktar kun korpi elektro-konduktiva; (b) separar korpi ek la elementi kemiala kun qui lu esas kombinita; (c) igar (ulu) vivar for lua kunhomi\lingui{DEFIRS}
\artiklo{izomera}%
\vorto{izomera}\fako{kemio} Qua konsistas ye atomi qualeso-sama e quanteso-sama\lingui{DEFIS}
\artiklo{izometra}%
\vorto{izometra}\fako{geom.) (pri perspektivo} Olta en qua la axi komparala esas inter-egala\lingui{DEFIRS}
\artiklo{izomorfa}%
\vorto{izomorfa}\fako{kemio} Di qua la formo kristala esas sama kam olta di korpo kemiala altra\lingui{DEFIRS}
\artiklo{izoperimetra}%
\vorto{izoperimetra}\fako{geom.} Dicesas pri la figuri di qui la perimetri esas inter-egala
\artiklo{izotermo}%
\vorto{izotermo}\fako{fiziko} Lineo, sur mapo, qua pasas tra omna punti di la Terglobo ube la temperaturo mezvalora dum la yaro esas inter-sama\lingui{DEFIRS}
\artiklo{izotonika}%
\vorto{izotonika}\fako{biol.} Dicesas pri la equilibro molekulala di du solvuri interseparita da membrano organika, e qui havas osmo-povo interegala
\artiklo{izotopa}%
\vorto{izotopa}\fako{kemio} Dicesas pri la kompozaji kemiala inter-identa, ma di qui la pezo atomala interdiferas\lingui{F}
\artiklo{izotropa}%
\vorto{izotropa} Qua havas la sama partikularaji fizikala omna-direcione
\end{multicols}
\sekciono{J}
\begin{multicols}{2}
\artiklo{ja}%
\vorto{ja} Sen vartar til kande fluabos multa tempo
\artiklo{jaboto}%
\vorto{jaboto} Ornivo ek muslino, ek dentelo, muntita an la aperturo di kamizo, sub la fauco, e qua su extensas sur la pektoro\lingui{DEFR}
\artiklo{jacar}%
\vorto{jacar}\fako{netrans.} Esar kushinta, sterninta, pozita\lingui{FIS}
\artiklo{jado}%
\vorto{jado} Lapido tre harda, bele olivea, sorto di silikato di alumino e kalko\lingui{DEFIS}
\artiklo{jaguaro}%
\vorto{jaguaro}\fako{zool.} Amerika-tigro, la maxim granda ek la felini pos la tigro e la leono, di qua la pilaro esas falva kun veinizuri ek nevusi nigra\latina{felix onca}\lingui{DEFIRS}
\artiklo{jako}%
\vorto{jako} Viro-vesto qua decensas kelke adsub la genuo, quan olim *weris la rurani, la plebeyi\lingui{DEFIRS}
\artiklo{jakobino}%
\vorto{jakobino} Partisano di la doktrino qua agnoskas ke la plebeyo esas irgo-povanta, irgo-darfanta\lingui{DEFIRS}
\artiklo{jakoneto}%
\vorto{jakoneto} Speco di muslino mi-diafana, de qua onu facas robi, kolumi, muliero-kofii\lingui{DEFIRS}
\artiklo{jalapo}%
\vorto{jalapo}\fako{bot.} Konvolvulo di qua la radiko esas purgiva\latina{convonvulus jalapa}\lingui{DEFIS}
\artiklo{jaluza}%
\vorto{jaluza} Di qua la amoro esas desquieta pro suspektemeso, desfidemeso\lingui{EFI}
\artiklo{janicharo}%
\vorto{janicharo} Elito-soldato ek la infantrio Turka, ek la guardo di la sultano\lingui{DEFIRS}
\artiklo{jansenismo}%
\vorto{jansenismo}\fako{religio kristana} Doktrino da Jansenius, ecese severa, kondamnita da la Roma-eklezio pro imputar tro multe a Graco, tro poke ad arbitrio, e qua duktas a la predestineso absoluta\lingui{DEFIRS}
\artiklo{januaro}%
\vorto{januaro} La monato unesma di la yaro (segun la maniero moderna)\lingui{DEFIRS}
\artiklo{jardiniero}%
\vorto{jardiniero}\senco{1}{Orno-moblo qua suportas kesto provizita ye flori}\senco{2}{Disho ek diversa sorti di legumi}\lingui{DEFIS}
\artiklo{jargonar}%
\vorto{jargonar}\fako{trans.}\senco{1}{Dicar ulo ye maniero nekomprenebla pro nekorekteso vortala, gramatikala, pronuncala, e c.}\senco{2}{Parolar idiomo stranjera koram ulu qua judikas olu kom barbara pro ne komprenar olu}\lingui{DEFIRS}
\artiklo{jaro}%
\vorto{jaro} Granda vazo de terakotajo, kun du ansi, en qua onu konservas aquo, oleo\lingui{EFIS}
\artiklo{jasmino}%
\vorto{jasmino}\fako{bot.} Arbusto sarmentoza, genero tipo di la familio « jasminei », di qua la flori esas odoroza\latina{jasminus}\lingui{DEFIRS}
\artiklo{jaspo}%
\vorto{jaspo} Lapido harda ed opaka, sama-natura kam agato, koloroza ye reda, flava o verda, uniforme, bendope o makulope\lingui{DEFIS}
\artiklo{javelino}%
\vorto{javelino} Armo dardo, longa e dina, uzita da la antiqui, qui lansis olu manue\lingui{EFIS}
\artiklo{jeleo}%
\vorto{jeleo}\fako{koquarto}\senco{1}{Suko de substanco animalala, qua solideskis per koldesko}\senco{2}{Suko konjelita de frukti}\lingui{DEFIS}
\artiklo{jelozio}%
\vorto{jelozio} Shutro, ma ek planketi paralela, horizontala, qui povas rotacar cirkum sua axo longesala, quan onu muntas avan fenestro, quan onu povas levar o deslevar per kordo o kordono; kande olu esas tote levita, la lati esas aplikita ici sur iti; e dop qua onu povas vidar sen videsar\lingui{DEFIS}
\artiklo{jemar}%
\vorto{jemar}\fako{netrans.} Agar lamento ne-artikulita\lingui{FIS}
\artiklo{jemelo}%
\vorto{jemelo}\senco{1}{Qua naskis sama-parture kam altra filio (o yuno)}\senco{2}{\textit{(plurale)} Kastor e Polux, un ek la 12 signi di zodiako}\lingui{EFIS}
\artiklo{jenar}%
\vorto{jenar}\fako{trans.} Igar nekomfortoza, nekomoda, per presar, klemar, opresar\lingui{DF}
\artiklo{jendarmo}%
\vorto{jendarmo} Soldato qua apartenas a korpo di qua la tasko konsistas ye sekurigar la vivo di la civitani\lingui{DEFIS}
\artiklo{jenio}%
\vorto{jenio}\fako{armeo} La korpo ek la oficiri e soldati qui aplikas la cienci a la fortifiko di la fortresi, a la konstrukto di navi, di ponti, e c.\lingui{DFI}
\artiklo{jenjivito}%
\vorto{jenjivito}\fako{patol.} Inflameso di la jenjivi\lingui{EFI}
\artiklo{jenjivo}%
\vorto{jenjivo} Tisuo karna qua garnisas la arkadi dentala e la perimetro di la denti per adherar a lia kolo\lingui{FI}
\artiklo{jenro-piktarto}%
\vorto{jenro-piktarto}\fako{arto} Qua traktas fantazio-temi (opoze a la pikto di temi historiala, peizajala)\lingui{DEF}
\artiklo{jentila}%
\vorto{jentila} Qua plezas delikate\lingui{EFIS}
\artiklo{jeo}%
\vorto{jeo}\fako{zool.} Ucelo, genero « paseri konobeka », vicina ye la korvi\latina{oarroloa ElandariuE}\lingui{EF}
\artiklo{jerar}%
\vorto{jerar}\fako{trans.} Administrar kom reprezentero di la proprietero (domeno, industrio-firmo, entraprezajo, e c.)\lingui{FIS}
\artiklo{jermo}%
\vorto{jermo}\senco{1}{Rudimento di la embriono, destinita a *reproduktar la planto, la animalo}\senco{2}{\textit{(metaf.)} Principo, elemento di la kresko di ulo}\lingui{EFIRS}
\artiklo{jerzeo}%
\vorto{jerzeo} Muliero-korsajo ek texajo elastika, qua tre strete fitas la busto\lingui{EF}
\artiklo{jetar}%
\vorto{jetar}\fako{trans.} Sendar aden spaco per impulso rapida\lingui{FI}
\artiklo{jezuito}%
\vorto{jezuito}\senco{1}{Nomo di la membri di la Iesu-ordeno quan fondis Ignatius de Loyola}\senco{2}{Hipokrito, quan oportas desfidar}\lingui{DEFIRS}
\artiklo{jibeto}%
\vorto{jibeto} Asembluro ek fosto vertikala, e trabo horizontala sur la libera extremajo di qua glitas la kordo a qua onu suspendos la persono pendeso-kondamnita\lingui{EF}
\artiklo{jigo}%
\vorto{jigo} Ario dansala, rapida-ritma\lingui{DEFIRS}
\artiklo{jileto}%
\vorto{jileto} Parto di la viro-vesto, qua kovras la torso, sen maniki, qua weresas sub la paltoto o la redingoto\lingui{F}
\artiklo{jilfloro}%
\vorto{jilfloro}\fako{bot.} Planto, genero ek la familio « kruciferi », di qua la flori esas blanka, flava o redatra, e di qua la odoro memorigas olta di la kariofilo\latina{cheiranthus cheiri}\lingui{E}
\artiklo{jineto}%
\vorto{jineto}\senco{1}{\textit{(zool.)} Mikra mamifero di la genero « karnivori digitigrada », di qua la felo uzesas kom furo}\senco{2}{\textit{(tekn.)} Morso di qua la kateno esas ringatra}\latina{viverra genetis}
\artiklo{jinjero}%
\vorto{jinjero}\fako{bot.} Planto de India, di qua la radiko aromata uzesas kondimente\latina{amomum zingiber}\lingui{EFIS}
\artiklo{jino}%
\vorto{jino} Brandio distilita sur junipero-beri\lingui{DEFIRS}
\artiklo{jirafo}%
\vorto{jirafo}\fako{zool.} Mamifero di qua la kolo esas tre longa, e la pilaro makuloza\latina{camelopardalis giraffa}\lingui{DEFIRS}
\artiklo{jirar}%
\vorto{jirar}\fako{netrans.} Parkurar cirklo-kurvo, o fragmenton di cirklo-kurvo\lingui{EFI}
\artiklo{jiravolto}%
\vorto{jiravolto} Jiro e retrojiro, agita rapide\lingui{FI}
\artiklo{jokar}%
\vorto{jokar}\fako{netrans.} Agar o dicar ulo quo distraktas ridigive\lingui{EI}
\artiklo{jokeo}%
\vorto{jokeo} Viro qua kavalkas kavalo en la kavalo-konkursi kurala, ed agas lo profesionale\lingui{DEFRS}
\artiklo{jonglar}%
\vorto{jonglar}\fako{netrans.} Lansar adciele e simultane plura buli, plura objekti, qui krucifas pasante de la manuo dextra a la manuo sinistra, od inverse\lingui{DEFR}
\artiklo{jonko}%
\vorto{jonko} Navo, di qua la pruo esas alta e krosatre kurvigita, uzata en Chinia e Japonia\lingui{DEFIS}
\artiklo{jonquilo}%
\vorto{jonquilo}\fako{bot.} Genero di narciso, planto kultivata en la gardeni pro la eleganteso di olua su-teno e pro la dolceso di olua odoro\latina{narcissus jonquilla}\lingui{DEFIRS}
\artiklo{jorno}%
\vorto{jorno} Klareso quan Suno produktas sur tero\lingui{FI}
\artiklo{jovdio}%
\vorto{jovdio} La dio kinesma di la semano\lingui{FIS}
\artiklo{joyar}%
\vorto{joyar}\fako{netrans., de, pro, pri} Tre plezur-impresesar\lingui{EFIS}
\artiklo{juar}%
\vorto{juar}\fako{trans.} Havar plezuro extreme granda de la posedo di ulo
\artiklo{jubeo}%
\vorto{jubeo}\fako{trans.} En kirko : quaza galerio inter la navo e la koreyo\lingui{EF}
\artiklo{jubilear}%
\vorto{jubilear}\fako{trans.} Celebrar festo, honore di la persono qua esas en sua ofico de kinadek yari\lingui{DEFIRS}
\artiklo{judiciar}%
\vorto{judiciar}\fako{trans.} Enuncar (pri ulu od ulo) decido, kom persono di qua la ofico konsistas ye aplikar la legi punisala, civila komercala\lingui{EFIS}
\artiklo{judikar}%
\vorto{judikar}\fako{trans.} Enuncar opiniono per qua onu aprobas o blamas : opinionar pri persono o kozo\lingui{EFIS}
\artiklo{jugulara}%
\vorto{jugulara}\fako{anat.} \textit{(veini)} qui apartenas a la fauco\lingui{DEFIS}
\artiklo{jujubo}%
\vorto{jujubo}\fako{bot.} Frukto (nukleo du-celula, kontenata da envelopilo pulpa, e di qua la suko uzesas farmaciale kom pektorala) di arboro ek la familio « ramnei », vicina ye la ilexo e ye la fuzeno, sinuoza-brancha, kun forta dorni, qua kreskas en la sudo-regioni di Europa\latina{zizyphus}\lingui{EFR}
\artiklo{julepo}%
\vorto{julepo}\fako{farmacio} Pociono dolcigiva, ek aquo e siropi, a qui onu adjuntas mikra dozo di opiumo, e quan onu donas drinkenda kom kalmigivo o dormigivo\lingui{DEFIS}
\artiklo{julieno}%
\vorto{julieno}\fako{koquarto} Supo ek plura sorti di herbi e di legumi (karoti, napi, celerio, kauli, pizi, e c.), hachita tenue e koquita en buliono\lingui{DEF}
\artiklo{julio}%
\vorto{julio} La monato sepesma di la yaro (segun la maniero moderna)\lingui{DEFIRS}
\artiklo{jungar}%
\vorto{jungar}\fako{trans.} Ligar (tir-bestii) a veturo, a plugilo; inteligar vagoni sur relvoyo\lingui{L}
\artiklo{junglo}%
\vorto{junglo} En India : plana lando, marshoza, kovrita ye kani e ye bushi dika e dornoza\lingui{EF}
\artiklo{junio}%
\vorto{junio} La monato sisesma di la yaro (segun la maniero moderna.)\lingui{DEFIRS}
\artiklo{juniora}%
\vorto{juniora}\fako{la persono} Qua naskis pos (frato)\lingui{EF}
\artiklo{junipero}%
\vorto{junipero}\fako{bot.} Arbusto ek la familio « cipresinei », di qua la beri esas aromatoza\latina{juniperus}\lingui{EFI}
\artiklo{junko}%
\vorto{junko}\fako{bot.} Genero tipa di la familio « junkacei », planto herbatra, rekta e flexebla, qua kreskas en la tereni humida\latina{juncus}\lingui{FIS}
\artiklo{juntar}%
\vorto{juntar}\fako{trans.}\senco{1}{Proximigar du kozi ye un ek olia extremaji respektiva, por ke oli interkontaktez od intertenez}\senco{2}{\textit{(aludante plura personi, trupi, e c.)} Irar unionar (su) ad (altra personi, altra trupi)}\lingui{EFIS}
\artiklo{jupo}%
\vorto{jupo} Parto di la muliero-vesto, qua kovras de la zono til la pedi\lingui{FR}
\artiklo{jurar}%
\vorto{jurar}\fako{trans.} Promisar, afirmar (ulo ad ulu) selektante (generale) deo kom testo pri lo\lingui{EFIS}
\artiklo{juraso}%
\vorto{juraso}\fako{geol.} Tereno ye formaceso sekundara inter la strato liasa e la tereno kretatra, quan onu renkontras precipue en la Jura-monti (Francia)\lingui{DEF}
\artiklo{jurio}%
\vorto{jurio}\fako{en la tribunalo kriminala} Ensemblo ek la civitani qui devas decidar ka la akuzato esas (o ne esas) kulpoza\lingui{DEFIS}
\artiklo{juriskonsulto}%
\vorto{juriskonsulto} La persono qua konsilas profesione pri la yuro-problemi\lingui{DEFIRS}
\artiklo{jurisprudenco}%
\vorto{jurisprudenco} Interpretado di la legi per la verdikti da la judiciisti : o, generale, maniero aplikar la legi a la kazi partikulara\lingui{DFIS}
\artiklo{jurnalo}%
\vorto{jurnalo} Edituro singladia o periodala, qua donas la niuzi politikala, literaturala, ciencala\lingui{DEFIS}
\artiklo{jus}%
\vorto{jus} Pasis nur kelka instanti de kande…\lingui{DEFIS}
\artiklo{justa}%
\vorto{justa}\senco{1}{Qua precize fitas, su adaptas a to por quo olu destinesis. (Dicesas pri mezuro, balanco, voco, vorto, penso)}\senco{2}{Qua agas o per quo onu agas perfekte lo exekutenda}\senco{3}{Qua judikas korekte (en la kazo konsiderata)}\lingui{EFIS}
\artiklo{justifikar}%
\vorto{justifikar}\fako{trans.}\senco{1}{Deklarar, demonstrar la inocenteso}\senco{2}{Deklarar, montrar la legitimeso}\lingui{DEFIS}
\artiklo{juto}%
\vorto{juto} India-kanabo, qua uzesas por facar fili e texuri ordinara\lingui{DEFIS}
\artiklo{juvelo}%
\vorto{juvelo} Ornamento tre valoroza, ek oro, arjento, lapidi, qua uzesas por la orno di la mulieri\lingui{DEFIRS}
\end{multicols}
\sekciono{K}
\begin{multicols}{2}
\artiklo{ka(d)}%
\vorto{ka(d)} Konjunciono questionala, sive por la questiono direta, sive por la nedireta, e qua sempre komencas la propoziciono
\artiklo{kabalo}%
\vorto{kabalo} Doktrino mistika, naskinta de la tradicionaji Biblala, qui modifikesis da la idei orientala\lingui{DEFIRS}
\artiklo{kabano}%
\vorto{kabano} Mikra habiteyo ek tero, ligno, e c.\lingui{EF}
\artiklo{kabino}%
\vorto{kabino}\senco{1}{\textit{(navig.)} Chambreto en navo}\senco{2}{Chambreto en qua onu su desvestizas jus ante balnar en maro od en rivero}\lingui{DEFIRS}
\artiklo{kablo}%
\vorto{kablo}\senco{1}{Tre grosa kordego de kanabo (uzata en la navi)}\senco{2}{Tre grosa kordego ek elektro-filo}\lingui{DEFIR}
\artiklo{kabo}%
\vorto{kabo}\fako{geogr.} Pinto ek tero, montetatra, qua salias aden maro\lingui{DEFIS}
\artiklo{kabrar}%
\vorto{kabrar}\fako{netrans.} \textit{(aludante kavalo qua pavoras e su defensas)} Elevar sua gambi avana, per erektar su sur sua gambi dopa\lingui{FS}
\artiklo{kabrioleto}%
\vorto{kabrioleto} Veturo lejera, un-kavala, kun aperturo avana e kun mi-tendo ledra, faldebla abanike\lingui{DEFIS}
\artiklo{kachuo}%
\vorto{kachuo} Substanco odoroza quan, en India, onu extraktas de la ligno, de la folii, de la sheli fresha di la akacio \textit{katechu}, e quan onu uzas en la tint-industrio, kom kolorizivo; en medicino, kom astriktivo e tonizivo\lingui{DEFIR}
\artiklo{kadastro}%
\vorto{kadastro} Registro publika, qua kontenas la relevo generala, la mezuro-rezultajo e la evaluo pekuniala di la domeni, ter-havaji, uzata kom fundamento por kalkular la ter-imposto\lingui{DEFIRS}
\artiklo{kadavro}%
\vorto{kadavro} Korpo mortinta, homala o bestiala\lingui{DEFIS}
\artiklo{kadencar}%
\vorto{kadencar}\fako{trans.}\senco{1}{Finar frazo, verso, hemistiko tale, ke la voco acentizas ta finalo per plufortigar olua pronunceso}\senco{2}{Produktar ritmo per acentizar simetrala ta finali, en muziko od en poezio}\senco{3}{Produktar movo ye intervali inter-egala ed uniforma, qua imitas la ritmo muzikala}\lingui{DEFIRS}
\artiklo{kadeto}%
\vorto{kadeto}\fako{armeo} Singla ek la membri di ta kompanii qui konsistis ye la familio-juniori volontaria\lingui{DEFIRS}
\artiklo{kadmio}%
\vorto{kadmio}\fako{kemio} Korpo simpla, forjebla, duktila, di qua la rupto produktas fibri, di qua la koloro e la brilo esas olti di stano\simbolo{Cd}\lingui{DEFIRS}
\artiklo{kadro}%
\vorto{kadro}\senco{1}{Bordumo irga-forma qua cirkondas pikturo, gravuro, spegulo}\senco{2}{Bordumo menuzura}\senco{3}{To quo limitizas ceno, temo; 2. \textit{(armeo)}. La korpo ek la oficiri generala, oficiri, sub-oficiri, di la divizioni, secioni, e c., imperote da qui la soldati su asemblas}\lingui{EFS}
\artiklo{kaduceo}%
\vorto{kaduceo} Emblemo di Merkurio, vergeto kronizita ye du mikra ali, cirkum qua esas du serpenti plektita – simbolo di komerco, eloquenteso, paco\lingui{DEFIRS}
\artiklo{kaduka}%
\vorto{kaduka}\senco{1}{Vicina ye sua krulo, dekado}\senco{2}{Quan ula defekto formala igas nihiligota tre balde}\lingui{EFIS}
\artiklo{kafeino}%
\vorto{kafeino} Alkaloido naturala, kristaligebla, quan onu extraktas de kafeo, teo, kakao\lingui{DEFIS}
\artiklo{kafeo}%
\vorto{kafeo}\fako{bot.}\senco{1}{Grano ek arboreto, sempre verda, ek Arabia, di qua la frukto esas grosa quale cerizo e konsistas ye shelo qua envelopas pulpo mucilajatra qua ipsa cirkondas du alveoli un-grana; ta grano, torefaktita, muelita ed infuzita, donas drinkajo agreabla, ecitiva e toniziva}\senco{2}{Infuzuro ek kafeo-grani torefaktita e muelita}\latina{coffeo arabica}\lingui{DEFIRS}
\artiklo{kaftano}%
\vorto{kaftano} Honor-peliso quan la sultani ofris a la oficiri precipua, a la ambasadisti stranjera, a personegi tre eminenta, e c. kom donaco\lingui{DEFIRS}
\artiklo{kagulo}%
\vorto{kagulo} Quaza froko sen-manika, qua kovris la kapo e la parti cetera di la korpo, e qua provizesis ye aperturi qui korespondis a la okuli ed a la boko\lingui{FS}
\artiklo{kaito}%
\vorto{kaito} Puero-ludilo, ek framo lejera qua esas kovrita ye papero tensita, quan onu igas acensar aden atmosfero per kurar kontre-vento e tenante olu per kordeto\lingui{E}
\artiklo{kajo}%
\vorto{kajo}\senco{1}{Quaza chambro, garnisita ye stangi ek ligno, ek fero, en qua onu enklozas animali ne-amansita por gardar li en menajerio o por transportar li}\senco{2}{Quaza chambreto portebla, garnisita ye latunfili, en qua onu retenas enklozite uceli}\lingui{EFI}
\artiklo{kakao}%
\vorto{kakao}\fako{bot.} Grano kontenata da la frukto di arbusto ek la familio « malvacei » e qua, grilite e triturite kun sukro, donas chokolado\latina{theobroma cacao}\lingui{DEFIRS}
\artiklo{kakar}%
\vorto{kakar}\fako{netrans.} Vakuigar sua ventro de la grosa exkrementi\lingui{DEFIS}
\artiklo{kakatuo}%
\vorto{kakatuo}\fako{zool.} Speco di papagayo penachoza, ucelo klimera, di qua la plumaro esas bele blanka\latina{piictolophus}\lingui{DEFIRS}
\artiklo{kakexio}%
\vorto{kakexio}\fako{patol.} Stando maladatra, quan karakterizas deperiso generala\lingui{DEFS}
\artiklo{kakofonio}%
\vorto{kakofonio}\senco{1}{Konsonanco qua lezas la oreli}\senco{2}{\textit{(muziko)} Asemblajo ek son-uni ne-akordanta}\lingui{DEFIS}
\artiklo{kaktuso}%
\vorto{kaktuso}\fako{bot.} Planto di la familio « kaktei », qua donas frukto laxigiva, figatra\latina{cactus}\lingui{DEFIRS}
\artiklo{kalamino}%
\vorto{kalamino}\fako{kemio} Nomo vulgara di la hidro-silikato di zinko\simbolo{SiO₅Z₂H₂}\lingui{DEFIS}
\artiklo{kalamitato}%
\vorto{kalamitato} Granda desfortuno qua unafoye frapas tre multa personi e mem populo\lingui{DEFIS}
\artiklo{kalamito}%
\vorto{kalamito} Rezino qualeso-inferiora, quan onu rekoltas de la kanostipi\latina{styrax calamita}\lingui{DEFIS}
\artiklo{kalamo}%
\vorto{kalamo} Fragmento de kano, quan la antiqui uzis por skribar\lingui{DFIS}
\artiklo{kalandrar}%
\vorto{kalandrar}\fako{trans.} Pasigar sub cilindro glata la stofi, drapi, teli, e c. quin onu volas briligar, muarizar, e c.\lingui{DEFS}
\artiklo{kalcedono}%
\vorto{kalcedono}\fako{mineral.} Varietato di agato, di qua la diafaneso esas laktatra\lingui{DEFIRS}
\artiklo{kalcinar}%
\vorto{kalcinar}\fako{trans.} Transformar metalo aden oxido, per submisar olu a la efiko da fairo intensa, en aero libera\lingui{DEFIS}
\artiklo{kalcio}%
\vorto{kalcio}\fako{kemio} Korpo simpla, metala, flave blanka\simbolo{Ca}\lingui{DEFIS}
\artiklo{kaldiero}%
\vorto{kaldiero} Vazo klozita (apartenanta, generale, a vapormashino) en qua onu boliigas aquo por obtenar vaporo\lingui{EFIS}
\artiklo{kaldrono}%
\vorto{kaldrono} Mikra recipiento, ordinare de kupro, en qua onu koquas la nutrivi o boliigas aquo\lingui{EFIS}
\artiklo{kalefaktar}%
\vorto{kalefaktar}\fako{netrans.) (fiziko} Aquo-guteti, disjetita sur metala plako kelke varmigita; vaporeskas tre rapide; ma se la surfaco esas tre varma, nula ebulio eventas, e la guteti « kalefaktas », t.e. li inter-asemblas e formacas globeto sfera\lingui{DEFIS}
\artiklo{kalemburo}%
\vorto{kalemburo} Vorto-ludo fondata sur la senco-difero inter vorti qui audale aspektas inter-sama\lingui{DFR}
\artiklo{kalendario}%
\vorto{kalendario}\senco{1}{Sistemo di divido di tempo, per yari, monati, e dii}\senco{2}{Tabelo ek la dii, monati, sezoni, pri singla yaro, qua generale indikas, pri singla dio, la nomo di santo}\lingui{DEFIRS}
\artiklo{kalendo}%
\vorto{kalendo} Nomo donita da la Romani antiqua a la dio unesma di singla monato\lingui{DEFIRS}
\artiklo{kalendulo}%
\vorto{kalendulo}\fako{bot.} Planto ek la familio « kompozaji », di qua la kapituli esas flava e radio-dispozita\latina{calendula}
\artiklo{kalenturo}%
\vorto{kalenturo}\fako{patol.} Kazo di febro ardorigiva, kun deliro, qua eventas ulafoye che la navani kande la navo trairas la zono torida\lingui{DEFIS}
\artiklo{kalesho}%
\vorto{kalesho} Veturo eleganta, quar-rota, maxim-multa-kaze apertita avane, e provizita (dope) ye tendo ek ledro qua esas desfaldebla e retrofaldebla segun-vole\lingui{DEFIR}
\artiklo{kalfatar}%
\vorto{kalfatar}\fako{trans.} Igar aquo-nepermeebla (barko, ramo-batelo), per stopar la junturi, la fenduri, la trui, per stupo gudroimpregnita\lingui{DFIS}
\artiklo{kalibro}%
\vorto{kalibro} Modelo sur qua esas trasita la konturi, la dimensioni, di la objekti fabrikenda\lingui{DEFIRS}
\artiklo{kalico}%
\vorto{kalico}\senco{1}{Vazo sakra en qua la sacerdoto sakrigas la vino di la eukaristio}\senco{2}{Envelopilo qua kovras la infra parto di la flor-korolo}\lingui{EFIS}
\artiklo{kalidoskopo}%
\vorto{kalidoskopo} Aparato qua kontenas mikra objekti multa-kolora qui su situas diverse ed intersimetrie kande onu agitas la instrumento\lingui{DEFIRS}
\artiklo{kalifo}%
\vorto{kalifo} Suvereno Mohamedista, qua unionas en su la povo politikala e la povo religiala\lingui{DEFIRS}
\artiklo{kaligrafar}%
\vorto{kaligrafar}\fako{trans.} Skribar (ulo) per literi eleganta, ornita\lingui{DEFS}
\artiklo{kaliko}%
\vorto{kaliko} Telo de kotono, min delikata kam perkalo\lingui{DEFIRS}
\artiklo{kalio}%
\vorto{kalio}\fako{kemio} Korpo simpla, metala, volatila, qua oxideskas per kontaktar kun aero humida\simbolo{K}\lingui{D}
\artiklo{kalkaneo}%
\vorto{kalkaneo}\fako{anat.} Osto di la tarso, qua formacas la talono\lingui{EF}
\artiklo{kalko}%
\vorto{kalko}\fako{kemio} Kali-oxido, alkalio minerala\lingui{DEFIRS}
\artiklo{kalkolo}%
\vorto{kalkolo}\fako{patol.} Konkreciono stonatra qua naskas acidente en ula organi (hepato, veziko, e c.)\lingui{EFIS}
\artiklo{kalkular}%
\vorto{kalkular}\fako{trans.}\senco{1}{Determinar, per operaci matematikala, ye nombri donita, nombron quan onu serchas}\senco{2}{\textit{(metaf.)} Agar sua aranji, kombinar la kozi, egarde di la skopo atingenda}\lingui{DEFIRS}
\artiklo{kalma}%
\vorto{kalma} Stando di to quo indijas agiteso\lingui{EFIS}
\artiklo{kalmaro}%
\vorto{kalmaro}\fako{zool.} Molusko ek la genero « sepii », qua sekrecas liquido nigra quan lu disjetas cirkum su kande onu atakas lu\latina{loligo}\lingui{DEFIRS}
\artiklo{kalo}%
\vorto{kalo}\fako{patol.} Dikesko od hardesko di la pelo, produktita da fricionado, ye la palmo, ye la plando, sur la pedo-fingri\lingui{EFIS}
\artiklo{kalomelo}%
\vorto{kalomelo}\fako{farmacio} Sub-kloruro di merkurio, korpo blanka, nedissolvebla en aquo, uzata medicine kom purgivo, e kom kontre-vermo\lingui{DEFIRS}
\artiklo{kaloriko}%
\vorto{kaloriko} Kaloro-principo quan kontenas la korpi\lingui{EFIS}
\artiklo{kalorimetrio}%
\vorto{kalorimetrio} La cienco qua traktas la mezuro di la kalorii\lingui{DEFIS}
\artiklo{kalorimetro}%
\vorto{kalorimetro} Fizik-instrumento, uzata por determinar la kaloro specifika di la korpi\lingui{DEFIS}
\artiklo{kalorio}%
\vorto{kalorio}\fako{fiziko} Mezur-unajo pri kaloro : la quanto de kaloro qua esas necesa por plualtigar ye un grado centigrada la temperaturo di un kilogramo de aquo\lingui{DEFIRS}
\artiklo{kaloro}%
\vorto{kaloro} Kauzo fizikala di ula modifiki molekulala di la korpo, vario da volumino – dilateso; vario di stando – liquidesko, vaporesko\lingui{EFIS}
\artiklo{kaloto}%
\vorto{kaloto}\senco{1}{Mikra boneto ek drapo, veluro, e c., ronda, qua kovras la suprajo di la kapo}\senco{2}{\textit{(geom.)} En sfero quan sekas plano : la surfaco di un ek la du parti, precipue di la plu mikra}\senco{3}{Quaza kapuco ek stano, ye qua onu envelopas la stopilo e la kolo di la flakoni, di la liquor-boteli, e c.}\lingui{DEFI}
\artiklo{kalquar}%
\vorto{kalquar}\fako{trans.} Kopiar desegnuro tra papero-folio diafana, per streki kontinua\lingui{DFIS}
\artiklo{kalsono}%
\vorto{kalsono} Suba vesto, quaza bracho, ek telo, flanelo, e c.\lingui{FRS}
\artiklo{kalumniar}%
\vorto{kalumniar}\fako{trans.} Dicar neexaktaji mala pri (ulu)\lingui{EFIS}
\artiklo{kalva}%
\vorto{kalva} Qua havas nur tre poka hari\lingui{DEFIS}
\artiklo{kalvario}%
\vorto{kalvario}\fako{kristanismo}\senco{1}{Kolino quan acensis Iesu Kristo, kande lu portis la kruco, e sur la somito di qua lu krucagesis}\senco{2}{Kolino, sur la suprajo di qua esas kruco, adube onu iras pilgrime, haltante e pregante ye la loki di la voyo, qui memorigas la cirkonstanci precipua di la morto di Iesu Kristo. – \textit{(anke metaf.)}}\lingui{DEFIS}
\artiklo{kalzo}%
\vorto{kalzo} Texuro ek lano, kotono, silko, *nilono, fitigita a la pedo ed a la gambo por kovrar olu\lingui{FIS}
\artiklo{kam}%
\vorto{kam} Partikulo uzata kom termino duesma en la gradi komparala : tam… kam; sama… kam; tale… kam; preferar… kam, e c.
\artiklo{kamalio}%
\vorto{kamalio}\senco{1}{En la Mezepoko, protektilo, konsistanta ye fera kaloto en qua esas muntita texuro ek mashi qua protektis la kolo e la shultri}\senco{2}{\textit{(religio katol.)} Quaza kapuco pelerina, quan la sacerdoti, la kantori metas sur la rocheto; mikra mantelo quan metas sur la rocheto la kardinali, la episkopi, la prelati, e c.}\senco{3}{Mikra mantelo mulierala, sen maniki, kun o sen kapuco, diversa-forma (segun la modo)}\lingui{DEF}
\artiklo{kamarado}%
\vorto{kamarado} La persono qua habitas normale en la sama loko kam ulu, e tale divenas, ad ica, kelke familiara\lingui{DEFIRS}
\artiklo{kamarilio}%
\vorto{kamarilio} Ensemblo ek la personi qui vivas intime kun princo, e probas influar lu politikale. (Vorto desprizala)\lingui{DEFIS}
\artiklo{kambiar}%
\vorto{kambiar}\fako{trans., po} Donar ad ulu (ulo) e recevar de lu (altra kozo kom equivaloza)\lingui{IS}
\artiklo{kambrar}%
\vorto{kambrar}\fako{trans.} Kurvigar arkatre\lingui{EF}
\artiklo{kameleono}%
\vorto{kameleono}\senco{1}{\textit{(zool.)} Reptero sauria, quaza lacerto kun grosa kapo, pelo shagrinatra, qua divenas diversa-kolora – ne (quankam onu kredis lo) pro reflektar la objekti vicina, ma segun ke la korpo esas plu o min plena ye aero}\senco{2}{\textit{(metaf.)} Persono qua modifikas sua opiniono, sua konduto, segun la okazioni, segun la cirkonstanci}\latina{chameleo vulgaris}\lingui{DEFIRS}
\artiklo{kamelio}%
\vorto{kamelio}\fako{bot.} Orno-planto, arbusto ek la familio « teacei »\latina{camellia}\lingui{DEFIRS}
\artiklo{kamelo}%
\vorto{kamelo}\fako{zool.} Quadripedo ruminera, longa-gamba, kun un o plura gibi sur la dorso, quan la Arabi, la Orientani, uzas kom charjo-bestio\latina{camelus bactrianus}\lingui{DEFIS}
\artiklo{kameno}%
\vorto{kameno} Konstrukturo masonura, aptigita por la kombusto, e qua konsistas esencale ye herdo qua komunikas kun la domexterajo per tubo, olqua produktas la aero-tiro ed esas la paseyo di la fumuri; la parto dom-extera di kameno-tubo o kameno-korpo\lingui{DFIRS}
\artiklo{kameo}%
\vorto{kameo} Lapido (onixo, sardonixo, e c.) skultita ad-reliefe, ube onu uzas la kolori diversa di la strati por nuancizar la parti diversa di la laboruro\lingui{DEFIRS}
\artiklo{kamerlingo}%
\vorto{kamerlingo}\fako{religio katol.} Kardinalo di qua la ofico konsistas, en la stati pontifikala, ye administrar la judiciofako e la trezoro-fako, qua prezidas la chambro apostolala, e guvernas dum la vaki da la Santa Sideyo\lingui{DEFIRS}
\artiklo{kamero}%
\vorto{kamero}\senco{1}{Chambro sen-fenestra}\senco{2}{\textit{(fotogr.)} Buxo kun aperturo truo-forma en qua esas inkastrita lenso qua igas la lumo-radii konvergar, quan la lumo-radii emaninta de la objekti extera trairas tale ke li formacas, sur la parieto funda, imajon reduktita de ta objekti}\lingui{DEFR}
\artiklo{kamforo}%
\vorto{kamforo} Substanco konkreta, blanka, diafana, qua saporas bitre e korodive, odoras penetrive, tre volatila, uzata kom kontre-spasmo ed antisepsivo, quan onu extraktas ek la Chinia-lauro\lingui{DEFIRS}
\artiklo{kamilo}%
\vorto{kamilo} Quaza bi-brankardo, sur qua onu transportas brakie kargaji\lingui{S}
\artiklo{kamiono}%
\vorto{kamiono} Veturo basa por la transporto di vari, di paki\lingui{FIS}
\artiklo{kamizo}%
\vorto{kamizo} Sub-vesto, ek telo, kaliko, lano, silko, *nilono, quan onu *weras nemediate sur la pelo\lingui{FIRS}
\artiklo{kamizolo}%
\vorto{kamizolo} Vesto kun maniki, ek telo, lano, e c., qua *weresas maxim-multa-kaze sur la kamizo\lingui{DEFIRS}
\artiklo{kamloto}%
\vorto{kamloto} Stofo grosiera quan onu fabrikis en la Oriento-landi, ek pili de kameli o de kapri\lingui{DEFIRS}
\artiklo{kamo}%
\vorto{kamo}\fako{tekn.} Saliajo kurvo-linea, muntita sur la perimetro di rotaco-axo, qua, renkontrante stango fixigita an la pistono-stango di pistilo, od an la extremajo di pisto-martelego, sublevas la pistilo, la martelego, e lasas lu retrofalar, tale transformante la movo cirklatra kontinua di la transmisarboro aden movo intermitoza, cirklatra o rekta-linea\lingui{DEFIS}
\artiklo{kamomilo}%
\vorto{kamomilo}\fako{bot.} Planto ek la familio « kompozaji », di qua la flori – uzata infuze – esas stomako-kuraciva e kontre-febra\latina{matricaria chamomilla}\lingui{DEFIS}
\artiklo{kampaniar}%
\vorto{kampaniar}\fako{netrans.}\senco{1}{\textit{(en la senco militala)} Movar, agar lo necesa, por sucesigar expediciono}\senco{2}{\textit{(anke metaf.)} Kampanio (tote pacoza) di navigado, di maristo o di navo; kampanii di explorado, di peskado, di laborado, di propago, di reklamo; kampanii « elektala » (politiko); kampanii di jurnali por o kontre ulu od ulo}\lingui{DEFIS}
\artiklo{kampanulo}%
\vorto{kampanulo}\fako{bot.} Planto di qua la flori, blua, violea, pendas quale klosheti\latina{campanula}\lingui{DEFIS}
\artiklo{kampar}%
\vorto{kampar}\fako{netrans.}\senco{1}{\textit{(armeo)} Establisar su, esar establisita plu o min duronte, sur tereno determinita}\senco{2}{Instalar su, esar instalita provizore}\lingui{DEFIRS}
\artiklo{kampesho}%
\vorto{kampesho}\fako{bot.} Arboro dornoza ek Mexikia, Antili, qua donas tintivo reda\latina{haematoxylon campechianum}\lingui{DEFIRS}
\artiklo{kanabino}%
\vorto{kanabino}\fako{zool.} Speco di mikra pasero griza, amansebla, e quan onu povas dresar por ke lu siflez arii\latina{fringilla cannabina}\lingui{L}
\artiklo{kanabo}%
\vorto{kanabo}\fako{bot.} Planto herbatra, di qua la stipo donas filamenti ek qui onu facas filo, telo, e c.\latina{canabia sativa}\lingui{FIS}
\artiklo{kanalio}%
\vorto{kanalio} Persono nehonesta, desestimenda\lingui{DEFIRS}
\artiklo{kanalo}%
\vorto{kanalo}\senco{1}{Quaza rivero artificala, destinita ad igar interkomunikar du baseni, du riveri o rivereti, o du parti ek un rivero o rivereto}\senco{2}{Aquo-voyo naturala od artificala en la enireyo di portuo, o paseyo navigebla qua duktas a ta enireyo}\lingui{DEFIRS}
\artiklo{kanapeo}%
\vorto{kanapeo} Sidilo di qua la dors-apogilo esas tale granda ke, sur lu, plura personi povas sidar ensemble, od un persono sternar su\lingui{DFIRS}
\artiklo{kanario}%
\vorto{kanario}\fako{zool.} Serino de Kanarii\lingui{DEFIRS}
\artiklo{kancelero}%
\vorto{kancelero} Gardero di la stato-siglili\lingui{DEFIRS}
\artiklo{kancero}%
\vorto{kancero}\fako{patol.} Nomo vulgara di la tumori qui rodas la tisui di ula parti di la korpo, di la organi\lingui{EFIS}
\artiklo{kande}%
\vorto{kande} En la epoko dum qua; ye la instanto dum qua…\lingui{FIS}
\artiklo{kandelabro}%
\vorto{kandelabro}\senco{1}{Kandeliero kun plura branchi por la bujii}\senco{2}{Granda suportilo qua portas un o plura lanterni gasala od elektrala por lumizar la voyi publika, la peristili, salonon e c.}\lingui{DEFIRS}
\artiklo{kandelo}%
\vorto{kandelo} Mecho ek kotono, cirkondita da sebo, da vaxo o da irgaltra materio grasa e kombustebla, qua, kande olu esas acendita, konsumesas lente, e donas flamo lumiziva\lingui{EFIRS}
\artiklo{kandida}%
\vorto{kandida} Qua esas sincera, pro ke lua pura anmo, psiko, koncienco havas nulo celenda\lingui{EFIS}
\artiklo{kandidato}%
\vorto{kandidato} La persono qua su propozas pri ofico o plaso qui vakas\lingui{DEFIRS}
\artiklo{kandio}%
\vorto{kandio} Sukro, kristaligita a peci\lingui{DEFIS}
\artiklo{kanelo}%
\vorto{kanelo} Quaza sulko longesala, mi-cilindra, qua alternas kun mikra muluro ronda, o kun aristo paralela\lingui{DEFI}
\artiklo{kanetilio}%
\vorto{kanetilio}\senco{1}{Filo ek oro od arjento, spulita sur longa agulo de fero, ed uzata por la fabriko di la texuri brodita ye oro od arjento}\senco{2}{Filo ek latuno arjentizita, qua, spiral-spulite, cirkum kordo ek tripo o cirkum metala kordo, donas la kordi maxim basa-sono di la violino o di la violoncelo}\senco{3}{Bendo streta, ek kotonajo, di qua la bordi esas garnisita ye latuna fili, uzata da la chapelifisti mulierala por muntar la formo di la ornivi chapelala}\lingui{DFIRS}
\artiklo{kanguruo}%
\vorto{kanguruo}\fako{zool.} Quadripedo ek Australia, ek la klaso « marsupiali », remarkinda pro la volumino di lua kaudo sur qua lu su apogas por saltar, e pro la tre granda longeso di lua membri dopa\latina{halmaturus}\lingui{DEFR}
\artiklo{kanibalo}%
\vorto{kanibalo}\senco{1}{La homo qua su nutras ek karno homala}\senco{2}{\textit{(metaf.)} Persono feroca quale homo qua su nutras ek karno homo-parta}\lingui{DEFIRS}
\artiklo{kanikulo}%
\vorto{kanikulo}\senco{1}{La stelo di Sirius o di la Hundo, qua su levas e su kushas sama-kloke kam Suno, de la 24esma di julio til la 24esma di agosto}\senco{2}{Periodo dum qua la vetero esas tre varma, e quan produktas ta serio de dii}\lingui{EFIS}
\artiklo{kanina}%
\vorto{kanina}\fako{anat.} \textit{(dento)} qua, che la homo, esas situita inter la incizivi e la molari\lingui{EFIS}
\artiklo{kankro}%
\vorto{kankro}\fako{zool.}\senco{1}{Krustaceo kun kin pari de gambi, di qua la gambi avana finas ye pinci, e qua marchas (egale-bone) avance o des-avance}\senco{2}{\textit{(astron.)} Stelaro, reprezentata da kankro, qua okupas la parto di zodiako, quan eniras Suno ye la solstico somerala}\latina{astacus}\lingui{EF}
\artiklo{kankroido}%
\vorto{kankroido}\fako{patol.} Kancero di la pelo e di la mukosi
\artiklo{kano}%
\vorto{kano}\fako{bot.}\senco{1}{Planto di qua la stipo, rekta, havas tuberi de qui spricas folii qui formacas gaino ye sua bazo}\senco{2}{Planto di qua la stipo esas glata, interne kava e plena ye medulo, genero ek la familio « graminei »}\latina{arundo, phragmites}\lingui{EFIS}
\artiklo{kanoniko}%
\vorto{kanoniko}\fako{eklezio katol.}\senco{1}{\textit{(olim)} Kleriko reguliera o sekulara, ek ordeno religiala od ek kirko katedrala o kolegiala}\senco{2}{\textit{(nun)} Kleriko sekulara qua, kun sua egali, deliberas pri la aferi di la kongregaciono}\senco{3}{Kleriko sekulara qua recevis de episkopo sua titulo honorizanta}\lingui{DEFIRS}
\artiklo{kanonizar}%
\vorto{kanonizar}\fako{trans.) (religio katol.}\senco{1}{Enskribar la nomo en la katalogo de la santi; plasizar inter la santi}\senco{2}{Deklarar konforma ye la kanoni ekleziala}\lingui{DEFIS}
\artiklo{kanono}%
\vorto{kanono}\senco{1}{\textit{(teol. katol.)} Eklezio-lego; decido, da la koncilo, pri la fido e pri la diciplino; 2. (meso)-kanono : parto di la mes-oficio, vorti, pregi sakramentala, quan la sacerdoto dicas nelaute, de la prefaco til la « pater »}\senco{2}{\textit{(armo)} Cilindro ek bronzo, giso-fero, stalo, borita ad-tubatre por kontenar la projektili (kuglegi, obusi, mitralio) quin lansas la explozo da kargajo ek pulvero quan onu acendas, muntita sur afusto}\lingui{DEFIRS}
\artiklo{kanoo}%
\vorto{kanoo} Kanoto pirogatra, kun extremaji pintatra, facita ek ligno dina, qua movesas ordinare per pagayi\lingui{DEFIS}
\artiklo{kanoto}%
\vorto{kanoto} Mikra batelo, sen-ferdeka, movata per segli, remili, motoro, ed ordinare uzata sur la mar-aquo\lingui{FIS}
\artiklo{kansono}%
\vorto{kansono} Ensemblo ek mikra versi, maxim-multa-kaze dividita ad kupleti kun refreno, kantata per ario populara\lingui{FIS}
\artiklo{kantabilo}%
\vorto{kantabilo}\fako{muziko} Kantajo ye movimento moderata\lingui{DEFI}
\artiklo{kantar}%
\vorto{kantar}\fako{trans.}\senco{1}{Produktar per la voco muzik-ario}\senco{2}{Celebrar per liriko-versi poeziala}\lingui{DEFIS}
\artiklo{kantarelo}%
\vorto{kantarelo}\fako{bot.} Speco de boleto, fungo manjebla\latina{cantharellus cibarius}\lingui{DEFSL}
\artiklo{kantarido}%
\vorto{kantarido}\fako{zool.} Insekto koleoptera, quan onu sekigas e reduktas a pulvero, uzata farmaciale en la preparaji vezikigiva\latina{cantharis vesicatoria}\lingui{DEFIRS}
\artiklo{kantato}%
\vorto{kantato} Ensemblo ek versi, qua konsistas (generale) ye recitativo e ye strofi, a qua onu povas adjuntar muzik-ario\lingui{DEFIRS}
\artiklo{kantiko}%
\vorto{kantiko}\fako{kristanismo}\senco{1}{Kantajo religiala ek la santa skriburi, por laudar, por dankar, deo}\senco{2}{Kantajo liturgiala}\senco{3}{Kantiko spiritala; kantajo religiala, en linguo vulgara, a qua onu adaptas arii, e quan kantas la pueri katekizata, la adolecantini di la kunfratari}\lingui{DEFIS}
\artiklo{kantileno}%
\vorto{kantileno} Speco di kansono, anciena e lamentala\lingui{DEFIRS}
\artiklo{kantino}%
\vorto{kantino} Taverno di regimento, di kazerno; taverno di fabrikeyo, di karcero, e c.\lingui{DEFIS}
\artiklo{kantono}%
\vorto{kantono}\fako{en Francia}\senco{1}{Dividuro teritoriala di distrikto}\senco{2}{\textit{(en Suisia)} Singla ek la 22 mikra stati, ye qui konsistas la federitaro (o : federuro) Suisa}\lingui{DEFIRS}
\artiklo{kantoro}%
\vorto{kantoro}\fako{kristanismo} La persono qua kantas lor la oficio en kirko, en templo\lingui{DEFIRS}
\artiklo{kanulo}%
\vorto{kanulo} Tubo, plu o min longa, rigida o flexebla, quan onu uzas en ula operaci kirurgiala, o por enduktar (aden la intestini o la vagino) ula liquidi\lingui{DEFIS}
\artiklo{kanvaso}%
\vorto{kanvaso}\senco{1}{Dika telo kruda, de qua onu facas la segli, la visho-tuki koqueyala}\senco{2}{Reto ek kotono-fili duopla, qui interkrucumas ajure, formacante quadrati, ye qua konsistas la fundo sur qua onu facas la agulo-tapeto}\lingui{DEFIRS}
\artiklo{kaolino}%
\vorto{kaolino} Arjilo blanka, tre pura, quan onu renkontras precipue en Chinia, ed ek qua onu facas la porcelano\lingui{DEFIRS}
\artiklo{kaoso}%
\vorto{kaoso}\senco{1}{Konfuzeso di la univers-elementi qui, segun ula teogonii, existis ante la organizeso di mondo}\senco{2}{\textit{(metaf.)} Konfuzeso e desordineso absoluta}\lingui{DEFIRS}
\artiklo{kapabla}%
\vorto{kapabla}\fako{la persono} Qua povas (mentale o korpale) exekutar bone to quon onu konfidas a lu kom agenda o facenda\lingui{DEFIS}
\artiklo{kapaca}%
\vorto{kapaca}\fako{extense di} Tale granda ke lu povas kontenar (ta o ca quanto de ulo)\lingui{DEFIS}
\artiklo{kapear}%
\vorto{kapear}\fako{netrans.} Navigar (dum vento forta) transverse en vento, sub seglaro mikra, la segli bursatrigite\lingui{FIS}
\artiklo{kapelino}%
\vorto{kapelino} Kofio por muliero, qua decensas til an-sur elua shultri\lingui{FIRS}
\artiklo{kapelo}%
\vorto{kapelo}\fako{kristanismo}\senco{1}{Loko sakrigita, ube onu konservas la mantelo, la restaji de santo}\senco{2}{Loko konsakrita a la kulto, en kastelo, kongregaciono, hospico, kolegio, e c.; la personi qui oficias en ta loko}\senco{3}{Parto di kirko qua kontenas altaro partikulara exter la koreyo}\senco{4}{Kirko qua ne havas ankore la titulo « parokiala »}\lingui{DEFIS}
\artiklo{kapero}%
\vorto{kapero}\fako{bot.} Yuna burjoneto de arboreto kultivata sude di Europa, qua, konfitite en vinagro, uzesas kom kondimento\latina{capparia spinosa}\lingui{DEFIRS}
\artiklo{kapilara}%
\vorto{kapilara} Di qua la diametro esas haro-diametra\lingui{DEFIRS}
\artiklo{kapilario}%
\vorto{kapilario}\fako{bot.} Filiko di qua la stipito e la foliaro esas tenua; e di qua la folii uzesas sive infuze, sive sirope\latina{adiantus capillus veneris}
\artiklo{kapistro}%
\vorto{kapistro}\fako{kirurg.} Bandajo per qua onu reduktas la rupturi e la luxacuri che la osti di la mandibulo\lingui{DFIS}
\artiklo{kapitalo}%
\vorto{kapitalo}\fako{financo} La pekunio quan ulu posedas (kontraste kun la revenuo, qua esas acesora). La pekunio, egardata kom produkto-moyeno\lingui{DEFIRS}
\artiklo{kapitano}%
\vorto{kapitano}\senco{1}{La persono qua komandas en regimento-kompanio}\senco{2}{La persono qua komandas sur milito-navo, kom chefo, oficiro, suprega}\lingui{DEFIRS}
\artiklo{kapitelo}%
\vorto{kapitelo}\fako{arkitekt.} Parto supra di kolono, di pilastro, qua esas quaze la krono di la fusto\lingui{DEFIRS}
\artiklo{kapitulacar}%
\vorto{kapitulacar}\fako{netrans.} Kontratar kun la enemiko, pri la kondicioni di la livro (forme di rezisto-ceso) di fortifikajo, fortreso, armeo\lingui{DEFIRS}
\artiklo{kapitulo}%
\vorto{kapitulo}\senco{1}{\textit{(bot.)} Infloresenco sama-forma kam kapo, spiko}\senco{2}{\textit{(liturgio katol.)} Kurta fragmento di la santa skriburi, qua relatas la oficio di la dio, quan onu recitas pos la psalmi, la precipua kloki di la breviario}\lingui{DEF}
\artiklo{kapo}%
\vorto{kapo}\fako{anat.}\senco{1}{Parto supra di la homo-korpo – parto avana di la animalo-korpo – qua kontenas la cerebro e la organi precipua di la sensi}\senco{2}{\textit{(metaf.)} Parto supra : 1. to quo direktas; 2. somito e (precipue) somito kalotatra; 3. \textit{(metaf.)} parto avana : to quo esas situita an-avan la objekto; 4. to quo komencas ulo (kom ex. : la kapo di relvoyo-lineo)}\lingui{ISL}
\artiklo{kaponiero}%
\vorto{kaponiero}\fako{milit-arto} Galerio konstruktita en sika fosato di fortreso, infre di la pafo-lineo, por komunikar de ta a ca fortifikajo, e por impedar la transiro di la fosato\lingui{DEFIRS}
\artiklo{kapono}%
\vorto{kapono} Hanulo yuna quan onu kastris por igar lu divenar plu grasoza\lingui{DEFIRS}
\artiklo{kaporalo}%
\vorto{kaporalo} Armeano di qua la grado esas la minim alta\lingui{DEFIR}
\artiklo{kapoto}%
\vorto{kapoto}\senco{1}{Longa redingoto por la viri, e precipue por la soldati}\senco{2}{Sur ula komerco-navo : bloko qua shirmas la enireyo di la pupo-eskalero; 2. sur la automobili, kovrilo metala qua shirmas la motoro}\lingui{DEFIRS}
\artiklo{kapreolo}%
\vorto{kapreolo}\fako{zool.} Speco di cervo, ma plu mikra, di qua la korni esas kurta, cilindra, ed havas nur un brancho\latina{cervus capreolus}\lingui{IL}
\artiklo{kaprico}%
\vorto{kaprico} Volo subita, variema, nejustifikata, fantazia\lingui{DEFIRS}
\artiklo{kaprifolio}%
\vorto{kaprifolio}\fako{bot.} Arboreto klimera, sarmentoza, di qua la flori odoras\latina{lonicera caprifolium}\lingui{DIR}
\artiklo{kaprikorno}%
\vorto{kaprikorno}\senco{1}{\textit{(zool.)} Insekto koleoptera, kun longa anteni artikoza, kurva quale korni}\senco{2}{Stelaro imaginata kom kaproforma, qua okupas ta parto di zodiako quan eniras Suno lor la solstico vintrala, e quan lu semblas transirar inter la 20esma di decembro e la 20esma di januaro}\latina{cerembyx}\lingui{EFIS}
\artiklo{kaprimulgo}%
\vorto{kaprimulgo}\fako{zool.} Genero di pasero kun beko fendita, e qua, flugante, havas sua beko apertita\latina{caprimulgus europeus}\lingui{IL}
\artiklo{kapriolar}%
\vorto{kapriolar}\fako{netrans.}\senco{1}{Saltar bruske e retrosaltar, lor folo-ludeto}\senco{2}{\textit{(aludante kavalo)} Saltar e, dum esar super sulo, kikar}\lingui{DEFIRS}
\artiklo{kapro}%
\vorto{kapro}\fako{zool.} Mamifero, speco ek la familio « rumineri », kun korni kava e vivaca, mentono kun barbo, e di qua la maskulo odoras multe e desagreable. (hircus)\latina{capra}\lingui{FIS}
\artiklo{kapstano}%
\vorto{kapstano} Vindilo vertikala, an-cirkum qua spulesas kablo\lingui{DERS}
\artiklo{kapsulo}%
\vorto{kapsulo}\senco{1}{\textit{(artilerio)} Alveolo de kupro, qua kontenas materio detoniva, e qua uzesas por komunikar la fairo a la kargajo de pulvero di la perkut-armi}\senco{2}{\textit{(farmacio)} Envelopilo solvebla, sen-sapora, aden qua onu inkluzas ula medikamenti por disimular lia saporo desagreabla}\lingui{DEFIRS}
\artiklo{kaptacar}%
\vorto{kaptacar}\fako{trans.} Obtenar per manovri ruzoza (donaco, legaco, e c.)\lingui{FS}
\artiklo{kaptar}%
\vorto{kaptar}\fako{trans.}\senco{1}{Sucesar pri sizar ulu, ulo, qua probas eskapo}\senco{2}{Sizar (la fluo di la liquido de fonto, por evitar la infiltri, la alteri, e c.)}\lingui{EFIS}
\artiklo{kapuchino}%
\vorto{kapuchino}\fako{bot.} Planto klimera, di qua la floro esas oranjo-flava, e qua havas, ye sua extremaji suba, apendico kapuco-forma\latina{tropaeolum maius}
\artiklo{kapucino}%
\vorto{kapucino}\fako{kristanismo} Reguliero ek un de la ordeni franciskana\lingui{DEFIRS}
\artiklo{kapuco}%
\vorto{kapuco}\senco{1}{Parto di mantelo, di froko, e c., forme di boneto tre ampla, qua esas charniratre adduktebla adsur la dorso}\senco{2}{Mantelo kapucoza, quan *weris la viri e la mulieri}\lingui{DFIRS}
\artiklo{kara}%
\vorto{kara}\senco{1}{\textit{(pri persono)} Qua inspiras multa tenereso (ad ulu). (anke familiare, ma kun senco tre febligita)}\senco{2}{\textit{(pri kozo)} Tre prizata pro olua charmo, pro olua importo}\lingui{FIS}
\artiklo{karabino}%
\vorto{karabino} Fusilo lejera, di qua la tubo havas (ordinare) surfaco interna qua esas striizita helicatre, ed extere surfaco taliita facioze\lingui{DEFIRS}
\artiklo{karabo}%
\vorto{karabo}\fako{zool.} Insekto karn-avida, genero ek la serio « koleopteri », e di qua ula speci lansas liquido fetida\latina{carabus}\lingui{EFIS}
\artiklo{karafo}%
\vorto{karafo} Vazo portebla, ordinare sen ansi, kun la ventro globoza, la kolo streta, ek kristalo od ek vitro sen-kolora, di qua la stopilo esas ek la sama materio, e qua generale destinesis a kontenor od a varsor la repast-aquo\lingui{DEFIS}
\artiklo{karakalo}%
\vorto{karakalo}\fako{zool.} Sorto di kato sovaja\latina{lynx caracal}\lingui{DEFIRS}
\artiklo{karako}%
\vorto{karako} Quaza kazako apertita, quan *weris olim la rurani, la soldati
\artiklo{karakolar}%
\vorto{karakolar}\fako{netrans.) (aludante kavalo}\senco{1}{Movar su turne, addextre ed adsinistre}\senco{2}{Irar kapricoze, per salti, per saluto-flexi}\lingui{DEFIS}
\artiklo{karakteristiko}%
\vorto{karakteristiko}\fako{matem.} La nombro integra di logaritmo\lingui{DEFIS}
\artiklo{karaktero}%
\vorto{karaktero}\senco{1}{Traito partikulara, eso-maniero propra, qua distingas kozo de altra kozo}\senco{2}{La eso-maniero morala, psikala, spiritala (kontraste kun : la kompreno-fakultato)}\lingui{DEFIRS}
\artiklo{karambolar}%
\vorto{karambolar}\fako{netrans.) (biliardo-ludo} Agar la stroko qua konsistas ye shokar du buli per la sua\lingui{DFIRS}
\artiklo{karamelo}%
\vorto{karamelo} Sukro fuzita e brunigita per la efiko da fairo\lingui{DEFRS}
\artiklo{karasino}%
\vorto{karasino}\fako{zool.} Fisho simila kam karpo\latina{carassius vulgaris}
\artiklo{karato}%
\vorto{karato} Pondero-unajo, uzata ankore da la juvelisti, por la diamanti, la lapidi, la perli, e qua equivalas cirkume 20 centigrami\lingui{DEFIRS}
\artiklo{karavano}%
\vorto{karavano}\senco{1}{En Avan-Azia, en Arabia, e c., trupo de komercisti, de voyajeri, de pilgrimanti, asemblita por trans-irar plu sekure la dezerti}\senco{2}{\textit{(metaf.)} Trupo de personi qui iras ensemble}\lingui{DEFIRS}
\artiklo{karavanserayo}%
\vorto{karavanserayo} En la Oriento-landi : granda edifico en qua la voyajeri povas shirmar e su ipsa e lia charjo-bestii\lingui{DEFIRS}
\artiklo{karavelo}%
\vorto{karavelo}\senco{1}{\textit{(olim)} Che la Turki, grosa milito-navo}\senco{2}{Mikra navo kun segli Latina, uzita precipue da la Portugalani}\senco{3}{Navo uzata por peskar la haringi}\lingui{F}
\artiklo{karbo}%
\vorto{karbo}\fako{kemio} Korpo simpla, metaloida, qua existas pluraforme : karbono, grafito, diamanto, e c.\simbolo{C}\lingui{DEFIS}
\artiklo{karbonario}%
\vorto{karbonario}\fako{zool.} Fisho ek la genero « gado »\latina{gadus carbonarius}
\artiklo{karbonaro}%
\vorto{karbonaro} Membro di societo revolucionista sekreta, fondita en Italia, komence di la 19esma yarcento\lingui{FI}
\artiklo{karbono}%
\vorto{karbono} Materio nigra, en qua dominacas la korpo simpla « karbo », e quan onu extraktas (ordinare) de subsulo per minar\lingui{EFIS}
\artiklo{karborundo}%
\vorto{karborundo}\fako{kemio} Siliko-karburo, preparita per submisar, a la kaloro de elektroforno, mixuro ek karbono e silicio : uzata por la fabriko di grindo-roti, e c.
\artiklo{karbunklo}%
\vorto{karbunklo}\senco{1}{\textit{(patol.)} Morbo qua desorganizas e nigrigas la tisui atakita}\senco{2}{Varietato di granato reda, di qua la brilo esas tre granda}\lingui{DEFIRS}
\artiklo{karburar}%
\vorto{karburar}\fako{netrans.) (tekn.} Preparar la mixuro gasa, explozala quan onu destinas a la alimento di la explozo-motori\lingui{EF}
\artiklo{karburatoro}%
\vorto{karburatoro}\fako{tekn.} Aparato qua produktas ta mixuro ek aero e vaporo de hidrokarburo o kombusto-gaso, quan onu uzas sive por lumizo, sive por alimentar explozo-motori\lingui{EF}
\artiklo{karcero}%
\vorto{karcero} Loko ube la arestito retenesas, malgre lua volo e, tale, privacesas de lua libereso\lingui{DEFIRS}
\artiklo{karcinomo}%
\vorto{karcinomo}\fako{patol.} Un ek la formi di la morbo « kancero »\lingui{DEFIRS}
\artiklo{kardamino}%
\vorto{kardamino}\fako{bot.} Planto krucifera\latina{cardamina}\lingui{DEFIS}
\artiklo{kardamomo}%
\vorto{kardamomo}\fako{bot.} Planto di India, genero « amomo »\latina{elitteria cardamomum}\lingui{DEFIRS}
\artiklo{kardano}%
\vorto{kardano}\fako{tekn.} « Artiko » di parti de mekanismi (stangi, e c.) qua posibligas movi irga-direcione\lingui{DEFIRS}
\artiklo{kardar}%
\vorto{kardar}\fako{trans.} Pektar, desintrikar (precipue lano e krino) per planketo provizita ye klovi ek latuno-filo\lingui{DEFIRS}
\artiklo{kardelo}%
\vorto{kardelo}\fako{zool.} Mikra ucelo kantera, de la serio « paseri », di qua la maskulo havas kapo reda e la ali markizita flave e brune\latina{garduelis elegans}\lingui{FIL}
\artiklo{kardinala}%
\vorto{kardinala}\senco{1}{\textit{(matem.)} Dicesas pri la nombri qui indikas quanta unaji sama-natura kontenesas en quanto determinita, e qui uzesas por formacar la cetera termini numerala}\senco{2}{Dicesas pri la quar horizonto-punti : nordo, sudo, esto, westo, segun qui onu determinas la situeso di loki cetera}\senco{3}{Dicesas pri la quar vertui precipua (prudenteso, yusteso, moderemeso, e psiko-forteso)}\lingui{DEFIS}
\artiklo{kardinalo}%
\vorto{kardinalo} Singla ek la sepa-dek prelati (episkopi, sacerdoti, diakoni) ye qui konsistas la Sakra Kolegio di la Eklezio katolika, qui votas en konklavo por elektar nova papo, o por elektesar, e qui havas kostumo reda\lingui{DEFIRS}
\artiklo{kardioido}%
\vorto{kardioido}\fako{geom.} Konkoido di la cirklo
\artiklo{kardono}%
\vorto{kardono}\fako{bot.} Planto ek la familio « kompozaji », di qua la folii e la kapituli esas dornoza\latina{carduus}\lingui{EFIS}
\artiklo{karduno}%
\vorto{karduno}\fako{bot.} Speco di artichoko, di qua onu manjas la kosto mediana di la folii, pos blankigir olu per kovrar lu ye palio\latina{cynara cardunculus}\lingui{DEFIS}
\artiklo{karear}%
\vorto{karear}\fako{trans.} Ne judikar su kom privacata de ulo, pro ke onu ne tre bezonas nek tre deziras olu\lingui{EFL}
\artiklo{karelo}%
\vorto{karelo}\senco{1}{\textit{(ludo-karto)} Quadrato reda, uzata kom la distingivo di grupo de karti, en qua la aso reprezentesas per quadrato meze di la karto}\senco{2}{Objekto di qua la surfaco esas quadrata o rektangula}\lingui{FL}
\artiklo{kareno}%
\vorto{kareno} La tota surfaco di navo, qua esas imersita\lingui{FIS}
\artiklo{karesmo}%
\vorto{karesmo} En la religio katolika, periodo ek 46 dii, de la cindro-merkurdio til la Pasko-sundio, dum qua la fideli devas abstinencar e fastar (ecepte dum la sundii)\lingui{FIS}
\artiklo{karezar}%
\vorto{karezar}\fako{trans.} Tushar lejere per sua manuo, sua labii, por signifikar sua afeciono, tenereso, amo, amoro – ed, en ula kazi, por produktar efekti erotika\lingui{EFIS}
\artiklo{kargar}%
\vorto{kargar}\fako{trans.} Igar (ulu, ulo) subisar la pezo da objekto qua esas portenda ad altra loko, transportenda\lingui{EFIS}
\artiklo{kariar}%
\vorto{kariar}\fako{netrans.) (patol.}\senco{1}{\textit{(aludante osto)} Domajesar (da morbo) forme di molesko di la tisuo osta kun pusifo}\senco{2}{\textit{(aludante la denti)} Domajesar (da morbo) forme di molesko lenta e gradoza, pri emalio ed ivoro}\lingui{DEFIS}
\artiklo{kariatido}%
\vorto{kariatido} Statuo de muliero, drapirita, quan onu reprezentas kom stacanta, e qua subtenas entablamento, kornico, vice kolono, pilastro\lingui{DEFIS}
\artiklo{karibuo}%
\vorto{karibuo}\fako{zool.} Rentiro di Kanada\latina{Rangifer caribus}\lingui{DEF}
\artiklo{karico}%
\vorto{karico}\fako{bot.} Planto ek la familio « ciperaci »\latina{carex}\lingui{EFIS}
\artiklo{kariero}%
\vorto{kariero} Profesiono ofico-hierarkioza\lingui{DEFIRS}
\artiklo{karikatar}%
\vorto{karikatar}\fako{trans.} Reprezentar persono, kozi, ma exajerante, groteske, por igar li rido-mokinda\lingui{DEFIS}
\artiklo{kariko}%
\vorto{kariko} Redingoto ampla, kun pelerino dispozita etajatre\lingui{DF}
\artiklo{kariocinezo}%
\vorto{kariocinezo}\fako{biol.} Divido di la celulo, lor qua la divido di la nuklei preiras olta di la celulo-korpo
\artiklo{kariofilo}%
\vorto{kariofilo}\fako{bot.} Burjono di la floro di arboro. di Moluko o di Antili, ek la familio « mirtacei », uzata kom spico, piklo\latina{caryophillus aromaticus}\lingui{I}
\artiklo{kariolizo}%
\vorto{kariolizo}\fako{bot.} Dissolvo di la celulo-nukleo
\artiklo{karitato}%
\vorto{karitato}\senco{1}{Amo kompatoza di la kunhomi}\senco{2}{\textit{(kristanismo)} Amo di la kunhomo, por deo}\lingui{EFIS}
\artiklo{karlino}%
\vorto{karlino}\fako{bot.} Planto di qua la radiko uzesis kom sudorifigivo\latina{carlina}\lingui{EFISL}
\artiklo{karmelito}%
\vorto{karmelito}\fako{religio katol.} Reguliero ek la ordeno « karmel » (ordeno fondita en 1451, da Jean Soreth)
\artiklo{karmezino}%
\vorto{karmezino} Koloro obskure-reda briloze\lingui{DEFIS}
\artiklo{karmino}%
\vorto{karmino} Koloro bele-reda, ma min brilo-reda e plu brune-reda kam cinabrea\lingui{DEFIRS}
\artiklo{karnaciono}%
\vorto{karnaciono} Koloro di la parti extera di la homo-korpo\lingui{DEFIS}
\artiklo{karnavalo}%
\vorto{karnavalo}\fako{religio katol.} La lasta karno-dii di la periodo de la Reji-dio til la Cindro-merkurdio, e dum qui eventas la lasta amuzi e distrakti di ta periodo; travestii, bali de maskiti, e c.\lingui{DEFIRS}
\artiklo{karnivoro}%
\vorto{karnivoro} Qua su nutras ye karno\lingui{DEFIS}
\artiklo{karno}%
\vorto{karno}\fako{anat.} La muskuli di la korpo di la homo o di animalo (kontraste kun la osti e la pelo)\lingui{DEFIS}
\artiklo{karonado}%
\vorto{karonado} Pafilo, ordinare ek fero, plu kurta kam kanono, qua uzesis, precipue sur la navi, e qua (quale la bombardilo) havas alveolo por la pulvero\lingui{DEFIS}
\artiklo{karoso}%
\vorto{karoso} Luxo-veturo, di qua la korpo su apogas sur suspenso-risorti, tektoza, kun quar roti\lingui{DFIS}
\artiklo{karotido}%
\vorto{karotido}\fako{anat.} Singla ek la du grosa arteri qui portas la sango a la kapo\lingui{DEFIS}
\artiklo{karotino}%
\vorto{karotino}\fako{kemio} Hidro-karburo, fuzebla ye 118° C., ye qua konsistas la materio flava di la karoti, e qua absorbas la oxo di aero
\artiklo{karoto}%
\vorto{karoto}\fako{bot.} Planto ek la familio « umbeliferi », di qua la radiko (pulpoza) esas manjebla\latina{daucus carota}\lingui{DEFI}
\artiklo{karpentar}%
\vorto{karpentar}\fako{trans.} Facar la asemblajo ek ligno-pecegi, ek trabi, qui uzesas en domo-konstrukturo por subtenar lua diversa parti\lingui{EFIS}
\artiklo{karpino}%
\vorto{karpino}\fako{bot.} Granda arboro qua kreskas en la foresti, e di qua la ligno, harda, kompakta e blanka, demandesas por la facaji da la tornisti, karpentisti, veturifisti, e c.\latina{carpinus betulus}\lingui{ISL}
\artiklo{karpo}%
\vorto{karpo}\senco{1}{\textit{(zool.)} Fisho qua vivas en la aquo sen-sala, ek la genero « ciprini »}\senco{2}{\textit{(anat.)} Parto di la membro, supra od avana, inter la manuo e la avan-brakio}\latina{cyprinus carpio}\lingui{DEF}
\artiklo{karpologio}%
\vorto{karpologio}\fako{bot.} Parto di botaniko, qua havas kom studiajo, la frukti
\artiklo{kartamino}%
\vorto{kartamino}\fako{kemio} Koloro-materio di la flori, genero « kartamo »\lingui{DE}
\artiklo{kartamo}%
\vorto{kartamo}\fako{bot.} Planto ek la familio « kompozaji », di qua la flori esas bele safrano-reda, uzata por tintar\latina{carthamus tinctorius}\lingui{FISL}
\artiklo{kartavar}%
\vorto{kartavar}\fako{netrans.} Pronuncar la litero « r » per la guturo (t.e. kelkete quale la « ch » Germana, la « j » Hispana)\lingui{R}
\artiklo{kartelo}%
\vorto{kartelo} Karto, papero-folio, sur qua onu sendis defio pri duelo\lingui{DEFS}
\artiklo{kartero}%
\vorto{kartero}\fako{tekn.} Envelopilo rigida qua protektas (kontre la enveno de korpi stranjera) la mashin-organi qui funcionas rapide\lingui{FIRS}
\artiklo{karteziana}%
\vorto{karteziana} Qua relatas Descartes (pr. Dekart') ed ilua doktrino\latina{Cartesius}\lingui{DEFIS}
\artiklo{kartilago}%
\vorto{kartilago}\fako{anat.} Tisuo animalala, flexebla, plastika, di qua la hardeso-grado fluktuas inter olta di la osti ed olta di la ligamenti\lingui{EFIS}
\artiklo{karto}%
\vorto{karto} Quadrato o rektangulo ek sorto di papero rezistiva, ma flexebla, facita ek plura papero-folii interglutinita\lingui{DEFIRS}
\artiklo{kartocho}%
\vorto{kartocho} Tubo kartona, cilindra, qua kontenas la pulvero di fusilo, di pistolo\lingui{DEFIRS}
\artiklo{*kartoflo}%
\vorto{*kartoflo}\fako{bot.} Solaneo tuberoza, di qua la nura tuberkulo esas manjebla (kom legumo), *reproduktebla per la plantaco di tuberkuli – prizentita en Francia da Parmentier. – \textit{(Nomo vulgara : terpomi, en Francia; ter-ovi, en Chinia)}
\artiklo{kartografio}%
\vorto{kartografio} La arto facar la mapi (geografiala, astronomiala, e c.)\lingui{DEFS}
\artiklo{kartomanciar}%
\vorto{kartomanciar}\fako{netrans.} Explorar la hazardo-kombinaji ek ludo-karti, por divinar la eventaji futura\lingui{DEFIS}
\artiklo{kartono}%
\vorto{kartono} Pasto ek papero, shifoni, e c., quan onu hardigis e foliigis; od asemblajo ek folii o karti interglutinita\lingui{DEFIRS}
\artiklo{kartusho}%
\vorto{kartusho} Kadro skultita, grabita, forme di karto di qua la bordumo, la angulo, aspektas quale volvajo, destinita por recevor enskriburo, devizo, moto, la monogramo di ulu, e quan onu plasizas sur edifico, infre ed en la angulo (sinistre o dextre), di tabelo, di pikturo, di graburo, di mapo. – Lineo-kadro elipsoida qua, en hieroglifo-texto, cirkondas la nomi di la deaji, di la reji, di la dinastii\lingui{DEFIRS}
\artiklo{kartuzio}%
\vorto{kartuzio}\fako{religio katol.} Monakeyo di la regulieri di la ordeno di « Santa Bruno », konstruktita (sen ecepto-kazi) en loko dezerta\lingui{DEFIS}
\artiklo{karubo}%
\vorto{karubo}\fako{bot.} La shelo pulpa di arboro, verda, ek la familio « leguminosi », e di qua la ligno esas tre harda\latina{ceratonia siliqua}\lingui{DEFIS}
\artiklo{karunklo}%
\vorto{karunklo}\senco{1}{\textit{(anat.)} Mikra peco karna (en la organismo); apendico karna an la fronto, an la pektoro di ula uceli (dindi, kazuari, e c.)}\senco{2}{\textit{(bot.)} Influro di la *funikulo qua kovras parto di la grano}\lingui{DEFIS}
\artiklo{karuselo}%
\vorto{karuselo} Quaza turniro, distraktivo ek konkursi, kur-konkursi, por recevor premie ringi fingrala, e c., agita da kavalokavalkanti dispozita quadrile\lingui{DEFIS}
\artiklo{kasacar}%
\vorto{kasacar}\fako{trans.) (yuro-cienco} Nihiligar yurale verdikto, akto, proceduro. – En Francia, la kasaco-korto esas la tribunalo suprega, qua (unika kom tala) darfas kasacar, nihiligar last-instance verdikto pro formo-defekto o lego-desegardo\lingui{DEFIRS}
\artiklo{kashaloto}%
\vorto{kashaloto}\fako{zool.} Mamifero ek la familio « cetacei », qua diferas de la baleno per to ke, vice barti, lu havas denti an la mandibulo ed, an la maxilo, alveoli en qui esas la denti\latina{physeter macrocephalus}\lingui{DEFIRS}
\artiklo{kashmiro}%
\vorto{kashmiro} Texajo tre delikata, facita ek la lano di la mutoni. ed ek la lanugo di la pektoro di la kaprini di Tibet\lingui{DEFIRS}
\artiklo{kasio}%
\vorto{kasio}\fako{bot.} Shelo (di qua la pulpo nigra uzesas kom laxigivo) di la arboro L. \textit{cassia}\latina{cassia}\lingui{DEFIS}
\artiklo{kasiso}%
\vorto{kasiso}\fako{bot.} Arboreto, analoga a la ribiero, di qua la frukti esas nigra, aromatoza e dispozita grape\latina{ribea nigrum}\lingui{FIS}
\artiklo{kaskado}%
\vorto{kaskado} Falo da mikra rivereto, de granda alteso-punto. – \textit{(metaf.)} Serio de agi qui eventas quaze salte\lingui{DEFIRS}
\artiklo{kasko}%
\vorto{kasko} Kapo-vesto (precipue : militistala) ek ledro, metalo, qua kovras e protektas la kapo\lingui{EFIRS}
\artiklo{kaso}%
\vorto{kaso}\senco{1}{Chambro ube eventas la inkasi e desinkasi da firmo komercala o da administrantaro}\senco{2}{Ensemblo ek la paperi ed akti financo-valoroza, ed ek la pekunio, quin firmo od administrantaro cirkuligas}\senco{3}{Establisuro financala qua recevas pekunio kom depozeyo, jeras lu ed esforcas augmentar lua quanto per la adjunto di revenui ek li}\lingui{DEFIR}
\artiklo{kasono}%
\vorto{kasono} Aparato quan onu uzas por konstruktar sub la surfaco di maro, di fluvio, e c., quaza pontono quan onu povas elevar od abasar, segun la naturo di la laborendi\lingui{DEFI}
\artiklo{kasqueto}%
\vorto{kasqueto} Kapo-vesto virala, kun viziero\lingui{F}
\artiklo{kasrolo}%
\vorto{kasrolo} Implemento koqueyala ek metalo, vazo cilindra kun mancho e fundo plata, en qua onu koquas nutrivi diversa\lingui{DEFIRS}
\artiklo{kastano}%
\vorto{kastano}\fako{bot.} Frukto di qua la shelo, dornoza e ledratra, kontenas un o plura « mandeli » farinoza qui donas nutrivo tre forta – quan produktas granda arboro ek la familio « amentacei », qua kreskas en la regioni tepida di Europa, Azia ed Amerika\latina{castanea sativa}\lingui{DFIRSL}
\artiklo{kastanyeto}%
\vorto{kastanyeto} Mikra peco ek buso-ligno od ek ivoro, kavigita skaliatre, juntita per kordeto an altra peco identa kontre qua onu igas lu frapar, por indikar forte la ritmo di dans-ario o kant-ario\lingui{DEFIS}
\artiklo{kastelo}%
\vorto{kastelo}\senco{1}{\textit{(en la mez-epoko)} Lojeyo siniorala, defensata per alta muregi, turmi, fosati, fortreso}\senco{2}{Rezideyo di rejo, di princo}\senco{3}{Granda distrakto-domo}\lingui{DEFIRS}
\artiklo{kasto}%
\vorto{kasto} En Egiptia antiqua ed en India : klaso qua esas un ek la dividuri hierarkiala di la socio. – \textit{(metaf.)} Socio-klaso, egardata kom havanta exkluzo-mento pri la personi di la klasi cetera\lingui{DEFIRS}
\artiklo{kastono}%
\vorto{kastono}\senco{1}{Kapo di fingro-ringo, parto salianta en qua esas inkastrita lapido}\senco{2}{Parto salianta, ek la sama metalo kam la fingro-ringo, sur qua onu grabas monogramo}\senco{3}{Singla ek la lapidi, ek la diamanti inkastrita, ye qui konsistas koliaro}\lingui{FI}
\artiklo{kastoreo}%
\vorto{kastoreo}\fako{farmacio} Substanco sekrecata da glandi qui esas sub la kaudo di la kastoro, e quan onu uzas kom kontre-spasmo
\artiklo{kastoro}%
\vorto{kastoro}\fako{zool.} Mamifero rodera, qua habitas somere tertrui quin lu exkavas rive di la fluvii, e vintro kabanin, quin lu konstruktas litoro di la fluvii o meze di aquo, e quin lu protektas per quaza digi\latina{castor faber}\lingui{EFIS}
\artiklo{kastrar}%
\vorto{kastrar}\fako{trans.} Igar sexuo-impotenta, per mutilar ed atrofiar la organi genitala\lingui{DEFIRS}
\artiklo{katafalko}%
\vorto{katafalko} Estrado dekorala, erektita meze di kirko, en qua la sarko esos enklozita dum la funero e simulita dum la oficii memorigala\lingui{DEFIRS}
\artiklo{katakaustika}%
\vorto{katakaustika}\fako{fiziko e geom.} Kurvo qua envelopas lumo-radii reflektita da spegulo konkava – la radii venante de un lumo-punto
\artiklo{kataklismo}%
\vorto{kataklismo} Subverso di la surfaco di la Terglobo, per inundo, sismo (tertremo), e c.\lingui{DEFIS}
\artiklo{katakombo}%
\vorto{katakombo} Chambrego subsula qua uzesis kom sepulteyo, osteyo; exkavajo adube onu asemblis osti\lingui{DEFIRS}
\artiklo{katakrezo}%
\vorto{katakrezo} Figuro retorikala, per qua la vorto qua indikas propre ta o ca objekto, uzesas por indikar altra objekto qua esas sur parto analoga ye la unesma\lingui{DEFIS}
\artiklo{katalazo}%
\vorto{katalazo}\fako{biol.} Fermento solvebla, deskovrita en la tabako-folii, e qua efikas reduktive
\artiklo{katalepsio}%
\vorto{katalepsio}\fako{patol.} Supreseso semblanta di vivo per la suspenseso di la sentiveso extera e di la moviveso volata kun rigidesko kadavrala\lingui{DEFIRS}
\artiklo{katalizar}%
\vorto{katalizar}\fako{trans.) (kemio} Efikigar ta korpi qui, per lia nura prezenteso (e sen esar modifikata) eventigas en altra korpi modifikesi\lingui{DEF}
\artiklo{katalogo}%
\vorto{katalogo} Listo, indexatra, ek la peci di kolektajo (libri, pikturi, graburi, medalii, e c.)\lingui{DEFIRS}
\artiklo{katalpo}%
\vorto{katalpo}\fako{bot.} Arboro ek la familio « bignonacei », di qua la folii esas larja ed ovala, la flori blanka e reda-puntizita, e quan onu kultivas en Europa kom plezuro-planto\latina{catalpa}\lingui{DEFIS}
\artiklo{kataplasmo}%
\vorto{kataplasmo} Topiko ek substanco moligiva, forme di pado kompakta\lingui{DEFIRS}
\artiklo{katapulto}%
\vorto{katapulto}\fako{che la antiqui} Milito-mashino, en qua trabo, per efikar quale risorto, lansis projektili tre pezoza\lingui{DEFIRS}
\artiklo{katarakto}%
\vorto{katarakto}\senco{1}{\textit{(geol.)} Sorto de aquo-fali qui interruptas la fluo di fluvio}\senco{2}{\textit{(oftalm.)} Opakeso di la kristalino, qua interceptas la lumo-radii}\lingui{DEFIRS}
\artiklo{kataro}%
\vorto{kataro}\fako{patol.} Inflameso di mukozo, kun sekreco\lingui{DEFIRS}
\artiklo{katastrofo}%
\vorto{katastrofo}\senco{1}{Inversesko bruska di fortuno}\senco{2}{Eventajo decidigiva qua solvas la problemo di poemo dramata, epika}\lingui{DEFIRS}
\artiklo{katedralo}%
\vorto{katedralo} Granda kirko di la arkitekturo kristana mezepokala\lingui{DEFIRS}
\artiklo{katedro}%
\vorto{katedro} Sidilo lokizita alte, e de qua (en kirko, templo, doceyo) onu su turnas docoze a sua askoltanti\lingui{DFIRS}
\artiklo{kategorika}%
\vorto{kategorika}\senco{1}{\textit{(segun la doktrino da Kant)} Lo absoluta e ne-kondicionata}\lingui{DEFIRS}\subvorto{}{imperativo kategorika}{Ago-regulo qua su impozas a la volado kom impero universala ed absoluta}{2}{}
\artiklo{kategorio}%
\vorto{kategorio}\senco{1}{\textit{(filoz.)} Singla ek la atributi generala di la ento, segun Aristoteles}\senco{2}{Singla ek la koncepti aprioria di la kompreno homala, segun qui lu konceptas necese la objekti di la experienco}\lingui{DEFIRS}
\artiklo{katekismo}%
\vorto{katekismo}\fako{religio kristana} Libro en qua esas rezumita la punti precipua di la docado religiala\lingui{DEFIRS}
\artiklo{katekizar}%
\vorto{katekizar}\fako{trans.} Instruktar voce (ulu) aden la religio kristana (precipue por preparar la infanti por la komunio unesma)\lingui{DEFIS}
\artiklo{katekumeno}%
\vorto{katekumeno}\fako{religio kristana} La homo qua recevas la doco religiala por preparar su a bapteso\lingui{DEFIRS}
\artiklo{kateno}%
\vorto{kateno}\senco{1}{Ligilo qua konsistas ye serio de ringi ek metalo, interplektita}\senco{2}{Longajo ek elementi interplektita, qui esas ica dop ita, od eventas ica pos ita, seninterrupto}\lingui{EIS}
\artiklo{katetero}%
\vorto{katetero}\fako{kirurg.} Sondilo metala, ordinare kurva, kun ranuro qua duktas la instrumento destinita a la exploro ed a la incizo di la veziko\lingui{DEFS}
\artiklo{kateto}%
\vorto{kateto}\fako{olima geometrio} Rekto, perpendikla ye altra\lingui{DEFS}
\artiklo{katguto}%
\vorto{katguto}\fako{kirurg.} Sorto di ligilo, facita ek intestini de mutono o de kato, uzata (aseptigite) por la suto di la plagi\lingui{DEFIRS}
\artiklo{kationo}%
\vorto{kationo}\fako{elektro} Iono qua iras a la katodo lor elektrolizo di solvuro\lingui{DEFIRS}
\artiklo{katisar}%
\vorto{katisar}\fako{trans.} Presar (drapo, lanaji) ed apretizar li por igar li ferma, briloza\lingui{DF}
\artiklo{kato}%
\vorto{kato}\fako{zool.} Animalo domestika, ek la genero « karnavidi », digitigrada, karnivora, kun unglegi retrotirebla (leono, tigro, pantero, linco, jaguaro)\latina{felis domestica}\lingui{DEFIS}
\artiklo{katodo}%
\vorto{katodo}\fako{elektro} Korpo konduktiva, uzata kom elektrodo en la parto liquida o gasa di elektro-cirkuito, e muntita an la polo negativa di la elektro-fluo-fonto\lingui{DEFIRS}
\artiklo{katolika}%
\vorto{katolika}\fako{religio kristana} Qua apartenas a la eklezio di Roma\lingui{DEFIRS}
\artiklo{katoptriko}%
\vorto{katoptriko} Parto di optiko, qua traktas la reflekto di la lumo-radii\lingui{DEFIRS}
\artiklo{kauchuko}%
\vorto{kauchuko} Rezino, suko koagulita, ek *heveo, e c. ed uzata da la industrii\latina{hevea}\lingui{DEFIRS}
\artiklo{kauciono}%
\vorto{kauciono} Pekunio-quanto quan onu depozas, o quan onu su engajas pagor, garantie di su-engajo agita da altru o da su ipsa\lingui{DEFIRS}
\artiklo{kaudo}%
\vorto{kaudo} Prolonguro plu o min oblonga, qua finas dope la torso di la vertebrozi. – (anke metaf.)\lingui{EFIS}
\artiklo{kaulo}%
\vorto{kaulo}\fako{bot.} Planto supala ek la familio « kruciferi »\latina{brassica oleracea}\lingui{DL}
\artiklo{kaurio}%
\vorto{kaurio} Konko uzata kom *koino en Bengal ed en la regioni centrala di Afrika\latina{cypraea}\lingui{FSL}
\artiklo{kaustika}%
\vorto{kaustika} Qua desorganizas, korodas la tisui di la animali e di la vejetanti\lingui{DEFIS}
\artiklo{kaustiko}%
\vorto{kaustiko}\fako{geom.} Kurvo qua envelopas lumo-radii refraktita da spegulo konkava, ta radii venante de un lumo-punto\lingui{DEFIS}
\artiklo{kautero}%
\vorto{kautero} Korpo varmega o kemiajo, quan onu uzas medicinale por desorganizar tisuo\lingui{DEFIS}
\artiklo{kauzo}%
\vorto{kauzo}\senco{1}{\textit{(filoz.)} To per quo kozo esas o divenas to quon lu esas}\senco{2}{To per quo ulo eventas; to pro quo onu agas o facas ulo}\lingui{DEFIS}
\artiklo{kava}%
\vorto{kava} Qua havas en su vakuo plu o min profunda
\artiklo{kavalkar}%
\vorto{kavalkar}\fako{netrans.} Sidar quale viro sur kavalo, sur biciklo\lingui{EFIRS}
\artiklo{kavalo}%
\vorto{kavalo}\fako{zool.} Animalo domestika, mamifera, ek la familio « solipedi », qua uzesas kom kavalkajo, kom tirbestio o kom charjo-bestio\latina{Caballus}\lingui{DEFIRS}
\artiklo{kavalrio}%
\vorto{kavalrio} La parto di armeo, qua konsistas ye viri sur kavali\lingui{DEFIRS}
\artiklo{kavatino}%
\vorto{kavatino}\fako{muziko} Kant-ario (« moderato ») di la operi italiana\lingui{DEFIRS}
\artiklo{kaverno}%
\vorto{kaverno} Kavajo naturala, sub sulo o sub roko, ed apta kom refujeyo\lingui{EFIRS}
\artiklo{kaviaro}%
\vorto{kaviaro} Disho acesora Rusa, ek sturg-ovi, forte presita e marinigita\lingui{DEFIRS}
\artiklo{kayero}%
\vorto{kayero}\senco{1}{Asemblajo ek plura papero-folii, generale intersutita unlatere}\senco{2}{Asemblajo ek la folii ye qui konsistas fragmento di verko skribala quan onu editas partope}\lingui{F}
\artiklo{kayo}%
\vorto{kayo}\senco{1}{Ter-amasajo, subtenata da muro ek quadro, e protektata da parapeto, facita alonge rivero por restriktar lu dum acenso, en la parto di lua fluo qua trairas urbo}\senco{2}{Rivo di portuo, ube onu kargas o deskargas la vari}\senco{3}{\textit{(metaf.)} kayo por la relvoyi}\lingui{DEF}
\artiklo{kazako}%
\vorto{kazako}\senco{1}{Mulier-vesto, surtuto di qua la maniki esas plu o min fitanta}\senco{2}{\textit{(kun maniki)} Mantelo quan olim *weris la musketieri Franca}\lingui{DFIRS}
\artiklo{kazeino}%
\vorto{kazeino} Substanco organika koagulebla, quan kontenas lakto\lingui{DEFIRS}
\artiklo{kazemato}%
\vorto{kazemato} Refujeyo subsula, vultoza, bomb-espruva, por shirmar la armeani, la municioni\lingui{DEFIRS}
\artiklo{kazeo}%
\vorto{kazeo} Fromajo ek nur lakto koagulita, gutifigita e muldita, e quan onu manjas pos adjuntir a lu kremo\lingui{DEFIS}
\artiklo{kazerno}%
\vorto{kazerno} Domego en qua onu lojas la armeani\lingui{DFIRS}
\artiklo{kazimiro}%
\vorto{kazimiro} Drapo dina e lejera, ma di qua la fil-interkrucumi esas plu komplikita, e la texajo plu mikra-masha, kam en la texuro simpla\lingui{DEFIRS}
\artiklo{kazino}%
\vorto{kazino} Loko di interrenkontro por karto-ludar, muzikar, lektar la jurnali, specale en la urbi mar-balnala o kuracaquala\lingui{DEFIRS}
\artiklo{kazo}%
\vorto{kazo}\senco{1}{Ek la eventaji quin povas produktar stando, konjunturo, la eventajo qua esas produktita, fakte o supozate o konjektate}\senco{2}{\textit{(gram.)} En ula lingui, dezinenco di la substantivo (e di la adjektivo e pronomo qui relatas lu), por indikar lua precipua roli en la frazo}\lingui{DEFIRS}
\artiklo{kazuala}%
\vorto{kazuala} Subordinita ad ula kazi.\lingui{DEFIS}\subvorto{}{kazualo}{Profito-quanto (qua varias segun la cirkonstanci) qua augmentas la revenuo singladia}{}{}
\artiklo{kazuaro}%
\vorto{kazuaro}\fako{zool.} Ucelo steltiera di la arkipelago Indiana e di Australia, simila a la strucho, ma di qua la ali esas plu kurta\latina{casuaria}\lingui{DEFIRSL}
\artiklo{kazuistiko}%
\vorto{kazuistiko}\senco{1}{Parto di teologio, qua traktas la solvo di la kazi di koncienco}\senco{2}{Tendenco a serchar subtilaji rezonala}\lingui{DEFIRS}
\artiklo{kazuisto}%
\vorto{kazuisto} Teologiisto qua su konsakris a studiar e solvar la kazi di koncienco\lingui{DEFIRS}
\artiklo{ke}%
\vorto{ke} Konjunciono subordinala\lingui{FIS}
\artiklo{keglo}%
\vorto{keglo} Singla ek la non koni ligna quin onu pozas sur sulo olia-baze, e quin onu probas abatar per ligno-bulo quan onu lansas de kelke fore\lingui{DFR}
\artiklo{kelero}%
\vorto{kelero} Spaco subsula, ordinare vultoza, rezervita sub la ter-etajo di la domi, ed apta (pro lua temperaturo sempre egala) konservar la vini e la diversa provizaji\lingui{DEF}
\artiklo{kelka}%
\vorto{kelka} Mikra quanto, mikra nombro. (\textit{Kelka} indikas quanto min granda kam multa, e plu granda kam \textit{plura}; \textit{plura} indikas, ula-cirkonstance, nur « du »)
\artiklo{kelo}%
\vorto{kelo}\fako{tekn.} Mikra stifto ek fero, ek ligno, muntita trans-verso aden la pinto di bolto, di charniro-cilindro, e c., kom haltigivo\lingui{D}
\artiklo{kemio}%
\vorto{kemio} La cienco qua studias la konsisto intima di la korpi diversa (simpla o komplexa), la *leyi (legi) segun qui interkombinas la elementi de qui oli naskas, ed olia qualesi propra\lingui{DEFIRS}
\artiklo{kemiotaktismo}%
\vorto{kemiotaktismo}\fako{biol.} Atrakto o repulso, che enti vivanta, da substanco kemiala
\artiklo{kemiotropismo}%
\vorto{kemiotropismo}\fako{biol.} Orientizo, direkto di la planti, per reaktivi kemiala
\artiklo{kenelo}%
\vorto{kenelo} Buleto oblonga de karno, de fisho-karno, uzata kom disho-garnituro\lingui{EF}
\artiklo{kepio}%
\vorto{kepio}\senco{1}{Kapo-vesto lejera, kun mikra viziero, quan *weras ordinare la armeani}\senco{2}{Kapo-vesto sama-forma di la skolani, kolegiani, liceani, e c.}\lingui{DEFIRS}
\artiklo{keratino}%
\vorto{keratino} Substanco albuminoida, tre sulfoza, qua esas elemento di la tisui sintezala (tegumenti, faneri)\lingui{FS}
\artiklo{kerlo}%
\vorto{kerlo} Viro vivaca e joyoza\lingui{D}
\artiklo{kermeso}%
\vorto{kermeso}\fako{zool.} Kochenilo qua vivas sur speco di querko, e di qua la femino facas, sur la folii, la stipi e la branchi di ta arboro, mikra sheli sferatra e reda\latina{coccus ilicia}\lingui{DEFIRS}
\artiklo{kerno}%
\vorto{kerno} Parto harda, ligna, qua esas interne di ula frukti, e qua kontenas la « mandelo », la grano\lingui{DE}
\artiklo{kerozeno}%
\vorto{kerozeno}\fako{kemio} Lum-oleo, obtenita per distilar petrol-olei kruda\lingui{EF}
\artiklo{kerubo}%
\vorto{kerubo}\senco{1}{En la Testamento Olda (biblo) : reprezenturo de la anjeli, qui ornis la tabernaklo; nomo di ula anjeli}\senco{2}{Che la kristani : anjelo qua, en la hierarkio cielala, esas situita pos la arkianjelo, e quan la pikturi e skulturi reprezentas ordinare ye kapo di infanto aloza}\lingui{DEFIRS}
\artiklo{kesto}%
\vorto{kesto} Quaza buxo por la transporto; asemblajo ek planki (ulakaze : ajuroza, kun fundo e kovrilo fixigita per klovi)
\artiklo{khan}%
\vorto{« khan »}\senco{1}{Sinioro (che la hordi nomada, Turka o Mongola)}\senco{2}{Suvereno (olim, che la Turki e la Tatari.)}\senco{3}{\textit{(en Persia)} Guvernisto di provinco, ed ulagrade oficiero statala}
\artiklo{khedivo}%
\vorto{« khedivo »} La suvereno di Egiptia, olim
\artiklo{kikar}%
\vorto{kikar}\fako{netrans.) (aludante kavalo, asno, e c} Lansar rapide addope sua membri dopa, apogante su sur la membri avana\lingui{E}
\artiklo{kilio}%
\vorto{kilio}\fako{navig.} Asemblajo longa e rezistiva de trabi, qua iras de la pruo a la pupo di navo, fundamento sur qua onu muntas la navo-guvernilo, e c.\lingui{DEFIS}
\artiklo{kiliogono}%
\vorto{kiliogono}\fako{geom.} Poligono mil-latera
\artiklo{kilo-}%
\vorto{kilo-} Prefixo ciencala : « … × 1000 »
\artiklo{kilogrametro}%
\vorto{kilogrametro}\fako{mekan.} Mezur-unajo por la efiko da la mashini e da la motori, forco qua esas necesa por sublevar 1 kilogramo ad 1 metro super sulo\lingui{DEFIRS}
\artiklo{kilogramo}%
\vorto{kilogramo} Pezo mil-grama. (En Britania ed Usa : cirkume 2.2 lbs.)
\artiklo{kilolitro}%
\vorto{kilolitro} Mezur-unajo pri kapaceso : 1000 litri\lingui{DEFIRS}
\artiklo{kilometro}%
\vorto{kilometro} Mezur-unajo pri longeso : l000metri\lingui{DEFIRS}
\artiklo{kilowatt}%
\vorto{« kilowatt »} 1000 « watt »
\artiklo{kimero}%
\vorto{kimero}\senco{1}{Monstro quan la antiqui reprezentis kom asembluro stranja ek parti di animali diversa}\senco{2}{\textit{(metaf.)} Kreajo imaginala da la spirito, quan onu aceptas kom realajo}\lingui{DEFIRS}
\artiklo{kimono}%
\vorto{kimono} Tuniko longa, *werata (en Japonia) da la viri e mulieri; lu karakterizesas da larja zono ye la tayo. – Jap
\artiklo{kin}%
\vorto{kin} 2 plus 3, o 1 plus 4\lingui{FIS}
\artiklo{kinazo}%
\vorto{kinazo}\fako{biol.} Fermento solvebla, apta transformar a fermento aktiva la profermenti
\artiklo{kiosko}%
\vorto{kiosko}\senco{1}{Belvedero, paviliono apertita omna-latere, situita en gardeno, sur teraso}\senco{2}{En la granda urbi : mikra paviliono en qua onu vendas jurnali, flori; o paviliono plu granda, en qua su instalas la muzikisti municipala qui venas ibe por plear publike arii por la muziko-prizanti}\lingui{DEFIRS}
\artiklo{kiragro}%
\vorto{kiragro}\fako{patol.} La formo di goto qua afektas la manui\lingui{DEFIRS}
\artiklo{kirko}%
\vorto{kirko} Edifico por la kulto katolika od ortodoxa (kristanismo)\lingui{DE}
\artiklo{kiromanciar}%
\vorto{kiromanciar}\fako{netrans.} \textit{(asertate)} anuncar la eventaji futura, divinar la karaktero di la homi, per la exploro di la manui di la konsultanto – exploro, precipue, di la linei sur olia palmi\lingui{DEFIRS}
\artiklo{kiroptero}%
\vorto{kiroptero}\fako{zool.} Serio de mamiferi, di qui la vespertilio esas la tipo\lingui{EFIL}
\artiklo{kirsho}%
\vorto{kirsho} Brandio de cerizi, obtenita per fermentacigar cerizi nigra ed olia kerni\lingui{DEFIRS}
\artiklo{kirurgio}%
\vorto{kirurgio} Parto di la medicin-arto qua traktas precipue la lezuri extera (vunduri, osto-rupturi, e c.) e la operaci quin igas necesan ula morbi interna\lingui{DEFIRS}
\artiklo{kisar}%
\vorto{kisar}\fako{trans.} Presar sua labii, manifeste di afeciono, di amo, di amoro, di respekto, adsur la vizajo, la manuo di persono\lingui{DE}
\artiklo{kismet}%
\vorto{« kismet »}\fako{che la Mohamedisti} Fato
\artiklo{kisto}%
\vorto{kisto}\fako{patol.} Quaza posho sen aperturo, qua kreskas en la tisui pro la dilateso di la kondukto-tubi extera di la glandi diversa di qui la orifico divenis obstruktita, e qua kontenas humori ed altra materii morbala\lingui{EFIRS}
\artiklo{kizelguro}%
\vorto{kizelguro} Materio qua konsistas precipue ye silico (til 90\%), generale verda, qua devenas de la deskompozeso di diatomei qui kaptabis silico; uzata por facar dinamito, smalio, cementi, e c.\lingui{IRS}
\artiklo{klado}%
\vorto{klado} Skriburo destinita a transskribesor, e quan onu emendas, surstrekizas, e c. ante la kopio neta e definitiva\lingui{D}
\artiklo{klakar}%
\vorto{klakar}\fako{netrans.} Bruisar pro la intershokeso (instantala) di du surfaci plu o min plata\lingui{EF}
\artiklo{klamar}%
\vorto{klamar}\fako{trans.} Dicar (ulo) per vorti pronuncata forte, rezonante\lingui{EFI}
\artiklo{klamido}%
\vorto{klamido}\fako{epoki antiqua} Mantelo Greka, Romana, di qua un ek la anguli elevesas til la shultro dextra, ube lu an-ligesas per agrafo\lingui{DEFIRS}
\artiklo{klano}%
\vorto{klano} Quaza tribuo, che la populi Kelta\lingui{DEFIRS}
\artiklo{klapo}%
\vorto{klapo}\senco{1}{\textit{(tekn.)} Quaza valvo qua apertesas e klozesas quale charniro-plako}\senco{2}{Peco quan onu povas quaze faldar sur altra peco (fixa)}\lingui{DF}
\artiklo{klara}%
\vorto{klara}\senco{1}{(a) Qua donas o recevas lumo-radii quin nulo interceptas; (b) Di qua nulo desbriligas la pureso}\senco{2}{\textit{(metaf.)} (a) Dicesas pri son-uno, fono, qua ne esas matida, qua resonas bone en la orelo; (b) Dicesas pri to quon perceptas la spirito – quo ne esas obskura, – quon la spirito komprenas bone}\lingui{DEFIRS}
\artiklo{klarineto}%
\vorto{klarineto}\fako{muziko} Sufl-instrumento, kun beko e langeto\lingui{DEFIS}
\artiklo{klasifikar}%
\vorto{klasifikar}\fako{trans.}\senco{1}{\textit{(pri la natur-cienci)} Repartisar la individui hierarkiale, segun speci, generi, familii, serii, klasi, e c.}\senco{2}{Distributar metodoze}\lingui{DEFIS}
\artiklo{klasika}%
\vorto{klasika} Egardata kom konforma a la reguli di la genero, ed uzebla kom modelo\lingui{DEFIRS}
\artiklo{klaso}%
\vorto{klaso}\senco{1}{Kategorio de civitani}\senco{2}{Kategorio de adolecanti, *konvokata singla-yare por lia instrukteso militala}\senco{3}{Kategorio de lernanti qui frequentas singla grado di lerno-kurso}\senco{4}{Omna kategorio de personi, de kozi, distributita segun la difero pri grado o mem (nur) pri naturo}\senco{5}{Singla ek la granda grupi de animali o de vejetanti}\lingui{DEFIRS}
\artiklo{klastika}%
\vorto{klastika}\fako{geol.} Qua devenas efike da la rodo mekanismala de la aquo movanta. (Ex. : la sedimento-roki)
\artiklo{klaudikar}%
\vorto{klaudikar}\fako{netrans.} Marchar, apogante ne-interegale la gambi sur sulo, sive pro ke onu doloras ye la gambo, sive pro ke un gambo esas plu kurta\lingui{F}
\artiklo{klauno}%
\vorto{klauno}\senco{1}{Protagonisto bufona di la farsi Angla}\senco{2}{Aktoro qua, en ula cirki, distraktas, amuzas, la spektanti per prodaji di flexebleso, akompanata da joki e moki}\lingui{DEFIRS}
\artiklo{klauzo}%
\vorto{klauzo} Preskriptajo specala di lego, kontrato, testamento\lingui{DEFIRS}
\artiklo{klavikordo}%
\vorto{klavikordo}\fako{muziko} Instrumento kun sono fixa, klavaro e kordi metala quin pinchas plumo-pinti, qua remplasesis da la piano\lingui{DEFIS}
\artiklo{klavikulo}%
\vorto{klavikulo}\fako{anat.} Osto qua funcionas kom apog-arko inter la du shultri\lingui{EFIS}
\artiklo{klavo}%
\vorto{klavo} Singla ek la tabuleti levera, blanka o nigra, quin la fingri abasas por plear orgeno, piano; \textit{(analoge)} la klavi di skribomashino, di komposto-mashino\lingui{DEFRS}
\artiklo{klefo}%
\vorto{klefo} Peco ek fero, quan onu enduktas aden la truo di seruro, e per qua onu funcionigas la mekanismo qua apertas o klozas ta seruro\lingui{F}
\artiklo{kleistogama}%
\vorto{kleistogama}\fako{bot.} Dicesas pri floro qua nulinstante apertas su, e di qua la fekundigo eventas dum olua klozeso
\artiklo{klemar}%
\vorto{klemar}\fako{trans.} Tenar tre strete : (a) per lokizar en spaco tre poka (un o plura kozi o personi); (b) per interproximigar per ligilo (plura kozi o parti di kozo); (c) per lokizar ici tre an iti (personi o kozi)\lingui{D}
\artiklo{klematido}%
\vorto{klematido}\fako{bot.} Planto klimera, familio « ranunkulacei », di qua la flori esas blanka ed odoroza\latina{clematis}\lingui{EFIS}
\artiklo{klementa}%
\vorto{klementa} Dicesas pri persono qua, autoritatoza pri punisar kulpozo, agas dolce ad ica, pardonante lu, od alejante lua puniseso\lingui{EFIS}
\artiklo{klepsidro}%
\vorto{klepsidro} Aquo-horlojo qua indikas la kloki per la fluo ne-interruptita da ula quanto de aquo en tempo-quanto predeterminita\lingui{DEFIS}
\artiklo{kleriko}%
\vorto{kleriko} Membro di la oficieraro di eklezio\lingui{DEFIRS}
\artiklo{klerko}%
\vorto{klerko} La persono qua laboras o yuro-studias en la kontoro di notario, di ushero, di « avoué »\lingui{EF}
\artiklo{kliento}%
\vorto{kliento}\senco{1}{\textit{(Roma, antiqua)} Plebeyo qua protektesis da patrico}\senco{2}{La persono qua konfidas kustumale sua interesti ad aferisto (advokato, notario, e c.), sua saneso a mediko; la persono di qua advokato pledas la afero}\senco{3}{\textit{(komerco)} La persono qua kompras kustumale de komercisto, qua frequentas butiko (restorerio, e c.)}\lingui{DEFIRS}
\artiklo{klifo}%
\vorto{klifo} Eskarpajo de tero o de roko qua garnisas la bordo di maro\lingui{E}
\artiklo{kliko}%
\vorto{kliko}\fako{tekn.} Peco movebla quan sublevas omna dento di ingrano roto kande ica turnas ica-sinse, e qua butas kontre la ingranajo se onu volas igar lu turnar inversa-sinse\lingui{DEF}
\artiklo{kliktar}%
\vorto{kliktar}\fako{netrans.} \textit{(aludante ula korpi sonora qui intershokas)} Bruisar ne-matide\lingui{DEF}
\artiklo{klimar}%
\vorto{klimar}\fako{netrans.} Elevar su per akrochar su a to quo povas helpar la acenso\lingui{DE}
\artiklo{klimato}%
\vorto{klimato} Ensemblo ek la cirkonstanci atmosferala, a qua regioni submisesas\lingui{DEFIRS}
\artiklo{klimatologio}%
\vorto{klimatologio} Studio di la klimati\lingui{DEFS}
\artiklo{klinar}%
\vorto{klinar}\fako{netrans.) (aludante du peci qui quaze interkavalkas} Interkrucumar\lingui{F}
\artiklo{kliniko}%
\vorto{kliniko} Docado medicinala, medikala, da maestro a sua dicipuli, an la lito di la maladi\lingui{DEFIRS}
\artiklo{klinko}%
\vorto{klinko} Klozilo di pordo, ek simpla fero-lamo quan onu abasas sur fero-peco muntita an la pordo-kadro, kande onu klozas la pordo e sublevas kande onu apertas lu\lingui{DF}
\artiklo{klinoida}%
\vorto{klinoida}\fako{anat.} Singla ek la quar apofizi, qui esas an la parto supra di la osto sfenoida
\artiklo{klinometro}%
\vorto{klinometro} Aparato uzata – en aernavigado, marnavigado – (an busolo) por mezurar la deviaco de la horizontalo\lingui{DFRS}
\artiklo{klipero}%
\vorto{klipero}\fako{navig.} Seglo-navo, multa-tuna, pasable rapida, uzata olim por la teo-komerco inter Britania e Chinia\lingui{DEFIRS}
\artiklo{klishar}%
\vorto{klishar}\fako{trans.}\senco{1}{\textit{(imprim-arto)} Reproduktar reliefe la traco de kompostajo ek tipi movebla, de la tipo originala di graburo, e c.}\senco{2}{\textit{(fotogr.)} Facar la vitro-plako negativa, de qua la fotografisto obtenos la pozitivo}\senco{3}{\textit{(hektogr.)} (a) Perforar la folio de materio vaxatra (per la tipi di skribo-mashino o per stilo se koncernesas manuo-skriburo o desegnuro) qua uzesos por la hektografilo; (b) skribar (o desegnar) sur la folio de materio paperatra, tre glata, per inko specala (hektografinko), to quo uzesos por la hektografilo}\lingui{DFRS}
\artiklo{klistero}%
\vorto{klistero} La medikamento liquida quan onu injektas, tra la anuso, aden la grosa intestino\lingui{DEFIRS}
\artiklo{klitorido}%
\vorto{klitorido}\fako{anat.} Mikra organo karnaja, enire di la vulvo\lingui{DEFIS}
\artiklo{klivar}%
\vorto{klivar}\fako{trans.} Dividar (diamanti, kristali, mikao-bloko) ad strati lamela\lingui{DEF}
\artiklo{kloako}%
\vorto{kloako}\senco{1}{Loko destinita ad esor kolektuyo di la rezidui}\senco{2}{\textit{(anat.)} Posho quan, en la korpo di la uceli, formacas la extremajo di la intestino}\lingui{DEFIRS}
\artiklo{kloko}%
\vorto{kloko} La 24-esma parto di dio, determinata per horlojo\lingui{E}
\artiklo{klonuso}%
\vorto{klonuso} Fenomeno qua rezultas de la exajereso di la reflexi\lingui{DEFIRS}
\artiklo{kloro}%
\vorto{kloro}\fako{kemio} Korpo simpla, gasa, flava verdatre, odorsufokiva, qua saporas kaustike, uzata kom senkolorigivo e kom desinfektivo\simbolo{Chl}\lingui{DEFIRS}
\artiklo{klorofilo}%
\vorto{klorofilo}\fako{bot.} Materio qua donas a la parti verda di la planto olia koloro\lingui{DEFS}
\artiklo{kloroformo}%
\vorto{kloroformo} Substanco anesteziala, liquida, oleatra, senkolora\lingui{DEFIRS}
\artiklo{kloroleucito}%
\vorto{kloroleucito}\fako{biol.} Granulo de protoplasmo, impregnita ye klorofilo
\artiklo{kloroplasto}%
\vorto{kloroplasto}\fako{biol.} Rolo quan pleas la kloroleucito la produkto di klorofilo
\artiklo{kloroso}%
\vorto{kloroso}\fako{patol.} Paleso verdatra. Etioleso di la adolecantini en qui la pubereso-kresko esas neregulala\lingui{DEFIRS}
\artiklo{klosho}%
\vorto{klosho}\senco{1}{Instrumento ek bronzo, forme di vazo inversigita, qua longe vibras efike da frapilo suspensita interne, o da martelo extera}\senco{2}{To quo aspektas quale klosho}\lingui{F}
\artiklo{klostro}%
\vorto{klostro}\fako{religio katol.}\senco{1}{Parto di kuvento quan klozilo separas de la parti cetera di la konstrukturo}\senco{2}{Portiko tektoza, cirkum la korto interna di monakeyo}\lingui{DEFIS}
\artiklo{klovo}%
\vorto{klovo} Mikra stango fera, pintoza, kun un kapo, quan onu uzas por fixigar, muntar o suspensar ulo\lingui{FIS}
\artiklo{klozar}%
\vorto{klozar}\fako{trans.} Dispozar (la klapo o klapi di pordo, di fenestro, barilo, parieto, muro, greto, kordono de soldati) tale ke (la chambro, la chambrego, la parko, la gardeno, e c.) esez *kluza. (Chambro esas *kluza se lua pordo e lua fenestro esas klozita)\lingui{EFI}
\artiklo{klubo}%
\vorto{klubo} Asociitaro ek personi qui pagas komuna-spense po la uzado di chambrego en qua li povas su asemblar por ludar (per karti, e c.), lektar la jurnali, la revui, diskutar pri politiko, e c.\lingui{DEFIRS}
\artiklo{klukar}%
\vorto{klukar}\fako{netrans.) (aludante la hanino, la ganso, e c} Audigar sua krio\lingui{DEFIRS}
\artiklo{*kluza}%
\vorto{*kluza} Dicesas pri loko qua esas cirkondita per barilo, parieti, muri, greti, kordono de soldati, qui impedas la aceso di olua internajo\lingui{EFI}
\artiklo{kluzo}%
\vorto{kluzo} Truo cirkla horizontala, borita en la parto avana di navo (dextre o sinistre di la pruo), tra qua pasas la kablo muntita an ankro\lingui{DR}
\artiklo{knaucio}%
\vorto{knaucio}\fako{bot.} Planto herbatra, kun flori violea, purpurea, ed anke blanka, quan onu uzis olim kom medikamento kontre skabio\lingui{L}
\artiklo{knikar}%
\vorto{knikar}\fako{trans.} Strenar korpo relative longa (kom ex. : stango) per preso, en la direciono longesala, sur la du extremaji\lingui{D}
\artiklo{knuto}%
\vorto{knuto} Instrumento di tormentado, che la Rusi caroza : flogilo qua konsistas ye plura ledro-bendi interplektita, qui, cirkum la extremajo, esas separita e provizita ye hoki ek fero\lingui{DEFIRS}
\artiklo{ko-}%
\vorto{ko-} Prefixo ciencala qua signifikas « komplemento di… » o « di komplemento »
\artiklo{koadjutoro}%
\vorto{koadjutoro}\fako{religio katol.} Kleriko adjuntita ad episkopo, ad arkiepiskopo, por helpar lu en lua funcioni episkopala e remplasar lu se la ofico vakas\lingui{DEFIRS}
\artiklo{koagular}%
\vorto{koagular}\fako{netrans.) (aludante liquido} Asemblar la parti solida quin lu kontenas interne-flotace
\artiklo{koaktar}%
\vorto{koaktar}\fako{trans., pri, ad} Reduktar (ulu) ad agar kontre lua volo\lingui{F}
\artiklo{koalisar}%
\vorto{koalisar}\fako{netrans.) (aludante plura populi, plura nacioni} Inter-unionar poke-duronte por atako kontre adverso komuna\lingui{DEFIRS}
\artiklo{koaltaro}%
\vorto{koaltaro} Gudro quan onu obtenas per distilar min-karbono\lingui{EF}
\artiklo{kobalto}%
\vorto{kobalto}\fako{kemio} Korpo simpla, metala, stalo-griza, di qua ula kombinuri (oxido, arseniato) uzesas por bluigar porcelano, vitro, e c.\simbolo{Co}\lingui{DEFIRS}
\artiklo{kobayo}%
\vorto{kobayo}\fako{zool.} Mamifero rodera, genero « subunglozi », kun gambi kurta, quar fingri ye la avan-gambi, tri fingri ye la dop-gambi, e qua vivas en la sudo-regioni di Amerika\latina{Cavis cobaya}\lingui{FS}
\artiklo{kobeo}%
\vorto{kobeo}\fako{bot.} Planto klimera, kun flori violea, kloshatra. (scandena)\latina{cobea}\lingui{DEFIRS}
\artiklo{koboldo}%
\vorto{koboldo} En la superstici plebeyala : mikra diablo malicoza qua tormentas la homi dum lia dormo\lingui{DEF}
\artiklo{kochenilo}%
\vorto{kochenilo}\fako{zool.} Insekto qua vivas sur la nopalo e donas bela tinto-materio reda\latina{coctus cacti}\lingui{DEFIRS}
\artiklo{kocigo}%
\vorto{kocigo}\fako{anat.} Mikra osto, en la extremajo infra di la spino, *artikumata a la sakrumo, che la homo e che la spinozi sen-kauda
\artiklo{kocinelo}%
\vorto{kocinelo}\fako{zool.} Mikra insekto koleoptera. (septempunctata)\latina{coccinela}\lingui{EFIS}
\artiklo{kodexo}%
\vorto{kodexo} Kolektajo ek la legi klasifikita tale ke olu reprezentas la ensemblo di la legi di lando pri ta o ca temo. – Anke pri la formuli farmaciala, la frazi konvencionala telegrafala\lingui{DEFIRS}
\artiklo{kodicilo}%
\vorto{kodicilo} Klauzo adjuntita – qua kompletigas, modifikas o nihiligas testamento\lingui{DEFIRS}
\artiklo{koeficiento}%
\vorto{koeficiento}\fako{matem.} Nombro qua, skribite avan quanto algebrala, indikas quantafoye olca esas uzenda\lingui{DEFIRS}
\artiklo{koercitiva}%
\vorto{koercitiva}\fako{elektro} Feldo koercitiva : valoro de magneto-feldo, necesa por retroduktar a zero la indukteso magnetala en ica od ita punto di korpo feromagnetala, pos submisir ica a sat multa cikli magnetala
\artiklo{kofio}%
\vorto{kofio} Kapo-vesto mulierala, ek stofo, trikoturo, furo, e c., sen bordo\lingui{EFS}
\artiklo{kofro}%
\vorto{kofro} Granda buxo ek ligno, ek metalo, rektangula, kun klozilo klapa, maxim-multa-kaze konvexa, *kluzigebla per seruro\lingui{DEFIRS}
\artiklo{koherar}%
\vorto{koherar}\fako{netrans.) (aludante la elementi di ensemblo} Havar sua parti strete unionita\lingui{DEFIS}
\artiklo{kohorto}%
\vorto{kohorto}\fako{epoki antiqua} Trupo qua konsistis ye la dekimo di legiono Romana\lingui{DEFIRS}
\artiklo{koincidar}%
\vorto{koincidar}\fako{netrans.} Interkontaktar precize, en omna punti\lingui{DEFIS}
\artiklo{*koino}%
\vorto{*koino} Moneto-peci ek arjento, kupro, bronzo, nikelo, qui valoras poke. (Kom ex. : Me kambiis moneto-peco 5-franka po koino.)\lingui{E}
\artiklo{koitar}%
\vorto{koitar}\fako{netrans., kun) (aludante homi} Unionar su sexuale, por juo o genito\lingui{DF}
\artiklo{kokaino}%
\vorto{kokaino}\fako{kemio} Alkaloido ek la kokao, quan onu uzas medicine por anesteziar lokale\lingui{DEFIRS}
\artiklo{kokao}%
\vorto{kokao}\fako{bot.} Arbusto aromatoza di Peru, di qua la folio kontenas elementi toniziva ed ecitigiva\latina{erythroxylum coca}\lingui{DEFIRS}
\artiklo{kokardo}%
\vorto{kokardo} Insigno, forme di mikra disko, qua *weresas an la chapelo\lingui{DEFIRS}
\artiklo{koketa}%
\vorto{koketa} Qua esforcas plezar a la personi di la altra sexuo\lingui{DEFIRS}
\artiklo{kokleo}%
\vorto{kokleo}\fako{zool.} Molusko gasteropoda, reptera, kun konko helicatra\latina{cochlea}\lingui{EFIS}
\artiklo{koklusho}%
\vorto{koklusho}\fako{patol.} Tusado konvulsatra, qua kaptas precipue la infanti\lingui{FR}
\artiklo{koko}%
\vorto{koko} Kombusteblo, reziduo de la min-karbono quan onu kalcinis en vazo *kluza, por separar de olu la hidro e la parti sulfa e bitumina\lingui{DEFIS}
\artiklo{kokono}%
\vorto{kokono} Envelopilo quan facas, ek sua fili, la maxim multa larvi (raupi) di la lepidoptari, ed en qua eventas lia lasta metamorfoso\lingui{DEFR}
\artiklo{kokoso}%
\vorto{kokoso}\fako{bot.} Frukto di granda arboro ek la familio « palmieri »\latina{cocos nucifera}\lingui{DEFIRS}
\artiklo{kolacionar}%
\vorto{kolacionar}\fako{trans., kun} Komparar ica kun ita, o kun la originalo, kopiurin, reprodukturin de texto, kom verifiko di lia exakteso\lingui{DEFI}
\artiklo{kolao}%
\vorto{kolao}\fako{bot.} Arboro de la regioni tropikala di Afrika, qua produktas « nuco » uzata kom tonizivo por la kordio e kom stimulivo nervala\latina{cola acuminata}\lingui{DEFI}
\artiklo{*kolapsar}%
\vorto{*kolapsar}\fako{netrans.) (aludante homo} Falar pro exhausteso, extenueso, febleso\lingui{E}
\artiklo{kolareto}%
\vorto{kolareto} Stofo-garnituro por la kolo, generale plisoza, quan *weris la viri lor la rejo Henriko IV di Francia (t. e. cirkum la yaro 1600) e la mulieri, divers-epoke\lingui{DEFI}
\artiklo{kolaterala}%
\vorto{kolaterala}\senco{1}{Situita laterale, relate ulo}\senco{2}{\textit{(metaf.)} Qua apartenas a la genealogio-lineo « frati, kuzi, nevi, onkli »}\lingui{DEFIS}
\artiklo{kolchiko}%
\vorto{kolchiko}\fako{bot.} Planto bulboza, kun flori violea, di qua la toxikeso uzesas por kuracar artritismo\latina{colchicum autumnale}\lingui{EFIS}
\artiklo{kolda}%
\vorto{kolda}\senco{1}{Qua kontenas poka kaloro}\senco{2}{\textit{(metaf.)} Qua ne esas ardoroza, nek pasionoza}\lingui{DE}
\artiklo{koldkremo}%
\vorto{koldkremo} Kosmetiko por la pelo, qua konsistas ye vaxo blanka, spermaceto, ed oleo ek mandeli\lingui{DEFIRS}
\artiklo{kolegio}%
\vorto{kolegio} Korpo ek personi submisata a regularo komuna : 1. \textit{(en Roma, antique)} Korporaciono; 2. la korpo ek elekteri politikala qui apartenas ad un distrikto votala; 3. edukerio publika – en Francia : edukerio municipala (kontraste kun la « liceo », edukerio statala)\lingui{DEFIRS}
\artiklo{kolego}%
\vorto{kolego} Singla de la homi qui havas ofico publika, civila od armeala, egardata relate a la ceteri\lingui{DEFIRS}
\artiklo{kolektar}%
\vorto{kolektar}\fako{trans.} Asemblar (por facar ek li ensemblo) kozi amasigita ek fonti diversa\lingui{DEFIRS}
\artiklo{kolektivismo}%
\vorto{kolektivismo}\fako{sociologio} Sistemo qua ofras, kom solvuro di la problemo « quale desaparigar povreso », la komunigo, profite da socio, di omna produktili\lingui{DEFIRS}
\artiklo{kolektoro}%
\vorto{kolektoro}\fako{tekn.) (elektro} Parto di \textit{elektro}-mashino, ek asemblitaro de lami ek metalo elektrokonduktiva (kupro, precipue), dispozita periferie di cilindro, ed inter-izolita per lami de mikao\lingui{DEFIRS}
\artiklo{koleoptero}%
\vorto{koleoptero}\fako{zool.} Insekto kun quar ali, de qui la du supera (la elitri), harda, dika, esas quaza etuyi por la du altri\lingui{DEFIRSL}
\artiklo{kolerino}%
\vorto{kolerino}\fako{patol.} Diareo fortega qua esas (ulakaze) simptomo qua anuncas kolero\lingui{DEFS}
\artiklo{kolero}%
\vorto{kolero}\fako{patol.} Morbo epidemiala, quan karakterizas grava perturbeso intestinala, kun koldesko interna e kontrakto spasmatra di la membri\lingui{DEFIRS}
\artiklo{koliar}%
\vorto{koliar}\fako{trans.} Deprenar de la stipo (floro, frukto, ramo)\lingui{FIS}
\artiklo{koliaro}%
\vorto{koliaro}\senco{1}{Ringo qua cirkondas la kolo : 1. Ornivo por la kolo}\senco{2}{Ringo ek metalo, ek ledro, quan onu metis cirkum la kolo di la sklavi, quan onu metas cirkum la kolo di la animali por ligar li ad ulo}\senco{3}{Kolo-harneso di la tiro-bestii}\senco{4}{Ringo qua cirkondas ula objekti}\lingui{DEFIS}
\artiklo{kolibrio}%
\vorto{kolibrio}\fako{zool.} Ucelo di la regioni tropikala di Amerika, kun beko arkatra, kolori briloza, la maxim mikra ek la uceli konocata\latina{trochilus}\lingui{DEFIRS}
\artiklo{kolikar}%
\vorto{kolikar}\fako{netrans.) (patol.} Dolorar en la intestini\lingui{DEFIRS}
\artiklo{kolimatoro}%
\vorto{kolimatoro}\fako{optiko} Aparato uzata en la artilrio por vizar lor la apunto di kanono\lingui{DEF}
\artiklo{kolimbo}%
\vorto{kolimbo}\fako{zool.} Ucelo aquala, qua povas permanar dum multa tempo sur la aquo-surfaco\latina{colymbus}\lingui{L}
\artiklo{kolino}%
\vorto{kolino} Ter-altajo (min kam 500-metro) qua dominacas la plana lando\lingui{FIS}
\artiklo{kolirio}%
\vorto{kolirio}\fako{medic.} Medikamento (aquo, pomado, e c.) destinata ad aplikesor sur la konjunktivo di la okulo\lingui{EF}
\artiklo{kolizionar}%
\vorto{kolizionar}\fako{netrans.) (aludante du korpi} Intershokar\lingui{DEFIRS}
\artiklo{kolkotaro}%
\vorto{kolkotaro}\fako{kemio} Peroxido reda de fero, quan onu obtenas per kalcinar la fer-sulfato, e quan onu uzas por polisar la speguli\lingui{DEFS}
\artiklo{kolmo}%
\vorto{kolmo} La grado maxim alta an qua arivas ulo
\artiklo{kolo}%
\vorto{kolo} Parto di la korpo di la homo, di animalo, qua unionas la kapo a la torso\lingui{FIS}
\artiklo{kolodiono}%
\vorto{kolodiono} Solvuro ek pulvero-kotono en etero, alkoholizita, uzata kirurgiale kom aglutinivo, e fotografale por igar adheriva, sur la vitro-plaki, la strato de iodido\lingui{DEFIRS}
\artiklo{kolofono}%
\vorto{kolofono} Materio rezinoza, ye qua onu fricionas la krini di la arketo por ke olu ne glitez sur la kordi di la violino, di la violoncelo, e c.\lingui{DEFIRS}
\artiklo{koloido}%
\vorto{koloido}\fako{kemio} Kompozajo ne-kristaligebla e ne-dializebla\lingui{DEF}
\artiklo{kolokar}%
\vorto{kolokar}\fako{trans.} Prestar (mediate o nemediate) pekunio por recevor de olu revenuo\lingui{IS}
\artiklo{kolombo}%
\vorto{kolombo}\fako{zool.} Ucelo, meza inter la « galinacei » e la « paseri », di qua la pektoro-plumi esas tre koloroza\latina{columba}
\artiklo{kolonelo}%
\vorto{kolonelo} La viro qua, en armeo, komandas regimento\lingui{DEFIS}
\artiklo{kolonio}%
\vorto{kolonio}\senco{1}{Establisuro fondita da naciono di qua ula quanto de habitanti iras aden regiono stranjera, ed ibe su instalas fixe, sen cesar apartenar a la metropolo}\senco{2}{Kolonio agrokultivala : establisuro rurala, destinita okupor laboro povri, yuna karcerani, ed emendar la sulo}\senco{3}{Grupo de individui ek un naciono, instalita fixe en lando stranjera, en rezido-regiono komuna}\lingui{DEFIRS}
\artiklo{kolono}%
\vorto{kolono}\senco{1}{\textit{(arkitekt.)} Pilastro ek quadro, ek marmoro, qua finas per kapitelo, uzata kom subtenilo o kom ornivo di edifico}\senco{2}{\textit{(anat.)} Parto di la grosa intestino, qua sequas la ceko e su extensas til la rektumo}\lingui{DEFIRS}
\artiklo{koloquinto}%
\vorto{koloquinto}\fako{bot.} Varietato di kukombro\latina{cucumis colycynthis}\lingui{DEFIRS}
\artiklo{kolorito}%
\vorto{kolorito} Efekto qua rezultas de la selekto, de la uzo di ula farbi en pikturo\lingui{DFIRS}
\artiklo{koloro}%
\vorto{koloro}\senco{1}{Impreso partikulara, quan agas, sur la vid-organo, ta o ca ek la elementi aden qui su deskompozas lumo-radio, forsendata da surfaco qua reflektas ta elemento absor-- ante la cetera}\senco{2}{Proprajo di ula korpi pri tale impresar}\senco{3}{\textit{(metaf.)} Karaktero semblanta di la kozi}\lingui{DEFIS}
\artiklo{koloso}%
\vorto{koloso}\senco{1}{Statuo di qua la dimensioni esas grandega}\senco{2}{Homo, animalo, di qua la dimensiono esas grandega}\senco{3}{\textit{(metaf.)} La homo qua esas grandega pri ca o ta ek lua qualesi}\lingui{DEFIRS}
\artiklo{kolportar}%
\vorto{kolportar}\fako{trans.} Portar (vari) aden diversa loki, regioni, e c., (por disvendar li en ca loki, regioni)\lingui{DEF}
\artiklo{kolubro}%
\vorto{kolubro}\fako{zool.} Reptero ne-venenoza, ek la familio « serpenti »\latina{coluber}\lingui{EFSL}
\artiklo{kolumbario}%
\vorto{kolumbario}\fako{en Roma, antique} Konstrukturo funerala qua kontenis nichi aden qui onu pozis la urni mortintala\lingui{DEFIS}
\artiklo{kolumno}%
\vorto{kolumno} Dividuro vertikala (kolonatra), inter-paralela, di la pagini de komposturo, de imprimuro (texto)\lingui{DE}
\artiklo{koluro}%
\vorto{koluro}\fako{astron.} Singla de la du granda cirkli di la Tero-sfero, qui intersekas ortangule ye la poli, e pasas, ica ye la punti equinoxala, ita ye la punti solsticala\lingui{DEFIS}
\artiklo{koluteo}%
\vorto{koluteo}\fako{bot.} Arboreto, ek la familio « papilionacei », di qua la folii olim uzesis kom purgivo, e di qua la sheli, pro explozar bruisoze, esas amuzivo por la infanti\latina{colutea arborescens}\lingui{IL}
\artiklo{koluzionar}%
\vorto{koluzionar}\fako{netrans.} \textit{(aludante du personi qui su prizentas a la judiciisto)} Interkonvencionar sekrete, detrimente di triesma koncernato di la proceso\lingui{DEFIS}
\artiklo{kolzo}%
\vorto{kolzo}\fako{bot.} Planto krucifera, di qua la grano donas oleo quan onu uzas precipue por la lumizado\latina{brassica napus oleifera}\lingui{EFIRS}
\artiklo{kom}%
\vorto{kom} « Havante la qualeso, la titulo, la ofico, la rolo, di »\lingui{F}
\artiklo{komandar}%
\vorto{komandar}\fako{trans.) (aludante armeo-chefo}\senco{1}{Signalar ke movo, manovro, esas agenda da la subordinito – impero extreme severa, rigoroza ed absoluta, qua supresas omna responso che la imperario}\senco{2}{\textit{(tekn.)} \textit{(aludante organo di mekanismo)} Komunikar sua movo ad altra organo per ingranajo, pulio, rimeno, e c.}\lingui{DEFIRS}
\artiklo{komanditar}%
\vorto{komanditar}\fako{trans.} Subtenar (entraprezajo o la persono qua entraprezas) kom nura pekunio-prestero, sen intervenar pri quale uzesas ica pekunio\lingui{DEFIS}
\artiklo{komandoro}%
\vorto{komandoro} La persono qua esas super la oficiro, en la ordeni honoriziva di la *chevalieraro\lingui{DEFIRS}
\artiklo{komao}%
\vorto{komao}\fako{muziko} Intervalo muzikala tre mikra, quan onu desegardas, ordinare – pro ke ol esas apene orelo-perceptebla\lingui{DEF}
\artiklo{komato}%
\vorto{komato}\fako{patol.} Dormesko morba qua igas la dormanto aspektar quale mortinto\lingui{DEFIRS}
\artiklo{kombatar}%
\vorto{kombatar}\senco{1}{\textit{(netrans.)} \textit{(aludante du adversi, du armei)} Interfrapadar. (Se la quanto de koncernati, o la rezultaji, esus granda, onu dicus « bataliar ».)}\senco{2}{\textit{(trans.)} Komencar la interfrapado}\lingui{EFIS}
\artiklo{kombinar}%
\vorto{kombinar}\fako{trans.} Asemblar (du o plura elementi) segun ordino o segun proporcioni determinita\lingui{DEFIRS}
\artiklo{kombinatoriko}%
\vorto{kombinatoriko}\fako{filoz.} Metodo da Raymond Lulle, qua konsistis ye kombinar la homala idei tale ke solvesas omna problemi posibla
\artiklo{kombustar}%
\vorto{kombustar}\fako{trans.} Konsumar, destruktar per la efiko da fairo\lingui{EFIS}
\artiklo{komedio}%
\vorto{komedio}\senco{1}{Teatrajo qua ecitas a rido, per intervenigar protagonisti quin karakterizas defekti psikala o korpala, e qui esas situita en cirkonstanci jokala}\senco{2}{La teatraji komika o la tipo de komedio-teatrajo, propra ye epoko, ye skriptisto}\lingui{DEFIRS}
\artiklo{komedono}%
\vorto{komedono} Stopilo kornatra qua obstruktas la orifico di glando sebala, e formacas la « punti nigra » di la vizajo\lingui{EF}
\artiklo{komencar}%
\vorto{komencar}\senco{1}{\textit{(trans.)} Agar, facar, la parto unesma di kozo}\senco{2}{\textit{(netrans.)} Enirar en sua parto unesma}\lingui{EFIS}
\artiklo{komendar}%
\vorto{komendar}\fako{trans.} Komisar (fabrikisto) ke lu facez ca o ta kozo : komisar (komercisto) ke lu furnisez ta o ca varo\lingui{F}
\artiklo{komensalo}%
\vorto{komensalo}\fako{biol.} Animalo qua partoprenas la nutrivi di sua hosto animala
\artiklo{komentar}%
\vorto{komentar}\fako{trans.} Explikar, notizar, por klarigar la loki obskura di texto\lingui{DEFIRS}
\artiklo{komercar}%
\vorto{komercar}\fako{netrans., ye} Relatar por la interkambio di vari\lingui{DEFIRS}
\artiklo{kometo}%
\vorto{kometo}\fako{astron.} Astro kun traco-lineo lumoza, forme di kaudo o di longa hararo, qua parkuras cirkli tre oblonga\lingui{DEFIRS}
\artiklo{komforto}%
\vorto{komforto} Estado en qua la vivanto-bezoni esas satisfacita\lingui{DEFIRS}
\artiklo{komico}%
\vorto{komico}\fako{en Roma, antique} Asemblitaro ek la populani, segun la kurii o la centurii\lingui{DEFIRS}
\artiklo{komika}%
\vorto{komika} Qua ridigas\lingui{DEFIRS}
\artiklo{komisar}%
\vorto{komisar}\fako{ulu, pri} Konfidar (ad ulu) la tasko qua konsistas ye agar o facar (to o co)\lingui{DEFIS}
\artiklo{komisario}%
\vorto{komisario} La persono qua esas delegita pri ula funcioni, pri ula ofico – provizore o longa-dure\lingui{DEFIRS}
\artiklo{komisiono}%
\vorto{komisiono} To quon privato demandas de ulu ke ica agez vice lu (kompro, quero, e c.)\lingui{DEFIRS}
\artiklo{komisuro}%
\vorto{komisuro}\fako{anat.} Singla ek la du anguli ube la labii interunionas\lingui{DEFIS}
\artiklo{komitato}%
\vorto{komitato} Grupo de personi, selektita ek korpo plu personoza, por divenar asemblitaro specala di qua la tasko esas precipue exekutigala\lingui{DEFIRS}
\artiklo{komizo}%
\vorto{komizo} Employato di administreyo, di komerco-firmo, di banko-firmo\lingui{DFI}
\artiklo{komo}%
\vorto{komo} *Puntuo-signo (,) qua indikas ke oportas (dicionale) separar ula membri di la frazo, sen haltar longe\lingui{DEF}
\artiklo{komocionar}%
\vorto{komocionar}\fako{trans.} Shanceligar subite\lingui{DEFIS}
\artiklo{komoda}%
\vorto{komoda} Quan karakterizas la uzo-facileso segun-deziro\lingui{DEFIS}
\artiklo{komodo}%
\vorto{komodo} Moblo, di qua la suprajo ne esas plu alta kam la adulto-kubiti, garnisita ye tirkesti en qui onu lokizas linjo, vesti, e c.\lingui{DEFIS}
\artiklo{komodoro}%
\vorto{komodoro}\senco{1}{\textit{(navaro Britaniana ed Usana)} La oficiro di qua la grado esas nemediate sub olta di sub-amiralo}\senco{2}{\textit{(navaro Nederlandana)} Navestro qua komandas eskadro}\lingui{DEFIS}
\artiklo{komono}%
\vorto{komono}\senco{1}{\textit{(lor feudismo)} Urbo liberigita de la yugo feudala, e di qua la borgezi su administris ipse, per oficistaro elektita}\senco{2}{En Francia : urbo, burgo, vilajo o grupo de domi, ye qua konsistas la unajo teritoriala quan administras « urbestro » kun « *konselio municipala »}\senco{3}{En Britania : urbo, burgo, qua nominas deputiti a la Parlamento}\lingui{DEFIRS}
\artiklo{kompakta}%
\vorto{kompakta}\fako{aludante korpo solida} Di qua la mas-uni esas tre inter-kompresanta\lingui{DEFIRS}
\artiklo{kompanio}%
\vorto{kompanio}\senco{1}{\textit{(komerco)} La asociito o la asociiti, aludita da la abreviuro « e Ko. »}\senco{2}{Parto di bataliono, quan komandas kapitano}\lingui{DEFIRS}
\artiklo{kompano}%
\vorto{kompano}\senco{1}{La persono qua vivas kustumale intima-relate kun ulu}\senco{2}{\textit{(olim)} Laboristo qua, ne-patrona, iris laborar che patrono kun altra laboristi}\lingui{EFR}
\artiklo{komparar}%
\vorto{komparar}\fako{trans.}\senco{1}{\textit{(komparar kun :)} Apudigar du o plura kozi por determinar lia punti di simileso e di nesimileso reciproka}\senco{2}{\textit{(komparar ad :)} Apudigar por igar relatar egalesale}\senco{3}{Apudigar to quon onu aludas, ad ulo analoga, por igar plu komprenebla, plu sentebla, lo dicata}\lingui{EFIS}
\artiklo{komparativo}%
\vorto{komparativo}\fako{gram.} Modifikuro di la adjektivo, qua indikas ke la qualeso expresata da lu pozitive esas egardata ye grado superiora\lingui{DEFIS}
\artiklo{kompaso}%
\vorto{kompaso} Instrumento ek du branchi de ligno o de metalo, asemblita ye un ek lia extremaji, tale ke li povas interforigar od inter-apudigar, por mezurar anguli, trasar cirkli dimensionale inter-distanta\lingui{EFIRS}
\artiklo{kompatar}%
\vorto{kompatar}\fako{trans., pri, pro} Partoprenar la sufro, la chagreno, da altru\lingui{EFIS}
\artiklo{kompendio}%
\vorto{kompendio} Rezumuro ek la ensemblo di cienco, di doktrino\lingui{DEFIRS}
\artiklo{kompensar}%
\vorto{kompensar}\fako{trans., per} Neutrigar (efekto) per efekto equivalanta, inversa-sinse\lingui{DEFIRS}
\artiklo{kompetenta}%
\vorto{kompetenta} A qua esas atribuenda la yuro decidar : 1. konseque de autoritato lego-grantita; 2. konseque de parkonoco, par-savo, de parposedo di la temo koncernata\lingui{DEFIRS}
\artiklo{kompilar}%
\vorto{kompilar}\fako{trans., ek, de} Asemblar aden un kolektajo, aden un korpo (texti pri un temo komuna, cherpita de fonti diversa)\lingui{DEFIRS}
\artiklo{komplemento}%
\vorto{komplemento} Lo adjuntenda a kozo por igar lu kompleta\lingui{DEFIRS}
\artiklo{kompleta}%
\vorto{kompleta} A qua mankas nulo ek la elementi ye qui lu konsistas normale\lingui{DEFIS}
\artiklo{kompletorio}%
\vorto{kompletorio}\fako{liturgio katol.} La lasta parto di la oficio, qua dicesas o kantesas pos la vespri
\artiklo{komplexa}%
\vorto{komplexa} Qua unionas en su plura elementi diversa\lingui{DEFIRS}
\artiklo{komplexiono}%
\vorto{komplexiono} Ensemblo ek la elementi ye qui konsistas la naturo korpala di individuo\lingui{EF}
\artiklo{komplezar}%
\vorto{komplezar}\fako{netrans.}\senco{1}{Plezar (ad ulu) per adaptar su a lua gusti, a lua sentimenti}\senco{2}{Inklinesar a perturbar su por plezar (ad ulu)}\lingui{EFIS}
\artiklo{komplico}%
\vorto{komplico} La persono qua helpas (ulu) kriminar, deliktar\lingui{DEFIS}
\artiklo{komplikar}%
\vorto{komplikar}\fako{trans., per} Inkombrar per elementi multa-sorta\lingui{DEFIS}
\artiklo{komplimentar}%
\vorto{komplimentar}\fako{trans., pri, pro}\senco{1}{Diskursar homaje di (ulu), respekt-exprese di (ulu), en okaziono solena}\senco{2}{Dicar (ad ulu) vorti qui laudas lu}\senco{3}{Dicar politeso-vorti}\lingui{DEFIRS}
\artiklo{komplotar}%
\vorto{komplotar}\fako{netrans.}\senco{1}{Projetar, kombine da plura personi, kontre la vivo, kontre la sekureso, di ulu}\senco{2}{\textit{(familiare)} Kombinar ulo kontre ulu}\lingui{DEFRS}
\artiklo{kompostar}%
\vorto{kompostar}\fako{trans.} Asemblar la tipi imprimala por *reproduktar la texto imprimenda\lingui{EFIS}
\artiklo{kompoto}%
\vorto{kompoto} Entremeso sukroza, ek frukti integra o quarima, koquita kun sukro, aquo o vino\lingui{DEFS}
\artiklo{kompozar}%
\vorto{kompozar}\fako{trans.}\senco{1}{Formacar (ensemblo) per la asemblo o la kombino di elementi diversa}\senco{2}{Facar verko skriptala, per kombinar olua elementi, olua parti}\lingui{DEFIS}
\artiklo{komprar}%
\vorto{komprar}\fako{trans., po, por, de} Aquirar po pekunio-quanto\lingui{IS}
\artiklo{komprenar}%
\vorto{komprenar}\fako{trans.} Embracar per la penso (la senco di ulo, la naturo di ulo, la motivo di ulo)\lingui{EFIS}
\artiklo{kompresar}%
\vorto{kompresar}\fako{trans.} Reduktar a volumino min granda, per presar o klemar\lingui{DEFIRS}
\artiklo{kompromisar}%
\vorto{kompromisar}\fako{trans.} Igar en situeso danjeroza (pri lua interesti, reputeso, digneso)\lingui{DEFIRS}
\artiklo{kompunciono}%
\vorto{kompunciono}\senco{1}{\textit{(teol. katol.)} Tristeso pia, qua devenas de la koncio di nia nedigneso e de la sufro ofensir Deo}\senco{2}{Mieno grava, rekolecala}\lingui{EFIS}
\artiklo{kompunda}%
\vorto{kompunda}\fako{tekn.} Dicesas pri ula organi od aparati qui esas « asociita », kombinita reciproke\lingui{EF}
\artiklo{komto}%
\vorto{komto}\senco{1}{En la epoki finala di la imperio Romana, e komence di la mez-epoko : oficiro di la suvereno-palaco, mayoro, guvernisto di teritorio-regiono}\senco{2}{Lor la feudismo : sinioro di feudo, qua esis, en la hierarkio, nemediate sub la dukio}\senco{3}{En la hierarkio ek la nobelo-tituli : la persono qua esas nemediate dop la markezo}\lingui{EFIS}
\artiklo{komuna}%
\vorto{komuna} Qua apartenas kune ad plura\lingui{DEFIRS}
\artiklo{komuniar}%
\vorto{komuniar}\fako{netrans.) (religio kristana} Recevar la sakramento « eukaristio »\lingui{DEFIS}
\artiklo{komunikar}%
\vorto{komunikar}\fako{trans., ad, kun}\senco{1}{Igar (kozo quan onu posedas) komuna ad altru per transmisar olu a lu}\senco{2}{Relateskar ideale, interestale e c. kun ulo}\senco{3}{Inter-relatar per paseyo}\lingui{DEFIS}
\artiklo{komuniono}%
\vorto{komuniono}\fako{religio} Unionitaro ek la personi qui profesas la sama deo-kredo\lingui{EFIS}
\artiklo{komutar}%
\vorto{komutar}\fako{trans.) (elektro} Cirkuligar od interruptar (elektro-fluo) o chanjar lua direciono\lingui{DEFIS}
\artiklo{koncentrar}%
\vorto{koncentrar}\fako{trans.}\senco{1}{Asemblar ad-cirkum centro komuna; asemblar to quo esas dispersita. \textit{(anke metaf.)}}\senco{2}{Retenar en spaco streta qua impedas la disperso}\lingui{DEFIRS}
\artiklo{konceptaklo}%
\vorto{konceptaklo}\fako{biol.} Quaza posho, qua kontenas la genitorgani di ula fungi
\artiklo{konceptar}%
\vorto{konceptar}\fako{trans.}\senco{1}{Formacar en su, per la fekundigo (la jermo di ento vivanta)}\senco{2}{\textit{(metaf.)} Formacar en sua spirito \textit{(ideo pri ulo)}; formacar en sua kordio (sentimento)}\lingui{DEFIRS}
\artiklo{koncernar}%
\vorto{koncernar}\fako{trans.} Relatar ad (ulu od ulo)\lingui{EFIS}
\artiklo{koncertar}%
\vorto{koncertar}\fako{netrans.} \textit{(dum *seanco muzikala)} exekutar ula quanto de kanto-peci o de muziko instrumentala\lingui{DEFIRS}
\artiklo{koncesar}%
\vorto{koncesar}\fako{trans.}\senco{1}{Abandonar a la uzado libera da ulu}\senco{2}{\textit{(metaf.)} Abandonar ad adverso, sen koaktesar a lo, un ek la punti diskutata}\lingui{DEFIRS}
\artiklo{koncesionar}%
\vorto{koncesionar}\fako{trans., ad) (aludante stato, urbo} Koncesar negratuite a privato, e lo esante valida dum tempo plu o min multa\lingui{DEFIS}
\artiklo{konciar}%
\vorto{konciar}\fako{trans.} Komprenar exakte la realaji\lingui{EFIS}
\artiklo{koncienco}%
\vorto{koncienco} Savo interna, da singlu, pri lo etiko-bona e lo etiko-mala\lingui{EFIS}
\artiklo{konciliar}%
\vorto{konciliar}\fako{trans.}\senco{1}{Adduktar ad interkonkordo pri litijo-punto}\senco{2}{Inter-harmoniigar kozi qui semblas interkontresar}\lingui{DEFIRS}
\artiklo{koncilo}%
\vorto{koncilo}\fako{eklezio katol.} Asemblitaro ek episkopi e teologio-doktori, por decidar pri ula problemi di doktrino, di diciplino ekleziala\lingui{DEFIRS}
\artiklo{konciza}%
\vorto{konciza}\fako{pri stilo} Sen to quo ne esas necesa por la senco\lingui{DEFIS}
\artiklo{kondamnar}%
\vorto{kondamnar}\fako{trans.}\senco{1}{Deklarar kulpoza, per verdikto : 1. frapar (ulu) per judicio qua deklaras ke lu kulpis per kriminir, deliktir, detrimentir; 2. \textit{(metaf.)}\textit{(pri libri, letri, e c.)} Koaktar ad ulo puniso}\senco{2}{Deklarar ke (ulu) esas kulpoza, blaminda, reprimandinda}\senco{3}{Deklarar ke (ulo) ne plus esas uzebla}\senco{4}{Deklarar kom ne plus uzenda (pordo, fenestro, quan onu klozas tale ke lu ne plus esas apertebla.)}\lingui{EFIS}
\artiklo{kondensar}%
\vorto{kondensar}\fako{trans.} Igar (gaso, vaporo) plu densa, per la pluproximigo di lua molekuli\lingui{DEFIS}
\artiklo{kondensatoro}%
\vorto{kondensatoro}\fako{elektro} Aparato quan onu uzas por akumular, sur surfaco, elektro pozitiva o negativa
\artiklo{kondicionar}%
\vorto{kondicionar}\fako{trans., per} Submisar a klauzo obliganta de qua dependas la valideso (di akto, e c.)\lingui{EFIS}
\artiklo{kondilo}%
\vorto{kondilo}\fako{anat.} Saliajo artikala, osta, rondigita ca-facie – platigita, ta-facie\lingui{EF}
\artiklo{kondilono}%
\vorto{kondilono}\fako{patol.} Surkreskajo morbala en la regiono di la anuso o di la perineo
\artiklo{kondimento}%
\vorto{kondimento} Substanco quan onu uzas por plufortigar la saporo di nutrivi, sive lor lia manjeso an la repasto-tablo, sive lor lia preparo\lingui{DEFIS}
\artiklo{kondolar}%
\vorto{kondolar}\fako{trans., pri, pro} Expresar sua partopreno di trauro, di desfortuno, qui afliktas (ulu)\lingui{DEFS}
\artiklo{kondominiar}%
\vorto{kondominiar}\fako{trans.} Kunposedar o kundominacar (lando)\lingui{EF}
\artiklo{kondoro}%
\vorto{kondoro}\fako{zool.} Granda vulturo, en la Andi\latina{sarcorhamphus gryphus}\lingui{DEFIRS}
\artiklo{konduktar}%
\vorto{konduktar}\fako{trans.) (fiziko} Propagar de ca a ta punto la manifesti elektrala, kalorala\lingui{EFIRS}
\artiklo{konduktoro}%
\vorto{konduktoro} \textit{(en vehilo publika : tramveturo, autobuso, treno)} La komizo di qua la rolo esas relatar kun la voyajeri, e c.\lingui{DEFRS}
\artiklo{kondutar}%
\vorto{kondutar}\fako{netrans., ad} Agar, en ta o ca cirkonstanco, en sua vivo, ye ta o ca maniero\lingui{EFIS}
\artiklo{konektar}%
\vorto{konektar}\fako{trans., kun} Ligar strete kun kozo sama-natura. – \textit{(elektro)}. Inter-relatigar, per elektro-konduktili, mashini od organi\lingui{EF}
\artiklo{konestablo}%
\vorto{konestablo}\fako{olim}\senco{1}{Oficiro di la reji di Francia}\senco{2}{Chef-mayoro di la armei rejala}\lingui{DEFIS}
\artiklo{konexo}%
\vorto{konexo}\fako{matem.} Equaciono algebrala homogena inter x, y, z e – x, y, z esante la koordinati homogena di punto — la koordinati homogena di rekto
\artiklo{konfecionar}%
\vorto{konfecionar}\fako{trans.} Facar (vesti) granda-serio, segun mezuri ne individuala, ma generala\lingui{DF}
\artiklo{konfekto}%
\vorto{konfekto} Bonbono qua su dissolvas en la boko, e qua kontenas liquoro o pasto sukroza e parfumoza\lingui{DEIS}
\artiklo{konferar}%
\vorto{konferar}\fako{netrans., kun} Konversar (kun ulu) quan onu konsultas pri temo konvencionita\lingui{DEFIRS}
\artiklo{konfervo}%
\vorto{konfervo}\fako{biol.} Speco di algo kun filamenti tubatra\latina{conferva}\lingui{DEFISL}
\artiklo{konfesar}%
\vorto{konfesar}\fako{trans., ad}\senco{1}{Agnoskar kom da su ulo quon onu agis, partikulare ulo etiko-mala}\senco{2}{\textit{(teol. katol.)} Deklarar volunte (sua peki) koram la tribunalo di penitenco}\lingui{EFIS}
\artiklo{konfesiono}%
\vorto{konfesiono}\fako{religio kristana} Deklaro koram publiko (*publico) pri sua acepto di la kredendaji di ta o ca eklezio\lingui{DEFIS}
\artiklo{konfetti}%
\vorto{« konfetti »} Drajei, buleti de gipsajo, dina rondaji de papero koloroza, quin onu interlansas dum karnavalo\lingui{FI}
\artiklo{konfidar}%
\vorto{konfidar}\fako{trans., ad ulu} Donar sekure (ulu od ulo) ad ulu quan onu komisas por sorgar lu\lingui{EFIS}
\artiklo{konfidencar}%
\vorto{konfidencar}\fako{trans., ad} Konfidar (ad ulu) ulo ne-materiala quon la konfidario devas ne divulgor – pri quo ta konfidario devas tacar (sekretajo)\lingui{DEFIRS}
\artiklo{konfirmar}%
\vorto{konfirmar}\fako{trans.}\senco{1}{Igar (ulu) mem plu ferma, plu certa (*certena) (pri). Igar plu certa, plu nedubitebla, per pruvo}\senco{2}{Adjuntar a la vigoro quan (lu) ja posedas, afirmon, pruvon, e c.}\senco{3}{\textit{(religio)} Donar ad (ulu) ta ek la sep sakramenti di la eklezio katolika, qua havas kom skopo igar la Santa Spirito decensar adsur la kristano por igar plu ferma la gracio de bapto}\senco{4}{\textit{(komerco)} Repetar (la texto di sua letro, telegramo, e c.) por ke la destinario esez mem plu *certena (certa) pri lo afirmita}\senco{5}{\textit{(pri informo qua ne esis absolute certa)} Afirmar ke ta informo esas nun certa, nedubitebla}\lingui{DEFIRS}
\artiklo{konfiskar}%
\vorto{konfiskar}\fako{trans.}\senco{1}{Atribuar a fisko, per dekreto, per lego (to quon proprietas ulu)}\senco{2}{Embargar ulo quan interdiktas la reguli di la skolo, di la liceo, e c.}\lingui{DEFIRS}
\artiklo{konfitar}%
\vorto{konfitar}\fako{trans.} Preparar (frukti) per igar li sejornar en liquoro, en siropo, qua penetras en li, qua konservas li\lingui{DFIS}
\artiklo{konfliktar}%
\vorto{konfliktar}\fako{netrans., kun}\senco{1}{Esar en interlukto, interkombato \textit{(aludante personi qui interbatas)}}\senco{2}{\textit{(metaf.)} \textit{(pri la interesti, la pasioni)}}\senco{3}{Interkontestar, (aludante \textit{du korpi politikala, judiciala, e c., singla de qui volas havar} \textit{kom suan, yuro, resortiso, fako, e c}.)}\lingui{DEFIRS}
\artiklo{konfluanta}%
\vorto{konfluanta} Di qua la elementi interjuntas. – (medic.) Pustuleti konfluanta : tante interproxima ke li intertushas ed intermixesas.\subvorto{bot.}{folii konfluanta}{Di qui la stipiti interjun tas ed intermixas}{}{}
\artiklo{konforma}%
\vorto{konforma}\fako{ad}\senco{1}{Di qua la formo koincidas precize kun olta di altra objekto quan onu egardas kom tipo}\senco{2}{Qua koincidas precize kun to quo imperesas da ulu, o kun to quo igesas necesa da ulo}\lingui{DEFIS}
\artiklo{konfrontar}%
\vorto{konfrontar}\fako{trans., ad}\senco{1}{Situar facio vers facio (du o plura personi da qui la afirmi esas interkontrea) por dicernar lo exakta}\senco{2}{Inter-apud-pozar (texti, dokumenti) por interkomparar li}\lingui{DEFIS}
\artiklo{konfundar}%
\vorto{konfundar}\fako{trans., ad, kun}\subvorto{}{konfundar kun}{Facar pelmelo pro distingo mentala qua esas mala o nesuficanta, e mixas ici kun iti (kom ex. « Il konfundas omno en sua odio, la kulpozi kun l'inocenti, la benigni kun la maligni »)}{1}{}\senco{2}{\textit{Konfundar ad} Transportar (mentale) erore la qualesi o la individueso di ulu od ulo, ad altra (kom ex. : « Vu konfundis manekino ad homo. Ne konfundez lampiro ad lanterno »)}\lingui{DEFIRS}
\artiklo{konfuza}%
\vorto{konfuza} Di qua la elementi esas intermixita tale ke li cesis esar dicernebla e distingebla\lingui{DEFIRS}
\artiklo{konglomerato}%
\vorto{konglomerato}\fako{mineral.} Bloko de fragmenti minerala aglomerita\lingui{DEFIRS}
\artiklo{kongregaciono}%
\vorto{kongregaciono}\senco{1}{(a) Kompanio de sacerdoti qui su obligis per vovi, per regulo komuna, agnoskita da la Santa Siejo; (b) Kompanio de kleriki qui, aganta segun la regulo komuna di la ordeno a qua li apartenas, su obligas per regulo specala}\senco{2}{Asociitaro de nekleriki, qui interunionas sub la nomo di santo, por la praktikado di pieso}\senco{3}{Komitato de kardinali, de kleriki, establisita da papo, en Roma, por studiar ula aferi}\lingui{DEFIRS}
\artiklo{kongresar}%
\vorto{kongresar}\fako{netrans.}\senco{1}{\textit{(aludante ciencisti, publicisti, politikisti, e c.)} *Seancar por la diskuto di temi ciencala, literaturala, politikala, e c.}\senco{2}{\textit{(aludante ministri de nacioni diversa, plenipotenciera)} *Seancar por diskutar ula temi internaciona}\lingui{DEFIRS}
\artiklo{kongro}%
\vorto{kongro}\fako{zool.} Marfisho ek la sama familio kam la mureno\latina{muraena conger}\lingui{EFIS}
\artiklo{kongruar}%
\vorto{kongruar}\fako{trans.}\senco{1}{Qua fitas remarkinde-bone ta o ca cirkonstanco, ta o ca kazo}\senco{2}{\textit{(matem.)} Nombri kongruanta : di qui la interdifero esas multoplo di altra nombro (modulo)}\lingui{DEFIS}
\artiklo{konidio}%
\vorto{konidio}\fako{bot.} Sporo di fungo, qua naskis sur organo sporifera acesora, che speci qui posedas altra reprodukt-organo
\artiklo{koniko}%
\vorto{koniko}\senco{1}{\textit{(geom.)} Singla de la kurvi (cirklo, elipso, hiperbolo, parabolo) quin onu obtenas per la seciono di kono per plani diverse situita relate la axo e la genitanto}\senco{2}{\textit{(geom. anat.)} Seciona plano di kono duesma-grada, e di qua la equaciono esas : Ax² + 2 Bxy + Cy² + 2 D x + 2 Ey + F = zero}\lingui{DEFIRS}
\artiklo{konio}%
\vorto{konio} Peco de ligno, de fero, angul-forma, quan onu stekas aden bloko de ligno por fendar ica; o quan onu lokizas (1) sub kanono, sub bombardilo, por elevar olu, o (2) sub la bazo di objekto por mantenar lu aplombe, e c.\lingui{F}
\artiklo{konivencar}%
\vorto{konivencar}\fako{netrans., kun}\senco{1}{Divenar, per klozar sua okuli por ne vidar, per silencar – quaze komplico etikala (pri la kulpo da ulu)}\senco{2}{Konsentar, konkordar, sekrete}\lingui{DEFIS}
\artiklo{konjektar}%
\vorto{konjektar}\fako{trans., de} Kredar per supozo agita por explikar qua-maniero fakto eventis od eventas, od eventos\lingui{DEFIRS}
\artiklo{konjelar}%
\vorto{konjelar}\fako{trans.} Transirigar aden la stando di solido (liquidon) per deprenar parto di lua kaloriko latenta\lingui{EFIS}
\artiklo{konjestionar}%
\vorto{konjestionar}\fako{trans.) (patol.} Igar la sango, acidente, fluar ecese aden la vaskuli di organo\lingui{DEFIS}
\artiklo{konjugar}%
\vorto{konjugar}\fako{trans.}\subvorto{gram.}{konjugar verbo}{Juntar intersucedante a la radiko di la verbo la afixi di qui la rolo esas expresar la varii di la personi, di la nombri, di la tempi, di la voci}{1}{}\subvorto{matem.}{imaginari konjugata}{Dicesas pri du imaginari a + b i e \textit{a' + b' i}, kande li havas la sama modulo}{2}{}\senco{3}{\textit{(bot.)} Dicesas pri la folii qui havas, sur petiolo komuna, un o plura pari de folieti inter-opozanta}\lingui{DEFIRS}
\artiklo{konjunciono}%
\vorto{konjunciono}\senco{1}{\textit{(gram.)} Vorto-speco qua juntas du propozicioni}\senco{2}{\textit{(astron.)} Situeso di du astri relate triesma, kande rekto, departanta de la centro di olca, pasas tra la centro di la du altri}\lingui{DEFIRS}
\artiklo{konjuntivito}%
\vorto{konjuntivito}\fako{patol.} Inflamuro di la konjuntivo di la okulo\lingui{DEFIS}
\artiklo{konjuntivo}%
\vorto{konjuntivo}\fako{anat.} Membrano mukoza qua kovras (quale tapeto kovras muro, interne di domo) la globo di la okulo ed unionas olca a la palpebro\lingui{DEFIS}
\artiklo{konjunturo}%
\vorto{konjunturo} Stando qua rezultas de la inter-kunefiko hazardala di eventaji, di cirkonstanci\lingui{DEFIS}
\artiklo{konjurar}%
\vorto{konjurar}\fako{trans.} Eskartar, per pregi religiala o per praktiki magiale (la influo da potenti maligna, nociva)\lingui{EFIS}
\artiklo{konkava}%
\vorto{konkava} Qua karakterizesas da kurveso sferala kava\lingui{DEFIS}
\artiklo{konklavo}%
\vorto{konklavo}\fako{religio katol.} Loko ube su inkluzas klefe la kardinali pos la morto di papo, por exekutar la elekto di lua sucedonto\lingui{DEFIS}
\artiklo{konkluzar}%
\vorto{konkluzar}\fako{trans.}\senco{1}{Finar, o finigar, per solvo definitiva}\senco{2}{\textit{(logiko)} Deduktar la konsequi de la premisi pozita}\lingui{EFIRS}
\artiklo{konko}%
\vorto{konko}\fako{zool.} Envelopilo kalka di ula moluski\lingui{EFIRS}
\artiklo{konkoido}%
\vorto{konkoido}\fako{geom.} Kurvo ek du branchi infinita, qui havas kom asimptoto rekto fixa, e simetra relate la perpendiklo trasita de punto fixa a ta rekto\lingui{DEFIS}
\artiklo{konkologio}%
\vorto{konkologio} La parto di zoologio qua koncernas la konko-moluski\lingui{EFIS}
\artiklo{konkordar}%
\vorto{konkordar}\fako{netrans., kun, pri, pro ke}\senco{1}{\textit{(pri plura personi)} Opinionar, sentar, intersame}\senco{2}{\textit{(pri kozi)} Inter-relatar harmonioze pro la existo inter li di raporto exakta}\lingui{DEFIRS}
\artiklo{konkordato}%
\vorto{konkordato}\senco{1}{Interkonsento da la papo kun suvereno, pri la yuri rispektiva di la Stato e di la Eklezio}\senco{2}{Kontrato per qua komercisto qua ne plus povas pagar sua debaji obtenas, de sua debarii, pago-fristo o la diminuto di parto de lua debajo}\lingui{DEFIRS}
\artiklo{konkreciono}%
\vorto{konkreciono}\fako{patol.} Agregajo solida, qua naskas acidente en ula tisui\lingui{DEFIS}
\artiklo{konkreta}%
\vorto{konkreta} Qua expresas ulo reala\lingui{DEFIRS}
\artiklo{konkubinato}%
\vorto{konkubinato}\fako{che la Romani, antique} Sorto di mariajo, inferiora ye la mariajo civila – e la unika pri qua darfis kontratar la ne-Romani\lingui{DEFIS}
\artiklo{konkubo}%
\vorto{konkubo}\senco{1}{Singla de la du personi qui kun-vivas spozatre, -quankam ne intermariajita}\senco{2}{\textit{(che la Romani, antique)} Singla de la personi qui interunionis per la konkubinato}\lingui{DEFIS}
\artiklo{konkupicenco}%
\vorto{konkupicenco}\fako{en la linguo kristanismala} Inklineso, da persono dekadinta, ad la plezuri sensala (precipue la sexuala)\lingui{DEFIS}
\artiklo{konkurencar}%
\vorto{konkurencar}\fako{trans.} \textit{(aludante personi qui persequas la sama skopo)} Inter-rivalesar pri sua interesti (precipue pri la personi di la sama komerco-fako, industrio-fako)\lingui{DFIRS}
\artiklo{konkursar}%
\vorto{konkursar}\fako{netrans., pri} Esar en la rangi kun altri, pri obtenor premio, nominesor, e c.\lingui{DFIRS}
\artiklo{konkusionar}%
\vorto{konkusionar}\fako{netrans.) (aludante stat-oficiero} Profitar (sen obtenir permiso pri o por lo) per la autoritato quan donas a lu lua ofico\lingui{DFIS}
\artiklo{kono}%
\vorto{kono}\senco{1}{\textit{(geom.)} Surfaco genitata da rekto moviva, qua pasas ye punto fixa (« somito ») e su apogas a kurvo fixa (« direktanto »). Solido formacita da ta surfaco restriktata ad seciono facita perpendikle ye olua axo}\senco{2}{Objekto konatra}\senco{3}{\textit{(bot.)} Frukto di ula arbori (« koniferi ») qua konsistas ye squami imbrikita cirkum axo komuna}\lingui{DEFIRS}
\artiklo{konocar}%
\vorto{konocar}\fako{trans.} Havar ideo plu o min kompleta (pri ulo, pri ulu)\lingui{FIS}
\artiklo{konoido}%
\vorto{konoido}\fako{matem.} Solido quan formacas la rotaco di koniko cirkum sua axo\lingui{DEFIRS}
\artiklo{konosmento}%
\vorto{konosmento}\fako{yuro-cienco navigala} Agnosko-dokumento pri vari embarkita e pri lia transporto-kondicioni – donata, a la armatoro di komerco-navo, da la navestro\lingui{DFRS}
\artiklo{konquasanta}%
\vorto{konquasanta}\fako{medic.} Qua quaze ruptas la muskuli, qua korpo-depresas, qua destruktas la vigoro. (Dicesas precipue pri la granda dolori parturala.)\lingui{DEFIRS}
\artiklo{konquestar}%
\vorto{konquestar}\fako{trans., per} Kaptar, sizar per violento (lando, provinco, e c.)\lingui{EFIS}
\artiklo{konsakrar}%
\vorto{konsakrar}\fako{trans., ad} Donar unike por persono, por ulo. (ex. : konsakrar sua pekunion, sua tempo, sua sorgi, ad ulu, ad ulo; konsakrar sua vivo a lernado)\lingui{EFIS}
\artiklo{konsanguina}%
\vorto{konsanguina} Qua naskis de la sama patrulo, ma ne de la sama patrino\lingui{F}
\artiklo{konsekracar}%
\vorto{konsekracar}\fako{trans.) (liturgio katol.}\senco{1}{Igar sakra (per konsakrar lu a la servo di Deo)}\senco{2}{Konsekracar hostio \textit{(aludante la sacerdoto qua celebras meso)} : Transformar la pano e la vino aden la korpo e la sango di Iesu Kristo}\lingui{DEFIS}
\artiklo{*konselio}%
\vorto{*konselio} Asemblitaro de personi komisita oficale pri guidar, per lia konsili, la chefo di stato, di administreri, di armeo, e c.; o pri kontrolar la administro, la jero, di kompanio industriala, komercala, financala\lingui{EFIS}
\artiklo{konsentar}%
\vorto{konsentar}\fako{trans.} Aceptar propra-vole e sua-parte (la opiniono o la satisfaco di la deziro da ulu)\lingui{DEFIRS}
\artiklo{konsequar}%
\vorto{konsequar}\fako{netrans., de, ek) (aludante principo enuncita} Eventar logikale pos (lu)\lingui{DEFIRS}
\artiklo{konservar}%
\vorto{konservar}\fako{trans.}\senco{1}{Sorgar (ulu, ulo) tale ke lu ne alteresas o destruktesas}\senco{2}{Ne lasar disipesar, destruktesar, desaparar}\lingui{DEFIRS}
\artiklo{konservatorio}%
\vorto{konservatorio} Establisuro publika, en qua onu konservas ula kolektaji\lingui{DEFIS}
\artiklo{konsiderar}%
\vorto{konsiderar}\fako{trans.} Atencar sorgoze, minucioze (ta o ca motivo decidigiva, rezolvigiva) por agar saje\lingui{DEFIRS}
\artiklo{konsignar}%
\vorto{konsignar}\fako{trans.} Donar (kozo konservenda kom garantiivo, kom depozajo) : (1) pekunio-quanto, ye la kazo di tribunalo; (2) vari, depoze che komisionisto, komercisto – bagaji depoze en relvoyo-staciono\lingui{DEFIRS}
\artiklo{konsilar}%
\vorto{konsilar}\fako{trans., ad} Guidar (ulu) per dicar a lu to, quon lu devas agar o facar se lu volas agar saje\lingui{EFIRS}
\artiklo{konsistar}%
\vorto{konsistar}\fako{netrans., ye, ek} Esar formacita (ek ula elementi)\lingui{EFIS}
\artiklo{konsistorio}%
\vorto{konsistorio}\senco{1}{\textit{(che la katoliki)} Asemblitaro de kardinali, prezidata da la papo, qua su okupas per la aferi generala di la eklezio katolika}\senco{2}{\textit{(che la protestantisti e la Izraelidi)} Asemblitaro ek kulto-oficieri e laiki, elektita por direktar la aferi di la kongregaciono}\lingui{DEFIRS}
\artiklo{konskriptar}%
\vorto{konskriptar}\fako{trans.} En-armeigar koakte la adolecanti qui jus atingis la evo-punto quan la lego determinis pri la milito-kapabligo\lingui{DEFIRS}
\artiklo{konsolacar}%
\vorto{konsolacar}\fako{trans., pri} Alejar la chagreno di (ulu)\lingui{EFIS}
\artiklo{konsoldo}%
\vorto{konsoldo}\fako{bot.} Planto qua formacas genero di la familio « boraginei », di qua la radiko uzesis olim kom astriktivo\latina{symphytum}
\artiklo{konsolidar}%
\vorto{konsolidar}\fako{trans., per} \textit{(financo)} \textit{(aludante la debajo publika)} Transformar la debajo rimborsebla di stato, aden debajo perpetua di qua nur la rento esos postulebla\lingui{DEFIRS}
\artiklo{konsolo}%
\vorto{konsolo}\senco{1}{\textit{(arkitekt.)} Saliajo fixigita an murego, por suportar balkono, kornico, e c., o kom piedestalo por statuo}\senco{2}{Orno-moblo pozita an muro, destinita a subtenar vazi, kupi, statueti, e c.}\lingui{DEF}
\artiklo{konsomeo}%
\vorto{konsomeo} Buliono qua, pro koqueso longa-tempa, absorbis la tota suko di la karno\lingui{DEFIS}
\artiklo{konsonancar}%
\vorto{konsonancar}\fako{netrans.}\senco{1}{\textit{(aludante foni muzikala)} Interakordar plu o min komplete}\senco{2}{\textit{(aludante silabi, e precipue finali di qui la termino-foni esas sama)} -Interproximesar}\lingui{DEFIRS}
\artiklo{konsonanto}%
\vorto{konsonanto} Artikulo-maniero di la fono, qua varias segun la movi da la lango o da la labii. (ex. : b, t, f esas konsonanti, kontre ke a, e, i esas vokali)\lingui{DEFIRS}
\artiklo{konsorto}%
\vorto{konsorto}\senco{1}{La persono di qua la estado sociala esas komuna kun olta di sua spozo. (Dicesas precipue pri la spozulo di la rejino di Britania.)}\senco{2}{Singla de la personi qui, en afero, havas interesto komuna}\lingui{DEFIRS}
\artiklo{konspirar}%
\vorto{konspirar}\fako{netrans., kontre, kun} Konvencionar sekrete por renversar la guvernistaro instalita, nuna\lingui{DEFIS}
\artiklo{konsputar}%
\vorto{konsputar}\fako{trans.) (aludante plura personi qui agas lo} Jetar (kontre ulu) desestimo-klami\lingui{F}
\artiklo{konstanta}%
\vorto{konstanta} Qua ne cesas esar sama.\lingui{DEFIS}\subvorto{matem.}{konstanto}{Quanto qua permanas sama, kompare ad altrai qui varias}{}{}
\artiklo{konstatar}%
\vorto{konstatar}\fako{trans.} Igar su *certena (certa) (per observo, verifiko) pri la exakteso di fakto\lingui{DF}
\artiklo{konsternar}%
\vorto{konsternar}\fako{trans.) (aludante, generale, katastrofo} Opresar, depresar (la psiko)\lingui{DEFIS}
\artiklo{konstipar}%
\vorto{konstipar}\fako{trans.} Impedar la kako, pro konstrikto intestinala\lingui{EFIRS}
\artiklo{konstitucar}%
\vorto{konstitucar}\fako{trans.}\senco{1}{Instalar (ulu) en estado legala}\senco{2}{Organizar esencale (ulo)}\senco{3}{Esar la elementi qui formacas}\lingui{DEFIRS}
\artiklo{konstituciono}%
\vorto{konstituciono} Estado anatomiala e fiziologiala di organismo vivoza\lingui{DEFIRS}
\artiklo{konstriktar}%
\vorto{konstriktar}\fako{trans.) (anat.} Kompresar, per klemar cirkume\lingui{DEFIRS}
\artiklo{konstruktar}%
\vorto{konstruktar}\fako{trans.}\senco{1}{Facar (domo, edifico, ponto, navo, mashino, e c.) segun programo, skisuro, desegnuro}\senco{2}{\textit{(pri frazo)} Ordinar la vorti}\senco{3}{Trasar (la ensemblo di la linei di mapo, di figuro geometriala)}\lingui{DEFIRS}
\artiklo{konsulo}%
\vorto{konsulo}\senco{1}{\textit{(lor la Roma-republiko antiqua)} Stat-oficiero, elektita singla-yare da la plebeyo, e qua partoprenis, kun un kolego, la autoritato lego-exekutigala}\senco{2}{\textit{(lor la Roma-imperio antiqua)} Stat-oficiero nominita da la senato, ma di qua la rolo esis nur honorizala}\senco{3}{Agento diplomacala, komisita da guvernistaro por protektor, en lando stranjera, la interesti di sua samalandani}\lingui{DEFIRS}
\artiklo{konsultar}%
\vorto{konsultar}\fako{trans., pri, por}\senco{1}{Questionar (ulu) por savar lua opiniono pri lo rezolvenda}\senco{2}{Serchar informeso (en libro, en spegulo)}\lingui{DEFIRS}
\artiklo{konsumar}%
\vorto{konsumar}\fako{trans.}\senco{1}{Destruktar gradoze (ulo) en lua substanco}\senco{2}{Igar (ulu) deperisar gradoze}\senco{3}{Prenar gradoze de olua substanco per uzar lu. (ex. konsumar vesto)}\lingui{DEFIRS}
\artiklo{kontagiar}%
\vorto{kontagiar}\fako{trans., ad} Transmisar, per kontakto, ulo nociva (anke la morbi qui venas de mikrobo, de miasmi)\lingui{DEFIS}
\artiklo{kontaktar}%
\vorto{kontaktar}\fako{trans.} Havar, kun altra korpo, du o plura surfacopunti komuna. \textit{(Anke metaf.)}\lingui{DEFIRS}
\artiklo{kontaminar}%
\vorto{kontaminar}\fako{trans., ye} Infektar ye morbo kontagiala\lingui{DEFIS}
\artiklo{kontar}%
\vorto{kontar}\fako{trans.) (aritm.} Determinar la relato da quanto ad unajo\lingui{DEFIRS}
\artiklo{kontemplar}%
\vorto{kontemplar}\fako{trans.} Regardar, par-sinkante en la vido di la objekto\lingui{EFIS}
\artiklo{kontenar}%
\vorto{kontenar}\fako{trans.}\senco{1}{Havar en su}\senco{2}{Povar havar en su}\lingui{EFIS}
\artiklo{kontenta}%
\vorto{kontenta} Qua ne deziras ulo plusa, ulo plubona, kam to quon lu havas\lingui{EFIS}
\artiklo{kontestar}%
\vorto{kontestar}\fako{trans.} Diskutar negante la yusteso di to quon ulu revendikas, o la justeso di to quon ulu afirmas\lingui{EFIS}
\artiklo{kontigua}%
\vorto{kontigua}\fako{aludante tereno, domo, edifico} Qua jacas an (altra tereno, domo, edifico)\lingui{DEFIS}
\artiklo{kontinencar}%
\vorto{kontinencar}\fako{netrans.} Abstenar de koito, de amor-plezuri\lingui{EFIS}
\artiklo{kontinento}%
\vorto{kontinento} Singla de la granda dividuri di la Terglobo quin la oceani separas, ica de ita, e quin onu povas parkurar tota-extense sen trairar maro; Europa, Azia, Afrika, Amerika, Oceania\lingui{DEFIRS}
\artiklo{kontingento}%
\vorto{kontingento}\senco{1}{\textit{(armeo)} Quanto de armeani qua esas furnisenda singla-yare a la rekruto-autoritatozi}\lingui{DEFIRS}\subvorto{financo}{moneto-kontingento}{Proporciono di la naturi diversa di monet-uni quin la Monetiferio devas fabrikar.}{2}{}\subvorto{filoz.}{la kontingenti}{To quo povas eventar, ma ne eventos necese}{}{}
\artiklo{kontinua}%
\vorto{kontinua} Qua konsistas ek elementi qui formacas serio ne-interruptita en spaco od en tempo\lingui{DEFIS}
\artiklo{kontoro}%
\vorto{kontoro}\senco{1}{Loko ube funcionas la oficieri e lia subordiniti qui administras, jeras, kom employati di la guvernistaro, di firmo komercala, industriala, di banko}\senco{2}{Loko ube asemblas su la membri di komisitaro, di komitato, exter la chambrego ube eventas la *seanci generala}\lingui{DEFR}
\artiklo{kontrabandar}%
\vorto{kontrabandar}\fako{trans.} Importacar celite, sekrete, vari qui esas interdiktita, e pri qui onu ne pagis a la doganeyo la imposto-pekunio\lingui{DEFIRS}
\artiklo{kontrabaso}%
\vorto{kontrabaso}\fako{muziko} La maxim fonogamo-basa ek la arketinstrumenti — plu granda kam violoncelo e samaforma, kun tri kordi\lingui{DEFIRS}
\artiklo{kontradanso}%
\vorto{kontradanso} Salono-danso, ek kin parti o figuri, exekutata sucedante da du o quar grupi qui stacas inter-regarde\lingui{DEFIS}
\artiklo{kontrafaktar}%
\vorto{kontrafaktar}\fako{trans.}\senco{1}{Imitar o *reproduktar kontre-lege (la verko di ulu, e lua-detrimente)}\senco{2}{Imitar fraude objekto di qua la karaktero esas autentika}\lingui{FI}
\artiklo{kontraforto}%
\vorto{kontraforto}\fako{arkitekt.} Pilastro, saliajo masonura, sur qua apogas su muro qua suportas kargajo\lingui{EFIRS}
\artiklo{kontraktar}%
\vorto{kontraktar}\fako{trans.}\senco{1}{\textit{(gram.)} Reduktar (du vokali o du silabi) ad un}\senco{2}{Reduktar a volumino min granda, sen diminuteso di la maso : (1) \textit{(anat.)} per la kurtigo di la fibri; (2) \textit{(fiziko)} per la pluapudigo di la molekuli}\lingui{DEFIS}
\artiklo{kontrakturo}%
\vorto{kontrakturo}\fako{patol.} Stretiguro, kun rigidesko, qua eventas en muskulo pos konvulso, nevralgio, rumatismi, e c.\lingui{DF}
\artiklo{kontralto}%
\vorto{kontralto}\fako{muziko} Che la mulieri, la voco maxim basa\lingui{DEFIRS}
\artiklo{kontrapunto}%
\vorto{kontrapunto}\fako{muziko} Formo di muziko klasika, en qua onu asemblas, cirkum frazo precipua quan onu selektis kom temo, parti acesora qui interunionas interseque en kombinuri harmoniala determinita, noto kontre noto\lingui{DEFS}
\artiklo{kontrastar}%
\vorto{kontrastar}\fako{netrans., kun) (aludante du kozi} Interopozesar tale ke lia interkompareso igas ica saliar\lingui{DEFIRS}
\artiklo{kontratar}%
\vorto{kontratar}\fako{netrans., pri}\senco{1}{Asumar, ad ulu (obligeso)}\senco{2}{Debatar, diskutar konvenciono solena inter du stati, determinante, klauzope, to quo esos stipulita da singla de la partopreninti}\senco{3}{\textit{(aludante plura personi)} Interobligar per akto autentika}\lingui{DEFIRS}
\artiklo{kontravencar}%
\vorto{kontravencar}\fako{netrans., pri} Agar kontre la preskripti da regularo, da lego\lingui{DEFIS}
\artiklo{kontre}%
\vorto{kontre} Opoze ad (ulu, ulo)
\artiklo{kontreskarpo}%
\vorto{kontreskarpo}\fako{tekn., milit-arto} Parto di la fosato di fortifikuro qua orientizesas kontre la eskarpajo
\artiklo{kontributar}%
\vorto{kontributar}\fako{trans., ad} Kunlaborar, sua-parte, a verko komuna. Adportar sua quoto a spenso komuna\lingui{DEFIRS}
\artiklo{kontricar}%
\vorto{kontricar}\fako{netrans., pri) (religio} Repentar pri la peko, quan eventigas la chagreno ofensir Deo\lingui{DEFIS}
\artiklo{kontrolar}%
\vorto{kontrolar}\fako{trans.}\senco{1}{Submisar a verifiko da la administrantaro}\senco{2}{Submisar ad exameno minucioza}\lingui{DEFIRS}
\artiklo{kontroversar}%
\vorto{kontroversar}\fako{trans.} Diskutar kun ulu pri punto di doktrino\lingui{DEFIS}
\artiklo{kontumaca}%
\vorto{kontumaca}\fako{yuro-cienco} Persono qua, akuzite, ne su prizentis avan la tribunalo kriminala\lingui{DEFIS}
\artiklo{konturo}%
\vorto{konturo} Surfaco, lineo qua limitizas extera objekto, e, partikulare, objekto ronda\lingui{DEFIRS}
\artiklo{kontuzar}%
\vorto{kontuzar}\fako{trans., ye} Lezar per ta altero plu o min profunda quan shoko o frapo produktas en la tisui, sen lacero di la pelo\lingui{DEFIRS}
\artiklo{konvalecar}%
\vorto{konvalecar}\fako{netrans.} Esar en la stando di persono qua jus esis malada e komencas risaneskar, ribonestandar\lingui{DEFIS}
\artiklo{konvektar}%
\vorto{konvektar}\fako{netrans.) (fiziko} Dicesas pri la fenomeno qua eventas kande korpo varma plunjesas aden korpo liquida o gasa, determinante vera fluo en la fluido, konseque de la varmesko kontakte di la korpo varma\lingui{DEFIRS}
\artiklo{konvenar}%
\vorto{konvenar}\fako{netrans., ad}\senco{1}{Esar exakte to quon postulas la cirkonstanci, la stando di ulu, di ulo}\senco{2}{Esar exakte to quon onu expektas de la karaktero, de la stando sociala di ulu}\senco{3}{Esar exakte to quon postulas la deziro, la prizo di ulu}\lingui{EFIS}
\artiklo{konvencionar}%
\vorto{konvencionar}\fako{netrans., pri} Interkonkordar explicite (e, maxim-multa-kaze, skribe) pri punto determinita ek ulo agota o facota\lingui{DEFIRS}
\artiklo{konvento}%
\vorto{konvento} Asemblitaro, (maxim-multa-kaze sekreta) de framasoni. *generala\lingui{DEFIS}
\artiklo{konvergar}%
\vorto{konvergar}\fako{netrans., ad}\senco{1}{\textit{(fiziko)} \textit{(aludante direcioni, o linei qui reprezentas li)} Departante de punti diversa, irar ad un punto, interproximeskante}\senco{2}{\textit{(matem.)} \textit{(aludante quanto)} Proximeskar sempre plu multe a quanto fixa, sen ul-instante atingar lu}\lingui{DEFIS}
\artiklo{konversar}%
\vorto{konversar}\fako{netrans., kun, pri) (aludante personi qui esas kune} Kambiar idei kun ulu\lingui{DEFIRS}
\artiklo{konversionar}%
\vorto{konversionar}\fako{netrans.) (strategio} Agar movo de cirkumiro per qua armeo-korpo prizentas altra facio\lingui{F}
\artiklo{konvertar}%
\vorto{konvertar}\fako{trans., ad}\senco{1}{Duktar a kredar ed adoptar religio, od altra doktrino, kam olta quan lu havis}\senco{2}{Chanjar kozo ad altra}\senco{3}{\textit{(financo)} Kambiar rento-akto qua super-iris « para », po altra di qua la koto esas min supra, per ofrar la rimborso ye « para » a la posedanti qui refuzas aceptar la kambio}\lingui{DEFIRS}
\artiklo{konvertoro}%
\vorto{konvertoro}\senco{1}{\textit{(gisado)} Aparato por konvertar gisfero aden stalo — retorto metalo-tola, provizita ye briki refraktara, en qua onu submisas gisfero a la efiko da kaloro e da forta aerofluo}\senco{2}{\textit{(elektrologio)} Aparato per qua onu chanjas la direcioni di *korento elektrala (elektro-fluo) sen diplasar la guidfilo di la baterio}\lingui{DEFIRS}
\artiklo{konvexa}%
\vorto{konvexa} Qua prizentas kurviguro sferala reliefa\lingui{DEFIRS}
\artiklo{konvinkar}%
\vorto{konvinkar}\fako{trans., pri} Duktar (ulu) ad agnoskar ulo kom exakta\lingui{EFIS}
\artiklo{konvokar}%
\vorto{konvokar}\fako{trans., ad} Pregar (ulu) ke lu voluntez venar a la asembleyo di lua simili\lingui{EFIS}
\artiklo{konvolvulo}%
\vorto{konvolvulo}\fako{bot.} Planto herbatra od arboretatra, kun flori dispozita funelatre, di qua la precipua speci reptas sur sulo o volvas cirkum la planti vicina\latina{convolvulus}\lingui{EFIS}
\artiklo{konvoyar}%
\vorto{konvoyar}\fako{trans.} Transportar soldati, municioni, viktualii per serio de navi, de veturi, de kamioni qui avancas linee (quale la vagoni di treno)\lingui{DEFIS}
\artiklo{konvulsar}%
\vorto{konvulsar}\fako{netrans.}\senco{1}{Tordar su konseque de ta kontrakto subita di la muskuli qua akompanas ula standi morbala di la nervaro. Afekto qua koncernas la infanti, karakterizata da kontrakti di ta o ca parte di la korpo qui povas produktar morto per konjestiono cerebrala, asfixieso}\senco{2}{Subisar ta violenta trepido di la traiti di la vizajo, di la korpo, quan produktas ula emoci}\lingui{DEFIRS}
\artiklo{konyako}%
\vorto{konyako} Brandio qua esas la produkturo naturala de la vini rekoltata e distilata en la regiono qua cirkondas la Franca urbo Cognac \textit{(pron. Konyak')}\lingui{DEFIRS}
\artiklo{kooperar}%
\vorto{kooperar}\fako{netrans.) (ekonomiko} Dicesas pri la ago da ta laboristi qui, vice donar a patrono sua laboro po salariopekunio, komunigas sua sparita pekunio e sua laboro, por produktar e vendar – singla de li partoprenonte la profito o la desprofito. – Dicesas anke pri la personi qui kompras komune sua bezonaji engrose, tale obtenante profito-quanto – quale komercisto endetalista – quan li distributas inter su proporcione a la valoro pekuniala di sua kompri\lingui{EF}
\artiklo{kooptar}%
\vorto{kooptar}\fako{trans. por} Admisar eceptale (membri qui ne satisfacas omna kondicioni postulata), pro lua autoritato (savo, erudeso, talento) partikulara\lingui{DEF}
\artiklo{koordinar}%
\vorto{koordinar}\fako{trans.} Ordinar segun ula relati (la parti di ensemblo) por facar ek lu ensemblo\lingui{DEFIRS}
\artiklo{kopaivo}%
\vorto{kopaivo} Substanco terpentinatra quan onu extraktas per incizo de la kortico di ula arbusto e quan onu uzas medicinale kontre ula morbi di la organi genitala ed urinifala\latina{copaifera}\lingui{DEFISL}
\artiklo{kopalo}%
\vorto{kopalo} Rezino flavatra, solida, nesolvebla en alkoholo, quan onu extraktas per incizo ek arbori rezinifera di India e di Brazilia, ed ek qua onu facas vernisi tre prizata\lingui{DEFIRS}
\artiklo{kopiar}%
\vorto{kopiar}\fako{trans.} Reproduktar (la texto di skriburo, desegnuro, e c.) un-exemplere o plur-exemplere\lingui{DEFIRS}
\artiklo{koprolito}%
\vorto{koprolito} Exkremento petreskinta di animali fosila\lingui{DEF}
\artiklo{kopso}%
\vorto{kopso} Mikra bosko ek arbori qui kreskis trunko-stumpe o sprose, e quin onu rekortas intermite\lingui{E}
\artiklo{kopulacar}%
\vorto{kopulacar}\fako{netrans., kun) (aludante maskulo e femino ne-homa} Interunionar sexuale\lingui{DEFIS}
\artiklo{kopulo}%
\vorto{kopulo}\fako{gram.} Termino qua unionas la subjekto di la propoziciono a la atributo\lingui{DEFIS}
\artiklo{koquar}%
\vorto{koquar}\fako{trans., en, per} Igar apta por manjeso, per la efiko da fairo (per boliigo, rosto, grilo, e c.)\lingui{EFIRS}
\artiklo{koralio}%
\vorto{koralio} Produkturo kalkoza ek ula polipi, plu o min dazlante-reda, kun la formo di rami, maxim-multa-kaze fixigita an roki submarsurfaca\lingui{DEFIRS}
\artiklo{koram}%
\vorto{koram} « … asistante », « … esante prezenta »
\artiklo{korano}%
\vorto{korano}\fako{religio} La libro sakra di la Mohamedisti, qua kontenas la naraco pri la misiono di Mohamed e la ensemblo ek la preskripti quin ilu donis a sua populo\lingui{DEFIRS}
\artiklo{korbelo}%
\vorto{korbelo}\fako{arkitekt.} Quadro plu o min salianta, origine taliita bizele, ulakaze ornita per skulturi, e qua uzesas por subtenar arkizuro, kornico, e c.\lingui{EF}
\artiklo{korbo}%
\vorto{korbo} Recevuyo ek oziero, ek junko, e c., di qua la dimensioni e la formi esas diversa, uzata por kontenar o transportar provizuri, viktualii, e c.\lingui{DF}
\artiklo{kordio}%
\vorto{kordio}\senco{1}{\textit{(anat.)} Vicero muskula, di qua la formo esas olta di kono subversata, e qua esas la centro, la agento precipua di la cirkulo di sango}\senco{2}{\textit{(metaf.)} Ta organo egardata kom la rezideyo di la afecioni, di la sentimento interna, di la deziro, di la sufro, di la joyo, di la sentimento pri etiko, pri koncienco}\senco{3}{\textit{(karto-ludo)} La grupo de la karti sur qui reprezentesis kordio desegnita skise}\lingui{DEFIS}
\artiklo{kordo}%
\vorto{kordo}\senco{1}{Asemblajo ek fili tre grosa, intertordita, di qui la groseso e la longeso esas plu o min granda}\senco{2}{Torduro ek tripi od ek fili metala, uzata por la faco di ula muzik-instrumenti e quan onu vibrigas per la fingri (gitaro, harpo, e c.) o per arketo (violino, violoncelo, e c.), o per marteleti qui korespondas kun klavi (piano)}\senco{3}{\textit{(geom.)} Rekto qua klozas arko}\lingui{EFIS}
\artiklo{kordono}%
\vorto{kordono}\senco{1}{Kordeto qua esas un ek la elementi di kordo}\senco{2}{Mikra treso ronda o plata, rubando plu o min larja, uzata por rolo determinita}\senco{3}{\textit{(anat.)} Branchifuro di la nervi}\senco{4}{Lineo kontinua quan formacas kozo, serio de kozi, alonge spaco determinita}
\artiklo{koregono}%
\vorto{koregono}\fako{zool.} Genero di fisho qua vivas en aquo sen-sala, familio di la salmonidei, qua vivas en la granda lagi di la centro di Europa (Suisia, Germania, nordo-regioni di Italia)
\artiklo{korekta}%
\vorto{korekta} Qua ne deviacas de la regulo, de la normo\lingui{DEFIRS}
\artiklo{korelatar}%
\vorto{korelatar}\fako{trans.} Esar ligita ad altra termino tale ke un ek la du advokas nemediate la altra\lingui{DEFIRS}
\artiklo{*korento}%
\vorto{*korento}\fako{elektro} Direciono determinita di la fluido elektra, magnetala\lingui{EFIS}
\artiklo{koreo}%
\vorto{koreo}\fako{patol.} Morbo nervala qua produktas movi sukusoza\lingui{DEFS}
\artiklo{korespondar}%
\vorto{korespondar}\fako{netrans., kun, pri}\senco{1}{Esar reciproke kon-forma (1) pro interkonveno di sua naturo, di sua proporcioni; (2) pro interkonkordo di sentimenti}\senco{2}{Esar reciproke komunikanta : (1) per paseyo komuna; (2) per kambio di letri}\lingui{DEFIRS}
\artiklo{koriandro}%
\vorto{koriandro}\senco{1}{\textit{(bot.)} Planto aromata, qua formacas genero ek la familio « ombeliferi »}\senco{2}{La semino sika di ta planto, uzata en la koquo, en la farmacio-medikamenti}\latina{coriandrum}\lingui{DEFIS}
\artiklo{koribanto}%
\vorto{koribanto}\fako{epoki antiqua} Sacerdoto di la diino Cibelo\lingui{DEFIRS}
\artiklo{koridoro}%
\vorto{koridoro}\senco{1}{\textit{(olim)} Voyo tektoza cirkonde di fortifikajo}\senco{2}{\textit{(nia-epoke)} Paseyo alonge plura chambri di un etajo, di un apartmento, per qua onu povas forirar de oli. Paseyo streta uzata kom forir-voyo, por irar de parto di edifico ad altra parto}\lingui{DEFIRS}
\artiklo{korifeo}%
\vorto{korifeo}\senco{1}{Chefo di la koro, en la tragedio e la komedio antiqua}\senco{2}{Chefo di koro kantala o dansala, en la opero moderna}\senco{3}{\textit{(metaf.)} \textit{(kun nuanco di ironio)} La persono qua maxime salias, unesma-ranga, en ulo}\lingui{DEFIRS}
\artiklo{korimbo}%
\vorto{korimbo}\fako{bot.} Infloresenco di qua la pedunkli acesora, departante de punti interdiferanta, elevas la flori a nivelo unika, tale ke formacesas quaza parasuno di qua la radii esas neinteregala\lingui{EFIS}
\artiklo{korinto}%
\vorto{korinto} Bitbero sika, tre mikra, preske nigra, qua devenas de Korinto, de la insuli Ioniana\lingui{DFRS}
\artiklo{koriono}%
\vorto{koriono}\fako{anat.} Envelopilo extera di la feto\lingui{DEFIS}
\artiklo{korizo}%
\vorto{korizo}\fako{biol.} Duimesko o duoplesko di ula parti, per formaco di organi supernombra
\artiklo{korko}%
\vorto{korko} Ta strato sponjatra e lejera di speco di querko, de qua onu facas botelo-stopili, flotaceri, e c.\latina{quercum suber}\lingui{DE}
\artiklo{kormorano}%
\vorto{kormorano}\fako{zool.} Ucelo qua formacas grupo ek la klaso « palmipedi », qua vivas rive di la aquo-loki (lagi, mari, riveri) e qua su nutras ye fishi. (carbo)\latina{phalacrocorax}\lingui{DEFI}
\artiklo{kornajar}%
\vorto{kornajar}\fako{netrans.) (aludante kavalo, asno} Produktar ta stertoro quan kreas morbo di la organi respirala\lingui{DEFIRS}
\artiklo{kornalino}%
\vorto{kornalino}\fako{mineral.} Agato mi-diafana, plu o min obskure-reda, de qua onu facas siglili, ringi fingrala, e c.\lingui{DEFIS}
\artiklo{kornamuzo}%
\vorto{kornamuzo}\fako{muziko} Muzik-instrumento rustika, di qua la sonkoloro esas nazo-sonanta, facita ek felto-sako, aero-tanko, quan onu inflas per la boko, mediace tubo, ed ek tri tubi muntita sur la felto-sako, de qui du emisas noto kontinua, e de qui la triesma esas provizita ye boko-peco langetoza aden qua suflas la pleanto, e borita ye trui sur qui lu promenigas sua fingri por produktar diversa noti\lingui{FIS}
\artiklo{kornelo}%
\vorto{kornelo}\fako{bot.} Frukto de arboro di qua la ligno esas tre harda, qua formacas la speco tipa di la familio « kornacei », e qua similesas mikra pruno reda\latina{cornus mas}
\artiklo{korneo}%
\vorto{korneo}\fako{anat.} Membrano qua tapisas la dopa parto di la okulo\lingui{DEFIS}
\artiklo{korneto}%
\vorto{korneto}\senco{1}{Mikra korno, rustika, korno ek terakoto, garnisita ye boko-peco, per qua la pueri su amuzas, pleante lu lor karnavalo. – Granda fluto quan onu uzas en la kori por sub-tenar la voci}\senco{2}{To quo esas kornatra}\lingui{DEFIS}
\artiklo{kornico}%
\vorto{kornico}\fako{arkitekt.} Ornamento salianta qua kronizas la cirklo di la entablamento di edifico. Ornamento analoga, ek petro, ligno, e c., qua existas cirkum plafono, chambro, qua kronizas moblo, pordo, e c.\lingui{EFIRS}
\artiklo{korniko}%
\vorto{korniko}\fako{zool.} Speco di korvo, mikra\latina{corvus cornix}\lingui{FIS}
\artiklo{korno}%
\vorto{korno}\senco{1}{Surkreskajo konatra, rekta o kurva, qua kreskas sur la fronto di la rumineri, ordinare che la maskulo, e formacas quaza prolonguro di la osto frontala}\senco{2}{\textit{(analogie)} Surkreskajo, kornatra, qua existas sur la kapo di ula animali. Anke : pedunkli retrotirebla, qua existas sur la kapo di la limaki e qua portas la organi vidala}\senco{3}{To quo finas per punto simila a korno}\senco{4}{La materio tenaca, lameloza, de qua facesas la korni di la rumineri}\senco{5}{Korno kavigita e truizita, quan uzis la muton-pastori, la muton-gardisti, la chaseri, e c. por signalar, advokar}\senco{6}{\textit{(muziko)} Instrumento suflala, ek kupro, tordita adspirale, di qua un extremajo recevas boko-peco (mikra kupeto ek metalo, ek ivoro, e c.) e di qua la altra extremajo formacas kono expansita quan onu nomizis « funelo »}\lingui{EFIS}
\artiklo{koro}%
\vorto{koro}\senco{1}{\textit{(epoki antiqua)} Asemblitaro ek viri e mulieri, qui dansas marchante segun kadenco, ye la soni de la voci, de la instrumenti muzikala}\senco{2}{\textit{(nia-epoke)} Asemblitaro de personi qui kantas ensemble}\lingui{DEFIRS}
\artiklo{korodar}%
\vorto{korodar}\fako{trans.} Rodar lente, per efikivo kemiala (maxim-multa-kaze)\lingui{DEFIS}
\artiklo{koroido}%
\vorto{koroido}\fako{anat.} Membrano interna, kovrita ye materio nigra (pigmento) qua tapisas la okulo
\artiklo{korolario}%
\vorto{korolario} Konsequo qua esas sequontade la konkluzo di demonstro\lingui{DEFIRS}
\artiklo{korolo}%
\vorto{korolo}\fako{bot.} Parto di la floro, formacita ek un o plura petali, qua envelopas la pistilo e la stamini\lingui{DEFIS}
\artiklo{koronala}%
\vorto{koronala}\fako{anat.} Dicesas pri la arterii e veini di la kordio e di la stomako, dispozita cirkle
\artiklo{koronero}%
\vorto{koronero} En Britania : oficiero judiciala di qua la tasko konsistas ye inquestar, helpate de jurio, pri la kazi di morto violentala, naufrajo, e la deskovri di trezori
\artiklo{korpo}%
\vorto{korpo}\senco{1}{Che la homi e la bestii : ensemblo ek la parti materiala qui kompozas organismo en qua rezidas la vivo animalala}\senco{2}{Omna agregajo ek molekuli materiala}\senco{3}{(1) Parto precipua di kozo; (2) ensemblo formacita ek kolektajo de personi, de kozi}\lingui{DEFIRS}
\artiklo{korporaciono}%
\vorto{korporaciono} Asociitaro de individui, interligita per reguli, obligesi, privileji komuna\lingui{DEFIRS}
\artiklo{korporalo}%
\vorto{korporalo}\fako{liturgio katol.} Linjo sakrigita quan la sacerdoto extensas sur la altaro, lor meso, e sur qua lu pozas la kalico e la hostio\lingui{DEFIS}
\artiklo{korpulenta}%
\vorto{korpulenta}\fako{pri homo} Karakterizata da ampleso, plu o min grandega, di sua korpo. – Di qua la tota korpo (anke la membri) esas grosa, dika; muskuloza forte, do robusta e sana\lingui{DEFIS}
\artiklo{korpuskulo}%
\vorto{korpuskulo} Materio-grano infinite mikra\lingui{DEFIS}
\artiklo{korsajo}%
\vorto{korsajo} Parto di la robo, qua kovras la busto (di muliero), la parto supra, de la shultri til la tayo (do, super la jupo)\lingui{EF}
\artiklo{korsaro}%
\vorto{korsaro} Viro qua parkuras la mari (ed anke, ma rare, landi) skope di spolio, rapto – ex-pirato qua obtenis (en Francia) de la rejo la titulo « korsaro »\lingui{DEFIRS}
\artiklo{korseto}%
\vorto{korseto} Mikra korsajo ek barti, destinita a subtenar la mami; a vidigar la konturi di la tayo di la mulieri; ad igar la mulieri repleta o grosa semblar preske gracila\lingui{DEFRS}
\artiklo{korso}%
\vorto{korso} Aleo uzata kom promeneyo publika, en urbo\lingui{DFI}
\artiklo{kortico}%
\vorto{kortico}\fako{bot.} Envelopilo, « shelo », di la stipo e di la branchi di ula planti\lingui{EFIS}
\artiklo{korto}%
\vorto{korto}\senco{1}{Rezideyo di suvereno e di lua cirkondanti}\senco{2}{\textit{(yuro-cienco)} Asemblitaro tribunala qua eventis, origine, en la lojeyo di la suvereno}\senco{3}{Parto di domeno, tereno sen-tekta, generale cirkondata da muri o da edifici, qua esas avan o dop la habiteyo precipua}\lingui{EFIS}
\artiklo{korundo}%
\vorto{korundo}\fako{mineral.} Lapido, preske tam harda kam diamanto, formacita ek alumino diverse-koloroza pro kontenar oxidi metalala\lingui{DEFIRS}
\artiklo{koruptar}%
\vorto{koruptar}\fako{trans., per}\senco{1}{Alterar per putro (substanco)}\senco{2}{(1) Alterar to quo esas sana, honesta, che homo; (2) Alterar che ulu, per donaci, promisi, la integreso di la koncienco; (3) Destruktar to quo esas justa, sana, pura en la maniero judicar, sentar, expresar su; (4) Destruktar la integreso di la senco, di la texto di ulo}\lingui{DEFIRS}
\artiklo{korvear}%
\vorto{korvear}\fako{netrans., por, pri}\senco{1}{\textit{(olim)} \textit{(aludante vasalo)} Laborar gratuite e koakte por sua sinioro, laboro konsistanta ye ulo agenda tota-jorne, da homo o da charjo-bestio}\senco{2}{Laborar por konservar la voyi}\senco{3}{Exekutar tasko tre desagreabla, de qua onu ne povas dispensar su}\lingui{EF}
\artiklo{korveto}%
\vorto{korveto}\fako{navig.} Milito-navo, di qua la tipo esas meza inter brigo e fregato\lingui{DEFIRS}
\artiklo{korvo}%
\vorto{korvo}\fako{zool.} Ucelo granda quale mikra hano, di qua la plumaro esas nigra, la krii stridanta (kroaso), qua su nutras ye frukti, ye mikra bestii, ma esas avida precipue ye la karno di la kadavri\latina{corvus corax}\lingui{I}
\artiklo{kosmetiko}%
\vorto{kosmetiko} La parto di higieno qua koncernas la kompozo e la uzo di la produkturi per qui onu sorgas la hari, la denti, la pelo (por lia beleso)\lingui{DEFIRS}
\artiklo{kosmo}%
\vorto{kosmo} Universo, egardata kom organizuro\lingui{DEFIRS}
\artiklo{kosmogonio}%
\vorto{kosmogonio} Sistemo pri la formaceso di universo\lingui{DEFIRS}
\artiklo{kosmografio}%
\vorto{kosmografio} Deskripto di la sistemo astronomiala di nia planeto e di universo\lingui{DEFIRS}
\artiklo{kosmologio}%
\vorto{kosmologio} La cienco pri la *leyi (legi) qui regnas universo\lingui{DEFIRS}
\artiklo{kosmopolita}%
\vorto{kosmopolita}\senco{1}{Qua egardas su kom civitano di mondo (kontraste kun civitano di ca o ta lando)}\senco{2}{Qua vivas lore en ca lando, lore en ita, lore en altra, senfine}\lingui{DEFIRS}
\artiklo{kosmoso}%
\vorto{kosmoso}\fako{bot.} Genero de planti ek la familio « kompozaji », qua konsistas ye dek speci, de qui un esas kultivata en la gardeni pro lua foliaro lejera e lua flori granda, blanka o reda\latina{cosmos bipinnatus}\lingui{DEFIRS}
\artiklo{kosto}%
\vorto{kosto}\fako{anat.} Singla ek la osti plata ed arkatra qui, departante de la spino, e venante rajuntar avane la sternumo, formacas la kavajo ostoza « pektoro »\lingui{EFIS}
\artiklo{kostumo}%
\vorto{kostumo} Maniero vestizar su, propra di populo, di epoko, di klaso, di cirkonstanco\lingui{DEFIR}
\artiklo{koteno}%
\vorto{koteno}\fako{patol.} Falsa membrano qua formacesas interne di la guturo, en ula kazi di angino\lingui{I}
\artiklo{kotiledono}%
\vorto{kotiledono}\senco{1}{\textit{(anat.)} Lobo karnoza di la placenta}\senco{2}{\textit{(bot.)} Lobo, folio embriona di la semino di vejetanto fanerogama}\lingui{DEFIS}
\artiklo{kotiliono}%
\vorto{kotiliono} Danso kun figuri qua finas la balo (dansata olim da 4 od 8 personi)\lingui{DEFIS}
\artiklo{kotiloido}%
\vorto{kotiloido}\fako{anat.} Di qua la formo esas olta di ta kavajo di osto qua recevas altra osto
\artiklo{kotleto}%
\vorto{kotleto}\fako{koquarto} Kosto de mutono, de bovyuno, de porko, e c., destinita a manjeso da homo\lingui{DEFIR}
\artiklo{koto}%
\vorto{koto}\senco{1}{Indiko pri la altitudo di la punti diversa di tereno, di aquo-fluo, e c., relate nivelo egardata kom bazo}\senco{2}{Indiko, per cifri, per literi, di singla serio di peci ek dokumentaro, ek arkiveyo}\lingui{F}
\artiklo{kotono}%
\vorto{kotono} Materio texebla, blanka, fina, qua kovras la semino di arboreto qua formacas genero ek la familio « malvacei », qua origine venis de India e quan onu kultivas nun anke en Egiptia ed en la regiono sud-esta di Usa\lingui{EFIS}
\artiklo{koturno}%
\vorto{koturno} Pedovesto quan *weris, che la populi antiqua, la viri e la mulieri di rango super-meza; sandalo kun suolo, fixigita an la pedo per ledro-bendi ligita super la maleolo e lasante la pedo-fingri nekovrata\lingui{DFIS}
\artiklo{kovar}%
\vorto{kovar}\fako{trans.) (aludante la femino di ucelo} Sternar su e permanar tale dum ula quanto de tempo, sur sua ovi, por konservar ad oli konstante la kaloro qua igos la uceleto ekirar kande parformacita\lingui{FI}
\artiklo{kovrar}%
\vorto{kovrar}\fako{trans., per, ye}\senco{1}{(1) Garnisar ye ulo quan onu pozas sur…; (2) Garnisar per esar pozita sur…}\senco{2}{(1) Dissemar ulo sur ulo, disjete; (2) Dissemesir disjete sur ulo}\senco{3}{(1) Protektar per ulo quan onu pozas sur…; (2) Protektar per esar sur, avan}\senco{4}{(1) Celar per pozar ulo sur, avan…; (2) Celar per esar sur, avan…}\lingui{EFIS}
\artiklo{koxalgio}%
\vorto{koxalgio}\fako{patol.} Artrito o tumoro blanka ye la hancho\lingui{DEF}
\artiklo{kozako}%
\vorto{kozako} Slavo, kelke mestica, qua habitas regioni di la Rusia di Europa e di Azia; kom kavalriano ecelanta, lu posedis privileji administreriala e fiskala, kondicionata per lua devo esar armeano dum duadek yari, e, pose, en la milico til la fino di sua valideso\lingui{DEFIRS}
\artiklo{kozo}%
\vorto{kozo} Omna realajo, konkreta od abstraktita, quan onu indikas ye maniero nedeterminita, nepreciza\lingui{FIS}
\artiklo{krabo}%
\vorto{krabo}\fako{zool.} Nomo vulgara di krustacei diversa, dekpeda, kurta-kauda, di qua la tipo esas la krabo ordinara, speco manjebla qua habitas en la sablo di maro\latina{cancer, carcinus}\lingui{DEFR}
\artiklo{krago}%
\vorto{krago} Peco de telo, de dentelo, qua, pos cirkumirir la kolo, retrovenas ad la supra parto di la pektoro\lingui{D}
\artiklo{krak!}%
\vorto{krak!} Interjeciono qua expresas evento subita, o la bruiso de krako
\artiklo{krakar}%
\vorto{krakar}\fako{netrans.} Produktar bruiso sika, per la krumplo di la parti\lingui{DEF}
\artiklo{kraknelo}%
\vorto{kraknelo} Kukajo harda, qua krakas sub la denti qui mastikas lu\lingui{EF}
\artiklo{krampo}%
\vorto{krampo}\fako{patol.} Kontrakto dolorigiva qua produktas rigideso kurta-dura di la muskuli brakiala, gambala, \textit{(analogie : stomakala)}. Kontrakto spasmala\lingui{DEF}
\artiklo{krampono}%
\vorto{krampono} La signo tipografiala
\artiklo{kranio}%
\vorto{kranio}\fako{anat.} Buxo osta qua kontenas e protektas la cerebro\lingui{EFIS}
\artiklo{krano}%
\vorto{krano}\fako{tekn.} Aparato por sublevar la kargaji e transportar li kurta-diste. Lu konsistas ek framo ek fero od ek ligno, muntita sur pivoto cirkum qua lu povas rotacar, ed ek du mekanismi : un mekanismo di sublevo (puliaro o vindilo) ed un mekanismo di rotacigo\lingui{DER}
\artiklo{krapular}%
\vorto{krapular}\fako{netrans.} Debochar grosiere; karakterizivo : eceso di vino-drinko\lingui{DEFIS}
\artiklo{kraso}%
\vorto{kraso} Strato de materio sordida qua formacesas ed akumulesas sur la pelo, sur linjo, sur la vesti quin onu *weras, sur la objekti\lingui{FS}
\artiklo{kratego}%
\vorto{kratego}\fako{bot.}\senco{1}{Arbusto dornoza, ek la familio « rozacei », kun mikra flori blanka}\senco{2}{La floro di ta arbusto}\latina{crataegus}\lingui{L}
\artiklo{kratero}%
\vorto{kratero}\senco{1}{\textit{(epoki antiqua)} Vazo larja, en qua onu mixis vino ed aquo por plenigar la kalici e la kupi}\senco{2}{\textit{(metaf.)} Aperturo formacita somite di volkano, per qua olca forjetas la materii inflamita : flami, lavao, cindri}\senco{3}{Aperturo facita ye la parto supera di fornego di vitrifisto}\lingui{DEFIRS}
\artiklo{kravacho}%
\vorto{kravacho} Vergo flexebla qua finas per mecho, quan la kavalkanti uzas kom flogilo\lingui{DF}
\artiklo{kravato}%
\vorto{kravato} Peco de stofo quan onu *weras cirkum sua kolo, kom ornivo, e qua (maxim-multa-kaze) pendas extremaje sur la pektoro\lingui{DEFI}
\artiklo{krayono}%
\vorto{krayono} Peco de erco (grafito, e c.), plu o min mola, quan onu uzas por trasar, desegnar, skribar, e c.\lingui{EF}
\artiklo{krazo}%
\vorto{krazo}\fako{gram.} Fuzuro de la vokalo o diftongo finala di vorto, kun la vokalo o diftongo komencala di la vorto sequanta\lingui{DEFIS}
\artiklo{krear}%
\vorto{krear}\fako{trans., de, ek, per}\senco{1}{\textit{(religio)} \textit{(aludante deo)} Facar (ulo) ek nulo, prenante lu de nihilo}\senco{2}{\textit{(aludante homo)} Kompozar (verko, kozo qua ne existis antee)}\lingui{EFIS}
\artiklo{krecento}%
\vorto{krecento}\senco{1}{(1) Formo sesguroza di Luno dum lua kresko o deskresko; (2) To quo havas ta formo}\senco{2}{Pano-bloko butropasta, kun formo di Luno-krecento; facata precipue en Paris, Francia}\lingui{EFIS}
\artiklo{kredar}%
\vorto{kredar}\fako{trans.}\senco{1}{(1) Admisar (ulo) kom exakta sen verifiko. (2) Egardar (ulu) kom dicanta lo exakta, sen verifikar lo}\senco{2}{\textit{(ye ulu, ye ulo)} (1) Fidar a la realeso di ulo, sen verifikar; fidar a la efikiveso di ulo, ad olua povo, sen verifikar; (2) Fidar a la verdico di ulu, sen verifiko; (3) Fidar a lua karaktero, sen verifiko}\senco{3}{\textit{(senco absoluta)} Fidar religiale}\lingui{EFIRS}
\artiklo{kredenco}%
\vorto{kredenco} Plado-moblo, konsolo o tablo sur qua onu pozas la pladi, pladeti, glasi, e c. por la servi di tablo lor re-pasto\lingui{DEFIS}
\artiklo{kreditar}%
\vorto{kreditar}\fako{trans., per ulo} Vendar, aceptante la ajorno di la pago da la kompranto, pro ke onu judikas lu solventa\lingui{DEFIRS}
\artiklo{kremacar}%
\vorto{kremacar}\fako{trans.} Reduktar a cindri (objekto o la korpo di homo mortinta)\lingui{EFIS}
\artiklo{kremo}%
\vorto{kremo}\senco{1}{La parto maxim densa di lakto, qua acensas a la surfaco kande onu lasas lu repozar, ed ek qua onu facas butro per quirlo. \textit{(metaf.)} La parto maxim bona de ulo}\senco{2}{Pelikulo qua eventas surface di lakto kande onu varmigas lu}\senco{3}{Preparajo qua aspektas quale kremo}\lingui{DEFIRS}
\artiklo{kremono}%
\vorto{kremono}\fako{tekn.} Ferajo uzata por klozar od apertar fenestro; stango de fero rekta, muntita an un ek la montanti, quan onu adsuprenirigas od adinfreirigas per rotacigar mancho tale ke onu enirigas singla ek la extremaji aden seruro-boko, od ekirigas li de olu\lingui{F}
\artiklo{krenelo}%
\vorto{krenelo} Aperturo facita ye intervali en la parapeto di remparo, di turmo, por lansar projektili sur la enemiko\lingui{EF}
\artiklo{kreolo}%
\vorto{kreolo} Persono di raso blanka, naskinta en la kolonii (olim) Hispana di Amerika. – Persono qua naskis en la kolonii Europana inter-tropike\lingui{DEFIRS}
\artiklo{kreozoto}%
\vorto{kreozoto}\fako{kemio} Liquido oleatra, korodiva ed antisepsiiva, quan onu obtenas per la distilo di gudro lignala\lingui{DEFIRS}
\artiklo{krepisar}%
\vorto{krepisar}\fako{trans.} Indutar (muro) per trulo o balayilo, ye strato di gipso, di mortero, quan onu lasas aspera vice planigar lu\lingui{F}
\artiklo{krepitar}%
\vorto{krepitar}\fako{netrans.} Produktar bruiso ek intersequo di mikra kraki\lingui{DEFIS}
\artiklo{krepo}%
\vorto{krepo} Stofo ek lano delikata, plu o min diafana, kun filo ondoza, quan onu preparas per kelke tordar la filo longesale (warpe), pose plunjante la texuro aden aquo e frotante lu per vaxo\lingui{DEFR}
\artiklo{krepono}%
\vorto{krepono} Stofo frizita quale krepo, ma di qua la texuro esas plu dika, por togi advokatala, sutani, e c.\lingui{DEFIRS}
\artiklo{krepuskulo}%
\vorto{krepuskulo} Lumo nepreciza quan donas la suno-radii kande li reflektesas da la strati supra di atmosfero, t.e. kande Suno esas ja sub horizonto, e qua deskreskas segun ke la astro desaparas\lingui{EFIS}
\artiklo{kreskar}%
\vorto{kreskar}\fako{netrans.}\senco{1}{(1) \textit{(aludante ento organizita, vivoza)} Divenar plu granda gradoze, til la fino di sua developo normala. (2) \textit{(metaf.)} Divenar plu granda, plu multa, sen-cese, gradoze}\senco{2}{\textit{(aludante vejetanti)} Developar su naturale en ula regioni}\lingui{FIS}
\artiklo{kreso}%
\vorto{kreso}\fako{bot.} Planto perena, qua formacas genero ek la familio « kruciferi », qua kreskas alonge la rivereti, an la tereni inundata\latina{lepidium nasturtium}\lingui{DEFIR}
\artiklo{krespo}%
\vorto{krespo} Rondi de pasto tre dina, ek farino dilutita kun lakto ed ovi, quan onu varsas, lejera-strate, aden padelo indutita ye grasajo, por koquar lu\lingui{F}
\artiklo{kresto}%
\vorto{kresto}\senco{1}{Surkreskajo karnoza, di qua la konturo esas sesgura dento-forme, qua salias sur la kapo di ula galinacei. \textit{(analogie : sur la kapo di ula repteri, fishi, uceli.)}}\senco{2}{Parto salianta, streta, sur la somito di monto, di ondo, la supra parto di tekto, di muro}\lingui{EFIS}
\artiklo{krestomatio}%
\vorto{krestomatio} Selektajo ek texto-peci extraktita de la verki da la autori klasika\lingui{DEFIRS}
\artiklo{kretino}%
\vorto{kretino}\fako{patol.}\senco{1}{Persono rakitika, idioto maxim-multakaze strumoza}\senco{2}{\textit{(metaf.)} Persono stupida quale kretino}\lingui{DEFIRS}
\artiklo{kreto}%
\vorto{kreto} Kalko-karbonato amorfa, quan onu renkontras forme di strati plu o min dika, en ula sulo-strati.\lingui{DEFIS}\subvorto{}{kreto-krayono}{Bastoneto ek ta substanco, per qua onu skribas sur tabelo nigra o obskure-verda (en skolo)}{}{}
\artiklo{kretono}%
\vorto{kretono} Telo rezistiva, kun warpo ek kanabo, e wefto ek lino\lingui{DEFRS}
\artiklo{krevar}%
\vorto{krevar}\senco{1}{\textit{(netrans.)} Apertar su, spliteskante, pro tenseso ecesanta}\senco{2}{\textit{(trans.)} Apertar, splitante}\lingui{FI}
\artiklo{kreveto}%
\vorto{kreveto}\fako{zool.} Mikra krustaceo dekapeda, manjebla, quan onu renkontras sur la litoro marala\latina{crangon vulgaris}\lingui{F}
\artiklo{krevisar}%
\vorto{krevisar}\senco{1}{\textit{(netrans.)} Subisar fendeso plu o min profunda \textit{(aludante la surfaco di korpo)}}\senco{2}{\textit{(trans.)} Fendar plu o min profunde (la surfaco di korpo)}\lingui{EFI}
\artiklo{kriar}%
\vorto{kriar}\fako{netrans.} Lansar voco-sonuni stridanta e neartikulita\lingui{EFIRS}
\artiklo{kriblo}%
\vorto{kriblo} Recevuyo di qua la kulo esas borita ye trui inter-egala, por separar la objekti di groseso neinteregala, lasanta trapasar ici, retenante iti\lingui{FIS}
\artiklo{krik!}%
\vorto{krik!} Interjeciono qua imitas la bruiso de kozo lacerata
\artiklo{krik-krak}%
\vorto{krik-krak} ! Interjeciono per qua onu probas *reproduktar la bruiso sika di kozo qua ruptesas, laceresas
\artiklo{kriketo}%
\vorto{kriketo} Ludilo, origine Angla, en qua la ludanti, dividita a du partisi, lansas singlu-suafoye baloneti grava, per kroso ligna\lingui{EF}
\artiklo{kriko}%
\vorto{kriko} Mashino por sublevar nealte kargajo (quadri, e c.) per dento-stango quan elevas roto dentoza, movata per manivelo\lingui{FIS}
\artiklo{kriminar}%
\vorto{kriminar}\fako{netrans.}\senco{1}{Agar tale ke la lego etikala esas ne-egardata maxim grave}\senco{2}{\textit{(yuro-cienco)} Exekutar ago pri qua la lego stipulas ke lu esas la kauzo di puniso korpala ed infamiganta}\senco{3}{\textit{(hiperbole)} Exekutar ago quan onu judikas tro severe pri lua posibla konsequi mala}\lingui{DEFIRS}
\artiklo{krino}%
\vorto{krino} Pilo aspera, longa e flexebla di ula bestii; longa pilo ek la kolo od ek la kaudo di kavalo\lingui{FIS}
\artiklo{krinolino}%
\vorto{krinolino}\senco{1}{Texuro di qua la warpo esas lina, e la wefto ek krino}\senco{2}{Subjupo ek texuro de krino, por subtenar la robo por ke lu aspektez balonatre}\senco{3}{Subjupo garnisita ek barti, ek ringegi stala, e c., sama-skope}\lingui{DEFIRS}
\artiklo{krioforo}%
\vorto{krioforo}\fako{fiziko} Instrumento per qua la aquo divenas konjelata per sua propra vaporesko\lingui{DEFIS}
\artiklo{kripla}%
\vorto{kripla} Privacita de la uzo di un o plura ek sua membri, pro vunduro, morbo, e c.
\artiklo{kripo}%
\vorto{kripo} Manjo-trogo por la bruti
\artiklo{kripto}%
\vorto{kripto}\senco{1}{Tombo masonita, subsula, en kirko, olim uzata kom sepulteyo por la martiri}\senco{2}{\textit{(anat.)} Folikulo, mikra glando sakoforma, en la pelo od en la membrani mukoza, e sekrecanta – per pori – liquidi adsurface}\lingui{DEFIRS}
\artiklo{kriptogama}%
\vorto{kriptogama}\fako{bot.} Di qua la organi fruktifala esas poke videbla.\lingui{DEFIS}\subvorto{}{kriptogami}{Planti qui formacas la lasta ek la 24 klasi da Linné}{}{}
\artiklo{kriptografar}%
\vorto{kriptografar}\fako{trans.} Skribar per signi celanta, per cifri, e c.\lingui{DEFIRS}
\artiklo{kriptomerio}%
\vorto{kriptomerio}\fako{bot.} Genero de arbori, ek la familio « koniferi », vicina ye la cipreso, qui devenas de Japonia, e quin onu kultivas en Francia
\artiklo{krismo}%
\vorto{krismo}\fako{liturgio katol. e Greka} Oleo sakra quan onu uzas por la uncioni en ula sakramenti di la Eklezio katolika e Greka\lingui{DEFIS}
\artiklo{krispa}%
\vorto{krispa} Tale frizita ke li formacas tufo densa\lingui{EI}
\artiklo{kristalino}%
\vorto{kristalino}\fako{anat.} Korpo lenso-forma, diafana, qua esas en la okulo-globo, dop la pupilo, destinita a refraktar la radii luma ed igar li konvergar a retino\lingui{DEFIRS}
\artiklo{kristalo}%
\vorto{kristalo}\senco{1}{\textit{(en la vortaro ordinara)} Materio sen-kolora, extreme diafana}\senco{2}{\textit{(en la vortaro ciencala)} Korpo qua, kande solida, havas la formo di poliedro di qua la anguli e la lateri esas egala}\lingui{DEFIRS}
\artiklo{kristaloido}%
\vorto{kristaloido}\fako{biol.} Nomo di masi albuminoida ek substanco protoplasma, kun formi geometriala analoga ad olti di la kristali
\artiklo{kristo}%
\vorto{kristo}\fako{kristanismo}\senco{1}{Mesio, redemtero}\senco{2}{Iesu Kristo}\senco{3}{Imajo de Iesu Kristo (kom ex. : kristo ek ebeno)}\lingui{DEFIRS}
\artiklo{kriterio}%
\vorto{kriterio} Karaktero per qua (per komparo) onu dicernas lo exakta de lo ne-exakta, lo vera de lo falsa\lingui{DEFIRS}
\artiklo{kriticismo}%
\vorto{kriticismo}\fako{filoz.} Doktrino qua igas diskutata la fidindeso di la idei di Raciono\lingui{DEFIRS}
\artiklo{kritikar}%
\vorto{kritikar}\fako{trans.}\senco{1}{Observar la qualesi, bona o mala, di verko (literaturaja, arto-produkturo, e c.)}\senco{2}{Igar saliar la defekti di la kozi, di la personi}\lingui{DEFIRS}
\artiklo{krizalido}%
\vorto{krizalido} Nimfo di lepidoptero e lua envelopilo\lingui{DEFIS}
\artiklo{krizantemo}%
\vorto{krizantemo}\fako{bot.} Arboreto qua formacas genero ek la familio « kompozaji », kultivata en la Europana gardeni pro lua flori, flava, blanka, rozea, violea\latina{chrysanthemum}\lingui{DEFIRSL}
\artiklo{krizo}%
\vorto{krizo}\senco{1}{Fazo danjeroza di maladeso-periodo, qua esos decidiganta, sive favoroze, sive perisige}\senco{2}{Fazo danjeroza quan transiras politiko, financo, industrio, komerco, e c.}\lingui{DEFIRS}
\artiklo{krizokalko}%
\vorto{krizokalko} Aloyuro de zinko e kupro, kompozajo qua simulas oro e quan onu uzas en la faco di juveli falsa\lingui{DFIS}
\artiklo{kroasar}%
\vorto{kroasar}\fako{netrans.) (aludante la korvi, la rani} Agar la krio qua karakterizas sua speco\lingui{F}
\artiklo{krochar}%
\vorto{krochar}\fako{trans.} Facar objekti ek mashi, de lano, filo, silko, e c., per mikra stango de fero o ligno, di qua un pinto havas hoketo\lingui{EFR}
\artiklo{kroketo}%
\vorto{kroketo} Ludo-speco, de origino Angla, en qua onu pulsas buli per frapi agita per ligna martelo, por igar ta buli transsubirar arketi stekita en la sulo\lingui{DEFIRS}
\artiklo{krokodilo}%
\vorto{krokodilo}\fako{zool.} Reptero sauria, amfibia, di la granda fluvii di India e di Egiptia, timinda pro lua grandeso e devoremeso\latina{crocodilus}\lingui{DEFIRS}
\artiklo{kromatiko}%
\vorto{kromatiko}\senco{1}{\textit{(muziko)} Qua konsistas ek intersequo de mi-toni (mi-toni kromatika), kontraste kun la mi-toni diatonika, plu mikra ye un komao}\senco{2}{Qua koncernas la kolori}\lingui{DEFIS}
\artiklo{kromatino}%
\vorto{kromatino}\fako{biol.} Substanco tre afina pri la kolorizivi bazala, e quan onu renkontras maxim-multa-kaze forme di grani kunjuntita a la filamento nukleala
\artiklo{kromatofila}%
\vorto{kromatofila}\fako{biol.} Dicesas pri la celuli, grani od *organiti, apta kaptar facile la kolorizivi
\artiklo{kromatoforo}%
\vorto{kromatoforo}\fako{biol.} Celulo di la tisuo konjuntiva, qua kontenas multa pigmenti, quan onu renkontras periferie di la neuraxo e di la membrani sensala (iriso, pelo, e c.)
\artiklo{kromatolizo}%
\vorto{kromatolizo}\fako{biol.} Dissolveso di la kromatino di la celulo-nukleo, qua karakterizas un ek la fazi di mitoso
\artiklo{kromidio}%
\vorto{kromidio}\fako{biol.} Nomo donita da Hertwig a la elementi di la nukleo nepreciza di la bakterii
\artiklo{kromo}%
\vorto{kromo}\fako{kemio} Korpo simpla, metala, di qua la kompozaji esas remarkinda pro lia kolori\simbolo{Kr}
\artiklo{kromolitografar}%
\vorto{kromolitografar}\fako{trans.} Agar imprimo litografiala koloroza
\artiklo{kromosomo}%
\vorto{kromosomo}\fako{biol.} Segmento di la filamento kromatika nuklea qua aparas dum mitoso
\artiklo{kromotipar}%
\vorto{kromotipar}\fako{trans.} Imprimar koloroze per la procedi tipografiala\lingui{DEFRS}
\artiklo{kron}%
\vorto{« kron »} Monet-uno, di qua la valoro varias segun la landi (Britania, Suedia, e c.)
\artiklo{kronglaso}%
\vorto{kronglaso} Vitro diafana di tre bela qualeso, qua kontenas multa potaso, tre facile gisebla (fuzebla), di qua la refraktiveso-grado esas tre alta, e qua uzesas por fabrikar lensi por la instrumenti optikala, ed anke prismi kromata\lingui{DEFIRS}
\artiklo{kronika}%
\vorto{kronika}\fako{aludante morbo}\senco{1}{Qua par-iras lente sua periodi}\senco{2}{Qua divenis forta segun ke tempo fluis e, konseque, ne plus esas akuta}\lingui{DEFIRS}
\artiklo{kroniko}%
\vorto{kroniko}\senco{1}{Kolektajo de fakti historiala, segun la ordino di lia evento}\senco{2}{Ensemblo ek la fakti, ek la informi qui esas mencionita pri la personi}\lingui{DEFIRS}
\artiklo{krono}%
\vorto{krono}\senco{1}{Ringo destinita a metesar cirkum la kapo : (1) Ringo de flori, de folii interplektita quan onu metas cirkum la kapo kom (a) ornivo; (b) distingivo armeala o civila. \textit{(metaf.)} Rekompenso cielala; (2) Ringo de metalo, ornita ye floroni, ye perli, e c. quan onu metas cirkum la kapo kom insigno di la suvereneso, di la nobeleso. La krono rejala, dukala, e c. qua esas super la stampuro blazonala}\senco{2}{To di quo la formo memorigas krono}\lingui{DEFIRS}\subvorto{}{dento-krono}{La parto supra qua salias de la jenjivo e kovresas ye emalio}{3}{}
\artiklo{kronologio}%
\vorto{kronologio} La cienco di qua la temo esas pruvar la dati di la eventaji.historiala\lingui{DEFIRS}
\artiklo{kronometro}%
\vorto{kronometro} Horlojeto di qua la pendulo esas kompensita, tale ke lu ne varias segun la temperaturo\lingui{DEFIRS}
\artiklo{kropo}%
\vorto{kropo}\fako{anat.} Sako o membrano quan la uceli havas sub la guturo, ed en qua la nutrivi sejornas ante pasar aden la stomako\lingui{DE}
\artiklo{kroso}%
\vorto{kroso}\senco{1}{Bastono eklezio-pastorala, signo di la episkopeso od abadeso, kun un extremajo kurvigita forme di voluto, quan la abadi kustumis portar kun la voluto adinterne, e la episkopi adextere}\senco{2}{Parto kurvigita per qua finas fusilo, pistolo, ye la parto opozita a la tubo}\lingui{F}
\artiklo{krotalo}%
\vorto{krotalo}\fako{zool.} Serpento venenoza, di qua la korneti, inter-inkastrita ye la extremajo di la kaudo, kliktas\latina{crotalus}\lingui{FISL}
\artiklo{krotono}%
\vorto{krotono}\fako{bot.} Planto qua formacas genero di la familio « euforbiacei », di qua un speco donas la tornasolo, e di qua un altra donas emetiko e purgivo, e grani de qui onu ekpresas oleo akra, uzata kom drastiko e kom revulsivo violenta\latina{croton}\lingui{DEFIRS}
\artiklo{krozar}%
\vorto{krozar}\fako{netrans.) (aludante navo} Trairar maro, en ula regioni, por surveyar li\lingui{DEFIRS}
\artiklo{krucho}%
\vorto{krucho} Vazo tera, terakota, larja-ventra, kun anso e kolo, destinita a kontenar liquidi\lingui{DFR}
\artiklo{krucifera}%
\vorto{krucifera}\fako{bot.} Di qua la petali esas krucoforma\lingui{DEFIRS}
\artiklo{krucifixo}%
\vorto{krucifixo} Kruco ek ligno, metalo, ivoro, e c., sur qua esas reprezentita Iesu Kristo krucagita\lingui{DEFIRS}
\artiklo{kruco}%
\vorto{kruco}\senco{1}{\textit{(epoki antiqua)} (1) Sorto di jibeto, maxim-multakaze trairata da traverso, sur qua onu igis mortar ula krimininti, surligita o surklovagita ye la extremaji di lia membri. (2) La jibeto sur qua Iesu Kristo esis klovagata ed ocidata \textit{(segun kristanismo)}. (3). \textit{Metaf}. : La religio di Iesu Kristo. (4) Objekto krucoforma, ek ligno, metalo, e c., simbolo di Iesu Kristo krucagita}\senco{2}{Distingivo honorizanta, ornivo ek kruco}\senco{3}{To quo esas krucoforma}\lingui{DEFIRS}
\artiklo{kruda}%
\vorto{kruda}\senco{1}{\textit{(aludante la substanci nutrivala)} Qua ne esas koquita o bakita}\senco{2}{\textit{(aludante materio quan onu submisas a preparo)} Qua ne subisis ta preparo}\senco{3}{\textit{(metaf.)} Quan nulo minfortigas (kom ex. : termino, dis-kurso, kolori qui ne fuzesas kun olti qui cirkondas li)}\lingui{EFIS}
\artiklo{kruela}%
\vorto{kruela}\senco{1}{Qua diletas per sufrigar}\senco{2}{Qua sufrigas}\lingui{DFIS}
\artiklo{krugo}%
\vorto{krugo} Vazo ansoza, kun ventro tre konvexa, kolo streta, boko expansita, ek ligno, kun cirkondo-bendi fera o stana, e qua kontenas inter 5 e 10 litri; uzata por rakar vino, precipue en vino-venderio\lingui{D}
\artiklo{krular}%
\vorto{krular}\fako{netrans.) (aludante edifico, propradice e metafore} Subite falar unbloke, e divenar fragmenti ed eskombri pro la shoko kontre sulo\lingui{FI}
\artiklo{krumo}%
\vorto{krumo} Parto interna di pano-bloko, qua, pro ne subisir la efiko da fairo, duras esar mola\lingui{DE}
\artiklo{krumplar}%
\vorto{krumplar}\fako{trans.} Desfreshigar (stofo) per quaza petriso qua igas lu aspektar quale shifono\lingui{DE}
\artiklo{krupiero}%
\vorto{krupiero} En hazardo-luderio, la employato qua direktas la *ruleto-ludo ye la nomo di la direktero di la firmo\lingui{DEFI}
\artiklo{krupo}%
\vorto{krupo}\fako{patol.} Laringito akuta, qua afektas la infanti, e karakterizata da tuso rauka e da la produkteso di falsa membrani qui obstruktas la respir-tubi\lingui{DEFIRS}
\artiklo{kruro}%
\vorto{kruro}\fako{anat.} Parto di la gambo, qua iras de la hancho til la genuo\lingui{DEFIS}
\artiklo{krustacei}%
\vorto{krustacei}\fako{zool.} Klaso ek animali artropoda, di qui la korpo esas envelopita da krusto kalka\lingui{DEFIS}
\artiklo{krusto}%
\vorto{krusto}\senco{1}{Parto extera di pano-bloko, hardigita ed igita krakanta per bako}\senco{2}{Parto bakita qua cirkondas pasteto, karno-torto o fisho-torto}\senco{3}{Parto solida, friabla, qua naskas periferie di ula kozi}\senco{4}{Envelopilo di la krustacei}\lingui{DEFIS}
\artiklo{*kruzar}%
\vorto{*kruzar}\fako{trans.} Pasar apud (ulu), irante inversa-sinse, quaze por atingar la loko de ube la altra persono departis
\artiklo{kruzelo}%
\vorto{kruzelo} Vazo ek materio fair-espruva, porcelano, metalo, maxim-multa-kaze plu-stretigita apud la fondo, e destinita a la giso o kalcino di ula substanci\lingui{EFIS}
\artiklo{kuafar}%
\vorto{kuafar}\fako{trans.} Ajustar la hararo per dispozar la hari ula-maniero\lingui{DEF}
\artiklo{kubebo}%
\vorto{kubebo}\fako{bot.} Sorto di pipriero. La frukto de ta arboreto uzesis olim, forme di pipro, en la kuraco di la blenoragio\latina{piper cubeba}\lingui{DEFIRS}
\artiklo{kubiko}%
\vorto{kubiko}\fako{matem.} Equaciono triesma-grada, en qua la nekonocato esas ye la tria potenco\lingui{DEFIRS}
\artiklo{kubito}%
\vorto{kubito}\fako{anat.} La plu grosa ek la osti di la avan-brakio, qua okupas olua parto interna e supra, e di qua la extremajo formacas la kudo per sua junto artike a la humero\lingui{EFIS}
\artiklo{kubo}%
\vorto{kubo}\senco{1}{Paralelepipedo sis-edra, qua formacas quadrati inter-egala}\senco{2}{La tria potenco di nombro, o la produto de la nombro, multiplikita per lua quadrato}\lingui{DEFIRS}
\artiklo{kuboido}%
\vorto{kuboido}\fako{anat.} La mikra osto kubatra, qua esas an la parto extera e supra di la tarso\lingui{EF}
\artiklo{kudo}%
\vorto{kudo}\senco{1}{\textit{(anat.)} Angulo salianta quan formacas la parto dopa di la artiko di la brakio kun la avanbrakio}\senco{2}{Loko di la maniko di vesto, qua korespondas a la kubito}\senco{3}{Angulo salianta (kom ex. : la rivero, la voyo, sinuifas kude)}\senco{4}{Parto di utensilo, di instrumento, qua formacas angulo salianta}\senco{5}{Peco de tubo, formacita ek du parti qui facas angulo}\lingui{FIS}
\artiklo{kuglo}%
\vorto{kuglo} Mikra globo ek metalo, uzata kom projektilo en la pafarmi portebla\lingui{D}
\artiklo{kuko}%
\vorto{kuko}\fako{koquarto} Mixuro ek farino, butro, ovi, reduktita a pasto, bakita, ed ulakaze sukrizita, aromatizita, garnisita ye kremo, e c.\lingui{DE}
\artiklo{kukombro}%
\vorto{kukombro}\fako{bot.} Planto leguma un-yar-vivanta, ek la familio « kukurbitacei ». Frukto di ta planto, oblonga, tre aquoza, quan onu manjas o koquita o salade\latina{cucumis sativus}\lingui{DEFS}
\artiklo{kukulo}%
\vorto{kukulo}\fako{zool.} Ucelo klimera, ek la genero « pigo »\latina{cuculus canorus}\lingui{DEFIRS}
\artiklo{kukurbitaceo}%
\vorto{kukurbitaceo}\fako{bot.} Familio de planti herbatra, kun stipito reptera, di qua la gurdego esas la tipo\lingui{EFIS}
\artiklo{kukurbito}%
\vorto{kukurbito}\fako{bot.} Varietato di gurdego. La frukto di ta varietato\latina{cucurbita}\lingui{EFIS}
\artiklo{kulato}%
\vorto{kulato} La fundo di la tubo di pafarmo, plu dika kam la cetero pro ke lu korespondas a la kargajo de pulvero\lingui{FIS}
\artiklo{kulbutar}%
\vorto{kulbutar}\fako{netrans.}\senco{1}{Bruske falar renverse}\senco{2}{Saltar rotacante, kapo adavane ed addope}\senco{3}{Renversesar rotace, tale ke la infrajo divenas la suprajo}\lingui{F}
\artiklo{kulco}%
\vorto{kulco}\fako{zool.} Mikra moskito\latina{culex}
\artiklo{kulebrino}%
\vorto{kulebrino} Sorto de kanono, longa e mikra-diametra, quan onu uzis olim\lingui{EFIS}
\artiklo{kuli}%
\vorto{« kuli »} En India, Indochinia, Insulo Mauricio : pen-laboristo, portisto
\artiklo{kuliero}%
\vorto{kuliero} Implemento di la tablo, di la koqueyo, e c. ek mancho, e parto konkava, ovala o cirklatra, per qua onu cherpas, de vazo, de plado, e c., nutrivi liquida o mi-liquida, por transvarsar li o portar li aden sua boko\lingui{F}
\artiklo{kuliso}%
\vorto{kuliso}\senco{1}{\textit{(tekn.)} Suportilo kun ranuro interna, alonge qua peco moviva povas irar e retrovenar sen obstaklo}\senco{2}{(1) Ranuro facita sur la *stejo di teatro, alonge qua onu movas la frami di la dekoruri. (2) Parto di la teatro qua esas dop la *stejo}\senco{3}{En borso financala : la loko ube operacas la kurtajisti qui ne havas permiso oficala}\senco{4}{Rifalduro en la supra parto di kurteno, en qua onu igas pasar stango – en la bordo di vesto, por igar pasar kordono qua posibligas lua min-longigo per plisuri}\lingui{DEFR}
\artiklo{kulminar}%
\vorto{kulminar}\fako{netrans.}\senco{1}{\textit{(astron.)} Atingar, pasante ye la meridiano, sua alteso maxima super horizonto}\senco{2}{\textit{(metaf.)} Esar (sociale) tre altagrade}\lingui{DEFIRS}
\artiklo{kulo}%
\vorto{kulo}\senco{1}{\textit{(anat.)} La dopa parto (ek la glutei e la supra parto di la kruri) di la korpo homa ed animala}\senco{2}{La fundo di ula kozi (botelo, plado, glaso, sako; en la artichoko : la parto pulpoza qua portas la feno e la foliilFIS)}
\artiklo{kulpar}%
\vorto{kulpar}\fako{netrans., pri, pro} Faliar pri agar sua devo, koncerne od etiko o prudenteso e habileso\lingui{EFIS}
\artiklo{kultar}%
\vorto{kultar}\fako{trans.}\senco{1}{Honorizar deo suprega}\senco{2}{\textit{(metaf.)} Adorar ulu, ulo}\senco{3}{Agar la ceremonii extera per qui la homi honorizas deo}\lingui{DEFIRS}
\artiklo{kultelo}%
\vorto{kultelo} Instrumento mikra-dimensiona, facita por tranchar o sekar, e qua konsistas ye lamo muntita an mancho\lingui{FI}
\artiklo{kultivar}%
\vorto{kultivar}\fako{trans.}\senco{1}{Submisar (tero) ad ula labori por igar lu fertila}\senco{2}{Submisar (la planti) a sorgi destinata a plubonigar la efiko da naturo}\senco{3}{Explotar (ula planti) por rekoltar lia produkturo}\senco{4}{\textit{(metafore)} Kultivar la arti pacala, kultivar la cienco, spirito, e c.}\lingui{DEFIRS}
\artiklo{kulturo}%
\vorto{kulturo} La rezultajo de la kultivo di la arti, di la cienci, di literaturo, di etiko\lingui{DEFIRS}
\artiklo{kumino}%
\vorto{kumino}\fako{bot.} Planto umbelifera, di qua la grani esas aromoza (li saporas quale anizo)\latina{carum carvi}\lingui{EFI}
\artiklo{kumis}%
\vorto{« kumis »} Drinkajo fermentacita quan la nomadi di la centroregioni di Azia e di la sudo-regioni di Rusia, fabrikas ek la lakto di la kavalini
\artiklo{kumpreso}%
\vorto{kumpreso} Peco de linjo, faldita plura-duople, quan onu aplikas sur korpo-parto malada, por ibe mantenar la pansuro, recevar la puso, e c.\lingui{DEFIRS}
\artiklo{kumular}%
\vorto{kumular}\fako{trans.}\senco{1}{\textit{(jurisprudenco)} Asemblar en sua persono (plura yuri, plura autoritati)}\senco{2}{Exercar plura funcioni nacionala salario-implikanta}\lingui{F}
\artiklo{kun}%
\vorto{kun}\senco{1}{Akompanata da…}\senco{2}{Karakterizata da… Juntita a…}
\artiklo{kuneiforma}%
\vorto{kuneiforma}\fako{arkeol.} Qua havas la formo di konio, di flecho-pinto (kom ex. la literi kuneiforma di la antiqua Asiriani, Medi e Persi)\lingui{DEFIRS}
\artiklo{kuniklo}%
\vorto{kuniklo}\fako{zool.} Quadripedo, ek la klaso « vageri », di qua la gambi dopa esas kelke plu longa kam le avana, e qua exkavas ter-truo quan lu habitas\latina{lepus cuniculus}\lingui{ISL}
\artiklo{kupelo}%
\vorto{kupelo}\fako{tekn.} Mikra vazo poroza, quan onu uzas por purigar oro ed arjento per litargiro, qua fortiras la aloyuro ed absorbesas kun lu da la parieti, dum ke la oro e la arjento restas maso en la fundo di la vazo\lingui{DEFIRS}
\artiklo{kupeo}%
\vorto{kupeo} Veturo *kluza (klozita), di qua la korpo havas la formo di berlino di qua la fako avana esabus tranchita\lingui{DEFIRS}
\artiklo{kuperoso}%
\vorto{kuperoso}\fako{patol.} Morbo di la pelo di la vizajo, karakterizata da makuli redatra\lingui{FIS}
\artiklo{kuplar}%
\vorto{kuplar}\fako{tekn.} Juntar du peci de fero per charniri, rivetaguri, e c.\lingui{DEFI}
\artiklo{kupleto}%
\vorto{kupleto} En kansono plura-strofa, omna parto quan finas refreno\lingui{EFIRS}
\artiklo{kupo}%
\vorto{kupo} Vazo drinkala, expansita, kun pedo, ek bronzo, arjento, oro, kristalo, e c.\lingui{EFIS}
\artiklo{kupolo}%
\vorto{kupolo}\fako{arkitekt.} Mi-sfero, di qua la formo esas olta di kupo subversita (kom ex. : la kupolo di Panthéon, en Paris; la kupolo di la katedralo Santa Paulo, en London; la kupolo di la Parlamento di Usa, en Washington)\lingui{DEFIRS}
\artiklo{kupono}%
\vorto{kupono} Parto de-tranchita, separita de ensemblo (ensemblo ek rent-akto, karto di porcionizo nutrivala, tabakala, dum indijo nacionala)\lingui{DEFIRS}
\artiklo{kupro}%
\vorto{kupro}\fako{kemio} Korpo simpla, metalo redatra, plu gisebla kam oro e kam arjento, tre duktila e martelagebla\simbolo{Cu}\lingui{DEFIS}
\artiklo{kupulo}%
\vorto{kupulo}\fako{bot.} Asembluro ek braktei, quaze soldita ye sua bazo, formacanta quaza kupo mikra, qua cirkondas la floro e salias pinto cirkum la frukto\lingui{DEFIS}
\artiklo{kuracao}%
\vorto{kuracao} Liquoro facita ek la shelo di oranji bitra\lingui{DEFIRS}
\artiklo{kuracar}%
\vorto{kuracar}\fako{trans.) (aludante mediko} Agar a malado ye la maniero quan lu judikas kom apta risanigor lu\lingui{DEFIRS}
\artiklo{kurajar}%
\vorto{kurajar}\fako{trans.}\senco{1}{Ne-angorar koram danjero}\senco{2}{Esar energioza koram obstaklo vinkenda}\senco{3}{Esar psike sat forta pri rezistar ulu, ulo}\lingui{DEFIRS}
\artiklo{kurar}%
\vorto{kurar}\fako{netrans.} \textit{(aludante homo, animalo)} Irar, per springi, plu rapide kam marche : I. esforcante devancar ulu; 2. esforcante atingar ulu; 3. esforcante forirar de ulu\lingui{FIS}
\artiklo{kuraro}%
\vorto{kuraro} Veneno ek planti, aden qua la Amerindiani di la sudo-regiono di Amerika trempas sua flechi\lingui{DEFIS}
\artiklo{kuraso}%
\vorto{kuraso}\senco{1}{Armo defensala ek ledro, metalo, e c., qua envolvas e protektas la pektoro e la dorso}\senco{2}{La squamaro qua envolvas ula fishi}\senco{3}{Kovrilo ek plaki de fero o de stalo qua protektas la kareno di ula milito-navi}\lingui{DEFIRS}
\artiklo{kuratar}%
\vorto{kuratar}\fako{yuro-cienco} Komisesar, segun la lego, pri defensor la havajo e la interesti sive di minoro, sive di majoro qua esas deklarita nekapabla administrar sua heredajo\lingui{DEFIRS}
\artiklo{kurbaturar}%
\vorto{kurbaturar}\fako{trans.) (aludante fatigo-kauzo ecesanta, o morbo} Igar extreme nekapabla pri esforco, pro la doloro specala qua afektas la muskuli pos granda esforco\lingui{F}
\artiklo{kurear}%
\vorto{kurear}\fako{netrans.) (aludante chaso kun hundi} Abandonar a la chaso-hundi parto de la bestio quan li kaptis\lingui{F}
\artiklo{kuretar}%
\vorto{kuretar}\fako{trans.} Deprenar la kalkoli de la veziko, lor la operaco « litotomio », per instrumento di qua un extremajo esas kulieratra, e qua uzesas por deprenar la materii qui normale ne esas en ta kavajo\lingui{F}
\artiklo{kuriero}%
\vorto{kuriero}\senco{1}{Persono qua vehas rapide, per kavalo, por arivar ante veturo o por preparar la relayi}\senco{2}{Portisto di depeshi qua cirkulas pede, kavale o bicikle}\senco{3}{Portisto di depeshi qua cirkulas per veturo ordinara, automobila, o per veturo postala}\lingui{DEFIRS}
\artiklo{kurio}%
\vorto{kurio}\fako{eklezio katol.} Ensemblo ek la tribunali, ek la organismi administrala, di la guverno-sistemo papala\lingui{DEFIRS}
\artiklo{kurioza}%
\vorto{kurioza} Qua deziras vidar, konocar, saveskar\lingui{DEFIRS}
\artiklo{kurkuliono}%
\vorto{kurkuliono}\fako{zool.} Insekto koleoptera, di qua un speco rodas frumento\latina{curculio}\lingui{ISL}
\artiklo{kurkumo}%
\vorto{kurkumo}\fako{bot.} Genero di planto herbacea, perena, quan onu renkontras en Azia, Afrika ed Amerika; la radiko di un speco di olu (la safrano di India) uzesas en la tinteyi, pro olua kolorizivo flava\latina{curcuma longa}\lingui{DEFIRS}
\artiklo{kurlio}%
\vorto{kurlio}\fako{zool.} Ucelo ek la familio « stiltieri » longa-beka\latina{numenius}\lingui{EFIS}
\artiklo{kursiva}%
\vorto{kursiva}\fako{pri skribo} Qua konsistas ek literi trasebla fluante, ed inklinata de dextre a sinistre. Dicesas anke pri la imprimo-tipo qua inklinesas same kam la skriburo ordinara\lingui{DEFIS}
\artiklo{kurso}%
\vorto{kurso}\senco{1}{Serio de lecioni qui formacas ensemblo docala}\senco{2}{Preco di kompro o di vendo, pri kambiajo}\senco{3}{Cirkulo regulala di kambiajo}\lingui{DEFIRS}
\artiklo{kursoro}%
\vorto{kursoro}\fako{tekn.} Bloko plata, mobila alonge linealo, limbo gradizita, e c., uzata kom indexo\lingui{EFIS}
\artiklo{kurta}%
\vorto{kurta}\senco{1}{Di qua la extenseso, de un de sua extremaji til la altra, sinse di longeso o di alteso, esas poka}\senco{2}{Di qua la duro esas poka}\lingui{DEFIS}
\artiklo{kurtajar}%
\vorto{kurtajar}\fako{netrans.} Agar kom mediacero, en komerco, inter la vendanto e la kompranto, po pago\lingui{FRS}
\artiklo{kurteno}%
\vorto{kurteno}\senco{1}{Peco de stofo tensita por celar o shirmar ulo od ulu}\senco{2}{\textit{(teatro)} Kurteno qua celas la *stejo kande ne pleesas}\senco{3}{Facio di la muro di remparo inter du bastioni}\lingui{EFIS}
\artiklo{kurtezar}%
\vorto{kurtezar}\fako{trans.} Prizentar (ad ulu) sua homaji por obtenar lua favoro; esforcar obtenar favoreso (da o da ulu, precipue da muliero) per anuncar o dicar a lu ulo qua interesas lu ed efikas a lua vanitato\lingui{DEFIS}
\artiklo{kurva}%
\vorto{kurva} Di qua omna elemento su eskartas, per deviaco preske neperceptebla, e segun *leyo (lego) fixa, de la rekto o de la plano. – \textit{(geom.)} Lineo qua esas nek rekta, nek kompozuro ek rekti\lingui{DEFIS}
\artiklo{kurvimetro}%
\vorto{kurvimetro} Instrumento por mezurar la longeso di la kurvi, sur papero\lingui{DEFIS}
\artiklo{kuseno}%
\vorto{kuseno} Mikra sako quadrata, ek stofo, ledro, polsterizita ye plumi, krini, e c., por ke onu povez apogar su an olu\lingui{EFIS}
\artiklo{kushar}%
\vorto{kushar}\fako{trans., sur}\senco{1}{Pozar aden lito, instalar en lito}\senco{2}{Pozar aden lito, instalar en lito, por repozar nokto-dorme}\lingui{F}
\artiklo{kuskuto}%
\vorto{kuskuto}\fako{bot.} Herbo parazita filoforma, qua plektokaptas la planti vicina e sufokigas li\latina{cuscuta}\lingui{FISL}
\artiklo{kuspido}%
\vorto{kuspido}\fako{geom.} Punto ube kurvo su dividas aden du bendi, kun tangento komuna a ta punto
\artiklo{kustar}%
\vorto{kustar}\fako{trans.}\senco{1}{\textit{(aludante objekto, laboro)} Igar necesa la pago di pekunio-quanto por obtenesar kambie}\senco{2}{\textit{(metaf.)} Igar necesa sakrifiko por obtenesor}\senco{3}{Havar kom konsequajo ulo quo impozas spenso}\lingui{DEFIS}
\artiklo{kustumar}%
\vorto{kustumar}\fako{trans.) (aludante populo, naciono od individuo} Agar konforme a la maniero establisita da longa experienco o praktiko\lingui{EFIS}
\artiklo{kutikulo}%
\vorto{kutikulo}\fako{anat.} La epidermo nevivanta qua cirkondas la bordo di la ungli di la fingri\lingui{EF}
\artiklo{kutino}%
\vorto{kutino}\fako{biol.} Substanco kemiala qua devenas de ula modifikeso di la celuloso ye la surfaco di la planto-korpo
\artiklo{kutio}%
\vorto{kutio} Telo densa, glata, simpla o fasonita, uzata por envelopar la matraci, o kom tapeto-stofo, o por facar somer-vesti\lingui{FS}
\artiklo{kuto}%
\vorto{kuto} Sorto di tuniko, vesto di homulo o di homino. – \textit{(olim)}.\lingui{DF}\subvorto{}{masho-kuto}{Tuniko ek felo sur qua esis suto-muntita dika bendi de ledro, lami, mashi de fero}{}{}
\artiklo{kutro}%
\vorto{kutro} Mikra milito-navo, kun nur un masto\lingui{DEFIS}
\artiklo{kuvento}%
\vorto{kuvento} Domo en qua su asemblas, por vivar konforme ad un regulo religiala, personi qui adheris ad ordeno\lingui{EFIS}
\artiklo{kuverto}%
\vorto{kuverto} Peco de papero faldita tale ke lu povas envelopar letro\lingui{DR}
\artiklo{kuvo}%
\vorto{kuvo} Granda vazo ek ligno, cirkla, facita ek duvi asemblita per cirkli fera, kun fundo\lingui{FS}
\artiklo{kuzo}%
\vorto{kuzo} Persono qua naskis de la onklulo, de la onklino, di ulu
\end{multicols}
\sekciono{L}
\begin{multicols}{2}
\artiklo{la}%
\vorto{la} Partikulo gramatikala qua, avan substantivo (expresata o tacata), indikas : (1) la tota speco, od omna individui di la speco; (2) un o plura individui determinata di la speco. Ta partikulo ne uzesas kun la nomi propra od di enti abstraktita, di qualesi, di vertui, di cienci, di astri, di sezoni, di monati
\artiklo{la}%
\vorto{« la »} Nomo di la noto sisesma di la « do »-gamo
\artiklo{labaro}%
\vorto{labaro} Standardo qua konsistas ek longa piquo trairata da bastono, qua subtenas banderolo de purpuro, e kronizita ye la aglo Romana qua portas, de lor la regno da la imperiestro Konstantino, la kruco e la monogramo di Iesu Kristo\lingui{DEFIS}
\artiklo{labio}%
\vorto{labio}\fako{anat.}\senco{1}{Parto karna qua garnisas extere la bordo di la boko}\senco{2}{\textit{(metaf.)} To di quo la formo aspektas kelke quale la labii}\lingui{DEFIS}
\artiklo{labirinto}%
\vorto{labirinto}\fako{epoki antiqua} Kluzajo kun cirkuiti komplikita, en qua la iranto desfacile *rekonocas sua voyo\lingui{DEFIRS}
\artiklo{labirintodonto}%
\vorto{labirintodonto}\fako{paleont.} Grupo de amfibii giganta, ek la klaso (qua ne plus existas) « stegocefali », kun vertebri tre ostoza, korpo quaze kurasa ek plaki derma, e di qui la denti havas kaneli dispozita labirinto-forme. (Onu renkontras tre multa de oli en la triaso di Britania)
\artiklo{laborar}%
\vorto{laborar}\senco{1}{\textit{(netrans.)} Esforcar (size, transporte, kombine, separe, e c.) igar ke la kozi naturala esez apta satisfacar la bezoni homala}\senco{2}{\textit{(trans.)} Transformar (generale per utensilo o mashino) la kozi naturala por igar ke li esez apta satisfacar la bezoni homala}\lingui{EFIS}
\artiklo{laboratorio}%
\vorto{laboratorio}\senco{1}{Chambro provizita ye la necesaji por la experimenti ciencala}\senco{2}{Chambro en qua la farmaciisti, la liquoro-distilisti, la bonbon-ifisti e c. preparas sua produkturi}\lingui{DEFIRS}
\artiklo{labro}%
\vorto{labro} Genero de fishi teleostea, familio « percidi », ek plura speci di la hemisfero boreala, kun labii dika\latina{labrus}
\artiklo{lacerar}%
\vorto{lacerar}\fako{trans.} Dispecigar arache, tale ke lu ne plus esos uzebla\lingui{EFIS}
\artiklo{lacerto}%
\vorto{lacerto}\fako{zool.} Reptero sauria, kun quar gambi kurta e tenua, kaudo longa ek ringi flexebla e facile ruptebla (ma ri-kreskas)\latina{lacerta}\lingui{EFIS}
\artiklo{laciva}%
\vorto{laciva} Inklinata a la plezuri amorala; qua su abandonas grosiere a la plezuri amorala, koitala\lingui{DEFIS}
\artiklo{laco}%
\vorto{laco} Kordono streta, di qua la extremaji esas provizita ye fero-pinto per qua onu asemblas du parti di korseto, di shuo, e c.\lingui{EFIS}
\artiklo{ladano}%
\vorto{ladano}\fako{farmacio} Gumo rezina aromata, quan donas diversa arbusti, precipue la cisto di Kretia\latina{ladanum}\lingui{DEFIRS}
\artiklo{lado}%
\vorto{lado} Fero-tolo dina, kovrita ek induturo de stano, qua protektas lu de rusto\lingui{IS}
\artiklo{lago}%
\vorto{lago} Mareto ek aquo stagnanta, interne di kontinento\lingui{EFIS}
\artiklo{laguno}%
\vorto{laguno}\fako{geogr.} Parto di maro, poke profunda, inter insuleti\lingui{DEFIRS}
\artiklo{laika}%
\vorto{laika} Qua ne resortisas ad eklezio\lingui{DEFIS}
\artiklo{lakeo}%
\vorto{lakeo} Servisto livreoza, employata ye akompanar seque sua mastro\lingui{DEFIRS}
\artiklo{lako}%
\vorto{lako}\senco{1}{Materio rezinoza, obtenata de ula arbori, mi-diafana, brune-reda, uzata por facar ula vaxi, tintivi, bela vernisi, e c.}\senco{2}{Verniso ek Chinia, tre prizata, nigra o reda, ornita ye figuri, arabeski, orizuri}\lingui{DEFIRS}
\artiklo{lakonika}%
\vorto{lakonika} Qua expresas lo pensata per tre poka vorti (segun la maniero di la habitanti di Lakonia, lor la epoki antiqua)\lingui{DEFIRS}
\artiklo{lakrimo}%
\vorto{lakrimo} Guto de humoro limpida, quan sekrecas glando di la okulo, efike da emoceso kordiala o da efekto kozala\lingui{EFIS}
\artiklo{laktazo}%
\vorto{laktazo}\fako{kemio} Diastazo, o fermento dissolvebla, sekrecata da kelka hefi, e qua transformas la laktoso a glikoso o galaktoso\lingui{DEFIRS}
\artiklo{lakto}%
\vorto{lakto} Liquido opaka, blanka, sukroza, sekrecata da la mamoglandi di la muliero e di la mamiferini, por la nutro di la jus-naskinti\lingui{EFIS}
\artiklo{laktometro}%
\vorto{laktometro} Areometro uzata por mezurar la denseso-grado di lakto\lingui{DEFIRS}
\artiklo{laktoso}%
\vorto{laktoso}\fako{kemio} Sukro C₁₂H₂₂O₁₁, en la lakto di la mamiferi\lingui{DEFIRS}
\artiklo{lakuno}%
\vorto{lakuno}\senco{1}{Spaco vakua en la kontinueso di korpo}\senco{2}{Vakuajo en texto, en kateniguro, en serio}\lingui{EFIS}
\artiklo{lamantino}%
\vorto{lamantino}\fako{zool.} Cetaceo herbivora\latina{manatus}\lingui{DEFI}
\artiklo{lamao}%
\vorto{lamao}\fako{zool.}\senco{1}{Quadripedo ruminera ek Peruvia, uzata kom charjo-bestio, e di qua la pili uzesas por facar texuri}\senco{2}{Sacerdoto di Buddha, en Tibet}\lingui{DEFIRS}
\artiklo{lamarckismo}%
\vorto{lamarckismo}\fako{biol.} Teorio (transformista) da Lamarck, qua asertas ke exkluzive la cirkonstanci extera efikas a la materio vivoza\lingui{DEFIRS}
\artiklo{lamelo}%
\vorto{lamelo}\senco{1}{\textit{(tekn.)} Mikra lamo de vitro sur qua onu pozas la objekto quan onu exploras per mikroskopo}\senco{2}{\textit{(bot.)} Apendico di la korolo ye qua konsistas la krono}\senco{3}{\textit{(historio naturala)} Mikra lamo ek bloko kayeroza, folioza}\lingui{DEFIS}
\artiklo{lamentar}%
\vorto{lamentar}\fako{netrans., pri, pro} Lasar explozar, lor la emoco de sufro, klameti iterata\lingui{DEFIS}
\artiklo{laminar}%
\vorto{laminar}\fako{trans.) (tekn.}\senco{1}{Transformar (bloko metala) a lami, per igar lu pasar inter du rotaco-cilindri, qui turnas inversa-sinse}\senco{2}{Igar cilindro metala tre pezoza rular sur (voyo) til kande ica divenabos glata de la ensuligeso ed aplasteso di la stoni o stoneti qui saliis sur lu}\lingui{EFIS}
\artiklo{lamo}%
\vorto{lamo}\senco{1}{Peco de metalo, de ligno, de ivoro, e c., plata, dina e streta}\senco{2}{Parto fera, stala, di instrumento, di utensilo, apta tranchar, taliar, e c.}\lingui{EFIS}
\artiklo{lampaso}%
\vorto{lampaso} Silko-stofo, quan onu obtenis olim de Chinia, kun granda desegnuri reliefa, maxim-multa-kaze sur fondo altrakolora\lingui{DEFIS}
\artiklo{lampiono}%
\vorto{lampiono} Kupeto qua kontenas materio kombustebla, kun mecho, uzata por la ilumini\lingui{FI}
\artiklo{lampiro}%
\vorto{lampiro}\fako{zool.} Insekto di qua la femino esas sen-ala e fosforecas en nia klimato, — aloza e fosforecoza pri la du sexui, en la landi tropikala\latina{lampyris}\lingui{EF}
\artiklo{lampo}%
\vorto{lampo} Utensilo konsistanta (origine) ye tanko qua kontenas liquido kombustebla, en qua esas imersita mecho quan onu acendas por produktar lumo. – Utensilo sama-skopo, en qua lumo produktesas da filo (metala o karba) inkandecoza per elektro\lingui{DEFIRS}
\artiklo{lampredo}%
\vorto{lampredo}\fako{zool.} Fisho ek la familio « ciklostomi », di qua la brankii esas sep-trua\latina{petromyzon}\lingui{DEFIS}
\artiklo{lancinar}%
\vorto{lancinar}\fako{netrans.) (pri doloro pikoza} Aparar bruske en parto di la korpo\lingui{FIS}
\artiklo{lanco}%
\vorto{lanco}\senco{1}{\textit{(epoki antiqua)} Armo, ek shafto finanta per ferajo quadriangula, pintoza, uzata sive quaze flecho quan onu lansas manue, sive quaze piquo per qua onu kombatas proxim la adverso}\senco{2}{\textit{(mezepoko)} Armo kun shafto longa e tre pezoza, finanta per ferajo pintoza, per qua la *chevaliero, su-precipitinte a sua adverso, probas renversar lu de lua kavalo, e trapikar lu}\senco{3}{\textit{(nia-epoke)} Armo kun longa shafto lejera, finanta per ferajo pintoza, ornita ye banderolo ek stofo, uzata da ula korpi di kavalrio (lancieri, dragoni)}\lingui{DEFIRS}
\artiklo{landau}%
\vorto{« landau »} Veturo quar-rota, kun suspenso-risorti e du tendi, forme di berlino, quan onu povas tegar o destegar segun-vole
\artiklo{landgravo}%
\vorto{landgravo} Titulo di ula suverena princi (olima) di Germania
\artiklo{lando}%
\vorto{lando} Dividuro de regiono geografiala, determinita per limiti naturala o politikala\lingui{DE}
\artiklo{lango}%
\vorto{lango}\fako{anat.}\senco{1}{Korpo longa, karnoza, moviva, qua esas en la kavajo boko}\senco{2}{\textit{(teknol.)} Nomo di diversa instrumenti langoforma}\lingui{FIS}
\artiklo{langorar}%
\vorto{langorar}\fako{netrans.} Esar kontinue depresita ed abatita (rispektive : psikale o korpale)\lingui{EFIS}
\artiklo{langusto}%
\vorto{langusto}\fako{zool.} Krustaceo analoga a la homardo, ma sen pinci e kun longa anteni\latina{palinurus vulgaris}\lingui{FS}
\artiklo{lanio}%
\vorto{lanio}\fako{zool.} Pasero kun beko kona e hokatra ye sua extremajo, qua kombatas uceli plu forta kam su e flugas lansante krii akuta\latina{lanius}\lingui{IL}
\artiklo{lanjo}%
\vorto{lanjo} Speco de stofo lana per qua onu envelopas la infanteti ankore kelka-monata\lingui{F}
\artiklo{lano}%
\vorto{lano} Pili dolca e tufatra, qui kreskas sur la pelo di la mutoni, e c.\lingui{FIS}
\artiklo{lansar}%
\vorto{lansar}\fako{trans.}\senco{1}{Igar departar tre rapide, aden direciono determinita}\senco{2}{\textit{(metaf.)} Igar funcionar, igar efikar}\lingui{DEFIS}
\artiklo{lantano}%
\vorto{lantano}\senco{1}{\textit{(bot.)} Planto exotika, ek la familio « verbenadei », kultivata precipue kom apartamento-planto}\senco{2}{\textit{(kemio)} Metalo, korpo simpla qua deskompozas aquo e qua divenas kombustebla en aero}\simbolo{La}\latina{lantana}\lingui{DEFIS}
\artiklo{lanterno}%
\vorto{lanterno} Sorto di buxo di qua la parieti esas garnisita ye materio plu o min diafana (korno, vitro, papero, e c.) en qua onu pozas lumo-fonto shirme de vento\lingui{DEFIS}
\artiklo{lanugo}%
\vorto{lanugo}\senco{1}{Mikra plumi dina e dolca qui kreskas (en la komenco di lua vivo) che la ucelo e formacas strato sub la plumi granda}\senco{2}{Mikra pili quin havas, lor sua nasko, la maxim multa quadripedi, e qui duras garnisar la korpo sub la pili plu forta qui formacas la pilaro.3 Sorto di kotono lejera qua tegas la stipo, la squami di la burjoneti, la folio, la frukto, di ula vejetanti}\senco{3}{Barbo naskanta, pilo tre tenua}\lingui{EFI}
\artiklo{lapar}%
\vorto{lapar}\fako{trans.) (aludante precipue la hundi} Drinkar per quaze pumpar la liquido per sua lango\lingui{DEF}
\artiklo{lapido}%
\vorto{lapido} Stono tre kustoza, kom ex. : rubino, agato, onixo, kornalino, esmeraldo, safiro, e c.\lingui{DEFS}
\artiklo{lapislazulo}%
\vorto{lapislazulo}\fako{mineral.} Stono opaka, blua, kun veini blanka, e kun punti ek piriti feroza qui cintilifas quale oro\lingui{DEFIRS}
\artiklo{lardo}%
\vorto{lardo} Graso qua garnisas la pelo di la porko e di kelka altra animali, precipue ye la ventro\lingui{FIS}
\artiklo{larico}%
\vorto{larico}\fako{bot.} Arboro konifera, de la trunko di qua dekulas rezino e di qua la folii sekrecas substanco viskozo\latina{larix}\lingui{DEFISL}
\artiklo{laringito}%
\vorto{laringito}\fako{patol.} Inflamuro di la laringo\lingui{DEFIS}
\artiklo{laringo}%
\vorto{laringo}\fako{anat.} Organo esencala di voco, formacita ek plura peci intermoviva, ye la parto supra di la trakeo\lingui{DEFIS}
\artiklo{laringoskopo}%
\vorto{laringoskopo}\fako{tekn.} Aparato per qua onu povas observar la parto interna di la laringo\lingui{DEFS}
\artiklo{laringotomio}%
\vorto{laringotomio}\fako{kirurg.} Operaco per qua onu facas aperturo en la laringo por preventar asfixio\lingui{DEFIS}
\artiklo{larja}%
\vorto{larja} Tre extensita en la sinso opozata a longeso\lingui{FI}
\artiklo{larvo}%
\vorto{larvo}\senco{1}{\textit{(epoki antiqua, en Roma)} Fantomo ledega, repugnanta}\senco{2}{\textit{(historio naturala)} Insekto vermoforma qua reprezentas la stando unesma di la metamorfosinsekti}\lingui{DEFIS}
\artiklo{lasar}%
\vorto{lasar}\fako{trans.}\senco{1}{Igar ulo laxa per cesar retenar lu}\senco{2}{Igar (ulu od ulo) ube lu esas – ne prenar lu kun su}\senco{3}{Igar ulo, gardata,- ne deprenar lu de ulu}\senco{4}{Igar tale ke ulu posedeskas ulo quon onu cesas gardar}\lingui{DFI}
\artiklo{lasho}%
\vorto{lasho}\senco{1}{\textit{(kirurg.)} Plako de ligno, de kartono, quan onu enpozas alonge membro ruptita, por mantenar la osti en plaso fixa}\senco{2}{\textit{(tekn.)} Plako de fero qua kunjuntas du reli}\lingui{D}
\artiklo{laskaro}%
\vorto{laskaro}\fako{navig.} Membro Indiana di la *kruo di navo\lingui{DEFIS}
\artiklo{lasta}%
\vorto{lasta}\senco{1}{Qua esas pos omni cetera, (1) pri tempo, (2) pri spaco, (3) pri rango en serio}\senco{2}{\textit{(*ultima)} Pos qua ne plus existas ceteri, (1) pri tempo, (2) pri spaco, (3) pri rango en serio}\lingui{DE}
\artiklo{latanio}%
\vorto{latanio}\fako{bot.} Palmiero kun larja folii dispozita abanike, quan onu vidas en Madagaskar ed en la insuli di Indonezia\latina{latania}\lingui{EFI}
\artiklo{latenta}%
\vorto{latenta} Qua ne su manifestas adextere.\lingui{DEFIS}\subvorto{}{kaloro latenta}{Kaloro quan absorbas korpo por liquideskar o vaporeskar, qua, pro ne manifestar su adextere, ne efikas a la termometro}{}{}
\artiklo{latero}%
\vorto{latero}\senco{1}{\textit{(anat.)} La tota parto, dextra o sinistra, di la korpo homala}\senco{2}{Facio, parto qua esas dextre o sinistre che objekto. Singla ek la facii interopozanta di objekto}\senco{3}{\textit{(geom.)} Linei qui limitizas surfaco}\lingui{EFIS}
\artiklo{latexo}%
\vorto{latexo}\fako{bot.} Suko propra di ula planti (euforbi, figieri, papaveri, latugi, e c.) e qua cirkulas en vaskuli partikulara\lingui{DEF}
\artiklo{laticifera}%
\vorto{laticifera}\fako{bot.} Qua produktas latexo\lingui{DEF}
\artiklo{latiro}%
\vorto{latiro}\fako{bot.} Planto klimera, di qua la flori esas odoroza\lingui{DEF}
\artiklo{latitudo}%
\vorto{latitudo}\senco{1}{\textit{(kosmogr.)} Disto de loko ad equatoro, mezurata per arko di la meridiano terala}\senco{2}{Regiono submisata a ta o ca rejimo atmosferala, segun ke lu esas plu o min for equatoro}\lingui{EFIRS}
\artiklo{lato}%
\vorto{lato} Peco de ligno fendita segun la longeso di la fibri, longa, dina e plata, uzata por facar parieti, plafoni, e c. e qua suportas la ardezo-plaki o la teguli di tekto\lingui{DEFS}
\artiklo{latrino}%
\vorto{latrino} Loko rezervita por kakar od urinifar\lingui{DEFIS}
\artiklo{latugo}%
\vorto{latugo}\fako{bot.} Planto leguma, qua kontenas suko laktatra, ek la tribuo « chikorei »\latina{lactuca}\lingui{IS}
\artiklo{latuno}%
\vorto{latuno} Metalo ek aloyuro de kupro e de zinko\lingui{DFRS}
\artiklo{laubo}%
\vorto{laubo} Treliso-muri paralela, kun mi-cilindro trelisa kom tekto, kovrita per planti klimera tre folioza\lingui{D}
\artiklo{laudano}%
\vorto{laudano}\fako{farmacio}\senco{1}{\textit{(olim)} Opiumo purigita, moligita en aquo e filtrita per preso}\senco{2}{Extraktajo de opiumo, preparita per vino, alkoholo, e c.}\lingui{DEFIS}
\artiklo{laudar}%
\vorto{laudar}\fako{trans.} Igar eminenta per la aprobo di lua meriti, di lua bona qualesi\lingui{EFIS}
\artiklo{laureato}%
\vorto{laureato}\senco{1}{Qua obtenis krono ek lauro}\senco{2}{Qua obtenis premio en konkurso}\lingui{DEFIRS}
\artiklo{lauro}%
\vorto{lauro}\senco{1}{\textit{(bot.)} Arboro aromatoza, sempre verda, kun folii permananta, lanco-forma, qua esas la tipo di la familio « laurinei », konsakrita ad Apolono da la Antiqui, qui facis ekodu palmi, kroni por la vinkinti lor la ludi poeziala, e kredis ke olua foliaro protektas de la fulmino-stroki}\lingui{DEFIRS}\subvorto{metaf.}{la lauri}{La glorio di la vinkinti}{2}{}
\artiklo{lauso}%
\vorto{lauso}\fako{zool.} Insekto parazita sen-ala, qua su akrochas a la hari ed a la korpo di la homo, a la pili di la animali\latina{pediculus}\lingui{DE}
\artiklo{lavao}%
\vorto{lavao}\senco{1}{Materio fuzanta qua eskapas de volkano}\senco{2}{Ta materio, koldeskinta e solideskinta}\lingui{DEFIRS}
\artiklo{lavar}%
\vorto{lavar}\fako{trans.} Deprenar (to quo sordidigas) per dilutar lu en aquo\lingui{DEFIS}
\artiklo{lavareto}%
\vorto{lavareto}\fako{zool.} Fisho qua vivas en aquo sen-sala, qua saporas delikate, analoga a truto
\artiklo{lavendo}%
\vorto{lavendo}\fako{bot.} Planto aromatoza ek la familio « labiei »\latina{lavendula}\lingui{DEFIR}
\artiklo{laxa}%
\vorto{laxa} Qua ne esas tensita\lingui{DEFIS}
\artiklo{lazareto}%
\vorto{lazareto} Edifico izolita ube, en ula portui, onu impozas quaranteno a la personi ed a la objekti qui devenas de regioni infektita de morbo kontagiala\lingui{DEFIS}
\artiklo{lazarone}%
\vorto{« lazarone »} Mendikanto di urbo Napoli (Italia)
\artiklo{lazo}%
\vorto{lazo} Ledro-bendo garnisita ye buli quan la habitanti di la landi olim Hispana, di Amerika lansas tale ke li plekte kaptas sua enemiko o la animalo quan li chasas\lingui{DEFIS}
\artiklo{le}%
\vorto{le} Pluralo di « la », kande ne existas altra indikivo pri la pluraleso en la propoziciono
\artiklo{lea}%
\vorto{lea} Qua rezolvis servor Deo : (1) Qua konsakris su erste tarde a la vivo monakeyala (opoze ad oblata); (2) Qua exekutas la taski muskulala en la kongregaciono di qua lu esas parto\lingui{EF}
\artiklo{leciono}%
\vorto{leciono} Exerco per qua docanto docas od igas lernar ta o ca parto di cienco, di arto\lingui{DEFIRS}
\artiklo{leda}%
\vorto{leda} Qua aspektas desagreabla pro ne-beleso. Qua desplezas etikale o morale\lingui{FI}
\artiklo{ledro}%
\vorto{ledro} Pelo dika, deprenita de la karno di ula animali (bovulo, bovino, bovyuno, bufalo, e c.), senpiligita per tanago, e preparita por uzi diversa\lingui{DE}
\artiklo{legacar}%
\vorto{legacar}\fako{trans., ad} Donacar (ad ulu) per testamento qua atribuas ad un od a plura personi parto di sua havaji o nur parto de olci\lingui{DEFIS}
\artiklo{legato}%
\vorto{legato}\senco{1}{Delegito di la imperiestri Romana, qua havis kom tasko reprezentar ici en la provinci}\senco{2}{\textit{(eklezio katol.)} Kardinalo delegita da papo por guvernar un ek la provinci di la eklezio katolika}\lingui{DEFIRS}
\artiklo{legendo}%
\vorto{legendo} Serio de rakonti populala, qui havas maxim-multa-kaze fundo exakta, ma developita, transformita da la tradicioni\lingui{DEFIRS}
\artiklo{legiono}%
\vorto{legiono}\senco{1}{(1)\textit{(epoki antiqua, en Roma)} Armeo-korpo qua konsistis ye infantrio e kavalrio. (2) \textit{En Francia}, lor la rejo Francisko I. (cirkum la yaro 1515) : armeo-korpo permananta}\senco{2}{\textit{(metaf.)} Grupo de individui qui esforcas atingar la sama skopo}\senco{3}{\textit{(en Francia)} Legiono di Honoro : meritordeno, civila e militala, institucita da Bonaparte, konsulo rang-unesma}\lingui{DEFIRS}
\artiklo{legitima}%
\vorto{legitima} Konsakrita da la lego
\artiklo{lego}%
\vorto{lego}\senco{1}{Ago-regulo impozata da autoritatozo supera}\senco{2}{Ago-regulo impozata a la homo da lua raciono, lua koncienco}\senco{3}{\textit{(*leyo)} Regulo konstanta, universala, a qua omna fenomeni di naturo esas obedienda. – Rezultajo de analizar exakte e precize la cirkonstanci qui esas elementi en la produkto di la fenomeni, ed interligar ta cirkonstanci per la relati normala di inter-sucedo e di inter-simileso. (Auguste Comte)}\lingui{DEFIRS}
\artiklo{leguminoso}%
\vorto{leguminoso}\fako{bot.} Di qua la frukto esas shelo qua kontenas maxim-multa-kaze legumo (pizo, fabo, fazeolo)\lingui{DEFIS}
\artiklo{legumo}%
\vorto{legumo} Parto di planto koliita por uzesar kom nutrivo : lensi, fazeoli, pizi, asparagi, e c.\lingui{DEFIS}
\artiklo{lejera}%
\vorto{lejera} Qua pezas nur poke\lingui{FIS}
\artiklo{lekar}%
\vorto{lekar}\fako{trans.} Pasigar sua lango sur (ulo)\lingui{DEFI}
\artiklo{lektar}%
\vorto{lektar}\fako{trans.}\senco{1}{Dicernar, en imprimuro o skriburo, la son-uni reprezentita per literi}\senco{2}{Konoceskar la kontenajo di skriburo, di libro}\senco{3}{Konocigar, savigar, da altra personi la kontenajo di skriburo, di libro, per pronuncar koram li to quo esas skribita, imprimita}\senco{4}{\textit{(metaf.)} Divinar la penso da ulu, lua sentimenti, ek lua konduto, fiziognomio, aspekto extera, e c.}\lingui{FIRS}
\artiklo{lektoro}%
\vorto{lektoro}\senco{1}{Sorto di profesoro adyunta en ula universitati}\senco{2}{\textit{(liturgio katolika)} Kleriko kun la duesma ek la quar ordeni minora, qua havis olim kom tasko lektar, kantar la lecioni, epistoli, e c., en la ceremonii kultala}\lingui{DEFIS}
\artiklo{lemniskato}%
\vorto{lemniskato}\fako{geom.} Kurvo ek la grado quaresma, qua havas la formo di la cifro « 8 », e qua esas la ligilo di omna punti tale ke la produto di lia disti a du punti fixa esas konstanta\lingui{DEFI}
\artiklo{lemo}%
\vorto{lemo}\fako{matem.} Propoziciono preliminara quan onu establisas por preparar la demonstro di altra propoziciono\lingui{DEFIRS}
\artiklo{lemuriano}%
\vorto{lemuriano}\fako{zool.} Familio de simii, di qui la tipo esas la makio\lingui{DEFIS}
\artiklo{lengo}%
\vorto{lengo}\fako{zool.} Speco di gado, di qua la korpo esas oblonga\latina{gadus molva}
\artiklo{lenio}%
\vorto{lenio} Peco di kalorizo-ligno, segita o tranchita tale ke onu povas pozar lu aden kameno, aden furnelo\lingui{IS}
\artiklo{lenso}%
\vorto{lenso}\senco{1}{\textit{(bot.)} Grano di leguminoso quan onu manjas kom legumo pos lua sikigeso}\senco{2}{\textit{(fiziko)} Vitro-bloko, taliita lenso-forme, finanta sive per du surfaci sferala, sive per surfaco plana, qua produktas deviaco regulala di la lumo-radii}\lingui{EF}
\artiklo{lenta}%
\vorto{lenta}\senco{1}{Qua ne iras rapide}\senco{2}{Qua trodurigas to quon lu esas aganta o facanta}\lingui{FIS}
\artiklo{lentigino}%
\vorto{lentigino}\fako{che ula personi} Altereso naturala, grano-forma, di la koloro di lia pelo\lingui{EFIS}
\artiklo{lentisko}%
\vorto{lentisko}\fako{bot.} Arboreto qua kreskas en la Franca provinco Provence ed en Italia, e di qua un speco donas la suko rezinoza de qua onu facas liquoro e mastico\latina{Pistacia lentiscus}\lingui{EFIS}
\artiklo{leono}%
\vorto{leono}\senco{1}{\textit{(zool.)} Mamifero karnivora, ek la familio « felini », kun pilkoloro flava, kun la shultri kovrata (che la maskulo) da dika krinaro}\senco{2}{\textit{(astrol.)} Signo di zodiako, reprezentata da leono}\senco{3}{\textit{(metaf.)} (1) viro brava, audacoza quale leon(in)o. (2) Personego tre segun-moda, e tre admirata, che la mondumani}\lingui{EFIS}
\artiklo{leopardo}%
\vorto{leopardo}\fako{zool.} Animalo karnavida, ek la genero « kato », kun pelo makuletoza\latina{felis leopardus}\lingui{DEFIRS}
\artiklo{lepidoptero}%
\vorto{lepidoptero}\fako{zool.} Familio ek insekti kun ali squamoza, nomata vulgare « papilioni », qui subisas metamorfosi kompleta, e di qui onu nomizas la larvo « raupo », e la nimfo « krizalido »\lingui{DEFIS}
\artiklo{leporo}%
\vorto{leporo}\fako{zool.} Quadripedo rodera, di qua la gambi avana esas plu kurta kam le dopa, e kun oreli longa\latina{lepus timidus}\lingui{FIS}
\artiklo{lepro}%
\vorto{lepro}\fako{patol.} Morbo di qua la fonto esas tuberkloso, qua invadas, rodas e repugnante sulkizas la pelo\lingui{DEFIS}
\artiklo{lernar}%
\vorto{lernar}\fako{trans.} Aquirar la savo (di ulo) per spirito-laboro\lingui{DE}
\artiklo{lesivar}%
\vorto{lesivar}\fako{trans.} Netigar (linjo) per lasar lu dum kelka tempo imersita en aquo varma a qua onu adjuntis sodo, natroxo o potaso, e qua esas apta dissolvar, ek la linjo sordida, la kraso e la nepuraji qui esas an lu, adherante\lingui{FI}
\artiklo{letargio}%
\vorto{letargio}\fako{patol.} Stando maladesala quan karakterizas dormo profunda e tre duroza, dum qua la fenomeni videbla di vivo esas interruptita, e qua prizentas la aspekto di morto\lingui{DEFIRS}
\artiklo{letro}%
\vorto{letro} Skriburo quan onu sendas a persono absenta, por komunikar a lu to quon onu ne povas dicar a lu voce\lingui{EFIS}
\artiklo{leucito}%
\vorto{leucito}\fako{kemio} Duopla silikato di alumino e di potaso, quan onu renkontras en la tereni volkan-agita\lingui{DEFIRS}
\artiklo{leukocito}%
\vorto{leukocito}\fako{anat.} Celulo senkolora di sango e di limfo, nukleoza, deformebla, e di qua la dimensioni esas plu granda kam olti di hematito\lingui{DEFIRS}
\artiklo{leukomo}%
\vorto{leukomo} Makulo blanka, opaka, qua naskas ulakaze sur la korneo di la okulo\lingui{DEFIRS}
\artiklo{levar}%
\vorto{levar}\fako{trans.} Igar movata, de infre ad supre\lingui{EFIS}
\artiklo{levero}%
\vorto{levero} Stango rigida, moviva sur apogo-punto, e quan onu uzas por levar, por igar objekto movar\lingui{EFI}
\artiklo{leviatano}%
\vorto{leviatano}\fako{biblo} Monstro marala\lingui{DEFIRS}
\artiklo{levito}%
\vorto{levito}\senco{1}{\textit{(biblo)} Izraelido (ek la tribuo di Levi) konsakrita a la oficio en la Templo}\senco{2}{\textit{(metaf.)} \textit{(poezio)} Sacerdoto}\lingui{DEFIRS}
\artiklo{levriero}%
\vorto{levriero}\fako{zool.} Hundo kun muzelo akuta, korpo tenua, gambi longa, di qua la aluro esas rapida, uzata partikulare por persequar la lepori\lingui{FIS}
\artiklo{lexiko}%
\vorto{lexiko} Vortolibro qua indikas lakonike (generale : per un vorto) la senco di la vorti di altra linguo\lingui{DEFIRS}
\artiklo{lexikografio}%
\vorto{lexikografio} Ensemblo ek la laboruro de lexikografo
\artiklo{lexikografo}%
\vorto{lexikografo} La persono qua exploras, studias, la origino e la senci di la vorti di linguo\lingui{DEFIS}
\artiklo{lexikologio}%
\vorto{lexikologio} Cienco pri la origino e la senci di la vorti di linguo
\artiklo{*leyo}%
\vorto{*leyo} Regulo konstanta, universala, a qua omna fenomeni di naturo esas obedienda. — Rezultajo de analizar exakte e precize la cirkonstanci qui esas elementi en la produkto di la fenomeni, ed interligar ta cirkonstanci per la relati normala di inter-sucedo e di inter-simileso. (Aug. Comte)
\artiklo{lezar}%
\vorto{lezar}\fako{trans.}\senco{1}{\textit{(fiziol.)} Igar ne-uzebla parte}\senco{2}{\textit{(yuro-cienco)} Detrimentar per desegardo di (la yuri)}\lingui{DEFIS}
\artiklo{li}%
\vorto{li}\fako{gram.} Pronomo personala, triesma persono, pluralo, formo epicena
\artiklo{liano}%
\vorto{liano}\fako{bot.} Planto sarmentoza, klimera, di la foresti di Amerika, qua plekte kaptas la arbori e formacas densaji nedesintrikebla\lingui{DEFIRS}
\artiklo{liaso}%
\vorto{liaso}\fako{geol.} Tereno kalkoza, marnoza, arjiloza\lingui{DEFIS}
\artiklo{libacionar}%
\vorto{libacionar}\fako{netrans.) (religio antiqua} Disvarsar adsule, honorize di deajo, mielo, lakto, oleo e precipue vino pura\lingui{DEFIS}
\artiklo{libelo}%
\vorto{libelo}\fako{tekn.} Instrumento uzata por trasar lineo, plano, horizontala, o por mezurar, relate plano horizontala, la elevaciono di lineo, di plano qua esas paralela a lu
\artiklo{libelulo}%
\vorto{libelulo}\fako{zool.} Insekto nevroptera, kun korpo dina ed ali diafana\latina{libellula}\lingui{DEFIS}
\artiklo{libera}%
\vorto{libera}\senco{1}{\textit{(aludante individuo)} Qua ne apartenas a mastro}\senco{2}{\textit{(aludante statano)} Qua ne dependas de autoritato arbitriala}\senco{3}{\textit{(aludante stato)} Qua ne esas sklavigita da lando stranjera}\senco{4}{\textit{(aludante psiko)} Di qua la volo povas determinar su sen subisar koakto}\senco{5}{\textit{(aludante la korpo, la spirito, la kordio di la homo)} Di qua la agado ne renkontras obstaklo}\senco{6}{\textit{(aludante la kozi)} Qua ne prizentas obstaklo}
\artiklo{libertina}%
\vorto{libertina}\senco{1}{\textit{(olim)} Qua su liberigis de la autoritato di la religio kristana, di la kredi, di la diciplino}\senco{2}{Qua su liberigis de omna regulo, de omna autoritato}\senco{3}{Di qua la mori sexuala esas sen-regulo}\lingui{DEFIS}
\artiklo{libraco}%
\vorto{libraco}\fako{astron.} Ocilo videbla di astro cirkum sua axo\lingui{DEFI}
\artiklo{libro}%
\vorto{libro}\senco{1}{Reprodukturo imprimura di la verko da autoro, ek paperfolii asemblita aden kayeri quin onu inter-asemblas}\senco{2}{La verko tale *reproduktita}\senco{3}{Singla ek la parti di ula verki}\lingui{EFIS}
\artiklo{licenco}%
\vorto{licenco}\senco{1}{\textit{(en Francia)} Diplomo universitatala, inter bachelereso-diplomo e doktoreso-diplomo, e qua yurizas pri docar en licei o kolegii, o pri pledar, e c.}\lingui{DEFIRS}\subvorto{}{licenco poeziala}{Su-liberigo, da poeto, pri la strikta reguli prozodiala}{2}{}
\artiklo{licensar}%
\vorto{licensar}\fako{trans.) (aludante stat-administrantaro} Koncesar partikulare la yuro pri debitar ula vari konsumebla, o pri peskar, chasar, e c., en ta o ca regiono\lingui{DEFIRS}
\artiklo{liceo}%
\vorto{liceo}\fako{en Francia} Edifico en qua agesas la doco sekundara, direktata da la stato (kontraste kun « kolegio »)\lingui{FI}
\artiklo{lico}%
\vorto{lico} Tereno *kluza per barili, en qua luktas la personi qui partoprenas turniri, konkursi, e c.\lingui{EFIS}
\artiklo{lido}%
\vorto{lido} Parto movebla quan onu pozas sur (o qua retrovenas ad) la aperturo di vazo, di kofro, e c., por kovrar o klozar lu\lingui{E}
\artiklo{lienterio}%
\vorto{lienterio}\fako{patol.} Diareo en qua la nutrivi esas forjetita mi-digestita
\artiklo{lietnanto}%
\vorto{lietnanto} Oficiro di qua la grado esas nemediate sub olta di kapitano (en armeo)\lingui{DEFIRS}
\artiklo{*lifto}%
\vorto{*lifto} Mashino qua elevas la personi a la etaji diversa di edifico\lingui{EF}
\artiklo{ligamento}%
\vorto{ligamento}\fako{anat.} Fasko de tisuo fibroza densa, poke extensebla, qua ligas osti, kartilagi, muskuli\lingui{DEFIS}
\artiklo{ligar}%
\vorto{ligar}\fako{trans.}\senco{1}{Cirkondar la parti di objekto o di plura objekti, por interjuntar li, per ulo flexebla, ulo longa (kordo, rubando, e c.)}\senco{2}{Juntar (plura kozi) per relato di sucedo, konsequo, e c.}\senco{3}{Unionar (plura personi) per relati di socio, di amikeso, e c.}\lingui{EFIRS}
\artiklo{lignino}%
\vorto{lignino} Substanco kemiala qua impregnas la elementi di ligno, e qua donas ad ica lua grado di fermeso, di solideso\lingui{DEFIS}
\artiklo{lignito}%
\vorto{lignito} Karbono fosila qua konservas en su, plu o min multe, traci di konstituciono vejetala\lingui{DEFIRS}
\artiklo{ligno}%
\vorto{ligno} Parto dura, ek la tisuo fibroza di arboro, di qua la pezo specifika e la dikeso-grado varias segun la speci e la evo\lingui{EFIS}
\artiklo{*liguo}%
\vorto{*liguo} Asociitaro formacita interne di stato por igar triumfar ula interesti politikala, religiala\lingui{EF}
\artiklo{ligustro}%
\vorto{ligustro}\fako{bot.} Arboreto ek la familio « oleacei », de qua onu formacas masivi\latina{ligustrum}\lingui{DISL}
\artiklo{likantropo}%
\vorto{likantropo}\fako{patol.} Persono qua afektesas da la morbo psikala pro qua lu imaginas transformesir a volfo\lingui{DEFIS}
\artiklo{likar}%
\vorto{likar}\fako{netrans.) (aludante generale liquido} Eskapar tra fenduro (de vazo, e c.)\lingui{DE}
\artiklo{likeno}%
\vorto{likeno}\fako{bot.} Planto kriptogama qua kreskas sur la kortico di la arbori, sur la roki, sur la petri humida, forme di krusto pulveratra\latina{cetraria}\lingui{DEFIS}
\artiklo{likopodio}%
\vorto{likopodio}\fako{bot.} Planto kriptogama, di qua la kapsuli kontenas pulvero inflamebla, uzata medicinale kom sikirivo\latina{lycopodium}\lingui{DEFIS}
\artiklo{liktoro}%
\vorto{liktoro}\fako{epoki antiqua, en Roma} Guardano qua iris avan la oficieri judiciala, portis hakilo plasizita meze di vergo-fasko ed havis kom tasko exekutar la verdikto qua kondamnis la krimininti a senkapigeso pos batesir per vergi\lingui{DEFIS}
\artiklo{lilaco}%
\vorto{lilaco}\fako{bot.} Arboreto di la genero « seringo », di qua la flori dispozesas grape\latina{Syringa}
\artiklo{liliaceo}%
\vorto{liliaceo}\fako{bot.} Di qua la naturo esas olta di lilio (planto monokotiledona)\lingui{DEFIS}
\artiklo{lilio}%
\vorto{lilio}\fako{bot.} Planto herbacea, bulboza kun stipo rekta, e di qua la varietato maxim konocata havas floro pure-blanka e klosho-forma\latina{lilium}\lingui{DEFIRS}
\artiklo{limako}%
\vorto{limako}\fako{zool.} Molusko gasteropoda, reptera, di qua la korpo esas oblonga, sen-konka, kovrata da humoro viskoza\latina{limax}\lingui{EF}
\artiklo{limando}%
\vorto{limando}\fako{zool.} Fisho plata, ek la genero « plii », di qua la pelo esas rugoza, aspera\latina{pleuronectes}\lingui{FIL}
\artiklo{limar}%
\vorto{limar}\fako{trans., per} Konsumar, planigar per froto-kontakto kun lamo de fero, de stalo, kun strii transversa qui aspektas quale rangi de denti\lingui{FS}
\artiklo{limbo}%
\vorto{limbo}\senco{1}{\textit{(bot.)} Parto supra, ordinare expansita, di la koroli monopetala, di la kalici monofila}\senco{2}{\textit{(tekn.)} Bordo gradizita, di la cirklo uzata por mezurar la anguli}\lingui{EFIRS}
\artiklo{limeto}%
\vorto{limeto}\fako{bot.} Frukto di arboro ek la genero « oranjieri », qua produktas mikra citrono dolca\latina{citrus limeta}\lingui{DEFISL}
\artiklo{limfatismo}%
\vorto{limfatismo}\fako{patol.} Ta temperamento quan karakterizas la moleso fiziologiala e la pelo preske sen-kolora\lingui{DEFIRS}
\artiklo{limfo}%
\vorto{limfo}\fako{anat.} Liquido blankatra qua cirkulas en sistemo de vaskuli partikulara, paralela a la sistemo por la cirkolo di sango, e di qua parto, pos la digesto, transportas la substanci asimilebla extraktita de la nutrivi\lingui{DEFIRS}
\artiklo{limito}%
\vorto{limito}\senco{1}{Parto extrema ube finas teritorio, domeno}\senco{2}{Punto ube cesas la apliko di povo, di fakultato, di yuro}\lingui{DEFIS}
\artiklo{limonado}%
\vorto{limonado} Drinkajo boko-refreshigiva, qua konsistas ye suko de citrono o de limono, en aquo\lingui{DEFIS}
\artiklo{limono}%
\vorto{limono}\fako{bot.} Frukto analoga a citrono, ma plu acida\latina{citrus limonus}\lingui{DEFIRS}
\artiklo{limpida}%
\vorto{limpida} Di qua nulo trublas la diafaneso\lingui{EFIS}
\artiklo{linchar}%
\vorto{linchar}\fako{trans.} Punisar, ocidar sen vartar la verdikto di la judicio-korto\lingui{DEFIRS}
\artiklo{linco}%
\vorto{linco}\fako{zool.} Quadripedo karnavida, sorto di kato sovaja, kun kaudo kurta, e pri qua la antiqui dicis ke lua vido esas tre granda-portea\latina{felis lynx}\lingui{EFISL}
\artiklo{linealo}%
\vorto{linealo} Instrumento longa e rekta, uzata por trasar rekti\lingui{DIR}
\artiklo{lineamento}%
\vorto{lineamento} Lineo elementala qua indikas, ye maniero generala, la formo, la konturo
\artiklo{lineara}%
\vorto{lineara}\fako{algebro} Dicesas pri equaciono unesma-grada\lingui{DEFIS}
\artiklo{lineo}%
\vorto{lineo}\senco{1}{Streko kontinua qua indikas direciono determinita}\senco{2}{\textit{(geom.)} Extensajo una-dimensiona}\senco{3}{Direciono kontinua ye sinso determinita}\senco{4}{Serio de vorti lokizita horizontale en paginedo skribura od imprimura}\senco{5}{Serio de faketi horizontala, vertikala o diagonala en damo-planko o shako-planko}\senco{6}{Serio de telegrafofili elektrala qui igas interkomunikar du loki}\senco{7}{Ensemblo ek la instaluri di konstrukturi ed edifici di relvoyo, tramvoyo, autobuso-exploto, aeroplano-exploto}\lingui{DEFIRS}
\artiklo{lingoto}%
\vorto{lingoto} Peco de metalo (precipue de metalo tre valoroza : oro, platino, arjento) qua ne ja laboresis ed havas ankore la formo di la mulduro en qua lu gisesis\lingui{EFIS}
\artiklo{linguistiko}%
\vorto{linguistiko} La cienco di la persono qua su konsakras a la studio di la lingui\lingui{DEFIRS}
\artiklo{linguo}%
\vorto{linguo} Ta maniero expresar sua penso per parolo o skribo, propra a ta o ca naciono\lingui{DEFIRS}
\artiklo{linimento}%
\vorto{linimento}\fako{farmacio} Medikamento grasatra, uncioniva, por la fricioni\lingui{DEFIS}
\artiklo{linjo}%
\vorto{linjo} Tuko ek filo o kotono, specaligita por diversa uzi hemala\lingui{F}
\artiklo{lino}%
\vorto{lino}\fako{bot.} Planto herbatra, kun floro blua, e di qua la stipo donas filaso uzata por facar filo, tre dina telo\latina{linum}\lingui{DEFIRS}
\artiklo{linoleum}%
\vorto{« linoleum »} Sorto di texajo nepermeebla, ek telo de juto, indutita ye lin-oleo e ye pudro de korko, uzata kom tapiso-materio\lingui{DEFIRS}
\artiklo{linono}%
\vorto{linono} Telo de lino, tre klara, kun apreto ferma, por robi, shultro-tuki, mulier-kamizi\lingui{DFIS}
\artiklo{lintelo}%
\vorto{lintelo} Traverso ye qua konsistas la parto supra di pordo, di fenestro kun aperturo rektangula, e subtenanta la masonuro\lingui{EFS}
\artiklo{lipitudo}%
\vorto{lipitudo}\fako{patol.} La stando di persono di qua la okuli esas seboza
\artiklo{lipomo}%
\vorto{lipomo}\fako{patol.} Tumoro grasoza\lingui{DEFIRS}
\artiklo{lipotimio}%
\vorto{lipotimio}\fako{patol.} Esvano
\artiklo{liquacar}%
\vorto{liquacar}\fako{teknol.} Separar un ek la metali kontenata en erco, en aloyuro\lingui{EFIS}
\artiklo{liquida}%
\vorto{liquida} Di qua la molekuli, poke interadheranta, fluas, obediante separe la legi gravitala kande li ne kontenesas da recipiento\lingui{EFIS}
\artiklo{liquidacar}%
\vorto{liquidacar}\fako{trans.} Konkluzar estado financala per pagar la pasivo e per atribuar a la yurozo la aktivo restanta\lingui{DEFIRS}
\artiklo{liquoro}%
\vorto{liquoro}\senco{1}{Drinkajo alkoholoza, obtenata per distilo}\senco{2}{Mixajo liquida uzata farmaciale, industriale, e c.}\lingui{DEFIRS}
\artiklo{lirika}%
\vorto{lirika}\senco{1}{Qua kontenas poemi drama, akompanata da kanti e muziko}\senco{2}{Qua konservis (quankam ne destinita a kanteso) la formo di la poezio lirika di la antiqui}\lingui{DEFIRS}
\artiklo{liro}%
\vorto{liro}\senco{1}{Muzik-instrumento kordoza, quan uzis la antiqui}\senco{2}{\textit{(astron.)} Stelaro en la hemisfero nordala}\senco{3}{\textit{(zool.)} Ucelo galinacea}\latina{lira}\lingui{DEFIRS}
\artiklo{listelo}%
\vorto{listelo}\fako{arkitekt.} Mikra muluro quadrata e plata, qua separas du altri plu granda e plu salianta – ed uzesas por reprezentar kanalo sur pilastro, o kolono\lingui{DEFIRS}
\artiklo{listo}%
\vorto{listo} Serio de nomi di personi o di kozi, de numeri, de nombri, e c., skribita ica dop ita\lingui{DEFIRS}
\artiklo{litanio}%
\vorto{litanio}\senco{1}{\textit{(religio katol.)} Prego forme di serio de advoki a Iesu Kristo, a la santa Virgino, e c., en qui la sama formulo di supliko repetesas multafoye. (Uzesas ordinare en pluralo.)}\senco{2}{Enumero tedoza}\lingui{DEFIRS}
\artiklo{litargiro}%
\vorto{litargiro}\fako{kemio} Protoxido di plombo fuzita, mi-vitrigita, uzata por preparar la plombo-sali qui esas un ek la kompozanti di kristalo e per qua onu falsigis ula vino-sorti por diminutar lia grado di acideso\lingui{DEFIRS}
\artiklo{literaturo}%
\vorto{literaturo} Ensemblo di la stilo-verki (en qui preponderas la beleso di kompozo, di expreso : poezio, eloquenteso, e c.) di yarcento, di naciono\lingui{DEFIRS}
\artiklo{litero}%
\vorto{litero} Signo alfabetala per qua onu reprezentas singla son-uno di linguo\lingui{DEFIRS}
\artiklo{litiero}%
\vorto{litiero} Lito ek palio, folii, filiko sika, por la animali, en la stabli\lingui{EFI}
\artiklo{litijo}%
\vorto{litijo}\fako{yuro-cienco} Kontesto-punto qua divenas la materio di proceso\lingui{DEFIS}
\artiklo{litio}%
\vorto{litio}\fako{kemio} Korpo simpla, la maxim lejera ek omna metali konocata\simbolo{Li}\lingui{DEFIRS}
\artiklo{lito}%
\vorto{lito}\senco{1}{Moblo destinita a recevor su-kushonti}\senco{2}{Parto kava, plu o min profunda, en qua kontenesas e fluas la aquo di fluvio, di rivero, di rivereto}\lingui{FI}
\artiklo{litografar}%
\vorto{litografar}\fako{trans.} Reproduktar per preso la skriburi, la desegnuri, per trasar li (per krayono od inko grasoza) sur petro kalkoza di qua la grano esas tre mikra\lingui{DEFIS}
\artiklo{litologio}%
\vorto{litologio} Ta parto di la historio naturala, qua koncernas la petri\lingui{DEFIRS}
\artiklo{litoro}%
\vorto{litoro} La regiono di qua ula parto di lua konturo esas an maro\lingui{DEFIS}
\artiklo{litoto}%
\vorto{litoto}\fako{retor.} Figuro per qua onu minfortigas la expreso di sua penso, por komprenigar la maxim multo per la minim multo\lingui{DEFIS}
\artiklo{litotomiar}%
\vorto{litotomiar} \textit{(trans.) (kirurg.)} Dividar ed extraktar (per operaco) kalkolo ek la veziko\lingui{DEFIRS}
\artiklo{litotriciar}%
\vorto{litotriciar}\fako{trans.) (kirurg.} Triturar (per operaco) la kalkolo qua esas en la veziko, enduktante instrumento specala en-alonge la uretro\lingui{DEFIRS}
\artiklo{litotripsiar}%
\vorto{litotripsiar}\fako{trans.) (kirurg.} Dividar per froto ed extraktar la kalkolo ek la veziko\lingui{DEFIRS}
\artiklo{litro}%
\vorto{litro}\senco{1}{\textit{(en la mezuro-sistemo metrala)} unajo di la mezuri pri kapaceso}\senco{2}{Botelo kapaca ye un litro}\lingui{DEFIRS}
\artiklo{liturgio}%
\vorto{liturgio}\fako{religio} Formo di la kulto; ordino di la ceremonii konsakrita\lingui{DEFIRS}
\artiklo{liuto}%
\vorto{liuto}\fako{muziko} Muzik-instrumento kun plura rangi de kordi muntita sur rezono-kesto sube-rondigita, e quan onu uzas per pinchar la kordi\lingui{DEFIRS}
\artiklo{livar}%
\vorto{livar}\fako{trans.} Departar de apud (lu) definitive\lingui{E}
\artiklo{livida}%
\vorto{livida}\fako{aludante la vizajo} Qua esas plombo-nigra, bluatra, malado-pala\lingui{DEFIS}
\artiklo{livrar}%
\vorto{livrar}\fako{trans., ad}\senco{1}{Igar ye la dispono di ulu}\senco{2}{Submisar a la ago da ulu}\lingui{EFS}
\artiklo{livreo}%
\vorto{livreo}\senco{1}{Vesto quan la reji, la siniori furnisas a la viri di sua eskorto, e di qua la koloro, la galoni, la butoni, e c. memorigas lia blazoni}\senco{2}{Distingo-kostumo (distingebla pro koloro, galoni, butoni, e c.) quan la hemestro igas *werar da sua servistuli}\lingui{DEFIRS}
\artiklo{lizimako}%
\vorto{lizimako}\fako{bot.} Planto qua formacas genero ek la familio « primulacei »\latina{lysimachia}
\artiklo{lizo}%
\vorto{lizo}\senco{1}{Sedimento quan la vino depozas sur la fondo di barelo}\senco{2}{Elemento di reziduo sociala}\lingui{EF}
\artiklo{lo}%
\vorto{lo} Pronomo qua : 1. reprezentas fakto, o propoziciono; 2. indikas kozo nedeterminita precize (= to quo, co)
\artiklo{lobelio}%
\vorto{lobelio}\fako{bot.} Genero de « lobeliacei », di qua la suko esas laktatra, akra, venena, e quan onu kultivas precipue kom planto orniva\latina{lobelia}\lingui{DEFS}
\artiklo{lobo}%
\vorto{lobo}\fako{anat.}\senco{1}{Parto di ula organi di qua la formo esas ronda}\senco{2}{Saliajo ronda per qua finas, ye sua infra parto, la funelo di la orelo}\lingui{EFIS}
\artiklo{locho}%
\vorto{locho}\fako{zool.} Fisho di aquo sen-sala, di qua la korpo esas tre oblonga\latina{cobitis barbatula}\lingui{EFIS}
\artiklo{locionar}%
\vorto{locionar}\fako{trans.} Fluigar liquido sur parto di la homo-korpo (o sur la korpo tota), sive per varsar olu, sive per ekirigar olu per preso ek sponjo, linjo humidigita, e c.\lingui{DEFIS}
\artiklo{logaritmo}%
\vorto{logaritmo}\fako{matem.} Exponento integra o fracionala di la potenco a qua oportas elevar nombro konstanta por *reproduktar ta o ca nombro\lingui{DEFIRS}
\artiklo{logiko}%
\vorto{logiko}\senco{1}{La cienco pri lo exakta}\senco{2}{La cienco pri la *leyi (legi) qui koncernas rezonado}\senco{3}{La maniero rezonar rigoroze}\senco{4}{Interligeso racionala di la kozi. (ex. : che \textit{Ido} : la arto kompozar la frazi o la vorti tale ke la adiciono di la senco di singla ek la elementi donas (sen arbitrialeso, sen deviaco de la normo) la senco preciza di la ensemblo di la frazo o di la vorto)}\lingui{DEFIRS}
\artiklo{logo}%
\vorto{logo}\fako{maro-navig.} Planketo balastizita, imersita ye la extremajo di kordo, ye la dopa parto di navo, por mezurar la rapideso-grado di la iro da la navo\lingui{DEFIR}
\artiklo{logogrifo}%
\vorto{logogrifo} Sorto di enigmato qua konsistas ye un vorto di qua la literi, kombinita diverse, formacas altra vorti qui anke esas divinenda\lingui{DEFIRS}
\artiklo{logomakio}%
\vorto{logomakio} Disputo pri vorti\lingui{DEFIRS}
\artiklo{lojar}%
\vorto{lojar}\fako{netrans., en, che} Esar instalita sub tekto, en domo od en parto di domo, sive kontinue, sive dum-kelka-tempe\lingui{DEFIRS}
\artiklo{lojio}%
\vorto{lojio}\senco{1}{Chambreto aden qua onu enklozas la dementi}\senco{2}{Mikra domo di domo-gardisto}\senco{3}{Chambreto aden qua onu enklozas la studenti pri estetiko (piktado, skultado, muzikado, e c.) admisita por konkursar pri la Roma-Premio, dum ke li traktas la temo propozita por la konkurso}\senco{4}{Singla ek la chambreti en qui la teatro-artisti chanjas sua pleo-vesti e fardizas su}\senco{5}{Singla ek la mikra dividuri di spektaklo-chambrego, qua kontenas tri, quar, kin o sis plasi}\senco{6}{*Kluzajo (klozajo) ube inter-asemblas grupo de framasoni, prezide da la veneracindo}\lingui{DEFR}
\artiklo{lokacar}%
\vorto{lokacar}\fako{trans.} Obtenar (de lua proprietanto) la yuro uzar, dum periodo konvencionata, ulo po la pago di la preco konvencionata\lingui{F}
\artiklo{lokativo}%
\vorto{lokativo}\fako{gram.} Kazo gramatikala qua, en ula lingui, expresas la loko, la destino\lingui{DEFIS}
\artiklo{loklo}%
\vorto{loklo} Volvuro (en hari, rubando) forme di ringo\lingui{DER}
\artiklo{loko}%
\vorto{loko}\senco{1}{Parto determinita ek spaco}\senco{2}{Plaso, parto di spaco, atribuita a kozo, a persono determinita}\senco{3}{Parto di la texto di libro en qua esas texto determinita}\senco{4}{Parto limitizita en lando, en regiono}\lingui{DEFIRS}
\artiklo{lokomobilo}%
\vorto{lokomobilo} Vapor-mashino transportebla\lingui{DEFIRS}
\artiklo{lokomocar}%
\vorto{lokomocar}\fako{netrans.} Agar la movo o la serio de movi per qui onu su transportas de loko ad altra loko\lingui{DEFIS}
\artiklo{lokomotivo}%
\vorto{lokomotivo}\senco{1}{\textit{(olim)} Vaporo-mashino portata per roti, ed uzata por tirar treni sur relvoyi}\senco{2}{\textit{(nun)} Vehilo analoga, ma movata da elektro o da Diesel-motoro}\lingui{DEFIRS}
\artiklo{lokusto}%
\vorto{lokusto}\fako{zool.} Insekto aloza, ek la familio « ortopteri », saltera\latina{locusta}\lingui{EFIS}
\artiklo{lolio}%
\vorto{lolio}\fako{bot.} Herbo di qua la grani esas nigra, qua kreskas sur frumento\latina{lolium temulentum}\lingui{EIL}
\artiklo{lombardo}%
\vorto{lombardo} Establisuro di presto po gajo (institucita kun skopo karitatala)\lingui{DEF}
\artiklo{lombriko}%
\vorto{lombriko}\fako{zool.} Anelideo, di qua la nomo esas anke tervermo\latina{lumbricus terrestris}\lingui{EFIS}
\artiklo{loncho}%
\vorto{loncho} Porciono plu o min dina di korpo, tranchita per de-seko tre egala, tale ke la surfaco esas tre plana
\artiklo{longa}%
\vorto{longa}\senco{1}{\textit{(spaco)} Qua esas extensita, de un de sua extremajo a la altra, plu multe pri ta dimensiono kam pri la cetera. – Qua esas (tante e tante) extensita pri ta dimensiono. (ex. Ta domo esas longa ye X metri)}\senco{2}{\textit{(tempo)} Qua esas tre extensita en tempo}\lingui{DEFI}
\artiklo{longitudo}%
\vorto{longitudo}\senco{1}{\textit{(geogr.)} Koordinato geografiala qua uzesas por determinar la situeso di loko sur la surfaco di la Ter-globo : disto de ta loko a la meridiano determinita}\senco{2}{\textit{(astron.)} Koordinato astronomiala uzata por determinar la situeso di astro sur la sfero cielala; disto de ta astro a la punto equinoxala di printempo}\lingui{EFIRS}
\artiklo{lor}%
\vorto{lor} En la epoko kande… En la epoko di… \textit{(lore… lore : ye ula instanto… ye altra…)}
\artiklo{lorno}%
\vorto{lorno} Aparato qua konsistas ye lensi quin onu interpozas inter la okulo regardanta e la objekto regardata, por vidar plu precize, plu dicerneble\lingui{FR}
\artiklo{lotao}%
\vorto{lotao}\fako{zool.} Fisho qua vivas en aquo sen-sala, ek la genero « gado »\latina{gadus lota}\lingui{EFSL}
\artiklo{loto}%
\vorto{loto} To quo atribuesas ad ulu lor partigo o disdono. To quo atribuesas da hazardo ad ulu lor lotrio\lingui{EF}
\artiklo{lotriar}%
\vorto{lotriar}\fako{trans.} Agar ta hazardo-ludo en qua onu distributas ula quanto de bilieti pagenda, qui havas numeri, ek qui uli, indikota da hazardo (fato), yurizas pri recevor pekunio-quanto od objekto\lingui{DEFIRS}
\artiklo{lotuso}%
\vorto{lotuso}\fako{bot.}\senco{1}{Speco di jujubiero. Che la antiqui : arboro di qua la frukto egardesis kom iganta la stranjeri qui konsumis de lu obliviar sua patrio}\senco{2}{\textit{(bot.)} Speco di nenufaro blua qua kreskas en Egiptia}\latina{nelumbium speciosum}\lingui{DEFIRS}
\artiklo{loxodromo}%
\vorto{loxodromo}\senco{1}{\textit{(geom.)} Kurvo trasita sur la surfaco di sfero e qua intersekas sam-angule omna meridiani}\senco{2}{\textit{(maro-navig.)} Kurvo quan trasas navo kande lu iras kontinue ad un punto di la busolo, en un rumbo di vento}\lingui{DEFIS}
\artiklo{loyala}%
\vorto{loyala} Skrupuloze fidela ye sua promisi, ye sua su-obligi\lingui{DEFIRS}
\artiklo{lu}%
\vorto{lu} Pronomo personala, triesma persono, singularo, epicena
\artiklo{lubrifikar}%
\vorto{lubrifikar}\fako{trans.}\senco{1}{Igar glitigiva (tisuo) por igar plu facila lua funciono}\senco{2}{Untizar ye graso (mashino, aparato)}\lingui{EFIS}
\artiklo{lucida}%
\vorto{lucida} Qua komprenas klare la kozi\lingui{EFIS}
\artiklo{lucio}%
\vorto{lucio}\fako{zool.} Fisho qua vivas en aquo sen-sala, tre devorera, ek la familio « ezoci », di qua la korpo esas funelatra, kun muzelo longa, boko tre fendita e garnisita ye granda quanto de denti, e di qua la karno tre prizesas\latina{esox lucius}\lingui{ISL}
\artiklo{luciolo}%
\vorto{luciolo}\fako{zool.} Lampiro aloza e fosforecanta\latina{luciola italica}\lingui{FIS}
\artiklo{ludar}%
\vorto{ludar}\fako{netrans.) (per}\senco{1}{Agar ulo por amuzar su}\senco{2}{Amuzar su per manuagar instrumento}\senco{3}{Amuzar su, kun altra persono o personi, per agi en qua uli ganas e la ceteri perdas}
\artiklo{ludiono}%
\vorto{ludiono}\fako{fiziko} Figureto, statueto, imajeto, ek vitro, garnisita ye la parto supera ye sfereto vakua qua flotacas segun la principo da Arkimedes, ed acensas o decensas (segun ke onu lo deziras) en aquo\lingui{FIS}
\artiklo{lugar}%
\vorto{lugar}\fako{trans., ad) (aludante proprietero} Cedar la uzo-yuro di kozo – cedo dum-frista, po preco konvencionata
\artiklo{lugro}%
\vorto{lugro}\fako{maro-navig.} Milito-naveto lejera, uzita precipue da la kontrabandisti e da la pirati\lingui{DEFIS}
\artiklo{lukano}%
\vorto{lukano}\fako{zool.} Insekto aloza, granda, kun mandibuli di qui la formo esas olta di korni dentoza, kelke simila a la kornaco di cervo\latina{lucanus}\lingui{FL}
\artiklo{luko}%
\vorto{luko}\fako{navo} Aperturo quadriangula en la ponto, por igar komunikar du etaji\lingui{DER}
\artiklo{luktar}%
\vorto{luktar}\fako{netrans., kontre, kun, por}\senco{1}{\textit{(aludante du individui)} Esforcar por renversar la altro, per embraco}\senco{2}{\textit{(aludante du individui, du populi, du partisi)} Esforcar por vinkar la altra, la adverso, kontestante pri sua superioreso}\senco{3}{\textit{(aludante la volo da la homo)} Esforcar por vinkar sua pasioni, por vinkar la obstakli extera}\senco{4}{\textit{(aludante forco)} Efikar ad altra forco, kontre-sinse}\lingui{FIS}
\artiklo{lular}%
\vorto{lular}\fako{trans.} Duktar (infanteto) ad dormo per kantar apud lu\lingui{DE}
\artiklo{lumar}%
\vorto{lumar}\fako{netrans.) (aludante ula korpi} Emisar ula radii qui igas videbla la objekti\lingui{EFIS}
\artiklo{lumbago}%
\vorto{lumbago}\fako{patol.} Doloro en la regiono lumbala, rumatismala, nevralgiala, o produktita da esforco tro granda\lingui{EFIS}
\artiklo{lumbo}%
\vorto{lumbo}\fako{anat.} Dopa parto di la torso, de la dorso til la hanchi, dextre e sinistre di la kolono vertebrala (spino)\lingui{EFIS}
\artiklo{lumpo}%
\vorto{lumpo}\fako{mineral.} Fragmento di roko ronda quan onu renkontras en ula strati mineralala\lingui{E}
\artiklo{lunatika}%
\vorto{lunatika} Submisita a la efiko (asertita) da Luno; dicesas precipue pri persono qua ne plus havas sua raciono o rezono-fakultato pro, asertite, subisir la efiko da Luno\lingui{DEFIRS}
\artiklo{lundio}%
\vorto{lundio} La dio duesma di la semano (olta qua sequas sundio)\lingui{FIS}
\artiklo{luno}%
\vorto{luno}\fako{astron.} Planeto qua jiras cirkum la Terglobo (di qua lu nomizesis la « satelito ») e qua lumizas lu dum nokto kande lu esas situita tale ke lu reflektas a la Terglobo la lumo-radii da Suno\lingui{DEFIRS}
\artiklo{lupanaro}%
\vorto{lupanaro} Domo en qua la mulieri su prostitucas\lingui{EFL}
\artiklo{lupino}%
\vorto{lupino}\fako{bot.} Planto leguminosa, di qua la shelo kontenas grani di qui la farino esas manjebla\latina{lupinus}\lingui{DEFIRS}
\artiklo{lupio}%
\vorto{lupio}\fako{patol.} Surkreskajo sub la pelo, tumoro sen-dolora, maxim-multa-kaze kisteskinta\lingui{FIS}
\artiklo{lupo}%
\vorto{lupo}\fako{optiko} Lenso bikonvexa, uzata por igar aspektar plu granda kam reale la objekto quan onu regardas tra ta lenso\lingui{DF}
\artiklo{lupulo}%
\vorto{lupulo}\fako{bot.} Planto ek la familio « urticei », di qua la floro esas uzata por facar biro\latina{humulus lupulus}\lingui{ISL}
\artiklo{lupuso}%
\vorto{lupuso}\fako{patol.} Pelo-morbo di naturo tuberklosala\lingui{DEFIS}
\artiklo{lurar}%
\vorto{lurar}\fako{trans., per}\senco{1}{Atraktar per ulo qua apetitigas}\senco{2}{\textit{(anke metaf.)}}\lingui{E}
\artiklo{lustracar}%
\vorto{lustracar}\fako{trans.) (en la epoki antiqua, en Roma} Purigar lor ceremonio ta-skopa\lingui{DEFIRS}
\artiklo{lustrino}%
\vorto{lustrino} Sorto di perkalino, tre apretita e briligita\lingui{DEFIS}
\artiklo{lustro}%
\vorto{lustro}\senco{1}{Lumizilo suspendita, plura-brancha, uzata por lumizar e dekorar la granda saloni, la kirki, la teatri}\senco{2}{Periodo de kin yari}\lingui{DEFIRS}
\artiklo{luto}%
\vorto{luto} Induturo uzata por klozar hermetike vazi, por tegar retorti, vitro-tubi o porcelano-tubi quin onu volas protektar kontre la tro-varmeso di fairo\lingui{EFIS}
\artiklo{lutro}%
\vorto{lutro}\fako{zool.} Quadripedo karnavida, ek la familio « martri », di qua la pilaro esas densa, qua vivas apud aquo, e qua su nutras ek fishi\latina{lutra}\lingui{FIS}
\artiklo{luxacar}%
\vorto{luxacar}\fako{trans.) (patol.} Igar (osto) ekirar de lua kavajo normala\lingui{DEFIS}
\artiklo{luxo}%
\vorto{luxo}\senco{1}{Richeso, brilo, ostentata en la kozi di la vivo}\senco{2}{\textit{(metaf.)} Kozo superflua}\lingui{DEFIS}
\artiklo{luxurio}%
\vorto{luxurio} Peko kontre chasteso\lingui{FIS}
\artiklo{luzerno}%
\vorto{luzerno}\fako{bot.} Planto leguminosa quan onu kultivas en prati artificala, kom nutrivo por la bestii\latina{medicago}\lingui{DEFIR}
\end{multicols}
\sekciono{M}
\begin{multicols}{2}
\artiklo{ma}%
\vorto{ma} Konjunciono qua indikas la interopozeso di la du propozicioni quin lu ligas, unionas : (1) per restriktar to quo jus dicesis; (2) per expresar ideo qua kontrastas kun olta qua jus expresesis; (3) per objecionar a to quo jus dicesis\lingui{FIS}
\artiklo{macedonio}%
\vorto{macedonio}\fako{koquarto} Disho en qua legumi diversa esas preparita kune. Entremeso sukroza en qua frukti diversa esas asemblita aden quaza jeleo\lingui{EFIS}
\artiklo{macerar}%
\vorto{macerar}\fako{trans. e netrans.} Trempar (substanco) aden liquido kolda, e lasar (ta substanco) en ta liquido, por deprenar la solveblaji\lingui{DEFIS}
\artiklo{macfarlane}%
\vorto{« macfarlane »} Mantelo sen maniki, kun aperturi tra qui pasas la brakii, sur qui kovre-pendas pelerino qua finas ye la tayo di la persono
\artiklo{mackintosh}%
\vorto{« mackintosh »} Sorto di muliero-mantelo di qua la stofo esas nepermeebla
\artiklo{madapolam}%
\vorto{« madapolam »} Texuro de kotono blanka, di qua la grano esas tre salianta e tre apretita – meze inter kaliko e perkalo, tre uzata por facar la avantali koqueyala
\artiklo{madeleno}%
\vorto{madeleno} Mikra kuko di qua la pasto esas kompakta\lingui{DEFIS}
\artiklo{madono}%
\vorto{madono} Statuo, pikturo, qua reprezentas (che la kristani) la virgino Maria\lingui{DEFIRS}
\artiklo{madras}%
\vorto{« madras »} Kapovesto ek fulardo de stofo di qua la fili warpa esas ek silko, e la wefto ek kotono – e qua, origine, fabrikesis en la urbo Madras (India)
\artiklo{madreporo}%
\vorto{madreporo} Genero de polipi agregita, asemblajo de celuli kalka qui interkomunikas, e kreskas til formacar strati, rifi, insuleti, e c.\latina{madrepora}\lingui{DEFIS}
\artiklo{madrigalo}%
\vorto{madrigalo} Kurta versajo, maxim-multa-kaze turnita a damo, e qua expresas penso tenera o galanta\lingui{DEFIRS}
\artiklo{maestro}%
\vorto{maestro} Eminento pri arto o cienco\lingui{DEFIRS}
\artiklo{magazino}%
\vorto{magazino} Chambro, loko ube onu depozas vari diversa, por konservar li ante lia uzeso o vendeso. (Ne konfundar a « butiko », « butikego »)\lingui{DEFIR}
\artiklo{magiar}%
\vorto{magiar}\fako{netrans.} Efikar (aserte) a la homi, a la elementi, per povo supernatura – en ula kazi : mediace diablo.— \textit{(metaf.)} Efikar extraordinare a la homi (per eloquenteso, per poezio, e c.)\lingui{DEFIS}
\artiklo{magnato}%
\vorto{magnato} Stato-grando, en Polonia ed en Hungaria, olim\lingui{DEFIRS}
\artiklo{magnetizar}%
\vorto{magnetizar}\fako{trans.} Submisar (ulu) a la efiko di magnetismo (animalala)\lingui{DEFIRS}
\artiklo{magneto}%
\vorto{magneto} Erco de fero oxidoza, di qua un ek la karakteri esas atraktar fero o stalo. – Ferbloko o stalobloko quan efikis magneto o koranto (elektro-fluo) o la magnetefekto di la Terglobo. – \textit{(metaf.)} To quo atraktas forte\lingui{DEFIRS}
\artiklo{magnezio}%
\vorto{magnezio}\fako{kemio} Korpo simpla, metala, qua esas kombustebla, en aero, kun lumo blanka dazliva, lasante reziduo ek polvo blanka qua esas magnez-oxido\simbolo{Mg}\lingui{DEFIRS}
\artiklo{magnolio}%
\vorto{magnolio}\fako{bot.} Arboro exotika, kun granda folii briloza, bele-verda, e flori odorifanta, pure-blanka\latina{magnolia}\lingui{DEFIS}
\artiklo{mago}%
\vorto{mago}\senco{1}{Sacerdoto di la religio di la Persi antiqua, adoranti di fairo}\senco{2}{\textit{(biblo)} Un ek le tri personegi – quin la tradiciono transformis a reji – qui venis adorar Iesu en lua kripo, en urbo Betleem}\lingui{DEFIS}
\artiklo{magra}%
\vorto{magra}\senco{1}{\textit{(aludante homo od animalo vivanta)} Di qua la korpo kontenas poka graso}\senco{2}{\textit{(aludante karno koquala)} En qua esas poka o nula graso}\senco{3}{\textit{(aludante nutrivi)} Qua kontenas nek karno, nek graso}\lingui{DEFIS}
\artiklo{mahagono}%
\vorto{mahagono}\fako{bot.} Arboro qua furnisas bela ligno uzata por la faco di mobli\latina{Swetenia mahogani}\lingui{DEI}
\artiklo{maifloro}%
\vorto{maifloro}\fako{bot.} Planto kun mikra flori blanka odorifanta\latina{convallaria majalis}\lingui{D}
\artiklo{maizo}%
\vorto{maizo}\fako{bot.} Planto cereala, deveninta de Amerika, di qua la grani esas manjebla\latina{zea mais}\lingui{DEFRS}
\artiklo{majesto}%
\vorto{majesto} Honor-titulo donata a la personi suverena, sive konversante kun li, sive parolante pri li. (« Vua majesto », « Lua majesto »)\lingui{DEFIS}
\artiklo{majora}%
\vorto{majora}\senco{1}{\textit{(yuro)} Qua evas lo determinita da la lego por ke lu povez posedar sua yuri}\senco{2}{\textit{(muziko)} (pri terco) Intervalo ek tri noti, kontenanta du toni. \textit{(pri modo)}. En qua la sexto super la noto tonala esas majora}\lingui{DEFIS}
\artiklo{majorano}%
\vorto{majorano}\fako{bot.} Planto aromatoza, de la familio « labiei »\latina{majorana}\lingui{DEFIS}
\artiklo{majorato}%
\vorto{majorato}\fako{yuro-cienco} Proprietajo imobla, ligita a nobeleso-titulo ne-alienebla, qua heredesas kun ta titulo, (en Francia, lor la rejiL)\lingui{EFIS}
\artiklo{majordomo}%
\vorto{majordomo} Jeranto di la servisti en familio princala di Italia e di Hispania. Chefo di la servistaro\lingui{DEFIRS}
\artiklo{majoritato}%
\vorto{majoritato} La maxim granda nombro di la voti, en asemblitaro deliberanta, pri un ek la punti submisita a voto.\lingui{DEFIS}\subvorto{}{la majoritato absoluta}{La duimo di la votanti + 1}{}{}
\artiklo{makadamo}%
\vorto{makadamo}\fako{tekn.} Quaza-petriguro di voyo, di choseo, per kovrar ica ye granito, silexo, triturita e pistita, e submisita a forta preso tale ke lu formacas surfaco kompakta\lingui{DEFIRS}
\artiklo{makako}%
\vorto{makako}\fako{zool.} Simio de Afrika, di qua la korpo esas grosa e kurta, e qua havas muzelo kurta e dika\latina{macacus}\lingui{DEFIS}
\artiklo{makaronika}%
\vorto{makaronika} Qua esas en linguo Latina burleska (en linguo moderna, vestizachita ye dezinenci Latina)\lingui{DEFIRS}
\artiklo{makaronio}%
\vorto{makaronio}\fako{koquarto} Nutro-pasto ek farino muldita aden longa cilindri kava\lingui{DEFIRS}
\artiklo{makarono}%
\vorto{makarono}\fako{koquarto} Pasto ek mandeli, sukro ed albumeno, fasonita aden mikra panbloki ronda\lingui{DEFR}
\artiklo{makiavelala}%
\vorto{makiavelala} Konforma a la doktrino asertita da Machiavelli, olqua, dum multa tempo, konsideresis kom perfekte kontre-etika (perfida, ruzoza, malicoza, sen-shama)\lingui{DEFIRS}
\artiklo{makio}%
\vorto{makio}\fako{zool.} Simio kun longa kaudo e muzelo longa\latina{Lemur}\lingui{DEFIRS}
\artiklo{maklo}%
\vorto{maklo}\fako{kristalografio} Ula grupi (qui havas la formo di genuo, di tekto-pluvkanalo, di kordio) de kristali di qui la naturo, la formo e la molekulo-strukturo esas perfekte intersimila\lingui{EFIS}
\artiklo{makrelo}%
\vorto{makrelo}\fako{zool.} Maro-fisho, kun makuleti diversa-kolora, qua, singla-yare, arivas trupe de la nordo-regioni\latina{scomber}\lingui{DEFR}
\artiklo{makronukleo}%
\vorto{makronukleo}\fako{biol.} La maxim grosa nukleo di la infuzorii
\artiklo{makrosporangio}%
\vorto{makrosporangio}\fako{biol.} Sporangio qua produktas la makrosporo
\artiklo{makrosporo}%
\vorto{makrosporo}\fako{biol.) (che ula algi} Gameto heterogama, konsiderata kom femina, opoze a la mikrosporo (o gameto maskula)\lingui{DEF}
\artiklo{makulaturo}%
\vorto{makulaturo}\fako{imprim-arto} Papero-folio qua uzesis por recevar la quanto ecesanta di la imprim-inko – o : imprimo-folio en qua la tipi esas desordinita pro mala imprimo\lingui{DEFIRS}
\artiklo{makulo}%
\vorto{makulo}\senco{1}{Marko qua sordidigis ulo}\senco{2}{Marko naturala sur la pelo di la animali}\senco{3}{Altereso, en un loko, di la koloro di la ensemblo. – \textit{(anke metaf.)}}\lingui{DEFIS}
\artiklo{mala}%
\vorto{mala}\senco{1}{Ye maniero regretinda, penigiva}\senco{2}{Ye maniero neperfekta, defektoza}\senco{3}{Ye maniero kontrea a la lego etikala, morala}\senco{4}{Kontrea a lo devata}\lingui{EFIS}
\artiklo{malada}%
\vorto{malada}\senco{1}{Qua subisis kelka altero di sua organi, kelka parturbo di sua funcioni organismala}\senco{2}{\textit{(metaf.)} Qua subisis kelka perturbo di sua fakultati spiritala, etikala}\lingui{EFI}
\artiklo{malakito}%
\vorto{malakito}\fako{mineral.} Lapido verda, ek kupro karbonatoza, apta divenar bele polisita, uzata da la juvelisti\lingui{DEFIRS}
\artiklo{malario}%
\vorto{malario}\fako{patol.} Marsho-febro
\artiklo{maldicensar}%
\vorto{maldicensar}\fako{netrans.) (pri} Dicar, pri ulu, lo mala quan onu savas pri lu\lingui{FIS}
\artiklo{malear}%
\vorto{malear}\fako{trans.} Extensar (ordinare : metalo) per martelagar (sen varmigo preparala) aden lami o folii plu o min dina\lingui{EFIS}
\artiklo{maledikar}%
\vorto{maledikar}\fako{trans.}\senco{1}{Advokar mala chanco kontre ulu}\senco{2}{\textit{(aludante patrulo)} Advokar la iraco deala (kontre filiulo kruelega, senkordia.)}\senco{3}{\textit{(aludante sacerdoto)} Advokar, per vorti, per ritui religiala, la blamo deala kontre ulu, kontre ulo}\senco{4}{\textit{(aludante Deo)} Destinar a la blamo eterna}\lingui{DEFIS}
\artiklo{maleolo}%
\vorto{maleolo}\fako{anat.} Saliajo osta, ye la infra extremajo di singla ek la gambi – interna-latere, forme di salio di la tibio; extera-latere, forme di salio di la extremajo di la peroneo\lingui{EFIS}
\artiklo{malgre}%
\vorto{malgre} Sen impedesar da la nevolo quan opozas ulu, da la obstaklo quan opozas ulo\lingui{FI}
\artiklo{malico}%
\vorto{malico}\senco{1}{Plezuro de agar ulo mala}\senco{2}{Su-amuzo de agar o dicar malignalaji negrava}\lingui{EF}
\artiklo{maligna}%
\vorto{maligna} Qua tendencas, inklinesas a, delektesar de la doloro o chagreno da kunhomo\lingui{EFIS}
\artiklo{malio}%
\vorto{malio} Ludo ye qua onu avancigas bulo ek busligno, per quaza martelo ek ligno harda, ferizita, kun mancho longa e flexebla\lingui{EFIS}
\artiklo{malto}%
\vorto{malto} Hordeo jermifinta e sika, quan onu uzas por fabrikar biro, pos deprenir la jermi\lingui{DEFIS}
\artiklo{malvarozo}%
\vorto{malvarozo}\fako{bot.} Planto ek la familio « malvacei », qua florigas abunde en la fino-parto di somero\latina{Althaea rosea}\lingui{IS}
\artiklo{malvavisko}%
\vorto{malvavisko}\fako{bot.} Planto mucilajatra ek la familio « malvacei »\latina{althaea officinalis}\lingui{IS}
\artiklo{malvo}%
\vorto{malvo}\fako{bot.} Planto apta moligar la tisui inflamita. Herbo kun folii alternanta, di qui la flori (sola o asemblita), dispozita spike o grape, donas frukto enkapsula\latina{malva}\lingui{DEFIRS}
\artiklo{malyoto}%
\vorto{malyoto} Sorto di kalsono strete-fitanta, por maro-balno, por baleto-danso\lingui{F}
\artiklo{mamifero}%
\vorto{mamifero}\fako{zool.} Animalo vivipara, kun mami qui sekrecas, che la femini, la lakto ye qua nutresas la jus-naskinti\lingui{EFIS}
\artiklo{mamilo}%
\vorto{mamilo} La pinto di la mamo
\artiklo{mamo}%
\vorto{mamo}\fako{anat.}\senco{1}{Saliajo mi-globo, dextre (o sinistre) di la pektoro di muliero, organo glandatra qua sekrecas la lakto ye qua nutresas la jus-naskinto – anke aludante la femino che la mamiferi}\senco{2}{Parto simila, ma rudimenta, e sen sekreco, che la maskulo}\lingui{EFI}
\artiklo{mamuto}%
\vorto{mamuto}\fako{paleont.} Animalo ek la genero « elefanto », ne plus existanta, quan onu renkontras kom fosilo, precipue en Siberia\lingui{DEFIRS}
\artiklo{manao}%
\vorto{manao}\fako{segun biblo} Nutrivo falinta de cielo, por nutrar la Hebrei en la dezerto. — \textit{(anke metaf.)}\lingui{DEFIRS}
\artiklo{mancho}%
\vorto{mancho} Parto muntita an utensilo od implemento, por tenar la utensilo od implemento en qua manuo kande onu uzas lu\lingui{FIS}
\artiklo{mandareno}%
\vorto{mandareno}\fako{olim} En Chinia, granda stat-employato, civila od armeana, selektita maxim-multa-kaze ek la klaso « literaturisti »\lingui{DEFIRS}
\artiklo{mandarino}%
\vorto{mandarino}\fako{bot.} Mikra oranjo aromatoza\latina{citrus nobilis}\lingui{DEFIRS}
\artiklo{mandato}%
\vorto{mandato}\senco{1}{Dokumento emisita en posto-kontoro, e po qua irga posto-kontoro di la lando pagos la pekunio-quanto indikita sur olu}\senco{2}{Dokumento donita ad ulu, da administro-fako statala, e po qua la Trezoro statala pagos la pekunio-quanto indikita sur olu}\senco{3}{\textit{(banko)} Dokumento per qua bankiero imperas a sua debanto ke ica pagez a ta o ca persono la pekunio-quanto indikita sur olu}\lingui{DFIRS}
\artiklo{mandelo}%
\vorto{mandelo}\fako{bot.} Frukto di arboro ek la familio « rozacei »\latina{amygdalus communis}\lingui{DFIRS}
\artiklo{mandibulo}%
\vorto{mandibulo}\fako{anat.} La parto infra ek la du parti osta di la boko, qui suportas la denti, che la animali vertebroza\lingui{DEFIS}
\artiklo{mandolino}%
\vorto{mandolino}\fako{muziko} Muzik-instrumento ek framo, konvexa sube, qua havas normale mancho kun quar kordi metala duopla, dispozita ed inter-akordita quale olti di violino, e quin onu vibrigas per peco de ucel-plumo o di skalio\lingui{DEFIRS}
\artiklo{mandoro}%
\vorto{mandoro}\fako{muziko} Muzik-instrumento (ne plus uzata) analoga a mandolino\lingui{DEFI}
\artiklo{mandragoro}%
\vorto{mandragoro}\fako{bot.} Planto ek la genero « solanei », kun radiko pulpoza, narkotigiva e purgiva, ed a qua onu atribuis efiki magiala\latina{mandragora officinarum}\lingui{EFIS}
\artiklo{mandrelo}%
\vorto{mandrelo}\fako{tekn.} Peco de fero fite a qua onu forjas peci quin onu volas igar kava\lingui{EFIS}
\artiklo{mandrino}%
\vorto{mandrino}\fako{teknol.} Peco quan onu visas sur tornilo horizontala e per qua la axo rotacanta rotacigas la objekto quan onu volas fasonar\lingui{FIS}
\artiklo{manejo}%
\vorto{manejo} Loko ube onu dresas la kavali, ube onu docas e praktikas la kavalk-arto\lingui{DEFIRS}
\artiklo{manekino}%
\vorto{manekino}\senco{1}{(1) Figuro ligna o vaxa, qua reprezentas viro o muliero, olquan uzas la piktisti, la statuifisti, por probar posturi, dispozar drapiraji, dum la absenteso di la homa modelo. (2) Figuro quan la taliori vestizas por uzesar kom modeli pri la faco di la kostumi}\senco{2}{Figuro realeso-granda, sur qua onu exercas la studenti di medicino, di kirurgio, pri la bandajizo o pri la movi quin postulas ula operaci}\lingui{DEFRS}
\artiklo{maneto}%
\vorto{maneto}\fako{teknol.} Mikra mancho quan onu sizas por movigar mekanismo\lingui{FI}
\artiklo{mangano}%
\vorto{mangano}\fako{kemio} Korpo simpla, metala, blanka grizatre, harda e facile ruptebla, qua oxidesas aero\simbolo{Mn}\lingui{DEFIRS}
\artiklo{manglo}%
\vorto{manglo}\fako{bot.} Frukto di arboro di la regiono intertropika, qua kreskas en maro, en tereni slamoza. – Mangliero\latina{Rhizophora mangle}\lingui{DEFIS}
\artiklo{mango}%
\vorto{mango}\fako{bot.} Frukto di granda arboro di India e di la sudo-regioni di Amerika, e qua kelke similesas persiko\latina{mango}\lingui{DEFIS}
\artiklo{mangusto}%
\vorto{mangusto}\fako{zool.} Rato di Egiptia\latina{Herpestes}\lingui{DEFIS}
\artiklo{maniero}%
\vorto{maniero} Sorto di ago o faco\lingui{DEFIRS}
\artiklo{manifestar}%
\vorto{manifestar}\senco{1}{\textit{(trans.)} Igar quaze palpebla, sentebla}\senco{2}{\textit{(netrans., pri, por)} Demonstrar kolektive, publike, favore di opiniono; expresar publike revendiko – multa-kaze per amaso-cirkular sur la voyi publika, en la stradi, avenui, bulvardi, placi di urbo}\lingui{DEFIRS}
\artiklo{maniko}%
\vorto{maniko} Parto di vesto qua kovras la brakio til la karpo, lasante deskovrata la manuo\lingui{FIS}
\artiklo{manio}%
\vorto{manio}\senco{1}{\textit{(patol.)} Alienaco kun deliro generala ed tendenco a furio}\senco{2}{Pasiono stranja}\lingui{DEFIRS}
\artiklo{manioko}%
\vorto{manioko}\fako{bot.} Arboreto di la sudo-regioni di Amerika, di qua la radiko, liberigite de la veneno quan olu kontenas, donas farino uzata por facar sorto di pano, e fekulo nutriva (tapiokao)\latina{manihot palmata}
\artiklo{maniplo}%
\vorto{maniplo}\senco{1}{\textit{(liturgio katol.)} Bendo stofo quan la sacerdoto katolika *weras sur sua brakio sinistra, kande lu oficias}\senco{2}{\textit{(che la Romani antiqua)} Kompanio ek cent (o duacent) infantriani}\lingui{DEFIRS}
\artiklo{manipular}%
\vorto{manipular}\fako{trans.}\senco{1}{\textit{(teknol.)} Manuagar (mixante li) ula substanci kemiala o farmaciala. – \textit{(anke metaf.)}}\senco{2}{\textit{(teknol.)} Igar instrumento funcionar per manuo-direktar lu}\lingui{DEFIRS}
\artiklo{manipulatoro}%
\vorto{manipulatoro}\fako{teknol.} Parto di la mekanismo telegrafala, quan, lor transmiso, onu funcionigas per la manuo\lingui{EFIRS}
\artiklo{manivelo}%
\vorto{manivelo}\fako{teknol.} Peco muntita ye la extremajo di la ax-arboro di mashino, di rotaxo, ed uzata por rotacigar lu\lingui{FIS}
\artiklo{manjar}%
\vorto{manjar}\fako{trans.} Mastikar e glutar nutrivo solida. \textit{(anke metafore)}\lingui{FIS}
\artiklo{mankar}%
\vorto{mankar}\fako{netrans.) (aludante kozo, persono} Esar absenta en la loko ube lu deziresas, ube lu esas necesa\lingui{FI}
\artiklo{mano}%
\vorto{mano}\fako{mitol. Romana} Fantomo di mortinto, qua, che la antiqui, kultesis\lingui{DEFIRS}
\artiklo{manometro}%
\vorto{manometro} Instrumento per qua onu mezuras la tenseso o la elastikeso di la gasi, di la vapori\lingui{DEFIRS}
\artiklo{manoto}%
\vorto{manoto} Ligilo korda o fera, per qua onu interligas la manui di homo kaptita por impedar lu eskapar od agar malaji\lingui{FIS}
\artiklo{manovrar}%
\vorto{manovrar}\senco{1}{\textit{(trans.)} Funcionigar normale, per manuo (instrumento, mashino)}\senco{2}{\textit{(netrans.)} (I) \textit{(aludante armei, navi)} movar, evolucionar. (2) \textit{(aludante la persono qua imperas ad armei)} Evolucionigar armeo, eskadro}\senco{3}{\textit{(netrans.)} Kombinar serio de agi, de demarshi por atingar skopo}\lingui{DEFIRS}
\artiklo{mansardo}%
\vorto{mansardo}\senco{1}{Tekto di qua la karpenturo esas quaze ruptita}\senco{2}{Chambro facita sub la quaze-ruptita karpenturo, e di qua la parto supra esas plu streta pro la inklineso di la tekto}\lingui{DFIR}
\artiklo{mansheto}%
\vorto{mansheto} Garnituro, sutita o muntita kom ornivo ye la extremajo di la maniki di kamizo, e qua salias de la vesto, de la robo, tale ke la karpo kovresas\lingui{DFIR}
\artiklo{mantear}%
\vorto{mantear}\fako{trans.} Molestar (ulu) per igar lua korpo saltar de lito-kovrilo quan sukusas quar personi\lingui{S}
\artiklo{mantelo}%
\vorto{mantelo} Vesto ampla e longa, sen maniki, quan la homi *weras sur lia vesti cetera\lingui{DEFIRS}
\artiklo{mantenar}%
\vorto{mantenar}\fako{trans.}\senco{1}{Tenar ye maniero tala ke nula movo lua esas posibla}\senco{2}{Tenar en stando tala ke nula chanjo o vario esas posibla}\lingui{EFIS}
\artiklo{mantilio}%
\vorto{mantilio} Peco de stofo, de dentelo nigra, quan la Hispanini *weras sur sua kapo e qua retrofalas adsur la shultro\lingui{DEFIRS}
\artiklo{mantiso}%
\vorto{mantiso}\fako{matem.} La parto decimala di logaritmo
\artiklo{manufakturo}%
\vorto{manufakturo} Granda laboreyo en qua fasonesas, manue, ula produkturi qui postulas kelka delikateso laborala. (La manufakturi sencese remplasesas da « fabrikeyi », segun ke la mashini substitucesas a la manuolaboro)\lingui{DEFRS}
\artiklo{manuo}%
\vorto{manuo}\fako{anat.} Organo di prenado e di tushado, ye la extremajo di la brakio, che la homo, e kun kin fingri de qui la kinesma (la polexo) povas opozar su a la quar cetera\lingui{EFIS}
\artiklo{manuskripto}%
\vorto{manuskripto} Verko skribita manue o per skribmashino, ankore ne imprimita\lingui{DEFIRS}
\artiklo{manzanilo}%
\vorto{manzanilo}\fako{bot.} Arboro di Antili, ek la familio « euforbiacei », di qua la stipo e la frukto kontenas veneno\latina{hippomano}\lingui{DEFIS}
\artiklo{mapo}%
\vorto{mapo} Folio ek papero rezistiva, ma flexebla (facita ek plura paperfolii kunglutinita) od ek papero, sur qua esas reprezentita la surfaco di la Terglobo o di un ek olua parti\lingui{E}
\artiklo{marabuto}%
\vorto{marabuto}\fako{religio di Islam} Mohamedisto konsakrita a la praktiko ed a la doco di sua religio\lingui{DEFIRS}
\artiklo{maraskino}%
\vorto{maraskino} Ratafio qua fabrikesas en Italia, ek cerizi acida\lingui{DEFIRS}
\artiklo{marasmo}%
\vorto{marasmo}\fako{patol.} Langoro qua akompanas la deperiso, la senkurajeso\lingui{DEFIRS}
\artiklo{marchandar}%
\vorto{marchandar}\fako{trans.} Probar obtenor po preco min granda\lingui{F}
\artiklo{marchar}%
\vorto{marchar}\fako{netrans.} Pozar un ek sua pedi, e pose la altra, e tale ri-agante lo, por irar ta-direcione o ca-direcione\lingui{DEFIRS}
\artiklo{marcipano}%
\vorto{marcipano} Kukajo ek mandeli pistita e sukro\lingui{DEFIRS}
\artiklo{mardio}%
\vorto{mardio} Dio triesma di la semano, qua sequas lundio\lingui{FIS}
\artiklo{mareo}%
\vorto{mareo} Movo alternanta (di depreseso e di plualteso) di maro qua, en 24 hori, retroiras pokope ed abandonas la litoro, e pose ri-venas kovrar la spaco quan lu lasabis sikeskar\lingui{FIS}
\artiklo{mareografo}%
\vorto{mareografo}\fako{fiziko} Instrumento por mezurar ed enrejistrar la varii di la marei, per flotacero, trasante kurvo.qua reprezentas la ampleso-grado di la mareo, sur paperbendo senfina quan desvolvas mekanismo horlojala\lingui{EF}
\artiklo{mareometro}%
\vorto{mareometro} Sinonimo di « mareografo »
\artiklo{margarino}%
\vorto{margarino}\fako{kemio} Kombinuro ek margar-acido e glicerino, uzata kom butro\lingui{DEFIS}
\artiklo{margravo}%
\vorto{margravo}\fako{historio} Nomo di ula (olima) princi suverena di Germania\lingui{DEFIS}
\artiklo{margrito}%
\vorto{margrito}\fako{bot.} Nomo di krizantemo kun kolori diversa, qua divenis duopla per kultivo – de la astero di Chinia. Margriteto\latina{chrysantemum leucanthemum; bellis}\lingui{EF}
\artiklo{mariabalno}%
\vorto{mariabalno} Aquo varma, en vazo, aden qua onu plunjas e mantenas korpo quan onu volas varmigar\lingui{DFIS}
\artiklo{mariajar}%
\vorto{mariajar}\fako{trans.}\senco{1}{Unionar legitime (viro e muliero), celebrante ta uniono ye la nomo di la lego o di religio}\senco{2}{Rezolvigar la mariajo}\lingui{EFI}
\artiklo{marinar}%
\vorto{marinar}\fako{netrans.) (aludante karno, fisho} Esar en vinagro od en vino kun herbi aromatoza, spici, dum kelka tempo ante koqueso\lingui{DEFIRS}
\artiklo{marioneto}%
\vorto{marioneto}\senco{1}{Figurino qua reprezentas persono (viro o muliero) quan onu igas agar, marchar, gestifar, e c. per la fili de qui lu pendas}\senco{2}{\textit{(metaf.)} Persono quan altru direktas, igas movar, segun sua arbitrio}\lingui{DEFIRS}
\artiklo{marjino}%
\vorto{marjino} Spaco sen-texta, lakuno, cirkum paginedo de texto skribura od imprimura\lingui{DEFIS}
\artiklo{markezo}%
\vorto{markezo}\senco{1}{\textit{(olim)} Chefo di domeno an la frontiero di stato}\senco{2}{En la hierarkio de la tituli di nobeleso : ta qua venas pos la duko ed avan la komto. (Ta vorto esas ligita a la nomo di la persono per la prepoziciono « de ». Ex. : markezo de Beaufront)}\lingui{DEFIRS}
\artiklo{marko}%
\vorto{marko}\senco{1}{Traco lasita da kozo e per qua onu povas *rekonocar lu}\senco{2}{Traco facita sur kozo por igar lu *rekonocebla}\senco{3}{Marko obliterala da la posto, qua indikas sur la letri la loko e la dio di la departo e di la arivo}\senco{4}{\textit{(*timbro)} Flug-sigluro quan onu glutinas sur letro (por afrankar lu), sur afisho, sur quitigo-akto, e c. (por pagar la imposto-pekunio) e qua indikas sua preco}\lingui{DEFIRS}
\artiklo{markotar}%
\vorto{markotar}\fako{trans.} Enterigar an-sub la surfaco di sulo brancho di la planto – brancho quan onu separas de la stipo precipua kande lu radifikabos (separo per trancho)\lingui{FI}
\artiklo{marmelado}%
\vorto{marmelado} Disho ek frukti qui aplastesis e cesis havar sua formo normala pro koqueso kun sukro. Anke : la stando di karno, di fisho, e c. hiperkoquita. \textit{(anke metaf.)}\lingui{DEFIRS}
\artiklo{marmito}%
\vorto{marmito} Implemento koquala – vazo terakota, o metala, kun o sen pedo o pedi, en qua onu koquas dishi\lingui{FIS}
\artiklo{marmoro}%
\vorto{marmoro} Petro kalkoza blanka o koloroza, veinoza, kun grano tenua, polisebla ad-bele, uzata en arkitekturo ed en la statuarto\lingui{DEFIRS}
\artiklo{marmoto}%
\vorto{marmoto}\fako{zool.} Quadripedo rodera, ek la klaso « gliri », amansebla, hibernanta\latina{arctomys alpina}\lingui{EFIS}
\artiklo{marno}%
\vorto{marno}\fako{geol.} Tero grasa, mixuro naturala ek argilo e kalko, apta emendar ula sorti di tereni\lingui{FI}
\artiklo{maro}%
\vorto{maro}\senco{1}{Extensajo vasta ek aquo saloza, qua okupas parto grandega di la surfaco di la Terglobo}\senco{2}{Parto di ta extensajo vasta ek aquo saloza, situita en ta o ca regiono}\senco{3}{\textit{(metaf.)} Extensajo vasta di liquido – ek aquo konjelita. (Uzata anke pri granda quanto de sucii, de desfacilaji en qui onu esas quaze imersita)}\lingui{DEFIRS}
\artiklo{marodar}%
\vorto{marodar}\fako{netrans.}\senco{1}{\textit{(aludante soldati qui marchas o kampanias)} Furtetar ek la agri, ek la farmaji, ek la vilaji}\senco{2}{Furtetar frukti, legumi, ek la agri senklozita}\lingui{DEFRS}
\artiklo{marokino}%
\vorto{marokino}\senco{1}{Felo de kapro plu o min granoza, apretita per galo o sumako}\senco{2}{Felo di mutono o de bovyuno apretita tale ke lu imitas marokino}\lingui{DFIS}
\artiklo{marono}%
\vorto{marono}\fako{bot.} Frukto, plu grosa kam kastaneo, qua havas nur un pulpo-bloko por okupar la shelo dornoza – produktita da varietato di kastaniero greftita\latina{aesculus hippocastanum}\lingui{EF}
\artiklo{marquizo}%
\vorto{marquizo} Avantekto, maxim-multa-kaze ek vitroplaki, muntita super la pordo di edifico por ke onu povez enirar od ekirar shirmate\lingui{DEFIS}
\artiklo{marshalo}%
\vorto{marshalo}\senco{1}{Oficiro generala qua preiris la armei por preparar la voyo e lo necesa por la kampo}\senco{2}{Oficiro generala di qua la grado korespondis ad olta di brigado-generalo}\senco{3}{Oficiro generala qua esis super la generali}\lingui{DEFIRS}
\artiklo{marsho}%
\vorto{marsho} Tereno dilutita da aquo qua povas nur desfacile for-fluar. – Tereno basa, humida, apta pri la kultivo di legumi\lingui{E}
\artiklo{marsuino}%
\vorto{marsuino}\fako{zool.} Cetaceo ek la klaso « delfini », kun muzelo obtuza e rondatra\latina{delphinus phocoena}\lingui{FIS}
\artiklo{marsupialo}%
\vorto{marsupialo}\fako{zool.} Quadripedo qua havas posho abdominala ube esas la mameli ed ube la yuni, naskinta premature, havas sua refujeyo dum la gravideso suplementa\lingui{EFIS}
\artiklo{martelo}%
\vorto{martelo} Utensilo di perkuto, qua konsistas ye mancho e fera bloko di qua un latero (kapo) esas ordinare rektangula, kelke konvexa, e la altra (dorso) dinigita bizelatre\lingui{FIS}
\artiklo{martinpeskero}%
\vorto{martinpeskero}\fako{zool.} Ucelo ek la klaso « paseri », kun plumaro dazlanta, qua vivas alonge la riveri e su nutras per fishi quin lu snapas per sua beko\latina{alcedo hispida}\lingui{F}
\artiklo{martiro}%
\vorto{martiro}\senco{1}{La persono qua tormentesis od ocidesis pro ne renegar sua profeso di fido a Iesu, en la historio di kristanismo}\senco{2}{La persono qua supliciesis sen agir por justifikar lo}\senco{3}{La persono qua sufris por permanar fidela a sua idealo}\lingui{DEFIS}
\artiklo{marto}%
\vorto{marto} La monato triesma di la yaro\lingui{DEFIRS}
\artiklo{martro}%
\vorto{martro}\fako{zool.} Mikra mamifero di qua la korpo esas longa e la felo tre prizata\latina{mustela martes}\lingui{DEFIS}
\artiklo{marvelo}%
\vorto{marvelo} To quo (o : ta qua) astonas pro lua beleso, lua grandeso\lingui{EFIS}
\artiklo{masajar}%
\vorto{masajar}\fako{trans.} Palpar, klemar, petrisar per manui (la parti diversa di la homo-korpo) por acelerar la sango-cirkulo, por igar plu flexebla la membri\lingui{DEFIRS}
\artiklo{masakrar}%
\vorto{masakrar}\fako{trans.} Ocidar per frapar obstinoze, e pelmele\lingui{DEFIS}
\artiklo{*mashelo}%
\vorto{*mashelo}\fako{anat.} Ensemblo ek la du parti osta di la boko (maxilo e mandibulo) qua esas garnisita ye denti che la vertebrozi
\artiklo{mashinacar}%
\vorto{mashinacar}\fako{trans.} Kombinar artificoze, ruzoze, ula moyeni por atingar skopo quan onu ne audacas agnoskar publike\lingui{DEFIS}
\artiklo{mashino}%
\vorto{mashino} Asembluro de peci, kombinita tale ke li produktas ula efekti\lingui{DEFIRS}
\artiklo{masho}%
\vorto{masho}\senco{1}{Singla ek la mikra slingi de filo ek lino, kotono, silko, kordo, di qui la interplekto formacas texuro laxa (reto, trikoturo, e c.)}\senco{2}{Spaco vakua, di qua ta slingo formacas la konturo}\lingui{DE}
\artiklo{masiva}%
\vorto{masiva} Qua esas maso kompakta\lingui{DEFIRS}
\artiklo{maskar}%
\vorto{maskar}\fako{trans.} Celar sub ta peco de veluro o de satino, di qua la formo quaze muldas la vizajo, kun trui qui korespondas a la du okuli ed a la boko – e quan olim *weris la mulieri por prezervar sua karnaciono – e quan eli *weras ankore nun en la bali de maskilieri\lingui{DEFIRS}
\artiklo{maskarono}%
\vorto{maskarono}\fako{arkitekt.} Figuro di fantazio, alta-reliefa, o basreliefa, per qua onu ornas entablamenti, klozo-petri\lingui{DEFRS}
\artiklo{maskulo}%
\vorto{maskulo}\senco{1}{Individuo ek la sexuo qua havas la fakultato fekundigar la altra sexuo}\senco{2}{\textit{Maskula} (gram.) : Qua apartenas a la genro quan onu uzas por indikar la enti \textit{maskula}, o ti qui (arbitriale) esas asimilita ad ilci}\lingui{DEFIS}
\artiklo{maso}%
\vorto{maso} Asemblajo grandega de kozi, o parti de kozi, qui formacas quaza ensemblo\lingui{DEFIRS}
\artiklo{masonar}%
\vorto{masonar}\fako{trans.} \textit{(en la konstrukto di edifico, di domo, di ponto, e c.)} Facar ta parto di la laboruro qua konsistas ye la asemblo di la quadri, di la kruda petri, di la briki, e c. e la junturi, la induturi ek mortero, gipso, e c.\lingui{EF}
\artiklo{mastico}%
\vorto{mastico} Nomo di substanci diversa, kompozuro, mi-solida, quin onu uzas por stopar trui, fenduri, junturi, e qui hardeskas segun ke tempo fluas\lingui{DEFI}
\artiklo{mastikar}%
\vorto{mastikar}\fako{trans.} Pecigar per la denti, per la movado di la maxilo sur la mandibulo, (la nutrivi solida) por igar li plu facile glutebla e digestebla\lingui{EFIS}
\artiklo{masto}%
\vorto{masto}\senco{1}{\textit{(navig.)} Longa peco de ligno, cilindra e rekta, muntita vertikale sur navo, barko, e destinita a suportar la yardi an qui esas ligita la segli}\senco{2}{Ligno-peco plu o min longa, muntita vertikale, en gimnastikeyo por exerco pri klimo}\lingui{DEFRS}
\artiklo{mastodonto}%
\vorto{mastodonto}\fako{paleont.} Mamifero fosila, proxima a la elefanto\lingui{DEFIRS}
\artiklo{mastro}%
\vorto{mastro}\senco{1}{La persono qua dominacas personi, kozi}\senco{2}{La persono qua dominacas en ta o ca domeno}\lingui{DEFIS}
\artiklo{matadoro}%
\vorto{matadoro}\fako{precipue en Hispania} La viro qua, en la tauro-kombati, havas kom tasko ocidar la animalo per espadostroko\lingui{DEFIRS}
\artiklo{matematiko}%
\vorto{matematiko} La cienco (aritmetiko, geometrio o mekaniko) di qua la skopo esas la mezuro di la *grandori\lingui{DEFIRS}
\artiklo{materio}%
\vorto{materio}\senco{1}{Substanco ye qua kozo konsistas}\senco{2}{To quo esas la objekto di la agado o facado da la homo}\senco{3}{\textit{(filoz.)} Substanco korpa. (antonimo di « spirito »)}\lingui{DEFIRS}
\artiklo{matida}%
\vorto{matida}\senco{1}{Qua ne brilas}\senco{2}{\textit{(bruiso)} Poke sonora}\senco{3}{\textit{(koloro)} Qua ne dazlas, nek mem nur brilas}\lingui{FS}
\artiklo{matino}%
\vorto{matino} La tempo qua fluas inter jornesko e dimezo\lingui{EFIS}
\artiklo{matloto}%
\vorto{matloto}\fako{koquarto} Disho ek plura sorti de fishi, pecigita, preparita per vino\lingui{F}
\artiklo{mato}%
\vorto{mato} Texuro de palio, de junko tresigita, per qua onu kovras la planko-sulo, la muri\lingui{DE}
\artiklo{matraco}%
\vorto{matraco} Longa kuseno, polsterizita, plenigita ye lano, buro, krino, e c., quan onu extensas sur lito e sur qua onu su kushas\lingui{DEFIR}
\artiklo{matraso}%
\vorto{matraso} Vazo de vitro, di qua la kolo esas longa e streta, quan onu uzas en farmacio, kemio\lingui{EFIS}
\artiklo{matrico}%
\vorto{matrico} Mulduro qua produktas reliefajo per perkuto, giso, e c.\lingui{DEFIRS}
\artiklo{matrikulo}%
\vorto{matrikulo} Listo, etato, en qua esas enskribita la nomi di omna individui di korpo, di societo, – e, specale, listo, etato, en qua onu enskribas la nomo, prenomo ed eniro-numero di omna soldati qui eniras regimento\lingui{DEFIRS}
\artiklo{matro}%
\vorto{matro} Patrino homa. (Equivalanto sentimentala di « patrino ».)
\artiklo{matrono}%
\vorto{matrono} Muliero respektinda pro lua matureso evala e lua bona reputeso\lingui{DEFIRS}
\artiklo{matura}%
\vorto{matura}\senco{1}{\textit{(aludante frukto, grano)} Qua esas komplete desvolvita, developita}\senco{2}{\textit{(metaf.)} Qua jus atingis la punto maxim favoroza (por), la maxim oportuna (por)}\lingui{EFIS}
\artiklo{matutino}%
\vorto{matutino}\fako{liturgio katol.} Parto unesma di la oficio, qua dicesis en la kloko unesma pos nokto-mezo – e dicesas nun (ecepte en ula monakeyi) en la vigilio\lingui{EFIS}
\artiklo{mauzoleo}%
\vorto{mauzoleo} Tombo monumenta\lingui{DEFIRS}
\artiklo{maxilo}%
\vorto{maxilo}\fako{anat.}\senco{1}{La parto supra di la du parti osta di la boko, qui suportas la denti che la animali vertebroza}\senco{2}{\textit{(metaf.)} Nomo di la du fero-peci qui uzesas por tenar fixe objekto per klemo}\lingui{EFI}
\artiklo{maxim}%
\vorto{maxim}\fako{gram.} Indikas la gradaro-somito pri la eso-maniero quan expresas adjektivo od adverbo\lingui{DEFIRS}
\artiklo{mayo}%
\vorto{mayo} La monato kinesma di la yaro\lingui{DEFIRS}
\artiklo{mayoliko}%
\vorto{mayoliko} Olima sorto di fayenco Hispana ed Italiana\lingui{DEFIRS}
\artiklo{mayonezo}%
\vorto{mayonezo}\fako{koquarto} Sauco kolda, ek oleo quan onu quirlas kun vitelo til kande lu divenis mi-ferma\lingui{DEFIRS}
\artiklo{mayoro}%
\vorto{mayoro}\fako{armeo} Oficiro generala qua, helpe di generalisimo, funcionas kom chefo di la stabo\lingui{DEFIRS}
\artiklo{mayuskulo}%
\vorto{mayuskulo} Litero plu granda, e di formo maxim-multa-kaze altra, kam la minuskulo korespondanta – uzata precipue kom la rang-unesma litero di la nomi propra, e pos punto\lingui{DEFIS}
\artiklo{mazo}%
\vorto{mazo}\fako{olim} Armo ek fero, kun kapo tre pezoza – en ula kazi : garnisita ye pinti, quan onu manuagis quale Herkules manuagis sua asom-bastono tuberoza, plu grosa ye un ek sua extremaji\lingui{EFIS}
\artiklo{mazurko}%
\vorto{mazurko} Danso segun mezuro tri-tempa, plu lenta kam valso, ed en qua la tempo triesma esas plufortigita\lingui{DEFIRS}
\artiklo{me}%
\vorto{me}\fako{gram.} Pronomo personala, qua reprezentas la persono qua parolas o skribas\lingui{DEFIRS}
\artiklo{meandro}%
\vorto{meandro}\fako{arkitekt. Greka} Ornamento-desegnuro, formacita ek linei e vergeti qui interkrucumas\lingui{F}
\artiklo{meceno}%
\vorto{meceno} Personego richa e socio-povoza, qua favoras pekuniale la literaturisti e la artisti\lingui{DEFIRS}
\artiklo{mecho}%
\vorto{mecho} Kordono, bendo, asembluro de kotono e kanabo, cirkumata da sebo o de vaxo, por facar kandeli, bujii, ceri, o -imbibite ye oleo, petrolo, e c. – por brular en lampo
\artiklo{medalio}%
\vorto{medalio}\senco{1}{Peco de metalo, estampita por perpetuigar la memoro di ago remarkinda, di personego famoza o glorioza, e qua reprezentas objekto o legendo qua koncernas lu}\senco{2}{Peco de metalo qua reprezentas Iesu, la Virgino, e c., od irgaltra objekto di santeso, quan la katoliki *weras sur su od igas pendar de sua chapleto}\senco{3}{Peco de metalo, donita premie a la laureato di ula konkursi – rekompenso a la personi qui agis prodi su-sakrifikala – memorigive, a ti qui partoprenis ula kampanio od ecelis per milit-agi}\senco{4}{Peco de metalo, kun numero, quan (en ula landi) devas *werar la personi qui esas portisti di kargaji, e c.}\lingui{DEFIRS}
\artiklo{medaliono}%
\vorto{medaliono} Juvelo plata, ronda od ovala, en qua onu enklozas portreto, haro-lokli, reliquii, e c.\lingui{DEFIRS}
\artiklo{mediacar}%
\vorto{mediacar}\fako{netrans., pri} Agar por konciliar personi qui deskonkordas (anke pri aferi komercala)\lingui{DEFIS}
\artiklo{mediano}%
\vorto{mediano}\fako{geom.) (pri figuro plana} Lineo rekta qua dividas aden du parti interegala, omna kordi, paralela ye un direciono. – Rekto qua juntas la somito di triangulo a la mezo di la latero opozita\lingui{DEFI}
\artiklo{mediata}%
\vorto{mediata}\senco{1}{\textit{(aludante la kozi)} Qua, situita inter du termini, efikas por pasigar de un de la termini a la altra}\senco{2}{\textit{(aludante la personi)} Dicesas pri la ago da persono quan onu uzas por atingar ta o ca rezultajo}\lingui{DEFIS}
\artiklo{medicino}%
\vorto{medicino} La arto risanigar, fondita sur la cienco pri la morbi e la medikamenti\lingui{DEFIRS}
\artiklo{medikamento}%
\vorto{medikamento} Substanco uzata, interne od extere, kom remediilo lor maladeso\lingui{DEFIS}
\artiklo{mediko}%
\vorto{mediko} Persono a qua esas grantita la diplomo « doktoro pri medicino »\lingui{DEFIS}
\artiklo{medio}%
\vorto{medio}\senco{1}{To quo cirkondas omna-sinse}\senco{2}{Ensemblo ek la socio, mori, eventaji, en qui vivas o vivis persono, od en qui eventas od eventis fakto}
\artiklo{meditar}%
\vorto{meditar}\senco{1}{\textit{(netrans., pri)} Pensar profunde pri temo}\senco{2}{\textit{(trans.)} Studiar profunde (temo)}\lingui{DEFIS}
\artiklo{mediumo}%
\vorto{mediumo} Persono qua mediacas en la operaci di spiritismo\lingui{DEFIRS}
\artiklo{medulo}%
\vorto{medulo}\senco{1}{\textit{(anat.)} Substanco mola, oleoza, flavatra, qua plenigas la kavaji e la areoli di la substanco sponjatra di la osti}\senco{2}{\textit{(metaf.)}(1) La parto maxim nutriva : (2) La parto maxim intima, di la substanco}\lingui{EFIS}
\artiklo{meduzo}%
\vorto{meduzo}\senco{1}{\textit{(mitol.)} Un ek la tri Gorgoni, qua unika esis mortiva lu senkapigesis da Perseus}\senco{2}{\textit{(historio naturala)} Zoofito ek substanco gelatinatra, sen-forma}\latina{medusa}\lingui{DEFIRS}
\artiklo{mefistofelala}%
\vorto{mefistofelala} Di qua la perfideso esas diablala\lingui{DEFI}
\artiklo{mefito}%
\vorto{mefito} Exhalajo qua kontenas gasi nekonvenanta por respireso\lingui{DEFIS}
\artiklo{megafono}%
\vorto{megafono}\fako{akust.} Aparato quan onu uzas por augmentar la audebleso di la soni de elektrofoni, radio- o televizionorecevili\lingui{DEFIRS}
\artiklo{megalito}%
\vorto{megalito} Petro grosega di la monumenti prehistoriala\lingui{DEFIRS}
\artiklo{megasporangio}%
\vorto{megasporangio}\fako{biol.} Sinonimo di « makrosporangio »\lingui{DEFIS}
\artiklo{megasporo}%
\vorto{megasporo}\fako{biol.} Sinonimo di « makrosporo »\lingui{DEFIS}
\artiklo{megaterio}%
\vorto{megaterio}\fako{paleont.} Granda mamifero fosila, ek la klaso « sen-denti »\lingui{DEFIRS}
\artiklo{megero}%
\vorto{megero} Muliero tre maligna, akra-humora\lingui{DFI}
\artiklo{mejisar}%
\vorto{mejisar} \textit{(trans.)} \textit{(teknol.)} Preparar (la feli di mutono, agnelo, kaprido) por la faco di ganti, futero, furvesti\lingui{F}
\artiklo{mekaniko}%
\vorto{mekaniko} La cienco di la movi\lingui{DEFIRS}
\artiklo{mekanismo}%
\vorto{mekanismo}\senco{1}{Aranjuro de peci, de risorti, qua eventigas ula movi}\senco{2}{\textit{(filoz.)} Doktrino segun qua materio konsideresas kom extensajo pasiva, movata da forco qua esas extera ye lu}\lingui{DEFIS}
\artiklo{mekanisto}%
\vorto{mekanisto} Partisano di la doktrino filozofiala « mekanismo »
\artiklo{mekanomorfoso}%
\vorto{mekanomorfoso}\fako{biol.} Ensemblo ek la influi quin la efiki mekanikala (shoki, presi, ocili) agas a la kresko di la embrio e di ula organi
\artiklo{melanismo}%
\vorto{melanismo}\fako{patol.} Stando di la pelo qua, loke o totakorpe, aquiras koloro brunatra o nigratra, kun intenseso-grado variiva, ma (omnakaze) ne precize inkluzita
\artiklo{melankolio}%
\vorto{melankolio} Tristeso austera e nepreciza. Stando maladatra, morbala, karakterizata da depreseso mentala, psikala e fiziologiala\lingui{DEFIRS}
\artiklo{melaso}%
\vorto{melaso} Reziduo siropatra de la kristalesko e rafineso di sukro\lingui{DEFIS}
\artiklo{meleo}%
\vorto{meleo}\fako{en kombato} Intermixeso sen-ordina, pelmele, di la kombatanti qui frapas reciproke\lingui{EF}
\artiklo{melinito}%
\vorto{melinito}\fako{kemio} Mixuro exploziva, di qua la bazo esas trinitrofenolo 2.4.6\lingui{DEFIRS}
\artiklo{meliso}%
\vorto{meliso}\fako{bot.} Planto aromata, ek la familio « labiozi », di qua la folii uzesas por la faco di aquo alkoholizita e stomako-medikamenta\lingui{DEFIRS}
\artiklo{melkar}%
\vorto{melkar}\fako{trans.} Extraktar la lakto di femino, per presar lua mamo\lingui{DE}
\artiklo{melodio}%
\vorto{melodio}\senco{1}{\textit{(muziko)} Intersequo de soni qui plezas a la audanti od askoltanti}\senco{2}{Intersequo de vorti, de frazi qui plezas a la audanti od askoltanti}\lingui{DEFIRS}
\artiklo{melodramo}%
\vorto{melodramo}\senco{1}{Dramato, akompanata da muziko}\senco{2}{Dramato populara, qua probas emocigar per la extremeso di la cirkonstanci e la exajero di la sentimenti}\lingui{DEFIRS}
\artiklo{melolonto}%
\vorto{melolonto}\fako{zool.} Insekto koleoptera, qua aparas cirkum la fino di aprilo, e di qua la larvo devoras la radiki di la planti\latina{melolontha vulgaris}\lingui{IL}
\artiklo{melono}%
\vorto{melono}\fako{bot.} Planto ek la familio « kukurbitacei », di qua la frukto esas sukoza e sukroza\latina{cucumis melo}\lingui{DEFS}
\artiklo{melopear}%
\vorto{melopear}\fako{netrans.} Kantar ritmoze e segun-mezure, por akompanar la deklamo parolata\lingui{EFIRS}
\artiklo{mem}%
\vorto{mem}\fako{gram.} Manier-adverbo uzata (1) avan komparativo (di supereso o di infreso) por plufortigar lu; (2) avan substantivo o verbo, por insistar\lingui{F}
\artiklo{membrano}%
\vorto{membrano}\senco{1}{Tisuo organala, che animalo o vejetanto, qua envolvas ula organi, absorbas, exhalas, sekrecas ula fluidi}\senco{2}{\textit{(teknol.)} Planko quan la bindisto pozas sur e sub kolono de kayeri quan lu volas pozor en presilo}\lingui{DEFIS}
\artiklo{membro}%
\vorto{membro}\senco{1}{Apendico unionita a la torso di la homo, di la animalo, per artiki}\senco{2}{Persono qua esas un ek la kompozanti di korporaciono politikala, literaturala – di asociitaro, di komuniono religiala}\senco{3}{Un ek la propozicioni ye qui konsistas frazo; un ek la frazi ye qui konsistas la ideo-developo}\senco{4}{Singla de la du parti di equaciono, quin separas la signo di inter-egaleso}\lingui{EFIS}
\artiklo{memento}%
\vorto{« memento »}\fako{liturgio katol.} Prego di la mes-kanono, en qua la sacerdoto memorigas la vivanti e, pose, la mortinti
\artiklo{memorando}%
\vorto{memorando} Noto diplomacala en qua memorigesas to quo eventis pri ta o ca afero. – Noto destinita memorigar ulo\lingui{DEF}
\artiklo{memorar}%
\vorto{memorar}\fako{trans.} Reprezentar, en sua spirito, kozo pasinta\lingui{EFIRS}
\artiklo{memorialo}%
\vorto{memorialo} Registro, libro, en qua esas konsignita la kozi historiala quin onu volas memoror\lingui{DEFIS}
\artiklo{menado}%
\vorto{menado}\fako{epoki antiqua} Muliero qua, lor la festi pri Bako, abandonis su a deliri furioza\lingui{F}
\artiklo{menajar}%
\vorto{menajar}\fako{netrans.} Okupesar da la administro, da la ekonomio, di sua hemo\lingui{F}
\artiklo{menajerio}%
\vorto{menajerio} Animali skarsa, stranja, remarkinda, ye qui konsistas la kolektajo di parko zoologiala\lingui{DEF}
\artiklo{mencionar}%
\vorto{mencionar}\fako{trans.} Igar remarkata per signo, nomo-dico, \textit{(pri libro)} per titulo-dico, ma sen cito\lingui{EFIS}
\artiklo{mendikar}%
\vorto{mendikar}\senco{1}{\textit{(netrans.)} Demandar almono}\senco{2}{\textit{(trans.)} Solicitar ulo kom almono o kom quaza almono}\lingui{EFIS}
\artiklo{menhiro}%
\vorto{menhiro}\fako{arkeol.} Nomo di granda bloki de petro, erektita izole, quin onu opinionis havar origino Keltala, ma qui aspektas evar de epoko mem plu antiqua\lingui{DEF}
\artiklo{meningito}%
\vorto{meningito}\fako{patol.} Inflamuro di la envelopilo cerebrala\lingui{EFIS}
\artiklo{meningo}%
\vorto{meningo}\fako{anat.} Singla de la tri membrani ye qui konsistas la envelopilo cerebrala\lingui{EFIS}
\artiklo{menisko}%
\vorto{menisko}\senco{1}{Korpo qua finas ye du surfaci sfera : ica, konvexa – ita, konkava}\senco{2}{Parto supra di la kolono de liquido, kontenata en tubo kapilara, konvexa o konkava (segun la naturo di la liquido)}\lingui{DEFIS}
\artiklo{menso}%
\vorto{menso}\fako{historio di la eklezio katol.} Revenuo atribuita a la kompri nutrala\lingui{F}
\artiklo{menstruar}%
\vorto{menstruar}\fako{netrans.} Esar la objekto di ta ekpulso di sango qua eventas cirkume singla-monate, dum kelka dii, de la korpo di la homino evo-mariajebla\lingui{DEFIRS}
\artiklo{mentiar}%
\vorto{mentiar}\fako{netrans.} Asertar voloze ulo kontrea ad exakteso, kontrea a la fakti\lingui{FIS}
\artiklo{mento}%
\vorto{mento} Maniero vidar, pensar, sentar : I. Impulso preponderanta, per qua onu agas; II. Stoko de idei, de sentimenti, qua preponderas en la ago-maniero kustumata di persono o di asemblitaro; III. Direciono generala di la intelekto, en la kozi a qui lu su aplikas; IV. Penso preponderanta di verko\lingui{EFIS}
\artiklo{mentono}%
\vorto{mentono}\fako{anat.} Saliajo di la mandibulo, sub la labio infra\lingui{FI}
\artiklo{mentoro}%
\vorto{mentoro} La persono qua esas selektita kom guidanto di ulu pri la agi di ica (generale puera od, admaxime, adolecanta)\lingui{DEFIRS}
\artiklo{menueto}%
\vorto{menueto}\senco{1}{Danso olima, segun mezuro tri-tempa, kun mikra pazi, di karaktero grava, en qua la akompanantulo e la siorino evolucionas, inter-reverencante ed intersalutante}\senco{2}{Muziko-peco, ordinare segun mezuro 3/4, qua sequas la « adagio », la « andante » di sonato, di simfonio}\lingui{DEFIRS}
\artiklo{menuo}%
\vorto{menuo} Listo de la dishi ye qui konsistas repasto, festino\lingui{DEFS}
\artiklo{menuzar}%
\vorto{menuzar}\fako{netrans.} Facar objekto ek ligno, ex. tabli, stuli, paneli, e c. (kontraste ye la faco di karpenturi)
\artiklo{mercenario}%
\vorto{mercenario} Soldato stranjera di qua onu pagas la servi\lingui{EFIS}
\artiklo{mercero}%
\vorto{mercero} Komerco di la persono qua vendas mikra vari por la sutisti, por la vesti\lingui{EFIS}
\artiklo{meridiano}%
\vorto{meridiano} Cirklo-kurvo di la sfero cielala e terglobala, qua pasas tra la du poli, perpendikle ye equatoro\lingui{DEFIRS}
\artiklo{meringo}%
\vorto{meringo} Kukajo lejera, ek la albumeno di ovi e sukro-pulvero\lingui{DEFIS}
\artiklo{merino}%
\vorto{merino} Mutono, di raso Hispana, di qua la lano esas tenua\lingui{DEFIRS}
\artiklo{meristemo}%
\vorto{meristemo}\fako{biol.} Tisuo yuna, di qua la celuli, parietizeskanta, interpresas tre forte ed havas la formo di poliedro\lingui{DEFIS}
\artiklo{meritar}%
\vorto{meritar}\fako{trans., pro} Havar yuro pri obtenor (ulo) pro sua konduto, maniero\lingui{EFIS}
\artiklo{merizo}%
\vorto{merizo}\fako{bot.} Frukto di la ceriziero sovaja\latina{cerasus avium}\lingui{F}
\artiklo{merkar}%
\vorto{merkar}\fako{trans.} Plasizar en sua stoko psikala de memoraji\lingui{DE}
\artiklo{merkato}%
\vorto{merkato}\senco{1}{Loko publika ube eventas la transakti komercala, ube vendesas la vari nutrala, e c.}\senco{2}{\textit{(ekonomiko)} Stando di ofro e demando en loko ube eventas la transakti komercala}\senco{3}{La komercistaro di urbo}\lingui{DEFIS}
\artiklo{merkurdio}%
\vorto{merkurdio} Dio di la semano, qua sequas mardio\lingui{FIS}
\artiklo{merkurio}%
\vorto{merkurio}\fako{kemio} Korpo simpla, metala, fluida ye la temperaturo ordinara\simbolo{Hg}\lingui{EFIS}
\artiklo{merlano}%
\vorto{merlano}\fako{zool.} Marfisho, ek la genero « gadi », di qua la karno esas nutrivo lejera\latina{gadus merlangus}\lingui{FR}
\artiklo{merlo}%
\vorto{merlo}\fako{zool.} Ucelo ek la klaso « paseri »\latina{turdus merula}\lingui{DFIS}
\artiklo{merlucho}%
\vorto{merlucho}\fako{zool.} Nomo di fishi diversa ek la genero « gadi », sikigita efike da suno-radii\latina{gadus merlucius}\lingui{EFIRS}
\artiklo{merogonio}%
\vorto{merogonio}\fako{biol.} Modo di fekundigo segun qua la spermatozoido unionas su a fragmento de protoplasmo feminala sen-nuklea\lingui{DEFIS}
\artiklo{meromorfa}%
\vorto{meromorfa}\fako{matem.) (pri funciono, interne di areo} Qua permanas holomorfa, ecepte pri ula punti ube lu divenas infinita\lingui{DEFIS}
\artiklo{merotomio}%
\vorto{merotomio}\fako{biol.} Operaco qua konsistas ye deprenar irga parto di organismo ed observar la konduto di ca parto\lingui{DEFIS}
\artiklo{mesajo}%
\vorto{mesajo} Komunikajo letra, oficala, da la chefo di la povo legaplikala, a la parlamento. – \textit{(anke metaf.)}\lingui{EFIS}
\artiklo{meso}%
\vorto{meso}\fako{religio katol.} Sakrifiko di la korpo e sango di Iesu, ofrata a Deo semble di pano e vino, por reprezentar e durar la sakrifiko di la kruco, ed aplikar a ni la frukti\lingui{DEFIRS}
\artiklo{mestiero}%
\vorto{mestiero}\senco{1}{Sorto di laboro manuala}\senco{2}{Okupivo quan onu komparas a laboro manuala}\lingui{FIR}
\artiklo{metabolio}%
\vorto{metabolio}\fako{biol.} Ta omna kambii qui eventas en organismo, de lua nasko til lua morto\lingui{DEFIS}
\artiklo{metacentro}%
\vorto{metacentro}\senco{1}{\textit{(fiziko)} Punto ube la perpendiklo ye la flotaco-plano, pasanta tra la gravito-centro e la impulseso-centro di korpo qua flotacas equilibroze, renkontras la vertikalo qua pasas tra la nova impulseso-centro, kande la equilibro di la korpo esas perturbita}\senco{2}{\textit{(navig.)} En navo qua ocilas cirkum axo horizontala, la centro di la kurvo trasata da la impulseso-centro di la fluido diplasita}\lingui{DEFIRS}
\artiklo{metacinezo}%
\vorto{metacinezo}\fako{biol.} Migro, dum mitoso, da du grupi de kromosomi ad lia poli respektiva\lingui{DEFIS}
\artiklo{metafazo}%
\vorto{metafazo}\fako{biol.} Fazo, o periodo, tre kurta di kariokinezo\lingui{DEFIS}
\artiklo{metafiziko}%
\vorto{metafiziko} La cienco pri la kozi qui transiras naturo, la domeno di la kauzi sekundara (qui recevas, de kauzo superiora, la povo produktar ula efekti)\lingui{DEFIRS}
\artiklo{metaforo}%
\vorto{metaforo}\fako{gram., retor.} Figuro qua konsistas ye indikar persono o kozo per termino qua implikas komparo tacita\lingui{DEFIRS}
\artiklo{metafrazar}%
\vorto{metafrazar}\fako{trans.} Interpretar verko skribura; agar traduko qua donas la senco di la originalo por igar lu komprenata, ma ne por igar perceptar lua belaji literaturala\lingui{DEFIRS}
\artiklo{metageometrio}%
\vorto{metageometrio} Omna geometrio qua esas plu generala kam la Euklid-geometrio, e di qua olca esas egardebla kom kazo partikulara
\artiklo{metakarpo}%
\vorto{metakarpo}\fako{anat.} Unionuro ek la kin osti paralela, ye qui konsistas la palmo di la manuo\lingui{EFIS}
\artiklo{metalo}%
\vorto{metalo} Korpo simpla, kun brilo partikulara, plu o min duktila e maleebla (forjebla o martelagebla)\lingui{DEFIRS}
\artiklo{metalografio}%
\vorto{metalografio} La cienco pri la metali\lingui{DEFS}
\artiklo{metaloido}%
\vorto{metaloido}\fako{kemio} Korpo simpla, solida, liquida o gasa, sen metal-brilo, mala konduktivo pri kaloro ed elektro, e qua, kande kombinita kun metalo, esas elektro-negativa\lingui{DEFS}
\artiklo{metalurgio}%
\vorto{metalurgio} La industrio qua extraktas la metali ek la mineyi ercala, e separas li de la materii ne-metala\lingui{DEFIS}
\artiklo{metamero}%
\vorto{metamero}\fako{biol.} Segmento primitiva quan onu renkontras sur la vertebrozi dum lua kresko, e qua, pose, remplasesas da la vertebro\lingui{DEFIS}
\artiklo{metamorfosar}%
\vorto{metamorfosar}\fako{trans.} Igar (ento) pasar de lua formo naturala a formo altra\lingui{DEFIRS}
\artiklo{metano}%
\vorto{metano}\fako{kemio} Termino unesma CH₄ ek la hidokarbidi saturita\lingui{DEFIS}
\artiklo{metar}%
\vorto{metar}\fako{trans.} Vestizar su per (co o to)\lingui{FI}
\artiklo{metastabila}%
\vorto{metastabila}\fako{fiziko} Dicesas pri stando fizikala partikulara – o di kristalo, o di mixuro – obtenita en kondicioni specala di temperaturo : stando nestabila teoriale, ma qua povas tamen persistar dum periodo tre longa, o mem senfine\lingui{DEF}
\artiklo{metatarso}%
\vorto{metatarso}\fako{anat.} Unionuro de la kin osti paralela ye qua konsistas la parto di la pedo qua esas inter tarso e fingri\lingui{EFIS}
\artiklo{metatezo}%
\vorto{metatezo}\fako{gram.} Permuto di literi en vorto. \textit{(ex. : prato vice parto)}\lingui{DEFIRS}
\artiklo{metempsikoso}%
\vorto{metempsikoso} Doktrino qua agnoskas existi intersequanta, ye qui la anmo pasas, de la korpo quan lu vivigis, aden olta di altra ento\lingui{DEFIRS}
\artiklo{meteoro}%
\vorto{meteoro}\fako{fiziko} Fenomeno qua eventas en atmosfero : cielarko, pol-auroro, pluvo, roso, nivo, vento, trombo, e c.\lingui{DEFIRS}
\artiklo{meteorologio}%
\vorto{meteorologio} Ta parto di fiziko qua traktas la fenomeni qui eventas en atmosfero\lingui{DEFIRS}
\artiklo{metileno}%
\vorto{metileno}\fako{kemio} Radiko divalenta CH₂, derivita de metano per la supreso di du atomi di hido\lingui{DEFIS}
\artiklo{metilo}%
\vorto{metilo}\fako{kemio} Radiko monovalenta CH₃, termino unesma di la serio de la radiki di la karbidi grasoza saturita, e qua derivesas de metil-alkoholo per supreso di oxihidrido\lingui{DEFIS}
\artiklo{metodo}%
\vorto{metodo}\senco{1}{Iro-modo segun-raciona quan onu praktikas por atingar sua skopo}\senco{2}{Iro-modo segun-raciona quan praktikas la spirito por atingar lo exakta, o por demonstrar ta exaktajo pos trovir lu}\lingui{DEFIRS}
\artiklo{metodologio}%
\vorto{metodologio} Ta parto di logiko qua studias « a posteriori » la metodi di la sorti diversa di konoci, di savi, e plu partikulare la metodi di la cienci diversa\lingui{DEFIS}
\artiklo{metonimio}%
\vorto{metonimio}\fako{gram.} Figuro di retoriko, uzata por expresar kozo per termino qua indikas altra kozo qua esas unionita a lu per relati necesa (kauzo po efekto, kontenanto po kontenato, e c.)\lingui{DEFIRS}
\artiklo{metonomazio}%
\vorto{metonomazio} Transformeso subisata da persono-nomo quan onu transportas aden linguo stranjera per tradukar lu \textit{(Dubois = Sylvius; Schwarzerd = Melanchton; Pesch = Persiko)}\lingui{DEFIRS}
\artiklo{metopo}%
\vorto{metopo}\fako{arkitekt.} Interspaco quadrata, maxim-multakaze ornita per skulturi – qua esas inter la triglifi di la frizo dorika\lingui{DEFIRS}
\artiklo{metriko}%
\vorto{metriko} Ensemblo ek la reguli qui koncernas la metro di la versi
\artiklo{metrito}%
\vorto{metrito}\fako{patol.} Inflamuro di la utero\lingui{F}
\artiklo{metro}%
\vorto{metro}\senco{1}{Unajo di longeso qua reprezentas la 1/10.000.000-ima parto di la quarimo di la meridiano Terglobala}\senco{2}{Nombro e dispozeso di la silabi ye qui konsistas singla verso, singla versajo, poemeto, e c.}\lingui{DEFIRS}
\artiklo{metrologio}%
\vorto{metrologio} La cienco di pondero e mezuro
\artiklo{metronomo}%
\vorto{metronomo}\fako{muziko} Instrumento, uzata por indikar la rapideso-gradi diversa di la movimento muzikala, per pendulo qua skandas plu o min rapide segun ke onu plualtigas o minaltigas pezajo movebla sur la stangeto gradizita\lingui{DEFIRS}
\artiklo{metropolito}%
\vorto{metropolito}\fako{eklezio Greka} Granda oficiero di la eklezio ortodoxa, di qua la rango hierarkiala esas inter la patriarko e la arkiepiskopi\lingui{DEFIRS}
\artiklo{metropolo}%
\vorto{metropolo}\senco{1}{Urbo, stato, konsiderata kompare a la kolonii qui naskis de lu}\senco{2}{\textit{(epoki antiqua)} Urbo chefala di provinco}\senco{3}{\textit{(historio religiala)} Urbo qua havas rezideyo arkiepiskopala}\lingui{DEFIRS}
\artiklo{mezentero}%
\vorto{mezentero}\fako{anat.} Rifalduro di la peritoneo, qua mantenas la parti diversa di la intestini\lingui{DEFIS}
\artiklo{mezereono}%
\vorto{mezereono}\fako{bot.} Arbusto ek la genero « dafnei », qua kreskas sur la monti arboroza, boskoza, di Francia e di la centroregioni di Europa\latina{daphne mezereum}
\artiklo{mezo}%
\vorto{mezo} Parto qua distas egale de la extremaji, en intervalo de loki o de epoki\lingui{I}
\artiklo{mezodermo}%
\vorto{mezodermo}\senco{1}{\textit{(embriol.)} Strato qua aparas, dum la su-developo prima, inter la du folieti di ektodermo ed endodermo}\senco{2}{\textit{(bot.)} Parto di la kortico, inter la strato korkatra e la envelopilo herbatra}\lingui{DEFIS}
\artiklo{mezokarpo}%
\vorto{mezokarpo}\fako{bot.} Strato di la perikarpo di frukto, inter la epidermo extera (epikarpo) e la epidermo interna (endokarpo) – pulpoza, che ula frukti (cerizi, abrikoti, e c.) – lignatra che altrai (bruo)\lingui{S}
\artiklo{mezolabo}%
\vorto{mezolabo} Olima instrumento di matematiki, qua konsistas ye tri paralelogrami, moviva sur kuliso, uzita por trovar mekanikale du meza proporcionali quin onu ne povis determinar geometriale en la problemo « duopligo di la kubo »\lingui{DEFIRS}
\artiklo{mezotinto}%
\vorto{mezotinto} Graburo por qua onu preparas, unesme, la grano sur la tota planko, til obtenar nigro perfekta, ed ube, pos dekalquir la desegnuro sur la kupro-plako, onu aplastas plu o min multe la grano per skrapilo por la loki klara, la minuanci, la mi-ombri\lingui{DEFI}
\artiklo{mezurar}%
\vorto{mezurar}\fako{trans., per} Evaluar (quanto) per lua raporto ye quanto determinita, sama-speca, selektita kom komparo-termino\lingui{EFI}
\artiklo{mi}%
\vorto{« mi »}\fako{muziko} Noto triesma di la gamo diatonika naturala, t.e. di la « do »-gamo
\artiklo{mi-}%
\vorto{mi-} Prefixo di qua la senco esas « duime »
\artiklo{miasmo}%
\vorto{miasmo} Emanajo des-salubra, nociva, de materii putranta\lingui{DEFIR}
\artiklo{miaular}%
\vorto{miaular}\fako{netrans.) (aludante kato} Agar la krio qua esas propra ye lua speco\lingui{DEFIS}
\artiklo{micelo}%
\vorto{micelo}\fako{biol.} La aglomerajo maxim mikra ek molekuli, qua havas la propraji fizikala di korpo
\artiklo{mielato}%
\vorto{mielato} Exuduro viskoza e sukroza, quan sekrecas, dum la periodi di aero-sikeso e di tero-sikeso, la folii di ula arbori
\artiklo{mielo}%
\vorto{mielo} Substanco siropatra, sukroza, quan la abeli elaboras ek la materii quin li rekoltis en la flori, e quin li igas ek-fluar aden la alveoli di lia abeluyo por la nutro di su e di la larvi\lingui{FIS}
\artiklo{mieno}%
\vorto{mieno} Aspekto di la vizajo\lingui{DEFR}
\artiklo{migrar}%
\vorto{migrar}\fako{netrans., ad, de} Diplasar su, livante ula lando por instalar su en altra. – Su-diplasar di ula sorti de animali qui iras ad altra klimati, sive periodale, segun la sezoni, sive pro cirkonstanci acidentala\lingui{EFIS}
\artiklo{migreno}%
\vorto{migreno} Doloro qua afektas, ordinare, nur parto di la kapo, partikulare la regiono temporala, ed akompanesas multa-kaze da perturbesi stomakala\lingui{DEFIR}
\artiklo{mikao}%
\vorto{mikao}\fako{mineral.} Silikato aluminoza, kun metal-brilo, quan onu povas separar facile ad lami diafana, per klivo\lingui{DEFIS}
\artiklo{mikologio}%
\vorto{mikologio} Ta parto di botaniko qua studias la fungi (champinioni)
\artiklo{mikologo}%
\vorto{mikologo} Autoro di traktato pri la fungi (champinioni)
\artiklo{mikorizo}%
\vorto{mikorizo} Fungo quan onu renkontras generale asociita kun radiki di vejetanti (precipue la radiki di la vejetanti montala : kastanieri, avelanieri, fagi, morusieri, monto-pini, e c.)\lingui{DEFIS}
\artiklo{mikra}%
\vorto{mikra}\senco{1}{Qua ne atingas la dimensioni ordinara, partikulare pri longeso, alteso, qualeso o quanteso}\senco{2}{Qua ne atingas la nivelo ordinara pri rango en la hierarkio, merito, la qualesi di kordio o di mento}\lingui{DEFIRS}
\artiklo{mikro-}%
\vorto{mikro-} Prefixo ciencala : la milionimo de…
\artiklo{mikrobo}%
\vorto{mikrobo} Celulo vivanta, di vejetanto od animalo, qua apartenas a la mondo di la infinite mikra, faktoro di ula fermentaci od ula morbi infektiva\lingui{DEFIRS}
\artiklo{mikrocentro}%
\vorto{mikrocentro}\fako{biol.} Parto centra di la sfero atraktiva, konsistanta ye plura kromosomi
\artiklo{mikrofono}%
\vorto{mikrofono}\fako{fiziko} Aparato akustikala, destinita ad ampligar la soni\lingui{DEFIRS}
\artiklo{mikrogameto}%
\vorto{mikrogameto}\fako{biol.} Elemento genetiva, mikra e moviva, kun karakteri maskulala
\artiklo{mikrografio}%
\vorto{mikrografio} La cienco qua traktas la preparo di la objekti por lia observado per mikroskopo, e la deskripto di olci pos lia exameneso per ta instrumento\lingui{DEFIRS}
\artiklo{mikrografo}%
\vorto{mikrografo} Sorto di pantografo, qua posibligas desegnar figuri extreme mikra
\artiklo{mikrometro}%
\vorto{mikrometro}\fako{metrologio} Instrumento destinita a la mezuro di mikra objekti e mikra imaji\lingui{DEFIRS}
\artiklo{mikrono}%
\vorto{mikrono}\fako{teknol.} La milima parto di la milimetro – unajo di mezuro qua adoptesis por la observi mikroskopala\lingui{DEFIRS}
\artiklo{mikronukleo}%
\vorto{mikronukleo}\fako{biol.} La nukleo maxim mikra che la infuzorii, opoze a la grosa nukleo \textit{(makronukleo)}
\artiklo{mikropilo}%
\vorto{mikropilo}\fako{biol.} Mikra orifico di la membrano vitelala di la ovulo, tra qua pasas (konjektate) la spermatozoido lor la fekundigo
\artiklo{mikroskopo}%
\vorto{mikroskopo} Instrumento optikala, ek lensi dispozita tale ke li igas la mikra objekti semblar plu grosa kam se regardata per okulo nuda\lingui{DEFIRS}
\artiklo{mikrosomo}%
\vorto{mikrosomo}\fako{biol.) (histol.} Granularo tenua quan onu renkontras en la parto figurala di citoplasmo\lingui{DEFIS}
\artiklo{mikrosporangio}%
\vorto{mikrosporangio} Organo, che algi diversa e kriptogami vaskuloza, en qua produktesas la mikrospori\lingui{DEFIS}
\artiklo{mikrosporo}%
\vorto{mikrosporo} Nomo di ta spori maxim mikra di la vejetanti, qui genitas la protali maskula\lingui{DEFIS}
\artiklo{mikrozoo}%
\vorto{mikrozoo} Ento vivanta qua esas videbla nur per mikroskopo\lingui{DEFIRS}
\artiklo{mil}%
\vorto{mil} 1.000
\artiklo{milano}%
\vorto{milano}\fako{zool.} Rapt-ucelo jornala, di qua la kaudo esas oblonga e forkatra, ye flugo rapida\latina{milvus ictinus}\lingui{FS}
\artiklo{milbo}%
\vorto{milbo}\fako{zool.} Insekto sen-ala, preske mikroskopala, qua prosperas en farino, fromajo, e c.
\artiklo{mili-}%
\vorto{mili-} Prefixo ciencala, qua signifikas « la milimo di… »
\artiklo{miliara}%
\vorto{miliara}\fako{medic.} Qua produktas, erupte, granda quanto de vezikuli tre mikra
\artiklo{miliardo}%
\vorto{miliardo} 1000 milioni
\artiklo{milibaro}%
\vorto{milibaro} Mezur-unajo di la preso atmosferala, qua equivalas 1/1000 di \textit{baro} (o : hektopiezo) od anke 1000 baryi
\artiklo{milico}%
\vorto{milico} Guardo nacionala, od institucuro analoga; trupo de soldati civitana, en la landi qui ne havas armeo permananta\lingui{DEFIRS}
\artiklo{miligramo}%
\vorto{miligramo} La milima parto di un gramo\lingui{DEFIRS}
\artiklo{milimetro}%
\vorto{milimetro} La milima parto di un metro\lingui{DEFIRS}
\artiklo{milimikrono}%
\vorto{milimikrono} La milionima parto di un milimetro\lingui{DEFIRS}
\artiklo{milio}%
\vorto{milio} Mezurunajo parkurala, ek mil pazi. (che la Romani antiqua 1472 metri; che la Britaniani ed Usani = 1609 metri; pri navigado = 1852 metri)\lingui{EF}
\artiklo{miliono}%
\vorto{miliono} 1.000.000. (10⁶)\lingui{DEFIRS}
\artiklo{militar}%
\vorto{militar}\fako{netrans., kontre} Agar lukto, per armi omnasorta (armi shokala, pafala, bombi, atomfuzei, e c.), kontre altra populo, o mem kontre la sua (milito civila)\lingui{DEFIS}
\artiklo{militarismo}%
\vorto{militarismo} Prepondero exajerita da la elemento « armeani », che naciono; doktrino di ti qui tendencas a ta prepondero\lingui{DEFIRS}
\artiklo{milyeto}%
\vorto{milyeto}\fako{bot.} Gramino qua produktas mikra grano, ye qua su nutras la uceli\latina{panicum miliaceum}
\artiklo{mimar}%
\vorto{mimar}\fako{trans.} Reprezentar per gesti, posturi, fizionomio-movi – sen uzar parolo\lingui{DEFIRS}
\artiklo{mimoso}%
\vorto{mimoso}\fako{bot.} Planto ek la familio « leguminosi », di qua ula speci kontraktas su kande onu tushas li\latina{mimosa}\lingui{DEFIRS}
\artiklo{min}%
\vorto{min} Adverbo qua indikas la inferioreso pri qualeso, quanteso, precozeso, extenseso, e c. = ne tam multe … -a\lingui{FIRS}
\artiklo{minacar}%
\vorto{minacar}\fako{trans., per}\senco{1}{Manifestar ad (ulu) sua intenco nocar lu. Avertar (ulu) ke lu timez ulo de altru}\senco{2}{\textit{(aludante kozo)} Manifestar, per kelka indici, ke, de lu, esas timenda ulo kom eventonta balde}\lingui{EFIS}
\artiklo{minar}%
\vorto{minar}\fako{trans.} Exkavar truo en qua onu pozas pulvero quan onu acendas por destruktar exploze to quo esas sur o super lu\lingui{DEFIRS}
\artiklo{minareto}%
\vorto{minareto} Turmo di moskeo\lingui{DEFIRS}
\artiklo{mineralo}%
\vorto{mineralo} Substanco qua konsistas ye materio kruda, neorganizita, neorganika\lingui{DEFIRS}
\artiklo{mineralogio}%
\vorto{mineralogio} Ta parto di la natur-historio qua traktas la minerali, la korpi ne-organika tala qualan kande onu renkontras li en tero o sur la surfaco di olca\lingui{DEFIRS}
\artiklo{minestrelo}%
\vorto{minestrelo}\fako{en la mezepoko} Poeto, muzikisto recitanta, qua kantis sua kompozuri ed olti di la poeti cetera\lingui{DEFIR}
\artiklo{miniaturo}%
\vorto{miniaturo}\senco{1}{Litero reda, trasita per miniumo, sur la manuskripti, la meso-libri, e c. por ornar la komenco-loko di la chapitri. – Litero ye kolori diversa}\senco{2}{Pikturi delikata, ek mikra temi, sur velino, pergamo, en la manuskripti, la meso-libri}\senco{3}{Pikturo tre delikata, tre mikra, facita per farbi dilutita en aquo e gumo, sur ivoro, sur velino}\senco{4}{Reprezenturo tre mikra-dimensiona}\lingui{DEFIRS}
\artiklo{minim}%
\vorto{minim}\fako{gram.} Superlativo absoluta di « pokeso »\lingui{DEFIRS}
\artiklo{miniona}%
\vorto{miniona} Di qua la mikreso esas charmoza\lingui{F}
\artiklo{ministerio}%
\vorto{ministerio} Laboreyo di ministro e di la oficieri ed komizi di lua departmenti. (EF.)
\artiklo{ministro}%
\vorto{ministro} Viro (o muliero) a qua la povo leg-aplikanta konfidis un ek la taski precipua di la guverno\lingui{DEFIRS}
\artiklo{miniumo}%
\vorto{miniumo}\senco{1}{\textit{(kemio)} Deutoxido di plumbo, substanco farba, reda}\senco{2}{Farbo oleoza, facita ek miniumo, uzata kom strato-unesma lor la induto di fero}\lingui{DEFIRS}
\artiklo{minora}%
\vorto{minora}\senco{1}{\textit{(yuro-cienco)} Qua ne atingis ankore la evo legostipulita koncerne la yuro agar libere pri sua persono e sua havaji}\senco{2}{\textit{(muziko)} Intervalo ek tri noti, qua kontenas un tono ed un mi-tono, o de sis noti, qua kontenas tri toni e du mi-toni}\senco{3}{\textit{(logiko)} \textit{(en silogismo)} La subjekto di la konkluzo (pro ke lu esas min ampla kam la atributo)}\lingui{DEFIS}
\artiklo{minoritato}%
\vorto{minoritato} La quanto minim granda de voci che asemblitaro qua esas votanta pri lo jus deliberita\lingui{DEFIS}
\artiklo{minto}%
\vorto{minto}\fako{bot.} Planto aromatoza, ek la familio « labiacei »\latina{mentha}\lingui{DEFIRS}
\artiklo{minucio}%
\vorto{minucio} Detalo mikrega qua ne meritas ke onu haltez por lu (pro lu)\lingui{DEFIS}
\artiklo{minus}%
\vorto{minus}\fako{aritm., algebro} Signo qua, avan quanto, signifikas ke olca esas sustraktenda\lingui{DEIRS}
\artiklo{minuskulo}%
\vorto{minuskulo}\fako{pri la literi di la alfabeto} Mikra litero (o tipo), kompare a mayuskulo, e ye qua konsistas la maxim multa literi di la vorti skribita\lingui{DEFIS}
\artiklo{minuto}%
\vorto{minuto}\senco{1}{Sisadekima parto di horo}\senco{2}{Sisadekima parto di un grado di cirklo}\lingui{DEFIRS}
\artiklo{mioceno}%
\vorto{mioceno}\fako{geol.} Un ek la quar granda faki di la ero terciara (olta qua sucedas la strati oligocena, e sur qua jacas la strati pliocena)\lingui{DEFIS}
\artiklo{miologio}%
\vorto{miologio} La anatomio di la muskuli\lingui{DEFIS}
\artiklo{miopa}%
\vorto{miopa}\fako{patol.} Di qua la vido-porteo esas kurta\lingui{DEFIS}
\artiklo{miozoto}%
\vorto{miozoto}\fako{bot.} Planto di qua la floro esas ciel-blua, rozea, e quan, en landi diversa di Europa, onu nomizas poezioze « Ne-obliviez-me »\latina{myosotis}\lingui{EFISL}
\artiklo{mirabelo}%
\vorto{mirabelo} Mikra pruno flava
\artiklo{mirajo}%
\vorto{mirajo} Optik-iluziono – refrakto qua, en la landi di qui la atmosfero esas varmega, igas vidar kom proxima la imaji di objekti tre fora. – \textit{(metaf. : Iluziono seduktoza)}\lingui{EFIR}
\artiklo{miraklo}%
\vorto{miraklo} Fakto supernatura, qua eventas violace di la *leyi (legi naturala)\lingui{DEFIR}
\artiklo{miriado}%
\vorto{miriado} 10.000
\artiklo{miriametro}%
\vorto{miriametro} 10.000 metri
\artiklo{miriapodo}%
\vorto{miriapodo}\fako{zool.} Animalo artikoza (artropodo) qua konsistas ye ringi, singla de qui havas paro de gambi\latina{myriapoda}\lingui{DEFIS}
\artiklo{mirmekofago}%
\vorto{mirmekofago}\fako{zool.} Mamifero ek la klaso « sen-denti », di qua la lango esas filo-forma, viskoza, e quan lu extensas sur la formikeyo por kaptar la formiki qui esas lua disho precipua\latina{myrmecophaga}\lingui{IL}
\artiklo{mirmekofila}%
\vorto{mirmekofila}\fako{bot.} Dicesas pri ula animali qui esas la gasti di la formiki – qui vivas kun ici, sive kom « invititi » (ex. : la afidii o planto-lausi), sive kom « neatenciti », sive kom vera paraziti\lingui{DEFIS}
\artiklo{mirmekoleono}%
\vorto{mirmekoleono}\fako{zool.} Insekto, analoga a libelulo, di qua la larvo jacas sur la fundo di funelo quan lu exkavis en tero e kaptas por manjar li la formiki od altra insekti qui falas aden lu\latina{myrmecoleon}\lingui{IL}
\artiklo{miro}%
\vorto{miro}\fako{bot.} Gumo-rezino aromoza\lingui{DEFIS}
\artiklo{mirtelo}%
\vorto{mirtelo}\fako{bot.} Genero de planti ek « vacinei » qua produktas beri nigratra, acideta, akreta\latina{vaccinium murtidus}\lingui{FIL}
\artiklo{mirto}%
\vorto{mirto}\fako{bot.} Arboreto kun flori blanka, quan la antiqui konsakrabis a Venus (Afrodito)\latina{myrtus}\lingui{DEFIRS}
\artiklo{mis-}%
\vorto{mis-} Prefixo qua signifikas « erore, nejuste »
\artiklo{misiono}%
\vorto{misiono} Senditaro (religiala, ciencala, diplomacala) di qua la tasko konsistas ye – rispektive – predikar, instruktar o docar, studiar, explorar, negociar, e c.\lingui{DEFIRS}
\artiklo{mispelo}%
\vorto{mispelo}\fako{bot.} Frukto qua havas plura kerni, qua manjesas erste kande lu oldeskis e komencas velkar\latina{mespilus}\lingui{DISL}
\artiklo{mistelo}%
\vorto{mistelo}\fako{bot.} Planto lignoza, qua vivas parazite sur ula arbori (querko, pomiero, fraxino.) L. \textit{viscum}\latina{viscum}\lingui{DE}
\artiklo{misterio}%
\vorto{misterio}\senco{1}{Rito sekreta di la politeismo antiqua, a qua admisesis nur poka iniciiti. – To quon onu konservas, mantenas, sekreta}\senco{2}{Dogmato kristana, revelita kom objekto di kredo, e qua preter-iras raciono}\lingui{DEFIRS}
\artiklo{mistifikar}%
\vorto{mistifikar}\fako{trans.} Trompar, erorigar (ulu), amuzante su de lua kredo\lingui{DEFIR}
\artiklo{mistiko}%
\vorto{mistiko} Doktrino filozofiala, religiala, qua konsideras kom la perfekteso sorto di kontemplo ed extazo qua elevas la homo, ja en la vivo sur-tera, ad uniono misterioza kun deo\lingui{DEFIRS}
\artiklo{miteno}%
\vorto{miteno} Ganto qua kovras la manuo, tote o duime, kun fingruyo nur por la polexo\lingui{EFS}
\artiklo{mito}%
\vorto{mito} Rakonto tradicionita qua atribuas, ad ula eventi o personegi, karaktero supernatura\lingui{DEFIRS}
\artiklo{mitokondrio}%
\vorto{mitokondrio}\fako{biol.} Grani o filamenti, tre mikra, kolorizebla per alizarino e la kristalo violea, qui existas ye quanti plu o min granda en omna celuli, di animali o di vejetanti, dum lia tota vivo\lingui{DEFIRS}
\artiklo{mitologio}%
\vorto{mitologio} Historio fablatra di la dei, di la enti supernatura\lingui{DEFIRS}
\artiklo{mitomo}%
\vorto{mitomo}\fako{biol.} Parto di la citoplasmo, ek filamenti e fibreti, branchigita od anastomozigita, ma ne retikuligita, qua segun-dice konsistas ye protoplasmo vera
\artiklo{mitoso}%
\vorto{mitoso}\fako{biol.} Partigeso nemediata di la celulo, en qua la divido di la nukleo eventas ante olta di la korpo celula
\artiklo{mitralio}%
\vorto{mitralio}\senco{1}{Peci de olda fero, kupro, bronzo}\senco{2}{\textit{(olim)} Olda ferajo per qua onu charjis la kanoni}\lingui{EFIS}
\artiklo{mitralioso}%
\vorto{mitralioso}\fako{milit.} Kanono ye pafado rapida ed iterata\lingui{DEFR}
\artiklo{mitro}%
\vorto{mitro}\senco{1}{\textit{(epoki antiqua)} Kapo-vesto Aziana, quan pose *weris la Romanini alta-ranga}\senco{2}{Kapo-vesto di la episkopi, kardinali ed ula abadi kande li oficias solene}\lingui{DEFIRS}
\artiklo{mixamebo}%
\vorto{mixamebo}\fako{biol.} Celulo qua devenas de la jermifo di la sporo che la fungi mixomiceta, qua konsistas ye amaso de protoplasmo, sen membrano celuloza, e ye movi amiboidala
\artiklo{mixar}%
\vorto{mixar}\fako{trans., kun}\senco{1}{Unionar ensemble (kozi diversa) tale ke la diverseso cesas esar videbla}\senco{2}{Enduktar aden medio (ulu od ulo stranjera)}\lingui{EFI}
\artiklo{mizantropo}%
\vorto{mizantropo} Individuo qua odias la genero « homo »\lingui{DEFIS}
\artiklo{mizerar}%
\vorto{mizerar}\fako{netrans.} Esar tante povra ke (lu) esas kompatinda\lingui{DEFIRS}
\artiklo{mizerikordio}%
\vorto{mizerikordio} Kompato qua pardonas la kulpinto\lingui{FIS}
\artiklo{mnemoniko}%
\vorto{mnemoniko} Arto helpar memoro per procedi qui igas la kozi plu facile retenebla, merkebla, ritrovebla\lingui{DEFIRS}
\artiklo{mobilizar}%
\vorto{mobilizar}\fako{milit.} Igar armeo, od armeo-fragmento, departopronta por militesko\lingui{DEFIRS}
\artiklo{moblo}%
\vorto{moblo} Objekto uzata por garnisar, ornar la interna parto di domo (stuli, tabli, ed objekti diversa altra-sorta : liti, e c.)\lingui{DFIRS}
\artiklo{modelo}%
\vorto{modelo}\senco{1}{To quo destinesas esor objekto di imito}\senco{2}{Homo qua esas koram artisto (piktisto, statuifisto) en la posturo quan ica volas *reproduktar}\senco{3}{Reprezenturo ek argilo, vaxo, e c. di subjekto destinata a reproduktesor ek marmoro, bronzo, e c.}\senco{4}{To quo destinesas esor objekto di imito, segun quo onu facas od agas ulo : (1) kompozo-maniero; (2) vivo-maniero}\lingui{DEFIRS}
\artiklo{moderar}%
\vorto{moderar}\fako{trans.} Forigar de eceso\lingui{DEFIS}
\artiklo{moderna}%
\vorto{moderna} Qua apartenas a nia epoko od ad epoko kompare recenta (opoze a la epoki olima od antiqua)\lingui{DEFIS}
\artiklo{modesta}%
\vorto{modesta}\senco{1}{For luxego, brilo, pompo}\senco{2}{Qua opinionas sua merito kom nur mezvalora, pasableta}\lingui{EFIS}
\artiklo{modifikar}%
\vorto{modifikar}\fako{trans.} Chanjar (kozo) sen alterar olua naturo esencala\lingui{DEFIRS}
\artiklo{modiliono}%
\vorto{modiliono}\fako{arkitekt.} Sorto di konsolo, qua havas la formo di voluto duopla, salianta, sub la pluv-gutifilo di la kornico\lingui{DEFIS}
\artiklo{modlar}%
\vorto{modlar}\fako{trans.} Fasonar ek argilo, vaxo, e c. (modelo qua exekutesos ye gipso, marmoro, bronzo, e c.). – \textit{(anke aludante la materio quan onu fasonas)}\lingui{DEFIRS}
\artiklo{modo}%
\vorto{modo}\senco{1}{Maniero vidar, agar, partikulare di persono, di lando, di regiono}\senco{2}{Kustumo efemera qua normizas (segun la gusto di la periodo) la maniero vivar, vestizar su, e c.}\senco{3}{Chapeli e kapo-vesti por la homini}\senco{4}{\textit{(filoz.)} Eso-maniero variiva di la substanco finita}\senco{5}{\textit{(logiko)} Kondiciono specala a qua subordinesas silogismo. Dispozeso-maniero di la tri propozicioni di silogismo segun lia quanteso o qualeso}\senco{6}{\textit{(yuro-cienco)} Klauzo qua subordinas la efiko da ago ad evento necerta qua dependas de la volo di ta qua profitos de olu}\senco{7}{\textit{(muziko)} Singla de la toni, che la muziko antiqua, quin distingas la diapazono plu o min alta e la plaso di la mi-toni en la gamo}\senco{8}{Singla de la toni di la plan-kanto quin distingas (ultre la diapazono plu o min alta e la plaso di la mi-toni en la gamo) la divideso, en singla gamo, di la oktavo aden quarto plasizita altre}\senco{9}{\textit{(gram.)} Formo di la verbo, qua indikas la manieri diversa ye qui onu afirmas la ago e la stando quin expresas la verbo (exter la cirkonstanci di tempo e di personi quin indikas altra formi)}\lingui{DEFIRS}
\artiklo{modulacar}%
\vorto{modulacar}\fako{trans.}\senco{1}{Reproduktar per flexi diversa di la voco}\senco{2}{\textit{(metaf.)} Expresar per soni poezioza}\lingui{DEFIRS}
\artiklo{modulo}%
\vorto{modulo}\senco{1}{\textit{(arkitekt.)} Mi-diametro di la infra parto di la kolono-fusto – unajo konvencionita per qua onu determinas la proporciono di la parti en ordino arkitekturala}\senco{2}{\textit{(matem.)} Nombro per qua onu multiplikas la logaritmo di sistemo por obtenar la logaritmi korespondanta di altra sistemo}\senco{3}{\textit{(fiziko)} Koeficiento di elastikeso, o pezo apta plulongigar per un unajo la longeso di prismo di qua la seciono normala esas la surfaco-unajo}\lingui{DEFIRS}
\artiklo{mohero}%
\vorto{mohero} Stofo ek pili di la Angora-kapro\lingui{DEF}
\artiklo{mokar}%
\vorto{mokar}\fako{trans., pro, pri} Igar (ulu, ulo) la objekto di desprizo-rido\lingui{DEFR}
\artiklo{mokasino}%
\vorto{mokasino} Pedo-vesto di la Amer-indiani, ek felo ne-tanagita, ma nur mejisita\lingui{DEFRS}
\artiklo{mola}%
\vorto{mola}\senco{1}{Qua cedas facile ye preso}\senco{2}{\textit{(metaf.)} Qua cedas facile (aludante persono)}\lingui{FI}
\artiklo{molara}%
\vorto{molara}\fako{anat.} Dento qua uzesas por triturar la manjaji\lingui{DEFIS}
\artiklo{moldar}%
\vorto{moldar}\fako{netrans.} Kovreskar ye vejetanti kriptogama qui alteras la substanco\lingui{DE}
\artiklo{molekulo}%
\vorto{molekulo}\fako{kemio} Parto di korpo, konceptata kom ne-dividebla da la agenti fizikala, ma deskompozebla kemiale aden partikuli plu mikra, la « atomi »\lingui{DEFIRS}
\artiklo{moleskino}%
\vorto{moleskino} Telo vernisizita, qua imitas ledro\lingui{DEFI}
\artiklo{molestar}%
\vorto{molestar}\fako{trans.} Tormentar per produktar desagreablaji\lingui{DEFIS}
\artiklo{molibdo}%
\vorto{molibdo}\fako{kemio} Korpo simpla, metala, blanka matide, maleebla preske ne-gisebla\simbolo{Mo}\lingui{DEFIRS}
\artiklo{molio}%
\vorto{molio} Masonuro, konstruktita ye la enireyo di portuo, ye la kapo di warfo, di digo, por ruptar la impetuo di la ondegi e shirmar la navi\lingui{DEFIRS}
\artiklo{moloso}%
\vorto{moloso}\senco{1}{Hundo ek la lando di la Molosi (Epiro, Grekia) quan onu uzis, en la epoki antiqua, por chasar o gardar la bestio-trupi. – Nun : grosa gardo-hundo, ek la raso « dogi »}\senco{2}{\textit{(prozodio)} Verso-pedo Greka o Latina, ek trisilabi longa}\lingui{DEFIRS}
\artiklo{moltono}%
\vorto{moltono} Stofo dika, kun pili longa e dolca, qui igas lu delikata e mola\lingui{DEFIS}
\artiklo{molusko}%
\vorto{molusko}\fako{zool.} Nomo di animali sen-vertebra, ye qui konsistas la brancho triesma di la animal-regno\lingui{DEFIRS}
\artiklo{momento}%
\vorto{momento}\fako{mekan.} Produto di forco per disto\lingui{DEFIS}
\artiklo{monado}%
\vorto{monado}\fako{filoz.}\senco{1}{Segun la doktrino da Pitagoro : unajo perfekta, principo nevidebla di la kozi spiritala e materiala}\senco{2}{Segun Leibnitz : substanco simpla, nedividebla, esencale agema}\senco{3}{\textit{(zool.)} Protozoo kun organo lokomocala flogilatra, ek la klaso « monatidi », karakterizata da un o du mikra flogili acesora, ye la bazo di la flogilo precipua}\lingui{DEFIS}
\artiklo{monadologio}%
\vorto{monadologio}\fako{filoz.} Teorio pri la monadi
\artiklo{monako}%
\vorto{monako}\fako{religio katol.} Monakulo enklostrigita qua vivas for la socio ordinara\lingui{DEFIRS}
\artiklo{monarkio}%
\vorto{monarkio} Stato guvernata da chefo unika\lingui{DEFIRS}
\artiklo{monarko}%
\vorto{monarko} Suvereno di stato qua guvernesas da chefo unika\lingui{DEFIRS}
\artiklo{monastero}%
\vorto{monastero}\fako{biol.} Komenco-fazo di la mitoso (prefazo) dum qua la centrosomo cirkondas su per filamenti protoplasma akromata, astroforma
\artiklo{monato}%
\vorto{monato} Tempo-quanto adoptita kom un ek la dekedu dividuri di yaro, e qua kontenas lore 30, lore 31 dii (ecepte la monato februaro, qua kontenas 28 dii dum la yari ordinara, e 29 dum la yari bisextila)\lingui{DE}
\artiklo{mondo}%
\vorto{mondo} Ensemblo ek lo existanta. Ensemblo ek Suno qua lumizas ni, Tero quan ni habitas, e la planeti cetera kun lia sateliti e kometi\lingui{EFIS}
\artiklo{monero}%
\vorto{monero}\fako{biol.} Grupo ek protozoi, kreita da Haeckel, qua igis ek li klaso aparta, karakterizata per (precipue) lia tota indijo di nukleo; li esus do celuli nekompleta\lingui{DEFIS}
\artiklo{moneto}%
\vorto{moneto} Peco de metalo qua uzesas lor kambii, kun estampuro legala qua atestas lua pezo e titruro e determinas lua valoro nominala\lingui{EFIRS}
\artiklo{monismo}%
\vorto{monismo}\fako{filoz.}\senco{1}{Sistemo filozofiala qua explikas universo per la existo di tipo unika di realajo}\senco{2}{Sistemo da Haeckel qua expresas universo per (unike) la efiko di la forci fizikala e kemiala}\lingui{DEFIRS}
\artiklo{monisto}%
\vorto{monisto} Adepto di monismo\lingui{DEFIRS}
\artiklo{monitoro}%
\vorto{monitoro} Milito-navo, di tun-konteno mezvalora, basa sur la aquo-nivelo, kurasizita tre forte, e provizita ye kanoni gros-modela\lingui{DEFIRS}
\artiklo{mono-}%
\vorto{mono-} Prefixo ciencala = un
\artiklo{monodroma}%
\vorto{monodroma}\fako{matem.} Funciono monodroma, interne di areo donita. Funciono \textit{F(z}) esas monodroma interne di areo S, kande, la punto z variante interne di la areo, la funciono \textit{F(z}) retroprenas sempre la sama valoro en sama punto\lingui{EFIS}
\artiklo{monogama}%
\vorto{monogama}\senco{1}{Qua esas mariajita kun nur un spozo}\senco{2}{Qua spozeskis nur unfoye}\lingui{DEFIRS}
\artiklo{monogena}%
\vorto{monogena}\fako{biol.} Qua agas ta sorto di genito en qua la enti vivanta su *reproduktas nemediate, e kun fazi inter-identa di su-developo, per ovo o per ovuli, opoze a la digeneso (genita alternanta)\lingui{DEFIS}
\artiklo{monografio}%
\vorto{monografio} Studiuro pri punto specala di historio, cienco, e c.\lingui{DEFIRS}
\artiklo{monogramo}%
\vorto{monogramo}\fako{skribarto}\senco{1}{Kombinuro ek plura literi, unionita tale ke un streko, un slingo, uzesas por du o tri literi interdiferanta}\senco{2}{Kombinuro de literi quin ula artisti stampas sur sua verki – ula kolekteri, sur la peci di lia kolektaji}\lingui{DEFIRS}
\artiklo{monoika}%
\vorto{monoika}\fako{bot.} Qua portas, sur un stipo, floruli e florini
\artiklo{monoklina}%
\vorto{monoklina}\fako{bot.} En qua la du sexui un-esas en la sama floro
\artiklo{monoklo}%
\vorto{monoklo} Mikra okulvitruno, por un ek la okuli\lingui{DEFIRS}
\artiklo{monokroma}%
\vorto{monokroma} Qua esas un-kolora\lingui{DEFIS}
\artiklo{monolito}%
\vorto{monolito} Monumento (kolono, obelisko, tombo, e c.) ek nur un petro\lingui{DEFIRS}
\artiklo{monologo}%
\vorto{monologo}\senco{1}{Ceno teatrala en qua un persono – qua esas, o qua kredas esar, sola – parolas a su ipsa}\senco{2}{Diskurso da persono qua ne lasas la kun-esanti parolar liafoye}\lingui{DEFIRS}
\artiklo{monomio}%
\vorto{monomio}\fako{algebro} Expresuro inter la parti di qua ne existas signo di adiciono o di sustraciono\lingui{DEFIS}
\artiklo{monopolo}%
\vorto{monopolo}\senco{1}{Privilejo exkluziva pri vendar ulo}\senco{2}{Yuro exkluziva, posedata da mikra quanto de homi}\lingui{DEFIRS}
\artiklo{monosperma}%
\vorto{monosperma}\fako{bot.} Di qua la frukto kontenas nur un grano (semino)\lingui{DEFIS}
\artiklo{monoteismo}%
\vorto{monoteismo} Kredo ye deo unika\lingui{DEFIRS}
\artiklo{monotona}%
\vorto{monotona}\senco{1}{Qua esas sempre sama-tona. Qua fatigas pro uzar konstante la sama tono}\senco{2}{Qua fatigas pro la repeto di la sama kozi, pro la itero di la sama ago}\lingui{DEFIRS}
\artiklo{monsinioro}%
\vorto{monsinioro} Titulo honorizanta donata a la princi di familio suverenala, a la kardinali, arkiepiskopi, episkopi e prelati\lingui{F}
\artiklo{monstro}%
\vorto{monstro}\senco{1}{Ento di qua la formo esas kontre-natura. (mitol.) Ento maligna, kun formo stranja}\senco{2}{Ento granda extraordinare}\senco{3}{\textit{(metaf.)} Persono qua havas defekto o vicio developita extraordinare}\lingui{DEFIRS}
\artiklo{montanto}%
\vorto{montanto} Peco de ligno, de petro, de fero, pozita vertikale (ex. : la montanti di skalo). – HIS
\artiklo{monto}%
\vorto{monto}\senco{1}{\textit{(geogr.)} Granda altajo ek tereno, qua rezultas de subleveso di sulo, ed havas sur sua flanki, ye altesogradi diversa, zoni de vejetanti interdiferanta}\senco{2}{Amasego}\senco{3}{\textit{(metaf.)} Granda desfacilajo vinkenda}\lingui{EFIS}
\artiklo{montrar}%
\vorto{montrar}\fako{trans., ad}\senco{1}{Vidigar}\senco{2}{\textit{(metaf.)} Lasar aparar}\lingui{FIS}
\artiklo{monumento}%
\vorto{monumento}\senco{1}{Masonuro erektita por perpetuigar la memoro di ulu, di ulo, memorinda}\senco{2}{Edifico remarkinda (palaco, kirko, kolono, statuo, e c.)}\lingui{DEFIRS}
\artiklo{mopso}%
\vorto{mopso}\fako{zool.} Mikra hundo rez-pila, kun muzelo nigra, nazo aplastita\lingui{DRS}
\artiklo{moqueto}%
\vorto{moqueto} Stofo por tapisi e mobli, veluratrajo ek lano di qua la wefto e la warpo esas ek lino-filo\lingui{DEFIS}
\artiklo{morbilo}%
\vorto{morbilo}\fako{patol.} Morbo eruptala, febrala, kontagiala, qua afektas precipue la infanti e pueri\lingui{IL}
\artiklo{morbo}%
\vorto{morbo} Speco medikala di ta o ca kelka altereso di la organi, o kelka perturbeso di la funcioni organismala\lingui{EFISL}
\artiklo{mordar}%
\vorto{mordar}\fako{trans.}\senco{1}{Entamar per sua denti}\senco{2}{Entamar per frotado, exkavo. Entamar per korodo}\senco{3}{\textit{(metaf.)} Atakar (ulo) ye maniero acerba}\lingui{EFIS}
\artiklo{morelo}%
\vorto{morelo}\fako{bot.} Champiniono manjebla, di qua la « chapelo » esas remarkinda pro la nekontebleso di la mikra konkavaji qui igas lu aspektar quale sponjo\latina{Morchella esculenta}\lingui{DEF}
\artiklo{moreno}%
\vorto{moreno}\fako{geogr.} Bendo de gravii, de fragmenti ek roko, qui glitis ed akumulesis an la pedo di glaciero od an lua flanki\lingui{DEF}
\artiklo{morfino}%
\vorto{morfino}\fako{kemio} Alkaloido di opiumo, C₁₇H₁₉NO₃ + H₂O; narkotigivo tre efikiva, uzata medicinale por mingravigar la sufri; drogo danjeroza pro kustumesko\lingui{DEFIRS}
\artiklo{morfogena}%
\vorto{morfogena}\fako{biol.} Dicesas pri la efiki e la agenti qui influas la formo e la strukturo di la organismi\lingui{DEFIS}
\artiklo{morfogenezo}%
\vorto{morfogenezo}\fako{biol.} Cienco di la kauzi qui determinas la formo e la strukturo di la organismi, e di la *leyi (legi) qui rezumas li\lingui{DEFIS}
\artiklo{morfologio}%
\vorto{morfologio}\senco{1}{\textit{(filoz.)} Cienco di la formi diversa di materio}\senco{2}{\textit{(gram.)} Cienco di la formi diversa gramatikala di la vorti}\lingui{DEFIRS}
\artiklo{morfoplasmo}%
\vorto{morfoplasmo}\fako{biol.} Un ek la substanci protoplasma di la celulo-korpo di la ovulo, qua efikas a la karakterizivi di la genitori
\artiklo{morfoso}%
\vorto{morfoso}\fako{biol.} Reflexo elementala, ne-nervala, qua esas respondo rapida ad ecito simpla di la substanco vivanta, e qua su manifestas per ula formo od equilibro-posturo di ta substanco vivanta en la medio cirkondanta\lingui{DEFIS}
\artiklo{morganatika}%
\vorto{morganatika}\fako{yuro-cienco} \textit{(pariajo)} quan princo o sinioro paktas kun persono di rango inferiora e qua ne havas omna efekti civila\lingui{DEFIRS}
\artiklo{morge}%
\vorto{morge} En la dio qua sequas nemediate olta pri qua onu parolas\lingui{DE}
\artiklo{morilio}%
\vorto{morilio} Fero-peco, di qua un extremajo kapo-forma esas ornita, e quan onu pozas ye singla latero di la herdo di kameno, por subtenar la kombusto-ligno\lingui{S}
\artiklo{moro}%
\vorto{moro} Un ek la kustumi di individuo, di populo, pri la praktiko di lo bona e di lo mala etikale\lingui{FL}
\artiklo{morso}%
\vorto{morso} Levero muntita en la boko di la kavalo, e qua, per efikar a la stangi, uzesas por direktar lu\lingui{FI}
\artiklo{mortadelo}%
\vorto{mortadelo} Grosa socisego ek hepati e sen-grasa karni (pistita)\lingui{DFIS}
\artiklo{mortar}%
\vorto{mortar}\fako{netrans.} Cesar vivar. \textit{(anke metaf.)}\lingui{DEFIS}
\artiklo{mortero}%
\vorto{mortero} Sablo o cimento mixita a kalko, e dilutita en aquo, quan onu uzas por kunjuntar la petri e quadri di konstrukturo, di edifico\lingui{DEFIS}
\artiklo{mortezo}%
\vorto{mortezo}\fako{teknol.} Entaliuro en peco de ligno o metalo, por recevar tenono facita en altra peco, tale ke formacesas asembluro\lingui{EFS}
\artiklo{mortifikar}%
\vorto{mortifikar}\fako{trans.} Igar (ulu) mortar ye su ipsa, per domtar sua impulsi karnala, e per humiligar la spirito per asketalaji\lingui{EFIRS}
\artiklo{moruo}%
\vorto{moruo}\fako{zool.} Marfisho ek la genero « gadi », de qui un speco, qua tre abundas en Nov-Lando, esas la objekto di komerco tre granda\latina{gadus morrhus}\lingui{FL}
\artiklo{moruso}%
\vorto{moruso}\fako{bot.} Frukto monoika, pulpoza, di arboro kun folii alternanta\latina{morus}\lingui{SL}
\artiklo{morvo}%
\vorto{morvo}\fako{patol.} Morbo kontagiala di la kavali, maxim-multa-kaze mortiganta, karakterizata da la inflameso e la ekfluado di la membrano pituitala\lingui{FS}
\artiklo{moskardino}%
\vorto{moskardino}\fako{zool.} Speco di gliro, di qua la odoro forta memorigas musko\latina{muscardinus avellanarius}
\artiklo{moskeo}%
\vorto{moskeo}\fako{religio} Templo di la Mohamedisti\lingui{DEFIRS}
\artiklo{moskito}%
\vorto{moskito}\fako{zool.} Insekto diptera, di qua la pikuro esas doloriganta\lingui{DEFRS}
\artiklo{mosko}%
\vorto{mosko} Substanco odoranta quan kontenas posho renkontrebla sur la ventro di la maskulo di sorto di kapreol-yuno\lingui{DEFIR}
\artiklo{moso}%
\vorto{moso}\fako{navig.} Adolecanto qua, sur navo, lernas la mestiero di maristo (*kruano)\lingui{FI}
\artiklo{mosto}%
\vorto{mosto} Suko di vitbero, pronta subisar la fermentaco alkoholifera\lingui{DEFIS}
\artiklo{motacilo}%
\vorto{motacilo}\fako{zool.} Mikra ucelo qua flugetas cirkum la bestio-trupi, ed apud la aqui fluanta o dormanta\latina{motacilla alba}\lingui{L}
\artiklo{moteto}%
\vorto{moteto}\senco{1}{\textit{(olim)} Kirko-kantajo, mi-sakra e mi-profana, plura-parta, tolerita apud la plan-kanto}\senco{2}{Muzikajo destinita a la kirko, e kompozita ek vorti Latina quin onu prenis exter la oficio, quale psalmo, himno, antifono, e c.}\lingui{DEFIRS}
\artiklo{motivo}%
\vorto{motivo}\senco{1}{La « pro quo » o « por quo » di ago o di faco}\senco{2}{En arto-verko : la temo developenda}\lingui{DEFIRS}
\artiklo{moto}%
\vorto{moto} Sentenco, proverbo preferata da ulu, od atribuita ad ulu, kom expresanta plu o min adequate lua maniero pensar, sentar\lingui{DEIR}
\artiklo{motoro}%
\vorto{motoro} Omna sistemo materiala destinita a transformar irga energio aden energio mekanikala\lingui{DEFIRS}
\artiklo{movar}%
\vorto{movar}\fako{trans.}\senco{1}{Igar funcionar. \textit{(anke metaf.)}}\senco{2}{\textit{(pri korpo o parto di korpo)} Diplasar su}\lingui{EFIS}
\artiklo{movimento}%
\vorto{movimento}\fako{muziko} Rapideso-grado ye qua frazo muzikala pleesas o kantesas\lingui{EFI}
\artiklo{moyeno}%
\vorto{moyeno} To quo uzesas por atingar la skopo\lingui{F}
\artiklo{mozaiko}%
\vorto{mozaiko} Laboruro ek peci inter-adjustita ek stono, petro, emalio, glaso, ligno, diversa-kolora, di qui la asembluro formacas desegnuri od ornamenti\lingui{DEFIRS}
\artiklo{muar}%
\vorto{muar}\fako{netrans.}\senco{1}{\textit{(aludante ula speci di animali)} Subisar, ye ula epoki determinita di lia vivo-duro, la falo o la rinovigo di epidermo, pili, plumi, korni, e c.}\senco{2}{\textit{(aludante la voco homala)} Subisar ula altero lor puberesko}\lingui{EFIS}
\artiklo{muaro}%
\vorto{muaro} Stofo di qua la reflekti esas kolor-chanjanta qui produktesas da la aplasto di la grani sub la cilindri interrotacanta\lingui{DEFIRS}
\artiklo{mucilajo}%
\vorto{mucilajo}\senco{1}{Substanco viskoza quan kontenas ula vejetanti}\senco{2}{\textit{(farmacio)} Liquido viskoza, ek la solvo di gumo en aquo}\lingui{EFIS}
\artiklo{muelar}%
\vorto{muelar}\fako{trans.} Triturar (grani) per pasigar inter grosa diski qui rotacas, ordinare ek petro harda – la supra frotante kontre la infra\lingui{FIRS}
\artiklo{muevo}%
\vorto{muevo}\fako{zool.} Mar-ucelo, palmipeda, long-ala\latina{larus}\lingui{DF}
\artiklo{muflo}%
\vorto{muflo}\senco{1}{\textit{(teknol.)} Tera vazo quan la kemiisti uzas por submisar korpo a la efiko da fairo sen kontakto kun la flamo}\senco{2}{Mikra forno quan onu pozas transverse en forno plu granda, e qua recevas la materii destinita a kupelago}\lingui{DEFIS}
\artiklo{muflono}%
\vorto{muflono}\fako{zool.} Sorto di mutono sovaja, ne-amansita (en Korsika, Sardinia, Grekia)\latina{ovis musimo}\lingui{DEFIS}
\artiklo{mufo}%
\vorto{mufo}\senco{1}{Furajo, gaino vatizita, apertita ye la du extremaji, en qua onu pozas sua manui por shirmar olci de la atmosfero kolda}\senco{2}{\textit{(teknol.)} Cilindro ek gisfero aden qua onu ajustas la extremaji di du tubi, di du ax-arbori, quin onu volas interjuntar}\lingui{DER}
\artiklo{mujar}%
\vorto{mujar}\fako{netrans.) (aludante vento, ondi} Produktar bruiso analoga a la klamo matida, nesonora e longa di la animali bova\lingui{FIS}
\artiklo{muko}%
\vorto{muko}\fako{fiziol.} Liquido viskoza quan sekrecas ula membrani\lingui{EFIS}
\artiklo{mulato}%
\vorto{mulato} Individuo qua naskis de la koito da pelblanko kun pelnigro\lingui{DEFIRS}
\artiklo{muldar}%
\vorto{muldar}\fako{trans.} Reproduktar modelo reliefa, aplikante sur lu materio mola, flexebla, qua prenas lua konturi. – \textit{(anke metaf.)}\lingui{EFIS}
\artiklo{muliero}%
\vorto{muliero} Homino adulta\lingui{FL}
\artiklo{mulo}%
\vorto{mulo}\fako{zool.} Produkturo maskula ek la kopulaco da asnulo kun kavalino, o da kavalulo kun asnino\lingui{EFIRS}
\artiklo{multa}%
\vorto{multa} Di qua la quanto o grandeso, se expresita per nombro, korespondas a nombro relative granda\lingui{EFIS}
\artiklo{multiforma}%
\vorto{multiforma}\fako{matem.}\lingui{DEFIS}\subvorto{}{funciono multiforma}{Qua ne esas monodroma, t.e. qua prenas valori diferanta pri la sama sistemo de valori di sua variebli}{}{}
\artiklo{multiplikar}%
\vorto{multiplikar}\fako{trans.) (aritm.} Iterar (un quanto fixa) n-foye\lingui{DEFIS}
\artiklo{muluro}%
\vorto{muluro}\fako{arkitekt.} Ornamento prolongebla, uniforma, salianta, en produkturo di arkitekturo o di menuzo
\artiklo{mumio}%
\vorto{mumio} Kadavro balzamizita e envelopita per bendeti, quan onu renkontras en la sepulturi di la Egiptiani antiqua\lingui{DEFIRS}
\artiklo{mungar}%
\vorto{mungar}\fako{trans.} Desembarasar de la humoro viskoza quan sekrecas la mukozo nazala\lingui{L}
\artiklo{municiono}%
\vorto{municiono}\fako{teknol. militala} Singla ek la kozi ye qui konsistas la ensemblo ek la moyeni di defenso, di su-nutro, ye qui onu provizas fortifikajo, fortreso, militisto-trupo. Anke : pulvero, kartochi, obusi, bombi, atomarmi\lingui{DEFIRS}
\artiklo{muntar}%
\vorto{muntar}\fako{trans.) (tekn.} Erektar (mashino, karpenturo, e c.) per konjuntar la parti\lingui{DEFIRS}
\artiklo{mureno}%
\vorto{mureno}\fako{zool.} Speco di kongro\latina{muraena helena}\lingui{DEFIRS}
\artiklo{murmurar}%
\vorto{murmurar}\fako{netrans.}\senco{1}{Exhalar quaza lamento matida, ne-sonora}\senco{2}{Exhalar bruisado matida, nesonora}\lingui{DEFIS}
\artiklo{muro}%
\vorto{muro} Panelo de masonuro, plu o min alta, uzata por cirkum-klozar o subtenar\lingui{EFIS}
\artiklo{musareno}%
\vorto{musareno}\fako{zool.} Mikra mamifero, karnavida, insektivora, kun odoro moskoza\latina{sorex}\lingui{FI}
\artiklo{muserono}%
\vorto{muserono}\fako{bot.} Speco di agariko, fungo manjebla
\artiklo{musho}%
\vorto{musho}\fako{zool.} Insekto diptera, ek la familio « muskidi »\latina{musca}\lingui{FIRS}
\artiklo{muskado}%
\vorto{muskado}\fako{bot.} Nuco aromatoza, quan produktas arboro exotika, uzata kom spico\latina{myristica fragrans}\lingui{DFIRS}
\artiklo{muskato}%
\vorto{muskato}\fako{bot.} Vitbero di qua la odoro similesas olta di musko\lingui{DEFIRS}
\artiklo{musketo}%
\vorto{musketo}\fako{olim} Pafarmo quan onu funcionigis per mecho acendita e qua remplasesis da fusilo\lingui{DEFIRS}
\artiklo{musketono}%
\vorto{musketono}\fako{olim}\senco{1}{Pafarmo plu grosa, ma min longa, kam musketo}\senco{2}{Speco di fusilo kurta, inter karabino e pistolo}
\artiklo{musko}%
\vorto{musko}\fako{bot.} Planto kriptogama celuloza, di qua la folieti tenua multigas su tale ke li tapetizas la loki ube lu kreskas\latina{muscua}\lingui{FI}
\artiklo{muskulo}%
\vorto{muskulo}\fako{anat.} Organo karnoza, ek fibri kontraktebla, qui produktas la movi (che la animali)\lingui{DEFIRS}
\artiklo{muslino}%
\vorto{muslino} Stofo de kotono tenua, e di qua la texuro esas diafana\lingui{DEFIRS}
\artiklo{muslo}%
\vorto{muslo}\fako{zool.} Molusko bivalva, manjebla, tipo di la tribuo « mitilinei »\latina{mytilus edulis}\lingui{DEF}
\artiklo{muso}%
\vorto{muso}\fako{zool.} Mikra mamifero, rodera, ek la genero « rato » kun pilaro nigra o griza, qua habitas la trui di la domi\latina{mus}\lingui{DERL}
\artiklo{musono}%
\vorto{musono}\fako{meteorol.} Vento qua aparas periodale en la oceano di India, de esto a westo, en la semestro unesma di la yaro, ed inversa-sense en la semestro duesma\lingui{DEFIRS}
\artiklo{mustar}%
\vorto{mustar}\fako{trans.} Koaktesar ne-eskapeble, ne-eviteble\lingui{DE}
\artiklo{mustardo}%
\vorto{mustardo} Kondimento ek grani di sinapo, triturita kun vino-mosto o vinagro\lingui{EFIS}
\artiklo{muta}%
\vorto{muta} Qua ne povas (fiziologiale) parolar\lingui{EFIS}
\artiklo{mutaco}%
\vorto{mutaco}\fako{biol.} Vario bruska e spontana qua eventas en decendanto di individuo normala, e qua, de lore, esas fixa e heredonta\lingui{DEFIS}
\artiklo{mutilar}%
\vorto{mutilar}\fako{trans.}\senco{1}{Alterar (la korpo) ye sua integreso per deprenar de lu membro od altra parto}\senco{2}{\textit{(metaf.)} Alterar (ensemblo) ye lua integreso, per deprenar parto esencala}\lingui{EFIS}
\artiklo{mutono}%
\vorto{mutono}\fako{zool.}\senco{1}{Mamifero ek la familio « rumineri », kun chanfreno konvexa, korni kava e permananta, pili lana, frizita}\senco{2}{\textit{(meteorol.)} Mikra nubi, ek ondi blanka, de vaporo, simila a lano-floki}\senco{3}{\textit{(relig.katol.)} Persono konfidita a la direkto de pastoro spiritala}\latina{ovis aries}\lingui{EFI}
\artiklo{mutuala}%
\vorto{mutuala} Dicesas pri societo di qua la membri intersukursas kontre ula riski – interhelpas\lingui{DEFIS}
\artiklo{*muvmento}%
\vorto{*muvmento} Ensemblo de la agi ed esforci, koordinata, da grupo de personi qui volas atingar skopo partikulara\lingui{EF}
\artiklo{muzelo}%
\vorto{muzelo}\senco{1}{Parto salianta, oblonga di la vizajo di ula mamiferi, di ula fishi}\senco{2}{\textit{(metaf.)} Grimaso per oblongigar la labii kom signo di deskontenteso o di moko}\lingui{EFI}
\artiklo{muzeo}%
\vorto{muzeo}\senco{1}{\textit{(epoki antiqua)} Establisuro ube kultivesis literaturo e la cienci}\senco{2}{Loko en qua asemblesis kolektaji de objekti di arto, cienco, industrio, e c.}
\artiklo{muzikar}%
\vorto{muzikar}\fako{netrans.} Ecitar la sentimenti per serii melodia o per kombinuri harmonioza de soni segun-mezuro e ritmoza\lingui{DEFIRS}
\artiklo{muzo}%
\vorto{muzo}\senco{1}{\textit{(mitol.)} Singla ek la non deini interfrata, qui prezidis la non arti liberala}\senco{2}{\textit{(metaf.)} Inspirofonto di la poeto}\lingui{DEFIRS}
\artiklo{myelino}%
\vorto{myelino}\fako{anat.} Substanco mola, refraktiva e komplexa, qua cirkondas la axono o cilindraxo, e kontenesas en la Schwann-gaino\lingui{DEF}
\artiklo{myelito}%
\vorto{myelito}\fako{patol.} Inflamuro di la myelo\lingui{DEFIS}
\artiklo{myelo}%
\vorto{myelo}\fako{anat.} Grosa kordono nerva, prolonguro di la cerebro, qua kontenesas en la kanalo vertebrala\lingui{GR}
\end{multicols}
\sekciono{N}
\begin{multicols}{2}
\artiklo{-n}%
\vorto{-n}\fako{gram.} Signo qua indikas ke o la komplemento direta o la atributo esas inversigita en la propoziciono – e, per lo, preventas kompren-erori che la lektanto o la askoltanto
\artiklo{nabeto}%
\vorto{nabeto}\fako{bot.} Varietato di napo, di qua la grani esas oleoza\latina{brassica rapa oleifera}\lingui{FS}
\artiklo{nabo}%
\vorto{nabo} Centro di roto, adube konvergas la radii e quan trairas la rotaxo\lingui{DE}
\artiklo{nabobo}%
\vorto{nabobo}\senco{1}{Princo mohamedista en India}\senco{2}{Personego qua retrovenis de India (o de altra landi exotika) kun havajo grandega}\lingui{DEFIRS}
\artiklo{naciono}%
\vorto{naciono}\senco{1}{Asemblitaro de homi ye qua konsistas socio politikala, qua esas guvernata, administrata o jerata da institucuri komuna}\senco{2}{Singla ek la dividuri en qui repartisesis olim la studenti di la universitato di la urbo Paris}\lingui{DEFIRS}
\artiklo{nadiro}%
\vorto{nadiro}\fako{astron.} Punto di la cielo-sfero ube abutas vertikalo qua departas de la loko ye qua onu esas, e pasante tra la centro di la Terglobo\lingui{DEFIRS}
\artiklo{naftalino}%
\vorto{naftalino}\fako{kemio} Hidrokarbido quan produktas la distilo di minkarbono, di gudro\lingui{DEFIRS}
\artiklo{nafto}%
\vorto{nafto}\fako{kemio} Hidrokarbido, sorto di bitumo liquida, volatila ed inflamebla\lingui{DEFIRS}
\artiklo{naiva}%
\vorto{naiva} Qua esas tam kredera kam mento sen-experienca\lingui{DEFR}
\artiklo{nakarato}%
\vorto{nakarato} Koloro reda oranjea\lingui{DEFS}
\artiklo{naktigalo}%
\vorto{naktigalo}\fako{zool.} Mikra ucelo, ek la klaso « paseri », di qua la plumaro esas griza, famoza pro la beleso di lua kantostilo\latina{sylvia luscinia}\lingui{DE}
\artiklo{nam}%
\vorto{nam} Konjunciono qua unionas, a propoziciono, propoziciono qua sequas olca e qua donas la « pro quo » di to quon afirmas la unesma\lingui{L}
\artiklo{nankino}%
\vorto{nankino} Telo de kotono, di qua la koloro esas uniforma, maxim-multa-kaze flava klare\lingui{DEFIR}
\artiklo{nano}%
\vorto{nano} Persono qua esas mikra extraordinare\lingui{FIS}
\artiklo{napo}%
\vorto{napo}\fako{bot.} Planto krucifera, kun radiko pulpoza e nutriva\latina{brassica napus esculenta}\lingui{FIS}
\artiklo{naracar}%
\vorto{naracar}\fako{trans.} Raportar pri fakto, enumerante detaloze lua cirkonstanci diversa\lingui{EFIS}
\artiklo{narciso}%
\vorto{narciso}\fako{bot.} Planto ek la familio « amarilidei »\latina{narcissus}\lingui{DEFIRS}
\artiklo{nardo}%
\vorto{nardo} Oleo parfumoza, quan, en la epoki antiqua, onu extraktis de planto aromatoza\lingui{DEFIS}
\artiklo{narkotar}%
\vorto{narkotar}\fako{netrans.} Esar en ta stando di stuporo o torporo produktita artificale per la uzo di ula substanci (kloroformo, e c.)\lingui{DEFIS}
\artiklo{narvalo}%
\vorto{narvalo}\fako{zool.} Cetaceo di qua la maxilo havas longa dentego defensala\latina{monodon monoceros}\lingui{DEFIS}
\artiklo{naskar}%
\vorto{naskar}\fako{netrans.} Ekirar la sino di la patrino\lingui{IS}
\artiklo{naso}%
\vorto{naso} Korbo ek oziero, kono-forma, a qua abutas ula reto ube la fishi arivanta kaptesas\lingui{FS}
\artiklo{natar}%
\vorto{natar}\fako{netrans.} Subtenar su ed avancar en aquo, per ula movi di la membri\lingui{EFIS}
\artiklo{natro}%
\vorto{natro}\fako{kemio} Korpo simpla, metala, qua esas elemento esencala di sodo\simbolo{Na}
\artiklo{*naturele}%
\vorto{*naturele}\fako{adverbo} Pro la naturo di la kozi
\artiklo{naturo}%
\vorto{naturo}\senco{1}{Ensemblo ek la kozi e forci qui existas eterne}\senco{2}{La esenco, la atributi qui esas propra ye ula ento}\senco{3}{La stando primitiva (di la homo), opoze a « civilizeso »}\lingui{DEFIRS}
\artiklo{naufrajar}%
\vorto{naufrajar}\fako{netrans.} Dicesas pri navo qua perisas kontre rifi o brizanti, o dum lua navigo\lingui{FIS}
\artiklo{nautilo}%
\vorto{nautilo}\fako{zool.} Molusko cefalopoda, kun konko extera parietoza\latina{nautilus pompilius}\lingui{DEFIS}
\artiklo{nauzear}%
\vorto{nauzear}\fako{netrans.} Tormentesar da la bezono vomar\lingui{EFIS}
\artiklo{navigar}%
\vorto{navigar}\fako{netrans.} Voyajar per irar sur od en aquo, altre kam per nato\lingui{DEFIRS}
\artiklo{navo}%
\vorto{navo}\senco{1}{Konstrukturo destinita a navigar sur od en mari ed oceani}\senco{2}{Parto di kirko, de la portalo a la koreyo, inter du rangi de pilastri, de koloni, qui subtenas la vulto}\lingui{EFIRS}
\artiklo{nayado}%
\vorto{nayado}\fako{mitol.} Deajo infra-ranga, femina, qua prezidis la fonti, riveri, e c.
\artiklo{nazo}%
\vorto{nazo} Organo qua salias meze di la vizajo, rezideyo di olfakto\lingui{DEFIRS}
\artiklo{ne}%
\vorto{ne} Adverbo di nego
\artiklo{nebular}%
\vorto{nebular}\fako{netrans.} Dicesas kande aquo-vaporo nekomplete kondensita flotacas en ta parto di atmosfero qua esas apud tero e trublas olua diafaneso\lingui{DEFIS}
\artiklo{nebuloso}%
\vorto{nebuloso}\fako{astron.} Aglomerajo ek steli qui esas tante multa e semble inter-apuda ke li ne plus esas dicernebla singla\lingui{DEFIRS}
\artiklo{necesa}%
\vorto{necesa}\senco{1}{Quan onu ne povas karear}\senco{2}{De qua onu ne povas dispensar su}\senco{3}{\textit{(filoz.)} Qua ne povas ne esar}\lingui{EFIS}
\artiklo{negar}%
\vorto{negar}\fako{trans.} Deklarar ke ulo ne esas\lingui{DEFIS}
\artiklo{negativa}%
\vorto{negativa}\senco{1}{\textit{(algebro ed elektr.)} Qua havas la signo minus}\senco{2}{Negativo \textit{(fotogr.)} : Ube la nigro di la modelo aparas blanka; e la blanko, nigra}\lingui{DEFIRS}
\artiklo{neglijar}%
\vorto{neglijar}\fako{trans.}\senco{1}{Lasar (ulu, ulo) sen sorgar lu}\senco{2}{Lasar sen egardar lu}\lingui{EFIS}
\artiklo{neglijeo}%
\vorto{neglijeo} Homo-vesto, quan onu *weras (generale) dum la matino\lingui{DEFI}
\artiklo{negociar}%
\vorto{negociar}\fako{trans.} Mediacar por konkluzar afero (mariajo, kontrato, paco)\lingui{EFIRS}
\artiklo{negro}%
\vorto{negro} Homo ek la raso pelnigra\lingui{DEFIS}
\artiklo{nek}%
\vorto{nek} Konjunciono qua separas propozicioni negala, o termini interdiferanta, di propoziciono negala\lingui{FIRS}
\artiklo{nekrobioso}%
\vorto{nekrobioso}\fako{patol.} Modifikeso strukturala di organo o tisuo en qua la sango-cirkulo cesis, generale pro infarkto\lingui{DEFIS}
\artiklo{nekrologo}%
\vorto{nekrologo} Notico pri un o plura personegi jus mortinta\lingui{DEFIRS}
\artiklo{nekromanciar}%
\vorto{nekromanciar}\fako{netrans.} Praktikar la arto (asertita) konjurar la mortinti per magio por obtenar de li la revelo di la eventontaji, di la kozi okulta\lingui{DEFIRS}
\artiklo{nekropolo}%
\vorto{nekropolo}\senco{1}{\textit{(epoki antiqua)} Sorto di tombeyo subtera che la Egiptiani ed altra populi antiqua}\senco{2}{Parto di urbo, qua esas destinita a la sepulturi}\lingui{DEFIS}
\artiklo{nekroso}%
\vorto{nekroso}\fako{patol.} Stando di osto, o di parto de osto, qua (che vivanto) cesis vivar\lingui{DEFIS}
\artiklo{nektario}%
\vorto{nektario}\fako{bot.} Organaro glanda, situita en la floro e qua kontenas suko mielatra
\artiklo{nektaro}%
\vorto{nektaro}\senco{1}{\textit{(mitol.)} Drinkajo delikatega di la dei}\senco{2}{Drinkajo bonega, delikatega, qua ecelas irga ceteri}\senco{3}{Liquido sukroza, quan sekrecas la folio o la floro, e qua atraktas la insekti}\lingui{DEFIRS}
\artiklo{nematoblasto}%
\vorto{nematoblasto}\fako{biol.} Celulo specala, qua karakterizas la celenterei, e qua kontenas organo urtikaganta
\artiklo{nenio}%
\vorto{nenio}\fako{epoki antiqua} Kanto funerala – da la ploristini, lor la sepulto di personegi\lingui{DEFI}
\artiklo{nenufaro}%
\vorto{nenufaro}\fako{bot.} Planto aquala, tipo di la familio « nimfeacei »\latina{Nuphar, Nymphaea}\lingui{EFIS}
\artiklo{neolamarckismo}%
\vorto{neolamarckismo}\fako{biol.} Teorio transformista qua admisas, por explikar la origino di la speci, la influo preponderanta, en la efekti da la uzado, di la efiki mekanismala\lingui{DEFIRS}
\artiklo{neologismo}%
\vorto{neologismo} Uzo di vorti nove-kreita, o di vorti anciena ma uzata nova-senco\lingui{DEFIS}
\artiklo{neono}%
\vorto{neono}\fako{kemio} Elemento qua esas gasa ye la temperaturo ordinara, ed existas en aero minima-quante\simbolo{Ne}\lingui{DEFIRS}
\artiklo{neovitalismo}%
\vorto{neovitalismo}\fako{biol.} Teorio segun qua la faktori fiziko-kemiala qui kondicionas la fenomeni vivala efikas per mekanismo partikulara e qui diferas de olta qua intervenas en la mondo materiala ordinara\lingui{DEFIRS}
\artiklo{neperala}%
\vorto{neperala}\fako{matem.} Dicesas pri la logaritmi qui inventesis da Neper, la matematikisto Skota
\artiklo{nepotismo}%
\vorto{nepotismo}\senco{1}{\textit{(historio eklez.)} Favoro quan juis, lor ula papi, la nevi e parenti di ilci}\senco{2}{Trouzo, da oficiero, di sua povo por promocar neyuste sua familiani o protektati}\lingui{DEFIS}
\artiklo{nepoto}%
\vorto{nepoto} La filio di onua filio\lingui{EFISL}
\artiklo{neptunala}%
\vorto{neptunala}\fako{geol.} Dicesas : 1. pri la tereni qui debas sua origino ad aquo; 2. pri sistemo qua explikis la formaco di omna roki per la efiko da la aqui
\artiklo{nereido}%
\vorto{nereido}\fako{mitol.} Nimfo di maro\lingui{DEFIRS}
\artiklo{nervaturo}%
\vorto{nervaturo}\senco{1}{\textit{(historio naturala)} Singla de la fileti salianta qui quaze branchifas en la limbo di la folio, che ula planti. (duras sur pag.393). – Tubo kornatra qua quaze branchifas en la alo di ula insekti}\senco{2}{\textit{(bind-arto)} Salianto quan formacas, sur la dorso di libro, la kordeti ye qui sutesis la kayeri}\senco{3}{\textit{(arkitekt.)} Muluro salianta di la aristi di vulto – di la kaneli di kolono, e c.}\lingui{IS}
\artiklo{nervo}%
\vorto{nervo}\fako{anat.} Singla de la filamenti qui igas la parti diversa di la korpo homala komunikar kun la cerebro e la myelo, transmisante ad olca la senti; ad olta, la impulsi movala\lingui{DEFIRS}
\artiklo{nesto}%
\vorto{nesto} Mikra korbo ronda, quan la uceli konstruktas ek stipeti de palio, junko, ligno, ed ek moso, lanugo, e c. por pondar en olu sua ovi, kovar li, ekirigar la yuneti ed edukar ici\lingui{DE}
\artiklo{neta}%
\vorto{neta}\senco{1}{Quan nula makulo igas sordida}\senco{2}{\textit{(pri preco)} De qua onu deduktis omno deduktenda : rabato, gelto, diskonto, e c.}\lingui{EFIS}
\artiklo{neurono}%
\vorto{neurono}\fako{histol.} Celulo nervala, ek celulo-korpo nukleoza e prolonguri, arboratraji protoplasmala ed axono
\artiklo{neutra}%
\vorto{neutra}\senco{1}{Qua ne adheras a la partiso di ulu en debato, en milito}\senco{2}{\textit{(kemio)} Dicesas pri korpo ek la kombinuro kemiala di du korpi (acido e bazo)}\lingui{DEFIRS}
\artiklo{neveo}%
\vorto{neveo} Strato de nivo-floki apud glaciero\lingui{EFIS}
\artiklo{nevo}%
\vorto{nevo} La filio di onua frato
\artiklo{nevralgio}%
\vorto{nevralgio}\fako{patol.} Doloro nervala\lingui{DEFIRS}
\artiklo{nevrastenio}%
\vorto{nevrastenio}\fako{patol.} Nevroso di qua la karaktero esencala esas la hipotonio neuro-muskulala e fatigeso\lingui{DEF}
\artiklo{nevrito}%
\vorto{nevrito}\fako{patol.} Lezuro inflamurala o degenero di la nervi\lingui{DEF}
\artiklo{nevrologio}%
\vorto{nevrologio}\fako{biol.} La cienco qua traktas la formaco, la strukturo, la roli e la morbi di la sistemo nervala e, specale, di la nervi\lingui{DEFIS}
\artiklo{nevroptero}%
\vorto{nevroptero}\fako{historio naturala} Insekto di qua la ali havas nervaturo dispozita ret-atre. (kom ex. : libelulo)\lingui{DEFIS}
\artiklo{nevroso}%
\vorto{nevroso}\fako{patol.} Stando maladatra, karakterizata da trubleso nervala\lingui{DEFIRS}
\artiklo{nevuso}%
\vorto{nevuso}\fako{patol.} Misformacuro di pelo, multa-kaze naskala, de origino embrionala, qua havas la formo di makulo o tumoro, ma qua povas eventar pos la nasko e mem che oldi\lingui{DEF}
\artiklo{*nexta}%
\vorto{*nexta} Qua arivos kom unesma (en serio od en spaco)\lingui{DE}
\artiklo{*nexta}%
\vorto{*nexta} Qua venas nemediate – spacale, tempale o rangale
\artiklo{ni}%
\vorto{ni} Ensemblo de la personi quin la parolanto aludas kom esanta en la sama grupo kam lu
\artiklo{nicho}%
\vorto{nicho}\fako{arkitekt.} Kavajo facita en muro por pozar ibe statuo, vazo, furnelo, e c. Mikra reduito en apartamento, ube onu lokizas lito\lingui{DEFIS}
\artiklo{nielo}%
\vorto{nielo}\fako{tekn.} Graburo kava di qua onu plenigas la streki ye sorto di emalio nigra\lingui{DEFIS}
\artiklo{nigelo}%
\vorto{nigelo}\fako{bot.} Morbo di la spiko (imputata a nebulo) qua transformas la grano aden amaso blanka ek multa-mil aguleti. – Morbo di la spiko, qua transformas la grano a polvo nigratra qua konsistas ye kriptogami\lingui{EFISL}
\artiklo{nigra}%
\vorto{nigra} Di qua la koloro, rezultante de la absorbeso di omna lumo-radii, produktas sur la vid-organo la impreso di obskureso kompleta, di tenebro\lingui{EFIS}
\artiklo{nihilo}%
\vorto{nihilo}\senco{1}{Absenteso di eso (ante la existo di naturo)}\senco{2}{Destrukteso absoluta di la ento}\lingui{DEFIRS}
\artiklo{nikelino}%
\vorto{nikelino}\fako{kemio} Aloyuro ek 80/100 kupro e 20/100 nikelo, uzata por la envelopilo di la fusil-kugli
\artiklo{nikelo}%
\vorto{nikelo}\fako{kemio} Korpo simpla, metala, grize-blanka, desfacile fuzebla, forjebla, maleebla, polisebla\simbolo{Ni}\lingui{DE}
\artiklo{nikotino}%
\vorto{nikotino}\fako{kemio} Alkaloido venena quan onu extraktas de tabako\lingui{DEFIRS}
\artiklo{niktalopo}%
\vorto{niktalopo}\fako{patol.} Persono qua dicernas la objekti maximbone en nokto\lingui{DEFIRS}
\artiklo{niktropismo}%
\vorto{niktropismo}\fako{biol.} Fenomeno per qua la organi di ula planti agas movo por haveskar posturo partikulara kande li esas en tenebro
\artiklo{nilgauo}%
\vorto{nilgauo}\fako{zool.} Antilopo di India, alta-statura, kun korni glata e kurvigita adavane\latina{boselaphus tragocamelus}
\artiklo{nimbo}%
\vorto{nimbo}\fako{meteorol.} Pluvo-nubo qua karakterizas la aero-fluo equatorala, sen formo preciza, kun periferio quaze lacerita\lingui{DEFIS}
\artiklo{nimfo}%
\vorto{nimfo}\fako{mitol.} Deajo di la fluvii, boski, monti, o di la eskorto di ula deini (Diana chasera, Tetis, e c.)\lingui{DEFIRS}
\artiklo{nipelo}%
\vorto{nipelo}\fako{tekn.} Junto-tubo kun fileto. ye lua extremaji, e muntebla interne\lingui{DE}
\artiklo{nirvana}%
\vorto{« nirvana »}\fako{filoz. Indiana} Ekiro di la transmigro per supreso di deziro e di nesavo
\artiklo{nitro}%
\vorto{nitro}\fako{kemio} Nitrato naturala di kalio, NO₃K
\artiklo{*niuzo}%
\vorto{*niuzo}\fako{tekn.} \textit{(aludante jurnali, inform-agenterii)} Sorto specala di informi : redakturo qua naracas fakti ne-expektita, di qui la exakteso esas verifikita rigoroze (quale agus lo historiisto), e qua eliminas la detali qui ne interesus la lektonti\lingui{E}
\artiklo{nivar}%
\vorto{nivar}\fako{netrans.} \textit{(aludante la aquo qua konjelesis en la supra regioni di atmosfero)} Falar adsule, forme di floki blanka\lingui{IS}
\artiklo{nivelo}%
\vorto{nivelo} Grado di alteso, kompare a plano horizontala, di lineo o di plano qua esas paralela ad olu\lingui{DFIRS}
\artiklo{niveolo}%
\vorto{niveolo}\fako{bot.} Genero de amarilidei, planto bulboza, di qua la folii esas lineala e la flori blanka\latina{galanthus nivalis}
\artiklo{nixo}%
\vorto{nixo}\fako{mitol., Germana e Skandinava} Feo qua habitas profunde en la aquo amasi\lingui{DE}
\artiklo{no}%
\vorto{no} Interjeciono uzata por enuncar ke kozo ne esas tala. (Se la adverso insistas pri la « eso », lu uzas la interjecioni « si »)\lingui{EFIS}
\artiklo{nobelo}%
\vorto{nobelo} Persono qua apartenas a klaso quan onu egardas kom superiora, e qua honorizesis per privileji e distingivi heredita\lingui{DEFIS}
\artiklo{nobla}%
\vorto{nobla}\senco{1}{Elevita, ek la ceteri, per digneso o merito}\senco{2}{Elevita super to quo esas ordinara}\lingui{DEFIS}
\artiklo{nocar}%
\vorto{nocar}\fako{trans., pri, pro, per} Produktar detrimento (ad ulu) o domajo (ad ulo)\lingui{EFI}
\artiklo{nocho}%
\vorto{nocho}\senco{1}{Entaliuro quan facas panvendisto, karnovendisto, sur bastono por indikar la quanto de pano, de karno, vendita kredite}\senco{2}{Entaliuro sur la fusto di arbalesto o sur la grosa extremajo di flecho, ube esas la kordo tensita di la arbalesto, di la arm-arko}\lingui{E}
\artiklo{nociono}%
\vorto{nociono} Savajo elementala\lingui{DEFIS}
\artiklo{nodo}%
\vorto{nodo} Krucumuro tre klemita qua formacas halto en la kontinueso di filo, di kordo, e c., o qua unionas strete du fili, du kordi, e c.\lingui{FIS}
\artiklo{noktar}%
\vorto{noktar}\fako{netrans.} Dicesas pri la parto di la dio dum qua la suno-radii cesis lumizar la loko ube onu esas\lingui{DEFIRS}
\artiklo{noktiluko}%
\vorto{noktiluko}\fako{zool.} Genero de protozoo flogilatra, ek la klaso « cistoflagelei » – di qua la « flogilumo », tre povoza, atingas un milimetro diametre kande lu esas adulta, fosforecanta\latina{noctiluca}\lingui{EFISL}
\artiklo{noktuo}%
\vorto{noktuo}\fako{zool.} Speco di strigo, rapt-ucelo noktala\latina{noctua}
\artiklo{nokturno}%
\vorto{nokturno}\senco{1}{\textit{(muziko)} Sorti di romanco du-voca, kun la karaktero di serenado}\senco{2}{\textit{(liturgio katol.)} Parto di la nokt-oficio, qua kontenas ula quanto de psalmi}\lingui{DEFIRS}
\artiklo{nomado}%
\vorto{nomado} Qua ne havas rezideyo fixa en teritorio determinita\lingui{DEFIRS}
\artiklo{nomar}%
\vorto{nomar}\fako{trans.} Dicar la nomo qua distingas (persono) ed indikas lu individuale\lingui{DEFIS}
\artiklo{nombro}%
\vorto{nombro}\fako{matem.} Vorto qua indikas quanta-foye la unajo kontenesas en quanto\lingui{EFIS}
\artiklo{nomenklaturo}%
\vorto{nomenklaturo} Ensemblo ek la termini, ek la nomi specala, qui uzesas da ca o ta cienco, arto, libro teknikala\lingui{DEFIS}
\artiklo{nominar}%
\vorto{nominar}\fako{trans.} Grantar titulo, ofico\lingui{EFIS}
\artiklo{nominativo}%
\vorto{nominativo}\fako{gram.}\senco{1}{Deklino-kazo qua indikas, ke la nomo esas la subjekto di la verbo}\senco{2}{La ipsa nomo qua esas la subjekto di la verbo}\lingui{DEFIS}
\artiklo{nomografio}%
\vorto{nomografio}\fako{matem.} Ensemblo ek la metodi qui havas kom skopo remplasar la kalkuli numerala per simpla lekto sur tabelo grafika di qua la nomo esas « abako »
\artiklo{non}%
\vorto{non}\fako{matem.} La cifro « 9 »
\artiklo{nona-}%
\vorto{nona-} Prefixo ciencala = non-…
\artiklo{noniliono}%
\vorto{noniliono}\fako{matem.} 10⁵⁴, od oktiliono × 1.000.000
\artiklo{nopalo}%
\vorto{nopalo}\fako{bot.} Arboro sur qua vivas la kochenilo, varietato di kaktuso\latina{opuntia cocoinillifera}\lingui{DEFIRS}
\artiklo{nordo}%
\vorto{nordo}\senco{1}{Un ek la quar punti kardinala, punto di horizonto situita en la direciono di la polo-stelo}\senco{2}{Parto di la Terglobo, regiono situita vicine a nordo}\lingui{DEFIRS}
\artiklo{norio}%
\vorto{norio}\fako{teknol.} Mashino hidraulikala, qua konsistas ye siteli ligita alonge kateno senfina, qua spulas sur tamburo, per rotaco, e per qua la siteli pleneskas ed acensas, ca-latere – vakueskas e decensas, ta-latere\lingui{DEF}
\artiklo{normo}%
\vorto{normo}\senco{1}{Regulo, principo di konduto}\senco{2}{\textit{(industrio)} Regulo qua fixigas la cirkonstanci di la faco di objekto o di la elaboro di produkturo di qua onu volas uniformigar la uzo o certigar la inter-identeso perfekta, do la interremplasableso sen desfacileso}\senco{3}{Tipo legala di la mezurili e ponderili lego-permisata}\lingui{DEFIRS}
\artiklo{nostalgio}%
\vorto{nostalgio} Langoro quan produktas la regretado di la nasko-lando\lingui{EFIRS}
\artiklo{nostoko}%
\vorto{nostoko}\fako{historio naturala} Algo filamentoza, envelopita da gaino gelatinatra\latina{nostoc}\lingui{DEFS}
\artiklo{notar}%
\vorto{notar}\fako{trans.} Skribar to di quo onu volas konservar la menciono, la indiko\lingui{DEFIRS}
\artiklo{notario}%
\vorto{notario} Oficiero publika, establisita por redaktar e recevar olti ek la akti, kontrati, e c., ad olqui la privati volas donar karaktero di autentikeso\lingui{DEFIRS}
\artiklo{notico}%
\vorto{notico} Skriburo destinata ad informar pri punto partikulara di historio, di cienco, e c.\lingui{DEFIS}
\artiklo{notifikar}%
\vorto{notifikar}\fako{trans.} Igar ulu savar, uzante por lo la formo oficala\lingui{DEFIS}
\artiklo{notora}%
\vorto{notora} Qua konocesas, savesas, da publiko\lingui{DEFIRS}
\artiklo{nova}%
\vorto{nova}\senco{1}{Qua aparas unesmafoye, od aparis erste de poka tempo}\senco{2}{Qua aparas pos altra, quan lu remplasas}\senco{3}{\textit{(*nuva)} Quan onu ne ja komencis konsumar, o quan onu jus komencis konsumar}\lingui{DEFIRS}
\artiklo{novelo}%
\vorto{novelo}\fako{literaturo} Naraco (kelka-pagina) pri aventuro interesiva\lingui{DEFIRS}
\artiklo{novembro}%
\vorto{novembro} La monato dek-ed-unesma di la yaro\lingui{DEFIRS}
\artiklo{novico}%
\vorto{novico} La persono, qua, jus diveninte reguliero (che la katoliki), subisas probado dum ula tempo ante profesar. – \textit{(anke metaf.)}\lingui{DEFIS}
\artiklo{nu!}%
\vorto{nu!} Interjeciono qua : (1) segun la pronunco-tono, expresas astoneso o konsento; (2) preiras afirmajo quan onu qualifikas kom stranja, neexpektita\lingui{DR}
\artiklo{nuanco}%
\vorto{nuanco}\senco{1}{Singla di la toni di un koloro, de la maxim klara til la maxim obskura}\senco{2}{\textit{(metaf.)} Difero di grado, preske nesentebla, neperceptebla, qua distingas kozo de olti qui esas lua-vicina}\lingui{DEFR}
\artiklo{nubo}%
\vorto{nubo}\fako{meteorol.} Amaso de aquo-vaporo, qua trublas la diafaneso di atmosfero, e finas per divenar pluvo-guti\lingui{EIS}
\artiklo{nucelo}%
\vorto{nucelo}\fako{biol.} Parto di la ovulo, qua envelopas la sako embrionala e qua tegesas da la tegumenti di la ovulo
\artiklo{nuco}%
\vorto{nuco}\fako{bot.} Frukto di arboro ek la familio « juglandei », di qua la ligno uzesas da la ebenisti, moblifisti\latina{juglans}
\artiklo{nuda}%
\vorto{nuda}\senco{1}{Qua ne *weras vesti}\senco{2}{Qua ne esas garnisita (ye ulo)}\lingui{EFIS}
\artiklo{nudlo}%
\vorto{nudlo}\fako{koquarto} Singla de la bendi streta qui konsistas ye pasto ek farino ed ovi, e tranchita por formacar ta bendi\lingui{DE}
\artiklo{nugato}%
\vorto{nugato} Pasto ek mandeli torefaktita e sukro\lingui{EF}
\artiklo{nukleino}%
\vorto{nukleino}\fako{biol.} Substanco ye qua konsistas la celulo-nukleo e qua kontenas til 5/100 fosforo\lingui{DEFIS}
\artiklo{nukleo}%
\vorto{nukleo}\senco{1}{\textit{(biol.)} Parto kompakta ye qua konsistas la centro di ula kozi (protoplasmo, e c.)}\senco{2}{\textit{(atomistiko)} Korpo elektro-pozitiva, konsistanta ye elektro-pozitiva protoni ed elektro-neutra neutroni, qua konstitucas la maso precipua di atomi}\senco{3}{\textit{(metaf.)} Ta personi qui su asemblis dat-unesme e cirkum qui plusi venas aglomerar su}\lingui{EFS}
\artiklo{nukleolo}%
\vorto{nukleolo}\fako{biol.} Mikra korpo ronda, unika o multopla, qua existas en la nukleo di la celuli
\artiklo{nukleoplasmo}%
\vorto{nukleoplasmo}\fako{biol.} Protoplasmo di la nukleo
\artiklo{nuko}%
\vorto{nuko}\fako{anat.} Parto dopa di kolo, ube ica sujuntas a la kapo\lingui{FIS}
\artiklo{nula}%
\vorto{nula} Ne mem nur un
\artiklo{numeratoro}%
\vorto{numeratoro}\fako{aritm.} Ta ek la du nombri di fraciono, qua indikas quanta partin inter-egala di la unajo lu kontenas\lingui{DEFIRS}
\artiklo{numero}%
\vorto{numero}\fako{aritm.} Vorto qua, en serio de termini, indikas la rango di un ek la termini relate la ceteri\lingui{DEFIRS}
\artiklo{numismatiko}%
\vorto{numismatiko} La cienco qua traktas la moneto-peci e medalii\lingui{DEFIRS}
\artiklo{numulario}%
\vorto{numulario}\fako{bot.} Planto, vicina ye la primrozo\latina{Lysimachia numularia}\lingui{FIS}
\artiklo{numulito}%
\vorto{numulito}\fako{paleont.} Genero de foraminiferi, tipo di la familio « numulitidei », kun konki ronda, disko-forma, e dividita ye mikra lojii multega\lingui{DEFIS}
\artiklo{nun}%
\vorto{nun} En la instanto, en la epoko, en la periodo, prezenta\lingui{GR}
\artiklo{nuncio}%
\vorto{nuncio} Ambasadisto di la papo\lingui{DEFIRS}
\artiklo{nur}%
\vorto{nur} Sen ulo od ulu plusa\lingui{D}
\artiklo{nutacar}%
\vorto{nutacar}\fako{netrans.) (biol. ed astron.} Ocilo-movo da organo vivoza o da objekto\lingui{DEFIS}
\artiklo{nutrar}%
\vorto{nutrar}\fako{trans.} Entratenar (la ento vivanta) per alimentar (lu) ye la kozi necesa por lua substanco\lingui{EFIS}
\end{multicols}
\sekciono{O}
\begin{multicols}{2}
\artiklo{o(d)}%
\vorto{o(d)} Konjunciono qua indikas alternativo
\artiklo{oaziso}%
\vorto{oaziso} Loko sola, kun vejetantaro, en dezerto de sablo\lingui{DEFIRS}
\artiklo{obcena}%
\vorto{obcena} Qua shokas pudoro\lingui{DEFIS}
\artiklo{obediar}%
\vorto{obediar}\fako{trans.}\senco{1}{Agar konforme a to quon imperas ulo}\senco{2}{\textit{(aludante kozo)} Subisar la ago da ulu o da ulo}\lingui{EFIS}
\artiklo{obelisko}%
\vorto{obelisko} Monumento (maxim-multa-kaze monolita) taliita forme di agulo quadrangula\lingui{DEFIRS}
\artiklo{obeza}%
\vorto{obeza} Tro repleta e, konseque, di qua la ventro (che la viri) o la hanchi e la gropo (che la mulieri) salias exajere pro akumulaji di graso sub-pele\lingui{EFIS}
\artiklo{objecionar}%
\vorto{objecionar}\fako{trans., ad ulu}\senco{1}{Opozar (ulo) ad afirmo por kombatar olu}\senco{2}{Opozar (ulo) a propozo, a projeto, por repulsar olu}\lingui{EFIS}
\artiklo{objekto}%
\vorto{objekto}\senco{1}{To quo esas vidata, rigardata}\senco{2}{To qua su prizentas en la spirito, en la penso}\senco{3}{\textit{(filoz.)} To quo pensesas (opoze a to qua pensas)}\senco{4}{\textit{(gram.)} Termino di la propoziciono a qua agas la subjekto}\senco{5}{To a quo iras la sentimenti di la anmo, di la psiko}\senco{6}{To a quo iras volo}\lingui{DEFIRS}
\artiklo{oblato}%
\vorto{oblato} Pano sen-hefa, tranchita ronde, uzata por siglar kuverti\lingui{DRS}
\artiklo{obleo}%
\vorto{obleo} Kukajo dina, (normale) volvita forme di cilindro vakua, o di korneto\lingui{FRS}
\artiklo{obligaciono}%
\vorto{obligaciono}\fako{financo} Akto qua reprezentas parto di la kapitalo pruntita da societo financala, industriala, e c. – qua yurizas pri interesto-procento fixa, e qua rimborsesos pos fristo determinita\lingui{DEFIRS}
\artiklo{obligar}%
\vorto{obligar}\fako{trans., ad ulo} Ligar (ulu) per lego, per konvenciono, quan lu mustas obediar\lingui{DEFIS}
\artiklo{obliqua}%
\vorto{obliqua}\senco{1}{Qua deviacas de vertikalo}\senco{2}{Qua agas per moyeni nedireta}\lingui{EFIS}
\artiklo{obliterar}%
\vorto{obliterar}\fako{trans.}\senco{1}{Igar nelektebla per efacar, makulizar, e c.}\senco{2}{\textit{(pri *timbro postala od ulo analoga)} Kovrar lu per makuli punta o linea, tale ke lu ne esas uzebla plusafoye}\senco{3}{\textit{(fiziol.)} Obstruktar (ula tubo organa) per igar adheranta lua parieto}\lingui{DEFIS}
\artiklo{obliviar}%
\vorto{obliviar}\fako{trans.}\senco{1}{Ne konservir en sua memoro}\senco{2}{Ne pensar (pri ulu, pri ulo) konseque de neglijo}\lingui{EFIS}
\artiklo{oblonga}%
\vorto{oblonga} Plu longa kam larja\lingui{DEFS}
\artiklo{obolo}%
\vorto{obolo}\senco{1}{En Grekia antiqua : Monetuno, pezunajo, qua equivalis la sisimo di drakmo}\senco{2}{\textit{(metaf.)} Quanto tre mikra de pekunio}\lingui{DEFIS}
\artiklo{obsedar}%
\vorto{obsedar}\fako{trans.}\senco{1}{Fatigar per demandi sencesa}\senco{2}{Tormentar per tenti sencesa}\lingui{EFIS}
\artiklo{observar}%
\vorto{observar}\fako{trans.} Regardar kun atenco kontinua\lingui{DEFIRS}
\artiklo{observatorio}%
\vorto{observatorio} Edifico provizita ye la kozi necesa por la observi astronomiala, meteorologiala, e c.\lingui{DEFIS}
\artiklo{obsidiano}%
\vorto{obsidiano}\fako{mineral.} Substanco vitratra, de origino volkanala, obskure-verda, tenaca, bele-polisebla\lingui{DEFIRS}
\artiklo{obskura}%
\vorto{obskura}\senco{1}{Qua esas sen-luma}\senco{2}{\textit{(metaf.)} (1) Ube lumo mankas. Qua ne esas klara ye spirito. (2) Qua ne esas eminenta. (3) \textit{(muziko)}. \textit{(pri voco)} Di qua la sono esas quaze veloza. (4) Dicesas pri la soni di ula vokali di la linguo Angla, qui kelke similesas la sonuno « e » di la linguo Franca}\lingui{DEFIS}
\artiklo{obskurantismo}%
\vorto{obskurantismo} Nomo quan (en Francia) la liberisti donis, lor la Restauro (en 1815) a la doktrini di la personi qui kombatis la idei naskinta de la agitado sociala di la yaro 1789\lingui{DEFIRS}
\artiklo{obsoleta}%
\vorto{obsoleta} Qua cesis esar segun-moda
\artiklo{obstaklo}%
\vorto{obstaklo}\senco{1}{To quo su opozas a la paso}\senco{2}{\textit{(metaf.)} To quo su opozas a la realigo di projetajo, a la obteno di dezirajo}\lingui{EFIS}
\artiklo{obstinar}%
\vorto{obstinar}\fako{netrans., pri} Perseverar tenace en rezolvajo, mem se nerealigebla, neatingebla\lingui{EFIS}
\artiklo{obstruktar}%
\vorto{obstruktar}\fako{trans.}\senco{1}{Embarasar (voyo, orifico) per obstaklo}\senco{2}{\textit{(metaf.)} \textit{(en asemblitaro deliberala)} Sistemoze opozar obstaklo kontre diskuto}\lingui{DEFIRS}
\artiklo{obtenar}%
\vorto{obtenar}\fako{trans., de}\senco{1}{Sucesar pri grantigar a su}\senco{2}{Sucesar pri haveskar, per favoro}\lingui{EFIS}
\artiklo{obturar}%
\vorto{obturar}\fako{trans., per) (tekn.} Stopar
\artiklo{obtuza}%
\vorto{obtuza}\senco{1}{Di qua la angulo esas rondatra, senpintigita}\senco{2}{\textit{(geom.)} Angulo plu granda kam orto}\senco{3}{\textit{(metaf.)} Spirito-nesagaca}\lingui{EFIS}
\artiklo{obuso}%
\vorto{obuso}\fako{tekniko milit.} Projektilo kava, oblonga, sen anso e sen kulo, qua spliteskas per perkuteso, lansate per kanono\lingui{FI}
\artiklo{oceano}%
\vorto{oceano} Vasta extensajo de aquo saloza, qua okupas granda parto di la Terglobo\lingui{DEFIRS}
\artiklo{oceloto}%
\vorto{oceloto}\fako{zool.} Kato-tigro di Amerika\latina{felis pardalia}\lingui{DEFIRS}
\artiklo{ociar}%
\vorto{ociar}\fako{netrans.} Agar nulo, facar nulo, pro repugneso da laboro\lingui{FIS}
\artiklo{ocidar}%
\vorto{ocidar}\fako{trans.} Mortigar vole e violente\lingui{EFI}
\artiklo{ocidento}%
\vorto{ocidento}\fako{literaturo}\senco{1}{Loko di la horizonto ube suno su kushas}\senco{2}{Parto di la Terglobo, regiono situita ne for la regiono ube suno semblas kushar su}\lingui{DEFIS}
\artiklo{ocilar}%
\vorto{ocilar}\fako{netrans.} Deviacar de sua baricentro, e retrovenar ad ica per movado alternanta konstanta\lingui{DEFIS}
\artiklo{ocipito}%
\vorto{ocipito}\fako{anat.} La infrajo di la dopa parto di la kapo\lingui{EFIS}
\artiklo{ocitar}%
\vorto{ocitar}\fako{netrans.} Apertar sua boko ye movo spasma di expiro, sequata da expiro plu longa, produktita da la bezono dormar, da hungro, da fatigeso, da enoyo, da tedeso\lingui{L}
\artiklo{od}%
\vorto{od} Konjunciono qua expresas alternativo
\artiklo{odalisko}%
\vorto{odalisko}\senco{1}{Sklavino en haremo, destinita por servar la spozini di la sultano}\senco{2}{Muliero ek haremo}\lingui{DEFIRS}
\artiklo{odeono}%
\vorto{« odeono »}\fako{en epoki antiqua} Edifico en qua onu repetis la muzikaji destinita a kantesor en la teatro\lingui{DEFIRS}
\artiklo{odiar}%
\vorto{odiar}\fako{trans.}\senco{1}{\textit{(ulu)} Malvolar profunde}\senco{2}{\textit{(ulo)} Repugnesar profunde da (lu)}\lingui{EFIS}
\artiklo{odiseo}%
\vorto{odiseo} Naraco di voyajo ed aventuri diversa\lingui{DEFIRS}
\artiklo{odo}%
\vorto{odo}\senco{1}{\textit{(en la epoki antiqua)} Poemo lirika, dividita aden strofo, antistrofo ed epodo, quan la koro kantis evolucionante}\senco{2}{Poemo lirika, dividita aden strofi intersimila pri nombro e la mezuro di la versi}\lingui{DEFIRS}
\artiklo{odontoblasto}%
\vorto{odontoblasto}\fako{biol.} Grosa celulo, kun nukleo tre voluminoza, quan onu renkontras en la parto periferia di la pulpo dentala, e di qua la rolo specala esas sekrecar ivoro
\artiklo{odorar}%
\vorto{odorar}\fako{netrans.) (aludante ula korpi} Exhalar ulo qua impresas partikulare la olfakt-organo\lingui{EFIS}
\artiklo{ofensar}%
\vorto{ofensar}\fako{trans.} Vundar (ulu) pri lua digneso, respektindeso\lingui{EFIS}
\artiklo{ofertorio}%
\vorto{ofertorio}\fako{liturgio katol.}\senco{1}{Deo-ofro di pano e vino en la meso-sakrifiko}\senco{2}{Prego, muziko, qui akompanas ta ofro}\lingui{DEFIS}
\artiklo{oficiar}%
\vorto{oficiar}\fako{netrans.) (eklezio katol.} Celebrar la deo-servo kun la ceremonii qui devas akompanar lu\lingui{EFIS}
\artiklo{oficiro}%
\vorto{oficiro}\fako{armeo} La persono qua havas ofico komandala en armeo (terala, marala, aerala)\lingui{DEFIRS}
\artiklo{ofico}%
\vorto{ofico} Funciono administrala o direktala, sive en stato, sive en societo o granda entraprezajo\lingui{EFIS}
\artiklo{ofikleido}%
\vorto{ofikleido}\fako{muziko} Sufl-instrumento kupra, kun boko-peco, qua ne esas altro kam la serpento-korno, quan onu uzis en la kirki, a qua onu adjuntis klavi\lingui{DEFIS}
\artiklo{ofito}%
\vorto{ofito}\fako{mineral.} Marmoro obskur-verda, kun strii flava, quan onu renkontras en Pirenei\lingui{DEF}
\artiklo{ofrar}%
\vorto{ofrar}\fako{trans., ad} Igar (ulo, o, metaf., ulu) disponebla da altru, sen ke ica demandis lo\lingui{DEFIS}
\artiklo{oftalmio}%
\vorto{oftalmio}\fako{patol.} Morbo inflamurala di la okulo\lingui{DEFIS}
\artiklo{oftalmologio}%
\vorto{oftalmologio} Cienco qua traktas la strukturo e la funciono di la okuli e la morbi qui afektas li
\artiklo{ofte}%
\vorto{ofte} Qua eventas granda-quanta-foye, ma ne ye intertempi inter-egala – nur intermite\lingui{DE}
\artiklo{ogivo}%
\vorto{ogivo}\senco{1}{\textit{(arkeol.)} Arko ruptita quan formacas, en la arkitekturo gotika, la nervaturi di la travei di la vulto, interkrucumante diagonale ye la kolmo}\senco{2}{\textit{(arkitekt.)} Arkado ek du arki qui interkrucumas formacante angulo kurvo-linea, akuta}\lingui{DEFIS}
\artiklo{oglar}%
\vorto{oglar}\fako{trans.} Regardar furtatre, quaze per la komisuri di singla okulo (ulo) quan onu deziregas\lingui{E}
\artiklo{ohm}%
\vorto{« ohm »}\fako{elektro} Unajo praktikala di la rezisto elektrala = 10⁹ unaji elektromagnetala C.G\lingui{DEFIRS}
\artiklo{-oid}%
\vorto{-oid} Sufixo ciencala : « qua havas la sama formo kam » (o formo analoga)
\artiklo{ok}%
\vorto{ok}\fako{aritm.} La cifro « 8 »
\artiklo{okarino}%
\vorto{okarino}\fako{muziko} Muzik-instrumento ek terakoto, truizita, qua havas la formo di grosa ovo\lingui{DEF}
\artiklo{okaziono}%
\vorto{okaziono}\senco{1}{Cirkonstanco qua eventas oportune}\senco{2}{Cirkonstanco qua igas rezolvar agor ulo}\lingui{EFIRS}
\artiklo{okro}%
\vorto{okro} Farbo, argilo friabla, flava o reda, (segun lua kompozanti kemiala)\lingui{DEFIRS}
\artiklo{okta-}%
\vorto{okta-} Prefixo ciencala : ok
\artiklo{oktaedro}%
\vorto{oktaedro}\fako{geom.} Qua havas ok edri\lingui{DEFIS}
\artiklo{oktanto}%
\vorto{oktanto}\senco{1}{\textit{(astron.)} Disto de 45 gradi (okima parto di la cirklokurvo) inter du astri}\senco{2}{Instrumento ek sektoro 45-grada, uzata por mezurar la altesi e la disti}\lingui{DEFIS}
\artiklo{oktavo}%
\vorto{oktavo}\senco{1}{\textit{(muziko)} Intertempi de ok gradi}\senco{2}{\textit{(liturgio katol.)} Intertempo de ok dii qua sequas granda festo-dio ekleziala e dum qua eventas la memorigo di ta festo. – La lasta dio ek ta serio, en qua la oficio esas maxim solena}\lingui{DEFIS}
\artiklo{oktiliono}%
\vorto{oktiliono}\fako{aritm.} 10⁴⁸
\artiklo{oktobro}%
\vorto{oktobro} La monato dekesma di la yaro\lingui{DEFIRS}
\artiklo{oktogono}%
\vorto{oktogono}\fako{geom.} Figuro qua havas ok anguli ed ok lateri\lingui{DEFIS}
\artiklo{oktopodo}%
\vorto{oktopodo}\fako{zool.} Sub-klaso de moluski cefalopoda du-brankia, qua kontenas la pulpi, la argonauti, e c., opoze a la subklaso de la dekapodi (kalmari, sepii. e c.)\latina{octo-}
\artiklo{okulo}%
\vorto{okulo}\fako{anat.} Organo di la vido-fakultato\lingui{DEFIRS}
\artiklo{okulta}%
\vorto{okulta} Di qua la kauzo permanas celata\lingui{DEFIS}
\artiklo{okultacar}%
\vorto{okultacar}\fako{trans.} \textit{(astron.)} \textit{(pri astro)} Celar per interpozar su\lingui{DEFIS}
\artiklo{okupar}%
\vorto{okupar}\fako{trans.}\senco{1}{Posedeskar loko, instalar su en lu, establisar su en lu}\senco{2}{Posedeskar (ulu) per absorbar lua spirito, lua kordio}\senco{3}{Konsumar la tempo (di ulu)}\lingui{DEFIS}
\artiklo{ol(u)}%
\vorto{ol(u)} Pronomo personala, triesma persono, singularo, qua reprezentas irgo altra kam persono
\artiklo{olda}%
\vorto{olda} Qua vivas, od existas, de multa tempo\lingui{E}
\artiklo{oleo}%
\vorto{oleo} Substanco grasa, liquida, obtenita de vejetanto, animalo, mineralo, uzata por nutro, medicino, lumizo, la arti, ed ula ceremonii religiala\lingui{DEFIS}
\artiklo{olfaktar}%
\vorto{olfaktar}\fako{trans.} Perceptar la odoro di (ulo, ulu)\lingui{EFIS}
\artiklo{olibano}%
\vorto{olibano}\fako{farmacio} Rezino qua donas incenso\lingui{EFIS}
\artiklo{oligarkio}%
\vorto{oligarkio} Formo di guvernado, en qua la autoritato imperala praktikesas da mikra quanto de civitani privilejiera\lingui{DEFIRS}
\artiklo{oligarko}%
\vorto{oligarko} Membro di guvernistaro oligarkiala\lingui{DEFIRS}
\artiklo{oligocena}%
\vorto{oligocena}\fako{geol.} Dicesas pri grupo di tereno-strati terciara (sur eoceno, sub-an mioceno)\lingui{DEFIS}
\artiklo{olim}%
\vorto{olim} En epoko qua forfluis de longe ante nun (quankam ne antiqua)\lingui{L}
\artiklo{olimpiado}%
\vorto{olimpiado}\fako{antique, en Grekia} Periodo de quar yari, de celebro di la Olimpio-ludi til la celebro *nexta\lingui{DEFIRS}
\artiklo{olivo}%
\vorto{olivo}\fako{bot.} Frukto kernoza, elipsoida, verdatra, quan produktas arboro ye foliaro nevelkebla, e de qua onu extraktas oleo manjebla\latina{olea europaea}\lingui{DEFIRS}
\artiklo{ombro}%
\vorto{ombro}\senco{1}{Parto di surfaco lumizita, ube la lumo-radii interceptesas da korpo opaka}\senco{2}{La parto nelumizita qua, en surfaco lumizata, *reproduktas la formo di korpo opaka qua interceptas la lumo-radii}\lingui{EFIS}
\artiklo{omento}%
\vorto{omento}\fako{kirurg.} Granda rifalduro di la peritonio, qua fluktuas avan la intestino tenua\lingui{DEFIRS}
\artiklo{omisar}%
\vorto{omisar}\fako{trans.} Ne inkluzar, sive vole, sive pro oblivio, (ulo quon onu devas dicar od agar)\lingui{EFIS}
\artiklo{omleto}%
\vorto{omleto}\fako{koquarto} Disho ek ovi plaudita e koquita en padelo, kun butro e kondimenti o spici\lingui{DEFR}
\artiklo{omna}%
\vorto{omna}\senco{1}{\textit{(avan substantivo plurala)} Ici ed iti, sen ecepto}\senco{2}{\textit{(avan substantivo singulara)} Ica samagrade kam ita, sen dicerno nek distingo}\lingui{EIL}
\artiklo{omnibuso}%
\vorto{omnibuso} Veturo publika, qua cirkulas en la quarteri diversa di urbo, e haltas dum sua parkuro por lasar enirar (embarkar) od ekirar (desembarkar) la voyajanti quin lu transportas bloke po preco mikra\lingui{DEFIRS}
\artiklo{omnivora}%
\vorto{omnivora} Qua su nutras ye nutrivi irgaspeca – devenanta de animali o de vejetanti\lingui{DEFIS}
\artiklo{on(u)}%
\vorto{on(u)} Pronomo personala nedefinita, uzata por indikar un o plura personi ye maniero tre generala\lingui{EF}
\artiklo{onagro}%
\vorto{onagro}\senco{1}{\textit{(zool.)} Asno sovaja}\senco{2}{\textit{(milit-arto)} *Asiejo-mashino, uzita da la Romani antiqua, qua esis quaza balisto per qua onu lansis aden la urbo *asiejata projektili diversa}\latina{asinus onager}\lingui{DEFIS}
\artiklo{onaniar}%
\vorto{onaniar}\fako{netrans.} Produktar la orgasmo venerala per manuo vice per la koito normala\lingui{DEFIRS}
\artiklo{ondo}%
\vorto{ondo} Amaso de aquo qua superstaceskas e retrofalas sur la surfaco di maro, lago, fluvio, agitata da vento\lingui{EFIS}
\artiklo{onixo}%
\vorto{onixo} Agato delikata, kun strati interparalela diversa-kolora\lingui{DEFIRS}
\artiklo{onklo}%
\vorto{onklo} Frato di onua patro\lingui{DEF}
\artiklo{onomatopeo}%
\vorto{onomatopeo}\fako{gram.} Vorto formacita per imito-harmonio. (ex. : gluglo, zumo)\lingui{DEFIS}
\artiklo{-ont-}%
\vorto{-ont-}\fako{gram.} Dezinenco di la participo futura aktiva
\artiklo{ontogenezo}%
\vorto{ontogenezo}\fako{biol.} Developo di la individuo, opoze a « filogenezo », o developo di la speco\lingui{DEFIS}
\artiklo{ontologio}%
\vorto{ontologio}\senco{1}{\textit{(biol.)} Parto di biologio, qua traktas la fenomeni individuala quin prizentas organismo dum lua developo, lua adulteso e lua oldeso}\senco{2}{\textit{(filoz.)} Cienco metafizikala qua retroiras til la esenco, la kauzo prima di la ento}\lingui{DEFIRS}
\artiklo{onyono}%
\vorto{onyono}\fako{bot.} Planto leguma, di qua la radiko bulba – qua odoras e saporas forte – uzesas en la preparaji nutrala\latina{allkum cepa}\lingui{EF}
\artiklo{oocito}%
\vorto{oocito}\fako{biol.} Ovulo nematura nek fekundigebla, qua produktas (per du dupartigi) la ovulo matura e fekundigebla, e la globuli polala\lingui{DEFIRS}
\artiklo{oogenezo}%
\vorto{oogenezo}\fako{biol.} Formaco di la ovulo matura e fekundigebla\lingui{DEFIRS}
\artiklo{oogonio}%
\vorto{oogonio}\fako{biol.} Formaco di la celulo en qua aparas la elementi feminala, o \textit{oosferi}, che multa vejetanti
\artiklo{oolito}%
\vorto{oolito}\fako{geol.} Kalkajo ek mikra grani ovoida qui kelke similesas fish-ovi\lingui{DEFIRS}
\artiklo{oosfero}%
\vorto{oosfero}\fako{biol.} Elemento feminala qua, pos fekundigesir da la elemento maskulala (polinido od anterozoido) genitas la ovo – departo-punto di nova ento vejetanta
\artiklo{oosporo}%
\vorto{oosporo}\fako{bot.} Nomo donita ulakaze a la ovo di la algi o di la fungi, precipue kande la organo feminala (oosfero) liberigesas ante fekundigeso, nam lore la ovo kelke similesas sporo\lingui{DEFIRS}
\artiklo{-op-}%
\vorto{-op-} Sufixo distributala, kun nomi di nombro o di quanto
\artiklo{opaka}%
\vorto{opaka} Qua ne lasas la lumo-radii pasar tra su\lingui{DEFIS}
\artiklo{opalecar}%
\vorto{opalecar}\fako{netrans.} Qua aquiras nuanco opalea\lingui{DEFIS}
\artiklo{opalo}%
\vorto{opalo}\fako{geol.} Lapido volkanala, lakto-blanka e bluatra, kun reflekti koloro-chanjanta\lingui{DEFIRS}
\artiklo{operacar}%
\vorto{operacar}\senco{1}{\textit{(trans.)} \textit{(kirurg.)} Submisar ad interveno kirurgiala}\senco{2}{\textit{(netrans.)} (1) \textit{(matem.)} Agar la kalkuli quin onu efektigas departante de quanti konocita, por renkontrar un o plura quanti nekonocita. (2) \textit{(milit-arto)} Agar la trupo-diplasi qui havas kom skopo sucesigar atako o defenso. – (3) \textit{(en borso financala o komercala)} Manovrar por sucesigar profitoze la kompro o la vendo di la komercajo}\lingui{DEFIRS}
\artiklo{operkulo}%
\vorto{operkulo}\fako{tekn.} Aparato destinita ad obturar ula orifici\lingui{EFIS}
\artiklo{opero}%
\vorto{opero} Verko teatrala e muzikala, qua konsistas ye recitativi, kanti helpata da orkestro, ed ulakaze kun dansi\lingui{DEFIRS}
\artiklo{opiato}%
\vorto{opiato}\senco{1}{Elektuario opiumoza}\senco{2}{Irga elektuario, pasto ek pudro dilutita en siropo}\lingui{DEFIS}
\artiklo{opinionar}%
\vorto{opinionar}\fako{netrans.} Enuncar to quon onu judikas (sen fundamento ciencale establisita) pri temo o persono\lingui{EFIS}
\artiklo{opiumo}%
\vorto{opiumo} Substanco narkotigiva, kalmigiva se uzata moderata-doze, ma venena granda-doze, quan onu extraktas de la suko di ula sorti de papaveri\lingui{DEFIRS}
\artiklo{-opl-}%
\vorto{-opl-} Sufixo pri multipliko
\artiklo{opodeldoko}%
\vorto{opodeldoko}\fako{farmacio} Medikamento qua havas kom fundamento solvuro alkoholoza ek sapono animalala o medicinala, a qua onu adjuntas kamforo, amoniako ed esenco de timiano e de rosmarino, uzata fricione (kontre la dolori). – DEFI?
\artiklo{oportar}%
\vorto{oportar}\fako{netrans.} Esar absolute necesa\lingui{L}
\artiklo{oportuna}%
\vorto{oportuna} Qua eventas precize kande lo esas bona por to quo koncernesas\lingui{DEFIS}
\artiklo{opozar}%
\vorto{opozar}\fako{trans.}\senco{1}{Lokizar (ulo) regarde ulo}\senco{2}{Lokizar (ulu, ulo) regarde altra, por esar obstaklo pri ica}\senco{3}{Alegar (ulo) kontre to quon dicas o pensas altru}\lingui{DEFIRS}
\artiklo{opresar}%
\vorto{opresar}\fako{trans.}\senco{1}{Jenar (la pektoro) quale per preso qua havas kom efekto, ke onu respiras penoze}\senco{2}{Koaktar per autoritato tiranala}\senco{3}{Igar prostracar korpale}\lingui{EFIS}
\artiklo{optativo}%
\vorto{optativo}\fako{gram.} Modo di la konjugo, qua expresas la deziro\lingui{DEFIS}
\artiklo{optiko}%
\vorto{optiko} Cienco pri la fenomeni di lumo e di vidado\lingui{DEFIRS}
\artiklo{optimismo}%
\vorto{optimismo}\senco{1}{\textit{(filoz.)} Doktrino segun qua deo kreis nia mondo kom la maxim bona ek le posibla}\senco{2}{Tendenco judikar ke irgo esas la maxim bona ek la posiblaji}\lingui{DEFIRS}
\artiklo{-or}%
\vorto{-or}\fako{gram.} Dezinenco di la infinitivo di futureso
\artiklo{or}%
\vorto{or}\fako{konjunciono}\senco{1}{En la punto quan atingis la rezono}\senco{2}{\textit{(logiko)} Konjunciono uzata por venigar la minoro, pos la majoro, di silogismo}\lingui{FIS}
\artiklo{oraklo}%
\vorto{oraklo}\senco{1}{\textit{(en la epoki antiqua)} Respondo agita, ye la nomo di la deajo, a la persono qua venis konsultar lu en lua templo}\senco{2}{Dico da deo, enuncita da la profeti, la apostoli, e c.}\senco{3}{\textit{(metaf.)} Respondo, decido, egardata kom neeroriva}\lingui{DEFIRS}
\artiklo{orangutano}%
\vorto{orangutano}\fako{zool.} Granda simio sen-kauda, qua similesas pasable la homo de la vidopunto « formo extera »\latina{simia satyrus}\lingui{DEFIS}
\artiklo{oranjo}%
\vorto{oranjo}\fako{bot.} Frukto de arbusto di qua la foliaro ne velkas sukoza, parfumoza, di qua la shelo esas flava oree\latina{citrus aurantium}\lingui{DEFIRS}
\artiklo{oratorio}%
\vorto{oratorio}\fako{muziko} Dramato religiala, akompanata da muzikajo, destinita a pleeso sen dekoruri nek kostumi teatrala\lingui{DEFIRS}
\artiklo{oratoro}%
\vorto{oratoro} Persono qua, profesione, agas diskursi\lingui{EFIRS}
\artiklo{orbito}%
\vorto{orbito}\senco{1}{\textit{(astron.)} Kurvo alonge qua iras, propra-move, ula korpi cielala (planeti, kometi, e c.)}\senco{2}{\textit{(anat.)} Ostokavajo en qua esas la okuli}\lingui{EFIRS}
\artiklo{ordeno}%
\vorto{ordeno} Singla de la grupi (religiala, *chevalierala, advokatala, medikala) importoza, ye qua konsistas klasifikuro\lingui{DEFIRS}
\artiklo{ordinacar}%
\vorto{ordinacar}\fako{trans.) (liturgio katol.} Elevar a sacerdoteso, diakoneso, e c., per konferar (a lu) ta ordo sakra\lingui{DEFIS}
\artiklo{ordinar}%
\vorto{ordinar}\fako{trans.} Dispozar – spacale, tempale o hierarkiale – segun ula regulo, ula normo, la kozi – ici relate iti\lingui{DEFIS}
\artiklo{ordinara}%
\vorto{ordinara} Qua esas konforma a la ordineso kustumala; qua esas ye la grado kustumala\lingui{EFIS}
\artiklo{ordo}%
\vorto{ordo}\fako{eklezio katol.} Sakramento qua grantas un ek la gradi en la hierarkio ekleziala
\artiklo{ordonanco}%
\vorto{ordonanco}\fako{armeo}\senco{1}{Servisto soldata di oficiro}\lingui{DFIRS}\subvorto{}{ordonanc-oficiro}{Kavaliero qua portas ad oficiro superiora la komandi de alta oficiero di la stato}{2}{}
\artiklo{oreado}%
\vorto{oreado}\fako{mitol.} Nimfo di la monti, di la boski\lingui{DEFS}
\artiklo{oreliono}%
\vorto{oreliono}\fako{patol.} Influro di la celulo-tisuo qua envelopas la glando parotida\lingui{F}
\artiklo{orelo}%
\vorto{orelo}\fako{anat.}\senco{1}{La organo di audado}\senco{2}{Parto extera di la aud-organo, ye singla latero di la kapo}\lingui{EFIS}
\artiklo{oremus}%
\vorto{« oremus »}\fako{liturgio katol.} Prego
\artiklo{orfano}%
\vorto{orfano} Infanto o puero di qua la patri mortis\lingui{EFIS}
\artiklo{organdio}%
\vorto{organdio} Muslino apretizita\lingui{DEFIS}
\artiklo{organicismo}%
\vorto{organicismo}\fako{filoz.} Doktrino qua admisas la ideo ke vivo rezultas, ne de forco qua « anmizas » la organi, ma de la organi ipsa\lingui{DEFIRS}
\artiklo{organika}%
\vorto{organika} Qua relatas la enti organizita.\lingui{DEFIRS}\subvorto{}{kemio organika}{Qua studias la produkturi di la enti organizita}{}{}
\artiklo{organismo}%
\vorto{organismo}\senco{1}{Ensemblo de la organi ye qui konsistas ento vivoza}\senco{2}{\textit{Anke metaf}}\lingui{DEFIRS}
\artiklo{organo}%
\vorto{organo}\senco{1}{En korpo vivoza : grupo de elementi, koordinita tale ke exekutesas funciono specala}\senco{2}{Anke pri mashino}\senco{3}{\textit{(metaf.)} To quo expresas la penso, la opiniono, la sentimenti di altru}\lingui{DEFIRS}
\artiklo{organzino}%
\vorto{organzino} Silko-filo qua pasis dufoye en la mashino qua tordas la silko kruda, e qua destinesas a formacar la warpo di stofo\lingui{DEFI}
\artiklo{orgeno}%
\vorto{orgeno}\fako{muziko} Sufl-instrumento muzikala, ek tubi di dimensioni diversa, quan onu resonigas, mediace klavaro, per enduktar en olu la aero quan furnisas suflilo\lingui{EFIRS}
\artiklo{orgio}%
\vorto{orgio}\senco{1}{\textit{(en la epoki antiqua)} Festo honorize di Bako (Bacchus)}\senco{2}{Krapulo pro troa manjo e troa drinko}\lingui{DEFIRS}
\artiklo{oriento}%
\vorto{oriento}\senco{1}{\textit{(literaturo)} Punto di la horizonto de ube suno semblas levar su}\senco{2}{Regiono di la terglobo, este di Europa (precipue Avan-Azia ed Egiptia)}\lingui{DEFIRS}\subvorto{framasonismo}{Granda Oriento}{Quaza lojio centra, en la urbo precipua di lando, qua konsistas ye la reprezenteri di la lojii framasonala di la provinci}{3}{}
\artiklo{orifico}%
\vorto{orifico} Aperturo qua formacas la enireyo di kavajo\lingui{EFIS}
\artiklo{oriflamo}%
\vorto{oriflamo} Mikra standardo di qua la parto qua fluktuas finas per du pinti\lingui{DEFIS}
\artiklo{origano}%
\vorto{origano}\fako{bot.} Planto aromatoza, ek la familio « labiei », majorano sovaja\latina{origanum}\lingui{DEFIS}
\artiklo{origino}%
\vorto{origino}\senco{1}{Departo-punto (loko, raso, e c.) di la nasko di individuo, familio, populo – di la formaco di ulo}\senco{2}{Instanto di la nasko}\senco{3}{\textit{(geom.)} Punto de qua onu kontas la koordinati}\lingui{DEFIRS}
\artiklo{orikalko}%
\vorto{orikalko}\fako{en la epoki antiqua} Kompozuro metala, analoga a latuno\lingui{DEFIS}
\artiklo{oriolo}%
\vorto{oriolo}\fako{zool.} Pasero di qua la plumaro esas flava\latina{oriolus}\lingui{EF}
\artiklo{oripelo}%
\vorto{oripelo}\senco{1}{Filo o folio de latuno, polisita, qua, de fore, brilas quale oro, e de qua onu facas ornamenti}\senco{2}{Folieto dina, oratra, quan onu uzas por oro-sembligar objekti diversa}\lingui{FI}
\artiklo{orjato}%
\vorto{orjato} Siropo facita, olim, ek dekokturo de hordeo – e, nun, ek emulsuro de mandeli\lingui{DEFIS}
\artiklo{orkestro}%
\vorto{orkestro}\senco{1}{(1) \textit{(en la epoki antiqua)} Parto di la teatro Greka avan e sub la ceneyo (« *stejo ») ube la koro stacis ed evolucionis. (2) Parto korespondanta di la teatro moderna, ube esas la muzikisti qui pleas muzikaji koncertala}\senco{2}{Ensemblo de ta muzikisti}\lingui{DEFIRS}
\artiklo{orkideo}%
\vorto{orkideo}\fako{bot.} Un ek la familii de planti monokotiledona, bulboza, de qua la « orkis » esas la tipo\latina{orchidea}\lingui{DEFIRS}
\artiklo{orlar}%
\vorto{orlar}\fako{trans.} Facar la bordo di texuro, per rifaldo e suto, por impedar la destexeso\lingui{FIS}
\artiklo{ornamento}%
\vorto{ornamento} Detalo adaptita ad ensemblo por igar lu plu bela\lingui{DEFIS}
\artiklo{ornar}%
\vorto{ornar}\fako{trans.} Kovrar per vesti eleganta, per ulo qua plubeligas, per ornamenti. – \textit{Anke metaf}\lingui{EFIS}
\artiklo{ornitologio}%
\vorto{ornitologio} Historio naturala pri la ucelo\lingui{DEFIRS}
\artiklo{ornitoptero}%
\vorto{ornitoptero} Papiliono belega, la maxim granda ek le jornala, qua vivas en la foresti humida, e flugas tre alte\lingui{DEFIRS}
\artiklo{oro}%
\vorto{oro}\fako{kemio} Metalo flava, tre precoza, tre densa, tre duktila, tre maleebla, nealterebla da aero nek da aquo, quan onu uzas (aloyita kun kelka kupro) por facar moneto-peci granda-valora, juveli, e c.\simbolo{Au}\lingui{FIS}
\artiklo{orografio}%
\vorto{orografio} Deskripto di la monti\lingui{DEFIRS}
\artiklo{orpimento}%
\vorto{orpimento} Sulfido flava di arseno, uzata kom farbo e por senpiligar la feli. Formulo kemiala : As₂S₃\lingui{DEFIS}
\artiklo{orselo}%
\vorto{orselo}\fako{bot.} Likeno qua kreskas precipue en Madagaskarinsulo, e donas bela farbo violea\latina{roccella tinctoria}
\artiklo{orta}%
\vorto{orta}\fako{geom.) (pri angulo} Di qua la aperturo kontenas 90 gradi di la cirklo-kurvo
\artiklo{ortocentro}%
\vorto{ortocentro}\fako{geom.} Nomo di la konvergo-punto di la tri altesi di triangulo
\artiklo{ortodoxa}%
\vorto{ortodoxa}\fako{teol. kristana} Konforma a la doktrino korekta. – \textit{(jokoze)}. Qua esas to quon lu devas esar\lingui{DEFIRS}
\artiklo{ortogenezo}%
\vorto{ortogenezo}\fako{biol.} Vario spontana, pro kauzo interna, qua pulsas la organismi a punto determinita, sen la interveno da adapto
\artiklo{ortogono}%
\vorto{ortogono}\fako{geom.} Qua formacas angulo orta – qua esas perpendikla
\artiklo{ortografika}%
\vorto{ortografika}\fako{matem.} Qua reprezentas objekto sur plano, lua omna punti esante projektata perpendikle sur la plano\lingui{DEFI}
\artiklo{ortografio}%
\vorto{ortografio}\fako{gram.} Maniero skribar linguo korekte. Maniero skribar la vorti\lingui{DEFIRS}
\artiklo{ortokrona}%
\vorto{ortokrona}\fako{fotogr.) (pri la plaki o filmi fotografala} Di qua la sentiveso pri la kolori diversa di la spektro igesis, per la adjunto di korpi kemiala, harmoniar kun la sentiveso di la hom-okulo
\artiklo{ortolano}%
\vorto{ortolano}\fako{zool.} Mikra ucelo migranta, qua vivas en Mezeuropa, ed en Sudeuropa, di qua la karno (manjebla) esas delikata\latina{emberiza hortulana}\lingui{DEFIS}
\artiklo{ortopedio}%
\vorto{ortopedio} La arto preventar o pokope supresar la deformuri di la korpo (precipue che la puerii), per aparato o kuracado\lingui{DEFIRS}
\artiklo{ortoptero}%
\vorto{ortoptero}\fako{zool.} Insekto qua havas quar ali, di qui la du infra esas faldita longeso-sinse\lingui{DEFIS}
\artiklo{ortoptika}%
\vorto{ortoptika}\fako{matem.) (pri kurvo} Loko geometriala di la punti de ube onu vidas kurvo plana donata ortangule
\artiklo{orveto}%
\vorto{orveto}\fako{zool.} Mikra reptero sauria, ne-nociva, tipo di la familio « anguidi », qua vivas en Europa ed en Azia orientala\latina{anguis fragilis}\lingui{EF}
\artiklo{orvietano}%
\vorto{orvietano} Elektuario (varianto nenociva di teriako) quan onu judikis, en la mezo di la 17esma yarcento, kom tre apta e tre efikiva\lingui{DEFIS}
\artiklo{-os}%
\vorto{-os}\fako{gram.} Che la verbi : dezinenco di indikativo, futuro
\artiklo{osifrago}%
\vorto{osifrago}\fako{zool.} Speco di raptucelo jornala, ek la familio « falkonidi », qua vivas en la nordo-regioni di Europa\latina{aquila ossifraga}\lingui{EFI}
\artiklo{oskular}%
\vorto{oskular}\fako{trans.) (geom.} Efektigar kontakto supra-klasa kun kurvo, surfaco, en un punto di altra kurvo, surfaco
\artiklo{osmio}%
\vorto{osmio}\fako{kemio} Korpo simpla, metala, quan onu renkontras en ula platin-erco\simbolo{Os}\lingui{DEFIRS}
\artiklo{osmoso}%
\vorto{osmoso} \textit{(biol. e kemio)} Fenomeno qua konsistas ye la displacesi qui eventas de ca a ta, ek du liquidi di qui la koncentreso-gradi interdiferas, tra parieto mipermeebla\lingui{DEFS}
\artiklo{ostensorio}%
\vorto{ostensorio}\fako{liturgio katol.} Suportilo ek arjento, arjento orizita, oro, en qua onu expozas, a la adoro da la religiani, la hostio sakrigita\lingui{DEFI}
\artiklo{ostentar}%
\vorto{ostentar}\fako{netrans., pri} Montrar afektacoze to quo esas apta flatar sua vanitato\lingui{EFIS}
\artiklo{osteoblasto}%
\vorto{osteoblasto}\fako{biol.} Celulo yuna, restinta e retroveninta a sua stando embrionala, quan onu vidas en la osto-medulo e la strato ostifera di periosto\lingui{DEFIRS}
\artiklo{osteologio}%
\vorto{osteologio} Parto di anatomio, qua traktas la osti\lingui{DEFIRS}
\artiklo{osteomo}%
\vorto{osteomo}\fako{patol.} Tumoro ek osto-tisuo\lingui{DE}
\artiklo{ostito}%
\vorto{ostito}\fako{patol.} Inflamuro di la osto-tisuo\lingui{DEF}
\artiklo{osto}%
\vorto{osto} Parto harda e solida, ye qua konsistas la skeleto che la animali vertebroza di la klasi supra\lingui{EFIS}
\artiklo{ostracismo}%
\vorto{ostracismo}\fako{en Grekia antiqua} Exilo, dek-yar-dura, di civitano qua divenis suspektata pro lua autoritato, lua povo, en la civito\lingui{DEFIRS}
\artiklo{ostro}%
\vorto{ostro}\fako{zool.} Genero de moluski lamelo-brankia, tipo di la familio « ostreidi », ek cirkume cent speci qui vivas en la mari varma e mi-varma, e kinacent speci fosila en la strati sekundara e terciara. – L. \textit{ostrea edulis} e \textit{gryphea angulata}\latina{ostrea edulis; gryphea angulata}\lingui{DEFIRS}
\artiklo{ostrogoto}%
\vorto{ostrogoto} Persono qua ne konformigas sua konduto a la kustumi, a la deci, qui regnas en la medio en qua lu vivas
\artiklo{-ot-}%
\vorto{-ot-}\fako{gram.} Verbo-dezinenco di la participo di futureso, pasiva
\artiklo{otalgio}%
\vorto{otalgio}\fako{patol.} Nevralgio di la orelo
\artiklo{otario}%
\vorto{otario}\fako{zool.} Speco di mar-hundo di qua la oreli salias\latina{otaria}\lingui{EFI}
\artiklo{otito}%
\vorto{otito}\fako{patol.} Inflamuro di la orelo
\artiklo{otolito}%
\vorto{otolito}\fako{biol.} Konkreciono tre mikra, kalkoza, qua existas en la aud-organaro di animali diversa
\artiklo{otomano}%
\vorto{otomano} Granda sidilo, sen dors-apogilo, sur qua onu povas repozar segun la posturo di la Orientani\lingui{DEFIRS}
\artiklo{otuso}%
\vorto{otuso}\fako{zool.} Raptucelo noktala\latina{otus}\lingui{L}
\artiklo{ovacionar}%
\vorto{ovacionar}\fako{trans.}\senco{1}{\textit{(en Roma, antique)} Ceremonio, min solena kam triumfo, dum qua onu imolis mutonino, e la generalo vinkinta avancis kavalkante sur kavalo}\senco{2}{Honoro agita a personego da la procesionani, eskortani}\lingui{DEFIRS}
\artiklo{ovario}%
\vorto{ovario}\senco{1}{\textit{(zool.)} Organo di la femino, qua kontenas la ovuli destinita produktar la feto, che la vivipari – la ovo, che la ovipari}\senco{2}{\textit{(bot.)} Parto di la pistilo, qua kontenas la ovuli destinita produktar la semini}\lingui{DEFIS}
\artiklo{ovipara}%
\vorto{ovipara} Qua produktas e fekundigas ovi de qui ekiros la yuni\lingui{DEFIS}
\artiklo{ovo}%
\vorto{ovo}\senco{1}{Maso globatra, quan pondas la ucelino, jermo di la ucelo futura, cirkondata da liquidi destinata nutror lu dum kelka tempo, ed envelopata da shelo kalkoza}\senco{2}{Jermo-ovulo che la animali vivipara qua, pos fekundigeso e developeso, pasas de la ovario aden la utero}\senco{3}{To quo havas la formo di ovo}\senco{4}{\textit{(tekn.)} Ornamento di arkitekturo, di juveli, e c. qua havas la formo di ovo}\lingui{DEFIRS}
\artiklo{ovogonio}%
\vorto{ovogonio}\fako{biol.} Celuli qui rezultas de la genito di la elementi jerma di la ovisako
\artiklo{ovoida}%
\vorto{ovoida} Di qua la formo aspektas quale ovo
\artiklo{ovulo}%
\vorto{ovulo}\senco{1}{\textit{(zool.)} Produkturo de la ovario, qua, pos fekundigeso, divenas la feto che la vivipari – e la ovo, che la ovipari}\senco{2}{\textit{(bot.)} Produkturo de la ovario qua, pos fekundigeso, divenos la semino}\lingui{EFIS}
\artiklo{oxalo}%
\vorto{oxalo}\fako{bot.} Planto supala, ek la familio « poligonei », qua saporas kelke acida, e quan onu uzas por facar supi o kom legumo\latina{rumex acetosa}\lingui{DEFIS}
\artiklo{oxidazo}%
\vorto{oxidazo}\fako{kemio biologiala} Fermento dissolvebla en aquo – kompozajo di qua un ek la kompozanti esas un molekulo metala qua igas lu apta kaptar oxo e transportar ica adsur substanco organika, oxidigebla facile
\artiklo{oxido}%
\vorto{oxido} Kompozuro de korpo simpla kun oxo, e qua ne su unionas kun la bazi\lingui{DEFIS}
\artiklo{oxigeno}%
\vorto{oxigeno}\fako{kemio} Oxo en la formo di gaso\lingui{DEFIS}
\artiklo{oxikoko}%
\vorto{oxikoko}\fako{bot.} Genero de erikacei – subarboreto reptera, di qua la frukti – beri globetatra – redatra, acido-saporoza, divenas kelke sukroza pos la komenco di la frostoperiodo, ed uzata por facar konfituro\latina{veccinium oxycoccos}
\artiklo{oxo}%
\vorto{oxo}\fako{kemio} Korpo simpla, ek la aero, ed uzata por respiro e kombusto\simbolo{O}
\artiklo{-oz-}%
\vorto{-oz-} Sufixo qua implikas la ideo : « qua naturale havas… qua kontenas naturale… » e « qua kontenas to quon indikas la radiko verba » (sen impliko di anteeso, prezenteso o futureso; ca nuanci indikesas per « -inta », « -anta », « -onta »)
\artiklo{oziero}%
\vorto{oziero}\fako{bot.} Mikra speco di saliko, di qua la rami, flexebla, uzesas por facar korbatraji\latina{salix viminalis}\lingui{EF}
\artiklo{ozono}%
\vorto{ozono}\fako{kemio} Oxigeno kondensita\simbolo{O₃ (oxigeno) O₂}
\end{multicols}
\sekciono{P}
\begin{multicols}{2}
\artiklo{pacar}%
\vorto{pacar}\fako{netrans.}\senco{1}{Ne militar}\senco{2}{\textit{(pri lando, familio, grupo de homi)} Esar en stando tala, ke ne eventas perturbesi interna – ke ne eventas deskonkordo violentoza}\lingui{DEFIS}
\artiklo{pachulio}%
\vorto{pachulio}\fako{bot.}\senco{1}{Planto labioza aromatoza, qua devenas de India, vicina (kom tipo) ye la minti}\senco{2}{La parfumo extraktita de ta planto}\latina{pogostemon patchouli}\lingui{DEFIS}
\artiklo{pacienta}%
\vorto{pacienta} Qua toleras la malaji, la varto di to quo esas *retarda (tarda), la duro di to quo plulongeskas tempale\lingui{EFIS}
\artiklo{padelo}%
\vorto{padelo} Implemento koqueyala, kun longa mancho, per qua onu fritas, frikasas, e c.\lingui{IS}
\artiklo{padloko}%
\vorto{padloko} Seruro movebla, quan onu akrochas per okulskrubi a pordo, a kesteto juvelala, por klozar olu\lingui{E}
\artiklo{pado}%
\vorto{pado} Kuseno plenigita ye buro\lingui{E}
\artiklo{paf!}%
\vorto{paf!} Interjeciono qua expresas frapo, stroko od acidento subita
\artiklo{pafar}%
\vorto{pafar}\fako{netrans.} Deskuplar la hano di fusilo, revolvero, kanono
\artiklo{pagano}%
\vorto{pagano} Qua adoras e kultas la dei quin la kristani qualifikas « falsa »\lingui{DEFIS}
\artiklo{pagar}%
\vorto{pagar}\fako{ulu, ye ulo} Quitigar su (ad ulu) pri to quon onu debas a lu. – \textit{(ulo, ad ulu)}. Livrar (ad ulu) to quon onu debas a lu (po ulo)\lingui{EFIS}
\artiklo{pagayo}%
\vorto{pagayo}\fako{navig.} Longa remilo pirogala, kun paleto simpla o duopla, quan onu uzas sen apogar lu sur la barko\lingui{DFIS}
\artiklo{pagino}%
\vorto{pagino} Singla ek la du lateri di folio sur qua onu skribas od imprimas\lingui{DEFIRS}
\artiklo{pagodo}%
\vorto{pagodo}\senco{1}{Templo konsakrita a la kulto di la idoli, en Azia}\senco{2}{Statueto Chiniana, figureto di etajero}\lingui{DEFIRS}
\artiklo{pajo}%
\vorto{pajo}\senco{1}{Puerulo qua servas princulo, princino, e c.}\senco{2}{\textit{(ludo-karti)} Figuro di qua la rango-valoro esas situita dop la rejulo e la damo}\lingui{DEFIRS}
\artiklo{paketboto}%
\vorto{paketboto}\fako{navig.} Navo extreme kapaca, movata per aquo-vaporo, petrolo (e mem per atom-forco) qua transportas, trans la oceani, pasajanti, vari, postalaji, e c.\lingui{DEFIRS}
\artiklo{pakidermo}%
\vorto{pakidermo}\fako{zool.} Klaso de mamiferi ne-ruminera, di qui la pelo esas tre dika (elefanto, rinoceroso, e c.)\lingui{DEFIS}
\artiklo{pako}%
\vorto{pako} Asemblajo ek plura kozi, interligita, envolvita ad-ensemble\lingui{DEFIRS}
\artiklo{pakotilio}%
\vorto{pakotilio}\senco{1}{Vari ne-grava quin la maristi di navo darfas kunportar por komercor, sen pagar freto-preco}\senco{2}{Sortaro de vari destinita a kambieso, a komerco, en la landi fora}\lingui{DFIS}
\artiklo{paktar}%
\vorto{paktar}\fako{netrans., kun} Agar konvenciono solena\lingui{DEFIS}
\artiklo{pala}%
\vorto{pala}\senco{1}{\textit{(aludante la vizajo)} Senkolora, terne blanka}\senco{2}{Qua esas poke lumoza, poke koloroza}\lingui{EFIS}
\artiklo{palaco}%
\vorto{palaco} Habiteyo luxoza di rejo, di princo. Domo splendida di richa privato\lingui{DEFIRS}
\artiklo{paladino}%
\vorto{paladino} Heroo di la *chevalier-epoko. – \textit{Anke metaf}\lingui{DEFIRS}
\artiklo{paladio}%
\vorto{paladio}\fako{en la epoki antiqua}\senco{1}{Statuo de Pallas, e gardita da la Troyani kom salvanto di lia urbo}\senco{2}{To quon populo egardas kom sekuriganta sua salveso}\lingui{DEFIRS}
\artiklo{palado}%
\vorto{palado}\fako{kemio} Metalo blanka, tre maleebla, poke gisebla\simbolo{Pd}\lingui{DP}
\artiklo{palankino}%
\vorto{palankino} En India, Chinia (ante Mao Tze Tung), e c., porto-lito quan servisti portas de sua shultro. (Anke en la Egiptia faraonala)\lingui{DEFIS}
\artiklo{palateno}%
\vorto{palateno}\fako{en la epoki feudala} Qua havas ofico en la palaco di suvereno\lingui{DEFIRS}
\artiklo{palatino}%
\vorto{palatino} Furo per qua la mulieri kovras sua shultri\lingui{DFIS}
\artiklo{palato}%
\vorto{palato}\fako{anat.} Parto supra di la boko-kavajo\lingui{EFIS}
\artiklo{paleografio}%
\vorto{paleografio} Cienco-fako di la persono qua su okupas dechifrar la surskriburi, la manuskripti antiqua o tre olda, la charti, la diplomi, e c.\lingui{DEFIRS}
\artiklo{paleontologio}%
\vorto{paleontologio} Historio naturala pri la animali e pri la vejetinti fosila\lingui{DEFIRS}
\artiklo{paleozoologio}%
\vorto{paleozoologio} Parto di paleontologio di qua la studiajo esas la animali fosila\lingui{DEFIS}
\artiklo{palestro}%
\vorto{palestro}\fako{en la epoki antiqua} Exerci korpala (lukto, boxo, salto, kuro, e c.)\lingui{DEFIRS}
\artiklo{paleto}%
\vorto{paleto} Parto plata di la remilo, qua frapas la aquo; o di roto navala, aquo-falala\lingui{F}
\artiklo{paliatar}%
\vorto{paliatar}\fako{trans.} Mingravigar (morbo) sen remediar lu komplete\lingui{DEFIRS}
\artiklo{palimpsesto}%
\vorto{palimpsesto}\fako{paleogr.} Manuskripto di qua la dat-unesma skriburo gratesis, skrapesis, por posibligar la skribo di nova texto (ulakaze : transverse)\lingui{DEFIRS}
\artiklo{palinodio}%
\vorto{palinodio}\fako{en la epoki antiqua} Poemo en qua onu retraktas to quon onu skribis en poemo antea\lingui{DEFIS}
\artiklo{palio}%
\vorto{palio} Stipito di planti cereala di qua onu deprenis la grani\lingui{FIS}
\artiklo{palisado}%
\vorto{palisado}\senco{1}{Klozilo ek planki, stangi, e c. – lineo de lignopeci triangula destinita ad obstruktar la fauco di konstrukturo militala – ad impedar la avanco di la enemiko ye la bordo di fosato}\senco{2}{\textit{(gardeno)} Lineo de arbori di qui onu lasas kreskar la branchi de la pedo, tale ke li formacas quaza hego}\lingui{DEFIS}
\artiklo{palisandro}%
\vorto{palisandro} Ligno violea, kun nuanci flava, nigra, quan onu povas polisar adbrile, uzata por la inkrusturi, por facar mobli\latina{dalbergia nigra, jacarania mimosifolia}\lingui{DEFIRS}
\artiklo{paliso}%
\vorto{paliso} Sproso ligna, di qua onu igis pintoza un extremajo, e de qua onu facas (ordinare) palisado\lingui{EFIS}
\artiklo{paliumo}%
\vorto{paliumo}\senco{1}{\textit{(en la epoki antiqua)} Nomo per qua la Romani signifikis la mantelo di la Greki}\senco{2}{\textit{(eklezio katol.)} Bendo de lano blanka, ornita ye quar kruci, qua cirkondas la shultri e pendas, avane e dope, quan la papo *weras o quan lu sendas a la arkiepiskopi, primati, e c. kom signo pri judicioyuro}\lingui{DEFIS}
\artiklo{palmata}%
\vorto{palmata}\fako{zool.} Di qua la fingri interligesas per membrano\lingui{EFIS}
\artiklo{palmeto}%
\vorto{palmeto} Mikra palmo, ornivo skultita, sur frisi, e c.
\artiklo{palmipedo}%
\vorto{palmipedo}\fako{zool.} Klaso de uceli di qui la pedi esas palmata (anado, ansero, e c.)\lingui{EFIS}
\artiklo{palmo}%
\vorto{palmo}\senco{1}{Brancho di palmiero, di qua la folii dispozesis quale manuo apertata}\senco{2}{La interna parto di la manuo}\lingui{DEFIRS}
\artiklo{palpar}%
\vorto{palpar}\fako{trans.} Explorar per la manuo, quale la blindi\lingui{EFIS}
\artiklo{palpebro}%
\vorto{palpebro}\fako{anat.} Singla ek la du membrani movebla qui, per interkontaktar, tegas la globo okulala\lingui{FIS}
\artiklo{palpitar}%
\vorto{palpitar}\fako{netrans.) (aludante parto di la homo-korpo} Anhelar konvulsatre. Pulsar konvulsatre\lingui{DEFIS}
\artiklo{paltoto}%
\vorto{paltoto} Vesto quan la viri *weras super frako o redingoto\lingui{DEFRS}
\artiklo{palumbo}%
\vorto{palumbo} Kolombo sovaja, qua lojas en nesto sur la arbori\lingui{FIS}
\artiklo{pamfleto}%
\vorto{pamfleto} Verketo violentoza en qua onu atakas ulu, ulo\lingui{DFR}
\artiklo{pampro}%
\vorto{pampro} La foliaro di vito
\artiklo{panaceo}%
\vorto{panaceo}\fako{medic.} Remediivo generala, unika, apta por irga morbo\lingui{DEFIRS}
\artiklo{panariso}%
\vorto{panariso}\fako{patol.} Inflamuro flegmona ye la extremajo di la fingri\lingui{EFIS}
\artiklo{pancho}%
\vorto{pancho} La maxim kapaca ek la stomaki di la rumineri\lingui{DEFIS}
\artiklo{panduro}%
\vorto{panduro} Soldato exter-armeana, privata\lingui{DEFIRS}
\artiklo{panear}%
\vorto{panear}\fako{netrans.} Ne plus povar avancar, funcionar\lingui{DFI}
\artiklo{panegiro}%
\vorto{panegiro} Diskurso publika laude di persono\lingui{DEFIRS}
\artiklo{panelo}%
\vorto{panelo}\senco{1}{\textit{(arkitekt.)} Parto plu o min granda e grosa di muro.- Ensemblo de karpenturo di qua onu plenigas la vakui ye masonuro}\senco{2}{Peco ek vesto, ek mantelo, ek robo, e c.}\lingui{DEFR}
\artiklo{pangenezo}%
\vorto{pangenezo}\fako{biol.} Teorio per qua Darwin explikas omna fenomeni biologiala, per la karakterizivi e per la migri da partikuli materia, quin lu nomizas « jemuli »\lingui{DEFIS}
\artiklo{pangeno}%
\vorto{pangeno}\fako{biol.} Segun la teorio di de Vries : mikra korpo, ek molekulo kemiala, ma dotita ye la karakterizivo di la materio vivanta : asimiliveso, kreskiveso, sureprodukto per su-divido\lingui{DEFIS}
\artiklo{pangolino}%
\vorto{pangolino}\fako{zool.} Sen-dento, di qua la korpo esas squamoza. ek Sud-Afrika ed ek India\latina{manis laticaudata}
\artiklo{paniklo}%
\vorto{paniklo} Graso qua garnisas la pelo di la porko e di kelka animali altra, precipue ye la ventro\lingui{EFIS}
\artiklo{paniko}%
\vorto{paniko} Teroro subita e sen justifikivo\lingui{DEFIRS}
\artiklo{panikulo}%
\vorto{panikulo}\fako{bot.} Sorto di infloresenco, en qua la flori di qui la pedunkli esas dividita inter-situesas ye alteso-punti inter-egala\lingui{EFIS}
\artiklo{panio}%
\vorto{panio} Peco de stofo per qua la negri mi-nuda tegas su de la talio til la genui\lingui{F}
\artiklo{pankraco}%
\vorto{pankraco}\fako{en la epoki antiqua} Exerco gimnastikala (lukto e boxo)\lingui{DEFIS}
\artiklo{pankreatito}%
\vorto{pankreatito}\fako{patol.} Inflamuro di la pankreato\lingui{EF}
\artiklo{pankreato}%
\vorto{pankreato}\fako{anat.} Glando di la abdomino, qua sekrecas suko qua esas faktoro di la digesto\lingui{EFIS}
\artiklo{panmixio}%
\vorto{panmixio}\fako{biol.} Partopreno da omna individui a la perpetuigo di la speco, kontraste kun la \textit{selekto}, segun qua nur la individui avantajoza su *reproduktas\lingui{DEFIS}
\artiklo{pano}%
\vorto{pano} Nutrivo ek maso de farino dilutita en aquo, petrisita e bakita\lingui{FIS}
\artiklo{panoramo}%
\vorto{panoramo} Pikturo sur muro cirklatra qua igas la spektanto qua esas en la centro vidar su iluzione cirkondata da objekti quin lu kredas reala. — anke metaf\lingui{DEFIRS}
\artiklo{pansar}%
\vorto{pansar}\fako{trans.) (kirurg.} Traktar vunduro per aplikar a lu la aparati necesa\lingui{F}
\artiklo{panseo}%
\vorto{panseo}\fako{bot.} Planto di la genero « violo »\latina{viola tricolor}\lingui{DFIRS}
\artiklo{pantalono}%
\vorto{pantalono} Bracho qua iras til la pedi\lingui{EFIRS}
\artiklo{panteismo}%
\vorto{panteismo} Sistemo filozofiala qua agnoskas nur un substanco di qua la esi diversa esas modi\lingui{DEFIRS}
\artiklo{panteisto}%
\vorto{panteisto} Persono qua agnoskas panteismo\lingui{DEFIRS}
\artiklo{pantero}%
\vorto{pantero}\fako{zool.} Karnivoro, ek la genero « kato », e di qua la felo esas makuletoza\latina{felix pardus}\lingui{DEFIRS}
\artiklo{pantoflo}%
\vorto{pantoflo} Pedo-vesto chambrala, di qua la peco dopa esas supresita, – lejera e komfortoza\lingui{DFS}
\artiklo{pantografo}%
\vorto{pantografo} Instrumento per qua onu povas mekanisme kopiar, reduktar, grandigar desegnuro, xilografuro, e c.\lingui{DEFIRS}
\artiklo{pantometro}%
\vorto{pantometro} Triopla linealo movebla, uzata por determinar la tri anguli di triangulo\lingui{DEFIS}
\artiklo{pantomimo}%
\vorto{pantomimo} Teatrajo en qua la roli expresas sua sentimenti nur per gesti\lingui{DEFIRS}
\artiklo{papagayo}%
\vorto{papagayo}\fako{zool.} Ucelo ek la klaso « klimeri », di qua la plumaro esas, che preske omni, bunte verda, reda, e c., kun beko harda e kurvatra, qua imitas facile la voco homala, la krii da la animali cetera\latina{psittacus}\lingui{DIRS}
\artiklo{papavero}%
\vorto{papavero}\fako{bot.} Planto narkotigiva\latina{papaver rhoeas}\lingui{EFIL}
\artiklo{papayo}%
\vorto{papayo}\fako{bot.} Frukto de arboro ek la regioni tropikala – qua aspektas e saporas kelke quale melono\latina{carica papaya}\lingui{DEFIRS}
\artiklo{papero}%
\vorto{papero} Folio por skribar od imprimar, por envelopar, e c., facita ek shifoni, papera rezidui, diversa materii vejetantala fibroza (esparto, maizo, e c.) quin onu pastigas, extensas, sekigas\lingui{DEFS}
\artiklo{papiliono}%
\vorto{papiliono}\fako{zool.} Insekto qua havas quar ali, kolorizita ye squami dina quale polvo\lingui{EFL}
\artiklo{papilo}%
\vorto{papilo}\fako{anat.} Mikra saliajo, sub epidermo o sur la surfaco di ula mukozi, ek branchifuri nerva, vaskula, qua produktas la sentiveso di la pelo\lingui{DEFIS}
\artiklo{papiro}%
\vorto{papiro} Folieto deprenata de sorto di kano qua kreskas en Egiptia, India, e c., e quan onu preparis (en la epoki antiqua) por la skribo\lingui{DEFIRS}
\artiklo{paplo}%
\vorto{paplo}\senco{1}{Farino o fekulo quan onu lasas boliar en lakto, aquo, por facar de lu nutrivo destinata precipue a la infanti}\senco{2}{\textit{(metaf.)} Shifoni boliita e reduktita a pasto, de qua facesos papero, kartono, e c.}\lingui{EIS}
\artiklo{papo}%
\vorto{papo}\fako{religio} Chefo suprega di la eklezio katolika\lingui{DEFIRS}
\artiklo{papriko}%
\vorto{papriko}\fako{bot.} Varietato di pimento, tre uzata en Hungaria\latina{capsicum}\lingui{DFL}
\artiklo{papulo}%
\vorto{papulo}\senco{1}{\textit{(patol.)} Lezuro di la pelo, karakterizata da saliajo (rozea, reda o bruna) di la pelo}\senco{2}{Protuberanco mola e plena ye liquido aquatra, sur la epidermo di ula planti}\lingui{EFIS}
\artiklo{par-}%
\vorto{par-} Prefixo qua signifikas : « komplete », kun la radiki di verbi
\artiklo{para-}%
\vorto{para-} Prefixo qua signifikas « qua protektas kontre », « shirmilo kontre »
\artiklo{parabolo}%
\vorto{parabolo}\senco{1}{\textit{(retor.)} Alegorio sub qua celesas docivo etikala, morala}\senco{2}{\textit{(geom.)} Kurvo plana di la grado duesma, qua rezultas de la seciono di kono da plano paralela ye la latero}\lingui{DEFIRS}
\artiklo{paraboloido}%
\vorto{paraboloido}\fako{geom.} Qua havas la formo di parabolo\lingui{DEFIS}
\artiklo{paradar}%
\vorto{paradar}\fako{netrans.) (aludante armeo-trupi} Defilar lor revueso\lingui{DEF}
\artiklo{paradigmo}%
\vorto{paradigmo}\fako{gram.} Tipo di deklino, di konjugo\lingui{DEFIRS}
\artiklo{paradizo}%
\vorto{paradizo}\senco{1}{\textit{(kristanismo)} Loko delicoza, gardeno ube Deo pozabis la viro dat-unesma}\senco{2}{\textit{(metaf.)} Remdsyo di la beateso cielala. — Loko, regiono, stando, extreme agreabla}\lingui{DEFIRS}
\artiklo{paradoxo}%
\vorto{paradoxo} Opiniono kontrea a la opiniono komuna\lingui{DEFIRS}
\artiklo{parafino}%
\vorto{parafino}\fako{meio} Reziduo de la distilo di la oleo minerala, substanco grasa, blanka, kristalatra, uzebla por lumifo\lingui{DEFIRS}
\artiklo{parafo}%
\vorto{parafo}\senco{1}{Plumo-stroko diversa-forma quan onu adjuntis a sua nomo por distingar sua signaturo}\senco{2}{Signaturo kurtigura quan onu trasas en la marjino di akti por aprobar la surstrekizuri}\lingui{DEFI}
\artiklo{parafrazar}%
\vorto{parafrazar}\fako{trans.} Developar texto por explikar lu\lingui{DEFIRS}
\artiklo{paragrafo}%
\vorto{paragrafo} Segmento di chapitro\lingui{DEFIRS}
\artiklo{paralaxo}%
\vorto{paralaxo}\fako{astron.} Angulo formacita, ye la centro di astro, da du rekti qui abutas a du observanti lokizita en du punti interdiferanta\lingui{DEFIRS}
\artiklo{paralela}%
\vorto{paralela}\fako{geom.) (aludante rekti, surfaci} Di qui la interdisto esas egala irge-punte\lingui{DEFIRS}
\artiklo{paralelepipedo}%
\vorto{paralelepipedo}\fako{geom.} Solido qua finas per sis paralelogrami inter-egala, qui inter-opozas duope\lingui{DEFIS}
\artiklo{paralelogramo}%
\vorto{paralelogramo}\fako{geom.} Quadrilatero di qua la lateri inter-opozata esas egala e paralela\lingui{DEFIRS}
\artiklo{paralizar}%
\vorto{paralizar}\fako{trans.}\senco{1}{\textit{(patol.)} Produktar la nepovo, tota o parta, pri sive movar e sentar, sive movar o sentar}\senco{2}{\textit{(metaf.)} Impedar la funciono normala}\lingui{DEFIRS}
\artiklo{paralogismo}%
\vorto{paralogismo} Rezono nekorekta\lingui{DEFIS}
\artiklo{paramento}%
\vorto{paramento} Speco de stofo, simpla od ornita, adjuntita a la extremajo di la maniko di vesto, uniformo, robo\lingui{FIS}
\artiklo{parametro}%
\vorto{parametro}\fako{matem.} Konstanto ek la equaciono, ek la konstrukto di kurvo, ek surfaco, qua posibligas la determino di lua dimensioni\lingui{DEFIS}
\artiklo{paranimfo}%
\vorto{paranimfo}\fako{en la epoki antiqua} Amiko di la quik-mariajonto, qua helpis lu pri la ceremoniaro di la mariajo\lingui{DEFIS}
\artiklo{paranukleino}%
\vorto{paranukleino}\fako{kemio biol.} Substanco ye qua konsistas la nukleoli di la celulo-nukleo
\artiklo{parapeto}%
\vorto{parapeto}\senco{1}{Muro, ye la nivelo di la kudo, alonge la bordo di ponto, kayo, teras-tekto, por impedar la neprudenti falar aden precipiso, abismo, maro, e c.}\senco{2}{Masivo ek tero o masonuro qua quaze kronizas remparo, e shirmas kontre la pafi de la enemiko}\lingui{DEFIRS}
\artiklo{paraplasmo}%
\vorto{paraplasmo}\fako{kemio biol.} Substanco vitratra, mi-fluida, qua konsistas ye globulino, ed en qua imersesas la elementi di citoplasmo\lingui{DEFIS}
\artiklo{parazito}%
\vorto{parazito}\senco{1}{\textit{(en la epoki antiqua)} Repasto-gasto singladia di richo, qua havis kom tasko amuzar ica}\senco{2}{Ta qua mestiere su invitigas por repastar}\senco{3}{Planto od insekto qua vivas ek altra animalo od altra planto}\lingui{DEFIRS}
\artiklo{parcelo}%
\vorto{parcelo} Ek ensemblo, mikra parto di qua la formo esas determinita precize\lingui{DEFIS}
\artiklo{pardonar}%
\vorto{pardonar}\fako{trans.}\senco{1}{\textit{(ulo ad ulu – o : ulu pri ulo)} Remisar (kulpo.)}\senco{2}{\textit{(en la formulo di politeso)} Indulgar}\lingui{EFIS}
\artiklo{parear}%
\vorto{parear}\fako{trans.} Evitar atingeso, per deviacigar (la frapostroko)\lingui{DEFIS}
\artiklo{parenkimo}%
\vorto{parenkimo}\fako{anat.}\senco{1}{Tisuo sponjatra, qua karakterizas la viceri}\senco{2}{\textit{(bot.)} Tisuo utrikula di la planti}\lingui{DEFIRS}
\artiklo{parentezo}%
\vorto{parentezo}\senco{1}{\textit{(gram.)} Frazo acesora, komplementa, e, kom tala, provizita ye signo partikulara qua izolas lu interne di la frazo precipua}\senco{2}{Singla de la signi : (), qui indikas parentezo}\lingui{DEFIS}
\artiklo{parento}%
\vorto{parento} Persono qua esas membro di la familio (di ulu) per nasko o mariajo\lingui{EFIS}
\artiklo{parfumo}%
\vorto{parfumo}\senco{1}{Odoro agreabla, naturala}\senco{2}{Substanco qua igas odorar tale}\lingui{DEFIRS}
\artiklo{parhelio}%
\vorto{parhelio}\fako{astron.} Fenomeno lunala qua produktesas da la reflekteso di lumo sur la kristali glacia qui suspendesas en atmosfero\lingui{DEFIS}
\artiklo{paria}%
\vorto{« paria »} En India, persono qua apartenas a la kasto maxim infra e qua, kom tala, esas la objekto di desestimo
\artiklo{pariar}%
\vorto{pariar}\fako{trans., pri, por, kontre} Konvencionar pri pekunio-quanto riskenda (generale, da singla koncernato) por ta qua esabos justa pri kozo kontestata\lingui{DFR}
\artiklo{parietario}%
\vorto{parietario}\fako{bot.} Planto ek la familio « urticei », qua kreskas sur la muregi\latina{parietaria}\lingui{DEFIS}
\artiklo{parieto}%
\vorto{parieto}\senco{1}{Separilo ek masonuro lejera, menuzuro, quan onu erektas en domo, apartamento, e c.}\senco{2}{To quo dividas kavajo aden du o plusa parti}\lingui{EFIS}
\artiklo{parko}%
\vorto{parko} Extensajo ek sulo, cirkondata da muro, greto, e c., qua dependas generale de kastelo, de domego, ed uzesas sive por konservar vildo, sive por plubeligar rezideyo, sive por esar promeneyo publika\lingui{DEFIRS}
\artiklo{parlamento}%
\vorto{parlamento}\senco{1}{Asemblitaro ek personi qui deliberas}\senco{2}{En Britania : Asemblitaro ek la baroni, prelati e deputiti di la urbi, qua helpas la rejo en lua praktiko di la povo suprega}\senco{3}{Nomo di la asemblitaro, di la « chambri », qui praktikas la povo suprega en la landi konstitucala}\lingui{DEFIRS}
\artiklo{parlementar}%
\vorto{parlementar}\fako{netrans.} Konferar, negociar, espere di trans-akto, koncilio\lingui{DEFIRS}
\artiklo{paro}%
\vorto{paro}\senco{1}{Duo de kozi od animali, sama-speca, destinita irar kune}\senco{2}{La maskulo e la femino di animali sama-speca, un de singlu}
\artiklo{parodiar}%
\vorto{parodiar}\fako{trans.} Imitar, deformante (lu) burleske, verko grava\lingui{DEFIRS}
\artiklo{paroko}%
\vorto{paroko}\fako{religio} Sacerdoto qua direktas distrikto ekleziala, che la katoliki\lingui{DEFIS}
\artiklo{parolar}%
\vorto{parolar}\senco{1}{\textit{(netrans.)} Uzar la artikulo-linguo}\senco{2}{Expresar sua penso per la artikulo-linguo}\senco{3}{\textit{(trans.)} Uzar, kom expresilo (ta o ca linguo)}\lingui{FIS}
\artiklo{paromo}%
\vorto{paromo} Batelo plata, uzata por transportar, sur rivero o lago o fluvio, de ca a ta rivo,-homi, animali, veturi, e quan onu glitigas alonge kablo tensita inter la du rivi, oblique ye la fluo\lingui{F}
\artiklo{paronimo}%
\vorto{paronimo}\fako{gram.} Preske homonimo\lingui{DEFIRS}
\artiklo{paronomiazo}%
\vorto{paronomiazo}\fako{retor.} Interproximigo, en frazo, di vorti qui pronuncesas inter-same, ma di qui la senci interdiferas
\artiklo{parotido}%
\vorto{parotido}\fako{anat.} Glando salivifera, situita apud la orelo\lingui{EFIS}
\artiklo{paroxismo}%
\vorto{paroxismo}\fako{patol.} Intenseso maxima di aceso\lingui{DEFIRS}
\artiklo{parqueto}%
\vorto{parqueto} Asemblajo ek lami de ligno harda, polisita, dispozita diverse, de qui onu facas la planko-suli eleganta\lingui{DEFR}
\artiklo{partenero}%
\vorto{partenero} La persono kun qua onu ludas kontre altra ludanti\lingui{DEFR}
\artiklo{partenogenezo}%
\vorto{partenogenezo} Modo di genito ye qua femino *reproduktas sua speco sen fekundigesir da maskulo\lingui{DEFIRS}
\artiklo{partenokarpio}%
\vorto{partenokarpio}\fako{biol.} Nomo quan Noll propozis en 1902 por indikar la modo di su-developo che ula frukti
\artiklo{partero}%
\vorto{partero}\senco{1}{Ta parto di gardeno ube la sulo esas tre plata e glata, sen-arbora, e dividita per bedi}\senco{2}{\textit{(teatro)} Ta parto di spektaklo-chambrego qua esas inter la orkestreyo e la fundo di la teatro ube, olim, la spektanti permanis stacanta}\lingui{DFIR}
\artiklo{particiono}%
\vorto{particiono}\fako{muziko} Asembluro de omna parti di kanto od orkestro\lingui{DEFIRS}
\artiklo{participo}%
\vorto{participo}\fako{gram.} Vortospeco qua esas parte verba e parte adjektiva – modo ne-personala di la verbo, qua expresas la ago forme di adjektivo\lingui{DEFIS}
\artiklo{partikulara}%
\vorto{partikulara} Qua esas la proprajo exkluziva di ulu, di ulo\lingui{DEFIRS}
\artiklo{partikulo}%
\vorto{partikulo}\senco{1}{Parto mikrega di korpo}\senco{2}{\textit{(gram.)} Mikra (kurta) vorto, generale du-silaba, nevarianta, qua uzesas por la kompozo di altra vorti}\senco{3}{La partikulo di nobeleso : la prepoziciono qua esas situita avan (ante) la nomo di multa familii nobela}\lingui{DEFIS}
\artiklo{partio}%
\vorto{partio} Inter-amuzo, inter-distrakto da grupo de personi\lingui{DEFS}
\artiklo{partiso}%
\vorto{partiso} Grupo de personi qui adoptis (od adheris a) la sama principi di konduto\lingui{DEFIRS}
\artiklo{partitivo}%
\vorto{partitivo}\fako{gram.} Qua indikas parto ek toto\lingui{DEFIS}
\artiklo{partizano}%
\vorto{partizano} Olima shaft-armo, di qua la lamo esis plu larja e plu akuta kam olta di la halbardo ordinara\lingui{DEFIS}
\artiklo{parto}%
\vorto{parto}\senco{1}{La porciono di kozo qua apartenas a singla persono kande onu dividas lu inter plura personi}\senco{2}{To quo (bona o mala) apartenas ad ulu, en kozo en qua li havas interesto}\senco{3}{Un de la elementi ye qui konsistas loko}\senco{4}{Elemento ek toto}\senco{5}{Elemento qua kontributas ad ensemblo}\senco{6}{\textit{(muziko)} To quon singla voco havas kantenda – singla instrumento, pleenda, en ensemblo}\lingui{DEFIS}
\artiklo{parturar}%
\vorto{parturar}\fako{trans.) (aludante patrino} Emisar de su (filio konceptita e qua atingis la nasko-instanto)\lingui{EFIS}
\artiklo{paruo}%
\vorto{paruo}\fako{zool.} Mikra ucelo, ek la klaso « paseri », gracila, di qua la bela plumaro esas bunta\latina{parus}\lingui{SL}
\artiklo{pasabla}%
\vorto{pasabla} Qua esas aceptebla o tolerebla, quankam minima qualesale o quantesale\lingui{DEFIS}
\artiklo{pasajar}%
\vorto{pasajar}\fako{netrans.} Voyajar per navo, pos lokacir sua plaso dure di la voyajo e pagir po lo\lingui{DEFIRS}
\artiklo{pasamano}%
\vorto{pasamano} Balustrado qua garnisas la bordo di eskalero, uzebla kom apogo-punto da la personi qui acensas o decensas\lingui{IS}
\artiklo{pasar}%
\vorto{pasar}\fako{netrans.}\senco{1}{Durar sua iro, sua fluo}\senco{2}{Venante de… abutar a…}\lingui{DEFIS}
\artiklo{pasero}%
\vorto{pasero} Ucelo kono-beka, kun plumaro griza, tre ordinara en west-Europa\latina{passer}\lingui{FIL}
\artiklo{pasha}%
\vorto{« pasha »}\fako{olim} Oficiero Turka, komisita por administrar provinco
\artiklo{pasifloro}%
\vorto{pasifloro}\fako{bot.} Planto de Amerika, di qua la semini saporas quale olti di grenado\lingui{DEFI}
\artiklo{pasiono}%
\vorto{pasiono}\senco{1}{Movo violentema, impetuoza, di la psiko ad to quon lu deziras}\senco{2}{Amoro violentema}\lingui{DEFIS}
\artiklo{pasivo}%
\vorto{pasivo}\senco{1}{\textit{(komerco)} To quon onu debas}\senco{2}{\textit{(gram.)} La voco qua indikas ke la subjekto subisas la ago quan expresas la verbo.}\lingui{DEFIRS}\subvorto{}{pasiva}{Qua subisas la ago (da ulu, da ulo)}{}{}
\artiklo{pasko}%
\vorto{pasko}\fako{liturgio katol.}\senco{1}{Festo solena, memorige di la rezurekto di Kristo}\senco{2}{\textit{(liturgio judala)} Festo solena, celebrata memorige di la ekiro de Egiptia}\lingui{EFIRS}
\artiklo{pasmento}%
\vorto{pasmento} Interplekturo de fili ora, arjenta, silka – facita sur kuseno, per spindeli e pingli, uzata kom bordumo o kom aplikajo sur vesti, mobli, e c.\lingui{DEFIRS}
\artiklo{paspelo}%
\vorto{paspelo} Bendo ek silko, drapo, per qua onu garnisas la bordo di ula parti di vesto, di uniformo\lingui{DF}
\artiklo{pasporto}%
\vorto{pasporto} Dokumento emisata da Stato qua igas sua regnati darfar voyajar libere, e sekurigas a li lia protekteso da la reprezentisti diplomacala di ta stato en la landi cetera\lingui{EFIRS}
\artiklo{pastelo}%
\vorto{pastelo} Krayono ek farbi triturita e pastigita per gum-aquo\lingui{DEFIRS}
\artiklo{pasterno}%
\vorto{pasterno} Infra parto di la gambo di kavalo o di bovo, ube onu fixigas lige la entravilo qua retenas lu kande lu pasturas\lingui{EFI}
\artiklo{pasteto}%
\vorto{pasteto} Kukajo qua kontenas karno, fisho, e c., koquita en pasto\lingui{DEFIRS}
\artiklo{pasteurisar}%
\vorto{pasteurisar}\fako{trans.} Steriligar (liquido) per varmigar lu til (segun lua naturo) 55o/75° C. dum kelka tempo, e koldigar lu bruske. (Onu pasteurisas lakto, kom ex.)\lingui{DEFIRS}
\artiklo{pastilo}%
\vorto{pastilo} Bonbono qua havas la formo di mikra disko, ek sukro aromatizita, ek chokolado, e c.\lingui{DEFIRS}
\artiklo{pastinako}%
\vorto{pastinako}\fako{bot.} Planto leguma, di qua la radiko-bulbo esas sukroza ed odoroza\latina{pastinaca}\lingui{DIRSL}
\artiklo{pasto}%
\vorto{pasto}\senco{1}{Farino dilutita e petrisita por facar pano, galeto, e c.}\senco{2}{Nomo di preparajo diversa en sukrajerii, apoteki, parfumifeyi, e c.}\lingui{DEFIS}
\artiklo{pastoro}%
\vorto{pastoro}\senco{1}{\textit{(olim)} La persono qua edukas e gardas bestio-trupi}\senco{2}{\textit{(religio)} Quaza sacerdoto di la religio protestantista}\lingui{DEFIRS}
\artiklo{pasturar}%
\vorto{pasturar}\fako{netrans.) (aludante bestii herbovira} Manjar herbo nefalchita, glani jus-falinta, e c.\lingui{EFIRS}
\artiklo{patar}%
\vorto{patar}\fako{shako-ludo} Agar ta stroko di la shak-ludo en qua la rejulo, koaktate ludar, povas lo nur per faliigar su\lingui{DF}
\artiklo{patato}%
\vorto{patato}\fako{bot.} Planto di qua la radiko tuberkula esas manjebla\latina{ipornoea batatas}\lingui{DEFIRSL}
\artiklo{patelo}%
\vorto{patelo}\fako{anat.} Mikra osto plata, di qua la anguli esas rondatra, qua formacas la parto avana di genuo\lingui{EL}
\artiklo{pateno}%
\vorto{pateno}\fako{liturgio katol.} Vazo sakra, qua havas la formo di pladego ronda, destinita kovrar la kalico, recevor la peceti di la hostio, e quan la oficianto prizentas por labio-froto a la personi qui iras por la donaco\lingui{DEFIS}
\artiklo{patento}%
\vorto{patento} Dokumento emisita da guvernistaro e qua grantas ula yuri.\lingui{DEFIRS}\subvorto{}{invento-patento}{Dokumento qua konstatas ke ulu revendikas la proprieto di inventuro, di perfektiguro, kom olua autoro}{}{}
\artiklo{pater}%
\vorto{« pater »}\fako{eklezio katolika} Prego sundiala qua komencas de ta vorto Latina
\artiklo{patero}%
\vorto{patero} Peco de metalo, de ligno, plu o min ornita, ye formo kurva, e muntita sur muro per lua quaza pedo, e de qua onu pendigas kurteni, vesti, e c.\lingui{EFIS}
\artiklo{patetika}%
\vorto{patetika} Qua agitas forte la sentimenti, la pasioni\lingui{DEFIRS}
\artiklo{patibulo}%
\vorto{patibulo} Asembluro de fosto vertikala e trabo horizontala, ye la extremajo di qua povas glitar la kordo destinita suspendor la homo qua kondamnesis a mortigeso per pendeso\lingui{F}
\artiklo{patino}%
\vorto{patino}\senco{1}{Oxidizo verdatra qua eventis, pos longa tempo, sur la statui, la vazi e la medalii antiqua de bronzo}\senco{2}{Nuanco uniforma, quaza polisuro, quan tempo donas a la statui, la pikturi, la objekti ivora, e c.}\lingui{DEFIS}
\artiklo{patogena}%
\vorto{patogena}\fako{patol.} Qua genitas morbo
\artiklo{patogenezo}%
\vorto{patogenezo}\fako{patol.} Exploro di la mekanismo per qua la kauzi morbifera efikas a la organismo vivoza por maladigar olca
\artiklo{patologio}%
\vorto{patologio} Parto di medicino, qua studias la morbi, traktata kom tala\lingui{DEFIRS}
\artiklo{patoso}%
\vorto{patoso}\fako{retor.} Uzo exajerita di la moyeni apta agitar forte la pasioni\lingui{DEFR}
\artiklo{patriarko}%
\vorto{patriarko}\senco{1}{\textit{(biblo)} Familio-chefo di qua la vivo esis tre longa}\senco{2}{\textit{(metaf.)} Oldo veneracinda, cirkondata da multa decendanti}\senco{3}{\textit{(historio di la eklezio kristana)} Episkopo di la episkopeyo komencala}\senco{4}{Chefo di la eklezio Greka}\lingui{DEFIRS}
\artiklo{patrico}%
\vorto{patrico}\fako{en Roma, antique} Oficiero alta-klasa, en la periodi *ultima di la imperio Romana\lingui{DEFIRS}
\artiklo{patrimonio}%
\vorto{patrimonio} Havajo, domeno, quan onu heredis de la ancestri patrulala o patrinala\lingui{DEFIRS}
\artiklo{patrio}%
\vorto{patrio} Lando ube onu naskis ed a qua onu apartenas kom civitano\lingui{FIS}
\artiklo{patriota}%
\vorto{patriota} Devota por sua patrio\lingui{DEFIRS}
\artiklo{patristiko}%
\vorto{patristiko}\fako{religio katolika} Cienco pri la doktrino di la patri ekleziala
\artiklo{patro}%
\vorto{patro}\senco{1}{La persono qua genitis filio}\senco{2}{\textit{(fig.)} Titulo ye veneraceso, respekteso}\senco{3}{Membro di kongregaciono religiala (katolika)}\lingui{EFIS}
\artiklo{patroliar}%
\vorto{patroliar}\fako{netrans.) (aludante detachmento} Cirkular por la sekureso di urbo, di kampeyo, por verifikar kad omno esas en bona stando\lingui{DEFIRS}
\artiklo{patrologio}%
\vorto{patrologio}\fako{eklezio katolika} Cienco qua traktas la verki, la biografio di la patri ekleziala\lingui{DEFIS}
\artiklo{patrono}%
\vorto{patrono}\senco{1}{\textit{(en Roma, dum la epoki antiqua)} Persono richa, potenta, cirkum qua su grupigis klienti}\senco{2}{La mastro di establisuro komercala, industriala, che qua salario-laboras kontor-employati, fabrik-employati}\lingui{DEFIRS}
\artiklo{pauperismo}%
\vorto{pauperismo} La kalamitato « povreso », en stato\lingui{DEFRS}
\artiklo{pauzar}%
\vorto{pauzar}\fako{netrans.}\senco{1}{Interruptar la agado}\senco{2}{\textit{(muziko)} Cesar la kanto, la pleo instrumentala, dure di noto kompleta}\lingui{DEFIRS}
\artiklo{pavano}%
\vorto{pavano} Olima (16esma-yarcenta) danso, segun ritmo binara (du o quar tempi) di qua la karaktero esis grava\lingui{DEFIS}
\artiklo{paviliono}%
\vorto{paviliono}\senco{1}{Edifico qua superstacas la konstrukturi cetera per la alteso di la tekto-karpenturo}\senco{2}{Sorto di tendo, ronda o quadrata, qua finas ye la somito per pinto, ed uzata por kampar}\lingui{DEFRS}
\artiklo{pavo}%
\vorto{pavo} Singla de la bloki de greso, granito, e c., kubigita per talio, quin onu asemblas por igar plu rezistiva la sulo di choseo, di voyo, di strado, di voyeto\lingui{EF}
\artiklo{pavono}%
\vorto{pavono}\fako{zool.} Ucelo, ek la klaso « galinacei », di qua la plumaro esas tre briloza, e la kaudo, ek longa plumi, qui havas, ye sua extremajo, makuli briloza okul-forma\lingui{FIRS}
\artiklo{pavorar}%
\vorto{pavorar}\fako{netrans.} Subisar ta emoco desagreabla maxime, angoriganta, quan produktas bruiso neexpektita, la aparo bruska di ento o di objekto timenda o nekonocata\lingui{FIS}
\artiklo{pazar}%
\vorto{pazar}\fako{netrans.} Avancar per la movo quan agas persono, animalo, qua pozas sua pedo ica, avan la altra\lingui{EFIS}
\artiklo{pazigrafiar}%
\vorto{pazigrafiar}\fako{trans.} Skribar per sistemo de signi komprenebla da omna personi, irgequan esas lia linguo\lingui{DEFIS}
\artiklo{peano}%
\vorto{peano}\fako{en la epoki antiqua, en Grekia}\senco{1}{Kanto solena, honore di Apolono o di altra deaji}\senco{2}{Kanto di milito, di vinko, di festo o di gayeso}\lingui{DEF}
\artiklo{peccavi}%
\vorto{« peccavi »}\fako{eklezio katolika} Konfeso di la peko agita
\artiklo{pecho}%
\vorto{pecho} Substanco rezinoza, bituminoza, devenanta generale de planti konifera\lingui{DEFIRS}
\artiklo{peco}%
\vorto{peco}\senco{1}{Parto di korpo solida ruptita, tranchita, e c.}\senco{2}{Parto extraktita de penso-verko}\lingui{EFIS}
\artiklo{pedagogio}%
\vorto{pedagogio} Cienco pri la eduko di la pueri\lingui{DEFIRS}
\artiklo{pedagogo}%
\vorto{pedagogo} La persono qua instruktas ed edukas pueri\lingui{DEFIRS}
\artiklo{pedanto}%
\vorto{pedanto} La persono qua ostentas sua savo\lingui{DEFIRS}
\artiklo{pederasto}%
\vorto{pederasto} Viro qua solicitas, por satisfacar sua bezono koitala, homo sama-sexua\lingui{DEFIRS}
\artiklo{pedikulo}%
\vorto{pedikulo}\senco{1}{\textit{(patol.)} Parto streta qua suportas ula tumori}\senco{2}{\textit{(historio naturala)} Suportilo longa, gracila (che la fungo, likeno, ovulo, e c.)}\lingui{EFIS}
\artiklo{pedikuloso}%
\vorto{pedikuloso}\fako{patol.} Ensemblo de la manifesti morbala, produktita da lausi sur korpo-parto, o mem sur la tota korpo, per granda quanto de lausi\lingui{DEFIS}
\artiklo{pedo}%
\vorto{pedo}\senco{1}{Parto artikoza, ye la extremajo di la gambo, che homo od animalo, qua, pozante su sur sulo, posibligas lua susteneso e marcho}\senco{2}{Fundamento per qua objekto es pozita sur sulo; singla de la suportili di moblo}\senco{3}{Mezuro prozodiala (mezur-unajo di la versi)}\lingui{DEFIRS}
\artiklo{pedonklo}%
\vorto{pedonklo}\senco{1}{\textit{(bot.)} Suportilo di floro. – Kaudo di frukto (kom ex. : che cerizo)}\senco{2}{\textit{(zool.)} Prolonguro di la bulbo spinala}\lingui{EFIS}
\artiklo{pegazo}%
\vorto{pegazo}\senco{1}{\textit{(mitol.)} Kavalo aloza qua spricigis de la monto Helikon (Grekia) la fonto di Hippokrenes ube onu cherpis la inspirado poeziala}\senco{2}{\textit{(metaf.)} Simbolo di la inspirado poeziala}\lingui{DEFIRS}
\artiklo{pego}%
\vorto{pego}\fako{zool.} Ucelo ek la klaso « klimeri », qua havas beko rekta, anguloza, e qua, per perforar la kortico di la arbori, kaptas e devoras la insekti qui minas ta arbori\latina{picus}
\artiklo{peizajo}%
\vorto{peizajo}\senco{1}{Parto di lando, plu o min vasta, e plu o min piktinda}\senco{2}{\textit{(metaf.)} Pikturi qua reprezentas tala parto}\lingui{FIRS}
\artiklo{pejorativa}%
\vorto{pejorativa}\fako{gram.} Qua donas senco desfavoroza, desprizigiva\lingui{DEFI}
\artiklo{pekar}%
\vorto{pekar}\fako{netrans.) (religio krist.} Kulpar kontre la lego de Deo, kontre la lego religiala\lingui{EFIS}
\artiklo{pekario}%
\vorto{pekario}\fako{zool.} Mamifero di Amerika, kelke simila a porko\latina{dicotylus torquatus}\lingui{DEFIS}
\artiklo{pekino}%
\vorto{pekino} Stofo de silko en qua strio satinatra alternas kun strio matida
\artiklo{pektar}%
\vorto{pektar}\fako{trans.} Desintrikar, netigar, akomodar (sua hararo, sua barbo) per instrumento qua havas denti de materio korna, skalia, ivora, metala\lingui{FIS}
\artiklo{pektino}%
\vorto{pektino}\fako{biol.} Substanco neutra, quan kontenas la frukti matura\lingui{DEFIRS}
\artiklo{pektoro}%
\vorto{pektoro}\fako{anat.} Parto di la korpo, qua kontenas la pulmoni e la kordio\lingui{EFIS}
\artiklo{pektoso}%
\vorto{pektoso}\fako{biol.} Kompozajo qua naskas en la tisui vejetantala, qua kontributas (kun kutoso e vaskuloso) ad aglutinar la celuloso-fibri\lingui{DEFIRS}
\artiklo{pekunio}%
\vorto{pekunio} Omna speco di moneto (metala, papera, e c.)\lingui{DEFIRS}
\artiklo{pelagika}%
\vorto{pelagika}\fako{zool.} Qua vivas en maro larja, litoro-fore\lingui{DEFIS}
\artiklo{pelagro}%
\vorto{pelagro}\fako{patol.} Morbo di pelo, qua divenas kakexio, e di qua la kauzo esas la indijo di vitamini en la nutrivi\lingui{DEFIS}
\artiklo{pelargonio}%
\vorto{pelargonio}\fako{bot.} Genero de « geraniacei »\latina{pelargonium}\lingui{DEFIS}
\artiklo{pelerino}%
\vorto{pelerino} Granda kolumo jacanta (manteleto) qua tegas la shultro e la pektoro\lingui{DEFIR}
\artiklo{pelikano}%
\vorto{pelikano}\fako{zool.} Ucelo aquala, kun beko larja, di qua la mandibulo esas garnisita ye granda posho membrana aden qua lu rezervas sua peskajo\latina{pelecanus}\lingui{DEFIRS}
\artiklo{pelikulo}%
\vorto{pelikulo}\senco{1}{Pelo, membrano, tre dina}\senco{2}{Membrano tenua, dina, qua formacesas sur la surfaco di ula substanci (jelei, lakto boliita)}\lingui{EFIS}
\artiklo{peliso}%
\vorto{peliso} Mantelo vatizita\lingui{EFIS}
\artiklo{pelmelo}%
\vorto{pelmelo} Desordino maxima, qua memorigas kaoso\lingui{EF}
\artiklo{pelo}%
\vorto{pelo}\fako{anat.} Membrano extera qua tegas omna parti di la korpo di homo o di animalo\lingui{DFIS}
\artiklo{pelorio}%
\vorto{pelorio}\fako{bot.} Anomalajo qua konsistas ye to, ke korolo normala ne-segunregula divenas segun-regula. (Digitalo esas exemplo pri lo)\lingui{DEFIS}
\artiklo{peloto}%
\vorto{peloto} Mikra amaso qua havas la formo di baloneto. (Peloto de nivofloki, de lino-filo, de lano-filo)\lingui{FS}
\artiklo{pelotono}%
\vorto{pelotono}\fako{armeo} Mikra trupo de soldati, e, partikulare, subdividuro di kompanio, di eskadrono\lingui{DEFIS}
\artiklo{pelvo}%
\vorto{pelvo}\fako{anat.} Quaza zono osta, qua konstitucas la bazo di la torso, e konsistas ye la du osti hanchala\lingui{EFIS}
\artiklo{pemikano}%
\vorto{pemikano} Preparajo de karno sikigita e kondensita\lingui{DEFIS}
\artiklo{penacho}%
\vorto{penacho} Fasko de plumi, (maxim-multa-kaze) koloroza, interklemita infre, fluktuanta supre\lingui{EFIS}
\artiklo{penar}%
\vorto{penar}\fako{netrans.} Facar granda esforco qua fatigas multe\lingui{EFIS}
\artiklo{penato}%
\vorto{penato}\fako{en Roma, en la epoki antiqua} Singla de la deaji qui protektis la hemo, e quin onu kultas che su\lingui{DEFIRS}
\artiklo{pendar}%
\vorto{pendar}\senco{1}{\textit{(trans.)} Ligar per la parto supra, diste de sulo}\senco{2}{\textit{(netrans.)} Ligar a jibeto, per kordo, cirkondante la kolo por mortigar per strangulo}\senco{3}{\textit{(netrans.)} Esar ligita per la parto supra, diste de sulo}\lingui{EFS}
\artiklo{pendentivo}%
\vorto{pendentivo}\fako{arkitekt.} Parto de vulto, inter la granda arki qui suportas kupolo\lingui{DEF}
\artiklo{pendulo}%
\vorto{pendulo} Stango tenua, o filo, finanta per maso, qua su movas ye plano vertikala, e di qua la ocili esas izokrona\lingui{DEFIRS}
\artiklo{penetrar}%
\vorto{penetrar}\senco{1}{\textit{(netrans.)} Enirar transirante to quo esas obstaklo}\senco{2}{\textit{(trans.)} Enirar kelke profunde aden (ulo)}\lingui{DEFIS}
\artiklo{peninsulo}%
\vorto{peninsulo}\fako{geogr.} Granda lando, cirkondata da aquo an omna lateri minus una per qua lu juntesas a la kontinento\lingui{DEFIS}
\artiklo{peniso}%
\vorto{peniso}\fako{anat.} Organo cilindroida, vaskuloza, rigideskiva, uzata por la kopulaco e koito che la mamiferi\lingui{DEFIS}
\artiklo{penitencar}%
\vorto{penitencar}\fako{netrans.) (religio}\senco{1}{Repentar pri pekir}\senco{2}{Expiacar pro pekir}\lingui{EFIS}
\artiklo{pensar}%
\vorto{pensar}\fako{trans. e netrans.) (ulo, ad ulo, pri ulu od ulo}\senco{1}{Uzar sua spirito por konceptar, judikar, ulo}\senco{2}{Uzar sua fakultato konceptar, judikar. Posedar ta fakultato}\lingui{EFIS}
\artiklo{pensionar}%
\vorto{pensionar}\fako{trans.} Pagar sen-intermanke, sive a stat-employato, sive por lojeyo, sive por nutreso, sive por edukeso od instrukteso\lingui{DEFIRS}
\artiklo{penta-}%
\vorto{penta-} Sufixo qua signifikas « qua havas kin… »
\artiklo{pentaedro}%
\vorto{pentaedro}\fako{geom.} Solido qua havas kin edri\lingui{DEFIS}
\artiklo{pentagono}%
\vorto{pentagono}\fako{geom.} Figuro geometriala qua havas kin anguli e kin lateri\lingui{DEFIS}
\artiklo{pentagramo}%
\vorto{pentagramo} Figuro talismana qua reprezentas stelo kin-pinta\lingui{DEF}
\artiklo{pentametro}%
\vorto{pentametro} Verso de kin pedi\lingui{DEFIRS}
\artiklo{pentatlo}%
\vorto{pentatlo}\fako{en la epoki antiqua} Asemblajo de la kin exerci atletala (lukto, kuro, salto, lanso di disko, lanso di javelino)\lingui{DEFIS}
\artiklo{pentekosto}%
\vorto{pentekosto}\fako{religio krist.}\senco{1}{\textit{(che la Judi)} Festo memorige di la lego donita sur la monto Sinai, celebrata sep semani pos la Pasko-dio duesma, e dum qua onu ofris a Deo la premici di la rekolto}\senco{2}{\textit{(che la kristani)} Festo celebrata en la sundio sepesma pos Pasko, memorige di la decenso di la Santa Spirito sur la apostoli}\lingui{EFIS}
\artiklo{pento}%
\vorto{pento} Direciono di plano, qua formacas angulo obliqua ye horizontalo\lingui{FIS}
\artiklo{peonio}%
\vorto{peonio}\fako{bot.} Planto ranunkulacea, kun granda flori intense reda, bunta o blanka\latina{paeonia}\lingui{DEFIS}
\artiklo{pepito}%
\vorto{pepito}\fako{mineral.} Mikra maso de metalo quan onu renkontras nature sen gango\lingui{DFIS}
\artiklo{pepsino}%
\vorto{pepsino}\fako{kemio} Substanco efikiva di la suko stomakala, qua dissolvas la materii albuminoida\lingui{DEFIRS}
\artiklo{peptono}%
\vorto{peptono}\fako{farmacio} Produkturo de la digesto artificala di karno da pepsino o da pankreatino, en kondicioni determinita di medio e di temperaturo\lingui{DEFIRS}
\artiklo{per}%
\vorto{per} Prepoziciono qua indikas la instrumento di la ago, to quo uzesas por produktar lu, la moyeno
\artiklo{perceptar}%
\vorto{perceptar}\fako{trans.} Kaptar, komprenar, la donataji de la sensi\lingui{DEFIS}
\artiklo{perchar}%
\vorto{perchar}\fako{netrans.) (aludante la uceli} Okupar kustume plaso sur brancho di arboro, sur stango, sur paliso\lingui{EFS}
\artiklo{perdar}%
\vorto{perdar}\fako{trans.}\senco{1}{Des-aquirar profitajo, havajo}\senco{2}{Cesar havar}\lingui{FIS}
\artiklo{perdriko}%
\vorto{perdriko}\fako{zool.} Ucelo galinacea, ne-domestika, qua vivas kun sua simili en la frument-agri, en la prati\lingui{FIS}
\artiklo{perena}%
\vorto{perena}\fako{bot.} Qua ri-kreskas dum plura yari intersequanta\lingui{DEFIS}
\artiklo{perfekta}%
\vorto{perfekta} Di qua la ecelo en sua genero esas absoluta\lingui{DEFIS}
\artiklo{perfekto}%
\vorto{perfekto}\fako{gram.} Tempo di verbo, qua indikas, ye maniero absoluta, ke ago eventis, sen dicar precize ye qua epoko, ye qua instanto\lingui{DEFIS}
\artiklo{perfida}%
\vorto{perfida} Qua agas quale trahizanto, desloyale, ad ulu qua fidis lu\lingui{DEFIS}
\artiklo{perforar}%
\vorto{perforar}\fako{trans.} Trairar facante truo, ma sen rotacar\lingui{DEFIS}
\artiklo{pergameno}%
\vorto{pergameno} Felo de mutono, preparita por recevar skriburo, imprimuro, pikturo, desegnuro\lingui{DEFIRS}
\artiklo{perianto}%
\vorto{perianto}\fako{bot.} Ensemblo ek la envelopili di floro, qua povas konsistar ye kalico e korolo\lingui{DEFIS}
\artiklo{periferio}%
\vorto{periferio}\senco{1}{\textit{(geom.)} Konturo di figuro kurvoforma}\senco{2}{Surfaco extera di korpo}\lingui{DEFIRS}
\artiklo{perifrazo}%
\vorto{perifrazo}\fako{retor.} Frazo-formo uzata kom equivalanto di la vorto apta, konvenanta. \textit{(perifrazar)}\lingui{DEFIRS}
\artiklo{perigeo}%
\vorto{perigeo}\fako{astron.} Punto ek la orbito di planeto ube olca esas maxime proxima a Tero\lingui{DEFIRS}
\artiklo{perihelio}%
\vorto{perihelio}\fako{astron.} Punto ek la orbito di planeto ube olca esas maxime proxima a Suno\lingui{DEFIS}
\artiklo{perikardio}%
\vorto{perikardio}\fako{anat.} Sako membrana qua envelopas la kordio\lingui{DEFIS}
\artiklo{perimetro}%
\vorto{perimetro} Konturo qua limitizas spaco-parto determinita\lingui{DEFIRS}
\artiklo{perineo}%
\vorto{perineo}\fako{anat.} Regiono inter la organi sexuala e la anuso\lingui{DEFIS}
\artiklo{periodo}%
\vorto{periodo} Tempo-quanto dum qua kozo parkuras la fazi di sua duro\lingui{DEFIRS}
\artiklo{periostito}%
\vorto{periostito}\fako{patol.} Inflamuro di la periosto\lingui{DEFIS}
\artiklo{periosto}%
\vorto{periosto}\fako{anat.} Membrano fibroza qua parkovras la osti\lingui{DEFIS}
\artiklo{peripatetiko}%
\vorto{peripatetiko} La persono qua adheris a la filozofio di Aristoteles\lingui{DEFIRS}
\artiklo{peripecio}%
\vorto{peripecio} Bruska iro-chanjo di la eventantajo ad altra\lingui{DEFIRS}
\artiklo{perisar}%
\vorto{perisar}\fako{netrans.} Supresesar per morto violentoza\lingui{EFIS}
\artiklo{periskopo}%
\vorto{periskopo} Optiko-tubo, uzata dat-unesme da la submarnavi, kom aparato qua posibligas vidar dum submerseso\lingui{DEFIS}
\artiklo{peristalta}%
\vorto{peristalta}\fako{fiziol.}\lingui{DEFIS}\subvorto{}{movo peristalta}{Kontrakto, de supre ad infre, di la fibri cirklatra di la digesto-tubo, qua havas kom efekto pulsar kontinue adavan la nutrivi}{}{}
\artiklo{peristilo}%
\vorto{peristilo}\fako{arkitekt.} Edifico di qua la periferio interna kontenas kolonaro\lingui{DEFIRS}
\artiklo{peritoneo}%
\vorto{peritoneo}\fako{anat.} Membrano seroza qua esas quaze la tapeto di la kavajo abdomina, ed envelopas la viceri quin ica kontenas\lingui{EFIS}
\artiklo{perjurar}%
\vorto{perjurar}\fako{netrans.}\senco{1}{Agar trompo-juro}\senco{2}{Violacar sua juro-promiso}\lingui{EFIS}
\artiklo{perkalino}%
\vorto{perkalino} Perkalo tintita e glata\lingui{DEFIRS}
\artiklo{perkalo}%
\vorto{perkalo} Kaliko dina, di qua la lanugo deprenesis\lingui{DEFIRS}
\artiklo{perko}%
\vorto{perko}\fako{zool.} Fisho qua vivas en aquo sen-sala, e havas flosi dornoza\latina{perca}\lingui{EFIS}
\artiklo{perkutar}%
\vorto{perkutar}\fako{trans.}\senco{1}{\textit{(medic.)} Explorar per la metodo qua konsistas ye frapar sur la parieto di korpo-kavajo por saveskar (segun la naturo di la bruiso) quan esas la stando di la organi en ta kavajo}\senco{2}{\textit{(armo)} \textit{(aludante fusilo)} Agar la frapo qua – per explozo di la pulvero – deskuplos la projektilo di la pafarmo}\lingui{DEFIS}
\artiklo{perlo}%
\vorto{perlo}\senco{1}{Globeto, maxim-multa-kaze arjento-blanka, kun reflekti ciel-arkea, qua formacesas en ula konki}\senco{2}{Mikra bulo ek metalo, emalio, vitracho, quan onu boras, e tra serio de qui onu pasigas filo por facar ek li koliari, pasmenti, e c.}\lingui{DEFIS}
\artiklo{perlomatro}%
\vorto{perlomatro} Materio blanka, kun reflekti ciel-arkea, ye qua konsistas la latero interna di ula konki, ed uzesas por tabulifado\lingui{DEFIRS}
\artiklo{permanar}%
\vorto{permanar}\fako{netrans.} Mantenar su, konservar su, sen interrupto. – Esar kontinue en la sama loko, sen absenteso nek intermito\lingui{DEFIRS}
\artiklo{permear}%
\vorto{permear}\fako{trans.) (aludante liquido, gaso} Pasar de ca a ta latero di parieto per trairar la pori di olca\lingui{EFIS}
\artiklo{permigar}%
\vorto{permigar}\fako{trans.) (ulo ad ulu} Grantar la darfo o la yuro (agor o facar ulo)\lingui{EFIS}
\artiklo{permutar}%
\vorto{permutar}\fako{trans.} Pozar ulo en la plaso di altra kozo, e reciproke\lingui{EFIS}
\artiklo{pero}%
\vorto{pero} Vasalo suprega-ranga, membro di la judicio-korto rejala\lingui{DEFIRS}
\artiklo{peroneo}%
\vorto{peroneo}\fako{anat.} Osto longa di la gambo, fixigita ye la latero extera di tibio\lingui{EFS}
\artiklo{perono}%
\vorto{perono}\fako{arkitekt.} Eskalero de kelka gradi, konstruktita avan la fasado di domo, e qua finas per platformo a qua abutas la eniro-koridoro precipua\lingui{EFR}
\artiklo{peronosporo}%
\vorto{peronosporo}\fako{bot.} Fungo, tipo di la familio « peronosporei »\latina{peronospora}
\artiklo{peroraciono}%
\vorto{peroraciono}\fako{retor.} La *ultima ek la parti ye qui konsistas diskurso\lingui{DEFIS}
\artiklo{peroxidazo}%
\vorto{peroxidazo} Fermento solvebla di la puso, apta deskompozar la aquo oxizita e transportar la oxo sur la tintivo di gayako, sur hidrokinono, sur pirogalo, e c.
\artiklo{perpendikla}%
\vorto{perpendikla}\fako{geom.} Qua produktas – kun un rekto, un plano – du anguli adjacanta\lingui{DEFIRS}
\artiklo{perpetua}%
\vorto{perpetua} Qua duras e duros konstante\lingui{DEFIS}
\artiklo{perplexa}%
\vorto{perplexa} Embarasata e nequieta, inter plura solvi, rezolvi, posibla\lingui{DEFIS}
\artiklo{perquisitar}%
\vorto{perquisitar}\fako{netrans.} \textit{(che ulu)} \textit{(aludante judiciisto inquestala)} : Agar sercho judiciala minucioza, en la domicilo di inkulpato, e che irgu, por embargar e forprenar omna dokumenti ed altra objekti utila por saveskar pri la fakti\lingui{DEFIS}
\artiklo{persekutar}%
\vorto{persekutar}\fako{trans.} Persequar sen-halte, sen-cese, e traktar kruele, violente\lingui{EFIS}
\artiklo{persequar}%
\vorto{persequar}\fako{trans.} Sequar proxime, esforcante por atingar\lingui{EFIS}
\artiklo{perseverar}%
\vorto{perseverar}\fako{netrans.} Permanar ferma e konstanta en sua maniero esar od agar\lingui{EFIS}
\artiklo{persieno}%
\vorto{persieno} Sorto di shutro ajurizita, ek (generale kelka) lami vertikala, faldebla, fixa o movebla, qua shirmas kontre la suno-radii\lingui{DEFI}
\artiklo{persikario}%
\vorto{persikario}\fako{bot.} Planto, ordinare en Francia, uzata kom ecitivo ed urinifigivo en la hidropsi\latina{polygonum persicaria}\lingui{EFIS}
\artiklo{persiko}%
\vorto{persiko}\fako{bot.} Frukto de arboro qua devenis de Persia, di qua la shelo esas veluratra, la pulpo saporoza, la kerno grosa e tre harda\latina{prunus persica}\lingui{DFRL}
\artiklo{persistar}%
\vorto{persistar}\fako{netrans.} Permanar ferma pri sua rezolvajo malgre la rezisti\lingui{EFIS}
\artiklo{persono}%
\vorto{persono}\senco{1}{Individuo ek la speco « homo ». To quo konstitucas la persono ipsa}\senco{2}{\textit{(gram.)} Formo di la konjugo, per qua distingesas la persono qua parolas, ta a qua onu parolas, ta pri qua onu parolas}\lingui{DEFIRS}
\artiklo{perspektiva}%
\vorto{perspektiva} Qua reprezentas objekti segun la diferi di aspekto quin produktas la foreso-grado e la dispozeso\lingui{DEFIRS}
\artiklo{persuadar}%
\vorto{persuadar}\fako{trans.) (ulu, pri}\senco{1}{Duktar (ulu) a kredar (ulo), sen pruvo}\senco{2}{Igar kredar (ulo), sen pruvo}\lingui{EFIS}
\artiklo{perturbar}%
\vorto{perturbar}\fako{trans.}\senco{1}{Desordinar la funciono di sistemo, di organo}\senco{2}{\textit{(astron.)} Deviacigar de lua iro normala, per la atrakto da astro vicina}\senco{3}{Trublar la stando di stato, di asemblitaro}\lingui{DEFIRS}
\artiklo{perucho}%
\vorto{perucho}\fako{zool.} Mikra papagayo di qua la plumaro esas bele-verda\latina{conorus}\lingui{FI}
\artiklo{peruko}%
\vorto{peruko} Hararo falsa\lingui{DEFIRS}
\artiklo{perversa}%
\vorto{perversa} Qua havas tendenci o gusto ne-normala, kontre-etika\lingui{DEFIS}
\artiklo{pervinko}%
\vorto{pervinko}\fako{bot.} Planto di qua la folii esas brile-verda, e la flori, blua, kin-festona\latina{vinca minor}\lingui{EFS}
\artiklo{pesario}%
\vorto{pesario} Bandajo destinita a mantenar la utero\lingui{DEFIS}
\artiklo{pesimismo}%
\vorto{pesimismo} Doktrino di ta qua kredas ke, pro ke lo mala esas la existo-*leyo (lego), onu mustas odiar vivo\lingui{DEFIRS}
\artiklo{peskar}%
\vorto{peskar}\fako{trans.} Probar kaptar fishi (e mem sucesar lo) per tirar li ek aquo per filo, o per reto, e c.\lingui{EFIS}
\artiklo{pesto}%
\vorto{pesto}\fako{patol.} Morbo epidemiala, kontagiala, qua mortigas personi multega\lingui{DEFIS}
\artiklo{petalo}%
\vorto{petalo}\fako{bot.} Parto di la korolo di floro\lingui{EFIS}
\artiklo{petardo}%
\vorto{petardo}\senco{1}{Explozivo por destruktar obstaklo}\senco{2}{Peco piroteknala qua explozas granda-bruisoze}\lingui{DEFIRS}
\artiklo{peticionar}%
\vorto{peticionar}\fako{netrans., pri ulo, ad ulu} Demandar skribe, a la reprezentanti di autoritatozaro, a granda korpi politikala; generale ta demando signatesas da granda nombro de personi\lingui{DEFIRS}
\artiklo{petiolo}%
\vorto{petiolo}\fako{bot.} Kaudo di la folio, qua unionas lu a la stipo\lingui{EFIS}
\artiklo{petonklo}%
\vorto{petonklo}\fako{zool.} Molusko sen-kapa, kun konko du-valva\latina{pectunculus}\lingui{FI}
\artiklo{petrelo}%
\vorto{petrelo}\fako{zool.} Mar-ucelo, kun ali longa, qua, dum sua flugo rapida, semblas frolar la aquo-surfaco\latina{procellaria}\lingui{EFS}
\artiklo{petrisar}%
\vorto{petrisar}\fako{trans.} Klemar, manuagar (substanco pasta)\lingui{F}
\artiklo{petro}%
\vorto{petro} Materio ne-organika, harda, solida, qua konsistas ye substanci diversa, agregita, difuzita surface ed interne di sulo\lingui{DEFIRS}
\artiklo{petrografio}%
\vorto{petrografio} Cienco di qua la studiajo esas la roki di la ter-kortico, pri lia kontenajo mineralogiala e kemiala, e lia genezo\lingui{DEFIRS}
\artiklo{petrolo}%
\vorto{petrolo} Oleo minerala, obtenata de fonti naturala, uzata por lumizo, varmigo di mashini, e c.\lingui{DEFIRS}
\artiklo{petroselo}%
\vorto{petroselo}\fako{bot.} Planto legum-gardenala, kun saporo pikanta, uzata precipue kom kondimento\latina{petroselinum}\lingui{DRL}
\artiklo{petular}%
\vorto{petular}\fako{netrans.} Manifestar sua vivaceso per agito bruisoza; folo-ludetar, movante sen su-dominaco
\artiklo{petunio}%
\vorto{petunio}\fako{bot.} Ornamentala planto, kun flori bunta ed bon-odoranta\latina{petunia}\lingui{DEFIS}
\artiklo{pezar}%
\vorto{pezar}\fako{aludante korpo}\senco{1}{Igar sentar sua preso sur sua suportanto, o sua tiro di to a qua lu suspendesas – preso e tiro plu o min granda segun lua maso}\senco{2}{\textit{(n grami, kilogrami)} Enuncar la quanto de mezurunaji qui korespondas a ta preso o tiro}\lingui{FIS}
\artiklo{pia}%
\vorto{pia}\fako{religio} Di qua la devoteso religiala esas vivaca, forta e sincera\lingui{EFIS}
\artiklo{piafar}%
\vorto{piafar}\fako{netrans.) (aludante kavalo} Frapar bruisoze sulo per pedo, levante alte la gambi avana\lingui{DEFS}
\artiklo{piamatro}%
\vorto{piamatro}\fako{anat.} Parto interna di la membrano triopla qua tegas la cerebro
\artiklo{piano}%
\vorto{piano}\fako{muziko} Instrumento kun klavaro, di qua la kordi ne tushesas per plumo-pinto (li tale tushesis pri klavikordo), ma frapesas per marteleti – procedo qua posibligas plufortigo o modero di la sono segun-vole, per presar plu o min forte sur la klavo qua movigas la marteleto\lingui{DEFIRS}
\artiklo{piceo}%
\vorto{piceo}\fako{bot.} Arboro konifera, kelke simila ad abieto\latina{picea excelsa}\lingui{FIL}
\artiklo{pichpino}%
\vorto{pichpino}\fako{bot.} Arboro (di Usa) di qua la bela ligno esas rede-flava, riche veinoza, plu harda e plu densa kam la pino-ligno ordinara, ma esas ruptebla e fendebla plu facile. (australis)\latina{pinus palustris}\lingui{DEFIR}
\artiklo{piedestalo}%
\vorto{piedestalo} Suportilo di kolono, di vazo, di statuo, establisita sur sulo\lingui{DEFIRS}
\artiklo{pierido}%
\vorto{pierido}\fako{zool.} Genero de insekti lepidoptera, ek la familio « pieridei »; un speco de oli produktas raupi qui devastas la kreskaji (kauli, e c.)\latina{pieris brassicae}
\artiklo{pietismo}%
\vorto{pietismo}\fako{religio krist.} Nomo qua karakterizas la mento qua impulsis, en la yarcento 17-esma, e mem ankore nun, grupi diversa de protestantisti Germana qui tendencas ad asketismo e proklamas la sacerdoteso universala di omna Deo-kredanti\lingui{DEFIRS}
\artiklo{pietisto}%
\vorto{pietisto} Adepto di pietismo\lingui{DEFIRS}
\artiklo{pigargo}%
\vorto{pigargo}\fako{zool.} Genero de raptuceli, ek la familio « falkonidei », qua kontenas cirkum dek speci, difuzita sur la Terglobo omna-regiono, ecepte en Sud-Amerika\latina{circus pugargus}\lingui{EF}
\artiklo{pigmento}%
\vorto{pigmento}\fako{anat.} Materio ek granularo diversa-kolora, quan onu trovas en la celuli di la muko-strato di epidermo, e qua produktas la nuanci interdiferanta di pelo che la homo e la animali\lingui{DEFIRS}
\artiklo{pigmeo}%
\vorto{pigmeo}\senco{1}{\textit{(mitol.)} Persono di la raso « nani »}\senco{2}{Persono sen merito nek valoro sociala}\lingui{DEFIRS}
\artiklo{pigo}%
\vorto{pigo}\fako{zool.} Ucelo babilera, ek la familio « korvi », kun plumaro blanka e nigra, e longa kaudo dispozita etajatre\latina{pica}\lingui{IS}
\artiklo{pikadoro}%
\vorto{pikadoro} Kavaliero qua, en tauro-kombato, atakas la animalo per piquo\lingui{DEFIRS}
\artiklo{pikar}%
\vorto{pikar}\fako{trans.}\senco{1}{Entamar neprofunde, per pinto}\senco{2}{Truizar per ulo quan onu igas penetrar per pinto}\senco{3}{Igar (ulo penetrar aden ulo, pinto-avane)}\senco{4}{Stimular vivace; frapar per sento vivaca, penetriva}\lingui{DEFS}
\artiklo{pikeo}%
\vorto{pikeo} Stofo de kotono, ek du texuri aplikita ica kontre ita, ed unionita per steburi qui formacas desegnuri\lingui{DEFIS}
\artiklo{piketo}%
\vorto{piketo}\senco{1}{\textit{(armeo)} Mikra trupo de kavalriani, pronta departar quik lor signalo}\senco{2}{Karto-ludo ek 32 karti, en qua la ludanto qua ganis 30 punti ante ke la adverso ganis un, kontas sua ganajo quale se lu atingabus 60}\lingui{DEFIRS}
\artiklo{piklo}%
\vorto{piklo} Singla ek la kondimenti vejetantala, konservita per vinagro\lingui{E}
\artiklo{piknikar}%
\vorto{piknikar}\fako{netrans.} Agar repasto en qua singla partoprenanto pagas sua quoto di la spenso, od adportas sua dishi\lingui{DEFRS}
\artiklo{pikono}%
\vorto{pikono} Utensilo, implemento, ek fero, pintoza, quan uzas la ministi karbonala, petrala, la terlaboristi, por exkavar roko o materii harda; la fera parto havas pinto perpendikla ye la mancho
\artiklo{piktar}%
\vorto{piktar}\fako{trans.} Reprezentar per farbi\lingui{EFIS}
\artiklo{pikverdo}%
\vorto{pikverdo}\fako{zool.} Ucelo kun plumaro flava e verda, ek la genero « pegi »\latina{picus viridis}\lingui{FIS}
\artiklo{pilastro}%
\vorto{pilastro}\fako{arkitekt.}\senco{1}{Kolono quadrata, maxim-multa-kaze sinkita en muro}\senco{2}{Montanto quan onu fixigas intervalope en greto, pasamano, balkono}\senco{3}{Kolono de petro-bloki, qua subtenas parti diversa en edifico, en ponto, e c.}\lingui{DEFIRS}
\artiklo{pilgrimar}%
\vorto{pilgrimar}\fako{netrans.} Voyajar pro devoceso, pieso, a loko sakrigita\lingui{DEFIR}
\artiklo{pilio}%
\vorto{pilio}\fako{elektro} Aparato (inventita da Volta, en la yaro 1800), qua transformas ad elektro-fluo la energio quan developas reakto kemiala\lingui{DEF}
\artiklo{pilo}%
\vorto{pilo}\fako{anat.}\senco{1}{Produkturo epidermala, qua havas la formo di filo tenua, qua kovras la korpo di preske omna mamiferi}\senco{2}{Ta produkturo che la homo, ube olu kreskas sur ula parti di la korpo (vangi, labii, mentono, axeli, pubio)}\lingui{EFIS}
\artiklo{pilono}%
\vorto{pilono}\fako{arkitekt.}\senco{1}{\textit{(en Egiptia, antique)} Sorto di avan-korpo qua havas la formo di piramido quadriangula trunka, e kun porda aperturo centrala, quan la Egiptiani antiqua kustumis erektar avan la enireyo di sua templi}\senco{2}{Dekor-motivo qua havas la formo di pilastro quadriangula, lokizita an la lateri di enireyo}\lingui{DEFIRS}
\artiklo{pilorio}%
\vorto{pilorio} Fosto an qua onu ligis per fero-ringo la persono qua esis kondamnita ad expozeso publika (kom puniso legala)
\artiklo{piloro}%
\vorto{piloro}\fako{anat.} Orifico infra di la stomako, tra qua la nutrivi pasas aden la intestini\lingui{EFIS}
\artiklo{pilotar}%
\vorto{pilotar}\fako{trans.} Duktar navo en la mar-paseyi desfacile trairebla, ye la enireyo od ekireyo di portuo\lingui{DEFIS}
\artiklo{pilulo}%
\vorto{pilulo}\fako{farmacio} Medikamento fasonita a mikra globo, maxim-multa-kaze envelopita ye folio dinega ek oro od arjento, qua impedas lua adhero e plufaciligas lua engluteso\lingui{DEFIRS}
\artiklo{pimento}%
\vorto{pimento}\senco{1}{Drinkajo ek mielo, spici e vino}\senco{2}{\textit{(bot.)} Bero konoforma, kun surfaco glata e briloza; komence bele-verda, lu divenas (lor lua matureso) brile-reda – uzata kom kondimento}\latina{mirtus pimenta}\lingui{DEFIS}
\artiklo{pimpinelo}%
\vorto{pimpinelo}\fako{bot.} Planto herbatra, aromata, ek la familio « rozacei », uzata kom kondimento\latina{pimpinella}
\artiklo{pinaklo}%
\vorto{pinaklo}\fako{arkitekt.} Tekto-karpenturo dekorita, e finanta per pinto\lingui{DEFIS}
\artiklo{pinastro}%
\vorto{pinastro}\fako{bot.} Pino sovaja (mar-litorala) de qua onu extraktas terpentino\latina{pinus pinaster}\lingui{EFIS}
\artiklo{pinchar}%
\vorto{pinchar}\fako{trans.} Klemar forte inter la fingri\lingui{EF}
\artiklo{pinco}%
\vorto{pinco}\fako{tekn.} Mikra tenalio, uzata diverse en la mestieri, en kirurgio. – « Pinco » diferas de « tenalio » per to ke lu esas utensilo sen charniro (la du parti ye qua lu konsistas esas asemblita per resorto o parto flexebla)\lingui{DEFIRS}
\artiklo{pinealo}%
\vorto{pinealo}\fako{anat.} Mikra korpo, avan la cerebelo\lingui{DEFIS}
\artiklo{pinglo}%
\vorto{pinglo} Mikra stango de latuno, pinta ye un extremajo, kapoza ye la altra, quan onu uzas por ligar\lingui{EF}
\artiklo{pinguino}%
\vorto{pinguino}\fako{zool.} Ucelo palmipeda di la Nordo-mari, kun ali kurta\latina{aptenodytes}\lingui{DEFIRS}
\artiklo{pinio}%
\vorto{pinio}\fako{bot.} Sorto di pino qua kreskas en la regioni sudala di Europa, kun kapo parasuno-forma\latina{pinus pinea}\lingui{DF}
\artiklo{piniono}%
\vorto{piniono}\fako{tekn.} Roto dentoza qua ingranas la denti di altra roto\lingui{EFS}
\artiklo{pino}%
\vorto{pino}\fako{bot.} Granda arboro rezinoza, sempre verda, ek la familio « koniferi »\latina{pinus sylvestris}\lingui{EFIS}
\artiklo{pinselo}%
\vorto{pinselo} Tufo de pili, ligita an la extremajo di bastono mancha, uzata por extensar farbi, gluo, e c.\lingui{DFS}
\artiklo{pintado}%
\vorto{pintado}\fako{zool.} Ucelo kriema, galinacea, kun kapo nuda, kaudo kurta, plumaro griza bluatre, sur qua dissemesas makuli blanka\latina{Numida meleagris}\lingui{EFS}
\artiklo{pinto}%
\vorto{pinto}\senco{1}{Extremajo tale tenuigita ke lu pikas. (Ex. : pinto di agulo, di espado, di kompaso)}\senco{2}{Ulo kelke pinto-forma. (ex. : pinto di monto)}\lingui{EFIS}
\artiklo{piocho}%
\vorto{piocho}\fako{tekn.} Utensilo di terlaboristo, ek mancho e fer-bloko, pikona ye un extremajo, e haua ye la altra
\artiklo{pioniro}%
\vorto{pioniro}\senco{1}{En armeo : laboristo quan onu uzas por voyifar, exkavar tranchei, e c.}\senco{2}{La persono qua kultiveskas tereni nekultivita til nun. \textit{(anke metaf.)}}\lingui{DEFIR}
\artiklo{piono}%
\vorto{piono} Singla de la disketi di la damo-ludo, e di la figurini di la shak-ludo\lingui{FS}
\artiklo{pipeto}%
\vorto{pipeto}\fako{tekn.} Tubo vitra, konkavigita ye un extremajo, tenuigita til pintesko ye la altra, quan onu povas stopar per un fingro, uzata por transvarsar liquido\lingui{DEFI}
\artiklo{pipiar}%
\vorto{pipiar}\fako{netrans.) (aludante la ucel-yuni} Kriar\lingui{DFS}
\artiklo{pipo}%
\vorto{pipo} Tubo qua finas per quaza furneleto en qua onu pozas tabako quan onu acendas por aspirar lua fumuro *smoke\lingui{EFIS}
\artiklo{pipro}%
\vorto{pipro}\fako{bot.} Mikra bero, qua odoras e saporas forte, e qua, sikigite, o muelite, uzesas kom spico, kondimento\lingui{EFI}
\artiklo{pipsar}%
\vorto{pipsar}\fako{netrans.) (aludante la uceli e precipue la hani} Maladesar pro la kresko sur la pinto di la lango, di pelikulo squamatra qua impedas drinkar\lingui{DEFIS}
\artiklo{piquo}%
\vorto{piquo}\senco{1}{Shaft-armo, di qua la shafto esas garnisita ye ferajo plata e pintoza, plu kurta kam lanco}\senco{2}{Un ek la du figuri nigra, sur ula ludo-karti, di qua la formo esas olta di piquo-ferajo}\lingui{F}
\artiklo{piramido}%
\vorto{piramido}\senco{1}{Monumento kun quar facii triangula e bazo quadri-angula, qua uzesis kom tombo por la olima reji di Egiptia (faraoni)}\senco{2}{\textit{(geom.)} Solido kun edri triangula unionita per somito komuna, kun bazo poligonala, e qua havas tam multa lateri kam edri}\lingui{DEFIRS}
\artiklo{pirato}%
\vorto{pirato} Persono qua cirkulas sur la mari por ibe kaptar, sizar, embargar, raptar, sen permiso da la rejo di sua lando\lingui{DEFIRS}
\artiklo{piretro}%
\vorto{piretro}\fako{bot.} Planto ek la familio « kompozaji ». di qua la kapituli, sikigite, divenas pulvero kontre-insekta\latina{pyrethrum}\lingui{EFI}
\artiklo{piriko}%
\vorto{piriko}\fako{prozodio antiqua} Pedo kompozita ek du vokali kurta, Onu renkontras lu ye la fino di ula hexametri da Homeros, vice spondeo o trokeo\lingui{DEFIS}
\artiklo{pirito}%
\vorto{pirito}\fako{mineral.} Sulfuro metala kruda, inflamebla\lingui{DEFIS}
\artiklo{piro}%
\vorto{piro}\fako{bot.} Frukto di arboro ek la familio « rozacei », oblonga, di qua la volumino deskreskas arivante proxim la kaudo\latina{pirus communis}\lingui{EFIS}
\artiklo{pirogalo}%
\vorto{pirogalo}\fako{kemio} 1-2-3 C₆H₃ (OH)₃
\artiklo{pirogo}%
\vorto{pirogo} Barko di sovaji, facita, en la kazi maxim multa, ek arboro-trunko kavigita\lingui{DEFIRS}
\artiklo{pirograbar}%
\vorto{pirograbar}\fako{trans.} Dekorar ligno-peco per grabar sur lu desegnuro per pintoza stango metala, varmigita til esar intense-reda\lingui{EFIS}
\artiklo{pirolo}%
\vorto{pirolo}\fako{zool.} Genero de pasero, kun beko konoforma, ek la familio « korvidei »\latina{oriolus galbula}\lingui{D}
\artiklo{pirometro}%
\vorto{pirometro}\fako{tekn.} Instrumento di fiziko, qua uzesas por mezurar la temperaturi extreme alta-grada\lingui{DEFIRS}
\artiklo{pironismo}%
\vorto{pironismo}\fako{filoz.} Doktrino da Pyrrhon, skeptiko di la Grekia antiqua, qua kombinas en su la filozofio di Demokrites, la dialektiko di la sofisti, e la asketismo di ula sekto di India (gimnosofista)\lingui{DEFIRS}
\artiklo{piroso}%
\vorto{piroso}\fako{patol.} Sento di bruluro departanta de la epigastro ed acensanta til ezofago e til la fauco, e qua akompanesas da rukti e da ri-engluti di liquido acida e brulanta\lingui{DEFIS}
\artiklo{pirotekno}%
\vorto{pirotekno} Kombinuri diversa di efekti de lumo, de koloro, produktata da pulveri specala quin onu inflamas\lingui{DEFIS}
\artiklo{piroxilo}%
\vorto{piroxilo}\fako{kemio} Produkturo obtenita per traktar, per nitracido, materio celulosa, exemple : ligno, papero, kanabo, kotono, palio, e c.
\artiklo{piruetar}%
\vorto{piruetar}\fako{netrans.) (aludante persono} Rotacar, nur un pedo tushante sulo\lingui{DEFIRS}
\artiklo{pistacho}%
\vorto{pistacho}\fako{bot.} Mandelo de frukto di arboro ek la familio « terebintacei », envolvita en pelikulo verda, uzata da la sukrajisti\latina{pistacia vera}\lingui{DEFIRS}
\artiklo{pistar}%
\vorto{pistar}\fako{trans.} Reduktar a mikra fragmenti, per shoki iterata\lingui{EFI}
\artiklo{pistolo}%
\vorto{pistolo} Pafarmo mikra, tre portebla, quan onu tenas en nur un ek la manui kande onu pafas\lingui{DEFIRS}
\artiklo{pistono}%
\vorto{pistono}\fako{tekn.} Cilindro, movanta en pumpilo per stango, retropulsante fluido o pulsata da lu\lingui{DEFIRS}
\artiklo{pitagorala}%
\vorto{pitagorala}\senco{1}{\textit{(filoz.)} Apartenanta a Pithagores, a lua penso-skolo od a lua doktrini}\senco{2}{\textit{(matem.)} Nombri pitagorala : nombri interligita per la relato a² = b² + c}
\artiklo{pitekantropo}%
\vorto{pitekantropo}\fako{paleont.} Genero de granda simii fosila, reprezentata da kranio-kaloto, kelka denti ed un femuro, trovita en la yaro 1891 en sulo-strato pliocena o pleistocena, en la insulo Java (Azia)\lingui{DEFIRS}
\artiklo{pitio}%
\vorto{pitio}\fako{epoki antiqua} Sacerdotino di la oraklo di Apolono, en Delphos\lingui{DEFIS}
\artiklo{pitoniso}%
\vorto{pitoniso}\fako{biblo} Muliero qua anuncas la eventontaji\lingui{DEFIS}
\artiklo{pitono}%
\vorto{pitono}\fako{zool.}\senco{1}{Sorto di boao ek India}\senco{2}{\textit{(mitol.)} Nomo di serpento monstra, mortigita da Apolono, apud Delphos}\latina{python}\lingui{DEFIS}
\artiklo{pitoreska}%
\vorto{pitoreska}\senco{1}{Qua, pro lua dispozeso, aspektas apta divenar piktajo}\senco{2}{Dicesas pri stilo qua quaze piktas la personi o la kozi}\lingui{DEFIS}
\artiklo{pituito}%
\vorto{pituito}\fako{patol.} Humoro viskoza sekrecata da ula mukosi di la nazo, di la bronkii, en stando morbala\lingui{EFIS}
\artiklo{pivoto}%
\vorto{pivoto}\fako{mekan.}\senco{1}{Suportilo di la axo cirkum qua rotacas korpo}\senco{2}{\textit{(milit-arto)} \textit{(aludante konversiono di trupo de soldati)} Punto cirkum qua ta konversiono eventas}\senco{3}{\textit{(metaf.)} Motivo precipua di la agi, che la homi}\lingui{EF}
\artiklo{pizo}%
\vorto{pizo}\fako{bot.} Singla de la grani de planto leguminosa, papilionacea, inkluzita en shelo verda\latina{pisus sativus}\lingui{EFI}
\artiklo{placento}%
\vorto{placento}\senco{1}{\textit{(anat.)} Organo karnoza, sponjatra, qua adheras a la utero, komunikante per la kordono umbilikala kun la feto, per la vaskuli di qua la sango patrinala arivas a la feto e retrovenas a la patrino}\senco{2}{\textit{(bot.)} Parto di la surfaco di karpelo, sur qua esas insertita la ovuli}\lingui{DEFIS}
\artiklo{placo}%
\vorto{placo} Loko publika, netektoza, perimetro-garnisata per do-mi o monumenti, kun, posible, monumento centrala\lingui{DEFIRS}
\artiklo{plado}%
\vorto{plado} Vazo di qua la fundo esas plata, quan onu pozas avan singla repastonto o festinonto por recevar la nutrivi\lingui{EFIS}
\artiklo{plafono}%
\vorto{plafono}\fako{arkitekt.} Surfaco ye qua konsistas la parto supra di chambro, salono, galerio, kirko, e c., super la kapi di la personi qui esas en lu\lingui{DFR}
\artiklo{plago}%
\vorto{plago}\senco{1}{\textit{(fiziol.)} Laceruro, extera od interna, kun desorganizeso di la tisui en parto di la korpo, produktita da vunduro o kauzo morbifera}\senco{2}{\textit{(metaf.)} Persono o kozo konsiderata kom instrumento per qua Deo frapas o kastigas la homi}\senco{3}{Kalamitato qua frapas populo}\senco{4}{To quo esas funesta a persono od a kozo}\lingui{DEFS}
\artiklo{plajiar}%
\vorto{plajiar}\fako{trans.) (literaturo} Furtar ek la verki da kunhomo\lingui{DEFIRS}
\artiklo{plajo}%
\vorto{plajo} Bendo, plu o min larja, de sulo alonge maro od oceano, quan ta maro o ta oceano alterne kovras e deskovras. (Existas plaji de sablo e plaji de rul-stoni)\lingui{FIS}
\artiklo{plako}%
\vorto{plako} Lamo de metalo – maxim-multa-kaze destinita aplikesor adsur surfaco plana\lingui{EFIR}
\artiklo{plando}%
\vorto{plando}\fako{anat.} Parto di la pedo qua tushas sulo kande onu stacas o kande onu marchas\lingui{FIS}
\artiklo{planeto}%
\vorto{planeto}\senco{1}{\textit{(astron. antiqua)} Astro vaganta (kontraste kun la steli fixa). \textit{(astrol.)} Astro a qua onu imputas efiko a la fato di la homo}\senco{2}{\textit{(astron. moderna)} Astro qua parkuras la spaco cirkum Suno (situita ad un ek la foki)}\lingui{DEFIRS}
\artiklo{planimetrio}%
\vorto{planimetrio} Parto di geometrio : la arto mezurar la surfaci plana\lingui{DEFIRS}
\artiklo{planimetro}%
\vorto{planimetro} Instrumento per qua onu povas evaluar la mezuri di areo plana\lingui{DEFIRS}
\artiklo{planisfero}%
\vorto{planisfero}\fako{geogr.} Mapo en qua la Ter-\textit{sfero} dividesas en du diametri, reprezentata da surfaco plana, egardinte la redukto demandata da perspektivo\lingui{DEFIRS}
\artiklo{planko}%
\vorto{planko} Ligno-peco, fendita longesale, plu longa kam larja, e nur poke dina\lingui{DEFIS}
\artiklo{plankton}%
\vorto{« plankton »}\senco{1}{Ensemblo ek omna vivanti (animali e vejetanti) qui fluktuas en la aquo di maro o di lago, inter la surfaco e la fundo}\senco{2}{Animali o vejetanti tre mikra di la aquo di maro o di lago}\lingui{DEFIRS}
\artiklo{plano}%
\vorto{plano}\senco{1}{Surfaco sen kurviguro nek ondaro}\senco{2}{\textit{(geom.)} Surfaco plana : sur qua rekto esas aplikebla omna-sinse.}\subvorto{}{angulo plana}{Intersekuro da du rekti sama-plana. Figuro plana : trasita sur surfaco plana}{}{}\senco{3}{Desegnuro qua reprezentas la dispozeso di la kompozanti di urbo, di regiono, di edifico, kun indiko di la proporcioni e di la plaso relativa di la parti diversa}\lingui{DEFIRS}
\artiklo{plantacar}%
\vorto{plantacar}\fako{trans.} Fixigar en sulo per radiki (vejetanto)\lingui{DEFIS}
\artiklo{plantago}%
\vorto{plantago}\fako{bot.} Planto herbatra di qua la spiko kovresas ye mikra semini – tipo di la familio « plantaginei »\latina{plantago}\lingui{EFI}
\artiklo{plantigrado}%
\vorto{plantigrado}\fako{zool.} Mamifero qua, marchante, pozas la plando di sua pedi sur sulo. (Ex. : urso, daxo, e c.)\lingui{DEFIS}
\artiklo{planto}%
\vorto{planto}\fako{bot.} Omno quo fixigesas en sulo per radiki\lingui{DEFIS}
\artiklo{plantono}%
\vorto{plantono}\fako{armeo} Soldato qua stacas avan la pordo di oficiro alta-grada, por transmisar komandi, e c.\lingui{F}
\artiklo{plasmo}%
\vorto{plasmo}\fako{biol.} Parto liquida di sango, en qua fluktuas la globuli di sango\lingui{DEFIS}
\artiklo{plasmodermo}%
\vorto{plasmodermo}\fako{biol.} Materio protoplasma qua formacas la komuniki intercelula quin onu vidas che la vejetanti e che ula animali\lingui{DEFIRS}
\artiklo{plasmodio}%
\vorto{plasmodio}\fako{parazitologio} Parazito di la genero « plasmodium », en qua asemblesas la tri speci de hematozoi di la marsho-febro amiboida, e qua vivas en la hematii di la homo e di la animali, e genitesas per skizogonio o sporogonio\lingui{DEFIRS}
\artiklo{plasmolizo}%
\vorto{plasmolizo}\fako{biol.} Fenomeno osmosala, qua abutas a la destrukto di la plasmo\lingui{DEFIRS}
\artiklo{plasmosomo}%
\vorto{plasmosomo}\fako{histol.} Nukleolo vera, nedependanta de la reto de linino, qua povas movar su sur la nukleo e kolorizesas intense da la kolorizivo di plasmo\lingui{DEFIRS}
\artiklo{plaso}%
\vorto{plaso}\senco{1}{Spaco, okupata od okupebla, da singla persono en irga loko (veturo, teatro, cinemo, e c.)}\senco{2}{Loko propra, rezervita o konvenanta por singla kozo}\senco{3}{Rolo atribuita ad employato en entraprezuro, komerco-kontori, fabrikeyo, e c.}\lingui{DEF}
\artiklo{*plastica}%
\vorto{*plastica} Quan onu povas formizar o modlar
\artiklo{plastido}%
\vorto{plastido}\fako{biol.} Ensemblo ek la substanci vivanta di la celulo, ecepte la rezervo-substanci, la ek-sekrecaji, la envelopo-membrano\lingui{DEFIS}
\artiklo{plastidulo}%
\vorto{plastidulo}\fako{biol.} Granularo dispozita ganso-marche, qua formacas la ensemblo de la filamenti di la sponjo-plasmo di la celulo\lingui{DEFIS}
\artiklo{plastiko}%
\vorto{plastiko}\fako{kemio} Singla de la materii qui esas kemio-stabila kande uzata normale, ma qui, ye ula instanti dum lia produkteso, esas plastika, ed esas formizita o muldebla per kaloro, preso, o per ta du faktori. La maxim multa plastiki esas polimera, ed esas klasifikebla a termoplastiki e termo-rezistivi\lingui{DEFIRS}
\artiklo{plastilino}%
\vorto{plastilino} Modlo-argilo
\artiklo{plastro}%
\vorto{plastro}\senco{1}{Sorto di gipso per qua la aktoro qua pleas la rolo di Petreto kovras sua vizajo en la pantomimi}\senco{2}{Blankajo quan la mulieri uzas kom fardo}\lingui{DEFIRS}
\artiklo{plastrono}%
\vorto{plastrono}\senco{1}{Peco pektorala di kuraso}\senco{2}{Peco pektorala polsterizita, di vesto di skermisto}\senco{3}{Parto pektorala di kamizo di viro}\lingui{EFI}
\artiklo{plata}%
\vorto{plata} Di qua la surfaco esas glata, sen saliajo\lingui{DEFIR}
\artiklo{platano}%
\vorto{platano}\fako{bot.} Arboro di qua la folii esas larja, simila a manuo apertita, tipo di la familio « platanei »\latina{platanus}\lingui{DEFIS}
\artiklo{plateno}%
\vorto{plateno}\fako{tipografio} En la imprimo-mashino, brakio-funcionanta, la peco de gisfero, quadrata o rektangula, perfekte glata, qua, kande abasata adsur la komposturo (per la stango) produktas la imprimuro
\artiklo{platformo}%
\vorto{platformo}\senco{1}{Kovrilo horizontala di domo sen tekto-karpenturo}\senco{2}{Maso tera, planigita, sur qua onu establisas la fundamento di edifico}\senco{3}{Planko-sulo qua kovras la kapi di la palisi di palisegaro sinkita en sulo}\senco{4}{Planko-sulo provizora sur qua la masonisto-helperi facas la mortero}\lingui{DEFIRS}
\artiklo{platino}%
\vorto{platino}\fako{kemio} Metalo precoza, griz-blanka, mola, duktila, plu pezoza kam oro, ne-alterebla en aero, tre tenaca. Atom-pezo : 194\simbolo{Pt}\lingui{DEFIS}
\artiklo{platonala}%
\vorto{platonala} Qua apartenas a la doktrino da Platono\lingui{DEFIRS}
\artiklo{plaudar}%
\vorto{plaudar}\fako{netrans.} Bruisar monotone quale aquo qua agitesas\lingui{L}
\artiklo{plear}%
\vorto{plear}\fako{trans.}\senco{1}{Amuzar su, o laborar, per funcionigar muzik-instrumento}\senco{2}{Reprezentar teatrajo, cinemo-filmo. \textit{(anke metaf.)}}
\artiklo{plebeyo}%
\vorto{plebeyo}\fako{historio di Roma, antique} Persono ek la infra strati di populo\lingui{DEFIRS}
\artiklo{plebicito}%
\vorto{plebicito} Voto di ula lego o rezolvo, da omna civitani di la stato (kom ex. : en Suisia)\lingui{DEFIRS}
\artiklo{pledar}%
\vorto{pledar}\fako{netrans.} Defensar (afero procesala) koram la judiciisti\lingui{DEFS}
\artiklo{plektar}%
\vorto{plektar}\fako{trans.} Fasonar texuro (ek fasko de hari, e c.) plata, rubando-forma, per interkrucumo\lingui{EL}
\artiklo{plektro}%
\vorto{plektro}\fako{epoki antiqua} Mikra vergo ek ivoro, quan onu uzis por igar vibrar la kordi di liro\lingui{DEFIS}
\artiklo{plena}%
\vorto{plena}\fako{de}\senco{1}{Qua kontenas la tota quanto quan lu povas kontenar}\senco{2}{Qua kontenas granda quanto \textit{(kom ex. : plenigar ye ta o ca kozo, per ta o ca instrumento)}}\senco{3}{\textit{(fiziko)} Plenigita tote, ye materio solida, liquida, o gasa}\lingui{EFIS}
\artiklo{plendar}%
\vorto{plendar}\fako{netrans., pri}\senco{1}{Expresar sua deskontenteso pri ulu od ulo}\senco{2}{Expresar, che autoritatozo policala, sua motivo di deskontenteso por obtenar protekto, indemno, e c.}\lingui{EFS}
\artiklo{plenipotenco}%
\vorto{plenipotenco} Yuro quan onu recevas de ulu pri agar ye lua nomo. (Dicesas precipue pri la yuro donita da suvereno, pri agar ye guvernistaro o korto stranjera.)\lingui{DEFIS}
\artiklo{pleomorfa}%
\vorto{pleomorfa}\fako{mikrobiol.} Di qua la formo varias efike da ula kondicioni determinita
\artiklo{pleonasmo}%
\vorto{pleonasmo}\senco{1}{\textit{(gram.)} Uzo di vorti qui esas superflua pri la senco}\senco{2}{\textit{(retor.)} Abundego di la termini por insistar pri la penso. (kom ex. : « Me vidis lu, ya, vidis lu per mea okuli propra »)}\lingui{DEFIRS}
\artiklo{plesiosauro}%
\vorto{plesiosauro}\fako{paleont.} Genero di reptero, tipo di la familio « plesiosauridei », granda animalo marala, 3/5 metri longa, remarkinda pro lua mikra lezardo-kapo, lua denti strioza, lua kolo longa nemezureble\latina{plesiosaurus}\lingui{DEFIS}
\artiklo{pleto}%
\vorto{pleto} Pladego de metalo, lako, porcelano, kristalo, sur qua onu servas teo, kafeo, o cetera drinkaji, e sur qua onu pozas la bako-kukaji\lingui{F}
\artiklo{pletoro}%
\vorto{pletoro}\senco{1}{\textit{(patol.)} Pleneso di la vaskuli pro abundego di sango, di humori}\senco{2}{\textit{(metaf.)} Tromulteso quaze morba}\lingui{DEFIS}
\artiklo{pleuronekto}%
\vorto{pleuronekto}\fako{zool.} Marfisho plata, di qua la korpo esas romboida, analoga a turboto, ma plu ovala\latina{pleuronectes}\lingui{L}
\artiklo{pleuronektoidi}%
\vorto{pleuronektoidi}\fako{zool.} Genero de fishi plata qui natas sur sua flanko\latina{pleuronectidae}\lingui{L}
\artiklo{plevrito}%
\vorto{plevrito}\fako{patol.} Inflamuro di la plevro\lingui{DEFIRS}
\artiklo{plevro}%
\vorto{plevro}\fako{anat.} Membrano seroza, qua tegas la pulmoni e kovras la parieti di la torako\lingui{EFIRS}
\artiklo{plexo}%
\vorto{plexo}\fako{anat.} Reto de fibreti nerva o de vaskuli interplektita\lingui{DEFIS}
\artiklo{pleyado}%
\vorto{pleyado}\senco{1}{\textit{(astron.)} Singla ek la sis steli (en la epoki antiqua, onu dicis : sep) qui formacas grupo en la Tauro-stelaro}\senco{2}{\textit{(literaturo)} Grupo de sep poeti (en Alexandria ed en la Francia di la yarcento 16-esma.)}\senco{3}{\textit{(filoz.)} Grupo ek la sep saji di la Grekia antiqua}\senco{4}{\textit{(en la domeno di la linguo universala, Ido)} Grupo ek la sep homa fonti di la linguo universala definitiva : Louis de Beaufront, Louis Couturat, Otto Jespersen, André Lalande, Wilhelm Ostwald, R. Lorenz e L. Pfaundler}\lingui{DEFIS}
\artiklo{plezar}%
\vorto{plezar}\fako{netrans., ad, per, pro) (aludante la kozi o la personi} Atraktar \textit{(metafore)} per sua charmo\lingui{EFI}
\artiklo{plezuro}%
\vorto{plezuro} Sentimento di komforto psikologiala, qua naskas de la satisfaceso di deziro, di aspiro\lingui{DEFIS}
\artiklo{pliko}%
\vorto{pliko}\fako{patol.} Exuduro ek la pelo haroza, qua aglutinas la hari aden fasketi\lingui{EFI}
\artiklo{plinto}%
\vorto{plinto}\fako{arkitekt.}\senco{1}{Tabulo quadrata qua formacas soklo}\senco{2}{Bendo plata qua jacas an la pedo di edifico, infre di muro di apartamento}\senco{3}{Bendo de ligno qua aplikesis cirkum panelo}\lingui{DEFIS}
\artiklo{plioceno}%
\vorto{plioceno}\fako{geol.} Strato supra di la sedimento terciara, qua kontenas la fosili maxim recenta\lingui{DEFIRS}
\artiklo{plisar}%
\vorto{plisar}\fako{trans.} Dispozar tale ke lu formacas plura falduri\lingui{DF}
\artiklo{*pliz}%
\vorto{*pliz} Uzata kande onu deziras demandar polite ad ulu, ke lu… -ez. (Plu konciza kam « Me pregas vu pri lo :… -ez » o « Voluntez esar tante komplezema ke vu… -us… ».)\lingui{E}
\artiklo{plombo}%
\vorto{plombo}\fako{kemio} Metalo bluatre-blanka, mola, tre pezoza\simbolo{Pb}\lingui{DEFIRS}
\artiklo{plorar}%
\vorto{plorar}\fako{netrans.}\senco{1}{Manifestar, expresar sua aflikteso}\senco{2}{Emisar lakrimi}\lingui{F}
\artiklo{plu}%
\vorto{plu} Indikas la augmento o la kresko di quanto o qualeso qui duras esar od egardesar kom identa a su ipsa\lingui{FI}
\artiklo{plugar}%
\vorto{plugar}\fako{trans.} Apertar kultivo-sulo, e renversar la glebi, per exkavar sulki en tero, uzante la soko di instrumento specala\lingui{DR}
\artiklo{plumo}%
\vorto{plumo}\senco{1}{\textit{(zool.)} Singla ek la tubi korno-harda, garnisita ye aristi e ye lanugo, qua kovras la korpo di la uceli}\senco{2}{Tubo di plumo, taliita por skribar per lu (pos trempir lua pinto aden inko)}\senco{3}{Mikra lamo de metalo, qua havas la formo di la pinto di plumo taliita, di qua onu glite muntas la extremajo sen-pinta an-en tubo ucel-plumatra}\lingui{EFIS}
\artiklo{plump!}%
\vorto{plump!} Interjeciono qua expresas la bruiso de falo, precipue di ulo tre pezoza, mola, inerta
\artiklo{plumpa}%
\vorto{plumpa}\fako{pri vivanto} Qua su movas quale se lu esus tre pezoza, pro la denseso di sua korpo\lingui{D}
\artiklo{plunjar}%
\vorto{plunjar}\fako{netrans.} Penetrar, sinkar profunde aden aquo, propra-vole\lingui{EF}
\artiklo{plura}%
\vorto{plura} Ula quanto (ne tre granda) de personi, de kozi – mem nur du\lingui{DEFIS}
\artiklo{plus}%
\vorto{plus}\senco{1}{Ultre to; adjuntenda a…}\senco{2}{\textit{(kun « ne » avan od ante lu)} – indikas la ceso di la ago, di la eso}
\artiklo{plusho}%
\vorto{plusho} Stofo simila a veluro, e kun longa pili e lanugo sur la averso\lingui{DEFR}
\artiklo{plusquamperfekto}%
\vorto{plusquamperfekto}\fako{gram.} Verbo-tempo qua indikas pasinto antea ye altra tempo pasinta (t. e. … -is… -inta)
\artiklo{plutokratio}%
\vorto{plutokratio} Guverno da la richi\lingui{DEFIRS}
\artiklo{plutonala}%
\vorto{plutonala}\fako{geol.) (pri roki} Produktita da la efiko di fairo sub-ter-sula naturala\lingui{DEFIRS}
\artiklo{pluvar}%
\vorto{pluvar}\fako{netrans.) (aludante la aquo-vaporo di la nubi} Kondensesar e falar gutope\lingui{FIS}
\artiklo{pluviometro}%
\vorto{pluviometro}\fako{fiziko} Instrumento quan onu uzas por mezurar la quanto de pluv-aquo qua falas en un loko dum ula quanto determinita de tempo\lingui{DEFIS}
\artiklo{pneumatika}%
\vorto{pneumatika} Dicesas pri mashino o pumpo per qua onu pumpas la aero di tanko por eventigar lua vakueso\lingui{DEFIS}
\artiklo{pneumokoko}%
\vorto{pneumokoko}\fako{patol.} Agento patogena, la maxim ordinara, di pulmonito, deskovrita da Louis Pasteur en la yaro 1881
\artiklo{po}%
\vorto{po} Kom preco di… – Kambie di…
\artiklo{pociono}%
\vorto{pociono} Medikamento quan onu glutas forme di drinkajo\lingui{EF}
\artiklo{podagro}%
\vorto{podagro}\fako{patol.} Morbo qua afektas, ordinar-kaze, la artiki, iras de ica ad ita, e quan onu imputis olim a la influo da guti de humoro
\artiklo{podoido}%
\vorto{podoido}\fako{matem.} Singla de la kurvi quin onu obtenas per prenar la simetrii di un punto Q, relate la tangenti di ta kurvo
\artiklo{poemo}%
\vorto{poemo} Versaro pasable longa\lingui{DEFIRS}
\artiklo{poetiko}%
\vorto{poetiko} Verko docala, en qua expozesas la reguli di la kompozo per versi\lingui{DEFIS}
\artiklo{poeto}%
\vorto{poeto} Persono qua su konsakras a poezio; persono qua esas dotita por poezio\lingui{DEFIRS}
\artiklo{poezio}%
\vorto{poezio}\senco{1}{La arto facar verki ek versi}\senco{2}{Beleso di la materio e di la expreso-maniero, maxim grava pri poezio, sendepende de la versifo}\lingui{DEFIRS}
\artiklo{poka}%
\vorto{poka} Mikra-quanta\lingui{FIS}
\artiklo{polarimetro}%
\vorto{polarimetro}\fako{fiziko} Instrumento por mezurar la proporciono di la lumo polarigita qua existas en lumo-radio o la rotaco di la polarigo-plano
\artiklo{polemikar}%
\vorto{polemikar}\fako{netrans., kun ulu, kontre ulu} Debatar skribe\lingui{DEFIRS}
\artiklo{poleno}%
\vorto{poleno}\fako{bot.} Polvo fekundigiva che la vejetanti\lingui{DEFIS}
\artiklo{polexo}%
\vorto{polexo}\fako{anat.} La fingro maxim grosa e maxim forta di la manui e di la pedi\lingui{EIL}
\artiklo{poli-}%
\vorto{poli-} Sufixo ciencala : « multa-… »
\artiklo{poliandra}%
\vorto{poliandra} Spozino qua havas plura spozuli\lingui{DEFIRS}
\artiklo{polico}%
\vorto{polico}\senco{1}{Regulari facita por mantenar la ordino e la sekureso statala, nacionala}\senco{2}{Surveyo quan onu agas por igar aplikata la regulari, por la manteno di la ordino e di la sekureso statala, nacionala}\lingui{DEFIRS}
\artiklo{polidaktila}%
\vorto{polidaktila}\fako{teratol.} Qua havas fingri supernombra
\artiklo{poliedro}%
\vorto{poliedro}\fako{geom.} Korpo solida qua havas plura edri\lingui{DEFIRS}
\artiklo{poliembriona}%
\vorto{poliembriona}\fako{biol.} Qua formacesis ek du o plura embrioni, da un ovulo
\artiklo{polifonio}%
\vorto{polifonio}\fako{muziko} Efekto muzikala qua rezultas de la ensemblo harmoniala di la instrumenti e di la voci\lingui{DEFIRS}
\artiklo{poligama}%
\vorto{poligama} Spozulo qua havas plura spozini\lingui{DEFIRS}
\artiklo{poliglota}%
\vorto{poliglota}\senco{1}{Redaktita ye plura lingui}\senco{2}{Qua savas plura lingui}\lingui{DEFIRS}
\artiklo{poligono}%
\vorto{poligono}\senco{1}{\textit{(geom.)} Figuro qua havas plura anguli e plura lateri}\senco{2}{\textit{(termino militala)}; (a) figuro qua determinas la trasuro di fortreso, di fortifikajo; (b) placo ube la artileriisti exercas su pri la pafo per kanoni}\lingui{DEFIRS}
\artiklo{poligrafo}%
\vorto{poligrafo} Autoro qua kompozas pri temi diversa\lingui{DEFIS}
\artiklo{polikroma}%
\vorto{polikroma} Qua esas plura-kolora\lingui{DEFIRS}
\artiklo{polimatio}%
\vorto{polimatio} Savo qua embracas plura cienci tre diversa\lingui{DEFIS}
\artiklo{polimorfa}%
\vorto{polimorfa} Qua chanjas ofte sua formo\lingui{DEFIS}
\artiklo{polinomio}%
\vorto{polinomio}\fako{matem.} Quanto algebrala, ek plura termini qui interligesas per la signi « plus » o « minus »\lingui{DEFIS}
\artiklo{polinuklea}%
\vorto{polinuklea}\fako{histol.} Dicesas pri celulo o pri plastido qua kontenas plura nuklei
\artiklo{polipo}%
\vorto{polipo}\senco{1}{\textit{(zool.)} Zoofito di qua la kapo havas periferie palpili}\senco{2}{\textit{(patol.)} Surkreskajo qua su developas en la kavaji qui havas membrano mukoza}\lingui{DEFIRS}
\artiklo{poliporo}%
\vorto{poliporo}\fako{bot.} Genero de fungi bazidiomiceta, karakterizata da chapelo lignatra, qua havas sube tubi sporoza, paralela
\artiklo{polisar}%
\vorto{polisar}\fako{trans.} Igar glata per froto, netigar (armi, koquilaro, e c.) per frotar per smerilo, greso, tri-palio\lingui{DEFIRS}
\artiklo{polispermio}%
\vorto{polispermio}\fako{biol.} Penetro en la ovulo matura di du o plura spermatozoidi
\artiklo{polita}%
\vorto{polita} Karakterizata da lua egardo, indulgo, respekto di sua kunhomi en mondo, en mondumo\lingui{EFS}
\artiklo{politeismo}%
\vorto{politeismo} Religio qua agnoskas plura dei\lingui{DEFIRS}
\artiklo{politeisto}%
\vorto{politeisto} La persono qua profesas politeismo\lingui{DEFIRS}
\artiklo{politeknika}%
\vorto{politeknika} Qua koncernas plura arti, plura cienci
\artiklo{politiko}%
\vorto{politiko}\senco{1}{La arto guvernar stato}\senco{2}{La maniero di guverno}\senco{3}{La maniero konceptar la direkto di la aferi nacionala, municipala}\senco{4}{La maniero prudenta e saja direktar afero komercala, industriala od agrokultivala}\lingui{DEFIRS}
\artiklo{politropo}%
\vorto{politropo}\fako{mineral.} Dicesas pri la kristali di qui la segmenti precipua di la lami inter-sucedante inklinesas ica sur ita, divers-angule\lingui{DEFIRS}
\artiklo{polkar}%
\vorto{polkar}\fako{netrans.} Agar ta danso du-tempa qua importacesis aden la westo-regioni di Europa cirkum la yaro 1845, de Polonia\lingui{DEFIRS}
\artiklo{polo}%
\vorto{polo}\senco{1}{\textit{(astron.)} Singla ek la extremaji di la axo di mondo, cirkum qua semblas movar su, en 24 hori, la sfero ciela}\senco{2}{\textit{(geom.)} Singla ek la extremaji di la axo cirkum qua, supozite, rotacas korpo sfera, elipsatra}\senco{3}{\textit{(geogr.)} Singla ek la du extremaji di la axo di Tero, cirkum qua la Terglobo semblas rotacar en 24 hori}\senco{4}{\textit{(fiziko)} (a) Singla ek la du punti inter-opozita di magneto ube eventas maxima-force la atrakti e la repulsi magnetala, (b) Singla ek la du extremaji inter-opozita di un fluo elektrala}\lingui{DEFIRS}
\artiklo{polodo}%
\vorto{polodo}\fako{geom.} Nomo kreita da Poinsot por indikar la kurvi quin parkuras, sur la elipsoido di inerteso di korpo solida qua povas rotacar libere cirkum punto fixa e qua ne subisas ula forco, la kontakto-punto di la elipsoido kun la plano fixa qua esas paralela ye la plano di la maximo di la arei sur qua lu rulas
\artiklo{polonezo}%
\vorto{polonezo} Danso nacionala di Polonia, nobla, kun marcho vice salteti, preske pompoza, segun ario tri-tempa, segun movimento mezvalora\lingui{DEFRS}
\artiklo{polpo}%
\vorto{polpo}\fako{zool.} Molusko cefalopoda, quaze sako viskoza kun grosa kapo de ube departas ok palpili\lingui{EFIS}
\artiklo{polstero}%
\vorto{polstero}\senco{1}{Amaso de pili, quin onu desligis de la felo di ula animali (bovyuni, bovini, kavali, e c.) e quin onu uzas por garnisar la seli, la tabureti, e c.}\senco{2}{La parto maxim grosiera di lano, quan onu ne filigas e ye qua onu garnisas la matraci, la kuseni}\lingui{DEIR}
\artiklo{poltrono}%
\vorto{poltrono} Persono qua pavoras exajere ed agas konseque\lingui{DEFI}
\artiklo{polutar}%
\vorto{polutar}\fako{netrans.) (patol.} Emisar nevole spermo, jorne o dumdorme, konseque de kontinenco o de onanio\lingui{DEFIRS}
\artiklo{polvo}%
\vorto{polvo}\senco{1}{Tero, reduktita ad esar partikuli tenua, quin aero, se agitata mem nur poke, sublevas e transportas}\senco{2}{Partikuli tenua}\lingui{IS}
\artiklo{pomado}%
\vorto{pomado}\senco{1}{Kosmetikalajo qua konsistis, origine, ye graso e pulpo de pomi; nun, ek grasajo a qua onu adjuntas mixe esenci, parfumi}\senco{2}{Preparajo farmaciala : grasajo a qua onu adjuntis mixe substanco medikamenta}\lingui{DEFIRS}
\artiklo{pomelo}%
\vorto{pomelo}\senco{1}{Kapo mi-sfer-forma di la mancho di spado}\senco{2}{Saliajo mi-sfer-forma, meze di sel-arko, avan ica}\lingui{EFIS}
\artiklo{pomerhundo}%
\vorto{pomerhundo}\fako{zool.} Hundo-raso de Pomerania
\artiklo{pomo}%
\vorto{pomo}\fako{bot.} Frukto sfera, kun pulpo abundanta, di qua la suko, kande fermantacinta, donas cidro\latina{pirus malus}\lingui{FI}
\artiklo{pompelmuso}%
\vorto{pompelmuso}\fako{bot.} Speco di « citrus » di qua la frukto, grosega, konsumesas da la Usani e da la Britaniani lor lia dejuneto. Lu kultivesas omnaloke en la regioni tropikala\latina{citrus decumana}\lingui{DFIS}
\artiklo{*pompiero}%
\vorto{*pompiero} La persono (uniformoza), ek soldato-korporaciono, di qua la mestiero konsistas ye (precipue) extingar la incendii, salvar la personi quin minacas la inund-aquo, od irga vivodanjero acidentala, e c.
\artiklo{pompo}%
\vorto{pompo}\senco{1}{Procesiono solena}\senco{2}{Expozo di ulo splendida, luxozega, quan onu vidigas pro ostento}\lingui{DEFIS}
\artiklo{pompono}%
\vorto{pompono} Tufo per qua onu ornas la kapvesti, la vestizuri hominala. Tufo de lano quan la soldati *weras somite di sua kepio\lingui{DEFS}
\artiklo{*pondar}%
\vorto{*pondar}\fako{aludante ucelino} Emisar sua ovi\lingui{F}
\artiklo{ponderar}%
\vorto{ponderar}\fako{trans.}\senco{1}{Mezurar la pezo di korpo, per komparar lu a pez-unajo}\senco{2}{\textit{(metaf.)} Judikar pri la valoro di kozo, per explorar la « por » e la « kontre »}\lingui{EFIS}
\artiklo{poneo}%
\vorto{poneo}\fako{zool.} Kavalo di raso karakterizata da staturo ne-alta\lingui{DEF}
\artiklo{poniardo}%
\vorto{poniardo} Armo atakala, kun lamo kurta, akuta, fixigita an mancho, quan onu tenas per sua pugno\lingui{EFIS}
\artiklo{ponsa}%
\vorto{ponsa} Koloro dazlante-reda, kelke simila a ta di papavereto\lingui{DFIRS}
\artiklo{pontifiko}%
\vorto{pontifiko}\senco{1}{Pastoro di religio}\senco{2}{Che la katoliki : episkopo}\lingui{DEFIRS}
\artiklo{ponto}%
\vorto{ponto} Konstrukturo quan onu uzas por pasar de ca a ta rivo di rivero, di fluvio, di fosato, di granda kavajo en la surfaco di sulo\lingui{FIS}
\artiklo{pontono}%
\vorto{pontono}\senco{1}{Ponto flotacanta, facita ek du bateli interjuntita per trabi}\senco{2}{Konstrukturo uzata kom staciono di departo o di arivo por la navi qui transportas pasajeri}\lingui{DEFIRS}
\artiklo{ponyeto}%
\vorto{ponyeto} Bendo a qua abutas la maniko di kamizo, di robo, di korsajo, en la loko ube lu kovras la karpo\lingui{F}
\artiklo{poplino}%
\vorto{poplino} Stofo di qua la warpo-fili esas silka, e la wefto-fili, lana\lingui{DEFIRS}
\artiklo{poplito}%
\vorto{poplito}\fako{anat.} Parto di la gambo, qua esas dop la genuo, opoze ad ica\lingui{ISL}
\artiklo{poplo}%
\vorto{poplo}\fako{bot.} Arboro ek la familio « saliki », kun stipo rekta, e folii alternanta\latina{populus}\lingui{DEFI}
\artiklo{popo}%
\vorto{popo} Sacerdoto di la rituo Greka, che la Rusi, la Serbi e la Bulgari
\artiklo{populara}%
\vorto{populara}\senco{1}{Qua esas komprenebla (kom texto) da la personi ordinara, ne-specalista, ne-ciencista}\senco{2}{Pri qua la populo opinionas favoroze}\lingui{DEFIRS}
\artiklo{populo}%
\vorto{populo}\senco{1}{Ensemblo ek la personi qui habitas la sama lando, esas (maxim-multa-kaze) sama-rasa, ed uzas la sama linguo}\senco{2}{La korpo di la naciono, opoze a la suvereno ed a la nobeli}\lingui{EFIS}
\artiklo{popurito}%
\vorto{popurito} Verko muzikala, literaturala, ek texti de origini diversa\lingui{DEF}
\artiklo{por}%
\vorto{por} Profite da… Destine di…
\artiklo{porcelano}%
\vorto{porcelano} Ceramikajo blanka, tre dina, importacita de la esto-regioni di Azia (precipue Chinia) aden Europa en la yarcento 16-esma\lingui{DEFIS}
\artiklo{porciono}%
\vorto{porciono} Parto qua destinesas a singlu, lor disdono o distributo\lingui{DEFIRS}
\artiklo{pordo}%
\vorto{pordo}\senco{1}{Aperturo facita en la muro di klozajo, di edifico, di chambro, por povar enirar od ekirar}\senco{2}{Klapo qua klozas ta aperturo e povas apertesar por paso}\lingui{EFIS}
\artiklo{porelo}%
\vorto{porelo}\fako{bot.} Supo-planto, di la genero « alio »\latina{allium porrum}\lingui{DFIRS}
\artiklo{porfiro}%
\vorto{porfiro}\fako{mineral.} Roko harda, obskure-reda, en qua dissemesas makuli blanka, quan la antiqui obtenis de la regioni alte-situita di Egiptia\lingui{DEFIS}
\artiklo{porismo}%
\vorto{porismo}\fako{geom.) (antique} Teoremo nekompleta, qua establisas proprajo komuna inter ula elementi quin la enunco determinas, ed altra teoremi qui rezultas de la hipotezo\lingui{DEFI}
\artiklo{porkespino}%
\vorto{porkespino}\fako{zool.} Mamifero rodera di qua la korpo havas kom armi pikili\latina{hystrix}
\artiklo{porko}%
\vorto{porko}\fako{zool.} Mamifero domestika quan onu repletigas por manjor lu. (scrofa domestica)\latina{sus}\lingui{EFIS}
\artiklo{pornografio}%
\vorto{pornografio}\senco{1}{Traktato pri prostituco}\senco{2}{Verki literatura, arta, quin – kontre-etike – onu igas servar libertinismo}
\artiklo{pornografo}%
\vorto{pornografo}\senco{1}{Autoro qua kompozas texti pri prostituco}\senco{2}{Autoro, piktisto, obcena}\lingui{DEFIS}
\artiklo{poro}%
\vorto{poro}\senco{1}{Singla ek la orifici preske neperceptebla di la pelo di animalo, per qui eventas la transpiro}\senco{2}{Singla ek la orifici analoga, ye qui truizesas la vejetanti}\senco{3}{Singla ek la interstici qui interseparas la molekuli di korpo ed igas ica plu o min permeebla}\lingui{DEFIRS}
\artiklo{portalo}%
\vorto{portalo}\senco{1}{Quaza vestibulo, multa-kaze ornita per koloni, en la enireyo di templo, di kirko, di palaco}\senco{2}{Fasado di kirko en qua esas la pordo precipua}\lingui{EF}
\artiklo{portar}%
\vorto{portar}\fako{trans.}\senco{1}{Tenar (inter sua manui, inter sua brakii, sur sua shultri) ulo pezoza}\senco{2}{Prenar kun su, e depozar en ta o ca loko}\lingui{EFIS}
\artiklo{portaveino}%
\vorto{portaveino}\fako{anat.} Veino qua recevas sango de la viceri abdominala (ventrala), quan lu distributas aden la hepato\lingui{EF}
\artiklo{porteo}%
\vorto{porteo} Punto til qua la fusilo povas portar la kuglo – la kanono, la obuso\lingui{F}
\artiklo{portfolio}%
\vorto{portfolio} Sorto di buxo kartona, kunportebla, destinita ye recevor paperfolii, desegnuro-folii, dokumenti\lingui{DEFI}
\artiklo{portiko}%
\vorto{portiko} Konstrukturo uzata por la exerci gimnastikala, ek fosti qui suportas trabo transversa a qua esas suspendita la rigi\lingui{DEFIRS}
\artiklo{portreto}%
\vorto{portreto} Reprezenturo de persono per pikto o desegno sur moneto-peco, medalio, telo, e c.\lingui{DEFR}
\artiklo{portulako}%
\vorto{portulako}\fako{bot.} Supo-planto di qua la folii esas dika\latina{portulaca}\lingui{DEIRL}
\artiklo{portuo}%
\vorto{portuo}\senco{1}{Konkavajo di maro aden tero – naturala od artificala, ube la navi povas shirmar su}\senco{2}{\textit{(metaf.)} Shirmeyo}\senco{3}{Loko ube onu embarkas o desembarkas la pasajeri e la vari transportita per aviacili (aeroplani, helikopteri, e c.) o navi}\lingui{EFIRS}
\artiklo{pos}%
\vorto{pos} En epoko, instanto, qua sequas la evento quan onu mencionas\lingui{FL}
\artiklo{posedar}%
\vorto{posedar}\fako{trans.} Havar la fakultato juar (ulo od ulu), mem sen proprietar lu\lingui{EFIS}
\artiklo{posho}%
\vorto{posho} Mikra sako de telo, de stofo, sutita en la dikajo di vesto, por pozar aden lu to quon onu volas portar kun su
\artiklo{posibla}%
\vorto{posibla} Qua povas eventar\lingui{EFIS}
\artiklo{posteno}%
\vorto{posteno}\senco{1}{\textit{(milit-arto)} Loko, spaco, okupata da trupo de militisti por operaco militala}\senco{2}{Plaso asignita a kombatanti}\lingui{DEFIRS}
\artiklo{posterno}%
\vorto{posterno}\fako{milit-arto} Galerio vultoza, subsula, en fortifikajo por komunikar kun la parto extera\lingui{DEFIS}
\artiklo{postigo}%
\vorto{postigo} Vitro-plako di fenestro, quan onu povas apertar sen apertar la tota fenestro, por parolar ad ulu, por lasar la aero enirar, e c.\lingui{S}
\artiklo{postiliono}%
\vorto{postiliono}\senco{1}{La duktisto di posto-veturo, di qua la vesti esis maxim-multa-kaze di koloro okul-frapanta, kun chapelo rubandoza, e qua kavalkis sur un de la kavali jungita}\senco{2}{Duesma duktisto di karoso tirata da 4 o 6 kavali – la viro qua duktas la kavali avana}\lingui{DEFIRS}
\artiklo{posto}%
\vorto{posto} Transporto nacionala di la letri, imprimuri, dokumenti, mikra paki\lingui{DEFIRS}
\artiklo{postpoziciono}%
\vorto{postpoziciono}\fako{gram.} Partikulo lokizita pos la vorto a qua lu relatas
\artiklo{postular}%
\vorto{postular}\fako{trans., de}\senco{1}{Demandar rigoroze (ulo) konseque de sua yuro, de sua autoritato, de sua forteso}\senco{2}{Igar nekareebla}\lingui{FL}
\artiklo{postulato}%
\vorto{postulato} Propoziciono quan onu prizentas kom egardenda, pro lua exakteso, sen demonstrar lu\lingui{DEFIRS}
\artiklo{postuma}%
\vorto{postuma} Naskinta pos la morto di sua patro. \textit{(anke metaf.)}\lingui{DEFIS}
\artiklo{posturar}%
\vorto{posturar}\fako{netrans.}\senco{1}{\textit{(\textit{aludante persono qua esas avan piktisto, fotografisto qui esas pronta operacar})} Movar sua korpo, sua muskuli, modifikar la traiti di sua vizajo, quaze agas lo modelo, segun la cirkonstanci}\senco{2}{Igar su aspektar plu valoroza kam onu esas reale}\lingui{DEFIS}
\artiklo{potaso}%
\vorto{potaso}\fako{kemio} Karbonato di kalio quan onu extraktas de la cindro de ula vejetanti\simbolo{2 KOH}\lingui{DEFIRS}
\artiklo{potencialo}%
\vorto{potencialo}\fako{mekan.}\senco{1}{Sumo de la forci qui esas necesa por startar sistemo}\senco{2}{\textit{(fiziko)} Intenseso-grado di la forco necesa por retroduktar molekulo deformita a lua formo prima}\lingui{DEFIS}
\artiklo{potenciometro}%
\vorto{potenciometro}\fako{elektro} Aparato, inventita da Clark, analoga a la reostato de Wheatstone, uzata por mezurar la difero di potenciali inter la poli di baterio
\artiklo{potenco}%
\vorto{potenco}\fako{matem.} Valoro di quanto nombrala, multiplikita un, du o tri-foye per lu ipsa\lingui{DEFIR}
\artiklo{potenta}%
\vorto{potenta}\senco{1}{Qua disponas la povo produktar granda efekti}\senco{2}{Qua povas impozar sua autoritato}\lingui{DEFIS}
\artiklo{poto}%
\vorto{poto} Menajo-vazo di qua la materio, la formo, la dimensioni varias segun la uzi a qua lu esas destinita\lingui{EFS}
\artiklo{povar}%
\vorto{povar}\fako{trans. e netrans.} Esar en la stando necesa por agar e facar ulo, pro havar la forteso e la habileso necesa\lingui{FIS}
\artiklo{povra}%
\vorto{povra} Qua ne havas lo necesa por suficar a su. (antonimo : richa)\lingui{EFIS}
\artiklo{pozar}%
\vorto{pozar}\fako{trans.}\senco{1}{Situar en ta o ca plaso}\senco{2}{Situar en ca o ta plaso, por permanar en ta loko}\lingui{FIS}
\artiklo{poziciono}%
\vorto{poziciono}\fako{milit-arto} Tereno okupata da militisto-trupo, skope di agor batalio\lingui{EF}
\artiklo{pozitiva}%
\vorto{pozitiva}\subvorto{algebro}{quanto pozitiva}{Supera ad zero, ed indikata per la signo « plus »}{1}{}\subvorto{fiziko}{elektro pozitiva}{Glas-elektro quan onu konsideris kom produktata da eceso di fluido}{2}{}\senco{3}{\textit{(gram.)} Qua esas en la grado unesma, ye qua la qualeso afirmesas simple (opoze a « komparativo », « superlativo »)}\lingui{DEFIRS}\subvorto{filoz.}{valoro pozitiva}{Doktrino filozofiala qua repulsas omna absolutajo, pro agnoskar exkluzive la fakti qui konstatesis per experimento}{4}{}
\artiklo{pozolano}%
\vorto{pozolano} Sablo vulkanala, redatra, qua uzesas por facar cemento\lingui{DEFIS}
\artiklo{pragmatiko}%
\vorto{pragmatiko} Akto solena qua regulizas la interesti di la familio di la monarko en la stato
\artiklo{pragmato}%
\vorto{pragmato} Fakto realigita, facita, pruvita nedubiteble\lingui{DEFIRS}
\artiklo{praktikar}%
\vorto{praktikar}\fako{trans.} Aplikar la reguli di sua mestiero, di sua arto, di sua cienco, e c.\lingui{DEFIS}
\artiklo{pralino}%
\vorto{pralino} Mandelo risoligita en sukro\lingui{DEF}
\artiklo{pramo}%
\vorto{pramo}\fako{navig.} Granda navo kun remili e segli, di qua la fundo esas plata, e kun nur un ferdeko\lingui{E}
\artiklo{prato}%
\vorto{prato} Tereno en qua esas semita diversa planti forejala quin onu falchas o quin onu igas la brutaro konsumar en la kreskeyo ipsa\lingui{FIS}
\artiklo{pre-}%
\vorto{pre-} Sufixo qua equivalas « ante-… », ma kun ideo implicita di kontakto kun to quon indikas la radiko
\artiklo{prebendo}%
\vorto{prebendo}\fako{eklezio katol.} Revenuo di kanoniko\lingui{DEFIS}
\artiklo{precedar}%
\vorto{precedar}\fako{trans.} Praktikar ante ulu (en ofico, profesiono, komerco-fako, e c.)\lingui{EFIS}
\artiklo{precedento}%
\vorto{precedento} Fakto eventinta, a qua onu referas su por justifikar su pri eventigar o justifikar fakto analoga\lingui{DEFIS}
\artiklo{precepto}%
\vorto{precepto} Formulo qua docas quale onu devas kondutar. Fakto etikala quan onu propozas kom regulo di konduto\lingui{EFIS}
\artiklo{preceptoro}%
\vorto{preceptoro} Persono a qua onu konfidis la eduko privata di infanto, di puero\lingui{DEFIS}
\artiklo{precesiono}%
\vorto{precesiono}\fako{astron.}\lingui{DEFIS}\subvorto{}{precesiono di la equinoxi}{Diplaso pokopa di la axo di Tero, pro quo la equinoxo printempala eventas ante la ri-eso di Tero en la sama punto di lua orbito}{}{}
\artiklo{precipiso}%
\vorto{precipiso}\senco{1}{Pento tre eskarpa}\senco{2}{\textit{(metaf.)} Kauzo di ruino}\lingui{EFIS}
\artiklo{precipitar}%
\vorto{precipitar}\fako{trans.}\senco{1}{Faligar bruske, de loko situita alte, ad fundo}\senco{2}{\textit{(metaf.)} (1) Faligar subite, de socio-grado alta; (b) faligar subite a situeso funesta}\senco{3}{\textit{(kemio)} Retrofaligar a la fundo di la vazo qua kontenas liquido, solidon qua esas dissolvita en lu}\senco{4}{Igar irar impetuoze}\lingui{DEFIS}
\artiklo{precipua}%
\vorto{precipua}\senco{1}{La maxim importoza, aludante (1) persono, (b) kozo}\lingui{L}\subvorto{gram.}{propoziciono precipua}{Opoze a la propoziciono incidenta}{2}{}
\artiklo{preciza}%
\vorto{preciza}\senco{1}{\textit{(imajo, ideo)} Strikte inkluzita, limitizita, determinita}\senco{2}{\textit{(tranchuro, sekuro, rupturo)} di qua nula saliajo igas nedicidita la konturi}\lingui{DEFIRS}
\artiklo{preco}%
\vorto{preco} Valor-quanto di kozo, di persono : (1) valoro komercala; (2) valoro etikala\lingui{DEFIS}
\artiklo{predikar}%
\vorto{predikar}\fako{trans.}\senco{1}{Anuncar lo dicita da Deo}\senco{2}{Instruktar (ulu) per docar lo dicita da Deo}\senco{3}{Diskursar, de katedro kristanismala, por la edifiko di la asistanti}\lingui{DFIS}
\artiklo{predikato}%
\vorto{predikato}\fako{logiko} Atributo di propoziciono\lingui{DEFIS}
\artiklo{prefaco}%
\vorto{prefaco} Expozo-texto preliminara quan autoro insertas ante la komenco-parto di sua verko ed en qua lu savigas olua skopo, olua karaktero, e c.\lingui{EFIS}
\artiklo{prefekto}%
\vorto{prefekto}\senco{1}{La persono qua esas la chefo di stat-administrofako en distrikto determinita}\senco{2}{La persono qua administras departmento, en Francia}\lingui{DEFIRS}
\artiklo{preferar}%
\vorto{preferar}\fako{trans., ulu, ulo, kam ulu, kam ulo} Amar, amorar, prizar, diletar plu multe (kam onu agas lo pri altru od altro)\lingui{EFIS}
\artiklo{prefixo}%
\vorto{prefixo}\fako{gram.} Partikulo qua, lokizite avan vorto, modifikas la senco di olca e formacas vorto nova\lingui{DEFIRS}
\artiklo{prefoliaciono}%
\vorto{prefoliaciono}\fako{bot.} Dispozeso-maniero di la folii en la burjono\lingui{DEFIRS}
\artiklo{pregar}%
\vorto{pregar}\fako{trans., ulu, pri o por}\senco{1}{Adorar la deajo, Deo, demandante a lu favoro od indulgo}\senco{2}{Urjar (ulu) pri ke lu grantez}\lingui{FI}
\artiklo{*prei}%
\vorto{*prei} Demando plu insistoza kam *pliz – preska supliko\lingui{E}
\artiklo{prekara}%
\vorto{prekara}\fako{yuro-cienco} Grantita segun-requizite, e sempre revokebla\lingui{DEFIS}
\artiklo{prekoca}%
\vorto{prekoca}\senco{1}{Di qua la matureso esas kresko-*anticipa}\senco{2}{\textit{(metaf.)} Plu developita kam esas normala segun lua evo}\lingui{EFIS}
\artiklo{prekursoro}%
\vorto{prekursoro} Ta qua preparas la veno di altru\lingui{EFIS}
\artiklo{prelato}%
\vorto{prelato}\fako{eklezio katol.}\senco{1}{Oficiero ekleziala alta-grada}\senco{2}{Oficiero ek la oficieraro papala, yurizita pri *werar vesti violea}\lingui{DEFIRS}
\artiklo{preliminaro}%
\vorto{preliminaro} Singla ek la elementi di to quo preiras e preparas la temo precipua\lingui{DEFIRS}
\artiklo{preludar}%
\vorto{preludar}\fako{netrans.) (muziko} Probar sua voco, sua muzik-instrumento per kantar, plear serio de noti\lingui{DEFIS}
\artiklo{premici}%
\vorto{premici} La frukti dat-unesma di tero, di brutaro, quin onu ofras a la deaji\lingui{FIS}
\artiklo{premio}%
\vorto{premio} To quo donacesas kom rekompenso a persono qua vinkas lor konkurso\lingui{DEFIRS}
\artiklo{premiso}%
\vorto{premiso}\fako{logiko} Singla ek la du propozicioni (majoro e minoro) di silogismo de qua onu inferas la konkluzo. (Uzesas precipue plurale)\lingui{DEFIS}
\artiklo{prenar}%
\vorto{prenar}\fako{trans., de, ek} En-manuigar por tenar od uzar\lingui{FI}
\artiklo{preparar}%
\vorto{preparar}\fako{trans., por} Igar apta ye facar od agar ta a quo lu esas destinita\lingui{DEFIRS}
\artiklo{preponderar}%
\vorto{preponderar}\fako{netrans., pri, en, che} Esar plu eminenta pri autoritato, influo; superesar (ulu, ulo) pri dominaco, vinko\lingui{DEFIS}
\artiklo{prepoziciono}%
\vorto{prepoziciono}\fako{gram.} Vorto-speco qua, lokizita avan nomo, avan propoziciono infinitivala, e c., ligas ici per raporto determinita a termino qua preiras\lingui{DEFIRS}
\artiklo{prepuco}%
\vorto{prepuco}\fako{anat.} Prolonguro di la penis-pelo qua tegas la glano\lingui{EFIS}
\artiklo{prerogativo}%
\vorto{prerogativo} Yuro atribuita ad ula rangi nobelesala privilejiera\lingui{DEFIRS}
\artiklo{presar}%
\vorto{presar}\fako{trans.} Klemar per apogar su pezoze (sur lu)\lingui{DEFIRS}
\artiklo{presbita}%
\vorto{presbita}\fako{patol.} \textit{(persono)} qua povas vidar nur to quo esas for lu, pro ke lua muskulo ciliara cesis donar a la kristalino la kurviguro necesa\lingui{DEFIS}
\artiklo{presbiterio}%
\vorto{presbiterio}\fako{religio} Asemblitaro ek la personi qui administras la proprietaji e la revenui di la parokio\lingui{DEFIRS}
\artiklo{presbiterismo}%
\vorto{presbiterismo}\fako{religio} Doktrino di ula sekto Kristana qua ne agnoskas la episkopi, e grantas la direkto di la eklezio a la sacerdoti o pastori\lingui{DEFIRS}
\artiklo{preske}%
\vorto{preske} « Poko mankas pri ke lo esez tale »
\artiklo{preskriptar}%
\vorto{preskriptar}\fako{trans., ulo ad ulu} Imperar per formulo preciza\lingui{EFIS}
\artiklo{preskritar}%
\vorto{preskritar}\fako{trans.) (yuro-cienco} Liberigar de debajo, de persequi judiciala, de servitudo, per la fakto ke quanto determinita de tempo fluis de lore\lingui{FIS}
\artiklo{prestaciono}%
\vorto{prestaciono} Laboro kelka-dia (redemtebla per pekunio) postulebla de la habitanti di komono, por la konservo en bona stando di la voyi publika qui apartenas a la komono\lingui{DEFIS}
\artiklo{prestar}%
\vorto{prestar}\fako{trans., ad} Igar (ulo) disponebla, posedata, da ulu dum ula quanto de tempo, pos qua lu devos retrodonar olu\lingui{FIS}
\artiklo{prestijo}%
\vorto{prestijo} To quo frapas, astonas, surprizas per sua brilo, e c. Autoritato od influo qua naskas de la aspekto extera, ed efikas a la opiniono ed a la imagino-fakultato di ula personi\lingui{DEFIRS}
\artiklo{pretendar}%
\vorto{pretendar}\fako{trans.} Persequar publike ulo pri quo onu dicas esar yuroza\lingui{DEFIRS}
\artiklo{preter}%
\vorto{preter} Pasante apud (ulu od ulo) ed irante plu ad-avane kam (li)\lingui{L}
\artiklo{preterito}%
\vorto{preterito}\fako{gram.} Tempo di verbo, qua expresas la pasinto\lingui{DEFIS}
\artiklo{pretextar}%
\vorto{pretextar}\fako{trans.} Alegar motivo qua aspektas exakta, por celar la motivo reala di ago\lingui{EFIS}
\artiklo{pretoro}%
\vorto{pretoro}\fako{en Roma, antique} Oficiero judiciala qua direktis la operaci militala, administris provinco, e c.\lingui{DEFIS}
\artiklo{prevarikar}%
\vorto{prevarikar}\fako{netrans.} Deviacar de la konduto yuro-konforma\lingui{FIS}
\artiklo{preventar}%
\vorto{preventar}\fako{trans.} Impedar la evento, la existesko\lingui{EFIS}
\artiklo{prevosto}%
\vorto{prevosto}\senco{1}{Oficiero di qua la tasko konsistas ye judiciar pri la kazi kriminala en la armei}\senco{2}{Chefo di la kanonikaro en kirko kolegiala, superioro en ula ordeni religiala}\lingui{EFI}
\artiklo{prezenta}%
\vorto{prezenta}\senco{1}{Qua esas en loko kande ulu arivas an olu, od esas en olu}\senco{2}{Qua eventas, qua existas en la periodo en qua onu esas}\lingui{DEFIS}\subvorto{gram.}{prezento}{Tempo di la verbo, qua expresas la tempo-punto ube onu esas}{3}{}
\artiklo{prezervar}%
\vorto{prezervar}\fako{trans., de, kontre} Shirmar, protektar, de atingeso da malajo\lingui{DEFIRS}
\artiklo{prezidar}%
\vorto{prezidar}\fako{trans.} Direktar la laboro, la debati da asemblitaro (chambro deliberala, *konselio, e c.)\lingui{DEFIRS}
\artiklo{pri}%
\vorto{pri} Koncerne, relate, segun quante koncernesas…
\artiklo{priklar}%
\vorto{priklar}\fako{netrans.} Sentar su quaze se pikata neprofunde ed iteroze\lingui{DE}
\artiklo{prima}%
\vorto{prima} Egardata kom origino\lingui{DEFIRS}
\artiklo{primadono}%
\vorto{primadono} Kantistino famoza, tre bone reputata pri lua talento\lingui{DEFIRS}
\artiklo{primara}%
\vorto{primara}\senco{1}{Qua apartenas a la grado unesma (departante de infre)}\lingui{DEFIS}\subvorto{}{kuzi primara}{Qui naskis de patri interfrata}{2}{}
\artiklo{primato}%
\vorto{primato}\fako{eklezio katol.} Arkiepiskopo qua dominacas fakte ed (od) honore la arkiepiskopi ed episkopi di regiono\lingui{DEFIS}
\artiklo{primino}%
\vorto{primino}\fako{bot.} Tegilo unesma di la ovulo (extera)
\artiklo{primitiva}%
\vorto{primitiva} Qua aparis ye la origino, e duras havar de ta origino, ta o ca karaktero\lingui{DEFIRS}
\artiklo{primrozo}%
\vorto{primrozo}\fako{bot.} Planto qua florifas en la komenco-dii di printempo\latina{primula grandiflora}\lingui{E}
\artiklo{primulo}%
\vorto{primulo}\fako{bot.} Narciso sovaja, kun mikra flori\latina{primula officinalis}\lingui{DEFIS}
\artiklo{principo}%
\vorto{principo} Exaktajo fundamenta, sur qua su apogas la rezono\lingui{DEFIRS}
\artiklo{princo}%
\vorto{princo}\senco{1}{Persono qua esas suvereno od apartenas a familio de suvereni}\senco{2}{Spozo di princo}\senco{3}{\textit{(metaf.)} Persono qua esas la unesma hierarkiale, la chefo (di la filozofi, di la oratori)}\lingui{DEFIRS}
\artiklo{printempo}%
\vorto{printempo} Sezono di la yaro, qua sequas nemediate vintro, dum qua la temperaturo di atmosfero divenas plu dolca e rinaskas la vejetantaro. (intertempo qua separas la equinoxo qua sequas nemediate vintro, de la solstico di somero). – \textit{(anke metaf.)}\lingui{EF}
\artiklo{priora}%
\vorto{priora} Qua havas la yuro venar ante od avan\lingui{DEFIS}
\artiklo{prioro}%
\vorto{prioro}\fako{eklezio katol.} Superioro di kuvento, di monakeyo\lingui{DEFIRS}
\artiklo{prismato}%
\vorto{prismato}\senco{1}{\textit{(geom.)} Poliedro di qua la bazo infra e la bazo supra esas poligoni paralela egala; e la edri laterala, paralelogrami qui interunionas la lateri homologa di la bazo}\senco{2}{\textit{(fiziko)} Prismo triangula de kristalo qua, kande lumo-radii trairas lu, refraktas ta radii e deskompozas li aden lia kolori prima}\lingui{DEFIRS}
\artiklo{pritaneo}%
\vorto{pritaneo}\fako{en la epoki antiqua} Edifico en qua esis lojigata ed entratenata, da la stato e lua-spense, la civitani qui tre meritis de la patrio\lingui{DEFIRS}
\artiklo{privacar}%
\vorto{privacar}\fako{trans., de} Afliktar (ulu) per igar lu ne plus havar, od indijar (ulo)\lingui{EFIS}
\artiklo{privata}%
\vorto{privata} Ne publika\lingui{DEFIR}
\artiklo{privilejo}%
\vorto{privilejo} Avantajo personala, individuala, qua grantesis kom ecepto ye la yuro komuna\lingui{DEFIRS}
\artiklo{prizar}%
\vorto{prizar}\fako{trans.} Judikar (ulo) kom agreabla a su\lingui{DEF}
\artiklo{prizentar}%
\vorto{prizentar}\fako{trans., ad}\senco{1}{Pozar (ulo) avan ulu por ke lu prenez la kozo}\senco{2}{Lokizar (ulu) avan ulu por ke ica videz lu}\lingui{DEFIRS}
\artiklo{pro}%
\vorto{pro} Efekte de, konseque de, efike da, kauze de
\artiklo{probabla}%
\vorto{probabla} Di qua la exakteso havas plu multa chanci « por » kam chanci « kontre »\lingui{EFIS}
\artiklo{probar}%
\vorto{probar}\fako{trans.} Aplikar, kom foyounesma, kozo a lua destineso, por povar judikar ka lu esas apta por lo\lingui{DFIS}
\artiklo{problemo}%
\vorto{problemo} Questiono desfacile solvebla. (Onu prizentas problemo)\lingui{DEFIRS}
\artiklo{procedar}%
\vorto{procedar}\fako{netrans., ad} Komencar la exekuto di ta kozo\lingui{EFIS}
\artiklo{proceduro}%
\vorto{proceduro}\fako{yuro-cienco} Formo segun qua onu devas procedar en korto judiciala\lingui{DEFIR}
\artiklo{procento}%
\vorto{procento} Raporto (proporciono) a cent unaji. (Onu parolas pri procenti en statistiki omna-speca : procento di la naski, di la morti, di la mariaji en populo od habitantaro. Do, ne la interesto o rento di kapitalo. Ex. : Me kolokis 1000 franki ye la procento di 4 \textit{po cent})\lingui{DEFIRS}
\artiklo{procesar}%
\vorto{procesar}\fako{netrans., kontre} Persequar (ulu) en korto judiciala\lingui{DEFIRS}
\artiklo{procesionar}%
\vorto{procesionar}\fako{netrans.}\senco{1}{\textit{(aludante la klerikaro od amaso de personi qui defilas ordinoze)} Marchar kolone en la kirko od extere, kantante psalmi, himni, litanii, okazione di ula festi religiala}\senco{2}{\textit{(metaf.)}\textit{(aludante longa serio de personi qui avancas ganso-marche)} Marchar ica dop ita, unope}\lingui{DEFIRS}
\artiklo{procidar}%
\vorto{procidar}\fako{netrans.) (teol.} Exekutar la ago per qua la Santa Spirito pasas de la Patro a la Filio (en la religio katolika)
\artiklo{proda}%
\vorto{proda} Qua agis plu kam brave\lingui{FIS}
\artiklo{prodigar}%
\vorto{prodigar}\fako{trans.}\senco{1}{Spensar sen modereso, sen limito}\senco{2}{\textit{(metaf.)} Donacar sen modereso, sen limito}\lingui{EF}
\artiklo{prodromo}%
\vorto{prodromo}\fako{medic.} Signo qua anuncas morbo, maladesko\lingui{DEFIS}
\artiklo{produktar}%
\vorto{produktar}\fako{trans.} Igar existar, igar naskar\lingui{DEFIRS}
\artiklo{produto}%
\vorto{produto}\fako{matem.} Rezultajo de la multipliko di nombro per altra\lingui{F}
\artiklo{profana}%
\vorto{profana}\senco{1}{Qua esas exter la kozi sakra}\senco{2}{\textit{(metaf.)} Qua esas exter ula cienco, ula arto}\lingui{DEFIS}
\artiklo{profanacar}%
\vorto{profanacar}\fako{trans., per} Violacar la santeso di la sakraji. – (metaf.) Violacar la respekto debata ad ulu od ulo\lingui{DEFIRS}
\artiklo{profazo}%
\vorto{profazo}\fako{biol.} Fazo unesma di mitoso
\artiklo{profermento}%
\vorto{profermento}\fako{biol.} Granularo fermentala qua, por divenar efikiva e plear sua rolo zimazala, bezonas la prezenteso di substanco preparanta
\artiklo{profesar}%
\vorto{profesar}\fako{trans.} Deklarar publike (sua fido religiala, sua ateisteso, sua aprobo di…)\lingui{EFIS}
\artiklo{profesiono}%
\vorto{profesiono} Sorto di agado per qua onu ganas sua vivo-necesaji\lingui{DEFIRS}
\artiklo{profesoro}%
\vorto{profesoro} Persono qua docas arto, cienco, en skolo superiora (liceo, fakultato di yuro-cienco, di medicino)\lingui{DEFIRS}
\artiklo{profeto}%
\vorto{profeto}\senco{1}{\textit{(religio)} La persono quan Deo inspiras, qua presavigas la eventontaji, o revelas fakti celita de la personi cetera}\senco{2}{La persono qua presavigas la eventontaji}\lingui{DEFIS}
\artiklo{profilaxio}%
\vorto{profilaxio}\fako{medic.} Rejimo qua prezervas kontre la malaji, kontre la morbi\lingui{DEFIRS}
\artiklo{profilo}%
\vorto{profilo}\senco{1}{Aspekto di vizajo qua videsas de latere}\senco{2}{Profilo di edifico, di fortreso, di tereno, di muluro : (a) lia aspekto kande vidata de latere; (b) lia seciono sur plano, segun perpendiklo}\lingui{DEFIRS}
\artiklo{profitar}%
\vorto{profitar}\fako{trans.}\senco{1}{Ektirar avantajo de ulo}\senco{2}{Obtenar gano de ulo}\lingui{DEFIR}
\artiklo{profunda}%
\vorto{profunda}\senco{1}{Di qua la fundo esas tre infra kompare kun la orifico, kun la surfaco}\senco{2}{Qua su extensas tre adfore, de lo avana til lo dopa}\senco{3}{\textit{(metaf.)} (1) Qua exploras absolute omna faktori di la kozi. (2) Qua esas extrema}\lingui{EFIS}
\artiklo{prognozar}%
\vorto{prognozar}\fako{trans.) (medic., meteorol.} Konjektar, de ula signi, to quo mustas eventor\lingui{DEFIRS}
\artiklo{programo}%
\vorto{programo}\senco{1}{Skriburo od imprimuro qua anuncas la detali di ceremonio, di festo publika, di koncerto, di reprezento teatrala, e c.}\senco{2}{Anunco di la temi quin profesoro traktos en sua kurso}\senco{3}{Expozo di la koncepti politikala di partiso, di individuo}\lingui{DEFIRS}
\artiklo{progresar}%
\vorto{progresar}\fako{netrans., ad} Avancar a grado superiora\lingui{DEFIRS}
\artiklo{progresiono}%
\vorto{progresiono}\fako{matem.} Intersequo de termini, inter qui la raporto aritmetikala (difero) o geometriala (quociento) esas konstanta\lingui{DEFIRS}
\artiklo{projektar}%
\vorto{projektar}\fako{trans.}\senco{1}{\textit{(geom.)} Trasar perpendikli de omna punti di solido adsur plano, por figurar lu sur plano}\senco{2}{\textit{(optiko)} Reflektar sur *\textit{skrino}, per radii qui emanas de lumo-foko}\lingui{DEFIRS}
\artiklo{projektilo}%
\vorto{projektilo}\senco{1}{Irga korpo qua esas lansita per la impulso da forco irgequala, por atingar ulu od ulo}\senco{2}{Korpo qua esas lansita per lans-armo o per pafarmo}\lingui{DEFIS}
\artiklo{projektoro}%
\vorto{projektoro}\fako{tekn.} Aparato uzata por reflektar, forme di un o plura lumo-faski, di qui la intenseso lumala esas tre granda, la lum-ondi de foko\lingui{DEFIRS}
\artiklo{projetar}%
\vorto{projetar}\fako{trans.} Pozar avane ideo kom realigenda\lingui{DEFIRS}
\artiklo{proklamar}%
\vorto{proklamar}\fako{trans.} Anuncar publike e kelke-solene\lingui{DEFIRS}
\artiklo{proklitiko}%
\vorto{proklitiko}\fako{gram.} Vorto qua, pro indijar acento tonika, su apogas sur la vorto qua sequas\lingui{DEFIS}
\artiklo{prokonsulo}%
\vorto{prokonsulo}\fako{en Roma, en la epoki antiqua} Konsulo qua jus ekiris de sua ofico-periodo ed a qua onu konfidis la komando di armeo, di provinco\lingui{DEFIRS}
\artiklo{prokuraco}%
\vorto{prokuraco} Darfo grantita legale da ulu a persono quan lu komisas por agar ye la nomo di la grantinto\lingui{DEFIRS}
\artiklo{prokurar}%
\vorto{prokurar}\fako{trans., ad} Igar obtenar\lingui{EFIS}
\artiklo{prokuratoro}%
\vorto{prokuratoro}\fako{judiciado} Judiciisto komisita da la stato por akuzar koram tribunalo\lingui{DEFIRS}
\artiklo{prolegomeno}%
\vorto{prolegomeno} Longa prefaco qua preparas la objekto precipua di la verko\lingui{DEFIS}
\artiklo{prolepso}%
\vorto{prolepso}\fako{retor.} Figuro qua konsistas ye preirar, preventar, objeciono\lingui{DEFIS}
\artiklo{proletario}%
\vorto{proletario}\senco{1}{\textit{(en Roma, antique)} Civitano ek la klaso maxim inferiora, imuna de imposti}\senco{2}{La persono qua vivas en stando di indijo sen-cesa, pri lua vivonecesaji}\lingui{DEFIRS}
\artiklo{proliferar}%
\vorto{proliferar}\fako{netrans.) (bot. e fiziol.} Produktar enti simila a su
\artiklo{prolixa}%
\vorto{prolixa} Qua (lor skribo, lor parolo) dilutas to quon lu havas dicenda. – (dicesas anke pri la stilo di ulu)\lingui{DEFIS}
\artiklo{prologo}%
\vorto{prologo}\senco{1}{En teatrajo : sorto di prefaco qua anuncas e preparas la serio de eventontaji pleala}\senco{2}{Quaza prefaco}\lingui{DEFIRS}
\artiklo{prolongar}%
\vorto{prolongar}\fako{trans.} Plulongigar \textit{en spaco} (kayo, trotuaro, muro). Adjuntar, ad ula kozo, ulo suplementa qua konstitucas integrajo kun olu e qua havas kom nomo « prolonguro ». Onu \textit{plulongigas} : vivo, milito, voyajo, la valideso di pasporto\lingui{DEFIS}
\artiklo{promenar}%
\vorto{promenar}\fako{netrans.} Irar marche, dise, por respirar en aero ne-enklozita, por exercar sua muskuli\lingui{DEFR}
\artiklo{promisar}%
\vorto{promisar}\fako{trans., ad ulu} Dicar, garantiante lo per sua honoro, ke onu agos o facos ulo\lingui{DEFIS}
\artiklo{promocar}%
\vorto{promocar}\fako{trans.} Elevar (ulu) ad ofico, a grado, plu alta en la hierarkio\lingui{DEFIS}
\artiklo{promontorio}%
\vorto{promontorio}\fako{geogr.} Triangulo de tero tre salianta qua quaze avancas aden maro\lingui{DEFIS}
\artiklo{promulgar}%
\vorto{promulgar}\fako{trans.} Publikigar (lego) segun la formi postulata por ke lu divenez obliganta\lingui{DEFIS}
\artiklo{pronacar}%
\vorto{pronacar}\fako{trans.) (anat.} Igar la karpo en posturo tala ke la palmo di la manuo « regardas » vers sulo. \textit{(antonimo : supinacar)}
\artiklo{pronar}%
\vorto{pronar}\fako{netrans.) (religio, aludante la paroko, la vikario} Donar instrucioni, lor la celebro di la meso parokiala, en la sundii e la festo-dii, adjuntante a li la rekomendi quin lu havas (eventuale) agenda a la ekleziani, la publikigi di la mariaji, e c.\lingui{F}
\artiklo{pronomo}%
\vorto{pronomo}\fako{gram.} Vorto-speco qua remplasas e pleas la rolo di la nomo tacita\lingui{DEFIS}
\artiklo{pronta}%
\vorto{pronta}\senco{1}{Tote preparita}\senco{2}{Qua esas quik… -onta}\lingui{IS}
\artiklo{pronukleo}%
\vorto{pronukleo}\fako{biol.} Parto nuklea di la ovulo e di la spermatozoido qui kunfuzesas lor la fekundigo
\artiklo{pronuncar}%
\vorto{pronuncar}\fako{trans.}\senco{1}{Igar audebla publike (pronuncar diskurso, laudo funerala, pledo)}\senco{2}{Enuncar, artikulante la silabi, la vorti}\lingui{EFIS}
\artiklo{propagar}%
\vorto{propagar}\fako{trans.} Diskonocigar proximope. – (\textit{metaf}. : genitar por plumultigar (sua) decendanti.)\lingui{DEFIRS}
\artiklo{proporcionar}%
\vorto{proporcionar}\fako{trans.} Konformigar a la legi segun qui(lu), kom parto di ensemblo, esas tale o tale granda, konsidere di la toto\lingui{DEFIRS}
\artiklo{proporciono}%
\vorto{proporciono}\senco{1}{En linguo vulgara : (a) grandeso di parto kompare a la toto od a la altra parti, en ensemblo; (b) grandeso di kozo kompare a kozo analoga, egardata kom tipo; (c) grandeso di kozo kompare ad altra kozo, la unesma kreskante o deskreskante segun ke la duesma kreskas o deskreskas}\senco{2}{En la linguo ciencala : (a)}\lingui{DEFIRS}\subvorto{matem.}{proporciono geometriala}{Inter-egaleso di du raporti per quociento.}{}{}\subvorto{}{proporciono-regulo}{Operaco per qua, donite un ek la du raporti ed un termino di la duesma, onu determinas la altra termino.}{}{}\subvorto{}{proporciono aritmetikala}{Interegaleso di du raporti per difero.}{}{}\subvorto{kemio}{proporciono definita}{Quantesi fixa, nevariiva, segun qui interunionas la korpi por formacar kombinuri kemiala}{}{}
\artiklo{propozar}%
\vorto{propozar}\fako{trans.) (ad}\senco{1}{Pozar (ulo) avan ulu kom objekto, kom temo, di studio, di exploro}\senco{2}{Pozar (ulo) avan ulu kom regulo, kom modelo a qui konformigor sua agado}\senco{3}{Pozar (ulu) avan (ulu, korporaciono) kom kandidato}\lingui{DEFIRS}
\artiklo{propoziciono}%
\vorto{propoziciono}\senco{1}{\textit{(logiko)} Enunco di judiko}\senco{2}{\textit{(gram.)} Asembluro ek objekto, verbo ed atributo}\senco{3}{\textit{(geom.)} Enunco di teoremo demonstrenda}\senco{4}{\textit{(teol.)} Afirmo dogmatala}\lingui{DEFIRS}
\artiklo{propra}%
\vorto{propra} Qua proprietesas da persono, da kozo, exkluze di omna altra\lingui{EFIS}
\artiklo{propretoro}%
\vorto{propretoro}\fako{en Roma, antique} Lietnanto (remplasanto) di la pretoro\lingui{DEFIRS}
\artiklo{proprietar}%
\vorto{proprietar}\fako{trans.} Posedar (ulo) kom sua, kom propra, komplete e legitime\lingui{EFIS}
\artiklo{propulsar}%
\vorto{propulsar}\fako{trans., ad} \textit{(mekan.)} (aludante helico, paletroto, e c.) Pulsar adavan, avance\lingui{DEFIS}
\artiklo{prorogar}%
\vorto{prorogar}\senco{1}{\textit{(yuro-cienco)} Ajornar oficale a dato plu o min fora}\senco{2}{\textit{(aludante asemblitaro)} Rezolvar cesar quik la *seanco, la *sesiono, til epoko determinata}\lingui{DEFIS}
\artiklo{prosektoro}%
\vorto{prosektoro}\fako{anat.} Helpisto qua preparas, por la profesoro di medicino, la disseko-peci
\artiklo{prosenkimo}%
\vorto{prosenkimo}\fako{bot.} Tisuo ek fibri
\artiklo{prosilogismo}%
\vorto{prosilogismo}\fako{logiko} Katenatra dispozeso di la silogismo, en qua la konkluzo di la unesma divenas la premiso di la duesma – e tale pluse
\artiklo{proskriptar}%
\vorto{proskriptar}\fako{trans.) (en Roma, antique}\senco{1}{Imperar la ocido (di ulu) sen formi judiciala, per nur enskribo sur afisho-texto publikigita}\senco{2}{Kondamnar ad ocideso granda-quante}\lingui{DEFIRS}
\artiklo{prospekto}%
\vorto{prospekto} Anuncilo, generale imprimura, pri verko qua balde publikigesos, o pri institucuro por publiko, o pri varo, e c., quan onu disdonas por atraktar la kompronti, la klienti, eventuala\lingui{DEFIRS}
\artiklo{prosperar}%
\vorto{prosperar}\fako{netrans.} Favoresar da fortuno, sucesar, standar plena-vigoro\lingui{DEFIS}
\artiklo{prostatito}%
\vorto{prostatito}\fako{patol.} Inflamuro di la prostato\lingui{DEFIS}
\artiklo{prostato}%
\vorto{prostato}\fako{anat.} Glando situita (che la viro) ye la juntopunto di la veziko-kolo kun la uretro\lingui{DEFIS}
\artiklo{prosternar}%
\vorto{prosternar}\fako{netrans.} Kushar su, kun sua ventro vers sulo, kun sua vizajo kontre sulo (od inklinar su til kande sua fronto tushas sulo) koram ulu, kom signo di respekto\lingui{EFIS}
\artiklo{prostezo}%
\vorto{prostezo}\fako{gram.} Adjunto di un litero, di un silabo, komence di vorto. (kom ex. : \textit{ica} vice \textit{ca}, en L. \textit{iscala} vice \textit{scala}; en F. \textit{lierre} vice \textit{ierre})\lingui{DEFIS}
\artiklo{prostitucar}%
\vorto{prostitucar}\fako{trans.}\senco{1}{\textit{(aludante muliero)} Koitar kun irgu po pekunio}\lingui{DEFIRS}\subvorto{metaf.}{prostitucar la judicio}{Sakrifikar lu pro la expekto di pekunio, di favori, e c.}{2}{}\subvorto{}{prostitucar sua plumo}{Vendar sua skribo, sua redakto, po pekunio, ofico, e c.}{}{}\subvorto{}{prostitucar amikeso}{Igar irgu sua amiko, expekte di profito}{}{}
\artiklo{prostracar}%
\vorto{prostracar}\fako{netrans.} Esar abatita extreme, mentale, spiritale, o korpale\lingui{DEFIRS}
\artiklo{protagonisto}%
\vorto{protagonisto} La persono qua, en teatrajo, pleas la rolo precipua. – \textit{Metaf}. La persono qua *iniciativas ed agas quaze kom chefo en ta o ca afero publik-utila, homaro-bona\lingui{DEFS}
\artiklo{protalo}%
\vorto{protalo}\fako{bot.} Mikra lamo verda, granda admaxime quale la unglo di la orel-fingro, kreskanta surface di tero, qua rezultas de la jermifo di la spori di filiki (o di la cetera kriptogami vaskuloza)
\artiklo{protamino}%
\vorto{protamino}\fako{kemio organika} Proteino kontenata da la spermatozoidi di la fishi
\artiklo{protandra}%
\vorto{protandra}\fako{bot.} Dicesas pri la flori hermafrodita che qui la stamini divenas matura ante la pistilo ed ube, konseque, la fekundigo povas eventar nur inter flori distingata
\artiklo{protecionar}%
\vorto{protecionar}\fako{trans.) (ekonomiko} Charjar la produkturi de la landi stranjera, per dogano-imposti, por (aserte) protektar la produkturi di sua propra lando, detrimente di la kompronti\lingui{DEFIRS}
\artiklo{proteino}%
\vorto{proteino}\fako{biol.} Nomo rezervata, prefere, por la albuminoidi simpla, opoze a la albuminoidi kompozura \textit{(proteidi)}\lingui{DEFIS}
\artiklo{protektar}%
\vorto{protektar}\fako{trans.}\senco{1}{Defensar kontre to quo minacas}\senco{2}{Helpar la suceso di ulu per sua famo, sua egardeso, sua moyeni financala od altra}\lingui{DEFIRS}
\artiklo{protektorato}%
\vorto{protektorato}\fako{politiko} Dependo impozata a stato, a teritorio pozita sub la protekto di stato plu potenta, qua administras lu\lingui{DEFIRS}
\artiklo{proteo}%
\vorto{proteo} Homo qua pleas omna roli\lingui{DEFIRS}
\artiklo{proteolitika}%
\vorto{proteolitika}\fako{biol.} Qua koncernas la disloko di la proteini aden grupi plu simpla, influe da la fermenti digestala o da la efiki fiziko-kemiala
\artiklo{protestar}%
\vorto{protestar}\fako{netrans.}\senco{1}{Agar reklamaco formala kontre ago, kontre rezolvo quin onu deklaras nelegitima}\senco{2}{\textit{(komerco)} (aludante pago-promiso, trato) Igar konstatar oficale ke lu ne esas pagita ye la pago-dato, e ke la persono qua devis pagar responsas pri omna spensi acesora e detrimenti}\lingui{DEFIRS}\subvorto{religio kristana}{protestanto}{Persono qua apartenas a la religio reformita da Luther, Calvin, e c. \textit{(antonimo : katoliko)}}{}{DEFIRS}
\artiklo{protezo}%
\vorto{protezo}\fako{kirurg.} Remplasigo, per artificalajo, di organo quan onu deprenis per ablaciono\lingui{DEFI}
\artiklo{protisto}%
\vorto{protisto}\fako{biol.} Organismo maxim simpla, qua havas karakteri qui esas komuna a la du regni i animalaro e plantaro\lingui{DEFIS}
\artiklo{protogina}%
\vorto{protogina}\fako{biol.} Floro che qua la pistilo matureskas ante la stamini. \textit{(antonimo : protandra)}\lingui{DEFIS}
\artiklo{protokolo}%
\vorto{protokolo}\senco{1}{Raporto pri la rezolvi da kongresintaro, da konferintaro diplomacala}\senco{2}{Raporto per qua persono yurizita pri lo, o persono autoritatoza legale, konstatas to quon lu agis, facis, vidis od audis}\senco{3}{Akto da oficiero judiciala, qua konstatas delikto o kontravenco}\senco{4}{Raporto pri to quo eventis en *seanco o *sesiono di asemblitaro}\lingui{DEFIRS}
\artiklo{protoplasmo}%
\vorto{protoplasmo}\fako{biol.} Substanco komplexa, pasable fluida e diafana, ye qua konsistas la korpo di la celulo vivanta, e qua kontenas parto diferenciita \textit{(nukleo)} e kelka organeti : mikrosomi, leciti, vakuoli, e c.\lingui{DEFIRS}
\artiklo{prototipo}%
\vorto{prototipo}\senco{1}{Tipo prima, segun qua onu exekutis la exempleri cetera}\senco{2}{\textit{(metaf.)} Modelo supera ad la ceteri}\lingui{DEFIS}
\artiklo{protozoo}%
\vorto{protozoo}\fako{zool.} Animaleto quan lua organismo elementala lokizas ye la grado maxim infra en la hierarkio animalala\lingui{DEFIS}
\artiklo{protuberanco}%
\vorto{protuberanco} Saliajo qua aspektas quale gibo, ye la surfaco di korpo\lingui{DEFIRS}
\artiklo{proverbo}%
\vorto{proverbo} Kurta precepto di sajeso praktikala, uzata da la turbo\lingui{DEFIS}
\artiklo{providenco}%
\vorto{providenco}\senco{1}{\textit{(religio)} La sajo dea qua guvernas omna existajo}\senco{2}{\textit{(religio)} Deo egardata kom guvernanto di omna existajo}\senco{3}{\textit{(metaf.)} Helpo, susteno, konstanta}\lingui{EFIRS}
\artiklo{provinco}%
\vorto{provinco}\senco{1}{\textit{(en Roma, antique)} Teritorio qua konquestesis exter Italia, ed administrata da guverniestro}\senco{2}{Dividuro teritoriala di stato, administrata ye la nomo di la suvereno da guverniestro partikulara}\senco{3}{Kontraste kun la urbo metropola : la cetero di la lando}\lingui{DEFIRS}
\artiklo{provizar}%
\vorto{provizar}\fako{trans., ye ulo} Igar havar to quo esas necesa\lingui{DEFIRS}
\artiklo{provizora}%
\vorto{provizora}\senco{1}{Qua supleas dum kelka tempo por la satisfaco di bezono, varte di to quo esos definitiva}\senco{2}{Qua funcionas en ofico, dum kelka tempo, sen havar la titulo qua inheras la ofico}\lingui{DEFIRS}
\artiklo{provokar}%
\vorto{provokar}\fako{trans., ad} Ecitar a lukto, pro agar atakeme\lingui{DEFIRS}
\artiklo{proxeneto}%
\vorto{proxeneto} Mediacanto en la intrigi di aventuro amorala\lingui{FIS}
\artiklo{proxim}%
\vorto{proxim} Poke distanta (tempale, spacale, gradale) de…\lingui{EFIRS}
\artiklo{prozelito}%
\vorto{prozelito} Nova adepto di religio\lingui{DEFIRS}
\artiklo{prozo}%
\vorto{prozo} Formo di la pens-expresuro qua ne submisesas a la mezuro ed a la ritmo di la verso\lingui{DEFIRS}
\artiklo{prozodio}%
\vorto{prozodio}\senco{1}{\textit{(gram.)} Pronunco regulkonforma di vorto, konforme a la acentizo}\senco{2}{Reguli pri la quanto de silabi e la mezuro di la versi diversa}\lingui{DEFIRS}
\artiklo{prozopopeo}%
\vorto{prozopopeo}\fako{retor.} Figuro qua konsistas ye egardigar kom vivanta, kozin, personin mortinta od absenta\lingui{DEFIS}
\artiklo{pruda}%
\vorto{pruda} Qua afektacas vertuo rigoroza\lingui{DEF}
\artiklo{prudenta}%
\vorto{prudenta} Qua praktikas ta sajeso qua regulas la konduto ed igas posibla la evito di kulpi, di erori\lingui{EFIS}
\artiklo{pruinar}%
\vorto{pruinar}\fako{netrans.} Kovreskar per roso konjelita, o per nebulo konjelita, lor la nokti kolda di autuno o di printempo
\artiklo{prunelo}%
\vorto{prunelo}\fako{bot.} Pruno sovaja, nigra, di qua la saporo esas akra\latina{prunus spinosa}\lingui{FI}
\artiklo{pruno}%
\vorto{pruno}\fako{bot.} Frukto kernoza, kun shelo glata\latina{prunus}\lingui{EFI}
\artiklo{pruntar}%
\vorto{pruntar}\fako{trans., de}\senco{1}{Obtenar ke ulu prestas}\senco{2}{Tirar (ulo) de kunhomo}\senco{3}{Rekursar ad helpo de altru}\lingui{F}
\artiklo{pruo}%
\vorto{pruo} Avana parto di navo, di batelo\lingui{EFIS}
\artiklo{pruritar}%
\vorto{pruritar}\fako{netrans.} Tale priklar ke divenas necesega skrapar la pelo\lingui{EFIS}
\artiklo{pruvar}%
\vorto{pruvar}\fako{trans.} Vidigar kom exakta, per rezono o per fakti\lingui{DEFIS}
\artiklo{psalmo}%
\vorto{psalmo}\fako{religio krist.} Kantiko religiala\lingui{DEFIS}
\artiklo{psalmodiar}%
\vorto{psalmodiar}\fako{trans.}\senco{1}{\textit{(kirko)} Kantar psalmi}\senco{2}{\textit{(metaf.)} Recitar monotone}\lingui{DEFIS}
\artiklo{psalterio}%
\vorto{psalterio}\fako{epoki antiqua} Muzik-instrumento kordoza, quan onu pleas per plektro\lingui{DEFIRS}
\artiklo{pseudo-}%
\vorto{pseudo-} Prefixo qua havas kom senco « ne-exakta, ne-reala »
\artiklo{pseudonima}%
\vorto{pseudonima} Qua publikigas (o publikigesas) kun autor-nomo ne-exakta.\lingui{DEFIRS}\subvorto{}{pseudonimo}{Nomo ne-exakta}{}{}
\artiklo{pseudosfero}%
\vorto{pseudosfero}\fako{geom.} Surfaco di rotaco, produktita da la traktorio qua rotacas cirkum sua asimptoto
\artiklo{psikiatrio}%
\vorto{psikiatrio} Kuracado di la morbi mentala\lingui{DEFIR}
\artiklo{psiko}%
\vorto{psiko} La koncio fenomenala, egardata kom responsiva; la studiajo di psikologio. (Ta termino esas plu generala kam « anmo ».)\lingui{DEFIRS}
\artiklo{psikologio}%
\vorto{psikologio} Ta parto di filozofio qua traktas la fakultati mentala e la agado di psiko\lingui{DEFIRS}
\artiklo{psikologo}%
\vorto{psikologo}\senco{1}{La persono qua su okupas pri psikologio, o kompozas verki pri psikologio}\senco{2}{Autoro, persono qua diletas per analizar la standi psikala}\lingui{DEFIRS}
\artiklo{psikoso}%
\vorto{psikoso}\fako{patol.} Irga ek la morbi mentala
\artiklo{psikrometro}%
\vorto{psikrometro} Aparato per qua onu mezuras, onu determinas, la tenseso higrometrala di atmosfero\lingui{DEFIRS}
\artiklo{psit!}%
\vorto{psit!} Interjeciono uzata por atraktar la atenco da persono quan onu ne ja rajuntis, e qua avancas avan la advokanto
\artiklo{psoaso}%
\vorto{psoaso}\fako{anat.} Nomo di la muskuli qui esas sur la lateri di la vertebri lumbala\lingui{DEFIS}
\artiklo{pterigoida}%
\vorto{pterigoida} Alo-forma
\artiklo{pteropodo}%
\vorto{pteropodo}\fako{zool.} Klaso de moluski di qui la pedo havas du larja apendici ek muskuli, o flosi\lingui{EFIS}
\artiklo{ptomaino}%
\vorto{ptomaino}\fako{kemio biol.} Substanco alkaloida, qua naskas lor la putro da la proteini, e quan onu renkontras, do, en la kadavri\lingui{DEFIS}
\artiklo{pubera}%
\vorto{pubera} La persono che qua aparis la genito-fakultato\lingui{EFIS}
\artiklo{pubio}%
\vorto{pubio}\fako{anat.} Osto, e parto, avana e supra, di la pelvo, proxim la regiono hipogastrala, qua divenas kovrata ye pili lor la puberesko\lingui{EFIS}
\artiklo{publicano}%
\vorto{publicano}\fako{biblo} Recevisto di la imposti o di la censo-revenuo por la Romani \textit{(antique)} e, pro lo, odiata e desestimata lore da la Judi\lingui{EFIS}
\artiklo{publicisto}%
\vorto{publicisto} Persono qua kompozas libri od artikli – od esas erudita – pri la yuri qui koncernas la relati da la guvernistaro di stato kun la civitani di ta lando, e pri politiko\lingui{DEFIRS}
\artiklo{*publicitear}%
\vorto{*publicitear}\fako{trans.} \textit{(aludante konsuminto od uzinto)} laudar, a komprero o konsumero eventuala, varo o servo da qua onu esas tre satisfacita – laudo absolute gratuita e generale spontana. \textit{(Antonimo : reklamo)}
\artiklo{*publico}%
\vorto{*publico}\fako{publiko} Irgu, sen distingo\lingui{DEFIRS}
\artiklo{publika}%
\vorto{publika} Uzebla, acesebla, da irgu, koram irgu (voyo, gardeno, lecioni, *seanco, diskuto – publika)\lingui{DEFIRS}
\artiklo{pudelo}%
\vorto{pudelo}\fako{zool.} Varietato de hundo, di qua la pili esas longa e frizita\lingui{DER}
\artiklo{pudlar}%
\vorto{pudlar}\fako{trans.) (tekn.} Rafinar gisfero per deprenar lua troa karbono, transformar lu a fero o stalo, per varmigar lua surfaci oxidigiva\lingui{DEFR}
\artiklo{pudorar}%
\vorto{pudorar} \textit{(netrans.)} \textit{(precipue pri la homini)} Shamar pri to quo koncernas sexuo\lingui{EFIS}
\artiklo{pudro}%
\vorto{pudro} Amilo, reduktita a partikuli extreme tenua, parfumizita, uzata por la karnaciono, la hari\lingui{DEFR}
\artiklo{puero}%
\vorto{puero} Homo, kande lu evas inter 7 e 15 yari\lingui{EFL}
\artiklo{pufo}%
\vorto{pufo} Grosa tabureto cilindra, polsterizita, sen apogilo\lingui{DER}
\artiklo{pugno}%
\vorto{pugno} La manuo kande klozita, do kun la fingri interklemanta\lingui{EFIS}
\artiklo{pulchinelo}%
\vorto{pulchinelo}\senco{1}{Rolo komika, reprezentata kun du gibi (un dorse, un ventre), nazo hokatra e koloroza, kostumo tre precoza, e mimaskilo nigra}\senco{2}{\textit{(metaf.)} Persono groteska, poke valoroza}\lingui{DEFIRS}
\artiklo{pulco}%
\vorto{pulco}\fako{zool.} Mikra insekto, parazito di la homo e di ula animali, qua su nutras ek lia sango\latina{pulex irritans}\lingui{FIS}
\artiklo{pulikario}%
\vorto{pulikario}\fako{bot.} Genero de kompozaji, tribuo de la radicei\latina{pulicaria}\lingui{L}
\artiklo{pulio}%
\vorto{pulio}\fako{tekn.} Roto, ligna o metala, qua rotacas per axo, e provizita ye lua cirklokurvo, ye ranuro en qua pasas kordo di qua la extremaji recevas : ica, la tiro-forco, ita la rezisto-forco multa-kaze, sur transmiso-pulio, rimeno remplasas la kordo, e la ranuri divenas neutila\lingui{EFIS}
\artiklo{pulmonario}%
\vorto{pulmonario}\fako{bot.} Nomo di planti diversa qui aspektas, plu o min multe, quale pulmono\latina{pulmonaria}
\artiklo{pulmonito}%
\vorto{pulmonito}\fako{patol.} Inflamuro di la parenkimo di la pulmono\lingui{FIS}
\artiklo{pulmono}%
\vorto{pulmono}\fako{anat.} Singla ek la du viceri sponjatra, qui esas intersimetre en la kavajo torakala, per qui agesas e la respiro, e la produkto di voco\lingui{EFIS}
\artiklo{pulo}%
\vorto{pulo} \textit{(lor ludo per karti, biliardo, en la parii pri kavalkuri)} Pario en qua la riskaji (pekunio) di omna ludanti esos tote por la ganonto\lingui{EF}
\artiklo{pulpo}%
\vorto{pulpo}\senco{1}{\textit{(bot.)} Parto quaze-karna di ula frukti, legumi, e c.}\senco{2}{\textit{(farmacio)} Pulpo pastigita por ula preparaji farmaciala}\lingui{DEFIS}
\artiklo{pulsar}%
\vorto{pulsar}\fako{trans.}\senco{1}{Diplasar per preso, stroko, shoko}\senco{2}{\textit{(pri kordio)} Pulsar sango per stroki iterata, sen-cesa}\lingui{DEFIRS}
\artiklo{pultro}%
\vorto{pultro} Ensemblo ek la uceli quin onu nutras ed edukas en farmodomo, ruro-domo, korto por la animali domestika (hani, gansi, anadi, e c.)\lingui{E}
\artiklo{pulular}%
\vorto{pulular}\fako{netrans.} Kreskar (nombrale) abundege e rapide\lingui{EFIS}
\artiklo{pulvero}%
\vorto{pulvero}\senco{1}{Substanco solida, medikamenta, reduktita a partikuli tre tenua}\senco{2}{\textit{(milit-arto)} Mixuro de sulfo, salpetro e karbono, reduktita a partikuli tenua, qua, kande acendata, esas extreme exploziva – ed uzesas en la pafarmi por lansar projektili}\lingui{DEFIS}
\artiklo{pumao}%
\vorto{pumao}\fako{zool.} Mamifero karnavida, ek la familio « felidi », qua vivas sur la kontinento Amerika\latina{felix concolor}\lingui{DEFIRS}
\artiklo{pumico}%
\vorto{pumico}\fako{geol.} Petro volkanala, poroza, lejera, uzata por frotar, polisar, e c.\lingui{EIS}
\artiklo{pumpar}%
\vorto{pumpar}\fako{trans.} Aspirar o pulsar (liquido, gaso) per aparato qua, ordinare, konsistas ye cilindro en qua su movas, rektalinee ed alterne, pistono\lingui{DEFIS}
\artiklo{puncar}%
\vorto{puncar}\fako{trans.}\senco{1}{Surpozar (quaze perforo-probe) marko konkava produktita per stalo-bloko kun, ye un ek lua extremaji pinta, graburo reliefa, e sur la altra extremajo di qua onu frapas per martelo}\senco{2}{Perforar (bilieto kartona o papera) per pinco por igar lu nevalida}\lingui{DEFIS}
\artiklo{puncho}%
\vorto{puncho} Drinkajo, mixuro ek aquo, od infuzuro de teo, kun brandio o rumo, citrono e sukro\lingui{DEFIRS}
\artiklo{puncionar}%
\vorto{puncionar}\fako{trans.) (kirurg.} Pikar kirurgiale por igar eskapar liquido qua amaseskis aden kavajo di la korpo (kavajo naturala od acidentala)\lingui{DEFIS}
\artiklo{punisar}%
\vorto{punisar}\fako{trans.) (pro, per} Aplikar (ad ulu) sufrigilo pro lua kulpo\lingui{EFIS}
\artiklo{punsar}%
\vorto{punsar}\fako{trans.} Kopiar desegnuro per pulvero quan onu pasigas tra la trui di papero perforita segun la linei di la desegnuro\lingui{EF}
\artiklo{punto}%
\vorto{punto}\senco{1}{\textit{(gram.)} Signo ortografiala super le « i »}\senco{2}{*Puntuo-signo : .;?! …}\senco{3}{\textit{(muziko)} Signo qua, apud noto, augmentas la duro-valoro di olca; qua indikas haltinstanto en la mezuro}\senco{4}{\textit{(geom.)} Elemento qua genitas la lineo, konsiderata per abstrakto, kom ulo qua esas nek longa, nek larja, nek profunda}\senco{5}{La subdividuro minima di la parti di la temo quan onu traktas}\senco{6}{\textit{(ludi : biliardo, karti, loto, e c.)} Marko quan onu facas sur paperfolio od ardezo-plako, pos singla gano-stroko}\senco{7}{\textit{(tipogr.)} Mezuro qua determinas la grandeso di la tipi}\lingui{DEFIRS}
\artiklo{*puntuar}%
\vorto{*puntuar}\fako{trans.} Lor-skribe, insertar, en la frazi, la signi ,; . :?!
\artiklo{pupeo}%
\vorto{pupeo} Mikra figuro homala, ek kartono, vaxo, ligno, uzata da la infantini kom ludilo\lingui{DEF}
\artiklo{pupilo}%
\vorto{pupilo}\fako{anat.} Aperturo, meze di iriso, tra qua la lumo-radii eniras la okulo\lingui{DEFIRS}
\artiklo{pupitro}%
\vorto{pupitro} Moblo en qua existas surfaco plana, inklinita, sur qua onu pozas libri, muziko-kayeri, papero-folii, por lektar, kantar, plear, skribar komfortoze\lingui{FIS}
\artiklo{pupo}%
\vorto{pupo}\fako{navig.} Dopa parto di navo, di batelo\lingui{EFIS}
\artiklo{pura}%
\vorto{pura}\senco{1}{Kun la naturo di qua ne esas mixita elemento stranjera.}\subvorto{}{la matematiki pura, la mekaniko pura}{Qui havas kom studiajo nur la teorio, sen apliki praktikala}{}{}\senco{2}{Di qua la naturo ne koruptesas da elemento stranjera}\lingui{DEFIS}
\artiklo{pureo}%
\vorto{pureo} Disho ek legumi reduktita a paplo\lingui{DEFIS}
\artiklo{purgar}%
\vorto{purgar}\fako{trans.}\senco{1}{Kuracar per eskapigar la feko quan kontenas la intestini, per absorbigar lino-semini, agar-agaro, ricinoleo, aquo minerala sulfoza, e c.}\senco{2}{\textit{(ulo)} Liberigar de to quo sordidigas, desparigas, koruptas, domajas (lu)}\lingui{DEFIS}
\artiklo{purgatorio}%
\vorto{purgatorio}\fako{religio katol.} Loko ube la homi qui mortas ye stando di graco, expiacas la peki pri qui li ne penitencis sat multe en ica mondo\lingui{DEFIS}
\artiklo{puritano}%
\vorto{puritano}\fako{kristanismo} Membro di sekto protestanta tre rigoroza, qua aspiras ad igar kristanismo ri-esar pura quale lor lua origino\lingui{DEFIRS}
\artiklo{purpuro}%
\vorto{purpuro} Kolorizivo bele-reda quan antique onu obtenis de L. \textit{murex trunculus} o murex brancaris, molusko gasteropoda, tipo di la tribuo « muricinei » qua vivas en la mari di qui la aquo esas kolda o tepida, o fosila de la epoko « kretacea »\latina{murex trunculus; murex brancaris}\lingui{DEFIRS}
\artiklo{puso}%
\vorto{puso}\fako{patol.} Humoro morbala, densa, quan produktas la korupteso di la tisui, lor lia inflameso\lingui{EFIS}
\artiklo{pustulo}%
\vorto{pustulo}\fako{patol.} Mikra tumoro qua pusifas\lingui{DEFIRS}
\artiklo{putativa}%
\vorto{putativa}\fako{yuro-cienco} Pri qua onu kredas ke lu esas valida legale\lingui{EFIS}
\artiklo{puteo}%
\vorto{puteo} Exkavajo facita en sulo por tirar la aquo quan donas, ye ta o ca profundeso-grado, fonti od infiltruri.\lingui{FIS}\subvorto{}{puteo arteza}{Boro-kanaleto, facita en sulo per sondilo, qua atingas (multa-kaze) profundeso-grado tre granda, til aquo-strato}{}{}
\artiklo{putoro}%
\vorto{putoro}\fako{zool.} Mamifero karnivora, ek la familio « martri », qua odoras fetide\latina{mustela putorius}\lingui{F}
\artiklo{putrar}%
\vorto{putrar}\fako{netrans.) (pri, precipue, animalo o planto mortinta} Deskompozar su\lingui{EFIS}
\end{multicols}
\sekciono{Q}
\begin{multicols}{2}
\artiklo{qua}%
\vorto{qua}\fako{gram.}\senco{1}{\textit{(pronomo relativa)} Ne havas ipse senco; lokizita kom vorto rang-unesma di propoziciono qua pluprecizigas antecedento tre nepreciza, e determinas lu : (a) se « qua » ne preiresas nemediate da komo, la antecedento esas la \textit{lasta} substantivo di la propoziciono preiranta; (b) se « qua » preiresas da komo, la antecedento esas la \textit{pre-lasta} substantivo di la propoziciono preiranta. (Ta regulo esas aplikita rigoroze en ica Radikaro)}\senco{2}{\textit{(pronomo questionala)} « Qua persono…? »}\senco{3}{\textit{(adjektivo questionala)} « Qua… ek la plura qui existas…? »}
\artiklo{quadrangula}%
\vorto{quadrangula}\fako{geom.} Qua havas quar anguli e quar lateri.\subvorto{}{prismo, piramido quadriangula}{Di qua quadrilatero formacas la bazo}{}{}
\artiklo{quadranto}%
\vorto{quadranto}\fako{geom.} Quarimo di la cirklo-kurvo, t.e. 90\lingui{DEFIRS}
\artiklo{quadratika}%
\vorto{quadratika}\fako{matem.}\subvorto{}{equaciono quadratika}{Equaciono algebrala di qua la grado egalesas 2-potenco e di qua la solvo esas reduktebla ad olta di nombro finita de equacioni di la grado duesma}{}{}
\artiklo{quadrato}%
\vorto{quadrato}\fako{geom.} Quadriangulo equilatera ed ortangula\lingui{DEFIRS}
\artiklo{quadraturo}%
\vorto{quadraturo}\fako{astron.} Aspekto di du astri qui interdistas per un quarimo di cirklokurvo\lingui{DEFIRS}
\artiklo{quadri-}%
\vorto{quadri-}\fako{prefixo ciencala} « Qua havas quar… -i »
\artiklo{quadrigo}%
\vorto{quadrigo}\fako{en la epoki antiqua} Charo adqua onu jungis quar kavali sama-fronte\lingui{DEFIS}
\artiklo{quadriko}%
\vorto{quadriko}\fako{geom.} Areo di la grado duesma, t.e. reprezentata da equaciono di la grado duesma\lingui{DEFIS}
\artiklo{quadriliono}%
\vorto{quadriliono}\fako{aritm.} Mil bilioni. (1000.)
\artiklo{quadrilo}%
\vorto{quadrilo}\senco{1}{(a) Trupo de kavalieri, dividita (ordinare) aden quar grupi, qua figurantesas en karuselo; (b) singla de ta grupi, distingita per kolori, per kostumi interdiferanta}\senco{2}{(a) Nombro para de dui de danseruli e de danserini qui figurantesas, ici regardante iti facie, en kontradanso; (b) la danso ipsa}\senco{3}{\textit{(teknol.)} To quo reprezentas quadrato, rombo, ajuroza o sen-ajura, en la desegnuri di stofo, gipuro, pasmento, e c.}\lingui{DEFIRS}
\artiklo{quadriremo}%
\vorto{quadriremo}\fako{navig.} Navo kun quar rangi de remeri, o kun quar remeri ye un remilo\lingui{DEFIRS}
\artiklo{quadro}%
\vorto{quadro} Petro-bloko qua esas taliita (o qua esos taliita) por esor uzebla kom parto di domo, di edifico\lingui{DI}
\artiklo{quago}%
\vorto{quago}\fako{zool.} Speco di zebro, nun desaparinta de la terglobo, qua vivis en la sud-regioni di Afrika : nur la duima parto (avana) di lua korpo esis blanke-strioza\latina{hippotigris quagga}
\artiklo{quak!}%
\vorto{quak!} Son-uno ne-akordanta, quan lasas eskapar acidente persono qua kantas o qua pleas langet-instrumento muzikala\lingui{EF}
\artiklo{quakero}%
\vorto{quakero}\fako{religio kristan.} Membro di sekto de teisti filantropa qui agnoskas revelo interna da Deo, e predikas la interfrateso di omna homi, e la paco universala\lingui{DEFIRS}
\artiklo{quala}%
\vorto{quala} Vorto qua indikas la naturo (bona o mala) di ulu, di ulo. (Korelativo di « tala » : tala… quala…)\lingui{DEFIS}
\artiklo{qualifikar}%
\vorto{qualifikar}\fako{trans.} Karakterizar per la atribuo di qualeso\lingui{DEFIS}
\artiklo{qualio}%
\vorto{qualio}\fako{zool.} Ucelo migrera, ek la genero « perdriko », ma plu mikra kam perdriko\latina{Tetrao coturnix}
\artiklo{quankam}%
\vorto{quankam} Konjunciono qua venigas la enunco di la obstaklo malgre qua ulo eventas\lingui{L}
\artiklo{quanta}%
\vorto{quanta} Vorto qua indikas la grado di multeso. (Korelativo di « tanta » : tanta… quanta…)\lingui{DEFIS}
\artiklo{quantika}%
\vorto{quantika}\lingui{DEFIRS}\subvorto{}{nombri quantika}{Ensemblo ek la quar nombri qui definas la karakteri di singla ek la elektroni planeta di atomo}{}{}
\artiklo{quar}%
\vorto{quar} La cifro « 4 »\lingui{EFIS}
\artiklo{quaranteno}%
\vorto{quaranteno} Sejorno quaradek-dia qua impozesas a la pasajinti, la vesti, la bagaji, la vari – sejorno sur la navo od en lazareto – qui arivas de lando ube prosperas morbo kontagiala\lingui{DEFIRS}
\artiklo{quarco}%
\vorto{quarco}\fako{mineral.} Silico pura\simbolo{SiO}\lingui{DEFIRS}
\artiklo{quartana}%
\vorto{quartana}\fako{patol.} Dicesas pri febro intermitanta qua riaparas singla-quardie\lingui{DEFIS}
\artiklo{quartermastro}%
\vorto{quartermastro}\fako{navig.} Grado unesma super olta di navano, che la navistari di la militistaro marala. (Lu korespondas a la grado « kaporalo » di la militistaro terala)\lingui{DEFIRS}
\artiklo{quartero}%
\vorto{quartero} Segmento administrala di urbo\lingui{DEFIRS}
\artiklo{quarterono}%
\vorto{quarterono} Homo qua devenis de la kopulaco da pelblanko kun mestico\lingui{DEFIRS}
\artiklo{quarteto}%
\vorto{quarteto}\fako{muziko}\senco{1}{Muzikajo kompozita por quar voci o por quar instrumenti (du violini, alto e violoncelo)}\senco{2}{\textit{(en orkestro)} Ensemblo ek la violini unesma e duesma, alti, violonceli (o kontrabasi) qui formacas quar parti}\lingui{DEFIRS}
\artiklo{quartiko}%
\vorto{quartiko}\fako{geom.} Kurvo di la klaso quaresma\lingui{EFIS}
\artiklo{quarto}%
\vorto{quarto}\senco{1}{\textit{(muziko)} Grado quaresma di la skalo diatonika}\senco{2}{\textit{(skerm.)} La sorto quaresma di la parei : pareo per situar la karpo vers-extere, kun la manuo-palmo vers cielo e la spado-punto alte}\senco{3}{\textit{(karto-ludo)} Eso, en manuo, di quar karti qui intersequas sen intermanko. (kom ex. : aso, rejulo, damo, pajo formacas quarto majora.)}\senco{4}{\textit{(navig.)} (sur navo) Singla de la periodi quar-hora dum qui agesas la observo surveya}\lingui{DEFIRS}
\artiklo{quasio}%
\vorto{quasio}\fako{bot.} Kortico de la arbusto L. \textit{quassia amara}, de-tranchita aden mikra spani, qua saporas tre bitre, ed uzata en medicino kom tonigivo pos macereso (5 grami singla-litre), sive en tintivo alkoholoza\latina{quassia amara}\lingui{DEFIS}
\artiklo{quaternara}%
\vorto{quaternara}\senco{1}{\textit{(matem.)} En qua quar unaji simpla formacas unajo ek la klaso quaresma, e tale pluse}\lingui{DEFIS}\subvorto{geol.}{strato quaternara}{La maxim recenta ek la strati di formaceso sedimentala}{2}{}
\artiklo{quaterniono}%
\vorto{quaterniono}\fako{matem.} Nomo donita da Hamilton ad ula expresuri imaginara por qui lu kreis kalkulo-speco partikulara qua posibligas la solvo facila di la problemi qui koncernas la geometrio tri-dimensiona\lingui{DEFIS}
\artiklo{quatreno}%
\vorto{quatreno} Stanco, parto di soneto, o mikra poeziajo quar-versa\lingui{DEFIS}
\artiklo{quaza}%
\vorto{quaza} Simila til ula grado, ma ne tote tala. (Indikas nur semblo o simileso, e ke la du objekti interdiferas per nur poko)\lingui{DEFIS}
\artiklo{quecho}%
\vorto{quecho}\fako{bot.} Sorto di grosa pruno oblonga, qua su *reproduktas per semuro\latina{prunus domestica}\lingui{F}
\artiklo{querar}%
\vorto{querar}\fako{trans.} Serchar, komisite od intencante adportor od venigor\lingui{FI}
\artiklo{quercitrono}%
\vorto{quercitrono}\fako{bot.} Granda querko sempre-verda de la nordo-regioni di Amerika, di qua la kortico donas materio tintala flava\latina{quercus tinctoria}\lingui{DEFI}
\artiklo{querko}%
\vorto{querko}\fako{bot.} Granda arboro ek la klaso « amentacei », di qua la ligno, tre harda, esas un ek le maxim uzata por menuzo e karpento, e di qua la kortico, pulverigite, donas tano\latina{quercus}\lingui{L}
\artiklo{questar}%
\vorto{questar}\fako{trans.} Demandar, rekoltar, (pekunio) por verki karitatala o piesala\lingui{FI}
\artiklo{questionar}%
\vorto{questionar}\fako{trans., pri} Interpelar per demando di informivi por igar su savoza pri ulo\lingui{EFIS}
\artiklo{questoro}%
\vorto{questoro}\senco{1}{\textit{(en la epoki antiqua, en Roma)} Oficiero qua administris financi publika}\senco{2}{\textit{(nun)} La persono qua, en ula korporacioni, esas komisita por surveyar la uzo di la pekunio korporacionala}\lingui{DEFIRS}
\artiklo{qui}%
\vorto{qui}\fako{gram.} Pluralo di « qua »
\artiklo{quieta}%
\vorto{quieta} Di qua la psiko ne agitesas da angori\lingui{EFIS}
\artiklo{quik}%
\vorto{quik} Sen ajorno, tardeso (*retardeso) o fristo, mem minima\lingui{E}
\artiklo{quilayo}%
\vorto{quilayo}\fako{bot.} Genero de rozacei, ek tri o quar speci de arbori, kun folii alternanta, flori dioika, qui kreskas en la sudo-regioni di Amerika, e di qui la pulvero, obtenata de lia kortico, uzesas por netigar la linji sordida de makuli grasa\latina{quillaja saponaria}
\artiklo{quin}%
\vorto{quin} Pluralo di « quan »
\artiklo{quinara}%
\vorto{quinara}\fako{matem.} En qua kin unaji simpla formacas unajo ek la klaso kinesma\lingui{DEFIS}
\artiklo{quingo}%
\vorto{quingo}\fako{bot.} Frukto flava, di qua la pulpo kontenas quaza stoneti, acerba ed astriktiva, – donata da arboro ek la familio « rozacei », kun folii di qui la suba facio aspektas kotonoza\latina{cydonia}
\artiklo{quinino}%
\vorto{quinino}\fako{kemio} Alkaloido de vejetanto, extraktita de quino\simbolo{C₂₀H₂₄O₂N₂}\lingui{DEFIRS}
\artiklo{quinkunco}%
\vorto{quinkunco} Asemblitaro de objekti (ordinare : arbori) dispozita kinope : quar qui formacas la anguli di quadrato, ed un en la centro\lingui{DEFI}
\artiklo{quino}%
\vorto{quino}\fako{bot.) (farmacio} Kortico kontre-febra, donata da la arbori exotika ek la genero « cinchona »\lingui{DEFIRS}
\artiklo{quinqua-}%
\vorto{quinqua-} Prefixo ciencala : qua havas kin…
\artiklo{quinquagezimo}%
\vorto{quinquagezimo}\fako{liturgio katol.} La kinadekesma dio ante la pasko-festo\lingui{DEFIS}
\artiklo{quintalo}%
\vorto{quintalo}\fako{mezur-sistemo metrala} Maso 100-kilograma\lingui{EFIS}
\artiklo{quintana}%
\vorto{quintana}\fako{patol.} Febro qua riaparas singla-kindie
\artiklo{quintesenco}%
\vorto{quintesenco}\senco{1}{(a) \textit{(olim)} Substanco eterala, egardata kom la elemento kinesma e maxim subtila. (b). \textit{(nun)}. La parto maxim subtila di substanco}\senco{2}{\textit{(metaf.)} Lo maxim bona, maxim precoza, maxim perfekta}\lingui{DEFIRS}
\artiklo{quinteto}%
\vorto{quinteto} Muzikajo de kin parti koncertala\lingui{DEFIS}
\artiklo{quintiko}%
\vorto{quintiko}\fako{matem.} Kurvo di la klaso kinesma\lingui{EFIS}
\artiklo{quintiliono}%
\vorto{quintiliono}\fako{matem.} Mil quadrilioni. (10³⁰)
\artiklo{quinto}%
\vorto{quinto}\senco{1}{\textit{(muziko)} Intervalo de kin noti intersequanta, la du extrema inkluzite. – Grado kinesma di la skalo diatonika}\senco{2}{\textit{(skerm.)} Pareo simpla, qua korespondas a la engajo (adinterne) di la lineo alta}\lingui{DEFIRS}
\artiklo{quiproquo}%
\vorto{quiproquo} Eroro konseque de qua onu egardas kozo o persono kom altra\lingui{DEFIS}
\artiklo{quirito}%
\vorto{quirito}\fako{che la Romani, antique} Civitano Romana qua rezidis en Roma, kontraste kun ti qui esis servanta en la armei\lingui{DEFIRS}
\artiklo{quirlar}%
\vorto{quirlar}\fako{trans.} Batar (liquido, ovi, kremo) per vergo por mixar, igar augmentar volumine, spumifar, e c.\lingui{DE}
\artiklo{quita}%
\vorto{quita} Tote liberigata de obligeso pekuniala od etikala\lingui{DEFS}
\artiklo{quo}%
\vorto{quo} Neutra formo di « qua »
\artiklo{quociento}%
\vorto{quociento}\fako{aritm.} Rezultajo de la divido da nombro per altra\lingui{DEFIS}
\artiklo{quodlibeto}%
\vorto{quodlibeto} Dico triviala, mokachanta, joko mala-mora, mala vorto-ludo\lingui{F}
\artiklo{*quoniam}%
\vorto{*quoniam} Ta vorto indikas la motivo, la kauzo, la justifikivo; de la instanto kande onu agnoskas kom valida ke…\lingui{L}
\artiklo{quoto}%
\vorto{quoto} La parto di singlu en la repartiso di pekunio-quanto pagenda o recevenda da ensemblo de personi\lingui{DEFIS}
\end{multicols}
\sekciono{R}
\begin{multicols}{2}
\artiklo{rabatar}%
\vorto{rabatar}\fako{trans., ek, de} \textit{(komerco)} \textit{(aludante furnisisto)} Grantar preco-diminuto a sua kompranto. Grantar konvencione (a kompranto) diminuto di la pekunio-quanto debata, kondicione ke ica pagesos ante la dato normala\lingui{DEFIRS}
\artiklo{rabino}%
\vorto{rabino}\senco{1}{Doktoro pri la lego religiala di la Judi}\senco{2}{Doktoro pri la kulto Judala, chefo di socio de Judi}\lingui{DEFIRS}
\artiklo{rabio}%
\vorto{rabio}\fako{patol.} Morbo virulenta qua aparas spontane che la hundo (plu rare che la kato, la volfo, e c.), komunikebla per mordo da ta animali, e finas ye acesi hororinda di konvulsi, e ye morto. (Morbo vinkita da Louis Pasteur)\lingui{EFIS}
\artiklo{raboto}%
\vorto{raboto}\fako{tekn.} Utensilo di menuzisto, cizelo ek stalo, muntita oblique interne di ligno-peco rektangula kavigita, qua lasas pasar la parto tranchiva; uzata por planigar, diminutar la areo di ligno-peco\lingui{F}
\artiklo{racemo}%
\vorto{racemo}\senco{1}{\textit{(biol.)} Nomo ciencala di la grapo}\senco{2}{\textit{(kemio)} Substanco qua ne efikas a la lumo polarigita, konsistanta ye kombinuro (kun molekuli inter-egala) di du inversaji optikala}\lingui{DEFIS}
\artiklo{raciono}%
\vorto{raciono} La povo dicernar la exaktaji : (a) la fakultato deskovrar lo exakta; (b)\textit{(filoz.)} la fakultato konceptar la exaktaji absoluta\lingui{DEFIRS}
\artiklo{radiar}%
\vorto{radiar}\fako{netrans.}\senco{1}{\textit{(aludante lumo-sprici rekta)} Lansesar quale dardo, de centro lumoza, diverge, omna-sinse}\senco{2}{\textit{(fiziko)} Lansar kaloro-sprici}\senco{3}{\textit{(geom.)} \textit{(pri rekti)} Irar de la centro di cirklo a la cirklo-kurvo, a la surfaco di la sfero}\lingui{DEFIRS}
\artiklo{radiatoro}%
\vorto{radiatoro}\fako{teknol.}\senco{1}{Organo di ensemblo mekanikala, destinita ad eventigar la koldesko di liquido o di gaso uzata, en la ensemblo, kom vehilo di kaloro}\senco{2}{Hem-aparato uzata por varmigar la atmosfero di la chambri per fluido (aero, aquo, vaporo) qua uzesas kom vehilo di kaloro por transportar olca kelka-diste}\lingui{F}
\artiklo{radiko}%
\vorto{radiko}\senco{1}{\textit{(bot.)} La parto infra di vejetanto, ordinare sinkanta aden tero, ube lu su developas inversa-sinse di la stipo, e di qua la rolo esas fixigar la planto en sulo e pumpar la elementi de qui nutresas la planto. – Parto di organo (dento haro, pilo, unglo) per qua lu su enplantacas en altra parto}\senco{2}{\textit{(metaf.)} (a) \textit{(gram.)} Vorto primitiva, de qua derivesas altra vorti. (b)}\lingui{EFIS}\subvorto{matem.}{radiki di equaciono}{Quanti quin onu povas substitucar a la nekonocaji, kom satisfacanta egale la donaji di la equaciono.}{}{}\subvorto{}{extraktar radiko, quadrata, kubala, di nombro}{Trovar la nombro qua, multiplikite unople o duople ye su ipsa, donis ta nombro}{}{}
\artiklo{radiografar}%
\vorto{radiografar}\fako{trans.} Fotografar la internajo di la homo-korpo per la X-radii (Roentgen-radii) – por skopi medicinala, kirurgiala, e c.\lingui{EFIRS}
\artiklo{radiometro}%
\vorto{radiometro} Aparato uzata por mezurar la intenseso-grado di la lumo-radii\lingui{EFIS}
\artiklo{radioskopar}%
\vorto{radioskopar}\fako{trans.) (medic.} Explorar la ombri quin produktas la X-radii (Roentgen-radii) sur *skrino fluorecanta ek platino-cianuro di bario o di tungsteno\lingui{EFIS}
\artiklo{radirar}%
\vorto{radirar}\fako{trans.} Grabar, facante la desegnuro per stal-instrumento stipeto-forma, kun punto tre akuta por povar trasar sur la varniso o la vaxo-strato linei tre dina e delikata\lingui{D}
\artiklo{radiumo}%
\vorto{radiumo}\fako{kemio} Elemento metala, tre skarsa, deskovrita en 1898 da Pierre e Marie Curie kun J. Bémont, en la rezidui barioza obtenita lor la trakto di pekblendo (oxido naturala di uranio)\simbolo{Ra}\lingui{DEFIRS}
\artiklo{radiuso}%
\vorto{radiuso}\fako{anat.} Ta, ek la du osti di la avanbrakio, qua esas sur la latero vers-extere\lingui{EFS}
\artiklo{rado}%
\vorto{rado} Parto de maro, qua penetras aden la tero ferma, e formacas quaza doko naturala ube la navi povas shirmar su\lingui{DFIRS}
\artiklo{radotar}%
\vorto{radotar}\fako{netrans., pri} Parolar nekonsequante, nekoherante, konseque de feblesko di la spirito\lingui{EF}
\artiklo{rafano}%
\vorto{rafano}\fako{bot.} Speco di kruciferi, di la genero « cochlearia » : planto perena, kun grosa radiko pulpoza, blanka, qua (freshe raspita) povas remplasar mustardo\latina{raphanus sativus}\lingui{I}
\artiklo{rafido}%
\vorto{rafido}\fako{biol.} Kristali aguloforma quin onu observas en ula celuli di la animali o di la vejetanti
\artiklo{rafinar}%
\vorto{rafinar}\fako{trans.}\senco{1}{Igar plu pura per deprenar la substanci stranjera quin (lu) kontenas}\senco{2}{\textit{(metaf.)} Igar plu delikata, plu subtila}\lingui{DEFIRS}
\artiklo{raflo}%
\vorto{raflo}\fako{bot.} Vito-grapo de qua onu deprenis lua beri\lingui{F}
\artiklo{rafto}%
\vorto{rafto} Asemblitaro de trabi (o de arbori abatita) interligita tale ke li formacas quaza planko-sulo flotaciva\lingui{E}
\artiklo{rago}%
\vorto{rago} Fragmento de stofo, telo, lano, e c., qua pendas, mi-lacerite, de olda vesto par-konsumita\lingui{E}
\artiklo{raguto}%
\vorto{raguto}\fako{koquarto} Karno-peco preparita per sauco\lingui{DEFI}
\artiklo{raitro}%
\vorto{raitro}\fako{olim} Kavalo-soldato de Germania, qua militoservis en Francia\lingui{DFIRS}
\artiklo{rajuntar}%
\vorto{rajuntar}\fako{trans.} Irar por juntar su itere ad (ulu) de qua onu esas separita\lingui{FIS}
\artiklo{rakahuto}%
\vorto{rakahuto} Mixuro de fekulo, de glani dolca, de rizo-farino, de tuberkuli de ter-mandeli, sukrizita ed aromatizita\lingui{DEFIRS}
\artiklo{rakar}%
\vorto{rakar}\fako{trans., de} Pasigar (liquido) de un vazo aden altra, cherpante sube (per robineto o sifono)\lingui{E}
\artiklo{raketo}%
\vorto{raketo} Planketo manuo-tenebla, kun reto de tripo-kordi, uzata por sendar la baloneto, en la baloneto-ludo – la flugbalono, en la ludo di ta nomo\lingui{DEFIRS}
\artiklo{rakito}%
\vorto{rakito}\fako{patol.} Halto di la kresko, manifestata (multa-kaze) da deformeso di la spino\lingui{DEFIRS}
\artiklo{rakontar}%
\vorto{rakontar}\fako{trans., ad} Dicar (ad ulu) narace (ulo fingita) por amuzar\lingui{EFI}
\artiklo{raliar}%
\vorto{raliar}\fako{trans.) (milit.} Asemblar sur batal-agro, la kombatanti qui esis dispersita\lingui{DEF}
\artiklo{ralo}%
\vorto{ralo}\fako{zool.} Genero de stiltieri, kun grandeso mezvalora, gambi pasable kurta ma tre-longa-fingra, tre apta pri kuro – vivanta en la plana landi ed en la marshi; fluganta tre poke; cirkulanta nur dum nokto; extreme karnavida\latina{rallus}\lingui{DEFI}
\artiklo{ramio}%
\vorto{ramio}\fako{bot.} Planto qua donas materio texebla, qua devenis de Oriento\latina{boehmeria}\lingui{DEFS}
\artiklo{ramno}%
\vorto{ramno}\fako{bot.} Arboreto di qua la beri (preske omna-kaze : nigra) esas purgiva, e di qua la kortico donas kolorizivo flava\latina{rhamnus}\lingui{DEFS}
\artiklo{ramo}%
\vorto{ramo}\fako{bot.}\senco{1}{Mikra brancho qua naskis de brancho precipua}\senco{2}{Subdividuro. \textit{(anat.)} Subdividuro di la vaskuli, di la nervi. – \textit{(geol.)} Masivo qua divergas de montaro, aden altra direciono. – \textit{(geneal.)} Subdividuro di brancho di familio, di arboro genealogiala}\lingui{DEFIS}
\artiklo{rampo}%
\vorto{rampo} Sur voyo, plano qua formacas angulo obliqua kun la horizontalo, e quan onu esas acensonta\lingui{DEF}
\artiklo{ranca}%
\vorto{ranca} Dicesas pri korpo grasa (lardo, oleo, e c.) qua aquiris odoro kelke mala, e saporo akra\lingui{DEFIS}
\artiklo{rango}%
\vorto{rango}\senco{1}{Singla de la linei sur qui kozi, o personi, qui intersucedas esas situita}\senco{2}{(a) En serio de personi, de kozi, la plaso qua apartenas a singla de li, avan o dop la ceteri. -b La plaso quan ulu okupas en la estimo da la cetera homi}\lingui{DEFR}
\artiklo{rankorar}%
\vorto{rankorar}\fako{netrans., ad, pro, pri} Memorar tenace la mala agi quin onu subisis de ulu\lingui{EFIS}
\artiklo{rano}%
\vorto{rano}\fako{zool.} Batrako sen-kauda qua vivas, prefere, apud aquo stagnanta, en la loki humida\latina{rana}\lingui{IS}
\artiklo{ransono}%
\vorto{ransono} Preco postulata po la livro di milito-kaptito, di karcerano\lingui{DEF}
\artiklo{ranunkulo}%
\vorto{ranunkulo}\fako{bot.} Planto herbatra, di qua un speco, kun flori dazlante flava, esas vulgara en la prati\latina{ranunculus}\lingui{DEFIRS}
\artiklo{ranuro}%
\vorto{ranuro}\fako{tekn.} Kanelo facita en ligno-bloko por uzesar kom kuliso o por recevar asemblo-peco\lingui{FS}
\artiklo{rapecar}%
\vorto{rapecar}\fako{trans.} Igar bona-standa itere (stofo) per pasigar interkruce fili en o sur la loko domajita (per truo, per lacereso partala)\lingui{FI}
\artiklo{rapida}%
\vorto{rapida}\senco{1}{Qua parkuras multa spaco en poka tempo}\senco{2}{Qua agas o facas en poka tempo (to pri quo onu esas parolanta)}\lingui{DEFIS}
\artiklo{rapiero}%
\vorto{rapiero}\fako{skerm-arto} Olim : Espado longa, akutigita, di qua la konko esis truoza (trui aden qui povis penetrar la espado di la adverso)\lingui{DEFR}
\artiklo{rapo}%
\vorto{rapo}\fako{bot.} Varietato di kaulo, kun radiko tre pulpoza e manjebla\latina{rapa}\lingui{FIR}
\artiklo{raportar}%
\vorto{raportar}\fako{netrans., pri}\senco{1}{Naracar ad ulu to quon onu vidis, audis, e c.}\lingui{DEFIRS}\subvorto{metaf.}{raporto}{Rezultajo de la interkomparo di du quanti}{2}{}
\artiklo{rapsodo}%
\vorto{rapsodo}\fako{en Grekia, antique} La recitisto di epikaji\lingui{DEFIRS}
\artiklo{raptar}%
\vorto{raptar}\fako{trans., de} Deprenar per violento (ulu, o to quo apartenas a kunhomo)\lingui{FIS}
\artiklo{rara}%
\vorto{rara} Tre ne-ofta\lingui{DEFIS}
\artiklo{raskalo}%
\vorto{raskalo} Persono qua nule skrupulas pri agar ne-honeste\lingui{E}
\artiklo{raso}%
\vorto{raso}\senco{1}{Ensemblo ek la ancestri ed ek la decendinti di un familio, di un populo}\senco{2}{Grupo di speco de animali o de planti di qua la karakterizivo esas konstanta e transmisesas per genitado}\lingui{DEFIRS}
\artiklo{raspar}%
\vorto{raspar}\fako{trans.} Transformar aden pulpo o pulvero, per ta implemento koqueyala qua konsistas ye plako de lado herisita ye asperaji sur qui onu frotas la substanco tale transformenda\lingui{DEFIRS}
\artiklo{rastaquero}%
\vorto{rastaquero} Exotiko qua ostentas per luxo o tituli suspektinda
\artiklo{rastar}%
\vorto{rastar}\fako{trans.} Amasigar, en gardeno, la folii defalinta en la alei, od, en prato, la feno tranchita – per implemento T-forma, di qua la stango transversa esas provizita ye denti ligna o fera\lingui{FIS}
\artiklo{rasteliero}%
\vorto{rasteliero}\senco{1}{Stablo-moblo fixigita super la manjo-trogo en qua la feno destinita a la bruti retenesas da stangi ligna, muntita quale la denti di rastilo}\senco{2}{(a) Sorto di etajero adsur qua la soldati lokizas sua fusili, en kazerno. (b) Mikra etajero destinita recevor pipi}\lingui{FIS}
\artiklo{ratafio}%
\vorto{ratafio} Liquoro qua konsistas ye brandio, sukro ed aromati o frukto-sukro\lingui{DEFIS}
\artiklo{rateno}%
\vorto{rateno}\fako{tekn.} Stofo ek lanofili kruce-dispozita, kun la pili tirita adextere e loklo-forme\lingui{DEFIS}
\artiklo{ratifikar}%
\vorto{ratifikar}\fako{trans.) (yuro-cienco} Konfirmar per akto specala (akti emanita de privati o de autoritato ek judicio-rango inferiora : od akti signatita plenipotenciere, ye la nomo di lando, ma qui, por divenar valida por exekuteso, esas aprobenda da parlamento)\lingui{DEFIRS}
\artiklo{rato}%
\vorto{rato}\fako{zool.} Mikra quadripedo, ek la familio « roderi », kun muzelo pintoforma, gambi kurta, kaudo longa, qua rodas la cerealo-grani, palio, e c.\latina{mus decumanus, mus rattus}\lingui{DEFIS}
\artiklo{ratono}%
\vorto{ratono}\fako{zool.} Genero de mamiferi karnavida, tipo di la familio « procionidei », vicina ye la urso, qua kontenas kin speci propra ye Amerika\latina{procyon lotor}\lingui{FS}
\artiklo{rauka}%
\vorto{rauka} Di qua la sono esas aspera, ruda, quale la voco di persono di qua la laringo esas inflamata\lingui{EFIS}
\artiklo{raupo}%
\vorto{raupo}\fako{zool.} Larvo di papilioni, kun korpo longa, qua konsistas ye serio de ringi, e normale viloza\lingui{D}
\artiklo{rauto}%
\vorto{rauto} Festo en salono, adube onu invitas mondumani\lingui{DEFR}
\artiklo{ravelino}%
\vorto{ravelino}\fako{milit-arto) (en la yarcenti 15esma e 16-esma} Fortifikajo pasable analoga a la nuna redani e mi-luni\lingui{DEFIRS}
\artiklo{ravino}%
\vorto{ravino} Voyo naturala, profunde exkavita da aquo qua fluis olim en ta loko\lingui{EF}
\artiklo{ravisar}%
\vorto{ravisar}\fako{trans., per} Exaltar, igar ne plus posedar su, per la impulso de entuziasmo\lingui{EFI}
\artiklo{rayo}%
\vorto{rayo}\fako{zool.} Mar-fisho, plata e kartilagoza\latina{raya}\lingui{EFIS}
\artiklo{razar}%
\vorto{razar}\fako{trans.} Reze-tondar la barbo-pili\lingui{DEFIS}
\artiklo{re}%
\vorto{re}\fako{muziko} Noto duesma di la gamo naturala di « do »
\artiklo{reaktar}%
\vorto{reaktar}\fako{netrans., ad, kontre}\senco{1}{Facar sur korpo efiko genitata da la efiko quan lu recevas de ica}\senco{2}{\textit{(kemio)} \textit{(aludante korpo)} Manifestar sua karakteri distinganta per la efekto di la efiko quan altra korpo facis a lu}\senco{3}{Esforcar por rezistar}\lingui{DEFIRS}
\artiklo{reala}%
\vorto{reala}\senco{1}{Di qua la existo esas fakta, e ne nur semblanta o posibla}\senco{2}{\textit{(metaf.)} Dicesas pri *grandoro, antonime ye \textit{imaginara}}\lingui{DEFIS}\subvorto{optiko}{foko reala}{Ube la lumo-radii reflektita interrenkontras \textit{(antonime ye virtuala : ube li interrenkontrus se li esus plulongigita)}}{3}{}
\artiklo{realgaro}%
\vorto{realgaro}\fako{kemio} Reda arseno-sulfajo\simbolo{AsS}
\artiklo{rebela}%
\vorto{rebela} Qua refuzas obediar la autoritati legitima\lingui{DEFIS}
\artiklo{rebordo}%
\vorto{rebordo}\fako{tekn.} Bordo igita plu alta kam la centro-parto\lingui{FIRS}
\artiklo{rebuso}%
\vorto{rebuso} Esprito-ludo qua konsistas ye reprezentar vorto, frazo per figuri di objekti di qui la nomi memorigas, per homonimeso, ta vorto, ta frazo\lingui{DEFIR}
\artiklo{recensar}%
\vorto{recensar}\fako{trans.) (pri libri, artikli di revuo, e c} Analizar \textit{(por *publico)} olia kontenajo, opinionar pri lia valoro spiritala, e c.\lingui{DER}
\artiklo{recenta}%
\vorto{recenta} Qua eventis, existeskis, nur poka tempo ante la instanto konsiderata\lingui{EFIS}
\artiklo{recepcionar}%
\vorto{recepcionar}\fako{trans.} Verifikar (pri qualeso e quanto) livrajo por saveskar ka lu esas absolute konforma a la klauzi di la kompro-kontrato\lingui{F}
\artiklo{receptaklo}%
\vorto{receptaklo}\fako{bot.} Parto di la floro ube su insertas la kalico, la korolo, la stamini e, normale, la karpelo (o karpeli) di la ovario\lingui{EFS}
\artiklo{recepto}%
\vorto{recepto}\senco{1}{Indiko quan onu recevas pri la preparo di ula medikamenti, di ula dishi}\senco{2}{\textit{(medic.)} Preskripto da mediko per qua lu donas la formulo di medikamento}\lingui{DEFIRS}
\artiklo{recesiva}%
\vorto{recesiva}\fako{biol.} Dicesas, pri la Mendel-lego \textit{(*leyo)}, por qualifikar la karakteri qui dominacesas – t.e. qui transmisesas, a la hibrido od a la mestico – da un ek la patri, ma sen aparar en la genituro\lingui{F}
\artiklo{recevar}%
\vorto{recevar}\fako{trans., de} Esar la sendario di sendajo, la donario di donajo, la donacario di donaco, e c.- esar la objekto di ago\lingui{EFIS}
\artiklo{recidivar}%
\vorto{recidivar}\fako{netrans., pri} Iterar krimino, delikto, kulpo pri qui lu ja kondamnesis e punisesis\lingui{DEFI}
\artiklo{recipiento}%
\vorto{recipiento}\fako{tekn.} Vazo adaptita ad ula aparati por recevar gasi, liquida, e c.\lingui{DEFIS}
\artiklo{reciproka}%
\vorto{reciproka}\senco{1}{Dicesas pri du termini, singla de qui agas a la altra ye maniero qua equivalas olta quan lu recevas}\senco{2}{\textit{(gram.)} Propozicioni reciproka : Di qui singla esas nur la altra renversita, singla di qui havas kom subjekto la atributo, e kom atributo la subjekto di la altra.}\subvorto{}{verbo reciproka}{Qua expresas la ago da du subjekti, ica ad ita, ed ita ad ica}{}{}\senco{3}{Fondita sur interkambio di agi, di sentimenti qui inter-korespondas}\lingui{DEFIRS}
\artiklo{recitar}%
\vorto{recitar}\fako{trans., ad, koram} Dicar voce, parole, memore, quale se onu lektus\lingui{DEFIS}
\artiklo{recitativo}%
\vorto{recitativo} \textit{(muziko)} En ensemblo, solo kantala od instrumentala. Kanto deklamata sen simetrio di mezuro e di ritmo, e kadencata segun la dividuro di la versi e di la flexi di la voco parolanta\lingui{DEFIRS}
\artiklo{reda}%
\vorto{reda}\senco{1}{Di qua la koloro esas olta di sango, di fairo, e c.}\lingui{E}\subvorto{metaf.}{la redi}{La revolucionisti qui havas kom emblemo standardo reda}{2}{}
\artiklo{redaktar}%
\vorto{redaktar}\fako{trans.} Kompozar segun la formo qua konvenas (libro, artiklo di jurnalo o di revuo, letro, kontrato, e c.)\lingui{DEFIRS}
\artiklo{redano}%
\vorto{redano}\fako{fortifik.} Fortifiko-masivo dento-forma, ek du facii qui intersekas, tale formacante angulo konvexa\lingui{DEF}
\artiklo{redemtar}%
\vorto{redemtar}\fako{trans.}\senco{1}{Obtenar la liberigo po ransono}\senco{2}{Liberigar de obligeso per pagar pekunio-quanto}\senco{3}{\textit{(metaf.)} \textit{(kristanismo)} Obtenar la salveso eterna po sua morto. \textit{(aludante Iesu)}}\lingui{EFIS}
\artiklo{redingoto}%
\vorto{redingoto} Vesto di qua la paneli envelopas la korpo\lingui{DF}
\artiklo{reduito}%
\vorto{reduito} Fortifikajo konstruktita dop altra por prolongar la defenso-fortifikaji se la unesma esus kaptita da la enemiko\lingui{DEFIRS}
\artiklo{reduktar}%
\vorto{reduktar}\fako{trans.}\senco{1}{Retroduktar a lua situeso naturala, normala (osto luxacita, ruptita)}\senco{2}{Retroduktar a formo plu elementa (kom ex. : reduktar grani a farino, ligno a cindri, substanco a pulvero, a pudro, e c.)}\lingui{DEFIRS}
\artiklo{reduto}%
\vorto{reduto}\fako{milit-arto} Fortifikajo aparta, tote klozita (*kluza), sen anguli konkava. – Anke : Fortifikajo flotacanta, garnisita ye bastioni, uzata por transportar trupi de soldati trans fluvio o rivero\lingui{DEFS}
\artiklo{referar}%
\vorto{referar}\fako{trans., ad}\senco{1}{Relatigar (ulo) a to quo explikas lu, konfirmas lu}\lingui{DEFIRS}\subvorto{}{referar su ad}{Rekursar a la autoritato di ulu}{2}{}
\artiklo{referendario}%
\vorto{referendario} Oficiero di kancelereyo qua havas kom taski gardar la rejo-siglilo e raportar pri la supliko-letri\lingui{DEF}
\artiklo{reflektar}%
\vorto{reflektar}\senco{1}{\textit{(trans.)} \textit{(aludante surfaco)} Retrosendar (en ula kazi : plu feble) to quo frapas lu}\senco{2}{\textit{(netrans., pri)} Retrovenar a sua pensajo por studiar lu plu profunde}\lingui{DEFIRS}
\artiklo{reflektoro}%
\vorto{reflektoro}\fako{tekn., optiko} Aparato qua konsistas ye surfaco de metalo polisita od ye spegulo, destinita reflektor la radii konsiderata segun la direciono selektita; li uzesas precipue por lumo, ma anke por soni, kaloro\lingui{DEFS}
\artiklo{reflexa}%
\vorto{reflexa}\fako{fiziol.} Produktita sen-vole da reakto organismala\lingui{DEFIRS}
\artiklo{refluxar}%
\vorto{refluxar}\fako{netrans.) (pri mareo} Decensar, pos acensir per fluxo\lingui{EFIS}
\artiklo{refo}%
\vorto{refo}\fako{navig.} Parto di seglo quan onu rifaldis por diminutar la surfaco opozata a vento
\artiklo{reformar}%
\vorto{reformar}\fako{trans.}\senco{1}{\textit{(eklezio kristan.)} Retroduktar a la formo originala}\senco{2}{Donar formo plu bona, qua igas plu apta por la skopo}\lingui{DEFIRS}
\artiklo{refraktar}%
\vorto{refraktar}\fako{trans.) (fiziko} Igar (lumo-radio) deviacar pro (per) lua paso (pasigo) en medio di qua la denseso-gradi rispektiva diferas\lingui{DEFIRS}
\artiklo{refraktara}%
\vorto{refraktara}\senco{1}{\textit{(tekn.)} Qua rezistas ula influi}\senco{2}{\textit{(milit.)} Qua eskapas la lego pri armeanigado}\lingui{DEFIRS}
\artiklo{refraktoro}%
\vorto{refraktoro}\fako{tekn.} Instrumento qua donas la indici di refrakto\lingui{DEFIRS}
\artiklo{refreno}%
\vorto{refreno}\senco{1}{Ri-aparo di verso en omna kupleto di kansono, di balado, di rondelo}\senco{2}{To quon onu repetas sencese}
\artiklo{refujar}%
\vorto{refujar}\fako{netrans., en, che} Irar rapide a loko ube onu povas vivar sekure\lingui{EFIS}
\artiklo{refutar}%
\vorto{refutar}\fako{trans.} Destruktar la valideso (di opiniono) per demonstrar ke lu esas ne-exakta\lingui{DEFIS}
\artiklo{refuzar}%
\vorto{refuzar}\fako{trans.}\senco{1}{Ne-grantar}\senco{2}{Ne-aceptar}\senco{3}{Ne-admisar}\lingui{EFIS}
\artiklo{regalar}%
\vorto{regalar}\fako{trans., per}\senco{1}{Restorar (ulu) per ulo bona drinkale o manjale}\senco{2}{Ofrar (a gasti) partio di plezuro}\lingui{DEFS}
\artiklo{regalio}%
\vorto{regalio}\fako{eklezio katol., en Francia} Yuro quan havis la rejo pri inkasar la revenui di la episkopeyi, di la abadeyi, vakanta, e nominar pri la benefici qui dependis de li\lingui{DEFIRS}
\artiklo{regardar}%
\vorto{regardar}\fako{trans.} Direktar sua okuli apertita (ad ulo od ulu)\lingui{EFI}
\artiklo{regato}%
\vorto{regato} Konkurso de bateli, qui avancas o per seglo, o per la agado da remeri\lingui{DEFIRS}
\artiklo{regenerar}%
\vorto{regenerar}\fako{trans.}\senco{1}{Reproduktar (parto destruktita)}\senco{2}{Ri-novigar \textit{(*ri-nuvigar)} etikale, morale}\lingui{DEFIRS}
\artiklo{regento}%
\vorto{regento} La persono qua guvernas dum la minoreso o dum la absenteso di la suvereno\lingui{DEFIRS}
\artiklo{regimento}%
\vorto{regimento}\fako{armeo} Trupo de plura batalioni od eskadroni, direkte da kolonelo\lingui{DEFIRS}
\artiklo{regio}%
\vorto{regio}\senco{1}{Exploto, administro, di domeno ye la nomo di altru e ye la profito da ica}\senco{2}{Imposto-inkaso agata nemediate da la agenti di la stato, profitote da ica (kontraste kun la sistemo qua konsistas ye luar la yuro inkasor la imposto-revenui, la lokacinto pagante a la stato pekunio-quanto konvencionita, e retenante por su la ecesajo)}\senco{3}{La employataro, la kontori, qui reletas ta percepto nemediata}\lingui{DEFIS}
\artiklo{regiono}%
\vorto{regiono}\senco{1}{Parto plu o min granda di lando qua submisesas a klimato komuna}\senco{2}{Parto determinita di korpo. \textit{(anke metaf.)}}\lingui{DEFIRS}
\artiklo{registro}%
\vorto{registro} Libro, kayero, sur la skriburo-virga pagini di qua onu notas senfalie la fakti pri qui onu volas memorar\lingui{DEFIRS}
\artiklo{regnar}%
\vorto{regnar}\fako{trans.} Aplikar sua povo rejala, sua autoritato suverenala\lingui{DEFIS}
\artiklo{regoliso}%
\vorto{regoliso}\fako{bot.} Planto di qua la radiko uzesas kom dolcigivo pri rumo, bronkito, e c.\latina{glycyr-rhiza}\lingui{IS}
\artiklo{regolo}%
\vorto{regolo}\fako{zool.} Mikra pasero kantera, ek la familio « silvidei »\latina{troglodytes parvulus}
\artiklo{regresar}%
\vorto{regresar}\fako{netrans.} Irar kontre progresado\lingui{EFIRS}
\artiklo{regretar}%
\vorto{regretar}\fako{trans.}\senco{1}{Chagrenar pro ne plus havar (ulu, ulo)}\senco{2}{Chagrenar pro agir (o pro ne agir) (ulo)}\lingui{EF}
\artiklo{regular}%
\vorto{regular}\fako{trans.) (tekn.} Submisar ad ordino rigoroza (la funciono di mekanismo, di mashino, e c.)\lingui{DEFIRS}
\artiklo{regulatoro}%
\vorto{regulatoro}\fako{tekn.} Organo qua regulas la funciono di mekanismo, di mashino\lingui{DEFIS}
\artiklo{rehabilitar}%
\vorto{rehabilitar}\fako{trans.}\senco{1}{Ri-establisar en lua stando originala, en lua yuri originala (persono quan onu deklarabis dekadinta de li)}\senco{2}{\textit{(metaf.)} Ri-estimigar da lua kun-homi}\lingui{DEFIRS}
\artiklo{reino}%
\vorto{reino} Singla ek la rimeni di la brido (o di la brideto) quin onu uzas por direktar kavalo\lingui{EFIS}
\artiklo{rejimo}%
\vorto{rejimo}\senco{1}{\textit{(politiko)} Formo di guverno}\senco{2}{Maniero administrar institucuro}\senco{3}{\textit{(medic.)} Maniero administrar sua higieno individuala, sua bona-stando}\senco{4}{\textit{(medic.)} Maniero vivar en qua onu surveyas su, precipue pri to quon onu manjas o drinkas}\senco{5}{\textit{(tekn.)} Maniero ye qua eventas la fluo di aquo nestagnanta}\lingui{DEFIRS}
\artiklo{rejo}%
\vorto{rejo}\senco{1}{Chefo suverena di ula stati}\senco{2}{\textit{(karto-ludo)} Karto qua reprezentas rejo en singla de la du kolori}\senco{3}{\textit{(shak-ludo)} Piono precipua di la ludo, quan onu ne povas prenar; la partio esas perdita kande la rejo esas prenebla}\lingui{EFIS}
\artiklo{rekalkar}%
\vorto{rekalkar}\fako{trans.) (solduro, e c} Kunpresar, per martelofrapi, la plombo, la stano, qua salias, transpasas\lingui{IS}
\artiklo{rekamar}%
\vorto{rekamar}\fako{trans.} Lor texo, pasigar, sur la fondo uniforma di stofo, fili qui formacas desegnuro\lingui{IS}
\artiklo{rekapitular}%
\vorto{rekapitular}\fako{trans.} Rezumar, kurte repetar, o memorigar, per enumero preciza e kompleta di omna punti, parti o chapitri\lingui{DEFIS}
\artiklo{reklamacar}%
\vorto{reklamacar}\senco{1}{\textit{(netrans.)} \textit{(ad, pri, por)} Protestar kontre to quon onu judikas kom neyusta, por obtenar la reparo di la domajo o lezo; desagnoskar to quon onu judikas kom neyusta, ne-equitatoza}\senco{2}{\textit{(trans.)} Postular (to quo debesas)}\senco{3}{Demandar (ulo) tre insistoze}\lingui{DEF}
\artiklo{reklamar}%
\vorto{reklamar}\fako{netrans., pri} Laudar (mem exajeroze) publike la komercajo, la varo, la objekto quin onu probas vendar\lingui{DFRS}
\artiklo{rekolecar}%
\vorto{rekolecar}\fako{netrans.}\senco{1}{Koncentrar su aden su, per izolar su de la kozi extera}\senco{2}{Koncentrar su ad pensi di pieso}\lingui{FIS}
\artiklo{rekoliar}%
\vorto{rekoliar}\fako{trans.} Deprenar, hike ed ibe, to quo falis, dispersite, sur sulo e retro-pozar lu a lua plaso\lingui{F}
\artiklo{rekoltar}%
\vorto{rekoltar}\fako{trans.} Rekoliar la produkturi da la tero kultivita. – \textit{(anke metaf.)}\lingui{FI}
\artiklo{rekomendar}%
\vorto{rekomendar}\fako{trans.}\senco{1}{\textit{(pri persono)} Indikar (lu) partikulare kom esonta la objekto di la bonvolo, di la protekto, da ulu}\senco{2}{\textit{(pri letro, pako, e c.) (en posto-kontoro)} , demandar ke la letro-livristo donez lu a la destinario ipsa, nemediate – po pago di afranko-suplemento}\lingui{DEFIRS}
\artiklo{rekompensar}%
\vorto{rekompensar}\fako{trans., pro, per} Gratifikar pro servo agita. quankam ne-debata\lingui{EFIRS}
\artiklo{*rekonocar}%
\vorto{*rekonocar}\fako{trans.} Renkontrar inter sua memoraji, kom ja konocita (persono, kozo)
\artiklo{rekordo}%
\vorto{rekordo} Prodajo sportala (od analoga), konstatita oficale, e qua ecelas omna agi antea di la sama genero\lingui{DEFR}
\artiklo{rekortar}%
\vorto{rekortar}\fako{trans.} Plukurtigar per tra-tranchar la extremajo (hari, kordo, e c.), per detranchar la bordo
\artiklo{rekrutar}%
\vorto{rekrutar}\fako{trans.) (ulu, ad}\senco{1}{Obtenar la su-engajo libera (di ulu) kom soldato}\senco{2}{Obtenar la su-engajo libera (di ulu) a trupo, a partiso, a *muvmento}\lingui{DEFIRS}
\artiklo{rekta}%
\vorto{rekta}\fako{geom.} Qua, en spaco, de sua komenco til sua fino, iras sen deviaco, t.e. segun la lineo maxim kurta qua esas trasebla de un punto ad altra punto\lingui{DEFIS}
\artiklo{rektangulo}%
\vorto{rektangulo}\fako{geom.} Qua havas un orto\lingui{EFIS}
\artiklo{rektifikar}%
\vorto{rektifikar}\fako{trans.}\senco{1}{\textit{(kemio)} Igar (liquido) pura per un distilo plusa}\senco{2}{\textit{(metaf.)} Igar exakta per emendo, korektigo}\lingui{DEFIRS}
\artiklo{rektolinea}%
\vorto{rektolinea}\fako{matem.} Qua konsistas ye linei rekta
\artiklo{rektoro}%
\vorto{rektoro} La persono qua direktas universitato, distrikto universitatala, o Jezuito-kolegio\lingui{DEFIRS}
\artiklo{rektumoptoso}%
\vorto{rektumoptoso}\fako{patol.} Laxesko di la muskuli qui subtenas la parieto di la rektumo
\artiklo{rekuperar}%
\vorto{rekuperar}\fako{trans., de} Riposedeskar (per agir por lo)\lingui{DEFIS}
\artiklo{rekurenta}%
\vorto{rekurenta}\senco{1}{\textit{(anat.)} \textit{(pri nervo)} Qua kurveskas, e retrovenas preske a sua departo-punto}\subvorto{algebro}{serio rekurenta}{Di qua singla ek la termini esas la sumo di ula quanto de termini qui preiras lu, e qui multiplikesas rispektive per expresuri nevarianta}{2}{}
\artiklo{rekursar}%
\vorto{rekursar}\fako{netrans., ad}\senco{1}{Irar demandar (ad ulu) helpeso quan onu ne recevis de altru}\senco{2}{\textit{(yuro-cienco)} Apelar ad (ulu, ula judicio-korto) kontre verdikto judiciala}\lingui{DEFIS}
\artiklo{rekuzar}%
\vorto{rekuzar}\fako{trans., ulu, su) (yuro-cienco}\senco{1}{Refuzar aceptar kom judicionto, kom juriano, kom testo, kom arbitronto}\senco{2}{Refuzar agnoskar la autoritato di ulu}\lingui{FIS}\subvorto{}{rekuzar su}{Refuzar savigar sua opiniono}{3}{}
\artiklo{relatar}%
\vorto{relatar}\fako{netrans. ad} \textit{(pri termini, muziko-soni, afecioni, interesti)} Inter-korespondar per adminime kelka traiti komuna\lingui{DEFIRS}\subvorto{}{relate}{Antonimo di « \textit{absolute} »}{}{}
\artiklo{relayo}%
\vorto{relayo}\senco{1}{(a) Loko ube hundi qui vartas, en la parkureyo di chasantaro, remplasos la hundi qui esos fatigita; (b) Loko ube kavali quin onu vartigas, ye intervali, en posto-voyo, remplasos la kavali qui esos fatigita}\senco{2}{\textit{(telegrafo)} Aparato di qua la rolo esas lasar la elektro-fluo pasar, od interruptar ica, kande ula cirkonstanci esos realigita en altra cirkuito}\lingui{DEFIRS}
\artiklo{relegar}%
\vorto{relegar}\fako{trans.}\senco{1}{\textit{(yuro-cienco, che la Romani antiqua)} Exilar aden loko determinita, sen privaco di la yuri civila e politikala}\senco{2}{\textit{(nun, en Francia)} Kondamnar ad enterneso en kolonio}\lingui{DEFIS}
\artiklo{*relektar}%
\vorto{*relektar}\fako{trans.} Rilektar (nur per regardo, se por su; laute, se por la jus-diktinto) to quon onu skribis, por korektigar olu
\artiklo{relevar}%
\vorto{relevar}\fako{trans.} Determinar la situeso topografiala di (ulo) ed indikar lu sur mapo\lingui{FIS}
\artiklo{reliefo}%
\vorto{reliefo} La parto qua salias ek la surfaco cirkondanta\lingui{DEFIRS}
\artiklo{religio}%
\vorto{religio} Ensemblo ek la kredi e praktiki e ritui, qui havas kom skopo homajar deo o dei\lingui{DEFIRS}
\artiklo{reliquio}%
\vorto{reliquio}\fako{religio kristan.} To quo restas de santo pos lua morto, ek lua korpo, ek to quon lu esis, ek la instrumenti di lua tormenteso\lingui{DEFIS}
\artiklo{relo}%
\vorto{relo}\fako{tekn.} Singla de la du bendi ek fero o stalo, polisita, sur qui la roti di treno-veturi rulas facile e mantenesas en la direciono volata; konvexan esas la parto qua kontaktas la roti. (Opoze a « tramo », ye qua ta parto esas konkava)\lingui{EFRS}
\artiklo{remar}%
\vorto{remar}\fako{netrans.} Uzar, por igar batelo (ed, olim, navo) avancar, longa ligno-stango, kun planketo ye la extremajo per qua onu quaze plaudas o batas la aquo e, tale, obtenas la pulso-forco\lingui{FIS}
\artiklo{remarkar}%
\vorto{remarkar}\fako{trans.} Atencar ne-vole (ulu, ulo) ne-ordinara, ne-kustumala, ne-expektita\lingui{EF}
\artiklo{remediar}%
\vorto{remediar}\fako{trans., per} Aplikar la kozo qua esas necesa por risanigar (fiziologiale o psikale). – \textit{(anke metaf.)}\lingui{EFIS}
\artiklo{reminicar}%
\vorto{reminicar}\fako{trans.} \textit{(lor kompozo literaturala, muzikala, artala)} Memorar fragmenti di ta o ca literaturajo, muzikajo, pikturo, statuo, quan onu imitas ne-vole, nekoncie\lingui{F}
\artiklo{remisar}%
\vorto{remisar}\senco{1}{\textit{(trans.)} Dispensar (ulu) de parto di lua debo, di lua puniseso-grado, di peko}\senco{2}{\textit{(netrans.)} \textit{(patol.)} (aludante la simptomi di morbo) Cesar, o diminutar, dum kelka tempo}\lingui{EFIS}
\artiklo{remonstrar}%
\vorto{remonstrar}\fako{netrans., pri}\senco{1}{Parolar, diskursar, por demonstrar (ad ulu) lua nejusteso, la detrimento quan lu agis}\senco{2}{\textit{(olim, en Francia)} Protestar (kom parlamento) a la rejo kontre to quon ta parlamento judikis kom misuzo o tro-uzo di lua autoritato rejala}\lingui{DEF}
\artiklo{remonto}%
\vorto{remonto}\fako{armeo} Segmento di la armei di qua la rolo esas furnisar a la militisto-korpi ed a la institucuri militala la kavali quin li bezonas\lingui{DEFIRS}
\artiklo{remorkar}%
\vorto{remorkar}\fako{trans.}\senco{1}{\textit{(aludante navo)} Tirar dop su altra navo per kablo, per kateno}\senco{2}{\textit{(aludante lokomotivo)} Tranar vagoni}\lingui{FIS}
\artiklo{remoro}%
\vorto{remoro}\fako{zool.} Ekeneido, fisho a qua la Romani e Greki antiqua atribuis la fakultato impedar absolute la avanco di la navi\latina{echeneis remora}
\artiklo{remorsar}%
\vorto{remorsar}\fako{netrans., pri, pro} Esar la objekto di reprocho de sua koncienco\lingui{EFIS}
\artiklo{remparo}%
\vorto{remparo}\fako{milit-arto} Cirkuito-murego, reliefajo de tero, e c., destinita protektor urbo, kastelo, kontre enemiko\lingui{EF}
\artiklo{remplasar}%
\vorto{remplasar}\fako{trans.} Prenar, havar la plaso di (ulu, ulo)\lingui{EF}
\artiklo{remulado}%
\vorto{remulado} Sauco lango-pikanta, ek mustardo quirlita kun oleo\lingui{DEF}
\artiklo{rendevuo}%
\vorto{rendevuo} Interkonvenciono da du o plura personi pri esor en ta o ca loko, en ta o ca dio, ye ta o ca kloko, quin li fixigas pri lia inter-renkontro\lingui{DEFR}
\artiklo{rendimento}%
\vorto{rendimento}\fako{tekn.} Efekto utila di mashino, di laboro homala\lingui{DFIS}
\artiklo{renegar}%
\vorto{renegar}\fako{trans.}\senco{1}{Ne plus agnoskar kom sua}\senco{2}{Deklarar (ne-exakte) ke onu ne konocas (ulu, ulo)}\lingui{DFIRS}
\artiklo{renesanco}%
\vorto{renesanco} Epoko (qua sequas la mezepoko ed anuncas la epoko moderna) dum qua rinaskis e prosperis literaturo e la arti en Italia, en Francia, e c.\lingui{DEFIRS}
\artiklo{renet-pomo}%
\vorto{renet-pomo}\fako{bot.} Pomo di qua la shelo, makul-buntita, aspektas kelke quale la kopso-rano\lingui{DEFIR}
\artiklo{renkontrar}%
\vorto{renkontrar}\fako{trans.} Trovar dum sua voy-iro – o pro hazardo, o pro ad-iro intencita, o pro rendevuo\lingui{DEFIS}
\artiklo{reno}%
\vorto{reno}\fako{anat.} Vicero duopla, situita en la regiono lombala, qua sekrecas urino\lingui{EFIS}
\artiklo{rentiro}%
\vorto{rentiro}\fako{zool.} Quadripedo di la polo-regioni, ek la genero « cervo », kun korni plata e quaze dentoza\latina{cervus tarandus}\lingui{DE}
\artiklo{rento}%
\vorto{rento}\fako{financo} Revenuo (yarala), ordinare personala, obtenata de sulo-proprieto, de domi lugata, o kapitalo prestita a la stato od a privato\lingui{DEFIRS}
\artiklo{renuncar}%
\vorto{renuncar}\fako{trans.} Abandonar definitive, cesar sua pretendo, sua aspiro (ad lu)\lingui{EFIRS}
\artiklo{renversar}%
\vorto{renversar}\fako{trans.} Faligar ulo, ulu, adsur sulo – generale per shoko, pulso subita e ne-expektita da la pulsito
\artiklo{reostato}%
\vorto{reostato}\fako{elektro} Aparato qua inkluzas rezistivo elektrala, regulebla segun-vole, e qua, interpozite en cirkuito, posibligas la variigo di la elektrofluo\lingui{DEFIRS}
\artiklo{reotaktismo}%
\vorto{reotaktismo}\fako{biol.} Karakterizivo di ula organismi : li su movas inverse di la aquo-fluo en qua li esas imersita
\artiklo{reotropismo}%
\vorto{reotropismo}\fako{biol.} Reotaktismo qua su manifestas che enti komplexa (metafiti o metazoi)
\artiklo{reparar}%
\vorto{reparar}\fako{trans.} Igar lo necesa por igar(lu) aden lua stando antea, pos lua domajiteso\lingui{DEFIRS}
\artiklo{repartisar}%
\vorto{repartisar}\fako{trans., inter} Partigar (ulo) ed atribuar a singlu la parto qua apartenas a lu\lingui{DEFIRS}
\artiklo{repastar}%
\vorto{repastar}\fako{netrans., ek} Manjar ye kloki determinita (dejunar, dinear, supear; matino-frue : dejunetar)\lingui{EFI}
\artiklo{repentar}%
\vorto{repentar}\fako{netrans., pri} Regretar sua kulpo, dezirante emendar o kompensar lu, e ne plus agar lu\lingui{EFIS}
\artiklo{reperar}%
\vorto{reperar}\fako{trans.}\senco{1}{Determinar la situeso topografiala preciza ed exakta}\senco{2}{\textit{(teknol.)} Markizar du bloki, du organi, qui devas funcionar ensemble, por posibligar la sameso absoluta di lia situeso rispektiva lor la ri-munto}\lingui{F}
\artiklo{reperkutar}%
\vorto{reperkutar}\fako{trans.} Produktar retro-shoko. – (La korpo shokata, shokas sua-foye la korpo quan lu kontaktas o preske kontaktas)\lingui{EFIS}
\artiklo{repertorio}%
\vorto{repertorio}\senco{1}{Sorto de tabelo, en qua la tituli, la nomi, e c. esas klasifikita segun ordino qua igas posibla trovar li facile}\senco{2}{Listo de la teatraji qui permanas en la teatro e qui forsan ri-pleesos}\lingui{DEFIRS}
\artiklo{repetar}%
\vorto{repetar}\fako{trans.} Dicar plura-foye, parole o skribe, to quon onu, od altru, ja dicis\lingui{DEFIRS}
\artiklo{repleta}%
\vorto{repleta}\fako{pri vivanto} Di qua la formo esas kelke ronda, la karno ferma, la pelo pasable tensita\lingui{EFS}
\artiklo{replikar}%
\vorto{replikar}\fako{trans., ad, pri}\senco{1}{(a) Respondar a to quon ulu respondis. (b)\textit{(yuro-cienco)} Respondar a la respondo da la adverso}\senco{2}{Respondar a to quo, aspekte, ne postulas respondo}\lingui{DEFIRS}
\artiklo{reportar}%
\vorto{reportar}\fako{trans.) (en borso} Plulongigar la vendo-kontrato kun pago-tempo, quan agis spekulisto, de liquidaco-dato a la *nexta, po pagar a la kapitaliero qua satisfacas provizore la engaji financala, la difero-quanto, plus, en ula kazi, premio\lingui{DFIR}
\artiklo{repozar}%
\vorto{repozar}\fako{netrans.}\senco{1}{Jacar en stando tala ke onu cesas movar su, por desaparigar fatigeso}\senco{2}{\textit{(aludante liquido ne-limpida)} Permanar sen-mova por ke la partikuli suspendita en olu povez falar adsur la fundo di la recipiento o tanko}\lingui{EFIS}
\artiklo{represar}%
\vorto{represar}\fako{trans.} Haltigar la ageso, la exekuteso (di to quo esas kondamninda)\lingui{EFIRS}
\artiklo{reprezalar}%
\vorto{reprezalar}\fako{netrans.) (kontre}\senco{1}{Agar a la enemiki malajo po malajo, domajo po domajo}\senco{2}{Agar ad ulu insulto po insulto, moko po moko}\lingui{DEFIRS}
\artiklo{reprezentar}%
\vorto{reprezentar}\senco{1}{\textit{(trans.)} Remplasar (ulu) en la praktiko di lua yuri}\subvorto{}{reprezentar firmo}{Voyajar komise da ta firmo, komercar profite da olu}{2}{}\senco{3}{\textit{(trans.)} Rilokizar avan la okuli di ulu : (a) per la imajo quan donas la impreso, la imagino-memoro; (b) per la imito quan produktas parolo, poemo, pikturo; (c) per la signo qua rivenigas la ideo}\lingui{DEFIRS}
\artiklo{reprimandar}%
\vorto{reprimandar}\fako{trans., pri} Agar (ad ulu) blamo pri lua kulpo\lingui{EFS}
\artiklo{reprochar}%
\vorto{reprochar}\fako{trans., ad} Imputar kom kulpo (ulo) ad ulu, turnante su a lu\lingui{EFS}
\artiklo{*reproduktar}%
\vorto{*reproduktar}\fako{trans.} Igar quaze ri-existar, forme di produkturo nova per kopio, per imito
\artiklo{repso}%
\vorto{repso} Stofo ek silko e lano, od ek lano e kotono\lingui{DEFIRS}
\artiklo{reptar}%
\vorto{reptar}\fako{netrans.}\senco{1}{\textit{(aludante ula animali)} Avancar per tranar su kun ventro an sulo}\senco{2}{\textit{(avan ulu)} Humileskar desestiminde}\lingui{DEFIRS}
\artiklo{republiko}%
\vorto{republiko}\senco{1}{Guverno-sistemo en qua la suvereneso apartenas ad asemblitaro ek la populo, od a senato, ed en qua la povo exekutigar la legi ne esas heredala}\lingui{DEFIRS}\subvorto{}{la republiko literaturala}{Ensemblo de la literaturisti, egardata kom formacanta quaza kolegaro}{2}{}
\artiklo{repudiar}%
\vorto{repudiar}\fako{trans.} Cesar agnoskar legale (kom sua spozino)\lingui{EFIS}
\artiklo{repugnar}%
\vorto{repugnar}\fako{trans.} Produktar antipatio nevinkebla, nerepresebla\lingui{EFIS}
\artiklo{repulsar}%
\vorto{repulsar}\fako{trans.}\senco{1}{Pulsar ad-retroe, igar *rekular (desavancar) per retro-pulso. \textit{(anke metaf.)}}\senco{2}{Ne-aceptar pro lua mala qualeso, pro lua ne-konformeso ye la komendo (varo jus livrita)}\lingui{DEFIS}
\artiklo{reputar}%
\vorto{reputar}\fako{trans., kom} Igar egardesar kom tala o tala\lingui{DEFIRS}
\artiklo{requiem}%
\vorto{« requiem »}\fako{liturgio katolika}\senco{1}{Prego por la mortinti}\senco{2}{Meso di « requiem » : Meso por la repozo di la anmo di mortinto}\senco{3}{Parti di la mortinto-meso, igita muzikajo}
\artiklo{requizitar}%
\vorto{requizitar}\fako{trans.) (aludante la stat-autoritatozi} Postular ke persono o kozi esez igata en la dispono di ta autoritatozi por servado od uzado publiko-destine \textit{(*publico-destine)}\lingui{DEFIRS}
\artiklo{resekar}%
\vorto{resekar}\fako{trans.) (kirurg.} Ablacionar un ek la extremaji di osto malada
\artiklo{resipicencar}%
\vorto{resipicencar}\fako{netrans.} Repentar ye maniero qua eventigas la ri-komenco agar etiko-konforme\lingui{DEFI}
\artiklo{reskripto}%
\vorto{reskripto}\senco{1}{\textit{(yuro-cienco, che la Romani antiqua)} Respondo, da la imperiestro, a la questioni agita a lu da la guberniestri, da la judiciisti, e c., pri ula problemi solvenda}\senco{2}{\textit{(en nia landi, nun)} Dekreto da la stato-chefo}\senco{3}{\textit{(Vatikano)} Averto papala, qua decidas pri ta o ca yuro-punto}\lingui{DEFIRS}
\artiklo{resonar}%
\vorto{resonar}\fako{netrans.}\senco{1}{Retrosendar la soni per reflekto}\senco{2}{Produktar soni qui igesas pluforta per la reflekto interne di klozajo (*kluzajo) di qua la kapaceso ne esas sat granda pri facar eko}\lingui{DEFIS}
\artiklo{resorbar}%
\vorto{resorbar}\fako{trans.) (medic.} Igar retroirar aden la cirkulado (liquido quan la vaskuli lasabis eskapar)\lingui{DEFIS}
\artiklo{resortisar}%
\vorto{resortisar}\fako{netrans., ad} Esar judicienda da ta o ca tribunalo\lingui{DEFRS}
\artiklo{resorto}%
\vorto{resorto}\senco{1}{Peco di mekanismo qua, per sua destenso, movigas peco vicina}\senco{2}{\textit{(metaf.)} To quo impulsas}\lingui{FRS}
\artiklo{respektar}%
\vorto{respektar}\fako{trans.}\senco{1}{Egardar kom kozo quan onu devas atencar pro lua importo}\senco{2}{Egardar kom kozo quan onu devas ne domajar, ne lezar, pro lua bonegeso}\senco{3}{Honorizar (ulu) per deferencar profunde lu, pro lua karaktero, pro lua alta situeso en la hierarkio}\lingui{DEFIRS}
\artiklo{respirar}%
\vorto{respirar}\fako{trans. e netrans.} Aspirar aero por regenerar la sango, e, pose, expirar lu\lingui{DEFIS}
\artiklo{respondar}%
\vorto{respondar}\fako{trans., ad} Sendar, direktar ad ulu di qua onu recevis demando, questiono, akuzo, e c. to quon onu havas dicenda kambie\lingui{EFIS}
\artiklo{responsar}%
\vorto{responsar}\fako{netrans., pri} Obligesar reparor, per indemno, la detrimento o domajo quin onu agis od agos a kunhomo\lingui{EFIS}
\artiklo{responsorio}%
\vorto{responsorio}\fako{liturgio katol.} Citaji ek biblo, dicata o kantata pos omna leciono di matutino\lingui{DEFIS}
\artiklo{restar}%
\vorto{restar}\fako{netrans.} Esar ankore prezenta pos ke onu deprenis un o plura de lua parti\lingui{DEFIS}
\artiklo{restaurar}%
\vorto{restaurar}\fako{trans.) (pri edifico, statuo, e c} Igar en bona stando itere\lingui{DEFIRS}
\artiklo{restitucar}%
\vorto{restitucar}\fako{trans., ad} Retrodonar ad ulu (to quon onu prenis de lu)\lingui{DEFIS}
\artiklo{restorar}%
\vorto{restorar}\fako{trans.} Donar (ad ulu) la manjebli e drinkebli necesa por ke lu ne plus hungrez nek durstez\lingui{DEFIRS}
\artiklo{restriktar}%
\vorto{restriktar}\fako{trans., ad} Retroduktar aden limiti plu streta\lingui{DEFIRS}
\artiklo{retablo}%
\vorto{retablo}\senco{1}{\textit{(tekn.)} Parto dopa di altaro (en kirko) qua supersalias la tablo}\senco{2}{Panelo kontre qua, en kirko, apogesas altaro e quan ornas, en la kazi maxim multa, pikturo, skulturo, en la parto qua superstacas la altaro}\lingui{FS}
\artiklo{*retarda}%
\vorto{*retarda} Qua eventas pos la instanto (1) stipulita o konvencionita – (2) ye qua lu eventas normale
\artiklo{retenar}%
\vorto{retenar}\fako{trans.}\senco{1}{Ne lasar de-irar (persono o kozo); impedar (persono) agar impulsate da pasiono; ne retrodonar (kozo), ne restitucar (kozo)}\senco{2}{Rezervigar por su (da ulu)}\lingui{EFIS}
\artiklo{retiario}%
\vorto{retiario}\fako{historio di la Roma antiqua} Gladiatoro di qua la armo esis reto per qua lu probis enplektar sua adverso\lingui{EFIS}
\artiklo{retikulo}%
\vorto{retikulo}\senco{1}{\textit{(optiko)} Fili interkrucumanta quin onu lokizas sur la objektalo di la teleskopi, por determinar ula punti di la feldo di ta aparati}\subvorto{biol.}{retikulo nukleala}{Parto diferenciita di la nukleo celula, qua konsistas ye substanco proteinala komplexa \textit{(linino)}, e qua rezultas de la pelotesko di filamento unika e kontinua}{2}{}
\artiklo{retino}%
\vorto{retino}\fako{anat.} Membrano quan formacas, sur la fundo di la okulo, la expanseso di la nervo vidala, sur qua venas facar su la imajo de la objekti\lingui{DEFIRS}
\artiklo{reto}%
\vorto{reto}\senco{1}{Texuro ek mashi interligita per nodi plu o min inter-distanta, facita ek kordeto, filo, e c., per bastono cilindra e naveto}\senco{2}{Produkturo reto-forma ek filo, silko, filo ora od arjenta}\senco{3}{Interplekturo de linei (voyi, kanali, relvoyi, e c.)}\lingui{EFIS}
\artiklo{retoro}%
\vorto{retoro}\senco{1}{\textit{(epoki antiqua, en Roma e Grekia)} Profesoro pri eloquenteso}\senco{2}{\textit{(nun)} Oratoro qua probas nur facar bela frazi}\lingui{DEFIS}
\artiklo{retorto}%
\vorto{retorto}\fako{kemio}\senco{1}{Recipiento vitra, gresa, kun kolo streta e kurvigita, uzata da la kemiisti}\senco{2}{\textit{(tekn.)} Tre granda recipiento ek tero refraktara, en qua onu varmigas min-karbono por la fabriko di lum-gaso}\lingui{DEIRS}
\artiklo{retraktar}%
\vorto{retraktar}\fako{trans.} Retroprenar, nihiligar, explicite, formale (to quon onu dicis, o plendo ad autoritato judiciala)\lingui{EFIS}
\artiklo{retretar}%
\vorto{retretar}\fako{netrans., de, ek}\senco{1}{\textit{(strategio)} Retro-iro \textit{(*rekulo)} da trupo de militisti qua koaktesas agar tale da la enemiko plu forta qua avancas}\senco{2}{Cesar sua vivo profesionala, stat-employatala, mestierala, mondumala}\lingui{DEFIRS}
\artiklo{retro-}%
\vorto{retro-} Prefixo qua indikas movo de \textit{avan} ad \textit{dop} – de ula epoko ad epoko antea\lingui{L}
\artiklo{retushar}%
\vorto{retushar}\fako{trans.} Ri-laboreskar (ulo) por korektigar o plubonigar ta o ca ek lua parti\lingui{DEFIRS}
\artiklo{reumatismo}%
\vorto{reumatismo}\fako{patol.} Dolorigo qua su manifestas sen inflamuro, en la muskuli, en la artiki, e di qua la kauzo ne esas determinita bone\lingui{DEFIRS}
\artiklo{revanchar}%
\vorto{revanchar}\fako{netrans., pri} Retroprenar de ulu la avantajo quan lu prenis de ni sen nia konsento\lingui{DEFI}
\artiklo{revar}%
\vorto{revar}\fako{netrans., pri} Lasar sua penso vagar hazarde\lingui{EF}
\artiklo{revelar}%
\vorto{revelar}\fako{trans., per, pri, ad}\senco{1}{Konocigar (to quo esis celita)}\senco{2}{\textit{(teol.)} \textit{(aludante Deo)} Konocigar a la homi per inspiro supernatura}\senco{3}{\textit{(fotogr.)} Eventigar reakto kemiala qua igas la imajo aparar (sur la vitro-plako o sur la filmo fotografala)}\lingui{EFIS}
\artiklo{revendikar}%
\vorto{revendikar}\fako{trans.} Postular la retrodono, kom apartenajo nia, di to quon kunhomo ditenas\lingui{DFIS}
\artiklo{revenuo}%
\vorto{revenuo}\fako{financo} To quon ulu obtenas, singla-yare, de sua kapitalo, de sua domeno, de pensiono, de rento\lingui{DEF}
\artiklo{reverberar}%
\vorto{reverberar}\fako{trans.} \textit{(fiziko)} \textit{(aludante surfaco)} Retrosendar la lumo-radii o la kaloro-radii qui frapas lu\lingui{DEFIS}
\artiklo{reverencar}%
\vorto{reverencar}\fako{trans.} Movar sua korpo por salutar ye maniero qua manifestas respekto profunda\lingui{DEFIRS}
\artiklo{reversionar}%
\vorto{reversionar}\fako{yuro-cienco} Dicesas pri havaji qui retroiras a la persono qua donacabis li a kunhomo, kaze ke ica mortas sen filii\lingui{EFIRS}
\artiklo{reverso}%
\vorto{reverso}\senco{1}{La facio, opozata a la facio quan, en kozo, onu regardas, o regardigas, prefere}\senco{2}{La facio di medalio, di moneto-peco, opozata a la facio sur qua esas reprezentita figuro}\lingui{DEFIS}
\artiklo{revizar}%
\vorto{revizar}\fako{trans.} Submisar itere (texto, proceso, kodexo konto) ad exploro por reformar, korektigar, se necesa\lingui{DEFIRS}
\artiklo{revokar}%
\vorto{revokar}\fako{trans.) (pro} Retrotirar (lu) de lua ofico \textit{(pro nekapableso, pro kulpo profesionala, e c.)}\lingui{EFIS}
\artiklo{revoltar}%
\vorto{revoltar}\fako{netrans., kontre, pri, pro} Refuzar la obedio di la imperi de la autoritatozi legitima, regnanta. – \textit{(anke metaf.)}\lingui{DEFIRS}
\artiklo{revolucionar}%
\vorto{revolucionar}\fako{netrans.) (aludante la populo di stato} Agar violente la chanjo di (sua) guvernistaro\lingui{DEFIRS}
\artiklo{revolvero}%
\vorto{revolvero} Pafarmo portebla, pistoleto qua posibligas pafar plura-foye sen interrupto, per provizesir ye plura kulati quin onu povas charjar anticipe, e quin resorto igas pasar singlope avan la tubo o la tubi di la armo\lingui{DEFIRS}
\artiklo{revuar}%
\vorto{revuar}\fako{trans.}\senco{1}{\textit{(armeo)} Manovro, defilo, di la armo-trupi, di qua la skopo esas certigar su \textit{(igar su *certena)} pri la korekteso di lia su-teno, instrukteso, precizeso di movi e gesti}\senco{2}{Explorar la aktualaji en edituro periodala}\lingui{DEFIS}
\artiklo{revulsar}%
\vorto{revulsar}\fako{trans.) (patol.} Igar inflamantajo, obstruktantajo, e c. deviacar de la loko malada, per atraktar lu ad altra parto di la korpo ube lu ne esas danjeroza\lingui{EF}
\artiklo{reza}%
\vorto{reza}\senco{1}{\textit{(pri pili, hari, e c.)} Tondita tale ke li ne salias de la surfaco cirkondanta}\senco{2}{Di qua la surfaco esas sen saliaji nek sulki}\senco{3}{\textit{(pri recipiento e lua kontenajo solida)} Qua ne salias super la bordo o bordi}\lingui{EF}
\artiklo{rezedo}%
\vorto{rezedo}\fako{bot.} Planto herbacea, kun flori odoranta, folii dispozita sur la stipo intervaloze, quan onu kultivas en la gardeni\latina{reseda}\lingui{DEFIRS}
\artiklo{rezervar}%
\vorto{rezervar}\fako{trans.}\senco{1}{Retenar, gardar, ula-destine}\senco{2}{Ne lasar sua impresesi, sua sentimenti, konocesar da kunhomi. – \textit{(Ca-kaze, uzesas nur ye la formo « esar rezervita »)}}\lingui{DEFIRS}
\artiklo{rezidar}%
\vorto{rezidar}\fako{netrans., en, che}\senco{1}{Esar establisita nun en ta o ca loko. Permanar en la loko ube onu havas sua ofico}\senco{2}{Havar sua *siejo (en, che)}\lingui{DEFIS}
\artiklo{rezidento}%
\vorto{rezidento} Agento diplomacala qua rezidas apud la guvernistaro di mikra lando, kun grado infra ye olta di ambasadisto o di plenipotenciero\lingui{DEFIS}
\artiklo{reziduo}%
\vorto{reziduo} La materio qua restas pos operaco kemiala, pos manipulo industriala; ta parto di la dishi qua ne manjesis dum la repasto; la restajo di objekto ruptita acidente\lingui{DEFIS}
\artiklo{rezignar}%
\vorto{rezignar}\fako{netrans., pri} Aceptar volunte desfeliceso, domajo, perdo, quan onu subisis e, generale, ne povis evitar\lingui{DEFIS}
\artiklo{rezino}%
\vorto{rezino} Materio inflamebla qua defluas, naturale o per incizo, de la pino, pinastro, abieto\lingui{EFIS}
\artiklo{rezistar}%
\vorto{rezistar}\fako{trans.}\senco{1}{Opozar, a la efiko da forco, forcon qua agas inversa-sinse}\senco{2}{Aplikar esforco kontre la uzo di forco, di violento}\senco{3}{Opozar sua volo ad impulso, a volo adversa}\lingui{EFIS}
\artiklo{rezolvar}%
\vorto{rezolvar}\fako{trans.} Determinar, pos decido, to quo esas agenda praktikale, aplikale\lingui{DEFIS}
\artiklo{rezonar}%
\vorto{rezonar}\fako{netrans., pri} Agar kateniguro ek propozicioni qui esas apta demonstrar ulo\lingui{DEFIRS}
\artiklo{rezultar}%
\vorto{rezultar}\fako{netrans., de} Eventar konseque de ago, de fakto\lingui{DEFIRS}
\artiklo{rezumar}%
\vorto{rezumar}\fako{trans.} Repetar to quo dicesis skribe o voce, ma agar lo plu kurte, per expresar nur lo esencala o quintesencala\lingui{DEFIS}
\artiklo{rezurektar}%
\vorto{rezurektar}\fako{netrans.) (religio kristana} Retrovenar de mortinteso a vivo, per miraklo
\artiklo{ri-}%
\vorto{ri-} Prefixo qua indikas ago quan onu iteras, iteris od iteros
\artiklo{ribo}%
\vorto{ribo}\fako{bot.} Frukto di arbusto ek la familio « grosulariei »\latina{ribes rubrum}\lingui{IL}
\artiklo{richa}%
\vorto{richa} Qua havas propriete multa havaji importoza. (richa de)\lingui{DEFIS}
\artiklo{ricino}%
\vorto{ricino}\fako{bot.} Genero de planti, ek la familio « euforbiacei », di qua un speco donas oleo purgiva\latina{ricinus}\lingui{DEFISL}
\artiklo{rictus}%
\vorto{« rictus »} Fenduro di la boko, qua lasas la denti videsar
\artiklo{ridar}%
\vorto{ridar}\fako{netrans., pri, pro, de} Sentar quaza expanseso di sua vizajo quan akompanas expiri sukusoza plu o min bruisanta, signi korpala di gayesko subita\lingui{FIS}
\artiklo{ridelo}%
\vorto{ridelo}\fako{tekn.} Ligno-peco horizontala qua, ye singla latero di chareto, formacas, per la ronda ligno-cilindri qui tra-iras lu, quaza rastelero destinita impedor la falo di la kargajo\lingui{FI}
\artiklo{rifo}%
\vorto{rifo}\fako{navig.} En maro, rokajo, grupo de rokaji o de sablo o de koralii, nur kelke salianta de la aquo-surfaco, kontre qua povas su ruptar navi\lingui{DER}
\artiklo{rigida}%
\vorto{rigida}\senco{1}{Qua ne cedas se onu probas flexar lu}\senco{2}{\textit{(metaf.)} Qua ne igas flexar la reguli; quan onu ne povas igar deviacar de la apliko di la reguli}\lingui{DEFIS}
\artiklo{riglo}%
\vorto{riglo} Mikra peco fera, fixigita sur pordo, sur fenestro, e qua, se pulsita aden seruro-boketo, impedas la aperto\lingui{D}
\artiklo{rigo}%
\vorto{rigo}\senco{1}{\textit{(navig.)} To quo uzesas por manovrar navo (kordegi, segli, e c.) od aerostato}\senco{2}{Singla de la kordegi, korda skali, ligna stangi qui pendas de la portiko di gimnastikeyo}\lingui{E}
\artiklo{rigoro}%
\vorto{rigoro} Precizeso, strikteso, extrema\lingui{DEFIS}
\artiklo{rikochar}%
\vorto{rikochar}\fako{netrans., sur} Dicesas pri la retro-salto de stono plata qua frapas la surfaco di aquo, o projektilo de pafarmo qua frapas sur sulo\lingui{DEFR}
\artiklo{rimborsar}%
\vorto{rimborsar}\fako{trans., ad} Donar (ad ulu) la equivalanto di la sumo de to quon lu preste spensis\lingui{DEFIS}
\artiklo{rimeno}%
\vorto{rimeno}\senco{1}{Bendo de ledro, longa e streta}\senco{2}{Bendo de ledro di qua un extremajo esas garnisita ye buklo, e la altra havas longeso-sinse, serio de trui aden qui eniras la buklo-pinto segun ke onu volas plektar plu o min strete}\senco{3}{Bendo kontinua de ledro, de kauchuko, di qua la extremaji interligesis tale ke formacesas cirklo-kurvo flexebla, uzata por transmisar la movo di mashino}\lingui{DR}
\artiklo{rimo}%
\vorto{rimo} Konsonanco di la lasta silabo acentoza di la vorto finala en du o plura versi\lingui{DEFIRS}
\artiklo{rinforcar}%
\vorto{rinforcar}\fako{trans., per}\senco{1}{\textit{(muziko)} Igar (sono) plu intensa, plu energioza}\senco{2}{\textit{(milit.)} Igar plu multa la efektivo de la militisti di garnizono}\lingui{EFIS}
\artiklo{ringo}%
\vorto{ringo}\senco{1}{Cirklokurvo de metalo o de altra materio rezistiva quan onu uzas, en la kazi maxim multa, por retenar, ligar}\senco{2}{To quo, sua-forme, similesas ringo}\senco{3}{\textit{(geom.)} Parto de plano, qua kontenesas inter du cirkli sama-centra}\lingui{DE}
\artiklo{rinocero}%
\vorto{rinocero}\fako{zool.} Granda mamifero, sovaja, ek la klaso « pakidermi », kun un o du korni sur la nazo\latina{rhinoceros}\lingui{DEFIS}
\artiklo{rinoplastiko}%
\vorto{rinoplastiko} Operaco kirurgiala qua konsistas ye rifacar la nazo per la grefto animalala, per pelo-peco deprenita de la fronto de la brakio, e c.
\artiklo{rinsar}%
\vorto{rinsar}\fako{trans., per} Pasar en aquo pura, to quo esas jus lavita, por igar lu absolute neta\lingui{EF}
\artiklo{ripar}%
\vorto{ripar}\fako{netrans.} Dicesas pri roto di biciklo o di automobilo qua glitas skrapante aden direciono qua esas obliqua ye lua plaso propra\lingui{F}
\artiklo{riskar}%
\vorto{riskar}\fako{trans.} Expozar (ulu, ulo) a chanco pri qua onu certa (*certena) ke lu esos bona\lingui{DEFIRS}
\artiklo{rismo}%
\vorto{rismo} Bloko de (en Britania) 480 od (en mondo, ecepte Britania) 500 folii de papero\lingui{DEIS}
\artiklo{risolar}%
\vorto{risolar}\fako{koquarto} Igar (kom exemplo : *kartofli – terpomi jus fritita) krakanta quale pano-krusto\lingui{FI}
\artiklo{rispektiva}%
\vorto{rispektiva}\fako{aludante plura personi, plura kozi} Qua koncernas singla de li, relate a le cetera\lingui{DEFIRS}
\artiklo{ristornar}%
\vorto{ristornar}\fako{trans.}\senco{1}{\textit{(navig.)} Abrogar, nihiligar (interkonsente) asekur-akto navala konseque de cirkonstanci qui eventis ante la departo di la navo}\senco{2}{Diminutar la pekunio-quanto ye qua onu asekuris la navo, kande la valoro di lua omna kargaji esas fakte min granda kam expektita}\lingui{DFIRS}
\artiklo{ritmo}%
\vorto{ritmo}\fako{muziko} Aparo inter-simetra di la tempi forta e di la tempi febla, qua rieventas periodope en frazo muzikala, en verso, en tamburo-frapado\lingui{DEFIRS}
\artiklo{ritornelo}%
\vorto{ritornelo}\fako{muziko} En la akompano di kanto, kurta frazo muzikala ye qua onu igas preirata o sequata singla kupleto\lingui{DEFIS}
\artiklo{rituo}%
\vorto{rituo}\fako{liturgio} La ceremonii di kulto\lingui{DEFIRS}
\artiklo{rivalo}%
\vorto{rivalo} La persono qua kontestas (ulo ad ulu); singla de la personi qui inter-luktas por havar ulo, ulu\lingui{DEFIS}
\artiklo{rivero}%
\vorto{rivero} Aquo-fluo inter du rivi, qua fluas aden fluvio\lingui{EF}
\artiklo{riveto}%
\vorto{riveto}\fako{tekn.} Klovo di qua la pinto esas abatebla ed aplastebla per martelago (pos igir lu trairar plako) por fixigar lu\lingui{EF}
\artiklo{rivo}%
\vorto{rivo} Parto di la tero-sulo qua esas quaze la bordo di aquoextensajo, olua limito\lingui{FS}
\artiklo{rizo}%
\vorto{rizo}\fako{bot.}\senco{1}{Planto cereala, graminea, kultivata en la regioni di qui la atmosfero esas tre varma}\senco{2}{La grani de ta planto}\latina{oryza sativa}\lingui{DEFIRS}
\artiklo{rizopodo}%
\vorto{rizopodo}\fako{zool.} Protozoo ek un celulo nuda, qua havas prolonguri (pseudopodi) quin lu uzas kom organi lokomocala e pren-ala\lingui{DEFIS}
\artiklo{robineto}%
\vorto{robineto}\fako{tekn.} Ajuto muntita an la extremajo extera di tubo di fonteno, di tanko, di cisterno, e trairata da klefo qua, segun la sinso di lua rotacigeso, retenas o lasas eskapar la liquido\lingui{FI}
\artiklo{robinio}%
\vorto{robinio}\fako{bot.} Genero de arbori ek la familio « leguminosi », di qui la speco maxim konocata esas granda arboro kun flori blanka, odoroza, qui kreskas grape, e di qua la ligno densa esas apta por facar chari\latina{acacia}
\artiklo{robo}%
\vorto{robo} Longa vesto qua decensas til la pedi e varias, pri lua formo e lua naturo, segun la persono qua *weras lu, o segun lua destineso\lingui{DEFS}
\artiklo{robro}%
\vorto{robro}\fako{en la wisto-ludo} Partio ligita\lingui{DEFR}
\artiklo{robusta}%
\vorto{robusta} Di qua la konstituciono esas tre forta, tre vigoroza, tre rezistiva\lingui{DEFIS}
\artiklo{rocheto}%
\vorto{rocheto}\fako{liturgio katol.} Surpliso kun maniki rekta, ek stofo blanka e tre dina, quan *weras la episkopi ed ula alt-oficieri di la eklezio katolika\lingui{DEFIS}
\artiklo{rodar}%
\vorto{rodar}\fako{trans.}\senco{1}{Entamar per mikra dento-mordi, beko-frapi}\senco{2}{Konsumar, destruktar (korpo) per efiko lenta e neperceptebla}\lingui{EFIS}
\artiklo{rodio}%
\vorto{rodio}\fako{kemio} Metalo harda e facile ruptebla, qua aspektas quale aluminio, esas duktila e maleebla kande inkandecanta til redeso; lu fuzas ye 2000° C\simbolo{Rh}
\artiklo{rododendro}%
\vorto{rododendro}\fako{bot.} Genero de planti ek la familio « erikacei », arboreto sempre verda, remarkinda pro la beleso di lua flori\latina{rhododendron}\lingui{DEFIS}
\artiklo{rodomonto}%
\vorto{rodomonto} Falsa bravo qua asertas ocidir multa enemiki ed esar kapabla trafendar irgu\lingui{EFI}
\artiklo{rogacioni}%
\vorto{rogacioni}\fako{religio katol.} Prego publika, agata procesione, dum la tri dii qui preiras la Acenso-festo, por obtenar bona rekolti
\artiklo{rogo}%
\vorto{rogo} Amaso de ligno brancha sur qua la antiqui kustumis kombustar sua mortinti\lingui{IL}
\artiklo{rokambolo}%
\vorto{rokambolo}\fako{bot.} Alio reda di Hispania, planto liliacea, perena, bulboza. Lua bulbi saporas kelke quale alio, ma plu feble, ed uzesas kom kondimento. ophioscorodon\latina{allium sativum var}\lingui{DEF}
\artiklo{roko}%
\vorto{roko}\senco{1}{Blokego de petro tre rezistiva qua retenesas da sulo}\senco{2}{Blokego de petro harda, eskarpa}\senco{3}{\textit{(metaf.)} Quan nulo povas emocigar}\senco{4}{\textit{(anat.)} Un ek la tri parti di la osto temporala, tale nomizita pro lua hardeso}\lingui{EFIS}
\artiklo{rokoko}%
\vorto{rokoko}\senco{1}{Stilo ornementala qua esis segun-moda en la yarcento 18-esma}\senco{2}{\textit{(metaf.)} Stilo ornementala obsoleta, antiquatra}\lingui{DEF}
\artiklo{rolero}%
\vorto{rolero} Cilindro, ligna, metala, e c., quan onu promenas sur agro por ruptar la glebi, por kunpresar tero; por planigar la gazon-agri; por extensar pasto por la kukaji; por extensar la inko sur la komposturi di imprimilo, e c.\lingui{EF}
\artiklo{rolo}%
\vorto{rolo}\senco{1}{\textit{(teatro)} (a) Transskriburo di to quon aktoro devas recitar ek teatrajo. – (b) To quon la aktoro devas recitar. (c) Persono quan reprezentas la aktoro}\senco{2}{\textit{(metaf.)} (a) Lo agenda kom funciono en ofico. – (b) Konduto-maniero di ulu en ta o ca cirkonstanco}\lingui{DEFR}
\artiklo{romanco}%
\vorto{romanco}\senco{1}{Mikra kansono di qua la karaktero esas sentimentala}\senco{2}{Melodio qua havas tala karaktero}\lingui{DFIRS}
\artiklo{romano}%
\vorto{romano}\senco{1}{Verko literaturala, fingura, en qua la autoro esforcas interesar la lektero, sive per la singulareso di la aventuri, sive per la deskripto di la mori, di la pasioni, e c.}\senco{2}{Aventuri romanatra extraordinara}\lingui{DEFIRS}
\artiklo{romantika}%
\vorto{romantika} Qua apartenas a la sistemo literaturala qua refuzas la imito klasika di la Greki e di la Romani antiqua, e volas liberigar su de la reguli kompozala qui konsequas de ta imito\lingui{DEFIRS}
\artiklo{rombo}%
\vorto{rombo}\fako{geom.} Quadrilatero, kun lateri inter-egala e kun anguli ne-orta
\artiklo{romboedro}%
\vorto{romboedro}\fako{geom.} Korpo solida, kristalajo, di qua la edri esas romba\lingui{DEF}
\artiklo{ronda}%
\vorto{ronda} Di qua la formo esas cirkla\lingui{DEFIS}
\artiklo{rondelo}%
\vorto{rondelo}\senco{1}{\textit{(poezio)} Mikra poemo Franca, ek dek-e-tri versi (ok kun ula rimo, e kin kun altra rimo) interseparita per pauzo ye la verso kinesma, e ye la okesma – la vorti versokomencala iteresas pos la verso okesma e la dek-e-triesma, sen esar parti de li}\senco{2}{\textit{(muziko)} Muzikajo di qua la temato rivenas plurafoye}\lingui{DEFIRS}
\artiklo{ronkar}%
\vorto{ronkar}\fako{netrans.} Emisar, dormante, respiro nazala tre bruisoza\lingui{FS}
\artiklo{ronronar}%
\vorto{ronronar}\fako{netrans.) (aludante kato} Agar mikra grondado de kontenteso\lingui{F}
\artiklo{roquar}%
\vorto{roquar}\fako{netrans.) (shako-ludo} Pozar la turmo apud la rejo, e la rejo trans la turmo kande nula piono stacas inter li, o dextre, o sinistre, kande onu judikas lo kom oportuna\lingui{F}
\artiklo{rosbifo}%
\vorto{rosbifo}\fako{koquarto} Peco de bovokarno rostita – maxim-multakaze pseudo-fileto\lingui{DEFIRS}
\artiklo{rosmarino}%
\vorto{rosmarino}\fako{bot.} Planto aromata, ek la familio « labiacei », qua kreskas precipue en maro\latina{rosmarinus}\lingui{DEFIR}
\artiklo{rosmaro}%
\vorto{rosmaro}\fako{zool.} Mamifero amfibia, karnavida, analoga a maro-hundo\latina{trichochus rosmarus}\lingui{L}
\artiklo{roso}%
\vorto{roso}\fako{meteorol.} Strato de aquo-guteti qua aparas quale sedimento, nokte, sur la surfaco di ula korpi qui expozesas en aero ne-inkluzita – aquo-vaporo de atmosfero, kondensita da la varmesko quan produktas la radiado noktala\lingui{FRS}
\artiklo{rospo}%
\vorto{rospo}\fako{zool.} Reptero batraka, kun gambi plu kurta kam olti di rano, korpo grosa e kurta, sur qua abundegas la veruki, de qui exudas odoro fetida\latina{bufo}\lingui{I}
\artiklo{rostar}%
\vorto{rostar}\fako{trans.} Koquar per fairo intensa (karno, pultruni) sur spiso, en forno, sur grililo\lingui{DEFI}
\artiklo{rostro}%
\vorto{rostro}\senco{1}{\textit{(che la antiqui)} Beko, quaze sporno, ye qua (kom armo) onu provizis la pruo di la navi, por plufaciligar la asaltabordo di la navi enemika}\senco{2}{\textit{(zool.)} Longa prolonguro di la nazo, che la elefanto, quan la animalo uzas quale manuon. (analoge pri la insekti)}\senco{3}{\textit{(arkitekt.)} Ornemento di arkitekturo o di skulturo, forme di beko, di sporno}
\artiklo{rotacar}%
\vorto{rotacar}\fako{netrans.) (aludante korpo} Movar su cirkum sua axo fixa\lingui{DEFIRS}
\artiklo{rotano}%
\vorto{rotano}\fako{bot.} Genero de palmiero, kun stipo longa e tenua\latina{calamus rotang}\lingui{DEFIRS}
\artiklo{rotifero}%
\vorto{rotifero}\fako{zool.} Genero de animali mikroskopala, qui divenas sen-mova kande vetero esas sika, e rimoveskas lor humidesko di atmosfero\lingui{EFIS}
\artiklo{roto}%
\vorto{roto}\fako{tekn.} Peco rigida, cirkla, qua rotacas cirkum axo qua trairas lua centro, e pleas la rolo di motoro, generale kun felgo e radii\lingui{FIS}
\artiklo{rotolar}%
\vorto{rotolar}\fako{netrans.} Produktar bruiso (quale per frapi qui eventas tre kurta-intertempe) kontinua, analoga ad olta de veturo qua rulas rapide sur voyo ne-glata
\artiklo{rotondo}%
\vorto{rotondo}\senco{1}{\textit{(edifico)} Edifico cirklatra, quaze kronizata ye kupolo}\senco{2}{\textit{(vesto)} Longa manteleto sen maniki qua envolvas la korpo}\lingui{DEFIRS}
\artiklo{rotoro}%
\vorto{rotoro}\fako{tekn.} Nomo di la parto nefixa, en la elektro-motori kun fluo alternanta – opoze statoro, qua esas la parto fixa
\artiklo{rotulo}%
\vorto{rotulo}\senco{1}{\textit{(anat.)} Mikra osto plata, kun anguli kelke rondigita, ye qua konsistas la parto avana di la genuo}\senco{2}{\textit{(tekn.)} Peco metala, sfero-forma, uzata kom artiko en la mashin-organi qui bezonas povar movar su omna-direcione}\lingui{FIS}
\artiklo{rovo}%
\vorto{rovo}\fako{bot.} Arbusto dornoza e reptera, ek la familio « rozacei », qua kreskas en la hegi, la boski, e donas frukto (rovbero)\latina{rubus caesius, rubus fructicosus}\lingui{I}
\artiklo{rovro}%
\vorto{rovro}\fako{bot.} Querko blanka, min granda kam la querko ordinara\latina{quercus robur}\lingui{FIS}
\artiklo{royalismo}%
\vorto{royalismo} Mento monarkiista. (« En rejio, royalismo esas o partiso, od opiniono ».)\lingui{DEFR}
\artiklo{rozario}%
\vorto{rozario} Granda chapleto, ek dek-e-kin deki de grani, sur singla de qui onu recitas (che la katoliki) un « Ave » – singla de qui pre-iresas da grano plu grosa sur qua onu recitas un « Pater »\lingui{DEFIS}
\artiklo{rozeolo}%
\vorto{rozeolo}\fako{patol.} Pel-erupto febla, karakterizata da mikra makuli rozea, kurta-dura, sen febro, qua eventas precipue che la infanteti e la infanti\lingui{EFIS}
\artiklo{rozlauro}%
\vorto{rozlauro}\fako{bot.} Arboreto ek la familio « apocinei »\latina{nerium oleander}
\artiklo{rozo}%
\vorto{rozo}\fako{bot.} Flori di arbusto ek la familio « rozacei », qua odoras dolcege, e di qua la tipo primitiva esas ye koloro delikate pale-reda.\latina{rosa}\lingui{DEFIRS}\subvorto{}{vento-rozo}{Plako qua subtenas « rozego » 32-folia, qua dividas la cirklokurvo ad 32 parti inter-egala, kun la indiko, en singla de ta dividuri, di la rumbo}{}{}
\artiklo{rubando}%
\vorto{rubando} Mikra bendo streta de stofo fila, lana, silka, velura, de celofano, de plastiko\lingui{EF}
\artiklo{rubarbo}%
\vorto{rubarbo}\fako{bot.} Planto ek la familio « poligonei », kun larja folii, di qua la radiki uzesas kom tonizivo e purgivo, e la pedunkli, por facar kompoti e tarti\latina{rheum rhubarba}\lingui{DEFIS}
\artiklo{rubidio}%
\vorto{rubidio}\fako{kemio} Metalo alkalioza, qua existas en ula aqui minerala, en ula vejetanti (betravo, tabako, teo, kafeo), en la lepidolito, bazalto, e c., e deskovrita samatempe kam cezio, da Kirchhoff e Bunsen\simbolo{Rb}
\artiklo{rubino}%
\vorto{rubino} Lapido, reda bele e diafane\lingui{DEFIRS}
\artiklo{rubio}%
\vorto{rubio}\fako{bot.} Planto ek la familio « rubimacei », di qua la radiko, kande sikigita e pulverigita, donas tintivo reda\latina{rubia tinctorum}\lingui{ISL}
\artiklo{rubriko}%
\vorto{rubriko} En revuo, en jurnalo : titulo generala (kom exemplo : « Progreso ») kun tituli specala (kom exemplo, ek « Progreso » – rubriki esas « Questioni linguala », « Kroniko », « Bibliografio », « Diversaji ».)\lingui{DEFIS}
\artiklo{ruda}%
\vorto{ruda} Di qua la acepto-maniero esas ne-delikata, kelke maletraktanta, mem se nur parole\lingui{DEFI}
\artiklo{rudimento}%
\vorto{rudimento}\senco{1}{La komenc-elementi di arto, di cienco}\senco{2}{Mikra libro qua kontenas la komenc-elementi di la gramatiko Latina}\senco{3}{Komenco-lineamento di la strukturo di organo}\lingui{DEFIRS}
\artiklo{rufa}%
\vorto{rufa} Reda kelke-flava\lingui{L}
\artiklo{rugo}%
\vorto{rugo} Plisuro di la pelo sur la fronto, la vizajo, la manui, quan produktas, che la oldi, la mult-evozeso. – \textit{(metaf. : Sulketo sur surfaco di irgo)}\lingui{FIS}
\artiklo{ruinar}%
\vorto{ruinar}\fako{trans.}\senco{1}{Igar (konstrukturo) krular}\senco{2}{Igar (persono) ne plus havar definitive sua havaji importoza}\lingui{DEFIRS}
\artiklo{ruko}%
\vorto{ruko} Stipo di kano, bastono di qua la extremajo supra garnisesis ye lino, kanabo, kotono, e c. por filifo\lingui{DEIS}
\artiklo{ruktar}%
\vorto{ruktar}\fako{netrans.) (aludante la stomako} Retrosendar bruske suflado de aero, produktante bruiso\lingui{EFIS}
\artiklo{rukular}%
\vorto{rukular}\fako{netrans.) (aludante la kolombi, la turturi} Emisar la murmuro karezanta qua karakterizas lia speco\lingui{F}
\artiklo{rulado}%
\vorto{rulado}\fako{muziko} Vokalizo qua konsistas ye intersucedo rapida e lejera de noti sur la sama silabo\lingui{EF}
\artiklo{rular}%
\vorto{rular}\senco{1}{\textit{(trans.)} Movar per igar rotacar sur su ipsa}\senco{2}{\textit{(netrans.)} Movar su per rotaco sur su ipsa}\senco{3}{\textit{(aludante navo)} Ocilar de dextre a sinistre, ed inverse}\lingui{DEF}
\artiklo{ruleto}%
\vorto{ruleto} Hazardo-ludo en qua mikra bulo de ivoro, per rular sur cirklo dividita aden 76 faketi numeroza (reda o nigra), determinas la gano o la perdo, segun ke ta bulo haltas sur la faketo\lingui{DEFIRS}
\artiklo{rumbo}%
\vorto{rumbo}\fako{navig.} Angulo-spaco qua inter-separas la 32 vent-arei di la busolo\lingui{DEFIS}
\artiklo{ruminar}%
\vorto{ruminar}\fako{trans. e netrans.) (aludante ula mamiferi herbivora} Igar la manjaji mi-mastikita retrovenar de la stomako aden la boko, ube li mastikesas itere ante esar englutata definitive\lingui{EFIS}
\artiklo{rumo}%
\vorto{rumo} Brandio facita ek la reziduo de melaso di kano-sukro\lingui{EFIR}
\artiklo{rumoro}%
\vorto{rumoro} Informo-vorti qui propagesas, difuzesas en publiko (*publico)\lingui{EFIS}
\artiklo{rumpo}%
\vorto{rumpo} Extremajo dopa di la korpo di la ucelo, qua suportas la kaudo\lingui{E}
\artiklo{rumsteko}%
\vorto{rumsteko} La parto maxim supra di la dopajo di bovo, an-sur la kaudo\lingui{DEF}
\artiklo{runo}%
\vorto{runo} Singla de la literi ek la alfabeti maxim antiqua di la populi Germana e Skandinaviana\lingui{DEFIRS}
\artiklo{ruptar}%
\vorto{ruptar}\fako{trans.} Pecigar, pecetigar, per shoko, frapo, preso tre forta. – Anke metafore\lingui{EFIS}
\artiklo{ruro}%
\vorto{ruro} Extensajo de sulo, situita exter la urbi, quan okupas la agri, la prati, e c.\lingui{EFIS}
\artiklo{rusho}%
\vorto{rusho} Bendo de stofo, de tulo, e c., plisita forme di mikra kupeti, di godroni, kelke simila a la alveoli di abeluyo, ed uzata por ornar la vesti hominala\lingui{DEFR}
\artiklo{rusko}%
\vorto{rusko}\fako{bot.} Genero de planti ek la familio « liliazei », arboreti sempre verda, simila ad la asparagi, kun folii rudimenta e rameti aplastita forme di folii verda e harda ledratre di qui la surfaco portas flori qui donas beri globatra\latina{ruscus}\lingui{ISL}
\artiklo{rustika}%
\vorto{rustika} Qua apartenas a la kozi di ruro. Dicesas anke pri la manieri di la rurani – pri la planti qui rezistas la domaji produktita da la mala vetero quale la planti sovaja – pri benko, kozo ek petro kruda od ek ligno ne-laborita\lingui{DEFIS}
\artiklo{rusto}%
\vorto{rusto} Fer-oxido, redatra\simbolo{Fe₂O₃ e Fe₃O₄}\lingui{DE}
\artiklo{ruteno}%
\vorto{ruteno}\fako{kemio} Metalo ek la grupo di platino, la maxim skarsa ek omni, deskovrita da Glaus (en la yaro 1843) en la osmio di iridio\simbolo{Ru}
\artiklo{rutino}%
\vorto{rutino} Kustumo su-inkrustinta agar o facar ulo segun ta o maniero. – \textit{Metaf}. : Kustumo a qua onu su konformas mekanismatre, sen-reflekte\lingui{DEFIRS}
\artiklo{ruto}%
\vorto{ruto}\fako{bot.} Planto medicinala, ek la familio « rutacei » di qui lu esas la tipo, qua saporas akre\latina{ruta}\lingui{DFIRS}
\artiklo{ruzar}%
\vorto{ruzar}\fako{netrans., per} Uzar, por trompar, moyeni qui igas falsajo aspektar quale exaktajo, quale verajo\lingui{DEF}
\end{multicols}
\sekciono{S}
\begin{multicols}{2}
\artiklo{sabato}%
\vorto{sabato}\fako{biblo} Repozo, preskriptita da lia religio, quan la Judi devas praktikar dum la tota dio sepesma di la semano\lingui{DEFIRS}
\artiklo{sabeliko}%
\vorto{sabeliko}\fako{bot.} Varietato di kaulo, di qua la folii esas tre frizita\latina{brassica oleracea Sabauda}\lingui{I}
\artiklo{sablo}%
\vorto{sablo} Substanco pulveratra, produktita da la desagregeso di la roki silicoza o quarcoza\lingui{FIS}
\artiklo{sabordo}%
\vorto{sabordo}\fako{navo} Aperturo quadri-angula, situita en la kareno di navo (specale por lasar pasar la boko di kanono)\lingui{F}
\artiklo{sabro}%
\vorto{sabro} Armo ne-pafala, kun pinto ed aristo tranchiva, la aristo tranchiva di la lamo esante konvexa\lingui{DEFIRS}
\artiklo{sacerdoto}%
\vorto{sacerdoto}\senco{1}{Ta qua prezidas la kulto deala, la ceremonii religiala}\senco{2}{\textit{(religio katol.)} Ta persono qua recevis de la episkopo, kun la ordo di sacerdoteso, la permiso celebrar la meso e donar la uncioni}\lingui{EFIS}
\artiklo{saciar}%
\vorto{saciar}\fako{trans.} Kontentigar per cesigar la bezono\lingui{FIS}
\artiklo{sadismo}%
\vorto{sadismo} Laciveso akompanata da krueleso\lingui{DEFIRS}
\artiklo{safena}%
\vorto{safena}\fako{anat.} Veino di la gambo, qua departas de la pedo-fingri
\artiklo{safiro}%
\vorto{safiro} Lapido briloza, bele blua\lingui{DEFIRS}
\artiklo{safrano}%
\vorto{safrano}\senco{1}{\textit{(bot.)} Planto bulboza di qua la floro, blua kun purpurajo havas stigmati flavatra qui, pulverigita, uzesas kom kolorizivo e kom kondimento}\senco{2}{La ipsa farbo}\latina{safranum}\lingui{DEFIRS}
\artiklo{sagaca}%
\vorto{sagaca} Karakterizata da instinto subtila pri deskovrar la kozi, pri penetrar til la fakti, qua komprenas la relati, mem minima, inter la kozi ipse, qua dicernas multa diferi ube altra personi vidas nur kozi uniforma\lingui{EFIS}
\artiklo{sagitario}%
\vorto{sagitario}\fako{bot.} Planto qua vivas en aquo, herbatra, perena, ula folii di qua esas flecho-forma e vertikala, la ceteri longa e flotacanta\latina{sagittaria}\lingui{FISL}
\artiklo{saguino}%
\vorto{saguino}\fako{zool.} Speco di mikra simio, kun kaudo longa\latina{rapale}\lingui{DEFIRS}
\artiklo{saguto}%
\vorto{saguto}\fako{bot.} Palmiero di qua la medulo donas fekulo amilatra, e la frukti donas, per distilo, liquoro ebriigiva\lingui{DEFIRS}
\artiklo{saimo}%
\vorto{saimo} Graso de la epiplono di la porko, fuzita e salizita, quan onu uzas koquarte, por fritar, e c.\lingui{I}
\artiklo{saja}%
\vorto{saja} Qua judikas yuste pri la kozi, dicernas lo vera de lo falsa, lo exakta de lo semblanta (segun ke ta fakultato esas komuna ye omna homi), agas segun raciono e modere – e di qua la konduto esas delikate regulkonforma\lingui{EFIS}
\artiklo{sakarimetro}%
\vorto{sakarimetro}\fako{kemio} Aparato qua mezuras, per la deviaco di la radio polarigita, la quanto de sukro quan kontenas liquido\lingui{DEFIS}
\artiklo{sakarino}%
\vorto{sakarino}\fako{kemio} Laktono fuzebla ye 100 gradi C., qua saporas bitre, obtenita per la efiko di kalko sur glukoso o levuloso\simbolo{C₆H₄ / COSO₂/ NH}\lingui{DEFIS}
\artiklo{sako}%
\vorto{sako}\senco{1}{Peco de ledro, de stofo faldita mediane; klozita (per suturo o per glutino) laterale; kun aperturo supre por posibligar la enpozo di objekti}\senco{2}{\textit{(fiziol.)} Kavajo en la korpo, cirkondata da parieto membrana}\lingui{DEFIS}
\artiklo{sakra}%
\vorto{sakra}\fako{religio}\senco{1}{Di qua la karaktero esas neviolacenda pro ke konsakrita a la oficio deala}\senco{2}{Egardata kom neviolacenda}\lingui{EFIS}
\artiklo{sakramento}%
\vorto{sakramento}\fako{religio katolika} Singla de la agi, institucita da Iesu Kristo, por la santigo di la anmi, signo sentebla di efekto spiritala, destinita naskigar vivo de gracio (bapto, penitenco) od augmentar lu (konfirmo, eukaristio, mariajo, ordo, unciono extrema)\lingui{DEFIRS}
\artiklo{sakrifikar}%
\vorto{sakrifikar}\senco{1}{\textit{(trans.)} \textit{(religio)}}\senco{2}{Kombustar (kozo quan onu egardas kom tre precoza) sur altaro konsakrita a deo, por honorizar o flexar lu}\senco{3}{Ofrar vivanto quan onu imolas por la deo}\senco{4}{Perdar volunte persono o kozo quan onu amas o prizas rispektive, po (ulu, ulo)}\lingui{EFIS}
\artiklo{sakrilejar}%
\vorto{sakrilejar}\fako{trans.) (religio} Violacar ulo sakra\lingui{DEFIS}
\artiklo{sakristo}%
\vorto{sakristo}\fako{katolikismo} La persono qua su okupas ye ta parto acesora di kirko ube esas depozita la vazi e la ornivi sakra, ed ube la oficionti metas sua vesti sacerdotala; e di ta parto di kirko, ube la jus-mariajiti e lia testi venas signatar sur la registri di la parokio, e gratulesar da la amiki qui asistis la mariajo\lingui{DEFIS}
\artiklo{sakrumo}%
\vorto{sakrumo}\fako{anat.} Osto di la pelvo, quan la antiqui ofris sakrifike\lingui{EFIS}
\artiklo{salado}%
\vorto{salado} Disho ek herbi o legumi kun (kom spici o kondimenti) salo, pipro, oleo e vinagro\lingui{DEFIRS}
\artiklo{salamandro}%
\vorto{salamandro}\fako{zool.} Batrako di qua onu judikis la morduro kom venenoza, ed a qua onu atribuis la proprajo povar vivar en fairo\latina{lacerta salamandra}\lingui{DEFIRS}
\artiklo{salariar}%
\vorto{salariar}\fako{trans., quante de} Pagar (ulo ad ulu) po lua laboro por la paganto\lingui{DEFIS}
\artiklo{saldar}%
\vorto{saldar}\fako{trans.) (komerco} Pagar la lasta (*ultima) parto di to quo esas ankore pagenda ye konto\lingui{DFI}
\artiklo{salepo}%
\vorto{salepo} Substanco nutriva, quan onu extraktas de la tuberkuli di speci diversa de orkidi\lingui{DEFIRS}
\artiklo{saliar}%
\vorto{saliar}\fako{netrans., de, ek} Esar plu alta kam la cirkondajo\lingui{EFIS}
\artiklo{salicino}%
\vorto{salicino}\fako{kemio} Glukosido C₆H₁₁O₅OC₆H₄ . CH₂OH, fuzebla ye 201° C., kontenata en la kortico di saliko
\artiklo{salika}%
\vorto{salika}\fako{historio} Legaro di la olima Franki Salika, qua preskriptis ke la heredo esas legitima exkluzive da e tra la maskuli (pri la trono)\lingui{DEFIRS}
\artiklo{salikario}%
\vorto{salikario}\fako{bot.} Planto di qua la flori esas reda, dispozita samanivele cirkum axo komuna, e kelke astriktiva\latina{lythrum salicaria}\lingui{F}
\artiklo{saliko}%
\vorto{saliko}\fako{bot.} Arboro qua kreskas, ordinare, an la litoro di la rivereti\latina{salix}\lingui{IL}
\artiklo{salino}%
\vorto{salino} Loko ube onu obtenas salo per la vaporesko di la mar-aquo\lingui{DEFIS}
\artiklo{salivo}%
\vorto{salivo}\fako{fiziol.} Liquido quan sekrecas ula glandi qui humidigas la boko, komencas la elaboro di la nutrivi manjata, e faciligas la digesto\lingui{EFIS}
\artiklo{salmono}%
\vorto{salmono}\fako{zool.} Mar-fisho qua retroiras la fluvio-fluo, e di qua la karno, redatra e lameloza, tre prizesas. (salar)\latina{salmo}\lingui{DEFIS}
\artiklo{salo}%
\vorto{salo}\senco{1}{Na Cl, substanco friabla, dissolvebla en aquo, qua saporas pikante, uzata kom kondimento, por konservar la nutrivi, por impedar lia putro}\senco{2}{\textit{(kemio)} Kombinajo de acido kun oxido}\lingui{DEFIRS}
\artiklo{salolo}%
\vorto{salolo}\fako{kemio} Salicilato di fenolo
\artiklo{salono}%
\vorto{salono} Chambro di apartamento, ordinare tre spacoza e kelke-luxoze mobloza, destinita por aceptar la vizitanti\lingui{DEFRS}
\artiklo{salpetro}%
\vorto{salpetro}\fako{kemio} Nomo vulgara di kalio-nitrato, qua formacesas sur la muri olda, sur la eskombri de gipso, e c.\lingui{DEFR}
\artiklo{salpo}%
\vorto{salpo}\fako{zool.} Animalo flotacanta, ek la brancho « prokondei », renkontrebla nur en maro diafana
\artiklo{salsifio}%
\vorto{salsifio}\fako{bot.} Planto di qua onu manjas la radiki kom legumo\latina{tragopogon}\lingui{EF}
\artiklo{saltar}%
\vorto{saltar}\fako{netrans.} Lansar su aden atmosfero, esforcanta retro-falar a la departo-punto, o por parkurar ula quanto de spaco\lingui{EFIS}
\artiklo{salubra}%
\vorto{salubra}\fako{pri aero, klimato, nutrivo} Qua efikas favoroze a la organismo homala\lingui{EFIS}
\artiklo{salutar}%
\vorto{salutar}\fako{trans.}\senco{1}{Uzar ta formulo klameta per qua onu sendas ad ulu sua deziro pri lua prosperor}\senco{2}{Honorizar, segun la reguli di politeso, (ulu) per parolo o gesto konvencionala, kande onu renkontras lu, od abordas lu, o departas de lu}\lingui{DEFIRS}
\artiklo{salvar}%
\vorto{salvar}\fako{trans., de, ek}\senco{1}{Igar eskapar mortor o ruinesor}\senco{2}{Igar eskapar destruktesor, o la raptemeso, la avideso di ulu}\lingui{DEFIS}
\artiklo{salvear}%
\vorto{salvear}\fako{netrans.} Honorizar ulu per la deskuplo di pafoserio\lingui{DEFIRS}
\artiklo{salvio}%
\vorto{salvio}\fako{bot.} Planto aromatoza, ek la familio « labiei »\latina{salvia}
\artiklo{sama}%
\vorto{sama}\senco{1}{Di qua la naturo esas absolute olta di altra kozo}\senco{2}{Di qua la naturo esas absolute ne-altreskinta, ne altrigita}\lingui{DER}
\artiklo{sambuko}%
\vorto{sambuko}\fako{bot.} Arbusto di qua la ligno, tre lejera, kontenas medulo-tubo tre longa, e di qua la floro, tre odorosa, uzesas por facar infuzuri kom sudorifigivi\latina{sambucua}\lingui{IL}
\artiklo{samovaro}%
\vorto{samovaro} Sorto di bolio-urno, uzata en Rusia, qua posible igas disponar irga-instante aquo bolianta\lingui{DEFR}
\artiklo{sana}%
\vorto{sana}\senco{1}{Di qua la organismo standas bone, morb-imuna}\senco{2}{Di qua la fakultati pensala, mentala, standas bone. e}\lingui{EFIS}
\artiklo{sanatorio}%
\vorto{sanatorio} Kuracerio, precipue higienala, por la maladi e la personi fatigita kronike, inter qui la tuberklosiki\lingui{DEFIRS}
\artiklo{sancionar}%
\vorto{sancionar}\fako{trans.} Igar (lego) aplikenda per la aprobo da la suvereno, da la chefo di la povo exekutigala. – Konfirmar (lego), igar (lu) aplikata per minaco ye punisi, o per promiso di rekompenso\lingui{DEFIRS}
\artiklo{sandalo}%
\vorto{sandalo} Pedo-vesto ek nura suolo, retenata da kordoni qui ligesas an la supra parto di la pedo\lingui{DEFIRS}
\artiklo{sandarako}%
\vorto{sandarako} Gumo rezinoza, pulvera, per qua onu frotas la paper-folio quan onu skrapis, por impedar lu « drinkar » (absorbar) la inko skribala\lingui{DEFIS}
\artiklo{sandro}%
\vorto{sandro}\fako{zool.} Genero de fishi teleostea, familio di la « perkidei », kelke simila a la perki\latina{sander lucioperca}
\artiklo{sandwich}%
\vorto{« sandwich »} Loncho dina de shinko, e c., inter du pano-lonchi dina, di qua la facio qua kontaktas kun la shinko esas butrizita\lingui{E}
\artiklo{sango}%
\vorto{sango}\fako{fiziol.} Liquido qua cirkulas, per la impulsi de kordio, e per arterii e veini, en la parti diversa di la korpo, por durigar la vivo di ica\lingui{EFIS}
\artiklo{sanguina}%
\vorto{sanguina}\fako{pri temperamento} En qua preponderas sango\lingui{DEFIRS}
\artiklo{sanguisugo}%
\vorto{sanguisugo}\senco{1}{\textit{(zool.)} Anelido qua sugas la sango di la animali, quan onu uzas medicinale por la sango-tiro de la vaskuli kapilara}\senco{2}{\textit{(metaf.)} (1) Persono qua richigas su, detrimente de sua kunhomi, per exhaustar li. (2) To quo exhaustas}\latina{hirudo medicinalis}\lingui{FIS}
\artiklo{sanio}%
\vorto{sanio}\fako{medic.} Materio pusoza, fetida, qua ekfluas de ula ulceri\lingui{DEFIS}
\artiklo{sanskrito}%
\vorto{sanskrito} Linguo sakra di la bramani, di qua la formo maxim arkaika esas renkontrebla en le \textit{Veda} (2000 yari ante I. K.)\lingui{DEFIRS}
\artiklo{santa}%
\vorto{santa}\senco{1}{\textit{(aludante Deo od irgo quo konsakresis a Deo)} Pura de irgequala makulo, de sordidajo, de desnoblajo, e c.}\senco{2}{Dicesas pri ula kreiti, elevita til pureso supernatura : e pri la personi quin la eklezio katolika agnoskas kom tala}\lingui{DEFIS}
\artiklo{santalo}%
\vorto{santalo}\fako{bot.} Substanco lignatra, kun odoro aromatoza\lingui{DEFR}
\artiklo{santolino}%
\vorto{santolino}\fako{bot.} Planto odoranta, ek la familio « kompozaji », uzata kom kontre-vermo\latina{santolina}\lingui{FISL}
\artiklo{santuario}%
\vorto{santuario}\fako{religio}\senco{1}{La loko maxim santa en templo, en kirko, interdiktita a la profani : (1) che la Judi, la parto sekreta di la templo ube gardesis la archo di kontrato; (2) che la katoliki : parto di la kirko, cirkondata, maxim-multa-kaze, per balustrado, ube esas la chefa altaro}\senco{2}{Loko quan onu devas respektar partikulare}\lingui{DEFIS}
\artiklo{sapajuo}%
\vorto{sapajuo}\fako{zool.} Simio di Amerika, di qua la kaudo esas preniva\latina{cebus}\lingui{DEFIR}
\artiklo{sapar}%
\vorto{sapar}\fako{trans.} Exkavar sub edifico, konstrukturo, por igar lu krular. – \textit{(anke metaf.)}\lingui{DEFIRS}
\artiklo{saponario}%
\vorto{saponario}\fako{bot.} Planto « kariofilea », sudorifigiva, di qua onu aquo-boliigas la folii por netigar la lanaji, e c.\latina{saponaria}\lingui{FISL}
\artiklo{saponino}%
\vorto{saponino} Glukosido C₃₂H₅₄O₁₈, kontenata en saponario precipue, e di qua la solvuro spumifas quale aquo saponoza
\artiklo{saponito}%
\vorto{saponito}\fako{mineral.} Silikato naturala amorfa, blankatra, vicina ye talko e magnezito, ma qua kontenas alumino
\artiklo{sapono}%
\vorto{sapono}\fako{kemio} Kombinuro ek alkalio kun la acido oleinala, margarinala, di graso, quan onu uzas por linjo-lavar, netigar, e c.\lingui{EFIS}
\artiklo{saporar}%
\vorto{saporar}\fako{netrans.) (aludante manjajo o drinkajo} Efikar, per sua propraji kemiala, agreable o desagreable a la papili di la lango lor mastiko o drinko\lingui{FIS}
\artiklo{saprofito}%
\vorto{saprofito}\fako{mikrobiol.} Nomo di la mikrobi qui, normale, ne povas developar su en organismo vivanta e, konseque, nutras su de materii mortinta
\artiklo{sapto}%
\vorto{sapto}\fako{bot.} Liquido ek la suki quin la planto cherpas ek tero per sua radiki, e qua, absorbite da sua parti diversa, nutras lu\lingui{DE}
\artiklo{sarabando}%
\vorto{sarabando} Olima danso tri-tempa, di qua la ritmo esis grava e lenta\lingui{DEFIRS}
\artiklo{saraceno}%
\vorto{saraceno}\fako{bot.} Genero de poligonacei di qui la flori, grapa, ne havas korolo. La farino ek la grani uzesas por facar galeti, paplo, krespi, e la grani, por nutrar la kavali, la pultro, e c.\latina{fagopyrus}\lingui{FIS}
\artiklo{sarcelo}%
\vorto{sarcelo}\fako{zool.} Aquo-ucelo analoga a la anado, ma plu mikra\latina{anas querquedula}\lingui{F}
\artiklo{sardino}%
\vorto{sardino}\fako{zool.} Mikra fisho quan onu manjas fresha, o konservita en aquo saloza, o macerita en oleo\latina{clupea sardina}\lingui{DEFIRS}
\artiklo{sardonika}%
\vorto{sardonika} Qua donas a la boko karaktero di moko malignala\lingui{DEFIRS}
\artiklo{sargaso}%
\vorto{sargaso}\fako{bot.} Algo tropikala\latina{sargassum}
\artiklo{sarigo}%
\vorto{sarigo}\fako{zool.} Mamifero ek la kategorio « marsupiali », di qua la femino havas posho ventrala ube lua yuni finas developar su\latina{didelphys virginiana}\lingui{FI}
\artiklo{sarkasmo}%
\vorto{sarkasmo} Moko, ironio mordanta\lingui{DEFIRS}
\artiklo{sarkino}%
\vorto{sarkino}\fako{mikrobiol.} Bakterio saprofita quan onu renkontras en la gangreno pulmonala ed ula morbi di la stomako
\artiklo{sarklar}%
\vorto{sarklar}\fako{trans.} Netigar (sulo-peco) per extirpar la herbi mala\lingui{FI}
\artiklo{sarko}%
\vorto{sarko}\senco{1}{Che la Greki e Romani antiqua : sorto de kofro senkovrita, en qua la mortinto portesis a la rogo od a la tombo}\senco{2}{Che la moderni : kofro de ligno, de plombo, e c., en qua onu enklozas la mortinti por sepultar li}\lingui{DE}
\artiklo{sarkofago}%
\vorto{sarkofago}\senco{1}{Sarko de petro en qua la antiqui depozis la korpi mortinta}\senco{2}{Parto di monumento funerala qua simulas petro sarko qua kontenas la kadavro}\lingui{DEFIRS}
\artiklo{sarkomo}%
\vorto{sarkomo}\fako{patol.} Surkreskajo de karno\lingui{DEFIS}
\artiklo{sarkoplasmo}%
\vorto{sarkoplasmo}\fako{histol.} Protoplasmo granuloza di la fibro muskulala, qua cirkondas omna nukleo, e su extensas, per traco-linei plu o min dina e longa, inter la grupi quin formacas la fibreti muskulala
\artiklo{sarmento}%
\vorto{sarmento}\fako{bot.} Ligno quan la vito produktas singlayare\lingui{EFIS}
\artiklo{sarsaparelo}%
\vorto{sarsaparelo}\fako{bot.} Planto di qua la radiko esas purigiva medicinale\latina{smilax}\lingui{DEFIRS}
\artiklo{sat(e)}%
\vorto{sat(e)} Tam (multa) kam bezonata\lingui{L}
\artiklo{satano}%
\vorto{satano}\fako{kristanismo} La nomo di la diablo\lingui{DEFIRS}
\artiklo{satelito}%
\vorto{satelito}\senco{1}{Viro armoza, dop chefo, por exekutar lua imperi. – \textit{(anke metaf.)}}\senco{2}{\textit{(astron.)} Planeto acesora qua jiras cirkum planeto e sequas ica en lua rotaco}\lingui{DEFIRS}
\artiklo{satino}%
\vorto{satino} Stofo ek silko, plana, en qua la wefto ne esas videbla sur la averso, e qua esas briligita per la preso da cilindro\lingui{DEFIR}
\artiklo{satiro}%
\vorto{satiro}\senco{1}{\textit{(literaturo Latina)} Verko ek prozo e versi, en qua la autoro blamegas la mori di publiko \textit{(*publico)}}\senco{2}{Verko ek versi ube la poeto atakas la vicii e la mokindaji di lua samepokani}\lingui{DEFIRS}
\artiklo{satiruso}%
\vorto{satiruso}\senco{1}{\textit{(mitol.)} Mi-deo qua habitas la boski, kun korpo longa pila, korni, gambi e pedi kaprulala}\senco{2}{\textit{(metaf.)} Viro laciva debochema}\lingui{DEFIRS}
\artiklo{satisfacar}%
\vorto{satisfacar}\fako{trans., per}\senco{1}{Igar (ulu) en stando agreabla per agar o facar por lu to quon lu expektas o deziras}\senco{2}{\textit{(ulo)} Realigar to quon postulas ulo}\lingui{DEFIS}
\artiklo{satisfecit}%
\vorto{« satisfecit »}\fako{vorto skolanala} Paperfolio per qua docisto atestas satisfacesir da la agi o faci da lua docario
\artiklo{satrapo}%
\vorto{satrapo}\senco{1}{\textit{(epoki antiqua)} Guvernisto qua darfis agar suverene en provinco ye la nomo di la rejulo di Persia}\senco{2}{\textit{(nia-epoke)} Persono qua, en sua ofico-resortiso, agas despote}\lingui{DEFIRS}
\artiklo{saturar}%
\vorto{saturar}\fako{trans.}\senco{1}{\textit{(ye)} Adjuntar mixe o kombine (korpo ad altra) tam granda-proporcione kam lu povas recevar de la adjuntajo}\senco{2}{\textit{(metaf.)} Satisfacar plu kam necesa la hungro od altra deziro, bezono}\lingui{DEFIS}
\artiklo{saturdio}%
\vorto{saturdio} Dio sepesma di la semano, quan la Judi konsakras a repozo\lingui{E}
\artiklo{saturnali}%
\vorto{saturnali}\fako{che la Romani, antique} Festo honore di Saturno, dum qua la sklavi, qui lore *weris la sama vesti kam sua mastri, sidis kun ici a la festino-tablo. – \textit{(nia-epoke)}. Periodi de tro granda libereso sexuala\lingui{DEFIRS}
\artiklo{sauco}%
\vorto{sauco} Liquido qua akompanas ula dishi, qua ordinare konsistas ye salo, spici, kondimenti\lingui{DEFIRS}
\artiklo{saurio}%
\vorto{saurio}\fako{zool.} Singla de la repteri ek la kategorio qua kontenas la lacerto, la krokodilo, la kameleono, e c.\lingui{DEFIS}
\artiklo{savano}%
\vorto{savano}\fako{geogr.} Granda sulo-planajo di Amerika ube kreskas la herbo qua manjesas da la bruti\lingui{DEFIS}
\artiklo{savar}%
\vorto{savar}\fako{trans.}\senco{1}{Posedar (ulo) en sua stoko de memoraji}\senco{2}{Posedar la cienco, la arto, la praktiko di ulo (kom exemplo : savar linguo, muziko, desegnar, e c.)}\lingui{FIS}
\artiklo{savorio}%
\vorto{savorio}\fako{bot.} Planto aromatoza, uzata kom kondimento\latina{satureia}
\artiklo{savurar}%
\vorto{savurar}\fako{trans.} Lor manjo, drinko, mastikar od englutar lente por igar plu duroza la juo. – \textit{(metaf.)} Juar lente de sua plezuro, de sua feliceso\lingui{EF}
\artiklo{saxhorno}%
\vorto{saxhorno}\fako{muziko} Instrumento di suflo-muziko, ek latuno, inventita da Adolf Sax en la yaro 1845\lingui{DEF}
\artiklo{saxifrago}%
\vorto{saxifrago}\fako{bot.} Planto herbacea, di qua la flori dispozesas grape o panikule, qua kreskas precipue inter la stoni, en la fenduri di la roki\latina{saxifraga}\lingui{EFISL}
\artiklo{saxofono}%
\vorto{saxofono}\fako{muziko} Instrumento latuna, kun klavi e boko-peco qua havas la formo di klarineto-beko\lingui{DEFIS}
\artiklo{sbiro}%
\vorto{sbiro}\fako{en la Italia olima} Agento di la polico\lingui{EFI}
\artiklo{se}%
\vorto{se} Kaze di, kaze ke – egardante la kazo eventuala (1) kom posibla, (2) kom nura supozo, (3) kom realigita\lingui{FIS}
\artiklo{*seancar}%
\vorto{*seancar}\fako{netrans.} Esar asemblita e permanar tale til normale la traktiteso di omna punti, di omna *\textit{itemi} di la deliberoprogramo, \textit{(aludante la membri di asemblitaro qui deliberas, kun-laboras)}\lingui{EF}
\artiklo{sebo}%
\vorto{sebo} Grasajo fuzita di la animali ruminera (bovi, mutoni, e c. de qua onu fabrikas kandeli e bujii stearinala)\lingui{EIS}
\artiklo{seciono}%
\vorto{seciono}\fako{geom.} Lineo, surfaco, segun qua su partigas du surfaci, du solidi, un solido ed un surfaco.\lingui{DEFIRS}\subvorto{}{secioni konala}{Secioni plana di kono rekta kun bazo cirkla (cirklo, elipso, parabolo, hiperbolo)}{}{}
\artiklo{sedentaria}%
\vorto{sedentaria}\senco{1}{Qua sidas maxim-multa-tempe}\senco{2}{Qua chanjas nur rare sua rezideyo}\lingui{DEFIS}
\artiklo{sediciar}%
\vorto{sediciar}\fako{netrans.} Rezistar publike la autoritatozi establisita\lingui{EFIS}
\artiklo{sedimento}%
\vorto{sedimento}\senco{1}{Depozajo di la materii solida qua esis suspendita en liquido}\senco{2}{\textit{(geol.)} Strato quan la aqui depozis lor sua departo}\lingui{DEFIS}
\artiklo{sedo}%
\vorto{sedo}\fako{anat.} La parto karnoza qua esas infra-rene, che la homo, – dope che la animali
\artiklo{seduktar}%
\vorto{seduktar}\fako{trans., per} Venigar (ulu) ad su per charmar (lu)\lingui{EFIS}
\artiklo{segar}%
\vorto{segar}\fako{trans.} Tranchar, o nur sekar, per instrumento qua konsistas ek lamo metala, dina, kun aristo dentoza, muntita an mancho, e qua, per movo alternanta o rotacala, entamas, dividas pokope la materii harda (ligno, quadro, petro, e c.)\lingui{DI}
\artiklo{seglo}%
\vorto{seglo}\fako{navig.} Peco de telo forta, qua konsistas (ordinare) ye plura peci asemblita, quan onu ligas an la yardi di la masti, e qua, desfaldite ed orientizite segun la vigoro e la direciono di vento, recevas de olca impulso qua igas la navo avancar ta-direcione o ca-direcione\lingui{DE}
\artiklo{segmento}%
\vorto{segmento} Parto determinita di objekto longa, lineatra, quan onu dividas, longesale : segmento di voyo, di fervoyo (inter du stacioni), di rivero, di kanalo (inter du sluzi), di filo telegrafala o telefonala (inter du stacioni o du posteni)\lingui{DEFIRS}
\artiklo{segregacar}%
\vorto{segregacar}\fako{trans.} Separar (parto) de la bloko, de la socio. – \textit{(antonimo : agregar)}\lingui{DEFIS}
\artiklo{segun}%
\vorto{segun}\senco{1}{Egardante (ulo quo dependas de la personi, de la cirkonstanci)}\senco{2}{Fidante su a (la arbitro da, la decido da, la selekto da, la dico da)}\lingui{IS}
\artiklo{seido}%
\vorto{seido} Ta qua devotesas fanatike ad ulu, qua esas pronta agar o facar irgo por obediar lu, mem kriminar\lingui{DFIS}
\artiklo{*seifo}%
\vorto{*seifo} Kofro ek metalo, kun sekureso-seruro, en qua onu klefe inkluzas pekunio o kozi altra (dokumenti, juveli) tre precoza, tre valoroza
\artiklo{sejornar}%
\vorto{sejornar}\fako{netrans., en, che} Esar dum kelka tempo, ma ne multa (en ta o ca loko, che ta o ca persono)\lingui{EF}
\artiklo{sekalo}%
\vorto{sekalo}\fako{bot.} Planto gramina, di qua la grano esas plu mikra e plu bruna kam olta di frumento\latina{secala cereale}\lingui{FIL}
\artiklo{sekar}%
\vorto{sekar}\fako{trans.) (kirurg.} Dividar, plu o min profunde, (korpo) per korpo plu harda qua esas taliita tale ke lu havas aristo tranchiva (skalpelo, e c.)\lingui{EFIS}
\artiklo{sekino}%
\vorto{sekino} Olima moneto-peco ora, en Venezia (Italia), qua aceptesis kom pago-moyeno en omna landi qui esis kontigua a la regiono estala di Mediteraneo\lingui{DEFIRS}
\artiklo{sekondar}%
\vorto{sekondar}\fako{trans.}\senco{1}{Helpar (ulu) per agar segun lua idei, lua plano, lua projeti}\senco{2}{\textit{(en la parlamento di ula landi, precipue Britaniana)} Expresar la aprobo (di propozo, e c.) kom preliminaro necesa pri la duro di la dis-kuto di ta propozo, e/o pri la voto pri olu}\lingui{DEFIR}
\artiklo{sekrecar}%
\vorto{sekrecar}\fako{trans.) (aludante tisui} Produktar e lasar eskapar ula humori\lingui{DEFIS}
\artiklo{sekreta}%
\vorto{sekreta} Celata de la homi cetera, ne-savata de la homi cetera, de publiko \textit{(*publico)}\lingui{DEFIRS}
\artiklo{sekretario}%
\vorto{sekretario}\senco{1}{La persono di qua la tasko, helpe di altru, konsistas ye redaktar – o skribar kom diktario – letri, memoriali, e c.}\senco{2}{La persono qua redaktas la protokoli di asemblitaro}\senco{3}{Oficiro komisita por facar ula akti}\senco{4}{Kontor-employato supra-ranga qua traktas la aferi dum la absenteso di sua chefo}\lingui{DEFIRS}
\artiklo{sekto}%
\vorto{sekto} Ensemblo ek la personi qui profesas doktrino partikulara, religiala, filozofiala, e c.\lingui{DEFIRS}
\artiklo{sektoro}%
\vorto{sektoro}\senco{1}{\textit{(geom.)} Parto de cirklo, qua inkluzesas inter arko di la cirklo-kurvo, e du radii trasita til la extremaji di la arko}\senco{2}{\textit{(fortifik.)} Singla de la parti di klozajo \textit{(*kluzajo)}}\lingui{DEFIRS}
\artiklo{sekulara}%
\vorto{sekulara}\fako{religio} Qua apartenas a la mondo de la laiki.\lingui{DEFIRS}\subvorto{}{klerikaro sekulara}{Qua vivas en la mondo de la laiki. \textit{(antonimo di : klerikaro reguliera, la monaki)}}{}{}
\artiklo{sekundara}%
\vorto{sekundara}\lingui{DEFIS}\subvorto{geol.}{periodo sekundara}{Ube la strati kontenas ne plu roki primitiva, ma nur depozaji ek materii transportita}{1}{}\subvorto{pedagogio}{doco sekundara}{La doco agata en la licei e la kolegii \textit{(kontraste kun la doco primara, agata en la skoli elementala)}}{2}{}\subvorto{parenteso}{kuzi sekundara}{Filii de kuzi primara}{3}{}
\artiklo{sekundo}%
\vorto{sekundo} La sisadekima parto di un minuto de tempo, o di arko (geometriala)\lingui{DEFIRS}
\artiklo{sekura}%
\vorto{sekura}\senco{1}{Qua shirmas kontre danjero}\senco{2}{Qua ne esas danjero, danjeroza}\lingui{EFIS}
\artiklo{selakto}%
\vorto{selakto} Parto seroza qua su separas de lakto kande ica koagulas
\artiklo{selektar}%
\vorto{selektar}\fako{trans.}\senco{1}{Prenar prefere}\senco{2}{Rezolvar por un ek du o plura kozi}\lingui{DEFRS}
\artiklo{selenio}%
\vorto{selenio}\fako{kemio} Korpo simpla, metaloido vicina ye sulfo. Atom-nombro : 34. Atom-pezo : 79,2\simbolo{Se}\lingui{DEFIRS}
\artiklo{selenografio}%
\vorto{selenografio} Deskripto astronomiala di Luno\lingui{DEFIRS}
\artiklo{selo}%
\vorto{selo} Parto di la harneso pozita sur la dorso di kavalo, quan la kavalkanto uzas kom sidilo\lingui{FIS}
\artiklo{semaforo}%
\vorto{semaforo} Aparato telegrafala, instalita vicine ye portuo, por signalar la adveno e la arivo di la navi, lia manovri vide de la litoro di portuo, e por korespondar kun li\lingui{DEFIRS}
\artiklo{semano}%
\vorto{semano} Periodo de sep dii, de sundio til saturdio (inkluzite)\lingui{FIS}
\artiklo{semantiko}%
\vorto{semantiko} Studio di la senco di la vorti e di lia varii – cienco pri la « vivo di la vorti »\lingui{DEFIRS}
\artiklo{semar}%
\vorto{semar}\fako{trans.} Disjetar, disenterigar, distope, en kultivo-sulo preparata (la semini di ta o ca planto), o disjetante, o per semo-mashino. – \textit{(anke metaf.)}\lingui{FIRS}
\artiklo{semblar}%
\vorto{semblar}\fako{netrans.} Aspektar quale se lu esus (ma ne esas reale)\lingui{EFI}
\artiklo{semestro}%
\vorto{semestro} Periodo de sis monati intersequanta\lingui{DEFIRS}
\artiklo{seminario}%
\vorto{seminario}\fako{religio} Domo religiala ube onu preparas, en singla diocezo (katolika) la kleriki aprentisa, qui aspiras a nominesor kleriki\lingui{DEFIRS}
\artiklo{semino}%
\vorto{semino}\fako{bot.} Parto di frukto (grano, kerno, kerneto) quan onu semas por produktar la planto quan lu kontenas jerme\lingui{EFIS}
\artiklo{semlo}%
\vorto{semlo} Pano-bloko mikra, generale bulatra, ek frumento delikata
\artiklo{semolo}%
\vorto{semolo} Gruelo de frumento bakita, e granetigita per pisto o trituro, quan onu uzas por facar supi ed entremesi\lingui{EFIS}
\artiklo{sempervivo}%
\vorto{sempervivo}\fako{bot.} Planto dika-folia, herbatra, ek la familio « kresulacei »\latina{sempervivum tectorum}\lingui{L}
\artiklo{sempre}%
\vorto{sempre}\senco{1}{« En la tempo futura, sen mem nur un ecepto-kazo »}\senco{2}{« En la tempo pasinta, sen mem nur un ecepto-kazo »}\senco{3}{« Ye irga epoki »}\lingui{EFIS}
\artiklo{sen}%
\vorto{sen} Prepoziciono qua indikas la absenteso, la manko, la indijo di persono, di kozo, di la eso-maniero o di la ago-maniero quan lu guvernas\lingui{FIS}
\artiklo{senato}%
\vorto{senato}\senco{1}{\textit{(en Roma, antique)} Unesma (hierarkiale) korpo politikala, ek la patrici}\senco{2}{\textit{(nun)} Unesma (hierarkiale) korpo politikala di ula stati moderna}\lingui{DEFIRS}
\artiklo{senatuskonsulto}%
\vorto{senatuskonsulto}\senco{1}{\textit{(en Roma, antique)} Rezolvo agita da la senato}\senco{2}{\textit{(en Francia)} Rezolvo agita da la senato di Francia lor la konsulo-rejimo e lor la imperio di Napoleono Iesma}\lingui{EF}
\artiklo{senco}%
\vorto{senco} Maniero ye qua kozo esas komprenenda. La senco konsistas ye la idei diversa quin lu expresas od implikas- kontre ke lua \textit{signifiko} konsistas ye la objekti diversa a qui lu esas aplikebla\lingui{EFIS}
\artiklo{sendar}%
\vorto{sendar}\fako{trans., ad} Igar (ulo, ulu) departar a destinario\lingui{DE}
\artiklo{senecio}%
\vorto{senecio}\fako{bot.} Planto ek la familio « kompozaji », di qua la semino uzesas por nutrar la uceli\latina{senecio}\lingui{FL}
\artiklo{seneshalo}%
\vorto{seneshalo}\fako{lor la feudismo-rejimo, en Francia} Oficiro feudala qua chefdirektis la domego, komandis la armeo, judiciis, e c.\lingui{DEFIRS}
\artiklo{senila}%
\vorto{senila}\fako{medic.} Di qua la korpo e la spirito divenis febla pro oldeso\lingui{DEFIS}
\artiklo{seniora}%
\vorto{seniora}\senco{1}{Qua naskis ante sua frato o frati}\lingui{DEF}\subvorto{}{senioro}{Persono plu evoza kam altra}{2}{}
\artiklo{sensaciono}%
\vorto{sensaciono} Intereseso o sentiment-eciteso tre forta, che turbo\lingui{DEFIS}
\artiklo{sensitivo}%
\vorto{sensitivo}\fako{bot.} Varietato di mimozo, di qua la folii su retrofaldas kande onu tushas li\latina{mimosa pudica}\lingui{EF}
\artiklo{senso}%
\vorto{senso} Fakultato (che la homo, che la animalo) perceptar sua impresesi de la objekti materiala (tusho, vido, audo, gusto, flaro)\lingui{DEFIRS}
\artiklo{sensuala}%
\vorto{sensuala} Qua ambicias o diletas la plezuro de la sensi\lingui{EFIS}
\artiklo{sentar}%
\vorto{sentar}\fako{trans.} Impresesar agreable o dolorigante (a) per sua sensi, (b) sen sua sensi\lingui{EFIS}
\artiklo{sentenco}%
\vorto{sentenco} Opiniono, expresata per formulo dogmatala\lingui{DEFIRS}
\artiklo{sentimento}%
\vorto{sentimento} Psiko-movo, psiko-reakto, en qua ne implikesas la sensi\lingui{EFIRS}
\artiklo{sentinelo}%
\vorto{sentinelo}\senco{1}{Soldato qua havas kom tasko surveyar avan posteno armeala (kazerno, fortreso, fortifikajo, kampeyo) por ke olca ne surprizesez}\senco{2}{\textit{(metaf.)} Sentinelo : La persono qua surveyas por evitar surprizo}\lingui{EFIS}
\artiklo{sep}%
\vorto{sep} Sis plus un (6 + 1, o 4 + 3)\lingui{DEFIS}
\artiklo{sepalo}%
\vorto{sepalo}\fako{bot.} Folieto di la kalico di floro\lingui{EFIS}
\artiklo{separar}%
\vorto{separar}\fako{trans., de}\senco{1}{Igar aparta, ici de iti, kozi, personi, qui inter-kunesis (esis asemblita, unionita)}\senco{2}{Esar situita inter kozi, personi, tale ke ici impedesas inter-kunesar}\lingui{DEFIRS}
\artiklo{sepio}%
\vorto{sepio}\fako{zool.} Molusko di qua la korpo esas karnoza e depresita – qua havas sur sua dorso lamelo kalkoza, friabla (sepiosto) e difuzas cirkum su, kande lu timas atakeso, liquido nigratra quan onu uzas por facar farbi (sepio-koloro)\latina{sepia}\lingui{DEFIRS}
\artiklo{septa}%
\vorto{septa}\fako{medic.} Qua putrigas\lingui{DEFIRS}
\artiklo{septembro}%
\vorto{septembro} La monato nonesma di la yaro\lingui{DEFIRS}
\artiklo{septeto}%
\vorto{septeto}\fako{muziko} Muziko koncertala por sep instrumenti o sep voci (kun o sen akompano)\lingui{DEFRS}
\artiklo{septiliono}%
\vorto{septiliono}\fako{matem.} 10⁴²
\artiklo{septimo}%
\vorto{septimo}\senco{1}{\textit{(skermo)} Un ek la fazi di atako; pareo di ta atako}\senco{2}{\textit{(muziko)} Intervalo ek sep noti, de « do » til « si », e qua preiras la oktavo}
\artiklo{septua-}%
\vorto{septua-} Prefixo ciencala : « qua havas sep… »
\artiklo{septuagesimo}%
\vorto{septuagesimo}\fako{liturgio katolika} La dio sepadekesma ante la oktavo di pasko
\artiklo{sepultar}%
\vorto{sepultar}\fako{trans., en} Envelopar (la kadavro di persono) per tuko, ante enterigar lu\lingui{EFIS}
\artiklo{sequar}%
\vorto{sequar}\fako{trans.}\senco{1}{Irar dop (ulu, ulo) – (a) Irar dop (ulu), akompanante; (b) Esar dop o pos (ulo)}\senco{2}{Direktar su konforme ad (ulo, ulu) – (a) Durar, en spaco, segun la direciono quan indikas ulu; \textit{(metaf.)} egardar, pri sua konduto, la precepto, la modereso-grado quan indikis ulu, ulo; (b) durar irar segun la direciono, kondutar segun sua selekto}\lingui{EFIS}
\artiklo{sequencio}%
\vorto{sequencio}\fako{liturgio katol.} Versajo ek versi segun-mezuro e ritmoza qua, en la mesi solena, sequas la graduelo (fragmento di la oficio, inter la epistolo e la prozo, ante Evangelio, quan onu pronuncis olim sur la gradi di la jubeo o di la ambono) e la « aleluya »
\artiklo{sequestrar}%
\vorto{sequestrar}\fako{trans.) (yuro-cienco} Depozar (la kozo di la litijo-punto) en la manui di altra persono qua gardos lu til kande decidesabos pri la persono a qua lu devas apartenar\lingui{DEFIRS}
\artiklo{sequoyo}%
\vorto{sequoyo}\fako{bot.} Genero de koniferi taxodinea, qua kreskas en Kalifornia (Usa) – arboro imensa di qua ula sorto esas alta ye 140 metri\latina{sequoja}\lingui{EF}
\artiklo{serafo}%
\vorto{serafo}\fako{teol.} Anjelo di la hierarkio unesma; spirito cielala, che la Judi e che la kristani\lingui{DEFIS}
\artiklo{serayo}%
\vorto{serayo} Palaco di la sultani, viziri, siniori Turka (en Turkia olima)\lingui{DEFIRS}
\artiklo{serbatano}%
\vorto{serbatano}\senco{1}{Tubo longa, kava, uzata por lansar, per homosuflo, pizi, mikra buli ek tero, e c.}\senco{2}{Tubo de fero, uzata por ensuflar vitro}\lingui{FIS}
\artiklo{serchar}%
\vorto{serchar}\fako{trans.} Probar deskovror, prokuror a su, renkontror (ulu, ulo)\lingui{EFI}
\artiklo{serena}%
\vorto{serena}\fako{aludante cielo, atmosfero, e c., o psiko} Pura e kalma, quieta\lingui{EFIS}
\artiklo{serenado}%
\vorto{serenado} Koncerto, donata vespere sub la fenestri di ulu (precipue en Hispania), olim\lingui{DEFIRS}
\artiklo{serfo}%
\vorto{serfo}\fako{feudismo-rejimo} Homo pri la persono o la havaji di qua la sinioro havis yuri qui variis segun la kustumi\lingui{EFIS}
\artiklo{seringo}%
\vorto{seringo}\fako{bot.} Arboreto tre rametoza, kun folii inter-opozanta, kun flori blanka qui dispozesas bukete, e tre odoras\latina{philadelphus}\lingui{EFIS}
\artiklo{serino}%
\vorto{serino}\fako{zool.} Pasero di qua la plumaro esas citron-flava che plura speci, e di qua la maskulo povas lernar siflo e kantar ula arii\latina{fringilla serinus}\lingui{EF}
\artiklo{serio}%
\vorto{serio}\senco{1}{Termini qui intersequas segun lego \textit{(*leyo)} determinita}\senco{2}{Intersequentaro kontinua}\lingui{DEFIRS}
\artiklo{serioza}%
\vorto{serioza} Qua egardas, che la kozi, lia gravajo, lia importozajo, prefere\lingui{EFIRS}
\artiklo{serjento}%
\vorto{serjento}\senco{1}{Oficiro qua resortisas a la autoritatozi statala, municipala od urbala, e qua havis plura roli pleenda}\senco{2}{\textit{(armeo)} Suboficiro di infantrio o di kavalrio}\lingui{DEFIRS}
\artiklo{serjo}%
\vorto{serjo} Stofo de silko, de filo lina, de lano, di qua la wefto esas min glata e min densa kam la fili longesala (warpa), qui esas apene videbla averse\lingui{DEF}
\artiklo{sero}%
\vorto{sero}\fako{patol.}\senco{1}{Liquido aquatra qua su separas de la parto koagulinta di sango (o di lakto)}\senco{2}{\textit{(fiziol.)} Humoro sekrecita (da plevro, peritoneo, e c.)}\senco{3}{\textit{(patol.)} Humoro ne-normala, en ula morbi}\lingui{DEFIRS}
\artiklo{serono}%
\vorto{serono} Pakego, envelopita per bovo-felo (e qua kontenas vari exotika)\lingui{DEFIS}
\artiklo{serpentino}%
\vorto{serpentino}\fako{mineral.} Silikato di magnezio, kun makuli verda\lingui{DEFIS}
\artiklo{serpento}%
\vorto{serpento}\fako{zool.} Reptero di qua la korpo esas tre oblonga, senmembra; ula speci, venenoza; qua su movas per la konkavaji quin lu facas sur sulo\latina{serpens}\lingui{DEFIS}
\artiklo{serpo}%
\vorto{serpo} Larja lamo krecento-forma, di qua la aristo konkava esas tranchiva, kun mancho kurta, uzata da la gardenisti, hakisti, por emundar la arbori\lingui{F}
\artiklo{serpolo}%
\vorto{serpolo}\fako{bot.} Planto labiea, aromatoza\latina{thymus serpillum}\lingui{EFIS}
\artiklo{serumo}%
\vorto{serumo}\fako{farmacio} Liquido aquatra, qua su separas de la parto koagulinta di la sango di animali diversa, e preparita por uzesor kom injektajo terapiala\lingui{DEFIS}
\artiklo{seruro}%
\vorto{seruro}\fako{tekn.} Aparato qua, esence, konsistas ye buxo fera (« palastro ») quan onu aplikas ad-an pordo, e kovrilo di kofro, e c.- peco movebla interne di ta buxo (« lango ») e peco kava (« boko ») qua recevas la lango e fixigesis ube la klapo di la pordo o di la kofro retrovenas por klozo; fine, klefo, por apertar o klozar la pordo, la kovrilo, e c. per igar la lango enirar od ekirar la boko\lingui{FIS}
\artiklo{servalo}%
\vorto{servalo}\fako{zool.} Tigro-kato di Afrika\latina{felis serva}\lingui{DEFRS}
\artiklo{servar}%
\vorto{servar}\fako{trans.}\senco{1}{Igar su utila ad (ulu) po salario-quanto}\senco{2}{\textit{(*serviar)} Prokurar (ad ulu) avantajo per sua interveno personala}\lingui{DEFIS}
\artiklo{*serviar}%
\vorto{*serviar}\fako{trans.} Prokurar (ad ulu), generale gratuite, avantajo per sua interveno propra e privata
\artiklo{servico}%
\vorto{servico} Ensemblo de la vazi (pladi, pladegi, glasi, tasi, boli, e c.) e manjili (forci, kulieri, kulteli, e c.) o di la linjo-tuki necesa por la repasto-tablo\lingui{DEFIRS}
\artiklo{servitudo}%
\vorto{servitudo}\fako{yuro-cienco} Obligeso, debo, a qua submisesas terproprietajo koncerne altra terproprietajo\lingui{F}
\artiklo{sesgar}%
\vorto{sesgar}\fako{trans.} Entamar per deprenar parto di la bordo\lingui{S}
\artiklo{*sesiono}%
\vorto{*sesiono} Parto di la yaro, dum qua *seancas kontinue deliberantaro, tribunalo nepermananta\lingui{DEFIRS}
\artiklo{setono}%
\vorto{setono}\fako{medic.} Mecho de kotono quan onu pasigas tra la pelo e la celularo por durigar la funciono di la moyeno per qua onu ekfluigas ulo mala e liberigas su de lu\lingui{EFI}
\artiklo{severa}%
\vorto{severa}\fako{ad, pri}\senco{1}{Sen indulgo pri la kulpi}\senco{2}{Qua obedias strikte la regulo}\lingui{EFIS}
\artiklo{seviciar}%
\vorto{seviciar}\fako{ad} \textit{(netrans.)}\textit{(yuro-cienco)} Agar violentoze : spozo, a sua spozo; patro, a sua filio; mastro, a sua servisto
\artiklo{sexa-}%
\vorto{sexa-} Prefixo ciencala : « qua havas sis… »
\artiklo{sexagesimo}%
\vorto{sexagesimo}\fako{liturgio katol.} La dio sisadekesma ante la oktavo di pasko
\artiklo{sextanto}%
\vorto{sextanto}\senco{1}{\textit{(geom.)} La parto sisima di cirklokurvo (arko de 60 gradi)}\senco{2}{Instrumento di qua la parto precipua esas limbo gradizita, longa ye la sisimo di cirklokurvo t.e. 60 gradi, e qua uzesas por mezurar la anguli til 60 gradi por determinar la disti di astri, la situeso di navo sur maro, e c.}\lingui{DEFIRS}
\artiklo{sexteto}%
\vorto{sexteto}\fako{muziko} Muzikajo por sis instrumenti o por sis voci (kun o sen akompano)\lingui{DEFIRS}
\artiklo{sextiliono}%
\vorto{sextiliono}\fako{matem.} 10³⁶
\artiklo{sexto}%
\vorto{sexto}\fako{muziko}\senco{1}{Intervalo qua kontenas sis noti – de « do » a « la »}\senco{2}{\textit{(skermo)} Pareo per la espado vers-supre, en la direciono adextere, kun la karpo vers-ciele}\lingui{DEFIRS}
\artiklo{sexuo}%
\vorto{sexuo}\senco{1}{Karaktero organala naturala, qua interdistingas la maskulo e la femino}\senco{2}{Ensemblo de la karakteri organala, pensala, morala, qui interdistingas la homulo e la homino}\senco{3}{La ensemblo de la individui di ta o ca sexuo}\lingui{DEFIRS}
\artiklo{sezamo}%
\vorto{sezamo}\fako{bot.} Planto oleifera ek la regioni estala\latina{sesamum}
\artiklo{sezono}%
\vorto{sezono} Singla ek la quar dividuri inter-egala di la yaro, determinita segun la quanto de tempo quan spensas Suno por irar de solstico ad equinoxo, e de equinoxo a solstico, e segun la varii temperaturala qui korespondas a li\lingui{DEFR}
\artiklo{sfacelo}%
\vorto{sfacelo}\fako{medic.} Gangreno qua atakas la tota parto dika di membro o di organo plura-tisua
\artiklo{sfenoida}%
\vorto{sfenoida}\fako{anat.} Osto, ye la bazo di la kranio, qua kontributas a la formaco di la orbiti, di la nazo-kavaji
\artiklo{sfero}%
\vorto{sfero}\fako{geom.} Solido qua havas kom limito surfaco kurva di qua omna punti distas egale de punto interna, la « centro »\lingui{DEFIS}
\artiklo{sferoido}%
\vorto{sferoido}\fako{geom.} Solido di qua la formo esas preske sfera\lingui{DEFIS}
\artiklo{sfintero}%
\vorto{sfintero}\fako{anat.} Muskulo cirklokurva qua klozas la orifico di ula kavaji naturala di la homo-korpo (anuso, veziko)\lingui{DEFIS}
\artiklo{sfinxo}%
\vorto{sfinxo}\senco{1}{\textit{(mitol.)} Monstro qua havas kapo e manui quale homo, korpo quale leono, ali quale aglo, e qua devoris la personi qui ne divinis la enigmati quin lu propozis}\senco{2}{\textit{(metaf.)} Persono qua ne lasas sua pensi divinesar}\lingui{DEFIRS}
\artiklo{sh-atra}%
\vorto{sh-atra} \textit{(pronuncar : she-atra.)} – \textit{(fonetiko)}. (litero) qua pronuncesas kun siflo palatala (kom ex. : sh, ch, shch, e c.)
\artiklo{sha}%
\vorto{« sha »} La suvereno di Iran
\artiklo{shablono}%
\vorto{shablono} Modelo segun qua onu facas ula objekti\lingui{DR}
\artiklo{shabrako}%
\vorto{shabrako}\fako{armeo} Kovrilo por la selo di kavaliero, e qua maxim-multa-kaze, garnisita ye muton-felo\lingui{DEFRS}
\artiklo{shafto}%
\vorto{shafto}\senco{1}{Longa mancho ligna sur qua onu muntas ferajo di lanco halbardo, partizano o standardo}\senco{2}{La parto ligna de fusilo di balayilo}\lingui{DE}
\artiklo{shagrino}%
\vorto{shagrino} Ledro ek la felo de asno, mulo, e c., quan onu laboris por igar lu granizita\lingui{DEFI}
\artiklo{shakalo}%
\vorto{shakalo}\fako{zool.} Karnavido de Afrika e de la esto-regioni (oriento) mediate ye la hundo e la foxo, qua chasas nokte, trupe, e devoras la mikra bruti, la kadavri\latina{canis aureus}\lingui{DEFIRS}
\artiklo{shakar}%
\vorto{shakar}\fako{trans.} Agar la stroko per qua onu darfas prenar la rejo (en la ludo en qua la du ludanti igas manovrar, ica kontre ita sur tabeleto dividita ye 64 faketi inter-sama, du serii de peci mikra-statuatra : rejo, damo, turmi, kavali, soldati, oficiri) ed en qua la ludanto qua sucesas shakar la rejo ganas la partio\lingui{DFIRS}
\artiklo{shakto}%
\vorto{shakto} Exkavajo facita por la exploto di mineyo (karbonala, petrola, e c.)\lingui{DE}
\artiklo{shalmo}%
\vorto{shalmo}\senco{1}{\textit{(muziko)} Tubo de kano, de palio, truizita — fluto rustika}\senco{2}{\textit{(tekn.)} Tubo de metalo, de vitro, quan onu uzas por direktar – per ecitar lu per aero-fluo, e c. – la flamo di lampo adsur objekto quan onu volas submisar a kaloro tre granda}\senco{3}{Aspiro-palio uzata por drinkar ek kon-forma glasi, mikra-quantope, liquidi tre kolda}\lingui{DEF}
\artiklo{shalo}%
\vorto{shalo}\senco{1}{Longa stofo-peco quan la Orientani *weras turban ye qua li drapiras sua shultri, e c.}\senco{2}{Granda peco de lanajo, de silkajo, e c., qua havas la formo di quadrato, o di quadrato longa, quan la mulieri uzas por drapirar sua shultri}\lingui{DEFIRS}
\artiklo{shaloto}%
\vorto{shaloto}\fako{bot.} Planto legum-gardenala, analoga ad alio, ma qua saporas min forte\latina{allium ascalonicum}\lingui{DEFRS}
\artiklo{shalupo}%
\vorto{shalupo} Barko segloza, kun remili, plu granda kam olti di la kanoti, quan onu uzas en la portui, en la radi, e quan la granda navi embarkas por sua servo. – (barko segloza o motoroza)\lingui{DEFIRS}
\artiklo{shamar}%
\vorto{shamar}\fako{netrans., pri, pro. de} Sentar su perturbata psike (til vizajo-redesko, en ula kazi) pro konciar sua nedignese o sua devo-neexekuto pro febleso karakterala, o pro moke\lingui{DE}
\artiklo{shamoto}%
\vorto{shamoto}\fako{mineral.} Argilo specala, ek silikato di alumino, qua debas sua karakteri a la maniero di sua molekul-agregeso, e de qua onu facas la kruzeli en qui onu submisas la metali ed altra substanci a la kaloro de furnazi
\artiklo{shampunar}%
\vorto{shampunar}\fako{trans.} Fricionar, por netigar lu, la hararo e la kranio-epidermo per lociono ek alkoholo e saponario, kun parfum-aquo\lingui{DEF}
\artiklo{shancelar}%
\vorto{shancelar}\fako{netrans., pro, de} Vacilar sur sua bazo, marchanto\lingui{F}
\artiklo{shapko}%
\vorto{shapko} Kapo-vesto armeala, uzata (olim) che la Poloniani, e quan *weris la lancieri di Napoleono IIIesma (Francia)
\artiklo{sharado}%
\vorto{sharado} Enigmato qua konsistas ye divinar vorto quan onu povas partigar e di qua singla parto formacas altra vorto, segun defino di la parto e di la toto\lingui{DEFIRS}
\artiklo{sharko}%
\vorto{sharko}\fako{zool.} Grosa fisho marala, tre devorema, ek la genero « squali »\latina{squalus carcharias}\lingui{E}
\artiklo{sharlatano}%
\vorto{sharlatano}\senco{1}{(a) Aventuristo qua iras de urbo ad urbo, frapante sur tamburego por atraktar la turbo, fanfaronante pri sua risanigi miraklatra, operacante e vendante sua drogi sen bubiko, en la ferii e sur la placi. (b) Mediko empirikala di qua la remediili esas, segun lua aserto, nofaliiva e preske panacea}\senco{2}{La persono qua explotas la kredemeso di publiko \textit{(*publico)} – qua esforcadas divenar famoza, populara}\lingui{DEFIRS}
\artiklo{sharpio}%
\vorto{sharpio} Amaso de fili quin onu detiris de linjo par-konsumita, e quan onu uzas por la pansi\lingui{DFR}
\artiklo{sharpo}%
\vorto{sharpo}\senco{1}{Larja bendo de stofo, *werata zone cirkum la korpo, cirkum la torso, de la shultro dextra a la hancho sinistra, o fixigita ye la tayo per nodo, kun insigno di ula ofici statala, municipala, urbala, e c.}\senco{2}{Bandajo, *werata bandoliere, o – subtenate da la kolo – falanta sur la pektoro, por subtenar la avanbrakio malada}\lingui{DFIR}
\artiklo{sheiko}%
\vorto{sheiko} Chefo di tribuo de Arabi\lingui{DEFIRS}
\artiklo{sheklo}%
\vorto{sheklo} Masho di kateno, di qua la formo ne esas olta di la mashi cetera, ed uzesas por skopo specala (fixigo di hoko, di sitelo en dragilo, e c.)\lingui{DE}
\artiklo{shelo}%
\vorto{shelo}\senco{1}{\textit{(zool.)} Envelopilo kalka di la ovo. Envelopilo kalka di ula moluski}\senco{2}{\textit{(bot.)} Envelopilo di ula frukti (nuci, e c.)}\senco{3}{\textit{(bot.)} Envelopilo quan onu deprenas de la frukti, de la legumi, kom preparo por lia konsumeso}\lingui{DER}
\artiklo{shibro}%
\vorto{shibro}\fako{tekn.} Valvo-pordo qua funcionas per glitar tale ke lu impedas (parte o tote) la paso di vaporo en vapor-mashino di aquo en kanalo, tubo, e c.\lingui{D}
\artiklo{shifono}%
\vorto{shifono} Peco de stofo, qua esas (maxim-multa-kaze) krumplita, ed egardata kom reziduo\lingui{EFR}
\artiklo{shikanar}%
\vorto{shikanar}\fako{trans., pri, pro} Eventigar kontesto, disputo (kun sua kunhomo) sen havar fundamento de justeso, yusteso o justifikeso pri lo\lingui{DEF}
\artiklo{shildo}%
\vorto{shildo}\fako{milit.} Plako di qua la dimensioni e la formo varias segun la kazi – maxim-multa-kaze konvexa de la centro radie – uzita da la militisti – dum la epoko antiqua, la mezepoko e che la populi sovaja – por protektar sua korpo de la frapi qui venas de enemiko\lingui{DE}
\artiklo{shindolo}%
\vorto{shindolo} Singla de la planki dina, uzata vice teguli, por tekto-kovrar kirko-klosheyi en ruro, vento-mueleyi, palio-dometi, e c.\lingui{DFS}
\artiklo{shinko}%
\vorto{shinko} Kruro o shultro de porko, de apro, de urso, e c., salagita, e, maxim-multa-kaze, fumur-agita por konserveso\lingui{D}
\artiklo{shirmar}%
\vorto{shirmar}\fako{trans., de} Pozar en loko ube lu protektesas kontre la domaji de vetero\lingui{D}
\artiklo{shokar}%
\vorto{shokar}\fako{trans.}\senco{1}{Renkontrar violentoze (ulu, ulo)}\senco{2}{(Ref. « Adjuntenda », fine di ca verko)}\lingui{DEFS}
\artiklo{shovar}%
\vorto{shovar}\fako{trans.} Glitigar per pulso\lingui{DE}
\artiklo{shovelo}%
\vorto{shovelo} Implemento ek plako de ligno o de fero, muntita ye un ek la extremaji di longa mancho, e qua uzesas por deprenar tero, sablo, karbono, e c.\lingui{DE}
\artiklo{shovina}%
\vorto{shovina} Qualeso di la persono qua admiras sua patrio blinde, sen-reflekte, sen-judike e sen-kompare\lingui{DEF}
\artiklo{shrapnelo}%
\vorto{shrapnelo}\fako{milit.} Obuso plena ye kugli\lingui{DEFR}
\artiklo{shultro}%
\vorto{shultro}\fako{anat.} Parto di la korpo homala, per qua la brakio unionas su an la torso\lingui{DE}
\artiklo{shunto}%
\vorto{shunto}\fako{elektro} Derivo ek cirkuito di elektro-fluo ye maniero tala ke nur parto di la fluo povas pasar en la cirkuito\lingui{DEFIR}
\artiklo{shuo}%
\vorto{shuo} Pedovesto, kun suolo ek ledro, qua kovras nur la pedo (do, ne la maleolo, nek la infra parto di la gambo)\lingui{DE}
\artiklo{shut!}%
\vorto{shut!} Interjeciono uzata por tacigar o silencigar
\artiklo{shutro}%
\vorto{shutro} Panelo de ligno o de fero-tolo, sen-trua, qua, muntite per charniro an la montanto di la framo di fenestro (dextre o sinistre, interne od externe), uzesas por protektar la klapo (dextra o sinistra) di ta fenestro\lingui{E}
\artiklo{si}%
\vorto{« si »}\fako{muziko} La noto sepesma di la « do » gamo
\artiklo{*si}%
\vorto{*si} Adverbo per qua onu afirmas la kontreajo di lo jus dicita. (ex. : « Me esas *certena ke vu ne venos. » – « Si, me venos. » – « No! » – « Si! »)\lingui{F}
\artiklo{sibarito}%
\vorto{sibarito} Persono qua esforcas obtenar komforto maxima por su ipsa, e di qua la sentemeso korpala esas rafinita\lingui{DEFIRS}
\artiklo{sibilo}%
\vorto{sibilo}\fako{en la epoki antiqua} Muliero pri qua onu asertis ke lu recevas de deo la talento dicar to quo esas eventonta\lingui{DEFIRS}
\artiklo{*siejo}%
\vorto{*siejo} Loko ube ulu, ulo, esas establisita- ube eventas lua operaci komercala, industriala, financala\lingui{F}
\artiklo{siena-tero}%
\vorto{siena-tero} Okro bruna, redatra, uzata por facar ula farbi
\artiklo{siestar}%
\vorto{siestar}\fako{netrans.} Dormar pos la repasto dimezala, dejuna\lingui{DEFIS}
\artiklo{sifiliso}%
\vorto{sifiliso}\fako{patol.} Morbo infektiva e kontagiala, transmisebla a la decedonti, e di qua la kauzo esas la mikrobo L. \textit{treponema pallidum} – morbo qua afektas precipue la organi sexuala\latina{treponema pallidum}\lingui{DEFIRS}
\artiklo{siflar}%
\vorto{siflar}\fako{trans. e netrans.} Produktar bruiso akuta per igar aero eskapar per aperturo streta. \textit{(Kuglo sisas)}\lingui{F}
\artiklo{sifono}%
\vorto{sifono}\fako{tekn.} Aparato U-forma, kun la U-branchi longa ne-inter-egale, uzata por transvarsar la liquido di vazo aden qua esas imersita la extremajo di la brancho kurta aden vazo ube abutas la brancho longa – la atmosferepreso esante plu granda sur la facio ube la kolono de liquido esas min alta kande la tubo esas amorcita\lingui{DEFIRS}
\artiklo{sigaro}%
\vorto{sigaro} Mikra rulo-bloko ek tabako-folii, di qua onu acendas un ek la extremaji, e de qua onu aspiras la fumuro per la altra extremajo quan onu pozas en sua boko\lingui{DEFIRS}
\artiklo{sigilario}%
\vorto{sigilario}\fako{bot.} Genero de kriptogamo vaskuloza, sorto di likopodo giganta, arboro ek la plantaro karbon-mineyala, di qua la staturo povas atingar 30 metri\lingui{DEFIRS}
\artiklo{siglar}%
\vorto{siglar}\fako{trans.} Klozar per vaxo markizita (kande lu esis ankore mola pro fuzeso) konkave o reliefe, per apliko di korpo harda\lingui{DIS}
\artiklo{signalar}%
\vorto{signalar}\fako{trans., ad} Indikar per signo ke ulo esas agenda ye instanto determinita\lingui{DEFIS}
\artiklo{signatar}%
\vorto{signatar}\fako{trans.} Skribar sua nomo sub la texto di letro, di akto, e c., kom atesto pri to ke lu ipsa skribis ta texto, o ke lu aprobas olca\lingui{DEFIS}
\artiklo{signifikar}%
\vorto{signifikar}\fako{trans.) (aludante senco} Enumerar la objekti a qui la senco di vorto o vorti esas aplikebla. Ex. : La vorti « la maxim granda de la nombri » havas senco komprenebla e mem klara, ma korespondas a nula objekto; lu signifikas nulo. « Blanka Monto » havas senco aplikebla ad omna blanka monti; ma lua signifiko esas restriktita ad un monto. (Cf. Progreso, marto 1913, pag. 80 e 81, artiklo da Louis Couturat)\lingui{EFIS}
\artiklo{signo}%
\vorto{signo}\senco{1}{Kozo perceptebla qua venigas la ideo di ento o di eso-maniero, konseque de interrelato naturala o konvencionala}\subvorto{}{la signi zodiakala}{La dek-e-du stelari qui formacas zodiako}{2}{}\senco{3}{Gesto per qua onu manifestas to quon onu sentas, pensas, od opinionas}\lingui{EFIS}
\artiklo{sika}%
\vorto{sika} Qua ne plus kontenas en sua korpo o sur sua surfaco partikuli aqua\lingui{EFIRS}
\artiklo{siklo}%
\vorto{siklo} Mancheto ligna kun ferajo mi-cirkla, tranchiva (kun o sen denti), quan onu uzas por mesonar la stipi di frumento, di sekalo, e c.\lingui{DE}
\artiklo{sikomoro}%
\vorto{sikomoro}\fako{bot.} Varietato di figiero di qua la folii pasable similesas olti di morusiero\latina{ficus sycomorus}\lingui{DEFIS}
\artiklo{silabo}%
\vorto{silabo}\fako{fonetiko} Son-uno artikulata, produktata per un voc-emiso, e qua konsistas ye o konsonanti e vokali, o vokalo sola, o diftongo\lingui{DEFIRS}
\artiklo{silencar}%
\vorto{silencar}\fako{netrans.} Nule bruisar\lingui{EFIS}
\artiklo{sileno}%
\vorto{sileno}\fako{bot.} Herbo un-yara, du-yara o perena, ed, en ula kazi, miarboreto di qua la folii interopozesas, e kun flori (multa-kaze) grupigita aden *kimi\latina{silenus}
\artiklo{silepso}%
\vorto{silepso}\fako{retor.} Figuro qua konsistas ye uzar vorto (che la lingui en qui existas genro qua ne relatas unike la sexuozi) sen egardar la genro, nek la nombro gramatikala, di altra vorto kun qua lu relatas – omisante la konkordo gramatikala pro egardar unike la ideo ipsa\lingui{DEFIRS}
\artiklo{silexo}%
\vorto{silexo}\fako{mineral.} Stono tre harda, qua konsistas esence ye silico, e qua produktas cintili se lu shokesas da stalo-peco\lingui{EFIS}
\artiklo{silfo}%
\vorto{silfo}\fako{mitol.} Ento fantastika, feo di aero\lingui{DEFIRS}
\artiklo{silico}%
\vorto{silico}\fako{kemio} Oxido di siliko, qua konstitucas silexo\simbolo{SiO₂}\lingui{EFIS}
\artiklo{siliko}%
\vorto{siliko}\fako{kemio} Korpo simpla, metala, qua, kombinite kun oxo – formacas silico\simbolo{Si}\lingui{DEFIS}
\artiklo{siliquo}%
\vorto{siliquo}\fako{bot.} Frukto longa, bivalva\latina{siliqua}\lingui{EFIS}
\artiklo{silko}%
\vorto{silko} Materio filamentoza, briloza, quan filifas speco di faleno di qua la nomo Latina esas \textit{bombyx}\lingui{ER}
\artiklo{silo}%
\vorto{silo} Exkavajo, facita en sulo, aden qua onu depozas la grani (di frumento, e c.) por konservar li\lingui{DEFIS}
\artiklo{silogismo}%
\vorto{silogismo}\fako{logiko} Rezono formala, qua pruvas ke propoziciono \textit{(konkluzo)} esas exakta, t.e. ke la atributo \textit{(termino majora)} di ta propoziciono konvenas a lua subjekto \textit{(termino minora)} per la termino meza qua, konvenante ad la termino majora, ed a la termino minora, su kombinas kun singla de li por formacar du propozicioni (majoro e minoro) – la « premisi » – di qui la konkluzo esas la konsequo\lingui{FIRS}
\artiklo{silueto}%
\vorto{silueto} Desegnuro en qua streko trasita cirkum la ombro de figuro, de korpo, reprezentas olua profilo, ed en qua to quon cirkondas la konturi esas nigrigita\lingui{DEFIRS}
\artiklo{siluriano}%
\vorto{siluriano}\fako{geol.} Nomo di la strati di tereno quin onu renkontras sur la strati di olda greso reda\lingui{DEFS}
\artiklo{siluro}%
\vorto{siluro}\fako{zool.} Genero de fishi teleostea, tipo di la familio « siluridi », qui vivas en Europa ed en la regioni Aziana di qui la klimato esas moderata – tre granda, kun kapo plata, pelo glata e mola\latina{siluris glanis}\lingui{EFIS}
\artiklo{silvano}%
\vorto{silvano}\fako{mitol.} Deo subalterna, deajo di la boski\lingui{F}
\artiklo{silvio}%
\vorto{silvio}\fako{zool.} Nomo vulgara di grupo de paseri, di qui la plumi esas falvatra, e la kanto agreabla\latina{silvia}
\artiklo{silvo}%
\vorto{silvo} Bosko en qua onu lasis la arbori atingar sua kreskopunto maxima\lingui{IL}
\artiklo{simaro}%
\vorto{simaro}\senco{1}{\textit{(en la epoki olima)} Longa robo di muliero}\senco{2}{Sutano quan ula judiciisti *weras sub sua robo oficala}\lingui{EFI}
\artiklo{simbioso}%
\vorto{simbioso}\fako{biol.} Asociitaro ek du o plura organismi, qui interdiferas, kom speci, quan li profitas samatempe. (Kom ex. : la asociitaro di algo kun fungo por formacar likeno)\lingui{DEFIRS}
\artiklo{simboliko}%
\vorto{simboliko} Cienco qua expozas ed interpretas la simboli qui koncernas ta o ca religio o populo
\artiklo{simbolo}%
\vorto{simbolo}\senco{1}{Objekto perceptebla, egardata kom la signo qua reprezentas kozo qua ne esas perceptebla, per ula analogajo qua frapis la imagino-fakultato}\senco{2}{\textit{(kemio)} Litero vorto-komencala, adoptita por determinar korpo simpla en la formuli kemiala}\senco{3}{\textit{(religio)} Formulo konsakrita, ek la kredendaji di religio}\lingui{DEFIRS}
\artiklo{simetra}%
\vorto{simetra} Karakterizata da la korespondo segun-norma (pri grandeso, figuro, situeso) inter korpo o parti di korpo. (a) \textit{(anat.)} Dispozeso di organi dua, situita, ica dextre, ita sinistre, di la lineo mediana (che la vertebrozi e la artropodi); (b) \textit{(geom.)} (figuro) di qua la elementi reciproka inter-egala esas situita inverse\lingui{DEFIRS}
\artiklo{simfizo}%
\vorto{simfizo}\fako{anat.} Artiko ek du osti, fixa o poke-moviva, per fibro kartilago\lingui{DEFIRS}
\artiklo{simfonio}%
\vorto{simfonio}\fako{muziko} Muzikajo por plura voco-toni od instrumenti koncertala : (a)\textit{(olim)} : Muzikajo orkestrala qua esas la uverturo di opero; (b)(de la yarcento 18-esma) Kompozajo muzikala por orkestro, segun la formulo di sonato, t.e. qua konsistas ye prefaco, adajio, skerco e finalo\lingui{DEFIRS}
\artiklo{simila}%
\vorto{simila} Di qua la eso-maniero diferas per nur tre poke de la eso-maniero di altra persono, di altra kozo\lingui{EFIS}
\artiklo{simio}%
\vorto{simio}\fako{zool.} Mamifero quadrimanua\latina{simia}\lingui{EFI}
\artiklo{simoniar}%
\vorto{simoniar}\fako{netrans.} Komercachar pri la kozi religiala (sakramenti, tituli ekleziala, e c.) quin onu donas od aquiras po pekunio o po avantajo sekulara di altra sorto\lingui{DEFIS}
\artiklo{simpatiar}%
\vorto{simpatiar}\fako{trans.}\senco{1}{Sentar su afina kun altru}\senco{2}{Atraktesar da altru per tendenco instintala}\lingui{DEFIRS}
\artiklo{simpatiko}%
\vorto{simpatiko}\fako{anat.}\senco{1}{Centro nervala aparta de cerebro}\senco{2}{Nervaro ganglionoza ek dua kordono, dextre e sinistre di la spino}
\artiklo{simpla}%
\vorto{simpla} Qua ne konsistas ek parti\lingui{EFIS}
\artiklo{simptomo}%
\vorto{simptomo}\senco{1}{\textit{(patol.)} Fenomeno qua karakterizas la existesko di ta o ca morbo}\senco{2}{\textit{(metaf.)} Signo karakteriziva}\lingui{DEFIRS}
\artiklo{simular}%
\vorto{simular}\fako{trans.}\senco{1}{Donar la semblo kom la realajo}\senco{2}{\textit{(yuro-cienco)} Facar (akto) qua eludas la lego per indiki falsa}\lingui{DEFIS}
\artiklo{simultana}%
\vorto{simultana} Qua valoras egale e kune, solidare. (kom ex. : equacioni simultana). – « Simultana » ne esas sinonimo di « sama-tempa », « sam-epoka », « sam-instanta »
\artiklo{sinagogo}%
\vorto{sinagogo}\senco{1}{\textit{(historio religiala)} Asemblitaro ek la religiani, ekleziani (lego da Mozes)}\senco{2}{Religio di la Judi (opoze a kristanismo)}\senco{3}{Templo religiala di la Judi}\lingui{DEFIRS}
\artiklo{sinapismo}%
\vorto{sinapismo}\fako{farmak.} Revulsivo, ek farino de sinapo-grani, quan onu aplikas sur la pelo\lingui{DEFIS}
\artiklo{sinapo}%
\vorto{sinapo}\fako{bot.} Planto krucifera; ek lua grani facesas mustardo\latina{sinapis}\lingui{FIL}
\artiklo{sinartroso}%
\vorto{sinartroso}\fako{anat.} Artiko fixa ek du osti
\artiklo{sincera}%
\vorto{sincera}\senco{1}{Qua igas konocar to quon lu sentas o pensas exakte, reale}\senco{2}{Qua esas sentata, pensata, exakte, reale}\lingui{EFIS}
\artiklo{sindikato}%
\vorto{sindikato}\senco{1}{Asemblitaro de kapitalieri qui spekulas por eventigar la acenso o la decenso di la borso-valoro di la acioni}\senco{2}{Asemblitaro de salariati, de komercisti, qua komisesis por sorgar la interesti di korporaciono}\senco{3}{Asociitaro de salariati qua esforcas defensar sua interesti, protektar sua kun-sindikatani}\lingui{DEFIRS}
\artiklo{sindiko}%
\vorto{sindiko} La persono qua komisesis por sorgar la interesti komuna di korporaciono, di kreditintaro (lor la falio komercala di la debanti, e c.)\lingui{DEFIRS}
\artiklo{sinekdoko}%
\vorto{sinekdoko}\fako{retor.} Figuro qua konsistas ye indikar kozo per la genero de qua lu esas parto, o per un ek lua parti esencala
\artiklo{sinektika}%
\vorto{sinektika}\fako{algebro} Funciono sinektika interne di konturo : nomo, donita da Cauchy, ad omna funciono qua ne cesas esar finita e kontinua, di qua la derivajo, anke, ne cesas esar finita e kontinua, e di qua la valoro esas sempre unika pri omna valori di la varieblo qua korespondas a la punti di la kontur-internajo
\artiklo{sinekuro}%
\vorto{sinekuro} Ofico en qua la oficiero recevas salario-quanto quankam lu ne havas tasko exekutenda\lingui{DEFIS}
\artiklo{sinergido}%
\vorto{sinergido}\fako{biol.} Singla de la energidi, calatere e talatere di oosfero, en la sako embrionala di la angiospermi
\artiklo{singla}%
\vorto{singla} Adjektivo distributala qua indikas ke la persono, la kozo, quan lu determinas esas parto di plurajo kolektiva. – Single : un po un, unope\lingui{E}
\artiklo{singlutar}%
\vorto{singlutar}\fako{netrans.} Subisar ta spasmo di la pektoro qua, kontrakte da doloro, lasas eskapar la voco konvulsatra, forme di son-uni intermitanta\lingui{FI}
\artiklo{singulara}%
\vorto{singulara}\senco{1}{Qua ne esas aplikebla a pluri}\senco{2}{\textit{(gram.)} Qua indikas nur un}\lingui{DEFIS}
\artiklo{sinioro}%
\vorto{sinioro} Titulo honorala quan onu donas a la personegi di la supra rangi en la hierarkio sociala\lingui{DEFIS}
\artiklo{sinistra}%
\vorto{sinistra} Situita sama-latere kam la pinto di la homo-kordio\lingui{IL}
\artiklo{sinkar}%
\vorto{sinkar}\senco{1}{\textit{(trans.)} Igar irar a la fundo od an la fundo}\senco{2}{\textit{(netrans.)} Irar a la fundo, od an la fundo}\lingui{DE}
\artiklo{sinkopar}%
\vorto{sinkopar}\fako{netrans.}\senco{1}{\textit{(patol.)} Subisar ceso subita, instantala, di la kordio-pulsado, qua eventigas falo pro feblesko maxima}\senco{2}{\textit{(gram.)} Depreno di un litero od un silabo en la korpo di vorto}\senco{3}{\textit{(muziko)} Artikular (noto) sur tempo febla, e durigar lu til sur tempo forta}\lingui{DEFIS}
\artiklo{sinkretismo}%
\vorto{sinkretismo} Probo kunfuzar plura doktrini filozofiala, medicinala\lingui{DEFIRS}
\artiklo{sinkrona}%
\vorto{sinkrona}\senco{1}{\textit{(matem.)} Dicesas pri la movi qui eventas en la sama tempo-parto}\senco{2}{\textit{(elektro)} \textit{(pri elektro-mashino)} Di qua la angulo-rapideso sempre egalesas la pulso-rapideso di la elektro-fluo alternanta quan lu recevas o produktas od esas, di olu, multoplo o multoplo-dividanto}\lingui{DEFIS}
\artiklo{sino}%
\vorto{sino}\senco{1}{\textit{(anat.)} Parto di la korpo di muliero, en qua elu facis sua filio}\senco{2}{\textit{(metaf.)} Medio}\lingui{FIS}
\artiklo{sinodika}%
\vorto{sinodika}\fako{astron.} Qua relatas la kunjunto di astri.\lingui{DEFIRS}\subvorto{}{jiro sinodika di Luno}{La tempo-quanto quan lu konsumas por ri-esar en la sama situeso relate Suno.}{}{}\subvorto{}{yaro sinodika}{Qua ri-esigas Tero ye la sama longitudo kam altra planeto}{}{}
\artiklo{sinodo}%
\vorto{sinodo}\fako{religio}\senco{1}{Asemblitaro de membri di la klerikaro}\senco{2}{Asemblitaro de la kleriki di diocezo, konvokita (advokita) da la episkopo}\senco{3}{Asemblitaro da la pastori protestantista}\lingui{DEFIRS}
\artiklo{sinonima}%
\vorto{sinonima} Dicesas pri vorto quan karakterizas lua analogeso generala, pri senco, kun altra vorto, ma diferas de ica per nuanco di signifiko\lingui{DEFIRS}
\artiklo{sinoptika}%
\vorto{sinoptika}\senco{1}{Qua igas videbla, per un regard-uno, omna parti di ensemblo. (Kom ex. La tabelo sinoptika di la \textit{Kurso di la filozofio pozitivista}, da Auguste Comte)}\lingui{DEFIRS}\subvorto{kristanismo}{la sinoptiki}{La evangelii da Metteus, da Markus, da Lukas e da Johannes, qui esas dispozita same en la ensemblo de la naraco}{2}{}
\artiklo{sinovio}%
\vorto{sinovio}\fako{anat.} Humoro di qua la rolo esas igar plu facila la interglito di la artiki moviva\lingui{DEFIRS}
\artiklo{sinso}%
\vorto{sinso} De dextre a sinistre, pasante avan la observanto, od inverse. – (sur la mapi, onu kustumas indikar la sinsi per flechi)
\artiklo{sintaxo}%
\vorto{sintaxo}\fako{gram.} Ensemblo ek la reguli qui regnas la dispozo di la vorti en la frazo, e la konstrukto di la propozicioni.\lingui{DEFIRS}\subvorto{}{sintaxo di koordino}{Qua regnas la konstrukto di la propozicioni inter-juntita.}{}{}\subvorto{}{sintaxo di subordino}{Qua regnas la konstrukto di la propozicioni inter-dependanta}{}{}
\artiklo{sintezar}%
\vorto{sintezar}\fako{trans.} Ri-kompozar la elementi di ensemblo qui esas separita; (a) \textit{(kemio)} Rikonstruktar korpo kompozura, di qua la elementi esas separita per analizo; (b) \textit{(kirurg.)} Riunionar la parti apartigita di vunduro, la fragmenti di osto ruptita; (c) \textit{(logiko)} Komencar per lo generala, de qua onu obtenas lo partikulara per nura dedukto logiko-konforma\lingui{DEFIRS}
\artiklo{sintona}%
\vorto{sintona}\fako{elektro} Apta emisar ocili elektrala, di frequeso intersama
\artiklo{sinui}%
\vorto{sinui}\fako{aludante voyo, rivero, fluvio, e c} Turni e jiri ne-inter-egala\lingui{EFIS}
\artiklo{sinuso}%
\vorto{sinuso}\senco{1}{\textit{(anat.)} Kavajo kavernetatra di ula osti, di ula vaskuli}\senco{2}{\textit{(geom.)} Perpendiklo trasita de un ek la extremaji di arko til la diametro trasita ye la altra extremajo}
\artiklo{sinusoido}%
\vorto{sinusoido}\fako{geom.} Kurvo plana qua havas kom equaciono, koordinati rektangula. : y = sin. x
\artiklo{sioro}%
\vorto{sioro} Titulo quan onu donas, pro politeso, ad omna adulti, sive turnante su a li, sive parolante pri li\lingui{FIS}
\artiklo{sireno}%
\vorto{sireno}\senco{1}{\textit{(mitol.)} Ento mi-muliera, mi-fisha, qua atraktis la mar-voyajanti til an la rifi, per la dolceso di sua kanti. (Existis anke sireni ucela)}\senco{2}{\textit{(metaf.)} Muliero tre seduktiva}\senco{3}{\textit{(tekn.)} Instrumento en qua sprico di vaporo, o di aero kompresita, produktas sono stridanta, qua uzesas kom signalo sur la navi}\lingui{DEFIRS}
\artiklo{sirfo}%
\vorto{sirfo}\fako{zool.} Genero de insekti diptera, tipo di la familio « sirfidei », qui vivas omnaloke sur la Terglobo\latina{syrphus}
\artiklo{siringo}%
\vorto{siringo} Mikra pumpilo kunportebla, uzata por klisteri od altra injekti\lingui{EFIS}
\artiklo{siroko}%
\vorto{siroko}\fako{navig.} Nomo donita, sur Mediteraneo, a vento varmega de sud-westo\lingui{DEFIS}
\artiklo{siropo}%
\vorto{siropo}\senco{1}{\textit{(farmac.)} Liquido viskoza, ek sukro koquita, unionita a ca o ta substanco di qua la rolo esas parfumizar lu per igar lu drinkajo agreabla}\senco{2}{Stando liquida di la konfituri ante kande li divenis kelke solida per koldesko}\senco{3}{Stando liquida di sukro qua esas suko densa ante lua solidesko per koldesko}\lingui{DEFIR}
\artiklo{sis}%
\vorto{sis} La cifro « 6 »\lingui{EF}
\artiklo{sisar}%
\vorto{sisar}\fako{netrans.} Produktar la son-uno \textit{sss} o \textit{fff}, t.e. bruiso ne-muzikala qua esas, fonetike, konsonanto. (Kom exemplo : serpento sisas; lokomotivo siflas per sua siflilo, e sisas pro emisar vaporo ek la kaldiero o la cilindri, tra robineti, ma sen siflar)
\artiklo{sismo}%
\vorto{sismo}\fako{geol.} Movo bruska di la Ter-kortico; sukuso qua shanceligas sulo sur plu o min granda areo, e qua esas apta eventigar deformuri permananta en la regiono koncernata\lingui{DEFIRS}
\artiklo{sismografo}%
\vorto{sismografo} Cienco pri la sismi e pri la movi di sulo
\artiklo{sismometro}%
\vorto{sismometro} Sinonimo di \textit{sismografo}\lingui{DEFIRS}
\artiklo{sismoskopo}%
\vorto{sismoskopo} Sinonimo di sismografo
\artiklo{sistematiko}%
\vorto{sistematiko} Parto di la natur-historio qua koncernas la klasifiko di la enti, e qua inkluzas la deskripto di olci\lingui{DEFIS}
\artiklo{sistemo}%
\vorto{sistemo}\senco{1}{Ensemblo di qua la parti esas koordinita segun lego (aludante enti, institucuri, doktrini, e c.)}\senco{2}{\textit{(anat.)} Ensemblo ek la organi qui esas uzata por ta o ca funciono di la ekonomio che la animalo. (Kom ex. : la sistemo digestala, respirala)}\lingui{DEFIRS}
\artiklo{sistolo}%
\vorto{sistolo}\fako{anat.} Movo per qua la fibri muskulala di la kordio kontraktesas por pulsar sango aden la arterii\lingui{DEFIS}
\artiklo{sistro}%
\vorto{sistro}\fako{antique} Muzik-instrumento di la Egiptiani antiqua, ek mikra ringego de metalo trairata da vergeti qui sonoris kande onu agitis li
\artiklo{sitelo}%
\vorto{sitelo} Vazo cilindra, ek ligno, zinko, e c., kun anso, kapaca de plura litri, qua uzesas por cherpar, transportar, aquo o liquido altra\lingui{I}
\artiklo{situar}%
\vorto{situar}\fako{trans.} Pozar en ta o ca loko (pri expozeso, aspekto, poziciono defensala, e c.)\lingui{DEFIS}
\artiklo{sive}%
\vorto{sive}\fako{gram.} Konjunciono qua indikas alternativo, uzata avan verbo expresita o tacita\lingui{L}
\artiklo{sivo}%
\vorto{sivo}\fako{tekn.} Cirklo-kurvo de ligno sur qua esas tensita telo kun mashi plu o min densa, uzata por kriblagar materii pulveratra (farino, gipso, e c.)\lingui{DE}
\artiklo{sizar}%
\vorto{sizar}\fako{trans.} Prenar, manuo-kaptar, rapide e vigoroze (ulo, ulu)\lingui{EF}
\artiklo{sizigio}%
\vorto{sizigio}\fako{astron.} Situeso di Suno e di Luno, kande ta du astri esas od en kunjunto (« nova Luno »), od en opozeso (« plena Luno »)\lingui{DEFIRS}
\artiklo{skabelo}%
\vorto{skabelo}\senco{1}{Sidilo ek ligno, poke alta, sen brakii nek dors-apogilo}\senco{2}{Quaza eskalereto movebla, kun du o tri gradi, sur qua onu acensas por atingar ulo}\lingui{FIS}
\artiklo{skabino}%
\vorto{skabino} Judiciisto municipala di ula urbi\lingui{FI}
\artiklo{skabio}%
\vorto{skabio}\fako{patol.} Morbo kontagiala di la pelo, quan produktas la prezenteso di la insekto « skaro »\latina{acarus}\lingui{EF}
\artiklo{skabioso}%
\vorto{skabioso}\fako{bot.} Planto herbatra, di qua la flori esas violea, purpurea od (en ula kazi) blanka\latina{scabiosa}\lingui{DEFIRS}
\artiklo{skafandro}%
\vorto{skafandro}\fako{tekn.} Aparato aptigita pri kontenar plunjisto qua, per lu, povas laborar sub la aquo-surfaco, per respirar la aero quan sendas a lu pumpilo situita exter la aquo\lingui{FIS}
\artiklo{skalaro}%
\vorto{skalaro}\fako{konkologio} Genero de moluski karnivora, gasteropoda, prosobrankia, tipo di la familio « skalaridei » – kun rostro retrotirebla, palpili longa, konko turmatra, solida, briloza, kun multa spiri kostoza\lingui{DEFIRS}
\artiklo{skaleno}%
\vorto{skaleno}\fako{geom.} Triangulo di qua la tri lateri esas ne-inter-egala
\artiklo{skalio}%
\vorto{skalio}\fako{zool.}\senco{1}{Krusto qua kovras la korpo di la ostri o di la tortugi}\senco{2}{Kuraso qua kovras la korpo di la tortugi}
\artiklo{skalo}%
\vorto{skalo}\senco{1}{Eskalero portebla, ek du montanti, an qui adjustesis traversi}\senco{2}{\textit{(metaf.)} (a) Serio acensanta o decensanta (kom ex. : la skalo de la enti, la skalo sociala, la skalo diatonika, la skalo de la kolori); (b) Lineo, dividita aden parti (gradi), e c., uzebla por mezurar}\lingui{DEFIRS}
\artiklo{skalpar}%
\vorto{skalpar}\fako{trans.} Deprenar la pelo di la kranio\lingui{DEFIR}
\artiklo{skalpelo}%
\vorto{skalpelo}\fako{kirurg.} Instrumento kun lamo fixa, pintoza, kun un o du aristi, quan onu uzas por dissekar\lingui{DEFIRS}
\artiklo{skamoneo}%
\vorto{skamoneo} Gumo rezina, purgiva, quan onu extraktas de la radiko di speci diversa de konvolvulo mikra\lingui{DEFIS}
\artiklo{skandalar}%
\vorto{skandalar}\fako{trans., per}\senco{1}{Igar sua kunhomi en danjero di etikofalio, di moro-falio, per lia tenteso imitar onua mala exemplo}\senco{2}{Shokar psikale per la prizento regretinda di mala exemplo}\lingui{DEFIRS}
\artiklo{skandar}%
\vorto{skandar}\fako{trans.}\senco{1}{Pronuncar (verso) acentizante forte, o deskompozante, ica pos ita, singla ek la pedi (metrala o silabala) ye qui lu konsistas}\senco{2}{Kantar, plear (frazo muzikala) acentizante forte la tempi di singla mezuro}\lingui{DEFIRS}
\artiklo{skapulario}%
\vorto{skapulario}\fako{religio kristan.}\senco{1}{Parto di la vesto di ula regulieri, qua kovras la robo avane e dope}\senco{2}{Asembluro ek du mikra peci de drapo sakrigita, sur qui esas brodita portreto de la santa Virgino (kelke simila ye la skapulario di la regulieri) quan ula personi *weras sub sua vesti}\lingui{DEFIRS}
\artiklo{skapulo}%
\vorto{skapulo}\fako{anat.} Osto triangula, larja e plata, ye qua konsistas parto dopa di la shultro, e kun qua la osto di la brakio formacas artiko\lingui{EL}
\artiklo{skarabeo}%
\vorto{skarabeo}\senco{1}{\textit{(zool.)} Nomo generala di ula insekti koleoptera, partikulare ula lamelikornieri}\senco{2}{\textit{(geom.)} Kurvo qua pensigas pri skarabeo (insekto) – loko di la pedi di la perpendikli trasita de punto ek la bisekanto di angulo orta til lineo konstante-longa, di qua la extremaji quaze glitas sur la lateri di la angulo}\lingui{EFIRS}
\artiklo{skaramucho}%
\vorto{skaramucho}\fako{teatro} Rolo komika, bufonala, en la olima komedio Italiana, di qua la vesto esis nigra\lingui{DEFI}
\artiklo{skarifikar}%
\vorto{skarifikar}\fako{trans.) (medic.} Incizar surfacale (la pelo) per lanceto, per bisturio, o per aparato ek kupra buxo di qua un ek la edri esas provizita ye fendureti tra qui ekiras, per la preso da kliko-resorto, plura lanceti qui esas muntita sur pivoto komuna\lingui{EFIRS}
\artiklo{skarlatino}%
\vorto{skarlatino}\fako{patol.} Influro di la pelo, kontagiala, karakterizata da makuli skarlata\lingui{DEFIRS}
\artiklo{skarlato}%
\vorto{skarlato} Koloro dazlante reda, obtenita de la kochenilo di nopalo, o de la kochenilo di querko, e qua uzesas por tintar\lingui{DEFIRS}
\artiklo{skarmuchar}%
\vorto{skarmuchar}\fako{trans. e netrans.} \textit{(lor milito)} (aludante detachmento, tiralieri di du armei) Atakar per kombateti\lingui{DEFIS}
\artiklo{skarsa}%
\vorto{skarsa} Qua ne abundas\lingui{EI}
\artiklo{skeleto}%
\vorto{skeleto}\fako{anat.} Osto karpenturo di animalo (vertebroza)\lingui{DEFIRS}
\artiklo{skemo}%
\vorto{skemo}\senco{1}{(a) Trasuro qua reprezentas, per la proporcioni e la relati di ula figuri, la vario-legi di ula sorti de fenomeni – en la domeni di fiziko, mekaniko, statistiko, e c.; (b) Trasuro en qua la planeti diversa esas reprezentita en lia loko ye ta o ca instanto; (c) Trasuro qua reprezentas la dispozeso di organo, di aparato, e c.}\senco{2}{\textit{(yuro ekleziala)} Propozajo, redaktita formale, por submiseso a koncilo}\senco{3}{\textit{(filoz.)} Formo ye qua ni reprezentas, en nia mento, koncepton intelektala}\lingui{DEFIRS}
\artiklo{skeno}%
\vorto{skeno} Asembluro ek fili de kanabo, de silko, de lano, e c., volvita ordinoze por ke li ne inter-intrikez\lingui{E}
\artiklo{skeptika}%
\vorto{skeptika} Qua profesas la dubito filozofiala, olqua konsistas ye egardar raciono kom neapta, pro olua naturo ipsa, konocar ulo certe\lingui{DEFIRS}
\artiklo{skerco}%
\vorto{skerco}\fako{muziko} Muzikajo ne-longa, segun stilo jokema, petulema lejera, lor la pleo di qua oportas, prefere, separar kam ligar la noti\lingui{DEFIR}
\artiklo{skermar}%
\vorto{skermar}\fako{netrans.} Praktikar ta exerco qua kustumigas pri uzar flereto, spado, sabro\lingui{FIS}
\artiklo{sketar}%
\vorto{sketar}\fako{netrans.} Glite kurar sur glacio-surfaco, *werante sorto de sandalo di qua la subajo esas provizita ye fero-lamo vertikala\lingui{E}
\artiklo{skio}%
\vorto{skio} Longa sketilo ek ligno, uzata por glitar sur nivo-surfaco\lingui{DEFIS}
\artiklo{skiro}%
\vorto{skiro}\fako{patol.} Tumoro sen-dolora, qua naskas precipue en la korpo-parti glandatra, ulakaze kancera
\artiklo{skisar}%
\vorto{skisar}\fako{trans.}\senco{1}{Trasar, per krayono, per kreto, e c. la linei precipua di desegnoto, di piktoto}\senco{2}{Facar la modelo vaxa, argilo, e c. di skulturo}\senco{3}{Trasar la « mapo » komencala di verko literaturala o ciencala}\senco{4}{Deskriptar per (lua) traiti precipua}\lingui{DFIRS}
\artiklo{skismo}%
\vorto{skismo}\senco{1}{En religio establisita : formaco di eklezio qua su separas de la eklezio til nun agnoskita, sen disidenteso kompleta pri la punti precipua di la dogmato o di la kulto}\senco{2}{Su-separo en la domeni « politiko », « arti », e c.}\lingui{DEFIRS}
\artiklo{skisto}%
\vorto{skisto}\fako{mineral.} Roko stratoza, exfoliebla ye lameli dina, e di qua la elementi precipua esas silico, argilo, alumino, o metal-oxidi diversa\lingui{EFIS}
\artiklo{sklavo}%
\vorto{sklavo}\senco{1}{La homo qua, en socio, ne esas en la strato de la liberi}\senco{2}{(a) Ta qua su submisas ad ulu, obedias ulu, sen rezisto, pro ambiciar ulo, por obtenar ulo; (b) Ta qua su submisas a la dominaco da muliero, pro amorar elu}\lingui{DEFIS}
\artiklo{sklerotika}%
\vorto{sklerotika}\fako{anat.} Membrano blanka ye qua konsistas parto tre granda di la exterajo di la okul-globo\lingui{DEFIS}
\artiklo{skolastiko}%
\vorto{skolastiko} La filozofio quan onu docis en la skoli di la mezepoko\lingui{DEFIRS}
\artiklo{skoliasto}%
\vorto{skoliasto} Komentero (antiqua od olima) di la texti Greka o Latina\lingui{DEFIRS}
\artiklo{skolio}%
\vorto{skolio}\senco{1}{\textit{(gram.)} Noturo da komentero antiqua od olima}\senco{2}{\textit{(geom.)} Expliko pri teoremo, por igar lua apliko, plu kompleta o limitoza}\lingui{DEFIS}
\artiklo{skolo}%
\vorto{skolo}\senco{1}{(a) Establisuro en qua onu docas a la pueri la rudimenti : lekto, skribo, kalkulo, gramatiko elementala, e c.; (b) Establisuro ube onu docas specalaji}\senco{2}{Ensemblo ek la dicipuli qui studias ta o ca doktrino (filozofiala, medicinala, artala, literaturala) – ed anke ta doktrino ipsa}\senco{3}{To quo esas apta igar posedar la experienco di ulo (la « skolo di vivo », e c.)}\lingui{DEFIRS}
\artiklo{skolopendro}%
\vorto{skolopendro}\fako{zool.} Genero de miriapodo, insekto\latina{scolopendra}\lingui{DEFIS}
\artiklo{skopo}%
\vorto{skopo}\senco{1}{To quon onu esforcas atingar}\senco{2}{To a quo tendencas ago}\senco{3}{To quon intencas kauzo inteligenta}\lingui{EI}
\artiklo{skopso}%
\vorto{skopso}\fako{zool.} Genero de rapt-ucelo noktala, ek la familio « bubonidi »\latina{scops aldrovandi}\lingui{EFL}
\artiklo{skorbuto}%
\vorto{skorbuto}\fako{patol.} Morbo quan karakterizas la sango-povresko, ulceri, makuli hemoragiala, la altereso di la jenjivi, la kario di la denti e di la osti maxilala, abateso, deperiso, produktita (normale) da la manjo tro ofta ed exkluziva, di salizaji, da la privaceso di legumi fresha e vitaminoza\lingui{DEFIRS}
\artiklo{skorio}%
\vorto{skorio} Substanco vitratra qua acensas a la surfaco di la metali fuzanta, e konsistas ye materii altra-sorta, o ye partikuli de metalo\lingui{EFIS}
\artiklo{skorpiono}%
\vorto{skorpiono}\fako{zool.} Animalo ek la klaso « araknidi pulmonoza », di qua la kaudo havas dardo qua komunikas kun glando venenifanta. – \textit{(astron.)} La signo okesma di zodiako\latina{euscorpius italicus}\lingui{DEFIRS}
\artiklo{skorzonero}%
\vorto{skorzonero}\fako{bot.} Genero di planto kompozaja, cikoriacea, quan karakterizas kapituli flava o purpurea, cirkondata da braktei dispozita plura-range, e frukti granda qui havas penacheto (senstipita) de pili plumatra\latina{scorzonera}\lingui{DEFIRS}
\artiklo{skotisho}%
\vorto{skotisho} Danso di qua la ritmo kelke similesas olta di polko, ma du-tempa (vice 2-4) e plu lenta\lingui{DEFIS}
\artiklo{skovelo}%
\vorto{skovelo}\senco{1}{Sorto di balayilo, ek linjo-peco ligita ye un ek la extremaji di longa bastono, por netigar la forno di la panobakeyi}\senco{2}{Sorto di balayilo o de brosilo, kun longa mancho, por netigar la parto interna di kanono pos pafo}
\artiklo{skrachar}%
\vorto{skrachar}\fako{trans., per} Lacerar ne-profunde la pelo per la ungli o per ulo pintoza\lingui{E}
\artiklo{skrapar}%
\vorto{skrapar}\fako{trans.} Frotar til deprenar parto di la surfaco\lingui{E}
\artiklo{skreno}%
\vorto{skreno}\senco{1}{Framo metala, kun pedi, kovrita per treliso densa, quan onu lokizas avan kameno-fairuyo, por protektar su kontre la ardoro tro granda di la fairo}\senco{2}{Framo ligna, kovrita per telo tensita, quan la piktisto, la desegnisto, lokizas avan la fenestro di sua chambro laborala por diminutar la lumo-dazlo tro granda.}\lingui{EFR}\subvorto{cinemo}{*skrino}{Telo blanka, tabelo blanka, sur qua la lumo-radii qui departas de filmo-projektoro, su reflektas por *reproduktar la imajo projektita. (analoge pri televiziono)}{}{}
\artiklo{skribar}%
\vorto{skribar}\fako{trans.} Trasar la grupi de karakteri konvencionita por reprezentar vorti\lingui{DEFIS}
\artiklo{skriptar}%
\vorto{skriptar}\fako{netrans.} Esar verkifisto di literaturaji
\artiklo{skrofulo}%
\vorto{skrofulo}\fako{patol.} Morbo quan karakterizas la infleso di la ganglioni limfala (precipue olti di la kolo), formacante tumoro qua divenas ulcero\lingui{DEFIS}
\artiklo{skrofuloso}%
\vorto{skrofuloso}\fako{patol.} Formo min grava di la tuberkloso ganglionala ed ostala, che la puerii
\artiklo{skroto}%
\vorto{skroto}\fako{anat.} Envelopilo pela di la du testikli\lingui{DEFIS}
\artiklo{skrubo}%
\vorto{skrubo}\fako{tekn.} Stangeto, cilindra o kona, ek metalo, ligno, kun heliko salianta (fileto) qua su fixigas en la korpo quan lu penetras per enirar en parto kun heliko kava a qua lu adheras forte\lingui{DE}
\artiklo{skrupulo}%
\vorto{skrupulo} Desquieteso di koncienco pri punto minucioza\lingui{DEFIRS}
\artiklo{skudo}%
\vorto{skudo}\senco{1}{\textit{(blazono)} Fundo qua havas la formo di shildo, sur qua reprezentesas la peci di la blazono}\senco{2}{\textit{(en Francia)} Olima moneto-peco (ora, arjenta) qua havis sur sua averso la skudo di Francia}\lingui{FIS}
\artiklo{skultar}%
\vorto{skultar}\fako{trans.} Fasonar per cizelo (marmoro, petro, ligno, metalo bloka) por produktar statuo\lingui{DEFIRS}
\artiklo{skunero}%
\vorto{skunero}\fako{navig.} Mikra navo du-masta, rigizita quale goeleto\lingui{DEFIR}
\artiklo{skurar}%
\vorto{skurar}\fako{trans.} Netigar recevuyo, recipiento, per deprenar la slamo, sordidaji, rezidui, sedimenta, lasita da liquido qua sejornis en olu, – generale per froto pumicage\lingui{DEFS}
\artiklo{skurelo}%
\vorto{skurelo}\fako{zool.} Mikra animalo, ek la familio « roderi », gracila, gracioza, di qua la kaudo, tufatra, povas su erektar penache\latina{scurius}\lingui{DE}
\artiklo{slamo}%
\vorto{slamo} Sedimento densa sur la fundo di aquo-quanto, precipue kande olca stagnas (lageti, fosati, riveri, maro)\lingui{DE}
\artiklo{slango}%
\vorto{slango} Linguo specala, konvencionala, e plu o min sekreta, propra ye profesiono ed anke ye societo extreme privata (kom ex.societi de studenti, de raptisti, e c.)\lingui{DE}
\artiklo{slingo}%
\vorto{slingo} To quo su volvas cirkum su, forme di ringi. (ex. : slingo di kordo, di litero, di imprimo-tipo)\lingui{DE}
\artiklo{slipo}%
\vorto{slipo} La retro-iro di heliko\lingui{DE}
\artiklo{slupo}%
\vorto{slupo}\fako{navig.} Mikra navo, kun un masto vertikala, topmasto, un seglo trapezoida ed un granda seglo quadrata por la kazo di vetero mala…\lingui{DEFIR}
\artiklo{sluzo}%
\vorto{sluzo} Quaza lageto quan formacas inkluzili movebla quin onu instalis inter du parti di rivero, di kanalo, ed en qua la aquo esas elevebla ad la nivelo supera, od abasebla a la nivelo suba per retenar lu ye un ek lua extremaji e lasar lu eskapar ye la altra, od inverse\lingui{DER}
\artiklo{smalto}%
\vorto{smalto}\fako{tekn.} Vitro blua, obtenata per la fuzo di materio vitrigebla kun kobalto\lingui{DEFIRS}
\artiklo{smeraldo}%
\vorto{smeraldo} Lapido diafana, verda\lingui{DEFIRS}
\artiklo{smerilo}%
\vorto{smerilo} Petro harda, ek alumino, silico e fer-oxido, pulverigita por polisar la lapidi, la petri, la metali, kristalo\lingui{FIS}
\artiklo{*smokar}%
\vorto{*smokar}\fako{trans.} Konsumar tabako forme di sigaro, sigareto, o per kombusto en pipo
\artiklo{snapar}%
\vorto{snapar}\fako{trans.} Sizar bruske, per un maxil-ago\lingui{DE}
\artiklo{sniflar}%
\vorto{sniflar}\fako{trans.} Aspirar ulo per sua nazo (tabako pulverigita, aquo, la humoro qua stagnas en la naztrui)\lingui{E}
\artiklo{snortar}%
\vorto{snortar}\fako{netrans.) (dicesas precipue pri kavalo} Sniflar tre bruisoze por manifestar sua repugneso; ekpulsar aero per sua naztrui tre bruisoze por manifestar sua pavoro, nepacienteso, e c.\lingui{E}
\artiklo{sobra}%
\vorto{sobra}\senco{1}{Qua su moderas pri manjo e drinko}\senco{2}{\textit{(metaf.)} Qua esas moderata pri sua \textit{agi}, ne exajerema. (Ex. : Lu esas sobra pri komplimentar)}\lingui{EFIS}
\artiklo{socio}%
\vorto{socio} Ensemblo, duroza, ek homi qui vivas obedie di legi komunana. (anke aludante ula sorti de animali qui vivas ensemble : abeli, formiki, e c.)\lingui{DEFIRS}
\artiklo{sociologio}%
\vorto{sociologio} Cienco di qua la studiajo esas la cirkonstanci ed existo-kondicioni di la socio homara\lingui{DEFIRS}
\artiklo{sociologo}%
\vorto{sociologo} Ciencisto qua su okupas pri sociologio\lingui{DEFIRS}
\artiklo{sociso}%
\vorto{sociso} Peco de la intestino di porko, di mutono, quan onu plenigis ye karno de porko, o de bovo, tre spicizita, triturita e pistita, e tre presita\lingui{EFIR}
\artiklo{sodo}%
\vorto{sodo}\fako{bot.} Planto ek la familio « kenopodiacei », propra ye la regioni klimato-moderata e subtropika; un-yar-dura; di qua la folii esas quaze pulpoza, lineala, e la flori mikra e verdatra; de qui onu extraktis olim natroxo. – L. salsola soda e salsola kali. – (kemio) Salo alkalioza quan onu extraktas de ta planto ed anke ek la fuki\simbolo{CO₃Na₂}\latina{salsola soda; salsola kali}\lingui{DEFIR}
\artiklo{sodomiar}%
\vorto{sodomiar}\fako{netrans.} Agar koito anusala, inter-vira, od inter viro e muliero, o kun bestioio\lingui{DEFIS}
\artiklo{sofao}%
\vorto{sofao} Sorto di repozo-lito, kun tri dors-apogili, uzata kom sidilo\lingui{DEFIS}
\artiklo{sofismo}%
\vorto{sofismo} Argumento trompera, seduktera\lingui{DEFIRS}
\artiklo{sofistiko}%
\vorto{sofistiko} La arto di la sofisto\lingui{DEFIRS}
\artiklo{sofisto}%
\vorto{sofisto}\senco{1}{\textit{(en la epoki antiqua, en Grekia)} Sorto de profesoro di filozofio, di retoriko, qua instruktis pri quale pledar « por » e « kontre » alterne, un tezo, sen egardo pri olua exakteso}\senco{2}{\textit{(nun)} Persono qua argumentas trompere, seduktera}\lingui{DEFIRS}
\artiklo{sofito}%
\vorto{sofito}\fako{arkitekt.} La suba parto di plafono, di arkitravo, e c., ornita ye rozfenestri, e c.\lingui{DEFIRS}
\artiklo{soklo}%
\vorto{soklo}\fako{tekn.} Suportilo, sub skultura busto, statueto, sub vazo, ornivo arkitekturala\lingui{DEFIRS}
\artiklo{soko}%
\vorto{soko}\fako{agrokultivo} Larja fero-peco triangula, tranchiva, muntita an la plugilo-parto qua glitas sur sulo, e qua, kande onu enterigas lu per preso, produktas la sulko, per fendar la ter-surfaco\lingui{EF}
\artiklo{sokursar}%
\vorto{sokursar}\fako{trans.} Helpar (ulu) eskapar danjero qua tre minacas lu\lingui{EFIS}
\artiklo{sol}%
\vorto{« sol »}\fako{muziko} Noto kinesma di la « do »-gamo
\artiklo{sola}%
\vorto{sola} Qua esas ne-akompanata da sua simili; qua ne esas kun sua simili\lingui{DEFIRS}
\artiklo{solano}%
\vorto{solano}\fako{bot.} Genero de planti ek la familio « solanei »\latina{morchella esculenta}\lingui{IL}
\artiklo{soldar}%
\vorto{soldar}\fako{trans.} Kunjuntar (speci de metalo) per kompozajo metala fuzebla (per shalmo). – \textit{(cf. weldar)}\lingui{EFIS}
\artiklo{soldato}%
\vorto{soldato}\fako{milit.} Armeano sen-grada\lingui{DEFIRS}
\artiklo{solecismo}%
\vorto{solecismo} Kulpo pri la reguli di la sintaxo di linguo\lingui{DEFIRS}
\artiklo{solena}%
\vorto{solena}\senco{1}{Celebrata singla-yare per ceremonio publika}\senco{2}{Akompanata da ceremonio, da agi publika, qui igas lu partikulare importoza ed impozanta}\senco{3}{\textit{(metaf.)} Grava til igar la kozo importoza ed impozanta}\lingui{DEFIS}
\artiklo{solenoido}%
\vorto{solenoido}\fako{elektro} Ensemblo de mikra ringi inter-egala, inter-paralela, tre inter-proxima, egal-distanta, perpendikla ye un rekto qua unionas lia centri, ed en qua fluas elektro sama-sinse\lingui{DEFIRS}
\artiklo{soleo}%
\vorto{soleo}\fako{zool.} Genero de fishi, teleostea, familio ek la pleuronekti, qui vivas en omna mari. – \textit{(anat.)} Un ek la muskuli di la gambo\latina{pleuronectes solea}\lingui{EFI}
\artiklo{solfejar}%
\vorto{solfejar}\fako{trans.) (muziko} Lektar (frazo muzikala) intonante e nomante singla noto, segun la movo-rapideso e la mezuro\lingui{DEFR}
\artiklo{solicitar}%
\vorto{solicitar}\fako{trans., pri}\senco{1}{Incitar (ulo) tre insistoze pri ulo}\senco{2}{Pregar (ulu) tre insistoze por obtenar (de lu) ulo}\lingui{DEFIS}
\artiklo{solida}%
\vorto{solida} Qua konsistas ye molekuli inter-agregita ye maniero koheranta, ferma\lingui{DEFIS}
\artiklo{solidara}%
\vorto{solidara}\senco{1}{\textit{(yuro-cienco)} Tale komuna ye pluri ke singla responsas pri la ensemblo}\senco{2}{Ligita a kunhomi per responso komuna}\lingui{DEFIRS}
\artiklo{solio}%
\vorto{solio} Peco de ligno o de quadro sur qua rezeskas tota-larjese la infra parto di pordo, quan onu transiras por enirar, e qua, lokizita super la nivelo di la sulo extera, posibligas – per impedar la eniro subporda di aquo, tero, e c. – la plu facila funciono di la pordo ed igas la klozo plu tota\lingui{FI}
\artiklo{solipeda}%
\vorto{solipeda}\fako{zool.} Qua havas nur un fingro hufoza ye singla ek sua pedi\lingui{DEFIS}
\artiklo{solitara}%
\vorto{solitara} Loko solitara : ube onu esas sola\lingui{EFIS}
\artiklo{solitero}%
\vorto{solitero}\senco{1}{Diamanto qua esas inkastrita sole}\senco{2}{Ludilo ek tabuleto qua kontenas 37 trui, ye qua onu ludas (sole) per 36 stifti}\lingui{DEFIRS}
\artiklo{solstico}%
\vorto{solstico}\fako{astron.} Singla ek la du epoki en qui Suno distas maxime de equatoro e semblas sen-mova dum kelka dii\lingui{DEFIRS}
\artiklo{solvar}%
\vorto{solvar}\fako{trans.}\senco{1}{Deskompozar e retro-igar ad stando elementala}\senco{2}{Desintrikar de to quo embarasas la spirito}\lingui{EFIS}
\artiklo{solventa}%
\vorto{solventa}\fako{komerco} Qua povas pagar sua debaji\lingui{DEFIRS}
\artiklo{somato}%
\vorto{somato}\fako{biol.} Ensemblo ek la celuli qui desaparas e mortas samatempe kam la individuo (kontraste kun la celuli jerma qui persistas sen-fine per genito)\lingui{DEFIS}
\artiklo{somero}%
\vorto{somero} Sezono kun vetero varma, qua sequas printempo e preiras autuno\lingui{DE}
\artiklo{somito}%
\vorto{somito} La parto maxim supra di to quo stacas sur sua bazo\lingui{EFI}
\artiklo{somnambula}%
\vorto{somnambula} Dicesas pri homo qua, en la stando ne-normala quan produktas lua tendenco naturala o la efiko di magnetizo, agas e parolas dum sua dormo\lingui{DEFIS}
\artiklo{somnolar}%
\vorto{somnolar}\fako{netrans.} Sentar sua bezono quik dormor\lingui{DEFIS}
\artiklo{sonalio}%
\vorto{sonalio} Mikra bulo, kava, ek metalo, qua kontenas frapilo moviva qua igas lu resonigar lor movo mem mikra\lingui{I}
\artiklo{sonar}%
\vorto{sonar}\fako{netrans.) (aludante objekto o materio} Emisar movi vibrala, propagate da medio aeroforma, liquida o solida, til la membrano di la orel-timpano, produktas sento partikulara\lingui{EFIS}
\artiklo{sonato}%
\vorto{sonato}\fako{muziko} Muzikajo klasika, por un o plura instrumenti, ordinare ek « prefaco », « adagio » o « andante », menueto o finalo\lingui{DEFIRS}
\artiklo{sondar}%
\vorto{sondar}\fako{trans.}\senco{1}{Determinar la profundeso-grado di maro, di fluvio, di rivero, e c., la naturo di la fundo, e c., per plumbo-bloko ligita a la extremajo libera di kordeto sur qua onu trasis repero-punti, interdiste di un braco}\senco{2}{Explorar la naturo di la tereni, o facar puteo arteza por sinkar sulo-trepano}\senco{3}{\textit{(kirurg.)} Explorar plago, veziko, per sinkar stangeto}\senco{4}{\textit{(metaf.)} Probar deskovrar la intenci (di ulu.)}\lingui{DEFRS}
\artiklo{soneto}%
\vorto{soneto} Poeziajo ek dek-e-quar versi, ek du quatreni (durima) sequata da du terceti\lingui{DEFIRS}
\artiklo{sonjar}%
\vorto{sonjar}\fako{netrans.) (aludante psiko} Dum dormo, facar kombinuri (stranja, en la kazi maxim multa) ek fakti, idei, senti imaginala\lingui{FIS}
\artiklo{sonko}%
\vorto{sonko}\fako{bot.} Planto di qua la suko esas laktatra, ek la tribuo « chikoracei »\latina{sonchus}\lingui{L}
\artiklo{sonora}%
\vorto{sonora} Qua resonas bone\lingui{DEFIS}
\artiklo{soporto}%
\vorto{soporto}\fako{mekan.} Organo di mashino, qua igas plufacila la movo horizontala di parto sur altra\lingui{F}
\artiklo{soprano}%
\vorto{soprano}\fako{muziko}\senco{1}{Voco qua povas atingar la grado supra sur la noto-skalo}\senco{2}{Persono qua havas tala voco}\lingui{DEFIRS}
\artiklo{soquo}%
\vorto{soquo} Pedo-vesto ek ligno, ledro, aden qua onu insertas sua pedo kun la pedovesto ordinara, por protektar su de humideso
\artiklo{sorbeto}%
\vorto{sorbeto} Kremo mi-konjelita a qua mixesis siropo de citroni, de fragi, kirsho-liquoro, – konsumata, kande fuzanta, kom drinkajo\lingui{DEFIRS}
\artiklo{sorbo}%
\vorto{sorbo}\fako{bot.} Frukto di arboro ek la familio « rozacei »\latina{sorbus}\lingui{DEFIS}
\artiklo{sorcar}%
\vorto{sorcar}\senco{1}{\textit{(trans.)} Praktikar la arto di la homo a qua onu imputas povo supernatura quan lu havas de pakto kun la spiriti infernala, kun diablo}\senco{2}{\textit{(trans.)} Submisa a la efiko da artifico di sorcisto}\lingui{EF}
\artiklo{sordida}%
\vorto{sordida} Di qua la neteso alteresas (da makuli, polvo, kraso, reziduo, exkrementi)\lingui{EFIS}
\artiklo{sordino}%
\vorto{sordino}\senco{1}{\textit{(muziko)} (a) Espineto, di qua la kordini vibrigas mikra peci de ligno envelopita en drapo, por moderar la soni; (b) Kono de kartono, kun truo ye la bazo, quan onu lokizas en la funelo di chaskorno – mikra tubo ek ligno quan onu lokizas en la parto infra di la trumpeto, por moderar la soni; (c) Mikra peco de ligno, de ivoro, e c., di qua la formo esas olta di pektilo kun tri denti kavigita, quan onu inkastras sur la tresto di la violini, alti, violonceli, kontrabasi, por diminutar la vibri di la framo e moderar la soni}\senco{2}{\textit{(metaf.)} Agar sen bruiso por ke nulu remarkez ta ago}\lingui{DEFIS}
\artiklo{sorgar}%
\vorto{sorgar}\fako{trans.} Asumar volunte peni e sucii por la komforto di ulu, por la prezervo di ulo\lingui{D}
\artiklo{sorgumo}%
\vorto{sorgumo}\fako{bot.} Planto ek la familio « graminei », kelke simila a maizo\latina{sorghum vulgare}
\artiklo{sorito}%
\vorto{sorito}\fako{logiko} Serio de propozicioni, interligita tale ke la atributo di la unesma divenas la subjekto di la duesma, e tale pluse – la lasta esas implicite inkluzita en la unesma se la rezono esas korekta\lingui{DEFIRS}
\artiklo{sorto}%
\vorto{sorto} Eso-maniero di persono, di kozo (kompare a to di quo la naturo esas analoga)\lingui{DEFIRS}
\artiklo{sospirar}%
\vorto{sospirar}\fako{netrans.} Agar ta expiro profunda e pasable duroza quan onu lasas eskapar kande onu opresesas da sento o da sentimento chagreniganta\lingui{FIS}
\artiklo{sovaja}%
\vorto{sovaja}\senco{1}{Qua vivas exter la socii civilizita}\senco{2}{\textit{(pri planto)} Qua kreskas sen kultiveso}\senco{3}{\textit{(pri loko)} Adube la homo ne venas, ne efikas por agar o facar}\senco{4}{\textit{(metaf.)} Qua ne volas relatar kun sua kunhomi}\lingui{EFIS}
\artiklo{sozio}%
\vorto{sozio} Persono quan onu konfundas ad altra, tante grandan esas lia intersimileso\lingui{DFI}
\artiklo{spaco}%
\vorto{spaco}\senco{1}{Extensajo de ca a ta punto, vakua ye korpi solida}\senco{2}{(1) Extensajo ideala, egardata quale se lu kontenus omna extensaji reala, omna korpi qui existas e quin la spirito konceptas kom posibla. (2) Parto di ta extensajo ideala}\lingui{EFIS}
\artiklo{spadico}%
\vorto{spadico}\fako{bot.} Infloresenco forme di kastono, od asembluro de folii sen-pedikula sur axo komuna, qua havas flori envelopita da granda brakteo (spato)
\artiklo{spado}%
\vorto{spado} Garden-implemento, ek fero-plako larja e tranchiva, muntita an la extremajo di mancho longa, quan onu uzas por fendar tero e kulbutigar la glebi\lingui{DE}
\artiklo{spanielo}%
\vorto{spanielo}\fako{zool.} Chaso-hundo de fonto Hispana, di qua la pili esas longa, e la oreli pendanta\lingui{DEF}
\artiklo{spano}%
\vorto{spano} Rekorturo plu o min dina quan onu deprenas (1) per raboto, cizelo, kande onu fasonas ligno-peco; (2) per utensilo tranchiva kande onu fasonas peco de stalo, de fero, de kupro, sur torno-mashino, frezo-mashino, boro-mashino\lingui{D}
\artiklo{spanto}%
\vorto{spanto}\fako{navig.} Parto di la korpo di navo, varango qua formacas un ek la kosti di la framo\lingui{DR}
\artiklo{spanyoleto}%
\vorto{spanyoleto} Fero-peco qua uzesas por klozar od apertar fenestro per mancho manovrebla rotace\lingui{DFI}
\artiklo{sparadrapo}%
\vorto{sparadrapo}\fako{medic.} Texuro di qua la averso provizesas ye strato de plastro
\artiklo{sparar}%
\vorto{sparar}\fako{trans.} Ne spensar la toto, por povar rezervar la restajo\lingui{D}
\artiklo{sparviero}%
\vorto{sparviero}\fako{zool.} Speco di falkono quan onu dresas por chaso\latina{falco nisus}\lingui{DFI}
\artiklo{spasmo}%
\vorto{spasmo}\fako{medic.} Kontrakteso bruska di ula organi (stomako, intestini, veziko, e c.)\lingui{EFIRS}
\artiklo{spato}%
\vorto{spato}\fako{mineral.} Nomo di plura substanci minerala lameloza e kolor-chanjanta\lingui{DFIRS}
\artiklo{spatulo}%
\vorto{spatulo} Instrumento de ligno, ivoro, metalo, e c., plu larja e plata ye un extremajo kam kuliero, quan uzas (1) la apotekisti por agitar la preparaji farmaciala, extensar la unguenti, plastri, e c.; (2) la piktisti, por dilutar e pistar la farbi; (c) la skultisti, por modlar vaxo, argilo; (d) la gisisti, por facar mulduri\lingui{DEFIS}
\artiklo{specifika}%
\vorto{specifika} Qua indikas ta o ca speco, exkluzante omna ceteri\lingui{DEFIS}
\artiklo{specimeno}%
\vorto{specimeno} Parto de ensemblo destinita a donar ideo pri la cetero\lingui{DEFS}
\artiklo{speco}%
\vorto{speco} Dividuro di genero\lingui{DEFIRS}
\artiklo{spegulo}%
\vorto{spegulo} Vitro-plako, polisita, di qua la reverso esas kovrita indute ye folio de stano merkuriizita, o metalo-plako polisita, ube onu povas vidar sua vizajo reflektita\lingui{DEI}
\artiklo{spektaklo}%
\vorto{spektaklo}\senco{1}{Reprezentajo teatrala}\senco{2}{Ensemblo quan embracas la regardo}\lingui{EFIRS}
\artiklo{spektar}%
\vorto{spektar}\fako{trans.} Vidar (ulo) eventar avan sua regardo\lingui{EFIS}
\artiklo{spektro}%
\vorto{spektro}\senco{1}{Ento nevidebla qua aparas, pavorigante, terorigive, lor halucino}\senco{2}{\textit{(optiko)} Imajo oblonga en qua reflektesas, forme di bendi interparalela, reda, oranjea, flava, verda, blua – la radii parta qui rezultas de la deskompozo di radio ek lumo sen-kolora trans prismo de vitro. – (en la senco I, preferindan esas uzar « fantomo »)}\lingui{DEFIRS}
\artiklo{spektroskopio}%
\vorto{spektroskopio}\fako{fiziko} Analizo studia di la lumo-spektri\lingui{DEFIRS}
\artiklo{spektroskopo}%
\vorto{spektroskopo}\fako{fiziko} Aparato quan onu uzas por studiar la lumo-spektri\lingui{DEFIRS}
\artiklo{spekular}%
\vorto{spekular}\fako{netrans., pri}\senco{1}{Agar explori teoriala}\senco{2}{\textit{(borso)} Facar operaci (kompri, vendi) riskoza, fundamento di chanci di kresko o di deskresko di la varo-preci, o di la aciono-preci}\lingui{DEFIRS}
\artiklo{spelto}%
\vorto{spelto}\fako{bot.} Speco di frumento di qua la brano duras adherar a la grano pos drasheso\latina{triticum spelta}\lingui{DEIS}
\artiklo{spencero}%
\vorto{spencero}\senco{1}{Frako sen baski}\senco{2}{Korsajo sen jupo}\lingui{DEFI}
\artiklo{spensar}%
\vorto{spensar}\fako{trans.}\senco{1}{Donar, uzar (pekunio) por prokurar a su ulo}\senco{2}{\textit{(metaf.)} Donar, uzar (sua tempo, esforci, fakultati) por obtenar la rezultajo quan onu deziras}\lingui{EFI}
\artiklo{spergulo}%
\vorto{spergulo}\fako{bot.} Planto ek la familio « kariofilei »\latina{spergula}\lingui{DEFIS}
\artiklo{spermaceto}%
\vorto{spermaceto} Materio perlomatrea, extraktita de oleo qua abundegas en la kavaji di la kapo di la kashaloto, ed uzata por facar bujii luxoza, koldkremo, e c.\lingui{DEFIR}
\artiklo{spermatido}%
\vorto{spermatido}\fako{biol.} Spermatozoido nova, derivita de la duopla du-partesko inter-sucedanta di la spermatociti
\artiklo{spermatocito}%
\vorto{spermatocito}\fako{biol.} Spermatogonio qua, cesante partigar su, komencas kreskar e donas, pose, per du duparteski intersucedanta, la spermatidi\lingui{DEFIS}
\artiklo{spermatoforo}%
\vorto{spermatoforo}\fako{biol.} Kapsulo qua kontenas pakaji de spermatozoidi e qua renkontresas precipue che la cefalopodi\lingui{DEFIS}
\artiklo{spermatogenezo}%
\vorto{spermatogenezo}\fako{biol.} Ensemblo ek la procesi evolucionala qua abutas a la formaco di la spermatozoidi\lingui{DEFIS}
\artiklo{spermatogonio}%
\vorto{spermatogonio}\fako{biol.} Celulo di la tubo seminifera qua su partigas multega-foye e deskreskas pri sua volumino : lu produktas pose la spermatociti\lingui{DEFIS}
\artiklo{spermatoreo}%
\vorto{spermatoreo}\fako{patol.} Emisi spermala ne-volata\lingui{DEFIS}
\artiklo{spermatozoido}%
\vorto{spermatozoido}\fako{biol.} Animaleto quan onu trovas en la spermo di la animali, ed en la poleno qua fekundigas la planti\lingui{DEFIS}
\artiklo{spermino}%
\vorto{spermino}\fako{farmacio} Bazo ne bone analizita qua esus la substanco agiva di la preparaji opoterapiala quin onu obtenas ek la testikli, ovarii, spleno, prostato, e c.\lingui{DEFIS}
\artiklo{spermo}%
\vorto{spermo}\fako{fiziol.} Liquoro fekundigiva de la maskulo\lingui{DEFIRS}
\artiklo{spermocentro}%
\vorto{spermocentro}\fako{biol.} Centrosomo di spermatozoido, pri qua Plättner dicis ke lu esas renkontrebla en la pinto antea di la kapo, kontraste kun la dico da Flemming e Hertwig, segun qui lu esus en la segmento meza
\artiklo{spico}%
\vorto{spico} Singla ek la substanci vejetala, aromata o pikanta, quin onu uzas por saporozigar disho : kariofilo, muskado, pipro, cinamo, jinjero pulverigita\lingui{DEFIS}
\artiklo{spiko}%
\vorto{spiko}\fako{bot.} Parto extrema di la stipo di la graminei, qua portas, asemblita cirkum axo, la grani di la planto\lingui{EIS}
\artiklo{spinato}%
\vorto{spinato}\fako{bot.} Planto ek la familio « kenopodiacei », herbo un-yar-dura, di qua la folii, pos koqueso, esas nutrivo bonega\lingui{DEIRS}
\artiklo{spindelo}%
\vorto{spindelo}\senco{1}{Stango ligna, di qua la extremaji esas dinigita, un de li pinta, la altra rondigita, quan onu uzas lor filifo per ruko, por tordar la filo e rular lu segun ke lu formacesas}\senco{2}{\textit{(roto di vehilo)} Parto di la rot-axo qua fricionas kontre la kuseneti}\lingui{DE}
\artiklo{spinelo}%
\vorto{spinelo}\fako{mineral.} Rubino pale-reda
\artiklo{spino}%
\vorto{spino}\fako{anat.) (che la homo ed ula animali altra} Ensemblo ek la vertebri dorsala, artikoza, ye qui konsistas la axo di la karpenturo osta\lingui{EFIS}
\artiklo{spionar}%
\vorto{spionar}\fako{trans.} Observar habile e sekrete ta (o ti) di qua (o di qui) onu profitos de saveskar lua (o lia) intenci – o to quo eventas che la enemiko\lingui{DEFIRS}
\artiklo{spiralo}%
\vorto{spiralo}\senco{1}{\textit{(geom.)} Kurvolineo, sur plano, qua trasas sama-sinse serio de rotaci cirkum punti fixa (polo) de qua lu deviacas gradoze segun *leyo (lego) determinita}\senco{2}{\textit{(horloj.)} Mikra resorto qua reglas la eskapo di la pendulo de horlojeto, e mantenas la izokroneso di lua ocili}\lingui{DEFIRS}
\artiklo{spiremo}%
\vorto{spiremo}\fako{biol.} Reto de nuklei qua, kande mitoso esas quik eventonta, divenas, komence, kordono longa, tenua, dornoza, kun falduri tre interproxima – e, pose, divenas plu kurta, plu dika, kun surfaco glata, e di qua la falduri minproximeskas\lingui{DEFIS}
\artiklo{spireo}%
\vorto{spireo}\fako{bot.} Genero de planti ek la familio « rozacei »\latina{spiraeum}\lingui{DEFIS}
\artiklo{spirito}%
\vorto{spirito}\fako{religio} Substanco nekorpala (kontraste kun materio, substanco korpala) a qua uli imputas la penso-fakultato\lingui{DEFIRS}
\artiklo{spiro}%
\vorto{spiro} Lineo qua volvas cirkum axo quale la fileto di skrubo, di helico\lingui{DEFIRS}
\artiklo{spiso}%
\vorto{spiso} Stango ek fero, di qua un extremajo esas pinta, quan onu pasigas tra grosa peco de karno, tra pultruno, por rostar li, e quan onu rotacigas, o manue, o per mekanismo\lingui{DES}
\artiklo{splenalgio}%
\vorto{splenalgio}\fako{patol.} Doloro ye la spleno\lingui{DEFIRS}
\artiklo{splendida}%
\vorto{splendida} Di qua la luxozeso brilanta esas plena ye grandeso\lingui{DEFIS}
\artiklo{splenito}%
\vorto{splenito}\fako{patol.} Inflamuro di la spleno\lingui{DEFI}
\artiklo{spleno}%
\vorto{spleno}\fako{anat.} Vicero sponjatra, saturita ye sango, en la hipokondro sinistra, di qua la rolo esas ne-bone savata – e quan, olim, onu egardis kom la fonto di melankolio\lingui{DEFRS}
\artiklo{splinar}%
\vorto{splinar}\fako{netrans.} Afektesar da ta hipokondrio qua karakterizis la Britaniani, e su manifestis per la enoyo pri omno\lingui{DEFRS}
\artiklo{splinto}%
\vorto{splinto}\fako{tekn.} Stifto qua asemblas la parti di charniri, di bukli, e c.\lingui{DR}
\artiklo{splisar}%
\vorto{splisar}\fako{trans.) (tekn.} Ligar (du kordi), extremajo an extremajo, interplektante la toroni\lingui{DEF}
\artiklo{splito}%
\vorto{splito}\senco{1}{Fragmento forlansita da korpo qua explozas}\senco{2}{Mikra peceto de ligno, de metalo, e c., tenua ed akuta, qua penetras la pelo acidente}\lingui{DE}
\artiklo{spoko}%
\vorto{spoko}\fako{tekn.} Singla de la parti di roto qui iras de la nabo a la felgo\lingui{E}
\artiklo{spoliar}%
\vorto{spoliar}\fako{trans.} Deprenar (de ulu), per ruzo o per violento, por proprigar lu a su (to quo apartenis a lu)\lingui{DEFIS}
\artiklo{spondeo}%
\vorto{spondeo}\fako{versif.} Pedo longa ye du silabi\lingui{DEFIRS}
\artiklo{sponjo}%
\vorto{sponjo}\senco{1}{\textit{(zool.)} Zoofito, qua konsistas ek amaso de tisui fibroza, plu o min densa, flexebla, elastika, tre lejera, imbibebla, e qua, kande vivanta, indutesas de materio gelatinatra mi-fluida, iritebla ed efemera}\senco{2}{La korpo di ta zoofito, uzata diversa-rolo, pro lua karakterizivo : absorbar la liquidi, e lasar li ekirar lor preseso mem tre poka}\latina{spongia}\lingui{EFIS}
\artiklo{spontana}%
\vorto{spontana}\senco{1}{Quan onu agas o facas ipse, propra-move, sen impulseso da forco aplikata de extere}\senco{2}{Qua aspektas agata o facata propra-move – di qua la kauzo ne esas perceptebla}\lingui{DEFIS}
\artiklo{sporadika}%
\vorto{sporadika}\senco{1}{\textit{(geol.)} Dispersita}\senco{2}{\textit{(patol.)} Qua afektas individui qui vivas ne-grupe – qua ne esas epidemia}\lingui{DEFIRS}
\artiklo{sporangio}%
\vorto{sporangio}\fako{bot.} Quaza sako mikra, qua kontenas la spori, che la kriptogami\lingui{DEFIS}
\artiklo{sporno}%
\vorto{sporno} Brancho ek metalo, qua esas muntita an la talono di la kavalkanto e finas per disko dentetoza e movebla, quan la kavalkanto sinkas aden la flanki di sua kavalo kom stimulo\lingui{DFI}
\artiklo{sporo}%
\vorto{sporo}\fako{bot.} Korpuskulo genitera (di ula planti kriptogama)\lingui{DEFIRS}
\artiklo{sporofito}%
\vorto{sporofito}\fako{biol.} Stando di la ovo di la « uridinei » lor la fekundigeso, kande la du pronuklei ne esas ja kunfuzita
\artiklo{sportar}%
\vorto{sportar}\fako{netrans.} Exercar su korpale, metodoze, por developar en su la bona qualesi di sua korpo ed ula bona qualesi di la psiko (loyaleso, energiozeso, perseveremeso) per pedo-balono, natado, skermado, aviacado, e c.\lingui{DEFIRS}
\artiklo{spoto}%
\vorto{spoto}\fako{fiziol.} Mikra saliajo, o mikra makulo sur la pelo, qua igas -ulakaze – plu « spicoza » la fizionomio (di muliero)\lingui{E}
\artiklo{spozo}%
\vorto{spozo} Persono unionita a persono altra-sexua per mariajo\lingui{EFIS}
\artiklo{sprato}%
\vorto{sprato}\fako{zool.} Nomo vulgara di mikra fisho, sorto di haringo, qua vivas proxim la litoro Franca di Atlantiko\latina{clupea sprattus}
\artiklo{spricar}%
\vorto{spricar}\fako{netrans.}\senco{1}{\textit{(aludante liquido, fluido)} Springar impetuoze, ek orifico relative mikra}\senco{2}{\textit{(metaf.)} Springar, des-intrikante su, ekirar subite, vivace}\lingui{D}
\artiklo{sproso}%
\vorto{sproso}\fako{bot.} Kresko komencala di burjono\lingui{D}
\artiklo{spular}%
\vorto{spular}\fako{trans.} Volvar (kozo) cirkum altra quaze filo sur bobino\lingui{DE}
\artiklo{spumo}%
\vorto{spumo}\senco{1}{Quaza moso qua naskas surface di liquido quan onu agitas, varmigas o qua fermentacas}\senco{2}{Bavajo mosatra qua naskas ye la boko di ula animali kande li esas ecitita od iritita}\lingui{EFIS}
\artiklo{sputar}%
\vorto{sputar}\fako{trans.} De-jetar (to quon onu havas en sua boko)\lingui{EFI}
\artiklo{squadro}%
\vorto{squadro} Instrumento ek ligno, metalo, quan onu uzas por trasar ortanguli o perpendikli\lingui{IS}
\artiklo{squalo}%
\vorto{squalo}\fako{zool.} Genero de fishi kartilagoza\latina{squalus}\lingui{FIS}
\artiklo{squamo}%
\vorto{squamo}\fako{zool.} Singla ek la plaki inter-apudpozita od imbrikita qui tegas la pelo di ula fishi, di ula repteri, e c.
\artiklo{squatar}%
\vorto{squatar}\fako{netrans.}\senco{1}{\textit{(aludante la animali)} Sidar sur sua gropo}\senco{2}{\textit{(aludante la homo)} Sidar tale ke la gambi esas rifaldita, la glutei tushante la taloni}\lingui{E}
\artiklo{squilo}%
\vorto{squilo}\fako{bot.} Genero de « liliacei »\latina{Scilla maritima}
\artiklo{st!}%
\vorto{st!} Vorto quan onu dicas ad ulu por ke lu tacez
\artiklo{stabila}%
\vorto{stabila} Qua tendencas permanar en la sama loko, en la sama situeso, en la sama stando\lingui{DEFIS}
\artiklo{stablo}%
\vorto{stablo} Klozajo muroza e tektoza ube onu lojigas la bestii, la bruti\lingui{DEFIS}
\artiklo{stabo}%
\vorto{stabo}\fako{armeo} Ensemblo de oficiri quin onu selektis por servar oficiro superiora qua komandas armeo, ed esar apud lu quaze kom konsilantaro qua transmisas ed obediigas lua komandi\lingui{DER}
\artiklo{stacar}%
\vorto{stacar}\fako{netrans.) (aludante homo} Havar sua korpo vertikala sur sua pedi, t.e. standar ne-sidante, ne-kushinte, ne-jacante, ne genu-pozinte, e ne-inklinite\lingui{I}
\artiklo{stacionara}%
\vorto{stacionara} Qua haltigas sua movo e permanas dum kelka tempo en la sama loko (o stando)\lingui{DEFIS}
\artiklo{staciono}%
\vorto{staciono}\senco{1}{Singla ek la loki ube haltas treno, por embarkar o desembarkar pasajeri o vari}\senco{2}{Loko ube esas, o haltas, la veturi publik-uzebla \textit{(autofiakri – F. taxi – autobusi, e c.)}}\lingui{DEFIS}
\artiklo{stadio}%
\vorto{stadio}\fako{antique} Areo klozita, longa ye 180 metri, ube onu interluktis por obtenar la premio pri homokuro. – \textit{(nun)}. \textit{Lico} por la konkursi atletala, automobilala, biciklala\lingui{DEFIRS}
\artiklo{stafilo}%
\vorto{stafilo}\fako{bot.} Arboreto, vicina ye la aceri, kun folii kompozita, kun aristi plumatra\latina{staphylea}
\artiklo{stafisagro}%
\vorto{stafisagro}\fako{bot.} Planto ek la familio « ranunkulacei »\latina{Delphinium staphis agria}
\artiklo{stagnar}%
\vorto{stagnar}\fako{netrans.} Qua ne fluas, qua ne cirkulas, qua ne rinoveskas e, pro lo, eventuale putreskas. – \textit{(anke metaf.)}\lingui{DEFIS}
\artiklo{stakiso}%
\vorto{stakiso}\fako{bot.} Planto ek la familio « labiei », deveninta de Japonia, di qua la *rizomi esas manjebla ed analoga a salsifio\latina{stachys}
\artiklo{stalagmito}%
\vorto{stalagmito} Agreguro, sur la sulo di kavajo subtera, per infiltruri qui falas de la vulto\lingui{DEFIRS}
\artiklo{stalaktito}%
\vorto{stalaktito} Agreguro an la vulto di kavajo subtera, per infiltruri qui kontenas dissolvite sali kalkoza, silicoza, e c.\lingui{DEFIRS}
\artiklo{stalo}%
\vorto{stalo} Metalo ek fero pura a qua esas enkorpigita ula quanto de karbo – plu harda, e di qua la grani esas plu tenua e plu briloza kam olti di fero\lingui{DER}
\artiklo{stamino}%
\vorto{stamino}\fako{bot.} Organo sexuala maskulo, che la vejetanti\lingui{EFIS}
\artiklo{stampar}%
\vorto{stampar}\fako{trans.} Provizar ye marko quan onu imprimas sur objekto por atestar lua autentikeso\lingui{DEFR}
\artiklo{stanco}%
\vorto{stanco} Quanto fixa de versi di qua la senco esas kompleta, dispozita segun ula ordino, e qua, repetate periodoze, formacas la dividuri di poemeto\lingui{DEFIRS}
\artiklo{standar}%
\vorto{standar}\fako{netrans.}\senco{1}{\textit{(aludante persono o kozo)} Esar tala en ica od en ita temp-uno}\senco{2}{\textit{(aludante persono)} Esar tala, en la hierarkio sociala}\lingui{DE}
\artiklo{standardo}%
\vorto{standardo} La insigno nacionala, di qua un ek la bordi vertikala muntesas sive an stango (armeo terala), sive an kordo (armeo navala)\lingui{DEFIRS}
\artiklo{stango}%
\vorto{stango} Peco de ligno, de metalo, rigida, longa e streta\lingui{DI}
\artiklo{stano}%
\vorto{stano}\fako{kemio} Korpo simpla, metalo grizatra, poke sonora, poke duktila, tre forjebla, maleebla, de qua onu facas implementi koqueyala\simbolo{Sn}\lingui{EFIS}
\artiklo{stapelio}%
\vorto{stapelio}\fako{bot.} Genero de « asklepidacei », ek planti xerofila, kun folii kaktusatra e granda flori. EFI?\latina{stapelia}
\artiklo{staplo}%
\vorto{staplo}\senco{1}{Butikego ube onu vendas, ye la nomo di la stato o di kompanio komercala}\senco{2}{Portuo, urbo komercala ube la vari depozesas varte di lia exportaceso pos-importaca}\lingui{DE}
\artiklo{starchar}%
\vorto{starchar}\fako{trans.}\senco{1}{Apretar per amilo}\lingui{DE}\subvorto{metaf.}{starchita}{Di qua la su-teno, la posturo, en mondumo, esas tam rigida kam kamizo-plastrono apretita per amilo}{2}{}
\artiklo{startar}%
\vorto{startar}\senco{1}{\textit{(trans.)} Igar funcioneskar}\senco{2}{\textit{(netrans.)} Funcioneskar}\lingui{DE}
\artiklo{startero}%
\vorto{startero}\fako{sur lico di keval-kuri} La persono qua emisas la departo-signalo, per inklinar la standardo quan lu tenas en sua manuo dextra\lingui{DEFIRS}
\artiklo{statiko}%
\vorto{statiko}\fako{fiziko} Parto di mekaniko qua traktas la kondicioni pri la equilibro di la forci\lingui{DEFIRS}
\artiklo{statistiko}%
\vorto{statistiko} Cienco di qua la skopo esas donar dokumenti exakta pri la parti diversa di la vivo di populo, per amasigar, klasifikar, kontar omna fakti ek ta o ca naturo (naski, morti, mariaji, produkturi, komerco, konsumo, kondamni, e c.) qui eventas en periodo determinita\lingui{DEFIRS}
\artiklo{stativo}%
\vorto{stativo}\fako{fotogr.} Sorto di suportilo kun tri pedi, ek tri parti faldebla, uzata por la fotografili di la fotografisti\lingui{D}
\artiklo{stato}%
\vorto{stato} La naciono egardata kom korpo politikala\lingui{DEFIRS}
\artiklo{statolito}%
\vorto{statolito}\fako{bot.} Nomo di la amilo-grani qui, per lia pezo, kontributas igar la vejetanti sentiva pri la efiko da gravito. \textit{(geotropismo)}
\artiklo{statoro}%
\vorto{statoro}\fako{elektro) (en motoro kun elektrofluo polifaza} Parto fixa di la mashino, opoze a \textit{rotoro}
\artiklo{statuo}%
\vorto{statuo} Figuro skultura, generale ek marmoro, ma anke ek metalo, ek ligno, qua reprezentas tota-reliefe homo od animalo\lingui{DEFIRS}
\artiklo{staturo}%
\vorto{staturo} Longeso di la korpo di homo, mezurata kande lu stacas\lingui{DEFIS}
\artiklo{statuto}%
\vorto{statuto}\fako{yuro-cienco} To quo rezolvesis kom aplikenda (pri la personi, la kozi, en ta o ca kazo)\lingui{DEFIRS}
\artiklo{stearino}%
\vorto{stearino}\fako{kemio} Materio ek lameli, perlomatrea, ek glicerino e CH₃.(CH₂)₁₆ CO₂H, quan onu extraktas de sebo\lingui{DEFIRS}
\artiklo{steatito}%
\vorto{steatito}\fako{mineral.} Silikato di magnezio, varietato di talko kompakta, granetatra\lingui{DEFIS}
\artiklo{stebar}%
\vorto{stebar}\fako{sutado} Trairar steke stofo per du fili quin onu sinkas, per la pinto di agulo suto-mashinala, aden trui di qui singla ek le sequanta situesas egaldiste de la pre-esanta
\artiklo{*stejo}%
\vorto{*stejo} La parto di la teatro ube pleas la aktori. \textit{(anke metaf.)}\lingui{E}
\artiklo{*stejo}%
\vorto{*stejo} Ta parto di teatro ube pleas la aktori
\artiklo{stekar}%
\vorto{stekar}\fako{trans.} Penetrigar per koakto la pinto di objekto (kom ex. : paliso aden tero, e c.)\lingui{DE}
\artiklo{stelario}%
\vorto{stelario}\fako{bot.} Planto unyardura, ek la familio « primulacei », quan konsumas ula uceli\latina{stellaria}
\artiklo{steleo}%
\vorto{steleo}\fako{arkitekt.) (epoki antiqua} Monolito di qua la formo esas olta di la fusto di kolono, di kolono ruptita o di cipo, destinita (maxim-multa-kaze) recevor epigrafo\lingui{DEFI}
\artiklo{stelionato}%
\vorto{stelionato}\fako{jurisprudenco} Fraudo qua konsistas ye prizentar kom libera, domenon hipotekizita – o ye vendar, hipotekizar domeno di qua onu ne esas la proprietero\lingui{DEFIS}
\artiklo{stelo}%
\vorto{stelo}\senco{1}{\textit{(linguo ne-ciencala)} Omna astro qua brilas en firmamento (ecepte Suno e Luno)}\senco{2}{\textit{(linguo ciencala)} Korpo cielala, lumifanta, sen movo videbla, e di qua la lumo quaze cintilifas}\senco{3}{\textit{(astrologio)} Astro qua, acensante en cielo lor la nasko di homo, egardesas kom influanta bone o male lua fato}\senco{4}{\textit{(teatro, cinemo)} Artisto qua brilas, unika o preske unika, en trupo de aktori}\senco{5}{\textit{(metaf.)} \textit{(pro analogio kun la radiifo di la steli)} (1) punto briloza; (2) objekto di qua la parti dispozesas quale la spoki di roto, e qua memorigas la formo di stelo reprezentita desegne}\lingui{EFIS}
\artiklo{stelto}%
\vorto{stelto} Stango, provizita ye quaza estribo ek ligno sur qua onu pozas sua pedo por plualtigar su (kom ex. : por marchar en marshi, e c.)\lingui{DE}
\artiklo{*stencilo}%
\vorto{*stencilo} Folio de papero impregnita ye vaxo, sur qua onu skribas (generale per skribo-mashino) tale ke la tipi perforas la folio; la inko imprimala pasas tra ta perforuri
\artiklo{stenografar}%
\vorto{stenografar}\fako{trans.} Enrejistrar vortope (per manuo-skribo, uzante abrevio-sistemo) omna vorti quin pronuncas ulu, mem se la ritmo esas rapidega – la limito esante la komprenebleso-grado di lo dicita (komprenebleso audala)\lingui{DEFIRS}
\artiklo{stenotipar}%
\vorto{stenotipar}\fako{trans.} Stenografar per mashino fingro-funcionanta, uzante la literi di la alfabeto ordinara, di qui onu povas skribar plura sama-stroke, e silabope\lingui{DEFIRS}
\artiklo{stentoro}%
\vorto{stentoro} Persono di qua la voco esas tre resonanta, soneganta\lingui{DEFS}
\artiklo{stepo}%
\vorto{stepo} Vasta regiono plana, ek tero, kovrata ye bushi, pastureyi\lingui{DEFIRS}
\artiklo{stereografio}%
\vorto{stereografio} La arto reprezentar la solidi\lingui{DEFIRS}
\artiklo{stereokamero}%
\vorto{stereokamero} Kamero di stereoskopo\lingui{DEFIS}
\artiklo{stereometrio}%
\vorto{stereometrio} Apliko-geometrio qua havas kom temo la mezuro di la korpi solida\lingui{DEFIRS}
\artiklo{stereoskopo}%
\vorto{stereoskopo} Instrumento qua, per du imaji plana di un objekto, kun la sama difero angulala kam la imaji produktita en singla ek la okuli di homo, da la objekto, igas ica aspektar reliefe kande onu regardas samatempe la du imaji, singla per un okulo\lingui{DEFIRS}
\artiklo{stereotipar}%
\vorto{stereotipar}\fako{trans.}\senco{1}{Fixigar (paginedo ek tipi asemblura) por imprimi futura, sive per soldar la tipi, sive per facar mulduro ek la bloko de oli, aden qua onu gisas aloyuro ek metali}\senco{2}{\textit{(metaf.)} Fixigar ferme e ne-chanjeble}\lingui{DEFIRS}
\artiklo{stereotomio}%
\vorto{stereotomio} Apliko-geometrio qua havas kom temo la seciono di la korpi solida\lingui{DEFIRS}
\artiklo{sterila}%
\vorto{sterila}\senco{1}{Dicesas pri sulo-peco qua ne produktas (la frukti obtenebla normale de tero).- \textit{anke metaf}}\senco{2}{Dicesas pri muliero o femino qua ne esas fekundigebla}\senco{3}{De qua nulo rezultas}\lingui{DEFIRS}
\artiklo{sterko}%
\vorto{sterko} Dungo ek la litiero di la animali domestika, a qua esas mixita lia exkrementi, e qua putras pro fermentaco\lingui{SL}
\artiklo{sterleto}%
\vorto{sterleto}\fako{zool.} Speco di sturgo, quan onu renkontras en la parto orientala di Europa ed en Azia, e qua aspektas kom la formo yuna di la sturgo. De olu onu obtenas kaviaro\latina{acipensar ruthenus}\lingui{DEFR}
\artiklo{sternar}%
\vorto{sternar}\fako{trans.} Extensar (persono, kozo, su) longesale sur ulo\lingui{L}
\artiklo{sternumo}%
\vorto{sternumo}\fako{anat.} Osto plata an-avan torako, meze di olca, e sur qua « muntesas » la kosti\lingui{DEFIS}
\artiklo{sternutar}%
\vorto{sternutar}\fako{netrans.} Expirar bruske, tra nazo e boko, per movo konvulsatra di diafragmo\lingui{EFIS}
\artiklo{stero}%
\vorto{stero} Mezur-unajo pri la lignajo karpentala, kalorizala = un metro³\lingui{DEFIRS}
\artiklo{stertorar}%
\vorto{stertorar}\fako{netrans.} Agar ta bruiso ne-normala quan produktas, en la respir-organi, la paso di aero tra mukosi quin sekrecis la bronki, la pulmoni\lingui{SL}
\artiklo{stetoskopo}%
\vorto{stetoskopo} Sorto di tubo akustikala, quan onu uzas por la auskulti medikala\lingui{DEFIS}
\artiklo{stifa}%
\vorto{stifa} Inter-parenta pro la mariajo dato-duesma di un ek la gepatri\lingui{D}
\artiklo{stifto}%
\vorto{stifto}\fako{tekn.} Kelo qua asemblas la peci ye qui konsistas charniro, buklo, e c.\lingui{DR}
\artiklo{stigmato}%
\vorto{stigmato}\senco{1}{Marko neefacebla quan lasis plago, bruluro, vunduro}\senco{2}{\textit{(religio kristana)} Marko ek la kin vunduri di la korpo di Iesu Kristo, imprimita mirakle sur la korpo di ula santi}\senco{3}{Marko imprimita, per fera stampilo redigita per varmigo, sur la shultro-pelo di ula krimininti}\senco{4}{\textit{(bot.)} Korpo glandatra ube finas la stilo di pistilo, e transmisas a lu, por igar olu pasar aden la ovario, la parto fekundigiva di poleno}\lingui{DEFIRS}
\artiklo{stileto}%
\vorto{stileto} Poniardo, kun lamo streta, secione triangula o quadri-angula\lingui{EF}
\artiklo{stilistiko}%
\vorto{stilistiko} Regularo pri stilo, qua esas situebla inter « gramatiko » e « retoriko ». – \textit{(segun Bally)}. Havante ideo klara o sentimento, la arto expresar li per la moyeni quin disponas linguo e quaze inter-sinonima\lingui{DF}
\artiklo{stilo}%
\vorto{stilo}\senco{1}{\textit{(epoki antiqua)} Puncilo di qua un extremajo pintoza uzesis por trasar la literi sur vaxo-plako, e di qua la altra, platigite, uzesis por efacar}\senco{2}{\textit{(literaturo)} Maniero expresar redakte sua pensi, sua idei}\senco{3}{\textit{(arto)} Karaktero generala di la verki da artisto, skolo, yarcento, naciono}\senco{4}{\textit{(bot.)} Prolonguro filoforma di la somito di la ovario, quan finas la stigmato}\lingui{DEFRS}
\artiklo{stilobato}%
\vorto{stilobato}\fako{arkitekt.} Basamento piedestalo-forma qua suportas koloni\lingui{DEFIS}
\artiklo{stimular}%
\vorto{stimular}\fako{trans., ad}\senco{1}{Augmentar la ardoro, la agereso di ulu}\senco{2}{\textit{(medic.)} Augmentar la agado di la funcioni organala}\lingui{DEFIRS}
\artiklo{stipendio}%
\vorto{stipendio} Pensiono gratuita, grantita a lernanto en edukerio (liceo, skolo administrata da la guvernistaro, seminario)\lingui{DEIR}
\artiklo{stipito}%
\vorto{stipito}\fako{bot.} Parto di la planto, qua departas vertikale de la radiko, e portas la folii, la flori e la frukti\lingui{EFS}
\artiklo{stipo}%
\vorto{stipo}\fako{bot.} Trunketo di la arbori monokotiledona, sen-rama, qua finas per fasko de folii\lingui{L}
\artiklo{stiptika}%
\vorto{stiptika}\fako{medic.} Di qua la efikiveso e la saporo esas astriktiva\lingui{EFIS}
\artiklo{stipular}%
\vorto{stipular}\fako{trans.) (yuro-cienco} Enuncar kom kondiciono, en kontrato\lingui{DEFIS}
\artiklo{stivar}%
\vorto{stivar}\fako{trans.} Lokizar la kompozanti di kargajo en la holdo di navo\lingui{EIS}
\artiklo{stofo}%
\vorto{stofo} Texuro ek lano, silko, filo, e c., de qua onu facas vesti, tapeti, e c.\lingui{DEFI}
\artiklo{stoika}%
\vorto{stoika} Qua savas tolerar sua sufro sen plendar – ferme\lingui{DEFIRS}
\artiklo{stoko}%
\vorto{stoko}\fako{komerco} Quanto de vari qua esas en magazino, staplo. – \textit{(anke metaf.)}\lingui{DEF}
\artiklo{stolo}%
\vorto{stolo}\fako{liturgio katol.} Larja bendo de stofo kun (kom ornivo) tri kruci, quan la sacerdoto pasigas cirkum sua kolo kande lu celebras la meso, kande lu donas ad ulu la unciono extrema\lingui{DEFIS}
\artiklo{stolono}%
\vorto{stolono}\fako{bot.} Singla de la radiki qui quaze reptas sur la surfaco di sulo, ek la stipiti reptera di ula vejetanti (kom ex. : fragiero)
\artiklo{stomako}%
\vorto{stomako}\fako{anat.} Vicero, en la supra regiono di abdomino, parto di la tubo nutrivala, posho-forma, en qua la nutrivi transformesas a chimo\lingui{EFIS}
\artiklo{stono}%
\vorto{stono} Fragmento de silexo, de quarco, o de roko irgaltra\lingui{DE}
\artiklo{stopar}%
\vorto{stopar}\fako{trans.} Plenigar (aperturo, truo, lakuno) per enduktar ulo aden lu\lingui{DE}
\artiklo{storaxo}%
\vorto{storaxo}\fako{bot.} Planto de qua onu extraktas balzamo per incizo\latina{styrax officinalis}\lingui{DEFIS}
\artiklo{storo}%
\vorto{storo} Peco de mato, o de stofo, movebla adsupre od adinfre, per resorto, avan la fenestri di chambro, di veturo, por shirmar kontre la suno-radii\lingui{FS}
\artiklo{stoterar}%
\vorto{stoterar}\fako{netrans.} Parolar hezitoze, emendante sua jusa dico, (a) pro defekto fiziologiala, (b) pro ke onu ne ja savas parolar, (c) pro emoco o timideso\lingui{DE}
\artiklo{straba}%
\vorto{straba}\fako{fiziol.} Di qua la du okuli ne regardas sama-direcione\lingui{EFIRS}
\artiklo{strado}%
\vorto{strado} Voyo kun bordumo ek domi, muri – en urbo, en vilajo \textit{(tale ke onu dicas : en strado, ma sur voyo)}\lingui{DEI}
\artiklo{strandar}%
\vorto{strandar}\fako{netrans.}\senco{1}{\textit{(aludante navo)} Venar, pro acidenteso o volo, aden loko ube ne plus existas sat multa aquo pri vehar normale, ed ube lu permanas kun plu o min multa aquo sub o cirkum su, o mem sen aquo}\senco{2}{\textit{(metaf.)} Lasar su falar (adsur ulo)}\lingui{DE}
\artiklo{strangular}%
\vorto{strangular}\fako{trans.) (per}\senco{1}{Ocidar, pro tote impedar la respiro, per kompreso ed obstrukto di la respir-organi, klemante pincoforme sua manui cirkum la kolo}\senco{2}{\textit{(metaf.)} \textit{(pri afero qua prosperas, pri revolto, e c.)} Haltigar la iro, la kresko, la suceso}\lingui{DEFIS}
\artiklo{stranja}%
\vorto{stranja} Di qua la nekustumeso kontenas ulo quo esis neexpektebla, neprevidebla, e qua desquietigas kelkete\lingui{EFIR}
\artiklo{stranjera}%
\vorto{stranjera}\senco{1}{Di altra lando, di altra familio, di altra loko}\senco{2}{Qua nule relatas, nule relatis, kun la koncernato}\senco{3}{\textit{(medic.)} Qua enduktesis acidente, o formacesis ne-normale, en organismo fiziologiala}\lingui{EFIS}
\artiklo{strapadar}%
\vorto{strapadar}\fako{trans.} Submisar a la kastigo (olim ordinara) qua konsistas ye elevar til ula alteso-grado la tormentato, per kordo, e lasar lu falar bruske, violentoze\lingui{EFIS}
\artiklo{straso}%
\vorto{straso} Kemialajo qua kontrafaktas diamanto e la lapidi\lingui{DEFIRS}
\artiklo{strategio}%
\vorto{strategio}\senco{1}{La arto direktar la ensemblo ek la operaci militala qui esas apta igar vinkar lor batalio o kampanio}\senco{2}{\textit{(metaf.)} La arto direktar ensemblo de operaci (kom ex. : komercala, politikala)}\lingui{DEFIRS}
\artiklo{stratego}%
\vorto{stratego}\fako{epoko antiqua, en Grekia} Chefo di armeo, generalo
\artiklo{strato}%
\vorto{strato}\senco{1}{\textit{(geol.)} Extensajo uniforma, de sulo sedimenta, o de aquo (eventuale subsula), sur ula quanto de spaco}\senco{2}{\textit{(arkitekt.)} Rango horizontala de quadri qua, en edifico, sustenesas o da sulo, o da rango plu infra}\lingui{EFIS}
\artiklo{streko}%
\vorto{streko} Lineo trasita sur surfaco\lingui{DE}
\artiklo{strenar}%
\vorto{strenar}\fako{trans.} Igar (ulo) subisar esforco di kompreso o di tenso o di tordo\lingui{E}
\artiklo{streptokoko}%
\vorto{streptokoko}\fako{patol.} Mikrobo, multa-kaze patogena, qua, kande lu grupeskas kun sua sami, formacas kolonii kateneto-forma\lingui{DEFIRS}
\artiklo{streta}%
\vorto{streta} Di qua la larjeso esas tre poka kompare kun lua longeso\lingui{EIS}
\artiklo{stridar}%
\vorto{stridar}\fako{netrans.} Sonar harde ed akre\lingui{EFIS}
\artiklo{striglo}%
\vorto{striglo} Plako ek fero kun rangi inter-paralela de denti, muntita ye mancho ligna, ed uzata por netigar la pili di la kavali, di la muli, e c.\lingui{DFI}
\artiklo{strigo}%
\vorto{strigo}\fako{zool.} Ucelo noktala, kun grosa okuli quin cirkondas cirklo de plumi tenua\latina{strix flammea}\lingui{IL}
\artiklo{strikar}%
\vorto{strikar}\fako{netrans.} Cesar bruske sua kontrato-laboro e durar sua nelaboro, por obtenar, de sua employanto, kelke plusa sua-profite\lingui{DE}
\artiklo{striknino}%
\vorto{striknino}\fako{kemio} Alkalio venena, extraktita de la grani di la vomiko-nuciero\lingui{DEFIRS}
\artiklo{strikta}%
\vorto{strikta} Qua ne lasas libera – nek cis, nek trans la limito determinita\lingui{DEFIS}
\artiklo{stringero}%
\vorto{stringero}\fako{arkitekt. navala} Super la holdo, singla ek la tenili qui mantenas la varangi\lingui{DER}
\artiklo{strio}%
\vorto{strio} Lineo trasita o gravita sur surfaco per akutajo; kanelo tre dina\lingui{EFIS}
\artiklo{stripo}%
\vorto{stripo} Klapo benda, ledra o stofo, qua pasas sub la shuo, ed esas ligita an la du lateri infra di pantalono, di getro tale ke lu impedas olca movar adsupre\lingui{DR}
\artiklo{strofo}%
\vorto{strofo} Singla ek la dividuri inter-egala di lirikajo\lingui{DEFIRS}
\artiklo{strofoido}%
\vorto{strofoido}\fako{matem.} Kurvo di la grado triesma, definita geometriale ye maniero simpla, e di qua ula quanto de karakterizivi esas remarkinda\lingui{DEFIS}
\artiklo{strokar}%
\vorto{strokar}\fako{netrans.} Efektigar un ek la agi di ensemblo\lingui{E}
\artiklo{stroncio}%
\vorto{stroncio}\fako{kemio} Korpo simpla, metala, di qua la atom-pezo esas 87,6\simbolo{Sr}\lingui{DEFIRS}
\artiklo{stropo}%
\vorto{stropo}\fako{navig.} Kordego, kun glito-nodo ye la du extremaji, e per qua onu embracas la kargaji por embarkar o des-embarkar li\lingui{DEFIRS}
\artiklo{strucho}%
\vorto{strucho}\fako{zool.} Stelt-ucelo, kun longa kolo, ali rudimenta, e la ucelo maxim granda, remarkinda pro lua devoremeso e la forteso di lua digesto-sistemo\latina{struthio camelus}\lingui{DEFI}
\artiklo{strukturo}%
\vorto{strukturo}\senco{1}{Maniero ye qua edifico esas konstruktita}\senco{2}{\textit{(metaf.)} Maniero ye qua esas dispozita la parti di ensemblo}\lingui{DEFIRS}
\artiklo{strumo}%
\vorto{strumo}\fako{patol.} Tumoro sen-dolora, ye la parto avana di la kolo\lingui{F}
\artiklo{studento}%
\vorto{studento} La adolecanto qua lernas en universitato, en fakultato\lingui{DEFIRS}
\artiklo{studiar}%
\vorto{studiar}\fako{trans. e netrans.}\senco{1}{Explorar atencoze, minucioze (ulo) por determinar olua karaktero, signifiko}\senco{2}{\textit{(metaf.)} Explorar (la cirkonstanci, la kunjunteso di la eventi) ante agar}\senco{3}{Explorar atencoze e minucioze (verko) por intertar lu, imitar lu, e c.}\lingui{DEFIS}
\artiklo{stufar}%
\vorto{stufar}\fako{trans.} Koquar (karno-bloko, pultruno) en recipiento klozita (*kluza) por impedar la eskapo di la vaporo e di la odoro – koquo kun acesoraji (fungi, keneli, e c.)\lingui{EFIS}
\artiklo{stuko}%
\vorto{stuko} Induturo qua imitas marmoro, quan onu aplikas sur muro, sur kolono, e quan onu polisas kande lu hardeskis per sikesko\lingui{DEFIRS}
\artiklo{stulo}%
\vorto{stulo} Sidilo, di qua la dors-apogilo esas min larja kam olta di fotelo, e sen brakii\lingui{DER}
\artiklo{stulta}%
\vorto{stulta} Rido-mokinda pro judiko-defekto quan lu unike ne perceptas, ne koncias\lingui{I}
\artiklo{stumpo}%
\vorto{stumpo}\senco{1}{\textit{(kirurg.)} La parto qua restas de membro amputita}\senco{2}{La centro di ula frukti, la pedo di ula legumi de qui onu deprenis la tota parto manjebla}\lingui{DE}
\artiklo{stuntar}%
\vorto{stuntar}\fako{trans.} Haltigar (lua) kresko normala\lingui{E}
\artiklo{stupida}%
\vorto{stupida} Di qua la spirito, la mento, la psiko esas quaze torporanta\lingui{DEFIS}
\artiklo{stupo}%
\vorto{stupo} La parto maxim grosiera ek la filamenti di kanabo o di lino\lingui{FIS}
\artiklo{stuporar}%
\vorto{stuporar}\fako{netrans.} Produktar sorto di torporo dum qua la pensofakultato aspektas quaze nihiligita\lingui{EFIS}
\artiklo{sturgo}%
\vorto{sturgo}\fako{zool.} Grosa fisho qua natas de maro aden la granda fluvii, e di qua la ovi uzesas por facar kaviaro\latina{acipenser sturio}\lingui{EFIS}
\artiklo{sturmar}%
\vorto{sturmar}\fako{netrans.} Verbo qua indikas ke atmosfero perturbesas da vento violentoza, pluvo o grelo, fulmini e tondri\lingui{DE}
\artiklo{sturno}%
\vorto{sturno}\fako{zool.} Pasero korno-rostra, kun plumaro nigra e blanka, qua vivas bande en la boski ed en la gardeni\latina{sturnus vulgaris}\lingui{FIS}
\artiklo{stuvo}%
\vorto{stuvo} Chambro, chambrego, en qua la temperaturo di la aero esas tre alta (multa-kaze : por desinfekto di la enklozaji)\lingui{EFIS}
\artiklo{su}%
\vorto{su}\fako{gram.} Pronomo reflektiva, qua referas la subjekto di la propoziciono di qua lu esas un ek la vorti, kande la subjekto esas « lu » o to quon « lu » reprezentas\lingui{DFIS}
\artiklo{sub}%
\vorto{sub}\fako{gram.}\senco{1}{Prepoziciono : 1. Indikas la situeso di kozo relate to quo esas plu adsupre, \textit{kontaktante} kun lu, en la sama direciono vertikala, (a) relate to quon lu suportas, (b) relate to quo kovras lu tote o parte. \textit{(anke metaf.)}}\senco{2}{Indikas la situeso di kozo relate to quo esas plu adsupre, \textit{sen kontaktar} lu, en la sama direciono vertikala. \textit{(anke metafore)}}\lingui{DEFI}
\artiklo{subalterna}%
\vorto{subalterna} Inferiora pri lua situeso en hierarkio\lingui{DEFIS}
\artiklo{subisar}%
\vorto{subisar}\fako{trans.} Tolerar malgre su\lingui{FI}
\artiklo{subita}%
\vorto{subita} Qua eventas neexpektite\lingui{FIS}
\artiklo{subjekto}%
\vorto{subjekto} \textit{(gram.)} e \textit{(logiko)}. Termino di la propoziciono di qua onu afirmas eso-maniero\lingui{DEFIRS}
\artiklo{subjuntivo}%
\vorto{subjuntivo}\fako{gram.) (en ula lingui nacionala} Modo di la konjugo-sistemo, qua reprezentas la stando o la ago expresata da verbo qua dependas de altra verbo\lingui{DEFIS}
\artiklo{sublima}%
\vorto{sublima}\senco{1}{Qua esas situita tre alte}\senco{2}{\textit{(metaf.)} Qua expresas lo bela segun lua formo maxim alta, partikulare en la domeni pensala, etikala o morala}\lingui{DEFIS}
\artiklo{submersar}%
\vorto{submersar}\fako{trans.} Plunjar komplete adsub la surfaco di aquo (di maro, di lago, di fluvio, di recipiento, e c.)\lingui{EFIS}
\artiklo{submisar}%
\vorto{submisar}\fako{trans.} Dependigar de la autoritato di ulu\lingui{EFIS}
\artiklo{submisionar}%
\vorto{submisionar}\fako{netrans., pri} Ofrar sua vari, sua laboro, a la administrantaro statala, municipala, e c., lor adjudiko-kompro\lingui{DF}
\artiklo{subordinar}%
\vorto{subordinar}\fako{trans., ad}\senco{1}{Dependigar de to quo esas en rango plu alta, superiora. – (gram.) Propoziciono subordinita : qua dependas, en la sintaxo, de altra propoziciono}\senco{2}{Dependigar de la ageso o faceso di ulo}\lingui{DEFIRS}
\artiklo{subornar}%
\vorto{subornar}\fako{trans.} Instigar ad ago kontre-etika, mala-mora, per lurivo, per charmivo\lingui{EFIS}
\artiklo{subrekargo}%
\vorto{subrekargo}\fako{komerco per navo-transporto} Prokuraciero di la persono o personi qua o qui proprietas la vari embarkita sur navo
\artiklo{subreptar}%
\vorto{subreptar}\fako{trans.) (yuro-cienco} Obtenar per moyeni frauda, kontrelega e trompa\lingui{EFIS}
\artiklo{subreto}%
\vorto{subreto}\fako{teatro} En komedio, la rolo (acesora) di chambro-servistino, di akompanantino, di qua la karaktero esas impetuoza, audacoza, tre alerta, tre inteligenta\lingui{DEFR}
\artiklo{subrogar}%
\vorto{subrogar}\fako{trans.) (yuro-cienco} Establisar en la rolo di ulu altru qua esis la tituliero normala\lingui{DEFIS}
\artiklo{subsidiar}%
\vorto{subsidiar}\fako{trans., pro. pri} Donacar pekunio ad ulu qua indijas extreme, e qua bezonas lu por exekutendi qui urjas. (Kom ex. Subsidiar navigado-kompanio pri la konstrukto di paketboti nova-tipa)\lingui{DEFIRS}
\artiklo{substanco}%
\vorto{substanco}\senco{1}{To ye quo konsistas la ento}\senco{2}{\textit{(filoz.)} To quo permanas kontinue en ento}\senco{3}{Materio ye qua kozo konsistas}\lingui{DEFIS}
\artiklo{substantivo}%
\vorto{substantivo}\fako{gram.} Vorto qua indikas la ento (kontraste kun la adjektivo, qua indikas la eso-maniero)\lingui{DEFIS}
\artiklo{substitucar}%
\vorto{substitucar}\fako{trans., ad} Pozar (persono, kozo) en la plaso di altra\lingui{DEFIS}
\artiklo{substituto}%
\vorto{substituto}\fako{yuro-cienco} Judiciisto, en tribunalo, qua komisesis por agar vice altra qua esas absenta od impedata\lingui{DEFIS}
\artiklo{subsumar}%
\vorto{subsumar}\fako{trans.) (filoz.}\senco{1}{Situar logikale individuo o speco aden la klaso o genero qua implikas li}\senco{2}{\textit{(che Kant)} Pensar intuice sub koncepto o kategorio di intelekto}\lingui{DEFIS}
\artiklo{subterfujo}%
\vorto{subterfujo} Moyeno nedireta di eskapar ulo embarasiva, evitar responso o justifikivo\lingui{EFIS}
\artiklo{subtila}%
\vorto{subtila}\senco{1}{Qua konsistas ye elementi extreme tenua}\senco{2}{\textit{(metaf.)} Di qua la nuanci pensala esas tenuega, e dicernebla nur desfacile}\senco{3}{Qua perceptas la nuanci pensala, mem tre desfacile dicernebla}\lingui{DEFIS}
\artiklo{subvencionar}%
\vorto{subvencionar}\fako{trans.} Grantar (ad ulu, ad ulo) parto ek la pekunio-stoko qua esas rezervita por spensi neexpektita, por helpar finance entraprezajo\lingui{DEFIS}
\artiklo{subversar}%
\vorto{subversar}\fako{trans.} Quaze renversar, en la spiriti, omna lego, omna regulo, omna principo, generale agnoskata kom bona\lingui{DEFIS}
\artiklo{sucedar}%
\vorto{sucedar}\fako{trans.} Venar sen interrupteso, sen diskontinueso, pos (dop) ulu, ulo\lingui{DEFIS}
\artiklo{sucesar}%
\vorto{sucesar}\fako{netrans., pri} Abutar fortunoze\lingui{EFIS}
\artiklo{suciar}%
\vorto{suciar}\fako{netrans., pri} Pensar tre ofte e kelke-desquiete pri la bona stando di (ulu, ulo)\lingui{F}
\artiklo{sucino}%
\vorto{sucino} Speco di rezino fosila, flava, diafana, odoroza, qua aquiras, per froto, ula karakterizivi elektrala\lingui{FI}
\artiklo{sudario}%
\vorto{sudario}\fako{epoki antiqua, en Roma} Tuko quan la Romani uzis por vishar sua vizajo sudorifanta\lingui{L}
\artiklo{sudo}%
\vorto{sudo} Un ek la quar punti kardinala, ta qua opozesas a nordo\lingui{DFIRS}
\artiklo{sudoro}%
\vorto{sudoro} Produkturo de la transpiro di la pelo, kondensita aden guteti\lingui{EFIS}
\artiklo{suficar}%
\vorto{suficar}\fako{netrans., a, pri} Furnisar sat multe (por satisfacar la bezoni, por saciar la deziri.)\lingui{EFIS}
\artiklo{sufixo}%
\vorto{sufixo}\fako{gram.} Afixo qua, en vorto, adjuntesas dop o pos la radiko, por determinar la naturo o senco di olca\lingui{DEFIRS}
\artiklo{suflar}%
\vorto{suflar}\senco{1}{\textit{(netrans.)} (a) Pulsar, sendar, aero per sua boko, (b) Igar ekirar de sua boko, per expiro, la aero respirita}\senco{2}{\textit{(trans.)} (a) Pulsar, sendar, aero adsur (ulo); (b) Pulsar, sendar, aero aden (ulo)}\senco{3}{\textit{(trans., ad ulu)} \textit{(teatro)} Dicar (ad ulu), tre deslaute, susure, kande lu falias memorar lo, to quon lu devas repetar laute}\lingui{DFI}
\artiklo{sufleo}%
\vorto{sufleo}\fako{koquarto} Mikra kuko ek pasto lejera, di qua la internajo esas vakua, koquita en mulduro, e qua manjesas kom entremeso\lingui{F}
\artiklo{sufokar}%
\vorto{sufokar}\senco{1}{\textit{(trans.)} Haltigar (dum tempo kelketa) la respiro (di ulu) per surprizo, dolorigo, iracigo, e c.}\senco{2}{\textit{(netrans.)} Koaktesar a cesar sua respiro (dum tempo kelketa) pro surprizeso, doloro, iraco, ridego, singluti, e c.}\lingui{EFIS}
\artiklo{sufragano}%
\vorto{sufragano}\fako{eklezio katol.} Episkopo subordinita ad arkiepiskopo\lingui{DEFIRS}
\artiklo{sufrajio}%
\vorto{sufrajio}\fako{liturgio katol.}\senco{1}{Prego, da eklezio, da la santi, por la religiani}\senco{2}{Prego, ye la fino di la laudi e di la vespri, por obtenar la interceso da la santi}
\artiklo{sufrar}%
\vorto{sufrar}\fako{netrans., de, pro} Tolerar sua doloro (fiziologiala o psikala) dum tempo plu o min multa\lingui{EFIS}
\artiklo{sufuziono}%
\vorto{sufuziono}\fako{medic.} Adfluego da sango sub la pelo
\artiklo{sugar}%
\vorto{sugar}\fako{trans.}\senco{1}{Aspirar, per la labii, la suko quan kontenas ula substanci}\senco{2}{Presar, per la labii, (a) (ulo, frukto, e c.) por extraktar olua suko; (b) \textit{(aludante infanteto)} la mamo di lua patrino, por extraktar lua lakto}\lingui{D}
\artiklo{sugestar}%
\vorto{sugestar}\fako{trans., ad} Venigar en la penso-domeno\lingui{DEFIS}
\artiklo{suko}%
\vorto{suko} Liquoro quan kontenas la substanco di frukti, karni, legumi, e quan onu povas extraktar per preso, macero, koquo, e c.\lingui{FIR}
\artiklo{sukombar}%
\vorto{sukombar}\fako{netrans., ad} Sinkar, falar, efike da kargajo tro pezoza, da koakto tro granda\lingui{EFIS}
\artiklo{sukro}%
\vorto{sukro}\senco{1}{Substanco qua saporas tre dolce, extraktata de plura vejetanti (precipue kano e betravo) e qua transformesas, efike da fairo, aden substanco kristalatra, solvebla en aquo}\senco{2}{\textit{(kemio)} Omna substanco qua esas apta fermentacar e transformesar aden alkoholo e CO₂}\lingui{DEFIRS}
\artiklo{sukubo}%
\vorto{sukubo}\fako{teol.} Demono noktala qua igas su aspektar quale muliero por koitar kun viro\lingui{DEFS}
\artiklo{sukusar}%
\vorto{sukusar}\fako{trans.}\senco{1}{Agitar bruske en omna lua parti}\senco{2}{\textit{(metaf.)} Agitar bruske (kozo) por desintrikar su de olu}\lingui{FI}
\artiklo{sulfo}%
\vorto{sulfo}\fako{kemio} Korpo simpla (metaloida), solida, sen-sapora, klare-flava, tre inflamebla, e qua, kande kombustata, exhalas odoro tre forta e sufokiva\simbolo{S}\lingui{EFIS}
\artiklo{sulko}%
\vorto{sulko}\fako{agrokultivo} Trancheo quan facas la soko di plugilo\lingui{IS}
\artiklo{sulo}%
\vorto{sulo} Extensajo sur qua su apogas la korpi, surface di Tero\lingui{EFIS}
\artiklo{sultano}%
\vorto{sultano}\fako{olim} La suvereno di Turkia\lingui{DEFIRS}
\artiklo{sumako}%
\vorto{sumako}\fako{bot.} Planto ek la familio di la terebinto, uzata da la mejisisti\latina{rhus}\lingui{DEFIRS}
\artiklo{sumaria}%
\vorto{sumaria} Rapid-agema; sen disipo di tempo e sen formalaji ne-utila\lingui{EF}
\artiklo{sumnar}%
\vorto{sumnar}\fako{trans., pri) (yuro-cienco} Imperar (ad ulu) ke lu agez to quon onu havas la yuro postular de lu, minacante lu ye la puniso a qua lu su expozus se lu refuzus\lingui{EF}
\artiklo{sumo}%
\vorto{sumo} Quanto ek du o plura quanti quin onu adicionis\lingui{DEFIRS}
\artiklo{sundio}%
\vorto{sundio} La dio unesma di la semano – konsakrita, da la eklezio kristana, a repozo ed a la kulto di la deo kristana, memorige di la rezurekto di Iesu Kristo\lingui{DE}
\artiklo{suno}%
\vorto{suno}\senco{1}{Astro lumifanta, qua esas la centro di la planetaro di qua Tero esas parto}\senco{2}{\textit{(metaf.)} To quo brilas quale Suno}\lingui{DE}
\artiklo{suolo}%
\vorto{suolo}\senco{1}{\textit{(anat.)} Suba parto di la hufo, che ula animali (kavalo, asno, mulo, cervo, e c.), plako korna inter la bordumo infra di la parieto, la forketo e la apogarki qui cirkondas lu}\senco{2}{Parto di la pedovesto, qua esas sub la plando, por protektar ica}\lingui{DEIS}
\artiklo{supear}%
\vorto{supear}\fako{netrans.} Repastar ye kloko (generale) kelke pos nokto-mezo, kande onu retrovenas de festo, de balo, e c.\lingui{DEF}
\artiklo{super}%
\vorto{super} Prepoziciono : Qua esas plu alta kam la suprajo e sen kontakto kun ica. – \textit{anke metaf}\lingui{EFIS}
\artiklo{superba}%
\vorto{superba} Qua estimas su tante ke lu situas su super sua kun-homi ed, pro lo, exajeras sua revendiki sociala, mondumala\lingui{FIS}
\artiklo{superfetaciono}%
\vorto{superfetaciono}\fako{fiziol.} Formaceso di feto duesma super feto unesma, konceptita antee\lingui{DEFIS}
\artiklo{superflua}%
\vorto{superflua} Qua esas ultre lo necesa, exter lo utila\lingui{EFIS}
\artiklo{superioro}%
\vorto{superioro} Persono qua esas super altra en hierarkio, e qua havas la yuro obediigar su da ta altra\lingui{DEFIS}
\artiklo{superlativo}%
\vorto{superlativo}\fako{gram.} La grado maxim alta di la eso-maniero quan expresas epiteto\lingui{DEFIS}
\artiklo{supernumerario}%
\vorto{supernumerario} Employato quan administranti, komerco-firmo, engajas suplemente, varte di la vako di ofico che su\lingui{DEFIS}
\artiklo{superstico}%
\vorto{superstico}\senco{1}{Kredo, praktiko a qua religio miskomprenata donas karaktero sakra}\senco{2}{Eceso di skrupulo pri la egardo di ula kozi}\lingui{EFIS}
\artiklo{supino}%
\vorto{supino}\fako{gram. Latina} Tempo di infinitivo qua esas analoga, pri la formo, a participo pasintala, e, pri la senco, a substantivo obtenita de la radiko di verbo\lingui{DEFIS}
\artiklo{suplantar}%
\vorto{suplantar}\fako{trans.} Eviktar ulu e substitucar su a lu\lingui{EFIS}
\artiklo{suplear}%
\vorto{suplear}\fako{trans.} Remplasar fristoze, pleante la sama rolo, (to quo ne povas, ta qua ne povas, plear sua rolo)\lingui{EFIS}
\artiklo{suplemento}%
\vorto{suplemento} To quo donesas o pagesas ultre lo debata\lingui{DEFIS}
\artiklo{suplikar}%
\vorto{suplikar}\fako{trans.) (pri, por} Pregar ye maniero tre humila e tre insistoza, ye la nomo di ulo kara e sakra\lingui{EFIS}
\artiklo{supo}%
\vorto{supo} Buliono quan onu varsas sur pano-lonchi dina, disho per qua komencas, ordinare, la repasto\lingui{DEFIRS}
\artiklo{suportar}%
\vorto{suportar}\fako{trans.}\senco{1}{(a) Havar sur su (kozo di qua onu portas la tota pezo) talo ke onu impedas olua falo; (b) \textit{(metaf.)} Havar sur su la tota kargajo, la tota embaraso di (ulo)}\senco{2}{(a) Recevar sur su sen shancelar de lo (la tota esforco … a ulo); (b) Recevar la tota efekto de ulo sen shancelar e lo}\lingui{EFIS}
\artiklo{supozar}%
\vorto{supozar}\fako{trans.} Admisar kom realigita en la tempo-punto di sua dico (to quon onu imaginas esar tala)\lingui{EFIS}
\artiklo{supozitorio}%
\vorto{supozitorio}\fako{medic.} Kono oblonga ek substanco medikamenta solida, quan onu igas penetrar aden la anuso por eventigar kaki\lingui{DEFIS}
\artiklo{supre}%
\vorto{supre} En objekto vertikala : en la parto qua distas maxime de lua pedo, de lua bazo. \textit{(anke metaf.)}\lingui{I}
\artiklo{supresar}%
\vorto{supresar}\fako{trans.} Desaparigar, partikulare per impedar lua manifesto\lingui{EFIS}
\artiklo{sur}%
\vorto{sur} Prepoziciono qua indikas la situeso di kozo relate a la kozo qua esas an lu, interkontakte, en sama direciono vertikala : (a) relate to quo portas lu; (b) relate to quon lu kovras o tegas
\artiklo{surda}%
\vorto{surda} Di qua la oreli ne perceptas, o perceptas nur tre desfacile, la soni, pro la desaparo, o la feblesko tre granda, di la audo-senso\lingui{EFIS}
\artiklo{surfaco}%
\vorto{surfaco}\senco{1}{Extensajo extera di korpo}\senco{2}{\textit{(geom.)} To quo limitizas korpo en spaco, ed havas nur du dimensioni : longeso e larjeso}\lingui{EFIS}
\artiklo{surjeto}%
\vorto{surjeto}\fako{sutado} Sut-uno ube la agulo trairas du stofo-peci pozite bordo-an-bordo, igante sempre la filo pasar super la du bordi unionita\lingui{FI}
\artiklo{surkruto}%
\vorto{surkruto} Kaulo hachita, qua fermentacis en sal-aquo\lingui{D}
\artiklo{surmenar}%
\vorto{surmenar}\fako{trans.} Fatigegar kavalko-bestio, charjo-bestio, e c. per igar lu irar tro rapide o dum tempo tro multa. \textit{(anke metafore)}\lingui{F}
\artiklo{surmuleto}%
\vorto{surmuleto}\fako{zool.} Maro-fisho\latina{Mullus surmuletus}\lingui{DEF}
\artiklo{suro}%
\vorto{suro}\fako{anat.} En la gambo, ta parto mola quan formacas, ye la extremajo, dope, la saliajo ek la muskuli jemela e la muskulo solea\lingui{L}
\artiklo{surogato}%
\vorto{surogato} To quo povas remplasar altra substanco (farmaciala, tintala, e c.) pro havar ula karakterizivi analoga, ula efekti analoga\lingui{DEIR}
\artiklo{surpliso}%
\vorto{surpliso}\fako{religio krist.} Vesto ek telo plisuroza quan *weras, sur sua sutano, la sacerdoti (en la konfeseyo ed en la procesioni funerala), la diakoni, la kantori\lingui{EFS}
\artiklo{surprizar}%
\vorto{surprizar}\fako{trans.}\senco{1}{Prenar (ulu, ulo) per arivar ne-expektite}\senco{2}{Trublar (ulu) per arivar ne-expektite}\senco{3}{Frapar la psiko per ulo ne-expektita}\lingui{EFIS}
\artiklo{surtuto}%
\vorto{surtuto} Vesto qua konservas la kaloro di la korpo, quan la viri *weras sur sua frako, redingoto\lingui{EFIS}
\artiklo{surveyar}%
\vorto{surveyar}\fako{trans.} Sorgar per regardar kun atenco sen-interrupta, time di defekto, di eroro, di neglijo, e c.\lingui{EFI}
\artiklo{sus!}%
\vorto{sus!} Interjeciono qua uzesas por ecitar ad inter-defenso, a persequo
\artiklo{suskriptar}%
\vorto{suskriptar}\fako{trans., pri} Engajar su pri pagor parto (di ulo)\lingui{DEFIS}
\artiklo{suspektar}%
\vorto{suspektar}\fako{trans., pri}\senco{1}{(a) Konjektar, de ula aspekti, ke (ulu) esas kulpozo; (b) Admisar la existo, pro ula aspekti, (di lo mala pri qua onu konjektas ke ulu esas la autoro)}\senco{2}{Konjektar (ulo) de ula aspekti}\senco{3}{\textit{(metaf. Dicernar instinte)}}\lingui{EFIS}
\artiklo{suspendar}%
\vorto{suspendar}\fako{trans.} Tenar super sulo (persono, kozo) tale ke lu pendas\lingui{DEFIS}
\artiklo{suspensar}%
\vorto{suspensar}\fako{trans.} Revokar, dure di nur kelka tempo, de lua oficiereso\lingui{FI}
\artiklo{suspensorio}%
\vorto{suspensorio} Bandajo qua sustenas la testikuli\lingui{DEFIS}
\artiklo{suspensoro}%
\vorto{suspensoro}\fako{biol.} Ligamento qua tenas suspendita
\artiklo{sustenar}%
\vorto{sustenar}\fako{trans.} Subtenar (ulo, ulu) per portar parto de lua pezo\lingui{EFIS}
\artiklo{sustracionar}%
\vorto{sustracionar}\fako{trans.) (aritm.} Deprenar nombro de altra, por determinar lia interdifero\lingui{DEFIS}
\artiklo{sustraktar}%
\vorto{sustraktar}\fako{trans., de}\senco{1}{Deprenar (ulo) de ulu, por ke lu ne povez uzar olu}\senco{2}{Deprenar (ulo qua apartenas ad ulu) cele, per habileso o fraudo}\lingui{FIS}
\artiklo{susurar}%
\vorto{susurar}\fako{trans., ulo, ad} Parolar deslaute, preske an la orelo (di ulu)\lingui{FI}
\artiklo{sutano}%
\vorto{sutano} Vesto butonagita de supre ad infre, qua kovras til la pedi, quan olim *weris la mediki, la sacerdoti, e quan nur ici duras *werar\lingui{DEFIS}
\artiklo{sutar}%
\vorto{sutar}\fako{trans.} Asemblar, ligar, per filo, kordeto, quan onu igas trapasar per agulo, per puncilo\lingui{EFI}
\artiklo{suverena}%
\vorto{suverena}\senco{1}{Qua havas la autoritato suprega en la stato}\lingui{DEFIS}\subvorto{politiko-yuro}{la suvereno}{La persono (personi) en qua (qui) rezidas la autoritato suprega leg-exekutigala (la rejo, en monarkio; la populo, en demokratio; la grandi, en oligarkio)}{2}{}
\artiklo{swabro}%
\vorto{swabro}\fako{navo} Fasko de fili kanaba, por lavar e vishar la ponto di navo\lingui{DER}
\artiklo{swichar}%
\vorto{swichar}\fako{trans.) (elektro} Chanjar la direciono di la \textit{elektro}-fluo, per aparato qua posibligas agar lo sen modifiko di la instaluro, di la filo-reto\lingui{E}
\artiklo{*swobo}%
\vorto{*swobo} Ligno-stipito, kartono-stipito, provizita, an lua extremaji, ye tufo piratra de kotono – uzata por skrapar ed uncionar la pelo di homo od animalo domestika\lingui{E}
\end{multicols}
\sekciono{T}
\begin{multicols}{2}
\artiklo{ta}%
\vorto{ta}\fako{gram.} Abreviuro di « ita »
\artiklo{tabako}%
\vorto{tabako}\senco{1}{\textit{(bot.)} Planto ek la familio « solanei »}\senco{2}{Folio de ta planto, preparita por snifleso, fumeso (*smokeso) o mastikeso}\latina{nicotiana}\lingui{DEFIRS}
\artiklo{tabano}%
\vorto{tabano}\fako{zool.} Insekto diptera, qua pikas la animali tante profunde ke lia sango ekfluas\latina{tabanus}\lingui{IL}
\artiklo{tabelo}%
\vorto{tabelo}\senco{1}{Panelo de ligno, indutita (maxim-multa-kaze) ye farbo nigra o obskur-verda, e sur qua onu skribas o trasas figuri per kreto-krayono}\senco{2}{Panelo ek ligno, kupro, e c., kadro sur qua tensesis telo – sur qua onu piktas}\senco{3}{Listo di qua la elementi (*itemi) esas ordinita metodoze (kom ex. : tabelo de logaritmi)}\lingui{DEFIRS}
\artiklo{tabernaklo}%
\vorto{tabernaklo}\senco{1}{\textit{(antique, che la Judi)} Tendo en qua esis enklozita la « archo di kontrato » til ante kande la templo esis konstruktita}\senco{2}{\textit{(liturgio katol.)} Recevuyo en qua esas enklozita la santa ciborio, super la tablo altara}\lingui{DEFIS}
\artiklo{tabeto}%
\vorto{tabeto}\fako{patol.} Morbo quan karakterizas precipue la ne-interkoordineso di la movi, kun reteno di la forteso di la muskuli
\artiklo{tablero}%
\vorto{tablero} FS. \textit{(en automobilo)}. Parieto, maxim-multa-kaze tola, qua separas la kapoto (o : spaco rezervita por la motoro) de la parto rezervita por la pasajeri. (Normale la automobilisto muntas sur la tablero la organi acesora : la exhaustilo, la pumpilo di lubrifiko, la klaxono, e c.)\lingui{FS}
\artiklo{tablo}%
\vorto{tablo} Areo plata ek ligno, petro, e c., subtenata da un o plura pedi, adsur qua onu pozas la objekti (por manjar, skribar, Laborar, ludar, e c.)\lingui{EFI}
\artiklo{tabulo}%
\vorto{tabulo} Planketo de ligno, en biblioteko, armoro, etajero, qua interseparas horizontale e superpoze la parti diversa di la kontenajo\lingui{FIS}
\artiklo{tabureto}%
\vorto{tabureto}\senco{1}{Sidilo kun quar pedi, sen brakii nek dors-apogilo}\senco{2}{Mikra suportilo, di qua la formo esas analoga, sur qua onu pozas sua pedi kande onu sidas sur stulo}\lingui{DFRS}
\artiklo{tacar}%
\vorto{tacar}\fako{netrans., pri}\senco{1}{Silencar, neparolar, pri (ulo, ulu)}\senco{2}{Permanar silencoza}\lingui{EFIS}
\artiklo{tafto}%
\vorto{tafto} Stofo ek silko, uniforma e briloza\lingui{DEFIRS}
\artiklo{tago}%
\vorto{tago} Peceto de metalo, ye la extremajo di laco, di kordono\lingui{E}
\artiklo{takigenezo}%
\vorto{takigenezo}\fako{biol.} Formo di la kresko en qua ula fazi embrionifala esas kondensita o supresita\lingui{DEFIS}
\artiklo{takigrafar}%
\vorto{takigrafar}\fako{trans.}\senco{1}{\textit{(en la epoki antiqua, en Roma)} Uzar, por la enrejistro di lo dicata, la sistemo de stenografo-signi quan inventis Tiro, la sekretario di Markus Tulius Kikero (Cicero)}\senco{2}{\textit{(nia-epoke)} Mi-stenografar, t.e. uzar abrevio-sistemo qua posibligas notar ye rapideso-grado pasable plu granda kam per la skribo-sistemo normala, ma ne tam rapide kam per stenografo}\lingui{DEFIRS}
\artiklo{takimetrio}%
\vorto{takimetrio} Metodo qua posibligas demonstrar rapide la teoremi geometriala od algebrala precipua, per konkretigar la operaci quin eventigas la demonstro\lingui{DEFIRS}
\artiklo{takimetro}%
\vorto{takimetro} Aparato uzata pri mekaniko, por mezurar la rapideso-gradi angulala, o grandeso irg-altra qua esas funciono lineara di ta rapideso\lingui{DEFIRS}
\artiklo{tako}%
\vorto{tako}\fako{tekn.} Mikra peco de ligno qua subtenas la extremajo di traverseto; qua mantenas angulo di moblo; qua garnisas la kuseneto di pulio por ke olca ne tendencez a sterno; quan onu sinkas aden tero por uzesar kom repero-punto lor la rektigo di voyo; qua startas la naveto « fluganta » di texo-mashino, e c.\lingui{FI}
\artiklo{taktiko}%
\vorto{taktiko}\senco{1}{\textit{(armeo)} La arto situar e movigar, sur tereno, la trupi di armo-speci diversa}\senco{2}{\textit{(metaf.)} La maniero ye qua onu agas por atingar ta o ca skopo, por obtenar ta o ca rezultajo}\lingui{DEFIRS}
\artiklo{taktismo}%
\vorto{taktismo}\fako{biol.} Influo quan agas ula substanci kemiala od ula formi di energio (koloro, kaloro, foto, elektro) a la protoplasmo e precipue a la direciono, la translaco di la celuli libera\lingui{DEFIS}
\artiklo{takto}%
\vorto{takto} Sento delikata pri to quo konvenas agar o ne agar por ne shokar o shameto-redigar la personi altra\lingui{DEFIRS}
\artiklo{tala}%
\vorto{tala} Sama kam kozo, persono, pri qua onu jus parolis o quik parolos
\artiklo{talento}%
\vorto{talento} Valoro, ecelo, naturala od aquirita per exercado, en irga ek la generi\lingui{DEFIRS}
\artiklo{taliar}%
\vorto{taliar}\fako{trans.} Tranchar (ulo) segun formo determinita\lingui{EFIS}
\artiklo{talio}%
\vorto{talio}\fako{kemio} Korpo simpla, metalo min blanka kam arjento\simbolo{Tl}\lingui{DEFIRS}
\artiklo{taliono}%
\vorto{taliono}\fako{yuro-cienco} Puniso qua konsistas ye aplikar a la kulpinto la ipsa trakto quan lu subisigis da sua kunhomo\lingui{F}
\artiklo{talioro}%
\vorto{talioro} La persono qua facas la vesti di viri e di mulieri\lingui{EFI}
\artiklo{talismano}%
\vorto{talismano} Lapido, ringo, e c., sur qua trasesis ula signi misterioza, e qua egardesas (da ula personi) kom apta efikar supernature pro lia relato kun la stelari sub qui li gravesis\lingui{DEFIRS}
\artiklo{talko}%
\vorto{talko} Substanco lameloza, griza verdatre, grasatra tushale\simbolo{H₂Mg₃Si₄O₁₂}\lingui{DEFIRS}
\artiklo{talo}%
\vorto{talo}\fako{bot.} La organaro vejetantala, plu o min komplexa, di la planti minim supra pri organizeso : algi, fungi, e c.\lingui{DEFIS}
\artiklo{talono}%
\vorto{talono}\senco{1}{\textit{(anat.)} La dopa parto di la pedo, che la homo}\senco{2}{Parto di la pedo-vesto qua inkluzas la dopa parto di la homo-pedo}\senco{3}{\textit{(metaf.)} (a) Extremajo infra di ula objekti; (b) Marjino di la folieti detranchebla, qua duras sur parto di la folieto fixigita ye la registro, ye la kayereto, e qua uzesas kom kontrolo-moyeno}\lingui{DFIS}
\artiklo{talpo}%
\vorto{talpo}\fako{zool.} Mikra mamifero karnavida, qua exkavas galerii sub la sulsurfaco, e di qua la okuli esas extreme mikra\latina{talpa}\lingui{FISL}
\artiklo{taluso}%
\vorto{taluso} Tereno, kun pento tre inklinita, qua rezultas ordinare de terlaboro, e qua garnisas la bordo di trancheo, di fosato, e c.\lingui{DEFS}
\artiklo{talvego}%
\vorto{talvego}\fako{geogr.} Lineo qua sequas la fundo di valo, di la lito di fluvio, di rivero – inter-sekuro di du penti laterala\lingui{DEFI}
\artiklo{tam}%
\vorto{tam} Adverbo : sama-grada pri qualeso
\artiklo{tamarindo}%
\vorto{tamarindo}\fako{bot.} Frukto de la arboreto L. \textit{tamarindus}, ek la familio « leguminosi », di qua la shelo kontenas pulpo laxigiva, mi-purgiva\latina{tamarindus}\lingui{DEFIRS}
\artiklo{tamarisko}%
\vorto{tamarisko}\fako{bot.} Arboreto di qua la folii esas tenua, e la flori spikoza\latina{tamarix}\lingui{DEFIRS}
\artiklo{tamburino}%
\vorto{tamburino} Tamburo plu basa kam la tamburo ordinara, quan onu igas resonar per nur un vergeto
\artiklo{tamburo}%
\vorto{tamburo} Kesto cilindra di qua la du fundi esas kovrita ye felo tensita, sur un ek qui onu frapas per vergeto por igar lu resonar\lingui{EFIRS}
\artiklo{tamen}%
\vorto{tamen}\fako{konjunciono} Opoze ad ico; malgre lo dicita; malgre ta obstaklo\lingui{L}
\artiklo{tampar}%
\vorto{tampar}\fako{trans.} Igar (pavo, paliso, e c.) irar ad, o til, la fundo per frapar (lu) per grava bloko ek ligno\lingui{E}
\artiklo{tampono}%
\vorto{tampono}\senco{1}{Mikra bloko, garnisita ye linjo, felto, e c., quan onu imbibas per substanco liquida quan onu volas extensar, aplikar sur surfaco}\senco{2}{\textit{(skerm-arto)} Rondo de felo quan onu muntas tampone ye la extremajo di flereto por igar ica nevundiva}\senco{3}{\textit{(tekn.)} Stifto ligna, ye qua onu garnisas truo borita en muro, por fixigar en lu skrubo, klovo}\lingui{DEFS}
\artiklo{tamtamo}%
\vorto{tamtamo} Perkut-instrumento (che la populi di la esto-regiono di Azia), disko metala, kun rebordo, qua sonegas kande onu frapas lu\lingui{DEFIR}
\artiklo{tandem}%
\vorto{tandem} Indikas ke lo eventis pos lasir (ulu) expektar ta evento dum multa tempo\lingui{L}
\artiklo{tandemo}%
\vorto{tandemo}\senco{1}{Veturo kun du kavali jungita an singla latero di timono}\senco{2}{Biciklo por du personi qui kavalkas lu, la duesma dop la unesma}\lingui{DEFI}
\artiklo{tangar}%
\vorto{tangar}\fako{netrans.} \textit{(navig.)} \textit{(aludante navo an aquo)} Ocilar altarne de la pruo a la pupo, de la pupo a la pruo\lingui{F}
\artiklo{tangenta}%
\vorto{tangenta}\fako{geom.} Qua tushas lineo, surfaco, ye nur un punto\lingui{DEFIRS}
\artiklo{tanino}%
\vorto{tanino}\fako{kemio} Substanco astriktiva, qua existas en la kortico di querko, sumako, kastaniero, ed igas olu apta preparar la ledri\lingui{DEFIRS}
\artiklo{tanko}%
\vorto{tanko} Granda recipiento metala, klozita e klozebla hermetike, por la liquidi\lingui{DEFS}
\artiklo{tano}%
\vorto{tano} Kortico di querko, di sumako, di kastaniero, pulverigita, quan onu uzas por la preparo di la ledri\lingui{EF}
\artiklo{tanta}%
\vorto{tanta} Indikas quanto, grado, nombro, nedefinita o tre granda\lingui{FIS}
\artiklo{tantalo}%
\vorto{tantalo}\fako{kemio} Metalo duktila, tre elastika, fasonebla por la faco di skribo-plumi, siringi por injekti, subderma, e c.\simbolo{Ta}\lingui{DEFIRS}
\artiklo{tapar}%
\vorto{tapar}\fako{trans.) (tekn.} Filetizar la parto interna di cilindro mikra-diametra, per utensilo di qua la formo esas olta di skrubo ek stalo, kun aristi tranchiva e quaza koridori por la eskapo adextera di la spani – la fileti unesma esante trunkigita por povar enirar la truo quan onu boris kom laborpreparuro\lingui{E}
\artiklo{tapeto}%
\vorto{tapeto} Peco de stofo per qua onu kovras la muri interna, e sur qua facesis, ek lano, silko, e c., figuri, ornivi diversa, sur fundo qua igas saliar lia beleso\lingui{DEFIS}
\artiklo{tapiokao}%
\vorto{tapiokao} Fekulo de radiko di manioko, lavita e sekigita\lingui{DEFIS}
\artiklo{tapiro}%
\vorto{tapiro}\fako{zool.} Mamifero pakiderma, de la sudo-regioni di Amerika, zoologie vicina a porko, asno-statura\latina{tapirus}\lingui{DEFIRS}
\artiklo{tapiso}%
\vorto{tapiso} Peco de stofo, de lano, de silko, e c. quan onu extensas sur tablo, sur parqueto, e c. – \textit{(anke metafore)}\lingui{DFIS}
\artiklo{tapotar}%
\vorto{tapotar}\fako{trans.}\senco{1}{\textit{(medic.)} Masajar per mikra frapi qui intersequas rapide, agita per la bordo kubitala di la manuo}\senco{2}{\textit{(tekn.)} Uniformigar la repartiseso (di pasto) en mulduro per mikra frapi repetita quin recevas la tablo qua suportas la mulduro}\lingui{EF}
\artiklo{tarantelo}%
\vorto{tarantelo}\fako{muziko} Dans-ario, danso agata en la sudo-regioni di Italia, segun la movimento 6/8 tre rapida\lingui{DEFIRS}
\artiklo{tarantulo}%
\vorto{tarantulo}\fako{zool.} Araneo-sorto, tre ordinara en la regiono qua cirkondas Tarente (Italia), e di qua la morduro produktas, che la homi, acidenti nervala\lingui{DEFIRS}
\artiklo{tarda}%
\vorto{tarda}\fako{*retarda} Qua eventas (a) pos la tempo-punto stipulita, (b) pos la tempo-punto ye qua lo eventas normale\lingui{EFIS}
\artiklo{*tarda}%
\vorto{*tarda} Preske ye la fino di la periodo aludata
\artiklo{tardigrado}%
\vorto{tardigrado}\fako{zool.} Familio de mamiferi ungloza, qui marchas lente\lingui{DEFIRS}
\artiklo{tarifo}%
\vorto{tarifo}\fako{komerco} Tabelo qua indikas la preci di vari, di taski exekutenda, di dogano-taxi, e c.\lingui{DEFIRS}
\artiklo{tarino}%
\vorto{tarino}\fako{zool.} Ucelo kantera, sorto di kardelo, kun beko konatra e plumaro verdatra\latina{fringilla spinus}
\artiklo{tarjo}%
\vorto{tarjo}\fako{armeo, olim} Shildo quadrata, di qua un ek la anguli esas sesgita por lasar la lanco pasar\lingui{DEFIRS}
\artiklo{tarlatano}%
\vorto{tarlatano} Muslino kun texuro poke densa, ordinare klar-kolora\lingui{DEFIRS}
\artiklo{taro}%
\vorto{taro} Pezo di la barelo, poti, tegili, kesti qui kontenas varo\lingui{DEFIRS}
\artiklo{taroko}%
\vorto{taroko} Ludo-karto di qua la reverso esas ornita ye faki de pikturi griza\lingui{DEFR}
\artiklo{tarso}%
\vorto{tarso}\fako{anat.} Parto dopa di la ped-osto, qua konsistas ye sep osti\lingui{EFIS}
\artiklo{tartano-navo}%
\vorto{tartano-navo}\fako{navig.} Mikra navo (sur Mediteraneo) pontoza, un-masta, e kun seglo triangula (kun un ortangulo)
\artiklo{tartano-stofo}%
\vorto{tartano-stofo} Stofo ek lano, kun figuri ek quadrato-frami, di qui la kolori alternas
\artiklo{tartaro}%
\vorto{tartaro}\fako{kemio}\senco{1}{Sedimento sala qua diferenciesas pokope ek la liquori vinatra, e adheras a la parieti di la bareli, forme di krusto}\senco{2}{Sedimento blankatra o flavatra qua amaseskas an la infra parto di la denti e karias pokope lia emalio}\lingui{EFIS}
\artiklo{tarto}%
\vorto{tarto} Kuko, ronda (maxim-multa-kaze), ek fundo plata de pasto, kun rebordo, e kovrita ye frukti, konfitaji, kremo, e c.\lingui{DEFIRS}
\artiklo{tasko}%
\vorto{tasko}\senco{1}{To quon onu obligesas laborar, facar, agar}\senco{2}{Quanto de laboro quan onu su obligis exekutor en fristo e po preco konvencionita}\senco{3}{Laboro skribala quan docisto donas kom exekutenda da lua docarii exter la doceyo}\lingui{EF}
\artiklo{taso}%
\vorto{taso} Mikra vazo kun anso, ek porcelano, metalo, vitro, quan onu uzas por drinkar (kafeo, teo, chokolado, e c.)\lingui{DFIS}
\artiklo{tastar}%
\vorto{tastar}\fako{netrans.}\senco{1}{Palpar per la manui, en tenebro, por avancar, por trovar to quon onu serchas}\senco{2}{Probar interdiferante, pro ne savar quale agar apte por atingar la skopo quan onu determinis}\lingui{DFI}
\artiklo{tatuar}%
\vorto{tatuar}\fako{trans.} Trasar (sur la pelo) desegnuri, per farbi quin onu koaktas – per la pinto di instrumento – penetrar sub la epidermo\lingui{DEFIR}
\artiklo{taumaturgio}%
\vorto{taumaturgio} La arto di la taumaturgo\lingui{DEFIS}
\artiklo{taumaturgo}%
\vorto{taumaturgo} Facisto di mirakli. \textit{(vorto pejorativa)}\lingui{DEFIS}
\artiklo{taverno}%
\vorto{taverno}\senco{1}{\textit{(olim)} Butiko en qua onu iris konsumar drinkaji plu o min alkoholoza, manjar, fumar (*smokar)}\senco{2}{\textit{(nun)} Kafe-drinkerio, restorerio plu o min luxoza, ed (en ula kazi) ornita da objekti artala (freski, pikturi, statui)}\lingui{EFIS}
\artiklo{taxar}%
\vorto{taxar}\fako{trans., ye} Enskribar (ulu) sur la registro de la impostebli kom debanta pagor pekunio-quanto (de…)\lingui{DEFR}
\artiklo{taximetro}%
\vorto{taximetro}\fako{automobilo} Kontilo specala, uzata sur la fiakri, e qua indikas la pekunio-quanto pagenda, kande la transporto jus finabos, proporcione ye la quanto de kilometri parkurita\lingui{DEFIS}
\artiklo{*taxio}%
\vorto{*taxio} Fiakro automobila, taximetroza\lingui{DEFIRS}
\artiklo{taxuso}%
\vorto{taxuso}\fako{bot.} Arboro verda, ek la familio « koniferi »\latina{taxus}\lingui{DL}
\artiklo{tayo}%
\vorto{tayo}\fako{anat.} Parto streta di la homo-korpo, inter hanchi e la infra parto di la busto\lingui{FIS}
\artiklo{teatro}%
\vorto{teatro}\senco{1}{Edifico ube reprezentesas verki qui havas la formo di konversi}\senco{2}{Ta edifico kun ti qui direktas lu, la aktori qui reprezentas la verki teatrala, la dekoruri, e c.}\lingui{DEFIRS}
\artiklo{tebaido}%
\vorto{tebaido} Retreteyo, refujeyo, solitara\lingui{DFIS}
\artiklo{tedar}%
\vorto{tedar}\fako{trans.} Fatigar per arivir desoportune, o pro tro repetar, tro iterar sua dici\lingui{EIS}
\artiklo{tegar}%
\vorto{tegar}\fako{trans.} Envelopar plu o min multa ek la parti (di objekto)\lingui{L}
\artiklo{tegulo}%
\vorto{tegulo} Plaketo de terakoto, uzata por facar la tekti\lingui{IL}
\artiklo{tegumento}%
\vorto{tegumento}\fako{anat.} Tisuo organika adheranta, qua tegas ula parti di la animali, di la vejetanti\lingui{DEFIS}
\artiklo{teismo}%
\vorto{teismo} Kredo pri la existo di deo\lingui{DEFIRS}
\artiklo{teisto}%
\vorto{teisto} Persono qua kredas ke existas deo\lingui{DEFIRS}
\artiklo{tekniko}%
\vorto{tekniko} Ensemblo ek la procedi di arto\lingui{DEFIRS}
\artiklo{teknologio}%
\vorto{teknologio} Expliko di la procedi, di la termini qui koncernas propre ta o ca arto o cienco\lingui{DEFIRS}
\artiklo{teko}%
\vorto{teko}\fako{bot.} Arboro di India, ek la familio « verbenacei », di qua la ligno, extreme-densa, e harda, uzesas por konstruktar domi, navi, e c.\latina{tectonia}\lingui{DEFIRS}
\artiklo{tekto}%
\vorto{tekto}\fako{arkitekt.} Parto di la karpenturo somite di edifico, di domo, ek teguli, ardezo-plaki, palio, e c., qua kovras ta konstrukturo\lingui{FIS}
\artiklo{telefonar}%
\vorto{telefonar}\fako{trans., ad} Transmisar voce (ulo ad ulu) per uzar instrumento ek plako dina de fero pura qua, kande la voco-soni produktesas, vibras e vibrigas unisone, per filo elektrala, lamo simila qua produktas, sur la membrano di la timpano, vibri korespondanta e qui genitas soni identa\lingui{DEFIRS}
\artiklo{telegonio}%
\vorto{telegonio}\fako{biol.} Transmiso, a la genituro ek mariajo dato-duesma, di karakteri qui devenas de la patrulo unesma-mariaja\lingui{DEFIS}
\artiklo{telegrafar}%
\vorto{telegrafar}\fako{trans., ad} Sendar (texto) per aparato qua transmisas la depeshi preske instantale\lingui{DEFIRS}
\artiklo{telegramo}%
\vorto{telegramo} Depesho transmisita per telegrafo\lingui{DEFIRS}
\artiklo{teleologio}%
\vorto{teleologio}\fako{biol.} Doktrino segun qua la enti esas organizita por la skopo a qua li destinesis\lingui{DEFIRS}
\artiklo{telepatio}%
\vorto{telepatio}\fako{metapsik.} Nomo per qua onu indikas ula senti (sen-sala o psikala) de persono, pri ulo reala quon agis sama-tempe altra persono, ma ye tanta disto ed en cirkonstanci tala, ke olua saveso da la persono unesme-mencionita aspektas neposibla\lingui{DEFIRS}
\artiklo{teleskopo}%
\vorto{teleskopo} Optik-instrumento, uzata por observar la objekti tre fora, per la reflektigo, sur speguli, di la imajo de ta objekti, e per plugrandigar ta reflekturo per lensi\lingui{DEFIRS}
\artiklo{telo}%
\vorto{telo} Texuro de fili ek lino o kanabo
\artiklo{telofazo}%
\vorto{telofazo}\fako{biol.} Fazo finala di la mitoso
\artiklo{telolecito}%
\vorto{telolecito}\fako{biol.} Dicesas pri la ovulo kun segmentesko partala (uceli, repteri, fishi) en qua la vitelo formaciva esas aparta de la vitelo nutriva
\artiklo{teluro}%
\vorto{teluro}\fako{kemio} Korpo simpla, metala, deskovrita en 1782 en la oro-mineyi di Transilvania\simbolo{Te}
\artiklo{temato}%
\vorto{temato}\fako{muziko}\senco{1}{Temo, motivo, melodiala di muzikajo}\senco{2}{\textit{(gram.)} Parto nevariiva di vorto, qua recevas la flexioni}\lingui{DEFIRS}
\artiklo{temerara}%
\vorto{temerara} Di qua la audaco, la sen-timeso, esas ne-reflektita\lingui{EFIS}
\artiklo{temo}%
\vorto{temo} Objekto di studio, propoziciono quan onu traktas, developas\lingui{DEFIRS}
\artiklo{temperamento}%
\vorto{temperamento}\fako{fiziol.} Konstituciono propra ye korpo organizita, qua dependas de la maniero ye qua la funcioni inter-equilibresas en lu\lingui{DEFIRS}
\artiklo{temperar}%
\vorto{temperar}\fako{trans.} Moderar per ta o ca pludolcigo, diminuto; igar min strikta, min rigoroza\lingui{DEFIS}
\artiklo{temperaturo}%
\vorto{temperaturo} Varmeso-grado di korpo, di loko, evaluebla per la senti quin lu produktas ye nia organi, o lua efiko a termometro\lingui{DEFIRS}
\artiklo{tempestar}%
\vorto{tempestar}\fako{netrans.} Furiar (aludante la elementi atmosferala quaze \textit{deskatenizita : uragano, fulmino, burasko}, e c.). – \textit{anke metafore}\lingui{EFIS}
\artiklo{templo}%
\vorto{templo} Edifico quan onu konsakris a la kulto publika di la deajo\lingui{DEFIS}
\artiklo{tempo}%
\vorto{tempo}\senco{1}{Duro di to quo havas komenco-punto e fino-punto}\senco{2}{\textit{(gram.)} Irga ek la flexioni diversa di un modo, qua indikas ke la ago o la subiso quan expresas la verbo esas prezenta, pasinta o futura}\senco{3}{\textit{(muziko)} Unajo dividala di la mezuro muzikala}\lingui{EFIS}
\artiklo{temporisar}%
\vorto{temporisar}\fako{netrans.} Ajornar, varte di tempo-punto oportuna\lingui{DEFIS}
\artiklo{temporo}%
\vorto{temporo}\fako{anat.} Regiono laterala di la kapo, inter la angulo di la okulo e la parto supra di la orelo\lingui{EFIS}
\artiklo{temprar}%
\vorto{temprar}\fako{trans.} Indutar ye farbo triturita en aquo e midilutita en gluo liquida, qua ne esas tam solida kam oleo-farbo\lingui{DEFIS}
\artiklo{tenaca}%
\vorto{tenaca}\senco{1}{Qua adheras forte ad ulo}\senco{2}{Di qua la parti inter-adheras forte}\lingui{EFIS}
\artiklo{tenalio}%
\vorto{tenalio} Utensilo fera, ek du branchi morsoza qui beas e, pose, interklemas por sizar e tenar forte\lingui{FIS}
\artiklo{tenar}%
\vorto{tenar}\fako{trans.} Havar en sua manui, inter sua brakii, e c., tale ke (lu) ne povas eskapar, forirar, falar, e c.\lingui{FIS}
\artiklo{tenaro}%
\vorto{tenaro}\fako{anat.} Saliajo ek la muskuli, ye la manuo-parto avana, extera, supra, apud la polexo\lingui{EFIS}
\artiklo{tencho}%
\vorto{tencho}\fako{zool.} Fisho qua vivas en la aquo sen-sala, zoologie vicina a karpo\latina{tinca}\lingui{EFIS}
\artiklo{tendencar}%
\vorto{tendencar}\fako{netrans., ad} Sentar en su ta eso-maniero naturale efike da qua ento atraktesas a skopo\lingui{DEFIRS}
\artiklo{tendino}%
\vorto{tendino}\fako{anat.} En la muskularo, fasko ek fibri, distingita de la muskulo per sua fibri qui esas ne-kontraktebla e sua koloro (brilante-blanka)\lingui{EFIS}
\artiklo{tendo}%
\vorto{tendo}\senco{1}{Paviliono ek stofo tensita, quan onu erektas en aero ne-enklozita, uzata kom shirmivo (precipue por kampo armeala)}\senco{2}{Kovrilo ek ledro, uzata sur kabrioleto od irgaltra veturo sen-tekta, automobila o kaval-tirata}\lingui{EFIS}
\artiklo{tendro}%
\vorto{tendro}\fako{teknol.} Veturo ligita a lokomotivo por transportar la aquo e la karbono ye qui ica alimentesas\lingui{DEFIRS}
\artiklo{tenebro}%
\vorto{tenebro}\senco{1}{Obskureso absolute kompleta, profunda}\senco{2}{\textit{(metaf.)} Eroro, ne-savo, qua impedas la vido di lo exakta, di la fakti}\lingui{EFIS}
\artiklo{tenera}%
\vorto{tenera} Qua havas por ulu afeciono qua abundas ye sentemeso e dolceso\lingui{EFIS}
\artiklo{tenio}%
\vorto{tenio} Vermo rimeno-forma, qua kreskas en la intestini di la homo e di ula animali\latina{taenia}\lingui{EFIS}
\artiklo{teniso}%
\vorto{teniso} Sporto quan onu praktikas per raketi e baloneti, sur tereno equipita specala\lingui{DEFIS}
\artiklo{tenono}%
\vorto{tenono}\fako{tekn.} Parto salianta, situita ye la extremajo di peco de ligno, de metalo, de marmoro, e c., por ajusteso a parto, a mortezo, korespondanta en altra peco qua esas asemblenda an la unesma\lingui{EF}
\artiklo{tenoro}%
\vorto{tenoro}\senco{1}{\textit{(muziko)} Voco virala qua iras de la unesma « do » di la alto, til la « sol » duesma di la violino}\senco{2}{Kantisto qua havas tala voco}\lingui{DEFIRS}
\artiklo{tensar}%
\vorto{tensar}\fako{trans.}\senco{1}{Igar rigida (korpo plu o min elastika) per igar lu plu longa, per tiro de un ek lua extremaji a la altra}\senco{2}{\textit{(pri horlojo murala)} Ri-suprigar la pezo-bloki kande li ter-venis}\lingui{EFIS}
\artiklo{tentakulo}%
\vorto{tentakulo}\fako{zool.} Apendico movebla ye qua esas provizita ula animali, e qua uzesas da li kom, precipue, organo di tushosenso\lingui{EFIS}
\artiklo{tentar}%
\vorto{tentar}\fako{trans.} Solicitar (ulu) ad agar ulo interdiktita, per olua atrakto\lingui{EFIS}
\artiklo{tenua}%
\vorto{tenua} Preske neprenebla pro lua mikregeso, dinegeso, delikateso\lingui{EFIS}
\artiklo{teo}%
\vorto{teo}\fako{bot.}\senco{1}{Folii di arboreto de Chinia, de Japonia, qua, per infuzo, donas drinkajo aromatoza, stimuliva}\senco{2}{Ta ipsa infuzuro}\latina{thea}\lingui{DEFIS}
\artiklo{teodiceo}%
\vorto{teodiceo}\fako{filoz.}\senco{1}{Traktato pri la yusteso di deo}\senco{2}{Parto di filozofio, qua traktas la existo e la atributi di deo}\lingui{DEFIRS}
\artiklo{teodolito}%
\vorto{teodolito}\fako{astron. e geom.} Instrumento di geodezio, quan onu uzas por relevar mapi, por mezurar la anguli reduktita a la horizonto, la disti zenitala e la azimuti\lingui{DEFIR}
\artiklo{teogonio}%
\vorto{teogonio} Genito di la dei; ensemblo de deaji di qui la kulto formacas la sistemo religiala di populo politeista\lingui{DEFIRS}
\artiklo{teokratio}%
\vorto{teokratio} Formo di guvernado segun qua la chefi di la naciono, de la kasto « sacerdoti », egardesas kom ministri di deo\lingui{DEFIRS}
\artiklo{teokrato}%
\vorto{teokrato} Singla ek la chefi di teokratio\lingui{DEFS}
\artiklo{teologala}%
\vorto{teologala} Dicesas pri singla de la tri vertui di qui deo esas la objekto ed un ek lua atributi kom motivo (kredo, espero e karitato)\lingui{FIS}
\artiklo{teologio}%
\vorto{teologio} Doktrino religiala pri la kozi deala\lingui{DEFIRS}
\artiklo{teorbo}%
\vorto{teorbo}\fako{muziko} Instrumento analoga a liuto, plu granda, kun du kapi : un por la kordi quin onu fingragas sur la mancho, la altra por la grosa kordi quin onu pinchas vakue\lingui{DEFIRS}
\artiklo{teoremo}%
\vorto{teoremo} Propoziciono ciencala quan demonstro igas evidenta\lingui{DEFIRS}
\artiklo{teorio}%
\vorto{teorio} Ensemblo ek la legi, la reguli atribuita a serio de fakti determinita, la principi quin onu propozas kom lia existo-kauzo\lingui{DEFIRS}
\artiklo{teozofio}%
\vorto{teozofio}\senco{1}{Sorto di iluminismo religiala}\senco{2}{Doktrino fondita sur intuico kontinue plu altagrada di la deajo, e qua posibligas efikar a la forci superiora}\lingui{DEFIRS}
\artiklo{tepalo}%
\vorto{tepalo}\fako{bot.} Singla ek la parti di la envelopilo florala
\artiklo{tepida}%
\vorto{tepida} Di qua la temperaturo esas inter kolda e varma\lingui{EFI}
\artiklo{teplico}%
\vorto{teplico} Chambrego vitro-mura, kalorizita, en qua onu shirmas la planti qui domajesus da temperaturo kolda\lingui{R}
\artiklo{terakoto}%
\vorto{terakoto} Argilo fasonita e pose bakita\lingui{DEFIRS}
\artiklo{terapio}%
\vorto{terapio}\fako{medic.} Parto di medicino qua havas kom objekto la kuraco e risanigo de la morbi\lingui{DEFIRS}
\artiklo{terario}%
\vorto{terario} Instaluro plu o min granda, en qua onu depozas tero, petri, stoni, planti, e c., por la eduko e la entrateno di kolektajo ek repteri, batraki, insekti, e c.\lingui{EF}
\artiklo{teraso}%
\vorto{teraso}\senco{1}{Platformo plu o min vasta, facita en gardeno, en parko, por regardar la peizajo til horizonto o por promenar – ek ter-amaso sustenata cirkonde da masonuro}\senco{2}{Platformo masonura, facita por regardar la cirkondajo o por promenar, sur un ek la etaji supra o sur la tekto-karpenturo di domo, di edifico}\lingui{DEF}
\artiklo{teratologio}%
\vorto{teratologio} Parto di la naturo-historio qua deskriptas la monstri\lingui{DEF}
\artiklo{tercerono}%
\vorto{tercerono} Persono qua naskis de pelblanko e de mestico
\artiklo{terceto}%
\vorto{terceto} Stanco de tri versi\lingui{DEFIR}
\artiklo{terciana}%
\vorto{terciana}\fako{patol.} Febro intermitanta, qua rieventas singla-tridio\lingui{DEFIS}
\artiklo{terciara}%
\vorto{terciara}\senco{1}{\textit{(geol.)} Dicesas pri un ek la granda faki ek la serio geologiala, qua inkluzas la strati eocena, oligocena, miocena e pliocena}\senco{2}{\textit{(patol.)} Di qua la simptomi aparas pos du etapi evidenta di morbo}\lingui{DEFIS}
\artiklo{terco}%
\vorto{terco}\senco{1}{\textit{(skermo)} Poziciono di la espado kande lu engajesas en la lineo extera, kun lua pinto vers-supre, la karpo rotacante de extere ad interne}\senco{2}{\textit{(muziko)} Intervalo inter la sekundo e la quarto}\senco{3}{\textit{(tempo)} La parto sisadekima di sekundo}\lingui{DEFIRS}
\artiklo{terebinto}%
\vorto{terebinto}\fako{bot.} Speco di pistachiero\latina{pistacia terebinthus}\lingui{DEFIRS}
\artiklo{tereno}%
\vorto{tereno} Extensajo de tero, de sulo, quan onu egardas kom apta por ta o ca uzeso\lingui{DEFIS}
\artiklo{teriako}%
\vorto{teriako}\fako{medic.} Elektuario kontre la morduro da la serpenti\lingui{DEFIRS}
\artiklo{teriero}%
\vorto{teriero}\fako{zool.} Mikra hundo quan olim onu uzis por chasar la animali qui rezidas en ter-trui\lingui{DEF}
\artiklo{terino}%
\vorto{terino} Vazo ek ligno, fayenco, porcelano, ceramiko vernisizita, ronda, un-peca, sen rebordo nek anso, nek mancho, kun fundo plata, uzata en la menaji\lingui{DEFI}
\artiklo{teritorio}%
\vorto{teritorio} Extensajo de lando qua formacas distrikto politikala\lingui{DEFIRS}
\artiklo{termi}%
\vorto{termi}\senco{1}{\textit{(epoki antiqua)} Establisuro publika por balni varm-aqua}\senco{2}{\textit{(nun)} Establisuro adube onu venas por kuracesar per aqui medicinala varma}\lingui{DEFIS}
\artiklo{termino}%
\vorto{termino}\senco{1}{Bloko petra qua indikas la limito di agro, di stadio, la frontiero}\senco{2}{\textit{(gram.)} En propoziciono : o la subjekto, o la verbo, o la atributo}\senco{3}{\textit{(logiko)} En silogismo : o la majoro, o la minoro, o la mezo}\senco{4}{\textit{(matem.)} Expresuro di quanto en raporto, en proporciono, en progresiono, e c.}\senco{5}{\textit{(mitol.)} La deajo – reprezentita per figuro di viro di qua la parto infra finas gaino-forme, skultita en bloko de petro – qua sekurigis la egardo di la limiti inter-agra, di la frontieri, e c.}\lingui{DEFIS}
\artiklo{termito}%
\vorto{termito}\fako{zool.} Nevroptero qua rodas interne la objekti ek ligno, lente e cele\latina{termes}\lingui{DEFRS}
\artiklo{termodinamiko}%
\vorto{termodinamiko} Parto di fiziko, qua havas kom studiajo la relati qui existas inter la fenomeni mekanikala e kalorifala\lingui{DEFIRS}
\artiklo{termoelektro}%
\vorto{termoelektro} Ensemblo ek la fenomeni elektrala e kalorala qui naskas de la varii od elektrala o kalorala\lingui{DEFIRS}
\artiklo{termografo}%
\vorto{termografo}\fako{fiziol.} Instrumento qua konsistas ye aero-termometro kun enrejistro-levero, ed uzata por mezurar la varii di temperaturo, e la duro di ta varii, en irga punto di la homo-korpo\lingui{DEIR}
\artiklo{termokemio}%
\vorto{termokemio}\fako{kemio} Fako di la kemio-cienco, qua traktas la determino di la quanto de kaloro qua pleas rolo en la kombini\lingui{DEFIRS}
\artiklo{termologio}%
\vorto{termologio} Fako di fiziko, qua traktas kaloro\lingui{DEFIS}
\artiklo{termometrio}%
\vorto{termometrio}\fako{fiziko} La mezurado di kaloro\lingui{DEFIRS}
\artiklo{termometro}%
\vorto{termometro} Aparato por la mezuro di la temperaturi, per la dilato o la kontrakto di merkuro, di alkoholo, e c. en tubo gradoza\lingui{DEFIRS}
\artiklo{termoskopo}%
\vorto{termoskopo}\fako{fiziko} Aparato por montrar la existo e la signo di difero pri la temperaturi reciproka di du medii, ma qua donas unike valori relativa pri la grandeso-grado di ta diferi\lingui{DEFIS}
\artiklo{termostato}%
\vorto{termostato}\fako{teknol.} Ensemblo teknikala qua posibligas la obteno, en klozajo (*kluzajo) di temperaturo konstanta
\artiklo{termotaktismo}%
\vorto{termotaktismo}\fako{biol.} Su-diplaso reaktala da protoplasmo, influe da kaloro\lingui{DEFIS}
\artiklo{termoterapio}%
\vorto{termoterapio}\fako{medic.} Kuraco di la morbi per kaloro\lingui{DEFIS}
\artiklo{termotropismo}%
\vorto{termotropismo}\fako{biol.} Influo, da kaloro, di la protoplasmo che la elementi tisua, e precipue di la direciono o la sinso di la kresko che la tisui qui kreskas\lingui{DEFIS}
\artiklo{ternara}%
\vorto{ternara}\fako{matem.) (pri sistemo di nombrifado} En qua tri unaji ek la klaso unesma formacas unajo ek la klaso duesma\lingui{DEFIS}
\artiklo{tero}%
\vorto{tero}\senco{1}{Parto solida di la globo quan ni habitas : (a) sulo qua suportas la homo, la animali, la planti, la urbi, e c.; (b) sulo qua produktas la vejetanti; (c) sulo aden qua onu sepultas la mortinti}\senco{2}{Materio solida ye qua konsistas nia globo; substanco tirita de ta materio ed uzata diversa-skope}\senco{3}{La ter-globo ipsa, kom planeto}\lingui{EFIS}
\artiklo{terorar}%
\vorto{terorar}\fako{netrans., pro, pri} Fremisar, parturbesar profunde en sua organismo, pro pavoro\lingui{DEFIS}
\artiklo{terpentino}%
\vorto{terpentino}\fako{kemio} Rezino liquida quan donas la arbori ek la familio « terebintacei », e di qua un elemento esas C₁₀H₁₆\lingui{DE}
\artiklo{terpleno}%
\vorto{terpleno}\fako{milit-arto} Platformo ek tero adjuntita, sustenata da parapeti masonura\lingui{F}
\artiklo{testaceo}%
\vorto{testaceo}\fako{zool.} Molusko di qua la korpo kovresas da konko\lingui{DEFIS}
\artiklo{testamentar}%
\vorto{testamentar}\fako{netrans., favore di}\senco{1}{Redaktar la akto per qua onu legacas sua havaji, e di qua la klauzi esas exekutenda pos la morto di la redaktanto}\senco{2}{\textit{(historio religiala)} Testamento : nomo di singla ek la du libri santa : che la Hebrei, la Olda Testamento; che la Kristani : la \textit{Nova Testamento}}\lingui{DEFIS}
\artiklo{testiero}%
\vorto{testiero} Harneso-peco sur la kapo di kavalo, an qua ligesas la morso e la brido
\artiklo{testikulo}%
\vorto{testikulo}\fako{anat.} Che la mamiferuli : korpo glandatra quan kontenas singla ek la burseti, qua sekrecas spermo\lingui{EFIS}
\artiklo{testo}%
\vorto{testo} Persono qua deklaras (maxim-multa-kaze : koram judiciisti) certifikar vidir od audir (ulo, ulu)\lingui{IS}
\artiklo{tetanar}%
\vorto{tetanar}\fako{netrans.} Subisar kontrakti konvulsala di uli ek sua muskuli – ed, en ula kazi, la rigidesko di omna sua muskuli – qui mortigas la subisanto ulakaze\lingui{DEFIS}
\artiklo{tetrado}%
\vorto{tetrado}\fako{biol.} Su-divido di la celulo patra di la poleno-grani aden quar celuli qui divenas quar poleno-grani inter-unionita en la kazi maxim multa
\artiklo{tetraedro}%
\vorto{tetraedro}\fako{geom.} Solido quar-edra\lingui{DEFIS}
\artiklo{tetraedroido}%
\vorto{tetraedroido} Solido di qua la formo esas kelke simila a tetraedro\lingui{DEFIS}
\artiklo{tetragono}%
\vorto{tetragono}\fako{bot.} Genero de « fikoidi », qua kontenas duadek speci en la regioni temperaturo-moderata, e quan onu manjas kom spinato-analoga o forme di supo\lingui{DEFIS}
\artiklo{tetralogio}%
\vorto{tetralogio}\senco{1}{\textit{(literaturo antiqua)} Ensemblo de quar peci : tri tragedii ed un dramato satirema, quan prizentis, lor la konkursi teatrala, la poeti tragediala di la Grekia antiqua}\senco{2}{Grupo de quar dialogi, che Platono}\senco{3}{\textit{(muziko)} Ensemblo de quar operi}\lingui{DEFIRS}
\artiklo{tetrarko}%
\vorto{tetrarko}\fako{epoki antiqua, en Grekia} Chefo di singla ek la quar gubernii di lando\lingui{DEFIRS}
\artiklo{tetraso}%
\vorto{tetraso}\fako{zool.} Granda urogalo\latina{tetrao urogallus}\lingui{FIL}
\artiklo{tetrasporo}%
\vorto{tetrasporo}\fako{biol.} Genero de algi verda, tipo di la familio « tetrasporacei », qui konsistas ye celuli qui permanas en grupi de du o quar
\artiklo{teurgio}%
\vorto{teurgio} Magio fondita sur relato asertita kun la spiriti cielala. \textit{(antonimo di nekromancio, magio infernala)}\lingui{DEFIRS}
\artiklo{texar}%
\vorto{texar}\fako{trans.} Transformar (substanco filigita : lano, kanabo, kotono, e c.) aden materio koheranta, flexebla e dina, per krucumar fili tensita horizontale (warpo-filo) kun fili quin onu igas pasar transverse (wefto)\lingui{EFIS}
\artiklo{texto}%
\vorto{texto}\senco{1}{La ipsa vorti di autoro, di lego, di akto (opoze a la komenti)}\senco{2}{Peci ek la libri santa quin la predikisti citas, komence di sermono, di homilio, kom aplikebla a la temo quan li intencas quik traktar}\lingui{DEFIRS}
\artiklo{tezo}%
\vorto{tezo} Propoziciono, asertata da ulu, e defensata da lu kontre la personi qui kontestas olua exakteso\lingui{DEFIRS}
\artiklo{tiaro}%
\vorto{tiaro}\senco{1}{Kapovesto alta, uzata olim che la Oriento-populi}\senco{2}{Kapovesto di la sacerdoto suprega, che la Hebrei}\senco{3}{\textit{(che la katoliki)} Kapovesto quan *weras la papo dum la ceremonii, boneto kun, kom ornivi, tri kroni, e, suprege, globo kun un krono}\lingui{DEFIRS}
\artiklo{tibio}%
\vorto{tibio}\fako{anat.} La plu grosa ek la du osti di la gambo, avan la peroneo\lingui{EFIS}
\artiklo{tifo}%
\vorto{tifo}\fako{patol.} Morbo kontagiema, quan karakterizas febro sen-cesa, la stando malada di la mukozi, e c.\lingui{DEFIRS}
\artiklo{tifoido}%
\vorto{tifoido}\fako{patol.} Morbo infektiva quan produktas mikrobo specifika (deskovrita da Eberth)\lingui{DEFIS}
\artiklo{tifono}%
\vorto{tifono}\fako{meteorol.} Tempesto-vento, qua forportas kun su – kun violento nerezistebla – to quon lu renkontras, sur-tere e sur-maro\lingui{EF}
\artiklo{tigmotropismo}%
\vorto{tigmotropismo}\fako{biol.} Influo, da la kontakto, di la direciono che la tisui kreskanta
\artiklo{tigro}%
\vorto{tigro}\fako{zool.} Animalo feroca, karnavida, ek la familio « felini », di qua la pilaro esas strioza e makuletoza\latina{felis tigris}\lingui{DEFIRS}
\artiklo{tikar}%
\vorto{tikar}\fako{netrans.}\senco{1}{\textit{(patol.)} \textit{(aludante la muskuli, precipue olti di la vizajo)} Kontraktesar konvulsatre, preske kustumale}\senco{2}{\textit{(metaf.)} Agar gesto kustumale}\lingui{F}
\artiklo{tiktakar}%
\vorto{tiktakar}\fako{netrans.) (aludante, precipue, horlojo} Produktar bruiseto bruska, ek du tempi intersucedanta qui iteresas uniforme\lingui{DEFS}
\artiklo{til}%
\vorto{til} Prepoziciono qua indikas la arivo (en tempo, en spaco) an termino quan onu ne transiras. Kande la limito esas ne punto, ma intervalo, lu sequesas da « exkluzite » od « inkluzite » (segun la nociono)\lingui{E}
\artiklo{tilbury}%
\vorto{« tilbury »} Sorto di kabrioleto lejera, sen tendo\lingui{EF}
\artiklo{tildo}%
\vorto{tildo}\fako{gram.} Signo diakritika, qua havas ica formo : \textasciitilde{} quan la Hispani lokizas sur la « n » kande li volas pronuncigar lu « ny' » – e quan la Portugalani lokizas sur « a » e « o » por pronuncigar li quale la F. « an » e « on », (t.e. la E. « ang, ong ») rispektive
\artiklo{tilio}%
\vorto{tilio}\fako{bot.} Arboro di qua uzesas : la floro, por facar infuzuri kalmigiva, e la kortico, por facar puteo-kordegi\latina{tilia}\lingui{FISL}
\artiklo{timar}%
\vorto{timar}\fako{trans.}\senco{1}{Tendencar evitor (ulu, ulo) (de qua onu expektas malajo)}\senco{2}{Judikar ulo (regretinda, desoportuna) kom eventonta plu o min balde}\lingui{EFIS}
\artiklo{timbalo}%
\vorto{timbalo} Mi-globo de kupro o latuno, kovrita per felo tensita quan onu klemas o desklemas per skrubi, por akordigar lu, e sur qua onu frapas per vergeti ek ligno harda o tegita ye felo-peco segun la sonoreso-grado quan onu deziras obtenor – muzik-instrumento qua uzesas en kavalrio ed en orkestro\lingui{EFIS}
\artiklo{timbro}%
\vorto{timbro}\fako{*tembro} Kaloto de metalo qua, frapate per marteleto, produktas sonuno pasable duroza\lingui{FS}
\artiklo{*timbro}%
\vorto{*timbro}\fako{postmarko} Peceto de papero imprimuroza (ordinare kun konturo dentetoza pro perforesir) quan onu glutinas sur la letro-kuverti por montrar ke la afranko-preco esas pagita. – Analoge pri la « timbri » fiskala, asekurala, e c.\lingui{FS}
\artiklo{timiano}%
\vorto{timiano}\fako{bot.} Planto odoroza, ek la familio « labiei », quan onu uzas en la faco di sauci\latina{thymus}\lingui{DEFIS}
\artiklo{timida}%
\vorto{timida} Qua indijas audaco pro desfidar su\lingui{EFIS}
\artiklo{timono}%
\vorto{timono} Longa stango de ligno, muntita ye la parto avana di veturo, di plugilo, ed an singla latero di qua tir-bestio (kavalo, bovo) esas jungita por apogar su ad olu por tirar, retenar, chanjar sua direciono, des-avancar\lingui{FIS}
\artiklo{timpano}%
\vorto{timpano}\senco{1}{\textit{(anat.)} Kavajo ye qua konsistas la meza parto di orelo, klozita, en la parto extera ed en la parto interna di la orelo, per membrani tensita}\senco{2}{\textit{(arkitekt.)} Spaco quan cirkondas la tri kornici di frontono, qua destinesas recevor basa reliefi, ornamenti}\senco{3}{\textit{(muziko)} Instrumento olima di muziko, garnisita ye kordi latuna, quan onu pleas per du vergeti ek ligno}\lingui{EFIRS}
\artiklo{timuso}%
\vorto{timuso}\fako{anat.} Korpo glandatra, biloba, situita dop la ster sternumo, an la parto infra di la kolo\lingui{DEFIS}
\artiklo{tindro}%
\vorto{tindro} Substanco sponjatra, extraktita de la agariko di la querko (fungo kun kortico lignatra) qua acendesas kontakte di cintilo\latina{ochropus fomentarius}\lingui{DE}
\artiklo{tineo}%
\vorto{tineo}\fako{zool.}\senco{1}{Mikra insekto lepidoptera, di qua la speco maxim difuzita esas ta qua devoras la grani}\senco{2}{Mikra insekto lepidoptera, noktala, di qua la larvo rodas la stofi la libri, e c.}\latina{tinea}\lingui{FISL}
\artiklo{tinio}%
\vorto{tinio}\fako{patol.} Morbo di la pelo haroza, che la homo\lingui{FISL}
\artiklo{tinklar}%
\vorto{tinklar}\fako{netrans.} \textit{(aludante klosho quan onu frapas sur nur un latero per la frapilo)} Produktar sonuni qui intersucedas lente\lingui{EFI}
\artiklo{tinlauro}%
\vorto{tinlauro}\fako{bot.} Arboreto ek la genero « viburni ». L. viburnum tinus\latina{viburnum tinus}
\artiklo{tino}%
\vorto{tino} Barelo di qua la fundo esas plu larja kam la parto supra uzata por la transporto di aquo, vito-beri lor la rekolto, e c.\lingui{DFIS}
\artiklo{tintar}%
\vorto{tintar}\fako{trans., ye, ad}\senco{1}{Impregnar (filo, lano, texuro, felo, e c.) ye substanco qua igas lu havar ta o ca koloro}\senco{2}{Igar stofo ja tintita, ma di qua la koloro esas velkinta, aquirar koloro fresha}\lingui{DEFIS}
\artiklo{tipo}%
\vorto{tipo}\senco{1}{Peco de ligno, de metalo, e c., qua havas stampuro qua esas destinita *reproduktar stampuro strikte sama}\senco{2}{Peco de gisuro di qua la stampuro formacas la literi imprimala}\senco{3}{Koncepto ideala di la traiti generala segun qua formacesis kozo}\lingui{DEFIRS}
\artiklo{tipografar}%
\vorto{tipografar}\fako{trans.} Imprimar per tipi\lingui{DEFIRS}
\artiklo{tipulo}%
\vorto{tipulo}\fako{zool.} Genero de insekti diptera, kun mikra korpo, ek la familio « tipulidi », kun longa gambi, tre dina e frajila\latina{tipula}\lingui{EFISL}
\artiklo{tirado}%
\vorto{tirado}\fako{teatro} To quon rolo recitas quaze « bloke »\lingui{DEFS}
\artiklo{tiraliar}%
\vorto{tiraliar}\fako{netrans.) (armeo} Pafar segun sua arbitrio, sen vartar la komando de sua chefi, kontre la enemiko\lingui{DF}
\artiklo{tirano}%
\vorto{tirano}\senco{1}{La persono qua, suprege-autoritatoza en la stato, uzas sua povo ye maniero opresanta}\senco{2}{\textit{(metaf.)} Persono, kozo, qua uzas opresante sua autoritato}\lingui{DEFIRS}
\artiklo{tirar}%
\vorto{tirar}\fako{trans., de, ek, per}\senco{1}{Venigar ula-direcione, per prenar (la persono, la kozo) ye un ek lua parti quan onu adduktas ad su : (a) esante sen-mova, (b) marchante ta-direcione}\senco{2}{Adduktar (persono, kozo) ek la loko ube lu esas inkluzita, retenata}\senco{3}{\textit{(pri kameno, furnelo, e c., alimentata per karbono o ligno)} Subisar en su ta aero-fluo per qua la aero necesa por la kombusto atraktesas a la herdo por remplasar la aero varmigita qua ekiras la tubo fluala, tranante kun su la fumuro}\lingui{FIS}
\artiklo{tiritano}%
\vorto{tiritano} Sorto de drapo, mi-lana e mi-lina\lingui{DFS}
\artiklo{tiroido}%
\vorto{tiroido}\fako{anat.} Glando an la parto infra di la laringo, sur la serio-unesma ringi di la trakeo\lingui{EFIS}
\artiklo{tirozimazo}%
\vorto{tirozimazo}\fako{biol.} Oxidazo, en la fungi di la genero L. \textit{russula}, qua povas oxidigar tirozino\latina{russula}
\artiklo{tirozino}%
\vorto{tirozino}\fako{biol.} HO. C₆H₄. CH₂. CH (NH₂) CO₂H, trovata en la hepato, en la spleno, en la pankreato
\artiklo{tirso}%
\vorto{tirso}\fako{mitol.} Emblemo di Bacchus, javelino an-cirkondata da hedero, pampri, e finanta forme di pino-kono, quan portis la bacchus-pastorini\lingui{DEFIRS}
\artiklo{tisuo}%
\vorto{tisuo}\fako{anat.} Materio koheranta ye qua konsistas la organi di la animali, di la vejetanti, asemblajo de fibri interplektita od inter-apud-pozita\lingui{EFI}
\artiklo{titano}%
\vorto{titano}\senco{1}{\textit{(mitol.)} Singla ek la deaji primitiva, qui preiris la Olimpani (t. e. Okeanos, Koyos, Krios, Hiperion, e c.)}\senco{2}{\textit{(kemio)} Metaloido tetratoma, di qua la karakteri situesas inter siliko e stano}\simbolo{Ti}\lingui{DEFIS}
\artiklo{titilar}%
\vorto{titilar}\fako{trans.} Submisar a tre lejera tusheti iterata, sur ula parto di la korpo\lingui{EFIS}
\artiklo{titrar}%
\vorto{titrar}\senco{1}{\textit{(teknol.)} Determinar la proporciono di oro, di arjento, unionita a la aloyuro che la metali precoza – la proporciono di silko, di lano, en texuro mixura}\senco{2}{\textit{(kemio)} Determinar la quanto (fixa) di la reaktivo quan kontenas dissolve ula volumino di liquoro}\lingui{DEF}
\artiklo{titulo}%
\vorto{titulo}\senco{1}{Nomo qua expresas distingo pri rango sociala, ofico}\senco{2}{Nomo qua expresas ofico, grado}\senco{3}{Indiko di la temo traktata en libro, maxim-multa-kaze kun la nomo di la autoro, skribita sur la kovrilo di libro broshita, sur la dorso, sur un ek la folii komencala}\senco{4}{Indiko, en la komenco-parto di chapitro, di dividuro di libro – di la parto di la temo qua traktesas en lu}\lingui{DEFRS}
\artiklo{tizano}%
\vorto{tizano} Infuzuro o dekokturo de substanci medikamenta\lingui{DEFIS}
\artiklo{to}%
\vorto{to} La kozo qua esas hike, ma kelke plu fore kam « co ». (abreviuro di « ito »)
\artiklo{tofo}%
\vorto{tofo} Petro tre poroza, sorto di agreguro kalka, quan onu renkontras strate en ula tereni, sub la tero vejetantala\lingui{DEFIRS}
\artiklo{togo}%
\vorto{togo} Kapo-vesto (de drapo, de veluro, de silko, e c.), ronda, sen rebordo (o kun rebordo nur tre mikra), kun la suprajo plata, plisoza periferie (maxim-multa-kaze)\lingui{EFIS}
\artiklo{toldo}%
\vorto{toldo} Kovrilo de dika telo gudrizita, od ek ledro, qua uzesas por shirmar la kargajo di la chareti, di la dilijenci, di la navi
\artiklo{tolerar}%
\vorto{tolerar}\fako{trans.} Permisar tacite, ne impedar, ne interdiktar, to quon onu desaprobas, o to per quo onu lezesas, detrimentesas, o per quo onu sufras\lingui{DEFIS}
\artiklo{tolo}%
\vorto{tolo}\fako{teknol.} Fero quan onu reduktis a lami per batado o per lamino\lingui{FIS}
\artiklo{tomato}%
\vorto{tomato}\fako{bot.}\senco{1}{Planto ek la familio « solanei », kun frukti brilante-reda}\senco{2}{La frukto di ta planto, uzata por facar dishi o salado-forme}\latina{solanum lycopersicum}\lingui{DEFRS}
\artiklo{tombako}%
\vorto{tombako} Aloyuro de 90\% latuno e 10\% zinko\lingui{DEFIRS}
\artiklo{tombo}%
\vorto{tombo}\senco{1}{Foso superkovrita de quaza tablo de petro, de marmoro, qua kontenas mortinto}\senco{2}{Monumento funerala super la foso qua kontenas mortinto}\lingui{EFS}
\artiklo{tomo}%
\vorto{tomo} Singla de la libri, broshura o bindura, ye qui konsistas verko imprimura\lingui{EFIRS}
\artiklo{tondar}%
\vorto{tondar}\fako{trans.} Tranchar reze (la lano, la pili)\lingui{FI}
\artiklo{tondrar}%
\vorto{tondrar}\fako{netrans.} \textit{(meteorol.)} \textit{(aludante fulmino-stroko)} Produktar bruiso resonanta, son-eganta\lingui{DEFI}
\artiklo{tonika}%
\vorto{tonika}\lingui{DEFIS}\subvorto{gram.}{acento tonika}{Plualtigo (che la Antiqui) o plu-intensigo (che la Moderni) di la sono lor la prononco di silabo determinita}{1}{}\subvorto{muziko}{noto tonika}{Noto qua karakterizas la tono, la gamo ye qua ario, melodio, muzikajo kompozesis – ta de qua lu komencas e qua donas a lu sua nomo}{2}{}
\artiklo{tono}%
\vorto{tono}\fako{psik.}\senco{1}{Flexo di la voco, qua indikas la diversa senti di la psiko}\senco{2}{\textit{(muziko)} (a) Intervalo inter du noti, (b) Singla de la gami en qui muzikajo esas kompozebla, indikata per la noto komencala qua determinas la repartiseso di la intervali e mi-intervali inter du noti}\senco{3}{\textit{(fotogr., pikt.)} Grado di la sama koloro segun ke lu esas plu o min obskura, plu o min briloza}\senco{4}{\textit{(fiziol.)} Kontrakturo mikra e permananta di ula muskuli, precipue olti di la sfinteri}\lingui{DEFIRS}
\artiklo{tonoplasto}%
\vorto{tonoplasto}\fako{biol.} Parto permananta e diferenciita di la protoplasmo, qua formacas la parieto di la vakuoli, e precipue di la vakuoli kontraktebla\lingui{DEFIS}
\artiklo{tonsilito}%
\vorto{tonsilito}\fako{patol.} Inflamuro di la tonsili
\artiklo{tonsilo}%
\vorto{tonsilo}\fako{anat.} Korpo glandatra, mandelo-forma, situita an singla latero di la istmo di la guturo, faringo, laringo\lingui{DEFIS}
\artiklo{tonsuro}%
\vorto{tonsuro}\fako{religio katolika}\senco{1}{Krono quan onu facas sur la kapo di la kleriki, diakoni, sacerdoti, monaki, e c. per razar parto di la hari}\senco{2}{Ceremonio dum qua la episkopo grantas ad ulu la grado unesma di la klerikeso, per razar la hari de-sur la somito di la kapo}\lingui{DEFIRS}
\artiklo{tontino}%
\vorto{tontino} Asociitaro de personi qui kolokas kapitalo komuna, por obtenar de lu, ye epoko determinita, rento dum-viva, qua reversionesos (pos singla morto) a la transvivinti, di qui la porciono kreskas proporcione\lingui{DEFIRS}
\artiklo{topazo}%
\vorto{topazo}\fako{mineral.} Lapido diafana, orea briloze\lingui{DEFIRS}
\artiklo{topika}%
\vorto{topika}\fako{medic.} Medikamento lokala, uzata extere (plastro, kataplasmo, e c.)\lingui{DEFIRS}
\artiklo{topinamburo}%
\vorto{topinamburo}\fako{bot.} Planto ek la familio « kompozaji », di qua la tuberkulo manjesas kom surogato di artichoko\latina{halianthus tuberosus}\lingui{DF}
\artiklo{topo}%
\vorto{topo}\fako{aludante navo} Platformo muntita salie cirkum, ed an la extremajo supra di basa masto qua trairas lu\lingui{E}
\artiklo{topografio}%
\vorto{topografio} La arto deskriptar e reprezentar mape la figuro di loko, di la obstakli di tereno\lingui{DEFIRS}
\artiklo{topologio}%
\vorto{topologio}\senco{1}{\textit{(geogr.)} Studio di la formi di la tereno, e di la legi naturala (*leyi) qui guvernas li}\senco{2}{\textit{(geom.)} Fako di la geometrio da Riemann, en qua onu studias la karakteri di spaco – ta karakteri esante nur qualesala, ed omna ideo pri mezuro eskartite}\lingui{DEFIS}
\artiklo{torako}%
\vorto{torako}\fako{anat.}\senco{1}{\textit{(che la vertebrozi)} Kavajo quan shirmas parieto osta, e qua kontenas la organi precipua di la sango-cirkulado e di la respirado}\senco{2}{\textit{(che la insekti)} Parto di la korpo inter la kapo e la abdomino, e sur qua insertesas la ali e la gambi}\lingui{DEFIS}
\artiklo{torcho}%
\vorto{torcho} Flamo-fonto grosiera, ek kordego tordita, indutita ye rezino o vaxo – od ek bastono ek ligno rezinoza, kovrita ye vaxo, sebo, quan onu portis lor la procesioni, la mariaji, la funeri – quan onu tenis en sua manuo lor sua penitenco publika o reparo-prego – e quan onu uzas ankore por lumizar defilo dum nokto\lingui{EFIS}
\artiklo{tordar}%
\vorto{tordar}\fako{trans.} Kontraktar violentoze (korpo flexebla) sur lu ipsa, per mantenar fixa un ek lua extremaji, en ica sinso, e rotacigar la altra inversa-sinse\lingui{EFIS}
\artiklo{toreadoro}%
\vorto{toreadoro} Lor tauro-kombato, ta omna personi qui kombatas kontre la tauro, en la areno : pikadoro, matadoro, e c.\lingui{DEFIRS}
\artiklo{torefaktar}%
\vorto{torefaktar}\fako{trans.} Submisar a la efiko da fairo intensa, qua produktas komenco di karbonigo\lingui{EFIS}
\artiklo{torento}%
\vorto{torento} Aquo-fluo, igita impetuoza, sive da la troa pento di la tereno, sive da acenso efemera\lingui{EFIS}
\artiklo{torfo}%
\vorto{torfo} Materio kombustebla, sponjatra, nigratra, qua naskas (maxim-multa-kaze) sub la aqui stagnanta, de la deskompozo di la rezidui de vejetanti\lingui{DEFIRS}
\artiklo{torida}%
\vorto{torida} Ube atmosfero esas extreme varma.\lingui{EFIS}\subvorto{}{la zono torida}{La parto di la Terglobo, inter la du tropiki}{}{}
\artiklo{torio}%
\vorto{torio}\fako{kemio} Metalo skarsa, ek la grupo de la metali teroza, quan Berzelius deskovris en 1828\simbolo{Th}\lingui{DEFIRS}
\artiklo{tormentar}%
\vorto{tormentar}\fako{trans.}\senco{1}{Submisar (kondamnito) a korpo-puniso tre granda, grava, kruela, netolerebla}\senco{2}{Probar arachar de pre-akuzato konfesi, per tordar, luxacar, dislokar, brular, kombustar, ca o ta parto di lua korpo}\lingui{EFIS}
\artiklo{tormentilo}%
\vorto{tormentilo}\fako{bot.} Planto ek la familio « rozacei », di qua la radiko esas astriktiva\latina{tormentilla}
\artiklo{tornar}%
\vorto{tornar}\fako{trans.} Fasonar per mashino sur qua onu muntis (vertikale od horizontale) peco de ligno, de metalo, de ivoro, e c., quan onu rotacigas, tale ke rivenas, pos intertempi inter-egala, ta o ca parto di la peco por subisar la efiko da cizelo od irg-altra-sorta utensilo konceptita por fasonar lu\lingui{FIS}
\artiklo{tornasolo}%
\vorto{tornasolo}\senco{1}{Materio kolorizanta, blua, quan onu extraktas de varietato di krotono (sunfloro)}\senco{2}{Solvuro de ta blua materio, quan la acidi igas redeskar}\latina{rocella lecanora}\lingui{EFS}
\artiklo{toro}%
\vorto{toro} \textit{(arkitekt.)} Muluro cilindra an la bazo di kolono, II. \textit{(geom.)} Surfaco produktita per la rotaco di cirklo cirkum rekto qua kontenesas en lua plano, ma qua ne pasas tra lua centro\lingui{DEFIS}
\artiklo{torono}%
\vorto{torono}\fako{teknol.} Asemblajo ek fili kuntordita por facar kordegi, kabli, e c., e di qua la nombro varias segun la groseso di la kordego, di la kablo\lingui{FI}
\artiklo{torpedo}%
\vorto{torpedo}\senco{1}{\textit{(zool.)} Fisho kartilagoza, kun korpo plata, disko-forma, qua havas, an singla latero di sua kapo, organo partikulara qua elektro-sukusas ti qui tushas ta fisho}\senco{2}{\textit{(milit-arto)} Mashino plenigita ye pulvero, kun amorcilo, qua, se shokata, explozas, e quan onu lansas kontre la navi enemikala, od imersas, por impedar la proximesko da ta navi a la litoro, e c.}\lingui{DEFIRS}
\artiklo{torporar}%
\vorto{torporar}\fako{netrans.} Afektesar da ta sorto di inerteso, di graveso quan produktas la ceso di la sentiveso e di la funciono, en un o plura parti di la homo-korpo, o di la psiko\lingui{EFIS}
\artiklo{torsado}%
\vorto{torsado}\senco{1}{Volvajo ek fili de silko, oro, e c., tordita spirale, de qui onu facas la kordeti di la franji, di la glani por kurteni, tapeti, e c.}\senco{2}{Volvajo ek fili de arjento, de oro, de qua onu facas la epoleti di la oficiri}\senco{3}{Volvajo de hari tordita spirale}\lingui{DF}
\artiklo{torso}%
\vorto{torso}\fako{anat.} Korpo di homo o di animalo, sen la kapo nek la membri\lingui{DEFIR}
\artiklo{tortikolo}%
\vorto{tortikolo}\fako{patol.} Doloro che la muskuli di la kolo, qua koaktas ad igar olca deviacar de vertikalo\lingui{DEFIS}
\artiklo{torto}%
\vorto{torto} Kukajo ronda, en qua onu pozas karno de bruti, de pultro, de fisho – e kaneli\lingui{DFIRS}
\artiklo{tortugo}%
\vorto{tortugo}\fako{zool.} Genero de repteri, amfibia, quar-peda, qua marchas lente, kun kuraso de materio korno-harda (skalio)\latina{chelonia}\lingui{EFI}
\artiklo{tostar}%
\vorto{tostar}\fako{netrans., ad} Propozar (a la personi prezenta) drinkor « por la duro di la saneso » (di ulu), akompanante ta propozo per diskurseto\lingui{DEFR}
\artiklo{tota}%
\vorto{tota}\senco{1}{De qua nulo deprenesis per separo}\senco{2}{(a) Qua nule diminutesis. (b) \textit{(metaf.)} Qua ne implikas restrikteso}\lingui{EFI}
\artiklo{totalizatoro}%
\vorto{totalizatoro}\fako{tekn.} Aparato qua donas la sumo di serio de operaci
\artiklo{toxiko}%
\vorto{toxiko} Substanco venenoza\lingui{DEFIRS}
\artiklo{toxikologio}%
\vorto{toxikologio} La cienco qua traktas la toxiki\lingui{DEFIS}
\artiklo{toxino}%
\vorto{toxino}\fako{kemio biol.} Veneno komplexa e solvebla, quan produktas la funciono di la mikrobi, sive en la organismi vivanta, sive en la medii artificala di kultivo\lingui{DEFIS}
\artiklo{tra}%
\vorto{tra} Irante de ca a ta extremajo, sinse di profundeso\lingui{FS}
\artiklo{trabekulo}%
\vorto{trabekulo}\fako{histol.} Dina prolonguro qua separas su de parieto di celulo e salias en kavajo, maxim-multa-kaze per kontraktar anastomosi kun la trabekuli vicina
\artiklo{trabo}%
\vorto{trabo}\fako{teknol.} Stangego de ligno, quadratigata, destinata a suportor ula parto di edifico (tekto, plankaro, e c.)\lingui{ISL}
\artiklo{traciono}%
\vorto{traciono}\fako{teknol.) (pri relvoyo} Ensemblo de la mashini qui tiras, tranas, la treni qui transportas la pasajeri o la vari\lingui{DEFIS}
\artiklo{traco}%
\vorto{traco}\senco{1}{Serio de marki lasita da korpo presita aden altra, per qui onu deskovras la direciono ye qua iris ulu, ulo}\senco{2}{Singla ek ta marki. \textit{(anke metafore)}}\senco{3}{Marko per qua onu deskovras la efiko da ulo, da ulu, en epoko pasinta}\lingui{EFI}
\artiklo{tradeskancio}%
\vorto{tradeskancio}\fako{bot.} Genero de « komelinacei », kun flori violea, purpurea o blanka, qua vivas en la regioni atmosfero-varma di la suda parto di Amerika\latina{tradescantia}
\artiklo{tradicionar}%
\vorto{tradicionar}\fako{trans.} Transmisar vocale la fakti historiala, la doktrini religiala, la kustumi\lingui{DEFIRS}
\artiklo{tradukar}%
\vorto{tradukar}\fako{trans., de, ad} Pasigar de linguo ad altra linguo\lingui{FIS}
\artiklo{trafiko}%
\vorto{trafiko} Importo, frequeso-grado di la paso di la treni, di la navi\lingui{EFIS}
\artiklo{tragakanto}%
\vorto{tragakanto}\fako{bot.} Arboreto ek la genero « astragalo », qua donas adraganto\latina{astragalus}\lingui{DEFR}
\artiklo{tragedio}%
\vorto{tragedio}\senco{1}{Teatrajo normala, di qua la roli luktas kontre sua pasioni o kontre desfortuno, e quin minacas e frapas katastrofi}\senco{2}{\textit{(metaf.)} Eventajo funesta}\lingui{DEFIRS}
\artiklo{tragikomedio}%
\vorto{tragikomedio} Teatrajo di qua la roli ne esas monarki, suvereni, quankam la temo esas tragediala\lingui{DEFIRS}
\artiklo{trahizar}%
\vorto{trahizar}\fako{trans.} Abandonar (ulu, a qua onu devas esar fidela), plu o min perfide\lingui{EFS}
\artiklo{traino}%
\vorto{traino}\fako{milit-arto} Korpo di qua la tasko konsistas ye duktar, transportar, la provizuri, la *asiejo-mashini e la municioni
\artiklo{traito}%
\vorto{traito}\senco{1}{Singla ek la linei di la vizajo}\senco{2}{Marko signifikiva di ulo (karaktero, prudenteso, jenerozeso, e c.)}\lingui{EFI}
\artiklo{trajektorio}%
\vorto{trajektorio}\senco{1}{Lineo quan parkuras la baricentro di korpo qua lansesis}\senco{2}{Kurvo qua krucumas, ye angulo donita, kurvi altra, infinite-multa}\lingui{DEFIRS}
\artiklo{trakar}%
\vorto{trakar}\fako{trans.}\senco{1}{Persequar (la vildo qua esas en bosko o foresto), cirkondante lu, diminutante gradoze la cirklo-diametro}\senco{2}{\textit{(metaf.)} Persequar (krimininti, deliktinti) akulante li gradoze a loko sen eskapo-voyo}\lingui{F}
\artiklo{trakeito}%
\vorto{trakeito}\fako{patol.} Inflamuro di la trakeo\lingui{DEFIRS}
\artiklo{trakeo}%
\vorto{trakeo}\fako{anat.} Tubo ek serio de ringi kartilagoza, qua venas de la laringo, decensas alonge la kolo, avan la ezofago, bifurkas ye lua parto infra, formacante du tubi (bronki) di qui singla komunikas kun un ek la du pulmoni, ed uzesas por la paso di aero\lingui{DEFIRS}
\artiklo{trakeotomio}%
\vorto{trakeotomio}\fako{medic.} Incizo di la trakeo, qua posibligas la eniro di la aero extera en la kazo di la afekti qui obstruktas la respiro-tubi (krupo, e c.)\lingui{DEFIRS}
\artiklo{trakito}%
\vorto{trakito}\fako{mineral.} Roko quan karakterizas la prezenteso (plugranda-skale) en olu di feldspato alkalioza, e la preska absenteso di quarco o nefelino\lingui{DEFIRS}
\artiklo{trakomo}%
\vorto{trakomo}\fako{patol.} Morbo infektiva di la okul-sistemo, di qua la jermo esas ankore nekonocata\lingui{F}
\artiklo{traktar}%
\vorto{traktar}\fako{trans.}\senco{1}{Agar ad (ulu) ye ta o ca maniero}\senco{2}{\textit{(aludante oratoro, literaturisto, piktisto, e c.)} Expozar, developar (tota temo)}\lingui{EFIRS}
\artiklo{traktato}%
\vorto{traktato}\fako{pri} Verko en qua onu traktas temo, explorante ye maniero ordinoza lua parti
\artiklo{traktoro}%
\vorto{traktoro}\fako{tekn.} Vehilo automobila di qua la roti esas provizita ye organi por igar lu adherar a sulo mem tre neplana (roti qui inter-relatas per bendo metala quaze-tapisa qua, lor la funciono, aspektas quale raupo movanta – o roti angul-pecoza) e qua posedas junto-sistemo por remorko (di plugilo, di falcho-mashini, e c.)\lingui{DEFIRS}
\artiklo{traktrico}%
\vorto{traktrico}\fako{geom.} Kurvo quan trasas kordo quan onu tir-esforcas – altra-vorto : di qua la tangento egalesas lineo konstanta\lingui{DEFIS}
\artiklo{tramo}%
\vorto{tramo}\fako{tekn.} Longa stango de fero o, maxim-multa-kaze, de stalo, en qua existas kanelo por recevar la parto salianta di la felgo di veturo-roto, e qua, instalite en sulo, ne salias de la surfaco di olca. (opoze a « relo »)\lingui{DEFIS}
\artiklo{tramplar}%
\vorto{tramplar}\fako{netrans.} Movar, agitar, sua pedi, sen avancar\lingui{DE}
\artiklo{trampolino}%
\vorto{trampolino}\fako{tekn.} Planajo pento-forma, ek planko instalita tale ke olu esas elastika, de-sur qua la saltanti ri-saltas, tale aquirinte springo plu granda\lingui{DFIRS}
\artiklo{tranar}%
\vorto{tranar}\fako{trans.}\senco{1}{Tirar dop su (ulo, ulu) malgre la rezisto de la pezo o de la froto di parto (de lu) sur sulo}\senco{2}{Koaktar (ulu) pri kun-venar}\lingui{EF}
\artiklo{trancar}%
\vorto{trancar}\fako{patol.} Nekonciar, dum epilepsio, dum hipnoto, e c.\lingui{DEFIRS}
\artiklo{trancendanta}%
\vorto{trancendanta}\fako{filoz.} La filozofio qua traktas exkluzive la konoco aprioria, la konoco pure racionala. (En la sistemo da Immanuel Kant, spaco e tempo esas koncepti trancendanta)\lingui{DEFIRS}
\artiklo{tranchar}%
\vorto{tranchar}\fako{trans.} Dividar par-separe (korpo) per korpo plu harda e taliita tale ke adminime un ek lua aristi esas akuta. (ex. : tranchar la kapo per gilotino, per hakilo-frapo)\lingui{EFIS}
\artiklo{trancheo}%
\vorto{trancheo}\fako{tekn.}\senco{1}{Exkavuro facita en sulo ula-longe (por kanalo, kanaleto, voyo, relvoyo, e c.)}\senco{2}{\textit{(milit-arto)} Fosato, shirme da qua onu povas proximigar su a fortifikajo, a fortreso, *asiejata (*siejata)}\lingui{DEFR}
\artiklo{tranquila}%
\vorto{tranquila} Qua ne esas agitema, bruisema\lingui{EFIS}
\artiklo{trans}%
\vorto{trans} Plu fore kam\lingui{L}
\artiklo{transaktar}%
\vorto{transaktar}\fako{netrans., pri} Abutar ad interkonsento per cedi reciproka pri la litijo-punto\lingui{EFIS}
\artiklo{transbordar}%
\vorto{transbordar}\fako{trans.} Igar (vari, pasajanti) pasar de ca a ta navo – o de loko ube la relvoyo esas interruptita, a la loko ube lu esas itere uzebla\lingui{FS}
\artiklo{transepto}%
\vorto{transepto}\fako{arkitekt.} Ta parto di kirko exter la navo, qua formacas la « brakii » di la kruco\lingui{DEFR}
\artiklo{transferar}%
\vorto{transferar}\fako{trans., de, ad}\senco{1}{Instalar, establisar, de ca loko ad ita}\senco{2}{\textit{(yuro-cienco)} Igar divenar la proprietajo di altru}\lingui{DEFIS}
\artiklo{transfigurar}%
\vorto{transfigurar}\fako{trans.} Transformar per altrigar la vizajo, la traiti extera\lingui{EFIS}
\artiklo{transfinita}%
\vorto{transfinita}\fako{matem.}\subvorto{}{nombri transfinita}{Nombri qui egardesis en la teorio pri la ensembli.}{}{}\subvorto{}{ensemblo matematikala}{Grupo determinita de nombri, objekti, e c.; la literi di la alfabeto, la nombri pozitiva ne-para, e c. esas ensembli}{}{}
\artiklo{transformar}%
\vorto{transformar}\fako{trans., ulo, ulu, ad}\senco{1}{Pasigar de ca formo ad ita}\senco{2}{\textit{(metaf.)} Par-modifikar}\lingui{DEFIS}
\artiklo{transformatoro}%
\vorto{transformatoro}\fako{tekn.) (elektro} Aparato qua, recevante energio elektrala di ta o ca speco, modifikas exkluzive lua formo\lingui{DEFIRS}
\artiklo{transitar}%
\vorto{transitar}\fako{netrans.} \textit{(komerco)} \textit{(aludante vari)} Transirar lando, urbo, sen pagar (pri la imposti doganala od acizala), pro yuro grantita pri lo en ula cirkonstanci\lingui{DEFIRS}
\artiklo{translacar}%
\vorto{translacar}\fako{netrans.) (aludante solido} Movar (o movesar) tale ke omna parti havas rapideso-grado e direciono konstanta\lingui{EFIS}
\artiklo{translucida}%
\vorto{translucida} Qua lasas lumo pasar tra su, ma tale difuzita ke onu ne povas dicernar precize la objekti tra lu\lingui{DEFIS}
\artiklo{transmisar}%
\vorto{transmisar}\fako{trans.}\senco{1}{\textit{(yuro-cienco)} Igar pasar ad ulu (ta o ca yuro, privilejo, quan onu posedas)}\senco{2}{\textit{(biol. Igar pasar a la decendanti)} \textit{(tekn.)} Komunikar la movo di organo ad altra organo (per ingranaji, kabli, rimeni, kateni, e c.)}\lingui{DEFIS}
\artiklo{transmutar}%
\vorto{transmutar}\fako{trans., ad) (alkemio} Chanjar la naturo di (metalo ordinara a metalo precoza)\lingui{EFIS}
\artiklo{transparar}%
\vorto{transparar}\fako{netrans.) (tra, sub} Esar videbla tra\lingui{DEFIS}
\artiklo{transparento}%
\vorto{transparento} Papero-folio lineizita horizontale, quan onu lokizas sub la papero-folio sur qua onu skribas, por ke la skriburo formacez linei rekta\lingui{DEFIRS}
\artiklo{transpirar}%
\vorto{transpirar}\fako{netrans.}\senco{1}{Pasar tra la pori di korpo, forme di guteti qui kelke aspektas quale sudoro}\senco{2}{\textit{(metaf.)} \textit{(pri informi)} Diskonocesar pro la nediskreteso di ula personi o di ula persono}\lingui{DEFIS}
\artiklo{transportar}%
\vorto{transportar}\fako{trans.} Portar de ca a ta loko\lingui{DEFIRS}
\artiklo{transubstanciar}%
\vorto{transubstanciar}\fako{teol. katol.} Agar la transmuto di la pano e di la vino a la korpo e la sango di Iesu Kristo\lingui{DEFIS}
\artiklo{transversa}%
\vorto{transversa}\fako{pri korpo} Qua, de un ek sua extremaji a la altra, trairas lineo, direciono, krucume\lingui{DEFIS}
\artiklo{trapezo}%
\vorto{trapezo}\senco{1}{\textit{(geom.)} Quadrilatero di qua du lateri esas inter-paralela, ma ne-interegala}\senco{2}{Implemento di gimnastiko ek bastono qua povas portar la pezo di persono, suspendita per du kordi ringoza a la portiko}\lingui{DEFIRS}
\artiklo{trapezoedro}%
\vorto{trapezoedro}\senco{1}{Solido di qua la edri esas trapezoida}\senco{2}{Solido ek duadek-e-quar edri quadrilatera intersimetre}\lingui{DEFIS}
\artiklo{trapezoido}%
\vorto{trapezoido}\senco{1}{\textit{(geom.)} Quadrilatero plana, di qua omna lateri esas inter-obliqua}\senco{2}{\textit{(anat.)} Muskulo sub-pela qua kovras preske tote la regiono dorsala}\lingui{DEFIS}
\artiklo{trapo}%
\vorto{trapo} Pordo horizontala qua klozas aperturo facita en la plankaro e quan onu levas od abasas segun-vole\lingui{EFIS}
\artiklo{trasar}%
\vorto{trasar}\fako{trans.} Indikar per streko, per lineo, la direciono, la formo\lingui{F}
\artiklo{tratar}%
\vorto{tratar}\fako{trans., ulu) (aludante komercisto} Emisar mandato di qua la pekunio-quanto stipulita esas pagenda da (la tratito) po la vari furnisita a lu\lingui{DFIRS}
\artiklo{traumato}%
\vorto{traumato}\fako{patol.} Ensemblo de la desordini fiziologiala quin produktis ago od efiko violentoza\lingui{DEFIRS}
\artiklo{traurar}%
\vorto{traurar}\fako{netrans.} Afliktesar da la morto di persono kara\lingui{DR}
\artiklo{traveo}%
\vorto{traveo}\fako{arkitekt.} La spaco inter du suportili (trabi, koloni, pilastri, e c.)\lingui{F}
\artiklo{traverso}%
\vorto{traverso}\fako{tekn.} Peco ek ligno o metalo perpendikla ye la elementi precipua di konstrukturo, destinita a mantenar la interdisto di ta elementi\lingui{EFIS}
\artiklo{travertino}%
\vorto{travertino}\fako{geol.} Tofo kalkoza, harda, griza, uzata por la konstrukto di domi, e c.
\artiklo{travestiar}%
\vorto{travestiar}\fako{trans.}\senco{1}{Metar (sur su, sur ulu) vesto, kostumo, qua ne esas olta di lua klaso sociala o di lua sexuo}\senco{2}{Metar (sur ulu, sur su) kostumo qua igas des-facila lua rekonoceso}\lingui{DEFI}
\artiklo{travo}%
\vorto{travo}\fako{tekn.} Mashino en qua onu ligas la kavali, la bovi, e c., sive por hufizar li kande li rezistas, sive por submisar li ad ula operaci\lingui{EFIS}
\artiklo{tre}%
\vorto{tre}\fako{adverbo} Alta-grade, pri lua qualeso\lingui{F}
\artiklo{trefo}%
\vorto{trefo}\fako{karto-ludo} Un ek la emblemi nigra-kolora qua reprezentas folio di trifolio\lingui{DFR}
\artiklo{treko}%
\vorto{treko} Mala voyeto quan facis nevole la personi, o trupi, o veturi, qui pasis ofte sur la traci di lia pre-irinto, en sulo naturala\lingui{E}
\artiklo{trelicho}%
\vorto{trelicho} Telo grosiera, por la faco di varo-saki, laborvesti\lingui{DFI}
\artiklo{treliso}%
\vorto{treliso}\fako{tekn.} Plako plu o min granda areale, ek fer-fili interplektita tale ke li formacas mashi plu o min larja, quan onu muntas avan fenestro o lokizas avan la fairuyo di kameno, od uzas (larja-bendo) kom klozilo di gardeno, di tereno, e c.\lingui{EF}
\artiklo{tremao}%
\vorto{tremao} Signo ek du punti inter-apuda, quan onu lokizas horizontale super vokalo por modifikar la pronunco di olca (en la lingui Germana o Germanida) o por indikar ke du vokali inter-apuda ne pronuncesas diftonge (en la lingui Franca, Angla, e c.) (kom ex. : F. noël, poële – A. coöperation, G. Brüder)\lingui{DFS}
\artiklo{tremar}%
\vorto{tremar}\fako{netrans.) (pro, de} Agiteso da serio de mikra ocili, quin produktas la koldeso di atmosfero, febro, pavoro, timo, e c. en la homo-korpo. – \textit{(anke metaf.)}\lingui{EFI}
\artiklo{tremolo}%
\vorto{tremolo}\fako{muziko} Tremo sencesa sur noto muzikala\lingui{DEFIRS}
\artiklo{trempar}%
\vorto{trempar}\fako{trans., en} Plunjar en liquido, por imbibar ye ta liquido (la plunjajo)\lingui{F}
\artiklo{tremulo}%
\vorto{tremulo}\fako{bot.} Speco di poplo, di qua la folii, lejera, tremas efike da vento mem febla\latina{populus tremula}
\artiklo{treno}%
\vorto{treno} Ensemblo de la vagoni quin tranas lokomotivo\lingui{EFIS}
\artiklo{trepano}%
\vorto{trepano}\senco{1}{\textit{(kirurg.)} Instrumento por borar la osti, precipue olti di la kranio, por posibligar la ekfluo de puso o la retrodukto di osto profunde sinkita a lua plaso normala}\senco{2}{\textit{(tekn.)} Granda mashino por perforar sulo roka, por facar putei, serche di petrolo, e c.}\lingui{DEFIRS}
\artiklo{trepidar}%
\vorto{trepidar}\fako{netrans.}\senco{1}{\textit{(aludante navo, automobilo, qua funcionas)} Tremar ye sukusi bruska quin produktas la vapormashino, la motoro}\senco{2}{Manifestar, malgre su, signi di sua nepacienteso}\lingui{EFIRS}
\artiklo{treso}%
\vorto{treso} Kordono plata, ek fili (lana, silka, e c.) od ek palio o hari interplektita\lingui{DEFIRS}
\artiklo{tresto}%
\vorto{tresto} Planko sur quar pedi, uzata por subtenar longa tablo, eshafodo, *stejo (ceneyo) feriala\lingui{EF}
\artiklo{trezoro}%
\vorto{trezoro}\senco{1}{Asemblajo ek moneto-peci ora, arjenta, ek objekti tre precoza – rezervita, celita od enterigita}\senco{2}{\textit{(yuro-cienco)} Omna kozo celita ed enterigita quan ulu deskovras hazarde, e pri qua nulu povas justifikar proprieto-yuro}\senco{3}{Ensemblo de la revenui di stato}\senco{4}{Kolektajo ek reliquii, vazi, ornivi tre precoza, qui Konservesas en kirko}\senco{5}{\textit{(metaf.)} Irgo quo egardesas kom tre precoza. \textit{(anke pri persono)}}\lingui{DEFIS}
\artiklo{tri}%
\vorto{tri}\fako{aritm.} La cifro « 3 »\lingui{DEFIRS}
\artiklo{tri-}%
\vorto{tri-} Prefixo ciencala : Qua havas \textit{tri}…
\artiklo{triangulo}%
\vorto{triangulo}\senco{1}{\textit{(geom.)} Figuro geometriala kun tri anguli – parto de planajo qua havas kom limiti tri rekti qui intersekas duope}\senco{2}{Objekto qua havas la sama formo kam triangulo : (a) \textit{(religio kristana)}. Simbolo di la Triuno, (b) Emblemo di la framasonaro; (c) \textit{(muziko)} Instrumento ek stalo, quan onu resonigas per vergeto stala}\lingui{DEFIRS}
\artiklo{triaso}%
\vorto{triaso}\fako{geol.} Sulo-strato ek marni, kalkajo, greso, qua sucedas liaso segun la decens-ordino
\artiklo{tribono}%
\vorto{tribono}\fako{en la Roma antiqua}\lingui{DEFIRS}\subvorto{}{tribono di la soldati}{Chefo qua komandis legiono, alterne kun kin altri}{1}{}\subvorto{}{tribono di la plebeyi}{Judiciisto qua defensis la interesti di la plebeyani}{2}{}
\artiklo{tribular}%
\vorto{tribular}\fako{trans., per} Tormentar (lu) psikale, mentale\lingui{DEFIS}
\artiklo{tribunalo}%
\vorto{tribunalo}\fako{yuro-cienco} Asemblitaro de judiciisti qui *seancas (kunsidas) ensemble\lingui{DEFIRS}
\artiklo{tribuno}%
\vorto{tribuno}\senco{1}{Loko situita kelke-alte, quaze estrado de ube la oratoro diskursis a la plebeyani sur placo (en la epoki antiqua) – de ube la oratoro diskursas en asemblitaro deliberanta (nia-epoke)}\senco{2}{Galerio ube onu rezervas plasi en kirko, en asemblitaro deliberanta}\lingui{DEFIRS}
\artiklo{tribuo}%
\vorto{tribuo} Grupo de familii qua formacas socio neperfekte organizita\lingui{DEFIS}
\artiklo{tributar}%
\vorto{tributar}\fako{netrans., ad) (aludante populo vinkita} Pagar la pekunio-quanto, livrar la vari o precozaji, quin la populo vinkinta impozis a sua dominacati\lingui{DEFIS}
\artiklo{tricepso}%
\vorto{tricepso}\fako{anat.} Muskulo ek tri faski\lingui{DEFS}
\artiklo{triciklo}%
\vorto{triciklo} Velocipedo kun tri roti, muntita triangule\lingui{EFIS}
\artiklo{trifolio}%
\vorto{trifolio}\fako{bot.} Leguminoso « papilionacea », qua havas, ordinare, nur tri folii\latina{trifolium}\lingui{DEFIS}
\artiklo{triftongo}%
\vorto{triftongo}\fako{gram.} Silabo ek tri vokali qui pronuncesas per un emiso di voco. \textit{(Ex., la F. eau qua pronuncesas « o »)}\lingui{DEIRS}
\artiklo{triglifo}%
\vorto{triglifo}\fako{arkitekt.} Saliajo kun ranuri vertikala, ye qua esas ornita singla latero di friso dorika, e qua reprezentas la extremaji di la trabeti pozita sur la arkitravo\lingui{DEFIRS}
\artiklo{trigonometrio}%
\vorto{trigonometrio}\fako{geom.} Parto di geometrio, di qua la temo esas mezurar la trianguli, per determinar, helpe da ula donataji, la lateri e la anguli\lingui{EF}
\artiklo{trigramo}%
\vorto{trigramo} Grupo de tri literi qui reprezentas un sono
\artiklo{trikino}%
\vorto{trikino} Vermo intestinala haro-tenua, qua vivas (en la kazi maxim multa) en la karno di la porko\latina{trichinella spiralis}\lingui{DEFIRS}
\artiklo{trikinoso}%
\vorto{trikinoso}\fako{patol.} Morbo produktata da la ingesto di porko-karno qua kontenas trikino\lingui{DEFIRS}
\artiklo{triklina}%
\vorto{triklina}\fako{kristalografio} Qua havas tri axi inklinita ica relate iti\lingui{DEFIS}
\artiklo{trikocisto}%
\vorto{trikocisto}\fako{biol.} Singla de la mikra organi akuta quin onu renkontras che ula infuzorii holotricha
\artiklo{trikogino}%
\vorto{trikogino}\fako{biol.} Prolonguro pilo-forma qua garnisas la bazo di la celulo feminala di la floridei\lingui{DEFIS}
\artiklo{trikotar}%
\vorto{trikotar}\fako{trans.} Facar manue la mashi di texuro (kotona, lana, e c.) sur stangeti stala, ligna, ivora, kun extremaji obtuza, quin onu movas rapide\lingui{F}
\artiklo{triktrako}%
\vorto{triktrako}\fako{ludilo} Sorto di damo-planko kun rebordi, dividita a du faki ube singla ek la ludanti situas sua dami segun la nombro de punti quin lu obtenas per la ludo-kubi quin lu lansas sur la planko\lingui{DEFI}
\artiklo{trilar}%
\vorto{trilar}\fako{netrans.) (muziko} Ornar ye ta acesorajo muzikala qua konsistas ye la frapo tre rapida e plu o min duroza di noto kun la noto qua esas super lu, sive diste de un tono, sive diste de du toni\lingui{DEFIS}
\artiklo{triliono}%
\vorto{triliono}\fako{aritm.} 1.000.000³\lingui{DEFIRS}
\artiklo{trilobito}%
\vorto{trilobito}\fako{paleont.} Granda grupo de krustacei, qua cesis existar, e qua karakterizas la sulo-strati sekundara ube li esas inkluzita\lingui{DEFIS}
\artiklo{trilogio}%
\vorto{trilogio}\senco{1}{\textit{(epoki antiqua)} Ensemblo de tri tragedii qui traktas (en la kazi maxim multa) tri parti intersequanta de temo komuna, quin onu reprezentis, ica pos ita, en la konkursi pri poeziaji qui eventis lor ula festii solena}\senco{2}{\textit{(nia-epoke)} Asemblajo de tri tragedii qui traktas tri parti intersequanta di temo komuna}\lingui{DEFIRS}
\artiklo{trimestro}%
\vorto{trimestro} Periodo de tri monati, intersequanta\lingui{DEFIRS}
\artiklo{trimetriko}%
\vorto{trimetriko}\fako{matem.} Qua koncernas tri mezuri interdiferanta
\artiklo{trinomio}%
\vorto{trinomio}\fako{algebro} Polinomio tri-termina\lingui{DEFIRS}
\artiklo{trioleto}%
\vorto{trioleto}\senco{1}{\textit{(poezio)} Kupleto de ok versi en qui la unesma repetesas pos la triesma, e la duesma pos la sisesma}\senco{2}{\textit{(muziko)} Grupo tri-elementa de noti en mezuro du-elementa}
\artiklo{triplikato}%
\vorto{triplikato} La triesma kopiuro de akto\lingui{DEFIS}
\artiklo{tripliko}%
\vorto{tripliko}\fako{yuro-cienco} Respondo a dupliko\lingui{F}
\artiklo{tripo}%
\vorto{tripo} Stomako di la rumineri (pancho, e c.), egardata kom nutrivo por la homi\lingui{EFIRS}
\artiklo{tripolio}%
\vorto{tripolio}\fako{tekn.} Tero silicoza, flava redatre, uzata por polisar la speguli, la metali, e c.\lingui{DEFIS}
\artiklo{triptiko}%
\vorto{triptiko} Pikturo o tabelo sur tri « shutri » faldebla, de qui du (ia dextra e la sinistra) kovras, pos faldo, la centrala\lingui{DEFIS}
\artiklo{triremo}%
\vorto{triremo}\fako{antique, en Roma} Navo kun tri rangi de remili\lingui{DEFIRS}
\artiklo{trismo}%
\vorto{trismo}\fako{patol.) (che tetananto} Kontrakturo che la muskuli di la maxilo e di la mandibulo
\artiklo{trista}%
\vorto{trista}\senco{1}{\textit{(pri persono)} Qua sufras psikale}\senco{2}{\textit{(pri kozo)} Qua produktas sufro psikala}\lingui{FIS}
\artiklo{tritono}%
\vorto{tritono}\senco{1}{\textit{(mitol.)} Maro-deajo qua, segun la reprezenti, havas vizajo homala e korpo qua finas fishatre}\senco{2}{\textit{(zool.)} Batrako squala, vicina ye salamandro}\latina{triton}\lingui{DEFIRS}
\artiklo{triturar}%
\vorto{triturar}\fako{trans.}\senco{1}{Igar pulvero, per frotado, per aplastar sen frapar}\senco{2}{Igar, de korpo harda, peceti extreme mikra per shoki o per la preso da korpo plu harda}\lingui{EFIS}
\artiklo{triumfar}%
\vorto{triumfar}\fako{netrans.}\senco{1}{\textit{(en la Roma antiqua)} Obtenar la honoro qua grantesis a la generalo qua agabis granda vinko, e qua konsistis ye enirar pompoze la urbo, sur charo, sequata da lua armeo, kun la kaptiti e la raptaji}\senco{2}{\textit{(netrans.)} (ye) Vinkar briloze, komplete}\lingui{DEFIRS}
\artiklo{triumviro}%
\vorto{triumviro}\fako{en la Roma antiqua} Stat-oficiero administrala, alta-ranga, qua, lor la republiko, funcionis kun du kolegi. – \textit{(anke metaf.)}\lingui{DEFIRS}
\artiklo{triviala}%
\vorto{triviala} Sen originaleso, infra-sorta e pasable grosiera\lingui{DEFIRS}
\artiklo{tro}%
\vorto{tro} Plu (multe) kam konvenas, kam necesa, kam bezonata, kam demandita\lingui{FI}
\artiklo{trofeo}%
\vorto{trofeo}\senco{1}{\textit{(epoki antiqua)} Asemblajo ek la spoliuro, la raptajo de enemiko vinkita, quin onu suspendis a la trunko di arboro, kom monumento pri sua vinko}\senco{2}{\textit{(nun-epoke)} Grupo de armi, de standardi, asalte konquestita de la enemiko, quan onu suspendas kom memorigivo pri sua vinko}\lingui{DEFIRS}
\artiklo{trofoplasmo}%
\vorto{trofoplasmo}\fako{histol.} Parto di la hialoplasmo, o citoplasmo, amorfa, qua pleas rolo pure nutrala – nomo di la substanco fibretoza fundamentala di la neurono\lingui{DEFIRS}
\artiklo{troglodito}%
\vorto{troglodito}\senco{1}{Habitanto di la kaverni subsula}\senco{2}{\textit{(zool.)} Ucelo insekto-manjanta, quan onu konfundas multa-kaze a la regolo}\latina{troglodytes}\lingui{DEFIS}
\artiklo{trogo}%
\vorto{trogo}\senco{1}{Vazo ligna quan uzas la masonisti por dilutar gipso, cemento}\senco{2}{Petro-bloko kavigita o ligno-vazo de qua drinkas la kavali, la bruti}\lingui{DE}
\artiklo{trokantero}%
\vorto{trokantero}\fako{anat.} Singla ek la du tuberi rondatra di la femuro ube (an la hancho) interjuntas la kolo di la femuro e la korpo\lingui{EF}
\artiklo{trokaro}%
\vorto{trokaro}\fako{kirurg.} Instrumento di qua la formo esas olta di puncilo cilindra muntita an mancho e kontenata en kanulo, quan onu uzas por puncionar\lingui{DEFIS}
\artiklo{trokeo}%
\vorto{trokeo}\fako{versif.} Pedo ek un silabo longa ed un silabo kurta\lingui{DEFIR}
\artiklo{trokoido}%
\vorto{trokoido}\fako{matem.} La figuro kurva quan prizentas kordo vibranta, fuzelo\lingui{DEFIS}
\artiklo{troleo}%
\vorto{troleo}\fako{tekn.} Che la vehili elektro-movata : organo por la preno di elektro-fluo, qua konsistas ye stango flexebla kun mikra roteto o kun kontaktilo glitanta, qua uzesas por transmisar la elektro-fluo de la kablo konduktiva til la motoro di la vehilo\lingui{DEFS}
\artiklo{trombo}%
\vorto{trombo} Kolono de aquo qua, levata da vento-vortico, jiradas, devastante omno ube lu pasas\lingui{DFIS}
\artiklo{trombono}%
\vorto{trombono}\fako{muziko} Instrumento ek du tubi latuna qui interinkastras, e glitas ita sur ica, tale ke varias la longeso di la tubo\lingui{DEFIRS}
\artiklo{trompar}%
\vorto{trompar}\fako{trans., per} Igar (ulu) erorar, per ruzo\lingui{F}
\artiklo{trono}%
\vorto{trono}\senco{1}{Sidilo situita plu alte kam la ceterai, sur qua sidas la suvereni dum la *seanci (kunsidi) solena}\senco{2}{\textit{(metaf.)} La autoritato suprega, en la stato monarkoza}\lingui{DEFIRS}
\artiklo{tropiko}%
\vorto{tropiko}\fako{geogr.} Paralelo terglobala, qua en la mi-sfero boreala ed en la mi-sfero australa, separas la zono torida de la zono varmeso-temperata\lingui{DEFIRS}
\artiklo{tropismo}%
\vorto{tropismo}\fako{biol.} Kresko di organo ta-direcione o ca-direcione, influe da stimulivi mekanikala, fizikala o kemiala\lingui{DEFIS}
\artiklo{tropo}%
\vorto{tropo}\fako{retor.} Figuro qua konsistas ye igar vorto, expresuro, deviacar de sua senco propra, por uzar lu kun senco figurala\lingui{DEFIRS}
\artiklo{trotar}%
\vorto{trotar}\fako{netrans.) (aludante kavalo od altra quadripedo} Irar ye aluro inter pazo nehastoza e galopo, en qua singla ek la du gambi diagonala (i. e. gambo avana dextra, e gambo dopa sinistra – o gambo dopa dextra e gambo avana sinistra) esas alterne movata adavan\lingui{DEFIS}
\artiklo{trotuaro}%
\vorto{trotuaro} Extensajo, petro-plakizita o bitumizita en la maxim multa kazi, situita an singla ek la lateri di la choseo di strado, o rive di relvoyo, por la cirkulo ed embarko di la ped-iranti\lingui{DFR}
\artiklo{trovar}%
\vorto{trovar}\fako{trans.} Vidar aparar (ulu, ulo) quan onu serchis, o ne serchis. (sercho-trovar e hazardo-trovar, risp.)
\artiklo{trubaduro}%
\vorto{trubaduro}\fako{historio literat.} Poeto qua kompozis per la linguo olima di la sudo-regioni di Francia\lingui{DEFIRS}
\artiklo{trublar}%
\vorto{trublar}\fako{trans.} Igar ulu interruptar sua iro quieta\lingui{EF}
\artiklo{truflo}%
\vorto{truflo}\fako{bot.} Sorto di fungo subsula, tubero-forma – tre prizata (kom disho) pro lua aromo\latina{tuber}\lingui{DEFRS}
\artiklo{truismo}%
\vorto{truismo} Dico di qua la exakteso esas tante evidenta ke, judikeble, esas neutila enuncar lu\lingui{EFR}
\artiklo{truko}%
\vorto{truko}\fako{tekn.} Vagono platforma per qua onu transportas la veturi, plankegi, trabi, arboro-trunki ed altra objekti inkombroza\lingui{DEF}
\artiklo{trulo}%
\vorto{trulo}\fako{tekn.} Utensilo spatulo-forma por dilutar ed aplikar gipso, mortero, cemento, e c.\lingui{EFS}
\artiklo{trumpeto}%
\vorto{trumpeto}\fako{muziko} Sufl-instrumento ek latuno od altra metalo, di qua la sono esas stridanta, boranta – uzata por la ario-soni armeala, por la proklami sur placi\lingui{DEFIS}
\artiklo{trumpo}%
\vorto{trumpo}\fako{karto-ludo} La emblemo qua preponderas omna ceterai, e quan onu determinas singlafoye, sive per vidigar (lor la distributo kun nur la reverso videbla) la averso di la distributo-lasta – sive per konvenciono irgaltra\lingui{DE}
\artiklo{trunko}%
\vorto{trunko}\senco{1}{Korpo di arboro, sen la branchi nek la radiki}\senco{2}{\textit{(geom.)} Segmento di korpo definita. (Ex. : prismotrunko, kono-trunko, e c.)}\lingui{EFIS}
\artiklo{trunkono}%
\vorto{trunkono}\senco{1}{\textit{(tekn.)} Elemento di la vapor-kaldieri, e di la tanki e cisterni metala, cilindra, qua konsistas ye ringego de fer-tolo}\senco{2}{\textit{(arkitekt.)} Parto infra di fusto}
\artiklo{truo}%
\vorto{truo} Aperturo qua tra-iras korpo, o qua penetras lu profunde\lingui{F}
\artiklo{trupo}%
\vorto{trupo}\senco{1}{Asemblitaro ek ula quanto de personi qui kuniras, agas ensemble, interkonsente}\senco{2}{Asemblitaro ek animali domestika quan onu edukas, nutras, duktas a la pastureyo kom ensemblo}\lingui{DEF}
\artiklo{trusar}%
\vorto{trusar}\fako{trans.} Eskartar de sulo per faldi (kozo qua pendas)\lingui{F}
\artiklo{trusto}%
\vorto{trusto}\fako{komerco} Sindikato de spekulisti, fondita por plugrandigar la preco di varo o di acioni, per akaparar li por pose vendar li granda-preco\lingui{DEFIS}
\artiklo{truto}%
\vorto{truto}\fako{zool.} Fisho ek la genero « salmoni », kun pelo makuletoza, di qua la karno esas tre delikate saporoza\latina{salmo fario}\lingui{EFIS}
\artiklo{truvero}%
\vorto{truvero}\fako{historio literat.} Poeto qua kompozis per la linguo olima di la nordo-regioni di Francia\lingui{DEFIS}
\artiklo{tu}%
\vorto{tu}\fako{gram.} Pronomo personala, vokativo, uzata kande onu parolas o skribas a persono kun qua onu relatas pasable intime (pro la ne-ceremonialeso di ta pronomo); a persono kun qua onu relatas intime e mem tre intime; a persono qua esas tre yuna – ed a persono quan onu desprizas\lingui{DEFIRS}
\artiklo{tualetar}%
\vorto{tualetar}\fako{netrans.} Lavar, razar, kuafar, vestizar ed ornar su, generale en kabineto\lingui{DEFIRS}
\artiklo{tuberkloso}%
\vorto{tuberkloso}\fako{patol.} Morbo infektiva e kontagiema, komuna a la homo ed a la animali, imputata a la bacilo quan deskovris Koch, ed a granulari partikulara\lingui{DEFIRS}
\artiklo{tuberkulino}%
\vorto{tuberkulino}\fako{medic.} Extraktajo glicerinoza – facita per varmigo e steriligita – de kultivaji ek tuberkloso-bacili de rasi inter-diferanta
\artiklo{tuberkulo}%
\vorto{tuberkulo}\fako{bot.}\senco{1}{Aglomerajo pulpoza, qua konsistas (en la maxim multa kazi) ye fekulo e qua kreskas, che ula planti, ye la extremajo di lia radiki o rami sub-sula (*kartofli – terpomi – patati, topinamburi, e c.)}\senco{2}{\textit{(anat.)} Protuberanco fibroza, kartilagoze, e c., di ula parti di la organismo homala}\lingui{DEFIRS}
\artiklo{tubero}%
\vorto{tubero}\fako{bot.} Protuberanco quan produktas, sur stipo, sur trunko, la interplekturo di la fibri e la influro di la tisuo celulara\lingui{EFIS}
\artiklo{tuberoso}%
\vorto{tuberoso}\fako{bot.} Planto kun floro blanka, tre odoroza\latina{polyanthus tuberosa}\lingui{DEFIRS}
\artiklo{tubo}%
\vorto{tubo}\fako{tekn.} Cilindro kava, naturala od artificala, quan onu uzas por pasigar, ula-direcione, liquidi, fluidi\lingui{EFIS}
\artiklo{tubulo}%
\vorto{tubulo}\fako{tekn.}\senco{1}{Tubaro qua duktas fluido a la organo qua uzos olca}\lingui{DEFS}\subvorto{en automobilo}{tubulo di alimento}{Tubo qua duktas la mixuro exploziva, de la karburatoro til la cilindri di la motoro}{2}{}
\artiklo{tuerta}%
\vorto{tuerta} Di qua nur un okulo esas vidiva\lingui{S}
\artiklo{tufo}%
\vorto{tufo} Fasko buketatra, densa, de hari, plumi, pili, flori, planti, e c.\lingui{EF}
\artiklo{tukano}%
\vorto{tukano}\fako{zool.} Grosa ucelo, en Brazilia, kun plumaro diversekolora\latina{ramphastus}\lingui{DEFIS}
\artiklo{tuko}%
\vorto{tuko} Peco de texuro (maxim-multa-kaze ek lino o kotono) uzata heme por vishar, envelopar, kovrar, e c.\lingui{D}
\artiklo{tulipo}%
\vorto{tulipo}\fako{bot.}\senco{1}{Planto ek la familio « liliacei », kun radiko bulbatra, floro ovoforma, diversa-kolora}\senco{2}{La floro de ta planto}\latina{tulipa}\lingui{DEFIRS}
\artiklo{tulo}%
\vorto{tulo}\fako{tekn.} Texuro qua konsistas ye reto tre desdensa, ek fili tre tenua, de silko, lino, kotono\lingui{DEFIRS}
\artiklo{tumefaktar}%
\vorto{tumefaktar}\fako{trans.) (patol.} Igar organo inflesar morbale\lingui{EFIS}
\artiklo{tumoro}%
\vorto{tumoro}\fako{patol.} Influro morbala en parto di la organismo\lingui{DEFIS}
\artiklo{tumulo}%
\vorto{tumulo}\senco{1}{Granda amaso artificala ek tero o petri, konoforma o piramido-forma, quan onu erektis olim sur sepulteyo}\senco{2}{Monteto de tero, sola, en lando plana, di qua la somito esas platforma}\lingui{EFIS}
\artiklo{tumultar}%
\vorto{tumultar}\fako{netrans.) (aludante personi asemblita} Agitar su, bruisoze, sen-ordine, generale pro granda nekontenteso e kom ensemblo. – \textit{(anke metaf.)}\lingui{DEFIS}
\artiklo{tunelo}%
\vorto{tunelo}\fako{tekn.} Paseyo subsula, facita sub rivero; sub fluvio; sub parto de maro, qua quaze, klemesas inter du terolangi; sub monto; tra monto, e c.\lingui{DEFIRS}
\artiklo{tuniko}%
\vorto{tuniko}\senco{1}{\textit{(epoki antiqua)} Subvesto}\senco{2}{\textit{(armeo)} Redingoto kurta di uniformo}\senco{3}{Nomo di ula peli tenua, di ula gaini membranatra (ex. : la tuniki di la okulo, di la arterii, di la onyono)}\lingui{DEFIRS}
\artiklo{tuno}%
\vorto{tuno} Mezur-unajo (1.000 kilogrami) por evaluar la kargajo tota di la vari sur navo, sur vagono, sur kamiono automobila\lingui{DEFIRS}
\artiklo{tupio}%
\vorto{tupio} Ludilo puerala, konatra, quan onu rotacigas sur lua pinto fera; onu startas lu per kordeto spulita an-cirkume – e quan onu desspulas tre rapide, per bruska tiro di la kordeto\lingui{F}
\artiklo{turbano}%
\vorto{turbano}\senco{1}{Kapo-vesto virala, che la Orientani – tuko spulita cirkum la kapo}\senco{2}{Kapo-vesto mulierala por la bali, qua imitas la turbano virala}\lingui{DEFIRS}
\artiklo{turbino}%
\vorto{turbino}\fako{tekn.} Roto trogoza od helicoza, qua kaptas la energio de aquo fluanta o falanta, de vaporo spricanta, de gaso, e qua transmisas ta energio per ax-arboro\lingui{DEFIRS}
\artiklo{turbo}%
\vorto{turbo} Quanto tre granda de personi ordinara\lingui{FIS}
\artiklo{turboto}%
\vorto{turboto}\fako{zool.} Marfisho, ek la familio « pleuronekti », di qua la karno saporas delikate\latina{rhombus maximus}\lingui{DEFR}
\artiklo{turdo}%
\vorto{turdo}\fako{zool.} Ucelo ek la klaso « paseri dentirostra », qua formacas speco ek la genero « merlo », kun plumaro blanka e bruna\latina{turdus}
\artiklo{turgecar}%
\vorto{turgecar}\fako{netrans.} \textit{(fiziol.)} \textit{(aludante ula organi, precipue le sexuala)} Inflesar e rigideskar da la adfluado di sango od altra fluidi\lingui{DEFIS}
\artiklo{turkezo}%
\vorto{turkezo} Lapido matida, bele blua\lingui{DEFIS}
\artiklo{turmalino}%
\vorto{turmalino}\fako{mineral.} Mineralo silicoza, qua polarigas lumo\lingui{DEFIRS}
\artiklo{turmo}%
\vorto{turmo}\senco{1}{Konstrukturo cilindra o plur-edra, tre alta kompare a la dimensioni di lua bazo, sola, o salianta ek edifici vicina, e qua dominacas urbo, kastelo. (ex. : la turmo da Eiffel, en Paris)}\senco{2}{\textit{(shako-ludo)} Piono turmo-forma, en la angulo di la shako-plako}\lingui{D}
\artiklo{turnar}%
\vorto{turnar}\fako{trans.} Modifikar la pozeso per kelka rotacigo, tale ke altra facio, altra edro, surfaco til-lore nevidebla, divenas videbla, efikaria\lingui{EF}
\artiklo{turnirar}%
\vorto{turnirar} \textit{(netrans.)} \textit{(olim)} \textit{(aludante la *chevalieri, en festo militala)} Interkombatar plee, per armi uzata polite, interemule, pri forteso, habileso, talento\lingui{DEFIRS}
\artiklo{turo}%
\vorto{turo} Voyajo segun parkuro plu o min cirklokurvatra, maxim-multa-kaze por sua plezuro\lingui{DEFR}
\artiklo{turto}%
\vorto{turto} Bloko ek la reziduo de la grani di planti oleoza, uzata por la nutro di la bruti\lingui{F}
\artiklo{turturo}%
\vorto{turturo}\fako{zool.} Speco di kolombo plu mikra, di qua la rukulo esas plendatra\latina{columba turtur}\lingui{DFIS}
\artiklo{tusar}%
\vorto{tusar}\fako{netrans.} Esforcar por des-obstruktar la respir-organi, per expiro-stroki konvulsatra, sukusoza\lingui{FIS}
\artiklo{tushar}%
\vorto{tushar}\fako{trans.} Kontakteskar, kontaktar, kun ulu, ulo (quaze por prenar, en ula kazi)\lingui{EFIS}
\artiklo{tusilajo}%
\vorto{tusilajo}\fako{bot.} Planto perena, kun stipo squamoza – folio di qua la suba facio esas blankatra; lua kapituli, flava, e la folii, uzesas kom tuso-cesigivo\latina{tussilago}\lingui{FSL}
\artiklo{tutelar}%
\vorto{tutelar}\fako{trans.} Protektar la persono, la havaji di minoro, di persono a qua la lego interdiktas la dispozo sen-limito di sua havaji (pro imbecileso, dementeso, furiikeso)\lingui{EFIS}
\artiklo{tuvar}%
\vorto{tuvar}\fako{trans.) (navig.} Avancigar (navo) per (de la navo) tirar sur kablo imersita en fluvio, rivero, kanalo, e di qua la altra extremajo esas fixigita per amaro an la punto quan onu volas atingar\lingui{F}
\artiklo{tuyero}%
\vorto{tuyero}\fako{tekn.} Vestizuro metala ye qua esas garnisita, por protektar la masonuro di la parieto, la aperturo ye la parto infra e laterala di fornego, por recevar la tubo o beko di la suflili\lingui{F}
\artiklo{tuyo}%
\vorto{tuyo}\fako{bot.} Arboro ek la familio « kupreinei », di qua la ligno, veinoza, uzesas da la moblifisti\latina{thuya}\lingui{DEFIS}
\end{multicols}
\sekciono{U}
\begin{multicols}{2}
\artiklo{ube}%
\vorto{ube}\fako{adverbo} En ta loko; en qua loko?\lingui{FL}
\artiklo{ubiqua}%
\vorto{ubiqua} Qua esas prezenta en omna loki\lingui{EFIS}
\artiklo{ucelo}%
\vorto{ucelo}\fako{zool.} Un de la animali ek la klaso de vertebrozi, ovipara, kun du pedi, la membri avana alo-forma, kovrita per plumi, beko korno-harda, sen denti\lingui{FI}
\artiklo{uf!}%
\vorto{uf!} Interjeciono qua expresas la alejeso, la sento di liberigeso pri ulo sufriganta, peniganta, tedanta
\artiklo{ukazo}%
\vorto{ukazo}\senco{1}{Edikto da la caro}\senco{2}{\textit{(metaf.)} Edikto di qua la autoritato esas egardenda absolute}\lingui{DEFIRS}
\artiklo{ula}%
\vorto{ula} Qua ne determinesas – nek mem esas determininda\lingui{L}
\artiklo{ulano}%
\vorto{ulano}\fako{olim} Soldato ek regimento de lancieri, en ula landi\lingui{DEFIRS}
\artiklo{ulcero}%
\vorto{ulcero}\fako{patol.} Lezuro di tisuo qua, pro inherar la konstituciono, tendencas adextensar su sencese\lingui{EFIS}
\artiklo{ulexo}%
\vorto{ulexo}\fako{bot.} Arbusto qua kreskas en la erikeyi\latina{ulex}
\artiklo{ulmo}%
\vorto{ulmo}\fako{bot.} Varietato di la fraxino\latina{ulmus}\lingui{DEFIS}
\artiklo{ulno}%
\vorto{ulno} Olima mezur-unajo por la texuro. (En Paris = 1,18 metro)\lingui{FIS}
\artiklo{-ulo}%
\vorto{-ulo} Sufixo qua indikas la maskulo di la animalo. (Ta sufixo uzesas nur kande lu esas nekareebla)
\artiklo{*ultima}%
\vorto{*ultima} Qua ne sequesos da *nexta; pos qua nula sama eventos\lingui{EFIS}
\artiklo{ultimato}%
\vorto{ultimato}\fako{diplomac.} Dico serio-lasta, *ultima, rezolvo quan onu deklaras ne-revokebla\lingui{DEFIRS}
\artiklo{ultramaro}%
\vorto{ultramaro} Farbo bele-azurea\lingui{DEFIRS}
\artiklo{ultramikroskopio}%
\vorto{ultramikroskopio} Ensemblo ek la procedi di explori pri la studio di la korpuskuli tante mikra-dimensiona ke li ne esas observebla per la mikroskopo ordinara\lingui{DEF}
\artiklo{ultramikroskopo}%
\vorto{ultramikroskopo}\fako{tekn.} Aranjuro partikulara di lumizo, adaptita a mikroskopo, qua posibligas vidar la existo e la movi di objekti d1 uui 1 ein i, kam olti observebla ye la plugrandigo maxima di mikroskopo ordinara\lingui{DEF}
\artiklo{ultramontana}%
\vorto{ultramontana} Qua preferas obediar la papo kam la autoritatozi civila di sua lando\lingui{DEFIRS}
\artiklo{ultre}%
\vorto{ultre} Adjunte ad\lingui{DEFIS}
\artiklo{ulular}%
\vorto{ulular}\fako{netrans.} \textit{(aludante la nokt-uceli)} Kriar jemante. – Kriar jemante, quale la nokt-uceli\lingui{DEFIS}
\artiklo{umbelifera}%
\vorto{umbelifera}\fako{bot.} Di qua la flori dispozesas umbele\lingui{DEFIS}
\artiklo{umbelo}%
\vorto{umbelo}\fako{bot.} Modo di infloresenco, forme di parasuno apertata\lingui{EFIS}
\artiklo{umbiliko}%
\vorto{umbiliko}\fako{anat.} Cikatruro rondatra quan havas la mamiferi ye la centra parto di abdomino ube pasis kordono nutrala che la feto\lingui{EFIS}
\artiklo{umbro}%
\vorto{umbro} Okro brune-redatra, uzata kom farbo\lingui{DEFIR}
\artiklo{-umo}%
\vorto{-umo} Sufixo di qua la senco esas nepreciza (quale la prepoziciono « ye »), rare uzata e, kande uzata, nur en ula kazi quin sancionis la Akademio di Ido
\artiklo{un}%
\vorto{un}\fako{matem.} La nombro kardinala : 1\lingui{EFIS}
\artiklo{-un-}%
\vorto{-un-} Sufixo qua anuncas : « unajo, fragmento »
\artiklo{unanima}%
\vorto{unanima} Qua expresas la interkonsento perfekta di plura personi\lingui{EFIS}
\artiklo{unciala}%
\vorto{unciala}\fako{paleogr.} Granda literi, uzata en la epigrafi, la tituli di libri, e c., analoga a la mayuskuli moderna, ma kun formo plu ronda\lingui{DEFIS}
\artiklo{uncionar}%
\vorto{uncionar}\fako{trans.) (religio katol.} Konsakrar per frotar ye oleo santa\lingui{EFIS}
\artiklo{unco}%
\vorto{unco}\fako{zool.} Varietato di pantero, di jaguaro\latina{felis onca}
\artiklo{unglo}%
\vorto{unglo} Lamo kornatra, sur la facio dorsa di la extremajo di la fingri, che la homo e che multa animali spinoza\lingui{EFIS}
\artiklo{unguento}%
\vorto{unguento} Substanco medikamenta, mola, por apliki extera\lingui{EFIS}
\artiklo{uniforma}%
\vorto{uniforma} Di qua la formo, la eso-maniero, esas identa en sua tota extensajo, dum sua tota duro.\lingui{DEFIS}\subvorto{}{uniformo}{Kostumo di qua la formo, la koloro, esas sama pri omna personi qui apartenas a ta o ca korpo}{}{}
\artiklo{unika}%
\vorto{unika} Sola ek sua genero, sorto, kategorio\lingui{EFIS}
\artiklo{unikorno}%
\vorto{unikorno}\senco{1}{\textit{(fablo)} Animalo kun korpo di kavalo, kapo di cervo e korno unika meze di lua fronto}\senco{2}{\textit{(astron.)} Stelaro equatorala, konsistante ye kelka steli nur poke briloza, situita inter la Granda Hundo, la Mikra Hundo e la Hidro}\lingui{EFS}
\artiklo{unikursala}%
\vorto{unikursala}\fako{geom.} Kurvo pri qua la koordinati esas expresebla racionale per parametro variiva
\artiklo{unionar}%
\vorto{unionar}\fako{trans.} Interligar ne-materiale – asociar per ligilo politikala, religiala, komercala, e c.\lingui{EFIS}
\artiklo{unipola}%
\vorto{unipola}\fako{elektro}\senco{1}{Qua havas nur un polo}\senco{2}{Qua duktas nur un elektro-fluo}
\artiklo{unisono}%
\vorto{unisono}\fako{muziko} Konsonanco ek la voci, ek la instrumenti, qui emisas sam-instante la sama noto. — \textit{anke metaf}\lingui{DEFIRS}
\artiklo{unitara}%
\vorto{unitara} Qua retroduktas la kozi ad uneso\lingui{DEFIRS}
\artiklo{univalva}%
\vorto{univalva}\fako{zool.) (pri molusko} Di qua la konko konsistas ye nur un shutro
\artiklo{universitato}%
\vorto{universitato} Doceyo supra-grada, ek grupi (fakultati) de qui singla grantas la gradi di bachelereso, licenciereso, doktoreso\lingui{DEFIRS}
\artiklo{universo}%
\vorto{universo} Ensemblo ek omna existanti\lingui{DEFIS}
\artiklo{unto}%
\vorto{unto} Materio per qua onu lubrifikas la rotaro, la ingranaji di mashino, la axi di la roti, quan igis nigran e densan la enmixo di polvo e di la parti metala quin desagregis la frotado\lingui{IS}
\artiklo{upupo}%
\vorto{upupo}\fako{zool.} Genero de paseri ek la familio « tenuirostri », qua havas, sur sua kapo, tufo ek duopla rango de plumi rufa kun bordo nigra\latina{upupa}\lingui{EFI}
\artiklo{uragano}%
\vorto{uragano} Tempesto quan genitas ventoflui inter-opozata, qui formacas vortici\lingui{EFIRS}
\artiklo{urano}%
\vorto{urano}\fako{kemio} Korpo simpla, metala, extraktita de U₃O atompezo 238,2 – materio uzata por la fabriko di la atombombi\simbolo{U}\lingui{DEFIRS}
\artiklo{uranografio}%
\vorto{uranografio} Deskripto di cielo\lingui{DEFIRS}
\artiklo{urato}%
\vorto{urato}\fako{kemio} Kombinuro ek urin-acido kun bazo
\artiklo{urbario}%
\vorto{urbario} Taxo feudala
\artiklo{urbo}%
\vorto{urbo} Agregajo grandega de domi habitata, dispozita alonge e borde di stradi, e qua havas (normale) kom limitilo cirkum-muro o remparo-cirkuito\lingui{EFIS}
\artiklo{urdar}%
\vorto{urdar}\fako{trans.} Preparar (la texo) per tensar la fili warpala di la texuro (futura)\lingui{FIS}
\artiklo{uremio}%
\vorto{uremio}\fako{patol.} Akumulajo de ureo en sango\lingui{DEFS}
\artiklo{ureo}%
\vorto{ureo}\fako{kemio} Substanco precipua di urino : NH₂.CO.NH₂\lingui{EFIS}
\artiklo{uretero}%
\vorto{uretero}\fako{anat.} Kanalo membrana qua duktas urino, de la reni aden la veziko
\artiklo{uretrito}%
\vorto{uretrito}\fako{patol.} Inflamuro di uretro
\artiklo{uretro}%
\vorto{uretro}\fako{anat.} Kanalo exkrecala di urino\lingui{DEFIS}
\artiklo{uretroskopo}%
\vorto{uretroskopo} Varietato di endoskopo, uzata por explorar la domajeso di la mukozo di uretro
\artiklo{uretrotomio}%
\vorto{uretrotomio} Incizo di la parti streta di uretro
\artiklo{urino}%
\vorto{urino} Liquido exkrementa quan sekrecas la reni, duktata da uretero aden la veziko de ube lu ekpulsesas per la kanalo uretra\lingui{DEFIRS}
\artiklo{urjar}%
\vorto{urjar}\fako{netrans.} Ne povar sen domajeso tardigesar od ajornesar\lingui{EFIS}
\artiklo{urno}%
\vorto{urno}\senco{1}{\textit{(en la epoki antiqua)} Granda vazo quan la Antiqui uzis por cherpar aquo, por enklozar la cindri de la kadavri homala, por rekoliar la voto-buli en la tribunali, en la asemblitari politikala, e c.}\senco{2}{\textit{(nun-epoke)} Vazo, buxo, aden qua onu rekolias la voto-bilieti}\lingui{DEFIRS}
\artiklo{-uro}%
\vorto{-uro} Sufixo qua indikas « la rezultajo (materiala) de la faco, de la ago »\lingui{EFL}
\artiklo{urogalo}%
\vorto{urogalo}\fako{zool.} Nomo vulgara di la tetraso\latina{tetrao urogalus}\lingui{I}
\artiklo{uroxo}%
\vorto{uroxo}\fako{zool.} Bovo sovaja, en la foresti di Lituania, la maxim granda quadripedo di Europa\latina{bos urus}\lingui{DEFS}
\artiklo{urso}%
\vorto{urso}\fako{zool.} Mamifero karnavida, plantigrada, di qua la pili esas densa\latina{ursus}\lingui{EFIS}
\artiklo{urtikario}%
\vorto{urtikario}\fako{patol.} Inflamuro di la pelo, qua eventas acese, forme di makuli – qua igas pruritar quale de urtiko-piki\lingui{EFIS}
\artiklo{urtiko}%
\vorto{urtiko}\fako{bot.} Planto sovaja di qua la folii, la stipo, pikas sente di bruleso\latina{urtica}\lingui{EFIS}
\artiklo{-us}%
\vorto{-us}\fako{gram.} Dezinenco di la verbo-kondicionalo
\artiklo{Usa}%
\vorto{Usa}\fako{geogr. politikala} La Stati Unionita di Nord-Amerika (di qui la urbo centrala esas Washington, D.C.)\lingui{DEFIS}
\artiklo{ushero}%
\vorto{ushero} Oficiero judicistala, di qua la tasko konsistas ye notifikar la akti procesala\lingui{EFS}
\artiklo{utardo}%
\vorto{utardo}\fako{zool.} Granda ucelo steltiera, formo-vicina a strucho\latina{otis}\lingui{FIS}
\artiklo{utensilo}%
\vorto{utensilo}\senco{1}{Labor-instrumento quan laboristo manuala uzas por la praktiko di sua mestiero}\senco{2}{Vazo, instrumento, por la hemo-labori diversa – implemento}\lingui{DEFIS}
\artiklo{utero}%
\vorto{utero}\fako{anat.} Organo di la filiifo o yunifo che la mamiferini – organo kava, kun parieto muskulatra, di qua la internajo komunikas kun vagino\lingui{DEFIS}
\artiklo{utila}%
\vorto{utila} Uzata efike por agar o facar ulo\lingui{EFIS}
\artiklo{utopio}%
\vorto{utopio}\senco{1}{Konceptajo imaginala pri guverno-sistemo ideala}\senco{2}{Konceptajo ideala, ne-realigebla}\lingui{DEFIRS}
\artiklo{utrikulo}%
\vorto{utrikulo}\fako{bot.} Singla ek mikra lojii en la tisuo celulara di la planti\lingui{DEFIS}
\artiklo{uveo}%
\vorto{uveo}\fako{anat.} Membrano interna, garnisita ye materio nigra (pigmento) qua kovras tapete la okulo\lingui{EFIS}
\artiklo{uvulo}%
\vorto{uvulo}\fako{anat.} Apendico karna qua pendas de la infrajo di la palato-velo, en la enireyo di la guturo\lingui{EFIS}
\artiklo{-uyo}%
\vorto{-uyo} Sufixo qua signifikas la « recipiento, kontenanto »
\artiklo{uzar}%
\vorto{uzar}\fako{trans.} Igar (la objekto) plear sua rolo por atingar la skopo vizata\lingui{DEFIS}
\artiklo{uzufruktar}%
\vorto{uzufruktar}\fako{trans.) (yuro-cienco} Havar la yuro uzar kozo quan altru proprietas, e recevar por su la frukti, la produkturi de lu, sen darfar alienar lu, nek destruktar lu\lingui{EFIS}
\artiklo{uzukaptar}%
\vorto{uzukaptar}\fako{trans.) (en la yuro-cienco di la Romani antiqua} Proprieteskar per nur posedo dum tempo pasable multa, per nur uzo\lingui{EFIS}
\artiklo{uzurar}%
\vorto{uzurar}\fako{netrans.) (pri presto} Postular interesti qui ecesas la procento legala, normala\lingui{EFIS}
\artiklo{uzurpar}%
\vorto{uzurpar}\fako{trans.} Proprigar a su, sen yuro (domeno, autoritato, titulo)\lingui{DEFIRS}
\end{multicols}
\sekciono{V}
\begin{multicols}{2}
\artiklo{vabo}%
\vorto{vabo} Bloko de vaxo, ek alveoli, aden qui la abeli depozas sua mielo\lingui{D}
\artiklo{vacilar}%
\vorto{vacilar}\fako{netrans.) (aludante la gambi, la kapo}\senco{1}{Agitesar da quaza tremo, ta-sinse o ca-sinse}\senco{2}{\textit{(metaf.)} Hezitar multe pri decidar o rezolvar}\lingui{EFIS}
\artiklo{vacinio}%
\vorto{vacinio}\fako{bot.} Genero de planti ek le « vacinei » – qua produktas beri nigratra, acideta\latina{vaccinium}\lingui{L}
\artiklo{vacino}%
\vorto{vacino}\fako{patol.} Morbo eruptala, kontagiala, propra ye la bovini\lingui{EFIR}
\artiklo{vadar}%
\vorto{vadar}\fako{netrans., tra} Trairar rivereto, rivero, ube la aquo esas tale basa ke onu povas tra-pedirar lu\lingui{EIS}
\artiklo{vaflo}%
\vorto{vaflo} Kuko facita en mulduro di qua la plaki esas dividita ye celuli alveolo-simila
\artiklo{vagar}%
\vorto{vagar}\fako{netrans.} Vivar aventure, irante sen skopo, hazarde\lingui{DEFIS}
\artiklo{vagino}%
\vorto{vagino}\fako{anat.} Kanalo qua departas de la vulvo ed abutas a la utero\lingui{DEFS}
\artiklo{vagono}%
\vorto{vagono} Veturo uzata sur la relvoyi por transportar pasajeri e vari; anke sur la tramvoyi\lingui{DEFRS}
\artiklo{vajisar}%
\vorto{vajisar}\fako{netrans.) (aludante infanteto jus-naskinta} Kriar\lingui{FIS}
\artiklo{vakancar}%
\vorto{vakancar}\fako{netrans., dure di} Juar la repozo-periodo, relative longa, dum qua interruptesas la lerno en la skoli, la laborado en la tribunali, kontori, fabrikeyi, e c.\lingui{DEFIRS}
\artiklo{vakar}%
\vorto{vakar}\fako{netrans.) (aludante ofico, habiteyo, plaso} Ne okupesar\lingui{DEFIRS}
\artiklo{vakua}%
\vorto{vakua} Di qua la spaco interna ne okupesas (da to quon lu normale kontenas)\lingui{DEFIS}
\artiklo{vakuolo}%
\vorto{vakuolo}\fako{biol.} Kavajo en la protoplasmo, per la prezenteso di multa liquidi qui ne povas mixesar kun la protoplasmo\lingui{DEFIS}
\artiklo{valenco}%
\vorto{valenco}\fako{kemio} La grado di saturebleso di atomo di korpo per hidro-atomi\lingui{EF}
\artiklo{valerianelo}%
\vorto{valerianelo}\fako{bot.} Planto herbatra, kun folii spatulo-forma, quan la homo manjas forme di salado\latina{valerianella olitoria}\lingui{EFI}
\artiklo{valeriano}%
\vorto{valeriano}\fako{bot.} Planto farmaciala, ek la brancho « dikotiledoni monopetala », kun flori rozea e flori preske kompozaja, di qua la radiko esas kontre-spasma, kontre-febra\latina{valeriana}\lingui{EFIRS}
\artiklo{valida}%
\vorto{valida}\senco{1}{\textit{(fiziol.)} Qua disponas tote sua vigoro}\senco{2}{\textit{(yuro-cienco)} Qua havas omno postulata da la legi por produktar sua efiko; qua havas la valoro postulata por admisesar legitime}\lingui{DEFIS}
\artiklo{valizo}%
\vorto{valizo} Mikra voyajo-kofro quan onu povas kunportar manue\lingui{EFIS}
\artiklo{valo}%
\vorto{valo}\fako{geogr.}\senco{1}{Fundo inter du o plura monti}\senco{2}{Baseno di fluvio o rivero}\lingui{EFIS}
\artiklo{valorar}%
\vorto{valorar}\fako{trans.}\senco{1}{Esar evaluita kom la equivalanto di ta o ca pekunio-quanto}\senco{2}{Esar evaluita kom havanta ta o ca qualeso, ta o ca merito}\lingui{EFIS}
\artiklo{valsar}%
\vorto{valsar}\fako{netrans.} Dansar jirante segun ario tri-tempa\lingui{DEFIRS}
\artiklo{valvo}%
\vorto{valvo}\senco{1}{\textit{(zool.)} Singla de la du shutri ye qua konsistas la konko di la moluski qui desklozesas o riklozesas}\senco{2}{\textit{(bot.)} Singla ek la peci di perikarpo sika}\senco{3}{\textit{(tekn.)} Klozilo moviva, baskuliva, quan preso-resorto klozas, per aplikar lu kontre aperturo, e quan forco plu granda, per efikar inverse, povas apertar dum kelketa tempo, ma qua cesas quik kande la forco cesas efikar}\lingui{EFIS}
\artiklo{vampiro}%
\vorto{vampiro}\senco{1}{Ento imaginala, quan onu kredas vidar ekirar de tombo por sugar la sango di la vivanti dum lia dormo}\senco{2}{\textit{(metaf.)} Persono qua richigas su detrimente, leze, di sua kunhomi}\lingui{DEFIRS}
\artiklo{vamso}%
\vorto{vamso} Olima vesto virala, manikoza, qua klemis la tayo e kovris til la genui\lingui{D}
\artiklo{vana}%
\vorto{vana}\senco{1}{Qua havas nur la aspekto, ma nula realeso nek valoro}\senco{2}{Qua ne sucesas efikar}\lingui{EFIS}
\artiklo{vanado}%
\vorto{vanado}\fako{kemio} Metalo blanka, kristalatra, denseso 5,7, atom-pezo 51, extraktita precipue de (V₂O₈Pb₂)₃ + PbCl₂
\artiklo{vandalo}%
\vorto{vandalo} Desestimero, destruktero di la monumenti di civilizeso, di kulturo\lingui{DEFIRS}
\artiklo{vanelo}%
\vorto{vanelo}\fako{zool.} Ucelo ek la klaso « steltieri », kun plumaro nigra\latina{vanellus}\lingui{FIL}
\artiklo{vango}%
\vorto{vango}\senco{1}{\textit{(anat.)} Che la homo : parto di la vizajo, singla-latere, de-sub la okulo til la mentono}\senco{2}{Parto korespondanta di la kapo, inter la nazo, la boko e la oreli (che la mamiferi); inter la bazo di la beko e la boko (che la uceli e la fishi); inter la okulo e la mandibulo (che la insekti)}\senco{3}{\textit{(tekn.)} Parto di objekto plu o min neprecise komparebla, per lua formo e simetreso, a la vangi}\lingui{D}
\artiklo{vanilo}%
\vorto{vanilo}\fako{bot.} Frukto di planto sarmentoza di Mexikia, di qua la shelo odoras aromatoze\latina{vanilla}\lingui{DEFIRS}
\artiklo{vanitato}%
\vorto{vanitato} La deziro ostentar\lingui{EFIS}
\artiklo{vanto}%
\vorto{vanto}\fako{navig.} Grosa kordego skalatra qua fixigas masto e mantenas lu vertikale\lingui{DR}
\artiklo{vaporo}%
\vorto{vaporo} Asemblajo ek la guteti, preske neperceptebla, qui su elevas de la surfaco di la liquidi, de la korpi humida, per la efiko da kaloro, da la diminuto di la preso atmosferala\lingui{EFIS}
\artiklo{varango}%
\vorto{varango}\fako{navig.} Ligno-peco kurvatra, fixigita per lua mezo, sur la kilio di navo, e qua formacas la fundamento di la spantaro ye qua konsistas la karpenturo di la navo\lingui{DFS}
\artiklo{variar}%
\vorto{variar}\fako{netrans.}\senco{1}{Submisar (serio de fakti) a chanji intersucedanta}\senco{2}{\textit{(netrans.)} Chanjar intersucedante. – Dicesas pri kozi qui, kun traiti komuna, interdiferas}\lingui{DEFIS}
\artiklo{varicelo}%
\vorto{varicelo}\fako{patol.} Morbo infektala, kontagiala, qua aparas precipue che la infanti e la pueri, e qua karakterizesas esencale da erupto bulo-forma, intervaloze, e qua desaparas pos kelka dii\lingui{DEFIS}
\artiklo{varietato}%
\vorto{varietato}\fako{zool. e bot.} Singla ek la grupi interdiferanta de individui ek un speco, qui havas ula traiti extera komuna\lingui{DEFIS}
\artiklo{variko}%
\vorto{variko}\fako{patol.} Dilaturo permananta, quan produktas akumulajo de sango en veino (en la gambi, maxim-multa-kaze)\lingui{EFIS}
\artiklo{variolo}%
\vorto{variolo}\fako{patol.} Morbo infektala, kontagiala ed epidemiala, quan genitas mikroparazito ankore nekonocata, e qua havas kom karakterizivo erupto postuletoza qua abutas a pusifo\lingui{EFIS}
\artiklo{varma}%
\vorto{varma}\senco{1}{Di qua la temperaturo esas alta}\senco{2}{Qua igas parto di la kapo sentar su kaloroza}\lingui{DE}
\artiklo{varo}%
\vorto{varo} To quo esas la objekto di komerco\lingui{DE}
\artiklo{varsar}%
\vorto{varsar}\fako{trans.}\senco{1}{Igar (liquido, sablo, grani) fluar, glitar, per inklinar la vazo qua kontenas li}\senco{2}{\textit{(financo)} Igar pekunio pasar, page, de kaso aden altra kaso}\lingui{FI}
\artiklo{vartar}%
\vorto{vartar}\fako{trans.} Permanar til la arivo (di ulu, di ulo)\lingui{DE}
\artiklo{vasalo}%
\vorto{vasalo}\fako{feudismo} La persono qua dependas de sinioro pro feudo\lingui{DEFIRS}
\artiklo{vasko}%
\vorto{vasko} Baseno, forme di kupo, qua recevas la aquo de fontano\lingui{FI}
\artiklo{vaskulo}%
\vorto{vaskulo}\fako{anat.} Kanalo en qua cirkulas la sango, la limfo, e c., che la animali — la sapo, che la vejetanti\lingui{DEFIS}
\artiklo{vasta}%
\vorto{vasta} Qua su extensas til for sua centro. Qua kovras multa spaco\lingui{EFIS}
\artiklo{vato}%
\vorto{vato} Kotono preparita por esor strato mola, dolca – uzata diversa-skope\lingui{DFIR}
\artiklo{vaxo}%
\vorto{vaxo} Substanco flavatra, mola, fuzebla, quan produktas la abeli, ed ek qua li facas alveoli di la abeluyo\lingui{DE}
\artiklo{vazelino}%
\vorto{vazelino}\fako{farmacio} Reziduo de la distilo di la petrolo de Amerika, konsistanta ye mixuro de hidrokarburi (CH₂)\lingui{DEFIRS}
\artiklo{vazo}%
\vorto{vazo} Recevuyo ek ligno, metalo, glaso, tero, porcelano, fayenco, e c., varie-granda, varie-forma, destinita kontenor omna speco de substanci, liquida o solida (liquori, nutrivi, frukti, flori, e c.)\lingui{DEFIRS}
\artiklo{ve!}%
\vorto{ve!} Klameto qua expresas la chagreno
\artiklo{vecho}%
\vorto{vecho}\fako{bot.} Genero de leguminosi, « papilionacei », perena, generale klimera, kun folii paripenea, stipitoza, flori axala, solitara o grapana – uzata kom forejo\latina{vicia}\lingui{EFI}
\artiklo{vedeto}%
\vorto{vedeto}\fako{milit.} Mikra milito-navo uzata kom observero. -Barkego movata per vaporo o petrolo, qua servas milito-navo\lingui{DEFI}
\artiklo{vegro}%
\vorto{vegro}\fako{navig.} Singla ek la planki uzata por kovrar la parti interna di la kareno di navo\lingui{DFS}
\artiklo{vehar}%
\vorto{vehar}\fako{netrans.} Irar per veturo, biciklo, navo, aeroplano, fuzeo, relvoyo-treno, e c.\lingui{EFI}
\artiklo{veino}%
\vorto{veino}\senco{1}{\textit{(anat.)} Vaskulo sangala, qua retroduktas la sango a la latero dextra di la kordio, de ube lu pulsesas a la pulmoni por rekuperar la oxigeno quan lu perdis, per la cirkulo arterala, en la parti diversa di la korpo}\senco{2}{\textit{(mineral.)} Strato di mineyo ube onu renkontras la erco}\senco{3}{Traito sinuoza, analoga a la veini quin onu vidas tra la pelo, sur marmoro, e c.}\lingui{EFIRS}
\artiklo{vejetar}%
\vorto{vejetar}\fako{netrans.}\senco{1}{Exekutar la funcioni ye qui konsistas la vivo di la planti}\senco{2}{\textit{(aludante la homi)} Vivar inerte, sen emoci, sen agado (same kam planto)}\lingui{DEFIRS}
\artiklo{vejetarar}%
\vorto{vejetarar}\fako{netrans.} Praktikar la nutro-sistemo (« da Pitagoro ») segun qua supresesas omna speci de karni o lia derivaji ne-mediata\lingui{DEFIRS}
\artiklo{vekar}%
\vorto{vekar}\fako{netrans.} Cesar sua dormo\lingui{DE}
\artiklo{vekto}%
\vorto{vekto} Levero di la balanco, qua ocilas sur sua axo e suportas la pleti\lingui{L}
\artiklo{vektoro}%
\vorto{vektoro}\fako{geom.} Radio qua duktesas (en kurvo fokoza) de foko ad irga punto di la kurvo\lingui{DEFIRS}
\artiklo{velino}%
\vorto{velino}\fako{tekn.} Felo de bovyuno, apretita quale pergameno, ma plu delikata e plu glata, uzata por piktar miniaturi, por skribar\lingui{DEFIR}
\artiklo{velito}%
\vorto{velito}\senco{1}{\textit{(che la Romani antiqua)} Infantriano kun armi lejera}\senco{2}{\textit{(armei di Napoleon Iesma)} Pedo-chasero}\lingui{EFIS}
\artiklo{velkar}%
\vorto{velkar}\fako{netrans.} Cesir esar fresha. \textit{(anke metaf.)}\lingui{D}
\artiklo{velo}%
\vorto{velo} Tuko de stofo per qua la mulieri kovras sua vizajo por shirmar su de la regardi, o por protektar su kontre atmosfero kolda, kontre vento, e c.\lingui{EFIS}
\artiklo{velocipedo}%
\vorto{velocipedo} Sidilo sur roti, quan onu movigas per presar pedalo\lingui{DEFIS}
\artiklo{velodromo}%
\vorto{velodromo} Lico por la konkursi di biciklisti\lingui{EF}
\artiklo{veluro}%
\vorto{veluro} Stofo kun du strati horizontala de fili : la suba (reverso) formacas texuro densa; la supra (averso) esas sole dika (per mikra slingi sekita ed igita interegale longa)\lingui{DEFIS}
\artiklo{venar}%
\vorto{venar}\fako{netrans.) (aludante la persono qua parolas} Vidar (persono, kozo) qua iras o portesas a la loko ube esas la parolante\lingui{FIS}
\artiklo{vendar}%
\vorto{vendar}\fako{trans., ad ulu} Cedar (ulo) (ad ulu) qua proprietos olu, po pekunio-quanto (o po valorajo altra-sorta)\lingui{EFIS}
\artiklo{veneno}%
\vorto{veneno}\senco{1}{Liquido quan sekrecas ula animali per glando specala, e qua, enduktite per morduro o pikuro aden la sango di altra animalo, divenas nocanta o mortiganta}\senco{2}{\textit{(metaf.)} Dico o doktrino danjeroza}\lingui{EFIS}
\artiklo{veneracar}%
\vorto{veneracar}\fako{trans.} Traktar kun respekto profunda, religiala, humila\lingui{EFIS}
\artiklo{venerala}%
\vorto{venerala} Qua koncernas la relati inter-sexua intima\lingui{DEFIRS}
\artiklo{venerdio}%
\vorto{venerdio} La dio kinesma di la semano
\artiklo{veniala}%
\vorto{veniala}\fako{teologio katol.; pri peko} Pardonebla, pardoninda\lingui{EFIS}
\artiklo{venjar}%
\vorto{venjar}\fako{trans., su, ulu, pri, pro} Punisar ulu (qua ofensis) por satisfacar sua rankoro, o la rankoro di ulu\lingui{EFIS}
\artiklo{ventar}%
\vorto{ventar}\fako{netrans.} Dicesas kande, en ta o ca parto di atmosfero, eventas diplaseso plu o min rapida di aeroquanto\lingui{DEFIS}
\artiklo{ventilar}%
\vorto{ventilar}\fako{trans.} Chanjar la aero di klozajo (*kluzajo), aerizar per eventigar fluo di aero\lingui{DEFIRS}
\artiklo{ventilatoro}%
\vorto{ventilatoro}\fako{tekn.} Aparato uzata por ventilar la chambrego di teatro, chambro di hospitalo, e c., o por alimentar fornego ye aero\lingui{EF}
\artiklo{ventrikulo}%
\vorto{ventrikulo}\fako{anat.} Kavajo qua sequas la parto supra di la dextrajo e sinistrajo di la kordio; la dextrala sendas la sango veinala a la pulmoni – la sinistrala, la sango arteriala, a la tota korpo\lingui{DEFIS}
\artiklo{ventro}%
\vorto{ventro}\fako{anat.} Parto di la korpo, kava, qua kontenas la stomako e la intestini. (sinonimo : abdomino)\lingui{EFIS}
\artiklo{ventuzo}%
\vorto{ventuzo}\fako{medic.} Mikra vazo klosho-forma, quan onu aplikas sur ta o ca parto di la korpo e de qua onu ekpulsas la aero, por sublevar la pelo e produktar iriteso lokala deturniva\lingui{FIRS}
\artiklo{Venus}%
\vorto{« Venus »}\senco{1}{\textit{(astron.)} Planeto di la sistemo Sunala, la disto-duesma, kontante de Suno de qua lu distas per (cirkume) 108 milion kilometri}\senco{2}{\textit{(mitol.)} (che la Latini) Deino di Amoro}\lingui{DEFIRS}
\artiklo{venusmonto}%
\vorto{venusmonto}\fako{anat.} Saliajo avan la pubio, super la vulvo, e qua pilozeskas lor pubereso\lingui{F}
\artiklo{vera}%
\vorto{vera} Konforma a to quo esas reala (kontraste kun : imituro, kontrafakturo, finguro)\lingui{EFIS}
\artiklo{verando}%
\vorto{verando}\fako{arkitekt.} Teraso tektoza qua formacas peristilo\lingui{DEFIR}
\artiklo{verbasko}%
\vorto{verbasko}\fako{bot.} Herboro moligiva, uzata por kataplasmi, alta-statura, kun folii simpla, alterna, kovrita ye pili kotonatra\latina{verbascum}\lingui{IL}
\artiklo{verbeno}%
\vorto{verbeno}\fako{bot.} Herboro di qua la stipito quadrata, la folii inter-opozata e dentoza, la infloresenco spikoforma, memorigas la labiacei – uzata por la faco di parfumi o … di ula liquori\latina{verbena}\lingui{DEFIRS}
\artiklo{verbo}%
\vorto{verbo}\senco{1}{\textit{(logiko)} Termino di la propoziciono, qua unionas la atributo a la subjekto}\senco{2}{\textit{(gram.)} Vorto-speco qua expresas la existo, la ago}\senco{3}{\textit{(teol.)} Persono duesma ek la Triuno kristana; la filiulo di Deo, egardata kom la sajo eterna}\lingui{DEFIS}
\artiklo{verda}%
\vorto{verda} Di qua la koloro esas olta di foliaro, di herbo, e c.\lingui{EFIS}
\artiklo{verdigriso}%
\vorto{verdigriso} Nomo vulgara di la kuproxido qua naskas sur la surfaco di la vazi de kupro\lingui{EF}
\artiklo{verdikto}%
\vorto{verdikto}\fako{yuro-cienco} Rezultajo de la deliberi da la jurio, proklamita publike\lingui{DEFIRS}
\artiklo{verdrono}%
\vorto{verdrono}\fako{zool.} Nomo vulgara di mikra pasero konirostra, kun plumaro verda, flava o makulo-nigra\latina{fringilla chloris}\lingui{F}
\artiklo{verdugo}%
\vorto{verdugo} La persono di qua la tasko konsistas ye exekutar la verdikto « kondamnita ad ocidesor », enuncita da la judiciisto, t.e. senkapigar o per gilotino, o per hakilo\lingui{S}
\artiklo{vergo}%
\vorto{vergo} Longa bastoneto, rekta e flexebla, uzata por kastigo-batar, o (kande ornita per ivoro-peci) por ula funcioni di la usheri
\artiklo{verifikar}%
\vorto{verifikar}\fako{trans.} Submisar (ulo) ad exploro, por certeskar (*certeneskar) pri to ke lu facesis od agesis korekte, rigoroze segun la preskriptaji\lingui{DEFIS}
\artiklo{verjuso}%
\vorto{verjuso} Suko acida de vitoberi nematura, uzata kom kondimento\lingui{EF}
\artiklo{verko}%
\vorto{verko} Rezultajo de laboro (pensala, artala, etikala, sociala)\lingui{DE}
\artiklo{vermicelo}%
\vorto{vermicelo}\fako{koquarto} Farino dilutita e petrisita, filigita, quan onu manjas supe\lingui{EFIR}
\artiklo{vermo}%
\vorto{vermo}\fako{zool.} Larvo di insekto, qua su developas en la karno putranta, e quan onu uzas kom lurivo di pesko\lingui{DEFI}
\artiklo{vermuto}%
\vorto{vermuto} Vino flava, en qua onu infuzis absinto-folii\lingui{DEFIRS}
\artiklo{verniero}%
\vorto{verniero}\fako{tekn.} Linealo movebla, per qua onu evaluas la fracioni di la dividuri di skalo gradizita\lingui{EF}
\artiklo{verniso}%
\vorto{verniso} Solvuro ek gumo-rezino en alkoholo, en petrol-esenco, qua formacas, per sikesko, strato glata e briloza, surface di la objekti qua esas indutita ye lu\lingui{EFI}
\artiklo{veroniko}%
\vorto{veroniko}\fako{bot.} Nomo di plura planti ek la grupo « personei », di qui la extremaji konsumesas infuzure\latina{veronica}\lingui{DEFIS}
\artiklo{vers}%
\vorto{vers}\fako{prepoziciono} \textit{(qua regardas)} direcione di… – o (qua quaze regardas) direcione di… (do : sen movo ad)\lingui{F}
\artiklo{versiono}%
\vorto{versiono} Maniero ye qua persono komprenas lekte, naracas, raportas, explikas ulo\lingui{DEFIS}
\artiklo{verso}%
\vorto{verso} Serio di vorti, ritmoza e mezurita segun la duro di la silabi (longi e kurti) o segun la nombro di la silabi\lingui{DEFIS}
\artiklo{versoro}%
\vorto{versoro}\fako{matem.} Lineo di qua la longeso e la situeso karakterizas nombro imaginara\lingui{EF}
\artiklo{vertebro}%
\vorto{vertebro}\fako{anat.} Singla ek la 24 osti ye qui konsistas la spino\lingui{DEFIS}
\artiklo{vertico}%
\vorto{vertico}\fako{anat.} La suprajo di la kapo\lingui{EFIS}
\artiklo{vertijar}%
\vorto{vertijar}\fako{netrans.) (patol.} Subisar ta sturdo kurta-dura dum qua semblas ke la objekti jiras o rotacas avan la okuli\lingui{EFIS}
\artiklo{vertikala}%
\vorto{vertikala}\fako{geom.} Perpendikla ye la plano horizontala\lingui{DEFIRS}
\artiklo{vertuo}%
\vorto{vertuo} Praktiko kustumata di lo etiko-bona\lingui{EFIS}
\artiklo{veruko}%
\vorto{veruko}\senco{1}{Mikra surkreskajo pelala, surfacala}\senco{2}{Protuberanco sur la kortico di arboro, sur folio}\lingui{FIS}
\artiklo{vervo}%
\vorto{vervo} Inspireso ardoroza, nur kelka-dura\lingui{DEF}
\artiklo{vespero}%
\vorto{vespero} La parto finala di la jorno-periodo\lingui{EFL}
\artiklo{vespertilio}%
\vorto{vespertilio}\fako{zool.} Mikra mamifero, karnavida, noktala, di qua la osti di la membri avana unionesas a la korpo di la membri dopa per membrano quaze ala\latina{vespertilio}\lingui{L}
\artiklo{vespo}%
\vorto{vespo}\fako{zool.} Insekto, genero ek la tribuo « himenopteri », di qua la femino (qua havas dardo longa, quale la abelo) konstruktas alveoli\latina{vespa}\lingui{DEIS}
\artiklo{vespri}%
\vorto{vespri}\fako{liturgio katol.} Recito di la oficio deala quan onu agis olim cirkum vespero, e nun posdimeze\lingui{DEFIS}
\artiklo{vestalo}%
\vorto{vestalo}\fako{historio di la Roma antiqua} Sacerdotino di Vesta, qua devis permanar virga vivo-dure, e di qua la rolo esis entratenar la fairo sakra. – \textit{(metaf.)} Muliero tre chasta\lingui{DEFIRS}
\artiklo{vestibulo}%
\vorto{vestibulo}\senco{1}{Chambro qua, situita ye la enireyo di edifico, di domo, igas acesar olca}\senco{2}{Placo avan la pordo precipua di kirko o katedralo}\lingui{DEFIS}
\artiklo{vestio}%
\vorto{vestio} Vesto sen baski\lingui{F}
\artiklo{vesto}%
\vorto{vesto} To quo uzesas por kovrar la korpo homala, skope di shirmar lu de la domaji da vetero mala, o cesigar lua nudeso\lingui{EFIS}
\artiklo{vestono}%
\vorto{vestono} Vesto virala, kelke plu kurta kam vestio\lingui{F}
\artiklo{vetar}%
\vorto{vetar}\fako{trans.} Refuzar sua sanciono di (ulo)\lingui{EF}
\artiklo{veterano}%
\vorto{veterano}\senco{1}{\textit{(olim, en Roma)} Soldato qua obtenabis sua konjedo definitiva, pos armeanesir dum multa tempo}\senco{2}{\textit{(nun)} Soldato qua esas armeano de multa tempo}\lingui{DEFIRS}
\artiklo{veterinaro}%
\vorto{veterinaro} Mediko qua kuracas la morbi di la kavali, hundi, e c.\lingui{DEFIRS}
\artiklo{vetero}%
\vorto{vetero}\fako{meteorol.} La stando di atmosfero\lingui{DE}
\artiklo{vetivero}%
\vorto{vetivero}\fako{bot.} Planto de India, uzata por facar parfumi\latina{andropogon vetiveris}\lingui{DEFS}
\artiklo{vetrino}%
\vorto{vetrino}\fako{tekn.}\senco{1}{La parto di la butiko qua separesas de la strado per nur granda panelo de vitro}\senco{2}{Armoro, tablo sur qua esas instalita kesto ek vitro, en qua onu expozas la objekti skarsa o frajila}\lingui{FI}
\artiklo{veturo}%
\vorto{veturo} Vehilo por transporti, ek framo muntita sur roti, quan tranas kavali, muli, e c. od homo\lingui{FI}
\artiklo{vexar}%
\vorto{vexar}\fako{trans., pri, pro} Igar (ulu) deskontenta, despitanta e pasable shamoza ed interne iracoza, kande ulu amuzas su per kontre-agar lu pri mikraji\lingui{DE}
\artiklo{vezikatorio}%
\vorto{vezikatorio} Substanco qua produktas veziketi sur la pelo\lingui{DEFIS}
\artiklo{veziko}%
\vorto{veziko}\senco{1}{\textit{(anat.)} Quaza sako, posho, en qua urino esas akumule, til kande la parieti kontraktas por ekpulsar lu}\senco{2}{Saketo en la abdomino di la maxim multa fishi, qua kontenas mixuro de azoto, oxigeno e karbo-bioxo, ed igas la fisho plu lejera}\senco{3}{Ta sako, tirita ek la korpo di la animalo, sikigita ed inflita da la aero quan onu ensuflis}\lingui{EFIL}
\artiklo{vi}%
\vorto{vi}\fako{gram.} Pluralo di « vu »
\artiklo{viadukto}%
\vorto{viadukto}\fako{tekn.} Ponto alta, kun, multa-kaze, plura rangi de arkadi, konstruktita super fluvio, rivero, valeto, e c.\lingui{DEFIRS}
\artiklo{viatiko}%
\vorto{viatiko}\senco{1}{Subsidio a persono por lua spensi dum lua voyajo}\senco{2}{\textit{(religio katol.)} Helpilo por acesar la vivo sequonta, la vivo *nexta, eukaristio portata a persono extreme malada por preparar lu pri morto bona}\lingui{DEFIS}
\artiklo{vibrar}%
\vorto{vibrar}\fako{netrans., de, pro} Agar ocilo duopla, movo alternanta di la molekuli di korpo, qua transmisesas proximope a la medio\lingui{DEFIS}
\artiklo{vibriono}%
\vorto{vibriono}\senco{1}{Ento mikroskopala, forme di filamento tenua}\senco{2}{\textit{(biol.) (opoze a « bakterio »)} Ento mikroskopala qua semblas avancar per movo vermo-forma}
\artiklo{viburno}%
\vorto{viburno}\fako{bot.} Arboreto kun folii blanka e beri reda, dispozita buketatre, ek la familio « kaprifoliacei »\latina{viburnus}\lingui{EFIS}
\artiklo{vice}%
\vorto{vice}\senco{1}{Qua pre-nominesis por helpar o suplear (sua superioro nemediata)}\senco{2}{Indikas « remplase di… »}\senco{3}{\textit{(kande sequata da infinitivo)} indikas ago kontrea ye olta quan onu jus parolis, aludis, mencionis. (ex. : vice laborar, lu ludis)}
\artiklo{vicero}%
\vorto{vicero}\fako{anat.} Omna organo vivala esencala quan kontenas un ek la tri kavaji : kranio, torako od abdomino\lingui{EFIS}
\artiklo{vicina}%
\vorto{vicina} Qua esas situita proxime\lingui{EFIS}
\artiklo{vicio}%
\vorto{vicio} Tendenco kustumala a lo etiko-mala\lingui{EFIS}
\artiklo{vidamo}%
\vorto{vidamo}\fako{feudismo} Reprezentanto di episkopo, di abado, di qua la tasko konsistis ye defensar lu ed, en la kazi maxim multa, feudizita da lu\lingui{EFIRS}
\artiklo{vidar}%
\vorto{vidar}\fako{trans.} Perceptar la imaji quin la lumo-radii qui departas de la objekti lumizata formacas sur la fundo di la okulo por konvergar adsur retino\lingui{FIR}
\artiklo{vidva}%
\vorto{vidva} Qua cesis havar spozo, pro la morto di ta spozo\lingui{DEFIRS}
\artiklo{vielo}%
\vorto{vielo}\fako{muziko} Muzikilo kordoza, quan onu pleas per klavi quin la manuo sinistra igas kontaktar rotaco-roto indutita ye kolofono qua pleas la rolo di arketo\lingui{DEF}
\artiklo{vigilar}%
\vorto{vigilar}\fako{netrans.} Ne dormar dum la tempo normale dormokonsakrita\lingui{EFIS}
\artiklo{vigilio}%
\vorto{vigilio}\fako{liturgio katol.} Hiero di granda festo, dum qua onu abstinencas e fastas\lingui{DEFIS}
\artiklo{vigoro}%
\vorto{vigoro} Forteso di la korpo, di la spirito o di la psiko, qua atingas sua maximo\lingui{EFIS}
\artiklo{vikario}%
\vorto{vikario}\fako{precipue pri religio} La persono qua, dum la absenteso di sua superioro, remplasas lu ed exekutas lua taski\lingui{DEFIRS}
\artiklo{viktimo}%
\vorto{viktimo}\senco{1}{Vivanto ofrata sakrifike a la deajo}\senco{2}{\textit{(metafore)} Persono qua sakrifikas vole sua vivo o sua feliceso ad ulu, ad ulo}\senco{3}{Persono qua sakrifikesas a la odio, a la venjo, di ulu}\lingui{EFIS}
\artiklo{viktualio}%
\vorto{viktualio} Singla unajo ek provizuro de nutrivi\lingui{DEFS}
\artiklo{vikuno}%
\vorto{vikuno}\fako{zool.} Ruminero sen korni, animalo lanoza ek la genero « lamao », devenanta de Peru\latina{auchenia vicunna}\lingui{DEFIRS}
\artiklo{vilajo}%
\vorto{vilajo} Mikra grupo de domi di rurani, qua ne havas stradi rekta-linea, nek limito cirkum-muro o ramparo-cirkuito\lingui{EFI}
\artiklo{vilanelo}%
\vorto{vilanelo} Kansono di vilajani, maxim-multa-kaze kun kupleti tri-versa e refreno, finanta per quatreno\lingui{DEFIRS}
\artiklo{vildo}%
\vorto{vildo} Animali quin onu kaptas chase (perdriki, qualii, lepori, kunikli, kapreoli)\lingui{D}
\artiklo{vilejar}%
\vorto{vilejar}\fako{netrans.} Sejornar en ruro dum la sezono bela\lingui{FI}
\artiklo{vilo}%
\vorto{vilo}\fako{anat.} Singla ek la pili sur la surfaco di ula membrani (mukozo di la stomako, e c.)\lingui{EFIS}
\artiklo{vilozito}%
\vorto{vilozito}\fako{anat.} Singla ek la mikra asperaji o saliaji qui kovras ula surfaci quin li igas aspektar piloza. (kom ex. : la viloziti intestinala)\lingui{EF}
\artiklo{vimplo}%
\vorto{vimplo}\senco{1}{Telo-tuko qua, en la kostumi di la regulierini, kovras la pektoro, la kolo, e cirkondas la infra parto di la vizajo}\senco{2}{Kamizeto sen maniki (ek tulo, dentelo, e c.) quan *weras la mulieri kun robo dekoltita}\lingui{E}
\artiklo{vinagro}%
\vorto{vinagro} Liquoro qua saporas acide e pikante, ek la acidofermentaco de vino\lingui{EFS}
\artiklo{vindar}%
\vorto{vindar}\fako{trans.} Elevar grand-esforce (kargajo) per mashino qua konsistas ye arboro o cilindro horizontala qua rotacas sur axo ed an-cirkum qua spulesas kordego an qua ligesas la kargajo\lingui{DE}
\artiklo{vinkar}%
\vorto{vinkar}\fako{trans.}\senco{1}{Igar su superesar (a) pri forteso, sua enemikon (en batalio, milit.), (b) pri habileso, talento, sua konkurencanton}\senco{2}{Superpasar to quo esas obstaklo}\lingui{EFIS}
\artiklo{vino}%
\vorto{vino} Suko fermentacita di vitoberi – drinkajo alkoholoza\lingui{DEFIRS}
\artiklo{vintro}%
\vorto{vintro} La sezono maxim kolda di la yaro, en qua la jorni esas maxime kurta, qua sequas autuno e pre-iras printempo\lingui{DE}
\artiklo{vinyeto}%
\vorto{vinyeto}\fako{tekn.} Ornamento qua reprezentas vito-branchi interplektita, supre di la pagino unesma di libro o di chapitro. Omna ornamento di la frontispico o di la pagini di libro. Graburo cirkondata da kartusho, Ornamento di la kovrilo di libro. Ornamento cirkum naztuko. Ornamento di paperfolii letrala\lingui{DEFIRS}
\artiklo{violacar}%
\vorto{violacar}\fako{trans.} Domajar, lezar (ulo) quan onu ya devas respektar. – Koitar muliero kontre elua volo\lingui{EFIS}
\artiklo{violentar}%
\vorto{violentar}\fako{trans.} Koaktar sovaje, per la uzo di vigoro sen-represa\lingui{EF}
\artiklo{violino}%
\vorto{violino}\fako{muziko} Instrumento kun kordi ed arketo, derivita de la viulo olima, igita plu mikra-dimensiona, e reduktita a quar kordi inter-akordita quintope (g, d, la, e)\lingui{DEFIS}
\artiklo{violo}%
\vorto{violo}\senco{1}{\textit{(bot.)} Planto ek la familio « violacei », kun mikra floro di qua la odoro esas dolca, e la koloro inter reda e blua…}\senco{2}{La floro di ta planto}\latina{viola}\lingui{EFIRS}
\artiklo{violoncelo}%
\vorto{violoncelo}\fako{muziko} Instrumento kun quar kordi ed arketo, ye un quinto sub alto, qua preske korespondas a la viulobaso olima, e quan onu pleas tenante lu inter sua genui\lingui{DEFIS}
\artiklo{vipero}%
\vorto{vipero}\fako{zool.} Reptero serpenta di qua la morduro esas venenoza\latina{vipera}\lingui{DEFIS}
\artiklo{virga}%
\vorto{virga}\senco{1}{\textit{(pri adolecantino homa)} Quan nula viro koitis}\senco{2}{\textit{(metaf.)} (a) Regiono virga : ne ja explorita; (b) Tero virga : ne ja kultivita; (c) Foresto virga : ne ja explotita; (d)}\lingui{EFIS}\subvorto{}{metalo virga}{Renkontrita pura en sulo}{}{}
\artiklo{viro}%
\vorto{viro} Adultulo homa, qua havas maxim-alta-grade la karakterizivi di la maskuli\lingui{EFIS}
\artiklo{virtuala}%
\vorto{virtuala}\senco{1}{\textit{(filoz.)} Qua povas efikar, ma ne ja agis lo}\senco{2}{\textit{(optiko)} La foko virtuala di spegulo : loko ube la lumo-radii diverganta quin reflektas la spegulo, se prolongite ideale, konvergus dop la spegulo}\lingui{DEFIRS}\subvorto{mekan.}{efiko virtuala}{Efiko quan agus moveblo se lu movesus ye quanto extreme mikra ye ta o ca instanto}{3}{}
\artiklo{virtuoza}%
\vorto{virtuoza} Di qua la talento igas (ta persono) ecepto-kazo\lingui{DEFIS}
\artiklo{virulenta}%
\vorto{virulenta}\fako{patol.} Qua kontenas la fonto di infekto e di transmiso morbala\lingui{DEFIS}
\artiklo{viruso}%
\vorto{viruso}\fako{patol.} Substanco qua esas la kauzo di infekto e di transmiso morbala\lingui{DEFIS}
\artiklo{visar}%
\vorto{visar}\fako{trans.) (tekn.} Rotacigar skrubo por sinkar lu (aden ulo)\lingui{F}
\artiklo{vishar}%
\vorto{vishar}\fako{trans.} Netigar (objekto) per pasigar sur olu ulo qua deprenas la strateto de polvo qua kovras lu; sekigar per froto\lingui{D}
\artiklo{viskomto}%
\vorto{viskomto}\senco{1}{Titulo siniorala, aplikata, en la komenco, a la lietnanto di la komto}\senco{2}{Nobeleso-titulo, sub komto e super barono}\lingui{DEFIRS}
\artiklo{viskoza}%
\vorto{viskoza} Di qua la molekuli esas tenaca, inter-adheras ed adheras a la korpo quan li tushas\lingui{DEFIS}
\artiklo{vistar}%
\vorto{vistar}\fako{trans.} Aplikar (sur akto) la formulo qua konstatas ke lu esas verifikita da la persono o da la organismo di qua la signaturo esas necesa por igar olu valida\lingui{EIS}
\artiklo{vitalismo}%
\vorto{vitalismo}\fako{filoz.} Doktrino segun qua la fenomeni fiziologiala havas kauzo specala, ne-sama kam olta di la fenomeni psikologiala\lingui{DEFIRS}
\artiklo{vitelo}%
\vorto{vitelo} La parto flava di la ovo\lingui{EFIS}
\artiklo{vito}%
\vorto{vito}\fako{bot.} Planto kun stipo lignatra, qua produktas la beri de qui facesas vino\latina{vitis vinifera}\lingui{EFIS}
\artiklo{vitralio}%
\vorto{vitralio} Panelo de vitro-plaki, asemblita fako-formacante, e (maxim-multa-kaze) dekorita ye pikturi\lingui{F}
\artiklo{vitriolo}%
\vorto{vitriolo} Nomo vulgara di SO₄H₂\lingui{DEFIS}
\artiklo{vitro}%
\vorto{vitro} Materio harda, frajila, diafana, quan onu obtenas per la kunfuzo di sablo silicoza kun potaso e natroxo\lingui{EFIS}
\artiklo{viulo}%
\vorto{viulo}\fako{olim} Nomo di muzikili diversa; kun kordi ed arketo, forme di granda violino kun kin, sis o sep kordi, ne-egale grosa\lingui{DEFIRS}
\artiklo{vivaca}%
\vorto{vivaca} Rapida ed animata pri lua ago-maniero\lingui{EFIS}
\artiklo{vivar}%
\vorto{vivar}\fako{netrans.}\senco{1}{\textit{(che la animalo)} Posedar e praktikar la funcioni « nutro », « genito »}\senco{2}{\textit{(che la homo)} Posedar e praktikar la funcioni « nutro », « genito », « raciono » ed « arbitrio »}\lingui{EFIS}
\artiklo{vivipara}%
\vorto{vivipara} Qua parturas ja vivanta sua yuni (kontraste kun « ovipara »)\lingui{EF}
\artiklo{vizajo}%
\vorto{vizajo} La facio di la kapo \textit{(che la homo)}\lingui{EFI}
\artiklo{vizar}%
\vorto{vizar}\fako{trans.} Regardar tre atencoze (lor sua lanso o jus ante sua lanso di objekto ad la punto quan onu volas atingar)
\artiklo{vizelo}%
\vorto{vizelo}\fako{zool.} Varietato ek la genero « martro », kun korpo svelta, longa, falva, falveso plu klara subventre, muzelo pinto-forma\latina{mustela vulgaris}\lingui{DE}
\artiklo{viziero}%
\vorto{viziero} Parto avana di la kasko, quan onu abasis por shirmar la vizajo, e tra qua onu povis vidar e respirar per greto streta. – Peco ek ledro muntita an-avan kasqueto por shirmar la okuli kontre la suno-radii e kontre la pluvoguti\lingui{DEFIS}
\artiklo{viziono}%
\vorto{viziono} Iluziono per qua onu kredas vidar ulu od ulo qua ne esas prezenta koram onua okuli\lingui{DEFIS}
\artiklo{viziro}%
\vorto{viziro}\fako{en Turkia olima} Oficiro ek la *konselio (konsilantaro) di la sultano\lingui{DEFIRS}
\artiklo{vizitar}%
\vorto{vizitar}\fako{trans.} Irar renkontre che ulu, pro devo di politeso, di deferenco\lingui{DEFIRS}
\artiklo{voco}%
\vorto{voco}\senco{1}{(a) Sono de la aero ekpulsita de la pulmoni, kande ta aero trairas laringo, qua igas ica vibrar. (b) Sono artikulata, parolo. (c) Sono modulata segun la skalo muzikala, aptigita pri kantado}\senco{2}{Expreso di la opiniono di persono lor voto}\senco{3}{\textit{(metaf.)} \textit{(gram.)} Formo di la konjugo, qua indikas se la efiko quan expresas la verbo esas agata o subisata da la subjekto}\lingui{EFIS}
\artiklo{vodevilo}%
\vorto{vodevilo} Teatrajo mixita kun kupleti\lingui{DEFIR}
\artiklo{vokalizar}%
\vorto{vokalizar}\fako{netrans.) (muziko}\senco{1}{Exercar su voce, per kantar la noto sen pronuncar lua nomo e sen vortizar la sonuno}\senco{2}{Adjuntar traiti, ornivi, en kantajo}\lingui{DEFIRS}
\artiklo{vokalo}%
\vorto{vokalo}\senco{1}{Sonuno simpla, emisata per lasar la aero ekirar de laringo e trairar la boko sen obstaklo}\senco{2}{Singla ek la sonuni diversa qui genitesas tale, segun la maniero quale onu apertas sua boko por lasar li ekirar. (la vokali : a, e, i, o, u)}\lingui{DEFIS}
\artiklo{vokar}%
\vorto{vokar}\fako{trans.} Invitar (ulu) pri venar, per pronuncar lua nomo, o per signalo\lingui{EIL}
\artiklo{vokativo}%
\vorto{vokativo}\fako{gram.} Un ek la kazi di la deklino \textit{(che ula lingui nacionala)} quan onu uzas kande onu su turnas a persono od a kozo quan onu egardas quale personon\lingui{DEFIS}
\artiklo{volar}%
\vorto{volar}\fako{trans.}\senco{1}{Rezolvar agor o ne agor (ulo)}\senco{2}{Praktikar plu o min energioze sua povo pri rezolvar koncerne agar (ulo) od igar ulu agar (ulo)}\senco{3}{Rezolvar ferme, obtenor (ulo) por su o por altru}\lingui{DEFIRS}
\artiklo{volatila}%
\vorto{volatila}\fako{aludante substanco solida o liquida} Qua vaporeskas facile\lingui{EFIS}
\artiklo{volfo}%
\vorto{volfo}\fako{zool.} Quadripedo karnavida, kun muzelo longa e nigra, oreli rekta, pilaro flavatra e nigre-buntita, kaudo tufatra\latina{canis lupus}\lingui{DE}
\artiklo{volkano}%
\vorto{volkano}\senco{1}{Abismo subtera, (normale) interne di monto, qua lansas, per aperturo (kratero), flami, fumuro, cindri, e materii fuzura}\senco{2}{\textit{(metaf.)} (a)}\lingui{DEFIRS}\subvorto{}{esar sur volkano}{En stando politikala quan karakterizas la minaco da revoluciono. (b)}{}{}\subvorto{}{viro volkana}{Di qua la karaktero esas impetuoza}{}{}
\artiklo{volontario}%
\vorto{volontario} Qua armeanesas, qua su engajis pri armeanesko, sen obligeso pri lo\lingui{DEFIRS}
\artiklo{voltar}%
\vorto{voltar}\fako{netrans.}\senco{1}{\textit{(manejo)} Igar (kavalo) jiro-movar plurafoye}\senco{2}{\textit{(skermo)} Chanjar sua plaso por evitar tusheso da la espado di sua adverso}\senco{3}{\textit{(kartoludo)} Stroko ye qua onu agas omna preni}\lingui{DEFIRS}
\artiklo{voltijar}%
\vorto{voltijar}\fako{netrans.}\senco{1}{\textit{(kavalk-arto)} Exercar su gimnastike por kustumigar su pri saltar sen estribi adsur kavalo qua avancas paze, trote, galope}\senco{2}{\textit{(akrobat-arto)} Exercar su pri dansar sur kordo laxa}\lingui{DEFIR}
\artiklo{voltmetro}%
\vorto{voltmetro}\fako{tekn.} Sorto di galvanometro, per qua onu mezuras la forco elektromovala o la falo di la potencialo inter du punti\lingui{DEFIRS}
\artiklo{volto}%
\vorto{volto}\fako{elektro} Unajo pri la interdifero di la potenciali elektrala = 108 unaji elektromagnetala C.G.S.\lingui{DEFIRS}\subvorto{}{volto internaciona}{Tensiono elektrala qua, aplikate a konduktivo di qua la rezistiveso egalesas un \textit{ohm} internaciona, genitas fluo de un \textit{ampero} internaciona}{}{}
\artiklo{volumino}%
\vorto{volumino}\fako{geom.} Parto di spaco, quan okupas korpo\lingui{DEFIS}
\artiklo{voluntar}%
\vorto{voluntar}\fako{trans.} Esar benigna, komplezema, servema (*serviema) til konsentar (agor co o to)\lingui{EFI}
\artiklo{volupto}%
\vorto{volupto} Plezuro di la sensi, qua igas juar forte\lingui{EFIS}
\artiklo{voluto}%
\vorto{voluto}\fako{arkitekt.} Ornamento spiralforma an la kapitelo di kolono (precipue en la klasifiko-grupo « ionika »), di konsolo, di modiliono\lingui{DEFIRS}
\artiklo{volvar}%
\vorto{volvar}\fako{trans.} Aplikar rule kozo an-cirkum altra\lingui{L}
\artiklo{vomar}%
\vorto{vomar}\fako{trans.}\senco{1}{Retrojetar konvulso, per sua boko (la materii solida o liquida qui esis en la stomako)}\senco{2}{\textit{(metaf.)} Pulsar ek su (ulo mala)}\lingui{EFIS}
\artiklo{vomero}%
\vorto{vomero}\fako{anat.} Osto sen-simetraja, dina, quadrilatera, anomala, ye qua konsistas la parto supra di la parieto di la naztrui
\artiklo{vomiknuco}%
\vorto{vomiknuco}\fako{bot.} Bero venenoza ek arboreto de India\latina{strychnos nux vomica}\lingui{EFIS}
\artiklo{vomiko}%
\vorto{vomiko}\fako{medic.} Amaso de materio pusoza qua naskas en la pulmoni (che la tuberklosiki) e forjetesas per quaza vomo
\artiklo{vomitorio}%
\vorto{vomitorio}\fako{en la Roma antiqua} Larja aperturo en la klozilo por lasar pasar la spekteri, en la amfiteatro qua destinesis a la ludi publika. (*publicala)\lingui{DEFIS}
\artiklo{vorticar}%
\vorto{vorticar}\fako{netrans.} Forportesar per jirado rapida\lingui{EI}
\artiklo{vorto}%
\vorto{vorto} Son-uno artikulata, ek un o plura silabi, uzata por reprezentar ento od eso-maniero. – Skriburo qua reprezentas ta son-uno\lingui{DE}
\artiklo{votar}%
\vorto{votar}\fako{netrans., pri, por, kontre} (\textit{aludante asemblitaro} \textit{en qua la rezolvi esas agenda segun la opiniono da la} majoritato – \textit{ed aludante singla membro}) Depozar bilieto, bulo, e c. qua indikas ke (ta membro) adoptas o ne-adoptas ta o ca propozajo\lingui{DEFIRS}
\artiklo{vovar}%
\vorto{vovar}\fako{trans.} Promisar a la deajo exekutor ta o ca tasko, verko meritoza, se lu exaucos la demando quan onu agas a lu\lingui{EF}
\artiklo{voyajar}%
\vorto{voyajar}\fako{netrans., de, ad} Diplasar su, parkurante voyo plu o min longa, por irar til altra urbo, lando o kontinento (o mem planeto)\lingui{EFIS}
\artiklo{voyo}%
\vorto{voyo}\senco{1}{La spaco parkurenda por irar de loko ad altra loko, equipita por la cirkulado di la veturi e di la pediranti}\senco{2}{La spaco parkurita da la vildo quan onu persequas e quan revelas la traci quin lu lasis}\lingui{F}
\artiklo{vu}%
\vorto{vu}\fako{gram.} Pronomo personala qua, dum nefamiliareso ye ulu, remplasas « tu »\lingui{FIR}
\artiklo{vulgara}%
\vorto{vulgara}\senco{1}{Qua admisesas, uzesas da la homi ordinara}\senco{2}{Quan nulo distingas de la homi ordinara}\lingui{DEFIRS}
\artiklo{vulgato}%
\vorto{vulgato}\fako{religio kristana} Tradukuro Latina de la texto Hebrea di biblo, agnoskata kom konforma a la kanoni di la eklezio katolika\lingui{DEFIRS}
\artiklo{vulkanala}%
\vorto{vulkanala}\fako{geol.} Qua imputas a fairo la formaco di la Terglobo
\artiklo{vulkanizar}%
\vorto{vulkanizar}\fako{trans.) (tekn.} Preparar (kauchuko) per imersar lu en sulfo-fuzita por igar lu ne apta influesar da la fluktui di temperaturo e di vetero\lingui{DEFIRS}
\artiklo{vulnerario}%
\vorto{vulnerario}\fako{bot.} Planto leguminosa, kun flori flava, reputata kom apta remediar la vunduri, la plagi\latina{vulneraria rustica}\lingui{FIS}
\artiklo{vulto}%
\vorto{vulto} Konstrukturo di qua la suprajo esas kurvatra, ek asembluro de quadri infre-kono-trunka qui inter-apogas\lingui{EFI}
\artiklo{vulturo}%
\vorto{vulturo}\senco{1}{\textit{(zool.)} Granda raptucelo, kun kapo e kolo nuda, garnisita ye nur lanugo, sen plumi}\senco{2}{\textit{(metaf.)} La persono da qua la febli esas viktima}\latina{vultur}\lingui{EFI}
\artiklo{vulvo}%
\vorto{vulvo}\fako{anat.} Orifico di la organi genitala che la femino\lingui{DEFIRS}
\artiklo{vundar}%
\vorto{vundar}\fako{trans.) (ye}\senco{1}{Frapar per stroko qua lezas, detrimentas la organismo}\senco{2}{\textit{(aludante objekti)} Lezar parto di korpo per frotado, aracho, e c.}\senco{3}{Afektar penigive (organo di la sensi)}\lingui{DE}
\end{multicols}
\sekciono{W}
\begin{multicols}{2}
\artiklo{waranto}%
\vorto{waranto}\fako{komerco} Recev-atesto, donata a komercisto qua depozas vari en depozeyo qua apartenas a la dogan-administrantaro\lingui{DEFIS}
\artiklo{warfo}%
\vorto{warfo} Amaso de petri, ronda stoni, sablo, e c., jetita an la enireyo di portuo, interligita forte, e subtenata per palisegaro sinkita en sulo – destinita ruptor la impetuo di la ondegi\lingui{DEF}
\artiklo{warpo}%
\vorto{warpo}\fako{tekn.} Asembluro ek fili longesala interparalela, tensita inter la du roleri di la texo-mashino, quin onu separas por lasar la fili transversa (para o ne-para) pasar, fili qui ri-nodigesas (forme di kateno) ye lia extremaji por evitar lia intermixo\lingui{E}
\artiklo{wato}%
\vorto{wato}\fako{elektro} La milima parto di la kilowato, unajo di elektropovo (sistemo M.T.S.)\lingui{DEF}
\artiklo{wefto}%
\vorto{wefto}\fako{tekn.) (aludante texo-mashino} La filo quan onu pasigas, per la naveto, tra la fili tensita longesale (warpo)
\artiklo{weldar}%
\vorto{weldar}\fako{trans.) (tekn.} Soldar per fuzar la metali asemblenda, opoze a la soldo per kunjunto di du korpi asemblenda per metalo od aloyuro di qua la fuzo-punto esas tam infra (sur la temperaturo-skalo) kam posibla\lingui{E}
\artiklo{*werar}%
\vorto{*werar}\fako{trans.} Havar kustumale (od, eceptale, dum ceremonio) sur sua korpo kom vesto; od an su, kom juvelo, ornivo\lingui{E}
\artiklo{westo}%
\vorto{westo} Ta ek la quar punti kardinala qua esas sur la latero ube Suno su kushas\lingui{DEFIRS}
\artiklo{*wishar}%
\vorto{*wishar}\fako{trans., ad, por} Expresar la deziro ke ulo eventez (por su, o por altra persono)\lingui{DE}
\artiklo{wiskio}%
\vorto{wiskio} Brandio ek cerealo-grani, de Skotia ed Irlando\lingui{DEFIRS}
\artiklo{wistitio}%
\vorto{wistitio}\fako{zool.} Mikra simio de Brazilia\latina{simia yacchus}\lingui{DEFIR}
\artiklo{wisto}%
\vorto{wisto} Ludo per 52 karti, da 2 personi kontre 2 personi (o da 3 personi kun un karto-grupo deskovrita (« mortinto ») qua uzesas kom partnero, per partii interfederita) maxim-multakaze 10-punte, ma en ula kazi nur 5-punte\lingui{DEFIRS}
\artiklo{wolframo}%
\vorto{wolframo}\fako{kemio} Korpo simpla, metalo obskur-griza, tre harda, deskovrita en 1781 da Scheele, atom-pezo 184,8 – uzata kom filamento por la ampuli elektrala\simbolo{Tu o Wo}
\end{multicols}
\sekciono{X}
\begin{multicols}{2}
\artiklo{xantino}%
\vorto{xantino}\fako{kemio} La C₅H₄N₄O₂ quan kontenas multa liquidi de animali : sango, urino, ed anke teo
\artiklo{xantofilo}%
\vorto{xantofilo}\fako{kemio} C₄₀H₅₆O₂, kristali flava qui fuzas ye 172° C., e renkontresas en la klorofilo di la vejetanto-folii
\artiklo{xenio}%
\vorto{xenio}\senco{1}{\textit{(versif.)} Epigramo satira (enduktita da Goethe e Schiller)}\senco{2}{\textit{(biol.)} Influo eventuala da la elemento maskula, en ula kazi, ad ula parti di la femino qui, pro lo, aquiras la karakteri di la maskulo fekundiginta. (Dicesas precipue pri la planti di qui la fekundeso devenas de poleno stranjera)}
\artiklo{xenono}%
\vorto{xenono}\fako{kemio} Un ek la gasi skarsa di atmosfero, monoatoma; atom-pezo 130,2\simbolo{X}
\artiklo{xerofito}%
\vorto{xerofito}\fako{biol.} Singla de la vejetanti quin karakterizas la hardeso di la organi diversa e la sekreco samtempa di olei eteroza abundanta
\artiklo{xifoida}%
\vorto{xifoida}\fako{anat.} Dicesas pri la apendico – prolonguro kartilaga, qua finas la extremajo infra di sternumo
\artiklo{xilofago}%
\vorto{xilofago} Singla de la insekti qui rodas ligno\lingui{EFIS}
\artiklo{xilofono}%
\vorto{xilofono}\fako{muziko} Instrumento ek ligno-tabuleti ne-interegale longa, subtenata da du suportili, e sur qui onu frapas per du vergeti ligna\lingui{DEFIS}
\artiklo{xilografar}%
\vorto{xilografar}\fako{trans.} Imprimar ek tipi ligna, od ek tabuleti ligna sur qui esas grabita vorti e figuri\lingui{DEFIRS}
\end{multicols}
\sekciono{Y}
\begin{multicols}{2}
\artiklo{ya}%
\vorto{ya}\fako{adverbo} Partikulo qua nulkaze uzesas sole, ma kun irga dico, e konfirmas ica emfazoze\lingui{D}
\artiklo{yako}%
\vorto{yako}\fako{zool.} Bufalo kun kaval-kaudo, devenanta de Tibet\latina{poephagus}\lingui{F}
\artiklo{yakto}%
\vorto{yakto} Exkurso-navo\lingui{DEFRS}
\artiklo{yardo}%
\vorto{yardo}\fako{navig.} Ligno-peco, dinigita ye lua extremaji, lokizita horizontale ye altesi diversa, sur la masti di navo, por subtenar la veli desvolvita o volvita\lingui{E}
\artiklo{yaro}%
\vorto{yaro}\senco{1}{Duro di un jiro da la Terglobo cirkum Suno (365 dii, 5 hori, 48 minuti, e 50 sekundi), egardata kom unajo di la mezuro di tempo}\senco{2}{Ta periodo determinata per ula quanto de fakti qui eventis dum olca}\senco{3}{Periodo de dek-e-du monati, egardata relate to quo eventas dum li, sen egardo pri la epoko kande li komencas}\senco{4}{Singla periodo de deke-du monati, kontita de la nasko-dio di persono}\lingui{DE}
\artiklo{yatagano}%
\vorto{yatagano} Sabro poniardo-forma, kun lamo kelke kurva proxim la pinto – armo Araba e Turka\lingui{DEFIS}
\artiklo{ye}%
\vorto{ye} Prepoziciono di qua la senco ne esas determinita, quan onu uzas nur en la kazi kande nula prepoziciono altra postulesas da la senco. (Lu indikas nome la loko o la dato exakta di evento, di fakto, la parto koncernata di la korpo : « ye la angulo di la strado », « ye la dekesma kilometro », « ye dimezo », « ye la lasta foyo », « me doloras ye la kapo », « ilu prenis elu ye la tayo », « il kaptis la kavalo per lazo ye la kolo », « plenigar ulo ye … per kuliero, spado, e c. »)
\artiklo{yen}%
\vorto{yen} Prepoziciono qua igas atencar la persono o la kozo pri qua onu quik parolos, o to quo quik eventos. \textit{(yen, avan nomo; yen ke, avan propoziciono)}
\artiklo{yes}%
\vorto{yes} Adverbo qua equivalas propoziciono per qua onu respondas afirme a questiono (expresita o tacita). \textit{(yes ya = *si)}\lingui{E}
\artiklo{yiterbo}%
\vorto{yiterbo}\fako{kemio} Substanco ne-izolita, di qua onu konocas nur la oxido, e qua ne esus altro kam mixuro de lutecio e di neoyiterbio\simbolo{Yb}
\artiklo{yitro}%
\vorto{yitro}\fako{kemio} Metalo skarsa, qua akompanas cerio en la maxim multa de lua erci. Atom-pezo : 88,92\simbolo{Y}
\artiklo{yodlar}%
\vorto{yodlar}\fako{netrans.} Kantar ta kansono montanala en qua onu guturagas stroke por pasar de la voco « pektorala » a la voco « kapala » (voco di qua la rezono eventas en la kavaji supra od infra, rispektive, di la organaro vocala), di qua la rezultajo esas quaza rukulo tre stranja\lingui{DEF}
\artiklo{yolo}%
\vorto{yolo}\fako{navig.} Mikra kanoto, streta e lejera, qua avancas rapide\lingui{DEFIR}
\artiklo{yufto}%
\vorto{yufto} Ledro traktita per birko-oleo\lingui{DR}
\artiklo{yugo}%
\vorto{yugo}\senco{1}{Ligno-peco quan onu muntas sur la kapo di la bovi e per qua onu jungas li duope}\senco{2}{\textit{(metaf.)} (a) Submiseso quan impozas mastro. (b) Koakteso quan impozas engajo, pasiono, e c.}\senco{3}{\textit{(en la Roma antiqua)} Piquo pozita horizontale sur du piqui vertikala, stekita en sulo, sub qua onu igis pasar la enemiki vinkita por humiligar e shamigar li}\lingui{DEFIS}
\artiklo{yuko}%
\vorto{yuko}\fako{bot.} Planto perena, ek la familio « labiacei », kun folii longa, streta e harda, qui finas per pinto akuta; kun flori blanka, dispozita panikule sur shafto longa\latina{yucca}\lingui{DEFIS}
\artiklo{*yulo}%
\vorto{*yulo} Festo-dio pri la vintro-solstico, konvencione situita ye la 25. di decembro – e qua tale koincidas kun la festo-dio Kristo-naskala di la kristani\lingui{DEFIRS}
\artiklo{yuna}%
\vorto{yuna} Qua ne evas multe\lingui{DEFI}
\artiklo{-yuno}%
\vorto{-yuno} Sufixo qua indikas la filio tre yuna di animalo
\artiklo{yuro}%
\vorto{yuro}\senco{1}{Lego-justifikeso pri postular ke onu retrodonez to quon onu debas a ni}\senco{2}{Ensemblo ek la legi qui determinas to quon singlu, segun sua situeso, darfas legitime agar (o ne-agar) o facar (o ne-facar) o postular de altru}\lingui{DEFIRS}
\artiklo{yusta}%
\vorto{yusta} Konforma a la yuro; qua respektas la yuro di la kun-homi\lingui{EF}
\end{multicols}
\sekciono{Z}
\begin{multicols}{2}
\artiklo{zagayo}%
\vorto{zagayo} Sorto di javelino, facata ed uzata da ula populi sovaja
\artiklo{zebro}%
\vorto{zebro}\fako{zool.} Animalo di Afrika, aspekto-vicina ye asno, kun pilaro flava e strii bruna\latina{equus zebra}\lingui{DEFIRS}
\artiklo{zebuo}%
\vorto{zebuo}\fako{zool.} Bovo kun un gibo, devenanta de India\latina{bos taurus indicus}\lingui{DEFIRS}
\artiklo{zefiro}%
\vorto{zefiro}\senco{1}{\textit{(en la epoki antiqua)} Vento de la ocidento equinoxala}\senco{2}{Irga vento lejera e dolca}\lingui{DEFIS}
\artiklo{zeko}%
\vorto{zeko}\fako{zool.} Insekto sen-ala qua su aplikas an la pelo di la bovi, mutoni, hundi, e sugas lia sango\latina{ixodes}\lingui{DI}
\artiklo{zelar}%
\vorto{zelar}\fako{netrans.) (por, pri} Agar propra-move, sen tardeso, kun komplezo e fervoro, por servar (o *serviar) ulu, ulo\lingui{EFIS}
\artiklo{zenito}%
\vorto{zenito}\senco{1}{\textit{(astron.)} La punto ube la vertikalo di loko iras renkontre di la sfero cielala}\senco{2}{\textit{(metaf.)} La kulmino-punto di kozo}\lingui{DEFIRS}
\artiklo{zeolito}%
\vorto{zeolito}\fako{mineral.} Silikato aluminoza hidrata\lingui{DEF}
\artiklo{zero}%
\vorto{zero}\fako{matem.}\senco{1}{Cifro, 0-forma, qua – sen valorar ipse – lokizite dextre di altra cifri, igas dekoplan la valoro di ici, per suplear la unaji decimala qui mankas}\senco{2}{\textit{(fiziko)} En la instrumenti uzata por mezurar la temperaturi, la preso da aero, e c., la departo-punto di la grado-konto}\lingui{EFIS}
\artiklo{zesto}%
\vorto{zesto}\fako{bot.}\senco{1}{Parieto membrana qua dividas quarime la internajo di nuco}\senco{2}{Shelo extera di oranjo e di citrono}\lingui{F}
\artiklo{zezear}%
\vorto{zezear}\fako{netrans.} Pro defekto, pronuncar le \textit{j} quale le \textit{z}, e \textit{sh} quale \textit{s}\lingui{FS}
\artiklo{zibelino}%
\vorto{zibelino}\senco{1}{\textit{(zool.)} Martro di Siberia, kun pili tenua}\senco{2}{La felo di ta animalo, uzata kom furo}\latina{martas zibellina}\lingui{FIS}
\artiklo{zigomato}%
\vorto{zigomato}\fako{anat.} Osto di la saliajo di la vango, sub la angulo extera di la okulo\lingui{EFIS}
\artiklo{zigomorfa}%
\vorto{zigomorfa}\fako{bot.} Dicesas pri la flori qui havas plano di simetreso (kom ex. : la violi, la muzelflori)\lingui{DEFIS}
\artiklo{zigosporo}%
\vorto{zigosporo}\fako{bot.} Organo genitala sexuoza qua rezultas de la fuzo di du gameti izogama o heterogama\lingui{DEFIS}
\artiklo{zigzago}%
\vorto{zigzago} Lineo ruptita, qua formacas anguli alterne salianta e konkava\lingui{DEFIRS}
\artiklo{zimazo}%
\vorto{zimazo}\fako{biol.} Substanco fermenta, azotoza, nesulfuroza, pasable parenta ye la albuminoidi, egardenda kom la reziduo di funcionado\lingui{DEFIS}
\artiklo{zimogeno}%
\vorto{zimogeno}\fako{biol.} Substanco quan sekrecas la celuli glandala e qua, mixite kun la aquo di la vakuoli, genitas fermento
\artiklo{zimologio}%
\vorto{zimologio}\fako{biol.} La parto di kemio qua traktas fermentaco\lingui{DEFIRS}
\artiklo{zimosimetro}%
\vorto{zimosimetro}\fako{fiziko} Instrumento uzata pri mezurar la grado di fermentaco di liquido\lingui{DEFIRS}
\artiklo{zimotekniko}%
\vorto{zimotekniko} La arto genitar e direktar la fermentaco\lingui{DEFIRS}
\artiklo{zinko}%
\vorto{zinko}\fako{kemio} Korpo simpla, metala, blue-blanka, kompozajo lameloza, duktila, fuzebla, e qua, pos lamino, uzesas por facar kovrili di tekti, pluvo-kanali tektala, tubi, balnokuvi, siteli, e c.\simbolo{Zn}
\artiklo{zinkografio}%
\vorto{zinkografio} Procedo analoga a litografado, ma en qua la petroplako litografala remplasesas da zinkoplako
\artiklo{zirkono}%
\vorto{zirkono}\fako{kemio} Korpo simpla, metala, kelke simila a titano. Atom-pezo : 90,6\simbolo{Zr}
\artiklo{zodiako}%
\vorto{zodiako}\fako{astron.} Zono cielala qua kovras 8 o 9 gradi ye singla latero di ekliptiko, ed embracas omna situeso di la planeti til nun deskovrita – e dividita aden 12 parti qui nomizesis segun la stelaro maxim vicina a li, e quin onu nomas la 12 signi di zodiako\lingui{DEFIRS}
\artiklo{zono}%
\vorto{zono}\senco{1}{\textit{(geom.)} Parto de la areo di sfero qua formacas quaza bendo kurva inter du plani interparalela}\senco{2}{\textit{(geog.)} Singla de la granda regioni di la Terglobo, quin separas cirkli imaginala, paralela ye equatoro, e di qui la klimato dependas de lia situeso kompare a Suno}\senco{3}{Bendo de stofo, de ledro, e c., destinata a klemar la vesti an la tayo}\senco{4}{To quo cirkondas quale vesto-zono}\lingui{DEFIRS}
\artiklo{zoocecidio}%
\vorto{zoocecidio} Deformuro genitata da animalo o planto vivanta
\artiklo{zoofito}%
\vorto{zoofito} Nomo di la animali inferiora, vicina ye la planti, qui formacas la brancho quaresma di la animal-regno (de qua onu deprenis nia-epoke la protozoi por formacar kinesma)\lingui{DEFIRS}
\artiklo{zoografio}%
\vorto{zoografio} Deskripto di la animali\lingui{EFIRS}
\artiklo{zoologio}%
\vorto{zoologio} Parto di la historio naturala qua traktas la animali\lingui{DEFIRS}
\artiklo{zoospermo}%
\vorto{zoospermo} (sinonimo di « spermatozoido »)\lingui{DEFIS}
\artiklo{zoosporo}%
\vorto{zoosporo}\fako{zool.} Sporo (korpuskulo genitala) quan karakterizas cilii vibranta, che ula algi\lingui{DEFIS}
\artiklo{zuavo}%
\vorto{zuavo}\fako{armeo} En Francia, soldato ek korpo de infantrio lejera, uzata en la armeo di Afrika\lingui{DEFIRS}
\artiklo{zumar}%
\vorto{zumar}\fako{netrans.}\senco{1}{Genitar bruiso kontinua, grava e matida, aludante insekto, o la murmuro konfuza ek la voci di plura personi asemblita}\senco{2}{Resonigar klosho per igar la frapilo shokar alterne dextre e sinistre, ma sen ociligar ta klosho}\senco{3}{\textit{(medic.)} Audar en sua oreli la sono quan genitas ibe la palpitado tro forta da la arterii, la stretesko di la tubo audala}
\end{multicols}
